% Generated mechanically from frozen input p01/ja/algebra.tex.
% Do not edit this generated file; edit the bound locale source instead.


% OK, start here.
%

\setcounter{chapter}{9}
\chapter{可換代数}



\phantomsection
\label{algebra-section-phantom}





\section{序論}
\label{algebra-section-introduction}

\noindent
本章では可換代数の基礎を説明する。参考文献として \cite{MatCA} がある。








\section{規約}
\label{algebra-section-conventions}

\noindent
環は単位元 $1$ をもつ可換環とする。零環も環である。実際、零環は素イデアルを
もたない唯一の環である。クロネッカー記号 $\delta_{ij}$ を用いる。$R \to S$ が
環準同型で、$\mathfrak q$ が $S$ の素イデアルであるとき、$R$ が $S$ の部分環で
なく、また $R \to S$ が単射でない場合でも、$R \to S$ による $\mathfrak q$ の
逆像である素イデアルを表すために「$\mathfrak p = R \cap \mathfrak q$」という
記法を用いる。






\section{基本概念}
\label{algebra-section-rings-basic}

\noindent
以下は可換代数の基本概念の一覧である。これらの概念の一部は後続の本文で
より詳しく論じ、一部はこの一覧中で定義するが、残りは基礎事項とみなして
定義しない。斜体で示した概念の大半になじみがない場合には、先に代数学の
入門書を参照することを勧める。

\begin{enumerate}
\item $R$ は{\it 環}である、
\label{algebra-item-ring}
\item $x\in R$ は{\it 冪零}である、
\label{algebra-item-ring-element-nilpotent}
\item $x\in R$ は{\it 零因子}である、
\label{algebra-item-ring-element-zerodivisor}
\item $x\in R$ は{\it 単元}である、
\label{algebra-item-ring-element-unit}
\item $e \in R$ は{\it 冪等元}である、
\label{algebra-item-ring-element-idempotent}
\item 冪等元 $e \in R$ が $e = 1$ または $e = 0$ を満たすとき、これを
{\it 自明}という、
\label{algebra-item-idempotent-trivial}
\item $\varphi : R_1 \to R_2$ は{\it 環準同型}である、
\label{algebra-item-ring-homomorphism}
\item
\label{algebra-item-ring-homomorphism-finite-presentation}
$\varphi : R_1 \to R_2$ は{\it 有限表示}である、または
{\it $R_2$ は有限表示 $R_1$-代数である}。定義 \ref{algebra-definition-finite-type}
を参照せよ、
\item
\label{algebra-item-ring-homomorphism-finite-type}
$\varphi : R_1 \to R_2$ は{\it 有限型}である、または
{\it $R_2$ は有限型 $R_1$-代数である}。定義 \ref{algebra-definition-finite-type}
を参照せよ、
\item
\label{algebra-item-ring-homomorphism-finite}
$\varphi : R_1 \to R_2$ は{\it 有限}である、または
{\it $R_2$ は有限 $R_1$-代数である}、
\item $R$ は{\it （整）域}である、
\label{algebra-item-ring-domain}
\item $R$ は{\it 被約}である、
\label{algebra-item-ring-reduced}
\item $R$ は{\it ネーター的}である、
\label{algebra-item-ring-Noetherian}
\item $R$ は{\it 単項イデアル整域}、すなわち{\it PID}である、
\label{algebra-item-ring-PID}
\item $R$ は{\it ユークリッド整域}である、
\label{algebra-item-ring-Euclidean}
\item $R$ は{\it 一意分解整域}、すなわち{\it UFD}である、
\label{algebra-item-ring-UFD}
\item $R$ は{\it 離散付値環}、すなわち{\it dvr}である、
\label{algebra-item-ring-dvr}
\item $K$ は{\it 体}である、
\label{algebra-item-field}
\item $L/K$ は{\it 体拡大}である、
\label{algebra-item-field-extension}
\item $L/K$ は{\it 代数的体拡大}である、
\label{algebra-item-field-extension-algebraic}
\item $\{t_i\}_{i\in I}$ は $K$ 上の $L$ の{\it 超越基底}である、
\label{algebra-item-transcendence-basis}
\item $K$ 上の $L$ の{\it 超越次数} $\text{trdeg}(L/K)$、
\label{algebra-item-transcendence-degree}
\item 体 $k$ は{\it 代数閉}である、
\label{algebra-item-algebraically-closed}
\item
\label{algebra-item-extend-into-algebraically-closed}
$L/K$ が代数的で、$\Omega/K$ が $\Omega$ を代数閉とする拡大ならば、
$K$ 上の写像を延長する環準同型 $L \to \Omega$ が存在する、
\item $I \subset R$ は{\it イデアル}である、
\label{algebra-item-ideal}
\item $I \subset R$ は{\it 根基的}である、
\label{algebra-item-ideal-radical}
\item $I$ がイデアルならば、その{\it 根基} $\sqrt{I}$ がある、
\label{algebra-item-radical-ideal}
\item
\label{algebra-item-ideal-nilpotent}
$I \subset R$ が{\it 冪零}であるとは、ある $n \in \mathbf{N}$ に対して
$I^n = 0$ であることをいう、
\item
\label{algebra-item-ideal-locally-nilpotent}
$I \subset R$ が{\it 局所冪零}であるとは、$I$ のすべての元が冪零であることをいう、
\item $\mathfrak p \subset R$ は{\it 素イデアル}である、
\label{algebra-item-prime-ideal}
\item
\label{algebra-item-prime-product-ideals}
$\mathfrak p \subset R$ が素イデアルで、$I, J \subset R$ がイデアルであり、
$IJ\subset \mathfrak p$ ならば、$I \subset \mathfrak p$ または
$J \subset \mathfrak p$ である、
\item $\mathfrak m \subset R$ は{\it 極大イデアル}である、
\label{algebra-item-maximal-ideal}
\item 任意の非零環は極大イデアルをもつ、
\label{algebra-item-exists-maximal-ideal}
\item
\label{algebra-item-jacobson-radical}
$R$ の{\it ジャコブソン根基}とは、$R$ のすべての極大イデアルの共通部分
$\text{rad}(R) = \bigcap_{\mathfrak m \subset R} \mathfrak m$ である、
\item 部分集合 $T \subset R$ によって{\it 生成される}イデアル $(T)$、
\label{algebra-item-ideal-generated-by}
\item {\it 商環} $R/I$、
\label{algebra-item-quotient-ring}
\item 環 $R$ のイデアル $I$ が素であるための必要十分条件は、$R/I$ が整域で
あることである、
\label{algebra-item-characterize-prime-ideal}
\item
\label{algebra-item-characterize-maximal-ideal}
環 $R$ のイデアル $I$ が極大であるための必要十分条件は、環 $R/I$ が体で
あることである、
\item
\label{algebra-item-inverse-image-ideal}
$\varphi : R_1 \to R_2$ が環準同型で、$I \subset R_2$ がイデアルならば、
$\varphi^{-1}(I)$ は $R_1$ のイデアルである、
\item
\label{algebra-item-image-ideal}
$\varphi : R_1 \to R_2$ が環準同型で、$I \subset R_1$ がイデアルならば、
$\varphi(I) \cdot R_2$（$I \cdot R_2$ または $IR_2$ と記すこともある）は
$\varphi(I)$ によって生成される $R_2$ のイデアルである、
\item
\label{algebra-item-inverse-image-prime}
$\varphi : R_1 \to R_2$ が環準同型で、$\mathfrak p \subset R_2$ が素イデアル
ならば、$\varphi^{-1}(\mathfrak p)$ は $R_1$ の素イデアルである、
\item $M$ は{\it $R$-加群}である、
\label{algebra-item-module}
\item
\label{algebra-item-annihilator}
$m \in M$ に対し、$m$ の $R$ における{\it 零化イデアル}は
$I = \{f \in R \mid fm = 0\}$ である、
\item $N \subset M$ は{\it $R$-部分加群}である、
\label{algebra-item-submodule}
\item $M$ は{\it ネーター的 $R$-加群}である、
\label{algebra-item-Noetherian-module}
\item $M$ は{\it 有限 $R$-加群}である、
\label{algebra-item-finite-module}
\item $M$ は{\it 有限生成 $R$-加群}である、
\label{algebra-item-finitely-generated-module}
\item $M$ は{\it 有限表示 $R$-加群}である、
\label{algebra-item-finitely-presented-module}
\item $M$ は{\it 自由 $R$-加群}である、
\label{algebra-item-free-module}
\item
\label{algebra-item-extension-free}
$0 \to K \to L \to M \to 0$ が $R$-加群の短完全列で、$K$, $M$ が自由ならば、
$L$ も自由である、
\item $N \subset M \subset L$ が $R$-加群ならば、$L/M = (L/N)/(M/N)$ である、
\label{algebra-item-isomorphism-theorem}
\item $S$ は{\it $R$ の乗法的集合}である、
\label{algebra-item-multiplicative-subset}
\item $R$ の{\it 局所化} $R \to S^{-1}R$、
\label{algebra-item-localization-ring}
\item
\label{algebra-item-localization-zero}
$R$ が環で $S$ が $R$ の乗法的集合ならば、$S^{-1}R$ が零環であるための
必要十分条件は $S$ が $0$ を含むことである、
\item
\label{algebra-item-localize-nonzerodivisors}
$R$ が環で、乗法的集合 $S$ の元がすべて非零因子ならば、$R \to S^{-1}R$
は単射である、
\item $\varphi : R_1 \to R_2$ が環準同型で、$S$ が $R_1$ の乗法的集合ならば、
$\varphi(S)$ は $R_2$ の乗法的集合である、
\item
\label{algebra-item-products-multiplicative-subsets}
$S$, $S'$ が $R$ の乗法的集合であり、$SS'$ が積の集合
$SS' = \{r \in R \mid \exists s\in S, \exists s' \in S', r = ss'\}$ を表すならば、
$SS'$ は $R$ の乗法的集合である、
\item
\label{algebra-item-localization-localization}
$S$, $S'$ が $R$ の乗法的集合で、$\overline{S}$ が $(S')^{-1}R$ における
$S$ の像を表すならば、$(SS')^{-1}R = \overline{S}^{-1}((S')^{-1}R)$ である、
\item $R$-加群 $M$ の{\it 局所化} $S^{-1}M$、
\label{algebra-item-localization-module}
\item
\label{algebra-item-localization-exact}
関手 $M \mapsto S^{-1}M$ は単射、全射、および完全性を保つ、
\item
\label{algebra-item-localization-localization-module}
$S$, $S'$ が $R$ の乗法的集合で、$M$ が $R$-加群ならば、
$(SS')^{-1}M = S^{-1}((S')^{-1}M)$ である、
\item
\label{algebra-item-localize-ideal}
$R$ が環、$I$ が $R$ のイデアル、$S$ が $R$ の乗法的集合ならば、
$S^{-1}I$ は $S^{-1}R$ のイデアルであり、
$S^{-1}R/S^{-1}I = \overline{S}^{-1}(R/I)$ が成り立つ。ここで
$\overline{S}$ は $R/I$ における $S$ の像である、
\item
\label{algebra-item-ideal-in-localization}
$R$ が環で $S$ が $R$ の乗法的集合ならば、$S^{-1}R$ の任意のイデアル
$I'$ は $S^{-1}I$ の形である。ここで $I$ として、$R$ における $I'$ の逆像を
とることができる、
\item
\label{algebra-item-submodule-in-localization}
$R$ が環、$M$ が $R$-加群、$S$ が $R$ の乗法的集合ならば、$S^{-1}M$ の
任意の部分加群 $N'$ は、ある部分加群 $N \subset M$ に対する $S^{-1}N$ の形で
ある。ここで $N$ として、$M$ における $N'$ の逆像をとることができる、
\item $S = \{1, f, f^2, \ldots\}$ ならば、$R_f = S^{-1}R$ および
$M_f = S^{-1}M$ である、
\label{algebra-item-localize-f}
\item
\label{algebra-item-localize-p}
ある素イデアル $\mathfrak p$ に対して
$S = R \setminus \mathfrak p = \{x\in R \mid x\not\in \mathfrak p\}$ ならば、
$R_{\mathfrak p} = S^{-1}R$ および $M_{\mathfrak p} = S^{-1}M$ と記すのが慣例である、
\item {\it 局所環}とは極大イデアルをちょうど一つもつ環である、
\label{algebra-item-local-ring}
\item {\it 半局所環}とは極大イデアルを有限個もつ環である、
\label{algebra-item-semi-local-ring}
\item
\label{algebra-item-localize-p-local-ring}
$\mathfrak p$ が $R$ の素イデアルならば、$R_{\mathfrak p}$ は極大イデアル
$\mathfrak p R_{\mathfrak p}$ をもつ局所環である、
\item
\label{algebra-item-residue-field}
環 $R$ の素イデアル $\mathfrak p$ の{\it 剰余体} $\kappa(\mathfrak p)$ とは、整域
$R/\mathfrak p$ の分数体である。これは
$R_\mathfrak p/\mathfrak pR_\mathfrak p = (R \setminus \mathfrak p)^{-1}R/\mathfrak p$
に等しい、
\item $R$ と $M_1$, $M_2$ が与えられたときの{\it テンソル積}
$M_1 \otimes_R M_2$、
\label{algebra-item-tensor-product}
\item
\label{algebra-item-cauchy-binet}
環 $R$ 上のそれぞれ $m \times n$ および $n \times m$ 行列 $A$, $B$ に対し、
$R$ において $\det(AB) = \sum \det(A_S)\det({}_SB)$ が成り立つ。ここで和は
大きさ $m$ の部分集合 $S \subset \{1, \ldots, n\}$ 全体を走り、$A_S$ は
$S$ に対応する列からなる $A$ の $m \times m$ 部分行列、${}_SB$ は $S$ に
対応する行からなる $B$ の $m \times m$ 部分行列である、
\item など。
\end{enumerate}




\section{蛇の補題}
\label{algebra-section-snake}

\noindent
蛇の補題とその変種は、アーベル圏の設定において
Homology, Section \ref{homology-section-abelian-categories} で論じる。

\begin{lemma}
\label{algebra-lemma-snake}
\begin{reference}
\cite[III, Lemma 3.3]{Cartan-Eilenberg}
\end{reference}
行が完全であるアーベル群の可換図式
$$
\xymatrix{
& X \ar[r] \ar[d]^\alpha &
Y \ar[r] \ar[d]^\beta &
Z \ar[r] \ar[d]^\gamma &
0 \\
0 \ar[r] & U \ar[r] & V \ar[r] & W
}
$$
が与えられたとき、標準的な完全列
$$
\Ker(\alpha) \to \Ker(\beta) \to \Ker(\gamma)
\to
\Coker(\alpha) \to \Coker(\beta) \to \Coker(\gamma)
$$
が存在する。さらに、$X \to Y$ が単射ならば最初の写像は単射であり、
$V \to W$ が全射ならば最後の写像は全射である。
\end{lemma}

\begin{proof}
写像 $\partial : \Ker(\gamma) \to \Coker(\alpha)$ を次のように定める。
$z \in \Ker(\gamma)$ をとり、$z$ に写る $y \in Y$ を選ぶ。このとき
$\beta(y) \in V$ は $W$ の零元に写る。したがって $\beta(y)$ はある
$u \in U$ の像である。$\partial z = \overline{u}$、すなわち $\alpha$ の余核に
おける $u$ の類と定める。完全性の証明は省略する。
\end{proof}
\section{有限加群と有限表示加群}
\label{algebra-section-module-finite-type}

\noindent
ここでは基本的な記法と補題をいくつか述べる。

\begin{definition}
\label{algebra-definition-module-finite-type}
$R$ を環とし、$M$ を $R$-加群とする。
\begin{enumerate}
\item ある $n \in \mathbf{N}$ と $x_1, \ldots, x_n \in M$ が存在し、$M$ の
すべての元が $x_i$ の $R$-線形結合であるとき、$M$ を{\it 有限 $R$-加群}
または{\it 有限生成 $R$-加群}という。同値な言い換えとして、ある
$n \in \mathbf{N}$ に対して全射 $R^{\oplus n} \to M$ が存在する。
\item ある整数 $n, m \in \mathbf{N}$ と完全列
$$
R^{\oplus m} \longrightarrow R^{\oplus n} \longrightarrow M \longrightarrow 0
$$
が存在するとき、$M$ を{\it 有限表示 $R$-加群}または{\it 有限表示の
$R$-加群}という。
\end{enumerate}
\end{definition}

\noindent
非形式的にいえば、$M$ が有限表示 $R$-加群であるための必要十分条件は、
それが有限生成であり、これらの生成元の間の関係式の加群も有限生成である
ことである。定義に現れる完全列の選択を $M$ の{\it 表示}という。

\begin{lemma}
\label{algebra-lemma-lift-map}
$R$ を環とする。$\alpha : R^{\oplus n} \to M$ と $\beta : N \to M$ を
加群準同型とする。$\Im(\alpha) \subset \Im(\beta)$ ならば、
$\alpha = \beta \circ \gamma$ を満たす $R$-加群準同型
$\gamma : R^{\oplus n} \to N$ が存在する。
\end{lemma}

\begin{proof}
$e_i = (0, \ldots, 0, 1, 0, \ldots, 0)$ を $R^{\oplus n}$ の第 $i$ 基底ベクトル
とする。仮定により、$\alpha(e_i) = \beta(x_i)$ を満たす元 $x_i \in N$ が存在する。
$\gamma(a_1, \ldots, a_n) = \sum a_i x_i$ と定める。構成から
$\alpha = \beta \circ \gamma$ である。
\end{proof}

\begin{lemma}
\label{algebra-lemma-extension}
$R$ を環とし、
$$
0 \to M_1 \to M_2 \to M_3 \to 0
$$
を $R$-加群の短完全列とする。
\begin{enumerate}
\item $M_1$ と $M_3$ が有限 $R$-加群ならば、$M_2$ は有限 $R$-加群である。
\item $M_1$ と $M_3$ が有限表示 $R$-加群ならば、$M_2$ は有限表示
$R$-加群である。
\item $M_2$ が有限 $R$-加群ならば、$M_3$ は有限 $R$-加群である。
\item $M_2$ が有限表示 $R$-加群で、$M_1$ が有限 $R$-加群ならば、$M_3$ は
有限表示 $R$-加群である。
\item $M_3$ が有限表示 $R$-加群で、$M_2$ が有限 $R$-加群ならば、$M_1$ は
有限 $R$-加群である。
\end{enumerate}
\end{lemma}

\begin{proof}
(1) を証明する。$x_1, \ldots, x_n$ を $M_1$ の生成元とし、
$y_1, \ldots, y_m \in M_2$ を、その $M_3$ における像が $M_3$ を生成する元とする。
このとき $x_1, \ldots, x_n, y_1, \ldots, y_m$ は $M_2$ を生成する。

\medskip\noindent
(3) は定義から直ちに従う。

\medskip\noindent
(5) を証明する。$M_3$ は有限表示、$M_2$ は有限であると仮定し、表示
$$
R^{\oplus m} \to R^{\oplus n} \to M_3 \to 0
$$
を選ぶ。補題 \ref{algebra-lemma-lift-map} により、実線からなる図式
$$
\xymatrix{
& R^{\oplus m} \ar[r] \ar@{..>}[d] & R^{\oplus n} \ar[r] \ar[d] &
M_3 \ar[r] \ar[d]^{\text{id}} & 0 \\
0 \ar[r] & M_1 \ar[r] & M_2 \ar[r] & M_3 \ar[r] & 0
}
$$
が可換となるような写像 $R^{\oplus n} \to M_2$ が存在する。これにより点線の
矢印が得られる。蛇の補題（補題 \ref{algebra-lemma-snake}）から同型
$$
\Coker(R^{\oplus m} \to M_1)
\cong
\Coker(R^{\oplus n} \to M_2)
$$
を得る。特に $\Coker(R^{\oplus m} \to M_1)$ は有限 $R$-加群である。
$\Im(R^{\oplus m} \to M_1)$ は (3) により有限なので、(1) から $M_1$ は有限である。

\medskip\noindent
(4) を証明する。$M_2$ は有限表示、$M_1$ は有限であると仮定する。表示
$R^{\oplus m} \to R^{\oplus n} \to M_2 \to 0$ を選び、全射
$R^{\oplus k} \to M_1$ を選ぶ。補題 \ref{algebra-lemma-lift-map} により、合成
$R^{\oplus k} \to M_1 \to M_2$ は $R^{\oplus k} \to R^{\oplus n} \to M_2$
と分解する。このとき $R^{\oplus k + m} \to R^{\oplus n} \to M_3 \to 0$
は表示である。

\medskip\noindent
(2) を証明する。$M_1$ と $M_3$ が有限表示であると仮定する。(1) の証明の
議論により、全射な垂直矢印をもつ可換図式
$$
\xymatrix{
0 \ar[r] & R^{\oplus n} \ar[d] \ar[r] & R^{\oplus n + m} \ar[d] \ar[r] &
R^{\oplus m} \ar[d] \ar[r] & 0 \\
0 \ar[r] & M_1 \ar[r] & M_2 \ar[r] & M_3 \ar[r] & 0
}
$$
が得られる。蛇の補題により短完全列
$$
0 \to \Ker(R^{\oplus n} \to M_1) \to
\Ker(R^{\oplus n + m} \to M_2) \to
\Ker(R^{\oplus m} \to M_3) \to 0
$$
を得る。(5) により両端の加群は有限である。したがって中央の加群も有限である。
(4) により $M_2$ は有限表示である。
\end{proof}

\begin{lemma}
\label{algebra-lemma-trivial-filter-finite-module}
\begin{slogan}
有限加群は、逐次商が巡回加群となるフィルトレーションをもつ。
\end{slogan}
$R$ を環、$M$ を有限 $R$-加群とする。有限 $R$-部分加群による
フィルトレーション
$$
0 = M_0 \subset M_1 \subset \ldots \subset M_n = M
$$
であって、各商 $M_i/M_{i - 1}$ が $R$ のあるイデアル $I_i$ に対する
$R/I_i$ と同型になるものが存在する。
\end{lemma}

\begin{proof}
$M$ の生成元の個数に関する帰納法を用いる。$x_1, \ldots, x_r \in M$ を
生成元とする。$M' = Rx_1 \subset M$ とおけば、$M/M'$ は $r - 1$ 個の生成元を
もち、帰納法の仮定を適用できる。また明らかに
$I_1 = \{f \in R \mid fx_1 = 0\}$ とおけば $M' \cong R/I_1$ である。
\end{proof}

\begin{lemma}
\label{algebra-lemma-finite-over-subring}
$R \to S$ を環準同型とし、$M$ を $S$-加群とする。$M$ が $R$-加群として
有限ならば、$M$ は $S$-加群として有限である。
\end{lemma}

\begin{proof}
実際、$M$ の任意の $R$-生成系は $M$ の $S$-生成系でもある。これは
$R$-加群構造が $S$ における $R$ の像から誘導されるためである。
\end{proof}



\section{有限型および有限表示の環準同型}
\label{algebra-section-finite-type}

\begin{definition}
\label{algebra-definition-finite-type}
$R \to S$ を環準同型とする。
\begin{enumerate}
% P01-E1272
\item ある $n \in \mathbf{N}$ と $R$-代数の全射
$R[x_1, \ldots, x_n] \to S$ が存在するとき、$R \to S$ は{\it 有限型}である、
または{\it $S$ は有限型 $R$-代数である}という。
\item ある整数 $n, m \in \mathbf{N}$ と多項式
$f_1, \ldots,\allowbreak f_m \in R[x_1, \ldots,\allowbreak x_n]$、および $R$-代数の同型
$R[x_1, \ldots,\allowbreak x_n]/(f_1, \ldots,\allowbreak f_m)\allowbreak \cong S$ が存在するとき、
$R \to S$ は{\it 有限表示}であるという。
\end{enumerate}
\end{definition}

\noindent
非形式的にいえば、$R \to S$ が有限表示であるための必要十分条件は、$S$ が
有限生成 $R$-代数であり、生成元の間の関係式のイデアルも有限生成である
ことである。定義に現れる全射 $R[x_1, \ldots, x_n] \to S$ の選択を
$S$ の{\it 表示}ということがある。

\begin{lemma}
\label{algebra-lemma-compose-finite-type}
有限型および有限表示という概念は、次の安定性をもつ。
\begin{enumerate}
\item 有限型の環準同型の合成は有限型である。
\item 有限表示の環準同型の合成は有限表示である。
\item $R \to S' \to S$ が与えられ、$R \to S$ が有限型ならば、
$S' \to S$ は有限型である。
\item $R \to S' \to S$ が与えられ、$R \to S$ が有限表示で、
$R \to S'$ が有限型ならば、$S' \to S$ は有限表示である。
\end{enumerate}
\end{lemma}

\begin{proof}
最後の主張だけを証明する。
$S = R[x_1, \ldots, x_n]/(f_1, \ldots, f_m)$ および
$S' = R[y_1, \ldots, y_a]/I$ と書く。$y_i$ の類 $\bar y_i$ が $S$ において
$h_i \bmod (f_1, \ldots, f_m)$ に写るとする。このとき明らかに
$S = S'[x_1, \ldots, x_n]/(f_1, \ldots, f_m,
h_1 - \bar y_1, \ldots, h_a - \bar y_a)$ である。
\end{proof}

\begin{lemma}
\label{algebra-lemma-finite-presentation-independent}
$R \to S$ を有限表示の環準同型とする。任意の全射
$\alpha : R[x_1, \ldots, x_n] \to S$ に対し、$\alpha$ の核は
$R[x_1, \ldots, x_n]$ の有限生成イデアルである。
\end{lemma}

\begin{proof}
$S = R[y_1, \ldots, y_m]/(f_1, \ldots, f_k)$ と書く。
$\alpha(x_i)$ の持ち上げとなる $g_i \in R[y_1, \ldots, y_m]$ を選ぶ。このとき
$S = R[x_i, y_j]/(f_l, x_i - g_i)$ である。$\alpha(h_j)$ が
$y_j \bmod (f_1, \ldots, f_k)$ に対応するような
$h_j \in R[x_1, \ldots, x_n]$ を選ぶ。写像
$\psi : R[x_i, y_j] \to R[x_i]$ を $x_i \mapsto x_i$, $y_j \mapsto h_j$ で
定める。このとき $\alpha$ の核は $\psi$ による $(f_l, x_i - g_i)$ の像であり、
主張が従う。
\end{proof}

\begin{lemma}
\label{algebra-lemma-finitely-presented-over-subring}
$R \to S$ を環準同型とし、$M$ を $S$-加群とする。$R \to S$ が有限型で、
$M$ が $R$-加群として有限表示ならば、$M$ は $S$-加群として有限表示である。
\end{lemma}

\begin{proof}
補題 \ref{algebra-lemma-compose-finite-type} の (4) の証明と同様である。
$S = R[x_1, \ldots, x_n]/J$ と仮定してよい。$M$ を $R$-加群として生成する
$y_1, \ldots, y_m \in M$ を選び、$R^{\oplus m} \to M$ の核を生成する関係式
$\sum a_{ij} y_j = 0$, $i = 1, \ldots, t$ を選ぶ。任意の
$i = 1, \ldots, n$ と $j = 1, \ldots, m$ に対して
$$
x_i y_j = \sum a_{ijk} y_k
$$
と書く。ただし $a_{ijk} \in R$ である。$y_1, \ldots, y_m$ によって生成され、
関係式 $\sum a_{ij} y_j = 0$, $i = 1, \ldots, t$ および
$x_i y_j = \sum a_{ijk} y_k$, $i = 1, \ldots, n$, $j = 1, \ldots, m$ を課した
$S$-加群を $N$ とする。このとき $N$ は表示
$$
S^{\oplus nm + t} \longrightarrow S^{\oplus m} \longrightarrow N
\longrightarrow 0
$$
をもつ。構成により全射 $\varphi : N \to M$ が存在する。証明を終えるには
$\varphi$ が単射であることを示せばよい。ある $b_j \in S$ に対し
$z = \sum b_j y_j \in N$ とする。$b_j$ を $R$ 係数の
$x_1, \ldots, x_n$ の多項式とみなしてよい。$x_i y_j = \sum a_{ijk} y_k$ の
形の関係式を適用すれば、多項式の次数を帰納的に下げられる。したがって
ある $c_j \in R$ に対して $z = \sum c_j y_j$ と書ける。よって
$\varphi(z) = 0$ ならば、ベクトル $(c_1, \ldots, c_m)$ はベクトル
$(a_{i1}, \ldots, a_{im})$ の $R$-線形結合である。したがって所望のとおり
$z = 0$ である。
\end{proof}









\section{有限環準同型}
\label{algebra-section-finite}

\noindent
まず定義を述べる。

\begin{definition}
\label{algebra-definition-finite-ring-map}
$\varphi : R \to S$ を環準同型とする。$S$ が $R$-加群として有限であるとき、
$\varphi : R \to S$ は{\it 有限}であるという。
\end{definition}

\begin{lemma}
\label{algebra-lemma-finite-module-over-finite-extension}
$R \to S$ を有限環準同型とし、$M$ を $S$-加群とする。$M$ が $R$-加群として
有限であるための必要十分条件は、$M$ が $S$-加群として有限であることである。
\end{lemma}

\begin{proof}
一方の含意は補題 \ref{algebra-lemma-finite-over-subring} から従う。逆を示すため、$M$ が
$S$-加群として有限であると仮定する。$R$-加群として $S$ を生成する
$x_1, \ldots, x_n \in S$ を選び、$S$-加群として $M$ を生成する
$y_1, \ldots, y_m \in M$ を選ぶ。このとき $x_i y_j$ は $R$-加群として
$M$ を生成する。
\end{proof}

\begin{lemma}
\label{algebra-lemma-finite-transitive}
$R \to S$ と $S \to T$ が有限環準同型ならば、$R \to T$ は有限である。
\end{lemma}

\begin{proof}
$t_i$ が $S$-加群として $T$ を生成し、$s_j$ が $R$-加群として $S$ を
生成するならば、$t_i s_j$ は $R$-加群として $T$ を生成する。
（補題 \ref{algebra-lemma-finite-module-over-finite-extension} からも従う。）
\end{proof}

\begin{lemma}
\label{algebra-lemma-finite-finite-type}
$\varphi : R \to S$ を環準同型とする。
\begin{enumerate}
\item $\varphi$ が有限ならば、$\varphi$ は有限型である。
\item $S$ が $R$-加群として有限表示ならば、$\varphi$ は有限表示である。
\end{enumerate}
\end{lemma}

\begin{proof}
(1) について、$x_1, \ldots, x_n \in S$ が $R$-加群として $S$ を生成するならば、
$x_1, \ldots, x_n$ は $R$-代数として $S$ を生成する。(2) について、$x_i$ を
$S$ の $R$-加群生成元とみたとき、その関係式の生成系を
$\sum r_j^ix_i = 0$, $j = 1, \ldots, m$ とする。さらに、ある $r_i \in R$ に
対して $1 = \sum r_ix_i$ と書き、ある $r_{ij}^k \in R$ に対して
$x_ix_j = \sum r_{ij}^k x_k$ と書く。このとき $R$-代数として
$$
S = R[t_1, \ldots, t_n]/
(\sum r_j^it_i,\ 1 - \sum r_it_i,\ t_it_j - \sum r_{ij}^k t_k)
$$
であり、(2) が従う。
\end{proof}

\noindent
有限環準同型の詳細については、Section \ref{algebra-section-finite-ring-extensions} を参照せよ。


\section{余極限}
\label{algebra-section-colimits}

\noindent
本節の内容の一部は、Categories, Sections \ref{categories-section-limits} --
\ref{categories-section-posets-limits} における余極限の一般論と重複する。
前順序集合の概念は Categories, Definition
\ref{categories-definition-directed-set} で定義する。これは半順序集合より
わずかに弱い概念である。

\begin{definition}
\label{algebra-definition-directed-system}
$(I, \leq)$ を前順序集合とする。{\it $I$ 上の $R$-加群の系
$(M_i, \mu_{ij})$}とは、$I$ で添字づけられた $R$-加群の族
$\{M_i\}_{i\in I}$ と、比較可能な各添字対に対する $R$-加群準同型の族
$\{\mu_{ij} : M_i \to M_j\}_{i \leq j}$ であって、すべての
$i \leq j \leq k$ に対して
$$
\mu_{ii} = \text{id}_{M_i}\quad
\mu_{ik} = \mu_{jk}\circ \mu_{ij}
$$
を満たすものをいう。$I$ が有向集合であるとき、$(M_i, \mu_{ij})$ を
{\it 有向系}という。
\end{definition}

\noindent
これは Categories, Definition \ref{categories-definition-system-over-poset} および
Section \ref{categories-section-posets-limits} で定義した概念と同じである。
任意の圏における図式または系の余極限の定義については、Categories,
Definition \ref{categories-definition-colimit} を参照せよ。

\begin{lemma}
\label{algebra-lemma-colimit}
$(M_i, \mu_{ij})$ を前順序集合 $I$ 上の $R$-加群の系とする。この系
$(M_i, \mu_{ij})$ の余極限は
商 $R$-加群 $(\bigoplus_{i\in I} M_i) /Q$ である。ここで $Q$ は、すべての元
$$
\iota_i(x_i) - \iota_j(\mu_{ij}(x_i))
$$
によって生成される $R$-部分加群であり、
$\iota_i : M_i \to \bigoplus_{i\in I} M_i$ は自然な包含である。余極限を
$M = \colim_i M_i$ と記す。射影を
$\pi : \bigoplus_{i\in I} M_i \to M$ と記し、
$\phi_i = \pi \circ \iota_i : M_i \to M$ とおく。
\end{lemma}

\begin{proof}
この補題は Categories, Lemma
\ref{categories-lemma-colimits-coproducts-coequalizers} の特別な場合であるが、
ここでは直接証明も与える。まず上の構成において
$\phi_i = \phi_j\circ \mu_{ij}$ である。対 $(M, \phi_i)$ が余極限であることを
示すには、普遍性を満たすことを示せばよい。すなわち、同じ性質をもつ任意の対
$(Y, \psi_i)$、ただし $\psi_i : M_i \to Y$ かつ
$\psi_i = \psi_j\circ \mu_{ij}$、に対して、次の図式を可換にする一意な
$R$-加群準同型 $g : M \to Y$ が存在することを示せばよい。
$$
\xymatrix{
M_i \ar[rr]^{\mu_{ij}} \ar[dr]^{\phi_i} \ar[ddr]_{\psi_i} & &
M_j\ar[dl]_{\phi_j} \ar[ddl]^{\psi_j} \\
& M \ar[d]^{g}\\
& Y
}
$$
これは、直和 $\bigoplus M_i$ の直和因子 $M_i$ 上で $\psi_i$ とすることにより
$g$ を定義できるので明らかである。
\end{proof}

\begin{lemma}
\label{algebra-lemma-directed-colimit}
$(M_i, \mu_{ij})$ を前順序集合 $I$ 上の $R$-加群の系とし、$I$ は有向であると
仮定する。この系 $(M_i, \mu_{ij})$ の余極限は、次のように定義される加群 $M$ と
標準的に同型である。
\begin{enumerate}
\item 集合として
$$
M = \left(\coprod\nolimits_{i \in I} M_i\right)/\sim
$$
とする。ここで $m \in M_i$ および $m' \in M_{i'}$ に対し、
$$
m \sim m' \Leftrightarrow
\mu_{ij}(m) = \mu_{i'j}(m')\text{ for some }j \geq i, i'
$$
と定める。
\item アーベル群として、$m \in M_i$ および $m' \in M_{i'}$ に対し、$M$ に
おける $m$ と $m'$ の類の和を $\mu_{ij}(m) + \mu_{i'j}(m')$ の類と定める。
ここで $j \in I$ は $i \leq j$ および $i' \leq j$ を満たす任意の添字である。
\item $R$-加群として、$m \in M_i$ および $x \in R$ に対し、$x$ と $M$ に
おける $m$ の類との積を $M$ における $xm$ の類と定める。
\end{enumerate}
標準写像 $\phi_i : M_i \to M$ は標準写像
$M_i \to \coprod_{i \in I} M_i$ から誘導される。
\end{lemma}

\begin{proof}
省略する。Categories, Section \ref{categories-section-directed-colimits} と比較せよ。
\end{proof}

\begin{lemma}
\label{algebra-lemma-zero-directed-limit}
$(M_i, \mu_{ij})$ を有向系とする。$M = \colim M_i$ とし、
$\mu_i : M_i \to M$ とする。このとき $x_i \in M_i$ に対して
$\mu_i(x_i) = 0$ であるための必要十分条件は、$\mu_{ij}(x_i) = 0$ を満たす
$j \geq i$ が存在することである。
\end{lemma}

\begin{proof}
補題 \ref{algebra-lemma-directed-colimit} における有向余極限の記述から明らかである。
\end{proof}

\begin{example}
\label{algebra-example-zero-colimit-different}
$a < b$ および $a < c$ のほかに真の不等式をもたない半順序集合
$I = \{a, b, c\}$ を考える。$I$ 上の系
$(M_a, M_b, M_c, \mu_{ab}, \mu_{ac})$ は、三つの $R$-加群 $M_a, M_b, M_c$ と、
二つの $R$-加群準同型 $\mu_{ab} : M_a \to M_b$ および
$\mu_{ac} : M_a \to M_c$ からなる。この系の余極限は
$$
M := \colim_{i \in I} M_i = \Coker(M_a \to M_b \oplus M_c)
$$
にほかならず、この写像は $\mu_{ab} \oplus -\mu_{ac}$ である。したがって標準写像
$M_a \to M$ の核は $\Ker(\mu_{ab}) + \Ker(\mu_{ac})$ である。また標準写像
$M_b \to M$ の核は、$\mu_{ab}$ による $\Ker(\mu_{ac})$ の像である。ゆえに
補題 \ref{algebra-lemma-zero-directed-limit} の結論は一般の系については偽である。
\end{example}

\begin{definition}
\label{algebra-definition-homomorphism-directed-systems}
$(M_i, \mu_{ij})$, $(N_i, \nu_{ij})$ を同じ前順序集合 $I$ 上の $R$-加群の系とする。
$(M_i, \mu_{ij})$ から $(N_i, \nu_{ij})$ への{\it 系の準同型} $\Phi$ とは、
すべての $i \leq j$ に対して $\phi_j \circ \mu_{ij} = \nu_{ij} \circ \phi_i$ を
満たす $R$-加群準同型の族 $\phi_i : M_i \to N_i$ のことである。
\end{definition}

\noindent
圏の言葉では、これは付随する図式 $M : I \to \text{Mod}_R$ と
$N : I \to \text{Mod}_R$ の間の関手の変換と同じ概念である。次の補題は
Categories, Lemma \ref{categories-lemma-functorial-colimit} の特別な場合である。

\begin{lemma}
\label{algebra-lemma-homomorphism-limit}
$(M_i, \mu_{ij})$, $(N_i, \nu_{ij})$ を同じ前順序集合上の $R$-加群の系とする。
$(M_i, \mu_{ij})$ から $(N_i, \nu_{ij})$ への系の射
$\Phi = (\phi_i)$ は、一意な準同型
$$
\colim \phi_i : \colim M_i \longrightarrow \colim N_i
$$
であって、すべての $i \in I$ に対して図式
$$
\xymatrix{
M_i \ar[r] \ar[d]_{\phi_i} & \colim M_i \ar[d]^{\colim \phi_i} \\
N_i \ar[r] & \colim N_i
}
$$
を可換にするものを誘導する。
\end{lemma}

\begin{proof}
$M = \colim M_i$, $N = \colim N_i$, $\phi = \colim \phi_i$ と書く（この写像は
これから構成する）。以下では補題 \ref{algebra-lemma-colimit} における $M$ と $N$ の
具体的記述を断りなく用いる。補題の条件は、図式
$$
\xymatrix{
\bigoplus_{i\in I} M_i \ar[r] \ar[d]_{\bigoplus\phi_i} & M \ar[d]^\phi \\
\bigoplus_{i\in I} N_i \ar[r] & N
}
$$
が可換であるという条件と同値である。したがって $\phi$ が存在すれば一意で
あることは明らかである。$\phi$ の存在を示すには、上の水平矢印の核が
$\bigoplus \phi_i$ によって下の水平矢印の核へ写ることを示せば十分である。
このため $j \leq k$ および $x_j \in M_j$ とすると、
$$
(\bigoplus \phi_i)(x_j - \mu_{jk}(x_j))
=
\phi_j(x_j) - \phi_k(\mu_{jk}(x_j))
=
\phi_j(x_j) - \nu_{jk}(\phi_j(x_j))
$$
であり、これは所望のとおり下の水平矢印の核に属する。
\end{proof}

\begin{lemma}
\label{algebra-lemma-directed-colimit-exact}
\begin{slogan}
フィルター余極限は完全である。有向余極限は完全である。
\end{slogan}
$I$ を有向集合とする。$(L_i, \lambda_{ij})$, $(M_i, \mu_{ij})$ および
$(N_i, \nu_{ij})$ を $I$ 上の $R$-加群の系とする。
$\varphi_i : L_i \to M_i$ および $\psi_i : M_i \to N_i$ を $I$ 上の系の射とする。
すべての $i \in I$ に対し、$R$-加群の列
$$
\xymatrix{
L_i \ar[r]^{\varphi_i} &
M_i \ar[r]^{\psi_i} &
N_i
}
$$
がホモロジー $H_i$ をもつ複体であると仮定する。このとき $R$-加群 $H_i$ は
$I$ 上の系をなし、$R$-加群の列
$$
\xymatrix{
\colim_i L_i \ar[r]^\varphi &
\colim_i M_i \ar[r]^\psi &
\colim_i N_i
}
$$
も複体である。そのホモロジーを $H$ と記すと、
$$
H = \colim_i H_i.
$$
である。
\end{lemma}

\begin{proof}
$
\xymatrix{
\colim_i L_i \ar[r]^\varphi &
\colim_i M_i \ar[r]^\psi &
\colim_i N_i
}
$
が複体であることは明らかである。各 $i \in I$ に対して標準的な
$R$-加群準同型 $H_i \to H$ がある。これは各
$[m] \in H_i = \Ker(\psi_i) / \Im(\varphi_i)$ を、$\colim_i M_i$ における
$m$ の像の $H = \Ker(\psi) / \Im(\varphi)$ における剰余類へ送る。この族から
準同型 $\colim_i H_i \to H$ が得られる。これが全射かつ単射であることを示せばよい。

\medskip\noindent
以下では補題 \ref{algebra-lemma-directed-colimit} による $I$ 上の余極限の記述を繰り返し
用いる。$h \in H$ とする。$H = \Ker(\psi)/\Im(\varphi)$ なので、$h$ は
$\Ker(\psi) \subset \colim_i M_i$ の元 $[m]$ の $\Im(\varphi)$ を法とする類である。
$[m]$ がある元 $m \in M_i$ から来るような $i$ を選ぶ。
$[m] \in \Ker(\psi)$ なので、$\nu_{ij}(\psi_i(m)) = 0$ を満たす $j \geq i$ を
選べる。$i$ を $j$ に、$m$ を $\mu_{ij}(m)$ に置き換えると、
$m \in \Ker(\psi_i)$ と仮定してよい。これにより写像
$\colim_i H_i \to H$ が全射であることが分かる。

\medskip\noindent
$h_i \in H_i$ の $H$ における像が零であると仮定する。
$H_i = \Ker(\psi_i)/\Im(\varphi_i)$ なので、$h_i$ は
$m \in \Ker(\psi_i) \subset M_i$ によって表せる。$H$ における $h_i$ の消滅は、
$\colim_i M_i$ における $m$ の類が $\varphi$ の像に属することを意味する。
したがって $\varphi_j(l) = \mu_{ij}(m)$ を満たす $j \geq i$ と $l \in L_j$ が
存在する。これは $H_j$ における $h_i$ の像が零であることを明らかに示す。
よって $\colim_i H_i \to H$ は単射である。
\end{proof}

\begin{example}
\label{algebra-example-colimit-not-exact}
一般には余極限をとる操作は完全ではない。例
\ref{algebra-example-zero-colimit-different} と同様に、$a < b$ および $a < c$ のほかに
真の不等式をもたない半順序集合 $I = \{a, b, c\}$ を考える。系の写像
$(0, \mathbf{Z}, \mathbf{Z}, 0, 0) \to
(\mathbf{Z}, \mathbf{Z}, \mathbf{Z}, 1, 1)$ を考える。例
\ref{algebra-example-zero-colimit-different} における余極限の記述から、この系の写像は
各対象上では単射であるにもかかわらず、付随する余極限の写像は単射でない。
したがって補題 \ref{algebra-lemma-directed-colimit-exact} の結論は一般の系については
偽である。
\end{example}

\begin{lemma}
\label{algebra-lemma-almost-directed-colimit-exact}
$\mathcal{I}$ を Categories, Lemma \ref{categories-lemma-split-into-directed} の仮定を
満たす添字圏とする。このとき $\mathcal{I}$ 上のアーベル群の図式の余極限を
とる操作は完全である（すなわち補題 \ref{algebra-lemma-directed-colimit-exact} の類似が
この状況で成り立つ）。
\end{lemma}

\begin{proof}
Categories, Lemma \ref{categories-lemma-split-into-directed} により、
$\mathcal{I} = \coprod_{j \in J} \mathcal{I}_j$ と書ける。ここで各
$\mathcal{I}_j$ はフィルター圏であり、$J$ は空でもよい。Categories, Lemma
\ref{categories-lemma-directed-category-system} により、添字圏 $\mathcal{I}_j$ 上の
余極限をとることは、ある有向集合上の余極限をとることと同じである。したがって
これらの余極限には補題 \ref{algebra-lemma-directed-colimit-exact} が適用できる。よって
問題は、集合 $J$ 上の $R$-加群の圏における余積が完全であることを示すことに
帰着する。言い換えれば、$R$ 加群の完全列 $L_j \to M_j \to N_j$ に対して、
% P01-E1273
$$
\bigoplus\nolimits_{j \in J} L_j
\longrightarrow
\bigoplus\nolimits_{j \in J} M_j
\longrightarrow
\bigoplus\nolimits_{j \in J} N_j
$$
が完全であることを示せばよい。これは直接確かめられ、$J$ が空の場合にも成り立つ。
\end{proof}






\section{局所化}
\label{algebra-section-localization}

\begin{definition}
\label{algebra-definition-multiplicative-subset}
$R$ を環、$S$ を $R$ の部分集合とする。$1\in S$ であり、$S$ が乗法について
閉じている、すなわち $s, s' \in S \Rightarrow ss' \in S$ であるとき、$S$ を
{\it $R$ の乗法的集合}という。
\end{definition}

\noindent
環 $A$ と乗法的集合 $S$ が与えられたとき、$A \times S$ 上の関係を次で定める。
$$
(x, s) \sim (y, t)
\Leftrightarrow
\exists u \in S \text{ such that } (xt-ys)u = 0
$$
これが同値関係であることは容易に確かめられる。$(x, s)$ の同値類を $x/s$
（または $\frac{x}{s}$）とし、すべての同値類の集合を $S^{-1}A$ とする。
$S^{-1}A$ における加法と乗法を次で定める。
$$
x/s + y/t = (xt + ys)/st, \quad
x/s \cdot y/t = xy/st
$$
これらの演算によって $S^{-1}A$ が環になることを確かめられる。

\begin{definition}
\label{algebra-definition-localization}
この環を{\it $S$ に関する $A$ の局所化}という。
\end{definition}

\noindent
$A$ からその局所化 $S^{-1}A$ への自然な環準同型
$$
A \longrightarrow S^{-1}A, \quad x \longmapsto x/1
$$
があり、これを{\it 局所化写像}ということがある。一般に局所化写像は単射では
ないが、$S$ が零因子を含まなければ単射である。実際 $x/1 = 0$ ならば、
$xu = 0$ を $A$ において満たす $u\in S$ が存在する。$S$ には零因子がないので
$x = 0$ である。環の局所化は次の普遍性をもつ。

\begin{proposition}
\label{algebra-proposition-universal-property-localization}
$f : A \to B$ を、$S$ のすべての元を $B$ の単元へ送る環準同型とする。このとき
次の図式を可換にする一意な準同型 $g : S^{-1}A \to B$ が存在する。
$$
\xymatrix{
A \ar[rr]^{f} \ar[dr] & & B \\
& S^{-1}A \ar[ur]_g
}
$$
\end{proposition}

\begin{proof}
存在を示す。写像 $g$ を次のように定める。$x/s\in S^{-1}A$ に対して
$g(x/s) = f(x)f(s)^{-1}\in B$ とおく。定義から、これは良定義な環準同型である
ことが容易に確かめられる。またこの写像が図式を可換にすることも明らかである。

\medskip\noindent
一意性を示す。$g' : S^{-1}A \to B$ が $g'(x/1) = f(x)$ を満たすならば
$g = g'$ であることを示す。図式の可換性から、$s \in S$ に対して
$f(s) = g'(s/1)$ である。一方 $B$ において $g'(1/s)f(s) = 1$ だから、
$g'(1/s) = f(s)^{-1}$ であり、したがって
$g'(x/s) = g'(x/1)g'(1/s) = f(x)f(s)^{-1} = g(x/s)$ である。
\end{proof}

\begin{lemma}
\label{algebra-lemma-localization-zero}
局所化 $S^{-1}A$ が零環であるための必要十分条件は $0\in S$ である。
\end{lemma}

\begin{proof}
$0\in S$ ならば、定義により任意の対 $(a, s)\sim (0, 1)$ である。
$0\not \in S$ ならば、明らかに $S^{-1}A$ において $1/1 \neq 0/1$ である。
\end{proof}

\begin{lemma}
\label{algebra-lemma-localization-and-modules}
$R$ を環とし、$S \subset R$ を乗法的集合とする。$S^{-1}R$-加群の圏は、
各 $s \in S$ が $N$ 上の自己同型として作用するという性質をもつ
$R$-加群 $N$ の圏と同値である。
\end{lemma}

\begin{proof}
同値を定める関手は、$S^{-1}R$-加群 $M$ に、局所化写像
$R \to S^{-1}R$ を介して $R$-加群とみなした同じ加群を対応させる。逆に、
$N$ が $R$-加群で、各 $s \in S$ が自己同型 $s_N$ として作用するならば、
$x/s$ を $x_N  \circ s_N^{-1}$ として作用させることにより、$N$ を
$S^{-1}R$-加群とみなせる。これら二つの関手が互いに準逆であることの確認は
省略する。
\end{proof}

\noindent
環の局所化という概念は加群の局所化へ一般化できる。$A$ を環、$S$ を $A$ の
乗法的集合、$M$ を $A$-加群とする。$M \times S$ 上の関係を次で定める。
$$
(m, s) \sim (n, t)
\Leftrightarrow
\exists u\in S \text{ such that } (mt-ns)u = 0
$$
これは明らかに同値関係である。$(m, s)$ の同値類を $m/s$（または
$\frac{m}{s}$）とし、すべての同値類の集合を $S^{-1}M$ とする。加法と
スカラー乗法を次で定める。
% P01-E1274
$$
m/s + n/t = (mt + ns)/st,\quad
m/s\cdot n/t = mn/st
$$
これにより $S^{-1}M$ が $S^{-1}A$-加群になることは明らかである。

\begin{definition}
\label{algebra-definition-localization-module}
$S^{-1}A$-加群 $S^{-1}M$ を、$S$ に関する $M$ の{\it 局所化}という。
\end{definition}

\noindent
$A$-加群準同型 $M \to S^{-1}M$, $m \mapsto m/1$ があることに注意せよ。
これを{\it 局所化写像}ということがある。これは次の普遍性を満たす。

\begin{lemma}
\label{algebra-lemma-universal-property-localization-module}
$R$ を環とする。$S \subset R$ を乗法的集合とする。$M$, $N$ を $R$-加群とする。
% P01-E1275
$S$ のすべての元が $N$ 上の自己同型として作用すると仮定する。このとき
局所化写像から誘導される標準写像
$$
\Hom_R(S^{-1}M, N) \longrightarrow \Hom_R(M, N)
$$
は同型である。
\end{lemma}

\begin{proof}
写像が良定義かつ $R$-線形であることは明らかである。単射性を示す。
$\alpha \in \Hom_R(S^{-1}M, N)$ とし、任意の元 $m/s \in S^{-1}M$ をとる。
$s \cdot \alpha(m/s) = \alpha(m/1)$ なので、
$ \alpha(m/s) =s^{-1}(\alpha (m/1))$ である。したがって $\alpha$ は、$M$ の
$S^{-1}M$ における像上での値によって完全に決定される。全射性を示す。
$\beta : M \rightarrow N$ を与えられた R-線形写像とする。これを $S^{-1}M$ へ
「延長」できることを示す必要がある。集合の写像
$$
M \times S \rightarrow N,\quad
(m,s) \mapsto s^{-1}\beta(m)
$$
を定める。この写像が上の同値関係を保つことは明らかなので、良定義な写像
$\alpha : S^{-1}M \rightarrow N$ に降りる。この写像が $R$-線形であることを
示せばよい。そこで $r, r' \in R$, $s, s' \in S$, $m, m' \in M$ とすると、
\begin{align*}
\alpha(r \cdot m/s + r' \cdot m' /s')
& =
\alpha((r \cdot s' \cdot m + r' \cdot s \cdot m') /(ss')) \\
& =
(ss')^{-1}\beta(r \cdot s' \cdot m + r' \cdot s \cdot m') \\
& =
(ss')^{-1} (r \cdot s' \beta (m) + r' \cdot s \beta (m')) \\
& =
r \alpha (m/s) + r' \alpha (m' /s')
\end{align*}
であり、主張が従う。
\end{proof}

\begin{example}
\label{algebra-example-localize-at-prime}
$A$ を環、$M$ を $A$-加群とする。局所化の重要な例を挙げる。
\begin{enumerate}
\item $A$ の素イデアル $\mathfrak p$ が与えられたとき、
$S = A\setminus\mathfrak p$ を考える。$S$ が乗法的集合であることは直ちに
確かめられる。この場合、$S$ に関する $A$ と $M$ の局所化をそれぞれ
$A_\mathfrak p$ および $M_\mathfrak p$ と記す。これらを
{\it $\mathfrak p$ における $A$, resp.\ $M$ の局所化}という。
\item $f\in A$ とし、$S = \{1, f, f^2, \ldots\}$ を考える。これは明らかに
$A$ の乗法的集合である。この場合、局所化 $S^{-1}A$（resp. $S^{-1}M$）を
$A_f$（resp. $M_f$）と記す。これを
{\it $f$ に関する $A$, resp.\ $M$ の局所化}という。
$A_f = 0$ であるための必要十分条件は、$f$ が $A$ において冪零であることである。
\item $S = \{f \in A \mid f \text{ is not a zerodivisor in }A\}$ とする。これは
$A$ の乗法的集合である。この場合、環 $Q(A) = S^{-1}A$ を $A$ の
{\it 全商環}または{\it 全分数環}という。
\item $A$ が整域ならば、全商環 $Q(A)$ は $A$ の分数体である。Fields, Example
\ref{fields-example-quotient-field} を参照せよ。
\end{enumerate}
\end{example}

\begin{lemma}
\label{algebra-lemma-localization-colimit}
$R$ を環、$S \subset R$ を乗法的集合、$M$ を $R$-加群とする。このとき
$$
S^{-1}M = \colim_{f \in S} M_f
$$
である。ここで $S$ 上の前順序は、$f \geq f' \Leftrightarrow f = f'f''$ と定める。
ただし $f'' \in R$ である。この場合、写像 $M_{f'} \to M_f$ は
$m/(f')^e \mapsto m(f'')^e/f^e$ で与えられる。
\end{lemma}

\begin{proof}
省略する。ヒントとして、補題
\ref{algebra-lemma-universal-property-localization-module} の普遍性を用いよ。
\end{proof}

\noindent
以下では、$A$ を環、$M, N$ を $A$ 上の加群とする。

\medskip\noindent
$S$ と $S'$ が $A$ の乗法的集合ならば、
$$
SS' = \{ss' : s\in S, \ s'\in S'\}
$$
も $A$ の乗法的集合であることは明らかである。このとき次が成り立つ。

\begin{proposition}
\label{algebra-proposition-localize-twice}
$\overline{S}$ を $S'^{-1}A$ における $S$ の像とする。このとき
$(SS')^{-1}A$ は $\overline{S}^{-1}(S'^{-1}A)$ と同型である。
\end{proposition}

\begin{proof}
$x\in A$ を $x/1\in (SS')^{-1}A$ へ送る写像は、普遍性により
$x/s\in S'^{-1}A$ を $x/s \in (SS')^{-1}A$ へ送る写像を誘導する。
$\overline{S}$ の元の像は $(SS')^{-1}A$ において可逆である。再び普遍性により、
$(x/s')/(s/1)$ を $x/ss'$ へ送る写像
$f : \overline{S}^{-1}(S'^{-1}A) \to (SS')^{-1}A$ を得る。

\medskip\noindent
一方、$x\in A$ を $(x/1)/(1/1)$ へ送る $A$ から
$\overline{S}^{-1}(S'^{-1}A)$ への写像も、普遍性により $x/ss'$ を
$(x/s')/(s/1)$ へ送る写像
$g : (SS')^{-1}A \to \overline{S}^{-1}(S'^{-1}A)$ を誘導する。
$f$ と $g$ が互いに逆であることは直ちに確かめられる。したがって両者は同型である。
\end{proof}

\noindent
加群 $M$ については次が成り立つ。

\begin{proposition}
\label{algebra-proposition-localize-twice-module}
$S'^{-1}M$ を $A$-加群とみなす。このとき $S^{-1}(S'^{-1}M)$ は
$(SS')^{-1}M$ と同型である。
\end{proposition}

\begin{proof}
与えられた $A$-加群 M に対して、$S^{-1}M$ の普遍性はまだ証明していない
ことに注意せよ。したがって直前の証明と同じ論法は使えず、同型を明示的に
構成しなければならない。

\medskip\noindent
写像を次のように定める。
\begin{align*}
& f : S^{-1}(S'^{-1}M) \longrightarrow (SS')^{-1}M, \quad \frac{x/s'}{s}\mapsto
x/ss'\\
& g : (SS')^{-1}M \longrightarrow S^{-1}(S'^{-1}M), \quad x/t\mapsto
\frac{x/s'}{s}\ \text{for some }s\in S, s'\in S', \text{ and }
t = ss'
\end{align*}
これらの準同型が良定義、すなわち分数の選択に依存しないことを確かめる必要が
ある。これは容易に確かめられ、また互いに逆であることも直ちに示せる。
\end{proof}

\noindent
$u : M \to N$ が $A$-加群準同型ならば、局所化は良定義な
$S^{-1}A$-加群準同型 $S^{-1}u : S^{-1}M \to S^{-1}N$ を実際に誘導し、これは
$x/s$ を $u(x)/s$ へ送る。この構成が関手的であることは直ちに確かめられるので、
$S^{-1}$ は $A$-加群の圏から $S^{-1}A$-加群の圏への関手である。さらにこの関手は
完全であることを次の命題で示す。

\begin{proposition}
\label{algebra-proposition-localization-exact}
\begin{slogan}
局所化は完全である。
\end{slogan}
$L\xrightarrow{u} M\xrightarrow{v} N$ を $R$-加群の完全列とする。このとき
$S^{-1}L \to S^{-1}M \to S^{-1}N$ も完全である。
\end{proposition}

\begin{proof}
局所化は関手なので、$S^{-1}L \to S^{-1}M \to S^{-1}N$ が複体であることは
明らかである。次に $x/s \in S^{-1}M$ であり、$x/s$ が $S^{-1}N$ の零元に
写ると仮定する。定義により、$N$ において
$v(xt) = v(x)t = 0$ を満たす $t\in S$ が存在する。したがって
$xt \in \Ker(v)$ である。$L \to M \to N$ の完全性により、ある $L$ の元 $y$ に
対して $xt = u(y)$ である。このとき $x/s$ は $y/st$ の像である。これで完全性が
証明された。
\end{proof}

\begin{lemma}
\label{algebra-lemma-localize-quotient-modules}
局所化は商を保つ。すなわち $N$ が $M$ の部分加群ならば、
$S^{-1}(M/N)\allowbreak\simeq (S^{-1}M)\allowbreak/(S^{-1}N)$ である。
\end{lemma}

\begin{proof}
完全列
$$
0 \longrightarrow N \longrightarrow M \longrightarrow M/N \longrightarrow 0
$$
から完全列
$$
0 \longrightarrow S^{-1}N \longrightarrow S^{-1}M
\longrightarrow S^{-1}(M/N) \longrightarrow 0
$$
を得る。したがって系が従う。
\end{proof}

\noindent
直前の系において、$I$ を $A$ のイデアルとして $N = I$, $M = A$ と
おけば、$A$-加群として $S^{-1}A/S^{-1}I \simeq S^{-1}(A/I)$ であることが分かる。
% P01-E1276
次の命題は、これらが環として同型であることを示す。

\begin{proposition}
\label{algebra-proposition-localize-quotient}
$I$ を $A$ のイデアル、$S$ を $A$ の乗法的集合とする。このとき $S^{-1}I$ は
$S^{-1}A$ のイデアルであり、$\overline{S}^{-1}(A/I)$ は
$S^{-1}A/S^{-1}I$ と同型である。ここで $\overline{S}$ は $A/I$ における $S$ の
像である。
\end{proposition}

\begin{proof}
$I$ 自身がイデアルなので、$S^{-1}I$ がイデアルであることは明らかである。写像
$$
f : S^{-1}A\longrightarrow \overline{S}^{-1}(A/I), \quad x/s\mapsto
\overline{x}/\overline{s}
$$
を定める。ここで $\overline{x}$ と $\overline{s}$ は $A/I$ における $x$ と $s$ の
像である。この証明では同様の記法を用いる。この写像は $S^{-1}A$ の普遍性により
良定義であり、その核は $S^{-1}I$ を含む。したがって写像
$$
\overline{f} : S^{-1}A/S^{-1}I \longrightarrow \overline{S}^{-1}(A/I), \quad
\overline{x/s}\mapsto \overline{x}/\overline{s}
$$
を誘導する。

\medskip\noindent
一方、$x$ を $\overline{x/1}$ へ送る写像 $A \to S^{-1}A/S^{-1}I$ は、
$\overline{x}$ を $\overline{x/1}$ へ送る写像 $A/I \to S^{-1}A/S^{-1}I$ を
誘導する。$\overline{S}$ の像は $S^{-1}A/S^{-1}I$ において可逆なので、普遍性により
写像
$$
g : \overline{S}^{-1}(A/I) \longrightarrow S^{-1}A/S^{-1}I, \quad
\frac{\overline{x}}{\overline{s}}\mapsto \overline{x/s}
$$
を誘導する。$\overline{f}$ と $g$ が互いに逆であることは明らかであり、したがって
両者は同型である。
\end{proof}

\noindent
次に、部分加群が局所化のもとでどのように振る舞うかを調べる。

\begin{lemma}
\label{algebra-lemma-submodule-localization}
$S^{-1}M$ の任意の部分加群 $N'$ は、ある $N\subset M$ に対して $S^{-1}N$ の形で
ある。実際、$N$ として $M$ における $N'$ の逆像をとることができる。
\end{lemma}

\begin{proof}
$N$ を $M$ における $N'$ の逆像とする。このとき $S^{-1}N\supset N'$ であることが
分かる。等号を示すため、$S^{-1}N$ の元 $x/s$ をとる。ただし $s\in S$ かつ
$x\in N$ である。このとき $x/1\in N'$ である。$N'$ は
$S^{-1}R$-部分加群なので、$x/s = x/1\cdot 1/s\in N'$ である。これで証明が
完了する。
\end{proof}

\noindent
$M = A$ とし、$N = I$ を $A$ のイデアルとすれば、命題
\ref{algebra-proposition-localize-quotient} の前半の逆とみなせる次の系を得る。

\begin{lemma}
\label{algebra-lemma-ideal-in-localization}
\begin{slogan}
環の局所化のイデアルは、イデアルの局所化である。
\end{slogan}
$S^{-1}A$ の各イデアル $I'$ は $S^{-1}I$ の形であり、$I$ として $A$ における
$I'$ の逆像をとることができる。
\end{lemma}

\begin{proof}
補題 \ref{algebra-lemma-submodule-localization} から直ちに従う。
\end{proof}










\section{内部 Hom}
\label{algebra-section-hom}

\noindent
$R$ が環で、$M$, $N$ が $R$-加群ならば、
$$
\Hom_R(M, N) = \{ \varphi : M \to N\}
$$
は $M$ から $N$ への $R$-線形写像の集合である。通常どおり
$(\varphi + \psi)(m) = \varphi(m) + \psi(m)$ と定めることにより、この集合は
アーベル群の構造をもつ。実際、$\Hom_R(M, N)$ は規則
$(x \varphi)(m) = x \varphi(m) = \varphi(xm)$ によって $R$-加群でもある。

\medskip\noindent
$R$-加群の写像 $a : M \to M'$ および $b : N \to N'$ が与えられたとき、
準同型に $a$ を前から、$b$ を後ろから合成できる。これにより可換図式
$$
\xymatrix{
\Hom_R(M', N) \ar[d]_{- \circ a} \ar[r]_{b \circ -} &
\Hom_R(M', N') \ar[d]^{- \circ a} \\
\Hom_R(M, N) \ar[r]^{b \circ -} &
\Hom_R(M, N')
}
$$
を得る。実際、この図式の写像は $R$-加群準同型である。したがって
$\Hom_R$ は加法的関手
$$
\text{Mod}_R^{opp} \times \text{Mod}_R \longrightarrow \text{Mod}_R, \quad
(M, N) \longmapsto \Hom_R(M, N)
$$
を定める。

\begin{lemma}
\label{algebra-lemma-hom-exact}
完全性と $\Hom_R$。$R$ を環とし、$M_1$, $M_2$, $M_3$ を $R$-加群とする。
$M_1 \to M_2$ と $M_2 \to M_3$ を $R$-加群準同型とする。
\begin{enumerate}
\item $M_1 \to M_2 \to M_3 \to 0$ が完全であるための必要十分条件は、すべての
$R$-加群 $N$ に対して
$0 \to \Hom_R(M_3, N) \to \Hom_R(M_2, N) \to \Hom_R(M_1, N)$ が完全である
ことである。
\item $0 \to M_1 \to M_2 \to M_3$ が完全であるための必要十分条件は、すべての
$R$-加群 $N$ に対して
$0 \to \Hom_R(N, M_1) \to \Hom_R(N, M_2) \to \Hom_R(N, M_3)$ が完全である
ことである。
\end{enumerate}
\end{lemma}

\begin{proof}
省略する。
\end{proof}

\begin{lemma}
\label{algebra-lemma-hom-from-finitely-presented}
$R$ を環、$M$ を有限表示 $R$-加群、$N$ を $R$-加群とする。
\begin{enumerate}
\item $f \in R$ に対して
$\Hom_R(M, N)_f = \Hom_{R_f}(M_f, N_f) = \Hom_R(M_f, N_f)$ である。
\item $R$ の乗法的集合 $S$ に対して
$$
S^{-1}\Hom_R(M, N) = \Hom_{S^{-1}R}(S^{-1}M, S^{-1}N) =
\Hom_R(S^{-1}M, S^{-1}N).
$$
\end{enumerate}
\end{lemma}

\begin{proof}
(1) は (2) の特別な場合である。(2) の第二の等号は補題
\ref{algebra-lemma-localization-and-modules} から従う。表示
$$
\bigoplus\nolimits_{j = 1, \ldots, m} R
\longrightarrow
\bigoplus\nolimits_{i = 1, \ldots, n} R
\to M \to 0.
$$
を選ぶ。補題 \ref{algebra-lemma-hom-exact} により、これは完全列
$$
0 \to
\Hom_R(M, N) \to
\bigoplus\nolimits_{i = 1, \ldots, n} N
\longrightarrow
\bigoplus\nolimits_{j = 1, \ldots, m} N.
$$
を与える。$S$ の元を可逆にし、命題 \ref{algebra-proposition-localization-exact} を用いると、
完全列
$$
0 \to
S^{-1}\Hom_R(M, N) \to
\bigoplus\nolimits_{i = 1, \ldots, n} S^{-1}N
\longrightarrow
\bigoplus\nolimits_{j = 1, \ldots, m} S^{-1}N
$$
を得る。一方 $S^{-1}M$ は完全列
$$
\bigoplus\nolimits_{j = 1, \ldots, m} S^{-1}R
\longrightarrow
\bigoplus\nolimits_{i = 1, \ldots, n} S^{-1}R \to S^{-1}M \to 0
$$
に属し、これは補題 \ref{algebra-lemma-hom-exact} により完全列
$$
0 \to
\Hom_{S^{-1}R}(S^{-1}M, S^{-1}N) \to
\bigoplus\nolimits_{i = 1, \ldots, n} S^{-1}N
\longrightarrow
\bigoplus\nolimits_{j = 1, \ldots, m} S^{-1}N
$$
を誘導する。これは上の完全列と同じなので、結論が従う。
\end{proof}




\section{有限加群および有限表示加群の特徴づけ}
\label{algebra-section-colim-and-hom}

\noindent
環 $R$ 上の加群 $N$ が有限加群または有限表示加群であるかどうかは、関手
$\Hom_R(N, -)$ によって特徴づけられる。

\begin{lemma}
\label{algebra-lemma-characterize-finite-module-hom}
$R$ を環、$N$ を $R$-加群とする。次は同値である。
\begin{enumerate}
\item $N$ は有限 $R$-加群である。
\item $R$-加群の任意のフィルター余極限 $M = \colim M_i$ に対して、写像
$\colim \Hom_R(N, M_i) \to \Hom_R(N, M)$ は単射である。
\end{enumerate}
\end{lemma}

\begin{proof}
(1) を仮定し、$N$ の生成元 $x_1, \ldots, x_m$ を選ぶ。$N \to M_i$ が加群準同型で、
合成 $N \to M_i \to M$ が零ならば、$M = \colim_{i' \geq i} M_{i'}$ なので、
各 $j \in \{1, \ldots, m\}$ に対し、$x_j$ が $M_{i'}$ の零元へ写るような
% P01-E1277
$i' \geq i$ を見つけられる。$x_j$ は有限個しかないので、すべてに対して働く
単一の $i'$ を見つけられる。このとき合成 $N \to M_i \to M_{i'}$ は零であり、
写像が単射、すなわち (2) が成り立つと結論できる。

\medskip\noindent
(2) を仮定する。有限部分集合 $E \subset N$ に対して、$E$ の元が生成する
$R$-部分加群を $N_E \subset N$ と記す。このとき
$0 = \colim N/N_E$ はフィルター余極限である。したがって
$\text{id} : N \to N$ はある $E$ に対して $N_E$ を経由する。すなわち $N$ は
有限生成である。
\end{proof}

\noindent
参照の便宜のため、加群の元の間の関係式という概念を定義する。

\begin{definition}
\label{algebra-definition-relation}
$R$ を環、$M$ を $R$-加群とする。$n \geq 0$ とし、
$i = 1, \ldots, n$ に対して $x_i \in M$ とする。$M$ における
$x_1, \ldots, x_n$ の間の{\it 関係式}とは、
$\sum_{i = 1, \ldots, n} f_i x_i = 0$ を満たす元の列
$f_1, \ldots, f_n \in R$ のことである。
\end{definition}

\begin{lemma}
\label{algebra-lemma-module-colimit-fp}
$R$ を環、$M$ を $R$-加群とする。このとき $M$ は、各 $M_i$ が有限表示
$R$-加群であるような $R$-加群の有向系 $(M_i, \mu_{ij})$ の余極限である。
\end{lemma}

\begin{proof}
任意の有限部分集合 $S \subset M$ と、$S$ の元の間の関係式の任意の有限集合
$E$ を考える。各 $s \in S$ は $x_s \in M$ に対応し、各 $e \in E$ は
$\sum f_{e, s} x_s = 0$ を満たす元 $f_{e, s} \in R$ のベクトルからなる。
$M_{S, E}$ を写像
$$
R^{\# E} \longrightarrow R^{\# S}, \quad
(g_e)_{e\in E} \longmapsto (\sum g_e f_{e, s})_{s\in S}.
$$
の余核とする。標準写像 $M_{S, E} \to M$ がある。$S \subset S'$ であり、
$E$ の元がこの写像を介して $E'$ の関係式に対応するならば、$M$ への写像と
可換する明らかな写像 $M_{S, E} \to M_{S', E'}$ がある。$I$ を、上の包含関係で
順序づけた対 $(S, E)$ の集合とする。この有向系の余極限が $M$ であることは
明らかである。
\end{proof}

\begin{lemma}
\label{algebra-lemma-characterize-finitely-presented-module-hom}
$R$ を環、$N$ を $R$-加群とする。次は同値である。
\begin{enumerate}
\item $N$ は有限表示 $R$-加群である。
\item $R$-加群の任意のフィルター余極限 $M = \colim M_i$ に対して、写像
$\colim \Hom_R(N, M_i) \to \Hom_R(N, M)$ は全単射である。
\end{enumerate}
\end{lemma}

\begin{proof}
(1) を仮定し、$F_i$ が有限自由である完全列 $F_{-1} \to F_0 \to N \to 0$ を
選ぶ。このとき $R$-加群 $M$ に関して関手的な完全列
$$
0 \to \Hom_R(N, M) \to \Hom_R(F_0, M) \to \Hom_R(F_{-1}, M)
$$
がある。$\Hom_R(R^{\oplus n}, M) = M^{\oplus n}$ なので、関手
$\Hom_R(F_i, M)$ はフィルター余極限と可換する。フィルター余極限は完全である
（補題 \ref{algebra-lemma-directed-colimit-exact}）から、(2) が成り立つ。

\medskip\noindent
(2) を仮定する。補題 \ref{algebra-lemma-module-colimit-fp} により、すべての $i$ に対して
$N_i$ が有限表示となるように、$N = \colim N_i$ をフィルター余極限として書ける。
したがって $\text{id}_N$ はある $i$ に対して $N_i$ を経由する。これは $N$ が
有限表示 $R$-加群（すなわち $N_i$）の直和因子であることを意味する。ゆえに
$N$ は有限表示である。
\end{proof}












\section{テンソル積}
\label{algebra-section-tensor-product}

\begin{definition}
\label{algebra-definition-bilinear}
$R$ を環、$M, N, P$ を三つの $R$-加群とする。写像
$f : M \times N \to P$（ここで $M \times N$ は二つの $R$-加群の直積集合として
のみ考える）が{\it $R$-双線形}であるとは、各 $x \in M$ に対して
$N$ から $P$ への写像 $y\mapsto f(x, y)$ が $R$-線形であり、かつ各
$y\in N$ に対して写像 $x\mapsto f(x, y)$ も $R$-線形であることをいう。
\end{definition}

\begin{lemma}
\label{algebra-lemma-tensor-product}
$M, N$ を $R$-加群とする。このとき、$R$-加群 $T$ と $R$-双線形写像
$g : M \times N \to T$ からなる対 $(T, g)$ であって、次の普遍性をもつものが
存在する。任意の $R$-加群 $P$ と任意の $R$-双線形写像
$f : M \times N \to P$ に対し、$f = \tilde{f} \circ g$ を満たす一意な
$R$-線形写像 $\tilde{f} : T \to P$ が存在する。言い換えれば、次の図式が可換である。
$$
\xymatrix{
M \times N \ar[rr]^f \ar[dr]_g & & P\\
& T \ar[ur]_{\tilde f}
}
$$
さらに、$(T, g)$ と $(T', g')$ がこの性質をもつ二つの対ならば、
$j\circ g = g'$ を満たす一意な同型 $j : T \to T'$ が存在する。
\end{lemma}

\noindent
上の普遍性を満たす $R$-加群 $T$ を、$R$-加群 $M$ と $N$ の{\it テンソル積}
といい、$M \otimes_R N$ と記す。

\begin{proof}
まず、このような $R$-加群 $T$ の存在を証明する。$M, N$ を $R$-加群とする。
$P$ を自由 $R$-加群 $R^{(M \times N)}$ とし、$Q$ を次の形のすべての元
（$x\in M, y\in N$）によって生成される $R$-加群として、$T$ を商加群 $P/Q$
とする。
\begin{align*}
(x + x', y) - (x, y) - (x', y), \\
(x, y + y') - (x, y) - (x, y'), \\
(ax, y) - a(x, y), \\
(x, ay) - a(x, y)
\end{align*}
自然な写像を $\pi : M \times N \to T$ と記す。双線形性の条件を確認すれば、
上の関係式からこの写像が $R$-双線形であることが分かる。像を
$\pi(x, y) = x \otimes y$ と記すと、これらの元は $T$ を生成する。ここで
$f : M \times N \to P$ を $R$-双線形写像とすると、
$f'(x \otimes y) = f(x, y)$ を延長することにより $f' : T \to P$ を定義できる。
明らかに $f = f'\circ \pi$ である。さらに $f'$ は生成系
$\{x \otimes y : x\in M, y\in N\}$ 上の値によって一意に決まる。

同じ性質を満たす別の対 $(T', g')$ があると仮定する。このとき
$g' = j\circ g$, $g = j'\circ g'$ を満たす一意な $j : T \to T'$ および
$j' : T' \to T$ が存在する。しかし写像 $(j\circ j') \circ g$ と $g$ はどちらも
% P01-E1278
普遍性を満たすので、一意性により等しく、したがって $j'\circ j$ は $T$ 上の
恒等写像である。同様に $(j'\circ j) \circ g' = g'$ であり、$j\circ j'$ は
$T'$ 上の恒等写像である。したがって $j$ は同型である。
\end{proof}

\begin{lemma}
\label{algebra-lemma-flip-tensor-product}
$M, N, P$ を $R$-加群とする。このとき双線形写像
\begin{align*}
(x, y) & \mapsto y \otimes x\\
(x + y, z) & \mapsto x \otimes z + y \otimes z\\
(r, x) & \mapsto rx
\end{align*}
% P01-E1279
は一意な同型
\begin{align*}
M \otimes_R N & \to N \otimes_R M, \\
(M\oplus N)\otimes_R P & \to (M \otimes_R P)\oplus(N \otimes_R P),  \\
R \otimes_R M & \to M
\end{align*}
を誘導する。
\end{lemma}

\begin{proof}
省略する。
\end{proof}

\noindent
二つの $R$-加群のテンソル積を有限個の $R$-加群へ一般化し、多重テンソル積と
多重線形写像との対応を設定できる。ほぼ同じ構成によって次を証明できる。

\begin{lemma}
\label{algebra-lemma-multilinear}
$M_1, \ldots, M_r$ を $R$-加群とする。このとき $R$-加群 T と $R$-多重線形写像
$g : M_1\times \ldots \times M_r \to T$ からなる対 $(T, g)$ であって、次の
普遍性をもつものが存在する。任意の $R$-多重線形写像
$f : M_1\times \ldots \times M_r \to P$ に対し、$f'\circ g = f$ を満たす一意な
$R$-加群準同型 $f' : T \to P$ が存在する。このような加群 $T$ は一意な同型を
除いて一意である。これを $M_1\otimes_R \ldots \otimes_R M_r$ と記し、普遍的な
多重線形写像を $(m_1, \ldots, m_r) \mapsto m_1 \otimes \ldots \otimes m_r$ と記す。
\end{lemma}

\begin{proof}
省略する。
\end{proof}

\begin{lemma}
\label{algebra-lemma-transitive}
準同型
$$
(M \otimes_R N)\otimes_R P \to
M \otimes_R N \otimes_R P \to
M \otimes_R (N \otimes_R P)
$$
であって、$f((x \otimes y)\otimes z) = x \otimes y \otimes z$ および
$g(x \otimes y \otimes z) = x \otimes (y \otimes z)$ を満たすものは良定義で
同型である。ただし $x\in M, y\in N, z\in P$ とする。
\end{lemma}

\begin{proof}
$f$ が良定義で同型であることを証明する。$g$ についても同様である。任意の
$z\in P$ を固定すると、写像 $(x, y)\mapsto x \otimes y \otimes z$ は
$x\in M, y\in N$ に関して $x$ と $y$ の $R$-双線形写像であり、したがって
$f_z(x \otimes y) = x \otimes y \otimes z$ とする準同型
$f_z : M \otimes N \to M \otimes N \otimes P$ を誘導する。次に
$(w, z)\mapsto f_z(w)$ で与えられる写像
$(M \otimes N)\times P \to M \otimes N \otimes P$ を考える。この写像は
$R$-双線形なので、$f((x \otimes y)\otimes z) = x \otimes y \otimes z$ を満たす
$f : (M \otimes_R N)\otimes_R P \to M \otimes_R N \otimes_R P$ を誘導する。
逆写像を構成するため、写像 $\pi : M \times N \times P \to (M \otimes N)\otimes P$
が $R$-三重線形であることに注意する。したがって普遍性と整合する $R$-線形写像
$h : M \otimes N \otimes P \to (M \otimes N)\otimes P$ を誘導する。ここで
$h(x \otimes y \otimes z) = (x \otimes y)\otimes z$ である。$f$ と $h$ の明示式から、
$f\circ h$ と $h\circ f$ はそれぞれ $M \otimes N \otimes P$ と
$(M \otimes N)\otimes P$ の恒等写像である。したがって $f$ は所望の同型である。
\end{proof}

\noindent
帰納法により、これが多重テンソル積へ拡張されることが分かる。補題
\ref{algebra-lemma-flip-tensor-product} と合わせると、$R$-加群の圏上のテンソル積演算は
結合的、可換、かつ分配的である。

\begin{definition}
\label{algebra-definition-bimodule}
アーベル群 $N$ が{\it $(A, B)$-双加群}であるとは、$A$-加群かつ $B$-加群であり、
すべての $a \in A$ と $b \in B$ に対して $a$ と $b$ による乗法が可換する、
すなわちすべての $n \in N$ に対して $b(an) = a(bn)$ となることをいう。この状況では
通常 $B$-作用を右側に書く。すなわち $b \in B$ と $n \in N$ に対し、$n$ に $b$ を
掛けた結果を $nb$ と記す。この規約のもとで、上の両立条件は、すべての
$a\in A, b\in B, x\in N$ に対して $(ax)b = a(xb)$ となることである。
$(A, B)$-双加群 $N$ を表す略記として $_AN_B$ を用いる。
\end{definition}

\begin{lemma}
\label{algebra-lemma-tensor-with-bimodule}
$A$-加群 $M$、$B$-加群 $P$、$(A, B)$-双加群 $N$ に対して、加群
$(M \otimes_A N)\otimes_B P$ と $M \otimes_A(N \otimes_B P)$ はどちらも
$(A, B)$-双加群構造をもつことができ、さらに
$$
(M \otimes_A N)\otimes_B P \cong M \otimes_A(N \otimes_B P).
$$
である。
\end{lemma}

\begin{proof}
先験的には $M \otimes_A N$ は $A$-加群であるが、
$$
(x \otimes y)b = x \otimes yb, \quad x\in M, y\in N, b\in B
$$
と定めることにより $B$-加群構造を与えられる。したがって $M \otimes_A N$ は
$(A, B)$-双加群となる。$N \otimes_B P$ についても同様であり、したがって
$(M \otimes_A N)\otimes_B P$ および $M \otimes_A(N \otimes_B P)$ についても
同様である。補題 \ref{algebra-lemma-transitive} により、これら二つの加群は同じ写像を
介して $A$-加群としても $B$-加群としても同型である。
\end{proof}

\begin{lemma}
\label{algebra-lemma-hom-from-tensor-product}
任意の三つの $R$-加群 $M, N, P$ に対して
$$
\Hom_R(M \otimes_R N, P) \cong \Hom_R(M, \Hom_R(N, P))
$$
である。
\end{lemma}

\begin{proof}
$R$-線形写像 $\hat{f}\in \Hom_R(M \otimes_R N, P)$ は $R$-双線形写像
$f : M \times N \to P$ に対応する。各 $x\in M$ に対し、写像
$y\mapsto f(x, y)$ は普遍性により $R$-線形である。したがって $f$ は写像
$\phi_f : M \to \Hom_R(N, P)$ に対応する。この写像は、
$$
\phi_f(ax + y)(z) =
f(ax + y, z) = af(x, z)+f(y, z) =
(a\phi_f(x)+\phi_f(y))(z),
$$
がすべての $a \in R$, $x \in M$, $y \in M$, $z \in N$ に対して成り立つので
$R$-線形である。逆に、任意の $f \in \Hom_R(M, \Hom_R(N, P))$ は
$(x, y)\mapsto f(x)(y)$ によって $R$-双線形写像 $M \times N \to P$ を定める。
したがってこれは二つの加群 $\Hom_R(M \otimes_R N, P)$ と
$\Hom_R(M, \Hom_R(N, P))$ の間の自然な一対一対応である。
\end{proof}

\begin{lemma}[テンソル積は余極限と可換する]
\label{algebra-lemma-tensor-products-commute-with-limits}
$(M_i, \mu_{ij})$ を前順序集合 $I$ 上の系とし、$N$ を $R$-加群とする。このとき
$$
\colim (M_i \otimes N) \cong (\colim M_i)\otimes N.
$$
である。さらに、この同型は、$M = \colim_i M_i$ として自然な写像
$\mu_i : M_i \to M$ をとったとき、準同型
$\mu_i \otimes 1: M_i \otimes N \to M \otimes N$ によって誘導される。
\end{lemma}

\begin{proof}
第一の証明。補題 \ref{algebra-lemma-hom-from-tensor-product} により、関手
$M' \mapsto M' \otimes_R N$ は関手 $N' \mapsto \Hom_R(N, N')$ の左随伴である。
したがって $M' \mapsto M' \otimes_R N$ はすべての余極限と可換する。Categories,
Lemma \ref{categories-lemma-adjoint-exact} を参照せよ。

\medskip\noindent
第二の直接証明。$P = \colim (M_i \otimes N)$ とし、余積射を
$\lambda_i : M_i \otimes N \to P$ とする。また $M = \colim M_i$ とし、余積射を
$\mu_i : M_i \to M$ とする。このときすべての $i\leq j$ に対して次の図式は可換である。
$$
\xymatrix{
M_i \otimes N \ar[r]_{\mu_i \otimes 1} \ar[d]_{\mu_{ij} \otimes 1} &
M \otimes N \ar[d]^{\text{id}} \\
M_j \otimes N \ar[r]^{\mu_j \otimes 1} &
M \otimes N
}
$$
補題 \ref{algebra-lemma-homomorphism-limit} により、これらの写像は
$\mu_i \otimes 1 = \psi \circ \lambda_i$ を満たす一意な準同型
$\psi : P \to M \otimes N$ を誘導する。

\medskip\noindent
逆写像を構成する。各 $i\in I$ に対し、標準的な $R$-双線形写像
$g_i : M_i \times N \to M_i \otimes N$ がある。これは
$\widehat{\phi} \circ (\mu_i \times 1) = \lambda_i \circ g_i$ を満たす一意な写像
$\widehat{\phi} : M \times N \to P$ を誘導する。この写像は $R$-双線形なので、
$R$-線形写像 $\phi : M \otimes N \to P$ を誘導する。次の可換図式から、
$$
\xymatrix{
M_i \times N \ar[r]^{g_i} \ar[d]^{\mu_i \times \text{id}} &
M_i \otimes N\ar[r]_{\text{id}} \ar[d]_{\lambda_i} &
M_i \otimes N \ar[d]_{\mu_i \otimes \text{id}} \ar[rd]^{\lambda_i} \\
M \times N \ar[r]^{\widehat{\phi}} &
P \ar[r]^{\psi} & M \otimes N \ar[r]^{\phi} & P
}
$$
$\psi\circ\widehat{\phi} = g$ が分かる。ここで
$g : M \times N \to M \otimes N$ は標準的な $R$-双線形写像である。したがって
$\psi\circ\phi$ は $M \otimes N$ 上の恒等写像である。右側の正方形と三角形から、
$\phi\circ\psi$ も $P$ 上の恒等写像である。
\end{proof}

\begin{lemma}
\label{algebra-lemma-tensor-product-exact}
\begin{align*}
M_1\xrightarrow{f} M_2\xrightarrow{g} M_3 \to 0
\end{align*}
を $R$-加群と準同型の完全列とし、$N$ を任意の $R$-加群とする。このとき列
\begin{equation}
\label{algebra-equation-2ndex}
M_1\otimes N\xrightarrow{f \otimes 1} M_2\otimes N \xrightarrow{g \otimes 1}
M_3\otimes N \to 0
\end{equation}
は完全である。言い換えれば、関手 $- \otimes_R N$ は{\it 右完全}である。
すなわち、もとの右完全列の各項とテンソル積をとると完全性が保たれる。
\end{lemma}

\begin{proof}
任意の $R$-加群 $P$ に対し、最初の完全列へ関手 $\Hom(-, \Hom(N, P))$ を
適用すると、補題 \ref{algebra-lemma-hom-exact} (1) により完全な列
$$
0 \to
\Hom(M_3, \Hom(N, P)) \to
\Hom(M_2, \Hom(N, P)) \to
\Hom(M_1, \Hom(N, P))
$$
を得る。補題 \ref{algebra-lemma-hom-from-tensor-product} により、これは列
$$
0 \to \Hom(M_3 \otimes N, P) \to
\Hom(M_2 \otimes N, P) \to \Hom(M_1 \otimes N, P)
$$
となり、したがってこれも完全である。補題 \ref{algebra-lemma-hom-exact} (1) を再び
用いると、所望の完全列を得る。
\end{proof}

\begin{remark}
\label{algebra-remark-tensor-product-not-exact}
しかし一般にテンソル積は完全列を保たない。言い換えれば、
$M_1 \to M_2 \to M_3$ が完全であっても、任意の $R$-加群 $N$ に対して
$M_1 \otimes N \to M_2 \otimes N \to M_3 \otimes N$ が完全であるとは限らない。
\end{remark}

\begin{example}
\label{algebra-example-tensor-product-not-exact}
$\mathbf{Z}$-加群の写像とみなした単射 $2 : \mathbf{Z}\to \mathbf{Z}$ を考える。
$N = \mathbf{Z}/2$ とする。このとき誘導される写像
$\mathbf{Z} \otimes \mathbf{Z}/2 \to \mathbf{Z} \otimes \mathbf{Z}/2$ は単射でない。
実際、$x \otimes y\in \mathbf{Z} \otimes \mathbf{Z}/2$ に対して
$$
(2 \otimes 1)(x \otimes y) = 2x \otimes y = x \otimes 2y = x \otimes 0 = 0
$$
である。したがって誘導写像は零写像だが、$\mathbf{Z} \otimes N\neq 0$ である。
\end{example}

\begin{remark}
\label{algebra-remark-flat-module}
$R$-加群 $N$ に対して、関手 $-\otimes_R N$ が完全、すなわち $N$ との
テンソル積がすべての完全列を保つならば、$N$ を{\it 平坦} $R$-加群という。
これについては後に Section \ref{algebra-section-flat} で論じる。
\end{remark}

\begin{lemma}
\label{algebra-lemma-tensor-finiteness}
$R$ を環、$M$ と $N$ を $R$-加群とする。
\begin{enumerate}
\item $N$ と $M$ が有限ならば、$M \otimes_R N$ も有限である。
\item $N$ と $M$ が有限表示ならば、$M \otimes_R N$ も有限表示である。
\end{enumerate}
\end{lemma}

\begin{proof}
$M$ が有限であると仮定し、表示 $0 \to K \to R^{\oplus n} \to M \to 0$ を選ぶ。
補題 \ref{algebra-lemma-tensor-product-exact} により、これは完全列
$K \otimes_R N \to N^{\oplus n} \to M \otimes_R N \to 0$ を与える。$N$ も有限ならば、
$M \otimes_R N$ は有限加群の商なので有限である。補題 \ref{algebra-lemma-extension} を
参照せよ。同様に $N$ と $M$ がともに有限表示ならば、$K$ は有限であり、
$M \otimes_R N$ は有限表示加群 $N^{\oplus n}$ を有限加群
$K \otimes_R N$ で割った商であることが分かる。したがって有限表示である。
補題 \ref{algebra-lemma-extension} を参照せよ。
\end{proof}

\begin{lemma}
\label{algebra-lemma-tensor-localization}
$M$ を $R$-加群とする。このとき $S^{-1}R$-加群 $S^{-1}M$ と
$S^{-1}R \otimes_R M$ は標準的に同型であり、標準同型
$f : S^{-1}R \otimes_R M \to S^{-1}M$ は
$$
f((a/s) \otimes m) = am/s, \forall a \in R, m \in M, s \in S
$$
で与えられる。
\end{lemma}

\begin{proof}
写像 $f' : S^{-1}R \times M \to S^{-1}M$ を $f'(a/s, m) = am/s$ で定めると、
これは明らかに双線形である。したがって普遍性により、この写像は補題の主張に
ある一意な $S^{-1}R$-加群準同型 $f : S^{-1}R \otimes_R M \to S^{-1}M$ を誘導する。
実際、$S^{-1}M$ のすべての元は $m/s$, $m\in M, s\in S$ の形であり、
$S^{-1}R \otimes_R M$ のすべての元は $1/s \otimes m$ の形である。後者を示すため、
$S^{-1}R \otimes_R M$ の元を
$$
\sum_k \frac{a_k}{s_k} \otimes m_k =
\sum_k \frac{a_k t_k}{s} \otimes m_k =
\frac{1}{s} \otimes \sum_k {a_k t_k}m_k = \frac{1}{s} \otimes m
$$
と書く。ここで $m = \sum_k {a_k t_k}m_k$ である。このとき $f$ が全射であることは
明らかである。また $f(\frac{1}{s} \otimes m) = m/s = 0$ ならば、$M$ において
$tm = 0$ を満たす $t \in S$ が存在する。このとき
$$
\frac{1}{s} \otimes m = \frac{1}{st} \otimes tm = \frac{1}{st} \otimes 0 = 0
$$
である。したがって $f$ は単射である。
\end{proof}

\begin{lemma}
\label{algebra-lemma-tensor-product-localization}
$M, N$ を $R$-加群とする。このとき
$$
f((m/s)\otimes(n/t)) = (m \otimes n)/st
$$
で与えられる標準的な $S^{-1}R$-加群同型
$f : S^{-1}M \otimes_{S^{-1}R}S^{-1}N \to S^{-1}(M \otimes_R N)$ が存在する。
\end{lemma}

\begin{proof}
補題 \ref{algebra-lemma-tensor-with-bimodule} と補題 \ref{algebra-lemma-tensor-localization} を繰り返し
用いると、$S^{-1}R$ が $(R, S^{-1}R)$-双加群であることから、これら二つの
$S^{-1}R$-加群が次のように同型であることが分かる。
\begin{align*}
S^{-1}(M \otimes_R N) & \cong S^{-1}R \otimes_R (M \otimes_R N)\\
 & \cong S^{-1}M \otimes_R N\\
 & \cong (S^{-1}M \otimes_{S^{-1}R}S^{-1}R)\otimes_R N\\
 & \cong S^{-1}M \otimes_{S^{-1}R}(S^{-1}R \otimes_R N)\\
 & \cong S^{-1}M \otimes_{S^{-1}R}S^{-1}N
\end{align*}
この同型が補題で述べたものに一致することは容易に分かる。
\end{proof}

\section{テンソル代数}
\label{algebra-section-tensor-algebra}

\noindent
$R$ を環、$M$ を $R$-加群とする。
$R$ 上の $M$ の {\it テンソル代数}を、非可換 $R$-代数
$$
\text{T}(M) = \text{T}_R(M) =
\bigoplus\nolimits_{n \geq 0} \text{T}^n(M)
$$
として定める。ここで
$\text{T}^0(M) = R$,
$\text{T}^1(M) = M$,
$\text{T}^2(M) = M \otimes_R M$,
$\text{T}^3(M) = M \otimes_R M \otimes_R M$ などとする。
積は、単純テンソル上で
$$
(x_1 \otimes x_2 \otimes \ldots \otimes x_n)
\cdot
(y_1 \otimes y_2 \otimes \ldots \otimes y_m)
=
x_1 \otimes x_2 \otimes \ldots \otimes x_n \otimes
y_1 \otimes y_2 \otimes \ldots \otimes y_m
$$
とし、これを線形に拡張することによって定める。

\medskip\noindent
$R$ 上の $M$ の {\it 外積代数 $\wedge(M)$} を、
元 $x \otimes x \in \text{T}^2(M)$ が生成する両側イデアルによる
$\text{T}(M)$ の商として定める。
$\wedge^n(M)$ における単純テンソル $x_1 \otimes \ldots \otimes x_n$
の像を $x_1 \wedge \ldots \wedge x_n$ と書く。
これらの元は $\wedge^n(M)$ を生成し、各 $x_i$ について $R$-線形であり、
二つの $x_i$ が等しいときは零である（すなわち
$x_1, x_2, \ldots, x_n$ の関数として交代的である）。
$\wedge(M)$ 上の積は次数付き可換である。すなわち、任意の
$x \in M$ と $y \in M$ は $x \wedge y = - y \wedge x$ を満たす。

\medskip\noindent
一例として、$M = Rx_1 \oplus \ldots \oplus Rx_n$ が有限自由加群である場合を考える。
このとき $\wedge(M)$ は $R$ 上自由であり、
$$
x_{i_1} \wedge \ldots \wedge x_{i_r}
$$
ただし $0 \leq r \leq n$ かつ
$1 \leq i_1 < i_2 < \ldots < i_r \leq n$ である元を基底にもつ。

\medskip\noindent
$R$ 上の $M$ の {\it 対称代数 $\text{Sym}(M)$} を、
元 $x \otimes y - y \otimes x \in \text{T}^2(M)$ が生成する両側イデアルによる
$\text{T}(M)$ の商として定める。
$\text{Sym}^n(M)$ における単純テンソル
$x_1 \otimes \ldots \otimes x_n$ の像を単に $x_1 \ldots x_n$ と書く。
これらの元は $\text{Sym}^n(M)$ を生成し、各 $x_i$ について $R$-線形である。
また、元の列 $x_1, \ldots, x_n$ が元の列
$x_1', \ldots, x_n'$ の置換ならば
$x_1 \ldots x_n = x_1' \ldots x_n'$ である。
したがって $\text{Sym}(M)$ は可換である。

\medskip\noindent
一例として、$M = Rx_1 \oplus \ldots \oplus Rx_n$ が有限自由加群である場合を考える。
このとき $\text{Sym}(M) = R[x_1, \ldots, x_n]$ は多項式代数である。

\begin{lemma}
\label{algebra-lemma-free-tensor-algebra}
$R$ を環、$M$ を $R$-加群とする。
$M$ が自由 $R$-加群ならば、各対称冪および外積冪も自由である。
\end{lemma}

\begin{proof}
省略する。ただし有限自由の場合については上記を参照せよ。
\end{proof}

\begin{lemma}
\label{algebra-lemma-presentation-sym-exterior}
$R$ を環とする。
$M_2 \to M_1 \to M \to 0$ を $R$-加群の完全列とする。
このとき完全列
$$
M_2 \otimes_R \text{Sym}^{n - 1}(M_1)
\to
\text{Sym}^n(M_1)
\to
\text{Sym}^n(M)
\to
0
$$
があり、同様に完全列
$$
M_2 \otimes_R \wedge^{n - 1}(M_1)
\to
\wedge^n(M_1)
\to
\wedge^n(M)
\to
0
$$
がある。
\end{lemma}

\begin{proof}
省略する。
\end{proof}

\begin{lemma}
\label{algebra-lemma-present-sym-wedge}
$R$ を環とする。
$M$ を $R$-加群とする。
$x_i$, $i \in I$ を、$R$-加群としての $M$ の所与の生成系とする。
$n \geq 2$ とする。このとき標準的な完全列
$$
\bigoplus_{1 \leq j_1 < j_2 \leq n}
\bigoplus_{i_1, i_2 \in I}
\text{T}^{n - 2}(M)
\oplus
\bigoplus_{1 \leq j_1 < j_2 \leq n}
\bigoplus_{i \in I}
\text{T}^{n - 2}(M)
\to
\text{T}^n(M)
\to
\wedge^n(M)
\to
0
$$
が存在する。第一の直和因子に属する単純テンソル
$m_1 \otimes \ldots \otimes m_{n - 2}$ は
\begin{align*}
\underbrace{
m_1 \otimes \ldots \otimes x_{i_1} \otimes \ldots
\otimes x_{i_2} \otimes \ldots \otimes m_{n - 2}
}_{\text{with } x_{i_1} \text{ and } x_{i_2}
\text{ occupying slots } j_1 \text{ and } j_2
\text{ in the tensor}} \\
+
\underbrace{
m_1 \otimes \ldots \otimes x_{i_2} \otimes \ldots
\otimes x_{i_1} \otimes \ldots \otimes m_{n - 2}
}_{\text{with } x_{i_2} \text{ and } x_{i_1}
\text{ occupying slots } j_1 \text{ and } j_2
\text{ in the tensor}}
\end{align*}
へ写り、第二の直和因子に属する
$m_1 \otimes \ldots \otimes m_{n - 2}$ は
$$
\underbrace{
m_1 \otimes \ldots \otimes x_i \otimes \ldots
\otimes x_i \otimes \ldots \otimes m_{n - 2}
}_{\text{with } x_{i} \text{ and } x_{i}
\text{ occupying slots } j_1 \text{ and } j_2
\text{ in the tensor}}
$$
へ写る。また、標準的な完全列
$$
\bigoplus_{1 \leq j_1 < j_2 \leq n}
\bigoplus_{i_1, i_2 \in I}
\text{T}^{n - 2}(M)
\to
\text{T}^n(M)
\to
\text{Sym}^n(M)
\to
0
$$
も存在し、単純テンソル $m_1 \otimes \ldots \otimes m_{n - 2}$ は
\begin{align*}
\underbrace{
m_1 \otimes \ldots \otimes x_{i_1} \otimes \ldots
\otimes x_{i_2} \otimes \ldots \otimes m_{n - 2}
}_{\text{with } x_{i_1} \text{ and } x_{i_2}
\text{ occupying slots } j_1 \text{ and } j_2
\text{ in the tensor}} \\
-
\underbrace{
m_1 \otimes \ldots \otimes x_{i_2} \otimes \ldots
\otimes x_{i_1} \otimes \ldots \otimes m_{n - 2}
}_{\text{with } x_{i_2} \text{ and } x_{i_1}
\text{ occupying slots } j_1 \text{ and } j_2
\text{ in the tensor}}
\end{align*}
へ写る。
\end{lemma}

\begin{proof}
省略する。
\end{proof}

\begin{lemma}
\label{algebra-lemma-present-wedge}
$A \to B$ を環準同型とし、$M$ を $B$-加群とする。
$n > 1$ とする。$A$-線形写像
$M \otimes_A \ldots \otimes_A M \to \wedge^n_B(M)$
の核は、$A$-加群として、$i \not = j$ かつ
$m_1, \ldots, m_n \in M$ を満たす $m_i = m_j$ なる元
$m_1 \otimes \ldots \otimes m_n$ と、$i \not = j$,
$m_1, \ldots, m_n \in M$, $b \in B$ を満たす元
$m_1 \otimes \ldots \otimes bm_i \otimes \ldots \otimes m_n -
m_1 \otimes \ldots \otimes bm_j \otimes \ldots \otimes m_n$
によって生成される。
\end{lemma}

\begin{proof}
省略する。
\end{proof}

\begin{lemma}
\label{algebra-lemma-colimit-tensor-algebra}
\begin{slogan}
テンソル代数をとる操作はフィルター余極限と可換する。
\end{slogan}
$R$ を環とし、$M_i$ を $R$-加群の有向系とする。このとき
$\colim_i \text{T}(M_i) = \text{T}(\colim_i M_i)$
であり、対称代数および外積代数についても同様である。
\end{lemma}

\begin{proof}
省略する。ヒント：補題 \ref{algebra-lemma-tensor-products-commute-with-limits} を適用せよ。
\end{proof}

\begin{lemma}
\label{algebra-lemma-tensor-algebra-localization}
$R$ を環、$S \subset R$ を乗法的集合とする。
このとき任意の $R$-加群 $M$ に対し
$S^{-1}T_R(M) = T_{S^{-1}R}(S^{-1}M)$ である。
対称代数および外積代数についても同様である。
\end{lemma}

\begin{proof}
省略する。ヒント：補題 \ref{algebra-lemma-tensor-product-localization} を適用せよ。
\end{proof}

\section{基底変換}
\label{algebra-section-base-change}

\noindent
代数における基底変換を次のように形式的に導入する。

\begin{definition}
\label{algebra-definition-base-change}
$\varphi : R \to S$ を環準同型、$M$ を $S$-加群とする。
$R \to R'$ を任意の環準同型とする。
$R \to R'$ による $\varphi$ の {\it 基底変換}とは、環準同型
$R' \to S \otimes_R R'$ のことである。この状況ではしばしば
$S' = S \otimes_R R'$ と書く。
$S$-加群 $M$ の {\it 基底変換}とは $S'$-加群
$M \otimes_R R'$ のことである。
\end{definition}

\noindent
ある変数族 $x_i$, $i \in I$ と多項式族
$f_j \in R[x_i]$, $j \in J$ に対して $S = R[x_i]/(f_j)$ ならば、
$S \otimes_R R' = R'[x_i]/(f'_j)$ である。ここで
$f'_j \in R'[x_i]$ は、$R \to R'$ が誘導する写像
$R[x_i] \to R'[x_i]$ による $f_j$ の像である。
この簡単な注意が基底変換を理解する鍵である。

\begin{lemma}
\label{algebra-lemma-base-change-finiteness}
$R \to S$ を環準同型、$M$ を $S$-加群とする。
$R \to R'$ を環準同型とし、$S' = S \otimes_R R'$ および
$M' = M \otimes_R R'$ をそれぞれの基底変換とする。
\begin{enumerate}
\item $M$ が有限 $S$-加群ならば、基底変換 $M'$ は有限 $S'$-加群である。
\item $M$ が有限表示 $S$-加群ならば、基底変換 $M'$ は有限表示
$S'$-加群である。
\item $R \to S$ が有限型ならば、基底変換 $R' \to S'$ も有限型である。
\item $R \to S$ が有限表示ならば、基底変換 $R' \to S'$ も有限表示である。
\end{enumerate}
\end{lemma}

\begin{proof}
(1) の証明。全射な $S$-線形写像
$S^{\oplus n} \to M \to 0$ をとる。
補題 \ref{algebra-lemma-flip-tensor-product} および
\ref{algebra-lemma-tensor-product-exact} により、$R^\prime$ とテンソル積を
とった結果は全射
${S^\prime}^{\oplus n} \to M^\prime \rightarrow 0$
となる。したがって $M^\prime$ は有限生成 $S^\prime$-加群である。
(2) の証明。表示
$S^{\oplus m} \to S^{\oplus n} \to M \to 0$ をとる。
補題 \ref{algebra-lemma-flip-tensor-product} および
\ref{algebra-lemma-tensor-product-exact} により、$R^\prime$ とテンソル積を
とった結果は $S^\prime$-加群 $M^\prime$ の有限表示
${S^\prime}^{\oplus m} \to {S^\prime}^{\oplus n} \to M^\prime \to 0$
を与える。(3) の証明。仮定により $I$ を有限にとれるので、補題の直前の
注意から従う。(4) の証明。仮定により $I$ と $J$ を有限にとれるので、
補題の直前の注意から従う。
\end{proof}

\noindent
$\varphi : R \to S$ を環準同型とする。$S$-加群 $N$ が与えられると、
規則 $r \cdot n = \varphi(r)n$ により $R$-加群 $N_R$ を得る。
これは $N$ の $R$ への {\it 制限}と呼ばれることがある。

\begin{lemma}
\label{algebra-lemma-adjoint-tensor-restrict}
$R \to S$ を環準同型とする。関手
$\text{Mod}_S \to \text{Mod}_R$, $N \mapsto N_R$（制限）と
$\text{Mod}_R \to \text{Mod}_S$, $M \mapsto M \otimes_R S$
（基底変換）は随伴関手である。式で書けば
$$
\Hom_R(M, N_R) = \Hom_S(M \otimes_R S, N)
$$
である。
\end{lemma}

\begin{proof}
$\alpha : M \to N_R$ が $R$-加群写像ならば、規則
$\alpha'(m \otimes s) = s\alpha(m)$ により
$\alpha' : M \otimes_R S \to N$ を定める。
$\beta : M \otimes_R S \to N$ が $S$-加群写像ならば、規則
$\beta'(m) = \beta(m \otimes 1)$ により $\beta' : M \to N_R$ を定める。
これらの構成が互いに逆であることの確認は省略する。
\end{proof}

\noindent
上の補題は、制限が左随伴、すなわち基底変換をもつことを述べている。
制限は右随伴ももつ。

\begin{lemma}
\label{algebra-lemma-adjoint-hom-restrict}
$R \to S$ を環準同型とする。関手
$\text{Mod}_S \to \text{Mod}_R$, $N \mapsto N_R$（制限）と
$\text{Mod}_R \to \text{Mod}_S$, $M \mapsto \Hom_R(S, M)$
は随伴関手である。式で書けば
$$
\Hom_R(N_R, M) = \Hom_S(N, \Hom_R(S, M))
$$
である。
\end{lemma}

\begin{proof}
$\alpha : N_R \to M$ が $R$-加群写像ならば、規則
$\alpha'(n) = (s \mapsto \alpha(sn))$ により
$\alpha' : N \to \Hom_R(S, M)$ を定める。
$\beta : N \to \Hom_R(S, M)$ が $S$-加群写像ならば、規則
$\beta'(n) = \beta(n)(1)$ により $\beta' : N_R \to M$ を定める。
これらの構成が互いに逆であることの確認は省略する。
\end{proof}

\begin{lemma}
\label{algebra-lemma-hom-from-tensor-product-variant}
$R \to S$ を環準同型とする。$S$-加群 $M, N$ と
$R$-加群 $P$ が与えられたとき、
$$
\Hom_R(M \otimes_S N, P) = \Hom_S(M, \Hom_R(N, P))
$$
が成り立つ。
\end{lemma}

\begin{proof}
これは直接証明できるが、補題 \ref{algebra-lemma-adjoint-hom-restrict} および
\ref{algebra-lemma-hom-from-tensor-product} の帰結でもある。実際、
\begin{align*}
\Hom_R(M \otimes_S N, P)
& =
\Hom_S(M \otimes_S N, \Hom_R(S, P)) \\
& =
\Hom_S(M, \Hom_S(N, \Hom_R(S, P))) \\
& =
\Hom_S(M, \Hom_R(N, P))
\end{align*}
となり、所望の結果を得る。
\end{proof}

\section{雑録}
\label{algebra-section-miscellany}

\noindent
この節の証明では、基本概念を扱う節、すなわち
Section \ref{algebra-section-rings-basic} の結果以外を参照してはならない。

\begin{lemma}
\label{algebra-lemma-product-ideals-in-prime}
$R$ を環、$I$ と $J$ を二つのイデアル、$\mathfrak p$ を積 $IJ$ を含む
素イデアルとする。このとき $\mathfrak{p}$ は $I$ または $J$ を含む。
\end{lemma}

\begin{proof}
そうでないと仮定し、$x \in I \setminus \mathfrak p$ および
$y \in J \setminus \mathfrak p$ をとる。その積は
$IJ \subset \mathfrak p$ の元であり、$\mathfrak p$ が素イデアルであるという
仮定に矛盾する。
\end{proof}

\begin{lemma}[素イデアル回避]
\label{algebra-lemma-silly}
\begin{slogan}
1. アフィンスキームの有限個の点がある開部分集合に含まれるならば、
それらはさらに小さい単項開部分集合に含まれる。
2. アフィンスキームの所与の有限集合の近傍の中で、アフィン開集合は共終である。
\end{slogan}
$R$ を環とする。$I_i \subset R$, $i = 1, \ldots, r$ および
$J \subset R$ をイデアルとする。次を仮定する。
\begin{enumerate}
\item $i = 1, \ldots, r$ に対して $J \not\subset I_i$ である。
\item $I_i$ のうち二つを除くすべてが素イデアルである。
\end{enumerate}
このとき、すべての $i$ に対して $x\not\in I_i$ を満たす
$x \in J$ が存在する。
\end{lemma}

\begin{proof}
$r = 1$ のとき結果は正しい。$r = 2$ ならば、
$x \not \in I_1$ および $y \not \in I_2$ を満たす $x, y \in J$ をとる。
$x \in I_2$ かつ $y \in I_1$ でない限り、これでよい。
その場合、元 $x + y$ は $I_1$ に属しえず（属すれば
$x + y - y \in I_1$ となる）、同様に $I_2$ にも属しえない。

\medskip\noindent
$r \geq 3$ とし、$r - 1$ に対して結果が成り立つと仮定する。
番号を付け替えて $I_r$ が素イデアルであるとしてよい。
また、$I_i$ の間に包含関係がないとしてよい。
すべての $i = 1, \ldots, r - 1$ に対して $x \not \in I_i$ を満たす
$x \in J$ をとる。$x \not\in I_r$ ならばこれでよいので、
$x \in I_r$ と仮定する。もし
$J I_1 \ldots I_{r - 1} \subset I_r$ ならば、補題
\ref{algebra-lemma-product-ideals-in-prime} により $J \subset I_r$ となり、矛盾する。
$y \not \in I_r$ を満たす
$y \in J I_1 \ldots I_{r - 1}$ をとれば、$x + y$ が所望の元である。
\end{proof}

\begin{lemma}
\label{algebra-lemma-silly-silly}
$R$ を環とする。$x \in R$、$I \subset R$ をイデアルとし、
$\mathfrak p_i$, $i = 1, \ldots, r$ を素イデアルとする。
$i = 1, \ldots, r$ に対して
$x + I \not \subset \mathfrak p_i$ であると仮定する。
このとき、すべての $i$ に対して
$x + y \not \in \mathfrak p_i$ を満たす $y \in I$ が存在する。
\end{lemma}

\begin{proof}
$\mathfrak p_i$ の間に包含関係がないとしてよい。
順序を付け替えて、$i < s$ では $x \not \in \mathfrak p_i$、
$i \geq s$ では $x \in \mathfrak p_i$ としてよい。
$s = r + 1$ ならばこれでよい。そうでなければ
$y \not \in \mathfrak p_s$ を満たす $y \in I$ をとれる。
$f \not \in \mathfrak p_s$ を満たす
$f \in \bigcap_{i < s} \mathfrak p_i$ を選ぶ。
このとき $x + fy$ は
$\mathfrak p_1, \ldots, \mathfrak p_s$ のいずれにも属さない。
したがって $s$ に関する帰納法により結論を得る。
\end{proof}

\begin{lemma}[中国剰余定理]
\label{algebra-lemma-chinese-remainder}
$R$ を環とする。
\begin{enumerate}
\item $I_1, \ldots, I_r$ が、$a \not = b$ のとき
$I_a + I_b = R$ を満たすイデアルならば、
$I_1 \cap \ldots \cap I_r = I_1I_2\ldots I_r$ であり、
$R/(I_1I_2\ldots I_r) \cong R/I_1 \times \ldots \times R/I_r$ である。
\item $\mathfrak m_1, \ldots, \mathfrak m_r$ が互いに異なる極大イデアルならば、
$a \not = b$ に対して $\mathfrak m_a + \mathfrak m_b = R$ であり、
上記が適用される。
\end{enumerate}
\end{lemma}

\begin{proof}
まず $I_1 \cap \ldots \cap I_r = I_1 \ldots I_r$ を示す。
これにより、誘導される環準同型
$R/(I_1 \ldots I_r) \rightarrow R/I_1 \times \ldots \times R/I_r$
の単射性も従う。
イデアルは任意の環の元との積について閉じているので、包含
$I_1 \cap \ldots \cap I_r \supset I_1 \ldots I_r$ は常に成り立つ。
逆の包含を示すため、イデアル
$$
I_1 \ldots \hat I_i \ldots I_r,\quad i = 1, \ldots, r
$$
が環 $R$ を生成することを主張する。これを $r$ に関する帰納法で示す。
$r = 2$ のときは成り立つ。$r > 2$ ならば、
$R$ はイデアル
$I_1 \ldots \hat I_i \ldots I_{r - 1}$, $i = 1, \ldots, r - 1$
の和である。したがって $I_r$ はイデアル
$I_1 \ldots \hat I_i \ldots I_r$, $i = 1, \ldots, r - 1$
の和である。イデアルの順序を逆にして同じ議論を適用すれば、
$I_1$ はイデアル
$I_1 \ldots \hat I_i \ldots I_r$, $i = 2, \ldots, r$
の和である。仮定により $R = I_1 + I_r$ なので、
$R$ は上に表示したイデアルの和である。
したがって、和が 1 となる元
$a_i \in I_1 \ldots \hat I_i \ldots I_r$
をとれる。この等式に $I_1 \cap \ldots \cap I_r$ の元を掛ければ
逆の包含を得る。
残るのは、標準写像
$R/(I_1 \ldots I_r) \rightarrow R/I_1 \times \ldots \times R/I_r$
が全射であることを示すことである。そのため、上で得た等式
$1 = \sum_{i=1}^r a_i$ に対するこの写像の作用を考える。
一方で環準同型は 1 を 1 へ写し、他方で任意の $a_i$ の像は
$j \neq i$ に対して $R/I_j$ で零である。
したがって $R/I_i$ における $a_i$ の像は単位元である。
ゆえに任意の元
$(\bar{b_1}, \ldots, \bar{b_r}) \in R/I_1 \times \ldots \times R/I_r$
に対し、元 $\sum_{i=1}^r a_i \cdot b_i$ が $R$ における逆像である。

\medskip\noindent
(2) を示す。互いに異なる極大イデアルの定義そのものから、
$a \neq b$ に対して
$\mathfrak{m}_a + \mathfrak{m}_b = R$ であるので、上記が適用される。
\end{proof}

\begin{lemma}
\label{algebra-lemma-matrix-left-inverse}
$R$ を環とする。$n \geq m$ とし、$A$ を $R$ に係数をもつ
$n \times m$ 行列とする。$J \subset R$ を $A$ の
$m \times m$ 小行列式が生成するイデアルとする。
\begin{enumerate}
\item 任意の $f \in J$ に対し、$BA = f 1_{m \times m}$ を満たす
$m \times n$ 行列 $B$ が存在する。
\item $f \in R$ であり、ある $m \times n$ 行列 $B$ に対して
$BA = f 1_{m \times m}$ ならば、$f^m \in J$ である。
\end{enumerate}
\end{lemma}

\begin{proof}
$|I| = m$ を満たす $I \subset \{1, \ldots, n\}$ に対し、
射影
$$
R^{\oplus n} = \bigoplus\nolimits_{i \in \{1, \ldots, n\}} R
\longrightarrow \bigoplus\nolimits_{i \in I} R
$$
の $m \times n$ 行列を $E_I$ と書き、
$A_I = E_I A$ とおく。すなわち $A_I$ は、$A$ のうち添字が
$I$ に属する行からなる $m \times m$ 行列である。
$B_I$ を $A_I$ の随伴行列（余因子行列の転置）、すなわち
$A_I B_I = B_I A_I = \det(A_I) 1_{m \times m}$ を満たす行列とする。
$A$ の $m \times m$ 小行列式は、
$|I| = m$ を満たすすべての $I \subset \{1, \ldots, n\}$ に対する
行列式 $\det A_I$ である。$f \in J$ ならば、ある $c_I \in R$ により
$f = \sum c_I \det(A_I)$ と書ける。
$B = \sum c_I B_I E_I$ とおけば (1) が成り立つことが分かる。

\medskip\noindent
$f 1_{m \times m} = BA$ ならば、コーシー・ビネの公式
(\ref{algebra-item-cauchy-binet}) により
$f^m = \sum b_I \det(A_I)$ である。ここで $b_I$ は、
$B$ のうち添字が $I$ に属する列からなる $m \times m$ 行列の行列式である。
\end{proof}

\begin{lemma}
\label{algebra-lemma-matrix-right-inverse}
$R$ を環とする。$n \geq m$ とし、$A = (a_{ij})$ を $R$ に係数をもつ
$n \times m$ 行列で、ブロック表示が
$$
A =
\left(
\begin{matrix}
A_1 \\
A_2
\end{matrix}
\right)
$$
であるものとする。ここで $A_1$ の大きさは $m \times m$ である。
$B$ を $A_1$ の随伴行列（余因子行列の転置）とする。このとき
$$
AB = 
\left(
\begin{matrix}
f 1_{m \times m} \\
C
\end{matrix}
\right)
$$
である。ここで $f = \det(A_1)$ であり、$c_{ij}$ は符号を除いて、
$A$ の行 $1, \ldots, \hat j, \ldots, m, i$ に対応する
$m \times m$ 小行列式である。
\end{lemma}

\begin{proof}
随伴行列は $A_1B = B A_1 = f$ を満たすので、$AB$ の表示の
第一ブロックは正しい。また
$$
c_{ij} = \sum\nolimits_k a_{ik}b_{kj} = \sum (-1)^{j + k}a_{ik} \det(A_1^{jk})
$$
である。ここで $A_1^{ij}$ は、$A_1$ から第 $j$ 行と第 $k$ 列を
取り除いた行列を表す。
% P01-E1280
最後の式は、補題の主張に現れる行列の行列式を行について展開したものである。
\end{proof}

\begin{lemma}
\label{algebra-lemma-map-cannot-be-injective}
\begin{slogan}
有限自由加群の間の写像は、始域の階数が終域の階数より大きければ単射になりえない。
\end{slogan}
$R$ を零でない環とする。$n \geq 1$ とし、$M$ を
$< n$ 個の元で生成される $R$-加群とする。このとき任意の
$R$-加群写像 $f : R^{\oplus n} \to M$ は零でない核をもつ。
\end{lemma}

\begin{proof}
全射 $R^{\oplus n - 1} \to M$ を選ぶ。
写像 $f$ を写像 $f' : R^{\oplus n} \to R^{\oplus n - 1}$ へ持ち上げられる
（補題 \ref{algebra-lemma-lift-map}）。
$f'$ が零でない核をもつことを示せば十分である。
写像 $f' : R^{\oplus n} \to R^{\oplus n - 1}$ は行列
$A = (a_{ij})$ で与えられる。$a_{ij}$ の一つ、たとえば
$a = a_{ij}$ が冪零でないならば、$R$ を局所化 $R_a$ で置き換えて、
$a_{ij}$ が単元であるとしてよい。実際、局所化は完全なので、
局所化後に零でない核が見つかれば、もとの写像も零でない核をもつ。
命題 \ref{algebra-proposition-localization-exact} を参照せよ。
この場合、$R^{\oplus n}$ と $R^{\oplus n - 1}$ の双方で基底を取り替え、
$$
A =
\left(
\begin{matrix}
1 & 0 & 0 & \ldots \\
0 & a_{22} & a_{23} & \ldots \\
0 & a_{32} & \ldots \\
\ldots & \ldots
\end{matrix}
\right)
$$
の場合に帰着できる。したがって、この場合は $n$ に関する帰納法で結論を得る。
そうでなければ各 $a_{ij}$ は冪零である。
$I = (a_{ij}) \subset R$ とおく。ある $m \geq 0$ に対して
$I^{m + 1} = 0$ であることに注意する。
$I^m \not = 0$ を満たす最大の整数 $m$ をとる。
すると $(I^m)^{\oplus n}$ が写像の核に含まれるので、結論を得る。
\end{proof}

\begin{lemma}
\label{algebra-lemma-rank}
\begin{slogan}
有限自由加群の階数は良定義である。
\end{slogan}
$R$ を零でない環とし、$n, m \geq 0$ を整数とする。
$R^{\oplus n}$ と $R^{\oplus m}$ が $R$-加群として同型ならば
$n = m$ である。
\end{lemma}

\begin{proof}
補題 \ref{algebra-lemma-map-cannot-be-injective} から直ちに従う。
\end{proof}

\section{ケイリー・ハミルトン}
\label{algebra-section-cayley-hamilton}

\begin{lemma}
\label{algebra-lemma-charpoly}
$R$ を環とする。$A = (a_{ij})$ を $R$ に係数をもつ
$n \times n$ 行列とする。$P(x) \in R[x]$ を $A$ の特性多項式
（$\det(x\text{id}_{n \times n} - A)$ として定義する）とする。
このとき $\text{Mat}(n \times n, R)$ において $P(A) = 0$ である。
\end{lemma}

\begin{proof}
いくつかの段階を経て、この問題を線形代数でよく知られた
ケイリー・ハミルトンの定理に帰着する。
\begin{enumerate}
\item $\phi :S \rightarrow R$ が環準同型であり、$b_{ij}$ がこの写像による
$a_{ij}$ の逆像ならば、$\phi$ は環準同型なので、$S$ と $(b_{ij})$ に
対して主張を示せば十分である。
\item $\psi :R \hookrightarrow S$ が単射な環準同型ならば、
$S$ の元とみなした $a_{ij}$ と $S$ に対して結果を示せば十分である。
\item したがって、まず多項式環の場合
$R = \mathbf{Z}[X_{ij}]$, $a_{ij} = X_{ij}$ に帰着し、さらに
$R = \mathbf{Q}(X_{ij})$ の場合へ帰着できる。最後の体上では
ケイリー・ハミルトンの定理を適用できる。
\end{enumerate}
\end{proof}

\begin{lemma}
\label{algebra-lemma-charpoly-module}
$R$ を環とする。
$M$ を有限 $R$-加群とする。
$\varphi : M \to M$ を自己準同型とする。
このとき、$M$ の自己準同型として $P(\varphi) = 0$ を満たす
モニック多項式 $P \in R[T]$ が存在する。
\end{lemma}

\begin{proof}
ある生成元 $x_i \in M$ に対し
$(a_1, \ldots, a_n) \mapsto \sum a_ix_i$ で与えられる全射な
$R$-加群写像 $R^{\oplus n} \to M$ を選ぶ。
$\varphi(x_i) = \sum a_{ij} x_j$ を満たす
$(a_{i1}, \ldots, a_{in}) \in R^{\oplus n}$ を選ぶ。
言い換えれば、$A = (a_{ij})$ として図式
$$
\xymatrix{
R^{\oplus n} \ar[d]_A \ar[r] & M \ar[d]^\varphi \\
R^{\oplus n} \ar[r] & M
}
$$
は可換である。補題 \ref{algebra-lemma-charpoly} により、
$P(A) = 0$ を満たすモニック多項式 $P$ が存在する。
したがって $P(\varphi) = 0$ である。
\end{proof}

\begin{lemma}
\label{algebra-lemma-charpoly-module-ideal}
$R$ を環、$I \subset R$ をイデアルとする。
$M$ を有限 $R$-加群とする。
$\varphi(M) \subset IM$ を満たす自己準同型
$\varphi : M \to M$ とする。
このとき、$a_j \in I^j$ であり、$M$ の自己準同型として
$P(\varphi) = 0$ を満たすモニック多項式
$P = t^n + a_1 t^{n - 1} + \ldots + a_n \in R[T]$
が存在する。
\end{lemma}

\begin{proof}
ある生成元 $x_i \in M$ に対し
$(a_1, \ldots, a_n) \mapsto \sum a_ix_i$ で与えられる全射な
$R$-加群写像 $R^{\oplus n} \to M$ を選ぶ。
$\varphi(x_i) = \sum a_{ij} x_j$ を満たす
$(a_{i1}, \ldots, a_{in}) \in I^{\oplus n}$ を選ぶ。
言い換えれば、$A = (a_{ij})$ として図式
$$
\xymatrix{
R^{\oplus n} \ar[d]_A \ar[r] & M \ar[d]^\varphi \\
I^{\oplus n} \ar[r] & M
}
$$
は可換である。補題 \ref{algebra-lemma-charpoly} により、多項式
$P(t) = \det(t\text{id}_{n \times n} - A)$ は所望の性質をすべてもつ。
\end{proof}

\noindent
楽しい応用例として、次の驚くべき補題を証明する。

\begin{lemma}
\label{algebra-lemma-fun}
$R$ を環とする。
$M$ を有限 $R$-加群とする。
$\varphi : M \to M$ を全射な $R$-加群写像とする。
このとき $\varphi$ は同型である。
\end{lemma}

\begin{proof}[第一の証明]
$R' = R[x]$ と書き、$x$ が $\varphi$ として作用することにより
$M$ を有限 $R'$-加群とみなす。$I = (x) \subset R'$ とおく。
$\varphi$ が全射であるという仮定から $IM = M$ である。
したがって、$R'$-加群としての $M$、イデアル $I$、自己準同型
$\text{id}_M$ に補題 \ref{algebra-lemma-charpoly-module-ideal} を適用できる。
その結果、$a_j \in I$ を満たす
$(1 + a_1 + \ldots + a_n)\text{id}_M = 0$ を得る。
ある $b_j(x) \in R[x]$ により $a_j = b_j(x)x$ と書く。
これを $\varphi$ の言葉に戻すと
$\text{id}_M = -(\sum_{j = 1, \ldots, n} b_j(\varphi)) \varphi$
となり、したがって $\varphi$ は可逆である。
\end{proof}

\begin{proof}[第二の証明]
$R$ 上の $M$ の生成元の個数について帰納法を行う。
$M$ が一つの元で生成されるならば、あるイデアル $I \subset R$ に対し
$M \cong R/I$ である。この場合、$R$ を $R/I$ で置き換えて
$M = R$ としてよい。このとき $\varphi : R \to R$ は
ある元 $r \in R$ による $M$ 上の乗法で与えられる。
$\varphi$ の全射性により、$\varphi$ は $1$ を像にもたなければならないので、
$r$ は可逆となり、したがって $\varphi$ は可逆である。

\medskip\noindent
$n - 1$ 個の元で生成される加群に対して補題を証明済みとし、
$n$ 個の元で生成される加群 $M$ を考える。$A$ を環 $R[t]$ とし、
$t$ が $\varphi$ として作用することにより $M$ を $A$-加群とみなす。
$M$ は $R$ 上有限なので $R[t]$ 上も有限である。
また、示そうとしている $\varphi$ の単射性は集合論的な性質なので、
$A$ 上の自己準同型 $t : M \to M$ の単射性を示してもよい。
これで、自己準同型が基礎環の元による乗法である場合に問題を帰着した。
$M$ の最初の $n - 1$ 個の生成元が生成する部分 $A$-加群を
$M' \subset M$ と書き、図式
$$
\xymatrix{
0 \ar[r] & M' \ar[r]\ar[d]^{\varphi\mid_{M'}} & M\ar[d]^\varphi \ar[r] &
M/M' \ar[d]^{\varphi \bmod M'} \ar[r] & 0 \\
0 \ar[r] & M' \ar[r]                                  & M \ar[r]
           & M/M' \ar[r]                                  & 0,
}
$$
を考える。$\varphi$ は基礎環の元による乗法であり、$M'$ と $M/M'$ は
それ自体 $A$-加群なので、$\varphi$ の $M'$ への制限と
$\varphi$ が商 $M/M'$ 上に誘導する写像はいずれも良定義である。
$n = 1$ の場合により、写像 $M/M' \to M/M'$ は同型である。
図式追跡により $\varphi|_{M'}$ は全射であり、帰納法により
$\varphi|_{M'}$ は同型である。したがって蛇の補題により中央の列も同型となる。
\end{proof}

\section{環のスペクトル}
\label{algebra-section-spectrum-ring}

\noindent
環のスペクトルを位相空間とみなしたものは代数の章に属し、
アフィンスキームはスキームの章に属するものと、ここでは便宜的に定める。

\begin{definition}
\label{algebra-definition-spectrum-ring}
$R$ を環とする。
\begin{enumerate}
\item $R$ の {\it スペクトル}とは、$R$ の素イデアル全体の集合である。
通常これを $\Spec(R)$ と書く。
\item 部分集合 $T \subset R$ が与えられたとき、
$V(T) \subset \Spec(R)$ を $T$ を含む素イデアル全体、すなわち
$V(T) = \{ \mathfrak p \in
\Spec(R) \mid \forall f\in T, f\in \mathfrak p\}$ とする。
\item 元 $f \in R$ が与えられたとき、
$D(f) \subset \Spec(R)$ を $f$ を含まない素イデアル全体とする。
\end{enumerate}
\end{definition}

\begin{lemma}
\label{algebra-lemma-Zariski-topology}
$R$ を環とする。
\begin{enumerate}
\item 環 $R$ のスペクトルが空であることと、$R$ が零環であることは同値である。
\item 零でない任意の環は極大イデアルをもつ。
\item 零でない任意の環は極小素イデアルをもつ。
\item イデアル $I \subset R$ と素イデアル
$I \subset \mathfrak p$ が与えられたとき、$\mathfrak q$ が
$I$ 上極小であるような素イデアル
$I \subset \mathfrak q \subset \mathfrak p$ が存在する。
\item $T \subset R$ であり、$(T)$ が $R$ において $T$ が生成する
イデアルならば、$V((T)) = V(T)$ である。
\item $I$ がイデアルで $\sqrt{I}$ がその根基ならば、基本概念
(\ref{algebra-item-radical-ideal}) を参照して、$V(I) = V(\sqrt{I})$ である。
\item $R$ のイデアル $I$ に対して
$\sqrt{I} = \bigcap_{I \subset \mathfrak p} \mathfrak p$ である。
\item $I$ がイデアルならば、$V(I) = \emptyset$ であることと
$I$ が単位イデアルであることは同値である。
\item $I$, $J$ が $R$ のイデアルならば
$V(I) \cup V(J) = V(I \cap J)$ である。
\item $(I_a)_{a\in A}$ が $R$ のイデアル族ならば
$\bigcap_{a\in A} V(I_a) = V(\bigcup_{a\in A} I_a)$ である。
\item $f \in R$ ならば $D(f) \amalg V(f) = \Spec(R)$ である。
\item $f \in R$ ならば、$D(f) = \emptyset$ であることと
$f$ が冪零であることは同値である。
\item ある単元 $u \in R$ に対して $f = u f'$ ならば
$D(f) = D(f')$ である。
\item $I \subset R$ がイデアルであり、$\mathfrak p$ が
$\mathfrak p \not\in V(I)$ を満たす $R$ の素イデアルならば、
$\mathfrak p \in D(f)$ かつ
$D(f) \cap V(I) = \emptyset$ を満たす $f \in R$ が存在する。
\item $f, g \in R$ ならば $D(fg) = D(f) \cap D(g)$ である。
\item $i \in I$ に対して $f_i \in R$ ならば、
$\bigcup_{i\in I} D(f_i)$ は $\Spec(R)$ における
$V(\{f_i \}_{i\in I})$ の補集合である。
\item $f \in R$ かつ $D(f) = \Spec(R)$ ならば、$f$ は単元である。
\end{enumerate}
\end{lemma}

\begin{proof}
各部分を対応する項目で扱う。
\begin{enumerate}
\item (2) または (3) から直ちに従う。
\item $\mathfrak{A}$ を $R$ の真のイデアル全体の集合とする。
この集合を包含関係で順序づける。$(0) \in \mathfrak{A}$ は真の
イデアルなので、この集合は空でない。$A$ を $\mathfrak A$ の全順序部分集合とする。
このとき $\bigcup_{I \in A} I$ は実際イデアルである。
すべての $I \in A$ に対して 1 $\notin I$ なので、この和集合は 1 を含まず、
したがって真である。
% P01-E1281
ゆえに $\bigcup_{I \in A} I$ は $\mathfrak{A}$ に属し、集合 $A$ の上界である。
よってツォルンの補題により $\mathfrak{A}$ は極大元をもち、これが求める
極大イデアルである。
\item $R$ は零でないので、素イデアルでもある極大イデアルをもつ。
したがって $R$ の素イデアル全体の集合 $\mathfrak{A}$ は空でない。
$\mathfrak{A}$ を逆包含で順序づける。$A$ を $\mathfrak{A}$ の
全順序部分集合とする。$J = \bigcap_{I \in A} I$ が実際イデアルであることは
ほぼ明らかである。しかし、これが素であることはそれほど明らかでない。
$xy \in J$ とする。このときすべての $I \in A$ に対して $xy \in I$ である。
$B = \{I \in A | y \in I\}$ とし、$K = \bigcap_{I \in B} I$ とおく。
$A$ は全順序づけられているので、$K = J$（このとき $y \in J$ なので結論を得る）
であるか、または $K \supset J$ であって、$K$ に真に含まれるすべての
$I \in A$ に対して $y \notin I$ である。後者の場合、それらは素なので、
すべてのそのような $I, x \in I$ である。したがって $x \in J$ である。
いずれの場合も所望どおり $J$ は素である。よってツォルンの補題から極大元を得るが、
この場合それは極小素イデアルである。
\item $\mathfrak{p}$ に含まれ、かつ $I$ を含む素イデアルだけを
考える点を除けば、(3) とまったく同じ議論である。
\item $(T)$ は $T$ を含む最小のイデアルである。したがって、あるイデアル
$I$ に対して $T \subset I$ ならば $(T) \subset I$ でもある。
ゆえに $I \in V(T)$ ならば $I \in V((T))$ でもある。
逆の包含は明らかである。
\item $I \subset \sqrt{I}, V(\sqrt{I}) \subset V(I)$ である。
ここで $\mathfrak{p} \in V(I)$ とし、$x \in \sqrt{I}$ とする。
ある $n$ に対して $x^n \in I$ である。したがって
$x^n \in \mathfrak{p}$ である。$\mathfrak{p}$ は素なので、単純な
帰納法により $x \in \mathfrak{p}$ を得る。ゆえに
$\sqrt{I} \subset \mathfrak{p}$ であり、
$\mathfrak{p} \in V(\sqrt{I})$ である。
\item $f \in R \setminus \sqrt{I}$ とする。このときすべての $n$ に対して
$f^n \notin I$ である。したがって
$S = \{1, f, f^2, \ldots\}$ は $0$ を含まない乗法的集合である。
$S^{-1}I$ を含む素イデアル
$\bar{\mathfrak{p}} \subset S^{-1}R$ をとる。
$\bar{\mathfrak{p}}$ の $R$ における引き戻し $\mathfrak{p}$ は、
$I$ を含み $S$ と交わらない素イデアルである。
これにより $\bigcap_{I \subset \mathfrak p} \mathfrak p \subset \sqrt{I}$ が分かる。
一方、$a \in \sqrt{I}$ ならば、ある $n$ に対して $a^n \in I$ である。
したがって $I \subset \mathfrak{p}$ ならば $a^n \in \mathfrak{p}$ である。
$\mathfrak{p}$ は素なので $a \in \mathfrak{p}$ である。以上で等式が示された。
\item $I$ が単位イデアルでないことと、$I$ がある極大イデアルに含まれることは
同値であり（これを見るには環 $R/I$ に (2) を適用せよ）、その極大イデアルは素である。
\item $\mathfrak{p} \in V(I) \cup V(J)$ ならば、
$I \subset \mathfrak{p}$ または $J \subset \mathfrak{p}$ であり、
したがって $I \cap J \subset \mathfrak{p}$ である。
逆に $I \cap J \subset \mathfrak{p}$ ならば
$IJ \subset \mathfrak{p}$ であり、$\mathfrak{p}$ は素なので
$I$ または $J$ が $\mathfrak{p}$ に含まれる。
% P01-E1282
\item $\mathfrak{p} \in \bigcap_{a \in A} V(I_a) \Leftrightarrow
I_a \subset \mathfrak{p}, \forall a \in A \Leftrightarrow
\mathfrak{p} \in V(\bigcup_{a\in A} I_a)$
\item $\mathfrak{p}$ が素イデアルで $f \in R$ ならば、
$f \in \mathfrak{p}$ または $f \notin \mathfrak{p}$ のいずれか一方が
成り立つ。これが直和記号の述べていることである。
\item $a \in R$ が冪零ならば、ある $n$ に対して $a^n = 0$ である。
したがって任意の素イデアルに対して $a^n \in \mathfrak{p}$ である。
帰納法により $a \in \mathfrak{p}$ であり、$D(a) = \emptyset$ である。
逆に (7) で示したとおり、$a \in R$ が冪零でなければ、
それを含まない素イデアルが存在する。
\item $u$ は可逆なので
$f \in \mathfrak{p} \Leftrightarrow uf \in \mathfrak{p}$ である。
\item $\mathfrak{p} \notin V(I)$ ならば
$\exists f \in I \setminus \mathfrak{p}$ である。
このとき $f \notin \mathfrak{p}$ なので $\mathfrak{p} \in D(f)$ である。
また $\mathfrak{q} \in D(f)$ ならば $f \notin \mathfrak{q}$ であり、
したがって $I$ は $\mathfrak{q}$ に含まれない。
ゆえに $D(f) \cap V(I) = \emptyset$ である。
\item $fg \in \mathfrak{p}$ ならば $f \in \mathfrak{p}$ または
$g \in \mathfrak{p}$ である。したがって
$f \notin \mathfrak{p}$ かつ $g \notin \mathfrak{p}$ ならば
$fg \notin \mathfrak{p}$ である。
$\mathfrak{p}$ はイデアルなので、$fg \notin \mathfrak{p}$ ならば
$f \notin \mathfrak{p}$ かつ $g \notin \mathfrak{p}$ である。
\item $\mathfrak{p} \in \bigcup_{i \in I} D(f_i) \Leftrightarrow \exists i \in
I, f_i \notin \mathfrak{p} \Leftrightarrow \mathfrak{p} \in \Spec(R)
\setminus V(\{f_i\}_{i \in I})$
\item $D(f) = \Spec(R)$ ならば $V(f) = \emptyset$ であり、
したがって $fR = R$ なので $f$ は単元である。
\end{enumerate}
\end{proof}

\noindent
この補題により、定義 \ref{algebra-definition-spectrum-ring} の部分集合
$V(T)$ は $\Spec(R)$ 上のある位相の閉集合系をなす。
また、集合 $D(f)$ は開であり、この位相の基底をなすことも分かる。

\begin{definition}
\label{algebra-definition-Zariski-topology}
$R$ を環とする。
閉集合が $V(T)$ という形の集合である $\Spec(R)$ 上の位相を
{\it ザリスキー}位相という。開部分集合 $D(f)$ を $\Spec(R)$ の
{\it 標準開集合}という。
\end{definition}

\noindent
$\Spec(R)$ を単なる集合とみなすのか、位相空間とみなすのかは、
文脈から明らかであろう。

\begin{lemma}
\label{algebra-lemma-spec-functorial}
\begin{slogan}
スペクトルの関手性
\end{slogan}
$\varphi : R \to R'$ を環準同型とする。誘導写像
$$
\Spec(\varphi) : \Spec(R') \longrightarrow \Spec(R),
\quad
\mathfrak p' \longmapsto \varphi^{-1}(\mathfrak p')
$$
はザリスキー位相について連続である。実際、任意の元
$f \in R$ に対して
$\Spec(\varphi)^{-1}(D(f)) = D(\varphi(f))$ である。
\end{lemma}

\begin{proof}
$\mathfrak p := \varphi^{-1}(\mathfrak p')$ が実際に $R$ の素イデアルで
あることは基本概念 (\ref{algebra-item-inverse-image-prime}) である。
補題の最後の主張は定義から直ちに従い、最初の主張を導く。
\end{proof}

\noindent
$\varphi' : R' \to R''$ を第二の環準同型とすると、合成
$$
\Spec(R'')
\longrightarrow
\Spec(R')
\longrightarrow
\Spec(R)
$$
は $\Spec(\varphi' \circ \varphi)$ に等しい。
言い換えれば、$\Spec$ は環の圏から位相空間の圏への反変関手である。

\begin{lemma}
\label{algebra-lemma-spec-localization}
$R$ を環、$S \subset R$ を乗法的集合とする。
写像 $R \to S^{-1}R$ は $\Spec$ の関手性により同相写像
$$
\Spec(S^{-1}R)
\longrightarrow
\{\mathfrak p \in \Spec(R) \mid S \cap \mathfrak p = \emptyset \}
$$
を誘導する。ここで右辺には $\Spec(R)$ のザリスキー位相から誘導される
位相を入れる。逆写像は
$\mathfrak p \mapsto S^{-1}\mathfrak p = \mathfrak p(S^{-1}R)$ で与えられる。
\end{lemma}

\begin{proof}
補題の矢印の右辺を $D$ と書く。
素イデアル $\mathfrak p' \subset S^{-1}R$ を選び、
$\mathfrak p$ を $\mathfrak p'$ の $R$ における逆像とする。
$\mathfrak p'$ は $1$ を含まないので、$\mathfrak p$ は $S$ の
どの元も含まない。したがって $\mathfrak p \in D$ であり、像は
$D$ に含まれる。$\mathfrak p \in D$ とする。
仮定により像 $\overline{S}$ は $0$ を含まない。
基本概念 (\ref{algebra-item-localization-zero}) により
$\overline{S}^{-1}(R/\mathfrak p)$ は零環でない。
基本概念 (\ref{algebra-item-localize-ideal}) により
$S^{-1}R / S^{-1}\mathfrak p = \overline{S}^{-1}(R/\mathfrak p)$
は整域なので、$S^{-1}\mathfrak p$ は素である。
この環の等式は、$S^{-1}\mathfrak p$ の $R$ における逆像が
$\mathfrak p$ に等しいことも示す。実際、基本概念
(\ref{algebra-item-localize-nonzerodivisors}) により
$R/\mathfrak p \to \overline{S}^{-1}(R/\mathfrak p)$ は単射である。
以上により、写像 $\Spec(S^{-1}R) \to \Spec(R)$ は $D$ の上への
全単射であり、逆写像は主張のとおりである。
補題 \ref{algebra-lemma-spec-functorial} によりこれは連続である。
最後に、$D(g) \subset \Spec(S^{-1}R)$ を標準開集合とする。
ある $h\in R$ と $s\in S$ により $g = h/s$ と書く。
$g$ と $h/1$ は単元倍だけ異なるので、
$\Spec(S^{-1}R)$ において $D(g) = D(h/1)$ である。
したがって補題 \ref{algebra-lemma-spec-functorial} と上の全単射性により、
$D(g) = D(h/1)$ の像は $D \cap D(h)$ である。
これで写像が開でもあることが分かる。
\end{proof}

\begin{lemma}
\label{algebra-lemma-standard-open}
$R$ を環、$f \in R$ とする。
写像 $R \to R_f$ は $\Spec$ の関手性により同相写像
$$
\Spec(R_f) \longrightarrow D(f) \subset \Spec(R).
$$
を誘導する。逆写像は
$\mathfrak p \mapsto \mathfrak p \cdot R_f$ で与えられる。
\end{lemma}

\begin{proof}
補題 \ref{algebra-lemma-spec-localization} の特別な場合である。
\end{proof}

\noindent
スペクトルのすべての「アフィン開集合」が標準開集合であるわけではない。
例 \ref{algebra-example-affine-open-not-standard} を参照せよ。

\begin{lemma}
\label{algebra-lemma-spec-closed}
$R$ を環、$I \subset R$ をイデアルとする。
写像 $R \to R/I$ は $\Spec$ の関手性により同相写像
$$
\Spec(R/I) \longrightarrow V(I) \subset \Spec(R).
$$
を誘導する。逆写像は
$\mathfrak p \mapsto \mathfrak p / I$ で与えられる。
\end{lemma}

\begin{proof}
像が $V(I)$ に含まれることは直ちに分かる。
一方、$\mathfrak p \in V(I)$ ならば $\mathfrak p \supset I$ であり、
イデアル $\mathfrak p /I \subset R/I$ を考えられる。
基本概念 (\ref{algebra-item-isomorphism-theorem}) を用いると
$(R/I)/(\mathfrak p/I) = R/\mathfrak p$ は整域であり、
したがって $\mathfrak p/I$ は素イデアルである。
さらに $D(f + I)$ の像が $D(f) \cap V(I)$ であることは直ちに分かるので、
この写像は同相写像である。
\end{proof}



\begin{lemma}
\label{algebra-lemma-quasi-compact}
\begin{slogan}
環のスペクトルは準コンパクトである
\end{slogan}
$R$ を環とする。空間 $\Spec(R)$ は準コンパクトである。
\end{lemma}

\begin{proof}
$\Spec(R)$ の標準開集合による任意の被覆が有限被覆に細分されることを
示せば十分である。そこで、$R$ の元の集合
$\{f_i\}_{i\in I}$ に対して
$\Spec(R) = \cup D(f_i)$ であると仮定する。
これは $\cap V(f_i) = \emptyset$ を意味する。
補題 \ref{algebra-lemma-Zariski-topology} によれば、これは
$V(\{f_i \}) = \emptyset$ を意味する。
同じ補題によれば、これは $f_i$ が生成するイデアルが $R$ の単位イデアルで
あることを意味する。したがって $1$ を {\it 有限}和
$1 = \sum_{i \in J} r_i f_i$ として書ける。ここで
$J \subset I$ は有限である。このとき
$\Spec(R) = \cup_{i \in J} D(f_i)$ が従う。
\end{proof}

\begin{lemma}
\label{algebra-lemma-topology-spec}
$R$ を環とする。
$X = \Spec(R)$ 上の位相は次の性質をもつ。
\begin{enumerate}
\item $X$ は準コンパクトである。
\item $X$ は準コンパクト開集合からなる位相の基底をもつ。
\item 任意の二つの準コンパクト開集合の共通部分は準コンパクトである。
\end{enumerate}
\end{lemma}

\begin{proof}
環のスペクトルは準コンパクトである。補題
\ref{algebra-lemma-quasi-compact} を参照せよ。
これは標準開集合
$D(f) = \Spec(R_f)$
からなる位相の基底をもち（補題 \ref{algebra-lemma-standard-open}）、
最初の注意によりそれらは準コンパクトである。
二つの標準開集合の共通部分は
$D(f) \cap D(g) = D(fg)$ なので準コンパクトである。
任意の二つの準コンパクト開集合 $U, V \subset X$ に対し、
$U = D(f_1) \cup \ldots \cup D(f_n)$ および
$V = D(g_1) \cup \ldots \cup D(g_m)$ と書ける。このとき
$U \cap V = \bigcup D(f_ig_j)$ は準コンパクトである。
\end{proof}

\section{局所環}
\label{algebra-section-local-rings}

\noindent
局所環は代数幾何学の最も基本的な道具である。

\begin{definition}
\label{algebra-definition-local-ring}
{\it 局所環}とは、極大イデアルをちょうど一つもつ環である。
$R$ が局所環ならば、その極大イデアルをしばしば
$\mathfrak m_R$ と書き、体 $R/\mathfrak m_R$ を局所環 $R$ の
{\it 剰余体}という。
$R$ が局所環、$\mathfrak m$ がその唯一の極大イデアルであることを表すため、
しばしば「$(R, \mathfrak m)$ を局所環とする」といい、
さらに $\kappa = R/\mathfrak m$ がその剰余体であることも表すため、
「$(R, \mathfrak m, \kappa)$ を局所環とする」という。
{\it 局所環の局所準同型}とは、$R$ と $S$ が局所環であり、
$\varphi(\mathfrak m_R) \subset \mathfrak m_S$ を満たす環準同型
$\varphi : R \to S$ のことである。
$R$ と $S$ が局所環であることが既知ならば、「{\it 局所環準同型
$\varphi : R \to S$}」という語は、$\varphi$ が局所環の局所準同型で
あることを意味する。
\end{definition}

\noindent
体は局所環である。体の間の任意の環準同型は局所環の局所準同型である。

\medskip\noindent
環 $R$ の素イデアル $\mathfrak p$ における局所化 $R_\mathfrak p$ は、
極大イデアル $\mathfrak p R_\mathfrak p$ をもつ局所環である。
実際、補題 \ref{algebra-lemma-spec-localization} により、$R_\mathfrak p$ の
すべての素イデアルは素イデアル $\mathfrak p R_\mathfrak p$ に含まれる
（したがってこれは極大イデアルであり、$R_\mathfrak p$ の唯一の極大イデアルである）。
$R_\mathfrak p$ の剰余体を $\kappa(\mathfrak p)$ と書き、
{\it $\mathfrak p$ の剰余体}という。命題
\ref{algebra-proposition-localize-quotient} により、
$\kappa(\mathfrak p)$ を整域 $R/\mathfrak p$ の分数体と同一視できる。
合成
$$
\Spec(\kappa(\mathfrak p)) \to \Spec(R_\mathfrak p) \to \Spec(R)
$$
により、始域の唯一の点は終域の点 $\mathfrak p$ へ写る。

\medskip\noindent
$\varphi : R \to S$ を環準同型とする。$\mathfrak q \subset S$ を
素イデアルとし、$R$ の素イデアル
$\mathfrak p = \varphi^{-1}(\mathfrak q)$ を考える。
$\varphi(\mathfrak p) \subset \mathfrak q$ なので、誘導環準同型
$$
R_\mathfrak p \to S_\mathfrak q,\quad r/g \mapsto \varphi(r)/\varphi(g)
$$
は局所環準同型であり、剰余体の誘導写像
$\kappa(\mathfrak p) \to \kappa(\mathfrak q)$ を得る。


\begin{example}
\label{algebra-example-not-local}
$R$ が局所環で $\mathfrak p \subset R$ が極大でない素イデアルならば、
$R \to R_\mathfrak p$ は局所準同型でない。
\end{example}

\begin{lemma}
\label{algebra-lemma-characterize-local-ring}
$R$ を環とする。次は同値である。
\begin{enumerate}
\item $R$ は局所環である。
\item $\Spec(R)$ は閉点をちょうど一つもつ。
\item $R$ は極大イデアル $\mathfrak m$ をもち、
$R \setminus \mathfrak m$ のすべての元が単元である。
\item $R$ は零環でなく、すべての $x \in R$ に対して
$x$ または $1 - x$ の少なくとも一方が可逆である。
\end{enumerate}
\end{lemma}

\begin{proof}
$R$ を環、$\mathfrak m$ を極大イデアルとする。
$x \in R \setminus \mathfrak m$ が単元でなければ、$x$ を含む極大イデアル
$\mathfrak m'$ が存在する。したがって $R$ は少なくとも二つの極大イデアルをもつ。
逆に、$\mathfrak m'$ が別の極大イデアルならば、
$x \in \mathfrak m'$, $x \not \in \mathfrak m$ を満たす $x$ を選ぶ。
明らかに $x$ は単元でない。これで (1) と (3) の同値性が示された。
(1) と (2) の同値性は同語反復である。
$R$ が局所環ならば、$x$ は $\mathfrak m$ に属するか属さないかの
いずれかなので (4) が成り立つ。
(4) が成り立ち、$\mathfrak m$, $\mathfrak m'$ が互いに異なる
極大イデアルならば、中国剰余定理
（補題 \ref{algebra-lemma-chinese-remainder}）により、
$x \bmod \mathfrak m' = 0$ かつ $x \bmod \mathfrak m = 1$ を満たす
$x \in R$ を選べる。この元 $x$ は可逆でなく、$1 - x$ も可逆でないので
矛盾する。したがって (4) と (1) は同値である。
\end{proof}

\begin{lemma}
\label{algebra-lemma-characterize-local-ring-map}
$\varphi : R \to S$ を環準同型とし、$R$ と $S$ が局所環であると仮定する。
次は同値である。
\begin{enumerate}
\item $\varphi$ は局所環準同型である。
\item $\varphi(\mathfrak m_R) \subset \mathfrak m_S$ である。
\item $\varphi^{-1}(\mathfrak m_S) = \mathfrak m_R$ である。
\item 任意の $x \in R$ に対し、$\varphi(x)$ が $S$ で可逆ならば、
$x$ は $R$ で可逆である。
\end{enumerate}
\end{lemma}

\begin{proof}
条件 (1) と (2) は定義により同値である。
(3) ならば (2) である。逆に (2) が成り立つならば、
$\varphi^{-1}(\mathfrak m_S)$ は極大イデアル $\mathfrak m_R$ を含む
素イデアルなので、$\varphi^{-1}(\mathfrak m_S) = \mathfrak m_R$ である。
最後に、補題 \ref{algebra-lemma-characterize-local-ring} により、
(4) は (2) の対偶である。
\end{proof}

\begin{remark}
\label{algebra-remark-fundamental-diagram}
環準同型 $\varphi : R \to S$ と素イデアル
$\mathfrak p \subset R$ に付随する基本的な可換図式は次である。
$$
\xymatrix{
\kappa(\mathfrak p) \otimes_R S =
S_{\mathfrak p}/{\mathfrak p}S_{\mathfrak p}
&
S_{\mathfrak p} \ar[l]
&
S \ar[r] \ar[l]
&
S/\mathfrak pS \ar[r]
&
(R \setminus \mathfrak p)^{-1}S/\mathfrak pS
\\
\kappa(\mathfrak p) =
R_{\mathfrak p}/{\mathfrak p}R_{\mathfrak p} \ar[u]
&
R_{\mathfrak p} \ar[u] \ar[l]
&
R \ar[u] \ar[r] \ar[l]
&
R/\mathfrak p \ar[u] \ar[r]
&
\kappa(\mathfrak p) \ar[u]
}
$$
この図式では最左列と最右列は同一である。
スペクトル上で水平写像はその像の上への同相写像を誘導し、各正方形は
位相空間のファイバー積の正方形を誘導する
（補題 \ref{algebra-lemma-spec-localization} および
\ref{algebra-lemma-spec-closed} を参照せよ）。
これにより、$\mathfrak p$ がスペクトル上の写像の像に属することと、
$S \otimes_R \kappa(\mathfrak p)$ が零環でないことは同値である。
$\mathfrak p$ の上にある素イデアル $\mathfrak q \subset S$、
すなわち $\mathfrak p = \varphi^{-1}(\mathfrak q)$ を満たすものが
存在するならば、図式を次の図式へ拡張できる。
$$
\xymatrix{
\kappa(\mathfrak q) = S_{\mathfrak q}/{\mathfrak q}S_{\mathfrak q}
&
S_{\mathfrak q} \ar[l]
&
S \ar[r] \ar[l]
&
S/\mathfrak q \ar[r]
&
\kappa(\mathfrak q)
\\
\kappa(\mathfrak p) \otimes_R S =
S_{\mathfrak p}/{\mathfrak p}S_{\mathfrak p} \ar[u]
&
S_{\mathfrak p} \ar[u] \ar[l]
&
S \ar[u] \ar[r] \ar[l]
&
S/\mathfrak pS \ar[u] \ar[r]
&
(R \setminus \mathfrak p)^{-1}S/\mathfrak pS \ar[u]
\\
\kappa(\mathfrak p) =
R_{\mathfrak p}/{\mathfrak p}R_{\mathfrak p} \ar[u]
&
R_{\mathfrak p} \ar[u] \ar[l]
&
R \ar[u] \ar[r] \ar[l]
&
R/\mathfrak p \ar[u] \ar[r]
&
\kappa(\mathfrak p) \ar[u]
}
$$
この図式でも、最左列と最右列は同一であり、スペクトル上で水平写像は
その像の上への同相写像を誘導する。
\end{remark}

\begin{lemma}
\label{algebra-lemma-in-image}
$\varphi : R \to S$ を環準同型とし、$\mathfrak p$ を $R$ の
素イデアルとする。次は同値である。
\begin{enumerate}
\item $\mathfrak p$ は $\Spec(S) \to \Spec(R)$ の像に属する。
\item $S \otimes_R \kappa(\mathfrak p) \not = 0$ である。
\item $S_{\mathfrak p}/\mathfrak p S_{\mathfrak p} \not = 0$ である。
\item $(S/\mathfrak pS)_{\mathfrak p} \not = 0$ である。
\item $\mathfrak p = \varphi^{-1}(\mathfrak pS)$ である。
\end{enumerate}
\end{lemma}

\begin{proof}
最初の二条件が同値であることは注意
\ref{algebra-remark-fundamental-diagram} ですでに見た。
残りはこの事実の言い換えにすぎない。
\end{proof}

\begin{remark}
\label{algebra-remark-local-ring-fibre}
$R \to S$ を環準同型とする。$\mathfrak q$ を、$R$ の素イデアル
$\mathfrak p$ の上にある $S$ の素イデアルとする。
注意 \ref{algebra-remark-fundamental-diagram} によれば、素イデアル
$\mathfrak q$ はファイバー環
$F = S \otimes_R \kappa(\mathfrak p)$ の一意な素イデアル
$\overline{\mathfrak q}$ に対応する。このとき
$$
F_{\overline{\mathfrak q}}
\cong S_\mathfrak q \otimes_{R_\mathfrak p} \kappa(\mathfrak p)
\cong S_\mathfrak q/\mathfrak p S_\mathfrak q
$$
である。実際、明らかな環準同型
$F \to S_\mathfrak q \otimes_{R_\mathfrak p} \kappa(\mathfrak p)$
があり、これは $F \to F_{\overline{\mathfrak q}}$ と同型であることが
容易に分かる。第二の同型は、$\kappa(\mathfrak p)$ が
$R_\mathfrak p$ の $\mathfrak pR_\mathfrak p$ による商であることから従う。
\end{remark}



\section{環のジャコブソン根基}
\label{algebra-section-radical}

\noindent
環 $R$ の {\it ジャコブソン根基} $\text{rad}(R)$ とは、$R$ の
すべての極大イデアルの共通部分のことである。$R$ が局所環ならば、
$\text{rad}(R)$ は $R$ の極大イデアルである。

\begin{lemma}
\label{algebra-lemma-contained-in-radical}
$R$ を、ジャコブソン根基が $\text{rad}(R)$ である環とする。
$I \subset R$ をイデアルとする。このとき、次は同値である。
\begin{enumerate}
\item $I \subset \text{rad}(R)$ である。
\item $1 + I$ のすべての元は $R$ の単元である。
\end{enumerate}
この場合、$R/I$ の単元に写る $R$ の任意の元は単元である。
\end{lemma}

\begin{proof}
$f \in \text{rad}(R)$ ならば、$R$ のすべての極大イデアル
$\mathfrak m$ に対して $f \in \mathfrak m$ である。したがって、
$R$ のすべての極大イデアル $\mathfrak m$ に対して
$1 + f \not \in \mathfrak m$ である。ゆえに $\Spec(R)$ の閉部分集合
$V(1 + f)$ は空である。補題 \ref{algebra-lemma-Zariski-topology} により、
これは $1 + f$ が単元であることを意味する。

\medskip\noindent
逆に、すべての $f \in I$ に対して $1 + f$ が単元であると仮定する。
$\mathfrak m$ が極大イデアルで $I \not \subset \mathfrak m$ ならば、
$I + \mathfrak m = R$ である。したがって、ある $g \in \mathfrak m$
および $f \in I$ に対して $1 = f + g$ と書ける。このとき
$g = 1 + (-f)$ は単元でなく、矛盾する。

\medskip\noindent
最後の主張について、$f \in R$ が $R/I$ の単元に写るとする。
$f \bmod I$ の乗法逆元に写る $g \in R$ を取ることができる。
すると $fg = 1 \bmod I$ である。したがって (2) により $fg$ は
$R$ の単元であり、ゆえに $f$ も単元である。
\end{proof}

\begin{lemma}
\label{algebra-lemma-surjective-on-spec-units}
$\varphi : R \to S$ を、誘導される写像
$\Spec(S) \to \Spec(R)$ が全射となる環準同型とする。このとき
$x \in R$ が単元であることと、$\varphi(x) \in S$ が単元であることは
同値である。
\end{lemma}

\begin{proof}
$x$ が単元ならば $\varphi(x)$ も単元である。逆に、$\varphi(x)$ が
単元ならば、すべての $\mathfrak q \in \Spec(S)$ に対して
$\varphi(x) \not \in \mathfrak q$ である。したがって、すべての
$\mathfrak q \in \Spec(S)$ に対して
$x \not \in \varphi^{-1}(\mathfrak q) = \Spec(\varphi)(\mathfrak q)$
である。$\Spec(\varphi)$ は全射なので、補題
\ref{algebra-lemma-Zariski-topology} の (17) により、$x$ は単元である。
\end{proof}




\section{中山の補題}
\label{algebra-section-nakayama}

\noindent
\cite{MatCA} から引用する。「この単純だが重要な補題は、
T.~Nakayama、G.~Azumaya および W.~Krull による。優先権は明らかでなく、
通常は中山の補題と呼ばれるものの、故 Nakayama 教授はこの名称を
好まなかった。」

\begin{lemma}[中山の補題]
\label{algebra-lemma-NAK}
\begin{reference}
\cite[1.M Lemma (NAK) page 11]{MatCA}
\end{reference}
\begin{history}
\cite{MatCA} から引用する。「この単純だが重要な補題は、
T.~Nakayama、G.~Azumaya および W.~Krull による。優先権は明らかでなく、
通常は中山の補題と呼ばれるものの、故 Nakayama 教授はこの名称を
好まなかった。」
\end{history}
$R$ を、ジャコブソン根基が $\text{rad}(R)$ である環とする。
$M$ を $R$-加群とし、$I \subset R$ をイデアルとする。
\begin{enumerate}
\item
\label{algebra-item-nakayama}
$IM = M$ かつ $M$ が有限ならば、$fM = 0$ を満たす
$f \in 1 + I$ が存在する。
\item $IM = M$ であり、$M$ が有限で、$I \subset \text{rad}(R)$
ならば $M = 0$ である。
\item $N, N' \subset M$ であり、$M = N + IN'$ かつ $N'$ が有限ならば、
$fM \subset N$ および $M_f = N_f$ を満たす $f \in 1 + I$ が存在する。
\item $N, N' \subset M$ であり、$M = N + IN'$、$N'$ が有限、かつ
$I \subset \text{rad}(R)$ ならば $M = N$ である。
\item $N \to M$ が加群準同型であり、$N/IN \to M/IM$ が全射で、
$M$ が有限ならば、$N_f \to M_f$ が全射となる $f \in 1 + I$ が存在する。
\item $N \to M$ が加群準同型であり、$N/IN \to M/IM$ が全射、
$M$ が有限、かつ $I \subset \text{rad}(R)$ ならば、
$N \to M$ は全射である。
\item $x_1, \ldots, x_n \in M$ が $M/IM$ を生成し、$M$ が有限ならば、
$x_1, \ldots, x_n$ が $R_f$ 上で $M_f$ を生成するような
$f \in 1 + I$ が存在する。
\item $x_1, \ldots, x_n \in M$ が $M/IM$ を生成し、$M$ が有限で、
$I \subset \text{rad}(R)$ ならば、$M$ は $x_1, \ldots, x_n$ により
生成される。
\item $IM = M$ で $I$ が冪零ならば、$M = 0$ である。
\item $N, N' \subset M$ であり、$M = N + IN'$ かつ $I$ が冪零ならば
$M = N$ である。
\item $N \to M$ が加群準同型で、$I$ が冪零、かつ
$N/IN \to M/IM$ が全射ならば、$N \to M$ は全射である。
\item $\{x_\alpha\}_{\alpha \in A}$ が $M$ の元の集合で $M/IM$ を生成し、
$I$ が冪零ならば、$M$ はこれらの $x_\alpha$ により生成される。
\end{enumerate}
\end{lemma}

\begin{proof}
(\ref{algebra-item-nakayama}) を証明する。$R$ 上の $M$ の生成元
$y_1, \ldots, y_m$ を選ぶ。各 $i$ に対して
$z_{ij} \in I$ を用いて $y_i = \sum z_{ij} y_j$ と書ける
（$M = IM$ だからである）。言い換えれば、
$\sum_j (\delta_{ij} - z_{ij})y_j = 0$ である。
$m \times m$ 行列 $A = (\delta_{ij} - z_{ij})$ の行列式を $f$ とする。
$f \in 1 + I$ であることに注意する（行列 $A$ は各成分について、
$I$ を法として $m \times m$ 単位行列と合同だからである）。
補題 \ref{algebra-lemma-matrix-left-inverse} の (1) により、
$BA = f 1_{m \times m}$ を満たす $m \times m$ 行列 $B$ が存在する。
成分で書けば、すべての $h$ と $j$ に対して
$\sum_{i} b_{hi} a_{ij} = f \delta_{hj}$ であり、したがって各 $h$ に対して
$\sum_{i, j} b_{hi} a_{ij} y_j
= \sum_{j} f \delta_{hj} y_j = f y_h$ である。言い換えると、各
$i$ について $\sum_j a_{ij} y_j = 0$ なので、すべての $h$ に対して
$0 = f y_h$ である。これは $f$ が $M$ を零化することを意味する。

\medskip\noindent
補題 \ref{algebra-lemma-contained-in-radical} により、$1 + \text{rad}(R)$ の元は
$R$ の可逆元である。したがって、(\ref{algebra-item-nakayama}) から (2) が従う。
(3) は、$N'$ が有限なので有限となる $M/N$ に (1) を適用して得られる。
(4) は、$N'$ が有限なので有限となる $M/N$ に (2) を適用して得られる。
(5) は、$M$ とその部分加群 $\Im(N \to M)$ および $M$ に (3) を
適用して得られる。(6) は、$M$ とその部分加群 $\Im(N \to M)$ および
$M$ に (4) を適用して得られる。
(7) は、写像 $R^{\oplus n} \to M$、
$(a_1, \ldots, a_n) \mapsto a_1x_1 + \ldots + a_nx_n$ に (5) を
適用して得られる。(8) は、写像 $R^{\oplus n} \to M$、
$(a_1, \ldots, a_n) \mapsto a_1x_1 + \ldots + a_nx_n$ に (6) を
適用して得られる。

\medskip\noindent
(9) が成り立つのは、$M = IM$ ならばすべての $n \geq 0$ に対して
$M = I^nM$ であり、$I$ が冪零であることから、ある $n \gg 0$ に対して
$I^n = 0$ となるためである。(10)、(11)、(12) は、上と同じ議論により
(9) から従う。
\end{proof}

\begin{lemma}
\label{algebra-lemma-NAK-localization}
$R$ を環、$S \subset R$ を乗法的部分集合、$I \subset R$ をイデアル、
$M$ を有限 $R$-加群とする。$x_1, \ldots, x_r \in M$ が
$S^{-1}(R/I)$-加群として $S^{-1}(M/IM)$ を生成するならば、
$x_1, \ldots, x_r$ が $R_f$-加群として $M_f$ を生成するような
$f \in S + I$ が存在する。\footnote{特別な場合：(I) $I = 0$。
この補題によれば、$x_1, \ldots, x_r$ が $S^{-1}M$ を生成するならば、
ある $f \in S$ に対して $x_1, \ldots, x_r$ は $M_f$ を生成する。
(II) $I = \mathfrak p$ が素イデアルで、$S = R \setminus \mathfrak p$
である場合。この補題によれば、$x_1, \ldots, x_r$ が
$M \otimes_R \kappa(\mathfrak p)$ を生成するならば、ある
$f \in R$、$f \not \in \mathfrak p$ に対して $x_1, \ldots, x_r$ は
$M_f$ を生成する。}
\end{lemma}

\begin{proof}
まず特別な場合 $I = 0$ を扱う。$R$ 上の $M$ の生成元を
$y_1, \ldots, y_s$ とする。$S^{-1}M$ は $x_1, \ldots, x_r$ により
生成されるので、各 $i$ に対し、ある $a_{ij} \in R$ および
$s_{ij} \in S$ を用いて、$S^{-1}M$ において
$y_i = \sum (a_{ij}/s_{ij})x_j$ と書ける。$s_{ij}$ の積
$s \in S$ を掛けると、ある $a'_{ij} \in R$ に対して
$S^{-1}M$ において $sy_i = \sum a'_{ij}x_j$ となる。
これはさらに、$M$ において
$t_isy_i = \sum t_ia'_{ij}x_j$ となるような $t_i \in S$ が
存在することを意味する。そこで $t \in S$ を $t_i$ の積とすれば、
$M_{st}$ の $x_1, \ldots, x_r$ が生成する $R_{st}$-部分加群に
$y_i$ が属することが分かる。したがって、$x_1, \ldots, x_r$ は
$M_{st}$ を生成する。

\medskip\noindent
一般の場合を考える。特別な場合により、$x_1, \ldots, x_r$ が
$(R/I)_s$ 上で $(M/IM)_s$ を生成するような $s \in S$ を取れる。
補題 \ref{algebra-lemma-NAK} により、$x_1, \ldots, x_r$ が $(R_s)_g$ 上で
$(M_s)_g$ を生成するような $g \in 1 + I_s \subset R_s$ を取れる。
$g = 1 + i/s'$ と書く。このとき $f = ss' + is$ とすればよい。
詳細は省略する。
\end{proof}

\begin{lemma}
\label{algebra-lemma-when-surjective-local}
$A \to B$ を局所環の局所準同型とする。次を仮定する。
\begin{enumerate}
\item $B$ は有限 $A$-加群である。
\item $\mathfrak m_B$ は有限生成イデアルである。
\item $A \to B$ は剰余体上の同型を誘導する。
\item $\mathfrak m_A/\mathfrak m_A^2 \to \mathfrak m_B/\mathfrak m_B^2$
は全射である。
\end{enumerate}
このとき $A \to B$ は全射である。
\end{lemma}

\begin{proof}
$A \to B$ が全射であることを示すため、これを $A$-加群の写像とみなし、
補題 \ref{algebra-lemma-NAK} の (6) を適用する。すると、
$A/\mathfrak m_A \to B/\mathfrak m_AB$ が全射であることを示せば十分である。
$A/\mathfrak m_A = B/\mathfrak m_B$ なので、
$\mathfrak m_AB \to \mathfrak m_B$ が全射であることを示せば十分である。
$\mathfrak m_AB \to \mathfrak m_B$ を $B$-加群の写像とみなし、
補題 \ref{algebra-lemma-NAK} の (6) を適用する。すると、
$\mathfrak m_AB/\mathfrak m_A\mathfrak m_B \to \mathfrak m_B/\mathfrak m_B^2$
が全射であることを示せば十分である。これは仮定 (4) から従う。
\end{proof}






\section{スペクトルの開閉部分集合}
\label{algebra-section-open-and-closed}

\noindent
スペクトルの開かつ閉な部分集合は、環の冪等元に対応することが分かる。

\begin{lemma}
\label{algebra-lemma-idempotent-spec}
$R$ を環とし、$e \in R$ を冪等元とする。この場合、
$$
\Spec(R) = D(e) \amalg D(1-e).
$$
\end{lemma}

\begin{proof}
整域の冪等元 $e$ は $1$ または $0$ のいずれかであることに注意する。
したがって、
\begin{eqnarray*}
D(e)
& = &
\{ \mathfrak p \in \Spec(R)
\mid
e \not\in \mathfrak p \} \\
& = &
\{ \mathfrak p \in \Spec(R)
\mid
e \not = 0\text{ in }\kappa(\mathfrak p) \} \\
& = &
\{ \mathfrak p \in \Spec(R)
\mid
e = 1\text{ in }\kappa(\mathfrak p) \}
\end{eqnarray*}
である。同様に、
\begin{eqnarray*}
D(1-e)
& = &
\{ \mathfrak p \in \Spec(R)
\mid
1 - e \not\in \mathfrak p \} \\
& = &
\{ \mathfrak p \in \Spec(R)
\mid
e \not = 1\text{ in }\kappa(\mathfrak p) \} \\
& = &
\{ \mathfrak p \in \Spec(R)
\mid
e = 0\text{ in }\kappa(\mathfrak p) \}
\end{eqnarray*}
を得る。任意の剰余体における $e$ の像は $1$ または $0$ のいずれかなので、
$D(e)$ と $D(1-e)$ は $\Spec(R)$ 全体を覆う。
\end{proof}

\begin{lemma}
\label{algebra-lemma-spec-product}
$R_1$ と $R_2$ を環とし、$R = R_1 \times R_2$ とおく。
写像 $R \to R_1$、$(x, y) \mapsto x$ および $R \to R_2$、
$(x, y) \mapsto y$ は、連続写像 $\Spec(R_1) \to \Spec(R)$ および
$\Spec(R_2) \to \Spec(R)$ を誘導する。誘導される写像
$$
\Spec(R_1) \amalg \Spec(R_2)
\longrightarrow
\Spec(R)
$$
は同相写像である。言い換えれば、$R = R_1\times R_2$ のスペクトルは
$R_1$ のスペクトルと $R_2$ のスペクトルの非交和である。
\end{lemma}

\begin{proof}
$e_1 = (1, 0)$ および $e_2 = (0, 1)$ として
$1 = e_1 + e_2$ と書く。$e_1$ と $e_2 = 1 - e_1$ は冪等元である。
$R_1 = R_{e_1}$ が $e_1$ における $R$ の局所化であることの確認は
読者に委ねる。$e_2$ についても同様である。したがって、この補題の主張は
補題 \ref{algebra-lemma-idempotent-spec} と補題 \ref{algebra-lemma-standard-open} から従う。
\end{proof}

\noindent
以下の補題は、関数の貼り合わせ補題を導入した後でもう一度証明する。
節 \ref{algebra-section-tilde-module-sheaf} を参照せよ。

\begin{lemma}
\label{algebra-lemma-disjoint-decomposition}
$R$ を環とする。開かつ閉な各 $U \subset \Spec(R)$ に対して、
$U = D(e)$ を満たす一意な冪等元 $e \in R$ が存在する。これにより、
開かつ閉な部分集合 $U \subset \Spec(R)$ と冪等元 $e \in R$ の間に
一対一対応が誘導される。
\end{lemma}

\begin{proof}
$U \subset \Spec(R)$ を開かつ閉とする。$U$ は閉なので、補題
\ref{algebra-lemma-quasi-compact} により準コンパクトであり、その補集合も同様である。
$U = \bigcup_{i = 1}^n D(f_i)$ を標準開集合の有限合併として書く。
同様に、$\Spec(R) \setminus U = \bigcup_{j = 1}^m D(g_j)$ を標準開集合の
有限合併として書く。$\emptyset = D(f_i) \cap D(g_j) = D(f_i g_j)$ なので、
補題 \ref{algebra-lemma-Zariski-topology} により $f_i g_j$ は冪零である。
$I = (f_1, \ldots, f_n) \subset R$ および
$J = (g_1, \ldots, g_m) \subset R$ とおく。$V(J)$ は $U$ に等しく、
$V(I)$ は $U$ の補集合に等しいので、$\Spec(R) = V(I) \amalg V(J)$ である。
上の冪零性に関する考察から、十分大きい整数 $N$ に対して
$(IJ)^N = (0)$ である。また、
$\bigcup D(f_i) \cup \bigcup D(g_j) = \Spec(R)$ なので、補題
\ref{algebra-lemma-Zariski-topology} により $I + J = R$ である。
この等式を $2N$ 乗すると $I^N + J^N = R$ を得る。
$x \in I^N$ および $y \in J^N$ を用いて $1 = x + y$ と書く。
$I^N J^N = (0)$ なので $0 = xy = x(1 - x)$ である。したがって
$x = x^2$ は冪等元であり、$I^N \subset I$ に含まれる。
冪等元 $y = 1 - x$ は $J^N \subset J$ に含まれる。
これにより、$x$ はすべての $\mathfrak p \in V(J)$ に対して剰余体
$\kappa(\mathfrak p)$ で $1$ に写り、すべての $\mathfrak p \in V(I)$ に
対して $x$ は $\kappa(\mathfrak p)$ で $0$ に写ることが分かる。

\medskip\noindent
一意性を示すため、$e_1, e_2$ が $R$ の相異なる冪等元であるとする。
$e_1 \in \mathfrak p$ かつ $e_2 \not \in \mathfrak p$、またはその逆を
満たす素イデアル $\mathfrak p$ が存在することを示せばよい。
$e_i' = 1 - e_i$ と書く。$e_1 \not = e_2$ ならば、
$0 \not = e_1 - e_2  = e_1(e_2 + e_2') - (e_1 + e_1')e_2
= e_1 e_2' - e_1' e_2$ である。したがって、冪等元
$e_1 e_2' \not = 0$ または $e_1' e_2 \not = 0$ のいずれかが成り立つ。
% P01-E0296: frozen English source assertion candidate; see ERRATA_CANDIDATES.jsonl.
冪等元は冪零ではないので、補題 \ref{algebra-lemma-Zariski-topology} により、
$e_1e_2' \not \in \mathfrak p$ または
$e_1'e_2 \not \in \mathfrak p$ を満たす素イデアル $\mathfrak p$ が見つかる。
これが求める素イデアルを与えることは容易に分かる。
\end{proof}

\begin{lemma}
\label{algebra-lemma-characterize-spec-connected}
$R$ を零でない環とする。このとき $\Spec(R)$ が連結であることと、
$R$ が非自明な冪等元をもたないことは同値である。
\end{lemma}

\begin{proof}
補題 \ref{algebra-lemma-disjoint-decomposition} と連結位相空間の定義から明らかである。
\end{proof}

\begin{lemma}
\label{algebra-lemma-ideal-is-squared-union-connected}
$R$ の有限生成イデアル $I \subset R$ が $I = I^2$ を満たすとする。このとき、
\begin{enumerate}
\item $I = (e)$ を満たす冪等元 $e \in R$ が存在する。
\item 冪等元 $e' = 1 - e \in R$ に対して $R/I \cong R_{e'}$ である。
\item $V(I)$ は $\Spec(R)$ の開かつ閉な部分集合である。
\end{enumerate}
\end{lemma}

\begin{proof}
中山の補題 \ref{algebra-lemma-NAK} により、$fI = 0$ を満たす元
$f = 1 + i$、$i \in I$ が存在する。このとき $f^2 = f + fi = f$ なので、
これは冪等元である。冪等元 $e = 1 - f = -i \in I$ を考える。
$j \in I$ に対して $ej = j - fj = j$ であるから $I = (e)$ である。
これで (1) が示された。

\medskip\noindent
(2) と (3) は (1) から従う。実際、
$V(I) = V(e) = \Spec(R) \setminus D(e)$ であり、これは補題
\ref{algebra-lemma-idempotent-spec} または補題
\ref{algebra-lemma-disjoint-decomposition} により開かつ閉である。これで (3) が示された。
(2) について、写像 $R \to R_{e'}$ は全射であることに注意する。実際、
$R_{e'}$ において $x/(e')^n = x/e' = xe'/(e')^2 = xe'/e' = x/1$ である。
写像 $R \to R_{e'}$ の核は、$e'$ の正の冪で零化される $R$ の元全体である。
$e'$ は冪等元なので、これは $e'$ により零化される元のイデアルであり、
$e + e' = 1$ は互いに直交する冪等元の対だから、そのイデアルは
$I = (e)$ である。これで (2) が示された。
\end{proof}






\section{スペクトルの連結成分}
\label{algebra-section-connected-components}

\noindent
スペクトルの連結成分は、初めに思うほど理解しやすくない。
これは、われわれが局所連結空間の位相に慣れている一方で、
環のスペクトルは一般には局所連結でないためである。

\begin{lemma}
\label{algebra-lemma-closed-union-connected-components}
$R$ を環とし、$T \subset \Spec(R)$ をスペクトルの部分集合とする。
次は同値である。
\begin{enumerate}
\item $T$ は閉で、$\Spec(R)$ の連結成分の合併である。
\item $T$ は $\Spec(R)$ の開かつ閉な部分集合の共通部分である。
\item $T = V(I)$ であり、$I \subset R$ は冪等元で生成されるイデアルである。
\end{enumerate}
さらに、(3) のイデアルが存在するならば、それは一意である。
\end{lemma}

\begin{proof}
補題 \ref{algebra-lemma-topology-spec} および
Topology, Lemma \ref{topology-lemma-closed-union-connected-components} により、
(1) と (2) は同値である。(2) を仮定し、
$U_\alpha \subset \Spec(R)$ を開かつ閉として
$T = \bigcap U_\alpha$ と書く。補題
\ref{algebra-lemma-disjoint-decomposition} により、ある冪等元
$e_\alpha \in R$ に対して $U_\alpha = D(e_\alpha)$ である。
そこで $I = (1 - e_\alpha)$ とおけば $T = V(I)$ となり、(3) が成り立つ。
最後に (3) を仮定する。$T = V(I)$ と書き、ある冪等元の族
$e_\alpha$ によって $I = (e_\alpha)$ と書く。このとき明らかに
$T = \bigcap V(e_\alpha) = \bigcap D(1 - e_\alpha)$ である。

\medskip\noindent
$I$ を冪等元で生成されるイデアルとする。$V(I) \subset V(e)$ を満たす
冪等元 $e \in R$ を取る。補題 \ref{algebra-lemma-Zariski-topology} により、
ある $n \geq 1$ に対して $e^n \in I$ である。$e$ は冪等元なので、
これは $e \in I$ を意味する。したがって $I$ は、ちょうど
$T \subset V(e)$ を満たす冪等元 $e$ たちにより生成される。
言い換えれば、イデアル $I$ は閉部分集合 $T$ によって完全に定まり、
一意性が示された。
\end{proof}

\begin{lemma}
\label{algebra-lemma-connected-component}
$R$ を環とする。$\Spec(R)$ の連結成分は $V(I)$ の形である。
ここで $I$ は冪等元で生成されるイデアルであり、$R$ の各冪等元は
$R/I$ において $0$ または $1$ のいずれかに写る。
\end{lemma}

\begin{proof}
$\mathfrak p$ を $R$ の素イデアルとする。補題
\ref{algebra-lemma-topology-spec} により、位相空間 $\Spec(R)$ は
% P01-E0297: frozen English source wording candidate; see ERRATA_CANDIDATES.jsonl.
Topology, Lemma \ref{topology-lemma-connected-component-intersection} の
仮定を満たす。したがって、$\Spec(R)$ における $\mathfrak p$ の連結成分は、
$\mathfrak p$ を含む $\Spec(R)$ の開かつ閉な部分集合すべての共通部分である。
ゆえに、それは $V(I)$ に等しい。ここで $I$ は、冪等元 $e \in R$ であって
$e$ が $\kappa(\mathfrak p)$ において $0$ に写るものにより生成される。
補題 \ref{algebra-lemma-disjoint-decomposition} を参照せよ。
この族に属さない任意の冪等元 $e$ は、明らかに $R/I$ において $1$ に写る。
\end{proof}









\section{貼り合わせの性質}
\label{algebra-section-more-glueing}

\noindent
この節では、「あることが標準開被覆のすべての要素上で成り立てば、
それは大域的にも成り立つ」という形の標準的な結果をいくつかまとめる。
実際、次の補題のように、局所環のレベルで確認すれば十分なことが多い。

\begin{lemma}
\label{algebra-lemma-characterize-zero-local}
$R$ を環とする。
\begin{enumerate}
\item $R$-加群 $M$ の元 $x$ に対して、次は同値である。
\begin{enumerate}
\item $x = 0$ である。
\item すべての $\mathfrak p \in \Spec(R)$ に対して、$x$ は
$M_\mathfrak p$ で零に写る。
\item $R$ のすべての極大イデアル $\mathfrak m$ に対して、$x$ は
$M_{\mathfrak m}$ で零に写る。
\end{enumerate}
言い換えれば、写像 $M \to \prod_{\mathfrak m} M_{\mathfrak m}$ は単射である。
\item $R$-加群 $M$ に対して、次は同値である。
\begin{enumerate}
\item $M$ は零である。
\item すべての $\mathfrak p \in \Spec(R)$ に対して
$M_{\mathfrak p}$ は零である。
\item $R$ のすべての極大イデアル $\mathfrak m$ に対して
$M_{\mathfrak m}$ は零である。
\end{enumerate}
\item $R$-加群の複体 $M_1 \to M_2 \to M_3$ に対して、次は同値である。
\begin{enumerate}
\item $M_1 \to M_2 \to M_3$ は完全である。
\item $R$ のすべての素イデアル $\mathfrak p$ に対して、局所化
$M_{1, \mathfrak p} \to M_{2, \mathfrak p} \to M_{3, \mathfrak p}$
は完全である。
\item $R$ のすべての極大イデアル $\mathfrak m$ に対して、局所化
$M_{1, \mathfrak m} \to M_{2, \mathfrak m} \to M_{3, \mathfrak m}$
は完全である。
\end{enumerate}
\item $R$-加群の写像 $f : M \to M'$ に対して、次は同値である。
\begin{enumerate}
\item $f$ は単射である。
\item $R$ のすべての素イデアル $\mathfrak p$ に対して
$f_{\mathfrak p} : M_\mathfrak p \to M'_\mathfrak p$ は単射である。
\item $R$ のすべての極大イデアル $\mathfrak m$ に対して
$f_{\mathfrak m} : M_\mathfrak m \to M'_\mathfrak m$ は単射である。
\end{enumerate}
\item $R$-加群の写像 $f : M \to M'$ に対して、次は同値である。
\begin{enumerate}
\item $f$ は全射である。
\item $R$ のすべての素イデアル $\mathfrak p$ に対して
$f_{\mathfrak p} : M_\mathfrak p \to M'_\mathfrak p$ は全射である。
\item $R$ のすべての極大イデアル $\mathfrak m$ に対して
$f_{\mathfrak m} : M_\mathfrak m \to M'_\mathfrak m$ は全射である。
\end{enumerate}
\item $R$-加群の写像 $f : M \to M'$ に対して、次は同値である。
\begin{enumerate}
\item $f$ は全単射である。
\item $R$ のすべての素イデアル $\mathfrak p$ に対して
$f_{\mathfrak p} : M_\mathfrak p \to M'_\mathfrak p$ は全単射である。
\item $R$ のすべての極大イデアル $\mathfrak m$ に対して
$f_{\mathfrak m} : M_\mathfrak m \to M'_\mathfrak m$ は全単射である。
\end{enumerate}
\end{enumerate}
\end{lemma}

\begin{proof}
(1) のように $x \in M$ とする。
$I = \{f \in R \mid fx = 0\}$ とおく。$I$ がイデアルであることは容易に
分かる（これは $x$ の零化イデアルである）。条件 (1)(c) は、すべての
極大イデアル $\mathfrak m$ に対して、$fx =0$ を満たす
$f \in R \setminus \mathfrak m$ が存在することを意味する。
言い換えれば、$V(I)$ は閉点を含まない。補題
\ref{algebra-lemma-Zariski-topology} により、$I$ は単位イデアルである。
したがって $x$ は零であり、すなわち (1)(a) が成り立つ。これで (1) が示された。

\medskip\noindent
(2) は、(1) を $M$ のすべての元に同時に適用すれば従う。

\medskip\noindent
(3) を証明する。$H$ をこの列のホモロジー、すなわち
$H = \Ker(M_2 \to M_3)/\Im(M_1 \to M_2)$ とする。命題
\ref{algebra-proposition-localization-exact} により、$H_\mathfrak p$ は列
$M_{1, \mathfrak p} \to M_{2, \mathfrak p} \to M_{3, \mathfrak p}$ の
ホモロジーである。したがって (3) は (2) から従う。

\medskip\noindent
(4) と (5) は (3) の特別な場合である。(6) は (4) と (5) を
組み合わせれば形式的に従う。
\end{proof}

\begin{lemma}
\label{algebra-lemma-cover}
\begin{slogan}
加群と代数のザリスキー局所的性質
\end{slogan}
$R$ を環、$M$ を $R$-加群、$S$ を $R$-代数とする。
$f_1, \ldots, f_n$ を、$\bigcup D(f_i) = \Spec(R)$、言い換えれば
$(f_1, \ldots, f_n) = R$ を満たす $R$ の元の有限列とする。
\begin{enumerate}
\item 各 $M_{f_i} = 0$ ならば $M = 0$ である。
\item 各 $M_{f_i}$ が有限 $R_{f_i}$-加群ならば、$M$ は有限 $R$-加群である。
\item 各 $M_{f_i}$ が有限表示 $R_{f_i}$-加群ならば、
$M$ は有限表示 $R$-加群である。
\item $M \to N$ を $R$-加群の写像とする。各 $i$ に対して
$M_{f_i} \to N_{f_i}$ が同型ならば、$M \to N$ は同型である。
\item $0 \to M'' \to M \to M' \to 0$ を $R$-加群の複体とする。
各 $i$ に対して $0 \to M''_{f_i} \to M_{f_i} \to M'_{f_i} \to 0$
が完全ならば、$0 \to M'' \to M \to M' \to 0$ は完全である。
\item 各 $R_{f_i}$ がネーター的ならば、$R$ はネーター的である。
\item 各 $S_{f_i}$ が有限型 $R_{f_i}$-代数ならば、
$S$ は有限型 $R$-代数である。
\item 各 $S_{f_i}$ が $R_{f_i}$ 上有限表示ならば、
$S$ は有限表示 $R$-代数である。
\end{enumerate}
\end{lemma}

\begin{proof}
各項を順に証明する。
\begin{enumerate}
\item 命題 \ref{algebra-proposition-localize-twice} により、これはすべての
$\mathfrak p \in \Spec(R)$ に対して $M_\mathfrak p = 0$ を意味するので、
補題 \ref{algebra-lemma-characterize-zero-local} から結論を得る。
\item 各 $i$ に対して、$M_{f_i}$ の有限生成集合 $X_i$ を取る。
一般性を失わず、$X_i$ の元は局所化写像 $M \rightarrow M_{f_i}$ の像に
属すると仮定できるので、$M$ における $X_i$ の元の逆像からなる有限集合
$Y_i$ を取る。これらの集合の合併を $Y$ とする。これは依然として有限集合である。
基底元 $e_y$ を $y$ に写す明らかな $R$-線形写像
$R^Y \rightarrow M$ を考える。仮定により、この写像は $R$ の任意の素イデアル
$\mathfrak p$ で局所化すると全射なので、補題
\ref{algebra-lemma-characterize-zero-local} により全射であり、$M$ は有限生成である。
\item (2) により、短完全列
% P01-E0298: frozen English source rank/index collision; see ERRATA_CANDIDATES.jsonl.
$$
0 \rightarrow K \rightarrow R^n \rightarrow M \rightarrow 0
$$
を得る。局所化は完全関手であり、$M_{f_i}$ は有限表示なので、補題
\ref{algebra-lemma-extension} により、すべての $1 \leq i \leq n$ に対して
$K_{f_i}$ は有限生成である。(2) により、$K$ は有限 $R$-加群であり、
したがって $M$ は有限表示である。
\item 命題 \ref{algebra-proposition-localize-twice} により、この仮定は、
すべての素イデアルでの局所化上に誘導される射が同型であることを意味する。
したがって、補題 \ref{algebra-lemma-characterize-zero-local} から結論を得る。
\item 命題 \ref{algebra-proposition-localize-twice} により、この仮定は、
すべての素イデアルでの局所化からなる誘導列が短完全であることを意味する。
したがって、補題 \ref{algebra-lemma-characterize-zero-local} から結論を得る。
\item $R$ の任意のイデアルが有限生成集合をもつことを示す。
そのため、$I \subset R$ を任意のイデアルとする。命題
\ref{algebra-proposition-localization-exact} により、各 $I_{f_i} \subset R_{f_i}$ は
イデアルである。仮定により、これらはすべて有限生成なので、(2) から結論を得る。
\item 各 $i$ に対して、$S_{f_i}$ の有限生成集合 $X_i$ を取る。
一般性を失わず、$X_i$ の元は局所化写像 $S \rightarrow S_{f_i}$ の像に
属すると仮定できるので、$S$ における $X_i$ の元の逆像からなる有限集合
$Y_i$ を取る。これらの集合の合併を $Y$ とする。これは依然として有限集合である。
$Y$ により誘導される代数準同型 $R[X_y]_{y \in Y} \rightarrow S$ を考える。
これは代数準同型なので、その像 $T$ は $R$-加群 $S$ の $R$-部分加群であり、
商加群 $S/T$ を考えることができる。仮定により、これは各 $f_i$ で
局所化すると零なので、(1) により零である。したがって $S$ は有限型
$R$-代数である。
\item 前項により、全射 $R$-代数準同型
$R[X_1, \ldots, X_n] \rightarrow S$ が存在する。この写像の核を $K$ とする。
これは $R[X_1, \ldots, X_n]$ のイデアルであり、各 $f_i$ での局所化において
有限生成である。$f_i$ は $R$ で単位イデアルを生成するので、
$R[X_1, \ldots, X_n]$ でも単位イデアルを生成する。したがって、
(2) を適用すれば証明が完了する。
\end{enumerate}
\end{proof}

\begin{lemma}
\label{algebra-lemma-cover-upstairs}
$R \to S$ を環準同型とする。$g_1, \ldots, g_n$ を、
$\bigcup D(g_i) = \Spec(S)$、言い換えれば $(g_1, \ldots, g_n) = S$ を
満たす $S$ の元の有限列とする。
\begin{enumerate}
\item 各 $S_{g_i}$ が $R$ 上有限型ならば、$S$ は $R$ 上有限型である。
\item 各 $S_{g_i}$ が $R$ 上有限表示ならば、$S$ は $R$ 上有限表示である。
\end{enumerate}
\end{lemma}

\begin{proof}
$\sum h_i g_i = 1$ を満たす $h_1, \ldots, h_n \in S$ を選ぶ。

\medskip\noindent
(1) を証明する。各 $i$ に対し、$R$-代数として $S_{g_i}$ を生成する元の
有限列 $x_{i, j} \in S_{g_i}$、$j = 1, \ldots, m_i$ を選ぶ。
ある $y_{i, j} \in S$ および $n_{i, j} \ge 0$ を用いて
$x_{i, j} = y_{i, j}/g_i^{n_{i, j}}$ と書く。
$g_1, \ldots, g_n$、$h_1, \ldots, h_n$、および
$y_{i, j}$、$i = 1, \ldots, n$、$j = 1, \ldots, m_i$ により生成される
$R$-部分代数 $S' \subset S$ を考える。局所化は完全なので（命題
\ref{algebra-proposition-localization-exact}）、$S'_{g_i} \to S_{g_i}$ は単射である。
一方、$y_{i, j}$ の選び方により、これは全射である。
$h_1, \ldots, h_n \in S'$ なので、$g_1, \ldots, g_n$ は $S'$ で
単位イデアルを生成する。したがって、$S' \to S$ を $S'$-加群の写像と
みなせば、補題 \ref{algebra-lemma-cover} により同型である。

\medskip\noindent
(2) を証明する。$S$ が有限型であることはすでに分かっている。
あるイデアル $J$ に対して $S = R[x_1, \ldots, x_m]/J$ と書く。
各 $i$ に対し、$g_i$ の持ち上げ $g'_i \in R[x_1, \ldots, x_m]$ と
$h_i$ の持ち上げ $h'_i \in R[x_1, \ldots, x_m]$ を選ぶ。このとき、
$$
S_{g_i} = R[x_1, \ldots, x_m, y_i]/(J_i + (1 - y_ig'_i))
$$
である。ここで $J_i$ は $J$ により生成される
$R[x_1, \ldots, x_m, y_i]$ のイデアルである。細部は省略する。
補題 \ref{algebra-lemma-finite-presentation-independent} により、
$J_i$ における $f_{i, j}$ の像と $1 - y_ig'_i$ がイデアル
$J_i + (1 - y_ig'_i)$ を生成するような有限個の元
$f_{i, j} \in J$、$j = 1, \ldots, m_i$ を選べる。次のようにおく。
$$
S' = R[x_1, \ldots, x_m]/\left(\sum h'_ig'_i - 1, f_{i, j}; 
i = 1, \ldots, n, j = 1, \ldots, m_i\right)
$$
全射 $R$-代数準同型 $S' \to S$ が存在する。$S'$ における
$g'_1, \ldots, g'_n$ の類は単位イデアルを生成し、構成により写像
$S'_{g'_i} \to S_{g_i}$ は単射である。したがって、(1) と同様に結論を得る。
\end{proof}





\section{関数の貼り合わせ}
\label{algebra-section-tilde-module-sheaf}

\noindent
この節では、標準開集合による開被覆
$$
\Spec(R) = \bigcup\nolimits_{i = 1}^n D(f_i)
$$
と、各 $i$ に対する元 $h_i \in R_{f_i}$ が与えられ、
$R_{f_i f_j}$ の元として $h_i = h_j$ であるとき、一意な $h \in R$ が存在し、
その $h$ の $R_{f_i}$ における像が $h_i$ となることを示す。この結果は二通りに
解釈できる。
\begin{enumerate}
\item 対応 $D(f) \mapsto R_f$ は標準開集合上の環の層である。
Sheaves, Section \ref{sheaves-section-bases} を参照せよ。
\item $R_f$ の元を $D(f)$ 上の「代数的」または「正則」関数と考えると、
これらは、位相多様体上の連続関数、あるいは可微分多様体上の可微分関数と
同じように貼り合わさる。
\end{enumerate}

\begin{lemma}
\label{algebra-lemma-cover-module}
$R$ を環とする。$f_1, \ldots, f_n$ を単位イデアルを生成する $R$ の元とし、
$M$ を $R$-加群とする。列
$$
0 \to
M \xrightarrow{\alpha}
\bigoplus\nolimits_{i = 1}^n M_{f_i} \xrightarrow{\beta}
\bigoplus\nolimits_{i, j = 1}^n M_{f_i f_j}
$$
は完全である。ここで
$\alpha(m) = (m/1, \ldots, m/1)$ および
$\beta(m_1/f_1^{e_1}, \ldots, m_n/f_n^{e_n})
= (m_i/f_i^{e_i} - m_j/f_j^{e_j})_{(i, j)}$ である。
\end{lemma}

\begin{proof}
任意の極大イデアル $\mathfrak m$ におけるこの列の局所化が完全であることを
示せば十分である。補題 \ref{algebra-lemma-characterize-zero-local} を参照せよ。
$f_1, \ldots, f_n$ は単位イデアルを生成するので、
$f_i \not \in \mathfrak m$ を満たす $i$ が存在する。番号を付け替えることにより、
$i = 1$ と仮定してよい。命題
\ref{algebra-proposition-localize-twice-module} により、
$(M_{f_i})_\mathfrak m = (M_\mathfrak m)_{f_i}$ および
$(M_{f_if_j})_\mathfrak m = (M_\mathfrak m)_{f_if_j}$ である。
特に、$f_1$ は単元なので、
$(M_{f_1})_\mathfrak m = M_\mathfrak m$ および
$(M_{f_1 f_i})_\mathfrak m = (M_\mathfrak m)_{f_i}$ である。
この列の写像は、補題
\ref{algebra-lemma-universal-property-localization-module} と $M$ 上の恒等写像から来る
標準的な写像である。以上を踏まえて、$R$ を $R_\mathfrak m$、
$M$ を $M_\mathfrak m$、$f_i$ を $R_\mathfrak m$ におけるその像、
$f_1$ を $1 \in R_\mathfrak m$ で置き換えれば、$f_1 = 1$ の場合に帰着する。

\medskip\noindent
$f_1 = 1$ と仮定する。このとき $\alpha$ の単射性は自明である。
$m = (m_i) \in \bigoplus_{i = 1}^n M_{f_i}$ を $\beta$ の核の元とする。
すると $m_1 \in M_{f_1} = M$ である。さらに、$\beta(m) = 0$ であることから、
$m_1$ と $m_i$ は $M_{f_1f_i} = M_{f_i}$ の同じ元に写る。
したがって $\alpha(m_1) = m$ であり、証明が完了する。
\end{proof}

\begin{lemma}
\label{algebra-lemma-standard-covering}
$R$ を環とし、$f_1, f_2, \ldots f_n\in R$ は $R$ で単位イデアルを
生成するとする。
% P01-E0299: frozen English source punctuation candidate; see ERRATA_CANDIDATES.jsonl.
このとき、次の列は完全である。
$$
0 \longrightarrow
R \longrightarrow
\bigoplus\nolimits_i R_{f_i} \longrightarrow
\bigoplus\nolimits_{i, j}R_{f_if_j}
$$
ここで写像 $\alpha : R \longrightarrow \bigoplus_i R_{f_i}$ および
$\beta : \bigoplus_i R_{f_i} \longrightarrow \bigoplus_{i, j} R_{f_if_j}$ は
次のように定義される。
$$
\alpha(x) = \left(\frac{x}{1}, \ldots, \frac{x}{1}\right)
\text{ and }
\beta\left(\frac{x_1}{f_1^{r_1}}, \ldots, \frac{x_n}{f_n^{r_n}}\right)
=
\left(\frac{x_i}{f_i^{r_i}}-\frac{x_j}{f_j^{r_j}}~\text{in}~R_{f_if_j}\right).
$$
\end{lemma}

\begin{proof}
補題 \ref{algebra-lemma-cover-module} の特別な場合である。
\end{proof}

\noindent
次の結果はすでに上で見たが、便宜のためここで明示的に述べる。

\begin{lemma}
\label{algebra-lemma-disjoint-implies-product}
$R$ を環とする。$U$ と $V$ がともに開で
$\Spec(R) = U \amalg V$ ならば、$R \cong R_1 \times R_2$ であり、補題
\ref{algebra-lemma-spec-product} の写像を通じて $U \cong \Spec(R_1)$ および
$V \cong \Spec(R_2)$ である。さらに、$R_1$ と $R_2$ はいずれも環 $R$ の
局所化であり、かつ商でもある。
\end{lemma}

\begin{proof}
補題 \ref{algebra-lemma-disjoint-decomposition} により、ある冪等元 $e$ に対して
$U = D(e)$ および $V = D(1-e)$ である。補題
\ref{algebra-lemma-standard-covering} により、
$R \cong R_e \times R_{1 - e}$ である（明らかに $R_{e(1-e)} = 0$ なので
貼り合わせ条件は自明である。もちろん、この場合には積分解を直接証明することも
自明である）。これで補題が従う。
\end{proof}

\begin{lemma}
\label{algebra-lemma-when-injective-covering}
$R$ を環、$f_1, \ldots, f_n \in R$ を元、$M$ を $R$-加群とする。
このとき $M \to \bigoplus M_{f_i}$ が単射であることと、
$$
M \longrightarrow \bigoplus\nolimits_{i = 1, \ldots, n} M, \quad
m \longmapsto (f_1m, \ldots, f_nm)
$$
が単射であることは同値である。
\end{lemma}

\begin{proof}
写像 $M \to \bigoplus M_{f_i}$ が単射であることは、
$f_i^{e_i}m = 0$、$i = 1, \ldots, n$ を満たすすべての
$m \in M$ および $e_1, \ldots, e_n \geq 1$ に対して
$m = 0$ となることと同値である。これは明らかに、表示された写像が
単射であることを意味する。逆に、表示された写像が単射であると仮定し、
$f_i^{e_i}m = 0$、$i = 1, \ldots, n$ を満たす
$m \in M$ および $e_1, \ldots, e_n \geq 1$ を取る。
すべての $i$ に対して $e_i = 1$ ならば、表示された写像の単射性から直ちに
$m = 0$ を得る。次に、この主張がこのような任意のデータに対して成り立つことを、
$e = \sum e_i$ に関する帰納法で示す。基底の場合は $e = n$ であり、これは
今しがた扱った。ある $e_i > 1$ に対して $m' = f_im$ とおく。
帰納法により $m' = 0$ である。したがって $f_i m = 0$、すなわち
$e_i = 1$ としてよく、これにより $e$ が減少するので結論を得る。
\end{proof}

\noindent
次の補題は、平坦降下というより一般的な文脈で述べて証明する方がよい。
しかし、上の結果とよく調和するので、ここで述べることにも意味がある。

\begin{lemma}
\label{algebra-lemma-glue-modules}
$R$ を環とし、$f_1, \ldots, f_n \in R$ とする。次のデータが与えられたとする。
\begin{enumerate}
\item 各 $i$ に対する $R_{f_i}$-加群 $M_i$。
\item 各組 $i, j$ に対する $R_{f_if_j}$-加群同型
$\psi_{ij} : (M_i)_{f_j} \to (M_j)_{f_i}$。
\end{enumerate}
これらは、すべての図式
$$
\xymatrix{
(M_i)_{f_jf_k}
\ar[rd]_{\psi_{ij}}
\ar[rr]^{\psi_{ik}}
& &
(M_k)_{f_if_j} \\
&
(M_j)_{f_if_k} \ar[ru]_{\psi_{jk}}
}
$$
が可換であるという「コサイクル条件」を満たすとする
（すべての三つ組 $i, j, k$ に対して）。このデータから
$$
M = \Ker\left(
\bigoplus\nolimits_{1 \leq i \leq n} M_i
\longrightarrow
\bigoplus\nolimits_{1 \leq i, j \leq n} (M_i)_{f_j}
\right)
$$
を定義する。ここで $(m_1, \ldots, m_n)$ は、その $(i, j)$ 成分が
$m_i/1 - \psi_{ji}(m_j/1)$ である元に写る。このとき自然な写像
$M \to M_i$ は同型 $M_{f_i} \to M_i$ を誘導する。さらに、明らかな記法のもとで、
すべての $m \in M$ に対して $\psi_{ij}(m/1) = m/1$ である。
\end{lemma}

\begin{proof}
$M_{f_1} \to M_1$ が同型であることを示すには、$R_{f_1}$ のすべての素イデアル
$\mathfrak p'$ でのその局所化が同型であることを示せば十分である。
補題 \ref{algebra-lemma-characterize-zero-local} を参照せよ。補題
\ref{algebra-lemma-standard-open} により、ある素イデアル $\mathfrak p \subset R$、
$f_1 \not \in \mathfrak p$ に対して
$\mathfrak p' = \mathfrak p R_{f_1}$ と書ける。局所化は完全なので
（命題 \ref{algebra-proposition-localization-exact}）、
\begin{align*}
(M_{f_1})_{\mathfrak p'} & =
M_\mathfrak p \\
& =
\Ker\left(
\bigoplus\nolimits_{1 \leq i \leq n} M_{i, \mathfrak p}
\longrightarrow
\bigoplus\nolimits_{1 \leq i, j \leq n} ((M_i)_{f_j})_\mathfrak p
\right) \\
& =
\Ker\left(
\bigoplus\nolimits_{1 \leq i \leq n} M_{i, \mathfrak p}
\longrightarrow
\bigoplus\nolimits_{1 \leq i, j \leq n} (M_{i, \mathfrak p})_{f_j}
\right)
\end{align*}
である。ここでは命題 \ref{algebra-proposition-localize-twice-module} も用いた。
$f_1$ は $R_\mathfrak p$ の単元なので、$R$ を $R_\mathfrak p$、
$f_i$ を $R_\mathfrak p$ における $f_i$ の像、$M$ を $M_\mathfrak p$、
$f_1$ を $1$ で置き換えることにより、$f_1 = 1$ の場合に帰着する。

\medskip\noindent
$f_1 = 1$ と仮定する。このとき
$\psi_{1j} : (M_1)_{f_j} \to M_j$ は $j = 2, \ldots, n$ に対して同型である。
これらの同型を用いて $M_j = (M_1)_{f_j}$ と同一視すると、
$\psi_{ij} : (M_1)_{f_if_j} \to (M_1)_{f_if_j}$ は標準的な同一視である。
したがって複体
$$
0 \to M_1 \to
\bigoplus\nolimits_{1 \leq i \leq n} (M_1)_{f_i}
\longrightarrow
\bigoplus\nolimits_{1 \leq i, j \leq n}
(M_1)_{f_if_j}
$$
は補題 \ref{algebra-lemma-cover-module} により完全である。したがって、この場合、
最初の写像は $M_1$ と $M$ を同一視し、すべてが明らかになる。
\end{proof}













\section{零因子と全商環}
\label{algebra-section-total-quotient-ring}

\noindent
極小素イデアルにおける局所環は、次の性質をもつ。

\begin{lemma}
\label{algebra-lemma-minimal-prime-reduced-ring}
$\mathfrak p$ を環 $R$ の極小素イデアルとする。$R_{\mathfrak p}$ の
極大イデアルのすべての元は冪零である。$R$ が被約ならば、
$R_{\mathfrak p}$ は体である。
\end{lemma}

\begin{proof}
${\mathfrak p}R_{\mathfrak p}$ のある元 $x$ が冪零でないならば、補題
\ref{algebra-lemma-Zariski-topology} により $D(x) \not = \emptyset$ である。
これは $\mathfrak p$ の極小性に反する。$R$ が被約ならば
${\mathfrak p}R_{\mathfrak p} = 0$ であり、したがってこれは体である。
\end{proof}

\begin{lemma}
\label{algebra-lemma-reduced-ring-sub-product-fields}
$R$ を被約環とする。このとき、
\begin{enumerate}
\item $R$ は体の積の部分環である。
\item $R \to \prod_{\mathfrak p\text{ minimal}} R_{\mathfrak p}$ は
体の積への埋め込みである。
\item $\bigcup_{\mathfrak p\text{ minimal}} \mathfrak p$ は $R$ の零因子全体である。
\end{enumerate}
\end{lemma}

\begin{proof}
補題 \ref{algebra-lemma-minimal-prime-reduced-ring} により、各環 $R_\mathfrak p$ は体である。
特に、環準同型 $R \to R_\mathfrak p$ の核は $\mathfrak p$ である。
補題 \ref{algebra-lemma-Zariski-topology} により
$\bigcap_{\mathfrak p} \mathfrak p = (0)$ である。したがって (2) と (1) が成り立つ。
$x y = 0$ かつ $y \not = 0$ ならば、ある極小素イデアル $\mathfrak p$ に対して
$y \not \in \mathfrak p$ である。したがって $x \in \mathfrak p$ である。
ゆえに $R$ のすべての零因子は
$\bigcup_{\mathfrak p\text{ minimal}} \mathfrak p$ に含まれる。
逆に、ある極小素イデアル $\mathfrak p$ に対して $x \in \mathfrak p$ とする。
すると $x$ は $R_\mathfrak p$ で零に写るので、$xy = 0$ を満たす
$y \in R$、$y \not \in \mathfrak p$ が存在する。言い換えれば、$x$ は零因子である。
これで (3) と補題の証明が完了する。
\end{proof}

\noindent
環 $R$ の全商環 $Q(R)$ は例 \ref{algebra-example-localize-at-prime} で導入した。

\begin{lemma}
\label{algebra-lemma-total-ring-fractions}
$R$ を環とし、$S \subset R$ を非零因子からなる乗法的部分集合とする。
このとき $Q(R) \cong Q(S^{-1}R)$ である。特に
$Q(R) \cong Q(Q(R))$ である。
\end{lemma}

\begin{proof}
$x \in S^{-1}R$ が非零因子で、ある $r \in R$、$f \in S$ に対して
$x = r/f$ ならば、$r$ は $R$ の非零因子である。これで補題が従う。
\end{proof}

\noindent
貼り合わせの結果を用いて、例 \ref{algebra-example-localize-at-prime} で導入した
全商環 $Q(R)$ についての結果を証明できる。

\begin{lemma}
\label{algebra-lemma-total-ring-fractions-no-embedded-points}
$R$ を環とする。$R$ は有限個の極小素イデアル
$\mathfrak q_1, \ldots, \mathfrak q_t$ をもち、
$\mathfrak q_1 \cup \ldots \cup \mathfrak q_t$ は $R$ の零因子全体であると仮定する。
このとき全商環 $Q(R)$ は
$R_{\mathfrak q_1} \times \ldots \times R_{\mathfrak q_t}$ に等しい。
\end{lemma}

\begin{proof}
% P01-E0300: frozen English source membership wording candidate; see ERRATA_CANDIDATES.jsonl.
任意の非零因子は $R \setminus \mathfrak q_i$ に属するので、自然な写像
$Q(R) \to R_{\mathfrak q_i}$ がある。したがって自然な写像
$Q(R) \to R_{\mathfrak q_1} \times \ldots \times R_{\mathfrak q_t}$ がある。
任意の極小でない素イデアル $\mathfrak p \subset R$ に対して、補題
\ref{algebra-lemma-silly} により
$\mathfrak p \not \subset \mathfrak q_1 \cup \ldots \cup \mathfrak q_t$ である。
したがって
$\Spec(Q(R)) = \{\mathfrak q_1, \ldots, \mathfrak q_t\}$ である
（$\Spec(R)$ の部分集合として。補題 \ref{algebra-lemma-spec-localization} を参照せよ）。
ゆえに $\Spec(Q(R))$ は有限離散集合であり、補題
\ref{algebra-lemma-disjoint-implies-product} により、
$\Spec(A_i) = \{\mathfrak{q}_i\}$ を満たす
$Q(R) = A_1 \times \ldots \times A_t$ が得られる。さらに $A_i$ は
$R$ の局所化である局所環である。したがって $A_i \cong R_{\mathfrak q_i}$ である。
\end{proof}

\section{スペクトルの既約成分}
\label{algebra-section-irreducible}

\noindent
環のスペクトルの既約成分が、その環の極小素イデアルに対応することを示す。

\begin{lemma}
\label{algebra-lemma-irreducible}
$R$ を環とする。
\begin{enumerate}
\item 素イデアル $\mathfrak p \subset R$ に対し、ザリスキー位相における
$\{\mathfrak p\}$ の閉包は $V(\mathfrak p)$ である。式で書けば
$\overline{\{\mathfrak p\}} = V(\mathfrak p)$ である。
\item $\Spec(R)$ の既約閉部分集合は、素イデアル
$\mathfrak p \subset R$ に対する部分集合 $V(\mathfrak p)$ にちょうど一致する。
\item $\Spec(R)$ の既約成分（Topology, Definition
\ref{topology-definition-irreducible-components} を参照）は、極小素イデアル
$\mathfrak p \subset R$ に対する部分集合 $V(\mathfrak p)$ にちょうど一致する。
\end{enumerate}
\end{lemma}

\begin{proof}
$ \mathfrak p \in V(I)$ ならば $I \subset \mathfrak p$ であることに注意する。
したがって、明らかに
$\overline{\{\mathfrak p\}} = V(\mathfrak p)$ である。特に
$V(\mathfrak p)$ は一点集合の閉包であり、それゆえ既約である。
第二の主張から第三の主張が従う。
第二の主張を示すため、$I$ を根基イデアルとして
$V(I) \subset \Spec(R)$ とする。$I$ が素イデアルでなければ、
$a, b\in R$、$a, b\not \in I$、$ab\in I$ を満たす元を選べる。
この場合 $V(I, a) \cup V(I, b) = V(I)$ だが、補題
\ref{algebra-lemma-Zariski-topology} により、
$V(I, b) = V(I)$ でも $V(I, a) = V(I)$ でもない。
したがって $V(I)$ は既約でない。
\end{proof}

\noindent
言い換えると、この補題は $\Spec(R)$ の任意の既約閉部分集合が、
ある素イデアル $\mathfrak p$ に対して $V(\mathfrak p)$ の形であることを示す。
$V(\mathfrak p) = \overline{\{\mathfrak p\}}$ なので、
各既約閉部分集合は一意な生成点をもつことが分かる。Topology, Definition
\ref{topology-definition-generic-point} を参照せよ。
特に、$\Spec(R)$ はソバー位相空間である。
この事実を次の補題として記録する。

\begin{lemma}
\label{algebra-lemma-spec-spectral}
環のスペクトルはスペクトル空間である。Topology, Definition
\ref{topology-definition-spectral-space} を参照せよ。
\end{lemma}

\begin{proof}
形式的には、これは補題 \ref{algebra-lemma-irreducible} と補題
\ref{algebra-lemma-topology-spec} から従う。上の議論も参照せよ。
\end{proof}

\begin{lemma}
\label{algebra-lemma-irreducible-components-containing-x}
$R$ を環とする。$\mathfrak p \subset R$ を素イデアルとする。
\begin{enumerate}
\item $\mathfrak p$ を通る $\Spec(R)$ の既約閉部分集合全体は、
その局所環の素イデアル $\mathfrak q \subset R_{\mathfrak p}$
全体と一対一に対応する。
\item $\mathfrak p$ を通る $\Spec(R)$ の既約成分全体は、
その局所環の極小素イデアル $\mathfrak q \subset R_{\mathfrak p}$
全体と一対一に対応する。
\end{enumerate}
\end{lemma}

\begin{proof}
補題 \ref{algebra-lemma-irreducible} と、$\Spec(R_\mathfrak p)$ が
$\mathfrak q \subset \mathfrak p$ を満たす $R$ の素イデアル
$\mathfrak q$ に対応することを示す補題
\ref{algebra-lemma-spec-localization} の $\Spec(R_\mathfrak p)$ の記述から従う。
\end{proof}

\begin{lemma}
\label{algebra-lemma-standard-open-containing-maximal-point}
$R$ を環とする。$\mathfrak p$ を $R$ の極小素イデアルとする。
$W \subset \Spec(R)$ を、点 $\mathfrak p$ を含まない準コンパクト開集合とする。
このとき、$D(f) \cap W = \emptyset$ を満たす
$f \in R$、$f \not \in \mathfrak p$ が存在する。
\end{lemma}

\begin{proof}
$W$ は準コンパクトなので、標準アファイン開集合 $D(g_i)$、
$i = 1, \ldots, n$ の有限和として書ける。
$\mathfrak p \not \in W$ なので、すべての $i$ に対して
$g_i \in \mathfrak p$ である。補題
\ref{algebra-lemma-minimal-prime-reduced-ring} により、各 $g_i$ は
$R_{\mathfrak p}$ において冪零である。したがって、すべての $i$ について、
ある $n_i > 0$ に対して $f g_i^{n_i} = 0$ となる
$f \in R$、$f \not \in \mathfrak p$ を取れる。この $D(f)$ が求めるものである。
\end{proof}

\begin{lemma}
\label{algebra-lemma-ring-with-only-minimal-primes}
$R$ を環とし、位相空間として $X = \Spec(R)$ とする。
次は同値である。
\begin{enumerate}
\item $X$ は副有限である。
\item $X$ はハウスドルフである。
\item $X$ は完全不連結である。
\item $X$ の任意の準コンパクト開集合は閉である。
\item その素イデアルの間に自明でない包含関係がない。
\item 任意の素イデアルは極大イデアルである。
\item 任意の素イデアルは極小イデアルである。
\item 任意の標準開集合 $D(f) \subset X$ は閉である。
% P01-E0301: frozen English source unfinished placeholder; see ERRATA_CANDIDATES.jsonl.
\item ここにさらに追加する。
\end{enumerate}
\end{lemma}

\begin{proof}
第一の証明。(5)、(6)、(7) が同値であることは明らかである。
任意の準コンパクト開集合は標準開集合の有限和なので、
(4) と (8) が同値であることも明らかである。
(7) $\Rightarrow$ (4) は補題
\ref{algebra-lemma-standard-open-containing-maximal-point} から従う。
(4) を仮定する。$\mathfrak p, \mathfrak p'$ を $R$ の相異なる
素イデアルとする。$f \in \mathfrak p'$、$f \not \in \mathfrak p$
を選ぶ（必要なら $\mathfrak p$ と $\mathfrak p'$ を入れ替える）。
すると $\mathfrak p' \not \in D(f)$ かつ $\mathfrak p \in D(f)$ である。
(4) により、開集合 $D(f)$ は閉でもある。
したがって $\mathfrak p$ と $\mathfrak p'$ は、和が $X$ となる
互いに交わらない開近傍に属する。ゆえに $X$ はハウスドルフかつ完全不連結である。
したがって (4) $\Rightarrow$ (2) および (3) である。
(3) が成り立てば、$\Spec(R)$ の点の間に特殊化は存在し得ず、
(5) が成り立つことが分かる。$X$ がハウスドルフならば各点は閉なので、
(2) から (6) が従う。したがって (2)、(3)、(4)、(5)、(6)、(7)、(8)
は同値である。任意の副有限空間はハウスドルフなので、(1) から (2) が従う。
$X$ が (2) と (3) を満たすならば、補題
\ref{algebra-lemma-quasi-compact} により $X$ は準コンパクトでもあるから、
Topology, Lemma \ref{topology-lemma-profinite} により副有限である。

\medskip\noindent
第二の証明。(4) と (8) の同値性を除けば、これは補題
\ref{algebra-lemma-spec-spectral} と純粋に位相的な事実から従う。Topology, Lemma
\ref{topology-lemma-characterize-profinite-spectral} を参照せよ。
\end{proof}













\section{環のスペクトルの例}
\label{algebra-section-examples-spectra}

\noindent
この節ではスペクトルの例をいくつか挙げる。

\begin{example}
\label{algebra-example-spec-Zxmodx2minus4}
この例では $X = \Spec(\mathbf{Z}[x]/(x^2 - 4))$ を記述する。
$\mathfrak{p}$ を $X$ の任意の素イデアルとする。
$\phi : \mathbf{Z} \to \mathbf{Z}[x]/(x^2 - 4)$ を自然な環写像とする。
すると $ \phi^{-1}(\mathfrak p)$ は $\mathbf{Z}$ の素イデアルである。
$ \phi^{-1}(\mathfrak p) = (2)$ ならば、$\mathfrak p$ は $2$ を含むので、
写像 $ \mathbf{Z}[x]/(x^2 - 4) \to  \mathbf{Z}[x]/(x^2 - 4, 2)$
を介して
$\mathbf{Z}[x]/(x^2 - 4, 2) \cong (\mathbf{Z}/2\mathbf{Z})[x]/(x^2)$
の素イデアルに対応する。
$(\mathbf{Z}/2\mathbf{Z})[x]/(x^2)$ の任意の素イデアルは、
$(x^2)$ を含む $(\mathbf{Z}/2\mathbf{Z})[x]$ の素イデアルに対応する。
したがって、そのような素イデアルは $x$ を含む。
$(\mathbf{Z}/2\mathbf{Z}) \cong (\mathbf{Z}/2\mathbf{Z})[x]/(x)$ は体なので、
$(x)$ は極大イデアルである。任意の素イデアルは $(x)$ を含み、
$(x)$ は極大なので、この環はただ一つの素イデアル $(x)$ しかもたない。
したがって、この場合 $\mathfrak p = (2, x)$ である。
次に $q > 2$ に対して $ \phi^{-1}(\mathfrak p) = (q)$ ならば、
$\mathfrak p$ は $q$ を含むので、写像
$ \mathbf{Z}[x]/(x^2 - 4) \to  \mathbf{Z}[x]/(x^2 - 4, q)$ を介して
$\mathbf{Z}[x]/(x^2 - 4, q) \cong (\mathbf{Z}/q\mathbf{Z})[x]/(x^2 - 4)$
の素イデアルに対応する。
$(\mathbf{Z}/q\mathbf{Z})[x]/(x^2 - 4)$ の任意の素イデアルは、
$(x^2 - 4) = (x -2)(x + 2)$ を含む
$(\mathbf{Z}/q\mathbf{Z})[x]$ の素イデアルに対応する。
したがって、これらの素イデアルは $x -2$ または $x + 2$ のいずれかを含む。
$(\mathbf{Z}/q\mathbf{Z})[x]$ は主イデアル整域なので、
すべての非零素イデアルは極大である。ゆえに
$(x-2)(x + 2)$ を含む $(\mathbf{Z}/q\mathbf{Z})[x]$ の素イデアルは
ちょうど二つ、すなわち $(x-2)$ と $(x + 2)$ である。
結論として、$2 \neq -2 \in \mathbf{Z}/(q)$ なので、
二つの素イデアル $(q, x-2)$ と $(q, x + 2)$ が存在する。
最後に $\phi^{-1}(\mathfrak p) = (0)$ の場合を扱う。
$\mathfrak p$ は $(x^2 - 4) = (x -2)(x + 2)$ を含む
$\mathbf{Z}[x]$ の素イデアルに対応することに注意する。
したがって $\mathfrak p$ は $(x-2)$ または $(x + 2)$ のいずれかを含む。
ゆえに、仮定により、$\mathfrak p$ は
$\mathbf{Z}$ と $0$ でしか交わらない
$\mathbf{Z}[x]/(x - 2)$ または $\mathbf{Z}[x]/(x + 2)$ の素イデアルに対応する。
$\mathbf{Z}[x]/(x - 2) \cong \mathbf{Z}$ かつ
$\mathbf{Z}[x]/(x + 2) \cong \mathbf{Z}$ なので、
これは $\mathfrak p$ がこれらの環の一方における $0$ に対応しなければならないことを意味する。
したがって、もとの環で
$\mathfrak p = (x - 2)$ または $\mathfrak p = (x + 2)$ である。
\end{example}

\begin{example}
\label{algebra-example-spec-Zx}
この例では $X = \Spec(\mathbf{Z}[x])$ を記述する。
$\mathfrak p \in X$ を固定する。
$\phi : \mathbf{Z} \to \mathbf{Z}[x]$ とし、
$\phi^{-1}(\mathfrak p) \in \Spec(\mathbf{Z})$ であることに注意する。
% P01-E0302: frozen English source positive-characteristic fibre classification candidate; see ERRATA_CANDIDATES.jsonl.
$q$ を素数 $q > 0$ として $\phi^{-1}(\mathfrak p) = (q)$ ならば、
$\mathfrak p$ は $(\mathbf{Z}/(q))[x]$ の素イデアルに対応し、
これは $(\mathbf{Z}/(q))[x]$ で既約な多項式により生成されなければならない。
% P01-E0303: frozen English source minimal-degree lift and malformed quotient-ring candidate; see ERRATA_CANDIDATES.jsonl.
この多項式の最小次数の代表元を選べば、それは $\mathbf{Z}[x]$ でも既約である。
したがって、この場合 $\mathfrak p = (q, f_q)$ である。ここで $f_q$ は、
$(\mathbf{Z}/(q) [x])$ で見ても既約な $\mathbf{Z}[x]$ の既約多項式である。
% P01-E0304: frozen English source generic-fibre zero-ideal omission candidate; see ERRATA_CANDIDATES.jsonl.
次に $\phi^{-1}(\mathfrak p) = (0)$ と仮定する。
この場合、$\mathfrak p$ は非定数多項式によって生成されなければならず、
$\mathfrak p$ は素イデアルなので、それらは $\mathbf{Z}[x]$ で既約と仮定してよい。
ガウスの補題により、これらの多項式は $\mathbf{Q}[x]$ でも既約である。
$\mathbf{Q}[x]$ はユークリッド整域なので、$\mathfrak p$ を生成する
相異なる既約元 $f, g$ が少なくとも二つあれば、
$a, b \in \mathbf{Q}[x]$ に対して $1 = af + bg$ となる。
共通分母を掛ければ、$\bar{a}, \bar{b} \in \mathbf{Z}[x]$ および
非零な $m \in \mathbf{Z}$ に対して
$m = \bar{a}f + \bar{b} g$ となることが分かる。これは矛盾である。
したがって、$\mathfrak p$ は $\mathbf{Z}[x]$ の一つの既約多項式で生成される。
\end{example}

\begin{example}
\label{algebra-example-spec-kxy}
この例では、$k$ を任意の体とするときの
$X = \Spec(k[x, y])$ を記述する。
明らかに $(0)$ は素イデアルであり、$k[x, y]$ は一意分解整域なので、
既約多項式で生成される任意の単項イデアルも素イデアルである。
ここで $\mathfrak p$ を、単項でない $X$ の元と仮定する。
$k[x, y]$ はネーター一意分解整域なので、素イデアル $\mathfrak p$ は
有限個の既約多項式 $(f_1, \ldots, f_n)$ で生成できる。
ここで、同伴でない $k[x, y]$ の既約多項式 $f, g$ に対して
$(f, g) \cap k[x] \neq 0$ であると主張する。これを示すには、
$k(x)[y]$ で見たときに $f$ と $g$ が互いに素であることを示せば十分である。
この場合、$k(x)[y]$ はユークリッド整域なので、ユークリッドの互除法を適用して
分母を払えば、
% P01-E0305: frozen English source coefficient-ring candidate; see ERRATA_CANDIDATES.jsonl.
$p, a, b \in k[x]$ に対して $p = af + bg$ を得る。
したがって、そうでない、すなわち、ある非単元 $h \in k(x)[y]$ が
$f$ と $g$ の両方を割ると仮定する。
% P01-E0306: frozen English source ambient-ring conflict candidate; see ERRATA_CANDIDATES.jsonl.
するとガウスの補題により、ある $a, b \in k(x)$ に対して、
$ah, bh \in k[x]$ であり、$ah | f$ および $bh | g$ となる。
既約性により、$ah = f$ かつ $bh = g$ である
（$h \notin k(x)$ だからである）。したがって $k(x)[y]$ に戻ると、
$\frac{a}{b} g = f$ なので $f, g $ は同伴である。
$k(x)$ は $k[x]$ の分数体なので、共通因子をもたない元
$r , s \in k[x]$ により $g = \frac{r}{s} f $ と書ける。
これは $k[x, y]$ において $sg = rf$ であることを意味し、
$k[x, y]$ は一意分解整域なので、$s$ は $f$ を割らなければならない。
したがって $s = 1$ または $s = f$ である。
$s = f$ ならば $r = g$ となり、$f, g \in k[x]$ を意味する。
したがって、それらは $k(x)$ の単元であり、
$k(x)[y]$ で互いに素でなければならず、仮定に矛盾する。
$s = 1$ ならば $g = rf$ となり、これも矛盾である。
ゆえに、ユークリッド整域 $k(x)[y]$ において
$f, g$ は互いに素でなければならない。
こうして、$\mathfrak p$ が
$p \in k[x] \subset k[x, y]$ なる既約多項式を含む場合に帰着した。
上の議論により、$\mathfrak p$ は環
$k[x, y]/(p) = k(\alpha)[y]$ の素イデアルに対応する。
ここで $\alpha$ は、最小多項式 $p$ をもち $k$ 上代数的な元である。
これは主イデアル整域なので、任意の素イデアルは
$(0)$ または $k(\alpha)[y]$ の既約多項式に対応する。
したがって、$\mathfrak p$ は $(p)$ または $(p, f)$ の形である。
ここで $f$ は、商 $k[x, y]/(p)$ で既約な $k[x, y]$ の多項式である。
\end{example}

\begin{example}
\label{algebra-example-affine-open-not-standard}
環
$$
R = \{ f \in \mathbf{Q}[z]\text{ with }f(0) = f(1) \}.
$$
を考える。写像
$$
\varphi : \mathbf{Q}[A, B] \to R
$$
を $\varphi(A) = z^2-z$ および $\varphi(B) = z^3-z^2$ により定める。
$(A^3 - B^2 + AB) \subset \Ker(\varphi)$ であり、
$A^3 - B^2 + AB$ が既約であることは容易に確かめられる。
$\varphi$ が全射であると仮定する。すると $R$ は整域
（整域の部分環である）なので、$\Ker(\varphi)$ は
$\mathbf{Q}[A, B]$ の素イデアルでなければならない。
% P01-E0307: frozen English source maximal-ideal classification candidate; see ERRATA_CANDIDATES.jsonl.
$(A^3-B^2 + AB)$ を含む素イデアルは、$(A^3-B^2 + AB)$ 自身と、
$f$ が $g$ を法として既約であるような
$f, g\in\mathbf{Q}[A, B]$ による任意の極大イデアル $(f, g)$ である。
しかし $R$ は体ではないので、核は $(A^3-B^2 + AB)$ でなければならない。
% P01-E0308: frozen English source induced-isomorphism direction candidate; see ERRATA_CANDIDATES.jsonl.
したがって $\varphi$ は同型
$R \to \mathbf{Q}[A, B]/(A^3-B^2 + AB)$ を与える。

\medskip\noindent
$\varphi$ が全射であることを確かめるには、任意の $f\in R$ を
$A(z) = z^2-z$ と $B(z) = z^3-z^2$ の $\mathbf{Q}$ 係数多項式として
表さなければならない。関係 $zA(z) = B(z)$ に注意する。
$a = f(0) = f(1)$ とする。このとき $z(z-1)$ は $f(z)-a$ を割るので、
$f(z) = z(z-1)g(z)+a = A(z)g(z)+a$ と書ける。
% P01-E0309: frozen English source polynomial-name discontinuity candidate; see ERRATA_CANDIDATES.jsonl.
$\deg(g) < 2$ ならば、
$h(z) = c_1z + c_0$ かつ
$f(z) = A(z)(c_1z + c_0)+a = c_1B(z)+c_0A(z)+a$ となるので完了である。
% P01-E0310: frozen English source dropped additive-constant candidate; see ERRATA_CANDIDATES.jsonl.
$\deg(g)\geq 2$ ならば、多項式の除法アルゴリズムにより
$g(z) = A(z)h(z)+b_1z + b_0$ と書ける
（$\deg(h)\leq\deg(g)-2$）。したがって
$f(z) = A(z)^2h(z)+b_1B(z)+b_0A(z)$ である。
$h(z)$ に除法を適用して反復すれば、$f(z)$ を $A(z)$ と $B(z)$ の
多項式として表す式を得る。ゆえに $\varphi$ は全射である。

\medskip\noindent
いま $a \in \mathbf{Q}$、$a \neq 0, \frac{1}{2}, 1$ とし、
$$
R_a = \{ f \in \mathbf{Q}[z, \frac{1}{z-a}]\text{ with }f(0) = f(1)
\}.
$$
を考える。これも有限生成 $\mathbf{Q}$-代数である。実際、関数
$z^2-z$、$z^3-z$、$\frac{a^2-a}{z-a}+z$ が
$R_a$ を $\mathbf{Q}$-代数として生成することは容易に確かめられる。
次の包含関係がある。
$$
R\subset R_a\subset\mathbf{Q}[z, \frac{1}{z-a}], \quad
R\subset\mathbf{Q}[z]\subset\mathbf{Q}[z, \frac{1}{z-a}].
$$
環 T と乗法的部分集合 $S\subset T$ に対して、環写像
$T \to S^{-1}T$ はスペクトル上の写像
$\Spec(S^{-1}T) \to \Spec(T)$ を誘導し、これは部分集合
$$
\{\mathfrak p \in \Spec(T) \mid S \cap \mathfrak p = \emptyset\}
\subset \Spec(T).
$$
の上への同相写像であることを思い出そう（補題
\ref{algebra-lemma-spec-localization}）。
$f\in T$ に対して $S = \{ 1, f, f^2, \ldots\}$ のとき、これは開集合
% P01-E0311: frozen English source ambient-space membership candidate; see ERRATA_CANDIDATES.jsonl.
$D(f)\subset T$ である。次に、環写像 $R \to R_a$ について対応する性質を確かめる。
包含 $R\subset R_a$ が誘導する写像
$\theta : \Spec(R_a) \to \Spec(R)$ について、$\theta$ が単射な局所同相写像であることを確かめ、
$\Spec(R)$ の開部分集合の上への同相写像であることを示す。
これを、これから記述する二つの標準開集合による
$\Spec(R_a)$ の開被覆に関して行う。
任意の $r\in\mathbf{Q}$ に対し、
$\text{ev}_r : R \to \mathbf{Q}$ を $r$ での評価により与えられる準同型とする。
$r = 0$ および $r = 1-a$ のとき、これは準同型
$\text{ev}_r' : R_a \to \mathbf{Q}$ に拡張できることに注意する
（後者については $a\neq\frac{1}{2}$ なので、
$z = 1-a$ で $\frac{1}{z-a}$ が定義されるためである）。
しかし $\text{ev}_a$ は $R_a$ に拡張されない。
$\mathfrak{m}_r = \Ker(\text{ev}_r)$ と書く。次が成り立つ。
$$
\mathfrak{m}_0 = (z^2-z, z^3-z),
$$
$$
\mathfrak{m}_a = ((z-1 + a)(z-a), (z^2-1 + a)(z-a)), \text{ and}
$$
$$
\mathfrak{m}_{1-a} = ((z-1 + a)(z-a), (z-1 + a)(z^2-a)).
$$
これを確かめるには、右辺が明らかに左辺に含まれることに注意する。
次に、生成元を $A$ と $B$ で表し、
$R$ を $\mathbf{Q}[A, B]/(A^3-B^2 + AB)$ と見て、
右辺が極大イデアルであることを確かめる。
$\mathfrak{m}_a$ は $\theta$ の像に含まれないことに注意する。実際、
$$
(z^2 - z)^2(z - a)\left(\frac{a^2 - a}{z - a} + z\right) =
(z^2 - z)^2(a^2 - a) + (z^2 - z)^2(z - a)z
$$
である。左辺は $\mathfrak m_a R_a$ に属する。なぜなら
$(z^2 - z)(z - a)$ は $\mathfrak m_a$ に属し、
$(z^2 - z)(\frac{a^2 - a}{z - a} + z)$ は $R_a$ に属するからである。
同様に、元 $(z^2 - z)^2(z - a)z$ は $\mathfrak m_a R_a$ に属する。
なぜなら $(z^2 - z)$ は $R_a$ に属し、
$(z^2 - z)(z - a)$ は $\mathfrak m_a$ に属するからである。
$a \not \in \{0, 1\}$ なので、$(z^2 - z)^2 \in \mathfrak m_a R_a$
と結論できる。したがって $R_a$ のどのイデアル $I$ も
$I \cap R = \mathfrak m_a$ を満たし得ない。というのも、そのような $I$ は
$(z^2 - z)^2$ を含まなければならないが、これは $R$ に属し
$\mathfrak m_a$ には属さないからである。
標準開集合 $D((z-1 + a)(z-a))\subset\Spec(R)$ は、閉集合
$\{\mathfrak{m}_a, \mathfrak{m}_{1-a}\}$ の補集合に等しい。
次に
$R_{(z-1 + a)(z-a)} = (R_a)_{(z-1 + a)(z-a)}$ であることを確かめる。
この局所化された環を $R'$ と呼ぶと、写像 $R \to R'$ は
$R \to R_a \to R'$ と分解される。すると補題
\ref{algebra-lemma-spec-localization} により、これらの写像は
$\Spec(R') \subset \Spec(R_a)$ および
$\Spec(R') \subset \Spec(R)$ を開部分集合として表す。
したがって、$\theta : \Spec(R_a) \to \Spec(R)$ を
$D((z-1 + a)(z-a))$ に制限すると、開部分集合の上への同相写像となる。
同様に、$\theta$ を
$D((z^2 + z + 2a-2)(z-a)) \subset \Spec(R_a)$ に制限すると、
開部分集合 $D((z^2 + z + 2a-2)(z-a)) \subset \Spec(R)$
の上への同相写像となる。
$z^2 + z + 2a-2$ が $\mathbf{Q}$ 上既約かどうかに応じて、
後者の標準開集合の補集合は、一つまたは二つの閉点と
閉点 $\mathfrak{m}_a$ からなる。さらに、
% P01-E0312: frozen English source changed generator candidate; see ERRATA_CANDIDATES.jsonl.
$R_a$ において元 $(z^2 + z + 2a-a)(z-a)$ と
$(z-1 + a)(z-a)$ が生成するイデアルは $R_a$ 全体なので、
これら二つの標準開集合は $\Spec(R_a)$ を被覆する。
したがって $\theta$ が
$\Spec(R)-\{\mathfrak{m}_a\}$ の上への同相写像であることを示すには、
これら一つまたは二つの点が $\mathfrak{m}_{1-a}$ と決して等しくならないことを
示せば十分である。実際そのとおりである。$1-a$ が
$z^2 + z + 2a-2$ の根であるのは、$a = 0$ または $a = 1$ のとき、
かつそのときに限るが、いずれも起こらないからである。

\medskip\noindent
この同相写像は $R$ の元による局所化の挙動を模倣しているが、
% P01-E0313: frozen English source localization-at-maximal-ideal candidate; see ERRATA_CANDIDATES.jsonl.
$\mathbf{Q}[z, \frac{1}{z-a}]$ が極大イデアル $(z-a)$ における
$\mathbf{Q}[z]$ の局所化である一方、環 $R_a$ は $R$ の局所化では
{\it ない}。
% P01-E0314: frozen English source overbroad localization-units candidate; see ERRATA_CANDIDATES.jsonl.
任意の局所化 $S^{-1}R$ では、もとの環 $R$ より単元が増えるからである。
$R$ の単元は、$\mathbf{Q}$ の単元 $\mathbf{Q}^\times$ である。
実際、$R_a$ の単元が $\mathbf{Q}^*$ であることは容易に分かる。
すなわち、$\mathbf{Q}[z, \frac{1}{z - a}]$ の単元は、
$c \in \mathbf{Q}^*$ および $n \in \mathbf{Z}$ に対する
$c (z - a)^n$ であり、これらが $R_a$ に属するのは
$n = 0$ のときに限ることは明らかである。
したがって $R_a$ は $R$ より多くの単元をもたず、
それゆえ $R$ の局所化ではあり得ない。

\medskip\noindent
% P01-E0315: frozen English source unbound-variable candidate; see ERRATA_CANDIDATES.jsonl.
$z = 0, 1$ で $\frac{1}{z-a}$ が意味をもつことを保証するために、
$a\neq 0, 1$ という事実を用いた。$a\neq 1/2$ という事実は数箇所で用いた。
(1) $R_a$ 上の $\text{ev}_{1-a}$ の核を論じられるようにするためである。
これにより $\mathfrak{m}_{1-a}$ が $R_a$ の点であること
（すなわち、$R_a$ が $R$ の点をちょうど一つだけ欠くこと）が保証される。
(2) 最後に、$(z-a)^{k + \ell}$ が $R$ に属し得るのは
$k = \ell = 0$ のときだけだと結論するためである。実際、
$a = 1/2$ ならば、$k + \ell$ が偶数である限りこれは $R$ に属する。
したがって、この場合 $R_a$ は実際に $R$ より多くの単元をもち、
$R_a$ は $R$ の局所化である可能性がある。
\end{example}

\section{素イデアルについてのメタ観察}
\label{algebra-section-oka-families}

\noindent
この節は CRing project から取られている。$R$ を環とし、
$S \subset R$ を乗法的部分集合とする。補題
\ref{algebra-lemma-spec-localization} の帰結として、
$S$ と交わらないという性質に関して極大なイデアル
$I \subset R$ は素イデアルである。その理由は、環 $S^{-1}R$ の
ある極大イデアル $\mathfrak m$ に対して
$I = R \cap \mathfrak m$ となるからである。
実は、イデアルの多くの性質について、その性質をもつ極大なものは素イデアルである。
これを見る一般的方法が \cite{Lam-Reyes} で展開された。
この節では本筋を離れ、この現象を説明する。

\medskip\noindent
$R$ を環とする。$I$ が $R$ のイデアルで $a \in R$ ならば、
$$
(I : a) = \left\{ x \in R \mid xa \in I\right\}.
$$
と定める。より一般に、$J \subset R$ がイデアルならば、
$$
(I : J) = \left\{ x \in R \mid xJ \subset I\right\}.
$$
と定める。

\begin{lemma}
\label{algebra-lemma-colon}
$R$ を環とする。単項イデアル $J \subset R$ と任意のイデアル
$I \subset J$ に対して $I = J (I : J)$ である。
\end{lemma}

\begin{proof}
$J = (a)$ とする。このとき $(I : J) = (I : a)$ である。
$I \subset J$ なので、任意の $y \in I$ は、ある
$x \in (I : a)$ によって $y = xa$ の形に書ける。
したがって $I \subset J (I : J)$ である。
逆に $x \in (I : a)$ ならば、$xJ = (xa) \subset I$ であり、
もう一方の包含が証明される。
\end{proof}

\noindent
$\mathcal{F}$ を $R$ のイデアルの集まりとする。
$\mathcal{F}$ の補集合の極大元が素イデアルとなることを保証する条件に関心がある。

\begin{definition}
\label{algebra-definition-oka-family}
$R$ を環とし、$\mathcal{F}$ を $R$ のイデアルの集合とする。
$R \in \mathcal{F}$ であり、イデアル $I \subset R$ と元 $a \in R$ に対して
$(I : a), (I, a) \in \mathcal{F}$ ならば常に
$I \in \mathcal{F}$ となるとき、$\mathcal{F}$ を {\it 岡族} という。
\end{definition}

\noindent
岡族の例をいくつか挙げよう。最初の例は、この節の導入で論じた基本例である。

\begin{example}
\label{algebra-example-oka-family-not-meet-multiplicative-set}
$R$ を環、$S$ を $R$ の乗法的部分集合とする。
$\mathcal{F} = \{I \subset R \mid I \cap S \not = \emptyset\}$ が
岡族であると主張する。実際、ある $a \in R$ に対して
$(I : a), (I, a) \in \mathcal{F}$ と仮定する。
$s \in (I, a) \cap S$ と $s' \in (I : a) \cap S$ を選ぶ。
すると $ss' \in I \cap S$ なので $I \in \mathcal{F}$ である。
したがって $\mathcal{F}$ は岡族である。
\end{example}

\begin{example}
\label{algebra-example-oka-family-finitely-generated}
$R$ を環、$I \subset R$ をイデアル、$a \in R$ とする。
$(I : a)$ が $a_1, \ldots, a_n$ で生成され、$(I, a)$ が
$b_1, \ldots, b_m \in I$ を用いて
$a, b_1, \ldots, b_m$ で生成されるならば、$I$ は
$aa_1, \ldots, aa_n, b_1, \ldots, b_m$ で生成される。
これを見るため、$x \in I$ ならば、$x \in (I, a)$ は
$a, b_1, \ldots, b_m$ の線形結合だが、
$a$ の係数は $(I:a)$ に属さなければならないことに注意する。
その結果、有限生成イデアルの族が岡族であることが分かる。
\end{example}

\begin{example}
\label{algebra-example-oka-family-principal}
環 $R$ の単項イデアルの族が岡族であることを示そう。
実際、$I \subset R$ をイデアル、$a \in R$ とし、
$(I, a)$ と $(I : a)$ が単項であると仮定する。
$(I : a) = (I : (I, a))$ であることに注意する。
$J = (I, a)$ と置くと、$J$ は単項であり $(I : J)$ も単項である。
補題 \ref{algebra-lemma-colon} により $I = J (I : J)$ である。
したがって、この状況では $J = (I, a)$ と $(I : J)$ が単項なので、
$I$ も単項である。
\end{example}

\begin{example}
\label{algebra-example-oka-family-bound-cardinality}
$R$ を環とする。$\kappa$ を無限基数とする。
高々 $\kappa$ 個の元で生成できるイデアルの族は岡族である。
議論は例 \ref{algebra-example-oka-family-finitely-generated} の議論と同様なので省略する。
\end{example}

\begin{example}
\label{algebra-example-oka-family-property-modules}
$A$ を環、$I \subset A$ をイデアル、$a \in A$ を元とする。
最初の矢印が $a$ による乗法で与えられる短完全列
$0 \to A/(I : a) \to A/I \to A/(I, a) \to 0$ がある。
したがって、$P$ が拡大に関して安定で $0$ に対して成り立つ
$A$-加群の性質ならば、$A/I$ が $P$ をもつようなイデアル $I$ の族は岡族である。
\end{example}

\begin{proposition}
\label{algebra-proposition-oka}
$\mathcal{F}$ がイデアルの岡族ならば、$\mathcal{F}$ の補集合の任意の極大元は
素イデアルである。
\end{proposition}

\begin{proof}
$I \not \in \mathcal{F}$ が $\mathcal{F}$ に属さないことに関して極大だが、
$I$ は素イデアルでないと仮定する。$R \in \mathcal{F}$ なので
$I \not = R$ であることに注意する。$I$ は素イデアルでないので、
$ab \in I$ を満たす $a, b \in R - I$ を取れる。
すると $(I, a) \neq I$ であり、$(I : a)$ は $b \not \in I$ を含むので
$(I : a) \neq I$ でもある。したがって $(I : a), (I, a)$ はいずれも
$I$ を真に含み、それゆえ $\mathcal{F}$ に属さなければならない。
岡条件により $I \in \mathcal{F}$ となり、矛盾である。
\end{proof}

\noindent
ここまでで、上の例の大部分を環の素イデアルに関する補題へ変換できる。

\begin{lemma}
\label{algebra-lemma-simple}
$R$ を環とし、$S$ を $R$ の乗法的部分集合とする。
$I \cap S = \emptyset$ という性質に関して極大なイデアル
$I \subset R$ は素イデアルである。
\end{lemma}

\begin{proof}
これはこの節の導入で論じた例である。別の証明として、例
\ref{algebra-example-oka-family-not-meet-multiplicative-set} と命題
\ref{algebra-proposition-oka} を組み合わせればよい。
\end{proof}

\begin{lemma}
\label{algebra-lemma-cohen}
$R$ を環とする。
\begin{enumerate}
\item 有限生成でないことに関して極大なイデアル $I \subset R$ は素イデアルである。
\item $R$ のすべての素イデアルが有限生成ならば、$R$ のすべてのイデアルが
有限生成である\footnote{後で $R$ はネーター的であるという。}。
\end{enumerate}
\end{lemma}

\begin{proof}
第一の主張は、例 \ref{algebra-example-oka-family-finitely-generated} と命題
\ref{algebra-proposition-oka} の直接の帰結である。第二について、有限生成でない
イデアル $I \subset R$ が存在すると仮定する。
有限生成でないイデアルの全順序鎖 $\left\{I_\alpha\right\}$ の和は
有限生成でない。実際、$I = \bigcup I_\alpha$ が
$a_1, \ldots, a_n$ で生成されるならば、すべての生成元はある
$I_\alpha $ に属し、したがってそれを生成することになる。
ツォルンの補題により、有限生成でないことに関して極大なイデアルが存在する。
第一の部分により、このイデアルは素イデアルである。
\end{proof}

\begin{lemma}
\label{algebra-lemma-primes-principal}
$R$ を環とする。
\begin{enumerate}
\item 単項でないことに関して極大なイデアル $I \subset R$ は素イデアルである。
\item $R$ のすべての素イデアルが単項ならば、$R$ のすべてのイデアルが単項である。
\end{enumerate}
\end{lemma}

\begin{proof}
第一の部分は例 \ref{algebra-example-oka-family-principal} と命題
\ref{algebra-proposition-oka} から従う。第二について、単項でないイデアル
$I \subset R$ が存在すると仮定する。
% P01-E1283: frozen English source missing-copula candidate; see ERRATA_CANDIDATES.jsonl.
単項でないイデアルの全順序鎖 $\left\{I_\alpha\right\}$ の和は単項でない。
実際、$I = \bigcup I_\alpha$ が $a$ で生成されるならば、
$a$ はある $I_\alpha $ に属し、$a$ がそれを生成することになる。
ツォルンの補題により、単項でないことに関して極大なイデアルが存在する。
第一の部分により、このイデアルは必然的に素イデアルである。
\end{proof}

\begin{lemma}
\label{algebra-lemma-characterize-domain}
$R$ を環とする。
\begin{enumerate}
\item 非零因子を含まないイデアルの中で極大なイデアルは素イデアルである。
\item $R$ が非零で、$R$ の任意の非零素イデアルが非零因子を含むならば、
$R$ は整域である。
\end{enumerate}
\end{lemma}

\begin{proof}
非零因子の集合 $S$ を考える。これは $R$ の乗法的部分集合である。
したがって $S$ と交わらないことに関して極大な任意のイデアルは素イデアルである。
補題 \ref{algebra-lemma-simple} を参照せよ。
ゆえに、任意の非零素イデアルが非零因子を含むならば、
$(0)$ は素イデアル、すなわち $R$ は整域である。
\end{proof}

\begin{remark}
\label{algebra-remark-cohen-bound-cardinality}
$R$ を環とし、$\kappa$ を無限基数とする。
例 \ref{algebra-example-oka-family-bound-cardinality} と命題
\ref{algebra-proposition-oka} を適用すると、
$\kappa$ 個の元で生成されないことに関して極大な任意のイデアルは
素イデアルであることが分かる。この結果はそれほど有用ではない。
$R$ の任意の素イデアルを $\aleph_0$ 個の元で生成できるが、
あるイデアルは生成できないような環が存在するからである。
実際、$k$ を体、$T$ を濃度が $\aleph_0$ より大きい集合とし、
$$
R = k[\{x_n\}_{n \geq 1}, \{z_{t, n}\}_{t \in T, n \geq 0}]/
(x_n^2, z_{t, n}^2, x_n z_{t, n} - z_{t, n - 1})
$$
とする。これは一意な素イデアル $\mathfrak m = (x_n)$ をもつ局所環である。
しかしイデアル $(z_{t, n})$ は可算個の元では生成できない。
\end{remark}

\begin{example}
\label{algebra-example-noetherian-topology-on-spec}
\begin{reference}
2020 年 11 月 12 日の Lukas Heger のコメント。
\end{reference}
$R$ を環、$X = \Spec(R)$ とする。
% P01-E1284: frozen English source “radical ideas” typo candidate; see ERRATA_CANDIDATES.jsonl.
$X$ の閉部分集合は $R$ の根基イデアルに対応するので
（補題 \ref{algebra-lemma-Zariski-topology}）、$X$ がネーター位相空間であることと、
根基イデアルについて上昇連鎖条件が成り立つことは同値である。
これは、任意の根基イデアルが有限生成イデアルの根基であることと同値である
（詳細は省略する）。次を置く。
$$
\mathcal{F} = \{I \subset R \mid
\sqrt{I} = \sqrt{(f_1, \ldots, f_n)}\text{ for some }n
\text{ and }f_1, \ldots, f_n \in R\}.
$$
任意のイデアル $I \subset R$ と任意の元 $a \in R$ に対して成り立つ等式
$$
\sqrt{I} = \sqrt{(I, a)(I : a)}
$$
を用いれば、$\mathcal{F}$ が岡族であることを読者は示せる。
一方、いずれも $\mathcal{F}$ に属さないイデアルの全順序鎖
$\{I_\alpha\}$ があれば、その和 $I = \bigcup I_\alpha$ も
$\mathcal{F}$ に属し得ない。そうでなければ、
$\sqrt{I} = \sqrt{(f_1, \ldots, f_n)}$ であり、
ある $e$ に対して $f_i^e \in I$、次に $i$ に依存しないある
$\alpha$ に対して $f_i^e \in I_\alpha$、さらに
$\sqrt{I_\alpha} = \sqrt{(f_1, \ldots, f_n)}$ となり、矛盾する。
したがって、$\mathcal{F}$ に属さないイデアルの集合が空でなければ、
それは極大元をもち、補題 \ref{algebra-lemma-cohen} とまったく同様に、
$X$ がネーター位相空間であることと、$R$ の任意の素イデアルが
ある $f_1, \ldots, f_n \in R$ に対して
$\sqrt{(f_1, \ldots, f_n)}$ に等しいことは同値であると結論できる。
この結果が必要になった場合には、ここで注意深く主張し証明することにする。
\end{example}








\section{有限表示環写像の像}
\label{algebra-section-images-finite-presentation}

\noindent
この節では、環写像 $R \to S$ が誘導する写像
$\Spec(S) \to \Spec(R)$ の位相、主としてシュヴァレーの定理について
いくつかの結果を証明する。そのために、Topology, Section
\ref{topology-section-constructible} で定義される構成可能集合、
準コンパクト集合、レトロコンパクト集合などの概念を用いる。

\begin{lemma}
\label{algebra-lemma-qc-open}
$U \subset \Spec(R)$ を開集合とする。次は同値である。
\begin{enumerate}
\item $U$ は $\Spec(R)$ においてレトロコンパクトである。
\item $U$ は準コンパクトである。
\item $U$ は標準開集合の有限和である。
\item ある有限生成イデアル $I \subset R$ が存在して
% P01-E0318: frozen English source undefined ambient symbol candidate; see ERRATA_CANDIDATES.jsonl.
$X \setminus V(I) = U$ となる。
\end{enumerate}
\end{lemma}

\begin{proof}
$\Spec(R)$ は準コンパクトなので (1) $\Rightarrow$ (2) である。
補題 \ref{algebra-lemma-quasi-compact} を参照せよ。
標準開集合は位相の基底をなすので (2) $\Rightarrow$ (3) である。
(3) $\Rightarrow$ (1) を証明する。
$U = \bigcup_{i = 1\ldots n} D(f_i)$ とする。
$U$ が $\Spec(R)$ においてレトロコンパクトであることを示すには、
$\Spec(R)$ の任意の準コンパクト開集合 $V$ に対して
$U \cap V$ が準コンパクトであることを示せば十分である。
(2) $\Rightarrow$ (3) により可能なので、
$V = \bigcup_{j = 1\ldots m} D(g_j)$ と書く。
各標準開集合はある環のスペクトルと同相であり、それゆえ準コンパクトである。
補題 \ref{algebra-lemma-standard-open} および
\ref{algebra-lemma-quasi-compact} を参照せよ。したがって
$U \cap V =
(\bigcup_{i = 1\ldots n} D(f_i)) \cap (\bigcup_{j = 1\ldots m} D(g_j))
= \bigcup_{i, j} D(f_i g_j)$ は準コンパクト開集合の有限和なので
準コンパクトである。証明を終えるには、補題
\ref{algebra-lemma-Zariski-topology} により (4) と (3) が同値であることに注意すればよい。
\end{proof}

\begin{lemma}
\label{algebra-lemma-affine-map-quasi-compact}
$\varphi : R \to S$ を環写像とする。誘導される連続写像
$f : \Spec(S) \to \Spec(R)$ は準コンパクトである。
任意の構成可能集合 $E \subset \Spec(R)$ に対して、逆像
$f^{-1}(E)$ は $\Spec(S)$ で構成可能である。
\end{lemma}

\begin{proof}
まず、任意の準コンパクト開集合 $U \subset \Spec(R)$ の逆像が
準コンパクトであることを示す。補題 \ref{algebra-lemma-qc-open} により、
% P01-E0319: frozen English source “finite open” noun-substitution candidate; see ERRATA_CANDIDATES.jsonl.
$U$ を標準開集合の有限和として書ける。したがって補題
\ref{algebra-lemma-spec-functorial} により、$f^{-1}(U)$ は標準開集合の有限和である。
ゆえに、再び補題 \ref{algebra-lemma-qc-open} によって $f^{-1}(U)$ は準コンパクトである。
第二の主張は Topology, Lemma
\ref{topology-lemma-inverse-images-constructibles} から従う。
\end{proof}

\begin{lemma}
\label{algebra-lemma-constructible}
$R$ を環とする。$\Spec(R)$ の部分集合が構成可能であるための必要十分条件は、
それが $f, g_1, \ldots, g_m \in R$ に対する
$D(f) \cap V(g_1, \ldots, g_m)$ の形の部分集合の有限和として書けることである。
\end{lemma}

\begin{proof}
補題 \ref{algebra-lemma-qc-open} により、部分集合 $D(f)$ と
$V(g_1, \ldots, g_m)$ の補集合はレトロコンパクト開集合である。
したがって $D(f) \cap V(g_1, \ldots, g_m)$ は構成可能な部分集合であり、
そのような集合の任意の有限和も構成可能である。
逆に、$T \subset \Spec(R)$ を構成可能とする。
Topology, Definition \ref{topology-definition-constructible} により、
$U, V \subset \Spec(R)$ をレトロコンパクト開集合として
$T = U \cap V^c$ と仮定してよい。補題 \ref{algebra-lemma-qc-open} により
$U = \bigcup_{i = 1, \ldots, n} D(f_i)$ および
$V = \bigcup_{j = 1, \ldots, m} D(g_j)$ と書ける。すると
$T = \bigcup_{i = 1, \ldots, n} \big(D(f_i) \cap V(g_1, \ldots, g_m)\big)$
である。
\end{proof}

\begin{lemma}
\label{algebra-lemma-constructible-is-image}
$R$ を環とし、$T \subset \Spec(R)$ を構成可能とする。
このとき、有限表示の環写像 $R \to S$ で、$\Spec(R)$ における
$\Spec(S)$ の像が $T$ となるものが存在する。
\end{lemma}

\begin{proof}
環の有限積のスペクトルはスペクトルの非交和である。補題
\ref{algebra-lemma-spec-product} を参照せよ。
したがって $T = T_1 \cup T_2$ であり、結果が $T_1$ と $T_2$
について成り立つならば、結果は $T$ についても成り立つ。
補題 \ref{algebra-lemma-constructible} により、
$T = D(f) \cap V(g_1, \ldots, g_m)$ と仮定してよい。
この場合、$T$ は写像
$\Spec((R/(g_1, \ldots, g_m))_f) \to \Spec(R)$ の像である。
補題 \ref{algebra-lemma-standard-open} および \ref{algebra-lemma-spec-closed} を参照せよ。
\end{proof}

\begin{lemma}
\label{algebra-lemma-open-fp}
$R$ を環とし、$f$ を $R$ の元とし、$S = R_f$ とする。
このとき $\Spec(S)$ の構成可能部分集合の像は $\Spec(R)$ で構成可能である。
\end{lemma}

\begin{proof}
以下、補題 \ref{algebra-lemma-qc-open} を断りなく繰り返し用いる。
$U, V$ を $\Spec(S)$ の準コンパクト開集合とする。
$U \cap V^c$ の像が構成可能であることを示す。
補題 \ref{algebra-lemma-standard-open} の同一視
$\Spec(S) = D(f)$ のもとで、集合 $U, V$ は
$\Spec(R)$ の準コンパクト開集合 $U', V'$ に対応する。
したがって、$\Spec(R)$ において $U' \cap (V')^c$ が構成可能であることを
示せば十分であり、これは明らかである。
\end{proof}

\begin{lemma}
\label{algebra-lemma-closed-fp}
$R$ を環とし、$I$ を $R$ の有限生成イデアルとし、$S = R/I$ とする。
このとき $\Spec(S)$ の構成可能部分集合の像は $\Spec(R)$ で構成可能である。
\end{lemma}

\begin{proof}
$I = (f_1, \ldots, f_m)$ ならば、$V(I)$ は
$\bigcup D(f_i)$ の補集合であることが分かる。補題
\ref{algebra-lemma-Zariski-topology} を参照せよ。
したがって補題 \ref{algebra-lemma-qc-open} により、それは構成可能である。
写像 $R \to S$ を $f \mapsto \overline{f}$ と書く。
$\overline{U}, \overline{V}$ が $\Spec(S)$ のレトロコンパクト開集合ならば、
$\Spec(R)$ における $\overline{U} \cap \overline{V}^c$ の像が
構成可能であることを示さなければならない。
補題 \ref{algebra-lemma-qc-open} により
$\overline{U} = \bigcup D(\overline{g_i})$ と書ける。
$U = \bigcup D({g_i})$ と置くと、$\overline{U}$ の像は
$U \cap V(I)$ であり、これは $\Spec(R)$ で構成可能である。
同様に $\overline{V}$ の像は、$\Spec(R)$ のあるレトロコンパクト開集合
$V$ に対する $V \cap V(I)$ に等しい。
したがって $\overline{U} \cap \overline{V}^c$ の像は
$U \cap V(I) \cap V^c$ に等しく、望む結果を得る。
\end{proof}

\begin{lemma}
\label{algebra-lemma-affineline-open}
$R$ を環とする。写像 $\Spec(R[x]) \to \Spec(R)$ は開写像であり、
任意の標準開集合の像は準コンパクト開集合である。
\end{lemma}

\begin{proof}
標準開集合 $D(f)$、$f\in R[x]$ の像が準コンパクト開集合であることを
示せば十分である。$D(f)$ の像は
$\Spec(R[x]_f) \to \Spec(R)$ の像である。
$\mathfrak p \subset R$ を素イデアルとし、$\overline{f}$ を
$\kappa(\mathfrak p)[x]$ における $f$ の像とする。
補題 \ref{algebra-lemma-in-image} を思い出すと、$\mathfrak p$ が像に属するための
必要十分条件は
$R[x]_f \otimes_R \kappa(\mathfrak p) =
\kappa(\mathfrak p)[x]_{\overline{f}}$ が零環でないことである。
これはちょうど $f$ が $\kappa(\mathfrak p)[x]$ で零に写らないという条件、
言い換えれば $f$ のある係数が $\mathfrak p$ に属さないという条件である。
したがって、$f = a_d x^d + \ldots + a_0$ ならば、
$D(f)$ の像は $D(a_d) \cup \ldots \cup D(a_0)$ である。
\end{proof}

\noindent
以下で用いる特性多項式の性質を証明する。

\begin{lemma}
\label{algebra-lemma-characteristic-polynomial-prime}
$R \to A$ を環準同型とする。$R$-加群として
$A \cong R^{\oplus n}$ であると仮定する。$f \in A$ とする。
乗法写像 $m_f: A \to A$ は $R$-線形なので、特性多項式
$P(T) = T^n + r_{n-1}T^{n-1} + \ldots + r_0 \in R[T]$ をもつ。
任意の素イデアル $\mathfrak{p} \in \Spec(R)$ に対し、
$f$ が $A \otimes_R \kappa(\mathfrak{p})$ 上冪零に作用するための必要十分条件は、
$\mathfrak p \in V(r_0, \ldots, r_{n-1})$ である。
\end{lemma}

\begin{proof}
乗法写像
$m_{\bar f}: A \otimes_R \kappa(\mathfrak p) \to
A \otimes_R \kappa(\mathfrak p)$ の特性多項式
$\bar P(T) \in \kappa(\mathfrak p)[T]$ が、単に $P(T)$ の
$\kappa(\mathfrak p)[T]$ における像であることを示せば、結果は容易に従う。
この写像は $A \otimes_R \kappa(\mathfrak p)$ の元に
% P01-E0320: frozen English source ambient-ring claim candidate; see ERRATA_CANDIDATES.jsonl.
$\bar f$、すなわち $\kappa(\mathfrak p)$ の元と見た $f$ の像を掛ける。
$R$-加群としての $A$ の基底 $e_1, \ldots, e_n$ を用い、
$R$ に成分をもつ写像 $m_f$ の行列を $(a_{ij})$ とする。
すると
$A \otimes_R \kappa(\mathfrak p) \cong (R \otimes_R
\kappa(\mathfrak p))^{\oplus n} = \kappa(\mathfrak p)^n$ であり、
これは基底 $e_1 \otimes 1, \ldots, e_n \otimes 1$ をもつ
$\kappa(\mathfrak p)$ 上の $n$ 次元ベクトル空間である。
像は $\bar f = f \otimes 1$ なので、乗法写像 $m_{\bar f}$ の行列は
$(a_{ij} \otimes 1)$ である。したがって、特性多項式はちょうど
$P(T)$ の像である。

\medskip\noindent
線形代数により、線形変換が $n$ 次元ベクトル空間上冪零に作用するための
必要十分条件は、特性多項式が $T^n$ であることである
（特性多項式は最小多項式のある冪を割るからである）。
したがって、$f$ が $A \otimes_R \kappa(\mathfrak p)$ 上
冪零に作用するための必要十分条件は $\bar P(T) = T^n$ である。
これは、すべての $0 \leq i \leq n - 1$ に対して
$r_i \in \mathfrak p$ となるとき、すなわち
$\mathfrak p \in V(r_0, \ldots, r_{n - 1}).$ となるとき、
かつそのときに限る。
\end{proof}

\begin{lemma}
\label{algebra-lemma-affineline-special}
$R$ を環とし、$f, g \in R[x]$ を多項式とする。
$g$ の最高次係数が $R$ の単元であると仮定する。
% P01-E0321: frozen English source subject-verb agreement candidate; see ERRATA_CANDIDATES.jsonl.
$i = 1\ldots, n$ に対する元 $r_i\in R$ で、
$\Spec(R)$ における $D(f) \cap V(g)$ の像が
$\bigcup_{i = 1, \ldots, n} D(r_i)$ となるものが存在する。
\end{lemma}

\begin{proof}
$g = ux^d + a_{d-1}x^{d-1} + \ldots + a_0$ と書く。
ここで $d$ は $g$ の次数であり、したがって $u \in R^*$ である。
環 $A = R[x]/(g)$ を考える。これは $R$-加群として、
$1, x, \ldots, x^{d-1}$ の像を基底とする有限自由加群である。
$A$ 上で $f$ の像による乗法を考える。これは $R$-加群写像である。
したがって、この写像の特性多項式を $P(T) \in R[T]$ と書ける。
$P(T) = T^d + r_{d-1} T^{d-1} + \ldots + r_0$ と書く。
$r_0, \ldots, r_{d-1}$ が求める性質をもつと主張する。
以下では、特性多項式の次の性質を用いる。
$$
\mathfrak p \in V(r_0, \ldots, r_{d-1})
\Leftrightarrow
\text{multiplication by }f\text{ is nilpotent on }
A \otimes_R \kappa(\mathfrak p).
$$
これは補題 \ref{algebra-lemma-characteristic-polynomial-prime} で証明した。

\medskip\noindent
$\mathfrak q\in D(f) \cap V(g)$ とし、
$\mathfrak p = \mathfrak q \cap R$ とする。このとき、$f$ による乗法と両立する
非零写像 $A \otimes_R \kappa(\mathfrak p) \to \kappa(\mathfrak q)$ がある。
そして $f$ は $\kappa(\mathfrak q)$ 上単元として作用する。
したがって $\mathfrak p \not \in  V(r_0, \ldots, r_{d-1})$ と結論する。

\medskip\noindent
逆に、$R$ のある素イデアル $\mathfrak p$ とある
$0 \leq i \leq d - 1$ に対して $r_i \not\in \mathfrak p$ と仮定する。
このとき $f$ による乗法は代数
$A \otimes_R \kappa(\mathfrak p)$ 上冪零でない。
したがって、$f$ の像を含まない素イデアル
$\overline{\mathfrak q} \subset A \otimes_R \kappa(\mathfrak p)$ が存在する。
$R[x]$ における $\overline{\mathfrak q}$ の逆像は、
$\mathfrak p$ に写る $D(f) \cap V(g)$ の元である。
\end{proof}

\begin{theorem}[シュヴァレーの定理]
\label{algebra-theorem-chevalley}
$R \to S$ が有限表示であると仮定する。
$\Spec(S)$ の構成可能部分集合の $\Spec(R)$ における像は構成可能である。
\end{theorem}

\begin{proof}
$S = R[x_1, \ldots, x_n]/(f_1, \ldots, f_m)$ と書く。
$R \to S$ を
$R \to R[x_1] \to R[x_1, x_2]
\to \ldots \to R[x_1, \ldots, x_{n-1}] \to S$ と分解できる。
したがって $S = R[x]/(f_1, \ldots, f_m)$ と仮定してよい。
この場合、写像を $R \to R[x] \to S$ と分解し、補題
\ref{algebra-lemma-closed-fp} により $S = R[x]$ の場合に帰着する。
% P01-E0322: frozen English source missing-subject candidate; see ERRATA_CANDIDATES.jsonl.
補題 \ref{algebra-lemma-qc-open} により、
$f_i , g_j \in R[x]$ に対して
$T = (\bigcup_{i = 1\ldots n} D(f_i)) \cap V(g_1, \ldots, g_m)$
であるとき、その $\Spec(R)$ における像が構成可能であることを
示せば十分である。構成可能集合の有限和は構成可能なので、
$n = 1$ の場合、すなわち
$T = D(f) \cap V(g_1, \ldots, g_m)$ の場合を扱えば十分である。

\medskip\noindent
$c \in R$ ならば、
$$
\Spec(R) =
V(c) \amalg D(c) =
\Spec(R/(c)) \amalg \Spec(R_c),
$$
であり、それに対応して
$\Spec(R[x]) =
V(c) \amalg D(c) = \Spec(R/(c)[x]) \amalg
\Spec(R_c[x])$ であることに注意する。
$T = D(f) \cap V(g_1, \ldots, g_m)$ と各部分との交わりも、
$f$、$g_i$ をそれぞれ $R/(c)[x]$、$R_c[x]$ における像で置き換えた
同じ形をしている。$\Spec(R)$ における $T$ の像は、
$T \cap V(c)$ の像と $T \cap D(c)$ の像の和であることに注意する。
補題 \ref{algebra-lemma-open-fp} と \ref{algebra-lemma-closed-fp} を用いると、
両方の部分の像がそれぞれ $\Spec(R/(c))$ と $\Spec(R_c)$ で
構成可能であることを証明すれば十分である。

\medskip\noindent
上のように $T = D(f) \cap V(g_1, \ldots, g_m)$ であり、
$\deg(g_1) \leq \deg(g_2) \leq \ldots \leq \deg(g_m)$ と仮定する。
$m$ と $g_i$ の次数に関する帰納法を用いる。
$d_1 = \deg(g_1)$、すなわち
$g_1 = c x^{d_1} + l.o.t$ とし、$c \in R$ は零でないとする。
$R$ を $R/(c)$ と $R_c$ の部分に分けると、$g_1$ の次数が下がる
（これは帰納法で扱われる）か、$c$ が可逆な場合に帰着する。
$c$ が可逆で $m > 1$ ならば、
$g_2 = c' x^{d_2} + l.o.t$ と書く。この場合
$g_2' = g_2 - (c'/c) x^{d_2 - d_1} g_1$ を考える。
イデアル $(g_1, g_2, \ldots, g_m)$ と
$(g_1, g_2', g_3, \ldots, g_m)$ は等しいので、
% P01-E0323: frozen English source missing list separator candidate; see ERRATA_CANDIDATES.jsonl.
$T = D(f) \cap V(g_1, g_2', g_3\ldots, g_m)$ であることが分かる。
しかしここで $g_2'$ の次数は $g_2$ の次数より真に小さいので、
この場合は帰納法で扱われる。

\medskip\noindent
% P01-E0324: frozen English source “bases case” number/form candidate; see ERRATA_CANDIDATES.jsonl.
上の帰納法の基底の場合は、(a) $g$ の最高次係数が可逆である
$T = D(f) \cap V(g)$ の場合と、(b) $T = D(f)$ の場合である。
これら二つの場合は補題 \ref{algebra-lemma-affineline-special} および
\ref{algebra-lemma-affineline-open} で扱われる。
\end{proof}
\section{像についてのさらなる結果}
\label{algebra-section-more-images}

\noindent
この節では、環写像に対する $\Spec$ 上の像に関する補題をいくつか集める。
例えば Section \ref{algebra-section-going-up} も参照せよ。

\begin{lemma}
\label{algebra-lemma-generic-finite-presentation}
$R \subset S$ を整域の包含とする。$R \to S$ が有限型であると仮定する。
このとき、非零元 $f \in R$ と非零元 $g \in S$ で、
$R_f \to S_{fg}$ が有限表示となるものが存在する。
\end{lemma}

\begin{proof}
$R$ 上の $S$ の生成元の個数に関する帰納法による。
証明中、ある非零元 $f \in R$ に対して、
$R$ を $R_f$ で、$S$ を $S_f$ で置き換えてよい。

\medskip\noindent
$S$ が $R$ 上一つの元で生成されると仮定する。
このとき、ある素イデアル $\mathfrak q \subset R[x]$ に対して
$S = R[x]/\mathfrak q$ である。$\mathfrak q = (0)$ ならば証明することはない。
$\mathfrak q \not = (0)$ ならば、$x$ に関する次数が最小の非零元
$h \in \mathfrak q$ を取る。$a_i \in R$ および $f \not = 0$ として
$h = f x^d + a_{d - 1} x^{d - 1} + \ldots + a_0$ と書く。
$R$ と $S$ で $f$ を逆にした後、$h$ がモニックであると仮定してよい。
全射な $R$-代数写像 $R[x]/(h) \to S$ を得る。
$R$-加群として
$R[x]/(h) = R \oplus Rx \oplus \ldots \oplus Rx^{d - 1}$ であり、
$d$ の最小性により $R[x]/(h)$ は $S$ に単射的に写ることが分かる。
したがって $R[x]/(h) \cong S$ は $R$ 上有限表示である。

\medskip\noindent
$S$ が $R$ 上 $n > 1$ 個の元で生成されると仮定する。
$x_1, \ldots, x_n \in S$ が $S$ を生成するとする。
$x_1, \ldots, x_{n-1}$ で生成される $R$-部分代数を $S' \subset S$ と書く。
% P01-E0325: frozen English source missing-article candidate; see ERRATA_CANDIDATES.jsonl.
帰納法の仮定により、非零元 $f\in R$ と $g \in S'$ で、
$R_f \to S'_{fg}$ が有限表示となるものが存在する。
次に帰納法の仮定を $S'_{fg} \to S_{fg}$ に適用すると、
$f' \in S'_{fg}$ と $g' \in S_{fg}$ で、
$S'_{fgf'} \to S_{fgf'g'}$ が有限表示となるものが存在することが分かる。
結論を導くことは読者に任せる。
\end{proof}

\begin{lemma}
\label{algebra-lemma-characterize-image-finite-type}
$R \to S$ を有限型環写像とする。
$X = \Spec(R)$ および $Y = \Spec(S)$ と書く。
% P01-E0326: frozen English source malformed map-naming construction candidate; see ERRATA_CANDIDATES.jsonl.
誘導されるスペクトルの写像を $f : Y \to X$ と書く。
$E \subset Y = \Spec(S)$ を構成可能集合とする。
点 $\xi \in X$ が $f(E)$ に属するならば、
$\overline{\{\xi\}} \cap f(E)$ は
$\overline{\{\xi\}}$ の開稠密部分集合を含む。
\end{lemma}

\begin{proof}
$\xi \in X$ を $f(E)$ の点とする。$\xi$ に写る点 $\eta \in E$ を選ぶ。
$\mathfrak p \subset R$ を $\xi$ に対応する素イデアルとし、
$\mathfrak q \subset S$ を $\eta$ に対応する素イデアルとする。
図式
$$
\xymatrix{
\eta \ar[r] \ar@{|->}[d] & E \cap Y' \ar[r] \ar[d] &
Y' = \Spec(S/\mathfrak q) \ar[r] \ar[d] &
Y \ar[d] \\
\xi \ar[r] & f(E) \cap X' \ar[r] &
X' = \Spec(R/\mathfrak p) \ar[r] &
X
}
$$
を考える。補題 \ref{algebra-lemma-affine-map-quasi-compact} により、
集合 $E \cap Y'$ は $Y'$ で構成可能である。
したがって $X$ を $X'$ で、$Y$ を $Y'$ で置き換えてよい。
ゆえに、$R \subset S$ は整域の包含、$\xi$ は $X$ の生成点、
$\eta$ は $Y$ の生成点であると仮定してよい。
補題 \ref{algebra-lemma-generic-finite-presentation} とシュヴァレーの定理
（定理 \ref{algebra-theorem-chevalley}）を組み合わせると、稠密開集合
$U \subset X$、$V \subset Y$ で、$f(V) \subset U$ かつ
$f : V \to U$ が構成可能集合を構成可能集合に写すものが存在することが分かる。
$E \cap V$ は $V$ で構成可能であることに注意する。Topology, Lemma
\ref{topology-lemma-open-immersion-constructible-inverse-image} を参照せよ。
したがって $f(E \cap V)$ は $U$ で構成可能であり、$\xi$ を含む。
Topology, Lemma \ref{topology-lemma-generic-point-in-constructible} により、
$f(E \cap V)$ は稠密開集合 $U' \subset U$ を含む。
\end{proof}

\noindent
この節の末尾では、構成可能集合とは関係のない、スペクトル上の写像の像に関する
結果をさらにいくつか示す。

\begin{lemma}
\label{algebra-lemma-surjective-spec-radical-ideal}
$\varphi : R \to S$ を環写像とする。次は同値である。
\begin{enumerate}
\item 写像 $\Spec(S) \to \Spec(R)$ は全射である。
\item 任意のイデアル $I \subset R$ に対して、
$R$ における $\sqrt{IS}$ の逆像は $\sqrt{I}$ に等しい。
\item 任意の根基イデアル $I \subset R$ に対して、
$R$ における $IS$ の逆像は $I$ に等しい。
\item $R$ の任意の素イデアル $\mathfrak p$ に対して、
$R$ における $\mathfrak p S$ の逆像は $\mathfrak p$ である。
\end{enumerate}
この場合、任意の基底変換の後でも同じことが成り立つ。
すなわち、環写像 $R \to R'$ が与えられると、環写像
$R' \to R' \otimes_R S$ も同値な性質 (1)、(2)、(3) をもつ。
\end{lemma}

\begin{proof}
$J \subset S$ がイデアルならば、
$\sqrt{\varphi^{-1}(J)} = \varphi^{-1}(\sqrt{J})$ である。
これは (2) と (3) が同値であることを示す。
(3) $\Rightarrow$ (4) は直ちに従う。
$I \subset R$ が根基イデアルならば、補題
\ref{algebra-lemma-Zariski-topology} により
$I = \bigcap_{I \subset \mathfrak p} \mathfrak p$ である。
したがって (4) $\Rightarrow$ (2) である。
補題 \ref{algebra-lemma-in-image} により、
$\mathfrak p = \varphi^{-1}(\mathfrak p S)$ であるための必要十分条件は、
$\mathfrak p$ が像に属することである。したがって
(1) $\Leftrightarrow$ (4) である。ゆえに (1)、(2)、(3)、(4) は同値である。

\medskip\noindent
(1) を仮定する。$R \to R'$ を環写像とする。
$\mathfrak p' \subset R'$ を $R$ の素イデアル $\mathfrak p$ の上にある
素イデアルとする。$\mathfrak p'$ が写像
$\Spec(R' \otimes_R S) \to \Spec(R')$ の像に属することを見るには、
補題 \ref{algebra-lemma-in-image} により
$(R' \otimes_R S) \otimes_{R'} \kappa(\mathfrak p')$ が零でないことを
示さなければならない。しかし
$$
(R' \otimes_R S) \otimes_{R'} \kappa(\mathfrak p') =
S \otimes_R \kappa(\mathfrak p)
\otimes_{\kappa(\mathfrak p)} \kappa(\mathfrak p')
$$
である。仮定により $S \otimes_R \kappa(\mathfrak p)$ は零でなく、
$\kappa(\mathfrak p) \to \kappa(\mathfrak p')$ は体の拡大なので、
これは零でない。
\end{proof}

\begin{lemma}
\label{algebra-lemma-domain-image-dense-set-points-generic-point}
$R$ を整域とし、$\varphi : R \to S$ を環写像とする。次は同値である。
\begin{enumerate}
\item 環写像 $R \to S$ は単射である。
\item 像 $\Spec(S) \to \Spec(R)$ は点の稠密集合を含む。
\item $R$ における逆像が $(0)$ となる素イデアル
$\mathfrak q \subset S$ が存在する。
\end{enumerate}
\end{lemma}

\begin{proof}
$K$ を整域 $R$ の分数体とする。$R \to S$ が単射であると仮定する。
局所化は完全なので、$K \to S \otimes_R K$ は単射であることが分かる。
したがって補題 \ref{algebra-lemma-in-image} により、$(0)$ に写る素イデアルが存在する。

\medskip\noindent
$(0)$ は $\Spec(R)$ で稠密であることに注意する。
したがって、最後の条件から第二の条件が従う。

\medskip\noindent
第二の条件が成り立つと仮定する。$f \in R$、$f \not = 0$ とする。
$R$ は整域なので、$V(f)$ は
% P01-E0327: frozen English source wrong ambient-space candidate; see ERRATA_CANDIDATES.jsonl.
$R$ の真の閉部分集合である。仮定により、$\varphi(f) \not \in \mathfrak q$
を満たす $S$ の素イデアル $\mathfrak q$ が存在する。
したがって $\varphi(f) \not = 0$ である。
ゆえに $R \to S$ は単射である。
\end{proof}

\begin{lemma}
\label{algebra-lemma-injective-minimal-primes-in-image}
$R \subset S$ を単射な環写像とする。このとき
$\Spec(S) \to \Spec(R)$ はすべての極小素イデアルに達する。
\end{lemma}

\begin{proof}
$\mathfrak p \subset R$ を極小素イデアルとする。
この場合 $R_{\mathfrak p}$ は一意な素イデアルをもつ。
したがって $S_{\mathfrak p}$ が零でないことを示せば十分である。
これは局所化が完全であることから従う。命題
\ref{algebra-proposition-localization-exact} を参照せよ。
\end{proof}

\begin{lemma}
\label{algebra-lemma-image-dense-generic-points}
$R \to S$ を環写像とする。次は同値である。
\begin{enumerate}
\item $R \to S$ の核は冪零元からなる。
\item $R$ の極小素イデアルは $\Spec(S) \to \Spec(R)$ の像に属する。
\item $\Spec(S) \to \Spec(R)$ の像は $\Spec(R)$ で稠密である。
\end{enumerate}
\end{lemma}

\begin{proof}
$I = \Ker(R \to S)$ とする。補題 \ref{algebra-lemma-Zariski-topology} により
$\sqrt{(0)} = \bigcap_{\mathfrak q \subset S} \mathfrak q$ であることに注意する。
したがって
$\sqrt{I} = \bigcap_{\mathfrak q \subset S} R \cap \mathfrak q$ である。
ゆえに $V(I) = V(\sqrt{I})$ は $\Spec(S) \to \Spec(R)$ の像の閉包である。
これは (1) と (3) が同値であることを示す。(2) から (3) が従うことは明らかである。
最後に (1) を仮定する。スペクトルの位相と写像に影響を与えずに、
$R$ を $R/I$ で、$S$ を $S/IS$ で置き換えてよい。
したがって (1) $\Rightarrow$ (2) は補題
\ref{algebra-lemma-injective-minimal-primes-in-image} から従う。
\end{proof}

\begin{lemma}
\label{algebra-lemma-minimal-prime-image-minimal-prime}
$R \to S$ を環写像とする。極小素イデアル $\mathfrak p \subset R$ が
$\Spec(S) \to \Spec(R)$ の像に属するならば、それは極小素イデアルの像である。
\end{lemma}

\begin{proof}
$\mathfrak p = \mathfrak q \cap R$ とする。
補題 \ref{algebra-lemma-Zariski-topology} により、$\mathfrak r \subset \mathfrak q$
を満たす極小素イデアル $\mathfrak r \subset S$ を選ぶ。
$\mathfrak p$ の極小性により
$\mathfrak p = \mathfrak r \cap R$ であることが分かる。
\end{proof}

\begin{lemma}
\label{algebra-lemma-intersection-not-zero}
$A \subset B$ を、分数体の代数拡大を誘導する整域の包含とする。
$J \subset B$ が非零イデアルならば、$A \cap J$ も非零である。
したがって $\Spec(B)$ の真の閉部分集合の像は $\Spec(A)$ で稠密でない。
\end{lemma}

\begin{proof}
$x \in J$ を非零元とする。$x$ は $A$ の分数体上代数的なので、
$a_0, a_d \not = 0$ を満たす $d \geq 1$ と
$a_0, \ldots, a_d \in A$ が存在して、
$B$ において
$a_d x^d + a_{d - 1} x^{d - 1} + \ldots + a_0 = 0$ となる。
すると $a_0 \in A \cap J$ である。
\end{proof}

\section{ネーター環}
\label{algebra-section-Noetherian}

\noindent
環 $R$ の任意のイデアルが有限生成であるとき、$R$ を {\it ネーター的} という。
これは明らかに $R$ のイデアルに対する昇鎖条件と同値である。
補題 \ref{algebra-lemma-cohen} により、$R$ の任意の素イデアルが有限生成であることを
確かめれば十分である。

\begin{lemma}
\label{algebra-lemma-Noetherian-permanence}
\begin{slogan}
ネーター性は有限型拡大および局所化への移行で保たれる。
\end{slogan}
ネーター環上の任意の有限生成環はネーター的である。
ネーター環の任意の局所化はネーター的である。
\end{lemma}

\begin{proof}
局所化に関する主張は、任意のイデアル $J \subset S^{-1}R$ が
$I \cdot S^{-1}R$ の形であることから従う。
ネーター環 $R$ の任意の商 $R/I$ はネーター的である。
実際、任意のイデアル $\overline{J} \subset R/I$ は、
あるイデアル $I \subset J \subset R$ に対する $J/I$ の形である。
したがって、$R$ がネーター的ならば $R[X]$ もそうであることを
示せば十分である。$J_1 \subset J_2 \subset \ldots$ を
$R[X]$ のイデアルの昇鎖とする。
$J_i$ の次数 $d$ の多項式の最高次係数として現れる $R$ の元のイデアルを
$I_{i, d}$ と定める。$i \leq i'$ かつ $d \leq d'$ ならば、明らかに
$I_{i, d} \subset I_{i', d'}$ である。
$R$ の昇鎖条件により、すべての $I_{i, d}$ の中には相異なるイデアルが
高々有限個しかない。
（ヒント：$\mathbf{N} \times \mathbf{N}$ の任意の無限部分集合は、
増加する無限列を含む。）
すべての $i \geq i_0$ とすべての $d$ に対して
$I_{i, d} = I_{i_0, d}$ となるほど十分大きい $i_0$ を取る。
ある $i \geq i_0$ に対して $f \in J_i$ と仮定する。
次数 $d = \deg(f)$ に関する帰納法により $f \in J_{i_0}$ を示す。
実際、次数が $d$ で $f$ と同じ最高次係数をもつ
$g\in J_{i_0}$ が存在する。帰納法により $f - g \in J_{i_0}$ となり、
証明が終わる。
\end{proof}

\begin{lemma}
\label{algebra-lemma-Noetherian-power-series}
$R$ がネーター環ならば、形式的冪級数環
$R[[x_1, \ldots, x_n]]$ もネーター的である。
\end{lemma}

\begin{proof}
$R[[x_1, \ldots, x_{n + 1}]] \cong R[[x_1, \ldots, x_n]][[x_{n + 1}]]$
なので、$R$ がネーター的ならば $R[[x]]$ がネーター的であるという主張を
証明すれば十分である。$I \subset R[[x]]$ をイデアルとする。
$I$ が有限生成イデアルであることを示さなければならない。
各整数 $d$ に対して
$I_d = \{a \in R \mid ax^d + \text{h.o.t.} \in I\}$ と書く。
$R$ はネーター的なので、$I_0 \subset I_1 \subset \ldots$ は安定する。
$I_{d_0} = I_{d_0 + 1} = \ldots$ となる $d_0$ を選ぶ。
各 $d \leq d_0$ に対して、$f_{d, j} \in I \cap (x^d)$、
$j = 1, \ldots, n_d$ なる元を選び、
$f_{d, j} = a_{d, j}x^d + \text{h.o.t}$ と書くと
$I_d = (a_{d, j})$ となるようにする。
$I' = (\{f_{d, j}\}_{d = 0, \ldots, d_0, j = 1, \ldots, n_d})$ と書く。
このとき $I' \subset I$ であることは明らかである。
$f \in I$ を取る。まず、次を満たす $c_{d, i} \in R$ を選べる。
$$
f - \sum c_{d, i} f_{d, i} \in (x^{d_0 + 1}) \cap I.
$$
次に、$i = 1, \ldots, n_{d_0}$ に対して
$c_{i, 1} \in R$ を選び、次が成り立つようにできる。
$$
f - \sum c_{d, i} f_{d, i} - \sum c_{i, 1}xf_{d_0, i} \in (x^{d_0 + 2}) \cap I.
$$
次に、$i = 1, \ldots, n_{d_0}$ に対して
$c_{i, 2} \in R$ を選び、次が成り立つようにできる。
$$
f - \sum c_{d, i} f_{d, i} - \sum c_{i, 1}xf_{d_0, i}
- \sum c_{i, 2}x^2f_{d_0, i}
\in (x^{d_0 + 3}) \cap I.
$$
以下同様である。最終的に
$$
f = \sum c_{d, i} f_{d, i} +
\sum\nolimits_i (\sum\nolimits_e c_{i, e} x^e)f_{d_0, i}
$$
% P01-E0328: frozen English source element/ideal containment wording candidate; see ERRATA_CANDIDATES.jsonl.
は望みどおり $I'$ に属することが分かる。
\end{proof}

\noindent
次の補題は容易だが、「絶対ネーター還元」と呼ばれる手法で
有限型 $\mathbf{Z}$-代数が頻繁に現れるため有用である。

\begin{lemma}
\label{algebra-lemma-obvious-Noetherian}
体上の任意の有限型代数はネーター的である。
$\mathbf{Z}$ 上の任意の有限型代数はネーター的である。
\end{lemma}

\begin{proof}
これは補題 \ref{algebra-lemma-Noetherian-permanence} と、
体がネーター環であること、および
% P01-E0329: frozen English source missing-article candidate; see ERRATA_CANDIDATES.jsonl.
$\mathbf{Z}$ がネーター環であること（主イデアル整域だからである）から直ちに従う。
\end{proof}

\begin{lemma}
\label{algebra-lemma-Noetherian-finite-type-is-finite-presentation}
$R$ をネーター環とする。
\begin{enumerate}
\item 任意の有限 $R$-加群は有限表示である。
\item 有限 $R$-加群の任意の部分加群は有限である。
\item 任意の有限型 $R$-代数は $R$ 上有限表示である。
\end{enumerate}
\end{lemma}

\begin{proof}
$M$ を有限 $R$-加群とする。補題
\ref{algebra-lemma-trivial-filter-finite-module} により、$M$ の有限フィルトレーションで、
その逐次商が $R/I$ の形となるものを取れる。
任意のイデアルは有限生成なので、各商 $R/I$ は有限表示である。
したがって補題 \ref{algebra-lemma-extension} により $M$ は有限表示である。
これで (1) が証明された。

\medskip\noindent
$N \subset M$ を部分加群とする。$M$ は有限なので、商 $M/N$ は有限である。
したがって (1) により $M/N$ は有限表示である。
ゆえに補題 \ref{algebra-lemma-extension} の (5) により $N$ は有限である。
これで (2) が証明された。

\medskip\noindent
(3) を見るには、補題 \ref{algebra-lemma-Noetherian-permanence} により
$R[x_1, \ldots, x_n]$ の任意のイデアルが有限生成であることに注意すればよい。
\end{proof}

\begin{lemma}
\label{algebra-lemma-Noetherian-topology}
$R$ がネーター環ならば、$\Spec(R)$ はネーター位相空間である。
Topology, Definition \ref{topology-definition-noetherian} を参照せよ。
\end{lemma}

\begin{proof}
これは、$\Spec(R)$ の任意の閉部分集合が、根基イデアル $I$ に対する
$V(I)$ の形に一意に書けるためである。補題
\ref{algebra-lemma-Zariski-topology} を参照せよ。
また、この対応は包含関係を逆転する。したがって結果は定義から従う。
\end{proof}

\begin{lemma}
\label{algebra-lemma-Noetherian-irreducible-components}
\begin{slogan}
ネーター・アファインスキームは有限個の生成点をもつ。
\end{slogan}
$R$ がネーター環ならば、$\Spec(R)$ は有限個の既約成分をもつ。
言い換えれば、$R$ は有限個の極小素イデアルをもつ。
\end{lemma}

\begin{proof}
補題 \ref{algebra-lemma-Noetherian-topology} と Topology, Lemma
\ref{topology-lemma-Noetherian} により、既約成分が有限個であることが分かる。
補題 \ref{algebra-lemma-irreducible} により、これらは $R$ の極小素イデアルに対応する。
\end{proof}

\begin{lemma}
\label{algebra-lemma-Noetherian-base-change-finite-type}
$R \to S$ を環写像とし、$R \to R'$ が有限型であるとする。
$S$ がネーター的ならば、基底変換
$S' = R' \otimes_R S$ はネーター的である。
\end{lemma}

\begin{proof}
補題 \ref{algebra-lemma-base-change-finiteness} により、有限型という性質は
基底変換で保たれる。したがって $S \to S'$ は有限型である。
$S$ はネーター的なので、補題 \ref{algebra-lemma-Noetherian-permanence} を適用できる。
\end{proof}

\begin{lemma}
\label{algebra-lemma-Noetherian-field-extension}
$k$ を体、$R$ をネーター的な $k$-代数とする。
$K/k$ が有限生成体拡大ならば、$K \otimes_k R$ はネーター的である。
\end{lemma}

\begin{proof}
$K/k$ は有限生成体拡大なので、有限生成 $k$-代数
$B \subset K$ で、$K$ が $B$ の分数体となるものが存在する。
言い換えれば、$S = B \setminus \{0\}$ として $K = S^{-1}B$ である。
すると $K \otimes_k R = S^{-1}(B \otimes_k R)$ である。
補題 \ref{algebra-lemma-Noetherian-base-change-finite-type} により
$B \otimes_k R$ はネーター的である。
最後に、補題 \ref{algebra-lemma-Noetherian-permanence} により
$K \otimes_k R = S^{-1}(B \otimes_k R)$ はネーター的である。
\end{proof}

\noindent
ときどき有用な、興味深い補題をいくつか挙げる。

\begin{lemma}
\label{algebra-lemma-subring-of-local-ring}
$R$ を環、$\mathfrak p \subset R$ を素イデアルとする。
次の各場合に、$R_f \to R_\mathfrak p$ が単射となる
$f \in R$、$f \not \in \mathfrak p$ が存在する。
\begin{enumerate}
\item $R$ は整域である。
\item $R$ はネーター的である。
\item $R$ は被約で、有限個の極小素イデアルをもつ。
\end{enumerate}
\end{lemma}

\begin{proof}
$R$ が整域ならば $R \subset R_\mathfrak p$ なので、$f = 1$ でよい。
$R$ がネーター的ならば、$R \to R_\mathfrak p$ の核 $I$ は有限生成イデアルであり、
$IR_f = 0$ を満たす $f \in R$、$f \not \in \mathfrak p$ を取れる。
この $f$ に対して写像 $R_f \to R_\mathfrak p$ は単射であり、$f$ は求めるものである。
$R$ が被約で、有限個の極小素イデアル
$\mathfrak p_1, \ldots, \mathfrak p_n$ をもつとする。
このとき
$f \in \bigcap_{\mathfrak p_i \not \subset \mathfrak p} \mathfrak p_i$、
$f \not \in \mathfrak p$ となる元を選べる。
実際、$\mathfrak{p}_i\not\subset \mathfrak{p}$ ならば
$f_i \in \mathfrak{p}_i$、$f_i \not\in \mathfrak{p}$ となる元が存在し、
$f = \prod f_i$ でよい。この $f$ に対して、
$R_f$ の極小素イデアルは $R_\mathfrak p$ の極小素イデアルに対応するので
$R_f \subset R_\mathfrak p$ であり、補題
\ref{algebra-lemma-reduced-ring-sub-product-fields} を適用できる
（いくつかの詳細は省略する）。
\end{proof}

\begin{lemma}
\label{algebra-lemma-surjective-endo-noetherian-ring-is-iso}
ネーター環の任意の全射自己準同型は同型である。
\end{lemma}

\begin{proof}
$f : R \to R$ がそのような自己準同型だが単射でないならば、
$$
\Ker(f) \subset \Ker(f \circ f) \subset
\Ker(f \circ f \circ f) \subset \ldots
$$
はイデアルの真の増加列となる。
\end{proof}

\section{局所冪零イデアル}
\label{algebra-section-locally-nilpotent}

\noindent
定義は次のとおりである。

\begin{definition}
\label{algebra-definition-locally-nilpotent-ideal}
$R$ を環、$I \subset R$ をイデアルとする。
すべての $x \in I$ に対して、$x^n = 0$ となる
$n \in \mathbf{N}$ が存在するとき、$I$ を {\it 局所冪零} という。
$I^n = 0$ となる $n \in \mathbf{N}$ が存在するとき、
$I$ を {\it 冪零} という。
\end{definition}

\begin{example}
\label{algebra-example-locally-nilpotent-not-nilpotent}
$R = k[x_n | n \in \mathbf{N}]$ を、体 $k$ 上の無限個の変数をもつ
多項式環とする。$I$ を $n \in \mathbf{N}$ に対する元 $x_n^n$ で
生成されるイデアルとし、$S = R/I$ とする。
$J \subset S$ を、$n \in \mathbf{N}$ に対する $x_n$ の像で
生成されるイデアルとする。このイデアルは局所冪零だが冪零でない。
実際、冪零元の $S$-線形結合は冪零なので、$J$ が局所冪零であることを
証明するには、すべての生成元が冪零であることを確認すれば十分である
（それらが冪零であることは明らかである）。
一方、各 $n \in \mathbf{N}$ に対して $x_{n + 1}^n \not \in I$ なので、
$J^n \not = 0$ である。したがって $J$ は冪零でない。
\end{example}

\begin{lemma}
\label{algebra-lemma-locally-nilpotent}
$R \to R'$ を環写像、$I \subset R$ を局所冪零イデアルとする。
このとき $IR'$ は $R'$ の局所冪零イデアルである。
\end{lemma}

\begin{proof}
これは、$x, y \in R'$ が冪零ならば $x + y$ も冪零であることから従う。
実際、$x^n = 0$ かつ $y^m = 0$ ならば
$(x + y)^{n + m - 1} = 0$ である。
\end{proof}

\begin{lemma}
\label{algebra-lemma-locally-nilpotent-unit}
$R$ を環、$I \subset R$ を局所冪零イデアルとする。
$R$ の元 $x$ が単元であるための必要十分条件は、
$R/I$ における $x$ の像が単元であることである。
\end{lemma}

\begin{proof}
$x$ が $R$ の単元ならば、その $R/I$ における像が単元であることは明らかである。
逆を証明すればよい。$R/I$ における $y \in R$ の像が
$x$ の像の逆元であると仮定する。このとき、ある $z \in I$ に対して
$xy = 1 - z$ である。これは $xR$ を法として $1\equiv z$ であることを意味する。
$z$ は局所冪零イデアル $I$ に属するので、十分大きいある $N$ に対して
$z^N = 0$ である。したがって、$xR$ を法として
$1 = 1^N \equiv z^N = 0$ である。
言い換えれば、$x$ は $1$ を割り、それゆえ単元である。
\end{proof}

\begin{lemma}
\label{algebra-lemma-Noetherian-power}
\begin{slogan}
ネーター環のイデアルは、その各元が冪零ならば冪零である。
\end{slogan}
$R$ をネーター環とし、$I, J$ を $R$ のイデアルとする。
$J \subset \sqrt{I}$ と仮定する。このとき、ある $n$ に対して
$J^n \subset I$ である。特に、ネーター環では
「局所冪零イデアル」と「冪零イデアル」の概念は一致する。
\end{lemma}

\begin{proof}
$J = (f_1, \ldots, f_s)$ とする。仮定により $f_i^{d_i} \in I$ である。
$n = d_1 + d_2 + \ldots + d_s + 1$ と取る。
\end{proof}

\begin{lemma}
\label{algebra-lemma-lift-idempotents}
$R$ を環、$I \subset R$ を局所冪零イデアルとする。
このとき $R \to R/I$ は冪等元の間の全単射を誘導する。
\end{lemma}

\begin{proof}[補題 \ref{algebra-lemma-lift-idempotents} の第一の証明]
$I$ は局所冪零なので、すべての素イデアルに含まれる。
したがって $\Spec(R/I) = V(I) = \Spec(R)$ である。
ゆえに補題 \ref{algebra-lemma-disjoint-decomposition} から結果が従う。
\end{proof}

\begin{proof}[補題 \ref{algebra-lemma-lift-idempotents} の第二の証明]
$\overline{e} \in R/I$ を冪等元とする。
$\overline{e}$ を $R$ の冪等元に持ち上げなければならない。

\medskip\noindent
まず、$\overline{e}$ の任意の持ち上げ $f \in R$ を選び、
$x = f^2 - f$ と置く。すると $x \in I$ なので、$x$ は冪零である
（$I$ が局所冪零だからである）。いま、$x$ で生成される
$R$ のイデアルを $J$ とする。このとき $J$ は、単に局所冪零なのではなく
冪零である。冪零元 $x$ で生成されるからである。

\medskip\noindent
ある $k \geq 1$ に対して $e^2 - e \in J^k$ となる
$\overline{e}$ の持ち上げ $e \in R$ が見つかったと仮定する。
$e' = e - (2e - 1)(e^2 - e) = 3e^2 - 2e^3$ と置く。
これは $\overline{e}$ の別の持ち上げである
（$\overline{e}$ の冪等性から $e^2 - e \in I$ となるためである）。
簡単な計算により
$$
(e')^2 - e' = (4e^2 - 4e - 3)(e^2 - e)^2 \in J^{2k}
$$
である。

\medskip\noindent
こうして、$e^2 - e \in J^k$ を満たす $\overline{e}$ の持ち上げ $e$ から出発し、
$(e')^2 - e' \in J^{2k}$ を満たす $\overline{e}$ の持ち上げ $e'$ を得た。
この方法で近似を逐次改善できる
（$k = 1$ の条件を満たす $e = f$ から始める）。
最終的に $J^k = 0$ となる段階に達し、その段階では
$e^2 - e \in J^k = 0$ を満たす $\overline{e}$ の持ち上げ $e$ が得られる。
すなわち、この $e$ は冪等元である。

\medskip\noindent
以上により、任意の冪等元 $\overline{e} \in R/I$ に対して、
$R$ の冪等元である $\overline{e}$ の持ち上げが存在することが分かった。
この持ち上げの一意性を証明すればよい。
実際、$e_1$ と $e_2$ をそのような二つの持ち上げとする。
$e_1 = e_2$ であることを示さなければならない。

\medskip\noindent
$e_1$ と $e_2$ の定義により、$e_1 \equiv e_2 \mod I$ であり、
$e_1$ と $e_2$ はともに冪等元である。
$e_1 \equiv e_2 \mod I$ から $e_1 - e_2 \in I$ が分かるので、
$e_1 - e_2$ は冪零である（$I$ が局所冪零だからである）。
$e_1$ と $e_2$ の冪等性を用いた直接の計算により、
$(e_1 - e_2)^3 = e_1 - e_2$ であることが分かる。
これと帰納法により、任意の正の奇数 $k$ に対して
$(e_1 - e_2)^k = e_1 - e_2$ を得る。
$e_1 - e_2$ は冪零なので、十分大きいすべての $k$ に対して
$(e_1 - e_2)^k = 0$ である。したがって $e_1 - e_2 = 0$、
すなわち $e_1 = e_2$ となり、証明が完了する。
\end{proof}

\begin{lemma}
\label{algebra-lemma-lift-idempotents-noncommutative}
$A$ を可換とは限らない代数とする。
$e \in A$ を、$x = e^2 - e$ が冪零となる元とする。
このとき、$a_{i, j} \in \mathbf{Z}$ を用いて
$e' = e + x(\sum a_{i, j}e^ix^j) \in A$ の形に書ける冪等元が存在する。
\end{lemma}

\begin{proof}
環 $R_n = \mathbf{Z}[e]/((e^2 - e)^n)$ を考える。
各 $R_n$ について結果を証明できれば補題が従うことは明らかである。
$R_n$ においてイデアル $I = (e^2 - e)$ を考え、補題
\ref{algebra-lemma-lift-idempotents} を適用する。
\end{proof}

\begin{lemma}
\label{algebra-lemma-lift-nth-roots}
$R$ を環、$I \subset R$ を局所冪零イデアルとする。
$n \geq 1$ を $R/I$ で可逆な整数とする。このとき
\begin{enumerate}
\item $n$ 乗写像 $1 + I \to 1 + I$、$1 + x \mapsto (1 + x)^n$ は全単射である。
% P01-E0330: frozen English source indefinite-article candidate; see ERRATA_CANDIDATES.jsonl.
\item $R$ の単元が $n$ 乗であるための必要十分条件は、その $R/I$ における像が
$n$ 乗であることである。
\end{enumerate}
\end{lemma}

\begin{proof}
$a \in R$ を単元とし、その $R/I$ における像が
$b \in R$ に対する $b^n$ の像と同じであるとする。
このとき $b$ は単元であり（補題 \ref{algebra-lemma-locally-nilpotent-unit}）、
ある $x \in I$ に対して $ab^{-n} = 1 + x$ である。
したがって (1) により $ab^{-n} = c^n$ である。
ゆえに (2) は (1) から従う。

\medskip\noindent
(1) の証明。写像 $1 + x \mapsto (1 + x)^n$ には逆写像があるため、
主張は成り立つ。すなわち、$1 + x$ を
\begin{align*}
(1 + x)^{1/n}
& =
1 + {1/n \choose 1}x +
{1/n \choose 2}x^2 +
{1/n \choose 3}x^3 + \ldots \\
& =
1 + \frac{1}{n} x + \frac{1 - n}{2n^2}x^2 +
\frac{(1 - n)(1 - 2n)}{6n^3}x^3 + \ldots
\end{align*}
に写す写像を、初等微積分と同様に考えればよい。
この級数は $k \gg 0$ に対して $x^k = 0$ となるため有限であり、
各係数 ${1/n \choose k} \in \mathbf{Z}[1/n]$ なので意味をもつ
（詳細は省略する。補題 \ref{algebra-lemma-locally-nilpotent-unit} により
$n$ が $R$ で可逆であることに注意せよ）。
\end{proof}

\section{余談}
\label{algebra-section-curiosity}

\noindent
補題 \ref{algebra-lemma-disjoint-implies-product} は、あるイデアル $I \subset R$ に対して
$V(I)$ が開である場合に何が起こるかを説明している。
では、$\Spec(S^{-1}R)$ が $\Spec(R)$ において閉である場合はどうだろうか。
次の二つの補題は、この問いに部分的な答えを与える。さらに詳しくは
節 \ref{algebra-section-pure-ideals} を参照せよ。

\begin{lemma}
\label{algebra-lemma-invert-closed-quotient}
$R$ を環、$S \subset R$ を乗法的部分集合とする。
写像 $\Spec(S^{-1}R) \to \Spec(R)$ の像が閉であると仮定する。
このとき、あるイデアル $I \subset R$ に対して $S^{-1}R \cong R/I$ である。
\end{lemma}

\begin{proof}
$I = \Ker(R \to S^{-1}R)$ と置くと、その像は $V(I)$ に含まれる。
その像が、あるイデアル $I' \subset R$ に対する閉部分集合
$V(I') \subset \Spec(R)$ であるとする。すると $V(I') \subset V(I)$ である。
$f \in I'$ に対し、$f/1 \in S^{-1}R$ はすべての素イデアルに属する。
% P01-E0331: frozen English source uses containment for element membership; see ERRATA_CANDIDATES.jsonl.
したがって、ある $n \geq 1$ に対して $f^n$ は $S^{-1}R$ で零に写る
（補題 \ref{algebra-lemma-Zariski-topology}）。
ゆえに $V(I') = V(I)$ である。
すると、すべての $g \in S$ は $I$ を法として可逆である。
したがって、環準同型 $R/I \to S^{-1}R$ および $S^{-1}R \to R/I$ を得る。
第一の写像は $I$ の選び方により単射である。
第二の写像は $S^{-1}R \to S^{-1}(R/I) = R/I$ であり、
局所化の完全性によりその核は $S^{-1}I$ である。
$S^{-1}I = 0$ なので、第二の写像も単射である。
したがって $S^{-1}R \cong R/I$ である。
\end{proof}

\begin{lemma}
\label{algebra-lemma-invert-closed-split}
$R$ を環、$S \subset R$ を乗法的部分集合とする。
写像 $\Spec(S^{-1}R) \to \Spec(R)$ の像が閉であると仮定する。
$R$ がネーター的であるか、$\Spec(R)$ がネーター位相空間であるか、
または $S$ がモノイドとして有限生成であるならば、
ある環 $R'$ に対して $R \cong S^{-1}R \times R'$ である。
\end{lemma}

\begin{proof}
補題 \ref{algebra-lemma-invert-closed-quotient} により、あるイデアル $I \subset R$ に対して
$S^{-1}R \cong R/I$ である。補題 \ref{algebra-lemma-disjoint-implies-product} により、
$V(I)$ が開であることを示せば十分である。
$R$ がネーター的ならば、補題 \ref{algebra-lemma-Noetherian-topology} により
$\Spec(R)$ はネーター位相空間である。
$\Spec(R)$ がネーター位相空間ならば、その補集合
$\Spec(R) \setminus V(I)$ は準コンパクトである。Topology の補題
\ref{topology-lemma-Noetherian-quasi-compact} を参照せよ。
したがって、$V(I) = V(f_1, \ldots, f_n)$ となる有限個の
$f_1, \ldots, f_n \in I$ が存在する。
各 $f_i$ は $S^{-1}R$ で零に写るので、$i = 1, \ldots, n$ に対して
$gf_i = 0$ となる $g \in S$ が存在する。
したがって、所望のとおり $D(g) = V(I)$ である。
$S$ がモノイドとして有限生成である場合、$S$ が
$g_1, \ldots, g_m$ で生成されるとすると、
$S^{-1}R \cong R_{g_1 \ldots g_m}$ であり、
$V(I) = D(g_1 \ldots g_m)$ と結論できる。
\end{proof}















\section{ヒルベルトの零点定理}
\label{algebra-section-nullstellensatz}

\begin{theorem}[ヒルベルトの零点定理]
\label{algebra-theorem-nullstellensatz}
$k$ を体とする。
\begin{enumerate}
\item
\label{algebra-item-finite-kappa}
任意の極大イデアル $\mathfrak m \subset k[x_1, \ldots, x_n]$ に対して、
体拡大 $\kappa(\mathfrak m)/k$ は有限である。
\item
\label{algebra-item-polynomial-ring-Jacobson}
任意の根基イデアル $I \subset k[x_1, \ldots, x_n]$ は、
それを含む極大イデアルの共通部分である。
\end{enumerate}
同じことは、任意の有限型 $k$-代数に対して成り立つ。
\end{theorem}

\begin{proof}
任意の有限生成 $k$-代数はそのような多項式代数の商だから、定理の
(\ref{algebra-item-finite-kappa}) を多項式代数 $k[x_1, \ldots, x_n]$ の場合に
証明すれば十分である。これを $n$ に関する帰納法で証明する。
$n = 0$ の場合は明らかである。
$\mathfrak m$ を $k[x_1, \ldots, x_n]$ の極大イデアルとする。
$\mathfrak p \subset k[x_n]$ を $\mathfrak m$ と $k[x_n]$ の共通部分とする。

\medskip\noindent
$\mathfrak p \not = (0)$ ならば、$\mathfrak p$ は極大であり、
既約モニック多項式 $P$ によって生成される
（$k[x_n]$ におけるユークリッド互除法による）。すると
$k' = k[x_n]/\mathfrak p$ は $k$ の有限体拡大であり、
$\kappa(\mathfrak m)$ に含まれる。この場合、全射
$$
k'[x_1, \ldots, x_{n-1}]
\to
k'[x_1, \ldots, x_n] =
k' \otimes_k k[x_1, \ldots, x_n]
\longrightarrow
\kappa(\mathfrak m)
$$
を得るので、帰納法の仮定により $\kappa(\mathfrak m)$ は $k'$ の有限拡大である。
したがって、$\kappa(\mathfrak m)$ は $k$ 上も有限である。

\medskip\noindent
$\mathfrak p = (0)$ ならば、環拡大
$k[x_n] \subset k[x_1, \ldots, x_n]/\mathfrak m$ を考える。
これは有限生成環拡大であり、したがって補題 \ref{algebra-lemma-obvious-Noetherian} および
\ref{algebra-lemma-Noetherian-finite-type-is-finite-presentation} により有限表示である。
ゆえに、$\Spec(k[x_1, \ldots, x_n]/\mathfrak m)$ の $\Spec(k[x_n])$ における像は、
定理 \ref{algebra-theorem-chevalley} により構成可能である。
この像は $(0)$ を含むので、ある非零な $f\in k[x_n]$ に対する標準開集合
$D(f)$ を含むと結論できる。明らかに $D(f)$ は無限なので、
$k[x_1, \ldots, x_n]/\mathfrak m$ が体である
（したがってそのスペクトルが一点からなる）という仮定に矛盾する。

\medskip\noindent
(\ref{algebra-item-polynomial-ring-Jacobson}) の証明。
$R$ を $k$ 上有限型とし、$I \subset R$ を根基イデアルとする。
$f \in R$, $f \not \in I$ とする。$I \subset \mathfrak m$ かつ
$f \not \in \mathfrak m$ を満たす極大イデアル $\mathfrak m \subset R$ を
見つけなければならない。環 $(R/I)_f$ は零でない。実際、この環で
$1 = 0$ ならば $f^n \in I$ であり、$I$ は根基的なので $f \in I$ となって、
$f$ に関する仮定に反する。したがって、補題 \ref{algebra-lemma-Zariski-topology} により、
$(R/I)_f$ の極大イデアル $\mathfrak m'$ を選べる。
$\mathfrak m \subset R$ を $R$ における $\mathfrak m'$ の逆像とする。
すると $I \subset \mathfrak m$ かつ $f \not \in \mathfrak m$ である。
$\mathfrak m$ が $R$ の極大イデアルであることを示せば、証明は終わる。
明らかに
$$
k \subset R/\mathfrak m \subset \kappa(\mathfrak m').
$$
である。(\ref{algebra-item-finite-kappa}) により、体拡大
$\kappa(\mathfrak m')/k$ は有限である。したがって Fields の補題
\ref{fields-lemma-subalgebra-algebraic-extension-field} により、
$R/\mathfrak m$ は体である。
ゆえに $\mathfrak m$ は極大であり、証明が完了する。
\end{proof}

\begin{lemma}
\label{algebra-lemma-field-finite-type-over-domain}
$R$ を環、$K$ を体とする。
$R \subset K$ かつ $K$ が $R$ 上有限型ならば、
$R_f$ が体となるような $f \in R$ が存在し、
$K/R_f$ は有限体拡大である。
\end{lemma}

\begin{proof}
補題 \ref{algebra-lemma-characterize-image-finite-type} により、
$\Spec(K) \to \Spec(R)$ の像 $\{(0)\}$ に含まれる空でない開集合
$U \subset \Spec(R)$ が存在する。
$D(f) \subset U$、すなわち $D(f) = \{(0)\}$ となるように、
$f \in R$, $f \not = 0$ を選ぶ。
すると $R_f$ は、そのスペクトルがちょうど一点からなる整域であり、$R_f$ は体である。
また、$K$ は体 $R_f$ 上の有限生成代数であり、したがってヒルベルトの零点定理
（定理 \ref{algebra-theorem-nullstellensatz}）により $R_f$ の有限体拡大である。
\end{proof}

\section{ジャコブソン環}
\label{algebra-section-ring-jacobson}

\noindent
$R$ を環とする。$\Spec(R)$ の閉点は $R$ の極大イデアルである。
代数幾何において自然に現れる環は、多くの場合、極大イデアルを豊富にもつ。
たとえば、体上または $\mathbf{Z}$ 上の有限型代数がそうである。
これらがジャコブソン環の例であることを示す。

\begin{definition}
\label{algebra-definition-ring-jacobson}
$R$ を環とする。すべての根基イデアル $I$ が、
それを含む極大イデアルの共通部分であるとき、$R$ を
{\it ジャコブソン環}という。
\end{definition}

\begin{lemma}
\label{algebra-lemma-finite-type-field-Jacobson}
体上の任意の有限型代数はジャコブソン環である。
\end{lemma}

\begin{proof}
定理 \ref{algebra-theorem-nullstellensatz} と定義
\ref{algebra-definition-ring-jacobson} から従う。
\end{proof}

\begin{lemma}
\label{algebra-lemma-jacobson-prime}
$R$ を環とする。$R$ のすべての素イデアルが、それを含む極大イデアルの
共通部分であるならば、$R$ はジャコブソン環である。
\end{lemma}

\begin{proof}
すべての根基イデアル $I \subset R$ が、それを含む素イデアルの
共通部分であることから直ちに明らかである。
補題 \ref{algebra-lemma-Zariski-topology} を参照せよ。
\end{proof}

\begin{lemma}
\label{algebra-lemma-jacobson}
環 $R$ がジャコブソン環であるための必要十分条件は、$\Spec(R)$ が
ジャコブソンであることである。Topology の定義
\ref{topology-definition-space-jacobson} を参照せよ。
\end{lemma}

\begin{proof}
$R$ がジャコブソン環であると仮定する。
$Z \subset \Spec(R)$ を閉部分集合とする。
$Z$ の閉点全体が $Z$ で稠密であることを示さなければならない。
$U \subset \Spec(R)$ を、$U \cap Z$ が空でないような開集合とする。
$Z \cap U$ が $\Spec(R)$ の閉点を含むことを示せばよい。
標準開集合は $\Spec(R)$ の位相の基底をなすので、
$U = D(f)$ と仮定してよい。
補題 \ref{algebra-lemma-Zariski-topology} により、$I$ を根基イデアルとして
$Z = V(I)$ と仮定してよい。また、$f \not \in I$ である。
仮定により、$I \subset \mathfrak m$ かつ $f \not\in \mathfrak m$ となる
極大イデアル $\mathfrak m \subset R$ が存在する。
したがって、所望のとおり
$\mathfrak m \in D(f) \cap V(I) = U \cap Z$ である。

\medskip\noindent
逆に、$\Spec(R)$ がジャコブソンであると仮定する。
$I \subset R$ を根基イデアルとする。
$J = \cap_{I \subset \mathfrak m} \mathfrak m$ を、$I$ を含む極大イデアルの
共通部分とする。明らかに $J$ は根基イデアルであり、
$V(J) \subset V(I)$ である。また $V(J)$ は、$V(I)$ のすべての閉点を含む
$V(I)$ の最小の閉部分集合である。仮定により $V(J) = V(I)$ である。
一方、補題 \ref{algebra-lemma-Zariski-topology} は、ザリスキー閉集合と
根基イデアルの間に全単射があることを示すので、所望のとおり $I = J$ である。
\end{proof}

\begin{lemma}
\label{algebra-lemma-characterize-jacobson}
$R$ を環とする。$R$ がジャコブソン環でないならば、次を満たす
素イデアル $\mathfrak p \subset R$ と元 $f \in R$ が存在する。
\begin{enumerate}
\item $\mathfrak p$ は極大イデアルでない。
\item $f \not \in \mathfrak p$ である。
\item $V(\mathfrak p) \cap D(f) = \{\mathfrak p\}$ である。
\item $(R/\mathfrak p)_f$ は体である。
\end{enumerate}
他方、$R$ がジャコブソン環ならば、(1) と (2) を満たす任意の組
$(\mathfrak p, f)$ に対して、集合 $V(\mathfrak p) \cap D(f)$ は無限である。
\end{lemma}

\begin{proof}
$R$ がジャコブソン環でないと仮定する。
補題 \ref{algebra-lemma-jacobson} により、これは、閉点全体の集合
$T_0 \subset T$ が $T$ で稠密でないような閉部分集合
$T \subset \Spec(R)$ が存在することを意味する。
$T_0 \subset V(f)$ だが $T \not \subset V(f)$ となる $f \in R$ を選ぶ。
$T = V(I)$ ならば、$T \cap D(f)$ は $\Spec((R/I)_f)$ と同相であることに注意する。
補題 \ref{algebra-lemma-spec-closed} および \ref{algebra-lemma-standard-open} を参照せよ。
任意の環は極大イデアルをもつので
（補題 \ref{algebra-lemma-Zariski-topology}）、空間 $T \cap D(f)$ の閉点 $t$ を選べる。
このとき $t$ は、極大でない素イデアル $\mathfrak p \subset R$ に対応する
（$t \not \in T_0$ だからである）。したがって (1) が成り立つ。
構成により $f \not \in \mathfrak p$ なので (2) が成り立つ。
$t$ は $T \cap D(f)$ の閉点なので、
$V(\mathfrak p) \cap D(f) = \{\mathfrak p\}$、すなわち (3) が成り立つ。
したがって、$(R/\mathfrak p)_f$ は、そのスペクトルが一点からなる整域であり、
(4) が成り立つと結論できる
（たとえば補題 \ref{algebra-lemma-characterize-local-ring} と
\ref{algebra-lemma-minimal-prime-reduced-ring} を組み合わせよ）。

\medskip\noindent
逆に、$R$ がジャコブソン環で、$(\mathfrak p, f)$ が (1) と (2) を満たすとする。
もし
$V(\mathfrak p) \cap D(f) =
\{\mathfrak p, \mathfrak q_1, \ldots, \mathfrak q_t\}$
ならば、$\mathfrak p \not = \mathfrak q_i$ なので、
$g \not \in \mathfrak p$ だがすべての $i$ に対して
$g \in \mathfrak q_i$ となる元 $g \in R$ が存在する。
したがって $V(\mathfrak p) \cap D(fg) = \{\mathfrak p\}$ である。
しかし、$\Spec(R)$ の各局所閉部分集合は少なくとも一つの閉点を含む。
これは $\Spec(R)$ がジャコブソン位相空間だからであり、矛盾する。
\end{proof}

\begin{lemma}
\label{algebra-lemma-pid-jacobson}
環 $\mathbf{Z}$ はジャコブソン環である。
より一般に、$R$ を次を満たす環とする。
\begin{enumerate}
\item $R$ は整域である。
\item $R$ はネーター的である。
\item 任意の非零素イデアルは極大イデアルである。
\item $R$ は無限個の極大イデアルをもつ。
\end{enumerate}
このとき $R$ はジャコブソン環である。
\end{lemma}

\begin{proof}
$R$ が (1)、(2)、(3)、(4) を満たすとする。
主張は $(0) = \bigcap_{\mathfrak m \subset R} \mathfrak m$ を意味する。
$R$ は無限個の極大イデアルをもつので、任意の非零元 $x \in R$ が
高々有限個の極大イデアルに属すること、言い換えれば $V(x)$ が有限であることを
示せば十分である。
% P01-E0332: frozen English source uses containment for element membership; see ERRATA_CANDIDATES.jsonl.
補題 \ref{algebra-lemma-spec-closed} により、$V(x)$ は $\Spec(R/xR)$ と同相である。
仮定 (3) により、$R/xR$ の各素イデアルは極小であり、したがって
$\Spec(R/xR)$ の既約成分に対応する
（補題 \ref{algebra-lemma-irreducible}）。
$R/xR$ はネーター的なので、位相空間 $\Spec(R/xR)$ はネーター的であり
（補題 \ref{algebra-lemma-Noetherian-topology}）、有限個の既約成分をもつ
（Topology の補題 \ref{topology-lemma-Noetherian}）。
したがって、所望のとおり $V(x)$ は有限である。
\end{proof}

\begin{example}
\label{algebra-example-infinite-product-fields-jacobson}
$A$ を無限集合とする。各 $\alpha \in A$ に対して、$k_\alpha$ を体とする。
$R = \prod_{\alpha\in A} k_\alpha$ はジャコブソン環であると主張する。
まず、任意の元 $f \in R$ は、$u \in R$ を単元、$e\in R$ を冪等元として
$f = ue$ の形をもつことに注意する（読者に委ねる）。
したがって $D(f) = D(e)$ であり、$R_f = R_e = R/(1-e)$ は $R$ の商である。
実際、この性質をもつ任意の環はジャコブソン環である。
すなわち、$\mathfrak p \subset R$ を素イデアル、
$f \in R$, $f \not \in \mathfrak p$ とする。
$\mathfrak p \subset \mathfrak m$ かつ $f \not\in \mathfrak m$ となる
$R$ の極大イデアル $\mathfrak m$ を見つけなければならない。
$R_f$ は $R$ の商なので、$R_f$ の各極大イデアルは、$f$ を含まない
$R$ の極大イデアルに対応する。したがって、$\mathfrak p R_f$ を含む
$R_f$ の極大イデアルを選べば結果が従う。
\end{example}

\begin{example}
\label{algebra-example-not-jacobson}
有限個の極大イデアル $\mathfrak m_i$, $i = 1, \ldots, n$ をもつ整域 $R$ は、
体である場合を除いてジャコブソン環でない。
実際、この場合、$(0)$ は極大イデアルの共通部分でない。すなわち
$(0) \not =
\mathfrak m_1 \cap \mathfrak m_2 \cap \ldots \cap \mathfrak m_n
\supset \mathfrak m_1 \cdot \mathfrak m_2 \cdot \ldots
\cdot \mathfrak m_n \not = 0$ である。
特に、離散付値環、または少なくとも二つの素イデアルをもつ任意の局所環は
ジャコブソン環でない。
\end{example}

\begin{lemma}
\label{algebra-lemma-finite-residue-extension-closed}
$R \to S$ を環準同型とする。
$\mathfrak m \subset R$ を極大イデアルとする。
$\mathfrak q \subset S$ を $\mathfrak m$ の上にある素イデアルで、
$\kappa(\mathfrak q)/\kappa(\mathfrak m)$ が代数的体拡大となるものとする。
このとき $\mathfrak q$ は $S$ の極大イデアルである。
\end{lemma}

\begin{proof}
図式
$$
\xymatrix{
S \ar[r] & S/\mathfrak q \ar[r] & \kappa(\mathfrak q) \\
R \ar[r] \ar[u] & R/\mathfrak m \ar[u]
}
$$
を考える。$\kappa(\mathfrak m) \subset S/\mathfrak q \subset
\kappa(\mathfrak q)$ である。
体拡大 $\kappa(\mathfrak m) \subset \kappa(\mathfrak q)$ は代数的なので、
$\kappa(\mathfrak m)$ と $\kappa(\mathfrak q)$ の間の任意の環は体である
（Fields の補題 \ref{fields-lemma-subalgebra-algebraic-extension-field}）。
したがって $S/\mathfrak q$ は体であり、従って $\kappa(\mathfrak q)$ に等しい。
\end{proof}

\begin{lemma}
\label{algebra-lemma-dimension}
$k$ を体、$V$ を $k$ 上の非零ベクトル空間とする。
$V$ の次元（これは基数である）が $k$ の濃度より小さいと仮定する。
このとき、任意の線形作用素 $T : V \to V$ に対して、
$P(T)$ が可逆でないようなモニック多項式 $P(t) \in k[t]$ が存在する。
\end{lemma}

\begin{proof}
そうでなければ、$V$ は体 $k(t)$ 上のベクトル空間の構造を受け継ぐ。
しかし、$k$ 上の $k(t)$ の次元は少なくとも $k$ の濃度である。
たとえば、元 $\frac{1}{t - \lambda}$ が $k$-線形独立だからである。
\end{proof}

\noindent
ヒルベルトの零点定理の別の形を述べる。

\begin{theorem}
\label{algebra-theorem-uncountable-nullstellensatz}
$k$ を体とする。$S$ を、元 $\{x_i\}_{i \in I}$ によって $k$ 上生成される
$k$-代数とする。$I$ の濃度が $k$ の濃度より小さいと仮定する。このとき
\begin{enumerate}
\item すべての極大イデアル $\mathfrak m \subset S$ に対して、
体拡大 $\kappa(\mathfrak m)/k$ は代数的である。
\item $S$ はジャコブソン環である。
\end{enumerate}
\end{theorem}

\begin{proof}
$I$ が有限ならば、結果はヒルベルトの零点定理、定理
\ref{algebra-theorem-nullstellensatz} から従う。
以下の証明では $I$ が無限であると仮定する。
$S$ はその商だから、変数 $x_i$ の多項式環における極大イデアル
$\mathfrak m \subset k[\{x_i\}_{i \in I}]$ に対して結果を証明すれば十分である。
$I$ は無限なので、単項式
$x_{i_1}^{e_1} \ldots x_{i_r}^{e_r}$、$i_1, \ldots, i_r \in I$、
$e_1, \ldots, e_r \geq 0$ 全体の濃度は、高々 $I$ の濃度に等しい。
実際、$I \times \ldots \times I$ の濃度は $I$ の濃度であり、
$\bigcup_{n \geq 0} I^n$ の濃度も同じである。
（$I$ が有限ならばこれは成り立たず、その場合、この証明が機能するのは
$k$ が非可算のときに限られる。）

\medskip\noindent
矛盾を導くため、$k$ 上超越的な $T \in \kappa(\mathfrak m)$ を選ぶ。
$T$ による乗法で与えられる $k$-線形写像
$T : \kappa(\mathfrak m) \to \kappa(\mathfrak m)$ は、すべてのモニック多項式
$P(t) \in k[t]$ に対して $P(T)$ が可逆であるという性質をもつことに注意する。
また、$\kappa(\mathfrak m)$ は $k$ 上のベクトル空間
$k[\{x_i\}_{i \in I}]$ の商であり、その次元は上で見たように $\# I$ なので、
その $k$ 上の次元は高々 $I$ の濃度である。
これは補題 \ref{algebra-lemma-dimension} により不可能である。

\medskip\noindent
$S$ がジャコブソン環であることを次のように示す。
そうでなければ、補題 \ref{algebra-lemma-characterize-jacobson} により、
素イデアル $\mathfrak q \subset S$ と元 $f \in S$,
$f \not \in \mathfrak q$ で、$\mathfrak q$ は極大でなく、
$(S/\mathfrak q)_f$ が体となるものが存在する。
しかし、$(S/\mathfrak q)_f$ は高々 $\# I + 1$ 個の元で生成されることに注意する。
したがって、体拡大 $(S/\mathfrak q)_f/k$ は代数的である
（証明の前半による）。これは $\kappa(\mathfrak q)$ が $k$ の代数拡大であることを
含意するので、補題 \ref{algebra-lemma-finite-residue-extension-closed} により
$\mathfrak q$ は極大である。この矛盾により証明が完了する。
\end{proof}

\begin{lemma}
\label{algebra-lemma-base-change-Jacobson}
$k$ を体、$S$ を $k$-代数とする。
$S$ の濃度より大きい濃度をもつ任意の体拡大 $K/k$ に対して、次が成り立つ。
\begin{enumerate}
\item $S_K$ の各極大イデアル $\mathfrak m$ に対して、体
$\kappa(\mathfrak m)$ は $K$ 上代数的である。
\item $S_K$ はジャコブソン環である。
\end{enumerate}
\end{lemma}

\begin{proof}
$K$ の濃度が $S$ の濃度より大きくなるように $k \subset K$ を選ぶ。
$S$ の元は $K$-代数 $S_K$ を生成するので、定理
\ref{algebra-theorem-uncountable-nullstellensatz} を適用できる。
\end{proof}

\begin{example}
\label{algebra-example-countable-trick-does-not-work}
定理 \ref{algebra-theorem-uncountable-nullstellensatz} の証明で用いた方法は、
$k$ が可算体で $I$ も可算である場合には実際に機能しない。
$k$ を可算体とする。$x$ を変数とし、$k(x)$ を $x$ の有理関数体とする。
多項式代数 $R = k[x, \{x_f\}_{f \in k[x]-\{0\}}]$ を考える。
$I = (\{fx_f - 1\}_{f\in k[x] - \{0\}})$ と置く。
$I$ は $R$ の真のイデアルであることに注意する。
極大イデアル $I \subset \mathfrak m$ を選ぶ。
すると $k \subset R/\mathfrak m$ は $k(x)$ と同型であり、$k$ 上代数的でない。
\end{example}

\begin{lemma}
\label{algebra-lemma-Jacobson-invert-element}
$R$ をジャコブソン環、$f \in R$ とする。環 $R_f$ はジャコブソン環であり、
$R_f$ の極大イデアルは、$f$ を含まない $R$ の極大イデアルに対応する。
\end{lemma}

\begin{proof}
Topology の補題 \ref{topology-lemma-jacobson-inherited} により、
$D(f) = \Spec(R_f)$ はジャコブソンであり、$D(f)$ の閉点は、
たまたま $D(f)$ に含まれる $\Spec(R)$ の閉点に対応する。
したがって、補題 \ref{algebra-lemma-jacobson} により $R_f$ はジャコブソン環である。
\end{proof}

\begin{example}
\label{algebra-example-localize-not-preserve-closed-points}
補題 \ref{algebra-lemma-Jacobson-invert-element} が $R$ がジャコブソン環でない場合に
成り立たないことを示す簡単な例を挙げる。
環 $R = \mathbf{Z}_{(2)}$、すなわち素イデアル $(2)$ における
$\mathbf{Z}$ の局所化を考える。元 $2$ における $R$ の局所化は
$\mathbf{Q}$ と同型である。式で書けば $R_2 \cong \mathbf{Q}$ である。
明らかに、写像 $R \to R_2$ は $\Spec(\mathbf{Q})$ の閉点を
$\Spec(R)$ の生成点に写す。
\end{example}

\begin{example}
\label{algebra-example-infinite-localize-not-preserve-closed-points}
$R$ がジャコブソン環であっても、無限個の元で局所化すると
補題 \ref{algebra-lemma-Jacobson-invert-element} が成り立たないことを示す簡単な例を挙げる。
すなわち、$R = \mathbf{Z}$ とし、すべての非零元からなる集合における $R$ の局所化
$(R \setminus \{0\})^{-1}R \cong \mathbf{Q}$ を考える。
明らかに、写像 $\mathbf{Z} \to \mathbf{Q}$ は $\Spec(\mathbf{Q})$ の閉点を
$\Spec(\mathbf{Z})$ の生成点に写す。
\end{example}

\begin{lemma}
\label{algebra-lemma-Jacobson-mod-ideal}
$R$ をジャコブソン環、$I \subset R$ をイデアルとする。
環 $R/I$ はジャコブソン環であり、$R/I$ の極大イデアルは、
$I$ を含む $R$ の極大イデアルに対応する。
\end{lemma}

\begin{proof}
証明は補題 \ref{algebra-lemma-Jacobson-invert-element} の証明と同じである。
\end{proof}

\begin{lemma}
\label{algebra-lemma-silly-jacobson}
$R$ をジャコブソン環、$K$ を体とする。
$R \subset K$ であり、$K$ が $R$ 上有限型であるとする。
% P01-E0333: frozen English source has a malformed joined hypothesis; see ERRATA_CANDIDATES.jsonl.
このとき $R$ は体であり、$K/R$ は有限体拡大である。
\end{lemma}

\begin{proof}
まず $R$ は整域であることに注意する。
補題 \ref{algebra-lemma-field-finite-type-over-domain} により、ある非零元 $f \in R$ に対して
$R_f$ は体であり、$K/R_f$ は有限体拡大である。
したがって $(0)$ は $R_f$ の極大イデアルであり、補題
\ref{algebra-lemma-Jacobson-invert-element} により、$(0)$ は $R$ の極大イデアルであると結論する。
\end{proof}

\begin{proposition}
\label{algebra-proposition-Jacobson-permanence}
$R$ をジャコブソン環とし、$R \to S$ を有限型環準同型とする。このとき
\begin{enumerate}
\item 環 $S$ はジャコブソン環である。
\item 写像 $\Spec(S) \to \Spec(R)$ は閉点を閉点に写す。
\item $\mathfrak m \subset R$ の上にある極大イデアル
$\mathfrak m' \subset S$ に対して、体拡大
$\kappa(\mathfrak m')/\kappa(\mathfrak m)$ は有限である。
\end{enumerate}
\end{proposition}


\begin{proof}
$\mathfrak m' \subset S$ を極大イデアルとし、
$R \cap \mathfrak m' = \mathfrak m$ とする。
$R/\mathfrak m \to S/\mathfrak m'$ は補題 \ref{algebra-lemma-silly-jacobson} の条件を、
補題 \ref{algebra-lemma-Jacobson-mod-ideal} により満たす。
したがって $R/\mathfrak m$ は体であり、$\mathfrak m$ は極大イデアルで、
誘導される剰余体拡大は有限である。これで (2) と (3) が示された。

\medskip\noindent
$S$ がジャコブソン環でないならば、補題
\ref{algebra-lemma-characterize-jacobson} により、$S$ の極大でない素イデアル
$\mathfrak q$ と元 $g \in S$, $g \not\in \mathfrak q$ で、
$(S/\mathfrak q)_g$ が体となるものが存在する。
矛盾を導くため、$\mathfrak q$ が極大イデアルであることを示す。
$\mathfrak p = \mathfrak q \cap R$ と置く。
$R/\mathfrak p \to (S/\mathfrak q)_g$ は補題
\ref{algebra-lemma-silly-jacobson} の条件を、補題
\ref{algebra-lemma-Jacobson-mod-ideal} により満たす。
したがって $R/\mathfrak p$ は体であり、体拡大
$\kappa(\mathfrak p) \to (S/\mathfrak q)_g = \kappa(\mathfrak q)$ は有限で、
従って代数的である。すると補題 \ref{algebra-lemma-finite-residue-extension-closed} により、
$\mathfrak q$ は $S$ の極大イデアルである。矛盾である。
\end{proof}

\begin{lemma}
\label{algebra-lemma-corollary-jacobson}
$\mathbf{Z}$ 上の任意の有限型代数はジャコブソン環である。
\end{lemma}

\begin{proof}
補題 \ref{algebra-lemma-pid-jacobson} と命題
\ref{algebra-proposition-Jacobson-permanence} を組み合わせればよい。
\end{proof}

\begin{lemma}
\label{algebra-lemma-image-finite-type-map-Jacobson-rings}
$R \to S$ をジャコブソン環の有限型環準同型とする。
$X = \Spec(R)$ および $Y = \Spec(S)$ と書く。
$f : Y \to X$ を誘導されるスペクトルの写像とする。
$E \subset Y = \Spec(S)$ を構成可能集合とする。
位相空間の閉点全体の集合を下添字 ${}_0$ で表す。
\begin{enumerate}
\item $f(E)_0 = f(E_0) = X_0 \cap f(E)$ である。
\item 点 $\xi \in X$ が $f(E)$ に属するための必要十分条件は、
$\overline{\{\xi\}} \cap f(E_0)$ が $\overline{\{\xi\}}$ で稠密であることである。
\end{enumerate}
\end{lemma}

\begin{proof}
連続写像の可換図式
$$
\xymatrix{
E \ar[r] \ar[d] & Y \ar[d] \\
f(E) \ar[r] & X
}
$$
をもつ。$x \in f(E)$ が $f(E)$ で閉であると仮定する。
このとき $f^{-1}(\{x\})\cap E$ は空でなく、$E$ で閉である。
二つの包含
$$
f^{-1}(\{x\}) \cap E \subset E \subset Y
$$
に Topology の補題 \ref{topology-lemma-jacobson-inherited} を適用すると、
$Y$ で閉である点 $y \in f^{-1}(\{x\}) \cap E$ が存在することが分かる。
言い換えれば、$x$ に写る $y \in Y_0$ かつ $y \in E_0$ が存在する。
したがって $x \in f(E_0)$ である。これで
$f(E)_0 \subset f(E_0)$ が示された。
命題 \ref{algebra-proposition-Jacobson-permanence} により
$f(E_0) \subset X_0 \cap f(E)$ である。
包含 $X_0 \cap f(E) \subset f(E)_0$ は自明である。これで第一の主張が示された。

\medskip\noindent
$\xi \in f(E)$ と仮定する。補題
\ref{algebra-lemma-characterize-image-finite-type} により、集合
$f(E) \cap \overline{\{\xi\}}$ は $\overline{\{\xi\}}$ の稠密開部分集合を含む。
$X$ はジャコブソンなので、$f(E) \cap \overline{\{\xi\}}$ は閉点の稠密集合を含む。
Topology の補題 \ref{topology-lemma-jacobson-inherited} を参照せよ。
補題の (1) により結論を得る。

\medskip\noindent
他方、$\overline{\{\xi\}} \cap f(E_0)$ が $\overline{\{\xi\}}$ で
稠密であると仮定する。補題 \ref{algebra-lemma-constructible-is-image} により、
$E$ が $Y' := \Spec(S') \to Y$ の像となるような有限表示環準同型
$S \to S'$ が存在する。環準同型 $S \to S'$ に補題の前半を適用すると、
$E_0$ は $Y'_0$ の像である。したがって、$S$ を $S'$ に置き換えることにより、
$E = Y$ と仮定してよい。$\xi$ が $\mathfrak p \subset R$ に対応するとする。
図式
$$
\xymatrix{
S \ar[r] & S/\mathfrak p S \\
R \ar[r] \ar[u] & R/\mathfrak p \ar[u]
}
$$
を考える。この図式と、$V(\mathfrak p)$ における
$f(Y_0) \cap V(\mathfrak p)$ の稠密性から、射
$R/\mathfrak p \to S/\mathfrak p S$ は補題
\ref{algebra-lemma-domain-image-dense-set-points-generic-point} の条件 (2) を満たす。
したがって、$(0)$ に写る素イデアル
$\overline{\mathfrak q} \subset S/\mathfrak pS$ が存在すると結論できる。
言い換えれば、$S$ における $\overline{\mathfrak q}$ の逆像
$\mathfrak q$ は所望のとおり $\mathfrak p$ に写る。
\end{proof}

\noindent
以上の補題の結論は、$f$ の像を、その像の閉点全体から読み取れるということである。
写像が有限表示である場合には、構成可能集合の像が構成可能であることが分かるので、
これは少し扱いやすくなる。
% P01-E1285: frozen English source has malformed singular/plural image wording; see ERRATA_CANDIDATES.jsonl.
これを述べる前に記法を導入する。
$\text{Constr}(X)$ を構成可能集合全体の集合とする。
% P01-E1286: frozen English source omits “by” in the notation introduction; see ERRATA_CANDIDATES.jsonl.
$R \to S$ を環準同型とする。
$X = \Spec(R)$ および $Y = \Spec(S)$ と書く。
$f : Y \to X$ を誘導されるスペクトルの写像とする。
位相空間の閉点全体の集合を下添字 ${}_0$ で表す。


\begin{lemma}
\label{algebra-lemma-conclude-jacobson-Noetherian}
以上の記法を用いる。$R$ をネーター的ジャコブソン環と仮定する。
さらに $R \to S$ が有限型であると仮定する。可換図式
$$
\xymatrix{
\text{Constr}(Y) \ar[r]^{E \mapsto E_0} \ar[d]^{E \mapsto f(E)} &
\text{Constr}(Y_0) \ar[d]^{E \mapsto f(E)} \\
\text{Constr}(X) \ar[r]^{E \mapsto E_0} &
\text{Constr}(X_0)
}
$$
が存在する。ここで横向きの矢印は Topology の補題
\ref{topology-lemma-jacobson-equivalent-constructible} の全単射である。
\end{lemma}

\begin{proof}
$R \to S$ は有限型なので有限表示である。補題
\ref{algebra-lemma-Noetherian-finite-type-is-finite-presentation} を参照せよ。
したがって、シュヴァレーの定理
（定理 \ref{algebra-theorem-chevalley}）により、$X$ の構成可能集合の像は
$Y$ で構成可能である。
% P01-E0334: frozen English source reverses the domain and codomain; preserve its wording and see ERRATA_CANDIDATES.jsonl.
これと補題 \ref{algebra-lemma-image-finite-type-map-Jacobson-rings} を組み合わせれば、
補題が従う。
\end{proof}

\noindent
ジャコブソン環の使い方を例示するため、次の二つの例を挙げる。

\begin{example}
\label{algebra-example-product-matrices-zero}
$k$ を体とする。空間 $\Spec(k[x, y]/(xy))$ は二つの既約成分、
すなわち $x$ 軸と $y$ 軸をもつ。一般化として、
$$
R = k[x_{11}, x_{12}, x_{21}, x_{22}, y_{11}, y_{12}, y_{21}, y_{22}]/
\mathfrak a,
$$
とする。ここで $\mathfrak a$ は、$k[x_{11}, x_{12}, x_{21}, x_{22}, y_{11}, y_{12}, y_{21}, y_{22}]$
において、$2 \times 2$ 積行列
$$
\left(
\begin{matrix}
x_{11} & x_{12}\\
x_{21} & x_{22}
\end{matrix}
\right)
 \left(
\begin{matrix}
y_{11} & y_{12}\\
y_{21} & y_{22}
\end{matrix}
\right).
$$
の成分によって生成されるイデアルである。この例では $\Spec(R)$ を記述する。

\medskip\noindent
$\Spec(k[x, y]/(xy))$ に関する主張を次のように証明する。
$\mathfrak p \subset k[x, y]$ が $xy$ を含む任意のイデアルならば、
$x$ または $y$ のいずれかが $\mathfrak p$ に含まれることになる。
% P01-E0335: frozen English source asserts the product implication for an arbitrary ideal; see ERRATA_CANDIDATES.jsonl.
したがって、そのような極小素イデアルは $(x)$ と $(y)$ だけである。
$k$ が代数閉の場合、これらの成分の $\text{max-Spec}$ は、
それぞれ $y$ 軸と $x$ 軸の点集合として可視化できる。

\medskip\noindent
一般化について、
$k[x_{11}, x_{12}, x_{21}, x_{22}, y_{11}, y_{12}, y_{21}, y_{22}])$
のスペクトルの閉点を、行列の空間
% P01-E0336: frozen English source has an unmatched closing parenthesis inside this formula; see ERRATA_CANDIDATES.jsonl.
$$
\left\{ (X, Y) \in \text{Mat}(2, k)\times \text{Mat}(2, k) \mid
X = \left(
\begin{matrix}
x_{11} & x_{12}\\
x_{21} & x_{22}
\end{matrix}
\right),
Y= \left(
\begin{matrix}
y_{11} & y_{12}\\
y_{21} & y_{22}
\end{matrix}
\right)
\right\}
$$
と同一視できることに注意する。少なくとも $k$ が代数閉ならばそうである。
ここで、行列の空間 $\{(X, Y)\}$ 上の
$\text{GL}(2, k)\times \text{GL}(2, k)\times \text{GL}(2, k)$ の群作用を
$$
(g_1, g_2, g_3) \times (X, Y) \mapsto ((g_1Xg_2^{-1}, g_2Yg_3^{-1})).
$$
で定める。また、代数的集合
$$
\text{GL}(2, k)\times \text{GL}(2, k)\times \text{GL}(2, k) \subset
\text{Mat}(2, k)\times \text{Mat}(2, k) \times \text{Mat}(2, k)
$$
は、整域
$$
k[x_{11}, x_{12}, \ldots, z_{21}, z_{22}, (x_{11}x_{22}-x_{12}x_{21})^{-1}
, (y_{11}y_{22}-y_{12}y_{21})^{-1}, (z_{11}z_{22}-z_{12}z_{21})^{-1}].
$$
の極大スペクトルなので既約であることにも注意する。
既約な代数的集合の像はなお既約だから、集合
$\{(X, Y)\in \text{Mat}(2, k)\times \text{Mat}(2, k)|XY = 0\}$ の軌道を分類し、
それらの閉包を取れば十分である。
% P01-E1287: frozen English source has malformed word order in the irreducible-image sentence; see ERRATA_CANDIDATES.jsonl.
標準的な線形代数により、次の三つの場合に帰着される。
\begin{enumerate}
\item $g_1Xg_2^{-1} = I_{2\times 2}$ となる $\exists (g_1, g_2)$ が存在する。
このとき $Y$ は必ず $0$ であり、この条件は代数的集合として群作用で不変である。
したがって、この軌道は素イデアル
$(y_{11}, y_{12}, y_{21}, y_{22})$ によって定まる既約代数的集合に含まれる。
閉包を取ると、$(y_{11}, y_{12}, y_{21}, y_{22})$ が実際に一つの成分であると分かる。
\item
$$
g_1Xg_2^{-1} = \left(
\begin{matrix}
1 & 0 \\
0 & 0
\end{matrix}
\right).
$$
となる $\exists (g_1, g_2)$ が存在する。
この場合が起こるための必要十分条件は、$X$ が階数 1 の行列であることである。
さらに、そのような $X$ によって $Y$ が消されるための必要十分条件は
$$
x_{11}y_{11}+x_{12}y_{21} = 0; \quad x_{11}y_{12}+x_{12}y_{22} = 0;
$$
$$
x_{21}y_{11}+x_{22}y_{21} = 0; \quad x_{21}y_{12}+x_{22}y_{22} = 0.
$$
である。階数 1 の $X$ を固定すると、上の方程式を満たす非零の $Y$ 全体は、
次の理由により既約代数的集合をなす
（$Y = 0$ は前の場合に含まれる）。
% P01-E1288: frozen English source has malformed punctuation and a missing preposition in this reason clause; see ERRATA_CANDIDATES.jsonl.
$0 = g_1Xg_2^{-1}g_2Y$ は
$$
g_2Y = \left(
\begin{matrix}
0 & 0 \\
y_{21}' & y_{22}'
\end{matrix}
\right).
$$
を含意する。右から $g_3$ によるさらなる $\text{GL}(2, k)$ 作用を施すと、
$g_2Y$ は
$$
g_2Yg_3^{-1} = \left(
\begin{matrix}
0 & 0 \\
0 & 1
\end{matrix}
\right),
$$
の形にできる。したがって、そのような $Y$ 全体は、この作用による
$\text{GL}(2, k)$ の像と同型な既約代数的集合をなす。
最後に、$X$ に対する「階数 1」という条件は、既約代数的集合
$\det X = x_{11}x_{22} - x_{12}x_{21} = 0$ の開稠密部分集合をなすことに注意する。
これより、五つの方程式すべては、非零行列の組の空間の開部分集合において、
既約成分
$(x_{11}y_{11}+x_{12}y_{21}, x_{11}y_{12}+x_{12}y_{22}, x_{21}y_{11}
+x_{22}y_{21}, x_{21}y_{12}+x_{22}y_{22}, x_{11}x_{22}-x_{12}x_{21})$
を定めることが分かる。
方程式の組 $\det X = 0$, $\det Y = 0$ は $\Spec(R)$ を、
上の軌跡を開稠密部分集合としてもつ既約成分に切り取ることが示せる。
\item $g_1Xg_2^{-1} = 0$ となる $\exists (g_1, g_2)$ が存在する。これは
$X = 0$ と同値である。このとき $Y$ は任意でよく、したがってこの成分は
$(x_{11}, x_{12}, x_{21}, x_{22})$ によって定まる。
\end{enumerate}
\end{example}


\begin{example}
\label{algebra-example-idempotent-matrices}
別の例として、$R = k[\{t_{ij}\}_{i, j = 1}^{n}]/\mathfrak a$ を考える。
ここで $\mathfrak a$ は、積行列 $T^2-T$、$T = (t_{ij})$ の成分によって
生成されるイデアルである。線形代数により、
$g, T \mapsto gTg^{-1}$ で定まる $GL(n, k)$ 作用のもとで、$T$ はその階数により
分類され、各 $T$ は、$r$ 個の 1 と $n-r$ 個の 0 をもつ
$\text{diag}(1, \ldots, 1, 0, \ldots, 0)$ のいずれかと共役である。
% P01-E0337: frozen English source duplicates the determiner in “by the its rank”; see ERRATA_CANDIDATES.jsonl.
したがって、このような $\text{diag}(1, \ldots, 1, 0, \ldots, 0)$ の
群作用による各軌道は既約成分をなし、各冪等行列はいずれか一つの軌道に含まれる。
% P01-E1289: frozen English source uses set containment for a matrix belonging to an orbit; see ERRATA_CANDIDATES.jsonl.
次に、異なる二つの軌道が必ず互いに素であることを示す。
このためには、異なる軌道上で異なる値を取る多項式関数を作ればよい。
標数 0 の場合、そのような関数として
$f(t_{ij}) = trace(T) = \sum_{i = 1}^nt_{ii}$ を取れる。
正標数の場合には、$T \neq 0$ でも $trace(T) = 0$ となりうるので、
事情は少し複雑である。たとえば $char = 3$ ならば
$$
trace\left(
\begin{matrix}
1 & & \\
& 1 & \\
& & 1
\end{matrix}
\right) = 3 = 0
$$
である。それでも、これらの成分は別の関数を用いて分離できる。
たとえば標数 3 の場合、$tr(\wedge^3T)$ は $diag(1, 1, 1)$ に対応する成分上で
値 1 を取り、他の成分上では値 0 を取る。
% P01-E0338: frozen English source's exterior-power trace claim fails for unrestricted matrix size; see ERRATA_CANDIDATES.jsonl.
\end{example}

\section{有限環拡大と整環拡大}
\label{algebra-section-finite-ring-extensions}

\noindent
有限環準同型と整環準同型に関する自明な補題を述べる。
定義を復習する。

\begin{definition}
\label{algebra-definition-integral-ring-map}
$\varphi : R \to S$ を環準同型とする。
\begin{enumerate}
\item 元 $s \in S$ に対して、$P^\varphi(s) = 0$ となるモニック多項式
$P(x) \in R[x]$ が存在するとき、この元は $R$ 上{\it 整}であるという。
ここで $P^\varphi(x) \in S[x]$ は、$\varphi : R[x] \to S[x]$ による
$P$ の像である。
\item すべての $s \in S$ が $R$ 上整であるとき、環準同型
$\varphi$ は{\it 整}であるという。
\end{enumerate}
\end{definition}

\begin{lemma}
\label{algebra-lemma-characterize-integral-element}
$\varphi : R \to S$ を環準同型、$y \in S$ とする。
$1 \in M$ かつ $yM \subset M$ となる $S$ の有限 $R$-部分加群 $M$ が
存在するならば、$y$ は $R$ 上整である。
\end{lemma}

\begin{proof}
写像 $\varphi : M \to M$, $x \mapsto y \cdot x$ を考える。
補題 \ref{algebra-lemma-charpoly-module} により、$P(\varphi) = 0$ となる
モニック多項式 $P \in R[T]$ が存在する。環 $S$ において
$P(y) = P(y) \cdot 1 = P(\varphi)(1) = 0$ を得る。
\end{proof}

\begin{lemma}
\label{algebra-lemma-finite-is-integral}
有限環準同型は整である。
\end{lemma}

\begin{proof}
$R \to S$ を有限とし、$y \in S$ とする。
$M = S$ に補題 \ref{algebra-lemma-characterize-integral-element} を適用すると、
$y$ が $R$ 上整であることが分かる。
\end{proof}

\begin{lemma}
\label{algebra-lemma-characterize-integral}
$\varphi : R \to S$ を環準同型とし、$s_1, \ldots, s_n$ を
$S$ の有限個の元とする。このとき、すべての $i = 1, \ldots, n$ に対して
$s_i$ が $R$ 上整であるための必要十分条件は、すべての $s_i$ を含み、
$R$ 上有限である $R$-部分代数 $S' \subset S$ が存在することである。
\end{lemma}

\begin{proof}
各 $s_i$ が整ならば、$\varphi(R)$ と $s_i$ で生成される部分代数は
$R$ 上有限である。実際、$s_i$ が $R$ 上次数 $d_i$ のモニック方程式を
満たすならば、この部分代数は、$0 \leq e_i \leq d_i - 1$ を満たす元
$s_1^{e_1} \ldots s_n^{e_n}$ によって $R$-加群として生成される。
逆に、すべての $s_i$ を含む有限 $R$-部分代数 $S'$ が与えられたとする。
すると補題 \ref{algebra-lemma-finite-is-integral} により、すべての $s_i$ は整である。
\end{proof}

\begin{lemma}
\label{algebra-lemma-characterize-finite-in-terms-of-integral}
$R \to S$ を環準同型とする。次は同値である。
\begin{enumerate}
\item $R \to S$ は有限である。
\item $R \to S$ は整かつ有限型である。
\item $R$ 上の代数として $S$ を生成する $x_1, \ldots, x_n \in S$ が存在し、
各 $x_i$ は $R$ 上整である。
\end{enumerate}
\end{lemma}

\begin{proof}
補題 \ref{algebra-lemma-characterize-integral} から明らかである。
\end{proof}

\begin{lemma}
\label{algebra-lemma-integral-transitive}
\begin{slogan}
整環準同型の合成は整である
\end{slogan}
$R \to S$ と $S \to T$ を整環準同型とする。このとき $R \to T$ は整である。
\end{lemma}

\begin{proof}
$t \in T$ とする。$P(t) = 0$ となるモニック多項式
$P(x) \in S[x]$ を取る。$P$ の有限個の係数に補題
\ref{algebra-lemma-characterize-integral} を適用する。
したがって $t$ は、$R$ 上有限なある部分代数 $S' \subset S$ 上整である。
補題 \ref{algebra-lemma-characterize-integral} をもう一度適用して、$S'$ 上有限で
$t$ を含む部分代数 $T' \subset T$ を得る。
$R \to S' \to T'$ に補題 \ref{algebra-lemma-finite-transitive} を適用すると、
$T'$ は $R$ 上有限であることが分かる。
すると補題 \ref{algebra-lemma-finite-is-integral} により、$t$ は $R$ 上整である。
\end{proof}

\begin{lemma}
\label{algebra-lemma-integral-closure-is-ring}
$R \to S$ を環準同型とする。集合
$$
S' = \{s \in S \mid s\text{ is integral over }R\}
$$
は $S$ の $R$-部分代数である。
\end{lemma}

\begin{proof}
補題 \ref{algebra-lemma-characterize-integral} と
\ref{algebra-lemma-finite-is-integral} から明らかである。
\end{proof}

\begin{lemma}
\label{algebra-lemma-finite-product-integral}
$R_i\to S_i$ を環準同型、$i = 1, \ldots, n$ とする。
$R$ と $S$ を、それぞれ $R_i$ と $S_i$ の積とする。
元 $s = (s_1, \ldots, s_n) \in S$ が $R$ 上整であるための必要十分条件は、
各 $s_i$ が $R_i$ 上整であることである。
\end{lemma}

\begin{proof}
省略する。
\end{proof}

\begin{definition}
\label{algebra-definition-integral-closure}
$R \to S$ を環準同型とする。$R$ 上整な元全体からなる環
$S' \subset S$ を、$S$ における $R$ の{\it 整閉包}という。
補題 \ref{algebra-lemma-integral-closure-is-ring} を参照せよ。
$R \subset S$ のとき、$R = S'$ ならば $R$ は $S$ において
{\it 整閉}であるという。
\end{definition}

\noindent
特に、$R \to S$ が整であるための必要十分条件は、
$S$ における $R$ の整閉包が $S$ 全体であることである。

\begin{lemma}
\label{algebra-lemma-finite-product-integral-closure}
$R_i\to S_i$ を環準同型、$i = 1, \ldots, n$ とする。
$S_i$ における $R_i$ の整閉包を $S'_i$ と書く。
さらに $R$ と $S$ を、それぞれ $R_i$ と $S_i$ の積とする。
このとき、$S$ における $R$ の整閉包は $S'_i$ の積である。
特に、$R \to S$ が整閉であるための必要十分条件は、
各 $R_i \to S_i$ が整閉であることである。
\end{lemma}

\begin{proof}
補題 \ref{algebra-lemma-finite-product-integral} から直ちに従う。
\end{proof}

\begin{lemma}
\label{algebra-lemma-integral-closure-localize}
整閉包は局所化と可換である。すなわち、$A \to B$ を環準同型、
$S \subset A$ を乗法的部分集合とすると、$S^{-1}B$ における
$S^{-1}A$ の整閉包は $S^{-1}B'$ である。ここで $B' \subset B$ は
$B$ における $A$ の整閉包である。
\end{lemma}

\begin{proof}
局所化は完全なので $S^{-1}B' \subset S^{-1}B$ である。
$x \in B'$ および $f \in S$ とする。このとき、ある $a_i \in A$ に対して
$B$ において
$x^d + \sum_{i = 1, \ldots, d} a_i x^{d - i} = 0$ である。
したがって $S^{-1}B$ においても
$$
(x/f)^d + \sum\nolimits_{i = 1, \ldots, d} a_i/f^i (x/f)^{d - i} = 0
$$
である。このようにして、$S^{-1}B'$ は $S^{-1}B$ における
$S^{-1}A$ の整閉包に含まれることが分かる。
逆に、$x/f \in S^{-1}B$ が $S^{-1}A$ 上整であるとする。
すると、ある $a_i \in A$ および $f_i \in S$ に対して、$S^{-1}B$ において
$$
(x/f)^d + \sum\nolimits_{i = 1, \ldots, d} (a_i/f_i) (x/f)^{d - i} = 0
$$
である。これは、適当な $f' \in S$ に対して
$$
(f'f_1 \ldots f_d x)^d +
\sum\nolimits_{i = 1, \ldots, d}
f^i(f')^if_1^i \ldots f_i^{i - 1} \ldots f_d^i a_i
(f'f_1 \ldots f_dx)^{d - i} = 0
$$
であることを意味する。したがって $f'f_1\ldots f_dx \in B'$ であり、
所望のとおり $x/f \in S^{-1}B'$ である。
\end{proof}

\begin{lemma}
\label{algebra-lemma-integral-closure-stalks}
\begin{slogan}
環上の代数の元がその環上整であるための必要十分条件は、
その環のすべての素イデアルで局所的に整であることである。
\end{slogan}
$\varphi : R \to S$ を環準同型、$x \in S$ とする。次は同値である。
\begin{enumerate}
\item $x$ は $R$ 上整である。
\item すべての素イデアル $\mathfrak p \subset R$ に対して、元
$x \in S_{\mathfrak p}$ は $R_{\mathfrak p}$ 上整である。
\end{enumerate}
\end{lemma}

\begin{proof}
(1) が (2) を含意することは明らかである。(2) を仮定する。
$S' \subset S$ を、$\varphi(R)$ と $x$ で生成される $R$-代数とする。
$\mathfrak p$ を $R$ の素イデアルとする。このとき、ある
$a_i \in R_{\mathfrak p}$ に対して、$S_{\mathfrak p}$ において
$x^d + \sum_{i = 1, \ldots, d} \varphi(a_i) x^{d - i} = 0$ である。
したがって、現れる分母を調べることにより、ある
$f \in R$, $f \not \in \mathfrak p$ に対して $a_i \in R_f$ であり、
$S_f$ において
$x^d + \sum_{i = 1, \ldots, d} \varphi(a_i) x^{d - i} = 0$ であると分かる。
これは $S'_f$ が $R_f$ 上有限であることを含意する。
$\mathfrak p$ は任意であり、$\Spec(R)$ は準コンパクトなので
（補題 \ref{algebra-lemma-quasi-compact}）、$R$ の単位イデアルを生成する有限個の元
$f_1, \ldots, f_n \in R$ で、各 $S'_{f_i}$ が $R_{f_i}$ 上有限となるものを
見つけられる。したがって補題 \ref{algebra-lemma-cover} により、$S'$ は $R$ 上有限である。
ゆえに補題 \ref{algebra-lemma-characterize-integral} により、$x$ は $R$ 上整である。
\end{proof}

\begin{lemma}
\label{algebra-lemma-base-change-integral}
\begin{slogan}
整性と有限性は基底変換で保たれる。
\end{slogan}
$R \to S$ と $R \to R'$ を環準同型とし、$S' = R' \otimes_R S$ と置く。
\begin{enumerate}
\item $R \to S$ が整ならば $R' \to S'$ も整である。
\item $R \to S$ が有限ならば $R' \to S'$ も有限である。
\end{enumerate}
\end{lemma}

\begin{proof}
(1) を証明する。$s_i \in S$ を $R$ 上の $S$ の生成元とする。
各元は $R$ 上モニック多項式方程式 $P_i$ を満たす。
したがって、元 $1 \otimes s_i \in S'$ は $R'$ 上 $S'$ を生成し、
$R'$ 上の対応する多項式 $P_i'$ を満たす。
これらの元は $R'$ 上 $S'$ を生成するので、$S'$ は $R'$ 上整である。
(2) の証明は省略する。
\end{proof}

\begin{lemma}
\label{algebra-lemma-integral-local}
$R \to S$ を環準同型とする。$f_1, \ldots, f_n \in R$ が
単位イデアルを生成するとする。
\begin{enumerate}
\item 各 $R_{f_i} \to S_{f_i}$ が整ならば、$R \to S$ も整である。
\item 各 $R_{f_i} \to S_{f_i}$ が有限ならば、$R \to S$ も有限である。
\end{enumerate}
\end{lemma}

\begin{proof}
(1) の証明。$s \in S$ とする。
$P(s) = 0$ となる多項式 $P$ 全体からなるイデアル
$I \subset R[x]$ を考える。$J \subset R$ を、$I$ の元の最高次係数全体からなる
イデアル (!) とする。仮定と分母の払出しにより、すべての $i$ とある
$n_i \geq 0$ に対して $f_i^{n_i} \in J$ である。
したがって $J$ は $1$ を含み、$s$ は $R$ 上整である。
(2) の証明は省略する。
\end{proof}

\begin{lemma}
\label{algebra-lemma-integral-permanence}
$A \to B \to C$ を環準同型とする。
\begin{enumerate}
\item $A \to C$ が整ならば $B \to C$ も整である。
\item $A \to C$ が有限ならば $B \to C$ も有限である。
\end{enumerate}
\end{lemma}

\begin{proof}
省略する。
\end{proof}

\begin{lemma}
\label{algebra-lemma-integral-closure-transitive}
$A \to B \to C$ を環準同型とする。
$B'$ を $B$ における $A$ の整閉包、$C'$ を $C$ における $B'$ の整閉包とする。
このとき $C'$ は $C$ における $A$ の整閉包である。
\end{lemma}

\begin{proof}
省略する。
\end{proof}

\begin{lemma}
\label{algebra-lemma-integral-overring-surjective}
$R \to S$ を $R \subset S$ である整環拡大とする。
このとき $\varphi : \Spec(S) \to \Spec(R)$ は全射である。
\end{lemma}

\begin{proof}
$\mathfrak p \subset R$ を素イデアルとする。
補題 \ref{algebra-lemma-in-image} により、
$\mathfrak pS_{\mathfrak p} \not = S_{\mathfrak p}$ を示さなければならない。
局所化 $R_{\mathfrak p} \to S_{\mathfrak p}$ は単射であり
（局所化は完全だからである）、補題 \ref{algebra-lemma-integral-closure-localize} または
\ref{algebra-lemma-base-change-integral} により整である。
したがって $R$, $S$ を $R_{\mathfrak p}$, $S_{\mathfrak p}$ に置き換え、
$R$ が極大イデアル $\mathfrak m$ をもつ局所環であると仮定してよい。
このとき $\mathfrak mS \not = S$ を示せば十分である。
矛盾を得るため、ある $f_i \in \mathfrak m$ と $s_i \in S$ に対して
$1 = \sum f_i s_i$ であると仮定する。
補題 \ref{algebra-lemma-characterize-integral} により、$R \to S'$ が有限で、
$s_i \in S'$ となる $R \subset S' \subset S$ を取る。
方程式 $1 = \sum f_i s_i$ は、有限 $R$-加群 $S'$ が
$S' = \mathfrak m S'$ を満たすことを含意する。
したがって中山の補題 \ref{algebra-lemma-NAK} により $S' = 0$ である。矛盾である。
\end{proof}

\begin{lemma}
\label{algebra-lemma-integral-under-field}
$R$ を環、$K$ を体とする。
$R \subset K$ かつ $K$ が $R$ 上整ならば、$R$ は体であり、
$K$ は代数拡大である。
$R \subset K$ かつ $K$ が $R$ 上有限ならば、$R$ は体であり、
$K$ は有限代数拡大である。
\end{lemma}

\begin{proof}
$R \subset K$ が整であると仮定する。
補題 \ref{algebra-lemma-integral-overring-surjective} により、$\Spec(R)$ は $1$ 点をもつ。
明らかに $R$ は整域なので、$R = R_{(0)}$ は体である
（補題 \ref{algebra-lemma-minimal-prime-reduced-ring}）。
他の主張はこれから直ちに従う。
\end{proof}

\begin{lemma}
\label{algebra-lemma-integral-over-field}
$k$ を体とし、$S$ を $k$ 上の $k$-代数とする。
\begin{enumerate}
\item $S$ が整域で、$k$ 上有限次元ならば、$S$ は体である。
\item $S$ が $k$ 上整な整域ならば、$S$ は体である。
\item $S$ が $k$ 上整ならば、$S$ の各素イデアルは極大イデアルである
（さらなる帰結については補題 \ref{algebra-lemma-ring-with-only-minimal-primes} を参照せよ）。
\end{enumerate}
\end{lemma}

\begin{proof}
素イデアルに関する主張は、「整 $+$ 整域 $\Rightarrow$ 体」という主張から従う。
$S$ が $k$ 上整で、$S$ が整域であるとする。$s \in S$ を取る。
% P01-E0339: frozen English source has a comma splice between the hypotheses and choice; see ERRATA_CANDIDATES.jsonl.
補題 \ref{algebra-lemma-characterize-integral} により、$s$ を含む有限次元 $k$-部分代数
$k \subset S' \subset S$ を見つけられる。
したがって、第一の主張を証明できれば $S$ は体である。
$S$ が $k$ 上有限次元の整域であると仮定する。$s\in S$ を取る。
% P01-E1290: frozen source omits the necessary nonzero hypothesis for this multiplication map; see ERRATA_CANDIDATES.jsonl.
$S$ は整域なので、乗法写像 $s : S \to S$ は次元の理由により全射である。
したがって $ss_1 = 1$ となる元 $s_1 \in S$ が存在する。
ゆえに $S$ は体である。
\end{proof}

\begin{lemma}
\label{algebra-lemma-integral-no-inclusion}
$R \to S$ が整であると仮定する。
$\mathfrak q, \mathfrak q' \in \Spec(S)$ を、$\Spec(R)$ において同じ像をもつ
相異なる素イデアルとする。このとき
$\mathfrak q \subset \mathfrak q'$ でも $\mathfrak q' \subset \mathfrak q$ でもない。
\end{lemma}

\begin{proof}
$\mathfrak p \subset R$ をその像とする。
注意 \ref{algebra-remark-fundamental-diagram} により、素イデアル
$\mathfrak q, \mathfrak q'$ は $S \otimes_R \kappa(\mathfrak p)$ のイデアルに対応する。
したがって補題 \ref{algebra-lemma-integral-over-field} から補題が従う。
\end{proof}

\begin{lemma}
\label{algebra-lemma-finite-finite-fibres}
$R \to S$ が有限であると仮定する。
このとき $\Spec(S) \to \Spec(R)$ のファイバーは有限である。
\end{lemma}

\begin{proof}
注意 \ref{algebra-remark-fundamental-diagram} の議論により、ファイバーは環
$S \otimes_R \kappa(\mathfrak p)$ のスペクトルである。
$R \to S$ は有限なので、これらのファイバー環は
$\kappa(\mathfrak p)$ 上有限であり、したがって補題
\ref{algebra-lemma-Noetherian-permanence} によりネーター的である。
補題 \ref{algebra-lemma-integral-no-inclusion} により、
$S \otimes_R \kappa(\mathfrak p)$ の各素イデアルは極小素イデアルである。
したがって補題 \ref{algebra-lemma-Noetherian-irreducible-components} により、
その個数は高々有限である。
\end{proof}

\begin{lemma}
\label{algebra-lemma-integral-going-up}
$R \to S$ を、$S$ が $R$ 上整であるような環準同型とする。
$\mathfrak p \subset \mathfrak p' \subset R$ を素イデアルとし、
$\mathfrak q$ を $\mathfrak p$ に写る $S$ の素イデアルとする。
このとき、$\mathfrak q \subset \mathfrak q'$ かつ $\mathfrak p'$ に写る
素イデアル $\mathfrak q'$ が存在する。
\end{lemma}

\begin{proof}
$R$ を $R/\mathfrak p$ に、$S$ を $S/\mathfrak q$ に置き換えてよい。
これにより、整域の整拡大 $R \subset S$ と素イデアル
$\mathfrak p' \subset R$ が与えられた状況に帰着する。
補題 \ref{algebra-lemma-integral-overring-surjective} により結論を得る。
\end{proof}

\noindent
以上の補題で表された性質を、環準同型 $R \to S$ の「上昇性」という。
定義 \ref{algebra-definition-going-up-down} を参照せよ。

\begin{lemma}
\label{algebra-lemma-finite-finitely-presented-extension}
$R \to S$ を有限かつ有限表示な環準同型とし、$M$ を $S$-加群とする。
このとき、$M$ が $R$-加群として有限表示であるための必要十分条件は、
$M$ が $S$-加群として有限表示であることである。
\end{lemma}

\begin{proof}
一方の含意は補題 \ref{algebra-lemma-finitely-presented-over-subring} から従う。
他方を示すため、$M$ が $S$-加群として有限表示であると仮定する。
表示
$$
S^{\oplus m} \longrightarrow
S^{\oplus n} \longrightarrow
M \longrightarrow 0
$$
を選ぶ。$S$ は $R$-加群として有限なので、$S^{\oplus n} \to M$ の核は
有限 $R$-加群である。したがって補題 \ref{algebra-lemma-extension} により、
$S$ が $R$-加群として有限表示であることを証明すれば十分である。

\medskip\noindent
$y_1, \ldots, y_n \in S$ を、$y_1, \ldots, y_n$ が $R$-加群として
$S$ を生成するように選ぶ。
補題 \ref{algebra-lemma-characterize-integral-element} により、各 $y_i$ は $R$ 上整である。
$P_i(y_i) = 0$ となるモニック多項式 $P_i(x) \in R[x]$ を選ぶ。
環
$$
S' = R[x_1, \ldots, x_n]/(P_1(x_1), \ldots, P_n(x_n))
$$
を考える。補題 \ref{algebra-lemma-compose-finite-type} により、$S$ は $S'$-代数として
有限表示である。また $S' \to S$ は全射なので、その核
$J = \Ker(S' \to S)$ は補題 \ref{algebra-lemma-finite-presentation-independent} により
イデアルとして有限生成である。したがって $J$ は有限 $S'$-加群である
（定義から直ちに従う）。
ゆえに $S = \Coker(J \to S')$ は $S'$-加群として有限表示である。
% P01-E0340: frozen English source has duplicate whitespace in this sentence; see ERRATA_CANDIDATES.jsonl.
補題 \ref{algebra-lemma-extension} を参照せよ。
したがって、第一段落と同様に論じると、$S'$ が $R$-加群として有限表示であることを
示せば十分である。実際、$S'$ は、$0 \leq e_i < \deg(P_i)$ に対する単項式
$x_1^{e_1} \ldots x_n^{e_n}$ を基底とする自由 $R$-加群である。
すなわち、$R \to S'$ を合成
$$
R \to R[x_1]/(P_1(x_1)) \to R[x_1, x_2]/(P_1(x_1), P_2(x_2)) \to
\ldots \to S'
$$
として書く。この列の $i$ 番目の環は、その $(i - 1)$ 番目の環上、
$1, x_i, \ldots, x_i^{\deg(P_i) - 1}$ を基底とする自由加群である。
これから帰納法により容易に結果が従う。いくつかの詳細は省略する。
\end{proof}

\begin{lemma}
\label{algebra-lemma-silly-normal}
$R$ を環、$x, y \in R$ を非零因子とする。
$R[x/y] \subset R_{xy}$ を $x/y$ で生成される $R$-部分代数とし、
部分代数 $R[y/x]$ および $R[x/y, y/x]$ も同様に定める。
$R$ が $R_x$ または $R_y$ において整閉ならば、列
$$
0 \to R \xrightarrow{(-1, 1)} R[x/y] \oplus R[y/x] \xrightarrow{(1, 1)}
R[x/y, y/x] \to 0
$$
は $R$-加群の短完全列である。
\end{lemma}

\begin{proof}
$x/y \cdot y/x = 1$ なので、写像
$R[x/y] \oplus R[y/x] \to R[x/y, y/x]$ が全射であることは明らかである。
$\alpha \in R[x/y] \cap R[y/x]$ とする。中間での完全性を示すには、
$\alpha \in R$ を証明しなければならない。仮定により、ある
$n, m \geq 0$ と $a_i, b_j \in R$ に対して
$$
\alpha = a_0 + a_1 x/y + \ldots + a_n (x/y)^n =
b_0 + b_1 y/x + \ldots + b_m(y/x)^m
$$
と書ける。$N > \max(n, m)$ となるものを選ぶ。
$R_{xy}$ の元
$$
(x/y)^N, (x/y)^{N - 1}, \ldots, x/y, 1, y/x, \ldots, (y/x)^{N - 1}, (y/x)^N
$$
で生成される有限 $R$-部分加群 $M$ を考える。
$\alpha M \subset M$ であると主張する。実際、
$0 \leq i \leq N$ に対して
$(x/y)^i (b_0 + b_1 y/x + \ldots + b_m(y/x)^m) \in M$ であり、
$0 \leq i \leq N$ に対して
$(y/x)^i (a_0 + a_1 x/y + \ldots + a_n(x/y)^n) \in M$ であることは明らかである。
したがって補題 \ref{algebra-lemma-characterize-integral-element} により、
$\alpha$ は $R$ 上整である。$\alpha \in R_x$ であることに注意する。
ゆえに、$R$ が $R_x$ において整閉ならば、所望のとおり $\alpha \in R$ である。
\end{proof}

\section{正規環}
\label{algebra-section-normal-rings}

\noindent
まず正規整域の概念を導入し、その後で（非常に一般的な）正規環の概念を導入する。

\begin{definition}
\label{algebra-definition-domain-normal}
整域 $R$ がその分数体において整閉であるとき、$R$ を{\it 正規}であるという。
\end{definition}

\begin{lemma}
\label{algebra-lemma-integral-closure-in-normal}
$R \to S$ を環準同型とする。
$S$ が正規整域ならば、$S$ における $R$ の整閉包は正規整域である。
\end{lemma}

\begin{proof}
省略する。
\end{proof}

\noindent
正規性を調べる際には、次の概念が役立つことがある。

\begin{definition}
\label{algebra-definition-almost-integral}
$R$ を整域とする。
\begin{enumerate}
\item $R$ の分数体の元 $g$ が{\it $R$ 上概整}であるとは、
ある元 $r \in R$, $r\not = 0$ が存在して、すべての $n \geq 0$ に対し
$rg^n \in R$ となることをいう。
\item 分数体の概整元がすべて $R$ に含まれるとき、整域 $R$ を
{\it 完整閉}であるという。
\end{enumerate}
\end{definition}

\noindent
次の補題は、ネーター整域が正規であることと完整閉であることが同値であることを示す。

\begin{lemma}
\label{algebra-lemma-almost-integral}
$R$ を分数体 $K$ をもつ整域とする。
$u, v \in K$ が $R$ 上概整ならば、$u + v$ と $uv$ もそうである。
$R$ 上整である任意の元 $g \in K$ は $R$ 上概整である。
$R$ がネーター的ならば、逆も成り立つ。
\end{lemma}

\begin{proof}
すべての $n \geq 0$ に対して $ru^n \in R$ であり、
すべての $n \geq 0$ に対して $v^nr' \in R$ ならば、
すべての $n \geq 0$ に対して $(uv)^nrr'$ と $(u + v)^nrr'$ は
$R$ に属する。したがって第一の主張が従う。
$g \in K$ が $R$ 上整であると仮定する。
% P01-E1291: frozen English source has the wrong article in “an d”; see ERRATA_CANDIDATES.jsonl.
このとき、環 $R[g]$ が $R$-加群として $1, g, \ldots, g^d$ で
生成されるような $d > 0$ が存在する。
$r \in R$ を $1, g, \ldots, g^d \in K$ の共通分母とする。
すると $rR[g] \subset R$ であり、したがって $g$ は $R$ 上概整である。

\medskip\noindent
$R$ がネーター的であり、$g \in K$ が $R$ 上概整であると仮定する。
$r \in R$, $r\not = 0$ を定義に現れる元とする。
このとき $R[g] \subset \frac{1}{r}R$ が $R$-加群として成り立つ。
$R$ はネーター的なので、これは $R[g]$ が $R$ 上有限であることを意味する。
したがって補題 \ref{algebra-lemma-finite-is-integral} により、$g$ は $R$ 上整である。
\end{proof}

\begin{lemma}
\label{algebra-lemma-localize-normal-domain}
正規整域の任意の局所化は正規である。
\end{lemma}

\begin{proof}
$R$ を正規整域、$S \subset R$ を乗法的部分集合とする。
$g$ を、$S^{-1}R$ 上整である $R$ の分数体の元とする。
$P = x^d + \sum_{j < d} a_j x^j$ を、$a_i \in S^{-1}R$ かつ
$P(g) = 0$ を満たす多項式とする。
% P01-E0341: frozen source changes the coefficient index from j to i and back; see ERRATA_CANDIDATES.jsonl.
すべての $i$ に対して $sa_i \in R$ となる $s \in S$ を選ぶ。
すると $sg$ は、係数 $s^{d-j}a_j$ が $R$ に属するモニック多項式
$x^d + \sum_{j < d} s^{d-j}a_j x^j$ を満たす。
$R$ は正規なので $sg \in R$ である。したがって $g \in S^{-1}R$ である。
\end{proof}

\begin{lemma}
\label{algebra-lemma-PID-normal}
単項イデアル整域は正規である。
\end{lemma}

\begin{proof}
$R$ を単項イデアル整域とする。
$g = a/b$ を、$R$ 上整である $R$ の分数体の元とする。
$R$ は単項イデアル整域なので、$a$ と $b$ の共通因子を除き、
$(a, b) = R$ と仮定してよい。このとき、$r_i \in R$ を係数とする任意の等式
$(a/b)^n + r_{n-1} (a/b)^{n-1} + \ldots + r_0 = 0$ は
$a^n \in (b)$ を含意する。これは $b$ が $R$ の単元でない限り、
$(a, b) = R$ に矛盾する。
\end{proof}

\begin{lemma}
\label{algebra-lemma-prepare-polynomial-ring-normal}
$R$ を分数体 $K$ をもつ整域とする。
$f = \sum \alpha_i x^i$ を $K[x]$ の元とする。
\begin{enumerate}
\item $f$ が $R[x]$ 上整ならば、すべての $\alpha_i$ は $R$ 上整であり、
\item $f$ が $R[x]$ 上概整ならば、すべての $\alpha_i$ は $R$ 上概整である。
\end{enumerate}
\end{lemma}

\begin{proof}
まず第二の主張を証明する。
$f = \alpha_0 + \alpha_1 x + \ldots + \alpha_r x^r$ と書き、
$\alpha_r \not = 0$ とする。
仮定により、$b_s \not = 0$ を満たす
$h = b_0 + b_1 x + \ldots + b_s x^s \in R[x]$ が存在し、
すべての $n \geq 0$ に対して $f^n h \in R[x]$ となる。
これは、すべての $n \geq 0$ に対して $b_s \alpha_r^n \in R$
となることを意味する。したがって $\alpha_r$ は $R$ 上概整である。
% P01-E1292: frozen English source has singular-subject agreement “the set ... form”; see ERRATA_CANDIDATES.jsonl.
概整元全体は部分環をなすので（補題 \ref{algebra-lemma-almost-integral}）、
$f - \alpha_r x^r = \alpha_0 + \alpha_1 x + \ldots + \alpha_{r - 1} x^{r - 1}$
は $R[x]$ 上概整である。$r$ に関する帰納法により結論を得る。

\medskip\noindent
第一の主張を証明するため、絶対ネーター還元を用いる。
すなわち $\alpha_i = a_i / b_i$ と書き、
$P(f) = 0$ を満たす、係数 $f_j \in R[x]$ の多項式
$P(t) = t^d + \sum_{j < d} f_j t^j$ をとる。
$f_j = \sum f_{ji}x^i$ と書く。$R_0 \subset R$ を、
$R$ の元 $a_i, b_i, f_{ji}$ の有限リストで生成される部分環とする。
これは整域であり、その分数体を $K_0$ とする。
$R_0$ は有限型 $\mathbf{Z}$-代数なので、補題
\ref{algebra-lemma-obvious-Noetherian} によりネーター的である。
$R$ における元 $a_i, b_i, f_{ji}$ の間のすべての恒等式は
$R_0$ においても成り立つので、依然として
$f \in K_0[x]$ は $R_0[x]$ 上整である。
補題 \ref{algebra-lemma-almost-integral} により、元 $f$ は $R_0[x]$ 上概整である。
補題の第二の主張により、元 $\alpha_i$ は $R_0$ 上概整である。
さらに $R_0$ はネーター的なので、補題 \ref{algebra-lemma-almost-integral} により
これらは $R_0$ 上整である。もちろん、するとこれらは $R$ 上整である。
\end{proof}

\begin{lemma}
\label{algebra-lemma-polynomial-domain-normal}
$R$ を正規整域とする。このとき $R[x]$ は正規整域である。
\end{lemma}

\begin{proof}
$R$ が体 $K$ ならば結果は成り立つ。実際、$K[x]$ はユークリッド整域であり、
したがって単項イデアル整域であり、補題 \ref{algebra-lemma-PID-normal} により正規である。
$g$ を、$R[x]$ 上整である $R[x]$ の分数体の元とする。
$R$ の分数体を $K$ とすれば、$g$ は $K[x]$ 上整なので、
$\alpha_i \in K$ を用いて
$g = \alpha_d x^d + \alpha_{d-1}x^{d-1} +
\ldots + \alpha_0$ と書ける。
補題 \ref{algebra-lemma-prepare-polynomial-ring-normal} により、元 $\alpha_i$ は
$R$ 上整であり、したがって $R$ に属する。
\end{proof}

\begin{lemma}
\label{algebra-lemma-power-series-over-Noetherian-normal-domain}
$R$ をネーター正規整域とする。このとき $R[[x]]$ は
ネーター正規整域である。
\end{lemma}

\begin{proof}
補題 \ref{algebra-lemma-Noetherian-power-series} により、冪級数環はネーター的である。
$f, g \in R[[x]]$ を非零元とし、$w = f/g$ が $R[[x]]$ 上整であるとする。
$R$ の分数体を $K$ とする。ローラン級数環
$K((x)) = K[[x]][1/x]$ は体なので、ある $n \in \mathbf{Z}$、
$a_i \in K$、$a_n \not = 0$ に対して
$w = a_n x^n + a_{n + 1} x^{n + 1} + \ldots$ と書ける。
補題 \ref{algebra-lemma-almost-integral} により、$b_m \not = 0$ を満たす
$R[[x]]$ の非零元
$h = b_m x^m + b_{m + 1} x^{m + 1} + \ldots$ が存在して、
すべての $e \geq 1$ に対して $w^e h \in R[[x]]$ となる。
このとき $n \geq 0$ であり、すべての $e \geq 1$ に対して
$b_m a_n^e \in R$ である。
$R$ はネーター的なので、同じ補題により $a_n \in R$ である。
ここで $a_n, a_{n + 1}, \ldots, a_{N - 1} \in R$ ならば、
同じ議論を
$w - a_n x^n - \ldots - a_{N - 1} x^{N - 1} = a_N x^N + \ldots$
に適用できる。このようにして、すべての $a_i \in R$ が分かり、
補題が証明される。
\end{proof}

\begin{lemma}
\label{algebra-lemma-normality-is-local}
$R$ を整域とする。次は同値である。
\begin{enumerate}
\item 整域 $R$ は正規整域である。
\item すべての素イデアル $\mathfrak p \subset R$ に対して、局所環
$R_{\mathfrak p}$ は正規整域である。
\item すべての極大イデアル $\mathfrak m$ に対して、環 $R_{\mathfrak m}$
は正規整域である。
\end{enumerate}
\end{lemma}

\begin{proof}
(1) $\Rightarrow$ (2) は補題 \ref{algebra-lemma-localize-normal-domain} から従う。
(2) $\Rightarrow$ (3) は直ちに従う。(3) $\Rightarrow$ (1) は、
任意の整域 $R$ に対し、$R$ の分数体の中で
$$
R = \bigcap\nolimits_{\mathfrak m} R_{\mathfrak m}
$$
となることから従う。実際、$g$ が右辺の元ならば、イデアル
$I = \{x \in R \mid xg \in R\}$ はどの極大イデアル $\mathfrak m$
にも含まれず、したがって $I = R$ である。
\end{proof}

\noindent
補題 \ref{algebra-lemma-normality-is-local} は、次の定義が定義
\ref{algebra-definition-domain-normal} と両立することを示す。
（これは EGA の定義である。すなわち \cite[IV, 5.13.5 and 0, 4.1.4]{EGA}
を参照せよ。）

\begin{definition}
\label{algebra-definition-ring-normal}
すべての素イデアル $\mathfrak p \subset R$ に対して局所化
$R_{\mathfrak p}$ が正規整域であるとき、環 $R$ を{\it 正規}であるという
（定義 \ref{algebra-definition-domain-normal} を参照）。
\end{definition}

\noindent
正規環は被約環であることに注意する。実際、$R$ はすべての素イデアルにおける
局所化の積の部分環である（例えば補題
\ref{algebra-lemma-characterize-zero-local} を参照）。

\begin{lemma}
\label{algebra-lemma-normal-ring-integrally-closed}
正規環はその全商環において整閉である。
\end{lemma}

\begin{proof}
$R$ を正規環とする。$x \in Q(R)$ を、$R$ 上整である $R$ の全商環の元とする。
$I = \{f \in R, fx \in R\}$ とおく。$\mathfrak p \subset R$ を素イデアルとする。
$R \to R_{\mathfrak p}$ は平坦なので、
$R_{\mathfrak p} \subset Q(R) \otimes_R R_{\mathfrak p}$ である。
$R_{\mathfrak p}$ は正規整域なので、$x \otimes 1$ は
$R_{\mathfrak p}$ の元である。したがって $f \not \in \mathfrak p$ を満たす
$a, f \in R$ を選んで、$x \otimes 1 = a \otimes 1/f$ とできる。
これは $fx - a$ が $Q(R) \otimes_R R_{\mathfrak p} = Q(R)_{\mathfrak p}$
において零に写ることを意味し、さらに $f' \not \in \mathfrak p$ を満たす
$f' \in R$ が存在して、$R$ において $f'fx = f'a$ となることを意味する。
言い換えると $ff' \in I$ である。ゆえに $I$ は $R$ のどの素イデアルにも
含まれないイデアルであり、すなわち $I = R$ であって $x \in R$ である。
\end{proof}

\begin{lemma}
\label{algebra-lemma-localization-normal-ring}
正規環の局所化は正規環である。
\end{lemma}

\begin{proof}
省略する。
\end{proof}

\begin{lemma}
\label{algebra-lemma-polynomial-ring-normal}
$R$ を正規環とする。このとき $R[x]$ は正規環である。
\end{lemma}

\begin{proof}
$\mathfrak q$ を $R[x]$ の素イデアルとする。
$\mathfrak p = R \cap \mathfrak q$ とおく。
補題 \ref{algebra-lemma-polynomial-domain-normal} により、
$R_{\mathfrak p}[x]$ は正規整域である。
したがって補題 \ref{algebra-lemma-localize-normal-domain} により、
$(R[x])_{\mathfrak q}$ は正規整域である。
\end{proof}

\begin{lemma}
\label{algebra-lemma-finite-product-normal}
正規環の有限積は正規である。
\end{lemma}

\begin{proof}
二つの正規環、例えば $R$ と $S$ の積が正規であることを示せば十分である。
補題 \ref{algebra-lemma-disjoint-decomposition} により、$R\times S$ の素イデアルは
$\mathfrak{p}\times S$ および $R\times
\mathfrak{q}$ の形である。ここで $\mathfrak{p}$ と $\mathfrak{q}$ はそれぞれ
$R$ と $S$ の素イデアルである。局所化すると
$(R\times S)_{\mathfrak{p}\times S}=R_{\mathfrak{p}}$ を得て、
これは仮定により正規整域である。$S$ についても同様である。
\end{proof}

\begin{lemma}
\label{algebra-lemma-characterize-reduced-ring-normal}
$R$ を環とする。$R$ が被約であり、極小素イデアルを有限個しかもたないと
仮定する。このとき次は同値である。
\begin{enumerate}
\item $R$ は正規環である。
\item $R$ はその全商環において整閉である。
\item $R$ は正規整域の有限積である。
\end{enumerate}
\end{lemma}

\begin{proof}
(1) $\Rightarrow$ (2) および (3) $\Rightarrow$ (1) は一般に成り立つ。
補題 \ref{algebra-lemma-normal-ring-integrally-closed} および
\ref{algebra-lemma-finite-product-normal} を参照せよ。

\medskip\noindent
$\mathfrak p_1, \ldots, \mathfrak p_n$ を $R$ の極小素イデアルとする。
補題 \ref{algebra-lemma-reduced-ring-sub-product-fields} および
\ref{algebra-lemma-total-ring-fractions-no-embedded-points} により
$Q(R) = R_{\mathfrak p_1} \times \ldots \times R_{\mathfrak p_n}$
であり、補題 \ref{algebra-lemma-minimal-prime-reduced-ring} により各因子は体である。
$Q(R)$ の第 $i$ 冪等元を
$e_i = (0, \ldots, 0, 1, 0, \ldots, 0)$ と書く。

\medskip\noindent
$R$ が $Q(R)$ において整閉ならば、とくに冪等元 $e_i$ を含む。
したがって $R$ は $n$ 個の整域の積である
（節 \ref{algebra-section-connected-components} および
\ref{algebra-section-tilde-module-sheaf} を参照）。
各因子は、分数体が $R_{\mathfrak p_i}$ である $R/\mathfrak p_i$ の形をしている。
補題 \ref{algebra-lemma-finite-product-integral-closure} により、各写像
$R/\mathfrak p_i \to R_{\mathfrak p_i}$ は整閉である。
したがって $R$ は正規整域の有限積である。
\end{proof}

\begin{lemma}
\label{algebra-lemma-colimit-normal-ring}
$(R_i, \varphi_{ii'})$ を環の有向系
（Categories, Definition \ref{algebra-definition-directed-system}）とする。
各 $R_i$ が正規環ならば、$R = \colim_i R_i$ も正規環である。
\end{lemma}

\begin{proof}
$\mathfrak p \subset R$ を素イデアルとする。
$\mathfrak p_i = R_i \cap \mathfrak p$ とおく（通常の記号の濫用である）。
すると
$R_{\mathfrak p} = \colim_i (R_i)_{\mathfrak p_i}$
である。各 $(R_i)_{\mathfrak p_i}$ は正規整域なので、正規整域について
補題の主張を証明することに帰着する。
$a, b \in R$ とし、$a/b$ がモニック多項式 $P(T) \in R[T]$ を満たすとする。
% P01-E1293: frozen source uses an undefined index set I; see ERRATA_CANDIDATES.jsonl.
このとき、$a, b$ が $R_i$ 上の対象 $a_i, b_i$ から来て、
$P$ がモニック多項式 $P_i\in R_i[T]$ から来ており、
$P_i(a_i/b_i)=0$ となる（十分大きい）$i \in I$ を見いだせる。
$R_i$ は正規なので $a_i/b_i \in R_i$ であり、したがって
$a/b \in R$ でもある。
\end{proof}

\section{正規環上の整拡大における下降}
\label{algebra-section-going-down-integral-over-normal}

\noindent
まず、イデアル上整な元という概念を少し調べ、その後で
節の表題に掲げた定理を証明する。

\begin{definition}
\label{algebra-definition-integral-over-ideal}
$\varphi : R \to S$ を環準同型とする。
$I \subset R$ をイデアルとする。
元 $g \in S$ が
{\it $I$ 上整}であるとは、係数が $a_j \in I^{d-j}$ である
モニック多項式
$P = x^d + \sum_{j < d} a_j x^j$
であって、$S$ において $P^\varphi(g) = 0$ となるものが存在することをいう。
\end{definition}

\noindent
これは主として $\varphi = \text{id}_R : R \to R$ の場合に用いる。
この場合、$I$ 上整な元からなる集合 $I'$ を
$I$ の{\it 整閉包}という。$I'$ が $R$ のイデアルであることを後で見る
（もちろん $I \subset I'$ である）。

\begin{lemma}
\label{algebra-lemma-characterize-integral-ideal}
$\varphi : R \to S$ を環準同型とする。
$I \subset R$ をイデアルとする。
$A = \sum I^nt^n \subset R[t]$ を、多項式環の部分集合
$R \oplus It \subset R[t]$ によって生成される部分環とする。
元 $s \in S$ が $I$ 上整であるための必要十分条件は、元
$st \in S[t]$ が $A$ 上整であることである。
\end{lemma}

\begin{proof}
$st$ が $A$ 上整であると仮定する。
係数が $A$ に属し $P^\varphi(st) = 0$ を満たすモニック多項式
$P = x^d + \sum_{j < d} a_j x^j$ をとる。
$a_j' \in A$ を $a_j$ の次数 $d-j$ の部分とする。言い換えると、
$a_j'' \in I^{d-j}$ を用いて $a_j' = a_j'' t^{d-j}$ と書ける。
次数を考えれば、依然として
$(st)^d + \sum_{j < d} \varphi(a_j'') t^{d-j} (st)^j = 0$ である。
したがって $s^d + \sum_{j < d} \varphi(a_j'') s^j = 0$ であり、
$s$ が $I$ 上整であることが分かる。

\medskip\noindent
$s$ が $I$ 上整であると仮定する。
$a_j \in I^{d-j}$ をもつ
$P = x^d + \sum_{j < d} a_j x^j$ があるとする。このとき、係数が
$A$ に属する多項式
$Q = x^d + \sum_{j < d} (a_j t^{d-j}) x^j$ が直ちに得られ、
$st$ が $A$ 上整であることが示される。
\end{proof}

\begin{lemma}
\label{algebra-lemma-integral-over-ideal-is-submodule}
$\varphi : R \to S$ を環準同型とする。
$I \subset R$ をイデアルとする。
% P01-E1294: frozen source has subject--verb and article errors; see ERRATA_CANDIDATES.jsonl.
$I$ 上整である $S$ の元全体は、$S$ の $R$-部分加群をなす。
さらに、$s \in S$ が $R$ 上整であり、$s'$ が $I$ 上整ならば、
$ss'$ は $I$ 上整である。
\end{lemma}

\begin{proof}
補題 \ref{algebra-lemma-integral-closure-is-ring} を、以下では断りなく用いる。
加法について閉じていることは、その記号を踏襲する補題
\ref{algebra-lemma-characterize-integral-ideal} の特徴づけから明らかである。
$R$ 上整な任意の元 $s \in S$ は、$S[t]$ の次数 $0$ の元 $s$ に対応し、
これは $A$ 上整である（$R \subset A$ だからである）。
% P01-E0342: frozen source proof ends here syntactically incomplete; see ERRATA_CANDIDATES.jsonl.
したがって、$S[t]$ における $s$ による乗法は、$A$ 上整であるという性質を保ち、
\end{proof}

\begin{lemma}
\label{algebra-lemma-integral-integral-over-ideal}
$\varphi : R \to S$ が整であると仮定する。
$I \subset R$ をイデアルとする。
このとき $IS$ のすべての元は $I$ 上整である。
\end{lemma}

\begin{proof}
補題 \ref{algebra-lemma-integral-over-ideal-is-submodule} から直ちに従う。
\end{proof}

\begin{lemma}
\label{algebra-lemma-polynomials-divide}
$K$ を体とする。$n, m \in \mathbf{N}$ および
$a_0, \ldots, a_{n - 1}, b_0, \ldots, b_{m - 1} \in K$ とする。
多項式 $x^n + a_{n - 1}x^{n - 1} + \ldots + a_0$ が
$K[x]$ において多項式
$x^m + b_{m - 1} x^{m - 1} + \ldots + b_0$ を割るならば、
\begin{enumerate}
\item $a_0, \ldots, a_{n - 1}$ は、元 $b_0, \ldots, b_{m - 1}$ を
含む $K$ の任意の部分環 $R_0$ 上整であり、
\item 元 $a_0,\allowbreak \ldots,\allowbreak a_{n - 1},\allowbreak
b_0,\allowbreak \ldots,\allowbreak b_{m - 1}$\hfill\break
を含む
任意の部分環 $R \subset K$ に対し、各 $a_i$ は
$\sqrt{(b_0, \ldots,\allowbreak b_{m-1})R}$ に属する。
\end{enumerate}
\end{lemma}

\begin{proof}
$x^m + b_{m - 1} x^{m - 1} + \ldots + b_0 =
\prod_{i = 1}^m (x - \beta_i)$ と $\beta_i \in L$ を用いて書けるような
体拡大 $L/K$ をとる。
Fields, Section \ref{fields-section-splitting-fieds} を参照せよ。
各 $\beta_i$ は $R_0$ 上整である。
各 $a_i$ は $\beta_1, \ldots, \beta_m$ の斉次多項式なので、
$a_i$ についても同じことが従う
（補題 \ref{algebra-lemma-integral-closure-is-ring} を用いよ）。これで (1) が示された。

\medskip\noindent
$R$ を (2) のとおりとする。次を満たす
$c_0, \ldots, c_{m - n - 1} \in K$ をとる。
$$
\begin{matrix}
x^m + b_{m - 1} x^{m - 1} + \ldots + b_0 =  \\
(x^n + a_{n - 1}x^{n - 1} + \ldots + a_0)
(x^{m - n} + c_{m - n - 1}x^{m - n - 1}+ \ldots + c_0).
\end{matrix}
$$
この等式は
\begin{align*}
c_{m - n - 1} & = b_{m - 1} - a_{n - 1}, \\
c_{m - n - 2} & = b_{m - 2} - a_{n - 2} - a_{n - 1}c_{m - n - 1}, \\
\ldots
\end{align*}
を含意する。したがって、すべての $j$ について $c_j \in R$ である。
根基 $\sqrt{(b_0, \ldots, b_{m - 1})}$ で割ると、被約環 $\overline{R}$ を得る。
像 $\overline{a}_i \in \overline{R}$ が零であることを示せばよい。そして
$\overline{R}[x]$ において関係式
$$
\begin{matrix}
x^m = x^m + \overline{b}_{m - 1} x^{m - 1} + \ldots + \overline{b}_0 = \\
(x^n + \overline{a}_{n - 1}x^{n - 1} + \ldots + \overline{a}_0)
(x^{m - n} + \overline{c}_{m - n - 1}x^{m - n - 1}+ \ldots + \overline{c}_0).
\end{matrix}
$$
が成り立つ。これからすべての $i$ について
$\overline{a}_i = 0$ となることは容易に分かる。実際、補題
\ref{algebra-lemma-minimal-prime-reduced-ring} により、極小素イデアル
$\mathfrak{p}$ における $\overline{R}$ の局所化は体であり、
$\overline{R}_{\mathfrak p}[x]$ は一意分解整域である。したがって
$f = x^n + \sum \overline{a}_i x^i$ は $x^n$ と同伴であり、
$f$ はモニックなので $\overline{R}_{\mathfrak p}[x]$ において $f = x^n$ である。
% P01-E1295: frozen source uses the wrong indefinite article before s; see ERRATA_CANDIDATES.jsonl.
すると、$s(f - x^n) = 0$ となる
$s \in \overline{R}$, $s \not\in \mathfrak p$ が存在する。
ゆえにすべての $\overline{a}_i$ は $\mathfrak p$ に属し、補題
\ref{algebra-lemma-reduced-ring-sub-product-fields} によって結論を得る。
\end{proof}

\begin{lemma}
\label{algebra-lemma-minimal-polynomial-normal-domain}
$R \subset S$ を整域の包含とする。
$R$ が正規であると仮定する。$g \in S$ を $R$ 上整とする。
このとき、$g$ の最小多項式の係数は $R$ に属する。
\end{lemma}

\begin{proof}
$P(g) = 0$ を満たす、係数が $R$ に属する多項式を
$P = x^m + b_{m-1} x^{m-1} + \ldots + b_0$ とする。
$R$ の分数体 $K$ 上での $g$ の最小多項式を
$Q = x^n + a_{n-1}x^{n-1} + \ldots + a_0$ とする。
このとき $Q$ は $K[x]$ において $P$ を割る。補題
\ref{algebra-lemma-polynomials-divide} により、$a_i$ は $R$ 上整である。
$R$ は正規なので、これはそれらが $R$ に属することを意味する。
\end{proof}

\begin{proposition}
\label{algebra-proposition-going-down-normal-integral}
$R \subset S$ を整域の包含とする。
$R$ は正規で、$S$ は $R$ 上整であると仮定する。
$\mathfrak p \subset \mathfrak p' \subset R$ を素イデアルとする。
$\mathfrak p' = R \cap \mathfrak q'$ を満たす $S$ の素イデアルを
$\mathfrak q'$ とする。このとき、$\mathfrak q \subset \mathfrak q'$ かつ
$\mathfrak p = R \cap \mathfrak q$ を満たす素イデアル
$\mathfrak q$ が存在する。言い換えると、$R \to S$ は下降性をもつ。
定義 \ref{algebra-definition-going-up-down} を参照せよ。
\end{proposition}

\begin{proof}
$\mathfrak p$, $\mathfrak p'$, $\mathfrak q'$ を主張のとおりとする。
$\mathfrak q \subset \mathfrak q'$ かつ
$R \cap \mathfrak q = \mathfrak p$ を満たす素イデアル
$\mathfrak q$ が存在することを示せばよい。これは、
$\mathfrak p$ に写る $S_{\mathfrak q'}$ の素イデアルを見いだすことと同じである。
補題 \ref{algebra-lemma-in-image} により、
$\mathfrak p S_{\mathfrak q'} \cap R = \mathfrak p$ を示せばよい。
$z \in \mathfrak p S_{\mathfrak q'} \cap R$ をとる。
$y \in \mathfrak pS$ および $g \in S$, $g \not\in \mathfrak q'$ を用いて
$z = y/g$ と書ける。別の書き方をすれば $zg = y$ である。

\medskip\noindent
補題 \ref{algebra-lemma-integral-integral-over-ideal} により、
$b_i \in \mathfrak p$ かつ $P(y) = 0$ となるモニック多項式
$P = x^m + b_{m-1} x^{m-1} + \ldots + b_0$ が存在する。

\medskip\noindent
補題 \ref{algebra-lemma-minimal-polynomial-normal-domain} により、
$K$ 上での $g$ の最小多項式の係数は $R$ に属する。それを
$Q = x^n + a_{n-1} x^{n-1} + \ldots + a_0$ と書く。
$a_i$, $i = n-1, \ldots, 0$ のすべてが $\mathfrak p$ に属するわけではない。
実際、もしすべてが属すれば、
% P01-E1296: frozen source omits the algebraically required minus sign; preserve canonical formula.
$g^n = \sum_{j < n} a_j g^j \in \mathfrak pS
\subset \mathfrak p'S
\subset \mathfrak q'$
となり、矛盾する。

\medskip\noindent
$y = zg$ なので、以上から
$Q' = x^n + za_{n-1} x^{n-1} + \ldots + z^{n}a_0$ が
$y$ の最小多項式であることが直ちに分かる。したがって
$Q'$ は $P$ を割り、補題 \ref{algebra-lemma-polynomials-divide} により
$z^ja_{n - j} \in \sqrt{(b_0, \ldots, b_{m-1})}
\subset \mathfrak p$, $j =  1, \ldots, n$ である。
$a_i$, $i = n-1, \ldots, 0$ のすべてが $\mathfrak p$ に属するわけではないので、
所望のとおり $z \in \mathfrak p$ と結論する。
\end{proof}

\section{平坦加群と平坦環準同型}
\label{algebra-section-flat}

\noindent
頻繁に用いられる結果の一つは、$M = \colim_{i\in \mathcal{I}} M_i$ が
$R$-加群の余極限であり、$N$ が $R$-加群ならば、
$$
M \otimes N
=
\colim_{i\in \mathcal{I}} M_i \otimes_R N,
$$
となることである。補題 \ref{algebra-lemma-tensor-products-commute-with-limits} を参照せよ。
この性質は通常、{\it $\otimes$ は余極限と可換である}と表現される。
もう一つ頻繁に用いられる結果は、
$0 \to N_1 \to N_2 \to N_3 \to 0$ が完全列で、$M$ が任意の
$R$-加群ならば、
$$
M \otimes_R N_1
\to
M \otimes_R N_2
\to
M \otimes_R N_3
\to
0
$$
も完全であることである。補題 \ref{algebra-lemma-tensor-product-exact} を参照せよ。
これら二つの性質はともに、関手
$N \mapsto M \otimes_R N$ が{\it 右完全}であることを意味する。
Categories, Section \ref{categories-section-exact-functor} および
Homology, Section \ref{homology-section-functors} を参照せよ。
$R$-加群 $M$ が平坦であるとは、$N \mapsto N \otimes_R M$ が左完全でもある、
すなわち完全であることをいう。正確な定義は次のとおりである。

\begin{definition}
\label{algebra-definition-flat}
$R$ を環とする。
\begin{enumerate}
\item $R$-加群 $M$ が{\it 平坦}であるとは、$R$-加群の完全列
$N_1 \to N_2 \to N_3$ が与えられるたびに、列
$M \otimes_R N_1 \to M \otimes_R N_2 \to M \otimes_R N_3$
も完全であることをいう。
\item $R$-加群 $M$ が{\it 忠実平坦}であるとは、$R$-加群の複体
$N_1 \to N_2 \to N_3$ が完全であるための必要十分条件が、列
$M \otimes_R N_1 \to M \otimes_R N_2 \to M \otimes_R N_3$
が完全であることとなる場合をいう。
\item 環準同型 $R \to S$ が{\it 平坦}であるとは、$S$ が
$R$-加群として平坦であることをいう。
\item 環準同型 $R \to S$ が{\it 忠実平坦}であるとは、$S$ が
$R$-加群として忠実平坦であることをいう。
\end{enumerate}
\end{definition}

\noindent
平坦性の条件の使い方の一例を挙げる。

\begin{lemma}
\label{algebra-lemma-flat-intersect-ideals}
$R$ を環とする。$I, J \subset R$ をイデアルとする。$M$ を平坦な
$R$-加群とする。このとき $IM \cap JM = (I \cap J)M$ である。
\end{lemma}

\begin{proof}
完全列 $0 \to I \cap J \to R \to R/I \oplus R/J$ を考える。
平坦加群 $M$ とテンソル積をとると、完全列
$$
0 \to (I \cap J) \otimes_R M \to M \to M/IM \oplus M/JM
$$
を得る。$M \to M/IM \oplus M/JM$ の核は $IM \cap JM$ に等しいので、
結論を得る。
\end{proof}


\begin{lemma}
\label{algebra-lemma-colimit-flat}
$R$ を環とする。$\{M_i, \varphi_{ii'}\}$ を平坦な $R$-加群の有向系とする。
このとき $\colim_i M_i$ は平坦な $R$-加群である。
\end{lemma}

\begin{proof}
$\otimes$ が余極限と可換であり、有向余極限が完全であることから従う。
補題 \ref{algebra-lemma-directed-colimit-exact} を参照せよ。
\end{proof}

\begin{lemma}
\label{algebra-lemma-composition-flat}
（忠実）平坦な環準同型の合成は（忠実）平坦である。
$R \to R'$ が（忠実）平坦で、$M'$ が（忠実）平坦な
$R'$-加群ならば、$M'$ は（忠実）平坦な $R$-加群である。
\end{lemma}

\begin{proof}
補題の第一の主張は第二の主張の特別な場合なので、後者を証明すれば十分である。
$R \to R'$ を平坦な環準同型、$M'$ を平坦な $R'$-加群とする。
$M'$ が平坦な $R$-加群であることを証明する必要がある。
$N_1 \to N_2 \to N_3$ を $R$-加群の完全な複体とする。
このとき、複体
$R' \otimes_R N_1 \to R' \otimes_R N_2 \to R' \otimes_R N_3$ は完全であり
（$R'$ が $R$-加群として平坦だからである）、したがって複体
$M' \otimes_{R'} \left(R' \otimes_R N_1\right)
\to M' \otimes_{R'} \left(R' \otimes_R N_2\right)
\to M' \otimes_{R'} \left(R' \otimes_R N_3\right)$ は完全である
（$M'$ が平坦な $R'$-加群だからである）。任意の $R$-加群 $N$ に対し、関手的に
$M' \otimes_{R'} \left(R' \otimes_R N\right)
\cong \left(M' \otimes_{R'} R'\right) \otimes_R N
\cong M' \otimes_R N$ であるから
（補題 \ref{algebra-lemma-tensor-with-bimodule} および
\ref{algebra-lemma-flip-tensor-product} による）、この複体は
$M' \otimes_R N_1 \to M' \otimes_R N_2 \to M' \otimes_R N_3$
と同型であり、したがって後者も完全である。これにより、$M'$ が平坦な
$R$-加群であることが示される。この議論を逆にたどれば、
$R \to R'$ が忠実平坦で、$M'$ が $R'$-加群として忠実平坦ならば、
$M'$ は $R$-加群として忠実平坦であることを示せる。
\end{proof}

\begin{lemma}
\label{algebra-lemma-flat}
$M$ を $R$-加群とする。次は同値である。
\begin{enumerate}
\item
\label{algebra-item-flat}
$M$ は $R$ 上平坦である。
\item
\label{algebra-item-injective}
$R$-加群の任意の単射 $N \subset N'$ に対し、写像
$N \otimes_R M \to N'\otimes_R M$ は単射である。
\item
\label{algebra-item-f-ideal}
任意のイデアル $I \subset R$ に対し、写像
$I \otimes_R M \to R \otimes_R M = M$ は単射である。
\item
\label{algebra-item-ffg-ideal}
任意の有限生成イデアル $I \subset R$ に対し、写像
$I \otimes_R M \to R \otimes_R M = M$ は単射である。
\end{enumerate}
\end{lemma}

\begin{proof}
（\ref{algebra-item-flat}）から（\ref{algebra-item-injective}）、
（\ref{algebra-item-f-ideal}）、（\ref{algebra-item-ffg-ideal}）へ至る含意の鎖はいずれも自明である。
したがって、（\ref{algebra-item-ffg-ideal}）ならば（\ref{algebra-item-flat}）を証明する。
$N_1 \to N_2 \to N_3$ が完全であると仮定する。
$K = \Ker(N_2 \to N_3)$ および $Q = \Im(N_2 \to N_3)$ とする。
すると写像
$$
N_1 \otimes_R M \to
K \otimes_R M \to
N_2 \otimes_R M \to
Q \otimes_R M \to
N_3 \otimes_R M
$$
を得る。第一および第三の矢印は全射であることに注意せよ。
したがって、第二および第四の矢印が単射であることを示せば十分である
\footnote{議論をさらに詳しく述べる。
第二および第四の矢印が単射であると分かっていると仮定する。
完全列 $K \to N_2 \to Q \to 0$ に補題
\ref{algebra-lemma-tensor-product-exact} を適用すると、列
$K \otimes_R M \to N_2 \otimes_R M \to
Q \otimes_R M \to 0$ は完全である。したがって、
$\Ker \left(N_2 \otimes_R M \to Q \otimes_R M\right)
= \Im \left(K \otimes_R M \to N_2 \otimes_R M\right)$ である。
また、$N_1 \otimes_R M \to K \otimes_R M$ が全射なので
$\Im \left(K \otimes_R M \to N_2 \otimes_R M\right)
= \Im \left(N_1 \otimes_R M \to N_2 \otimes_R M\right)$ であり、
$Q \otimes_R M \to N_3 \otimes_R M$ が単射なので
$\Ker \left(N_2 \otimes_R M \to Q \otimes_R M\right)
= \Ker \left(N_2 \otimes_R M \to N_3 \otimes_R M\right)$ である。
ゆえにこれは
$\Ker \left(N_2 \otimes_R M \to N_3 \otimes_R M\right)
= \Im \left(N_1 \otimes_R M \to N_2 \otimes_R M\right)$
となる。したがって関手 $- \otimes_R M$ は完全であり、$M$ は平坦である。}。
よって、$- \otimes_R M$ が $R$-加群の単射を $R$-加群の単射に移すことを示せば十分である。

\medskip\noindent
$K \to N$ を $R$-加群の単射とし、
$x \in \Ker(K \otimes_R M \to N \otimes_R M)$ とする。
$x$ が零であることを示さなければならない。
$R$-加群 $K$ は、その有限生成 $R$-部分加群の合併である。
したがって $K \otimes_R M$ は、$K$ のすべての有限生成
$R$-部分加群を $K_i$ が走るときの $R$-加群
$K_i \otimes_R M$ の余極限である
（テンソル積は余極限と可換だからである）。
ゆえに、ある $i$ に対して、$x$ は元
$x_i \in K_i \otimes_R M$ から来る。したがって、$K$ が有限生成
$R$-加群であると仮定してよい。そう仮定する。
単射 $K \to N$ を包含とみなし、$K \subset N$ とする。

\medskip\noindent
$R$-加群 $N$ は、$K$ を含む有限生成 $R$-部分加群の合併である。
したがって $N \otimes_R M$ は、$K$ を含む $N$ のすべての有限生成
$R$-部分加群を $N_i$ が走るときの $R$-加群
$N_i \otimes_R M$ の余極限である
（ここでもテンソル積が余極限と可換であることを用いた）。
これは有向系上の余極限であることに注意せよ
（$N$ の二つの有限生成部分加群の和は再び有限生成だからである）。
したがって、補題 \ref{algebra-lemma-zero-directed-limit} により、元
$x \in K \otimes_R M$ は、これらの $R$-加群の少なくとも一つ
$N_i \otimes_R M$ において零に写る
（$x$ は $N \otimes_R M$ において零に写るからである）。
ゆえに、$N$ は有限生成 $R$-加群であると仮定してよい。

\medskip\noindent
$N$ が有限生成 $R$-加群であると仮定する。ある
$L \subset L' \subset R^{\oplus n}$ に対して
$N = R^{\oplus n}/L$ および $K = L'/L$ と書く。
任意の $R$-部分加群 $G \subset R^{\oplus n}$ に対し、合成
$G \otimes_R M \to R^n \otimes_R M = M^{\oplus n}$ によって得られる
標準写像 $G \otimes_R M \to M^{\oplus n}$ がある。
$L \otimes_R M \to M^{\oplus n}$ および
$L' \otimes_R M \to M^{\oplus n}$ が単射であることを証明すれば十分である。
実際、そうならば、
$K \otimes_R M = L' \otimes_R M/L \otimes_R M \to M^{\oplus n}/L \otimes_R M$
も単射であることが分かる\footnote{$L' \otimes_R M$ および
$L \otimes_R M$ を $M^{\oplus n}$ の部分加群と同一視すれば、これは明らかになる
（写像 $L \otimes_R M \to M^{\oplus n}$ および
$L' \otimes_R M \to M^{\oplus n}$ が単射であり、
% P01-E1297: frozen source reverses the inclusion-induced tensor map; see ERRATA_CANDIDATES.jsonl.
明らかな写像 $L' \otimes_R M \to L \otimes_R M$ と可換なので、この同一視は正当である）。}。

\medskip\noindent
したがって、$L \subset R^{\oplus n}$ が $R$-部分加群であるとき、
$L \otimes_R M \to M^{\oplus n}$ が単射であることを示せば十分である。
$n$ に関する帰納法でこれを示す。基底の場合 $n = 1$ は後で扱う。
帰納段階として $n > 1$ と仮定し、
$L' = L \cap R \oplus 0^{\oplus n - 1}$ とおく。このとき
$L'' = L/L'$ は $R^{\oplus n - 1}$ の部分加群である。図式
$$
\xymatrix{
&
L' \otimes_R M \ar[r] \ar[d] &
L \otimes_R M \ar[r] \ar[d] &
L'' \otimes_R M \ar[r] \ar[d] &
0 \\
0 \ar[r] &
M \ar[r] &
M^{\oplus n} \ar[r] &
M^{\oplus n - 1} \ar[r] & 0
}
$$
を得る。帰納法の仮定および基底の場合により、左と右の垂直矢印は単射である。
各行は完全である。したがって中央の垂直矢印も単射である。

\medskip\noindent
以上の帰納法の基底の場合は、$L \subset R$ がイデアルである場合である。
言い換えると、$R$ の任意のイデアル $I$ に対して
$I \otimes_R M \to M$ が単射であることを示さなければならない。
$I$ が有限生成ならば、これは真であると分かっている。しかし、
$I = \bigcup I_\alpha$ は、その中に含まれる有限生成イデアル
$I_\alpha$ の合併である。言い換えると、$I = \colim I_\alpha$ である。
$\otimes$ は余極限と可換なので、
$I \otimes_R M = \colim I_\alpha \otimes_R M$ である。
仮定によりすべての射 $I_\alpha \otimes_R M \to M$ は単射だから、
$I \otimes_R M \to M$ についても同じことが成り立つ。
\end{proof}

\begin{lemma}
\label{algebra-lemma-colimit-rings-flat}
$\{R_i, \varphi_{ii'}\}$ を有向集合 $I$ 上の環の系とする。
$R = \colim_i R_i$ とおく。
\begin{enumerate}
\item $M$ を $R$-加群とし、すべての $i$ について $M$ が
$R_i$-加群として平坦ならば、$M$ は $R$-加群として平坦である。
\item $i \in I$ に対して $M_i$ を平坦な $R_i$-加群とし、
$i' \geq i$ に対して $f_{ii'} : M_i \to M_{i'}$ を
$f_{i' i''} \circ f_{i i'} = f_{i i''}$ を満たす $\varphi_{ii'}$-線形写像とする。
このとき $M = \colim_{i \in I} M_i$ は平坦な $R$-加群である。
\end{enumerate}
\end{lemma}

\begin{proof}
(1) は、すべての $i$ に対して $M_i = M$、$f_{i i'} = \text{id}_M$ とした
(2) の特別な場合である。(2) を証明する。
$\mathfrak a \subset R$ を有限生成イデアルとする。補題 \ref{algebra-lemma-flat} により、
$\mathfrak a \otimes_R M \to M$ が単射であることを示せば十分である。
$\mathfrak a = \mathfrak a'R$ を満たす $i \in I$ と有限生成イデアル
$\mathfrak a' \subset R_i$ を見いだせる。
このとき $\mathfrak a = \colim_{i' \geq i} \mathfrak a'R_{i'}$ である。
$\otimes$ は余極限と可換なので、写像
$\mathfrak a \otimes_R M \to M$ は写像
$$
\mathfrak a'R_{i'} \otimes_{R_{i'}} M_{i'} \longrightarrow M_{i'}
$$
の余極限である。仮定により、これらの写像はすべて単射である。
補題 \ref{algebra-lemma-directed-colimit-exact} により $I$ 上の余極限は完全なので、
結論を得る。
\end{proof}

\begin{lemma}
\label{algebra-lemma-flat-base-change}
$M$ が $R$ 上（忠実）平坦で、$R \to R'$ が環準同型であると仮定する。
このとき $M \otimes_R R'$ は $R'$ 上（忠実）平坦である。
\end{lemma}

\begin{proof}
任意の $R'$-加群 $N$ に対し、標準同型
$N \otimes_{R'} (R'\otimes_R M) = N \otimes_R M$ がある。
したがって、関手 $-\otimes_{R'}(R'\otimes_R M)$ に望まれる完全性は、
関手 $-\otimes_R M$ の対応する完全性から従う。
\end{proof}

\begin{lemma}
\label{algebra-lemma-flatness-descends}
$R \to R'$ を忠実平坦な環準同型とする。
$M$ を $R$ 上の加群とし、$M' = R' \otimes_R M$ とおく。
$M$ が $R$ 上平坦であるための必要十分条件は、$M'$ が $R'$ 上平坦であることである。
\end{lemma}

\begin{proof}
補題 \ref{algebra-lemma-flat-base-change} により、$M$ が平坦ならば $M'$ も平坦である。
逆に、$M'$ が平坦であると仮定する。
$N_1 \to N_2 \to N_3$ を $R$-加群の完全列とする。
$N_1 \otimes_R M \to N_2 \otimes_R M \to N_3 \otimes_R M$
が完全であることを示したい。$R \to R'$ が平坦なので、
$N_1 \otimes_R R' \to N_2 \otimes_R R' \to N_3 \otimes_R R'$ は完全である。
$M'$ の平坦性により、
$N_1 \otimes_R R' \otimes_{R'} M'
\to N_2 \otimes_R R' \otimes_{R'} M'
\to N_3 \otimes_R R' \otimes_{R'} M'$ は完全である。
これは
$N_1 \otimes_R M \otimes_R R'
\to N_2 \otimes_R M \otimes_R R'
\to N_3 \otimes_R M \otimes_R R'$ と書ける。
最後に、忠実平坦性により
$N_1 \otimes_R M \to N_2 \otimes_R M \to N_3 \otimes_R M$
は完全である。
\end{proof}

\begin{lemma}
\label{algebra-lemma-flatness-descends-more-general}
$R$ を環とする。$S \to S'$ を $R$-代数の平坦な準同型とする。
$M$ を $S$ 上の加群とし、$M' = S' \otimes_S M$ とおく。
\begin{enumerate}
\item $M$ が $R$ 上平坦ならば、$M'$ は $R$ 上平坦である。
\item $S \to S'$ が忠実平坦ならば、$M$ が $R$ 上平坦であるための
必要十分条件は、$M'$ が $R$ 上平坦であることである。
\end{enumerate}
\end{lemma}

\begin{proof}
$N \to N'$ を $R$-加群の単射とする。$S \to S'$ の平坦性により
$$
\Ker(N \otimes_R M \to N' \otimes_R M) \otimes_S S'
=
\Ker(N \otimes_R M' \to N' \otimes_R M')
$$
である。$M$ が $R$ 上平坦ならば左辺は零であり、補題
\ref{algebra-lemma-flat} における平坦性の第二の特徴づけにより、$M'$ が
$R$ 上平坦であることが分かる。$M'$ が $R$ 上平坦ならば右辺は零であり、
さらに $S \to S'$ が忠実平坦ならば、これは
$\Ker(N \otimes_R M \to N' \otimes_R M)$ が零であることを含意する。
したがって $M$ は $R$ 上平坦である。
\end{proof}

\begin{lemma}
\label{algebra-lemma-flat-permanence}
$R \to S$ を環準同型とする。$M$ を $S$-加群とする。
$M$ が $R$-加群として平坦で、$S$-加群として忠実平坦ならば、
$R \to S$ は平坦である。
\end{lemma}

\begin{proof}
$N_1 \to N_2 \to N_3$ を $R$-加群の完全列とする。
仮定により、$N_1 \otimes_R M \to N_2 \otimes_R M \to N_3 \otimes_R M$
は完全である。これは
$$
N_1 \otimes_R S \otimes_S M
\to
N_2 \otimes_R S \otimes_S M
\to
N_3 \otimes_R S \otimes_S M.
$$
と書ける。$S$ 上の $M$ の忠実平坦性により、
$N_1 \otimes_R S \to N_2 \otimes_R S \to N_3 \otimes_R S$ は完全である。
したがって $R \to S$ は平坦である。
\end{proof}

\noindent
$R$ を環とする。
$M$ を $R$-加群とする。
$\sum f_i x_i = 0$ を $M$ における関係式とする。
関係式 $\sum f_i x_i$ が{\it 自明}であるとは、整数 $m \geq 0$、
元 $y_j \in M$, $j = 1, \ldots, m$、および元 $a_{ij} \in R$,
$i = 1, \ldots, n$, $j = 1, \ldots, m$ であって、
$$
x_i = \sum\nolimits_j a_{ij} y_j, \forall i,
\quad\text{and}\quad
0 = \sum\nolimits_i f_ia_{ij}, \forall j.
$$
を満たすものが存在することをいう。

\begin{lemma}[平坦性の等式判定法]
\label{algebra-lemma-flat-eq}
$R$ 上の加群 $M$ が平坦であるための必要十分条件は、
$M$ におけるすべての関係式が自明であることである。
\end{lemma}

\begin{proof}
$M$ が平坦であると仮定し、$\sum f_i x_i = 0$ を $M$ における関係式とする。
$I = (f_1, \ldots, f_n)$ とおき、
$K = \Ker(R^n \to I, (a_1, \ldots, a_n) \mapsto \sum_i a_i f_i)$ とする。
すると短完全列 $0 \to K \to R^n \to I \to 0$ がある。
$\sum f_i \otimes x_i$ は $I \otimes_R M$ の元であり、
$R \otimes_R M = M$ において零に写る。平坦性により、
$\sum f_i \otimes x_i$ は $I \otimes_R M$ において零である。
したがって、$\sum e_i \otimes x_i \in R^n \otimes_R M$ に写る元が
$K \otimes_R M$ に存在する。ここで $e_i$ は $R^n$ の第 $i$ 基底元である。
この元を $\sum k_j \otimes y_j$ と書き、さらに $R^n$ における
$k_j$ の像を $\sum a_{ij} e_i$ と書けば、結論を得る。

\medskip\noindent
すべての関係式が自明であると仮定し、$I$ を有限生成イデアルとし、
$x = \sum f_i \otimes x_i$ を $I \otimes_R M$ の元で
$R \otimes_R M = M$ において零に写るものとする。
これはまさに $\sum f_i x_i$ が $M$ における関係式であることを意味する。
それが自明であることから、$x$ が零であることが容易に従う。実際、
$$
x
=
\sum f_i \otimes x_i
=
\sum f_i \otimes \left(\sum a_{ij}y_j\right)
=
\sum \left(\sum f_i a_{ij}\right) \otimes y_j
=
0
$$
である。
\end{proof}

\begin{lemma}
\label{algebra-lemma-flat-tor-zero}
$R$ を環、$0 \to M'' \to M' \to M \to 0$ を短完全列、
$N$ を $R$-加群とする。$M$ が平坦ならば、
$N \otimes_R M'' \to N \otimes_R M'$ は単射である。すなわち、列
$$
0 \to N \otimes_R M'' \to N \otimes_R M' \to N \otimes_R M \to 0
$$
は短完全列である。
\end{lemma}

\begin{proof}
自由加群から $N$ への全射 $R^{(I)} \to N$ を、その核を $K$ としてとる。
次の図式に蛇の補題を適用すれば、結果が従う。
$$
\begin{matrix}
 & & 0 & & 0 & & 0 & & \\
 & & \uparrow & & \uparrow & & \uparrow & & \\
 & & M''\otimes_R N & \to & M' \otimes_R N & \to & M \otimes_R N & \to & 0 \\
 & & \uparrow & & \uparrow & & \uparrow & & \\
0 & \to & (M'')^{(I)} & \to & (M')^{(I)} & \to & M^{(I)} & \to & 0 \\
 & & \uparrow & & \uparrow & & \uparrow & & \\
 & & M''\otimes_R K & \to & M' \otimes_R K & \to & M \otimes_R K & \to & 0 \\
 & & & & & & \uparrow & & \\
 & & & & & & 0 & &
\end{matrix}
$$
ここで、各行と各列は完全である。自由加群 $R^{(I)}$ とのテンソル積をとることは
完全なので、中央の行は完全である。
\end{proof}


\begin{lemma}
\label{algebra-lemma-flat-ses}
$0 \to M' \to M \to M'' \to 0$ を $R$-加群の短完全列とする。
$M'$ と $M''$ が平坦ならば $M$ も平坦である。
$M$ と $M''$ が平坦ならば $M'$ も平坦である。
\end{lemma}

\begin{proof}
加群 $N$ が平坦であることの判定法として、任意のイデアル $I \subset R$ に対し
写像 $N \otimes_R I \to N$ が単射であることを用いる。
補題 \ref{algebra-lemma-flat} を参照せよ。
イデアル $I \subset R$ を考える。図式
$$
\begin{matrix}
0 & \to & M' & \to & M & \to & M'' & \to & 0 \\
& & \uparrow & & \uparrow & & \uparrow & & \\
& & M'\otimes_R I & \to & M \otimes_R I & \to & M''\otimes_R I & \to & 0
\end{matrix}
$$
を考える。各行は完全である。これにより第一の主張が直ちに示される。
第二の主張は、$M''$ が平坦ならば、補題 \ref{algebra-lemma-flat-tor-zero} により
下段左の水平矢印が単射であることから従う。
\end{proof}

\begin{lemma}
\label{algebra-lemma-easy-ff}
$R$ を環とする。
$M$ を $R$-加群とする。
次は同値である。
\begin{enumerate}
\item $M$ は忠実平坦である。
\item $M$ は平坦であり、すべての $R$-加群準同型 $\alpha : N \to N'$ に対して、
$\alpha = 0$ であるための必要十分条件は $\alpha \otimes \text{id}_M = 0$ である。
\end{enumerate}
\end{lemma}

\begin{proof}
$M$ が忠実平坦ならば、$0 \to \Ker(\alpha) \to N \to N'$ が完全であるための
必要十分条件は、$M$ とテンソル積をとった後にも同じことが成り立つことである。
これで (1) ならば (2) が示される。逆に (2) を仮定する。
$N_1 \to N_2 \to N_3$ を複体とし、複体
$N_1 \otimes_R M \to N_2 \otimes_R M \to N_3\otimes_R M$
が完全であると仮定する。$x \in \Ker(N_2 \to N_3)$ をとり、写像
$\alpha : R \to N_2/\Im(N_1)$, $r \mapsto rx + \Im(N_1)$ を考える。
複体 $-\otimes_R M$ の完全性により、$\alpha \otimes
\text{id}_M$ は零である。仮定により $\alpha$ は零である。
したがって $x $ は $N_1 \to N_2$ の像に属する。
\end{proof}

\begin{lemma}
\label{algebra-lemma-ff}
\begin{slogan}
平坦加群が忠実平坦であるための必要十分条件は、そのファイバーが零でないことである。
\end{slogan}
$M$ を平坦な $R$-加群とする。
次は同値である。
\begin{enumerate}
\item $M$ は忠実平坦である。
% P01-E1298: frozen source says "then tensor product" instead of "the tensor product".
\item 任意の零でない $R$-加群 $N$ に対し、テンソル積 $M \otimes_R N$ は零でない。
\item すべての $\mathfrak p \in \Spec(R)$ に対し、テンソル積
$M \otimes_R \kappa(\mathfrak p)$ は零でない。
\item $R$ のすべての極大イデアル $\mathfrak m$ に対し、テンソル積
$M \otimes_R \kappa(\mathfrak m) = M/{\mathfrak m}M$ は零でない。
\end{enumerate}
\end{lemma}

\begin{proof}
$M$ が忠実平坦で、$N \not = 0$ と仮定する。補題 \ref{algebra-lemma-easy-ff} により、
零でない写像 $1 : N \to N$ は零でない写像
$M \otimes_R N \to M \otimes_R N$ を誘導する。したがって
$M \otimes_R N \not = 0$ である。これで (1) ならば (2) が示された。
(2) $\Rightarrow$ (3) $\Rightarrow$ (4) は直ちに従う。

\medskip\noindent
(4) を仮定する。$N_1 \to N_2 \to N_3$ を複体とし、
$N_1 \otimes_R M \to N_2\otimes_R M \to N_3\otimes_R M$ が完全であると仮定する。
$H$ をこの複体のコホモロジー、すなわち
$H = \Ker(N_2 \to N_3)/\Im(N_1 \to N_2)$ とする。
証明を終えるため、$H = 0$ を示す。平坦性により $H \otimes_R M = 0$ である。
$x \in H$ をとり、その零化イデアルを
$I = \{f \in R \mid fx = 0 \}$ とする。$R/I \subset H$ なので、
$M$ の平坦性により $M/IM \subset H \otimes_R M = 0$ を得る。
$I \not =  R$ ならば、極大イデアル $I \subset \mathfrak m \subset R$ を選べる。
これは直ちに矛盾を与える。
\end{proof}

\begin{lemma}
\label{algebra-lemma-ff-rings}
$R \to S$ を平坦な環準同型とする。
次は同値である。
\begin{enumerate}
\item $R \to S$ は忠実平坦である。
\item 誘導される $\Spec$ 上の写像は全射である。
\item 任意の閉点 $x \in \Spec(R)$ は写像
$\Spec(S) \to \Spec(R)$ の像に属する。
\end{enumerate}
\end{lemma}

\begin{proof}
補題 \ref{algebra-lemma-ff} から直ちに従う。実際、注意
\ref{algebra-remark-fundamental-diagram} で見たように、$\mathfrak p$ が像に属するための
必要十分条件は、環 $S \otimes_R \kappa(\mathfrak p)$ が零でないことである。
\end{proof}

\begin{lemma}
\label{algebra-lemma-local-flat-ff}
局所環の平坦な局所環準同型は忠実平坦である。
\end{lemma}

\begin{proof}
補題 \ref{algebra-lemma-ff-rings} から直ちに従う。
\end{proof}

\noindent
平坦性は局所化とよく整合する。

\begin{lemma}
\label{algebra-lemma-flat-localization}
$R$ を環とする。$S \subset R$ を積閉部分集合とする。
\begin{enumerate}
\item 局所化 $S^{-1}R$ は平坦な $R$-代数である。
\item $M$ を $S^{-1}R$-加群とする。$M$ が平坦な $R$-加群であるための
必要十分条件は、$M$ が平坦な $S^{-1}R$-加群であることである。
\item $M$ を $R$-加群とする。$M$ が平坦な $R$-加群であるための
必要十分条件は、$R$ のすべての素イデアル $\mathfrak p$ に対し
$M_{\mathfrak p}$ が平坦な $R_{\mathfrak p}$-加群であることである。
\item $M$ を $R$-加群とする。$M$ が平坦な $R$-加群であるための
必要十分条件は、$R$ のすべての極大イデアル $\mathfrak m$ に対し
$M_{\mathfrak m}$ が平坦な $R_{\mathfrak m}$-加群であることである。
\item $R \to A$ を環準同型、$M$ を $A$-加群とし、
$g_1, \ldots, g_m \in A$ を $A$ の単位イデアルを生成する元とする。
このとき $M$ が $R$ 上平坦であるための必要十分条件は、各局所化
$M_{g_i}$ が $R$ 上平坦であることである。
\item $R \to A$ を環準同型、$M$ を $A$-加群とする。
$M$ が平坦な $R$-加群であるための必要十分条件は、$A$ のすべての素イデアル
$\mathfrak q$ に対し、局所化 $M_{\mathfrak q}$ が平坦な
$R_{\mathfrak p}$-加群であることである
（$\mathfrak p$ は $\mathfrak q$ の下にある $R$ の素イデアルとする）。
\item $R \to A$ を環準同型、$M$ を $A$-加群とする。
$M$ が平坦な $R$-加群であるための必要十分条件は、$A$ のすべての極大イデアル
$\mathfrak m$ に対し、局所化 $M_{\mathfrak m}$ が平坦な
$R_{\mathfrak p}$-加群であることである
（$\mathfrak p = R \cap \mathfrak m$ とする）。
\end{enumerate}
\end{lemma}

\begin{proof}
補題の最後の主張を証明する。
証明では、局所化が完全でテンソル積と可換であることを繰り返し用いる。
Sections \ref{algebra-section-localization} および \ref{algebra-section-tensor-product} を参照せよ。

\medskip\noindent
$R \to A$ を環準同型、$M$ を $A$-加群とする。
$A$ のすべての極大イデアル $\mathfrak m$ に対し
$M_{\mathfrak m}$ が平坦な $R_{\mathfrak p}$-加群であると仮定する
（$\mathfrak p = R \cap \mathfrak m$ とする）。$I \subset R$ をイデアルとする。
写像 $I \otimes_R M \to M$ が単射であることを示さなければならない。
これを $A$-加群の写像と考えられる。仮定により、局所化
$(I \otimes_R M)_{\mathfrak m} \to M_{\mathfrak m}$ は単射である。なぜなら
$(I \otimes_R M)_{\mathfrak m} =
I_{\mathfrak p} \otimes_{R_{\mathfrak p}} M_{\mathfrak m}$ だからである。
したがって、補題 \ref{algebra-lemma-characterize-zero-local} により、
$I \otimes_R M \to M$ の核は零である。ゆえに $M$ は $R$ 上平坦である。

\medskip\noindent
逆に、$M$ が $R$ 上平坦であると仮定する。$R$ の素イデアル
$\mathfrak p$ の上にある $A$ の素イデアル $\mathfrak q$ をとる。
$I \subset R_{\mathfrak p}$ をイデアルとする。
$I \otimes_{R_{\mathfrak p}} M_{\mathfrak q} \to M_{\mathfrak q}$
が単射であることを示さなければならない。あるイデアル $J \subset R$ により
$I = J_{\mathfrak p}$ と書ける。このとき写像
$I \otimes_{R_{\mathfrak p}} M_{\mathfrak q} \to M_{\mathfrak q}$ は、
単射 $J \otimes_R M \to M$ の（$\mathfrak q$ における）局所化にほかならない。
局所化は完全なので、$M_{\mathfrak q}$ は平坦な
$R_{\mathfrak p}$-加群であることが分かる。

\medskip\noindent
これで (7) と (6) が示された。他の主張は最後の主張から直接に従う
（証明は省略する）。
\end{proof}

\begin{lemma}
\label{algebra-lemma-flat-going-down}
$R \to S$ を平坦とする。$\mathfrak p \subset \mathfrak p'$ を $R$ の素イデアルとする。
$\mathfrak q' \subset S$ を $\mathfrak p'$ に写る $S$ の素イデアルとする。
このとき、$\mathfrak p$ に写る素イデアル
$\mathfrak q \subset \mathfrak q'$ が存在する。
\end{lemma}

\begin{proof}
補題 \ref{algebra-lemma-flat-localization} により、局所環準同型
$R_{\mathfrak p'} \to S_{\mathfrak q'}$ は平坦である。
補題 \ref{algebra-lemma-local-flat-ff} により、この局所環準同型は忠実平坦である。
補題 \ref{algebra-lemma-ff-rings} により、$\mathfrak p R_{\mathfrak p'}$ に写る素イデアルがある。
この素イデアルの $S$ における逆像が所望のものである。
\end{proof}

\noindent
補題に記述された $R \to S$ の性質を「下降性」という。
定義 \ref{algebra-definition-going-up-down} を参照せよ。

\begin{lemma}
\label{algebra-lemma-colimit-faithfully-flat}
$R$ を環とする。$\{S_i, \varphi_{ii'}\}$ を忠実平坦な $R$-代数の有向系とする。
このとき $S = \colim_i S_i$ は忠実平坦な $R$-代数である。
\end{lemma}

\begin{proof}
補題 \ref{algebra-lemma-colimit-flat} により $S$ は平坦である。
$\mathfrak m \subset R$ を極大イデアルとする。補題
\ref{algebra-lemma-ff-rings} により、環 $S_i/\mathfrak m S_i$ はいずれも零環ではない。
したがって $1$ は零に等しくないので、
$S/\mathfrak mS = \colim S_i/\mathfrak mS_i$ も零環ではない。
ゆえに $\Spec(S) \to \Spec(R)$ の像は $\mathfrak m$ を含み、補題
\ref{algebra-lemma-ff-rings} により $R \to S$ は忠実平坦である。
\end{proof}

\section{台と零化イデアル}
\label{algebra-section-supp-and-ann}

\noindent
ごく基本的な定義と補題をいくつか述べる。

\begin{definition}
\label{algebra-definition-support-module}
$R$ を環、$M$ を $R$-加群とする。
$M$ の{\it 台}とは、集合
$$
\text{Supp}(M)
=
\{
\mathfrak p \in \Spec(R)
\mid
M_{\mathfrak p} \not = 0
\}
$$
をいう。
\end{definition}

\begin{lemma}
\label{algebra-lemma-support-zero}
\begin{slogan}
環上の加群の台が空であるための必要十分条件は、その加群が零加群であることである。
\end{slogan}
$R$ を環とする。$M$ を $R$-加群とする。このとき
$$
M = (0) \Leftrightarrow \text{Supp}(M) = \emptyset.
$$
\end{lemma}

\begin{proof}
実際、補題 \ref{algebra-lemma-characterize-zero-local} は、$M$ が零でなければ
$\text{Supp}(M)$ が常に極大イデアルを含むことさえ示している。
\end{proof}

\begin{definition}
\label{algebra-definition-annihilator}
$R$ を環とする。$M$ を $R$-加群とする。
\begin{enumerate}
\item 元 $m \in M$ が与えられたとき、$m$ の{\it 零化イデアル}とは、イデアル
$$
\text{Ann}_R(m) = \text{Ann}(m) = \{f \in R \mid fm = 0\}.
$$
をいう。
\item $M$ の{\it 零化イデアル}とは、イデアル
$$
\text{Ann}_R(M) = \text{Ann}(M) = \{f \in R \mid fm = 0\ \forall m \in M\}.
$$
をいう。
\end{enumerate}
\end{definition}

\begin{lemma}
\label{algebra-lemma-annihilator-flat-base-change}
$R \to S$ を平坦な環準同型とする。$M$ を $R$-加群とし、$m \in M$ とする。
このとき $\text{Ann}_R(m) S = \text{Ann}_S(m \otimes 1)$ である。
$M$ が有限 $R$-加群ならば、
$\text{Ann}_R(M) S = \text{Ann}_S(M \otimes_R S)$ である。
\end{lemma}

\begin{proof}
$I = \text{Ann}_R(m)$ とおく。定義により、$f$ を $fm$ に写す写像
$R \to M$ をもつ完全列 $0 \to I \to R \to M$ がある。
平坦性を用いると完全列
$0 \to I \otimes_R S \to S \to M \otimes_R S$ を得て、第一の主張が示される。
$m_1, \ldots, m_n$ が $M$ の生成元の集合ならば、
$\text{Ann}_R(M) = \bigcap \text{Ann}_R(m_i)$ である。同様に
$\text{Ann}_S(M \otimes_R S) = \bigcap \text{Ann}_S(m_i \otimes 1)$ である。
$I_i = \text{Ann}_R(m_i)$ とおく。このとき、
$\bigcap_{i = 1, \ldots, n} (I_i S) = (\bigcap_{i = 1, \ldots, n} I_i)S$
を示せば十分である。これは補題 \ref{algebra-lemma-flat-intersect-ideals} である。
\end{proof}

\begin{lemma}
\label{algebra-lemma-support-closed}
$R$ を環、$M$ を $R$-加群とする。
$M$ が有限ならば、$\text{Supp}(M)$ は閉である。
より正確には、$I = \text{Ann}(M)$ が $M$ の零化イデアルならば、
$V(I) = \text{Supp}(M)$ である。
\end{lemma}

\begin{proof}
$V(I) = \text{Supp}(M)$ を示す。

\medskip\noindent
$\mathfrak p \in \text{Supp}(M)$ と仮定する。このとき $M_{\mathfrak p} \not = 0$ である。
$M_\mathfrak p$ における像が零でない元 $m \in M$ を選ぶ。
局所化 $M_\mathfrak p$ の構成により、$m$ の零化イデアルは
$\mathfrak p$ に含まれる。したがって、なおさら
$I = \text{Ann}(M)$ は $\mathfrak p$ に含まれなければならない。

\medskip\noindent
逆に、$\mathfrak p \not \in \text{Supp}(M)$ と仮定する。
このとき $M_{\mathfrak p} = 0$ である。
$x_1, \ldots, x_r \in M$ を生成元とする。
補題 \ref{algebra-lemma-localization-colimit} により、
$x_i/1 = 0$ が $M_f$ において成り立つような
$f \in R$, $f\not\in \mathfrak p$ が存在する。
したがって、ある $n_i \geq 1$ に対し $f^{n_i} x_i = 0$ である。
ゆえに、所望のとおり $n = \max\{n_i\}$ に対し $f^nM = 0$ である。
\end{proof}

\begin{lemma}
\label{algebra-lemma-support-base-change}
$R \to R'$ を環準同型、$M$ を有限 $R$-加群とする。
このとき $\text{Supp}(M \otimes_R R')$ は $\text{Supp}(M)$ の逆像である。
\end{lemma}

\begin{proof}
$\mathfrak p \in \text{Supp}(M)$ とする。中山の補題
（補題 \ref{algebra-lemma-NAK}）により、
$$
M \otimes_R \kappa(\mathfrak p) = M_\mathfrak p/\mathfrak p M_\mathfrak p
$$
は零でない $\kappa(\mathfrak p)$-ベクトル空間である。
したがって、$\mathfrak p$ の上にあるすべての素イデアル
$\mathfrak p' \subset R'$ に対し、
$$
(M \otimes_R R')_{\mathfrak p'}/\mathfrak p' (M \otimes_R R')_{\mathfrak p'} =
(M \otimes_R R') \otimes_{R'} \kappa(\mathfrak p') =
M \otimes_R \kappa(\mathfrak p) \otimes_{\kappa(\mathfrak p)}
\kappa(\mathfrak p')
$$
は零でない。これは $\mathfrak p' \in \text{Supp}(M \otimes_R R')$ を含意する。
逆に、$\mathfrak p' \subset R'$ が任意の素イデアル
$\mathfrak p \subset R$ の上にあるならば、
$$
(M \otimes_R R')_{\mathfrak p'} =
M_\mathfrak p \otimes_{R_\mathfrak p} R'_{\mathfrak p'}.
$$
したがって、$\mathfrak p \subset R$ の上にある
$\mathfrak p' \in \text{Supp}(M \otimes_R R')$ に対し、
$\mathfrak p \in \text{Supp}(M)$ である。
\end{proof}

\begin{lemma}
\label{algebra-lemma-support-element}
$R$ を環、$M$ を $R$-加群、$m \in M$ とする。
$\mathfrak p \in V(\text{Ann}(m))$ であるための必要十分条件は、
$m$ が $M_\mathfrak p$ において零に写らないことである。
\end{lemma}

\begin{proof}
$M$ を $Rm \subset M$ に置き換えてよい。このとき
(1) $\text{Ann}(m) = \text{Ann}(M)$ であり、(2) $m$ が
$M_\mathfrak p$ において零に写らないための必要十分条件は
$\mathfrak p \in \text{Supp}(M)$ であることである。
したがって、補題 \ref{algebra-lemma-support-closed} から結果が従う。
\end{proof}

\begin{lemma}
\label{algebra-lemma-support-finite-presentation-constructible}
$R$ を環、$M$ を $R$-加群とする。
$M$ が有限表示 $R$-加群ならば、$\text{Supp}(M)$ は $\Spec(R)$ の閉部分集合で、
その補集合は準コンパクトである。
\end{lemma}

\begin{proof}
表示
$$
R^{\oplus m} \longrightarrow R^{\oplus n} \longrightarrow M \to 0
$$
を選ぶ。第一の写像の行列を $A \in \text{Mat}(n \times m, R)$ とする。
中山の補題 \ref{algebra-lemma-NAK} により、
$$
M_{\mathfrak p} \not = 0 \Leftrightarrow
M \otimes \kappa(\mathfrak p) \not = 0 \Leftrightarrow
\text{rank}(A \bmod \mathfrak p) < n.
$$
である。したがって、$I$ を $A$ の $n \times n$ 小行列式が生成する
$R$ のイデアルとすれば、$\text{Supp}(M) = V(I)$ である。
$I$ は有限生成なので、$I = (f_1, \ldots, f_t)$ と書ける。
すると $\Spec(R) \setminus V(I)$ は標準開集合 $D(f_i)$ の有限合併であり、
したがって準コンパクトである。
\end{proof}

\begin{lemma}
\label{algebra-lemma-support-quotient}
$R$ を環、$M$ を $R$-加群とする。
\begin{enumerate}
\item $M$ が有限ならば、$M/IM$ の台は $\text{Supp}(M) \cap V(I)$ である。
\item $N \subset M$ ならば、$\text{Supp}(N) \subset
\text{Supp}(M)$ である。
\item $Q$ が $M$ の商加群ならば、$\text{Supp}(Q) \subset
\text{Supp}(M)$ である。
\item $0 \to N \to M \to Q \to 0$ が短完全列ならば、
$\text{Supp}(M) = \text{Supp}(Q) \cup \text{Supp}(N)$ である。
\end{enumerate}
\end{lemma}

\begin{proof}
関手 $M \mapsto M_{\mathfrak p}$ は完全である。これにより、第一の主張以外は
直ちに従う。第一の主張については、$M_\mathfrak p \not = 0$ かつ
$I \subset \mathfrak p$ ならば
$(M/IM)_{\mathfrak p} = M_\mathfrak p/IM_\mathfrak p \not = 0$
となることを示す必要がある。これは中山の補題 \ref{algebra-lemma-NAK} から従う。
\end{proof}

\section{上昇と下降}
\label{algebra-section-going-up}

\noindent
$\mathfrak p$, $\mathfrak p'$ を環 $R$ の素イデアルとする。
$X = \Spec(R)$ にザリスキー位相を入れる。
% P01-E1299: frozen source omits "by" in both denote constructions.
$\mathfrak p$ に対応する点を $x \in X$、$\mathfrak p'$ に対応する点を
$x' \in X$ と表す。このとき
$$
x' \leadsto x \Leftrightarrow \mathfrak p' \subset \mathfrak p.
$$
である。言い換えると、$x$ が $x'$ の特殊化であるための必要十分条件は
$\mathfrak p' \subset \mathfrak p$ であることである。
用語と記法については Topology, Section
\ref{topology-section-specialization} を参照せよ。

\begin{definition}
\label{algebra-definition-going-up-down}
$\varphi : R \to S$ を環準同型とする。
\begin{enumerate}
\item $R$ の素イデアル $\mathfrak p \subset \mathfrak p'$ と
$\mathfrak p$ の上にある $S$ の素イデアル $\mathfrak q$ が与えられたとき、
(a) $\mathfrak q \subset \mathfrak q'$ かつ (b) $\mathfrak q'$ が
$\mathfrak p'$ の上にあるような $S$ の素イデアル $\mathfrak q'$ が存在するならば、
$\varphi : R \to S$ は{\it 上昇性}を満たすという。
\item $R$ の素イデアル $\mathfrak p \subset \mathfrak p'$ と
$\mathfrak p'$ の上にある $S$ の素イデアル $\mathfrak q'$ が与えられたとき、
(a) $\mathfrak q \subset \mathfrak q'$ かつ (b) $\mathfrak q$ が
$\mathfrak p$ の上にあるような $S$ の素イデアル $\mathfrak q$ が存在するならば、
$\varphi : R \to S$ は{\it 下降性}を満たすという。
\end{enumerate}
\end{definition}

\noindent
% P01-E1300: frozen source has "have see" instead of "have seen".
これまでに、この性質が成り立つ次の場合を見た。
\begin{enumerate}
\item 整な環準同型は上昇性を満たす。補題 \ref{algebra-lemma-integral-going-up} を参照せよ。
\item 特に、有限な環準同型は上昇性を満たす。
\item 特に、商写像 $R \to R/I$ は上昇性を満たす。
\item 平坦な環準同型は下降性を満たす。補題 \ref{algebra-lemma-flat-going-down} を参照せよ。
\item 特に、任意の局所化は下降性を満たす。
\item 整域の拡大 $R \subset S$ で、$R$ が正規、$S$ が $R$ 上整ならば
下降性を満たす。命題 \ref{algebra-proposition-going-down-normal-integral} を参照せよ。
\end{enumerate}
下降性が成り立つもう一つの場合を挙げる。

\begin{lemma}
\label{algebra-lemma-open-going-down}
$R \to S$ を環準同型とする。誘導される写像
$\varphi : \Spec(S) \to \Spec(R)$ が開ならば、$R \to S$ は下降性を満たす。
\end{lemma}

\begin{proof}
$\mathfrak p \subset \mathfrak p' \subset R$ とし、
$\mathfrak q' \subset S$ が $\mathfrak p'$ の上にあると仮定する。
$\varphi$ は開なので、すべての $g \in S$, $g \not \in \mathfrak q'$ に対し、
$\mathfrak p$ は $D(g) \subset \Spec(S)$ の像に属する。
言い換えると、$S_g \otimes_R \kappa(\mathfrak p)$ は零でない。
$S_{\mathfrak q'}$ はこれらの $S_g$ の有向余極限なので、これは
$S_{\mathfrak q'} \otimes_R \kappa(\mathfrak p)$ が零でないことを含意する。
補題 \ref{algebra-lemma-localization-colimit} および
\ref{algebra-lemma-tensor-products-commute-with-limits} を参照せよ。
したがって、所望のとおり $\mathfrak p$ は
$\Spec(S_{\mathfrak q'}) \to \Spec(R)$ の像に属する。
\end{proof}

\begin{lemma}
\label{algebra-lemma-going-up-down-specialization}
$R \to S$ を環準同型とする。
\begin{enumerate}
\item $R \to S$ が下降性を満たすための必要十分条件は、写像
$\Spec(S) \to \Spec(R)$ に沿って一般化が持ち上がることである。
Topology, Definition \ref{topology-definition-lift-specializations} を参照せよ。
\item $R \to S$ が上昇性を満たすための必要十分条件は、写像
$\Spec(S) \to \Spec(R)$ に沿って特殊化が持ち上がることである。
Topology, Definition \ref{topology-definition-lift-specializations} を参照せよ。
\end{enumerate}
\end{lemma}

\begin{proof}
省略する。
\end{proof}

\begin{lemma}
\label{algebra-lemma-going-up-down-composition}
$R \to S$ と $S \to T$ が下降性を満たす環準同型であると仮定する。
このとき $R \to T$ も下降性を満たす。上昇性についても同様である。
\end{lemma}

\begin{proof}
補題 \ref{algebra-lemma-going-up-down-specialization} により、Topology, Lemma
\ref{topology-lemma-lift-specialization-composition} から従う。
\end{proof}

\begin{lemma}
\label{algebra-lemma-image-stable-specialization-closed}
$R \to S$ を環準同型とする。$T \subset \Spec(R)$ を
$\Spec(S)$ の像とする。$T$ が特殊化について安定ならば、$T$ は閉である。
\end{lemma}

\begin{proof}
二つの証明を与える。

\medskip\noindent
第一の証明。$\mathfrak p \subset R$ を、その $\Spec(R)$ における対応点が
$T$ の閉包に属する素イデアルとする。これは、すべての
$f \in R$, $f \not \in \mathfrak p$ に対し
$D(f) \cap T \not = \emptyset$ となることを意味する。
$D(f) \cap T$ は $\Spec(S_f)$ の $\Spec(R)$ における像であることに注意せよ。
したがって $S_f \not = 0$ と結論する。言い換えると、環 $S_f$ において
$1 \not = 0$ である。$S_{\mathfrak p}$ は環 $S_f$ の有向余極限なので、
$S_{\mathfrak p}$ において $1 \not = 0$ と結論する。
言い換えると $S_{\mathfrak p} \not = 0$ であり、
$\Spec(S_{\mathfrak p}) \to \Spec(S) \to \Spec(R)$ の像を考えると、
$\mathfrak p' \subset \mathfrak p$ となる $\mathfrak p' \in T$ が存在することが分かる。
$T$ は特殊化について閉じていると仮定したので、$\mathfrak p$ が
所望のとおり $T$ の点であると結論する。

\medskip\noindent
第二の証明。$I = \Ker(R \to S)$ とおく。$R$ を $R/I$ に置き換えてよい。
この場合、環準同型 $R \to S$ は単射である。
補題 \ref{algebra-lemma-injective-minimal-primes-in-image} により、$R$ のすべての極小素イデアルは
像 $T$ に含まれる。したがって、$T$ が特殊化について安定ならば、
すべての素イデアルを含む。
\end{proof}

\begin{lemma}
\label{algebra-lemma-going-up-closed}
$R \to S$ を環準同型とする。次は同値である。
\begin{enumerate}
\item $R \to S$ に対して上昇性が成り立つ。
\item 写像 $\Spec(S) \to \Spec(R)$ は閉である。
\end{enumerate}
\end{lemma}

\begin{proof}
位相空間の閉写像に沿って特殊化が持ち上がることは一般的な事実である。
Topology, Lemma \ref{topology-lemma-closed-open-map-specialization} を参照せよ。
したがって第二の条件は第一の条件を含意する。

\medskip\noindent
$R \to S$ に対して上昇性が成り立つと仮定する。
$V(I) \subset \Spec(S)$ を閉集合とする。$V(I)$ の $\Spec(R)$ における像が
閉であることを示したい。環準同型 $S \to S/I$ は明らかに上昇性を満たす。
したがって、補題 \ref{algebra-lemma-going-up-down-composition} により
$R \to S \to S/I$ は上昇性を満たす。
$S$ を $S/I$ に置き換えると、$\Spec(S)$ の $\Spec(R)$ における像 $T$ が
閉であることを示せば十分である。Topology, Lemmas
\ref{topology-lemma-open-closed-specialization} および
\ref{topology-lemma-lift-specializations-images} により、この像は特殊化について安定である。
したがって補題 \ref{algebra-lemma-image-stable-specialization-closed} から結果が従う。
\end{proof}

\begin{lemma}
\label{algebra-lemma-constructible-stable-specialization-closed}
$R$ を環とする。$E \subset \Spec(R)$ を構成可能部分集合とする。
\begin{enumerate}
\item $E$ が特殊化について安定ならば、$E$ は閉である。
\item $E$ が一般化について安定ならば、$E$ は開である。
\end{enumerate}
\end{lemma}

\begin{proof}
第一の証明。第一の主張は、補題 \ref{algebra-lemma-image-stable-specialization-closed} と
補題 \ref{algebra-lemma-constructible-is-image} を組み合わせると従う。
第二の主張は、構成可能集合の補集合が構成可能であること
（Topology, Lemma \ref{topology-lemma-constructible} を参照せよ）、
補題の第一の部分、および Topology, Lemma
\ref{topology-lemma-open-closed-specialization} から従う。

\medskip\noindent
第二の証明。補題 \ref{algebra-lemma-spec-spectral} により $\Spec(R)$ はスペクトル空間なので、
これは Topology, Lemma
\ref{topology-lemma-constructible-stable-specialization-closed} の特別な場合である。
\end{proof}

\begin{proposition}
\label{algebra-proposition-fppf-open}
$R \to S$ を平坦かつ有限表示とする。
このとき $\Spec(S) \to \Spec(R)$ は開である。
より一般に、下降性を満たす任意の有限表示環準同型 $R \to S$ について同じことが成り立つ。
\end{proposition}

\begin{proof}
$R \to S$ が平坦ならば、補題 \ref{algebra-lemma-flat-going-down} により
$R \to S$ は下降性を満たす。したがって、
% P01-E1302: frozen source calls this proposition a lemma.
$R \to S$ が有限表示で下降性を満たすと仮定して証明すればよい。

\medskip\noindent
標準開集合 $D(g) \subset \Spec(S)$, $g \in S$ は位相の基底をなすので、
$D(g)$ の像が開であることを証明すれば十分である。
$\Spec(S_g) \to \Spec(S)$ は $\Spec(S_g)$ から $D(g)$ への同相写像であることを
思い出そう（補題 \ref{algebra-lemma-standard-open}）。
$S \to S_g$ は下降性を満たすので（上を参照）、補題
\ref{algebra-lemma-going-up-down-composition} により $R \to S_g$ は下降性を満たす。
したがって $S$ を $S_g$ に置き換えれば、像が開であることを証明すれば十分である。
Chevalley の定理（定理 \ref{algebra-theorem-chevalley}）により、像は構成可能集合 $E$ である。
$R \to S$ は下降性を満たすので、$E$ は一般化について安定である。
Topology, Lemmas \ref{topology-lemma-open-closed-specialization} および
\ref{topology-lemma-lift-specializations-images} を参照せよ。
したがって、補題 \ref{algebra-lemma-constructible-stable-specialization-closed} により
$E$ は開である。
\end{proof}

\begin{lemma}
\label{algebra-lemma-same-image}
$k$ を体とし、$R$, $S$ を $k$-代数とする。
$S' \subset S$ を $k$-部分代数とし、$f \in S' \otimes_k R$ とする。
可換図式
$$
\xymatrix{
\Spec((S \otimes_k R)_f) \ar[rd] \ar[rr] & &
\Spec((S' \otimes_k R)_f) \ar[ld] \\
& \Spec(R) &
}
$$
において、二つの斜めの矢印の像は同じである。
\end{lemma}

\begin{proof}
$\mathfrak p \subset R$ が南西向きの矢印の像に属するとする。
これは補題 \ref{algebra-lemma-in-image} により
$$
(S' \otimes_k R)_f \otimes_R \kappa(\mathfrak p)
=
(S' \otimes_k \kappa(\mathfrak p))_f
$$
が零環ではないことを意味する。すなわち、
$S' \otimes_k \kappa(\mathfrak p)$ は零環ではなく、その中での $f$ の像は
冪零ではない。環準同型
$S' \otimes_k \kappa(\mathfrak p) \to S \otimes_k \kappa(\mathfrak p)$
は単射である。したがって $S \otimes_k \kappa(\mathfrak p)$ も零環ではなく、
その中での $f$ の像も冪零ではない。
ゆえに $(S \otimes_k R)_f \otimes_R \kappa(\mathfrak p)$ は零環ではない。
したがって補題 \ref{algebra-lemma-in-image} により、所望のとおり
$\mathfrak p$ は南東向きの矢印の像に属する。
\end{proof}

\begin{lemma}
\label{algebra-lemma-map-into-tensor-algebra-open}
$k$ を体とする。
$R$ と $S$ を $k$-代数とする。
写像 $\Spec(S \otimes_k R) \to \Spec(R)$ は開である。
\end{lemma}

\begin{proof}
$f \in S \otimes_k R$ とする。標準開集合 $D(f)$ の像が開であることを
証明すれば十分である。$f \in S' \otimes_k R$ となる有限型
$k$-部分代数 $S' \subset S$ をとる。写像 $R \to S' \otimes_k R$ は
平坦かつ有限表示なので、命題 \ref{algebra-proposition-fppf-open} により
$\Spec((S' \otimes_k R)_f) \to \Spec(R)$ の像 $U$ は開である。
補題 \ref{algebra-lemma-same-image} により、これは $D(f)$ の像でもあり、結論を得る。
\end{proof}

\noindent
時に有用となる、少し技巧的な補題を挙げる。

\begin{lemma}
\label{algebra-lemma-unique-prime-over-localize-below}
$R \to S$ を環準同型とする。
$\mathfrak p \subset R$ を素イデアルとする。
次を仮定する。
\begin{enumerate}
\item $\mathfrak p$ の上にある素イデアル $\mathfrak q \subset S$ がただ一つ存在する。
\item 次のいずれかが成り立つ。
\begin{enumerate}
\item $R \to S$ に対して上昇性が成り立つ。
\item $R \to S$ に対して下降性が成り立ち、$R$ の各素イデアルの上にある
$S$ の素イデアルは高々一つである。
\end{enumerate}
\end{enumerate}
このとき $S_{\mathfrak p} = S_{\mathfrak q}$ である。
\end{lemma}

\begin{proof}
$\Spec(S_{\mathfrak p})$ の点に対応する任意の素イデアル
$\mathfrak q' \subset S$ を考える。これは
$\mathfrak p' = R \cap \mathfrak q'$ が $\mathfrak p$ に含まれることを意味する。
図で表すと
$$
\xymatrix{
\mathfrak q' \ar@{-}[d] \ar@{-}[r] & ? \ar@{-}[r] \ar@{-}[d] & S \ar@{-}[d] \\
\mathfrak p' \ar@{-}[r] & \mathfrak p \ar@{-}[r] & R
}
$$
である。(1) と (2)(a) を仮定する。
上昇性により、$\mathfrak q' \subset \mathfrak q''$ かつ
$\mathfrak q''$ が $\mathfrak p$ の上にあるような素イデアル
$\mathfrak q'' \subset S$ が存在する。$\mathfrak q$ の一意性により、
$\mathfrak q'' = \mathfrak q$ と結論する。言い換えると、
$\mathfrak q'$ は $\Spec(S_{\mathfrak q})$ の点を定める。

\medskip\noindent
(1) と (2)(b) を仮定する。
下降性により、$\mathfrak p'$ の上にある素イデアル
$\mathfrak q'' \subset \mathfrak q$ が存在する。
$\mathfrak p'$ の上にある素イデアルの一意性により
$\mathfrak q' = \mathfrak q''$ と分かる。言い換えると、
$\mathfrak q'$ は $\Spec(S_{\mathfrak q})$ の点を定める。

\medskip\noindent
いずれの場合にも、写像
$\Spec(S_{\mathfrak q}) \to \Spec(S_{\mathfrak p})$ は全単射であると結論する。
これは明らかに $S - \mathfrak q$ のすべての元が
$S_{\mathfrak p}$ において可逆であること、言い換えると
$S_{\mathfrak p} = S_{\mathfrak q}$ であることを意味する。
\end{proof}

\noindent
次の補題は、平坦な環準同型に対する下降性の一般化である。

\begin{lemma}
\label{algebra-lemma-going-down-flat-module}
$R \to S$ を環準同型とする。$N$ を $R$ 上平坦な有限 $S$-加群とする。
$\text{Supp}(N) \subset \Spec(S)$ に誘導位相を入れる。
このとき $\text{Supp}(N) \to \Spec(R)$ に沿って一般化が持ち上がる。
\end{lemma}

\begin{proof}
主張の意味は次のとおりである。$\mathfrak p \subset \mathfrak p' \subset R$ を
素イデアルとする。
% P01-E1301: frozen source omits "with" and terminal punctuation in this hypothesis.
$\mathfrak q' \subset S$ を $\mathfrak q' \in \text{Supp}(N)$ である素イデアルとする。
このとき、$\mathfrak p$ の上にあり、$\mathfrak q \in \text{Supp}(N)$ である
素イデアル $\mathfrak q \subset \mathfrak q'$ が存在する。
$N$ は $R$ 上平坦なので、補題 \ref{algebra-lemma-flat-localization} により
$N_{\mathfrak q'}$ は $R_{\mathfrak p'}$ 上平坦である。
$N_{\mathfrak q'}$ は $S_{\mathfrak q'}$ 上有限であり、
$\mathfrak q' \in \text{Supp}(N)$ なので零でない。したがって中山の補題
\ref{algebra-lemma-NAK} により
$N_{\mathfrak q'} \otimes_{S_{\mathfrak q'}} \kappa(\mathfrak q')$ は零でない。
ゆえに $N_{\mathfrak q'} \otimes_{R_{\mathfrak p'}} \kappa(\mathfrak p')$ も零でない。
補題 \ref{algebra-lemma-ff} により、
$N_{\mathfrak q'} \otimes_{R_{\mathfrak p'}} \kappa(\mathfrak p)$ は零でないと結論する。
$J \subset S_{\mathfrak q'} \otimes_{R_{\mathfrak p'}} \kappa(\mathfrak p)$ を、
零でない有限加群
$N_{\mathfrak q'} \otimes_{R_{\mathfrak p'}} \kappa(\mathfrak p)$ の零化イデアルとする。
$J$ は真のイデアルなので、
$S_{\mathfrak q'} \otimes_{R_{\mathfrak p'}} \kappa(\mathfrak p)/J$
の素イデアルに対応する素イデアル $\mathfrak q \subset S$ を選べる。
この素イデアルは $N$ の台に属し、$\mathfrak p$ の上にあり、
所望のとおり $\mathfrak q'$ に含まれる。
\end{proof}

\section{分離拡大}
\label{algebra-section-separability}

\noindent
この節では、代数的とは限らない体拡大の分離性を扱う。
これは幾何学的被約代数の概念と密接に関係する。定義
\ref{algebra-definition-geometrically-reduced} を参照せよ。

\begin{definition}
\label{algebra-definition-separable-field-extension}
$K/k$ を体拡大とする。
\begin{enumerate}
\item $K/k$ の超越基底 $\{x_i; i \in I\}$ であって、拡大
$K/k(x_i; i \in I)$ が分離代数拡大となるものが存在するとき、
$K$ は $k$ 上{\it 分離生成}であるという。
\item $k \subset K' \subset K$ かつ $K'$ が $k$ 上有限生成である
すべての中間拡大に対して、拡大 $K'/k$ が分離生成であるとき、
$K$ は $k$ 上{\it 分離的}であるという。
\end{enumerate}
\end{definition}

\noindent
この扱いにくい定義からは、分離生成体拡大そのものが分離的であることは明らかでない。
実際にはそうなることが後で分かる。補題
\ref{algebra-lemma-separably-generated-separable} を参照せよ。

\begin{lemma}
\label{algebra-lemma-subextensions-are-separable}
$K/k$ を分離体拡大とする。
任意の中間拡大 $K/K'/k$ に対して、体拡大 $K'/k$ は分離的である。
\end{lemma}

\begin{proof}
定義から直ちに従う。
\end{proof}

\begin{lemma}
\label{algebra-lemma-generating-finitely-generated-separable-field-extensions}
% P01-E1303: frozen source has a spurious comma between coordinated hypotheses.
$K/k$ を分離生成かつ有限生成な体拡大とする。
$r = \text{trdeg}_k(K)$ とおく。このとき、$K$ の元
$x_1, \ldots, x_{r + 1}$ で次を満たすものが存在する。
\begin{enumerate}
\item $x_1, \ldots, x_r$ は $k$ 上の $K$ の超越基底である。
\item $K = k(x_1, \ldots, x_{r + 1})$ である。
\item $x_{r + 1}$ は $k(x_1, \ldots, x_r)$ 上分離的である。
\end{enumerate}
\end{lemma}

\begin{proof}
定義と Fields, Lemma \ref{fields-lemma-primitive-element} を組み合わせればよい。
\end{proof}

\begin{lemma}
\label{algebra-lemma-make-separably-generated}
$K/k$ を有限生成体拡大とする。
図式
$$
\xymatrix{
K \ar[r] & K' \\
k \ar[u] \ar[r] & k' \ar[u]
}
$$
であって、$k'/k$, $K'/K$ が有限純非分離体拡大、かつ
$K'/k'$ が分離生成体拡大となるものが存在する。
\end{lemma}

\begin{proof}
この補題が興味深いのは、$k$ の標数が $p > 0$ の場合だけである。
% P01-E1304: frozen source omits "to be" in the choice construction.
$x_1, \ldots, x_r$ を $k$ 上の $K$ の超越基底として選ぶ。
$K$ は $k$ 上有限生成なので、拡大
$k(x_1, \ldots, x_r) \subset K$ は有限である。
Fields, Lemma \ref{fields-lemma-separable-first} で得られる中間拡大を
$K/K_{sep}/k(x_1, \ldots, x_r)$ とする。
$K = K_{sep}$ ならば証明は終わる。
$d = [K : K_{sep}]$ に関する帰納法を用いる。

\medskip\noindent
$d > 1$ と仮定する。$\alpha = \beta^p \in K_{sep}$ かつ
$\beta \not \in K_{sep}$ となる $\beta \in K$ を選ぶ。
$P = T^n + a_1T^{n - 1} + \ldots + a_n$ を、
$k(x_1, \ldots, x_r)$ 上の $\alpha$ の最小多項式とする。
各 $a_i$ が $k'(x_1^{1/p}, \ldots, x_r^{1/p})$ における
$p$ 乗となるように $p$ 乗根を添加して得られる有限純非分離拡大を
$k'/k$ とする。そのような拡大は存在する。詳細は省略する。
図式
$$
\xymatrix{
K \ar[r] & L \\
k(x_1, \ldots, x_r) \ar[u] \ar[r] & k'(x_1^{1/p}, \ldots, x_r^{1/p}) \ar[u]
}
$$
に入る体を $L$ とする。$L$ は $K$ と
$k'(x_1^{1/p}, \ldots, x_r^{1/p})$ の合成体であると仮定してよく、実際そう仮定する。
Fields, Lemma \ref{fields-lemma-separable-first} で得られる中間拡大を
$L/L_{sep}/k'(x_1^{1/p}, \ldots, x_r^{1/p})$ とする。
このとき $L_{sep}$ は $K_{sep}$ と
$k'(x_1^{1/p}, \ldots, x_r^{1/p})$ の合成体である。
元 $\alpha \in L_{sep}$ は多項式 $P$ の零点であり、その係数はすべて
$k'(x_1^{1/p}, \ldots, x_r^{1/p})$ における $p$ 乗で、その根は互いに異なる。
Fields, Lemma \ref{fields-lemma-pth-root} により、ある
$\alpha' \in L_{sep}$ に対して $\alpha = (\alpha')^p$ であることが分かる。
明らかに、これは $\beta$ が $\alpha' \in L_{sep}$ に写ることを意味する。
言い換えると、体の塔
$$
\xymatrix{
K \ar[r] & L \\
K_{sep}(\beta) \ar[r] \ar[u] & L_{sep} \ar[u] \\
K_{sep} \ar[r] \ar[u] & L_{sep} \ar@{=}[u] \\
k(x_1, \ldots, x_r) \ar[u] \ar[r] & k'(x_1^{1/p}, \ldots, x_r^{1/p}) \ar[u] \\
k \ar[r] \ar[u] & k' \ar[u]
}
$$
を得る。したがって、この構成から
$[L : L_{sep}] < [K : K_{sep}]$ となる新たな状況を得る。
帰納法により、拡大 $L/k'$ に対して補題のとおりとなる
$k' \subset k''$ と $L \subset L'$ を見いだせる。
このとき拡大 $k''/k$ と $L'/K$ が拡大 $K/k$ に対して機能する。
これで補題が示された。
\end{proof}






























\section{幾何学的被約代数}
\label{algebra-section-geometrically-reduced}

\noindent
幾何学的被約代数に関する主結果は補題
\ref{algebra-lemma-geometrically-reduced-finite-purely-inseparable-extension} である。
定義を読んだ後、この補題まで読み飛ばすことを読者に勧める。

\begin{definition}
\label{algebra-definition-geometrically-reduced}
$k$ を体とする。$S$ を $k$-代数とする。
任意の体拡大 $K/k$ に対して $K$-代数 $K \otimes_k S$ が被約であるとき、
$S$ は $k$ 上{\it 幾何学的被約}であるという。
\end{definition}

\noindent
$k$ を体とし、$S$ を被約 $k$-代数とする。
$S$ が幾何学的被約であることを確認するには、
（$\overline{k}$ は $k$ の代数閉包を表す）$\overline{k} \otimes_k S$ が
被約であることを確認すれば十分である。
実際、有限純非分離体拡大 $k'/k$ についてこれを確認するだけで十分である。
補題 \ref{algebra-lemma-geometrically-reduced-finite-purely-inseparable-extension} を参照せよ。

\begin{lemma}
\label{algebra-lemma-subalgebra-separable}
幾何学的被約代数の基本的性質。
$k$ を体とする。$S$ を $k$-代数とする。
\begin{enumerate}
\item $S$ が $k$ 上幾何学的被約ならば、すべての $k$-部分代数もそうである。
\item $S$ のすべての有限生成 $k$-部分代数が幾何学的被約ならば、$S$ は幾何学的被約である。
\item 幾何学的被約 $k$-代数の有向余極限は幾何学的被約である。
\item $S$ が $k$ 上幾何学的被約ならば、$S$ の任意の局所化は $k$ 上幾何学的被約である。
\end{enumerate}
\end{lemma}

\begin{proof}
省略する。第二および第三の性質は、テンソル積が余極限と可換することから従う。
\end{proof}

\begin{lemma}
\label{algebra-lemma-geometrically-reduced-permanence}
$k$ を体とする。
$R$ が $k$ 上幾何学的被約で、$S \subset R$ が乗法的部分集合ならば、局所化
$S^{-1}R$ は $k$ 上幾何学的被約である。
$R$ が $k$ 上幾何学的被約ならば、$R[x]$ は $k$ 上幾何学的被約である。
\end{lemma}

\begin{proof}
省略する。被約環の局所化は被約であり、局所化はテンソル積と可換することがヒントである。
\end{proof}

\noindent
以下の補題の証明では、次の観察を繰り返し用いる。
$R' \subset R$ および $S' \subset S$ が $k$-代数の包含だとする。
このとき写像
$R' \otimes_k S' \to R \otimes_k S$
は単射である。

\begin{lemma}
\label{algebra-lemma-limit-argument}
$k$ を体とする。$R$, $S$ を $k$-代数とする。
\begin{enumerate}
\item $R \otimes_k S$ が被約でないならば、$R' \otimes_k S'$ が被約でないような
有限生成部分代数 $R' \subset R$, $S' \subset S$ が存在する。
\item $R \otimes_k S$ が非零の零因子を含むならば、$R' \otimes_k S'$ が非零の零因子を
含むような有限生成部分代数 $R' \subset R$, $S' \subset S$ が存在する。
\item $R \otimes_k S$ が非自明な冪等元を含むならば、$R' \otimes_k S'$ が非自明な冪等元を
含むような有限生成部分代数 $R' \subset R$, $S' \subset S$ が存在する。
\end{enumerate}
\end{lemma}

\begin{proof}
$z \in R \otimes_k S$ を冪零元とする。$z = \sum_{i = 1, \ldots, n} x_i \otimes y_i$
と書ける。したがって、$x_i$ たちが生成する $k$-部分代数を $R'$ とし、
$y_i$ たちが生成する $k$-部分代数を $S'$ と取ればよい。
第二および第三の主張も同じ方法で証明される。
\end{proof}
\begin{lemma}
\label{algebra-lemma-geometrically-reduced-any-reduced-base-change}
$k$ を体とする。
$S$ を幾何学的被約 $k$-代数とする。
$R$ を任意の被約 $k$-代数とする。
このとき $R \otimes_k S$ は被約である。
\end{lemma}

\begin{proof}
補題 \ref{algebra-lemma-limit-argument} により、$R$ は $k$ 上有限型であると仮定してよい。
すると $R$ は被約ネーター環なので、有限個の体の積に埋め込まれる
（補題 \ref{algebra-lemma-total-ring-fractions-no-embedded-points},
\ref{algebra-lemma-Noetherian-irreducible-components}, および
\ref{algebra-lemma-minimal-prime-reduced-ring} を参照）。
したがって $R$ は体の有限積であると仮定してよい。
この場合、定義 \ref{algebra-definition-geometrically-reduced} から
$R \otimes_k S$ が被約であることが従う。
\end{proof}

\begin{lemma}
\label{algebra-lemma-separable-extension-preserves-reducedness}
$k$ を体とする。
$S$ を被約 $k$-代数とする。
$K/k$ が分離体拡大、または分離生成体拡大のいずれかであるとする。
このとき $K \otimes_k S$ は被約である。
\end{lemma}

\begin{proof}
$k \subset K$ が分離的であると仮定する。
補題 \ref{algebra-lemma-limit-argument} により、$S$ は $k$ 上有限型で、$K$ は $k$ 上有限生成であると
仮定してよい。
すると $S$ は有限個の体の積、すなわちその全商環に埋め込まれる
（補題 \ref{algebra-lemma-minimal-prime-reduced-ring} および
\ref{algebra-lemma-total-ring-fractions-no-embedded-points} を参照）。
よって実際には $S$ が域であると仮定してよい。
補題 \ref{algebra-lemma-generating-finitely-generated-separable-field-extensions} にあるとおり、
$x_1, \ldots, x_{r + 1} \in K$ を選ぶ。
$P \in k(x_1, \ldots, x_r)[T]$ を $x_{r + 1}$ の最小多項式とする。
これは分離多項式である。
容易に
$k[x_1, \ldots, x_r] \otimes_k S = S[x_1, \ldots, x_r]$ が域であることが分かる。
これは $k(x_1, \ldots, x_r) \otimes_k S$ が $S[x_1, \ldots, x_r]$ の局所化なので域であることを
意味する。
環拡大
$k(x_1, \ldots, x_r) \otimes_k S \subset K \otimes_k S$
は単一元 $x_{r + 1}$ と一つの方程式、すなわち $P$ によって生成される。
したがって $K \otimes_k S$ は、$F$ を
$k(x_1, \ldots, x_r) \otimes_k S$ の分数体とすると、$F[T]/(P)$ に埋め込まれる。
$P$ は分離的なので、これは体の有限積であり、証明が終わる。

\medskip\noindent
この時点では、分離生成体拡大が分離的であることはまだ分かっていないので、
この場合にも補題を証明しなければならない。
そのため、$\{x_i\}_{i \in I}$ が $K$ の $k$ 上の分離超越基底だとする。
 $\lambda_j \in K$ の任意の有限集合に対して、ある有限部分集合 $T \subset I$ が存在し、
$k(\{x_i\}_{i\in T}) \subset k(\{x_i\}_{i \in T} \cup \{\lambda_j\})$
が有限分離拡大となる。
したがって $K$ は $k$ の有限生成かつ分離生成な拡大の有向余極限である。
よって、前段落の議論がこの場合にも適用される。
\end{proof}

\begin{lemma}
\label{algebra-lemma-generic-points-geometrically-reduced}
$k$ を体とし、$S$ を $k$-代数とする。
$S$ が被約であり、$S_{\mathfrak p}$ が $S$ のすべての極小素イデアル
$\mathfrak p$ に対して幾何学的被約であると仮定する。
このとき $S$ は幾何学的被約である。
\end{lemma}

\begin{proof}
$S$ は被約なので、写像
$S \to \prod_{\mathfrak p\text{ minimal}} S_{\mathfrak p}$
は単射である。補題 \ref{algebra-lemma-reduced-ring-sub-product-fields} を参照せよ。
$K/k$ が体拡大ならば、写像
$$
S \otimes_k K \to (\prod S_\mathfrak p) \otimes_k K \to
\prod S_\mathfrak p \otimes_k K
$$
は単射である。第一の写像は $k \to K$ が平坦であることによるものであり、
第二の写像は $K$ が自由 $k$-加群であることから直ちに分かる。
$S_\mathfrak p$ は幾何学的被約なので、右辺の環は被約である。
したがって $S \otimes_k K$ は被約環の部分環として被約である。
\end{proof}

\begin{lemma}
\label{algebra-lemma-separable-algebraic-diagonal}
$k'/k$ を分離代数体拡大とする。
乗法写像 $k' \otimes_k k' \to k'$ が
$k' \otimes_k k' \to S^{-1}(k' \otimes_k k')$ と同一視されるような
乗法的部分集合 $S \subset k' \otimes_k k'$ が存在する。
\end{lemma}

\begin{proof}
まず $k'/k$ が有限分離であると仮定する。このとき
$k' = k(\alpha)$ である（Fields, 補題 \ref{fields-lemma-primitive-element}）。
$P \in k[x]$ を $k$ 上の $\alpha$ の最小多項式とする。
すると $P$ は既約かつ分離的なモニック多項式である（Fields, 節
\ref{fields-section-separable-extensions}）。
すると $k'[x]/(P) \to k' \otimes_k k'$、
$\sum \alpha_i x^i \mapsto \alpha_i \otimes \alpha^i$ は同型である。
% P01-E1305: frozen source omits the summation on the displayed image of the polynomial map.
$P = (x - \alpha) Q$ を $k'[x]$ において因数分解できる。
$P$ は分離的なので $Q(\alpha) \not = 0$ である。
したがって、$k'[x]/(P)$ における $Q$ が生成する乗法的集合 $S'$ が機能すること、すなわち
$k' = (S')^{-1}(k'[x]/(P))$ は明らかである。
構造の移送により、$k' \otimes_k k'$ における $S'$ の像 $S$ も機能する。

\medskip\noindent
一般の場合、$k' = \bigcup k_i$ を $k$ 上の有限部分体拡大の合併と書く。
各 $i$ に対して、$S_i \subset k_i \otimes_k k_i$ であって
$k_i = S_i^{-1}(k_i \otimes_k k_i)$ を満たす乗法的部分集合が存在する。
$S$ を $\bigcup S_i \subset k' \otimes_k k'$ が生成する乗法的閉包とする。
% P01-E1306: frozen source has singular-subject/plural-verb agreement error in "multiplication maps sends".
乗法写像は $S$ の各元を $k'$ の可逆元に写すことは明らかである。
$k' \otimes_k k'$ が環 $k_i \otimes_k k_i$ たちの合併であることから、
乗法写像の核の各元は $S^{-1}(k' \otimes_k k')$ で零に写る。
局所化の完全性を用いれば結果が従う。
\end{proof}
\begin{lemma}
\label{algebra-lemma-geometrically-reduced-over-separable-algebraic}
$k'/k$ を分離代数体拡大とする。
$A$ を $k'$ 上の代数とする。このとき、$A$ が $k$ 上幾何学的被約であることと、
$k'$ 上幾何学的被約であることは同値である。
\end{lemma}

\begin{proof}
$A$ が $k'$ 上幾何学的被約であると仮定する。
$K/k$ を体拡大とする。このとき、補題
\ref{algebra-lemma-separable-extension-preserves-reducedness} により
$K \otimes_k k'$ は被約環である。
したがって補題 \ref{algebra-lemma-geometrically-reduced-any-reduced-base-change} により
$K \otimes_k A = (K \otimes_k k') \otimes_{k'} A$ が被約であることが分かる。

\medskip\noindent
$A$ が $k$ 上幾何学的被約であると仮定する。$K/k'$ を体拡大とする。このとき
$$
K \otimes_{k'} A = (K \otimes_k A) \otimes_{(k' \otimes_k k')} k'
$$
である。
$k' \otimes_k k' \to k'$ は補題 \ref{algebra-lemma-separable-algebraic-diagonal} により局所化なので、
$K \otimes_{k'} A$ は被約代数の局所化であり、したがって被約である。
\end{proof}
\section{分離拡大（続き）}
\label{algebra-section-separability-continued}

\noindent
この節では、節 \ref{algebra-section-separability} で始めた議論を続ける。

\begin{lemma}
\label{algebra-lemma-mini-separability}
標数 $p > 1$ の体 $k$ を取る。$K/k$ を、次を満たす
$x_1, \ldots, x_{n + 1} \in K$ によって生成される体拡大とする。
\begin{enumerate}
\item $\{x_1, \ldots, x_n\}$ は $K/k$ の超越基底である。
\item $K$ の任意の $k$-線形独立部分集合 $\{a_1, \ldots, a_m\}$ に対して、集合
 $\{a^p_1, \ldots, a_m^p\}$ は $k$-線形独立である。
\end{enumerate}
このとき、ある $1 \leq j \leq n+1$ が存在し、
$\{ x_1, \ldots, \widehat{x}_j, \ldots, x_{n+1}\}$ は
$K / k$ の分離超越基底となる。
\end{lemma}

\begin{proof}
仮定により $x_{n + 1}$ は $k(x_1, \ldots, x_n)$ 上代数的なので、
$F(x_1, \ldots, x_{n+1}) = 0$ を満たす非零多項式
$F \in k[X_1, \ldots, X_{n + 1}]$ が存在する。$F$ は全次数が最小となるように選ぶ。
すると $F$ は既約である。なぜなら、既約因子の少なくとも一つが同じ性質をもつからである。

\medskip\noindent
ある $i$ について、$F$ に現れる $X_i$ のすべての冪が $p$ の倍数ではないことを主張する。
矛盾を導くため、$F$ に現れるすべての指数が $p$ の倍数だと仮定する。このとき、
$$
\{x_1^{\alpha_1} \ldots x^{\alpha_{n+1}}_{n+1} \mid \lambda_\alpha \neq 0\}
\subset
K
$$
は $k$-線形従属である。ここで $\lambda_\alpha$ は $F$ の係数である。
仮定 (2) により、
$$
\{x_1^{\alpha_1 / p} \ldots x^{\alpha_{n+1} / p}_{n+1}
\mid \lambda_\alpha \neq 0 \}
$$
も $k$-線形従属であると結論でき、これは $\deg(F)$ の最小性に反する。

\medskip\noindent
$F$ に $X_i$ の $p$ 乗でない冪が現れるような $i$ を選ぶ。
すると $x_i$ は
$L = k(x_1, \ldots, x_{i - 1}, x_{i + 1}, \ldots, x_{n+1})$
上代数的であることが分かる。
Fields, 補題 \ref{fields-lemma-transcendence-degree} により、
$x_1, \ldots, x_{i - 1}, x_{i + 1}, \ldots, x_{n+1}$ は
$K/k$ の超越基底である。したがって $L$ は、
$x_1, \ldots, x_{i - 1}, x_{i + 1}, \ldots, x_{n + 1}$ における
$k$ 上の多項式環の分数体である。
Gauss の補題により、
$$
P(T) =
F(x_1, \ldots, x_{i - 1}, T, x_{i + 1}, \ldots, x_{n + 1}) \in L[T]
$$
は既約であると結論する。構成により $P(T)$ は $L[T^p]$ に含まれない。
したがって必要なとおり $K/L$ は分離的である。
\end{proof}
\noindent
$p$ を素数とし、$k$ を標数 $p$ の体とする。
このとき、$k$ のすべての元の $p$ 乗根を $k$ に添加して得られる拡大について
$k^{1/p}$ と書く。
（言い換えれば、これは $k$ の代数閉包の中で、$k$ の元の $p$ 乗根によって生成される $k$-部分体である。）

\begin{lemma}
\label{algebra-lemma-characterize-separable-field-extensions}
標数 $p > 0$ の体 $k$ を取る。
$K/k$ を体拡大とする。
次は同値である。
\begin{enumerate}
\item $K$ は $k$ 上分離的である。
\item $K$ の任意の $k$-線形独立部分集合 $\{a_1, \ldots, a_m\}$ に対して、
集合 $\{a^p_1, \ldots, a_m^p\}$ は $k$-線形独立である。
\item 環 $K \otimes_k k^{1/p}$ は被約であり、
\item $K$ は $k$ 上幾何学的被約である。
\end{enumerate}
\end{lemma}

\begin{proof}
(1) $\Rightarrow$ (4) は補題
\ref{algebra-lemma-separable-extension-preserves-reducedness} による。
(4) $\Rightarrow$ (3) は明らかである。

\medskip\noindent
(3) を仮定する。写像
$m : K \otimes_k k^{1/p} \rightarrow K$ で定まる環準同型を考える。
$$
\lambda \otimes \mu \rightarrow \lambda^p \mu^p
$$
この環のすべての $x \in K \otimes_k k^{1/p}$ に対して
$x^p = m(x) \otimes 1$ であることに注意する。
$K \otimes_k k^{1/p}$ は被約なので、$m$ は単射であることが分かる。
$\{a_1, \ldots, a_m\} \subset K$ が $k$-線形独立ならば、
$\{a_1 \otimes 1, \ldots, a_m \otimes 1\}$ は $k^{1/p}$-線形独立である。
$m$ の単射性から、$a_1^p, \ldots, a_m^p$ の非自明な $k$-線形結合は
零ではないと結論する。したがって (3) は (2) を含意する。

\medskip\noindent
(2) を仮定する。(1) を証明するには、$K$ が $k$ 上有限生成であると仮定してよく、
$K$ が $k$ 上分離生成であることを示さなければならない。
$\{x_1, \ldots, x_d\}$ を $K/k$ の超越基底とする。
Fields, 補題 \ref{fields-lemma-algebraic-finitely-generated} により、
$K' = k(x_1, \ldots, x_d)$ とおくと $[K : K'] < \infty$ である。
非分離次数 $[K : K']_i$ が最小になるように超越基底を選ぶ。
$K / K'$ が分離的ならば証明は終わる。
そうでないと仮定して矛盾を導く。すると、$K'$ 上分離的でない
$x_{d + 1} \in K$ が存在し、特に $[K'(x_{d+1}) : K']_i > 1$ である。
補題 \ref{algebra-lemma-mini-separability} により、
% P01-E1307: frozen source uses n+1 in this bound although the active generators are indexed through d+1.
ある $1 \leq j \leq n + 1$ が存在し、$K'' = k(x_1, \ldots, \widehat{x}_j, \ldots, x_{d+1})$
は $[K'(x_{d+1}) : K'']_i = 1$ を満たす。
次数の乗法性により $[K : K'']_i < [K : K']_i$ となり、矛盾を得る。
\end{proof}
\begin{lemma}
\label{algebra-lemma-separably-generated-separable}
分離生成体拡大は分離的である。
\end{lemma}

\begin{proof}
補題 \ref{algebra-lemma-separable-extension-preserves-reducedness} と補題
\ref{algebra-lemma-characterize-separable-field-extensions} を組み合わせよ。
\end{proof}

\noindent
次の補題では、定義 \ref{algebra-definition-perfection} で定義される完全閉包の概念を用いる。

\begin{lemma}
\label{algebra-lemma-geometrically-reduced-finite-purely-inseparable-extension}
$k$ を体とする。$S$ を $k$-代数とする。
次は同値である。
\begin{enumerate}
\item $k$ の任意の有限純非分離拡大 $k'$ に対して、$k' \otimes_k S$ は被約である。
\item $k^{1/p} \otimes_k S$ は被約である。
\item $k^{perf} \otimes_k S$ は被約である。ここで $k^{perf}$ は $k$ の完全閉包である。
\item $\overline{k} \otimes_k S$ は被約である。ここで $\overline{k}$ は $k$ の代数閉包であり、
\item $S$ は $k$ 上幾何学的被約である。
\end{enumerate}
\end{lemma}

\begin{proof}
任意の有限純非分離拡大 $k'/k$ は $k^{perf}$ に埋め込まれることに注意する。
さらに $k^{1/p}$ は $k^{perf}$ に埋め込まれ、それは $\overline{k}$ に埋め込まれる。
したがって
(5) $\Rightarrow$ (4) $\Rightarrow$ (3) $\Rightarrow$ (2) および
(3) $\Rightarrow$ (1) は明らかである。

\medskip\noindent
(1) $\Rightarrow$ (5) を証明する。
$k$ の有限純非分離拡大 $k'$ すべてに対して $k' \otimes_k S$ が被約であると仮定する。
$K/k$ を体の拡大とする。$K \otimes_k S$ が被約であることを示さなければならない。
補題 \ref{algebra-lemma-limit-argument} により、$K/k$ が有限生成体拡大である場合に帰着する。
補題 \ref{algebra-lemma-make-separably-generated} にある図式
$$
\xymatrix{
K \ar[r] & K' \\
k \ar[u] \ar[r] & k' \ar[u]
}
$$
を選ぶ。
仮定により $k' \otimes_k S$ は被約である。
補題 \ref{algebra-lemma-separable-extension-preserves-reducedness} により
$K' \otimes_k S$ も被約である。したがって、求めるとおり
$K \otimes_k S$ は被約であると結論する。

\medskip\noindent
最後に (2) $\Rightarrow$ (5) を証明する。
$k^{1/p} \otimes_k S$ が被約であると仮定する。このとき $S$ は被約である。
さらに、極小素イデアル $\mathfrak p$ における各局所化 $S_{\mathfrak p}$ について、環
$k^{1/p}\otimes_k S_{\mathfrak p}$ は $k^{1/p} \otimes_k S$ の局所化なので被約である。
しかし補題 \ref{algebra-lemma-minimal-prime-reduced-ring} により $S_{\mathfrak p}$ は体であり、
したがって補題 \ref{algebra-lemma-characterize-separable-field-extensions} により
$S_{\mathfrak p}$ は幾何学的被約である。
補題 \ref{algebra-lemma-generic-points-geometrically-reduced} から、$S$ が幾何学的被約であることが従う。
\end{proof}

\section{完全体}
\label{algebra-section-perfect-fields}

\noindent
ここに定義を与える。

\begin{definition}
\label{algebra-definition-perfect}
体 $k$ が {\it 完全} であるとは、$k$ のすべての体拡大が
$k$ 上分離的であることをいう。
\end{definition}

\begin{lemma}
\label{algebra-lemma-perfect}
体 $k$ が完全であることと、$k$ が標数 $0$ の体であるか、または
標数 $p > 0$ で各元が $p$ 乗根をもつ体であることは同値である。
\end{lemma}

\begin{proof}
標数 0 の場合は明らかである。
$k$ の標数が $p > 0$ であると仮定する。
$k$ が完全なら、$k$ の元の $p$ 乗根を添加するすべての体拡大は
自明でなければならない。したがって $k$ のすべての元は
$p$ 乗根をもつ。逆に、すべての元が $p$ 乗根をもつならば、
$k = k^{1/p}$ であり、補題
\ref{algebra-lemma-characterize-separable-field-extensions} により
$k$ のすべての体拡大は分離的である。
\end{proof}

\begin{lemma}
\label{algebra-lemma-make-separable}
$K/k$ を有限生成体拡大とする。次の図式が存在する。
$$
\xymatrix{
K \ar[r] & K' \\
k \ar[u] \ar[r] & k' \ar[u]
}
$$
ここで $k'/k$, $K'/K$ は有限純非分離体拡大であり、$K'/k'$ は
分離体拡大である。この状況では $K' = k'K$ を合成体と仮定でき、
さらに $K' = (k' \otimes_k K)_{red}$ と仮定できる。
\end{lemma}

\begin{proof}
補題 \ref{algebra-lemma-make-separably-generated} により、
$K'/k'$ が分離生成となるような図式を見いだせる。補題
\ref{algebra-lemma-separably-generated-separable} により、これは $K'$ が
$k'$ 上分離的であることを意味する。
合成体 $k'K$ は $K'/k'$ の部分拡大であり、したがって補題
\ref{algebra-lemma-subextensions-are-separable} により $k' \subset k'K$ は
分離的である。環 $(k' \otimes_k K)_{red}$ は領域である。
ある $n \gg 0$ に対し写像 $x \mapsto x^{p^n}$ がこれを $K$ に写すからである。
したがって補題 \ref{algebra-lemma-integral-over-field} により体である。
よって $(k' \otimes_k K)_{red} \to K'$ はこれを $k'K$ へ同型的に写す。
\end{proof}

\begin{lemma}
\label{algebra-lemma-perfection}
\begin{slogan}
すべての体は一意な完全閉包をもつ。
\end{slogan}
任意の体 $k$ に対して、$k'/k$ が純非分離拡大で $k'$ が完全となるものが
存在する。体拡大 $k'/k$ は一意同型を除いて一意である。
\end{lemma}

\begin{proof}
$k$ の標数が 0 なら、$k' = k$ が唯一の選択である。
$k$ の標数が $p > 0$ であると仮定する。
任意の $n > 0$ に対して、次の (a), (b) を満たす一意な代数拡大
$k \subset k^{1/p^n}$ が存在する。(a) の条件は、任意の元
$\lambda \in k$ が $k^{1/p^n}$ 内に $p^n$ 乗根をもつことであり、
(b) の条件は、任意の元 $\mu \in k^{1/p^n}$ に対して
$\mu^{p^n} \in k$ となることである。
具体的には、環準同型 $k \to k^{1/p^n} = k$,
$x \mapsto x^{p^n}$ を考える。これは単射であり、(a), (b) を満たす。
$k$ の拡大として、写像 $y \mapsto y^p$ を介して
$k^{1/p^n} \subset k^{1/p^{n + 1}}$ であることは明らかである。
したがって $k' = \bigcup k^{1/p^n}$ と取れる。詳細は省略する。
\end{proof}

\begin{definition}
\label{algebra-definition-perfection}
体 $k$ に対し、補題 \ref{algebra-lemma-perfection} の体拡大 $k'/k$ を
$k$ の {\it 完全閉包} と呼ぶ。記法は $k^{perf}/k$ とする。
\end{definition}

\noindent
$k'/k$ が任意の代数的純非分離拡大なら、$k'$ は $k^{perf}$ の部分拡大、
すなわち $k^{perf}/k'/k$ であることに注意する。実際、補題
\ref{algebra-lemma-perfection} の一意性により $(k')^{perf}$ は $k^{perf}$ と同型である。

\begin{lemma}
\label{algebra-lemma-perfect-reduced}
完全体 $k$ を取る。任意の被約 $k$-代数は $k$ 上幾何学的被約である。
$k$-代数 $R$, $S$ を取る。$R$ と $S$ がともに被約であると仮定する。
このとき $k$-代数 $R \otimes_k S$ は被約である。
\end{lemma}

\begin{proof}
最初の主張は補題
\ref{algebra-lemma-geometrically-reduced-finite-purely-inseparable-extension} から従う。
第二の主張には、最初の主張と補題
\ref{algebra-lemma-geometrically-reduced-any-reduced-base-change} を用いる。
\end{proof}


\section{普遍同相写像}
\label{algebra-section-universal-homeomorphism}

\noindent
$k'/k$ を代数的純非分離体拡大とする。このとき任意の $k$-代数
$R$ に対して、環準同型 $R \to k' \otimes_k R$ はスペクトルの同相写像を
誘導する。その理由は、以下のより一般的な補題
\ref{algebra-lemma-p-ring-map} にある。

\begin{lemma}
\label{algebra-lemma-surjective-locally-nilpotent-kernel}
$\varphi : R \to S$ を、核が局所冪零である全射とする。このとき
$\varphi$ はスペクトルの同相写像と剰余体の同型を誘導する。
任意の環準同型 $R \to R'$ に対して、環準同型
$R' \to R' \otimes_R S$ も、核が局所冪零である全射である。
\end{lemma}

\begin{proof}
補題 \ref{algebra-lemma-spec-closed} により、写像
$\Spec(S) \to \Spec(R)$ は閉部分集合 $V(\Ker(\varphi))$ 上への
同相写像である。もちろん、$R$ のすべての素イデアルは $R$ の
すべての冪零元を含むので、$V(\Ker(\varphi)) = \Spec(R)$ である。
これにより剰余体についての主張も従う。テンソル積の右完全性により、
$\Ker(\varphi)R'$ は全射
$R' \to R' \otimes_R S$ の核であることが分かる。
したがって最後の主張は補題 \ref{algebra-lemma-locally-nilpotent} による。
\end{proof}

\begin{lemma}
\label{algebra-lemma-powers-field}
\begin{reference}
\cite[Lemma 3.1.6]{Alper-adequate}
\end{reference}
$k'/k$ を体拡大とする。次は同値である。
\begin{enumerate}
\item 各 $x \in k'$ に対して、ある $n > 0$ が存在して $x^n \in k$ となり、
\item $k' = k$、または $k$ と $k'$ の標数が $p > 0$ であり、$k'/k$ が純非分離体拡大であるか、または
$k$ と $k'$ が $\mathbf{F}_p$ の代数拡大である。
\end{enumerate}
\end{lemma}

\begin{proof}
(2) に列挙した各可能性が (1) を満たすことに注意する。
そこで、$k'/k$ が (1) を満たすと仮定し、(2) の場合のいずれかに
該当することを示す。$k = k'$ の場合を除外すれば、$k' \not = k$ と
仮定してよい。$k'/k$ が代数的であることは明らかである。
したがって、$k'/k$ が自明でない有限拡大であると仮定してよい。
Fields, 補題 \ref{fields-lemma-separable-first} で得られる分離部分拡大を
$k'/k'_{sep}/k$ と書く。
$k = k'_{sep}$ または $k$ が $\mathbf{F}_p$ 上代数的であることを
示さなければならない。したがって、$k'/k$ が自明でない有限分離拡大であり、
$k$ が $\mathbf{F}_p$ 上代数的であることを示すと仮定してよい。

\medskip\noindent
$x \in k'$ かつ $x \not \in k$ となる元を取る。$x^n \in k$ かつ
$(x + 1)^m \in k$ となる $n, m > 0$ を取る。$\overline{k}$ を $k$ の
代数閉包とする。$\sigma(x) \not = \tau(x)$ を満たす埋め込み
$\sigma, \tau : k' \to \overline{k}$ を選ぶことができる。
これは Fields, 節 \ref{fields-section-separable-extensions} の議論から従う。
より正確には、$k'$ を $x$ によって生成される $k$-拡大で置き換えれば、
Fields, 補題 \ref{fields-lemma-count-embeddings} から従う。
すると、$\overline{k}$ のある $n$ 乗根 $\zeta$ に対して
$\sigma(x) = \zeta \tau(x)$ であることが分かる。
同様に、ある $m$ 乗根 $\zeta' \in \overline{k}$ に対して
$\sigma(x + 1) = \zeta' \tau(x + 1)$ であることが分かる。
$\sigma(x + 1) \not = \tau(x + 1)$ なので $\zeta' \not = 1$ である。
すると
$$
\zeta' (\tau(x) + 1) =
\zeta' \tau(x + 1) =
\sigma(x + 1) =
\sigma(x) + 1 =
\zeta \tau(x) + 1
$$
となるので、
$$
\tau(x) (\zeta' - \zeta) = 1 - \zeta'
$$
を得る。したがって $\zeta' \not = \zeta$ であり、
$$
\tau(x) = (1 - \zeta')/(\zeta' - \zeta)
$$
である。
したがって、$k'$ の $k$ に属さないすべての元は素体上代数的である。
$k'$ は、素体上、$k'$ の $k$ に属さない元たちによって生成されるので、
$k'$（したがって $k$）は素体上代数的であると結論する。

\medskip\noindent
最後に、$k$ の標数が $0$ なら、上の議論は次のように矛盾を導く
（読者自身の証明を見つけることを勧める）。
各有理数 $y$ に対して同様に、ある 1 の冪根 $\zeta_y$ が存在して
$\sigma(x + y) = \zeta_y\tau(x + y)$ となる。
すると
$$
\zeta \tau(x) + y = \zeta_y(\tau(x) + y)
$$
を得る。上の $\tau(x)$ の公式により
$\zeta_y \in \mathbf{Q}(\zeta, \zeta')$ と結論できる。
数体 $\mathbf{Q}(\zeta, \zeta')$ は有限個の 1 の冪根しか含まないので、
$\zeta_y = \zeta_{y'}$ となる相異なる有理数 $y, y'$ が存在する。
したがって
$$
y - y' =
\sigma(x + y) - \sigma(x + y') =
\zeta_y(\tau(x + y)) - \zeta_{y'}\tau(x + y') = \zeta_y(y - y')
$$
となる。これは $\zeta_y = 1$ を意味し、矛盾である。
\end{proof}

\begin{lemma}
\label{algebra-lemma-powers}
環準同型 $\varphi : R \to S$ が次を満たすとする。
\begin{enumerate}
\item 任意の $x \in S$ に対して、ある $n > 0$ が存在し、$x^n$ は
$\varphi$ の像に属し、
\item $\Ker(\varphi)$ は局所冪零である。
\end{enumerate}
このとき $\varphi$ はスペクトル上の同相写像を誘導し、さらに剰余体の拡大は
補題 \ref{algebra-lemma-powers-field} の同値な条件を満たす。
\end{lemma}

\begin{proof}
(1) と (2) を仮定する。$R$ の同じ素イデアル $\mathfrak p$ の上にある
$S$ の素イデアル $\mathfrak q, \mathfrak q'$ を取る。
$x \in S$ が $x \in \mathfrak q$ かつ $x \not \in \mathfrak q'$ を満たすと仮定する。
このとき、すべての $n > 0$ に対して $x^n \in \mathfrak q$ かつ
$x^n \not \in \mathfrak q'$ である。ある $n > 0$ に対して
$x^n = \varphi(y)$ かつ $y \in R$ ならば、
$$
x^n \in \mathfrak q \Rightarrow y \in \mathfrak p \Rightarrow
x^n \in \mathfrak q'
$$
となり、矛盾である。したがって上のような $x$ は存在せず、
$\mathfrak q = \mathfrak q'$、すなわちスペクトル上の写像は単射である。
仮定 (2) により核 $I = \Ker(\varphi)$ はすべての素イデアルに含まれるので、
位相空間として $\Spec(R) = \Spec(R/I)$ である。
仮定 (1) により誘導される写像 $R/I \to S$ は整である。
補題 \ref{algebra-lemma-integral-overring-surjective} により
$\Spec(S) \to \Spec(R/I)$ は全射である。以上を合わせると、
$\Spec(S) \to \Spec(R)$ は全単射である。
任意の $x \in S$ を取り、ある $n > 0$ に対して
$\varphi(y) = x^n$ となる $y \in R$ を選ぶ。このとき開集合
$D(x) \subset \Spec(S)$ は、上の全単射を介して開集合
$D(y) \subset \Spec(R)$ に対応する。したがって
$\Spec(S) \to \Spec(R)$ は同相写像である。

\medskip\noindent
剰余体についての主張を見るため、素イデアル $\mathfrak q \subset S$ を取り、
これが素イデアル $\mathfrak p \subset R$ の上にあるとする。$x \in \kappa(\mathfrak q)$ とする。
$\kappa(\mathfrak q)$ を局所環 $S_\mathfrak q$ の剰余体と考えると、
$x$ は $y/z \in S_\mathfrak q$ の像であることが分かる。
ここで $y \in S$, $z \in S$, $z \not \in \mathfrak q$ である。
$y^n, z^m$ が $\varphi$ の像に属するような $n, m > 0$ を選ぶ。
このとき $x^{nm}$ は
$(y/z)^{nm} = (y^n)^m/(z^m)^n$ の剰余であり、
$R_\mathfrak p \to S_\mathfrak q$ の像に属する。
したがって $x^{nm}$ は
$\kappa(\mathfrak p) \to \kappa(\mathfrak q)$ の像に属する。
\end{proof}

\begin{lemma}
\label{algebra-lemma-2-3-ring-map}
$\varphi : R \to S$ を環準同型とする。次を仮定する。
\begin{enumerate}
\item[(a)] $S$ は $R$-代数として、$x^2, x^3 \in \varphi(R)$ を満たす
元 $x$ によって生成される。
\item[(b)] $\Ker(\varphi)$ は局所冪零である。
\end{enumerate}
このとき $\varphi$ は剰余体上の同型とスペクトルの同相写像を誘導する。
任意の環準同型 $R \to R'$ に対して、環準同型
$R' \to R' \otimes_R S$ も (a) と (b) を満たす。
\end{lemma}

\begin{proof}
(a) と (b) を仮定する。$S$ は $R$ 上整であるから、補題
\ref{algebra-lemma-going-up-closed} および
\ref{algebra-lemma-integral-going-up} により、スペクトル上の写像は閉写像である。
補題 \ref{algebra-lemma-image-dense-generic-points} により像は稠密である。
したがって $\Spec(S) \to \Spec(R)$ は全射である。
$\mathfrak p \subset R$ の上にある素イデアル $\mathfrak q \subset S$ を取ると、
体拡大 $\kappa(\mathfrak q)/\kappa(\mathfrak p)$ は、平方と立方が
$\kappa(\mathfrak p)$ に属する元 $\alpha \in \kappa(\mathfrak q)$
によって生成される。したがって明らかに $\alpha \in \kappa(\mathfrak p)$ であり、
$\kappa(\mathfrak q) = \kappa(\mathfrak p)$ を得る。
$\mathfrak p$ の上に相異なる二つの素イデアル $\mathfrak q, \mathfrak q'$ が
存在したとする。このとき (a) における $S$ の生成元 $x$ の少なくとも一つは、
$\kappa(\mathfrak q) = \kappa(\mathfrak p)$ と
$\kappa(\mathfrak q') = \kappa(\mathfrak p)$ において異なる像を持つはずである。
これは $x^2$ と $x^3$ がいずれも同じ像を持つことに矛盾する。
よって $\Spec(S) \to \Spec(R)$ は単射であり、すでに示したことと合わせて
同相写像である。

\medskip\noindent
$\varphi$ はスペクトル上の同相写像を誘導するので、特にスペクトル上全射である。
これは任意の基底変換で保たれる性質である。補題
\ref{algebra-lemma-surjective-spec-radical-ideal} を参照せよ。
したがって任意の $R \to R'$ に対し、環準同型
$R' \to R' \otimes_R S$ の核は冪零元からなる。補題
\ref{algebra-lemma-image-dense-generic-points} を参照せよ。
言い換えれば、$R' \to R' \otimes_R S$ に対して (b) が成り立つ。
(a) が基底変換で保たれることは明らかである。
\end{proof}

\begin{lemma}
\label{algebra-lemma-help-with-powers}
$p$ を素数とする。$n, m > 0$ を二つの整数とする。このとき
$(x + y)^{p^a}, p^a(x + y) \in \mathbf{Z}[x^{p^n}, p^nx, y^{p^m}, p^my]$
を満たす整数 $a$ が存在する。
\end{lemma}

\begin{proof}
$a \geq n, m$ ならば $p^a(x + y)$ については明らかである。
実際、$a \gg n, m$ を取る。次のように書く。
$$
(x + y)^{p^a}  = \sum\nolimits_{i, j \geq 0, i + j = p^a}
{p^a \choose i, j} x^iy^j
$$
$i + j = p^a$ を満たす各 $i, j \geq 0$ に対して、
$r \in \{0, \ldots, p^n - 1\}$ を用いて $i = q p^n + r$ と書き、
$r' \in \{0, \ldots, p^m - 1\}$ を用いて $j = q' p^m + r'$ と書く。
条件 $(x + y)^{p^a} \in \mathbf{Z}[x^{p^n}, p^nx, y^{p^m}, p^my]$ は、
$$
p^{nr + mr'} \text{ divides } {p^a \choose i, j}
$$
ならば成り立つ。$r = r' = 0$ ならば、この整除性は成り立つ。
$r \not = 0$ ならば、次のように書く。
$$
{p^a \choose i, j} = \frac{p^a}{i} {p^a - 1 \choose i - 1, j}
$$
$r \not = 0$ であるから、有理数 $p^a/i$ の $p$ 進付値は少なくとも
$a - (n - 1)$ である（$i$ は $p^n$ で割り切れないためである）。
したがってこの場合、${p^a \choose i, j}$ は $p^{a - n + 1}$ で割り切れる。
同様に、$r' \not = 0$ ならば ${p^a \choose i, j}$ は
$p^{a - m + 1}$ で割り切れる。$a = np^n + mp^m + n + m$ と取ればよい。
\end{proof}

\begin{lemma}
\label{algebra-lemma-p-ring-map-field}
$k'/k$ を体拡大とする。$p$ を素数とする。次は同値である。
\begin{enumerate}
\item $k'$ は $k$ の体拡大として、ある $n > 0$ に対して
$x^{p^n} \in k$ かつ $p^nx \in k$ を満たす元 $x$ によって生成される。
\item $k = k'$ であるか、または $k$ と $k'$ の標数が $p$ であり、
$k'/k$ は純非分離である。
\end{enumerate}
\end{lemma}

\begin{proof}
$x \in k'$ とする。ある $n > 0$ に対して $x^{p^n} \in k$ かつ
$p^nx \in k$ であり、標数が $p$ でなければ $x \in k$ である。
標数が $p$ ならば $x^{p^n} \in k$ を得るので、$x$ は $k$ 上純非分離である。
\end{proof}

\begin{lemma}
\label{algebra-lemma-p-ring-map}
$\varphi : R \to S$ を環準同型とする。$p$ を素数とする。次を仮定する。
\begin{enumerate}
\item[(a)] $S$ は $R$-代数として、ある $n > 0$ に対して
$x^{p^n} \in \varphi(R)$ かつ $p^nx \in \varphi(R)$ を満たす
元 $x$ によって生成される。
\item[(b)] $\Ker(\varphi)$ は局所冪零である。
\end{enumerate}
このとき $\varphi$ はスペクトルの同相写像を誘導し、補題
\ref{algebra-lemma-p-ring-map-field} の同値な条件を満たす剰余体拡大を誘導する。
任意の環準同型 $R \to R'$ に対して、環準同型
$R' \to R' \otimes_R S$ も (a) と (b) を満たす。
\end{lemma}

\begin{proof}
(a) と (b) を仮定する。(b) は補題 \ref{algebra-lemma-powers} の条件 (2) と同値である。
$T \subset S$ を、ある整数 $n > 0$ に対して
$x^{p^n} , p^n x \in \varphi(R)$ を満たす元 $x \in S$ の集合とする。
$T = S$ であることを主張する。これにより補題 \ref{algebra-lemma-powers} の条件 (1) が
成り立つことが分かり、したがって $\varphi$ はスペクトル上の同相写像を誘導する。
仮定 (a) により、$T \subset S$ が $R$-部分代数であることを示せば十分である。
$x \in T$ かつ $y \in R$ ならば、$yx \in T$ であることは明らかである。
$x, y \in T$ とし、$x^{p^n}, y^{p^m}, p^n x, p^m y \in \varphi(R)$ を
満たす $n, m > 0$ を取る。このとき
$(xy)^{p^{n + m}}, p^{n + m}xy \in \varphi(R)$ であるから $xy \in T$ である。
補題 \ref{algebra-lemma-help-with-powers} により $x + y \in T$ であり、主張が従う。

\medskip\noindent
$\varphi$ はスペクトル上の同相写像を誘導するので、特にスペクトル上全射である。
これは任意の基底変換で保たれる性質である。補題
\ref{algebra-lemma-surjective-spec-radical-ideal} を参照せよ。
したがって任意の $R \to R'$ に対し、環準同型
$R' \to R' \otimes_R S$ の核は冪零元からなる。補題
\ref{algebra-lemma-image-dense-generic-points} を参照せよ。
言い換えれば、$R' \to R' \otimes_R S$ に対して (b) が成り立つ。
(a) が基底変換で保たれることは明らかである。
最後に、剰余体に関する条件は (a) から従う。実際、$R$-代数としての
$S$ の生成元は、剰余体拡大の生成元へ写る。
\end{proof}

\begin{lemma}
\label{algebra-lemma-radicial}
$\varphi : R \to S$ を環準同型とする。次を仮定する。
\begin{enumerate}
\item $\varphi$ はスペクトル上の単射を誘導する。
\item $\varphi$ は純非分離な剰余体拡大を誘導する。
\end{enumerate}
このとき任意の環準同型 $R \to R'$ に対して、性質 (1) と (2) は
$R' \to R' \otimes_R S$ についても成り立つ。
\end{lemma}

\begin{proof}
$S' = R' \otimes_R S$ と置くと、環のスペクトルの連続写像からなる可換図式
$$
\xymatrix{
\Spec(S') \ar[r] \ar[d] & \Spec(S) \ar[d] \\
\Spec(R') \ar[r] & \Spec(R)
}
$$
を得る。$\mathfrak p \subset R$ の上にある素イデアル
$\mathfrak p' \subset R'$ を取る。$\mathfrak p$ の上にある
$S$ の素イデアルが存在しなければ、$\mathfrak p'$ の上にある
$S'$ の素イデアルも存在しない。そうでなければ、注意
\ref{algebra-remark-fundamental-diagram} により
$F = S \otimes_R \kappa(\mathfrak p)$ の素イデアル $\mathfrak r$ で、
その剰余体が $\kappa(\mathfrak p)$ 上純非分離であるものが一意に存在する。
環準同型
$$
\kappa(\mathfrak p) \to F \to \kappa(\mathfrak r)
$$
を考える。補題 \ref{algebra-lemma-minimal-prime-reduced-ring} により、イデアル
$\mathfrak r \subset F$ は局所冪零である。したがって環準同型
$F \to \kappa(\mathfrak r)$ に補題
\ref{algebra-lemma-surjective-locally-nilpotent-kernel} を適用できる。
環準同型 $\kappa(\mathfrak p) \to \kappa(\mathfrak r)$ には補題
\ref{algebra-lemma-p-ring-map} を適用できる。したがって、次の写像における合成写像と
第二の射は、スペクトル上の全単射と純非分離な剰余体拡大を誘導する。
$$
\kappa(\mathfrak p') \to
\kappa(\mathfrak p') \otimes_{\kappa(\mathfrak p)} F \to
\kappa(\mathfrak p') \otimes_{\kappa(\mathfrak p)} \kappa(\mathfrak r)
$$
よって第一の写像についても同じことが成り立つ。さらに
$$
\kappa(\mathfrak p') \otimes_{\kappa(\mathfrak p)} F =
\kappa(\mathfrak p') \otimes_{\kappa(\mathfrak p)}
\kappa(\mathfrak p) \otimes_R S =
\kappa(\mathfrak p') \otimes_R S =
\kappa(\mathfrak p') \otimes_{R'} R' \otimes_R S
$$
であるから、注意 \ref{algebra-remark-fundamental-diagram} の議論により結論を得る。
\end{proof}

\begin{lemma}
\label{algebra-lemma-radicial-integral}
$\varphi : R \to S$ を環準同型とする。次を仮定する。
\begin{enumerate}
\item $\varphi$ は整である。
\item $\varphi$ はスペクトル上の単射を誘導する。
\item $\varphi$ は純非分離な剰余体拡大を誘導する。
\end{enumerate}
このとき $\varphi$ は $\Spec(S)$ から $\Spec(R)$ の閉部分集合への
同相写像を誘導する。また任意の環準同型 $R \to R'$ に対して、性質
(1), (2), (3) は $R' \to R' \otimes_R S$ についても成り立つ。
\end{lemma}

\begin{proof}
補題 \ref{algebra-lemma-going-up-closed} および \ref{algebra-lemma-integral-going-up} により、
スペクトル上の写像は閉写像である。補題 \ref{algebra-lemma-radicial} および
\ref{algebra-lemma-base-change-integral} により、これらの性質は基底変換で保たれる。
\end{proof}

\begin{lemma}
\label{algebra-lemma-radicial-integral-bijective}
% P01-E0344: frozen source has the article error "an bijective map".
$\varphi : R \to S$ を環準同型とする。次を仮定する。
\begin{enumerate}
\item $\varphi$ は整である。
\item $\varphi$ はスペクトル上の全単射を誘導する。
\item $\varphi$ は純非分離な剰余体拡大を誘導する。
\end{enumerate}
このとき $\varphi$ はスペクトル上の同相写像を誘導する。また任意の環準同型
$R \to R'$ に対して、性質 (1), (2), (3) は
$R' \to R' \otimes_R S$ についても成り立つ。
\end{lemma}

\begin{proof}
補題 \ref{algebra-lemma-radicial-integral} および
\ref{algebra-lemma-surjective-spec-radical-ideal} から従う。
\end{proof}

\begin{lemma}
\label{algebra-lemma-universally-bijective}
$\varphi : R \to S$ を次を満たす環準同型とする。
\begin{enumerate}
\item $\varphi$ の核は局所冪零である。
\item $S$ は $R$-代数として、ある $n > 0$ と多項式
$P(T) \in R[T]$ が存在し、その $S[T]$ における像が $(T - x)^n$ となるような
元 $x$ によって生成される。
\end{enumerate}
このとき $\Spec(S) \to \Spec(R)$ は同相写像であり、$R \to S$ は
純非分離な剰余体拡大を誘導する。さらに、条件 (1) と (2) は任意の基底変換で
引き続き成り立つ。
\end{lemma}

\begin{proof}
補題 \ref{algebra-lemma-surjective-locally-nilpotent-kernel} により、$R$ を
$R/\Ker(\varphi)$ で置き換えてよい。仮定 (2) により、$S$ は $R$ 上整な元によって
$R$ 上生成される。したがって $R \subset S$ は整である
（補題 \ref{algebra-lemma-integral-closure-is-ring}）。特に
$\Spec(S) \to \Spec(R)$ は全射かつ閉写像である（補題
\ref{algebra-lemma-integral-overring-surjective},
\ref{algebra-lemma-going-up-closed}, および
\ref{algebra-lemma-integral-going-up}）。

\medskip\noindent
% P01-E0345: frozen source has the malformed phrase "there exists an n > 0 be such that".
$x \in S$ を (2) の生成元の一つとする。すなわち、
$(T - x)^n \in R[T]$ を満たす $n > 0$ が存在する。
$\mathfrak p \subset R$ を素イデアルとする。上で示したことと補題
\ref{algebra-lemma-in-image} により、環 $\kappa(\mathfrak p) \otimes_R S$ は零環でない。
$\kappa(\mathfrak p)$ の標数が零ならば、$nx \in R$ であることから
$1 \otimes x$ は $\kappa(\mathfrak p) \to \kappa(\mathfrak p) \otimes_R S$ の
像に属することが分かる。したがって
$\kappa(\mathfrak p) \to \kappa(\mathfrak p) \otimes_R S$ は同型である。
$\kappa(\mathfrak p)$ の標数が $p > 0$ ならば、$p$ と互いに素な $m$ を用いて
$n = p^k m$ と書く。$\kappa(\mathfrak p) \otimes_R S[T]$ において
$$
(T - 1 \otimes x)^n = ((T - 1 \otimes x)^{p^k})^m =
(T^{p^k} - 1 \otimes x^{p^k})^m
$$
が成り立ち、$mx^{p^k} \in R$ であることが分かる。したがって
$1 \otimes x^{p^k}$ は
$\kappa(\mathfrak p) \to \kappa(\mathfrak p) \otimes_R S$ の像に属する。
よって環準同型 $\kappa(\mathfrak p) \to \kappa(\mathfrak p) \otimes_R S$ に
補題 \ref{algebra-lemma-p-ring-map} を適用できる。いずれの場合にも、
$\kappa(\mathfrak p) \otimes_R S$ は一意な素イデアルを持ち、その剰余体は
$\kappa(\mathfrak p)$ 上純非分離であると結論できる。注意
\ref{algebra-remark-fundamental-diagram} により、$\varphi$ はスペクトル上全単射である。

\medskip\noindent
基底変換についての主張は直ちに従う。
\end{proof}










\section{幾何学的既約代数}
\label{algebra-section-algebras-over-fields}

\noindent
体 $k$ 上の代数 $S$ が幾何学的既約であるとは、任意の体拡大 $k'/k$ に対して
代数 $S \otimes_k k'$ が一意な極小素イデアルを持つことをいう。
本節ではこの概念に関連する理論を少し展開する。

\begin{lemma}
\label{algebra-lemma-flat-fibres-irreducible}
$R \to S$ を環準同型とする。次を仮定する。
\begin{enumerate}
\item[(a)] $\Spec(R)$ は既約である。
\item[(b)] $R \to S$ は平坦である。
\item[(c)] $R \to S$ は有限表示である。
\item[(d)] $R$ の素イデアル $\mathfrak p$ の稠密な集合に対して、
ファイバー環 $S \otimes_R \kappa(\mathfrak p)$ のスペクトルは既約である。
\end{enumerate}
このとき $\Spec(S)$ は既約である。さらに一般に、(b) $+$ (c) を
「写像 $\Spec(S) \to \Spec(R)$ は開写像である」で置き換えても成り立つ。
\end{lemma}

\begin{proof}
仮定 (b) と (c) により、スペクトル上の写像は開写像である。命題
\ref{algebra-proposition-fppf-open} を参照せよ。したがって補題は
Topology, Lemma \ref{topology-lemma-irreducible-on-top} から従う。
\end{proof}

\begin{lemma}
\label{algebra-lemma-separably-closed-irreducible}
$k$ を分離閉体とする。$R$, $S$ を $k$-代数とする。
$R$, $S$ がそれぞれ一意な極小素イデアルを持つならば、
$R \otimes_k S$ も一意な極小素イデアルを持つ。
\end{lemma}

\begin{proof}
$k \subset \overline{k}$ を完全閉包とする。定義
\ref{algebra-definition-perfection} を参照せよ。仮定により $\overline{k}$ は代数閉である。
環準同型 $R \to R \otimes_k \overline{k}$ と
$S \to S \otimes_k \overline{k}$ および
$R \otimes_k S \to (R \otimes_k S) \otimes_k \overline{k}
= (R \otimes_k \overline{k}) \otimes_{\overline{k}} (S \otimes_k \overline{k})$
は補題 \ref{algebra-lemma-p-ring-map} の仮定を満たす。
したがって $k$ は代数閉であると仮定してよい。

\medskip\noindent
$R$ と $S$ をそれぞれの被約化で置き換えてよい。
したがって $R$ と $S$ は整域であると仮定してよい。
補題 \ref{algebra-lemma-perfect-reduced} により $R \otimes_k S$ は被約である。
よってそのスペクトルが可約であることと、非零零因子を含むことは同値である。
補題 \ref{algebra-lemma-limit-argument} により、$R$ と $S$ が代数閉体 $k$ 上
有限型の整域である場合に帰着する。

\medskip\noindent
環準同型 $R \to R \otimes_k S$ は有限表示かつ平坦である。
さらに、$R$ の任意の極大イデアル $\mathfrak m$ に対して
$(R \otimes_k S) \otimes_R R/\mathfrak m \cong S$ が成り立つ。実際、
Hilbert の零点定理 \ref{algebra-theorem-nullstellensatz} により
$k \cong R/\mathfrak m$ である。また補題
\ref{algebra-lemma-finite-type-field-Jacobson} により $\Spec(R)$ は Jacobson であるから、
極大イデアルの集合は $R$ のスペクトルで稠密である。
したがって環準同型 $R \to R \otimes_k S$ に補題
\ref{algebra-lemma-flat-fibres-irreducible} を適用でき、望みどおり
$R \otimes_k S$ のスペクトルは既約であると結論できる。
\end{proof}

\begin{lemma}
\label{algebra-lemma-geometrically-irreducible}
$k$ を体とする。$R$ を $k$-代数とする。次は同値である。
\begin{enumerate}
\item 任意の体拡大 $k'/k$ に対して、$R \otimes_k k'$ のスペクトルは既約である。
\item 任意の有限分離体拡大 $k'/k$ に対して、$R \otimes_k k'$ のスペクトルは既約である。
\item $\overline{k}$ を $k$ の分離閉包とすると、
$R \otimes_k \overline{k}$ のスペクトルは既約である。
\item $\overline{k}$ を $k$ の代数閉包とすると、
$R \otimes_k \overline{k}$ のスペクトルは既約である。
\end{enumerate}
\end{lemma}

\begin{proof}
(1) が (2) を含意することは明らかである。

\medskip\noindent
(2) を仮定し、$\overline{k}$ を $k$ の分離閉包とする。
% P01-E0346: frozen source has the malformed construction "let ... is".
$\mathfrak q_i \subset R \otimes_k \overline{k}$, $i = 1, 2$ を二つの
極小素イデアルとする。任意の有限部分拡大 $\overline{k}/k'/k$ に対して、
拡大 $k'/k$ は分離的であり、環準同型
$R \otimes_k k' \to R \otimes_k \overline{k}$ は平坦である。
したがって $\mathfrak p_i = (R \otimes_k k') \cap \mathfrak q_i$ は
極小素イデアルである（補題 \ref{algebra-lemma-flat-going-down} により、
平坦な環準同型に対して下降定理が成り立つ）。よって仮定 (2) により
$\mathfrak p_1 = \mathfrak p_2$ である。
$\overline{k} = \bigcup k'$ であるから $\mathfrak q_1 = \mathfrak q_2$ と結論できる。
したがって $\Spec(R \otimes_k \overline{k})$ は既約である。

\medskip\noindent
(3) を仮定し、$\overline{k}$ を $k$ の代数閉包とする。
$\overline{k}/\overline{k}'/k$ を対応する $k$ の分離閉包とする。
このとき $\overline{k}/\overline{k}'$ は（正標数では）純非分離であるか、
または自明である。したがって、例えば補題 \ref{algebra-lemma-p-ring-map} により、
$R \otimes_k \overline{k}' \to R \otimes_k \overline{k}$ は
スペクトル上の同相写像を誘導する。よって (4) を得る。

\medskip\noindent
(4) を仮定する。$k'/k$ を任意の体拡大とし、$\overline{k}$ を
$k$ の代数閉包とする。$k'$ と $\overline{k}$ の双方が $F$ の部分体と
同型になるような体 $F$ を選べる。このとき
$$
R \otimes_k F = (R \otimes_k \overline{k}) \otimes_{\overline{k}} F
$$
である。したがって補題 \ref{algebra-lemma-separably-closed-irreducible} により、
$R \otimes_k F$ は一意な極小素イデアルを持つ。最後に、環準同型
$R \otimes_k k' \to R \otimes_k F$ は平坦かつ単射である。よって
$R \otimes_k k'$ の任意の極小素イデアルは $R \otimes_k F$ の
極小素イデアルの像である（補題
\ref{algebra-lemma-injective-minimal-primes-in-image} と下降定理による）。
したがってそのような極小素イデアルは一つしかなく、証明が完了する。
\end{proof}

\begin{definition}
\label{algebra-definition-geometrically-irreducible}
$k$ を体とし、$S$ を $k$-代数とする。任意の体拡大 $k'/k$ に対して
$S \otimes_k k'$ のスペクトルが既約であるとき、$S$ は
{\it $k$ 上幾何学的既約} であるという\footnote{既約空間は空でない。}。
\end{definition}

\noindent
補題 \ref{algebra-lemma-geometrically-irreducible} により、これは有限分離体拡大
$k'/k$ について、または $k'$ が $k$ の分離閉包に等しい場合について
確認すれば十分である。

\begin{lemma}
\label{algebra-lemma-separably-closed-irreducible-implies-geometric}
$k$ を体とし、$R$ を $k$-代数とする。$k$ が分離閉ならば、
$R$ が $k$ 上幾何学的既約であることと、$R$ のスペクトルが既約であることは同値である。
\end{lemma}

\begin{proof}
定義 \ref{algebra-definition-geometrically-irreducible} の直後の注意から直ちに従う。
\end{proof}

\begin{lemma}
\label{algebra-lemma-subalgebra-geometrically-irreducible}
$k$ を体とし、$S$ を $k$-代数とする。
\begin{enumerate}
\item $S$ が $k$ 上幾何学的既約ならば、任意の $k$-部分代数もそうである。
\item $S$ のすべての有限生成 $k$-部分代数が幾何学的既約ならば、
$S$ は幾何学的既約である。
\item 幾何学的既約な $k$-代数の有向余極限は幾何学的既約である。
\end{enumerate}
\end{lemma}

\begin{proof}
$S' \subset S$ を部分代数とする。このとき任意の拡大 $k'/k$ に対して、
環準同型 $S' \otimes_k k' \to S \otimes_k k'$ も単射である。
したがって (1) は補題 \ref{algebra-lemma-injective-minimal-primes-in-image} から従う
（既約空間の連続写像による像が既約であることも用いる）。
% P01-E0348: frozen source uses singular "property" for properties (2) and (3).
第二および第三の性質は、テンソル積が余極限と可換であることから従う。
\end{proof}

\begin{lemma}
\label{algebra-lemma-geometrically-irreducible-any-base-change}
$k$ を体とする。$S$ を幾何学的既約な $k$-代数とする。
$R$ を任意の $k$-代数とする。写像
$$
\Spec(R \otimes_k S) \longrightarrow \Spec(R)
$$
は既約成分上の全単射を誘導する。
\end{lemma}

\begin{proof}
既約成分が極小素イデアルに対応することを思い出そう
（補題 \ref{algebra-lemma-irreducible}）。$R \to R \otimes_k S$ は平坦であるから、
下降定理（補題 \ref{algebra-lemma-flat-going-down}）により、
$R \otimes_k S$ の任意の極小素イデアルは $R$ の極小素イデアルの上にある。
逆に、$\mathfrak p \subset R$ を（極小）素イデアルとすると、
$$
R \otimes_k S/\mathfrak p(R \otimes_k S)
=
(R/\mathfrak p) \otimes_k S
\subset
\kappa(\mathfrak p) \otimes_k S
$$
が $R \to R \otimes_k S$ の平坦性により成り立つ。
仮定により、環 $\kappa(\mathfrak p) \otimes_k S$ のスペクトルは既約である。
したがって $R \otimes_k S/\mathfrak p(R \otimes_k S)$ は一意な
極小素イデアルを持つ（補題 \ref{algebra-lemma-injective-minimal-primes-in-image}）。
言い換えれば、既約集合 $V(\mathfrak p)$ の逆像は既約である。
よって補題が従う。
\end{proof}

\noindent
幾何学的既約な体拡大という概念について、いくつか注意を述べよう。

\begin{lemma}
\label{algebra-lemma-field-extension-geometrically-irreducible}
$K/k$ を体拡大とする。$k$ が $K$ の中で代数閉ならば、
$K$ は $k$ 上幾何学的既約である。
\end{lemma}

\begin{proof}
$k$ が $K$ の中で代数閉であると仮定する。定義
\ref{algebra-definition-geometrically-irreducible} と補題
\ref{algebra-lemma-geometrically-irreducible} により、任意の有限分離拡大 $k'/k$ に対して
$K \otimes_k k'$ のスペクトルが既約であることを示せば十分である。
Fields, Lemma \ref{fields-lemma-primitive-element} により、$k'$ は $k$ 上
$\alpha \in k'$ によって生成されるとする。$\alpha$ の最小多項式を
$P = T^d + a_1 T^{d - 1} + \ldots + a_d \in k[T]$ とする。このとき
$K \otimes_k k' \cong K[T]/(P)$ である。
$K[T]/(P)$ のスペクトルが可約になり得るのは、$P$ が $K[T]$ で
可約な場合だけである。矛盾を導くため、$K[T]$ において
$P = P_1 P_2$ が非自明な因数分解であると仮定する。
補題 \ref{algebra-lemma-polynomials-divide} により、$P_1$ と $P_2$ の係数は
$k$ 上代数的である。仮定により $P_1$ と $P_2$ の係数は $k$ に属し、
これは $P$ が $k$ 上既約であることに矛盾する。
\end{proof}

\begin{lemma}
\label{algebra-lemma-geometrically-irreducible-transitive}
$K/k$ を幾何学的既約な体拡大とする。$S$ を幾何学的既約な
$K$-代数とする。このとき $S$ は $k$ 上幾何学的既約である。
\end{lemma}

\begin{proof}
定義 \ref{algebra-definition-geometrically-irreducible} と補題
\ref{algebra-lemma-geometrically-irreducible} により、任意の有限分離拡大 $k'/k$ に対して
$S \otimes_k k'$ のスペクトルが既約であることを示せば十分である。
$K$ は $k$ 上幾何学的既約であるから、$K' = K \otimes_k k'$ は
$K$ の有限分離体拡大である。$S$ は $K$ 上幾何学的既約であると仮定したので、
$S \otimes_k k' = S \otimes_K K'$ のスペクトルは既約である。
\end{proof}

\begin{lemma}
\label{algebra-lemma-geometrically-irreducible-base-change-transcendental}
$K/k$ を体拡大とする。次は同値である。
\begin{enumerate}
\item $K$ は $k$ 上幾何学的既約である。
\item 純超越拡大から誘導される拡大 $K(t)/k(t)$ は幾何学的既約である。
\end{enumerate}
\end{lemma}

\begin{proof}
(1) を仮定する。
% P01-E0349: frozen source omits "by" in "Denote Omega".
$\Omega$ を $k(t)$ の代数閉包とする。定義
\ref{algebra-definition-geometrically-irreducible} により
$$
K \otimes_k \Omega = K \otimes_k k(t) \otimes_{k(t)} \Omega
$$
のスペクトルは既約である。$K(t)$ は $K \otimes_k k(T)$ の局所化であるから、
% P01-E0350: frozen source changes the indeterminate from t to T in this localization.
$K(t) \otimes_{k(t)} \Omega$ のスペクトルは既約であると結論できる。
したがって補題 \ref{algebra-lemma-geometrically-irreducible} により、
$K(t)/k(t)$ は幾何学的既約である。

\medskip\noindent
(2) を仮定する。$k'/k$ を体拡大とする。
$K \otimes_k k'$ が一意な極小素イデアルを持つことを示さなければならない。
次のスペクトルが既約、すなわち一意な極小素イデアルを持つことは分かっている。
$$
K(t) \otimes_{k(t)} k'(t)
$$
単射 $K \otimes_k k' \to K(t) \otimes_{k(t)} k'(t)$ が存在するので
（詳細は省略する）、補題 \ref{algebra-lemma-injective-minimal-primes-in-image} および
\ref{algebra-lemma-minimal-prime-image-minimal-prime} により結論を得る。
\end{proof}

\begin{lemma}
\label{algebra-lemma-geometrically-irreducible-add-transcendental}
$K/L/M$ を体の塔とし、$L/M$ は幾何学的既約であるとする。
$x \in K$ は $L$ 上超越的であるとする。このとき $L(x)/M(x)$ は
幾何学的既約である。
\end{lemma}

\begin{proof}
補題 \ref{algebra-lemma-geometrically-irreducible-base-change-transcendental} から従う。
実際、体 $L(x)$ と $M(x)$ はそれぞれ $L$ と $M$ の純超越拡大である。
\end{proof}

\begin{lemma}
\label{algebra-lemma-geometrically-irreducible-separable-elements}
$K/k$ を体拡大とする。次は同値である。
\begin{enumerate}
\item $K/k$ は幾何学的既約である。
\item $k$ 上分離代数的な任意の元 $\alpha \in K$ は $k$ に属する。
\end{enumerate}
\end{lemma}

\begin{proof}
(1) を仮定し、$\alpha \in K$ は $k$ 上分離代数的であるとする。
このとき $k' = k(\alpha)$ は $k$ の有限分離拡大であり、$K$ に含まれる。
補題 \ref{algebra-lemma-subalgebra-geometrically-irreducible} により、拡大 $k'/k$ は
幾何学的既約である。特に $k' \otimes_k \overline{k}$ のスペクトルは既約である
（したがって、それが体の積ならば因子はちょうど一つである）。
Fields, Lemma \ref{fields-lemma-finite-separable-tensor-alg-closed} により
$\Hom_k(k', \overline{k})$ は一つの元を持つ。すると Fields, Lemma
\ref{fields-lemma-separable-equality} により $k' = k$ である。よって (2) が成り立つ。

\medskip\noindent
(2) を仮定する。$k$ 上代数的な元からなる部分体を $k' \subset K$ とする。
補題 \ref{algebra-lemma-field-extension-geometrically-irreducible} により、拡大 $K/k'$ は
幾何学的既約である。仮定により $k'/k$ は純非分離拡大である。
補題 \ref{algebra-lemma-p-ring-map} により拡大 $k'/k$ は幾何学的既約である。
したがって補題 \ref{algebra-lemma-geometrically-irreducible-transitive} により、
$K/k$ は幾何学的既約である。
\end{proof}

\begin{lemma}
\label{algebra-lemma-make-geometrically-irreducible}
$K/k$ を体拡大とする。$k$ 上分離代数的な元からなる部分拡大
$K/k'/k$ を考える。このとき $K$ は $k'$ 上幾何学的既約である。
$K/k$ が有限生成体拡大ならば、$[k' : k] < \infty$ である。
\end{lemma}

\begin{proof}
第一の主張は、補題 \ref{algebra-lemma-geometrically-irreducible-separable-elements} と、
分離代数拡大の推移性により $k'$ 上分離代数的な元は $k'$ に属することから
直ちに従う。Fields, Lemma \ref{fields-lemma-separable-permanence} を参照せよ。
$K/k$ が有限生成ならば、Fields, Lemma
\ref{fields-lemma-algebraic-closure-in-finitely-generated} により
$k'$ は $k$ 上有限である。
\end{proof}

\begin{lemma}
\label{algebra-lemma-Galois-orbit}
$K/k$ を体拡大とする。$\overline{k}/k$ を分離閉包とする。
このとき $\text{Gal}(\overline{k}/k)$ は
$\overline{k} \otimes_k K$ の素イデアルに推移的に作用する。
\end{lemma}

\begin{proof}
補題 \ref{algebra-lemma-make-geometrically-irreducible} で得られる部分拡大を
$K/k'/k$ とする。$k \subset \overline{k}$ は整であるから、
$\overline{k} \otimes_k K$ と $\overline{k} \otimes_k k'$ のすべての
素イデアルは極大であることに注意する。補題
\ref{algebra-lemma-integral-no-inclusion} を参照せよ。補題
\ref{algebra-lemma-geometrically-irreducible-any-base-change} により、写像
$$
\Spec(\overline{k} \otimes_k K) \to \Spec(\overline{k} \otimes_k k')
$$
は全単射である。実際、(1) すべての素イデアルは極小素イデアルであり、
(2) $\overline{k} \otimes_k K = (\overline{k} \otimes_k k') \otimes_{k'} K$ であり、
(3) $K$ は $k'$ 上幾何学的既約である。
したがって $\overline{k} \otimes_k k'$ の素イデアルへの
$\text{Gal}(\overline{k}/k)$ の作用について補題を証明すれば十分である。

\medskip\noindent
$\overline{k} \otimes_k k'$ のすべての素イデアルは極大であるから、
剰余体は $\overline{k}$ と同型である。したがって
$\overline{k} \otimes_k k'$ の素イデアルは $\Hom_k(k', \overline{k})$ の元と
一対一に対応する。ここで $\sigma \in \Hom_k(k', \overline{k})$ は
$1 \otimes \sigma : \overline{k} \otimes_k k' \to \overline{k}$ の核
$\mathfrak p_\sigma$ に対応する。特に $\text{Gal}(\overline{k}/k)$ は
この集合に推移的に作用し、望む結論を得る。
\end{proof}




\section{幾何学的連結代数}
\label{algebra-section-geometrically-connected}

\begin{lemma}
\label{algebra-lemma-separably-closed-connected}
% P01-E0351: frozen source uses the nonstandard phrase "separably algebraically closed" twice in this section.
% P01-E0352: frozen source has a spurious comma between the coordinated spectra.
$k$ を分離閉体とする。$R$, $S$ を $k$-代数とする。
$\Spec(R)$ と $\Spec(S)$ が連結ならば、$\Spec(R \otimes_k S)$ も連結である。
\end{lemma}

\begin{proof}
$\Spec(R)$ が連結であることと、$R$ が非自明な冪等元を持たないことは同値である。
補題 \ref{algebra-lemma-characterize-spec-connected} を参照せよ。したがって補題
\ref{algebra-lemma-limit-argument} により、$R$ と $S$ は $k$ 上有限型であると仮定してよい。
この場合 $R$ と $S$ はネーターであり、有限個の極小素イデアルを持つ。
補題 \ref{algebra-lemma-Noetherian-irreducible-components} を参照せよ。
よって $n + m$ に関する帰納法を用いる。ここで $n$, resp.\ $m$ は
$\Spec(R)$, resp.\ $\Spec(S)$ の既約成分の個数である。
$n$ または $m$ が零の場合はもちろん自明である。$n = m = 1$、すなわち
$\Spec(R)$ と $\Spec(S)$ がともに一つの既約成分を持つ場合、
補題 \ref{algebra-lemma-separably-closed-irreducible} により結果が成り立つ。
$n > 1$ とする。既約閉部分集合 $T \subset \Spec(R)$ に対応する極小素イデアルを
$\mathfrak p \subset R$ とする。$T' \subset \Spec(R)$ を、残りの
$n - 1$ 個の既約成分の合併とする。$T' = V(I) = \Spec(R/I)$ を満たす
イデアル $I \subset R$ を選ぶ（補題 \ref{algebra-lemma-spec-closed}）。
極小素イデアルを注意深く選ぶことにより、さらに $T'$ が連結になるようにできる。
Topology, Lemma \ref{topology-lemma-remove-irreducible-connected} を参照せよ。
このとき $T \cup T' = \Spec(R)$ であり、
$\Spec(R)$ は連結であると仮定したので
$T \cap T' = V(\mathfrak p + I) = \Spec(R/(\mathfrak p + I))$ は空でない。
$\Spec(R \otimes_k S)$ における $T$ の逆像は
$\Spec(R/\mathfrak p \otimes_k S)$ であり、$\Spec(R \otimes_k S)$ における
$T'$ の逆像は
% P01-E0353: frozen source omits "image" in the second inverse-image phrase.
$\Spec(R/I \otimes_k S)$ である。帰納法により、これらはいずれも連結である。
$T \cap T'$ の逆像は $\Spec(R/(\mathfrak p + I) \otimes_k S)$ であり、空でない。
したがって $\Spec(R \otimes_k S)$ は連結である。
\end{proof}

\begin{lemma}
\label{algebra-lemma-geometrically-connected}
$k$ を体とし、$R$ を $k$-代数とする。次は同値である。
\begin{enumerate}
\item 任意の体拡大 $k'/k$ に対して、$R \otimes_k k'$ のスペクトルは連結である。
\item 任意の有限分離体拡大 $k'/k$ に対して、$R \otimes_k k'$ のスペクトルは連結である。
\end{enumerate}
\end{lemma}

\begin{proof}
任意の体拡大 $k'/k$ に対し、$R \otimes_k k'$ のスペクトルが連結であることと、
$R \otimes_k k'$ が非自明な冪等元を持たないことは同値である。
補題 \ref{algebra-lemma-characterize-spec-connected} を参照せよ。(2) を仮定する。
$k \subset \overline{k}$ を $k$ の分離閉包とする。補題
\ref{algebra-lemma-limit-argument} により、(2) は $R \otimes_k \overline{k}$ が
非自明な冪等元を持たないことと同値である。任意の体拡大 $k'/k$ に対して、
$k' \subset \overline{k}'$ を満たす体拡大 $\overline{k}'/\overline{k}$ が存在する。
補題 \ref{algebra-lemma-separably-closed-connected} により、
$R \otimes_k \overline{k}'$ は非自明な冪等元を持たない。
$R \otimes_k k'$ が非自明な冪等元を持てば、$R \otimes_k \overline{k}'$ も
非自明な冪等元を持つことになり、矛盾である。
\end{proof}

\begin{definition}
\label{algebra-definition-geometrically-connected}
$k$ を体とし、$S$ を $k$-代数とする。任意の体拡大 $k'/k$ に対して
$S \otimes_k k'$ のスペクトルが連結であるとき、$S$ は
{\it $k$ 上幾何学的連結} であるという。
\end{definition}

\noindent
補題 \ref{algebra-lemma-geometrically-connected} により、これは有限分離体拡大
$k'/k$ について確認すれば十分である。

\begin{lemma}
\label{algebra-lemma-separably-closed-connected-implies-geometric}
$k$ を体とし、$R$ を $k$-代数とする。$k$ が分離閉ならば、
$R$ が $k$ 上幾何学的連結であることと、$R$ のスペクトルが連結であることは同値である。
\end{lemma}

\begin{proof}
定義 \ref{algebra-definition-geometrically-connected} の直後の注意から直ちに従う。
\end{proof}

\begin{lemma}
\label{algebra-lemma-subalgebra-geometrically-connected}
$k$ を体とし、$S$ を $k$-代数とする。
\begin{enumerate}
\item $S$ が $k$ 上幾何学的連結ならば、任意の $k$-部分代数もそうである。
\item $S$ のすべての有限生成 $k$-部分代数が幾何学的連結ならば、
$S$ は幾何学的連結である。
\item 幾何学的連結な $k$-代数の有向余極限は幾何学的連結である。
\end{enumerate}
\end{lemma}

\begin{proof}
これは、非自明な冪等元が存在しないという連結性の特徴づけから従う。
% P01-E0354: frozen source uses singular "property" for properties (2) and (3).
第二および第三の性質は、テンソル積が余極限と可換であることから従う。
\end{proof}

\noindent
次の補題は、より一般的な Varieties, Lemma
\ref{varieties-lemma-bijection-connected-components} によって後に置き換えられる。

\begin{lemma}
\label{algebra-lemma-geometrically-connected-any-base-change}
$k$ を体とする。$S$ を幾何学的連結な $k$-代数とし、
$R$ を任意の $k$-代数とする。写像
$$
R \longrightarrow R \otimes_k S
$$
は冪等元上の全単射を誘導し、写像
$$
\Spec(R \otimes_k S) \longrightarrow \Spec(R)
$$
は連結成分上の全単射を誘導する。
\end{lemma}

\begin{proof}
第二の主張は、第一の主張と補題 \ref{algebra-lemma-connected-component} から従う。
補題 \ref{algebra-lemma-subalgebra-geometrically-connected} および
\ref{algebra-lemma-limit-argument} により、$R$ と $S$ は $k$ 上有限型であると仮定してよい。
すると $R \otimes_k S$ も $k$ 上有限型である。この場合、すべての環は
ネーターであり、したがってそれらのスペクトルは有限個の連結成分を持つことに注意する
（補題 \ref{algebra-lemma-Noetherian-irreducible-components} により、有限個の既約成分を持つためである）。
特に、問題となるすべての連結成分は開である。したがって補題
\ref{algebra-lemma-disjoint-implies-product} により、この場合の補題の第一の主張は
第二の主張と同値である。これを証明しよう。
代数 $S$ は幾何学的連結かつ零でないので、
$X = \Spec(R \otimes_k S) \to \Spec(R) = Y$ のすべてのファイバーは連結かつ空でない。
また $R \to R \otimes_k S$ は平坦かつ有限表示であるから、写像
$X \to Y$ は開写像である（命題 \ref{algebra-proposition-fppf-open}）。
% P01-E0355: frozen source omits the article before "bijection".
Topology, Lemma \ref{topology-lemma-connected-fibres-connected-components} により、
$X \to Y$ は連結成分上の全単射を誘導する。
\end{proof}




\section{幾何学的整代数}
\label{algebra-section-geometrically-integral}

\noindent
定義は次のとおりである。

\begin{definition}
\label{algebra-definition-geometrically-integral}
$k$ を体とし、$S$ を $k$-代数とする。任意の体拡大 $k'/k$ に対して
$S \otimes_k k'$ が整域であるとき、$S$ は {\it $k$ 上幾何学的整} であるという。
% P01-E0356: frozen source has the extraneous preposition in "the ring of ... is a domain".
\end{definition}

\noindent
% P01-E0357: frozen source uses "translated in" instead of "translated into".
幾何学的整代数に関する任意の問題は、幾何学的被約かつ幾何学的既約な代数に
関する問題へと言い換えられる。

\begin{lemma}
\label{algebra-lemma-geometrically-integral}
$k$ を体とし、$S$ を $k$-代数とする。このとき $S$ が $k$ 上幾何学的整で
あることと、$S$ が $k$ 上幾何学的既約かつ幾何学的被約であることは同値である。
\end{lemma}

\begin{proof}
省略する。
\end{proof}

\begin{lemma}
\label{algebra-lemma-characterize-geometrically-integral}
$k$ を体とし、$S$ を $k$-代数とする。次は同値である。
\begin{enumerate}
\item $S$ は $k$ 上幾何学的整である。
\item 任意の有限体拡大 $k'/k$ に対して、環 $S \otimes_k k'$ は整域である。
\item $\overline{k}$ を $k$ の代数閉包とすると、$S \otimes_k \overline{k}$ は整域である。
\end{enumerate}
\end{lemma}

\begin{proof}
補題 \ref{algebra-lemma-geometrically-integral},
\ref{algebra-lemma-geometrically-reduced-finite-purely-inseparable-extension}, および
\ref{algebra-lemma-geometrically-irreducible} から従う。
\end{proof}

\begin{lemma}
\label{algebra-lemma-geometrically-integral-any-integral-base-change}
$k$ を体とする。$S$ を幾何学的整な $k$-代数とする。
$R$ を $k$-代数かつ整域とする。このとき $R \otimes_k S$ は整域である。
\end{lemma}

\begin{proof}
補題 \ref{algebra-lemma-geometrically-reduced-any-reduced-base-change} により
環 $R \otimes_k S$ は被約であり、補題
\ref{algebra-lemma-geometrically-irreducible-any-base-change} により環
$R \otimes_k S$ は既約である（スペクトルは既約成分を一つだけ持つ）。
したがって $R \otimes_k S$ は整域である。
\end{proof}




\section{付値環}
\label{algebra-section-valuation-rings}

\noindent
いくつかの定義を述べる。

\begin{definition}
\label{algebra-definition-valuation-ring}
付値環について次を定める。
\begin{enumerate}
\item $K$ を体とし、$A$, $B$ を $K$ に含まれる局所環とする。
$A \subset B$ かつ $\mathfrak m_A = A \cap \mathfrak m_B$ であるとき、
$B$ は $A$ を {\it 支配する} という。
\item $A$ を環とする。$A$ が局所整域であり、かつ $A$ の分数体に含まれる
局所環の間で $A$ が支配関係について極大であるとき、$A$ を {\it 付値環} という。
\item $A$ を分数体 $K$ を持つ付値環とする。$R \subset K$ が $K$ の部分環ならば、
$R \subset A$ であるとき $A$ は $R$ 上に {\it 中心をもつ} という。
\end{enumerate}
\end{definition}

\noindent
この定義では、体も付値環である。

\begin{lemma}
\label{algebra-lemma-dominate}
$K$ を体とし、$A \subset K$ を局所部分環とする。このとき分数体が $K$ であり、
$A$ を支配する付値環が存在する。
\end{lemma}

\begin{proof}
$K$ の局所部分環全体を、支配関係によって半順序集合とみなす。
$\{A_i\}_{i \in I}$ を $K$ の局所部分環からなる全順序族とする。このとき
$B = \bigcup A_i$ はすべての $A_i$ を支配する局所部分環である。
したがって Zorn の補題により、$A \subset K$ が分数体を $K$ としない局所環ならば、
$A$ を支配する局所環 $B \subset K$, $B \not = A$ が存在することを示せば十分である。

\medskip\noindent
$A$ の分数体に属さない $t \in K$ を取る。$t$ が $A$ 上超越的ならば、
$A[t] \subset K$ であり、したがって $A[t]_{(t, \mathfrak m)} \subset K$ は
$A$ と異なり $A$ を支配する局所環である。$t$ が $A$ 上代数的であるとする。
このとき、ある非零元 $a \in A$ に対して $at$ は $A$ 上整である。
$A$ と $ta$ により生成される部分環 $A' \subset K$ は $A$ 上有限である。
補題 \ref{algebra-lemma-integral-overring-surjective} により、$\mathfrak m$ の上にある
素イデアル $\mathfrak m' \subset A'$ が存在する。このとき
$A'_{\mathfrak m'}$ は $A$ を支配する。$A = A'_{\mathfrak m'}$ ならば
$t$ は $A$ の分数体に属することになり、仮定に反する。
したがって望みどおり $A \not = A'_{\mathfrak m'}$ である。
\end{proof}

\begin{lemma}
\label{algebra-lemma-valuation-ring-normal}
$A$ を付値環とする。このとき $A$ は正規整域である。
\end{lemma}

\begin{proof}
$x$ が $A$ の分数体に属し、$A$ 上整であるとする。
% P01-E0358: frozen proof uses K without introducing it as the fraction field of A.
$A$ と $x$ により生成される $K$ の部分環を $A'$ とする。
$A\subset A'$ は整拡大であるから、補題
\ref{algebra-lemma-integral-overring-surjective} により $\mathfrak m$ の上にある
素イデアル $\mathfrak m' \subset A'$ が存在する。このとき
$A'_{\mathfrak m'}$ は $A$ を支配する。$A$ は付値環であるから
$A=A'_{\mathfrak m'}$ と結論でき、したがって $x\in A$ である。
\end{proof}


\begin{lemma}
\label{algebra-lemma-valuation-ring-x-or-x-inverse}
$A$ を極大イデアル $\mathfrak m$ と分数体 $K$ を持つ付値環とする。
$x \in K$ とする。このとき $x \in A$ または $x^{-1} \in A$、あるいはその両方である。
\end{lemma}

\begin{proof}
$x$ は $A$ に属さないと仮定する。$A$ と $x$ により生成される
$K$ の部分環を $A'$ とする。$A$ は付値環であるから、$\mathfrak m$ の上にある
$A'$ の素イデアルは存在しない。$\mathfrak m$ は極大であるから
$V(\mathfrak m A') = \emptyset$ である。補題 \ref{algebra-lemma-Zariski-topology} により
$\mathfrak m A' = A'$ である。したがって $t_i \in \mathfrak m$ を用いて
$1 = \sum_{i = 0}^d t_i x^i$ と書ける。これは
% P01-E0359: frozen transformed polynomial omits the index and bounds of the second sum.
$(1 - t_0) (x^{-1})^d - \sum t_i (x^{-1})^{d - i} = 0$ を含意する。
特に $x^{-1}$ は $A$ 上整であるから、補題
\ref{algebra-lemma-valuation-ring-normal} により $x^{-1} \in A$ である。
\end{proof}

\begin{lemma}
\label{algebra-lemma-x-or-x-inverse-valuation-ring}
$A \subset K$ を体 $K$ の部分環とし、すべての $x \in K$ に対して
$x \in A$ または $x^{-1} \in A$、あるいはその両方が成り立つとする。
このとき $A$ は分数体 $K$ を持つ付値環である。
\end{lemma}

\begin{proof}
$A$ が $K$ でなければ、$A$ は体でなく、非零極大イデアル $\mathfrak m$ が存在する。
$\mathfrak m'$ を第二の極大イデアルとすると、
$x \in \mathfrak m$, $y \not \in \mathfrak m$,
$x \not \in \mathfrak m'$, $y \in \mathfrak m'$ を満たす $x, y \in A$ を選べる。
このとき $x/y \in A$ でも $y/x \in A$ でもなく、補題の仮定に矛盾する。
したがって $A$ は局所環である。$A'$ を $K$ に含まれ $A$ を支配する局所環とする。
$x \in A'$ を取り、$x \in A$ を示す必要がある。そうでなければ
$x^{-1} \in A$ であり、もちろん $x^{-1} \in \mathfrak m_A$ である。
すると $x^{-1} \in \mathfrak m_{A'}$ となり、$x \in A'$ に矛盾する。
\end{proof}

\begin{lemma}
\label{algebra-lemma-colimit-valuation-rings}
\begin{slogan}
付値環はフィルター付き直極限のもとで安定である
\end{slogan}
$I$ を有向集合とし、$(A_i, \varphi_{ij})$ を $I$ 上の付値環の系とする。
このとき $A = \colim A_i$ は付値環である。
\end{lemma}

\begin{proof}
$A$ が整域であることは明らかである。$a, b \in A$ とする。
補題 \ref{algebra-lemma-x-or-x-inverse-valuation-ring} により、$A$ において
$a | b$ または $b | a$ が成り立つことを示せばよい。
$a, b$ へ写る $a_i, b_i \in A_i$ が存在するほど大きい $i$ を選ぶ。
$A_i$ における $a_i, b_i$ に補題 \ref{algebra-lemma-valuation-ring-x-or-x-inverse} を
適用すれば、$A$ における $a, b$ について結果が従う。
\end{proof}

\begin{lemma}
\label{algebra-lemma-valuation-ring-cap-field}
$L/K$ を体拡大とする。$B \subset L$ が付値環ならば、
$A = K \cap B$ は付値環である。
\end{lemma}

\begin{proof}
$L$ を $B$ の分数体 $F$ で、$K$ を $K \cap F$ で置き換えてよい。
すると補題 \ref{algebra-lemma-valuation-ring-x-or-x-inverse} と
\ref{algebra-lemma-x-or-x-inverse-valuation-ring} を組み合わせれば従う。
\end{proof}

\begin{lemma}
\label{algebra-lemma-valuation-ring-cap-field-finite}
$L/K$ を代数体拡大とする。$B \subset L$ が分数体 $L$ を持つ付値環であり、
体でないならば、$A = K \cap B$ は体でない付値環である。
\end{lemma}

\begin{proof}
補題 \ref{algebra-lemma-valuation-ring-cap-field} により環 $A$ は付値環である。
$A$ が体ならば $A = K$ である。このとき $A = K \subset B$ は整拡大であるから、
$B$ の素イデアルの間に真の包含関係は存在しない
（補題 \ref{algebra-lemma-integral-no-inclusion}）。これは $B$ が体でない局所整域であるという
仮定に矛盾する。
\end{proof}

\begin{lemma}
\label{algebra-lemma-make-valuation-rings}
$A$ を付値環とする。任意の素イデアル $\mathfrak p \subset A$ に対して
商 $A/\mathfrak p$ は付値環である。局所化 $A_\mathfrak p$ についても同様であり、
実際 $A$ の任意の局所化について同じことが成り立つ。
\end{lemma}

\begin{proof}
補題 \ref{algebra-lemma-x-or-x-inverse-valuation-ring} で与えられた付値環の特徴づけを用いる。
\end{proof}

\begin{lemma}
\label{algebra-lemma-stack-valuation-rings}
$A'$ を剰余体 $K$ をもつ付値環とする。
$A$ を分数体 $K$ をもつ付値環とする。このとき
$C = \{\lambda \in A' \mid \lambda \bmod \mathfrak m_{A'} \in A\}$
は付値環である。
\end{lemma}

\begin{proof}
$\mathfrak m_{A'} \subset C$ かつ $C/\mathfrak m_{A'} = A$ であることに注意する。
特に、$C$ の分数体は $A'$ の分数体に等しい。補題
\ref{algebra-lemma-x-or-x-inverse-valuation-ring} の判定基準を用いて補題を示す。
$x$ を $C$ の分数体の元とする。補題により $x \in A'$ と仮定してよい。
$x \in \mathfrak m_{A'}$ ならば $x \in C$ である。そうでなければ $x$ は
$A'$ の単元であり、$x^{-1} \in A'$ でもある。したがって再び補題により、
$x$ または $x^{-1}$ のいずれかが $A$ の元へ写る。
\end{proof}

\begin{lemma}
\label{algebra-lemma-find-valuation-rings}
$A$ を分数体 $K$ をもつ正規整域とする。
\begin{enumerate}
\item 任意の $x \in K$, $x \not \in A$ に対して、分数体 $K$ をもち
$x \not \in V$ を満たす付値環 $A \subset V \subset K$ が存在する。
\item $A$ が局所環ならば、さらに $A$ を支配する $V$ を選べる。
\end{enumerate}
言い換えると、$A$ は $A$ を含む $K$ 内のすべての付値環の共通部分であり、
$A$ が局所環ならば、$A$ は $A$ を支配する $K$ 内のすべての付値環の共通部分である。
\end{lemma}

\begin{proof}
$x \in K$, $x \not \in A$ とする。$B = A[x^{-1}]$ を考える。
このとき $x \not \in B$ である。実際、
$x = a_0 + a_1x^{-1} + \ldots + a_d x^{-d}$ ならば
$x^{d + 1} - a_0x^d - \ldots - a_d = 0$ となり、$x$ は $A$ 上整となるが、
これは $A$ が正規であることに反する。したがって $x^{-1}$ は $B$ の単元でない。
ゆえに $V(x^{-1}) \subset \Spec(B)$ は空でなく
（補題 \ref{algebra-lemma-Zariski-topology}）、$x^{-1} \in \mathfrak p$ を満たす
素イデアル $\mathfrak p \subset B$ を選べる。$B_\mathfrak p$ を支配する
付値環 $V \subset K$ を選ぶ（補題 \ref{algebra-lemma-dominate}）。
$x^{-1} \in \mathfrak m_V$ なので $x \not \in V$ である。

\medskip\noindent
$A$ が局所環であるとする。このとき
$x^{-1} B + \mathfrak m_A B \not = B$ であると主張する。実際、
$a_i \in A$ および $a'_i \in \mathfrak m_A$ によって
$1 = (a_0 + a_1x^{-1} + \ldots + a_d x^{-d})x^{-1} +
a'_0 + \ldots + a'_d x^{-d}$ と書けるならば、
$$
(1 - a'_0) x^{d + 1} - (a_0 + a'_1) x^d - \ldots - a_d = 0
$$
を得る。$a'_0 \in \mathfrak m_A$ なので $1 - a'_0$ は $A$ の単元であり、
$x$ は $A$ 上整となるが、これは先ほどと同じく矛盾である。
% P01-E0360: frozen sentence runs together the prime and valuation-ring choices.
次に $\mathfrak p \supset x^{-1} B + \mathfrak m_A B$ を満たす素イデアルを選び、
その局所化を支配する付値環を取れば、$A$ を支配する $V$ が得られる。
\end{proof}

\noindent
{\it 全順序アーベル群}とは、アーベル群 $\Gamma$ と全順序 $\geq$ の組
$(\Gamma, \geq)$ であって、任意の
$\gamma, \gamma', \gamma'' \in \Gamma$ に対し
$\gamma \geq \gamma' \Rightarrow
\gamma + \gamma'' \geq \gamma' + \gamma''$ を満たすものをいう。

\begin{lemma}
\label{algebra-lemma-valuation-group}
$A$ を分数体 $K$ をもつ付値環とする。
$\Gamma = K^*/A^*$ と置く（群演算は加法的に書く）。
$\gamma, \gamma' \in \Gamma$ に対し、$\gamma - \gamma'$ が
$A - \{0\} \to \Gamma$ の像に属するとき、かつそのときに限り
$\gamma \geq \gamma'$ と定める。このとき $(\Gamma, \geq)$ は全順序アーベル群である。
\end{lemma}

\begin{proof}
省略するが、補題 \ref{algebra-lemma-valuation-ring-x-or-x-inverse} から容易に従う。
$A = K$ の場合には、唯一の全順序を備えた零群 $\Gamma = \{0\}$ が得られることに注意する。
\end{proof}

\begin{definition}
\label{algebra-definition-value-group}
$A$ を付値環とする。
\begin{enumerate}
\item 補題 \ref{algebra-lemma-valuation-group} の全順序アーベル群 $(\Gamma, \geq)$ を、
付値環 $A$ の {\it 値群}という。
\item 写像 $v : A - \{0\} \to \Gamma$ および $v : K^* \to \Gamma$ を、
$A$ に付随する {\it 付値}という。
\item $\Gamma \cong \mathbf{Z}$ であるとき、付値環 $A$ を
{\it 離散付値環}という。
\end{enumerate}
\end{definition}

\noindent
$\Gamma \cong \mathbf{Z}$ ならば、$1 \geq 0$ を満たすそのような同型は一意である。
このように同型を選ぶと、順序は整数の通常の順序になることに注意する。

\begin{lemma}
\label{algebra-lemma-properties-valuation}
$A$ を付値環とする。付値 $v : A -\{0\} \to \Gamma_{\geq 0}$ は
次の性質をもつ。
\begin{enumerate}
\item $v(a) = 0 \Leftrightarrow a \in A^*$,
\item $v(ab) = v(a) + v(b)$,
\item $a + b \not = 0$ ならば $v(a + b) \geq \min(v(a), v(b))$.
\end{enumerate}
\end{lemma}

\begin{proof}
省略する。
\end{proof}

\begin{lemma}
\label{algebra-lemma-characterize-valuation-ring}
環 $A$ に対して次は同値である。
\begin{enumerate}
\item $A$ は付値環である。
\item $A$ は局所整域であり、$A$ の任意の有限生成イデアルは単項である。
\end{enumerate}
\end{lemma}

\begin{proof}
% P01-E0361: frozen proof omits zero cases before applying the valuation and forming ratios.
$A$ を付値環とし、$f_1, \ldots, f_n \in A$ とする。
$v(f_j)$ のうち $v(f_i)$ が最小となる $i$ を選ぶ。
すると $(f_i) = (f_1, \ldots, f_n)$ である。逆に、$A$ が局所整域であり、
$A$ の任意の有限生成イデアルが単項であるとする。
$f, g \in A$ を取り、$(f, g) = (h)$ と書く。このとき、ある
$a, b, c, d \in A$ に対して $f = ah$, $g = bh$ および $h = cf + dg$ である。
したがって $ac + bd = 1$ であり、$a$ または $b$ のいずれかは単元である。
すなわち $g/f$ または $f/g$ のいずれかは $A$ の元である。
補題 \ref{algebra-lemma-x-or-x-inverse-valuation-ring} により、これは $A$ が付値環であることを示す。
\end{proof}

\begin{lemma}
\label{algebra-lemma-valuation-valuation-ring}
$(\Gamma, \geq)$ を全順序アーベル群とし、$K$ を体とする。
$v : K^* \to \Gamma$ をアーベル群の準同型であって、非零な
$a, b, a + b$ を満たす $a, b \in K$ に対して
$v(a + b) \geq \min(v(a), v(b))$ となるものとする。このとき
$$
A =
\{
x \in K \mid x = 0 \text{ or } v(x) \geq 0
\}
$$
は値群 $\Im(v) \subset \Gamma$ をもつ付値環であり、その極大イデアルは
$$
\mathfrak m =
\{
x \in K \mid x = 0 \text{ or } v(x) > 0
\}
$$
であり、単元群は
$$
A^* =
\{
x \in K^* \mid v(x) = 0
\}.
$$
である。
\end{lemma}

\begin{proof}
省略する。
\end{proof}

\noindent
$(\Gamma, \geq)$ を全順序アーベル群とする。
{\it $\Gamma$ のイデアル}とは、$I \subset \Gamma$ であって、
$I$ のすべての元が $\geq 0$ であり、$\gamma \in I$ かつ
$\gamma' \geq \gamma$ ならば $\gamma' \in I$ となるものをいう。
さらに、$0 \not \in I$ であり、
$\gamma + \gamma' \in I, \gamma, \gamma' \geq 0
\Rightarrow \gamma \in I \text{ or } \gamma' \in I$ ならば、
このイデアルを {\it 素}であるという。

\begin{lemma}
\label{algebra-lemma-ideals-valuation-ring}
$A$ を付値環とする。$A$ のイデアルと $\Gamma$ のイデアルとは
$1 - 1$ に対応する。この全単射は包含を保ち、素イデアルを素イデアルへ写す。
\end{lemma}

\begin{proof}
省略する。
\end{proof}

\begin{lemma}
\label{algebra-lemma-valuation-ring-Noetherian-discrete}
付値環がネーター環であるための必要十分条件は、
それが離散付値環または体であることである。
\end{lemma}

\begin{proof}
$A$ を離散付値環とし、付値 $v : A \setminus \{0\} \to \mathbf{Z}$ を
$\Im(v) = \mathbf{Z}_{\geq 0}$ となるように正規化する。
% P01-E1308: the frozen list of all ideals omits the zero ideal.
補題 \ref{algebra-lemma-ideals-valuation-ring} により、$A$ のイデアルは部分集合
$I_n = \{0\} \cup v^{-1}(\mathbf{Z}_{\geq n})$ である。
$v(x) = n$ を満たす任意の元 $x \in A$ が $I_n$ を生成することは明らかである。
したがって $A$ は単項イデアル整域であり、特にネーター環である。

\medskip\noindent
$A$ を値群 $\Gamma$ をもつネーター付値環とする。
補題 \ref{algebra-lemma-ideals-valuation-ring} により、$\Gamma$ のイデアルに対して
昇鎖条件が成り立つ。$A$ は体でない、すなわち
$\gamma > 0$ を満たす $\gamma \in \Gamma$ が存在すると仮定してよい。
$\gamma > 0$ に対する部分集合 $\gamma + \Gamma_{\geq 0}$ に昇鎖条件を適用すると、
$0$ より大きい最小の元 $\gamma_0$ が存在することが分かる。
$\gamma \in \Gamma$ を $\gamma > 0$ である元とする。
$\gamma$, $\gamma - \gamma_0$, $\gamma - 2\gamma_0$, などの列を考える。
昇鎖条件により、これらがすべて $> 0$ であることはできない。
$\gamma - n \gamma_0$ を最後の $\geq 0$ である項とする。
$\gamma_0$ の最小性により $0 = \gamma - n \gamma_0$ である。
したがって $\Gamma$ は望みどおり巡回群である。
\end{proof}


\section{ネーター環についての続論}
\label{algebra-section-Noetherian-again}


\begin{lemma}
\label{algebra-lemma-Noetherian-basic}
$R$ をネーター環とする。
任意の有限 $R$-加群は有限表示である。
有限 $R$-加群の任意の部分加群は有限である。
有限 $R$-加群の $R$-部分加群に対して昇鎖条件が成り立つ。
\end{lemma}

\begin{proof}
まず、有限 $R$-加群 $M$ の任意の部分加群 $N$ が有限であることを示す。
$M$ の生成元の個数に関する帰納法を用いる。この個数が $1$ ならば、
あるイデアル $I \subset J \subset R$ に対して
$N = J/I \subset M = R/I$ と書けるので、ネーター性の定義から結果が従う。
$M$ の生成元の個数が $1$ より大きいならば、
$M'$ と $M''$ の生成元の個数がより少ない短完全列
$0 \to M' \to M \to M'' \to 0$ を取れる。
$N' = M' \cap N$, $N'' = \Im(N \to M'')$ と置くと、
$N$ に対して同様の短完全列が得られることに注意する。
$N$ の生成元の個数は $N'$ の生成元の個数と $N''$ の生成元の個数との和以下なので、
帰納法の仮定から結果が従う。

\medskip\noindent
$M$ が有限表示であることを示すには、表示 $R^n \to M$ の核に
上の結果を適用すればよい。

\medskip\noindent
$M$ の $R$-部分加群に対する昇鎖条件と、
$M$ の任意の部分加群が有限 $R$-加群であるという条件が同値であることは
よく知られており、容易に示される。証明は省略する。
\end{proof}

\begin{lemma}[Artin-Rees]
\label{algebra-lemma-Artin-Rees}
% P01-E0362: the frozen second hypothesis lacks a finite verb.
$R$ をネーター環とし、$I \subset R$ をイデアルとする。
$N \subset M$ を有限 $R$-加群とする。このとき定数 $c > 0$ が存在して、
すべての $n \geq c$ に対して
$I^n M \cap N  =  I^{n-c}(I^cM \cap N)$ が成り立つ。
\end{lemma}

\begin{proof}
環 $S = R \oplus I \oplus I^2 \oplus \ldots
= \bigoplus_{n \geq 0} I^n$ を考える。規約として $I^0 = R$ とする。
積は $R$ における積によって $I^n \times I^m$ を $I^{n + m}$ へ写す。
$I = (f_1, \ldots, f_t)$ ならば、$S$ はネーター環
$R[X_1, \ldots, X_t]$ の商であることに注意する。
その写像は単項式 $X_1^{e_1}\ldots X_t^{e_t}$ を
$f_1^{e_1}\ldots f_t^{e_t}$ へ写すだけである。したがって $S$ はネーター環である。
同様に、加群 $M \oplus IM \oplus I^2M \oplus \ldots
= \bigoplus_{n \geq 0} I^nM$ を考える。これは有限生成 $S$-加群である。
実際、$x_1, \ldots, x_r$ が $R$ 上 $M$ を生成するならば、
これらは $S$ 上 $\bigoplus_{n \geq 0} I^nM$ も生成する。
次に部分加群 $\bigoplus_{n \geq 0} I^nM \cap N$ を考える。
容易に確かめられるように、これは $S$-部分加群である。
補題 \ref{algebra-lemma-Noetherian-basic} により、これは $S$-加群として有限生成であり、
生成元を $\xi_j \in \bigoplus_{n \geq 0} I^nM \cap N$,
$j = 1, \ldots, s$ とする。各 $\xi_j$ を斉次成分に分解することにより、
各 $\xi_j \in I^{d_j}M \cap N$ がある $d_j$ に対して成り立つと仮定してよい。
% P01-E1309: the frozen maximum need not be positive and is undefined for an empty generating set.
$c = \max\{d_j\}$ と置く。するとすべての $n \geq c$ に対して、
$I^nM \cap N$ の各元は $h_j \in I^{n - d_j}$ によって
$\sum h_j \xi_j$ の形に書ける。これと自明な包含関係
$I^{n-d_j}(I^{d_j}M \cap N) \subset I^{n-c}(I^cM \cap N)$ から補題が従う。
\end{proof}

\begin{lemma}
\label{algebra-lemma-map-AR}
$R$ をネーター環とし、
$0 \to K \to M \xrightarrow{f} N$ をその上の有限生成加群の完全列とする。
$I \subset R$ をイデアルとする。このとき、ある $c$ が存在して、
すべての $n \geq c$ に対して
$$
f^{-1}(I^nN) = K + I^{n-c}f^{-1}(I^cN)
\quad\text{and}\quad
f(M) \cap I^nN \subset f(I^{n - c}M)
$$
が成り立つ。
\end{lemma}

\begin{proof}
$\Im(f) \subset N$ に補題 \ref{algebra-lemma-Artin-Rees} を適用し、
$f : I^{n-c}M \to I^{n-c}f(M)$ が全射であることに注意する。
\end{proof}

\begin{lemma}[Krull の共通部分定理]
\label{algebra-lemma-intersect-powers-ideal-module-zero}
$R$ をネーター局所環とする。$I \subset R$ を真のイデアルとする。
$M$ を有限 $R$-加群とする。このとき $\bigcap_{n \geq 0} I^nM = 0$ である。
\end{lemma}

\begin{proof}
$N = \bigcap_{n \geq 0} I^nM$ と置く。
するとすべての $n \geq 0$ に対して $N = I^nM \cap N$ である。
Artin-Rees の補題 \ref{algebra-lemma-Artin-Rees} により、十分大きなある $n$ に対して
$N = I^nM \cap N \subset IN$ となる。
中山の補題 \ref{algebra-lemma-NAK} により $N = 0$ を得る。
\end{proof}

\begin{lemma}
\label{algebra-lemma-intersection-powers-ideal-module}
$R$ をネーター環とする。$I \subset R$ をイデアルとする。
$M$ を有限 $R$-加群とする。$N = \bigcap_n I^n M$ と置く。
\begin{enumerate}
\item
% P01-E0363: the frozen containment condition is not grammatically attached to the quantified prime.
$I \subset \mathfrak p$ を満たす任意の素イデアル $\mathfrak p$ に対して、
$f \in R$, $f \not \in \mathfrak p$ であって $N_f = 0$ となるものが存在する。
\item $I$ が $R$ の Jacobson 根基に含まれるならば $N = 0$ である。
\end{enumerate}
\end{lemma}

\begin{proof}
(1) の証明。加群 $N$ の生成元を $x_1, \ldots, x_n$ とする
（補題 \ref{algebra-lemma-Noetherian-basic} 参照）。
$I \subset \mathfrak p$ を満たす任意の素イデアル $\mathfrak p$ に対し、
補題 \ref{algebra-lemma-intersect-powers-ideal-module-zero} により、
局所化 $M_{\mathfrak p}$ における $N$ の像は零である。
したがって $x_i$ が $N_{g_i}$ で零に写るような
$g_i \in R$, $g_i \not \in \mathfrak p$ を取れる。
ゆえに $N_{g_1g_2\ldots g_n} = 0$ である。

\medskip\noindent
(2) は (1) と補題 \ref{algebra-lemma-characterize-zero-local} から従う。
\end{proof}

\begin{remark}
\label{algebra-remark-intersection-powers-ideal}
補題 \ref{algebra-lemma-intersect-powers-ideal-module-zero} は特に、
ネーター局所環 $R$ において $I \subset R$ が単位イデアルでなければ
$\bigcap_n I^n = (0)$ であることを意味する。
より一般に、$R$ をネーター環、$I \subset R$ をイデアルとする。
$f \in \bigcap_{n \in \mathbf{N}} I^n$ と仮定する。
すると補題 \ref{algebra-lemma-intersection-powers-ideal-module} により、
$I \subset \mathfrak p$ を満たす任意の素イデアルに対し、
$g \in R$, $g \not \in \mathfrak p$ であって、
$f$ が $R_g$ において零に写るものが存在する。
代数幾何ではこれを「$f$ は $\Spec(R)$ の閉集合 $V(I)$ のある開近傍で零である」
と表現する。
\end{remark}

\begin{lemma}[Artin-Tate]
\label{algebra-lemma-Artin-Tate}
$R$ をネーター環とし、$S$ を有限生成 $R$-代数とする。
$T \subset S$ が $R$-部分代数であって、$S$ が $T$-加群として有限生成ならば、
$T$ は $R$ 上有限型である。
\end{lemma}

\begin{proof}
$S$ を $R$-代数として生成する元 $x_1, \ldots, x_n \in S$ を選ぶ。
$S$ を $T$-加群として生成する $S$ の元 $y_1, \ldots, y_m$ を選ぶ。
したがって $x_i = \sum a_{ij} y_j$ を満たす $a_{ij} \in T$ が存在する。
また、$y_i y_j = \sum b_{ijk} y_k$ を満たす $b_{ijk} \in T$ も存在する。
$a_{ij}$ と $b_{ijk}$ が生成する $R$-部分代数を $T' \subset T$ とする。
これは有限生成 $R$-代数であり、したがってネーター環である。代数
$$
S' = T'[Y_1, \ldots, Y_m]/(Y_i Y_j - \sum b_{ijk} Y_k).
$$
を考える。$S'$ は $T'$ 上有限である。実際、$T'$-加群として
$1, Y_1, \ldots, Y_m$ の類により生成される。
$Y_i$ を $y_i$ へ写す $T'$-代数準同型 $S' \to S$ を考える。
% P01-E0364: the frozen conclusion changes the generator index from i to j.
$a_{ij} \in T'$ なので $x_j$ はこの写像の像に属する。
したがって $S' \to S$ は全射である。ゆえに $S$ も $T'$ 上有限である。
$T'$ はネーター環なので $T \subset S$ は $T'$ 上有限であり、
望みの結論を得る。
\end{proof}


\section{長さ}
\label{algebra-section-length}

\begin{definition}
\label{algebra-definition-length}
$R$ を環とする。任意の $R$-加群 $M$ に対し、$R$ 上の $M$ の
{\it 長さ}を次の式で定める。
% P01-E0365: the frozen adjacent-step condition omits the range of i.
$$
\text{length}_R(M)
=
\sup
\{
n
\mid
\exists\ 0 = M_0 \subset M_1 \subset \ldots \subset M_n = M,
\text{ }M_i \not = M_{i + 1}
\}.
$$
\end{definition}

\noindent
言い換えると、これは部分加群の鎖の長さの上限である。
部分加群の鎖が別の鎖の細分であるという概念は明らかである。
これによって部分加群のすべての鎖からなる集まりに半順序が入り、
$M$ が零でないとき、最小の鎖は $0 = M_0 \subset M_1 = M$ の形である。
$M$ の長さが有限ならば、任意の鎖を極大な鎖に細分できるという明らかな事実に注意する。
しかし、すべての極大な鎖が同じ長さをもつことは、それほど明らかではない
（このことは後で示す）。

\begin{lemma}
\label{algebra-lemma-finite-length-finite}
\begin{slogan}
有限の長さをもつ加群は有限加群である。
\end{slogan}
$R$ を環とする。$M$ を $R$-加群とする。
$\text{length}_R(M) < \infty$ ならば $M$ は有限 $R$-加群である。
\end{lemma}

\begin{proof}
省略する。
\end{proof}

\begin{lemma}
\label{algebra-lemma-length-additive}
\begin{slogan}
長さは短完全列に対して加法的である。
\end{slogan}
$0 \to M' \to M \to M'' \to 0$ が $R$ 上の加群の短完全列ならば、
$M$ の長さは $M'$ と $M''$ の長さの和である。
\end{lemma}

\begin{proof}
$M'$ と $M''$ の長さ $n', n''$ のフィルトレーションが与えられると、
それらに対応する $M$ の長さ $n' + n''$ のフィルトレーションを容易に作れる。
したがって $\text{length}_R M
\geq \text{length}_R M' + \text{length}_R M''$ である。
逆に、$M$ のフィルトレーション
$M_0 \subset M_1 \subset \ldots \subset M_n$ が与えられたとき、
誘導されるフィルトレーション
$M_i' = M_i \cap M'$, $M_i'' = \Im(M_i \to M'')$ を考える。
$\{M'_i\}$（それぞれ $\{M''_i\}$）のフィルトレーションの段数を
$n'$（それぞれ $n''$）とする。
$M_i' = M_{i + 1}'$ かつ $M_i'' = M_{i + 1}''$ ならば
$M_i = M_{i + 1}$ である。よって $n' + n'' \geq n$ となる。
先ほどの結果と合わせれば結論を得る。
\end{proof}

\begin{lemma}
\label{algebra-lemma-length-infinite}
$R$ を極大イデアル $\mathfrak m$ をもつ局所環とする。
$M$ が $R$-加群であり、すべての $n \geq 0$ に対して
$\mathfrak m^n M \not = 0$ ならば、$\text{length}_R(M) = \infty$ である。
言い換えると、$M$ の長さが有限ならば、ある $n$ に対して
$\mathfrak m^nM = 0$ である。
\end{lemma}

\begin{proof}
すべての $n\geq 0$ に対して $\mathfrak m^n M \not = 0$ と仮定する。
$f_1f_2 \ldots f_n x \not = 0$ となるように
$x \in M$ と $f_1, \ldots, f_n \in \mathfrak m$ を選ぶ。
フィルトレーション
$$
0 \subset R f_1 \ldots f_n x \subset R f_1 \ldots f_{n - 1} x
\subset \ldots \subset R x \subset M
$$
の最初の $n$ 段は互いに異なる。例えば
$R f_1 x = R f_1 f_2 x$ ならば、ある $g$ に対して
$f_1 x = g f_1 f_2 x$ となり、したがって $(1 - gf_2) f_1 x = 0$ である。
$1 - gf_2$ は単元なので $f_1 x = 0$ となるが、
これは $x$ および $f_1, \ldots, f_n$ の選び方に矛盾する。
したがって長さは無限である。
\end{proof}

\begin{lemma}
\label{algebra-lemma-length-independent}
$R \to S$ を環準同型とする。$M$ を $S$-加群とする。
常に $\text{length}_R(M) \geq \text{length}_S(M)$ が成り立つ。
$R \to S$ が全射ならば等号が成り立つ。
\end{lemma}

\begin{proof}
% P01-E0366: the frozen proof omits "to" and a finite verb in its concluding sentence.
$M$ の $S$-部分加群によるフィルトレーションは、
$M$ の $R$-部分加群によるフィルトレーションを与える。これで不等式が示された。
さらに $R \to S$ が全射ならば、$M$ の任意の $R$-部分加群は
自動的に $S$-部分加群でもある。したがってこの場合には等号が成り立つ。
\end{proof}

\begin{lemma}
\label{algebra-lemma-dimension-is-length}
$R$ を極大イデアル $\mathfrak m$ をもつ環とする。
$M$ を $\mathfrak m M  =  0$ を満たす $R$-加群とする。
このとき $M$ の $R$-加群としての長さは、
$M$ の $R/\mathfrak m$-ベクトル空間としての次元に等しい。
長さが有限であるための必要十分条件は、$M$ が有限 $R$-加群であることである。
\end{lemma}

\begin{proof}
最初の主張は補題 \ref{algebra-lemma-length-independent} の特別な場合である。
したがって長さが有限であることは、$M$ が $R/\mathfrak m$-ベクトル空間として
有限な基底をもつことと同値であり、さらに $M$ が $R$-加群として
有限個の生成元をもつことと同値である。
\end{proof}

\begin{lemma}
\label{algebra-lemma-length-localize}
$R$ を環とする。$M$ を $R$-加群とする。
$S \subset R$ を積閉集合とする。このとき
$\text{length}_R(M) \geq \text{length}_{S^{-1}R}(S^{-1}M)$ である。
\end{lemma}

\begin{proof}
補題 \ref{algebra-lemma-submodule-localization} により、
任意の部分加群 $N' \subset S^{-1}M$ は、ある $R$-部分加群
$N \subset M$ に対して $S^{-1}N$ の形である。これより補題が従う。
\end{proof}

\begin{lemma}
\label{algebra-lemma-length-finite}
$R$ を有限生成な極大イデアル $\mathfrak m$ をもつ環とする
（例えば $R$ がネーター環の場合）。
$M$ を有限 $R$-加群とし、ある $n$ に対して
$\mathfrak m^n M  =  0$ と仮定する。
このとき $\text{length}_R(M) < \infty$ である。
\end{lemma}

\begin{proof}
フィルトレーション
$0 = \mathfrak m^n M \subset
\mathfrak m^{n-1} M \subset
\ldots \subset \mathfrak m M \subset M$ を考える。
その各部分商は、補題 \ref{algebra-lemma-dimension-is-length} が適用できる
有限生成 $R$-加群である。加法性によって結論を得る
（補題 \ref{algebra-lemma-length-additive} 参照）。
\end{proof}

\begin{definition}
\label{algebra-definition-simple-module}
$R$ を環とする。$M$ を $R$-加群とする。
$M \not = 0$ であり、$M$ の任意の部分加群が $M$ または $0$ に等しいとき、
$M$ は {\it 単純}であるという。
\end{definition}

\begin{lemma}
\label{algebra-lemma-characterize-length-1}
$R$ を環とする。$M$ を $R$-加群とする。次は同値である。
\begin{enumerate}
\item $M$ は単純である。
\item $\text{length}_R(M) = 1$ である。
\item ある極大イデアル $\mathfrak m \subset R$ に対して
$M \cong R/\mathfrak m$ である。
\end{enumerate}
\end{lemma}

\begin{proof}
$\mathfrak m$ を $R$ の極大イデアルとする。
補題 \ref{algebra-lemma-dimension-is-length} により、加群 $R/\mathfrak m$ の長さは $1$ である。
最初の二つの主張の同値性は定義から直ちに従う。
$M$ を単純加群とする。$x \in M$, $x \not = 0$ を選ぶ。
$M$ は単純なので $M = R \cdot x$ である。
$I \subset R$ を $x$ の零化イデアル、すなわち
$I = \{f \in R \mid fx = 0\}$ とする。
写像 $R/I \to M$, $f \bmod I \mapsto fx$ は同型なので、
$R/I$ は単純 $R$-加群である。$R/I \not = 0$ なので $I \not = R$ である。
% P01-E0367: the frozen choice misbinds the maximal ideal; its containment formula is retained.
$I$ を含む極大イデアルを $I \subset \mathfrak m$ となるように選ぶ。
$I \not = \mathfrak m$ ならば、$\mathfrak m /I \subset R/I$ は
非自明な部分加群となり、$R/I$ の単純性に矛盾する。
したがって望みどおり $I = \mathfrak m$ である。
\end{proof}

\begin{lemma}
\label{algebra-lemma-simple-pieces}
$R$ を環とする。$M$ を有限の長さをもつ $R$-加群とする。
$M_i \not = M_{i-1}$, $i = 1, \ldots, n$ を満たす
任意の極大な部分加群の鎖
$$
0 = M_0 \subset M_1 \subset M_2 \subset \ldots \subset M_n = M
$$
を選ぶ。このとき
\begin{enumerate}
\item $n = \text{length}_R(M)$ である。
\item 各 $M_i/M_{i-1}$ は単純である。
\item 各 $M_i/M_{i-1}$ は、ある極大イデアル $\mathfrak m_i$ に対して
$R/\mathfrak m_i$ の形である。
\item 極大イデアル $\mathfrak m \subset R$ が与えられたとき、
$$
\# \{i \mid \mathfrak m_i = \mathfrak m\}
=
\text{length}_{R_{\mathfrak m}} (M_{\mathfrak m}).
$$
が成り立つ。
\end{enumerate}
\end{lemma}

\begin{proof}
$M_i/M_{i-1}$ が単純でなければフィルトレーションを細分できるので、
もとのフィルトレーションは極大でない。したがって $M_i/M_{i-1}$ は単純である。
補題 \ref{algebra-lemma-characterize-length-1} により、各加群 $M_i/M_{i-1}$ の長さは $1$ であり、
ある極大イデアル $\mathfrak m_i$ に対して $R/\mathfrak m_i$ の形をもつ。
長さの加法性（補題 \ref{algebra-lemma-length-additive}）により
$n = \text{length}_R(M)$ を得る。局所化は完全なので、
$$
0 = (M_0)_{\mathfrak m}
\subset (M_1)_{\mathfrak m}
\subset (M_2)_{\mathfrak m}
\subset \ldots
\subset (M_n)_{\mathfrak m} = M_{\mathfrak m}
$$
は $M_{\mathfrak m}$ のフィルトレーションであり、その逐次商は
$(M_i/M_{i-1})_{\mathfrak m}$ である。
したがって最後の主張は、$R$ の極大イデアル $\mathfrak m$, $\mathfrak m'$ に対して
$$
(R/\mathfrak m')_{\mathfrak m}
\cong
\left\{
\begin{matrix}
0 &
\text{if } \mathfrak m \not = \mathfrak m', \\
R_{\mathfrak m}/\mathfrak m R_{\mathfrak m} &
\text{if } \mathfrak m  = \mathfrak m'
\end{matrix}
\right.
$$
であるという事実から直接従う。この確認は読者に委ねる。
\end{proof}

\begin{lemma}
\label{algebra-lemma-pushdown-module}
$A$ を極大イデアル $\mathfrak m$ をもつ局所環とする。
$B$ を極大イデアル $\mathfrak m_i$, $i = 1, \ldots, n$ をもつ半局所環とする。
$A \to B$ を準同型とし、各 $\mathfrak m_i$ が $\mathfrak m$ の上にあり、
$$
[\kappa(\mathfrak m_i) : \kappa(\mathfrak m)] < \infty.
$$
を満たすと仮定する。$M$ を有限の長さをもつ $B$-加群とする。このとき
$$
\text{length}_A(M) = \sum\nolimits_{i = 1, \ldots, n}
[\kappa(\mathfrak m_i) : \kappa(\mathfrak m)]
\text{length}_{B_{\mathfrak m_i}}(M_{\mathfrak m_i}),
$$
特に $\text{length}_A(M) < \infty$ である。
\end{lemma}

\begin{proof}
補題 \ref{algebra-lemma-simple-pieces} のように、$B$-部分加群による極大な鎖
$$
0 = M_0
\subset M_1
\subset M_2
\subset \ldots
\subset M_m = M
$$
を選ぶ。すると各商 $M_j/M_{j - 1}$ は、
ある $i(j) \in \{1, \ldots, n\}$ に対して $\kappa(\mathfrak m_{i(j)})$ と同型である。
さらに補題 \ref{algebra-lemma-dimension-is-length} により
$\text{length}_A(\kappa(\mathfrak m_i)) =
[\kappa(\mathfrak m_i) : \kappa(\mathfrak m)]$ である。
長さの加法性（補題 \ref{algebra-lemma-length-additive}）から補題が従う。
\end{proof}

\begin{lemma}
\label{algebra-lemma-pullback-module}
$A \to B$ を局所環の間の平坦な局所準同型とする。
このとき任意の $A$-加群 $M$ に対して
$$
\text{length}_A(M) \text{length}_B(B/\mathfrak m_AB)
=
\text{length}_B(M \otimes_A B).
$$
が成り立つ。特に $\text{length}_B(B/\mathfrak m_AB) < \infty$ ならば、
$M$ が有限の長さをもつことと $M \otimes_A B$ が有限の長さをもつことは同値である。
\end{lemma}

\begin{proof}
補題 \ref{algebra-lemma-local-flat-ff} により、環準同型 $A \to B$ は忠実平坦である。
したがって $0 = M_0 \subset M_1 \subset \ldots \subset M_n = M$ が
$M$ における長さ $n$ の鎖ならば、対応する鎖
$0 = M_0 \otimes_A B \subset M_1 \otimes_A B \subset
\ldots \subset M_n \otimes_A B = M \otimes_A B$ の長さも $n$ である。
これにより
$\text{length}_A(M) = \infty \Rightarrow
\text{length}_B(M \otimes_A B) = \infty$ が示された。
次に $\text{length}_A(M) < \infty$ と仮定する。
この場合、$M$ は長さ $\ell = \text{length}_A(M)$ のフィルトレーションをもち、
その各商は $A/\mathfrak m_A$ である。
上と同様に論じると、$M \otimes_A B$ は長さ $\ell$ のフィルトレーションをもち、
その各商は $B \otimes_A A/\mathfrak m_A =  B/\mathfrak m_AB$ と同型である。
これより補題が従う。
\end{proof}

\begin{lemma}
\label{algebra-lemma-pullback-transitive}
$A \to B \to C$ を局所環の間の平坦な局所準同型とする。このとき
$$
\text{length}_B(B/\mathfrak m_A B)
\text{length}_C(C/\mathfrak m_B C)
=
\text{length}_C(C/\mathfrak m_A C)
$$
\end{lemma}

\begin{proof}
環準同型 $B \to C$ と $B$-加群 $M = B/\mathfrak m_A B$ に
補題 \ref{algebra-lemma-pullback-module} を適用すれば従う。
\end{proof}


\section{アルティン環}
\label{algebra-section-artinian}

\noindent
アルティン環、特に局所アルティン環は、代数幾何において、
例えば変形理論で重要な役割を果たす。

\begin{definition}
\label{algebra-definition-artinian}
環 $R$ がイデアルに対する降鎖条件を満たすとき、これを {\it アルティン環}という。
\end{definition}

\begin{lemma}
\label{algebra-lemma-finite-dimensional-algebra}
$R$ を体上の有限次元代数とする。このとき $R$ はアルティン環である。
\end{lemma}

\begin{proof}
イデアルに対する降鎖条件は明らかに成り立つ。
\end{proof}

\begin{lemma}
\label{algebra-lemma-artinian-finite-nr-max}
$R$ がアルティン環ならば、$R$ の極大イデアルは有限個しかない。
\end{lemma}

\begin{proof}
$\mathfrak m_i$, $i = 1, 2, 3, \ldots$ が相異なる極大イデアルであるとする。
すると
$\mathfrak m_1 \supset \mathfrak m_1\cap \mathfrak m_2
\supset \mathfrak m_1 \cap \mathfrak m_2 \cap \mathfrak m_3 \supset \ldots$
は無限降鎖である
（中国剰余定理により、すべての写像
$R \to \oplus_{i = 1}^n R/\mathfrak m_i$ は全射だからである）。
\end{proof}

\begin{lemma}
\label{algebra-lemma-artinian-radical-nilpotent}
$R$ をアルティン環とする。$R$ の Jacobson 根基は冪零イデアルである。
\end{lemma}

\begin{proof}
$I \subset R$ を Jacobson 根基とする。
$I \supset I^2 \supset I^3 \supset \ldots$ は降鎖なので、
ある $n$ に対して $I^n = I^{n + 1}$ となる。
$J = \{ x\in R \mid xI^n = 0\}$ と置く。$J = R$ を示す必要がある。
そうでなければ、$J' \not = J$, $J \subset J'$ を満たす極小のイデアルを選ぶ
（降鎖条件により可能である）。すると $J'/J$ は単純 $R$-加群であり、
ある極大イデアル $\mathfrak m$ に対して $R/\mathfrak m$ と同型である
（補題 \ref{algebra-lemma-characterize-length-1} 参照）。
このとき $\mathfrak m I^n$ は $J'$ を零化する。
$I \subset \mathfrak m$ なので $I^{n + 1} = I^n$ も $J'$ を零化する。
ゆえに $J' = J$ となり、矛盾である。
\end{proof}

\begin{lemma}
\label{algebra-lemma-product-local}
極大イデアルが有限個で Jacobson 根基が局所冪零である任意の環は、
その極大イデアルにおける局所化の直積である。また、すべての素イデアルは極大である。
\end{lemma}

\begin{proof}
$R$ を極大イデアル $\mathfrak m_1, \ldots, \mathfrak m_n$ をもつ環とする。
$I = \bigcap_{i = 1}^n \mathfrak m_i$ を $R$ の Jacobson 根基とし、
$I$ は局所冪零と仮定する。$\mathfrak p$ を $R$ の素イデアルとする。
すべての素イデアルは $R$ のすべての冪零元を含むので、
$ \mathfrak p \supset \mathfrak m_1 \cap \ldots \cap \mathfrak m_n$ である。
$\mathfrak m_1 \cap \ldots \cap \mathfrak m_n \supset
\mathfrak m_1 \ldots \mathfrak m_n$ なので、
$\mathfrak p \supset \mathfrak m_1 \ldots \mathfrak m_n$ を得る。
したがってある $i$ に対して $\mathfrak p \supset \mathfrak m_i$ であり、
$\mathfrak p = \mathfrak m_i$ となる。ゆえに $R$ のスペクトルは離散位相空間
$\{\mathfrak m_1, \ldots, \mathfrak m_n\}$ である。
補題 \ref{algebra-lemma-disjoint-implies-product} を $n - 1$ 回適用すると、
$R = R_1 \times \ldots \times R_n$ と書け、各 $i$ に対して
$R_i$ のスペクトルは一点集合となる。
したがって $R_i$ は局所環であり、同補題により $R$ の局所化でもあるから、
望みどおり $R$ の局所環の一つである。
\end{proof}

\begin{lemma}
\label{algebra-lemma-artinian-finite-length}
環 $R$ がアルティン環であるための必要十分条件は、
それ自身の上の加群として有限の長さをもつことである。
このような環 $R$ はアルティン環かつネーター環であり、
$R$ の任意の素イデアルは極大イデアルであり、
$R$ はその極大イデアルにおける局所化の（有限）直積に等しい。
\end{lemma}

\begin{proof}
$R$ がそれ自身の上で有限の長さをもつならば、
イデアルに対する昇鎖条件と降鎖条件をともに満たす。
したがってネーター環かつアルティン環である。
% P01-E0368: the frozen product sentence omits its article.
補題 \ref{algebra-lemma-artinian-finite-nr-max},
\ref{algebra-lemma-artinian-radical-nilpotent}, \ref{algebra-lemma-product-local} により、
任意のアルティン環はその極大イデアルにおける局所化の直積に等しい。

\medskip\noindent
$R$ をアルティン環とする。$R$ がそれ自身の上で有限の長さをもつことを示す。
逐次商が有限の長さをもつ部分加群の鎖を与えればよい。
上で述べたことにより、$R$ は極大イデアル $\mathfrak m$ をもつ局所環と仮定してよい。
補題 \ref{algebra-lemma-artinian-radical-nilpotent} により、
ある $n$ に対して $\mathfrak m^n =0$ である。鎖
$0 = \mathfrak m^n \subset \mathfrak m^{n-1} \subset
\ldots \subset \mathfrak m \subset R$ を考える。
補題 \ref{algebra-lemma-dimension-is-length} により、各部分商
$\mathfrak m^j/\mathfrak m^{j + 1}$ の長さは、
その $\kappa(\mathfrak m)$ 上のベクトル空間としての次元である。
% P01-E0369: the frozen phrase "sub vector spaces" means vector subspaces.
この次元は有限でなければならない。そうでなければベクトル部分空間の無限降鎖が存在し、
それは $R$ のイデアルの無限降鎖に対応するからである。
\end{proof}


\section{本質的に有限型な準同型}
\label{algebra-section-essentially-of-finite-type}

\noindent
有限型環準同型の局所化について、いくつかの簡単な注意を述べる。

\begin{definition}
\label{algebra-definition-essentially-finite-p-t}
$R \to S$ を環準同型とする。
\begin{enumerate}
\item $S$ が有限型 $R$-代数の局所化であるとき、
$R \to S$ は {\it 本質的に有限型}であるという。
\item $S$ が有限表示 $R$-代数の局所化であるとき、
$R \to S$ は {\it 本質的に有限表示}であるという。
\end{enumerate}
\end{definition}

\begin{lemma}
\label{algebra-lemma-composition-essentially-of-finite-type}
本質的に有限型な環準同型のクラスは合成で保たれる。
本質的に有限表示な環準同型についても同様である。
\end{lemma}

\begin{proof}
省略する。
\end{proof}

\begin{lemma}
\label{algebra-lemma-base-change-essentially-of-finite-type}
本質的に有限型な環準同型のクラスは基底変換で保たれる。
本質的に有限表示な環準同型についても同様である。
\end{lemma}

\begin{proof}
省略する。
\end{proof}

\begin{lemma}
\label{algebra-lemma-essentially-of-finite-type-into-artinian-local}
$R \to S$ を環準同型とする。$S$ は極大イデアル $\mathfrak m$ をもつ
局所アルティン環であると仮定する。このとき
\begin{enumerate}
\item $R \to S$ が有限であることと $R \to S/\mathfrak m$ が有限であることは同値である。
\item $R \to S$ が有限型であることと $R \to S/\mathfrak m$ が有限型であることは同値である。
\item $R \to S$ が本質的に有限型であることと、
合成 $R \to S/\mathfrak m$ が本質的に有限型であることは同値である。
\end{enumerate}
\end{lemma}

\begin{proof}
$R \to S$ が有限ならば、補題 \ref{algebra-lemma-finite-transitive} により
$R \to S/\mathfrak m$ は有限である。
逆に $R \to S/\mathfrak m$ が有限であると仮定する。
$S$ はそれ自身の上で有限の長さをもつので
（補題 \ref{algebra-lemma-artinian-finite-length}）、
イデアルによるフィルトレーション
$$
0 \subset I_1 \subset \ldots \subset I_n = S
$$
であって、$S$-加群として $I_i/I_{i - 1} \cong S/\mathfrak m$ となるものを選べる。
したがって $S$ は $R$-部分加群 $I_i$ によるフィルトレーションをもち、
その各逐次商は有限 $R$-加群である。ゆえに補題 \ref{algebra-lemma-extension} により、
$S$ は有限 $R$-加群である。

\medskip\noindent
$R \to S$ が有限型ならば、補題 \ref{algebra-lemma-compose-finite-type} により
$R \to S/\mathfrak m$ は有限型である。
逆に $R \to S/\mathfrak m$ が有限型であると仮定する。
$S/\mathfrak m$ の生成元へ写る $f_1, \ldots, f_n \in S$ を選ぶ。
すると $A = R[x_1, \ldots, x_n] \to S$, $x_i \mapsto f_i$ は環準同型であり、
$A \to S/\mathfrak m$ は全射（特に有限）である。
したがって (1) により $A \to S$ は有限であり、
補題 \ref{algebra-lemma-compose-finite-type} により $R \to S$ は有限型である。

\medskip\noindent
$R \to S$ が本質的に有限型ならば、
補題 \ref{algebra-lemma-composition-essentially-of-finite-type} により
$R \to S/\mathfrak m$ は本質的に有限型である。
逆に $R \to S/\mathfrak m$ が本質的に有限型であると仮定する。
$S/\mathfrak m$ が $R[x_1, \ldots, x_n]/I$ の局所化であるとする。
$\mathfrak m$ を法とする剰余類が、$I$ を法とする
$x_1, \ldots, x_n$ の剰余類に対応するような $f_1, \ldots, f_n \in S$ を選ぶ。
写像 $R[x_1, \ldots, x_n] \to S$, $x_i \mapsto f_i$ を考え、その核を $J$ とする。
$A = R[x_1, \ldots, x_n]/J \subset S$ および
$\mathfrak p = A \cap \mathfrak m$ と置く。
$A/\mathfrak p \subset S/\mathfrak m$ は、$S/\mathfrak m$ における
$R[x_1, \ldots, x_n]/I$ の像に等しいことに注意する。
したがって $\kappa(\mathfrak p) = S/\mathfrak m$ である。
ゆえに (1) により $A_\mathfrak p \to S$ は有限である。
補題 \ref{algebra-lemma-composition-essentially-of-finite-type} により、
$S$ は本質的に有限型であると結論できる。
\end{proof}

\noindent
次の補題は、ここで与える代数的な議論の代わりに、
射影空間の固有性を用いて示すこともできる。

\begin{lemma}
\label{algebra-lemma-localization-at-closed-point-special-fibre}
$\varphi : R \to S$ を本質的に有限型とし、$R$ と $S$ は局所環とする
（ただし $\varphi$ は局所準同型とは限らない）。
このとき、ある $n$ と、$\mathfrak m_R$ の上にある極大イデアル
$\mathfrak m \subset R[x_1, \ldots, x_n]$ が存在して、
$S$ は $R[x_1, \ldots, x_n]_\mathfrak m$ の商の局所化となる。
\end{lemma}

\begin{proof}
$S$ は $R[x_1, \ldots, x_n]$ の商の局所化として書ける。
したがって、ある素イデアル $\mathfrak q \subset R[x_1, \ldots, x_n]$ に対して
$S = R[x_1, \ldots, x_n]_\mathfrak q$ である場合に補題を示せばよい。
$\mathfrak q + \mathfrak m_R R[x_1, \ldots, x_n] \not = R[x_1, \ldots, x_n]$
ならば、補題の主張にあるような極大イデアル $\mathfrak m$ で
$\mathfrak q \subset \mathfrak m$ を満たすものを取れ、結果は明らかである。

\medskip\noindent
$R \to \kappa(\mathfrak q)$ の像を支配する付値環
$A \subset \kappa(\mathfrak q)$ を選ぶ（補題 \ref{algebra-lemma-dominate}）。
% P01-E0370: the frozen antecedent does not explicitly quantify all indexed images.
$x_i$ の像 $\lambda_i \in \kappa(\mathfrak q)$ が $A$ に含まれるならば、
$\mathfrak q$ は $R[x_1, \ldots, x_n] \to A$ による
$\mathfrak m_A$ の逆像に含まれるので、前の場合に帰着する。
したがって、$\lambda_i^{-1} \in A$ であり、かつすべての
$j = 1, \ldots, n$ に対して $\lambda_j/\lambda_i \in A$ となる $i$ が存在する
（$A$ の値群は全順序群であるからである。補題 \ref{algebra-lemma-valuation-group} 参照）。
そこで写像
$$
R[y_0, y_1, \ldots, \hat{y_i}, \ldots, y_n]
\to R[x_1, \ldots, x_n]_\mathfrak q,\quad
y_0 \mapsto 1/x_i,\quad y_j \mapsto x_j/x_i
$$
を考える。$\mathfrak q$ の逆像を
$\mathfrak q' \subset R[y_0, \ldots, \hat{y_i}, \ldots, y_n]$ とする。
$y_0 \not \in \mathfrak q'$ なので、表示した矢印が局所化の間の同型を定めることは容易に分かる。
他方、$y_j$ は $A$ の元へ写るので、最初の段落の結果を
$R[y_0, \ldots, \hat{y_i}, \ldots, y_n]$ に適用できる。
これで証明が終わる。
\end{proof}


\section{K 群}
\label{algebra-section-K-groups}

\noindent
$R$ を環とする。$R$ に付随する二つのアーベル群を導入する。
一つ目を $K'_0(R)$ と書き、これは次の性質をもつ
\footnote{定義は任意の環に対して意味をもつが、$R$ がネーター環でない場合には
ほとんど用いられない。}。
\begin{enumerate}
\item 任意の有限 $R$-加群 $M$ に対して、$K'_0(R)$ の元 $[M]$ が与えられている。
\item 有限 $R$-加群の任意の短完全列 $0 \to M' \to M \to M'' \to 0$ に対して、
関係式 $[M] = [M'] + [M'']$ が成り立つ。
\item 群 $K'_0(R)$ は元 $[M]$ によって生成される。
\item $K'_0(R)$ における生成元 $[M]$ の間のすべての関係式は、
上のような完全列から得られる関係式の $\mathbf{Z}$-線形結合である。
\end{enumerate}
実際の構成には少し手間がかかる。有限生成 $R$-加群全体は真のクラスをなすことに
注意しなければならないからである。
% P01-E0371: the frozen generator construction lets n range over all integers, including negative ones.
しかし、群 $K'_0(R)$ の生成元集合として、$n$ がすべての整数を、
$K$ がすべての部分加群 $K \subset R^n$ を動くときの元 $[R^n/K]$ を取れば、
この問題を克服できる。これらの元に課す関係式の部分群の生成元は、
各項が $R^n/K$ の形をもつ短完全列から得られる関係式とする。
元 $[M]$ は、$M \cong R^n/K$ となる $n$ と $K$ を選び、
$[M] = [R^n/K]$ と置いて定める。詳細は読者に委ねる。

\begin{lemma}
\label{algebra-lemma-length-K0}
$R$ が局所アルティン環ならば、長さの関数は自然なアーベル群準同型
$\text{length}_R : K'_0(R) \to \mathbf{Z}$ を定める。
\end{lemma}

\begin{proof}
任意の有限 $R$-加群の長さは有限である。実際、それは $R^n$ の商であり、
補題 \ref{algebra-lemma-artinian-finite-length} により後者の長さは有限である。
また、長さの関数は加法的である（補題 \ref{algebra-lemma-length-additive} 参照）。
\end{proof}

\noindent
二つ目を $K_0(R)$ と書き、これは次の性質をもつ。
\begin{enumerate}
\item 任意の有限射影 $R$-加群 $M$ に対して、$K_0(R)$ の元 $[M]$ が与えられている。
\item 有限射影 $R$-加群の任意の短完全列 $0 \to M' \to M \to M'' \to 0$ に対して、
関係式 $[M] = [M'] + [M'']$ が成り立つ。
\item 群 $K_0(R)$ は元 $[M]$ によって生成される。
\item $K_0(R)$ におけるすべての関係式は、上のような完全列から得られる関係式の
$\mathbf{Z}$-線形結合である。
\end{enumerate}
この群は上と同じ方法で構成される。

\medskip\noindent
明らかな写像 $K_0(R) \to K'_0(R)$ が存在するが、
一般には同型でないことに注意する。

\begin{example}
\label{algebra-example-K0-field}
$R = k$ が体ならば、明らかに $K_0(k) = K'_0(k) \cong \mathbf{Z}$ であり、
この同型は次元の関数（これは長さの関数でもある）によって与えられる。
\end{example}

\begin{example}
\label{algebra-example-K0-PID}
$R$ を単項イデアル整域とする。$K_0(R) = K'_0(R) = \mathbf{Z}$ であると主張する。
実際、任意の有限射影 $R$-加群は有限自由加群である。
補題 \ref{algebra-lemma-rank} により、有限自由加群の階数は矛盾なく定まる。
有限自由加群の短完全列
$$
0 \to M' \to M \to M'' \to 0
$$
が与えられると、$\text{rank}(M) = \text{rank}(M') + \text{rank}(M'')$ である。
% P01-E0372: the frozen splitting repeats M' as the second summand.
これは、この場合 $M \cong M' \oplus M'$ となるからである
（例えば補題 \ref{algebra-lemma-lift-map} により分裂が存在する）。
よって $K_0(R) = \mathbf{Z}$ を得る。

\medskip\noindent
単項イデアル整域上の加群の構造定理によれば、
任意の有限生成 $R$-加群は
$M = R^{\oplus r} \oplus R/(d_1) \oplus \ldots \oplus R/(d_k)$ の形である。
短完全列
$$
0 \to (d_i) \to R \to R/(d_i) \to 0
$$
を考える。イデアル $(d_i)$ は加群として $R$ と同型である
（$d_i$ を生成元とする自由加群である）から、
$K'_0(R)$ において $[(d_i)] = [R]$ である。
したがって $[R/(d_i)] = [(d_i)]-[R] = 0$ である。
これより、捩れ加群の $K'_0(R)$ における類は零である。
自由部分の階数を用いると $K'_0(R) = \mathbf{Z}$ と同一視でき、
標準準同型 $K_0(R) \to K'_0(R)$ は同型となる。
\end{example}

\begin{example}
\label{algebra-example-K0-polynomial-ring}
$k$ を体とする。このとき $K_0(k[x]) = K'_0(k[x]) = \mathbf{Z}$ である。
$R = k[x]$ は単項イデアル整域なので、これは例 \ref{algebra-example-K0-PID} から従う。
\end{example}

\begin{example}
\label{algebra-example-K0-node}
$k$ を体とし、$R = \{f \in k[x] \mid f(0) = f(1)\}$ と置く
（例 \ref{algebra-example-affine-open-not-standard} と比較せよ）。
この場合、$K_0(R) \cong k^* \oplus \mathbf{Z}$ であるが、
$K'_0(R) = \mathbf{Z}$ である。
\end{example}

\begin{lemma}
\label{algebra-lemma-K0-product}
$R = R_1 \times R_2$ とする。このとき
$K_0(R) = K_0(R_1) \times K_0(R_2)$ および
$K'_0(R) = K'_0(R_1) \times K'_0(R_2)$ である。
\end{lemma}

\begin{proof}
省略する。
\end{proof}

\begin{lemma}
\label{algebra-lemma-K0prime-Artinian}
$R$ を局所アルティン環とする。
補題 \ref{algebra-lemma-length-K0} の写像
$\text{length}_R : K'_0(R) \to \mathbf{Z}$ は同型である。
\end{lemma}

\begin{proof}
省略する。
\end{proof}

\begin{lemma}
\label{algebra-lemma-K0-local}
$(R, \mathfrak m)$ を局所環とする。
任意の有限射影 $R$-加群は有限自由加群である。
$[M] \to \text{rank}_R(M)$ によって定められる写像
$\text{rank}_R : K_0(R) \to \mathbf{Z}$ は矛盾なく定義され、同型である。
\end{lemma}

\begin{proof}
$P$ を有限射影 $R$-加群とする。
$P/\mathfrak m P$ の基底へ写る元 $x_1, \ldots, x_n \in P$ を選ぶ。
中山の補題 \ref{algebra-lemma-NAK} により、これらの元は $P$ を生成する。
$P$ は射影的なので、対応する全射 $u : R^{\oplus n} \to P$ は分裂する。
したがって $Q = \Ker(u)$ と置くと $R^{\oplus n} = P \oplus Q$ である。
これより $Q/\mathfrak m Q = 0$ となり、中山の補題により $Q$ は零である。
このように、任意の有限射影 $R$-加群は有限自由加群であることが分かる。
補題 \ref{algebra-lemma-rank} により、有限自由加群の階数は矛盾なく定まる。
有限自由 $R$-加群の短完全列
$$
0 \to M' \to M \to M'' \to 0
$$
が与えられると、$\text{rank}(M) = \text{rank}(M') + \text{rank}(M'')$ である。
% P01-E0373: the frozen local-ring splitting repeats M' as the second summand.
これは、この場合 $M \cong M' \oplus M'$ となるからである
（例えば補題 \ref{algebra-lemma-lift-map} により分裂が存在する）。
よって $K_0(R) = \mathbf{Z}$ を得る。
\end{proof}

\begin{lemma}
\label{algebra-lemma-K0-and-K0prime-Artinian-local}
$R$ を局所アルティン環とする。可換図式
$$
\xymatrix{
K_0(R) \ar[rr] \ar[d]_{\text{rank}_R} & &
K'_0(R) \ar[d]^{\text{length}_R} \\
\mathbf{Z} \ar[rr]^{\text{length}_R(R)} & &
\mathbf{Z}
}
$$
が存在し、垂直の写像は補題 \ref{algebra-lemma-K0prime-Artinian} と
\ref{algebra-lemma-K0-local} により同型である。
\end{lemma}

\begin{proof}
$P$ を有限射影 $R$-加群とする。
$\text{length}_R(P) = \text{rank}_R(P) \text{length}_R(R)$ を示す必要がある。
補題 \ref{algebra-lemma-K0-local} により、加群 $P$ は有限自由加群である。
したがってある $n \geq 0$ に対して $P \cong R^{\oplus n}$ である。
このとき $\text{rank}_R(P) = n$ であり、長さの加法性
（補題 \ref{algebra-lemma-length-additive}）により
$\text{length}_R(R^{\oplus n}) = n \text{length}_R(R)$ である。
したがって結果が成り立つ。
\end{proof}


\section{次数付き環}
\label{algebra-section-graded}

\noindent
ここでは、{\it 次数付き環}とは、その基礎となるアーベル群の直和分解
$S = \bigoplus_{d \geq 0} S_d$ を備え、
$S_d \cdot S_e \subset S_{d + e}$ を満たす環 $S$ をいう。
負の次数の非零元は認めないことに注意する。
{\it 無縁イデアル}とは、イデアル $S_{+} = \bigoplus_{d > 0} S_d$ のことである。
{\it 次数付き加群}とは、基礎となるアーベル群の直和分解
$M = \bigoplus_{n\in \mathbf{Z}} M_n$ を備え、
$S_d \cdot M_e \subset M_{d + e}$ を満たす $S$-加群 $M$ をいう。
加群については負の次数の非零元も認めることに注意する。
$k > 0$ に対して $S_{-k} = (0)$ と置くことにより、
$S$ を次数付き $S$-加群とみなす。
$M$（それぞれ $S$）の元 $x$（それぞれ $f$）が、
ある $d$ に対して $x \in M_d$（それぞれ $f \in S_d$）を満たすとき、
{\it 斉次}であるという。
{\it 次数付き $S$-加群の写像}とは、$S$-加群の写像
$\varphi : M \to M'$ であって $\varphi(M_d) \subset M'_d$ を満たすものをいう。
次数をずらす写像は認めない。
$M$ から $N$ への次数付き加群の準同型からなる $S_0$-加群を
$\text{GrHom}_0(M, N)$ と書く。

\medskip\noindent
ここで次数付きイデアル、次数付き商環、次数付き部分加群、次数付き商加群、
次数付きテンソル積などの概念が得られる。関連する定義と補題は読者に委ねる。
例えば、次数付き加群の短完全列は各次数において短完全である。

\medskip\noindent
次数付き環 $S$、次数付き $S$-加群 $M$ および $n \in \mathbf{Z}$ に対して、
$M(n)_d = M_{n + d}$ となる次数付き $S$-加群を $M(n)$ と書く。
これを {\it $M$ の $n$ による捻り}という。
特に、射影スキームの研究で重要な役割を果たす加群 $S(n)$,
$n \in \mathbf{Z}$ が得られる。次のような明らかな関手的同型がある。
$(M \oplus N)(n) = M(n) \oplus N(n)$,
$(M \otimes_S N)(n) = M \otimes_S N(n) = M(n) \otimes_S N$.
さらに、$S_0$-加群
$$
\text{GrHom}(M, N) =
\bigoplus\nolimits_{n \in \mathbf{Z}} \text{GrHom}_n(M, N),
\quad
\text{GrHom}_n(M, N) = \text{GrHom}_0(M, N(n)).
$$
に次数付き $S$-加群の構造を定めることができる。
% P01-E0374: the frozen text calls the S-module operation multiplication.
この乗法の定義は省略する。

\begin{lemma}
\label{algebra-lemma-graded-NAK}
$S$ を次数付き環とする。$M$ を次数付き $S$-加群とする。
\begin{enumerate}
\item $S_+M = M$ であり $M$ が有限ならば $M = 0$ である。
\item $N, N' \subset M$ が次数付き部分加群であり、
$M = N + S_+N'$ であって $N'$ が有限ならば $M = N$ である。
\item $N \to M$ が次数付き加群の写像で、
$N/S_+N \to M/S_+M$ が全射であり $M$ が有限ならば、$N \to M$ は全射である。
\item $x_1, \ldots, x_n \in M$ が斉次であり $M/S_+M$ を生成し、
$M$ が有限ならば、$x_1, \ldots, x_n$ は $M$ を生成する。
\end{enumerate}
\end{lemma}

\begin{proof}
(1) の証明。$S$ 上の $M$ の生成元 $y_1, \ldots, y_r$ を選ぶ。
$y_i$ は次数 $d_i$ の斉次元と仮定してよい。
番号を付け替えて $d_r = \min(d_i)$ と仮定してよい。
すると条件 $S_+M = M$ から $y_r = 0$ が従う。
したがって $r$ に関する帰納法により $M = 0$ である。
(2) は $M/N$ に (1) を適用すれば従う。
(3) は部分加群 $\Im(N \to M)$ と $M$ に (2) を適用すれば従う。
(4) は加群の写像
$\bigoplus S(-d_i) \to M$, $(s_1, \ldots, s_n) \mapsto \sum s_i x_i$
に (3) を適用すれば従う。
\end{proof}

\noindent
$S$ を次数付き環とする。$d \geq 1$ を整数とする。
$S^{(d)} = \bigoplus_{n \geq 0} S_{nd}$ と置く。
$S^{(d)}$ は、次数 $n$ の直和因子を
$(S^{(d)})_n = S_{nd}$ とする次数付き環とみなす。
次数付き $S$-加群 $M$ に対して、同様に
$M^{(d)} = \bigoplus_{n \in \mathbf{Z}} M_{nd}$ を考えることができ、
これは次数付き $S^{(d)}$-加群である。

\begin{lemma}
\label{algebra-lemma-uple-generated-degree-1}
$S$ を $S_0$ 上有限生成な次数付き環とする。
このとき、十分に割り切れるすべての $d$ に対し、
代数 $S^{(d)}$ は $S_0$ 上次数 $1$ の元で生成される。
\end{lemma}

\begin{proof}
$S$ は $S_0$ 上 $f_1, \ldots, f_r \in S$ によって生成されるとする。
$f_i$ をその斉次成分で置き換えて、$f_i$ は次数 $d_i > 0$ の斉次元と仮定してよい。
すると $S_n$ の任意の元は、$\sum e_i d_i = n$ を満たす単項式
$f_1^{e_1} \ldots f_r^{e_r}$ の、$S_0$ に係数をもつ線形結合である。
$m$ を $\text{lcm}(d_i)$ の倍数とする。
任意の $N \geq r$ に対して
$$
\sum e_i d_i = N m
$$
ならば、初等的な議論によって、ある $i$ に対して $e_i \geq m/d_i$ である。
したがって次数 $N m$ の任意の単項式は、次数 $m$ の単項式、
すなわち $f_i^{m/d_i}$ と、次数 $(N - 1)m$ の単項式との積である。
これより、$n \geq 2$ に対して次数 $nrm$ の任意の単項式は、
次数 $rm$ の単項式の積である。
ゆえに $S^{(rm)}$ は $S_0$ 上次数 $1$ の元で生成される。
\end{proof}

\begin{lemma}
\label{algebra-lemma-integral-closure-graded}
\begin{reference}
\cite[Theorem 2.3.2]{Huneke-Swanson}
\end{reference}
$R \to S$ を次数付き環の準同型とする。
$S' \subset S$ を $S$ における $R$ の整閉包とする。このとき
$$
S' = \bigoplus\nolimits_{d \geq 0} S' \cap S_d,
$$
すなわち $S'$ は $S$ の次数付き $R$-部分代数である。
\end{lemma}

\begin{proof}
次を示す必要がある。
$s = s_n + s_{n + 1} + \ldots + s_m \in S'$ ならば、
各斉次成分について $s_j \in S'$ が成り立つ。
次数付き環のすべての準同型 $R \to S$ を対象として、
$m - n$ に関する帰納法でこれを示す。
まず、環準同型 $R \to R_0$ と両立する環準同型 $S \to S_0$ があるので、
$s_0$ が $R_0$ 上（したがって $R$ 上）整であることは直ちに分かる。
よって $s$ を $s - s_0$ で置き換え、$n > 0$ と仮定してよい。
$t$ の次数を $0$ とした次数付き環の拡大
$R[t, t^{-1}] \to S[t, t^{-1}]$ を考える。可換図式
$$
\xymatrix{
S[t, t^{-1}] \ar[rr]_{s \mapsto t^{\deg(s)}s} & & S[t, t^{-1}] \\
R[t, t^{-1}] \ar[u] \ar[rr]^{r \mapsto t^{\deg(r)}r} & &  R[t, t^{-1}] \ar[u]
}
$$
があり、水平の写像は環の自己同型である。
したがって $S[t, t^{-1}]$ における $R[t, t^{-1}]$ の整閉包 $C$ はそれ自身へ写る。
ゆえに
$$
t^m(s_n + s_{n + 1} + \ldots + s_m) -
(t^ns_n + t^{n + 1}s_{n + 1} + \ldots + t^ms_m) \in C
$$
であり、帰納法の仮定により $i = n, \ldots, m - 1$ に対して
各 $(t^m - t^i)s_i \in C$ が従う。
任意の環 $A$ と $m > i \geq n > 0$ に対して、
$A[t, t^{-1}]/(t^m - t^i - 1) \cong A[t]/(t^m - t^i - 1) \supset A$
であることに注意する。これは $A[t]/(t^m - t^i - 1)$ において
$t(t^{m - 1} - t^{i - 1}) = 1$ であるからである。
$t^m - t^i$ は $1$ へ写るので、
$i = n, \ldots, m - 1$ に対して、環 $S[t]/(t^m - t^i - 1)$ における
$s_i$ の像は $R[t]/(t^m - t^i - 1)$ 上整である。
$R \to R[t]/(t^m - t^i - 1)$ は有限なので、
推移性により $s_i$ は $R$ 上整である（補題 \ref{algebra-lemma-integral-transitive} 参照）。
最後に、$s_m = s - \sum_{i = n, \ldots, m - 1} s_i$ も $R$ 上整であると結論できる。
\end{proof}


\section{次数付き環の Proj}
\label{algebra-section-proj}

\noindent
$S$ を次数付き環とする。{\it 斉次イデアル}とは、単に
$S$ の次数付き部分加群でもあるイデアル $I \subset S$ のことである。
同値な言い方をすれば、斉次元によって生成されるイデアルである。
さらに同値な条件として、$f \in I$ であり、
$$
f = f_0 + f_1 + \ldots + f_n
$$
が $S$ における $f$ の斉次成分への分解ならば、
各 $i$ に対して $f_i \in I$ となる。
斉次イデアル $\mathfrak p$ が素であることを確かめるには、
$a, b$ が斉次であり $ab \in \mathfrak p$ ならば、
$a \in \mathfrak p$ または $b \in \mathfrak p$ となることを確かめればよい。

\begin{definition}
\label{algebra-definition-proj}
$S$ を次数付き環とする。
$\text{Proj}(S)$ を、$S_{+} \not \subset \mathfrak p$ を満たす
$S$ の斉次素イデアル $\mathfrak p$ の集合と定める。
集合 $\text{Proj}(S)$ は $\Spec(S)$ の部分集合であり、これに誘導位相を入れる。
位相空間 $\text{Proj}(S)$ を、次数付き環 $S$ の {\it 斉次スペクトル}という。
\end{definition}

\noindent
構成から連続写像
$$
\text{Proj}(S) \longrightarrow \Spec(S_0).
$$
が存在することに注意する。

\medskip\noindent
$S = \oplus_{d \geq 0} S_d$ を次数付き環とする。
$f\in S_d$ とし、$d \geq 1$ と仮定する。
$S_{(f)}$ を、$r$ が斉次で $\deg(r) = nd$ を満たすときの
$r/f^n$ の形の元からなる $S_f$ の部分環と定める。
$M$ が次数付き $S$-加群ならば、$S_{(f)}$-加群 $M_{(f)}$ を、
$x$ が次数 $nd$ の斉次元であるときの $x/f^n$ の形の元からなる
$M_f$ の部分加群と定める。

\begin{lemma}
\label{algebra-lemma-Z-graded}
$S$ を、正の次数の可逆な斉次元を含む $\mathbf{Z}$-次数付き環とする。
このとき、$S$ の $\mathbf{Z}$-次数付き素イデアルの集合
$G \subset \Spec(S)$ は、誘導位相に関して $\Spec(S_0)$ へ同相に写る。
\end{lemma}

\begin{proof}
まず逆写像を構成して、この写像が全単射であることを示す。
$f \in S_d$, $d > 0$ を $S$ の可逆元とする。
$\mathfrak p_0$ が $S_0$ の素イデアルならば、
$\mathfrak p_0S$ は $S$ の $\mathbf{Z}$-次数付きイデアルであり、
$\mathfrak p_0S \cap S_0 = \mathfrak p_0$ である。
さらに $a$, $b$ が斉次で $ab \in \mathfrak p_0S$ ならば、
$a^db^d/f^{\deg(a) + \deg(b)} \in \mathfrak p_0$ である。
したがって $a^d/f^{\deg(a)} \in \mathfrak p_0$ または
$b^d/f^{\deg(b)} \in \mathfrak p_0$ であり、言い換えると
$a^d \in \mathfrak p_0S$ または $b^d \in \mathfrak p_0S$ である。
これより $\sqrt{\mathfrak p_0S}$ は $S$ の $\mathbf{Z}$-次数付き素イデアルであり、
$S_0$ との共通部分は $\mathfrak p_0$ となる。

\medskip\noindent
この写像が同相写像であることを示すため、$G \cap D(g)$ の像が開であることを示す。
$g_i \in S_i$ により $g = \sum g_i$ と書くと、
上の議論から $G \cap D(g)$ は開集合 $\bigcup D(g_i^d/f^i)$ の上に写る。
\end{proof}

\noindent
次数 $> 0$ の斉次元 $f \in S$ に対して、
$$
D_{+}(f) = \{ \mathfrak p \in \text{Proj}(S) \mid f \not\in \mathfrak p \}.
$$
と定める。最後に、斉次イデアル $I \subset S$ に対して
$$
V_{+}(I) = \{ \mathfrak p \in \text{Proj}(S) \mid I \subset \mathfrak p \}.
$$
と定める。より一般に、斉次元の任意の集合 $E$、$E \subset S$ に対しても
$V_{+}(E)$ という記号を用いる。

\begin{lemma}[Proj 上の位相]
\label{algebra-lemma-topology-proj}
$S = \oplus_{d \geq 0} S_d$ を次数付き環とする。
\begin{enumerate}
\item 集合 $D_{+}(f)$ は $\text{Proj}(S)$ において開である。
\item $D_{+}(ff') = D_{+}(f) \cap D_{+}(f')$ が成り立つ。
\item $g = g_0 + \ldots + g_m$ を $S$ の元とし、$g_i \in S_i$ とする。このとき
$$
D(g) \cap \text{Proj}(S) =
(D(g_0) \cap \text{Proj}(S))
\cup
\bigcup\nolimits_{i \geq 1} D_{+}(g_i).
$$
\item
$g_0\in S_0$ を次数 $0$ の斉次元とする。このとき
$$
D(g_0) \cap \text{Proj}(S)
=
\bigcup\nolimits_{f \in S_d, \ d\geq 1} D_{+}(g_0 f).
$$
\item 開集合 $D_{+}(f)$ は $\text{Proj}(S)$ の位相の基底をなす。
\item $f \in S$ を正の次数の斉次元とする。
環 $S_f$ は自然な $\mathbf{Z}$-次数付けをもつ。
環準同型 $S \to S_f \leftarrow S_{(f)}$ は同相写像
$$
D_{+}(f)
\leftarrow
\{\mathbf{Z}\text{-graded primes of }S_f\}
\to
\Spec(S_{(f)}).
$$
を誘導する。
\item $\text{Proj}(S)$ が準コンパクトでないような $S$ が存在する。
\item 集合 $V_{+}(I)$ は閉である。
\item 任意の閉部分集合 $T \subset \text{Proj}(S)$ は、
ある斉次イデアル $I \subset S$ に対して $V_{+}(I)$ の形である。
\item 任意の次数付きイデアル $I \subset S$ に対して、
$V_{+}(I) = \emptyset$ であることと $S_{+} \subset \sqrt{I}$ は同値である。
\end{enumerate}
\end{lemma}

\begin{proof}
$D_{+}(f) = \text{Proj}(S) \cap D(f)$ なので、これらの集合は開である。
これで (1) が示された。また、$D(ff') = D(f) \cap D(f')$ より (2) が従う。
同様に、集合 $V_{+}(I) = \text{Proj}(S) \cap V(I)$ は閉である。
これで (8) が示された。

\medskip\noindent
$T \subset \text{Proj}(S)$ を閉集合とする。
すると、あるイデアル $J \subset S$ に対して
$T = \text{Proj}(S) \cap V(J)$ と書ける。
斉次イデアルの定義により、$g \in J$ が $g_d \in S_d$ によって
$g = g_0 + \ldots + g_m$ と書けるならば、
すべての $\mathfrak p \in T$ に対して $g_d \in \mathfrak p$ である。
したがって $J$ の元の斉次成分によって生成されるイデアルを
$I \subset S$ とすれば、$T = V_{+}(I)$ となる。これで (9) が示された。

\medskip\noindent
$g \in S$ に対する $\text{Proj}(S) \cap D(g)$ の式は定義から直ちに従う。
これで (3) が示された。
$\text{Proj}(S) \cap D(g_0)$ の式を考える。右辺が左辺に含まれることは明らかである。
逆の包含については、$\mathfrak p \in \text{Proj}(S)$ で
$g_0 \not \in \mathfrak p$ とする。
正の次数のすべての斉次元 $f$ に対して $g_0f \in \mathfrak p$ ならば、
$S_{+} \subset \mathfrak p$ となり、矛盾である。
これより逆の包含を得る。これで (4) が示された。

\medskip\noindent
標準開集合 $D(g) \subset \Spec(S)$ は $\Spec(S)$ の位相の基底をなすので、
開集合 $D(g) \cap \text{Proj}(S)$ の集まりは位相の基底をなす。
上の式により、$D(g) \cap \text{Proj}(S)$ は開集合 $D_{+}(f)$ の和集合として表せる。
したがって開集合 $D_{+}(f)$ の集まりも位相の基底をなす。これで (5) が示された。

\medskip\noindent
(6) の証明。まず、補題 \ref{algebra-lemma-standard-open} により、
$D_{+}(f)$ を $D(f) = \Spec(S_f)$ の部分集合と誘導位相込みで同一視できる。
環 $S_f$ は $\mathbf{Z}$-次数付けをもつことに注意する。
斉次元は斉次な $r \in S$ による $r/f^n$ の形であり、その次数は
$\deg(r/f^n) = \deg(r) - n\deg(f)$ である。
部分集合 $D_{+}(f)$ は、$\mathbf{Z}$-次数付きイデアル
（すなわち斉次元によって生成されるイデアル）である素イデアル
$\mathfrak p \subset S_f$ にちょうど対応する。
したがって、$S_f$ の $\mathbf{Z}$-次数付き素イデアルの集合が
$\Spec(S_{(f)})$ へ同相に写ることを示せばよい。
これは補題 \ref{algebra-lemma-Z-graded} から従う。

\medskip\noindent
$S = \mathbf{Z}[X_1, X_2, X_3, \ldots]$ とし、
各 $X_i$ の次数を $1$ とする次数付けを入れる。このとき
$$
\text{Proj}(S) = \bigcup\nolimits_{i = 1}^\infty D_{+}(X_i)
$$
が有限な細分をもたないことは容易に分かる。これで (7) が示された。

\medskip\noindent
$I \subset S$ を次数付きイデアルとする。
$\sqrt{I} \supset S_{+}$ ならば $V_{+}(I) = \emptyset$ である。
実際、任意の素イデアル $\mathfrak p \in \text{Proj}(S)$ は、
定義により $S_{+}$ を含まない。
逆に $S_{+} \not \subset \sqrt{I}$ と仮定する。
このとき、元 $f \in S_{+}$ で、$f$ が $I$ を法として冪零でないものが存在する。
明らかに、これは $f$ のある斉次成分が $I$ を法として冪零でないことを意味する。
言い換えると、$f$ は斉次と仮定してよく、以下そうする。
すると $I S_f \not = S_f$ であり、言い換えると $(S/I)_f$ は零でない。
したがって $(S/I)_{(f)} \not = 0$ である。
これは $(S/I)_f$ へ写る環だからである。
素イデアル $\mathfrak q \subset (S/I)_{(f)}$ を選ぶ。
これは、無縁イデアル $(S/I)_{+}$ を含まない $S/I$ の次数付き素イデアルに対応する。
さらにこれは、望みどおり $I$ を含み $S_{+}$ を含まない
$S$ の次数付き素イデアル $\mathfrak p$ に対応する。
これで (10) が示され、証明が終わる。
\end{proof}

\begin{example}
\label{algebra-example-proj-polynomial-ring-1-variable}
$R$ を環とする。$S = R[X]$, $\deg(X) = 1$ ならば、自然な写像
$\text{Proj}(S) \to \Spec(R)$ は全単射であり、実際には同相写像である。
実際、$\mathfrak p \in \text{Proj}(S)$ とすると、
$S_{+} \not \subset \mathfrak p$ より $X \not \in \mathfrak p$ である。
したがって $a \in R$, $n > 0$ に対して $aX^n \in \mathfrak p$ ならば
$a \in \mathfrak p$ である。これより
$\mathfrak p_0 = \mathfrak p \cap R$ と置けば $\mathfrak p = \mathfrak p_0S$ である。
\end{example}

\noindent
$\mathfrak p \in \text{Proj}(S)$ に対し、
$S_{(\mathfrak p)}$ を次のような環と定める。
その元は、$r, f \in S$ が同じ次数の斉次元であり $f \not\in \mathfrak p$ を満たすときの
分数 $r/f$ である。通常どおり、$r/f = r'/f'$ であるとは、
ある斉次元 $f'' \in S$, $f'' \not \in \mathfrak p$ が存在して
$f''(rf' - r'f) = 0$ となることをいう。
次数付き $S$-加群 $M$ に対し、$M_{(\mathfrak p)}$ を次の
$S_{(\mathfrak p)}$-加群とする。
その元は、$x \in M$ と $f \in S$ が同じ次数の斉次元であり
$f \not \in \mathfrak p$ を満たすときの分数 $x/f$ である。
$x/f = x'/f'$ であるとは、ある斉次元 $f'' \in S$,
$f'' \not \in \mathfrak p$ が存在して $f''(xf' - x'f) = 0$ となることをいう。

\begin{lemma}
\label{algebra-lemma-proj-prime}
$S$ を次数付き環とする。$M$ を次数付き $S$-加群とする。
$\mathfrak p$ を $\text{Proj}(S)$ の元とする。
$f \in S$ を正の次数の斉次元で、$f \not \in \mathfrak p$、
すなわち $\mathfrak p \in D_{+}(f)$ を満たすものとする。
補題 \ref{algebra-lemma-topology-proj} のように、
$\mathfrak p$ に対応する $\Spec(S_{(f)})$ の元を
$\mathfrak p' \subset S_{(f)}$ とする。このとき
$S_{(\mathfrak p)} = (S_{(f)})_{\mathfrak p'}$ であり、これと両立して
$M_{(\mathfrak p)} = (M_{(f)})_{\mathfrak p'}$ である。
\end{lemma}

\begin{proof}
写像 $\psi : M_{(\mathfrak p)} \to (M_{(f)})_{\mathfrak p'}$ を定める。
$x/g \in M_{(\mathfrak p)}$ に対して
$$
\psi(x/g) = (x g^{\deg(f) - 1}/f^{\deg(x)})/(g^{\deg(f)}/f^{\deg(g)}).
$$
と置く。$\deg(x) = \deg(g)$ かつ
$g^{\deg(f)}/f^{\deg(g)} \not \in \mathfrak p'$ なので、この式は意味をもつ。
$\psi$ が矛盾なく定義され、加群の写像であり、同型であることの確認は省略する。
ヒント：逆写像は $(x/f^n)/(g/f^m)$ を $(xf^m)/(g f^n)$ へ写す。
\end{proof}

\noindent
次は補題 \ref{algebra-lemma-silly} の次数付き版である。

\begin{lemma}
\label{algebra-lemma-graded-silly}
$S$ を次数付き環、$\mathfrak p_i$, $i = 1, \ldots, r$ を斉次素イデアル、
$I \subset S_{+}$ を次数付きイデアルとする。
すべての $i$ に対して $I \not\subset \mathfrak p_i$ と仮定する。
このとき、正の次数の斉次元 $x\in I$ で、
すべての $i$ に対して $x\not\in \mathfrak p_i$ となるものが存在する。
\end{lemma}

\begin{proof}
$\mathfrak p_i$ の間に包含関係はないと仮定してよい。
$r = 1$ では結果は成り立つ。$r - 1$ について結果が成り立つと仮定する。
すべての $i = 1, \ldots, r - 1$ に対して $x \not \in \mathfrak p_i$ となる、
正の次数の斉次元 $x \in I$ を選ぶ。
$x \not\in \mathfrak p_r$ ならば終わりである。そこで $x \in \mathfrak p_r$ と仮定する。
$I \mathfrak p_1 \ldots \mathfrak p_{r-1} \subset \mathfrak p_r$ ならば
$I \subset \mathfrak p_r$ となり、矛盾である。
$y \not \in \mathfrak p_r$ となる斉次元
$y \in I\mathfrak p_1 \ldots \mathfrak p_{r-1}$ を選ぶ。
すると $x^{\deg(y)} + y^{\deg(x)}$ が求める元である。
\end{proof}

\begin{lemma}
\label{algebra-lemma-smear-out}
$S$ を次数付き環とする。$\mathfrak p \subset S$ を素イデアルとする。
$\mathfrak p$ の斉次元によって生成される $S$ の斉次イデアルを $\mathfrak q$ とする。
このとき $\mathfrak q$ は $S$ の素イデアルである。
\end{lemma}

\begin{proof}
$\mathfrak q$ が素であることを示すには、$f, g \in S$ が斉次で
$fg \in \mathfrak q$ ならば $f$ または $g$ が $\mathfrak q$ に属することを
確かめればよい。
% P01-E0375: the frozen proof reverses the containment between p and q.
$\mathfrak p \subset \mathfrak q$ なので $fg \in \mathfrak p$ である。
$\mathfrak p$ は素であるから、$f \in \mathfrak p$ または $g \in \mathfrak p$ である。
$f$ と $g$ は斉次なので、$f \in \mathfrak q$ または $g \in \mathfrak q$ であることは明らかである。
\end{proof}

\begin{lemma}
\label{algebra-lemma-graded-ring-minimal-prime}
$S$ を次数付き環とする。
\begin{enumerate}
\item $S$ の任意の極小素イデアルは $S$ の斉次イデアルである。
\item 斉次イデアル $I \subset S$ が与えられたとき、
$I$ の上の任意の極小素イデアルは斉次である。
\end{enumerate}
\end{lemma}

\begin{proof}
最初の主張は、補題 \ref{algebra-lemma-smear-out} で構成した素イデアル $\mathfrak q$ が
$\mathfrak q \subset \mathfrak p$ を満たすことから従う。
二つ目は、$S/I$ を考えて最初の主張を適用すれば従う。
\end{proof}

\begin{lemma}
\label{algebra-lemma-dehomogenize-finite-type}
$R$ を環とする。$S$ を次数付き $R$-代数とする。
$f \in S_{+}$ を斉次元とする。$S$ は $R$ 上有限型と仮定する。このとき
\begin{enumerate}
\item 環 $S_{(f)}$ は $R$ 上有限型である。
\item 任意の有限次数付き $S$-加群 $M$ に対して、
加群 $M_{(f)}$ は有限 $S_{(f)}$-加群である。
\end{enumerate}
\end{lemma}

\begin{proof}
$S$ を $R$-代数として生成する $f_1, \ldots, f_n \in S$ を選ぶ。
各 $f_i$ をその斉次成分に分解することにより、各 $f_i$ は斉次と仮定してよい。
$S_{(f)}$ の元は、$\lambda_{e_1 \ldots e_n} \in R$ を係数とする
$$
\sum\nolimits_{e\deg(f) =
\sum e_i\deg(f_i)} \lambda_{e_1 \ldots e_n} f_1^{e_1} \ldots f_n^{e_n}/f^e
$$
の形の和である。したがって $S_{(f)}$ は $R$-代数として、
$e\deg(f) = \sum e_i\deg(f_i)$ を満たす
$f_1^{e_1} \ldots\allowbreak f_n^{e_n} /f^e$ によって生成される。
$e_i \geq \deg(f)$ ならば、これを
$$
f_1^{e_1} \ldots f_n^{e_n}/f^e =
f_i^{\deg(f)}/f^{\deg(f_i)} \cdot
f_1^{e_1} \ldots f_i^{e_i - \deg(f)} \ldots f_n^{e_n}/f^{e - \deg(f_i)}
$$
と書ける。したがって必要なのは、元 $f_i^{\deg(f)}/f^{\deg(f_i)}$ と、
$e \deg(f) = \sum e_i \deg(f_i)$ および $e_i < \deg(f)$ を満たす元
$f_1^{e_1} \ldots f_n^{e_n} /f^e$ だけである。
これは有限なリストなので (1) が成り立つ。

\medskip\noindent
(2) を示すため、$M$ が斉次元 $x_1, \ldots, x_m$ で生成されるとする。
上と同様に論じると、$M_{(f)}$ は $S_{(f)}$-加群として、
$e \deg(f) = \sum e_i \deg(f_i) + \deg(x_j)$ および
$e_i < \deg(f)$ を満たす $f_1^{e_1} \ldots f_n^{e_n} x_j /f^e$ の形の
有限個の元によって生成される。
\end{proof}

\begin{lemma}
\label{algebra-lemma-homogenize}
$R$ を環とする。
$R'$ を有限型 $R$-代数、$M$ を有限 $R'$-加群とする。
次数付き $R$-代数 $S$、次数付き $S$-加群 $N$、
および次数 $1$ の斉次元 $f \in S$ で、次を満たすものが存在する。
\begin{enumerate}
\item $R' \cong S_{(f)}$ かつ（加群として）$M \cong N_{(f)}$ である。
\item $S_0 = R$ であり、$S$ は $R$ 上次数 $1$ の有限個の元で生成される。
\item $N$ は有限 $S$-加群である。
\end{enumerate}
\end{lemma}

\begin{proof}
あるイデアル $I$ に対して $R' = R[x_1, \ldots, x_n]/I$ と書ける。
% P01-E1310: minimal-degree homogenization of constant relations can change S_0.
元 $g \in R[x_1, \ldots, x_n]$ に対し、
$g = \tilde g(1, x_1, \ldots, x_n)$ を満たす最小次数の斉次元を
$\tilde g \in R[X_0, \ldots, X_n]$ と書く。
すべての元 $\tilde g$, $g \in I$ によって生成されるイデアルを
$\tilde I \subset R[X_0, \ldots, X_n]$ とする。
$S = R[X_0, \ldots, X_n]/\tilde I$ と置き、$S$ における $X_0$ の像を $f$ と書く。
構成により同型
$$
S_{(f)} \longrightarrow R', \quad
X_i/X_0 \longmapsto x_i.
$$
が得られる。加群 $M$ に対して同様のことをするため、表示
$$
M = (R')^{\oplus r}/\sum\nolimits_{j \in J} R'k_j
$$
を選び、$k_j = (k_{1j}, \ldots, k_{rj})$ とする。
% P01-E0376: polynomial lifts of the quotient-algebra matrix entries are not specified.
$d_{ij} = \deg(\tilde k_{ij})$ と置き、
$d_j = \max\{d_{ij}\}$ と置く。
$K_{ij} = X_0^{d_j - d_{ij}}\tilde k_{ij}$ と置けば、これは次数 $d_j$ の斉次元である。
この記号のもとで
$$
N = \Coker\Big(
\bigoplus\nolimits_{j \in J} S(-d_j) \xrightarrow{(K_{ij})} S^{\oplus r}
\Big)
$$
と置けばよい。いくつかの詳細は省略する。
\end{proof}


\section{ネーター次数付き環}
\label{algebra-section-noetherian-graded}

\noindent
ヒルベルト多項式に関する事項を含め、ネーター次数付き環の理論を少し述べる。

\begin{lemma}
\label{algebra-lemma-S-plus-generated}
$S$ を次数付き環とする。斉次元 $f_i \in S_{+}$ の集合が
$S$ を $S_0$ 上の代数として生成することと、
$S_{+}$ を $S$ のイデアルとして生成することは同値である。
\end{lemma}

\begin{proof}
$f_i$ が $S$ を $S_0$ 上の代数として生成するならば、
$S_{+}$ の任意の元は $f_i$ に関する定数項のない多項式である。
したがって $S_{+}$ はイデアルとして $f_i$ によって生成される。
逆に $S_{+} = \sum Sf_i$ と仮定する。
$S$ の任意の元 $f$ が、$S_0$ に係数をもつ $f_i$ の多項式として
書けることを示す。斉次元について示せば十分である。
$f$ の次数を $d$ とする。$d$ に関する帰納法を用いる。
$d = 0$ の場合は明らかである。$d > 0$ ならば $f \in S_{+}$ なので、
ある $g_i \in S$ により $f = \sum g_i f_i$ と書ける。
$S$ は次数付きなので、$g_i$ を次数 $d - \deg(f_i)$ の斉次成分で置き換えてよい。
帰納法により各 $g_i$ は $f_i$ の多項式であり、主張が従う。
\end{proof}

\begin{lemma}
\label{algebra-lemma-graded-Noetherian}
次数付き環 $S$ がネーターであることと、
$S_0$ がネーターであり $S_{+}$ が $S$ のイデアルとして有限生成であることは同値である。
\end{lemma}

\begin{proof}
$S$ がネーターならば $S_0 = S/S_{+}$ はネーターであり、
$S_{+}$ は有限生成であることは明らかである。
逆に、$S_0$ はネーターであり、$S_{+}$ は $S$ のイデアルとして有限生成と仮定する。
生成元 $S_{+} = (f_1, \ldots, f_n)$ を選ぶ。
$f_i$ を斉次成分に分解することにより、各 $f_i$ は斉次と仮定してよい。
補題 \ref{algebra-lemma-S-plus-generated} により、
% P01-E0377: the frozen polynomial-variable list omits the comma after ldots.
$X_i$ を $f_i$ へ写す $S_0[X_1, \ldots X_n] \to S$ は全射である。
したがって補題 \ref{algebra-lemma-Noetherian-permanence} により $S$ はネーターである。
\end{proof}

\begin{definition}
\label{algebra-definition-numerical-polynomial}
$A$ をアーベル群とする。
十分大きいすべての整数 $n$ に対して定義された関数
$f : n \mapsto f(n) \in A$ が {\it 数値的多項式}であるとは、
$r \geq 0$ および元 $a_0, \ldots, a_r\in A$ が存在して、
すべての $n \gg 0$ に対して
$$
f(n) = \sum\nolimits_{i = 0}^r \binom{n}{i} a_i
$$
となることをいう。
\end{definition}

\noindent
二項係数を用いる理由は、整数点での値がすべて整数である任意の多項式
$P \in \mathbf{Q}[T]$ が、$a_i \in \mathbf{Z}$ によって
$P(T) = \sum a_i \binom{T}{i}$ の形の和に等しくなるという初等的な事実にある。
特に、式 $\binom{T + 1}{i + 1}$ もこの形であることに注意する。

\begin{lemma}
\label{algebra-lemma-numerical-polynomial-functorial}
$A \to A'$ がアーベル群の準同型であり、
$f : n \mapsto f(n) \in A$ が数値的多項式ならば、その合成も数値的多項式である。
\end{lemma}

\begin{proof}
これは定義から直ちに従う。
\end{proof}

\begin{lemma}
\label{algebra-lemma-numerical-polynomial}
$f: n \mapsto f(n) \in A$ が十分大きいすべての $n$ に対して定義され、
$n \mapsto f(n) - f(n-1)$ が数値的多項式であると仮定する。
このとき $f$ は数値的多項式である。
\end{lemma}

\begin{proof}
すべての $n \gg 0$ に対して
$f(n) - f(n-1) = \sum\nolimits_{i = 0}^r \binom{n}{i} a_i$ とする。
$g(n) = f(n) - \sum\nolimits_{i = 0}^r \binom{n + 1}{i + 1} a_i$ と置く。
すると、すべての $n \gg 0$ に対して $g(n) - g(n-1) = 0$ である。
したがって $g$ は最終的に定数となり、その値を $a_{-1}$ とする。
$a_{-1} + \sum\nolimits_{i = 0}^r \binom{n + 1}{i + 1} a_i$
が必要な形をもつことの確認は読者に委ねる（補題の前の注意を参照）。
\end{proof}

\begin{lemma}
\label{algebra-lemma-graded-module-fg}
$M$ が有限生成次数付き $S$-加群であり、$S$ が $S_0$ 上有限生成ならば、
各 $M_n$ は有限 $S_0$-加群である。
\end{lemma}

\begin{proof}
$M$ の生成元を $m_i$、$S$ の生成元を $f_i$ とする。
斉次成分を取ることにより、$m_i$ と $f_i$ は斉次と仮定でき、
さらに $f_i \in S_{+}$ と仮定してよい。
この場合、各 $M_n$ は、次数が $n$ である「単項式」
$\prod f_i^{e_i} m_j$ によって $S_0$ 上生成されることが明らかである。
\end{proof}

\begin{proposition}
\label{algebra-proposition-graded-hilbert-polynomial}
$S$ をネーター次数付き環、$M$ を有限次数付き $S$-加群とする。
関数
$$
\mathbf{Z} \longrightarrow K'_0(S_0), \quad
n \longmapsto [M_n]
$$
を考える（補題 \ref{algebra-lemma-graded-module-fg} 参照）。
$S_{+}$ が次数 $1$ の元で生成されるならば、この関数は数値的多項式である。
\end{proposition}

\begin{proof}
$S_1$ の生成元の最小個数に関する帰納法で示す。
この個数が $0$ ならば、すべての $n \gg 0$ に対して $M_n = 0$ であり、結果は成り立つ。
帰納法の段階を示すため、最小生成系の一つの元 $x\in S_1$ を取り、
次数付き環 $S/(x)$ には帰納法の仮定が適用できるものとする。

\medskip\noindent
まず、$x$ が $M$ 上冪零である場合に結果が成り立つことを示す。
これは $x^r M  = 0$ となる最小の整数 $r$ に関する帰納法で行う。
$r = 1$ ならば $M$ は $S/xS$ 上の加群なので、
結果はもう一方の帰納法の仮定によって成り立つ。
% P01-E0378: the frozen proof leaves the submodules and their exponents implicit.
$r > 1$ ならば、整数 $r', r''$ が $r$ より真に小さくなるような
短完全列 $0 \to M' \to M \to M'' \to 0$ を見つけられる。
したがって $M''$ と $M'$ については結果が分かっている。
ゆえに $K'_0(S_0)$ における関係
$
[M_d]  = [M'_d] + [M''_d]
$
から $M$ についても結果を得る。

\medskip\noindent
$x$ が $M$ 上冪零でないならば、$x$ が冪零となる最大の部分加群を
$M' \subset M$ とする。完全列 $0 \to M' \to M \to M/M' \to 0$ を考えると、
やはり $M/M'$ について結果を示せば十分であることが分かる。
言い換えると、$x$ による乗法が単射であると仮定してよい。

\medskip\noindent
$\overline{M} = M/xM$ と置く。
写像 $x : M \to M$ は $M_d$ を $M_d$ へ写さないので、
次数付き $S$-加群の写像{\it ではない}ことに注意する。
すなわち各 $d$ に対して、次の短完全列がある。
$$
0 \to M_d \xrightarrow{x} M_{d + 1} \to \overline{M}_{d + 1} \to 0
$$
これより $[M_{d + 1}] - [M_d] = [\overline{M}_{d + 1}]$ が従う。
したがって補題 \ref{algebra-lemma-numerical-polynomial} により結論を得る。
\end{proof}

\begin{remark}
\label{algebra-remark-period-polynomial}
$S$ は依然としてネーターだが、$S$ が次数 $1$ の元で生成されない場合、
次数付き $S$-加群に付随する関数は周期的多項式となる
（すなわち、ある $n$ に対し、整数の $n$ を法とする各合同類上で数値的多項式となる）。
\end{remark}

\begin{example}
\label{algebra-example-hilbert-function}
$S = k[X_1, \ldots, X_d]$ とする。
例 \ref{algebra-example-K0-field} により $K_0(k) = K'_0(k) = \mathbf{Z}$ と同一視できる。
したがって任意の有限生成次数付き $k[X_1, \ldots, X_d]$-加群から、
数値的多項式 $n \mapsto \dim_k(M_n)$ が得られる。
\end{example}

\begin{lemma}
\label{algebra-lemma-quotient-smaller-d}
$k$ を体とする。
$I \subset k[X_1, \ldots, X_d]$ を非零の次数付きイデアルとする。
$M = k[X_1, \ldots, X_d]/I$ と置く。
このとき数値的多項式 $n \mapsto \dim_k(M_n)$
（例 \ref{algebra-example-hilbert-function} 参照）の次数は $ < d - 1$ である
（$d = 1$ ならば零である）。
\end{lemma}

\begin{proof}
次数付き加群 $k[X_1, \ldots, X_d]$ に付随する数値的多項式は
$n \mapsto \binom{n - 1 + d}{d - 1}$ である。
% P01-E0379: the frozen source uses literal >> rather than gg.
次数 $e$ の任意の非零斉次元 $f \in I$ と任意の次数 $n >> e$ に対して
$I_n \supset f \cdot k[X_1, \ldots, X_d]_{n-e}$ であり、
したがって $\dim_k(I_n) \geq \binom{n - e - 1 + d}{d - 1}$ である。
ゆえに
$\dim_k(M_n) \leq \binom{n - 1 + d}{d - 1} - \binom{n - e - 1 + d}{d - 1}$
を得る。最後の式の次数は $ < d - 1$ である
（$d = 1$ ならば零である）から、主張が従う。
\end{proof}


\section{ネーター局所環}
\label{algebra-section-Noetherian-local}

\noindent
この節を通して $(R, \mathfrak m, \kappa)$ はネーター局所環とする。
この節では加群のヒルベルト関数の理論を展開する。
$M$ を有限 $R$-加群とする。$M$ の {\it ヒルベルト関数}を、
すべての整数 $n \geq 0$ に対して定義される関数
$$
\varphi_M : n
\longmapsto
\text{length}_R(\mathfrak m^nM/{\mathfrak m}^{n + 1}M)
$$
と定める。もう一つの重要な不変量は、すべての整数 $n \geq 0$ に対して
定義される関数
$$
\chi_M : n
\longmapsto
\text{length}_R(M/{\mathfrak m}^{n + 1}M)
$$
である。補題 \ref{algebra-lemma-length-additive} により
$$
\chi_M(n) = \sum\nolimits_{i = 0}^n \varphi_M(i).
$$
が成り立つことに注意する。
この構成には、定義イデアルを用いる変種がある。

\begin{definition}
\label{algebra-definition-ideal-definition}
$(R, \mathfrak m)$ をネーター局所環とする。
$\sqrt{I} = \mathfrak m$ を満たすイデアル $I \subset R$ を
{\it $R$ の定義イデアル}という。
\end{definition}

\noindent
$I \subset R$ を定義イデアルとする。
$R$ はネーターなので、これはある $r$ に対して
$\mathfrak m^r \subset I$ となることを意味する
（補題 \ref{algebra-lemma-Noetherian-power} 参照）。
% P01-E0380: the source's indefinite article before finite length is nonstandard.
したがって $I$ の冪によって零化される任意の有限 $R$-加群は有限の長さをもつ
（補題 \ref{algebra-lemma-length-finite} 参照）。
ゆえに、すべての $n \geq 0$ に対して
$$
\varphi_{I, M}(n) = \text{length}_R(I^nM/I^{n + 1}M)
\quad\text{and}\quad
\chi_{I, M}(n) = \text{length}_R(M/I^{n + 1}M)
$$
と定めることができる。ここでも
$$
\chi_{I, M}(n) = \sum\nolimits_{i = 0}^n \varphi_{I, M}(i).
$$
が成り立つ。

\begin{lemma}
\label{algebra-lemma-differ-finite}
$M' \subset M$ を、商が有限の長さをもつ有限 $R$-加群とする。
% P01-E0381: the frozen English has singular/plural disagreement in the constants.
このとき定数 $c_1, c_2$ が存在し、すべての $n \geq c_2$ に対して
$$
c_1 + \chi_{I, M'}(n - c_2) \leq \chi_{I, M}(n) \leq
c_1 + \chi_{I, M'}(n)
$$
が成り立つ。
\end{lemma}

\begin{proof}
$M/M'$ は有限の長さをもつので、$I^{c_2}M \subset M'$ となる
$c_2 \geq 0$ が存在する。$c_1 = \text{length}_R(M/M')$ と置く。
$n \geq c_2$ に対して
\begin{eqnarray*}
\chi_{I, M}(n)
& = &
\text{length}_R(M/I^{n + 1}M) \\
& = &
c_1 + \text{length}_R(M'/I^{n + 1}M) \\
& \leq &
c_1 + \text{length}_R(M'/I^{n + 1}M') \\
& = &
c_1 + \chi_{I, M'}(n)
\end{eqnarray*}
である。一方、$I^{c_2}M \subset M'$ なので、
$n \geq c_2$ に対して $I^nM \subset I^{n - c_2}M'$ である。
したがって $n \geq c_2$ に対して
\begin{eqnarray*}
\chi_{I, M}(n)
& = &
\text{length}_R(M/I^{n + 1}M) \\
& = &
c_1 + \text{length}_R(M'/I^{n + 1}M) \\
& \geq &
c_1 + \text{length}_R(M'/I^{n + 1 - c_2}M') \\
& = &
c_1 + \chi_{I, M'}(n - c_2)
\end{eqnarray*}
を得て、証明が終わる。
\end{proof}

\begin{lemma}
\label{algebra-lemma-hilbert-ses}
$0 \to M' \to M \to M'' \to 0$ を有限 $R$-加群の短完全列とする。
このとき、有限の余長 $l$ をもつ部分加群 $N \subset M'$ と
$c \geq 0$ が存在して、すべての $n \geq c$ に対して
$$
\chi_{I, M}(n) = \chi_{I, M''}(n) + \chi_{I, N}(n - c) + l
$$
および
$$
\varphi_{I, M}(n) = \varphi_{I, M''}(n) + \varphi_{I, N}(n - c)
$$
が成り立つ。
\end{lemma}

\begin{proof}
$M/I^nM \to M''/I^nM''$ は全射であり、その核は
$M' / M' \cap I^nM$ であることに注意する。
アルティン--リースの補題 \ref{algebra-lemma-Artin-Rees} により、
$M' \cap I^nM = I^{n - c}(M' \cap I^cM)$ となる定数 $c$ が存在する。
$N = M' \cap I^cM$ と書く。$I^c M' \subset N \subset M'$ に注意すると、
$n \geq c$ に対して
$\text{length}_R(M' / M' \cap I^nM)
= \text{length}_R(M'/N) + \text{length}_R(N/I^{n - c}N)$
である。短完全列
$$
0 \to M' / M' \cap I^nM \to M/I^nM \to M''/I^nM'' \to 0
$$
と長さの加法性（補題 \ref{algebra-lemma-length-additive}）から、
$n \geq c$ に対して等式
$$
\chi_{I, M}(n - 1)
=
\chi_{I, M''}(n - 1)
+
\chi_{I, N}(n - c - 1)
+
\text{length}_R(M'/N)
$$
を得る。
$\varphi_{I, M}(n) = \chi_{I, M}(n) - \chi_{I, M}(n - 1)$
であり、加群 $M''$ と $N$ についても同様である。
したがって $n \geq c$ に対して
$\varphi_{I, M}(n) = \varphi_{I, M''}(n) + \varphi_{I, N}(n-c)$ を得る。
\end{proof}

\begin{lemma}
\label{algebra-lemma-hilbert-change-I}
$I$, $I'$ をネーター局所環 $R$ の二つの定義イデアルとする。
$M$ を有限 $R$-加群とする。このとき定数 $a$ が存在し、
$n \geq 1$ に対して $\chi_{I, M}(n) \leq \chi_{I', M}(an)$ が成り立つ。
\end{lemma}

\begin{proof}
$(I')^c \subset I$ となる整数 $c \geq 1$ が存在する。
したがって全射 $M/(I')^{c(n + 1)}M \to M/I^{n + 1}M$ を得る。
ゆえに $a = 2c - 1$ とすれば結果が従う。
\end{proof}

\begin{proposition}
\label{algebra-proposition-hilbert-function-polynomial}
$R$ をネーター局所環とする。$M$ を有限 $R$-加群とする。
$I \subset R$ を定義イデアルとする。
ヒルベルト関数 $\varphi_{I, M}$ と関数 $\chi_{I, M}$ は数値的多項式である。
\end{proposition}

\begin{proof}
次数付き環 $S = R/I \oplus I/I^2 \oplus I^2/I^3 \oplus
\ldots = \bigoplus_{d \geq 0} I^d/I^{d + 1}$ と、
次数付き $S$-加群 $N = M/IM \oplus IM/I^2M \oplus \ldots =
\bigoplus_{d \geq 0} I^dM/I^{d + 1}M$ を考える。
この対 $(S, N)$ は命題 \ref{algebra-proposition-graded-hilbert-polynomial} の仮定を満たす。
したがって $\varphi_{I, M}$ についての結果は、その命題と
補題 \ref{algebra-lemma-length-K0} から従う。
$\chi_{I, M}$ についての結果は、これと補題 \ref{algebra-lemma-numerical-polynomial} から従う。
\end{proof}

\begin{definition}
\label{algebra-definition-hilbert-polynomial}
$R$ をネーター局所環とする。$M$ を有限 $R$-加群とする。
$R$ 上の $M$ の {\it ヒルベルト多項式}とは、
$n \gg 0$ に対して $P(n) = \varphi_M(n)$ となる元
$P(t) \in \mathbf{Q}[t]$ のことである。
\end{definition}

\noindent
命題 \ref{algebra-proposition-hilbert-function-polynomial} により、
ヒルベルト多項式が存在することが分かる。

\begin{lemma}
\label{algebra-lemma-d-independent}
$R$ をネーター局所環とする。$M$ を有限 $R$-加群とする。
\begin{enumerate}
\item 数値的多項式 $\varphi_{I, M}$ の次数は、定義イデアル $I$ に依存しない。
\item 数値的多項式 $\chi_{I, M}$ の次数は、定義イデアル $I$ に依存しない。
\end{enumerate}
\end{lemma}

\begin{proof}
(2) は補題 \ref{algebra-lemma-hilbert-change-I} から直ちに従う。
$n \geq 1$ に対して
$\varphi_{I, M}(n) = \chi_{I, M}(n) - \chi_{I, M}(n - 1)$
なので、(1) は (2) から従う。
\end{proof}

\begin{definition}
\label{algebra-definition-d}
$R$ をネーター局所環、$M$ を有限 $R$-加群とする。
次のように定められる $\{-\infty, 0, 1,\allowbreak 2, \ldots \}$ の元を
{\it $d(M)$} と書く。
\begin{enumerate}
\item $M = 0$ ならば $d(M) = -\infty$ と置く。
\item $M \not = 0$ ならば $d(M)$ は数値的多項式 $\chi_M$ の次数である。
\end{enumerate}
\end{definition}

\noindent
すべての $n$ に対して $\mathfrak m^nM \not = 0$ ならば、
$d(M)$ は $M$ のヒルベルト多項式の次数に $+1$ したものとなる。

\begin{lemma}
\label{algebra-lemma-differ-finite-chi}
$R$ をネーター局所環とする。$I \subset R$ を定義イデアルとする。
$M$ を、有限の長さをもたない有限 $R$-加群とする。
% P01-E0382: the source's degree-comparison wording is grammatically incomplete.
$M' \subset M$ が有限の余長をもつ部分加群ならば、
$\chi_{I, M} - \chi_{I, M'}$ は、その次数がいずれの多項式の次数よりも真に小さい（$<$）多項式である。
\end{lemma}

\begin{proof}
補題 \ref{algebra-lemma-differ-finite} から初等的な計算により従う。
\end{proof}

\begin{lemma}
\label{algebra-lemma-hilbert-ses-chi}
$R$ をネーター局所環とする。$I \subset R$ を定義イデアルとする。
$0 \to M' \to M \to M'' \to 0$ を有限 $R$-加群の短完全列とする。
このとき
\begin{enumerate}
\item $M'$ が有限の長さをもたないならば、
$\chi_{I, M} - \chi_{I, M''} - \chi_{I, M'}$ は、
次数が $\chi_{I, M'}$ の次数より真に小さい（$<$）数値的多項式である。
\item $\max\{ \deg(\chi_{I, M'}), \deg(\chi_{I, M''}) \} = \deg(\chi_{I, M})$ である。
\item $\max\{d(M'), d(M'')\} = d(M)$ である。
\end{enumerate}
\end{lemma}

\begin{proof}
まず (1) を示す。$N \subset M'$ を補題 \ref{algebra-lemma-hilbert-ses} のように取る。
補題 \ref{algebra-lemma-differ-finite-chi} により、
数値的多項式 $\chi_{I, M'} - \chi_{I, N}$ の次数は、
$\chi_{I, M'}$ と $\chi_{I, N}$ の共通の次数より真に小さい（$<$）。
補題 \ref{algebra-lemma-hilbert-ses} により、差
$$
\chi_{I, M}(n) - \chi_{I, M''}(n) - \chi_{I, N}(n - c)
$$
は $n \gg 0$ に対して定数である。
初等的な計算により、差 $\chi_{I, N}(n) - \chi_{I, N}(n - c)$ の次数は
$\chi_{I, N}$ の次数より真に小さく（$<$）、この後者の次数は零より大きい（上を参照）。
これらを合わせると (1) を得る。

\medskip\noindent
$\chi_{I, M'}$ と $\chi_{I, M''}$ の最高次係数は非負であることに注意する。
したがって $\chi_{I, M'} + \chi_{I, M''}$ の次数は、それぞれの次数の最大値に等しい。
ゆえに $M'$ が有限の長さをもたないならば、(2) は (1) から従う。
$M'$ が有限の長さをもつならば、アルティン--リース
（補題 \ref{algebra-lemma-Artin-Rees}）により、すべての $n \gg 0$ に対して
$I^nM \to I^nM''$ は同型である。
したがって $n \gg 0$ に対して $M/I^nM \to M''/I^nM''$ は
核が $M'$ である全射となり、
すべての $n \gg 0$ に対して
$\chi_{I, M}(n) - \chi_{I, M''}(n) = \text{length}(M')$ が分かる。
ゆえに、この場合も (2) は成り立つ。

\medskip\noindent
(3) の証明。$M$、$M'$、$M''$ のいずれかが零である場合を除き、
これは (2) から従う。これらの特別な場合の証明は省略する。
\end{proof}


\section{次元}
\label{algebra-section-dimension}

\noindent
「位相」節 \ref{topology-section-krull-dimension} と比較されたい。

\begin{definition}
\label{algebra-definition-chain-primes}
$R$ を環とする。{\it 素イデアルの鎖}とは、
$i = 0, \ldots, n - 1$ に対して $\mathfrak p_i \not = \mathfrak p_{i + 1}$
を満たす $R$ の素イデアルの列
$\mathfrak p_0 \subset \mathfrak p_1 \subset \ldots \subset \mathfrak p_n$
をいう。この素イデアルの鎖の {\it 長さ}は $n$ である。
\end{definition}

\noindent
環 $R$ の素イデアルと $\Spec(R)$ の既約閉部分集合の間には、
包含関係を逆転させる全単射があることを思い出そう
（補題 \ref{algebra-lemma-irreducible} 参照）。

\begin{definition}
\label{algebra-definition-Krull}
環 $R$ の {\it クルル次元}とは、位相空間 $\Spec(R)$ のクルル次元である
（「位相」定義 \ref{topology-definition-Krull} 参照）。
言い換えると、$R$ が長さ $n$ の素イデアルの鎖
$$
\mathfrak p_0
\subset
\mathfrak p_1
\subset
\ldots
\subset
\mathfrak p_n, \quad
\mathfrak p_i \not = \mathfrak p_{i + 1}.
$$
をもつような整数 $n\geq 0$ の上限である。
\end{definition}

\begin{definition}
\label{algebra-definition-height}
環 $R$ の素イデアル $\mathfrak p$ の {\it 高さ}とは、
局所環 $R_{\mathfrak p}$ の次元である。
\end{definition}

\begin{lemma}
\label{algebra-lemma-dimension-height}
$R$ のクルル次元は、その（極大）素イデアルの高さの上限である。
\end{lemma}

\begin{proof}
素イデアルの鎖の末尾には常に極大イデアルを加えることができるので、主張が従う。
\end{proof}

\begin{lemma}
\label{algebra-lemma-Noetherian-dimension-0}
次元 $0$ のネーター環はアルティン環である。
逆に、任意のアルティン環は次元零のネーター環である。
\end{lemma}

\begin{proof}
$R$ を次元 $0$ のネーター環とする。
補題 \ref{algebra-lemma-Noetherian-topology} により、空間 $\Spec(R)$ はネーターである。
「位相」補題 \ref{topology-lemma-Noetherian} により、
$\Spec(R)$ は有限個の既約成分をもち、
$\Spec(R) = Z_1 \cup \ldots \cup Z_r$ と書ける。
補題 \ref{algebra-lemma-irreducible} によれば、
% P01-E0383: the frozen source says minimal ideal, not minimal prime ideal.
各 $Z_i = V(\mathfrak p_i)$ において $\mathfrak p_i$ は極小イデアルである。
次元が $0$ なので、これらの $\mathfrak p_i$ は極大でもある。
したがって $\Spec(R)$ は、元が $\mathfrak p_i$ である離散位相空間である。
ジャコブソン根基 $\bigcap \mathfrak p_i$ のすべての元 $f$ は冪零である。
そうでなければ $R_f$ は零環でなくなり、別の素イデアルが存在するからである。
補題 \ref{algebra-lemma-product-local} により $R$ は $\prod R_{\mathfrak p_i}$ に等しい。
$R_{\mathfrak p_i}$ も次元 $0$ のネーター環なので、
先の議論から、その根基 $\mathfrak p_iR_{\mathfrak p_i}$ は局所冪零である。
補題 \ref{algebra-lemma-Noetherian-power} により、ある $n \geq 1$ に対して
$\mathfrak p_i^nR_{\mathfrak p_i} = 0$ である。
補題 \ref{algebra-lemma-length-finite} により、$R_{\mathfrak p_i}$ は $R$ 上有限の長さをもつと結論できる。
したがって補題 \ref{algebra-lemma-artinian-finite-length} により、$R$ はアルティン環である。

\medskip\noindent
$R$ がアルティン環ならば、補題 \ref{algebra-lemma-artinian-finite-length} によりネーターである。
補題 \ref{algebra-lemma-artinian-finite-nr-max}、\ref{algebra-lemma-artinian-radical-nilpotent}、
\ref{algebra-lemma-product-local} を合わせると、その素イデアルはすべて極大である。
\end{proof}

\noindent
以下では、定義 \ref{algebra-definition-d} で定めた不変量 $d(-)$ を用いる。
まず準備として次の補題を述べる。

\begin{lemma}
\label{algebra-lemma-dimension-0-d-0}
$R$ をネーター局所環とする。
このとき $\dim(R) = 0 \Leftrightarrow d(R) = 0$ である。
\end{lemma}

\begin{proof}
これは、$d(R) = 0$ であることと $R$ が $R$-加群として有限の長さをもつことが
同値だからである。補題 \ref{algebra-lemma-artinian-finite-length} を参照。
\end{proof}

\begin{proposition}
\label{algebra-proposition-dimension-zero-ring}
$R$ を環とする。次は同値である。
\begin{enumerate}
\item $R$ はアルティン環である。
\item $R$ はネーターであり $\dim(R) = 0$ である。
\item $R$ は自身の上の加群として有限の長さをもつ。
\item $R$ は有限個のアルティン局所環の直積である。
\item $R$ はネーターであり $\Spec(R)$ は有限離散位相空間である。
\item $R$ は次元 $0$ の有限個のネーター局所環の直積である。
\item $R$ は $d(R_i) = 0$ を満たす有限個のネーター局所環 $R_i$ の直積である。
\item $R$ は極大イデアルが冪零である有限個のネーター局所環 $R_i$ の直積である。
\item $R$ はネーターであり、極大イデアルを有限個しかもたず、そのジャコブソン根基イデアルは冪零である。
\item $R$ はネーターであり、その素イデアルの間に真の包含関係はない。
\end{enumerate}
\end{proposition}

\begin{proof}
これは補題 \ref{algebra-lemma-product-local}、
\ref{algebra-lemma-artinian-finite-length}、
\ref{algebra-lemma-Noetherian-dimension-0}、
\ref{algebra-lemma-dimension-0-d-0} を合わせたものである。
\end{proof}

\begin{lemma}
\label{algebra-lemma-height-1}
$R$ をネーター局所環とする。次は同値である。
\begin{enumerate}
\item
\label{algebra-item-dim-1}
$\dim(R) = 1$ である。
\item
\label{algebra-item-d-1}
$d(R) = 1$ である。
\item
\label{algebra-item-Vx}
$x \in \mathfrak m$ で、$x$ が冪零でなく $V(x) = \{\mathfrak m\}$ となるものが存在する。
\item
\label{algebra-item-x}
$x \in \mathfrak m$ で、$x$ が冪零でなく $\mathfrak m = \sqrt{(x)}$ となるものが存在する。
\item
\label{algebra-item-ideal-1}
$1$ 個の元で生成される定義イデアルが存在し、
$0$ 個の元で生成される定義イデアルは存在しない。
\end{enumerate}
\end{lemma}

\begin{proof}
まず $\dim(R) = 1$ と仮定する。
$R$ の極小素イデアルを $\mathfrak p_i$ とする。
次元が $1$ なので、$R$ のそれ以外の素イデアルは $\mathfrak m$ だけである。
補題 \ref{algebra-lemma-Noetherian-irreducible-components} により、極小素イデアルは有限個である。
したがって $x \in \mathfrak m$, $x \not \in \mathfrak p_i$ となる元を選べる
（補題 \ref{algebra-lemma-silly} 参照）。
ゆえに $x$ を含む素イデアルは $\mathfrak m$ だけであり、
(\ref{algebra-item-Vx}) を得る。

\medskip\noindent
(\ref{algebra-item-Vx}) が成り立つならば、
補題 \ref{algebra-lemma-Zariski-topology} により $\mathfrak m = \sqrt{(x)}$ であり、
したがって (\ref{algebra-item-x}) が成り立つ。逆も明らかである。
% P01-E0385: the frozen English misplaces the adverb directly.
(\ref{algebra-item-x}) と (\ref{algebra-item-ideal-1}) の同値性は定義から直ちに従う。

\medskip\noindent
(\ref{algebra-item-ideal-1}) を仮定する。$I = (x)$ を定義イデアルとする。
$x^n$ による乗法により $I^n/I^{n + 1}$ は $R/I$ の商であり、
したがって $\text{length}_R(I^n/I^{n + 1})$ は有界であることに注意する。
ゆえに $d(R) = 0$ または $d(R) = 1$ だが、
$0$ が定義イデアルでないという仮定により $d(R) = 0$ は除外される。

\medskip\noindent
(\ref{algebra-item-d-1}) を仮定する。矛盾を導くため、
両方の包含が真である素イデアル $\mathfrak p \subset \mathfrak q \subset \mathfrak m$
が存在すると仮定する。定義イデアル $I \subset R$ を一つ選ぶ。
補題 \ref{algebra-lemma-hilbert-ses-chi} を繰り返し用いる。
まず、完全列 $0 \to \mathfrak p \to R \to R/\mathfrak p \to 0$ によって
$d(R/\mathfrak p) \leq 1$ が従う。しかし、これは明らかに零ではあり得ない。
$x\in \mathfrak q$, $x\not \in \mathfrak p$ を選ぶ。短完全列
% P01-E0386: the frozen sequence does not label the multiplication-by-x arrow.
$$
0 \to R/\mathfrak p \to R/\mathfrak p \to R/(xR + \mathfrak p) \to 0.
$$
を考える。これより
$\chi_{I, R/\mathfrak p} - \chi_{I, R/\mathfrak p}
- \chi_{I, R/(xR + \mathfrak p)} = - \chi_{I, R/(xR + \mathfrak p)}$
の次数は $ < 1$ である。
言い換えると $d(R/(xR + \mathfrak p)) = 0$ であり、
補題 \ref{algebra-lemma-dimension-0-d-0} により $\dim(R/(xR + \mathfrak p)) = 0$ となる。
しかし $R/(xR + \mathfrak p)$ は相異なる素イデアル
$\mathfrak q/(xR + \mathfrak p)$ と $\mathfrak m/(xR + \mathfrak p)$ をもち、
求める矛盾を得る。
\end{proof}

\begin{proposition}
\label{algebra-proposition-dimension}
$R$ をネーター局所環とする。$d \geq 0$ を整数とする。次は同値である。
\begin{enumerate}
\item
\label{algebra-item-dim-d}
$\dim(R) = d$ である。
\item
\label{algebra-item-d-d}
$d(R) = d$ である。
\item
\label{algebra-item-ideal-d}
$d$ 個の元で生成される定義イデアルが存在し、
$d$ 個未満の元で生成される定義イデアルは存在しない。
\end{enumerate}
\end{proposition}

\begin{proof}
この証明は実質的に補題 \ref{algebra-lemma-height-1} の証明と同じである。
$d$ に関する帰納法で命題を示す。
補題 \ref{algebra-lemma-dimension-0-d-0} と \ref{algebra-lemma-height-1} により、$d > 1$ と仮定してよい。
$R$ の定義イデアルの生成元の最小個数を $d'(R)$ と書く。
不等式 $\dim(R) \geq d'(R) \geq d(R) \geq \dim(R)$ を示す。
これにより、これらはすべて等しくなる。

\medskip\noindent
まず $\dim(R) = d$ と仮定する。
$R$ の極小素イデアルを $\mathfrak p_i$ とする。
補題 \ref{algebra-lemma-Noetherian-irreducible-components} により、これらは有限個である。
したがって $x \in \mathfrak m$, $x \not \in \mathfrak p_i$ となる元を選べる
（補題 \ref{algebra-lemma-silly} 参照）。
素イデアルの任意の極大鎖は、ある $\mathfrak p_i$ から始まることに注意する。
ゆえに $R/xR$ の次元は高々 $d-1$ である。
帰納法により、$R/xR$ の定義イデアルを生成する $x_2, \ldots, x_d$ が存在する。
したがって $R$ は（高々）$d$ 個の元で生成される定義イデアルをもつ。

\medskip\noindent
$d'(R) = d$ と仮定する。$I = (x_1, \ldots, x_d)$ を定義イデアルとする。
$x_1, \ldots, x_d$ の次数 $n$ のすべての単項式による乗法により、
$I^n/I^{n + 1}$ は $\binom{d + n - 1}{d - 1}$ 個の $R/I$ の直和の商となることに注意する。
したがって $\text{length}_R(I^n/I^{n + 1})$ は
次数 $d-1$ の多項式によって上から抑えられる。ゆえに $d(R) \leq d$ である。

\medskip\noindent
$d(R) = d$ と仮定する。
すべての包含が真である素イデアルの鎖
$\mathfrak p \subset \mathfrak q \subset
\mathfrak q_2 \subset \ldots \subset \mathfrak q_e = \mathfrak m$ を考え、$e \geq 2$ とする。
定義イデアル $I \subset R$ を一つ選ぶ。
補題 \ref{algebra-lemma-hilbert-ses-chi} を繰り返し用いる。
まず、完全列 $0 \to \mathfrak p \to R \to R/\mathfrak p \to 0$ によって
$d(R/\mathfrak p) \leq d$ が従う。しかし、これは明らかに零ではあり得ない。
$x\in \mathfrak q$, $x\not \in \mathfrak p$ を選び、短完全列
$$
0 \to R/\mathfrak p \to R/\mathfrak p \to R/(xR + \mathfrak p) \to 0.
$$
を考える。これより
$\chi_{I, R/\mathfrak p} - \chi_{I, R/\mathfrak p}
- \chi_{I, R/(xR + \mathfrak p)} = - \chi_{I, R/(xR + \mathfrak p)}$
の次数は $ < d$ である。
言い換えると $d(R/(xR + \mathfrak p)) \leq d - 1$ であり、
帰納法により $\dim(R/(xR + \mathfrak p)) \leq d - 1$ である。
ここで $R/(xR + \mathfrak p)$ は素イデアルの鎖
$\mathfrak q/(xR + \mathfrak p) \subset \mathfrak q_2/(xR + \mathfrak p)
\subset \ldots \subset \mathfrak q_e/(xR + \mathfrak p)$ をもつので、
$e - 1 \leq d - 1$ を得る。
任意の素イデアルの鎖から始めたので、これで $\dim(R) \leq d(R)$ が示された。

\medskip\noindent
読み返せば、求めていた循環する一連の不等式が示されたことが分かる。
\end{proof}

\noindent
$(R, \mathfrak m)$ をネーター局所環とする。
% P01-E0387: the frozen text says variables rather than elements or generators.
上の議論から、$\mathfrak m$ は $\dim(R)$ 個未満の変数では生成されないことが明らかである。
中山の補題 \ref{algebra-lemma-NAK} により、$\mathfrak m$ の生成元の最小個数は
$\dim_{\kappa(\mathfrak m)} \mathfrak m/\mathfrak m^2$ に等しい。
したがって次の基本不等式を得る。
$$
\dim(R) \leq \dim_{\kappa(\mathfrak m)} \mathfrak m/\mathfrak m^2.
$$
等号が成り立つ環は、多くのよい性質をもつことが分かる。
これらを正則局所環という。

\begin{definition}
\label{algebra-definition-regular-local}
$(R, \mathfrak m)$ を次元 $d$ のネーター局所環とする。
\begin{enumerate}
\item {\it $R$ のパラメータ系}とは、$R$ の定義イデアルを生成する
元の列 $x_1, \ldots, x_d \in \mathfrak m$ をいう。
\item $\mathfrak m = (x_1, \ldots, x_d)$ を満たす
$x_1, \ldots, x_d \in \mathfrak m$ が存在するとき、
$R$ を {\it 正則局所環}と呼び、$x_1, \ldots, x_d$ を
{\it 正則パラメータ系}と呼ぶ。
\end{enumerate}
\end{definition}

\noindent
以下の補題は、上の補題と命題の証明から明らかであるが、
参照しやすいように明記しておく。

\begin{lemma}
\label{algebra-lemma-minimal-over-1}
$R$ をネーター環とする。$x \in R$ とする。
\begin{enumerate}
\item $\mathfrak p$ が $(x)$ の上で極小ならば、
$\mathfrak p$ の高さは $0$ または $1$ である。
\item $\mathfrak p, \mathfrak q \in \Spec(R)$ であり、
$\mathfrak q$ が $(\mathfrak p, x)$ の上で極小ならば、
$\mathfrak p$ と $\mathfrak q$ の真の中間に素イデアルは存在しない。
\end{enumerate}
\end{lemma}

\begin{proof}
(1) の証明。$\mathfrak p$ が $x$ の上で極小ならば、
$R_\mathfrak p$ の $x$ を含む素イデアルは極大イデアル $\mathfrak p R_\mathfrak p$ だけである。
これは $R_\mathfrak p$ の素イデアルが、$\mathfrak p$ に含まれる
$R$ の素イデアルと $1$ 対 $1$ に対応するからである
（補題 \ref{algebra-lemma-spec-localization} 参照）。
ゆえに、$x$ が $R_\mathfrak p$ において冪零でないならば、
補題 \ref{algebra-lemma-height-1} により $\dim(R_\mathfrak p) = 1$ である。
もちろん、$x$ が $R_\mathfrak p$ において冪零ならば、
この議論から $\mathfrak pR_\mathfrak p$ が唯一の素イデアルとなり、高さは $0$ である。

\medskip\noindent
(2) の証明。(1) により、$\mathfrak q/\mathfrak p$ は
$R/\mathfrak p$ において高さ $1$ または $0$ の素イデアルである。
これから直ちに、$\mathfrak p$ と $\mathfrak q$ の真の中間には
素イデアルが存在し得ないことが従う。
\end{proof}

\begin{lemma}
\label{algebra-lemma-minimal-over-r}
$R$ をネーター環とする。$f_1, \ldots, f_r \in R$ とする。
\begin{enumerate}
\item $\mathfrak p$ が $(f_1, \ldots, f_r)$ の上で極小ならば、
$\mathfrak p$ の高さは $\leq r$ である。
\item $\mathfrak p, \mathfrak q \in \Spec(R)$ であり、
$\mathfrak q$ が $(\mathfrak p, f_1, \ldots, f_r)$ の上で極小ならば、
$\mathfrak p$ と $\mathfrak q$ の間の任意の素イデアルの鎖の長さは高々 $r$ である。
\end{enumerate}
\end{lemma}

\begin{proof}
(1) の証明。$\mathfrak p$ が $f_1, \ldots, f_r$ の上で極小ならば、
$R_\mathfrak p$ の $f_1, \ldots, f_r$ を含む素イデアルは
極大イデアル $\mathfrak p R_\mathfrak p$ だけである。
これは $R_\mathfrak p$ の素イデアルが、$\mathfrak p$ に含まれる
$R$ の素イデアルと $1$ 対 $1$ に対応するからである
（補題 \ref{algebra-lemma-spec-localization} 参照）。
ゆえに命題 \ref{algebra-proposition-dimension} により $\dim(R_\mathfrak p) \leq r$ である。

\medskip\noindent
(2) の証明。(1) により、$\mathfrak q/\mathfrak p$ は高さ $\leq r$ の素イデアルである。
これから直ちに、$\mathfrak p$ と $\mathfrak q$ の間の素イデアルの鎖に関する主張が従う。
\end{proof}

\begin{lemma}
\label{algebra-lemma-one-equation}
$R$ をネーター局所環とし、$x\in \mathfrak m$ をその極大イデアルの元とする。
このとき $\dim R \leq \dim R/xR + 1$ である。
$x$ が $R$ のどの極小素イデアルにも含まれなければ等号が成り立つ。
（例えば $x$ が非零因子ならばそうである。）
\end{lemma}

\begin{proof}
$x_1, \ldots, x_{\dim R/xR} \in R$ の像が $R/xR$ の定義イデアルを生成する
$R/xR$ の元であるならば、$x, x_1, \ldots,
x_{\dim R/xR}$ は $R$ の定義イデアルを生成する。
したがって命題 \ref{algebra-proposition-dimension} により不等式を得る。
一方、$x$ が $R$ のどの極小素イデアルにも含まれなければ、
$R/xR$ の素イデアルの鎖はすべて、$R$ において
極大な鎖となるまでに少なくとももう一段階を要する鎖を与える。
\end{proof}

\begin{lemma}
\label{algebra-lemma-elements-generate-ideal-definition}
$(R, \mathfrak m)$ をネーター局所環とする。
$x_1, \ldots, x_d \in \mathfrak m$ が定義イデアルを生成し、
$d = \dim(R)$ と仮定する。このとき、すべての $i = 1, \ldots, d$ に対して
$\dim(R/(x_1, \ldots, x_i)) = d - i$ が成り立つ。
\end{lemma}

\begin{proof}
命題 \ref{algebra-proposition-dimension} の証明から従う。
または、$d$ に関する帰納法と補題 \ref{algebra-lemma-one-equation} を用いてもよい。
\end{proof}


\section{次元論の応用}
\label{algebra-section-applications-dimension-theory}

\noindent
次元に関する結果を用いて、ある種の環のスペクトルが無限であることを示し、
さらに多くのヤコブソン環を構成できる。

\begin{lemma}
\label{algebra-lemma-Noetherian-local-domain-dim-2-infinite-opens}
$R$ を次元 $\geq 2$ のネーター局所整域とする。
空でない開部分集合 $U \subset \Spec(R)$ は無限である。
\end{lemma}

\begin{proof}
矛盾を導くため、$U \subset \Spec(R)$ が有限であると仮定する。
このとき $(0) \in U$ であり、$\{(0)\}$ は $U$ の開部分集合である
（$\{(0)\}$ の補集合は、ほかの点の閉包の合併だからである）。
したがって $U = \{(0)\}$ と仮定してよい。
$\mathfrak m \subset R$ を極大イデアルとする。
$V(x) \cup U = \Spec(R)$ を満たす
$x \in \mathfrak m$, $x \not = 0$ を見つけることができる。
言い換えると $D(x) = \{(0)\}$ である。特に
$\dim(R/xR) = \dim(R) - 1 \geq 1$ であることが分かる
（補題 \ref{algebra-lemma-one-equation} 参照）。
$\overline{y}_2, \ldots, \overline{y}_{\dim(R)} \in R/xR$ を、
$R/xR$ の定義イデアルを生成する元とする
（命題 \ref{algebra-proposition-dimension} 参照）。
持ち上げ $y_2, \ldots, y_{\dim(R)} \in R$ を選び、
$x, y_2, \ldots, y_{\dim(R)}$ が $R$ の定義イデアルを生成するようにする。
これより $\dim(R/(y_2)) = \dim(R) - 1$ および
$\dim(R/(y_2, x)) = \dim(R) - 2$ が従う
（補題 \ref{algebra-lemma-elements-generate-ideal-definition} 参照）。
したがって、$y_2$ を含むが $x$ を含まない素イデアル
$\mathfrak p$ が存在する。これは $D(x) = \{(0)\}$ に矛盾する。
\end{proof}

\noindent
% P01-E0388: the frozen example says p-adic numbers rather than p-adic integers.
$k$ を体とするときの環 $k[[t]]$、および $p$ 進数の環は、
素イデアルを有限個しかもたない次元 $1$ のネーター環である。
この現象が起こり得る次元として、これは最大である。

\begin{lemma}
\label{algebra-lemma-Noetherian-finite-nr-primes}
素イデアルを有限個しかもたないネーター環の次元は $\leq 1$ である。
\end{lemma}

\begin{proof}
$R$ を素イデアルを有限個しかもたないネーター環とする。
$R$ が局所整域ならば、補題
\ref{algebra-lemma-Noetherian-local-domain-dim-2-infinite-opens} から主張が従う。
$R$ が整域ならば、局所の場合により、すべての極大イデアル
$\mathfrak m$ に対して $R_\mathfrak m$ の次元は $\leq 1$ である。
したがって補題 \ref{algebra-lemma-dimension-height} により $\dim(R) \leq 1$ である。
$R$ が一般の場合、$R$ の各極小素イデアル $\mathfrak q$ に対して
$\dim(R/\mathfrak q) \leq 1$ である。
各素イデアルは極小素イデアルを含むので
（補題 \ref{algebra-lemma-Zariski-topology}）、これより $\dim(R) \leq 1$ が従う。
\end{proof}

\begin{lemma}
\label{algebra-lemma-finite-type-algebra-finite-nr-primes}
$S$ を体 $k$ 上の零でない有限型代数とする。
次の条件は同値である。
\begin{enumerate}
\item $\dim(S) = 0$,
\item $S$ は素イデアルを有限個しかもたない,
\item $S$ は極大イデアルを有限個しかもたない,
\item $\Spec(S)$ は補題 \ref{algebra-lemma-ring-with-only-minimal-primes} の
同値な条件の一つを満たす,
\item $\dim_k(S) < \infty$,
\item $S$ はアルティン環である,
\item $\Spec(S)$ は離散位相空間である,
% P01-E0389: the frozen eighth item is an unfinished editorial placeholder.
\item ここにさらに項目を追加する。
\end{enumerate}
\end{lemma}

\begin{proof}
定義から、補題 \ref{algebra-lemma-ring-with-only-minimal-primes} の (5) を見れば、
(1) と (4) が同値であることは直ちに分かる。
$\Spec(S)$ はソバー、ネーター、かつヤコブソンであることを思い出そう。
補題 \ref{algebra-lemma-spec-spectral}, \ref{algebra-lemma-Noetherian-topology},
\ref{algebra-lemma-finite-type-field-Jacobson}, \ref{algebra-lemma-jacobson} を参照されたい。
$S$ の次元が $0$ ならば、各点が一つの既約成分を定め、
既約成分は有限個しかない
（Topology, Lemma \ref{topology-lemma-Noetherian}）。
したがって (1) は (2) を含意する。(2) が (3) を含意するのは自明である。
(3) が成り立つならば、Topology, Lemma
\ref{topology-lemma-finite-jacobson} により $\Spec(S)$ は離散的であり、
したがって $S$ の次元は $0$ である。

\medskip\noindent
ここまでで (1) -- (4) が同値であることが分かった。
含意 (5) $\Rightarrow$ (6) は補題 \ref{algebra-lemma-finite-dimensional-algebra} である。
含意 (6) $\Rightarrow$ (7) は命題
\ref{algebra-proposition-dimension-zero-ring} から従う。
含意 (7) $\Rightarrow$ (4) は直ちに分かる。
逆に $S$ が (1) -- (4) を満たすならば、$S$ は有限個の素イデアル
$\mathfrak m_1, \ldots, \mathfrak m_r$ をもち、それらはすべて極大である。
ヒルベルトの零点定理（定理 \ref{algebra-theorem-nullstellensatz}）により、
$\kappa(\mathfrak m_i)$ は $k$ の有限次拡大である。
命題 \ref{algebra-proposition-dimension-zero-ring} により、$S$ がアルティン環であることも分かる。
次に補題 \ref{algebra-lemma-artinian-finite-length} から
$\text{length}_S(S) < \infty$ を得る。したがって補題
\ref{algebra-lemma-pushdown-module} により $\dim_k(S) < \infty$ である。
以上より (1) -- (7) は同値である。
（注意：$\dim(S) = 0 \Rightarrow \dim_k(S) < \infty$ を示す別の、
より標準的な方法はネーター正規化を用いることであるが、
この段階ではまだ利用できない。）
\end{proof}

\begin{lemma}
\label{algebra-lemma-noetherian-dim-1-Jacobson}
ネーター・ヤコブソン環について、次が成り立つ。
\begin{enumerate}
\item 次元 $1$ で素イデアルを無限個もつ任意のネーター整域 $R$ は
ヤコブソン環である。
\item 各素イデアル $\mathfrak p$ が、極大であるか、または無限個の
素イデアルに含まれるような任意のネーター環はヤコブソン環である。
\end{enumerate}
\end{lemma}

\begin{proof}
(1) は補題 \ref{algebra-lemma-pid-jacobson} の言い換えである。

\medskip\noindent
$R$ を、各極大でない素イデアル $\mathfrak p$ が無限個の素イデアルに
含まれるようなネーター環とする。
矛盾を導くため $\Spec(R)$ がヤコブソンでないと仮定する。
補題 \ref{algebra-lemma-irreducible} と \ref{algebra-lemma-Noetherian-topology} により、
$\Spec(R)$ はソバーなネーター位相空間である。
Topology, Lemma \ref{topology-lemma-non-jacobson-Noetherian-characterize}
により、$\{\mathfrak p\}$ が $\Spec(R)$ の局所閉部分集合となるような
% P01-E0390: the frozen text says only non-maximal ideal although p is a spectrum point.
極大でないイデアル $\mathfrak p \subset R$ が存在する。
言い換えると、$\mathfrak p$ は極大でなく、
$\{\mathfrak p\}$ は $V(\mathfrak p)$ の開部分集合である。
% P01-E0391: the frozen subset relations fail to exclude q=p.
$\mathfrak p \subset \mathfrak q$ を満たす素イデアル
$\mathfrak q \subset R$ を考える。
$(R/\mathfrak p)_{\mathfrak q} = R_{\mathfrak q}/\mathfrak pR_{\mathfrak q}$
のスペクトルの位相は $\Spec(R)$ の位相から誘導されることを思い出そう。
補題 \ref{algebra-lemma-spec-localization} と \ref{algebra-lemma-spec-closed} を参照されたい。
したがって $\{(0)\}$ は
$\Spec((R/\mathfrak p)_{\mathfrak q})$ の局所閉部分集合である。
補題 \ref{algebra-lemma-Noetherian-local-domain-dim-2-infinite-opens} により
$\dim((R/\mathfrak p)_{\mathfrak q}) = 1$ と結論できる。
これがすべての $\mathfrak q \supset \mathfrak p$ に対して成り立つので、
$\dim(R/\mathfrak p) = 1$ と結論できる。
ここで $\mathfrak p$ が無限個の素イデアルに含まれるという仮定を用いると、
$\Spec(R/\mathfrak p)$ が無限であることが分かる。
したがって補題の (1) により
$V(\mathfrak p) \cong \Spec(R/\mathfrak p)$ は閉点の閉包である。
これは $\{\mathfrak p\} \subset V(\mathfrak p)$ が開ではあり得ないことを
意味するので、求めていた矛盾である。
\end{proof}


\section{加群の台と次元}
\label{algebra-section-support}

\noindent
加群の台と次元に関するいくつかの基本的な結果を述べる。

\begin{lemma}
\label{algebra-lemma-filter-Noetherian-module}
$R$ をネーター環とし、$M$ を有限 $R$-加群とする。
次のような $R$-部分加群によるフィルトレーションが存在する。
$$
0 = M_0 \subset M_1 \subset \ldots \subset M_n = M
$$
ここで各商 $M_i/M_{i-1}$ は、$R$ のある素イデアル
$\mathfrak p_i$ に対する $R/\mathfrak p_i$ と同型である。
\end{lemma}

\begin{proof}[第一の証明]
補題 \ref{algebra-lemma-trivial-filter-finite-module} により、あるイデアル $I$ に対する
$M = R/I$ の場合を示せば十分である。
加群 $R/J$ に対して補題が成り立たないようなイデアル $J$ の集合を $S$ とし、
包含関係によって順序を入れる。
矛盾を導くため、$S$ が空でないと仮定する。
$R$ はネーター環なので、$S$ は極大元 $J$ をもつ。
$S$ の定義により、イデアル $J$ は素イデアルではあり得ない。
$ab \in J$ だが $a \in J$ でも $b\in J$ でもないような
$a, b\in R$ を選ぶ。フィルトレーション
$0 \subset aR/(J \cap aR) \subset R/J$ を考える。
部分加群 $aR/(J \cap aR)$ と商加群
$(R/J)/(aR/(J \cap aR))$ はともに巡回加群であることに注意する。
これらを $R/J'$ および $R/J''$ と書けば、短完全列
$0 \to R/J' \to R/J \to R/J'' \to 0$ を得る。
$b \in J'$ なので包含 $J \subset J'$ は真であり、
$a \in J''$ なので包含 $J \subset J''$ も真である。
したがって $J$ の極大性により、$R/J'$ と $R/J''$ はともに
上のようなフィルトレーションをもち、それゆえ $R/J$ ももつ。これは矛盾である。
\end{proof}

\begin{proof}[第二の証明]
$R$-加群 $M$ に対し、補題の主張にあるようなフィルトレーションが
存在するとき $P(M)$ が成り立つということにする。
$P$ は拡大について安定であり、$0$ に対して成り立つことに注意する。
補題 \ref{algebra-lemma-trivial-filter-finite-module} により、すべてのイデアル $I$ に対して
$P(R/I)$ が成り立つことを示せば十分である。
もしそうでなければ、$R$ はネーター環なので、極大な
% P01-E0392: the frozen English writes the noun as two words.
反例 $J$ が存在する。例 \ref{algebra-example-oka-family-property-modules} と
命題 \ref{algebra-proposition-oka} により、イデアル $J$ は素イデアルであり、矛盾する。
\end{proof}

\begin{lemma}
\label{algebra-lemma-filter-primes-in-support}
$R$, $M$, $M_i$, $\mathfrak p_i$ を補題
\ref{algebra-lemma-filter-Noetherian-module} のとおりとする。
このとき $\text{Supp}(M) = \bigcup V(\mathfrak p_i)$ であり、
特に $\mathfrak p_i \in \text{Supp}(M)$ である。
\end{lemma}

\begin{proof}
補題 \ref{algebra-lemma-support-closed} と \ref{algebra-lemma-support-quotient} から従う。
\end{proof}

\begin{lemma}
\label{algebra-lemma-support-point}
$R$ を極大イデアル $\mathfrak m$ をもつネーター局所環とする。
$M$ を零でない有限 $R$-加群とする。
このとき $\text{Supp}(M) = \{ \mathfrak m\}$ であることと、
$M$ が $R$ 上有限長をもつことは同値である。
\end{lemma}

\begin{proof}
$\text{Supp}(M) = \{ \mathfrak m\}$ と仮定する。
補題 \ref{algebra-lemma-filter-Noetherian-module} のフィルトレーションに現れる
すべての素イデアル $\mathfrak p_i$ が極大イデアルであることを示せば十分である。
これは補題 \ref{algebra-lemma-filter-primes-in-support} から明らかである。

\medskip\noindent
$M$ が $R$ 上有限長をもつと仮定する。
補題 \ref{algebra-lemma-length-infinite} により $\mathfrak m^n M = 0$ である。
$R$ の任意の素イデアル $\mathfrak p \not = \mathfrak m$ に対し、
$\mathfrak m$ のある元は $R_{\mathfrak p}$ で単元になるので、
$M_{\mathfrak p} = 0$ である。
\end{proof}

\begin{lemma}
\label{algebra-lemma-Noetherian-power-ideal-kills-module}
$R$ をネーター環とする。
$I \subset R$ をイデアルとする。
$M$ を有限 $R$-加群とする。
このとき、ある $n \geq 0$ に対して $I^nM = 0$ であることと
$\text{Supp}(M) \subset V(I)$ であることは同値である。
\end{lemma}

\begin{proof}
実際、$I^nM = 0$ であることは $I^n \subset \text{Ann}(M)$ と同値である。
$R$ はネーター環なので、これは
$I \subset \sqrt{\text{Ann}(M)}$ と同値である
（補題 \ref{algebra-lemma-Noetherian-power} 参照）。
さらにこれは $V(I) \supset V(\text{Ann}(M))$ と同値である
（補題 \ref{algebra-lemma-Zariski-topology} 参照）。
補題 \ref{algebra-lemma-support-closed} により、これは
$V(I) \supset \text{Supp}(M)$ と同値である。
\end{proof}

\begin{lemma}
\label{algebra-lemma-filter-minimal-primes-in-support}
$R$, $M$, $M_i$, $\mathfrak p_i$ を補題
\ref{algebra-lemma-filter-Noetherian-module} のとおりとする。
集合 $\{\mathfrak p_i\}$ の極小元は $\text{Supp}(M)$ の極小元である。
極小素イデアル $\mathfrak p$ が現れる回数は
$$
\#\{i \mid \mathfrak p_i = \mathfrak p\}
=
\text{length}_{R_\mathfrak p} M_{\mathfrak p}.
$$
である。
\end{lemma}

\begin{proof}
第一の主張は
$\text{Supp}(M) = \bigcup V(\mathfrak p_i)$ であることから従う
（補題 \ref{algebra-lemma-filter-primes-in-support} 参照）。
$\mathfrak p \in \text{Supp}(M)$ を極小とする。
$M_{\mathfrak p}$ の台は極大イデアル
$\mathfrak p R_{\mathfrak p}$ だけからなる集合である。
したがって補題 \ref{algebra-lemma-support-point} により、
$M_{\mathfrak p}$ の長さは有限かつ $> 0$ である。
次に、$M_{\mathfrak p}$ は次の部分商をもつフィルトレーションをもつことに注意する。
$
(R/\mathfrak p_i)_{\mathfrak p}
=
R_{\mathfrak p}/{\mathfrak p_i}R_{\mathfrak p}
$.
これは $\mathfrak p_i \not \subset \mathfrak p$ のとき零であり、
$\mathfrak p_i \subset \mathfrak p$ のとき $\kappa(\mathfrak p)$ に等しい。
後者の場合、$\mathfrak p$ の極小性により
$\mathfrak p_i = \mathfrak p$ だからである。
$\kappa(\mathfrak p)$ の長さは $1$ なので、結果が従う。
\end{proof}

\begin{lemma}
\label{algebra-lemma-support-dimension-d}
$R$ をネーター局所環とする。
$M$ を有限 $R$-加群とする。
このとき $d(M) = \dim(\text{Supp}(M))$ である。
ここで $d(M)$ は定義 \ref{algebra-definition-d} のものとする。
\end{lemma}

\begin{proof}
$M_i, \mathfrak p_i$ を補題 \ref{algebra-lemma-filter-Noetherian-module} のとおりとする。
補題 \ref{algebra-lemma-hilbert-ses-chi} により
$d(M) = \max \{ d(R/\mathfrak p_i) \}$ を得る。
命題 \ref{algebra-proposition-dimension} により
$d(R/\mathfrak p_i) = \dim(R/\mathfrak p_i)$ である。
$\dim(R/\mathfrak p_i) = \dim V(\mathfrak p_i)$ は自明である。
$\text{Supp}(M)$ のすべての極小素イデアルは
$\mathfrak p_i$ の中に現れるので
（補題 \ref{algebra-lemma-filter-minimal-primes-in-support}）、主張が従う。
\end{proof}

\begin{lemma}
\label{algebra-lemma-ses-dimension}
$R$ をネーター環とする。有限 $R$-加群の短完全列
$0 \to M' \to M \to M'' \to 0$ をとる。このとき
$\max\{\dim(\text{Supp}(M')), \dim(\text{Supp}(M''))\} =
\dim(\text{Supp}(M))$ である。
\end{lemma}

\begin{proof}
$R$ が局所環ならば、補題 \ref{algebra-lemma-support-dimension-d} と
\ref{algebra-lemma-hilbert-ses-chi} から直ちに従う。
$R$ が局所環でない場合にも通用する、より初等的な議論としては、
$\text{Supp}(M')$, $\text{Supp}(M'')$, $\text{Supp}(M)$ が閉であること
（補題 \ref{algebra-lemma-support-closed}）と、
$\text{Supp}(M) = \text{Supp}(M') \cup \text{Supp}(M'')$ であること
（補題 \ref{algebra-lemma-support-quotient}）を用いればよい。
\end{proof}


\section{随伴素イデアル}
\label{algebra-section-ass}

\noindent
標準的な定義を述べる。非ネーター環や有限でない加群については、
\ref{algebra-section-weakly-ass} 節の定義を用いる方が適切なこともある。

\begin{definition}
\label{algebra-definition-associated}
$R$ を環とし、$M$ を $R$-加群とする。
$R$ の素イデアル $\mathfrak p$ が $M$ に {\it 随伴する}とは、
零化イデアルが $\mathfrak p$ である元 $m \in M$ が存在することをいう。
このような素イデアルすべての集合を $\text{Ass}_R(M)$ または
$\text{Ass}(M)$ と表す。
\end{definition}

\begin{lemma}
\label{algebra-lemma-ass-support}
$R$ を環とし、$M$ を $R$-加群とする。
このとき $\text{Ass}(M) \subset \text{Supp}(M)$ である。
\end{lemma}

\begin{proof}
$m \in M$ の零化イデアルが $\mathfrak p$ ならば、
特に $R \setminus \mathfrak p$ のどの元も $m$ を零化しない。
したがって $m$ は $M_{\mathfrak p}$ の零でない元であり、
すなわち $\mathfrak p \in \text{Supp}(M)$ である。
\end{proof}

\begin{lemma}
\label{algebra-lemma-ass}
$R$ を環とし、$R$-加群の短完全列
$0 \to M' \to M \to M'' \to 0$ をとる。このとき
$\text{Ass}(M') \subset \text{Ass}(M)$ および
$\text{Ass}(M) \subset \text{Ass}(M') \cup \text{Ass}(M'')$ が成り立つ。
また $\text{Ass}(M' \oplus M'') = \text{Ass}(M') \cup \text{Ass}(M'')$ である。
\end{lemma}

\begin{proof}
$m' \in M'$ ならば、$M'$ の元として見た $m'$ の零化イデアルは、
$M$ の元として見た $m'$ の零化イデアルと同じである。
したがって包含 $\text{Ass}(M') \subset \text{Ass}(M)$ を得る。
$m \in M$ を、零化イデアルが素イデアル $\mathfrak p$ である元とする。
$g \in R$, $g \not \in \mathfrak p$ であって
$m' = gm \in M'$ となるものが存在すれば、$m'$ の零化イデアルは
$\mathfrak p$ である。そのような $g \in R$, $g \not \in \mathfrak p$ で
$gm \in M'$ となるものが存在しなければ、
% P01-E0393: the frozen English misspells annihilator.
商加群における $m$ の像 $m'' \in M''$ の零化イデアルは $\mathfrak p$ である。
これで包含
$\text{Ass}(M) \subset \text{Ass}(M') \cup \text{Ass}(M'')$
が示された。最後の主張の証明は省略する。
\end{proof}

\begin{lemma}
\label{algebra-lemma-ass-filter}
$R$ を環とし、$M$ を $R$-加群とする。
次のような $R$-部分加群によるフィルトレーションが存在すると仮定する。
$$
0 = M_0 \subset M_1 \subset \ldots \subset M_n = M
$$
ここで各商 $M_i/M_{i-1}$ は、$R$ のある素イデアル
$\mathfrak p_i$ に対する $R/\mathfrak p_i$ と同型である。
このとき
$\text{Ass}(M) \subset \{\mathfrak p_1, \ldots, \mathfrak p_n\}$ である。
\end{lemma}

\begin{proof}
フィルトレーション $\{ M_i \}$ の長さ $n$ に関する帰納法による。
零化イデアルが素イデアル $\mathfrak p$ である元 $m \in M$ を選ぶ。
$m \in M_{n-1}$ ならば帰納法によりよい。そうでなければ、
$m$ は $M/M_{n-1} \cong R/\mathfrak p_n$ の零でない元に写る。
したがって $\mathfrak p \subset \mathfrak p_n$ である。
等号が成り立たなければ、$f \in \mathfrak p_n$,
$f \not\in \mathfrak p$ となる元 を見つけることができる。
このとき $fm$ の零化イデアルは依然として $\mathfrak p$ であり、
$fm \in M_{n-1}$ である。したがって帰納法により主張を得る。
\end{proof}

\begin{lemma}
\label{algebra-lemma-finite-ass}
$R$ をネーター環とする。
$M$ を有限 $R$-加群とする。
このとき $\text{Ass}(M)$ は有限である。
\end{lemma}

\begin{proof}
補題 \ref{algebra-lemma-ass-filter} と補題
\ref{algebra-lemma-filter-Noetherian-module} から直ちに従う。
\end{proof}

\begin{proposition}
\label{algebra-proposition-minimal-primes-associated-primes}
$R$ をネーター環とする。
$M$ を有限 $R$-加群とする。
次の素イデアルの集合は同じである。
\begin{enumerate}
\item $M$ の台における極小素イデアル。
\item $\text{Ass}(M)$ における極小素イデアル。
\item $M_i/M_{i-1} \cong R/\mathfrak p_i$ を満たす任意のフィルトレーション
$0 = M_0 \subset M_1 \subset \ldots
\subset M_{n-1} \subset M_n = M$ に対する、
集合 $\{\mathfrak p_i\}$ の極小素イデアル。
\end{enumerate}
\end{proposition}

\begin{proof}
(3) のようなフィルトレーションを一つ選ぶ。
補題 \ref{algebra-lemma-filter-minimal-primes-in-support} で、
(1) と (3) の集合が等しいことを示した。

\medskip\noindent
$\mathfrak p$ を集合 $\{\mathfrak p_i\}$ の極小元とする。
$\mathfrak p = \mathfrak p_i$ となる最小の $i$ をとる。
$m \in M_i$, $m \not \in M_{i-1}$ を選ぶ。
$m$ の零化イデアルは $\mathfrak p_i = \mathfrak p$ に含まれ、
$\mathfrak p_1 \mathfrak p_2 \ldots \mathfrak p_i$ を含む。
$i$ と $\mathfrak p$ の選び方により、$j < i$ に対して
$\mathfrak p_j \not \subset \mathfrak p$ であり、したがって
$\mathfrak p_1 \mathfrak p_2 \ldots \mathfrak p_{i - 1}
\not \subset \mathfrak p_i$ である。
$f \in \mathfrak p_1 \mathfrak p_2 \ldots \mathfrak p_{i - 1}$,
$f \not \in \mathfrak p$ を選ぶ。このとき $fm$ の零化イデアルは
$\mathfrak p$ である。このようにして、$\mathfrak p$ が $M$ の随伴素イデアル
であることが分かる。補題 \ref{algebra-lemma-ass-support} により
$\text{Ass}(M) \subset \text{Supp}(M)$ であるから、
$\mathfrak p$ は $\text{Ass}(M)$ において極小である。
したがって (1) の素イデアルの集合は (2) の素イデアルの集合に含まれる。

\medskip\noindent
$\mathfrak p$ を $\text{Ass}(M)$ の極小元とする。
$\text{Ass}(M) \subset \text{Supp}(M)$ なので、
$\mathfrak q \subset \mathfrak p$ を満たす
$\text{Supp}(M)$ の極小元 $\mathfrak q$ が存在する。
上で $\mathfrak q \in \text{Ass}(M)$ を示した。
したがって $\mathfrak p$ の極小性により
$\mathfrak q = \mathfrak p$ である。
ゆえに (2) の素イデアルの集合は (1) の素イデアルの集合に含まれる。
\end{proof}

\begin{lemma}
\label{algebra-lemma-ass-zero}
\begin{slogan}
ネーター環上では、零でない各加群は随伴素イデアルをもつ。
\end{slogan}
$R$ をネーター環とし、$M$ を $R$-加群とする。
このとき
$$
M = (0) \Leftrightarrow \text{Ass}(M) = \emptyset.
$$
\end{lemma}

\begin{proof}
$M = (0)$ ならば、定義により $\text{Ass}(M) = \emptyset$ である。
$M \not = 0$ ならば、零でない有限生成部分加群
$M' \subset M$ を任意に一つ選ぶ。例えば、零でない一元で生成される部分加群を選べばよい。
補題 \ref{algebra-lemma-support-zero} により、$\text{Supp}(M')$ は空でない。
命題 \ref{algebra-proposition-minimal-primes-associated-primes} により、
これは $\text{Ass}(M')$ が空でないことを含意する。
補題 \ref{algebra-lemma-ass} により、これは
$\text{Ass}(M) \not = \emptyset$ を含意する。
\end{proof}

\begin{lemma}
\label{algebra-lemma-ass-minimal-prime-support}
$R$ をネーター環とする。
$M$ を $R$-加群とする。
$\text{Supp}(M)$ の元のうちで極小である任意の
$\mathfrak p \in \text{Supp}(M)$ は $\text{Ass}(M)$ の元である。
\end{lemma}

\begin{proof}
$M$ が有限 $R$-加群ならば、命題
\ref{algebra-proposition-minimal-primes-associated-primes} から従う。
一般には、$M = \bigcup M_\lambda$ と有限部分加群の合併として書き、
$\text{Supp}(M) = \bigcup \text{Supp}(M_\lambda)$ および
$\text{Ass}(M) = \bigcup \text{Ass}(M_\lambda)$ であることを用いる。
\end{proof}

\begin{lemma}
\label{algebra-lemma-ass-zero-divisors}
$R$ をネーター環とする。
$M$ を $R$-加群とする。
合併 $\bigcup_{\mathfrak q \in \text{Ass}(M)} \mathfrak q$ は、
$M$ 上で零因子となる $R$ の元全体の集合である。
\end{lemma}

\begin{proof}
任意の随伴素イデアルに属する元は、明らかに $M$ 上の零因子である。
逆に、$x \in R$ が $M$ 上の零因子であると仮定する。
部分加群 $N = \{m \in M \mid xm = 0\}$ を考える。
$N$ は零でないので、補題 \ref{algebra-lemma-ass-zero} により
随伴素イデアル $\mathfrak q$ をもつ。
このとき $x \in \mathfrak q$ であり、補題 \ref{algebra-lemma-ass} により
$\mathfrak q$ は $M$ の随伴素イデアルである。
\end{proof}

\begin{lemma}
\label{algebra-lemma-one-equation-module}
% P01-E0394: the frozen English says Let R is rather than Let R be.
$R$ をネーター局所環、$M$ を有限 $R$-加群、
$f \in \mathfrak m$ を $R$ の極大イデアルの元とする。このとき
$$
\dim(\text{Supp}(M/fM)) \leq
\dim(\text{Supp}(M)) \leq
\dim(\text{Supp}(M/fM)) + 1
$$
$M$ の台のどの極小素イデアルにも $f$ が含まれなければ
（例えば $f$ が $M$ 上の非零因子ならば）、右側の不等式では等号が成り立つ。
\end{lemma}

\begin{proof}
（括弧内の主張は補題 \ref{algebra-lemma-ass-zero-divisors} から従う。）
第一の不等式は
$\text{Supp}(M/fM) \subset \text{Supp}(M)$ から従う
（補題 \ref{algebra-lemma-support-quotient} 参照）。
第二の不等式については、
$\text{Supp}(M/fM) = \text{Supp}(M) \cap V(f)$ であることに注意する
（補題 \ref{algebra-lemma-support-quotient} 参照）。
例えば補題 \ref{algebra-lemma-filter-primes-in-support} と次元の初等的な性質から、
$R$ の素イデアル $\mathfrak p$ に対して
$\dim V(\mathfrak p) \leq \dim (V(\mathfrak p) \cap V(f)) + 1$
を示せば十分であることが従う。これは補題
\ref{algebra-lemma-one-equation} の帰結である。
最後に、$M$ の台のどの極小素イデアルにも $f$ が含まれなければ、
$\text{Supp}(M/fM)$ の素イデアルの鎖はいずれも、
$\text{Supp}(M)$ において極大になるまでに少なくともあと一段階を要する鎖を与える。
\end{proof}

\begin{lemma}
\label{algebra-lemma-ass-functorial}
$\varphi : R \to S$ を環写像とする。
$M$ を $S$-加群とする。
このとき
$\Spec(\varphi)(\text{Ass}_S(M)) \subset \text{Ass}_R(M)$ である。
\end{lemma}

\begin{proof}
$\mathfrak q \in \text{Ass}_S(M)$ ならば、$S$ における $m$ の零化イデアルが
$\mathfrak q$ であるような $M$ の元 $m$ が存在する。
このとき $R$ における $m$ の零化イデアルは $\mathfrak q \cap R$ である。
\end{proof}


\begin{remark}
\label{algebra-remark-ass-reverse-functorial}
$\varphi : R \to S$ を環写像とする。
$M$ を $S$-加群とする。このとき常に
$\Spec(\varphi)(\text{Ass}_S(M)) \supset \text{Ass}_R(M)$
となるとは限らない。例えば環写像
$R = k \to S = k[x_1, x_2, x_3, \ldots]/(x_i^2)$ と $M = S$ を考える。
このとき $\text{Ass}_R(M)$ は空でないが、
$\text{Ass}_S(S)$ は空である。
\end{remark}

\begin{lemma}
\label{algebra-lemma-ass-functorial-Noetherian}
$\varphi : R \to S$ を環写像とする。
$M$ を $S$-加群とする。$S$ がネーター環ならば、
$\Spec(\varphi)(\text{Ass}_S(M)) = \text{Ass}_R(M)$ である。
\end{lemma}

\begin{proof}
補題 \ref{algebra-lemma-ass-functorial} で既に
$\Spec(\varphi)(\text{Ass}_S(M)) \subset \text{Ass}_R(M)$
を示した。逆の包含について、
素イデアル $\mathfrak p \in \text{Ass}_R(M)$ を選ぶ。
$R$ における $m$ の零化イデアルが $\mathfrak p$ となる元
$m \in M$ をとる。
$I = \{g \in S \mid gm = 0\}$ を $S$ における $m$ の零化イデアルとする。
このとき $R/\mathfrak p \subset S/I$ は単射である。
補題 \ref{algebra-lemma-injective-minimal-primes-in-image} と
\ref{algebra-lemma-minimal-prime-image-minimal-prime} を組み合わせると、
$I$ の上で極小であり $\mathfrak p$ に写る素イデアル
$\mathfrak q \subset S$ が存在することが分かる。
命題 \ref{algebra-proposition-minimal-primes-associated-primes} により、
$\mathfrak q$ は $S/I$ の随伴素イデアルである。
したがって補題 \ref{algebra-lemma-ass} により
$\mathfrak q$ は $M$ の随伴素イデアルであり、主張を得る。
\end{proof}

\begin{lemma}
\label{algebra-lemma-ass-quotient-ring}
$R$ を環とする。
$I$ をイデアルとする。
$M$ を $R/I$-加群とする。
標準的な単射
$\Spec(R/I) \to \Spec(R)$
を通じて $\text{Ass}_{R/I}(M) = \text{Ass}_R(M)$ である。
\end{lemma}

\begin{proof}
省略する。
\end{proof}

\begin{lemma}
\label{algebra-lemma-associated-primes-localize}
$R$ を環とする。
$M$ を $R$-加群とする。
$\mathfrak p \subset R$ を素イデアルとする。
\begin{enumerate}
\item $\mathfrak p \in \text{Ass}(M)$ ならば
$\mathfrak pR_{\mathfrak p} \in \text{Ass}(M_{\mathfrak p})$ である。
\item $\mathfrak p$ が有限生成ならば、逆も成り立つ。
\end{enumerate}
\end{lemma}

\begin{proof}
$\mathfrak p \in \text{Ass}(M)$ ならば、零化イデアルが
$\mathfrak p$ である元 $m \in M$ が存在する。
局所化は完全なので
（命題 \ref{algebra-proposition-localization-exact}）、
$M_{\mathfrak p}$ における $m/1$ の零化イデアルは
$\mathfrak pR_{\mathfrak p}$ であり、したがって (1) が成り立つ。
$\mathfrak pR_{\mathfrak p} \in \text{Ass}(M_{\mathfrak p})$ および
$\mathfrak p = (f_1, \ldots, f_n)$ を仮定する。
$m/g$ を、零化イデアルが $\mathfrak pR_{\mathfrak p}$ である
$M_{\mathfrak p}$ の元とする。
これは $m$ の零化イデアルが $\mathfrak p$ に含まれることを含意する。
$M_{\mathfrak p}$ で $f_i m/g = 0$ なので、
$g_i \in R$, $g_i \not \in \mathfrak p$ であって
$M$ で $g_i f_i m = 0$ となるものが存在する。
合わせると、$g_1\ldots g_nm$ の零化イデアルは $\mathfrak p$ である。
したがって $\mathfrak p \in \text{Ass}(M)$ である。
\end{proof}

\begin{lemma}
\label{algebra-lemma-localize-ass}
$R$ を環とし、$M$ を $R$-加群とする。
$S \subset R$ を積閉部分集合とする。
標準的な単射 $\Spec(S^{-1}R) \to \Spec(R)$ を通じて、次が成り立つ。
\begin{enumerate}
\item $\text{Ass}_R(S^{-1}M) = \text{Ass}_{S^{-1}R}(S^{-1}M)$,
\item
$\text{Ass}_R(M) \cap \Spec(S^{-1}R) \subset \text{Ass}_R(S^{-1}M)$,
\item $R$ がネーター環ならば、この包含は等号である。
\end{enumerate}
\end{lemma}

\begin{proof}
$m \in S^{-1}M$ に対し、$I \subset R$ と $J \subset S^{-1}R$ を
$m$ の零化イデアルとする。このとき $I$ は写像
$R \to S^{-1}R$ による $J$ の逆像であり、$J = S^{-1}I$ である。
(1) の等式は、補題 \ref{algebra-lemma-spec-localization} における写像
$\Spec(S^{-1}R) \to \Spec(R)$ の記述から従う。
$m \in M$ に対し、$I \subset R$ を $R$ における $m$ の零化イデアル、
$J \subset S^{-1}R$ を $m/1 \in S^{-1}M$ の零化イデアルとする。
$J = S^{-1}I$ であり、これから (2) が従う。
ネーターの場合の等号は補題
\ref{algebra-lemma-associated-primes-localize} から従う。
実際、
% P01-E0395: the frozen formula types the prime ideal as an element of R.
$\mathfrak p \in R$, $S \cap \mathfrak p = \emptyset$ ならば
$M_{\mathfrak p} = (S^{-1}M)_{S^{-1}\mathfrak p}$ である。
\end{proof}

\begin{lemma}
\label{algebra-lemma-localize-ass-nonzero-divisors}
$R$ を環とし、$M$ を $R$-加群とする。
$S \subset R$ を積閉部分集合とする。
すべての $s \in S$ が $M$ 上の非零因子であると仮定する。
このとき
$$
\text{Ass}_R(M) = \text{Ass}_R(S^{-1}M).
$$
\end{lemma}

\begin{proof}
仮定により $M \subset S^{-1}M$ なので、補題 \ref{algebra-lemma-ass} から包含
$\text{Ass}(M) \subset \text{Ass}(S^{-1}M)$ を得る。
逆に、$n/s \in S^{-1}M$ を、零化イデアルが素イデアル
$\mathfrak p$ である元とする。このとき
$n \in M$ の零化イデアルも $\mathfrak p$ である。
\end{proof}

\begin{lemma}
\label{algebra-lemma-ideal-nonzerodivisor}
$R$ を極大イデアル $\mathfrak m$ をもつネーター局所環とする。
$I \subset \mathfrak m$ をイデアルとし、$M$ を有限 $R$-加群とする。
次の条件は同値である。
\begin{enumerate}
\item $M$ 上で零因子でない $x \in I$ が存在する。
\item すべての $\mathfrak q \in \text{Ass}(M)$ に対して
$I \not \subset \mathfrak q$ である。
\end{enumerate}
\end{lemma}

\begin{proof}
非零因子 $x$ が $I$ に存在すれば、
$x$ は明らかに $M$ のどの随伴素イデアルにも属し得ない。
逆に、すべての $\mathfrak q \in \text{Ass}(M)$ に対して
$I \not \subset \mathfrak q$ であると仮定する。
このとき補題 \ref{algebra-lemma-finite-ass} と \ref{algebra-lemma-silly} により、
すべての $\mathfrak q \in \text{Ass}(M)$ に対して
$x \not \in \mathfrak q$ となる $x \in I$ を選ぶことができる。
補題 \ref{algebra-lemma-ass-zero-divisors} により、元 $x$ は
$M$ 上の零因子ではない。
\end{proof}

\begin{lemma}
\label{algebra-lemma-zero-at-ass-zero}
$R$ を環とし、$M$ を $R$-加群とする。$R$ がネーター環ならば、写像
$$
M
\longrightarrow
\prod\nolimits_{\mathfrak p \in \text{Ass}(M)} M_{\mathfrak p}
$$
は単射である。
\end{lemma}

\begin{proof}
$x \in M$ を写像の核の元とする。
$\mathfrak p$ が $Rx \subset M$ の随伴素イデアルならば、一方で
$\mathfrak p \in \text{Ass}(M)$ であり
（補題 \ref{algebra-lemma-ass}）、他方で
$(Rx)_{\mathfrak p} \subset M_{\mathfrak p}$ は零でない。
この矛盾から $\text{Ass}(Rx) = \emptyset$ が分かる。
したがって補題 \ref{algebra-lemma-ass-zero} により $Rx = 0$ である。
\end{proof}

\noindent
% P01-E0396: the frozen source renders this internal placement note to readers.
この補題はおそらく別の場所に置くべきである。

\begin{lemma}
\label{algebra-lemma-dim-not-zero-exists-nonzerodivisor-nonunit}
$k$ を体とし、$S$ を有限型 $k$-代数とする。
$\dim(S) > 0$ ならば、非零因子かつ非単元である
$f \in S$ が存在する。
\end{lemma}

\begin{proof}
補題 \ref{algebra-lemma-finite-ass} により、環 $S$ の随伴素イデアルは有限個である。
補題 \ref{algebra-lemma-finite-type-algebra-finite-nr-primes} により、
環 $S$ は無限個の極大イデアルをもつ。
したがって $S$ の随伴素イデアルではない極大イデアル
$\mathfrak m \subset S$ を選べる。
素イデアル回避（補題 \ref{algebra-lemma-silly}）により、
$S$ のどの随伴素イデアルにも含まれない
零でない $f \in \mathfrak m$ を選べる。
補題 \ref{algebra-lemma-ass-zero-divisors} により元 $f$ は非零因子であり、
$f \in \mathfrak m$ なので $f$ は単元でない。
\end{proof}


\section{記号冪}
\label{algebra-section-symbolic-power}

\noindent
定義を述べる。

\begin{definition}
\label{algebra-definition-symbolic-power}
$R$ を環とし、$\mathfrak p$ を素イデアルとする。$n \geq 0$ に対し、
$\mathfrak p$ の $n$ 次 {\it 記号冪}とは、イデアル
$\mathfrak p^{(n)} = \Ker(R \to R_\mathfrak p/\mathfrak p^nR_\mathfrak p)$
をいう。
\end{definition}

\noindent
$\mathfrak p^n \subset \mathfrak p^{(n)}$ であるが、
等号は常に成り立つとは限らないことに注意する。

\begin{lemma}
\label{algebra-lemma-symbolic-power-associated}
$R$ をネーター環とする。
$\mathfrak p$ を素イデアルとする。
$n > 0$ とする。このとき
$\text{Ass}(R/\mathfrak p^{(n)}) = \{\mathfrak p\}$ である。
\end{lemma}

\begin{proof}
$\mathfrak q$ が $R/\mathfrak p^{(n)}$ の随伴素イデアルならば、
明らかに $\mathfrak p \subset \mathfrak q$ である。
一方、任意の $x \in R$, $x \not \in \mathfrak p$ は
$R/\mathfrak p^{(n)}$ 上の非零因子である。
実際、$y \in R$ かつ
$xy \in \mathfrak p^{(n)} = R \cap \mathfrak p^nR_{\mathfrak p}$
ならば $y \in \mathfrak p^nR_{\mathfrak p}$ であり、
したがって $y \in \mathfrak p^{(n)}$ である。これで補題が従う。
\end{proof}

\begin{lemma}
\label{algebra-lemma-symbolic-power-flat-extension}
% P01-E0397: the frozen English is missing the article before flat ring map.
$R \to S$ を平坦な環写像とする。
$\mathfrak p \subset R$ を、$\mathfrak q = \mathfrak p S$ が
$S$ の素イデアルとなるような素イデアルとする。
このとき $\mathfrak p^{(n)} S = \mathfrak q^{(n)}$ である。
\end{lemma}

\begin{proof}
$\mathfrak p^{(n)} = \Ker(R \to R_\mathfrak p/\mathfrak p^nR_\mathfrak p)$
なので、平坦性を用いると、$\mathfrak p^{(n)} S$ は写像
$S \to S_\mathfrak p/\mathfrak p^nS_\mathfrak p$ の核であることが分かる。
一方、$\mathfrak q^{(n)}$ は写像
$S \to S_\mathfrak q/\mathfrak q^nS_\mathfrak q =
S_\mathfrak q/\mathfrak p^nS_\mathfrak q$ の核である。
したがって
$$
S_\mathfrak p/\mathfrak p^nS_\mathfrak p
\longrightarrow
S_\mathfrak q/\mathfrak p^nS_\mathfrak q
$$
が単射であることを示せば十分である。
右辺の加群は、左辺の加群を
$f \in S$, $f \not \in \mathfrak q$ となる元で局所化したものであることに注意する。
したがって、これらの元が
$S_\mathfrak p/\mathfrak p^nS_\mathfrak p$ 上の非零因子であることを示せば十分である。
平坦性により、加群
$S_\mathfrak p/\mathfrak p^nS_\mathfrak p$ は、次の部分商をもつ
有限フィルトレーションをもつ。
$$
\mathfrak p^iS_\mathfrak p/\mathfrak p^{i + 1}S_\mathfrak p
\cong \mathfrak p^iR_\mathfrak p/\mathfrak p^{i + 1}R_\mathfrak p
\otimes_{R_\mathfrak p} S_\mathfrak p \cong
V \otimes_{\kappa(\mathfrak p)} (S/\mathfrak q)_\mathfrak p
$$
ここで $V$ は $\kappa(\mathfrak p)$-ベクトル空間である。
したがって
% P01-E0398: the frozen proof says invertibly where it establishes injectively.
$f$ は望みどおり可逆に作用する。
\end{proof}


\section{相対随伴素イデアル集合}
\label{algebra-section-relative-assassin}

\noindent
相対随伴素イデアル集合について論じる。
$R \to S$ を環写像とし、$N$ を $S$-加群とする。
この状況で、$S$ の素イデアル $\mathfrak q$ の次の集合を導入できる。
\begin{enumerate}
\item $A$: $\mathfrak p = R \cap \mathfrak q$ とおくと
$\mathfrak q \in \text{Ass}_S(N \otimes_R \kappa(\mathfrak p))$ となるもの,
\item $A'$: $\mathfrak p = R \cap \mathfrak q$ とおくと、
% P01-E0400: the frozen fibre-ring subscript omits the base R.
$\mathfrak q$ が、標準写像
$\Spec(S \otimes_R \kappa(\mathfrak p)) \to \Spec(S)$ による
$\text{Ass}_{S \otimes \kappa(\mathfrak p)}(N \otimes_R \kappa(\mathfrak p))$
の像に属するもの,
\item $A_{fin}$: $\mathfrak p = R \cap \mathfrak q$ とおくと
$\mathfrak q \in \text{Ass}_S(N/\mathfrak pN)$ となるもの,
\item $A'_{fin}$: ある素イデアル $\mathfrak p' \subset R$ に対して
$\mathfrak q \in \text{Ass}_S(N/\mathfrak p'N)$ となるもの,
\item $B$: ある $R$-加群 $M$ に対して
$\mathfrak q \in \text{Ass}_S(N \otimes_R M)$ となるもの,
\item $B_{fin}$: ある有限 $R$-加群 $M$ に対して
$\mathfrak q \in \text{Ass}_S(N \otimes_R M)$ となるもの。
\end{enumerate}
これらの集合の間の関係のいくつかを決定しよう。

\begin{lemma}
\label{algebra-lemma-compare-relative-assassins}
$R \to S$ を環写像とし、$N$ を $S$-加群とする。
% P01-E0399: the frozen declaration omits A'_fin from its five-name list.
$A$, $A'$, $A_{fin}$, $B$, $B_{fin}$ を、
上で導入した $\Spec(S)$ の部分集合とする。
\begin{enumerate}
\item 常に $A = A'$ である。
\item 常に $A_{fin} \subset A$,
$B_{fin} \subset B$, $A_{fin} \subset A'_{fin} \subset B_{fin}$
および $A \subset B$ である。
\item $S$ がネーター環ならば、$A = A_{fin}$ および $B = B_{fin}$ である。
\item $N$ が $R$ 上平坦ならば、
$A = A_{fin} = A'_{fin}$ および $B = B_{fin}$ である。
\item $R$ がネーター環で $N$ が $R$ 上平坦ならば、すべての集合は等しい。
すなわち $A = A' = A_{fin} = A'_{fin} = B = B_{fin}$ である。
\end{enumerate}
\end{lemma}

\begin{proof}
証明の議論の一部は、後の、これより精密な補題の証明でも繰り返される
（それらはすべての加群や素イデアルの集まりではなく、与えられた加群 $M$ または
与えられた素イデアル $\mathfrak p$ を扱うからである）。

\medskip\noindent
(1) の証明。$\mathfrak p$ を $R$ の素イデアルとする。このとき
$$
\text{Ass}_S(N \otimes_R \kappa(\mathfrak p)) =
\text{Ass}_{S/\mathfrak pS}(N \otimes_R \kappa(\mathfrak p)) =
\text{Ass}_{S \otimes_R \kappa(\mathfrak p)}(N \otimes_R \kappa(\mathfrak p))
$$
である。第一の等式は補題 \ref{algebra-lemma-ass-quotient-ring}、
第二の等式は補題 \ref{algebra-lemma-localize-ass} の (1) による。
% P01-E0401: the frozen English has subject-verb disagreement here.
これで $A = A'$ が示された。
包含 $A_{fin} \subset A'_{fin}$ は明らかである。

\medskip\noindent
(2) の証明。各包含は定義から直ちに従う。ただし
$A_{fin} \subset A$ だけは、補題 \ref{algebra-lemma-localize-ass} と、
$A_{fin}$ の定式化で $\mathfrak p = R \cap \mathfrak q$ を要求したことから従う。

\medskip\noindent
(3) の証明。$S$ がネーター環ならば、等式
$A = A_{fin}$ は補題 \ref{algebra-lemma-localize-ass} の (3) から従う。
$\mathfrak q = (g_1, \ldots, g_m)$ を $S$ の有限生成素イデアルとする。
$z \in N \otimes_R M$ を、零化イデアルが $\mathfrak q$ である元とする。
$z$ が $z' \in N \otimes_R M'$ の像となるような有限部分加群
$M' \subset M$ を選べる。このとき
$\text{Ann}_S(z') \subset \mathfrak q = \text{Ann}_S(z)$ である。
$N \otimes_R -$ は余極限と可換し、$M$ は有限 $R$-加群の有向余極限なので、
$M' \subset M'' \subset M$ であって、像
$z'' \in N \otimes_R M''$ が $g_1, \ldots, g_m$ により
零化されるものを見つけられる。
したがって $\text{Ann}_S(z'') = \mathfrak q$ である。
これで、$S$ がネーター環ならば $B = B_{fin}$ であることが示された。

\medskip\noindent
(4) の証明。$N$ が平坦ならば、関手 $N \otimes_R -$ は完全である。
特に $M' \subset M$ ならば
$N \otimes_R M' \subset N \otimes_R M$ である。
したがって $z \in N \otimes_R M$ の零化イデアル
$\mathfrak q = \text{Ann}_S(z)$ が素イデアルならば、
$z \in N \otimes_R M'$ となる任意の有限 $R$-部分加群
$M' \subset M$ を選ぶことができ、$N \otimes_R M'$ の元としての
$z$ の零化イデアルも $\mathfrak q$ に等しい。
したがって $B = B_{fin}$ である。
$\mathfrak p'$ を $R$ の素イデアルとし、$\mathfrak q$ を
$N/\mathfrak p'N$ の随伴素イデアルであるような $S$ の素イデアルとする。
これは $\mathfrak p'S \subset \mathfrak q$ を含意する。
$N$ は $R$ 上平坦なので、$N/\mathfrak p'N$ は整域
$R/\mathfrak p'$ 上平坦である。
したがって $R/\mathfrak p'$ の各非零元は、
% P01-E0402: the frozen module quotient omits the final N.
$N/\mathfrak p'$ 上の非零因子である。
よってこれらの元の像はどれも $\mathfrak q$ に属し得ず、
$\mathfrak p' = R \cap \mathfrak q$ と結論できる。
したがって $A_{fin} = A'_{fin}$ である。
最後に補題 \ref{algebra-lemma-localize-ass-nonzero-divisors} により
$\text{Ass}_S(N/\mathfrak p'N) =
\text{Ass}_S(N \otimes_R \kappa(\mathfrak p'))$ であり、
すなわち $A'_{fin} = A$ である。

\medskip\noindent
(5) の証明。(4) でほかの等式を示したので、
$A'_{fin} = B_{fin}$ だけを示せばよい。
$M$ を有限 $R$-加群とする。
補題 \ref{algebra-lemma-filter-Noetherian-module} により、
次のような $R$-部分加群によるフィルトレーションが存在する。
$$
0 = M_0 \subset M_1 \subset \ldots \subset M_n = M
$$
ここで各商 $M_i/M_{i-1}$ は、$R$ のある素イデアル
$\mathfrak p_i$ に対する $R/\mathfrak p_i$ と同型である。
$N$ は平坦なので、次のような $S$-部分加群によるフィルトレーションを得る。
$$
0 = N \otimes_R M_0 \subset N \otimes_R M_1 \subset \ldots \subset
N \otimes_R M_n = N \otimes_R M
$$
各部分商は $N/\mathfrak p_iN$ と同型である。
補題 \ref{algebra-lemma-ass} により
$\text{Ass}_S(N \otimes_R M) \subset \bigcup \text{Ass}_S(N/\mathfrak p_iN)$
と結論できる。したがって $B_{fin} \subset A'_{fin}$ である。
逆の包含は (2) に含まれるので、主張を得る。
\end{proof}

\noindent
$S/R$ 上の $N$ の相対随伴素イデアル集合を、上の集合
$A = A'$ と定義する。動機として、これはファイバー環上のファイバー加群
$N \otimes_R \kappa(\mathfrak p)$ だけに依存することを指摘しておく。
加群の随伴素イデアルの場合と同様、この概念はファイバー環
$S \otimes_R \kappa(\mathfrak p)$ がネーターであるとき、
例えば $R \to S$ が有限型であるときに最もよく機能することを注意しておく。

\begin{definition}
\label{algebra-definition-relative-assassin}
$R \to S$ を環写像とし、$N$ を $S$-加群とする。
{\it $S/R$ 上の $N$ の相対随伴素イデアル集合}とは、集合
$$
\text{Ass}_{S/R}(N)
=
\{ \mathfrak q \subset S \mid
\mathfrak q \in \text{Ass}_S(N \otimes_R \kappa(\mathfrak p))
\text{ with }\mathfrak p = R \cap \mathfrak q\}.
$$
をいう。これは補題 \ref{algebra-lemma-compare-relative-assassins} で
$A$ と名づけた集合である。
\end{definition}


\noindent
次のいくつかの結果は、一見しただけではそう見えないこともあるが、
相対随伴素イデアル集合に関するものである。

\begin{lemma}
\label{algebra-lemma-bourbaki}
$R \to S$ を環写像とする。
$M$ を $R$-加群、$N$ を $S$-加群とする。
$N$ が $R$-加群として平坦ならば
$$
\text{Ass}_S(M \otimes_R N)
\supset
\bigcup\nolimits_{\mathfrak p \in \text{Ass}_R(M)} \text{Ass}_S(N/\mathfrak pN)
$$
であり、$R$ がネーター環ならば等号が成り立つ。
\end{lemma}

\begin{proof}
$\mathfrak p \in \text{Ass}_R(M)$ ならば、単射
$R/\mathfrak p \to M$ が存在する。
$N$ は $R$ 上平坦なので、単射
$R/\mathfrak p \otimes_R N \to M \otimes_R N$ を得る。
$R/\mathfrak p \otimes_R N = N/\mathfrak pN$ なので、
補題 \ref{algebra-lemma-ass} により
$\text{Ass}_S(N/\mathfrak pN) \subset \text{Ass}_S(M \otimes_R N)$
と結論できる。したがって右辺は左辺に含まれる。

\medskip\noindent
$M = \bigcup M_\lambda$ と有限生成 $R$-部分加群の合併として書く。
このとき $N$ は $R$-平坦なので、
$N \otimes_R M = \bigcup N \otimes_R M_\lambda$ でもある。
随伴素イデアルの定義により
$\text{Ass}_S(N \otimes_R M) = \bigcup \text{Ass}_S(N \otimes_R M_\lambda)$
および $\text{Ass}_R(M) = \bigcup \text{Ass}(M_\lambda)$ である。
したがって $M$ が有限生成であると仮定してよい。

\medskip\noindent
$\mathfrak q \in \text{Ass}_S(M \otimes_R N)$ とし、
$R$ がネーター環で $M$ が有限 $R$-加群であると仮定する。
証明を終えるには、$\mathfrak q$ が右辺の元であることを示さなければならない。
まず、補題 \ref{algebra-lemma-associated-primes-localize} により
$\mathfrak qS_{\mathfrak q} \in
\text{Ass}_{S_{\mathfrak q}}((M \otimes_R N)_{\mathfrak q})$
であることに注意する。
$\mathfrak p$ を $R$ の対応する素イデアルとする。次の等式に注意する。
$$
(M \otimes_R N)_{\mathfrak q} = M \otimes_R N_{\mathfrak q}
= M_{\mathfrak p} \otimes_{R_{\mathfrak p}} N_{\mathfrak q}
$$
もし
$\mathfrak pR_{\mathfrak p} \not \in
\text{Ass}_{R_{\mathfrak p}}(M_{\mathfrak p})$
ならば、$M_{\mathfrak p}$ 上の非零因子である元
$x \in \mathfrak pR_{\mathfrak p}$ が存在する
（補題 \ref{algebra-lemma-ideal-nonzerodivisor} 参照）。
$N_{\mathfrak q}$ は $R_{\mathfrak p}$ 上平坦なので、
$\mathfrak qS_{\mathfrak q}$ における $x$ の像は
$(M \otimes_R N)_{\mathfrak q}$ 上の非零因子である。
これは
% P01-E0403: the frozen contradiction formula uses Ass_S instead of Ass_{S_q}.
$\mathfrak qS_{\mathfrak q} \in \text{Ass}_S((M \otimes_R N)_{\mathfrak q})$
という仮定に矛盾する。
したがって $\mathfrak p$ は $M$ の随伴素イデアルの一つである。

\medskip\noindent
議論を続け、次のようなフィルトレーションを選ぶ。
$$
0 = M_0 \subset M_1 \subset \ldots \subset M_n = M
$$
ここで各商 $M_i/M_{i-1}$ は、$R$ のある素イデアル
$\mathfrak p_i$ に対する $R/\mathfrak p_i$ と同型である
（補題 \ref{algebra-lemma-filter-Noetherian-module} 参照）。
（補題 \ref{algebra-lemma-ass-filter} により、少なくとも一つの $i$ に対して
$\mathfrak p_i = \mathfrak p$ である。）
これから次のフィルトレーションを得る。
$$
0 = M_0 \otimes_R N \subset M_1 \otimes_R N \subset \ldots
\subset M_n \otimes_R N = M \otimes_R N
$$
その部分商は $N/\mathfrak p_iN$ と同型である。
$\mathfrak p_i \not = \mathfrak p$ ならば、
前段落の結果により、$\mathfrak q$ は加群
$N/\mathfrak p_iN$ に随伴し得ない
（$\text{Ass}_R(R/\mathfrak p_i) = \{\mathfrak p_i\}$ だからである）。
したがって望みどおり、$\mathfrak q$ は
$N/\mathfrak pN$ に随伴すると結論できる。
\end{proof}

\begin{lemma}
\label{algebra-lemma-post-bourbaki}
$R \to S$ を環写像とする。
$N$ を $S$-加群とする。
$N$ が $R$-加群として平坦であり、$R$ が分数体 $K$ をもつ整域であると仮定する。
このとき
$$
\text{Ass}_S(N) =
\text{Ass}_S(N \otimes_R K) =
\text{Ass}_{S \otimes_R K}(N \otimes_R K)
$$
である。ただし、標準的な包含
$\Spec(S \otimes_R K) \subset \Spec(S)$ を用いている。
\end{lemma}

\begin{proof}
$S \otimes_R K = (R \setminus \{0\})^{-1}S$ および
$N \otimes_R K = (R \setminus \{0\})^{-1}N$ であることに注意する。
零でない任意の $x \in R$ に対し、$N$ は $R$ 上平坦なので、
$N$ 上の $x$ による乗法は単射である。
したがって、補題 \ref{algebra-lemma-localize-ass-nonzero-divisors} と
補題 \ref{algebra-lemma-localize-ass} の (1) を組み合わせれば主張が従う。
\end{proof}

\begin{lemma}
\label{algebra-lemma-bourbaki-fibres}
$R \to S$ を環写像とする。
$M$ を $R$-加群、$N$ を $S$-加群とする。
$N$ が $R$-加群として平坦であると仮定する。このとき
$$
\text{Ass}_S(M \otimes_R N)
\supset
\bigcup\nolimits_{\mathfrak p \in \text{Ass}_R(M)}
\text{Ass}_{S \otimes_R \kappa(\mathfrak p)}(N \otimes_R \kappa(\mathfrak p))
$$
である。ここでは注意 \ref{algebra-remark-fundamental-diagram} を用いて、
ファイバー環のスペクトルを $\Spec(S)$ の部分集合とみなす。
$R$ がネーター環ならば、この包含は等号である。
\end{lemma}

\begin{proof}
これは補題 \ref{algebra-lemma-bourbaki} と同値である。
そのことは補題 \ref{algebra-lemma-ass-quotient-ring},
\ref{algebra-lemma-flat-base-change}, \ref{algebra-lemma-post-bourbaki} による。
\end{proof}

\begin{remark}
\label{algebra-remark-bourbaki}
$R \to S$ を環写像とし、$N$ を $S$-加群とする。
$\mathfrak p$ を $R$ の素イデアルとする。このとき
$$
\text{Ass}_S(N \otimes_R \kappa(\mathfrak p)) =
\text{Ass}_{S/\mathfrak pS}(N \otimes_R \kappa(\mathfrak p)) =
\text{Ass}_{S \otimes_R \kappa(\mathfrak p)}(N \otimes_R \kappa(\mathfrak p)).
$$
% P01-E0404: the frozen explanatory sentence is a finite-verb fragment.
第一の等式は補題 \ref{algebra-lemma-ass-quotient-ring} により、
第二の等式は補題 \ref{algebra-lemma-localize-ass} の (1) により成り立つ。
\end{remark}

\section{弱随伴素イデアル}
\label{algebra-section-weakly-ass}

\noindent
これは随伴素イデアルの概念の変種であり、非ネーター環や有限でない加群に対して有用である。

\begin{definition}
\label{algebra-definition-weakly-associated}
$R$ を環とし、$M$ を $R$-加群とする。
$R$ の素イデアル $\mathfrak p$ が $M$ に {\it 弱随伴する} とは、
ある元 $m \in M$ が存在して、$\mathfrak p$ が零化イデアル
$\text{Ann}(m) = \{f \in R \mid fm = 0\}$
を含む素イデアルのうち極小であることをいう。
このような素イデアル全体の集合を $\text{WeakAss}_R(M)$
または $\text{WeakAss}(M)$ と表す。
\end{definition}

\noindent
したがって、随伴素イデアルは弱随伴素イデアルである。
素イデアルにおける局所化を用いた特徴づけを次に与える。

\begin{lemma}
\label{algebra-lemma-weakly-ass-local}
$R$ を環、$M$ を $R$-加群とし、
$\mathfrak p$ を $R$ の素イデアルとする。
次の条件は同値である。
\begin{enumerate}
\item $\mathfrak p$ は $M$ に弱随伴する。
\item $\mathfrak pR_{\mathfrak p}$ は $M_{\mathfrak p}$ に弱随伴する。
\item $M_{\mathfrak p}$ は、零化イデアルの根基が
$\mathfrak pR_{\mathfrak p}$ に等しい元を含む。
\end{enumerate}
\end{lemma}

\begin{proof}
(1) を仮定する。このとき、ある元 $m \in M$ が存在して、
$\mathfrak p$ は $m$ の零化イデアル
$I = \{x \in R \mid xm = 0\}$ を含む素イデアルのうち極小である。
局所化は完全なので、$M_{\mathfrak p}$ における $m$ の零化イデアルは
$I_{\mathfrak p}$ である。
したがって $\mathfrak pR_{\mathfrak p}$ は、
$M_{\mathfrak p}$ における $m$ の零化イデアル
$I_{\mathfrak p}$ を含む $R_{\mathfrak p}$ の極小素イデアルである。
ゆえに (2) が成り立ち、さらに
$\sqrt{I_{\mathfrak p}} = \mathfrak pR_{\mathfrak p}$
であるから (3) も成り立つ。

\medskip\noindent
$R_{\mathfrak p}$ 上の $M_{\mathfrak p}$ に (1) $\Rightarrow$ (3)
を適用すれば、(2) $\Rightarrow$ (3) が従う。

\medskip\noindent
最後に (3) を仮定する。すなわち、ある元
$m/f \in M_{\mathfrak p}$ が存在し、その零化イデアルの根基が
$\mathfrak pR_{\mathfrak p}$ に等しい。
$M$ における $m$ の零化イデアル
$I = \{x \in R \mid xm = 0\}$ は
$\sqrt{I_{\mathfrak p}} = \mathfrak pR_{\mathfrak p}$ を満たす。
これは明らかに、$\mathfrak p$ が $I$ を含み、かつ
$I$ を含む素イデアルのうち極小であることを意味する。すなわち (1) が成り立つ。
\end{proof}

\begin{lemma}
\label{algebra-lemma-reduced-weakly-ass-minimal}
被約環の弱随伴素イデアルは、その極小素イデアルである。
\end{lemma}

\begin{proof}
$(R, \mathfrak m)$ を被約局所環とする。
零化イデアルの根基が $\mathfrak m$ である元 $x \in R$ を考える。
$\mathfrak m \not = 0$ ならば $x$ は単元ではあり得ないので、
$x \in \mathfrak m$ である。したがって特に、ある $n \geq 0$ に対して
$x^{1 + n} = 0$ である。ゆえに $x = 0$ である。
% P01-E0405: the frozen English continues with a fragment beginning “Which”.
これは $x$ の零化イデアルが $\mathfrak m$ に含まれるという仮定に矛盾する。
したがって $\mathfrak m = 0$、すなわち $R$ は体である。
補題 \ref{algebra-lemma-weakly-ass-local} により、これから補題の主張が従う。
\end{proof}

\begin{lemma}
\label{algebra-lemma-weakly-ass}
$R$ を環とし、
$0 \to M' \to M \to M'' \to 0$ を $R$-加群の短完全列とする。
このとき $\text{WeakAss}(M') \subset \text{WeakAss}(M)$ かつ
$\text{WeakAss}(M) \subset \text{WeakAss}(M') \cup \text{WeakAss}(M'')$
である。
\end{lemma}

\begin{proof}
補題 \ref{algebra-lemma-weakly-ass-local} による弱随伴素イデアルの特徴づけを用いる。
$\mathfrak p$ を $R$ の素イデアルとする。
局所化は完全なので、短完全列
$0 \to M'_{\mathfrak p} \to M_{\mathfrak p} \to M''_{\mathfrak p} \to 0$
を得る。
$m \in M_{\mathfrak p}$ の零化イデアルの根基が
$\mathfrak pR_{\mathfrak p}$ であると仮定する。
$M''_{\mathfrak p}$ における $m$ の像 $\overline{m}$ が零であって
$m \in M'_{\mathfrak p}$ であるか、または $\overline{m}$ の
零化イデアルの根基が $\mathfrak pR_{\mathfrak p}$ である。
これにより
$\text{WeakAss}(M) \subset \text{WeakAss}(M') \cup \text{WeakAss}(M'')$
が示される。包含
$\text{WeakAss}(M') \subset \text{WeakAss}(M)$ は定義から直ちに従う。
\end{proof}

\begin{lemma}
\label{algebra-lemma-weakly-ass-zero}
\begin{slogan}
零でない任意の加群は弱随伴素イデアルをもつ。
\end{slogan}
$R$ を環、$M$ を $R$-加群とする。このとき
$$
M = (0) \Leftrightarrow \text{WeakAss}(M) = \emptyset
$$
である。
\end{lemma}

\begin{proof}
$M = (0)$ ならば、定義により
$\text{WeakAss}(M) = \emptyset$ である。
逆に $M \not = 0$ と仮定し、零でない元 $m \in M$ を選ぶ。
$m$ の零化イデアルを
$I = \{x \in R \mid xm = 0\}$ と書く。このとき $R/I \subset M$ である。
したがって補題 \ref{algebra-lemma-weakly-ass} により
$\text{WeakAss}(R/I) \subset \text{WeakAss}(M)$ である。
一方 $I \not = R$ なので、$V(I) = \Spec(R/I)$ は極小素イデアルを含む。
補題 \ref{algebra-lemma-Zariski-topology} および
\ref{algebra-lemma-spec-closed} を参照せよ。これで証明が完了する。
\end{proof}

\begin{lemma}
\label{algebra-lemma-weakly-ass-support}
$R$ を環、$M$ を $R$-加群とする。このとき
$$
\text{Ass}(M) \subset \text{WeakAss}(M) \subset \text{Supp}(M).
$$
\end{lemma}

\begin{proof}
第一の包含は定義から直ちに従う。
$\mathfrak p \in \text{WeakAss}(M)$ ならば、
補題 \ref{algebra-lemma-weakly-ass-local} により
$M_{\mathfrak p} \not = 0$ であるから、
$\mathfrak p \in \text{Supp}(M)$ である。
\end{proof}

\begin{lemma}
\label{algebra-lemma-weakly-ass-zero-divisors}
$R$ を環、$M$ を $R$-加群とする。
和集合 $\bigcup_{\mathfrak q \in \text{WeakAss}(M)} \mathfrak q$
は、$M$ 上の零因子である $R$ の元全体の集合である。
\end{lemma}

\begin{proof}
$f \in \mathfrak q \in \text{WeakAss}(M)$ と仮定する。
このとき、ある元 $m \in M$ が存在して、
$\mathfrak q$ は $I = \{x \in R \mid xm = 0\}$ 上極小である。
したがって、$g \in R$、$g \not \in \mathfrak q$ および $n > 0$
が存在して $f^ngm = 0$ となる。
$g \not \in I$ なので $gm \not = 0$ であることに注意する。
上のような $n$ を最小に取れば、
$f (f^{n - 1}gm) = 0$ かつ $f^{n - 1}gm \not = 0$ であるから、
$f$ は $M$ 上の零因子である。
逆に $f \in R$ が $M$ 上の零因子であると仮定する。
部分加群 $N = \{m \in M \mid fm = 0\}$ を考える。
$N$ は零でないので、補題 \ref{algebra-lemma-weakly-ass-zero} により
弱随伴素イデアル $\mathfrak q$ をもつ。
明らかに $f \in \mathfrak q$ であり、補題
\ref{algebra-lemma-weakly-ass} により
$\mathfrak q$ は $M$ の弱随伴素イデアルである。
\end{proof}

\begin{lemma}
\label{algebra-lemma-weakly-ass-minimal-prime-support}
$R$ を環、$M$ を $R$-加群とする。
$\text{Supp}(M)$ の元のうち極小である任意の
$\mathfrak p \in \text{Supp}(M)$ は
$\text{WeakAss}(M)$ の元である。
\end{lemma}

\begin{proof}
$\Spec(R_{\mathfrak p})$ において
$\text{Supp}(M_{\mathfrak p}) = \{\mathfrak pR_{\mathfrak p}\}$
であることに注意する。特に $M_{\mathfrak p}$ は零でないので、
補題 \ref{algebra-lemma-weakly-ass-zero} により
$\text{WeakAss}(M_{\mathfrak p}) \not = \emptyset$ である。
補題 \ref{algebra-lemma-weakly-ass-support} により
$\text{WeakAss}(M_{\mathfrak p}) \subset \text{Supp}(M_{\mathfrak p})$
だから、
$\text{WeakAss}(M_{\mathfrak p}) = \{\mathfrak pR_{\mathfrak p}\}$
と結論できる。したがって補題 \ref{algebra-lemma-weakly-ass-local} により
$\mathfrak p \in \text{WeakAss}(M)$ である。
\end{proof}

\begin{lemma}
\label{algebra-lemma-ass-weakly-ass}
$R$ を環、$M$ を $R$-加群とする。
$\mathfrak p$ を有限生成な $R$ の素イデアルとする。このとき
$$
\mathfrak p \in \text{Ass}(M) \Leftrightarrow
\mathfrak p \in \text{WeakAss}(M).
$$
特に $R$ がネーター環ならば
$\text{Ass}(M) = \text{WeakAss}(M)$ である。
\end{lemma}

\begin{proof}
ある $g_i \in R$ によって
$\mathfrak p = (g_1, \ldots, g_n)$ と書く。
% P01-E0406: the frozen English has “enough the prove” instead of “enough to prove”.
逆向きの含意は一般に成り立つので、含意
「$\Leftarrow$」を示せば十分である。補題
\ref{algebra-lemma-weakly-ass-support} を参照せよ。
$\mathfrak p \in \text{WeakAss}(M)$ と仮定する。
補題 \ref{algebra-lemma-weakly-ass-local} により、ある元
$m \in M_{\mathfrak p}$ が存在して、
$I = \{x \in R_{\mathfrak p} \mid xm = 0\}$ の根基は
$\mathfrak pR_{\mathfrak p}$ である。
したがって各 $i$ に対して、$M_{\mathfrak p}$ において
$g_i^{e_i}m = 0$ となる最小の $e_i > 0$ が存在する。
ある $i$ に対して $e_i > 1$ ならば、$m$ を
$g_i^{e_i - 1} m \not = 0$ で置き換えることにより
$\sum e_i$ を減らせる。
ゆえに $m \in M_{\mathfrak p}$ の零化イデアルを
$(g_1, \ldots, g_n)R_{\mathfrak p} = \mathfrak p R_{\mathfrak p}$
としてよい。補題 \ref{algebra-lemma-associated-primes-localize} により
$\mathfrak p \in \text{Ass}(M)$ が従う。
\end{proof}

\begin{remark}
\label{algebra-remark-weakly-ass-not-functorial}
$\varphi : R \to S$ を環写像とし、$M$ を $S$-加群とする。
随伴素イデアルの場合とは異なり（補題
\ref{algebra-lemma-ass-functorial} 参照）、必ずしも
$\Spec(\varphi)(\text{WeakAss}_S(M)) \subset \text{WeakAss}_R(M)$
とはならない。
例として、環写像
$$
R = k[x_1, x_2, x_3, \ldots] \to
S = k[x_1, x_2, x_3, \ldots, y_1, y_2, y_3, \ldots]/
(x_1y_1, x_2y_2, x_3y_3, \ldots)
$$
と $M = S$ を考える。この場合
$\mathfrak q = \sum x_iS$ は $S$ の極小素イデアルであり、
したがって $M = S$ の弱随伴素イデアルである
（補題 \ref{algebra-lemma-weakly-ass-minimal-prime-support} 参照）。
しかし一方で、$S$ の任意の零でない元の $R$ における零化イデアルは
有限生成であり、したがって、その根基が
$R \cap \mathfrak q = (x_1, x_2, x_3, \ldots)$
に等しくなることはない（詳細は省略する）。
\end{remark}

\begin{lemma}
\label{algebra-lemma-weakly-ass-reverse-functorial}
$\varphi : R \to S$ を環写像とし、$M$ を $S$-加群とする。このとき
$\Spec(\varphi)(\text{WeakAss}_S(M)) \supset \text{WeakAss}_R(M)$
である。
\end{lemma}

\begin{proof}
$\mathfrak p$ を $\text{WeakAss}_R(M)$ の元とする。
このとき、ある $m \in M_{\mathfrak p}$ が存在して、その零化イデアル
$I = \{x \in R_{\mathfrak p} \mid xm = 0\}$ の根基は
$\mathfrak pR_{\mathfrak p}$ である。
$S_{\mathfrak p}$ における $m$ の零化イデアル
$J = \{x \in S_{\mathfrak p} \mid xm = 0 \}$ を考える。
$IS_{\mathfrak p} \subset J$ なので、$J$ 上の任意の極小素イデアル
$\mathfrak q \subset S_{\mathfrak p}$ は $\mathfrak p$ 上にある。
さらに、例えば補題 \ref{algebra-lemma-weakly-ass-local} により、
このような $\mathfrak q$ は $M$ の弱随伴素イデアルに対応する。
\end{proof}

\begin{remark}
\label{algebra-remark-ass-functorial}
$\varphi : R \to S$ を環写像、$M$ を $S$-加群とする。
% P01-E0407: both frozen English naming clauses omit “by” after “Denote”.
スペクトル上の対応する写像を $f : \Spec(S) \to \Spec(R)$ と表す。
このとき
$$
f(\text{Ass}_S(M)) \subset
\text{Ass}_R(M) \subset
\text{WeakAss}_R(M) \subset
f(\text{WeakAss}_S(M))
$$
である。補題 \ref{algebra-lemma-ass-functorial}、
\ref{algebra-lemma-weakly-ass-reverse-functorial}、および
\ref{algebra-lemma-weakly-ass-support} を参照せよ。
一般にはすべての包含が真に狭くなり得る。注意
\ref{algebra-remark-ass-reverse-functorial} および
\ref{algebra-remark-weakly-ass-not-functorial} を参照せよ。
$S$ がネーター環ならば、補題 \ref{algebra-lemma-ass-weakly-ass} により
外側の二つが等しいので、すべての包含は等号となる。
\end{remark}

\begin{lemma}
\label{algebra-lemma-weakly-ass-finite-ring-map}
$\varphi : R \to S$ を環写像、$M$ を $S$-加群とする。
% P01-E0407: second occurrence of the same frozen English naming defect.
スペクトル上の対応する写像を $f : \Spec(S) \to \Spec(R)$ と表す。
$\varphi$ が有限環写像ならば
$$
\text{WeakAss}_R(M) = f(\text{WeakAss}_S(M)).
$$
\end{lemma}

\begin{proof}
一方の包含はすでに示されている。注意
\ref{algebra-remark-ass-functorial} を参照せよ。
他方を示すため、$\mathfrak q \in \text{WeakAss}_S(M)$ と仮定し、
$\mathfrak p$ を対応する $R$ の素イデアルとする。
ある元 $m \in M$ が存在して、$\mathfrak q$ は
$J = \{g \in S \mid gm = 0\}$ 上の極小素イデアルであるとする。
したがって $JS_{\mathfrak q}$ の根基は
$\mathfrak qS_{\mathfrak q}$ である。
$R \to S$ は有限なので、$\mathfrak p$ 上には有限個の素イデアル
$\mathfrak q = \mathfrak q_1, \mathfrak q_2, \ldots, \mathfrak q_l$
がある。補題 \ref{algebra-lemma-finite-finite-fibres} を参照せよ。
補題 \ref{algebra-lemma-silly} により、
$i > 1$ に対して $x \not \in \mathfrak q_i$ となる
$x \in \mathfrak q$ を選ぶ。
上記により、ある元 $y \in S$、$y \not \in \mathfrak q$ と
整数 $t > 0$ が存在して $y x^t m = 0$ となる。
したがって元 $ym \in M$ は $x^t$ により零化されるので、
$ym$ は $M_{\mathfrak q_i}$、$i = 2, \ldots, l$ において零に写る。
% P01-E0408: the frozen source writes S_q although ym is a module element.
なお、$ym$ は $S_{\mathfrak q}$ において零には写らない。

\medskip\noindent
有限環写像に対する上昇定理により、環 $S_{\mathfrak p}$ は
極大イデアル $\mathfrak q_i S_{\mathfrak p}$ をもつ半局所環である。
補題 \ref{algebra-lemma-integral-going-up} を参照せよ。
$f \in \mathfrak pR_{\mathfrak p}$ ならば、$f$ のある冪は
$JS_{\mathfrak q}$ に入る。したがって、ある $t > 0$ に対して
$f^t ym$ は $M_{\mathfrak q}$ において零に写る。
$ym$ は $S_{\mathfrak p}$ の他の極大イデアルで消えるので、
補題 \ref{algebra-lemma-characterize-zero-local} により
$f^t ym$ は $M_{\mathfrak p}$ において零である。
このようにして、$\mathfrak p$ は $R$ における $ym$ の
零化イデアル上の極小素イデアルであることが分かり、証明が完了する。
\end{proof}

\begin{lemma}
\label{algebra-lemma-weakly-ass-quotient-ring}
$R$ を環、$I$ をイデアル、$M$ を $R/I$-加群とする。
標準的な単射
$\Spec(R/I) \to \Spec(R)$
を介して
$\text{WeakAss}_{R/I}(M) = \text{WeakAss}_R(M)$
である。
\end{lemma}

\begin{proof}
補題 \ref{algebra-lemma-weakly-ass-finite-ring-map} の特別な場合である。
\end{proof}

\begin{lemma}
\label{algebra-lemma-localize-weakly-ass}
$R$ を環、$M$ を $R$-加群とし、
$S \subset R$ を乗法的部分集合とする。
標準的な単射 $\Spec(S^{-1}R) \to \Spec(R)$ を介して
$\text{WeakAss}_R(S^{-1}M) = \text{WeakAss}_{S^{-1}R}(S^{-1}M)$
であり、さらに
$$
\text{WeakAss}(M) \cap \Spec(S^{-1}R) = \text{WeakAss}(S^{-1}M).
$$
\end{lemma}

\begin{proof}
$m \in S^{-1}M$ とする。
$I = \{x \in R \mid xm = 0\}$ および
$I' = \{x' \in S^{-1}R \mid x'm = 0\}$ とおく。
このとき $I' = S^{-1}I$ であり、$I = R$ でない限り
$I \cap S = \emptyset$ である（確認は省略する）。
したがって、$I'$ 上の $S^{-1}R$ の極小素イデアルは、
$I$ 上の $R$ の極小素イデアルで $S$ を避けるものと一対一に対応する。
これにより等式
$\text{WeakAss}_R(S^{-1}M) = \text{WeakAss}_{S^{-1}R}(S^{-1}M)$
が示される。
補題 \ref{algebra-lemma-weakly-ass-local} と、
% P01-E0409: the frozen source types the prime ideal as p in R.
$\mathfrak p \in R$、$S \cap \mathfrak p = \emptyset$ に対する
$M_{\mathfrak p} = (S^{-1}M)_{S^{-1}\mathfrak p}$ から
第二の等式が従う。
\end{proof}

\begin{lemma}
\label{algebra-lemma-localize-weakly-ass-nonzero-divisors}
$R$ を環、$M$ を $R$-加群とし、
$S \subset R$ を乗法的部分集合とする。
すべての $s \in S$ が $M$ 上の非零因子であると仮定する。このとき
$$
\text{WeakAss}(M) = \text{WeakAss}(S^{-1}M).
$$
\end{lemma}

\begin{proof}
仮定により $M \subset S^{-1}M$ なので、補題
\ref{algebra-lemma-weakly-ass} から
$\text{WeakAss}(M) \subset \text{WeakAss}(S^{-1}M)$
を得る。
逆に、$n/s \in S^{-1}M$ が零化イデアル $I$ をもち、
$\mathfrak p$ が $I$ 上極小な素イデアルであると仮定する。
このとき $n \in M$ の零化イデアルは $I$ であり、
$\mathfrak p$ は $I$ 上極小な素イデアルである。
\end{proof}

\begin{lemma}
\label{algebra-lemma-zero-at-weakly-ass-zero}
$R$ を環、$M$ を $R$-加群とする。写像
$$
M
\longrightarrow
\prod\nolimits_{\mathfrak p \in \text{WeakAss}(M)} M_{\mathfrak p}
$$
は単射である。
\end{lemma}

\begin{proof}
$x \in M$ をこの写像の核の元とし、
$N = Rx \subset M$ とおく。
$\mathfrak p$ が $N$ の弱随伴素イデアルならば、一方で
$\mathfrak p \in \text{WeakAss}(M)$
（補題 \ref{algebra-lemma-weakly-ass}）であり、他方で
$N_{\mathfrak p} \subset M_{\mathfrak p}$ は零でない。
これは矛盾なので $\text{WeakAss}(N) = \emptyset$ である。
したがって補題 \ref{algebra-lemma-weakly-ass-zero} により
$N = 0$、すなわち $x = 0$ である。
\end{proof}

\begin{lemma}
\label{algebra-lemma-weak-post-bourbaki}
$R \to S$ を環写像とし、$N$ を $S$-加群とする。
$N$ が $R$-加群として平坦であり、$R$ が分数体 $K$ をもつ
整域であると仮定する。このとき
$$
\text{WeakAss}_S(N) = \text{WeakAss}_{S \otimes_R K}(N \otimes_R K)
$$
である。ただし、標準的な包含
$\Spec(S \otimes_R K) \subset \Spec(S)$
を用いている。
\end{lemma}

\begin{proof}
$S \otimes_R K = (R \setminus \{0\})^{-1}S$ および
$N \otimes_R K = (R \setminus \{0\})^{-1}N$
であることに注意する。
零でない任意の $x \in R$ に対し、$N$ は $R$ 上平坦なので、
$N$ 上の $x$ による乗法は単射である。
したがって補題
\ref{algebra-lemma-localize-weakly-ass-nonzero-divisors} から主張が従う。
\end{proof}

\begin{lemma}
\label{algebra-lemma-weakly-ass-change-fields}
$K/k$ を体拡大、$R$ を $k$-代数、$M$ を $R$-加群とする。
$\mathfrak q \subset R \otimes_k K$ を
$\mathfrak p \subset R$ 上にある素イデアルとする。
$\mathfrak q$ が $M \otimes_k K$ に弱随伴するならば、
$\mathfrak p$ は $M$ に弱随伴する。
\end{lemma}

\begin{proof}
$z \in M \otimes_k K$ を、$\mathfrak q$ が $z$ の零化イデアル
$J \subset R \otimes_k K$ 上極小となるような元とする。
$z \in M \otimes_k L$ となる有限生成部分拡大 $K/L/k$ を選ぶ。
$R \otimes_k L \to R \otimes_k K$ は平坦なので、
$J = I(R \otimes_k K)$ である。ここで
$I \subset R \otimes_k L$ は小さい方の環における $z$ の
零化イデアルである
（補題 \ref{algebra-lemma-annihilator-flat-base-change}）。
したがって平坦降下により
$\mathfrak q \cap (R \otimes_k L)$ は $I$ 上極小である
（補題 \ref{algebra-lemma-flat-going-down}）。
このようにして、次段落で述べる場合に帰着する。

\medskip\noindent
$K/k$ が有限生成体拡大であると仮定する。
Fields, Section \ref{fields-section-transcendence} を参照し、
$x_1, \ldots, x_r \in K$ を $k$ 上の $K$ の超越基底とする。
$L = k(x_1, \ldots, x_r)$ とおき、$[K : L] = n$ とする。
このとき $R \otimes_k L \to R \otimes_k K$ は有限環写像である。
したがって補題 \ref{algebra-lemma-weakly-ass-finite-ring-map} により
$\mathfrak q \cap (R \otimes_k L)$ は、
$R \otimes_k L$-加群とみなした $M \otimes_k K$ の
弱随伴素イデアルである。
$R \otimes_k L$-加群として
$M \otimes_k K \cong (M \otimes_k L)^{\oplus n}$
なので、$\mathfrak q \cap (R \otimes_k L)$ は
$M \otimes_k L$ の弱随伴素イデアルである
（例えば補題 \ref{algebra-lemma-weakly-ass} と帰納法を用いる）。
このようにして、次段落で論じる場合に帰着する。

\medskip\noindent
$K = k(x_1, \ldots, x_r)$ が純超越体拡大であると仮定する。
$R$ を $R_\mathfrak p$ で、$M$ を $M_\mathfrak p$ で、
$\mathfrak q$ を $\mathfrak q(R_\mathfrak p \otimes_k K)$
で置き換えてよい。補題 \ref{algebra-lemma-localize-weakly-ass} を参照せよ。
このようにして、次段落で論じる場合に帰着する。

\medskip\noindent
$K = k(x_1, \ldots, x_r)$ が純超越体拡大であり、
$R$ が極大イデアル $\mathfrak p$ をもつ局所環であると仮定する。
$f \in R \otimes_k K$、$f \not \in \mathfrak p(R \otimes_k K)$
ならば、この元は $M \otimes_k K$ 上の非零因子であると主張する。
実際、$z \in M \otimes_k K$ を元とする。
ある有限 $R$-部分加群 $M' \subset M$ が存在して、
$z \in M' \otimes_k K$ であり、かつ $M'$ はこの性質に関して極小である。
これを見るには、$K$ の $k$-ベクトル空間としての基底
$\{t_\alpha\}$ を選び、$z = \sum m_\alpha \otimes t_\alpha$ と書き、
$m_\alpha$ が生成する $R$-部分加群を $M'$ とすればよい。
$z \in \mathfrak p(M' \otimes_k K) = \mathfrak p M' \otimes_k K$
ならば、$\mathfrak pM' = M'$ であり、補題
\ref{algebra-lemma-NAK} により $M' = 0$ となって矛盾する。
したがって $z$ は
$M'/\mathfrak p M' \otimes_k K$ において非零の像
$\overline{z}$ をもつ。
% P01-E0410: the frozen polynomial-ring display retains the stale index n.
しかし $R/\mathfrak p \otimes_k K$ は
$\kappa(\mathfrak p)[x_1, \ldots, x_n]$ の局所化として整域であり、
$M'/\mathfrak p M' \otimes_k K$ は自由加群であるから、
$f\overline{z} \not = 0$ である。これで主張が示された。

\medskip\noindent
最後に、$\mathfrak q$ が $z$ の零化イデアル
$J \subset R \otimes_k K$ 上極小となるような
$z \in M \otimes_k K$ を選ぶ。
$f \in \mathfrak p$ に対し、ある $n \geq 1$ と
$g \in R \otimes_k K$、$g \not \in \mathfrak q$ が存在して
% P01-E0411: frozen membership is ill-typed; the sidecar proposes g f^n in J.
$g f^n z \in J$、すなわち $g f^n z = 0$ となる。
（これは、$\mathfrak q$ が $\mathfrak p$ 上にあり、
$\mathfrak q$ が $J$ 上極小だからである。）
上で $g$ は非零因子であることを見たので $f^n z = 0$ である。
これは、$R$-加群とみなした $M \otimes_k K$ に
$\mathfrak p$ が弱随伴することを意味する。
$M \otimes_k K$ は $M$ のコピーの直和なので、
前と同様に $\mathfrak p$ は $M$ に弱随伴すると結論する。
\end{proof}

\section{埋没素イデアル}
\label{algebra-section-embedded-primes}

\noindent
まず定義を与える。

\begin{definition}
\label{algebra-definition-embedded-primes}
$R$ を環、$M$ を $R$-加群とする。
\begin{enumerate}
\item $M$ の随伴素イデアルのうち、$M$ の随伴素イデアル全体の中で
極小でないものを、$M$ の {\it 埋没随伴素イデアル} という。
\item {\it $R$ の埋没素イデアル} とは、$R$ を $R$-加群とみなしたときの
埋没随伴素イデアルをいう。
\end{enumerate}
\end{definition}

\noindent
これらを取り除く方法を次に述べる。

\begin{lemma}
\label{algebra-lemma-remove-embedded-primes}
$R$ をネーター環、$M$ を有限 $R$-加群とする。
$R$-部分加群の集合
$$
\{
K \subset M
\mid
\text{Supp}(K)
\text{ nowhere dense in }
\text{Supp}(M)
\}.
$$
を考える。この集合は最大元 $K$ をもち、商
$M' = M/K$ は次の性質をもつ。
\begin{enumerate}
\item $\text{Supp}(M) = \text{Supp}(M')$。
\item $M'$ は埋没随伴素イデアルをもたない。
\item $M$ のすべての埋没随伴素イデアルに含まれる任意の
$f \in R$ に対して $M_f \cong M'_f$ である。
\end{enumerate}
\end{lemma}

\begin{proof}
以下、補題 \ref{algebra-lemma-finite-ass} と命題
\ref{algebra-proposition-minimal-primes-associated-primes} を断りなく用いる。
$\mathfrak q_1, \ldots, \mathfrak q_t$ を $M$ の台の極小素イデアル、
$\mathfrak p_1, \ldots, \mathfrak p_s$ を $M$ の
埋没随伴素イデアルとする。このとき
$\text{Ass}(M) = \{\mathfrak q_j, \mathfrak p_i\}$ である。
次のようにおく。
$$
K = \{m \in M \mid \text{Supp}(Rm) \subset \bigcup V(\mathfrak p_i)\}
$$
これは直ちに部分加群であると分かる。
$M$ はネーター環上有限なので $K$ も有限である。
したがって $\text{Supp}(K)$ は $\text{Supp}(M)$ において
稠密な部分を含まない。
$K' \subset M$ を、台が $\text{Supp}(M)$ において
稠密な部分を含まない別の部分加群とする。これは
% P01-E0412: the frozen localization drops the prime mark from K'.
$K_{\mathfrak q_j} = 0$ を意味する。
したがって $m \in K'$ ならば、$m$ は
$M_{\mathfrak q_j}$ において零に写り、その結果
$(Rm)_{\mathfrak q_j} = 0$ となる。
一方 $\text{Ass}(Rm) \subset \text{Ass}(M)$ である。
ゆえに $Rm$ の台は $\bigcup V(\mathfrak p_i)$ に含まれる。
したがって $m \in K$ であり、$m$ は $K'$ の任意の元だったので
$K' \subset K$ である。

\medskip\noindent
$M' = M/K$ とおく。$K_{\mathfrak q_j}=0$ なので、すべての $j$ に対して
$M'_{\mathfrak q_j} = M_{\mathfrak q_j}$ である。
したがって $M$ と $M'$ は同じ台をもつ。

\medskip\noindent
$\mathfrak q = \text{Ann}(\overline{m}) \in \text{Ass}(M')$
と仮定する。ここで $\overline{m} \in M'$ は $m \in M$ の像である。
このとき $m \not \in K$ なので、$Rm$ の台は
$\mathfrak q_j$ の一つを含まなければならない。
$M_{\mathfrak q_j} = M'_{\mathfrak q_j}$ であるから、
$\overline{m}$ は $M'_{\mathfrak q_j}$ において零には写らない。
したがって $\mathfrak q \subset \mathfrak q_j$
（実際には等号が成り立つ）である。
これは、$M'$ のすべての随伴素イデアルが埋没していないことを意味する。

\medskip\noindent
$f$ をすべての $\mathfrak p_i$ に含まれる元とする。このとき
% P01-E0413: the frozen formula uses lowercase supp and compares with 0.
$D(f) \cap \text{supp}(K) = 0$ である。
ゆえに $K_f = 0$ なので $M_f = M'_f$ である。
\end{proof}

\begin{lemma}
\label{algebra-lemma-remove-embedded-primes-localize}
$R$ をネーター環、$M$ を有限 $R$-加群とする。
任意の $f \in R$ に対して $(M')_f = (M_f)'$ である。
ここで $M \to M'$ および $M_f \to (M_f)'$ は、
補題 \ref{algebra-lemma-remove-embedded-primes} で構成した商である。
\end{lemma}

\begin{proof}
省略する。
\end{proof}

\begin{lemma}
\label{algebra-lemma-no-embedded-primes-endos}
$R$ をネーター環とし、$M$ を埋没随伴素イデアルをもたない
有限 $R$-加群とする。
$I = \{x \in R \mid xM = 0\}$ とおく。このとき環 $R/I$ は
埋没素イデアルをもたない。
\end{lemma}

\begin{proof}
$R$ を $R/I$ で置き換えてよい。
したがって、$R$ のすべての零でない元が $M$ に非自明に作用すると
仮定してよい。
補題 \ref{algebra-lemma-support-closed} により、これは
$\Spec(R)$ が $M$ の台に等しいことを意味する。
$\mathfrak p$ が $R$ の埋没素イデアルであると仮定する。
零化イデアルが $\mathfrak p$ である元 $x \in R$ を取る。
零でない加群 $N = xM \subset M$ を考える。
これは $\mathfrak p$ により零化される。
したがって $N$ の任意の随伴素イデアル $\mathfrak q$ は
$\mathfrak p$ を含み、かつ $M$ の随伴素イデアルでもある。
すると $\mathfrak q$ は $M$ の埋没随伴素イデアルとなり、
補題の仮定に矛盾する。
\end{proof}

\section{正則列}
\label{algebra-section-regular-sequences}

\noindent
この節では正則列の基本的な性質をいくつか展開する。

\begin{definition}
\label{algebra-definition-regular-sequence}
$R$ を環、$M$ を $R$-加群とする。
$R$ の元の列 $f_1, \ldots, f_r$ が {\it $M$-正則列} であるとは、
次の条件が成り立つことをいう。
\begin{enumerate}
\item 各 $i = 1, \ldots, r$ に対して、$f_i$ は
$M/(f_1, \ldots, f_{i - 1})M$ 上の非零因子である。
\item 加群 $M/(f_1, \ldots, f_r)M$ は零でない。
\end{enumerate}
$I$ が $R$ のイデアルで $f_1, \ldots, f_r \in I$ ならば、
$f_1, \ldots, f_r$ を {\it $I$ 内の $M$-正則列} という。
$M = R$ のときは、$f_1, \ldots, f_r$ を単に
（$I$ 内の）{\it 正則列} という。
\end{definition}

\noindent
この定義が元 $f_1, \ldots, f_r$ の順序に依存することに注意されたい
（下の例を参照）。
論文や書籍によっては、加群 $M/(f_1, \ldots, f_r)M$ が
零でないという要件を課さない。この流儀には、正則列であるという性質が
局所化で保たれるという利点がある。
しかし本書では主として、$R$ が局所環の場合に加群の深さを定義するために
この定義を用いる。その場合 $f_i$ は極大イデアルに属することを要求されるが、
この条件は $R$ から局所化 $R_\mathfrak p$ に移る際には保たれない。

\begin{example}
\label{algebra-example-global-regular}
$k$ を体とする。環 $k[x, y, z]$ において
列 $x, y(1-x), z(1-x)$ は正則であるが、
列 $y(1-x), z(1-x), x$ は正則でない。
\end{example}

\begin{example}
\label{algebra-example-local-regular}
$k$ を体とする。環
$k[x, y, w_0, w_1, w_2, \ldots]/I$
を考える。ここで $I$ は
$yw_i$、$i = 0, 1, 2, \ldots$ と
$w_i - xw_{i + 1}$、$i = 0, 1, 2, \ldots$
によって生成される。
列 $x, y$ は正則であるが、$y$ は零因子である。
さらに、極大イデアル $(x, y, w_i)$ で局所化しても例が得られる。
\end{example}

\begin{lemma}
\label{algebra-lemma-permute-xi}
$R$ をネーター局所環、$M$ を有限 $R$-加群とする。
$x_1, \ldots, x_c$ を $M$-正則列とする。このとき
% P01-E0414: the frozen conclusion drops the module qualifier M.
$x_i$ の任意の置換も正則列である。
\end{lemma}

\begin{proof}
まず $c = 2$ の場合を扱う。
$K \subset M$ を $x_2 : M \to M$ の核とする。
任意の $z \in K$ に対し、$x_2$ は $M/x_1M$ 上の非零因子なので、
ある $z' \in M$ によって $z = x_1 z'$ と書ける。
$x_1$ は $M$ 上の非零因子だから $x_2 z' = 0$ でもある。
したがって $x_1 : K \to K$ は全射である。
ゆえに中山の補題 \ref{algebra-lemma-NAK} により $K = 0$ である。
次に $M/x_2M$ 上の $x_1$ による乗法を考える。
$z \in M$ の像 $\overline{z} \in M/x_2M$ がこの写像の核に入るならば、
ある $y \in M$ に対して $x_1 z = x_2 y$ である。
しかし $x_1, x_2$ は正則列なので、ある $y' \in M$ に対して
$y = x_1 y'$ である。したがって
$x_1 ( z - x_2 y' ) =0$、ゆえに $z = x_2 y'$ であり、
望みどおり $\overline{z} = 0$ となる。

\medskip\noindent
一般の場合、任意の置換は隣り合う添字の互換の合成であることに注意する。
したがって
$$
x_1, \ldots, x_{i-2}, x_i, x_{i-1}, x_{i + 1}, \ldots, x_c
$$
が $M$-正則列であることを示せば十分である。
これは、加群
% P01-E0415: the frozen quotient omits the final module factor M.
$M/(x_1, \ldots, x_{i-2})$
と長さ $2$ の正則列 $x_{i-1}, x_i$ に、先ほど示した場合を
適用すれば従う。
\end{proof}

\begin{lemma}
\label{algebra-lemma-flat-increases-depth}
\begin{slogan}
平坦な局所環準同型は正則列を保存し、かつ反映する。
\end{slogan}
$R, S$ を局所環、$R \to S$ を平坦な局所環準同型とする。
$x_1, \ldots, x_r$ を $R$ 内の列、$M$ を $R$-加群とする。
次の条件は同値である。
\begin{enumerate}
\item $x_1, \ldots, x_r$ は $R$ 内の $M$-正則列である。
\item $S$ における $x_1, \ldots, x_r$ の像は
$M \otimes_R S$-正則列をなす。
\end{enumerate}
\end{lemma}

\begin{proof}
補題 \ref{algebra-lemma-local-flat-ff} により
$R \to S$ は忠実平坦だからである。
\end{proof}

\begin{lemma}
\label{algebra-lemma-regular-sequence-in-neighbourhood}
$R$ をネーター環、$M$ を有限 $R$-加群とする。
$\mathfrak p$ を素イデアルとする。
$x_1, \ldots, x_r$ を $R$ 内の列とし、その $R_{\mathfrak p}$ における像が
$M_{\mathfrak p}$-正則列をなすと仮定する。
このとき、ある $g \in R$、$g \not \in \mathfrak p$ が存在して、
$R_g$ における $x_1, \ldots, x_r$ の像は
$M_g$-正則列をなす。
\end{lemma}

\begin{proof}
次のようにおく。
$$
K_i = \Ker\left(x_i : M/(x_1, \ldots, x_{i - 1})M \to
M/(x_1, \ldots, x_{i - 1})M\right).
$$
これは有限 $R$-加群であり、仮定により $\mathfrak p$ における局所化は
零である。したがって、ある $g \in R$、$g \not \in \mathfrak p$ が存在して、
すべての $i = 1, \ldots, r$ に対して $(K_i)_g = 0$ となる。
この $g$ が求めるものである。
\end{proof}

\begin{lemma}
\label{algebra-lemma-join-regular-sequences}
$A$ を環とし、$I$ を $A$ 内の正則列
$f_1, \ldots, f_n$ が生成するイデアルとする。
$g_1, \ldots, g_m \in A$ を、その像
$\overline{g}_1, \ldots, \overline{g}_m$ が
$A/I$ 内で正則列をなすような元とする。
このとき $f_1, \ldots, f_n, g_1, \ldots, g_m$ は
$A$ 内の正則列である。
\end{lemma}

\begin{proof}
定義から直ちに従う。
\end{proof}

\begin{lemma}
\label{algebra-lemma-regular-sequence-short-exact-sequence}
$R$ を環とし、$0 \to M_1 \to M_2 \to M_3 \to 0$
を $R$-加群の短完全列とする。
$f_1, \ldots, f_r \in R$ とする。
$f_1, \ldots, f_r$ が $M_1$-正則かつ $M_3$-正則ならば、
$f_1, \ldots, f_r$ は $M_2$-正則である。
\end{lemma}

\begin{proof}
補題 \ref{algebra-lemma-snake} により、
$f_1 : M_1 \to M_1$ と $f_1 : M_3 \to M_3$ が単射ならば
$f_1 : M_2 \to M_2$ も単射であり、短完全列
$$
0 \to M_1/f_1M_1 \to M_2/f_1M_2 \to M_3/f_1M_3 \to 0
$$
を得る。これと $r$ に関する帰納法から補題が従う。
細部は省略する。
\end{proof}

\begin{lemma}
\label{algebra-lemma-regular-sequence-powers}
$R$ を環、$M$ を $R$-加群とする。
$f_1, \ldots, f_r \in R$ とし、
$e_1, \ldots, e_r > 0$ を整数とする。
このとき $f_1, \ldots, f_r$ が $M$-正則列であることと、
$f_1^{e_1}, \ldots, f_r^{e_r}$ が $M$-正則列であることは同値である。
\end{lemma}

\begin{proof}
$r$ に関する帰納法で示す。$r = 1$ の場合は、次の二つの簡単な事実から従う。
(a) $M$ 上の非零因子の冪は $M$ 上の非零因子であり、
(b) $M$ 上の非零因子の約元は $M$ 上の非零因子である。
$r > 1$ とする。$M/f_1M$ に帰納法を適用すれば、
$f_1, f_2, \ldots, f_r$ が $M$-正則列であることと
$f_1, f_2^{e_2}, \ldots, f_r^{e_r}$ が $M$-正則列であることは同値である。
したがって、$e > 0$ が与えられたとき、
$f_1^e, f_2, \ldots, f_r$ が $M$-正則列であることと
$f_1, \ldots, f_r$ が $M$-正則列であることが同値だと示せば十分である。
これを $e$ に関する帰納法で示す。$e = 1$ の場合は自明である。
いずれの仮定の下でも $f_1$ は非零因子なので
（$r = 1$ の場合による）、短完全列
$$
0 \to M/f_1M \xrightarrow{f_1^{e - 1}} M/f_1^eM \to M/f_1^{e - 1}M \to 0
$$
を得る。
$f_1, f_2, \ldots, f_r$ が $M$-正則列であると仮定する。
帰納法により、元 $f_2, \ldots, f_r$ は
$M/f_1M$-正則列かつ $M/f_1^{e - 1}M$-正則列である。
補題 \ref{algebra-lemma-regular-sequence-short-exact-sequence} により、
$f_2, \ldots, f_r$ は $M/f_1^eM$-正則である。
したがって $f_1^e, f_2, \ldots, f_r$ は $M$-正則である。
逆に、$f_1^e, f_2, \ldots, f_r$ が $M$-正則列であると仮定する。
このとき $f_2 : M/f_1^eM \to M/f_1^eM$ は単射なので、
$f_2 : M/f_1M \to M/f_1M$ は単射である。
ゆえに帰納法（！）により
$f_2 : M/f_1^{e - 1}M \to M/f_1^{e - 1}M$ は単射であり、
$$
0 \to
M/(f_1, f_2)M \xrightarrow{f_1^{e - 1}}
M/(f_1^e, f_2)M \to
M/(f_1^{e - 1}, f_2)M \to 0
$$
は補題 \ref{algebra-lemma-snake} により短完全列である。
これで $r = 2$ の場合の逆が示された。
$r > 2$ ならば、
$f_3 : M/(f_1^e, f_2)M \to M/(f_1^e, f_2)M$ は単射であり、
したがって
$f_3 : M/(f_1, f_2)M \to M/(f_1, f_2)M$ は単射である。
以下も同様である。細部は省略する。
\end{proof}

\begin{lemma}
\label{algebra-lemma-regular-sequence-in-polynomial-ring}
$R$ を環とし、$f_1, \ldots, f_r \in R$ は単位イデアルを
生成しないものとする。次の条件は同値である。
\begin{enumerate}
\item $f_1, \ldots, f_r$ の任意の置換は正則列である。
\item $f_1, \ldots, f_r$ の任意の部分列は
（与えられた順序で）正則列である。
\item $f_1x_1, \ldots, f_rx_r$ は多項式環
$R[x_1, \ldots, x_r]$ 内の正則列である。
\end{enumerate}
\end{lemma}

\begin{proof}
(1) が (2) を含意することは明らかである。
(2) が (1) を含意することを $r$ に関する帰納法で示す。
$r = 1$ の場合は自明である。
$r = 2$ の場合に示すべきことは、
$a, b \in R$ が正則列で $b$ が非零因子ならば、
$b, a$ が正則列であるということである。
$a$ と $b$ がともに非零因子ならば、
$a : R/(b) \to R/(b)$ の核は
$b : R/(a) \to R/(a)$ の核と同型なので、これは明らかである。
$r > 2$ の場合を考える。(2) が成り立つと仮定し、ある置換
$\sigma$ に対して $f_{\sigma(1)}, \ldots, f_{\sigma(r)}$
が正則列であることを示す。
帰納法により
$f_{\sigma(1)}, \ldots, f_{\sigma(r - 1)}$ が正則列であることは
すでに分かっている。
したがって $s = \sigma(r)$ とするとき、$f_s$ が
$f_1, \ldots, \hat f_s, \ldots, f_r$ を法として
非零因子であることを示せば十分である。
$s = r$ ならば証明は終わる。
$s < r$ ならば、帰納法の仮定を再び用いることにより、
$f_s$ と $f_r$ はともに環
$R/(f_1, \ldots, \hat f_s, \ldots, f_{r - 1})$
における非零因子である。
この環において $f_s, f_r$ は正則列だと分かっているので、
長さ $2$ の場合により $f_r, f_s$ も正則列である。

\medskip\noindent
$R$-加群として
$R[x_1, \ldots, x_r]/(f_1x_1, \ldots, f_ix_i)$
は、複数添字 $E = (e_1, \ldots, e_r)$ で添字づけられた加群
$$
R/I_E \cdot x_1^{e_1} \ldots x_r^{e_r}
$$
の直和であることに注意する。ここで $I_E$ は、
$1 \leq j \leq i$ かつ $e_j > 0$ である添字に対する
$f_j$ によって生成されるイデアルである。
したがって
% P01-E0416: the frozen next polynomial factor uses x_i rather than x_{i+1}.
$f_{i + 1}x_i$ がこれの上の非零因子であることと、
すべての $E$ に対して $f_{i + 1}$ が $R/I_E$ 上の
非零因子であることは同値である。
すべての成分が正である $E$ を取れば、
$f_{i + 1}$ が $R/(f_1, \ldots, f_i)$ 上の非零因子だと分かる。
したがって (3) は (2) を含意する。
逆に (2) が成り立つならば、
$f_1, \ldots, f_i, f_{i + 1}$ の任意の部分列は正則列であり、
特に $f_{i + 1}$ はすべての $R/I_E$ 上の非零因子である。
このようにして (2) が (3) を含意することが分かる。
\end{proof}

\section{準正則列}
\label{algebra-section-quasi-regular}

\noindent
正則列よりわずかに弱く、扱いやすい準正則列の概念を導入する。
$R$ を環とし、$f_1, \ldots, f_c \in R$ とする。
$J = (f_1, \ldots, f_c)$ とおき、$M$ を $R$-加群とする。
このとき標準写像
\begin{equation}
\label{algebra-equation-quasi-regular}
M/JM \otimes_{R/J} R/J[X_1, \ldots, X_c]
\longrightarrow
\bigoplus\nolimits_{n \geq 0} J^nM/J^{n + 1}M
\end{equation}
があり、これは次数付き $R/J[X_1, \ldots, X_c]$-加群の写像で、規則
$$
\overline{m} \otimes X_1^{e_1} \ldots X_c^{e_c} \longmapsto
f_1^{e_1} \ldots f_c^{e_c} m \bmod J^{e_1 + \ldots + e_c + 1}M.
$$
によって定義される。
(\ref{algebra-equation-quasi-regular}) は常に全射であることに注意する。

\begin{definition}
\label{algebra-definition-quasi-regular-sequence}
$R$ を環、$M$ を $R$-加群とする。
$R$ の元の列 $f_1, \ldots, f_c$ が
{\it $M$-準正則} であるとは、
(\ref{algebra-equation-quasi-regular}) が同型であることをいう。
$M = R$ のときは、$f_1, \ldots, f_c$ を単に
{\it 準正則列} という。
\end{definition}

\noindent
したがって $f_1, \ldots, f_c$ が準正則列ならば
$$
R/J[X_1, \ldots, X_c] = \bigoplus\nolimits_{n \geq 0} J^n/J^{n + 1}
$$
である。ここで $J = (f_1, \ldots, f_c)$ である。
準正則列であることが $f_1, \ldots, f_c$ の順序に依存しないのは
明らかである。

\begin{lemma}
\label{algebra-lemma-regular-quasi-regular}
$R$ を環とする。
\begin{enumerate}
\item $R$ の正則列 $f_1, \ldots, f_c$ は準正則列である。
\item $M$ を $R$-加群とし、$f_1, \ldots, f_c$ を
$M$-正則列とする。このとき $f_1, \ldots, f_c$ は
$M$-準正則列である。
\end{enumerate}
\end{lemma}

\begin{proof}
$J = (f_1, \ldots, f_c)$ とおく。
第一の主張を $c$ に関する帰納法で示す。
$a_I \in R$ を満たす任意の関係
$\sum_{|I| = n} a_I f^I \in J^{n + 1}$ が与えられたとき、
すべての複数添字 $I$ に対して実際に $a_I \in J$ であることを
示さなければならない。
$J^{n + 1}$ の任意の元は $b_I \in J$ を用いて
$\sum_{|I| = n} b_I f^I$ の形に書けるので、
$a_I$ を $a_I - b_I$ で置き換えることにより、関係を
$\sum_{|I| = n} a_I f^I = 0$ と仮定してよい。
これは次のように書き直せる。
$$
\sum\nolimits_{e = 0}^n
\left(
\sum\nolimits_{|I'| = n - e}
a_{I', e} f^{I'}
\right)
f_c^e
=
0
$$
以下、「プライム付き」の複数添字 $I'$ は
$I' = (i_1, \ldots, i_{c - 1}, 0)$ の形であるものとする。
$l \in \{0, \ldots, n\}$ に関する帰納法により、関係
$$
\sum\nolimits_{e = 0}^l
\left(
\sum\nolimits_{|I'| = n - e}
a_{I', e} f^{I'}
\right)
f_c^e
=
0
$$
があれば、すべての $I', e$ に対して $a_{I', e} \in J$
となることを示す。
% P01-E0418: the frozen induction does not separate l=0 before l-1 and l-2 appear.
実際、$J' = (f_1, \ldots, f_{c-1})$ とおく。
$\sum\nolimits_{|I'| = n - l} a_{I', l} f^{I'}$ は
$f_c^{l}$ により $(J')^{n - l + 1}$ に写されることに注意する。
帰納法の仮定（$c$ に関する帰納法）により
$f_c^l a_{I', l} \in J'$ である。
$f_1, \ldots, f_c$ は正則列なので、$f_c$ は
$R/J'$ 上の零因子でない。したがって
$a_{I', l} \in J'$ と結論できる。
これにより、項
$(\sum\nolimits_{|I'| = n - l} a_{I', l} f^{I'})f_c^l$
を
$(\sum\nolimits_{|I'| = n - l + 1} f_c b_{I', l - 1}
f^{I'})f_c^{l-1}$
の形に書き直せる。したがって、次の形の新しい関係を得る。
$$
\left(\sum\nolimits_{|I'| = n - l + 1}
(a_{I', l-1} + f_c b_{I', l - 1}) f^{I'}\right)f_c^{l-1}
+
\sum\nolimits_{e = 0}^{l - 2}
\left(
\sum\nolimits_{|I'| = n - e}
a_{I', e} f^{I'}
\right)
f_c^e
=
0
$$
今度は $l$ に関する帰納法の仮定により、すべての
$a_{I', l-1} + f_c b_{I', l - 1} \in J$ および
$e \leq l - 2$ に対するすべての $a_{I', e} \in J$ を得る。
これと、上で示した $a_{I', l} \in J' \subset J$ を合わせれば、
帰納段階の証明が完了する。

\medskip\noindent
第二の主張は、$m_I \in M$ を満たす任意の形式的表式
$F = \sum_{|I| = n} m_I X^I$ に対し、
$\sum m_I f^I \in J^{n + 1}M$ ならば、すべての係数
% P01-E0417: the frozen source places module coefficients in the ring ideal J.
$m_I$ が $J$ に属するということである。
これは、上で第一の主張に対して対応する結果を示したのと
まったく同じ方法で証明される。
\end{proof}

\begin{lemma}
\label{algebra-lemma-flat-base-change-quasi-regular}
$R \to R'$ を平坦な環写像、$M$ を $R$-加群とする。
$f_1, \ldots, f_r \in R$ が $M$-準正則列をなすと仮定する。
このとき $R'$ における $f_1, \ldots, f_r$ の像は
$M \otimes_R R'$-準正則列をなす。
\end{lemma}

\begin{proof}
$J = (f_1, \ldots, f_r)$、$J' = JR'$、
$M' = M \otimes_R R'$ とおく。
標準写像
$\mu : R'/J'[X_1, \ldots X_r] \otimes_{R'/J'} M'/J'M' \to
\bigoplus (J')^nM'/(J')^{n + 1}M'$
が同型であることを示さなければならない。
$R \to R'$ は平坦なので、列
$0 \to J^nM \to M$ と
$0 \to J^{n + 1}M \to J^nM \to J^nM/J^{n + 1}M \to 0$
は $R'$ とテンソル積を取っても完全である。
前者からまず
$J^nM \otimes_R R' = (J')^nM'$ が従い、次いで後者から
$(J')^nM'/(J')^{n + 1}M' = J^nM/J^{n + 1}M \otimes_R R'$
が従う。
したがって $\mu$ は、仮定により同型である
(\ref{algebra-equation-quasi-regular}) と $\text{id}_{R'}$ との
テンソル積である。ゆえに結論を得る。
\end{proof}

\begin{lemma}
\label{algebra-lemma-quasi-regular-sequence-in-neighbourhood}
$R$ をネーター環、$M$ を有限 $R$-加群とする。
$\mathfrak p$ を素イデアルとする。
$x_1, \ldots, x_c$ を $R$ 内の列とし、その $R_{\mathfrak p}$
における像が $M_{\mathfrak p}$-準正則列をなすと仮定する。
このとき、ある $g \in R$、$g \not \in \mathfrak p$ が存在して、
$R_g$ における $x_1, \ldots, x_c$ の像は
$M_g$-準正則列をなす。
\end{lemma}

\begin{proof}
写像 (\ref{algebra-equation-quasi-regular}) の核 $K$ を考える。
$M/JM \otimes_{R/J} R/J[X_1, \ldots, X_c]$ は
有限 $R/J[X_1, \ldots, X_c]$-加群であり、
$R/J[X_1, \ldots, X_c]$ はネーター環なので、
$K$ も有限 $R/J[X_1, \ldots, X_c]$-加群である。
斉次生成元 $k_1, \ldots, k_t \in K$ を選ぶ。
仮定により、各 $i = 1, \ldots, t$ に対して
ある $g_i \in R$、$g_i \not \in \mathfrak p$ が存在して
$g_i k_i = 0$ となる。
したがって $g = g_1 \ldots g_t$ が求めるものである。
\end{proof}

\begin{lemma}
\label{algebra-lemma-truncate-quasi-regular}
$R$ を環、$M$ を $R$-加群とする。
$f_1, \ldots, f_c \in R$ を $M$-準正則列とする。
任意の $i$ に対し、$\overline{R} = R/(f_1, \ldots, f_i)$ の列
$\overline{f}_{i + 1}, \ldots, \overline{f}_c$ は
$\overline{M} = M/(f_1, \ldots, f_i)M$-準正則列である。
\end{lemma}

\begin{proof}
$i = 1$ の場合を示せば十分である。
$\overline{J} = (\overline{f}_2, \ldots, \overline{f}_c)
\subset \overline{R}$ とおく。このとき
\begin{align*}
\overline{J}^n\overline{M}/\overline{J}^{n + 1}\overline{M}
& =
(J^nM + f_1M)/(J^{n + 1}M + f_1M) \\
& = J^nM / (J^{n + 1}M + J^nM \cap f_1M).
\end{align*}
したがって補題を示すには
% P01-E0419: the frozen formula uses J^{n-1} without separating degree zero.
$J^{n + 1}M + J^nM \cap f_1M = J^{n + 1}M + f_1J^{n - 1}M$
を示せば十分である。実際、これにより
$\bigoplus_{n \geq 0}
\overline{J}^n\overline{M}/\overline{J}^{n + 1}\overline{M}$
は
% P01-E0420: the frozen source calls the coefficient-module expression an algebra.
$\bigoplus_{n \geq 0} J^nM/J^{n + 1}M \cong M/JM[X_1, \ldots, X_c]$
を $X_1$ で割った商であると分かる。
実際 $J^nM \cap f_1M = f_1J^{n - 1}M$ である。
なぜなら $m \not \in J^{n - 1}M$ ならば、
仮定により $\bigoplus J^nM/J^{n + 1}M$ は多項式代数
$M/J[X_1, \ldots, X_c]$ だから、
$f_1m \not \in J^nM$ となるからである。
\end{proof}

\begin{lemma}
\label{algebra-lemma-quasi-regular-regular}
$(R, \mathfrak m)$ をネーター局所環、
$M$ を零でない有限 $R$-加群とする。
$f_1, \ldots, f_c \in \mathfrak m$ を $M$-準正則列とする。
このとき $f_1, \ldots, f_c$ は $M$-正則列である。
\end{lemma}

\begin{proof}
$J = (f_1, \ldots, f_c)$ とおく。
$f_1$ が $M$ 上の非零因子であることを示そう。
$x \in M$ は零でないと仮定する。
クルルの交叉定理により、$x \in J^rM$ かつ
$x \not \in J^{r + 1}M$ となる整数 $r$ が存在する。
補題 \ref{algebra-lemma-intersect-powers-ideal-module-zero} を参照せよ。
このとき $f_1 x \in J^{r + 1}M$ の
$J^{r + 1}M/J^{r + 2}M$ における類は、
$\bigoplus J^nM/J^{n + 1}M$ の仮定された構造により零でない。
したがって $f_1x \not = 0$ である。

\medskip\noindent
ここで補題 \ref{algebra-lemma-truncate-quasi-regular} を用い、
$c$ に関する帰納法によって証明を終えられる。
\end{proof}

\begin{remark}[他の種類の正則列]
\label{algebra-remark-koszul-regular}
論文 \cite{Kabele} では、環 $R$ の元の列
$x_1, \ldots, x_r$ に対する、さらに二つの正則性条件が論じられている。
すなわち、$K_{\bullet}(R, x_{\bullet})$ をコズル複体とするとき、
$i \geq 1$ に対して $H_i(K_{\bullet}(R, x_{\bullet})) = 0$ ならば、
この列は {\it コズル正則} であるという。
$H_1(K_{\bullet}(R, x_{\bullet})) = 0$ ならば、
この列は {\it $H_1$-正則} であるという。
正則 $\Rightarrow$ コズル正則 $\Rightarrow$
$H_1$-正則 $\Rightarrow$ 準正則という含意がある。
著者は例によって、$R$ が（非ネーター）局所環で列が $R$ の
極大イデアルを生成する場合でさえ、これらの含意は一般には逆にできないことを示している。
これらの概念は More on Algebra, Section
\ref{more-algebra-section-koszul-regular} でより詳しく導入する。
\end{remark}

\begin{remark}
\label{algebra-remark-join-quasi-regular-sequences}
$k$ を体とし、環
$$
A = k[x, y, w, z_0, z_1, z_2, \ldots]/
(y^2z_0 - wx, z_0 - yz_1, z_1 - yz_2, \ldots)
$$
を考える。
この環において $x$ は非零因子であり、$A/xA$ における
$y$ の像は準正則列を与える。
しかし $x, y$ は $A$ 内の準正則列ではない。実際、
$(x, y)/(x, y)^2$ は $A/(x, y)$ 上階数二の自由加群ではない。
なぜなら $(x, y)/(x, y)^2$ において $wx = 0$ である一方、
$A/(x, y)$ において $w$ は零でないからである。
したがって、補題 \ref{algebra-lemma-join-regular-sequences} の類似は
準正則列については成り立たない。
\end{remark}

\begin{lemma}
\label{algebra-lemma-quasi-regular-on-quotient}
$R$ を環、$J = (f_1, \ldots, f_r)$ を $R$ のイデアル、
$M$ を $R$-加群とする。
$\overline{R} = R/\bigcap_{n \geq 0} J^n$、
$\overline{M} = M/\bigcap_{n \geq 0} J^nM$ とおき、
$\overline{f}_i$ を $\overline{R}$ における $f_i$ の像とする。
このとき $f_1, \ldots, f_r$ が $M$-準正則であることと、
$\overline{f}_1, \ldots, \overline{f}_r$ が
$\overline{M}$-準正則であることは同値である。
\end{lemma}

\begin{proof}
% P01-E0421: the frozen proof uses overline J without defining it.
これは
$J^nM/J^{n + 1}M \cong
\overline{J}^n\overline{M}/\overline{J}^{n + 1}\overline{M}$
だからである。
\end{proof}

\section{ブローアップ代数}
\label{algebra-section-blow-up}

\noindent
この節ではブローアップについていくつかの初等的な考察を行う。

\begin{definition}
\label{algebra-definition-blow-up}
$R$ を環、$I \subset R$ をイデアルとする。
\begin{enumerate}
\item 対 $(R, I)$ に付随する {\it ブローアップ代数}、または
{\it Rees 代数} とは、次数付き $R$-代数
$$
\text{Bl}_I(R) =
\bigoplus\nolimits_{n \geq 0} I^n =
R \oplus I \oplus I^2 \oplus \ldots
$$
のことである。ここで直和因子 $I^n$ は次数 $n$ に置かれる。
\item $a \in I$ を元とする。
Rees 代数の次数 $1$ の元とみなした $a$ を $a^{(1)}$ と表す。
このとき {\it アフィン・ブローアップ代数}
$R[\frac{I}{a}]$ とは、Section \ref{algebra-section-proj} で構成した代数
$(\text{Bl}_I(R))_{(a^{(1)})}$ のことである。
\end{enumerate}
\end{definition}

\noindent
言い換えると、$R[\frac{I}{a}]$ の元は
$x \in I^n$ を用いて $x/a^n$ の形の式で表される。
二つの代表元 $x/a^n$ と $y/a^m$ が同じ元を定めるための
必要十分条件は、ある $k \geq 0$ に対して
$a^k(a^mx - a^ny) = 0$ となることである。

\begin{lemma}
\label{algebra-lemma-affine-blowup}
$R$ を環、$I \subset R$ をイデアル、$a \in I$ とする。
$R' = R[\frac{I}{a}]$ をアフィン・ブローアップ代数とする。このとき
\begin{enumerate}
\item $R'$ における $a$ の像は非零因子である。
\item $IR' = aR'$ である。
\item $(R')_a = R_a$ である。
\end{enumerate}
\end{lemma}

\begin{proof}
上の $R[\frac{I}{a}]$ の記述から直ちに従う。
\end{proof}

\begin{lemma}
\label{algebra-lemma-blowup-base-change}
$R \to S$ を環写像とする。
$I \subset R$ をイデアル、$a \in I$ とする。
$J = IS$ とおき、$b \in J$ を $a$ の像とする。
このとき $S[\frac{J}{b}]$ は
$S \otimes_R R[\frac{I}{a}]$ を、$b$ のある冪により
零化される元のイデアルで割った商である。
\end{lemma}

\begin{proof}
$S'$ を $S \otimes_R R[\frac{I}{a}]$ の
$b$-冪捩れ元による商とする。環写像
$$
S \otimes_R R[\textstyle{\frac{I}{a}}]
\longrightarrow
S[\textstyle{\frac{J}{b}}]
$$
は全射であり、$S[\frac{J}{b}]$ において $b$ は非零因子なので、
% P01-E0424: the frozen proof switches from b-power to a-power torsion.
$a$-冪捩れを零化する。
したがって全射 $S' \to S[\frac{J}{b}]$ を得る。
核が自明であることを見るため、逆写像を構成する。
すなわち、$z = y/b^n$ を $S[\frac{J}{b}]$ の元とする。
つまり $y \in J^n$ である。
$x_i \in I^n$、$s_i \in S$ を用いて
$y = \sum x_is_i$ と書く。
$z$ を $S'$ における $\sum s_i \otimes x_i/a^n$ の類に写す。
これは well-defined である。実際、写像
$S \otimes_R I^n \to J^n$ の核の元は
% P01-E0424: the same frozen a/b image switch is retained at a^n.
$a^n$ により零化されるので、$S'$ において零に写る。
\end{proof}

\begin{example}
\label{algebra-example-rees-algebra-polynomial}
$R$ を環とし、$P = R[t_1, \ldots, t_n]$ を多項式代数とする。
$I = (t_1, \ldots, t_n) \subset P$ とする。
定義 \ref{algebra-definition-blow-up} の記法の下で、
$T_i$ を $t_i^{(1)}$ に写す同型
$$
P[T_1, \ldots, T_n]/(t_iT_j - t_jT_i) \longrightarrow \text{Bl}_I(P)
$$
がある。この写像が well-defined であることの確認は読者に任せる。
$I$ は $t_1, \ldots, t_n$ により生成されるので、この写像は全射である。
この写像が単射であることを見るには、各 $e \geq 1$ に対して
$P$-加群 $I^e$ が、次数 $|E| = e$ の複数添字
$E = (e_1, \ldots, e_n)$ に対する単項式
% P01-E0422: the frozen monomial changes the last variable from t_n to x_n.
$t^E = t_1^{e_1} \ldots x_n^{e_n}$
によって生成され、関係が
$|E| = |E'| = e$ かつ
$e_a + \delta_{a i} = e'_a + \delta_{a j},\ a = 1, \ldots, n$
のときの $t_i t^E = t_j t^{E'}$ のみであることを示す必要がある
（クロネッカーのデルタ）。細部は省略する。
\end{example}

\begin{example}
\label{algebra-example-affine-blowup-algebra-polynomial}
$R$ を環とし、$P = R[t_1, \ldots, t_n]$ を多項式代数とする。
$I = (t_1, \ldots, t_n) \subset P$ とし、$a = t_1$ とする。
定義 \ref{algebra-definition-blow-up} の記法の下で、
$x_i$ を $t_i/t_1$ に写す同型
$$
P[x_2, \ldots, x_n]/(t_1x_2 - t_2, \ldots, t_1x_n - t_n)
\longrightarrow
\textstyle{P[\frac{I}{a}] = P[\frac{I}{t_1}]}
$$
がある。この写像が well-defined であることの確認は読者に任せる。
$I$ は $t_1, \ldots, t_n$ により生成されるので、この写像は全射である。
この写像が単射であることを見るには、写像の始域と終域が
$t_1$-捩れをもたず、$t_1$ を可逆にした後で写像が同型になることを
用いればよい。補題 \ref{algebra-lemma-affine-blowup} を参照せよ。
別の方法として、例 \ref{algebra-example-rees-algebra-polynomial} の
Rees 代数の記述を用いてもよい。
細部は省略する。
\end{example}

\begin{lemma}
\label{algebra-lemma-affine-blowup-quotient-description}
$R$ を環、$I = (a_1, \ldots, a_n)$ を $R$ のイデアルとし、
$a = a_1$ とする。このとき全射
$$
R[x_2, \ldots, x_n]/(a x_2 - a_2, \ldots, a x_n - a_n)
\longrightarrow
\textstyle{R[\frac{I}{a}]}
$$
があり、その核は始域の $a$-冪捩れである。
\end{lemma}

\begin{proof}
$t_i$ を $a_i$ に写す環写像
$P = \mathbf{Z}[t_1, \ldots, t_n] \to R$ を考える。
$J = (t_1, \ldots, t_n)$ とおく。
例 \ref{algebra-example-affine-blowup-algebra-polynomial} により
$P[\frac{J}{t_1}] =
P[x_2, \ldots, x_n]/(t_1 x_2 - t_2, \ldots, t_1 x_n - t_n)$
である。
% P01-E0423: the frozen proof applies base change to undefined target A.
写像 $P \to A$ に補題 \ref{algebra-lemma-blowup-base-change} を適用すれば
結論を得る。
\end{proof}

\begin{lemma}
\label{algebra-lemma-blowup-in-principal}
$R$ を環、$I \subset R$ をイデアル、$a \in I$ とし、
$R' = R[\frac{I}{a}]$ とおく。
$f \in R$ が $V(f) = V(I)$ を満たすならば、
$f$ は $R'$ において非零因子に写り、
$R'_f = R'_a = R_a$ である。
\end{lemma}

\begin{proof}
以下、補題 \ref{algebra-lemma-affine-blowup} の結果を断りなく用いる。
仮定 $V(f) = V(I)$ は
$V(fR') = V(IR') = V(aR')$ を含意する。
したがって、ある $b, c \in R'$ に対して
$a^n = fb$ かつ $f^m = ac$ である。これから補題が従う。
\end{proof}

\begin{lemma}
\label{algebra-lemma-blowup-add-principal}
$R$ を環、$I \subset R$ をイデアル、
$a \in I$、$f \in R$ とする。
$R' = R[\frac{I}{a}]$ および
$R'' = R[\frac{fI}{fa}]$ とおく。
このとき全射な $R$-代数写像 $R' \to R''$ があり、
その核は $R'$ の $f$-冪捩れ元全体の集合である。
\end{lemma}

\begin{proof}
この写像は、$x \in I^n$ に対する $x/a^n$ を
$f^nx/(fa)^n$ に写すことで与えられる。
この写像が well-defined かつ全射であることは直ちに確認できる。
$R''$ において $af$ は非零因子なので
（補題 \ref{algebra-lemma-affine-blowup}）、
$f$-冪捩れ元全体の集合は零に写る。
逆に $x \in R'$ が、すべての $n > 0$ に対して
$f^n x \not = 0$ を満たすと仮定する。
$R'$ において $a$ は非零因子なので、すべての $n$ に対して
$(af)^n x \not = 0$ である。
定義 \ref{algebra-definition-blow-up} の後にある $R''$ の記述から、
$R''$ における $x$ の像は零でないと分かる。
\end{proof}

\begin{lemma}
\label{algebra-lemma-blowup-reduced}
\begin{slogan}
被約性はブローアップで不変である
\end{slogan}
$R$ が被約ならば、$R$ のすべての（アフィン）ブローアップ代数は
被約である。
\end{lemma}

\begin{proof}
$I \subset R$ をイデアル、$a \in I$ とする。
$x \in I^n$ に対する $x/a^n$ が
$R[\frac{I}{a}]$ の冪零元であると仮定する。
このとき $(x/a^n)^m = 0$ である。
したがって、ある $N \geq 0$ に対して $R$ において
$a^N x^m = 0$ である。
必要なら $N$ を増やすことにより、ある $e \geq 0$ に対して
$N = me$ と仮定してよい。
すると $(a^e x)^m = 0$ であり、$R$ は被約なので
$a^e x = 0$ となる。
これは $R[\frac{I}{a}]$ において $x/a^n = 0$ であることを意味する。
\end{proof}

\begin{lemma}
\label{algebra-lemma-blowup-domain}
$R$ を整域、$I \subset R$ をイデアルとし、
$a \in I$ を零でない元とする。
このときアフィン・ブローアップ代数 $R[\frac{I}{a}]$ は整域である。
\end{lemma}

\begin{proof}
$x \in I^n$、$y \in I^m$ に対する
$x/a^n$、$y/a^m$ が $R[\frac{I}{a}]$ の元で、その積が零だと仮定する。
このとき $R$ において $a^N x y = 0$ である。
$R$ は整域なので、$x = 0$ または $y = 0$ と結論する。
\end{proof}

\begin{lemma}
\label{algebra-lemma-blowup-dominant}
$R$ を環、$I \subset R$ をイデアル、$a \in I$ とする。
$a$ が $R$ のどの極小素イデアルにも含まれないならば、
$\Spec(R[\frac{I}{a}]) \to \Spec(R)$ の像は稠密である。
\end{lemma}

\begin{proof}
$x \in R$ に対して $a^k x = 0$ ならば、
$x$ は $R$ のすべての極小素イデアルに含まれ、したがって冪零である。
補題 \ref{algebra-lemma-Zariski-topology} を参照せよ。
したがって $R \to R[\frac{I}{a}]$ の核は冪零元からなる。
ゆえに補題 \ref{algebra-lemma-image-dense-generic-points} から結果が従う。
\end{proof}

\begin{lemma}
\label{algebra-lemma-valuation-ring-colimit-affine-blowups}
$(R, \mathfrak m)$ を分数体 $K$ をもつ局所整域とする。
$R \subset A \subset K$ を $R$ を支配する付値環とする。このとき
$$
A = \colim R[\textstyle{\frac{I}{a}}]
$$
は、次の性質をもつアフィン・ブローアップ
$R \to R[\frac{I}{a}]$ の有向余極限である。
\begin{enumerate}
\item $a \in I \subset \mathfrak m$。
\item $I$ は有限生成である。
\item $\mathfrak m$ における
$R \to R[\frac{I}{a}]$ のファイバー環は零でない。
\end{enumerate}
\end{lemma}

\begin{proof}
補題 \ref{algebra-lemma-blowup-domain} を参照すれば、任意のブローアップ代数
$R[\frac{I}{a}]$ は $K$ に含まれる整域である。
% P01-E0425: the unrestricted frozen statement includes the field case.
補題が述べているのは単に、$A$ が、(1)、(2)、(3) の性質をもつ
$a \in I$ に対応するものの有向和だということである。
$R[\frac{I}{a}] \subset A$ かつ
$R[\frac{J}{b}] \subset A$ ならば
$$
R[\textstyle{\frac{I}{a}}] \cup R[\textstyle{\frac{J}{b}}]  \subset
R[\textstyle{\frac{IJ}{ab}}] \subset A
$$
である。
第一の包含は $x/a^n = b^nx/(ab)^n$ により従う。
第二の包含については、$z \in (IJ)^n$ ならば
$x_i \in I^n$、$y_i \in J^n$ を用いて
$z = \sum x_iy_i$ と書けるので、
$z/(ab)^n = \sum (x_i/a^n)(y_i/b^n)$ は $A$ に含まれる。

\medskip\noindent
有限部分集合 $E \subset A$ を考え、
$E = \{e_1, \ldots, e_n\}$ とする。
すべての $i = 1, \ldots, n$ に対して $e_i = f_i/a$ と書けるような
零でない $a \in R$ を選ぶ。
$I = (f_1, \ldots, f_n, a)$ とおく。
% P01-E0426: the frozen proof does not establish I subset m or the nonzero fibre.
$R[\frac{I}{a}] \subset A$ と主張する。
$R[\frac{I}{a}]$ の元は $e_i$ の多項式として表せるので、これは明らかである。
この考察から補題は直ちに従う。
\end{proof}

\section{Ext 群}
\label{algebra-section-ext}

\noindent
この節では、ネーター局所環上の深さの基本的な性質を確立するため、
ホモロジー代数を少しだけ扱う。

\begin{lemma}
\label{algebra-lemma-resolution-by-finite-free}
$R$ を環、$M$ を $R$-加群とする。
\begin{enumerate}
\item 完全複体
$$
\ldots \to F_2 \to F_1 \to F_0 \to M \to 0.
$$
で、各 $F_i$ が自由 $R$-加群であるものが存在する。
\item $R$ がネーター環で $M$ が有限 $R$-加群ならば、
各 $F_i$ が有限自由となるようにこの複体を選べる。
言い換えると、完全複体
$$
\ldots \to R^{\oplus n_2} \to R^{\oplus n_1} \to R^{\oplus n_0} \to M \to 0.
$$
を取れる。
\end{enumerate}
\end{lemma}

\begin{proof}
ネーターの場合だけ説明する。
最初に全射 $R^{n_0} \to M$ を選ぶ。
長さ $e$ の完全複体を構成したならば、単に全射
% P01-E0430: the frozen recursion does not define the first e=0 kernel.
$R^{n_{e + 1}} \to
\Ker(R^{n_e} \to R^{n_{e-1}})$
を選ぶ。$R$ はネーター環なので、これは可能である。
\end{proof}

\begin{definition}
\label{algebra-definition-finite-free-resolution}
$R$ を環、$M$ を $R$-加群とする。
\begin{enumerate}
\item $M$ の（左）{\it 分解} $F_\bullet \to M$ とは、
$R$-加群の完全複体
$$
\ldots \to F_2 \to F_1 \to F_0 \to M \to 0
$$
のことである。
\item {\it 自由 $R$-加群による $M$ の分解} とは、
各 $F_i$ が自由 $R$-加群である分解 $F_\bullet \to M$ のことである。
\item {\it 有限自由 $R$-加群による $M$ の分解} とは、
各 $F_i$ が有限自由 $R$-加群である分解
$F_\bullet \to M$ のことである。
\end{enumerate}
\end{definition}

\noindent
$R$-加群の複体
$$
\ldots \to F_i \to F_{i-1} \to \ldots
$$
を表すため、しばしば記法 $F_{\bullet}$ を用いる。
この場合、写像 $F_i \to F_{i-1}$ を $d_i$ または
$d_{F, i}$ と表すことが多い。
この節では常に、$F_0$ が複体の最後の非零項であると仮定する。
複体 $F_{\bullet}$ の {\it 第 $i$ ホモロジー群} とは
群 $H_i = \Ker(d_{F, i})/\Im(d_{F, i + 1})$ のことである。
{\it 複体の写像 $\alpha : F_{\bullet} \to G_{\bullet}$} は、
% P01-E0427: the frozen chain-map identity retains the ill-typed d_{G,i-1}.
$\alpha_{i-1} \circ d_{F, i} = d_{G, i-1} \circ \alpha_i$
を満たす写像 $\alpha_i : F_i \to G_i$ によって与えられる。
このような写像はホモロジー上の写像
$H_i(\alpha) :
H_i(F_{\bullet}) \to H_i(G_{\bullet})$
を誘導する。
$\alpha, \beta
:  F_{\bullet} \to G_{\bullet}$ が複体の写像ならば、
$\alpha$ と $\beta$ の間の {\it ホモトピー} とは、
$\alpha_i - \beta_i = d_{G, i + 1} \circ h_i +
h_{i-1} \circ d_{F, i}$
を満たす写像の族 $h_i : F_i \to G_{i + 1}$ のことである。
$\alpha$ と $\beta$ の間にホモトピーが存在するとき、
二つの写像 $\alpha, \beta : F_{\bullet} \to G_{\bullet}$ は
{\it ホモトピック} であるという。

\medskip\noindent
次の形をした複体 $F^{\bullet}$ についても、非常によく似た記法を用いる。
$$
\ldots \to F^i \xrightarrow{d^i} F^{i + 1} \to \ldots
$$
複体の写像、ホモトピーなどが同様に定義される。
この場合
$H^i(F^{\bullet}) =
\Ker(d^i)/\Im(d^{i - 1})$
とおき、これを {\it 第 $i$ コホモロジー群} という。

\begin{lemma}
\label{algebra-lemma-homotopic-equal-homology}
ホモトピックな任意の二つの複体の写像は、（コ）ホモロジー群上に
同じ写像を誘導する。
\end{lemma}

\begin{proof}
省略する。
\end{proof}

\begin{lemma}
\label{algebra-lemma-compare-resolutions}
$R$ を環とし、$M \to N$ を $R$-加群の写像とする。
$N_\bullet \to N$ を任意の分解とする。
各 $F_i$ が自由 $R$-加群である $R$-加群の複体
$$
\ldots \to F_2 \to F_1 \to F_0 \to M
$$
を取る。このとき
\begin{enumerate}
\item 図式
$$
\xymatrix{
F_0 \ar[r] \ar[d] & M \ar[d] \\
N_0 \ar[r] & N
}
$$
を可換にする複体の写像 $F_\bullet \to N_\bullet$ が存在する。
\item (1) のような任意の二つの写像
$\alpha, \beta : F_\bullet \to N_\bullet$ はホモトピックである。
\end{enumerate}
\end{lemma}

\begin{proof}
(1) の証明。$F_0$ は自由なので、写像 $F_0 \to M \to N$ を
持ち上げる写像 $F_0 \to N_0$ を取れる。
これにより写像 $F_1 \to F_0 \to N_0$ が誘導され、その像は
$N_1 \to N_0$ の像に入る。
$F_1$ は自由なので、これを写像 $F_1 \to N_1$ に持ち上げられる。
さらにこれは写像 $F_2 \to F_1 \to N_1$ を誘導し、
その $N_0$ における像は零である。
$N_\bullet$ は完全なので、この写像の像は
$N_2 \to N_1$ の像に含まれる。
したがって持ち上げて写像 $F_2 \to N_2$ を得られる。
これを繰り返す。

\medskip\noindent
(2) の証明。$\alpha, \beta$ がホモトピックであることを示すには、
差 $\gamma = \alpha - \beta$ が零写像とホモトピックであることを
示せば十分である。
$\gamma_0 : F_0 \to N_0$ の像は $N_1 \to N_0$ の像に含まれる。
したがって $\gamma_0$ を写像 $h_0 : F_0 \to N_1$ に持ち上げられる。
写像 $\gamma_1' = \gamma_1 - h_0 \circ d_{F, 1}$ を考える。
$h_0$ の選び方により、$\gamma_1'$ の像は
$N_1 \to N_0$ の核に含まれる。
$N_\bullet$ は完全なので、$\gamma_1'$ を
写像 $h_1 : F_1 \to N_2$ に持ち上げられる。
この段階で
$\gamma_1 = h_0 \circ d_{F, 1}
+ d_{N, 2} \circ h_1$
を得る。これを繰り返す。
\end{proof}

\noindent
これで群 $\Ext^i_R(M, N)$ を定義する準備ができた。
すなわち、補題 \ref{algebra-lemma-resolution-by-finite-free} を参照して、
自由 $R$-加群による $M$ の分解 $F_{\bullet}$ を選ぶ。
（コホモロジー的）複体
$$
\Hom_R(F_\bullet, N) :
\Hom_R(F_0, N) \to
\Hom_R(F_1, N) \to
\Hom_R(F_2, N) \to \ldots
$$
を考える。
$i \geq 0$ に対し、$\Ext^i_R(M, N)$ をこの複体の第 $i$
コホモロジー群として定義する\footnote{この段階では、Ext 群「そのもの」
ではなく Ext 群「の一つ」と言う方が適切かもしれない。}。
$i < 0$ に対しては $\Ext^i_R(M, N) = 0$ とおく。
先へ進む前に、
$$
\Ext^0_R(M, N) = \Ker(\Hom_R(F_0, N) \to \Hom_R(F_1, N)) =
\Hom_R(M, N)
$$
であることを指摘しておく。
これは完全列 $F_1 \to F_0 \to M \to 0$ に
補題 \ref{algebra-lemma-hom-exact} の (1) を適用できるからである。
次の補題は、この定義がどの意味で well-defined であるかを説明する。

\begin{lemma}
\label{algebra-lemma-ext-welldefined}
$R$ を環とし、$M_1, M_2, N$ を $R$-加群とする。
$F_{\bullet}$ を加群 $M_1$ の自由分解、
$G_{\bullet}$ を加群 $M_2$ の自由分解とする。
$\varphi : M_1 \to M_2$ を加群写像とする。
$\alpha : F_{\bullet} \to G_{\bullet}$ を、
$M_1 = \Coker(d_{F, 1}) \to M_2 = \Coker(d_{G, 1})$
上に $\varphi$ を誘導する複体の写像とする。
補題 \ref{algebra-lemma-compare-resolutions} を参照せよ。
このとき誘導される写像
% P01-E0428: the frozen cohomology arrow treats Hom(-,N) covariantly.
$$
H^i(\alpha) :
H^i(\Hom_R(F_{\bullet}, N))
\longrightarrow
H^i(\Hom_R(G_{\bullet}, N))
$$
は $\alpha$ の選び方に依存しない。
$\varphi$ が同型ならば、すべての写像
$H^i(\alpha)$ も同型である。
$M_1 = M_2$、$F_\bullet = G_\bullet$ で、
$\varphi$ が恒等写像ならば、
% P01-E0429: the frozen statement switches from cohomology H^i to H_i.
すべての写像 $H_i(\alpha)$ も恒等写像である。
\end{lemma}

\begin{proof}
$\varphi$ を誘導する別の写像
$\beta : F_{\bullet} \to G_{\bullet}$ は、
補題 \ref{algebra-lemma-compare-resolutions} により $\alpha$ とホモトピックである。
したがって写像
% P01-E0428: the frozen proof retains the same contravariant direction defect.
$\Hom_R(F_\bullet, N) \to
\Hom_R(G_\bullet, N)$
はホモトピックである。
ゆえに補題 \ref{algebra-lemma-homotopic-equal-homology} から
独立性の主張が従う。

\medskip\noindent
$\varphi$ が同型であると仮定し、
$\psi : M_2 \to M_1$ をその逆写像とする。
$\psi :
M_2 = \Coker(d_{G, 1}) \to M_1 = \Coker(d_{F, 1})$
を誘導する写像
$\beta : G_{\bullet} \to F_{\bullet}$ を選ぶ。
補題 \ref{algebra-lemma-compare-resolutions} を参照せよ。
さて、写像
$H^i(\alpha) \circ H^i(\beta) =
H^i(\alpha \circ \beta)$
を考える。
上記により、写像 $H^i(\alpha \circ \beta)$ は
$H^i(\text{id}_{G_{\bullet}}) = \text{id}$ と
{\it 同じ} である。
合成 $H^i(\beta) \circ H^i(\alpha)$ についても同様である。
したがって $H^i(\alpha)$ と $H^i(\beta)$ は互いに逆である。
\end{proof}

\begin{lemma}
\label{algebra-lemma-long-exact-seq-ext}
$R$ を環、$M$ を $R$-加群とする。
$0 \to N' \to N \to N'' \to 0$ を短完全列とする。
このとき長完全列
$$
\begin{matrix}
0
\to \Hom_R(M, N')
\to \Hom_R(M, N)
\to \Hom_R(M, N'')
\\
\phantom{0\ }
\to \Ext^1_R(M, N')
\to \Ext^1_R(M, N)
\to \Ext^1_R(M, N'')
\to \ldots
\end{matrix}
$$
を得る。
\end{lemma}

\begin{proof}
自由分解 $F_{\bullet} \to M$ を選ぶ。
各 $F_i$ は自由なので、複体の短完全列
$$
0 \to
\Hom_R(F_{\bullet}, N') \to
\Hom_R(F_{\bullet}, N) \to
\Hom_R(F_{\bullet}, N'') \to
0
$$
を得る。
これに蛇の補題を適用して長完全列を得る。
\end{proof}

\begin{lemma}
\label{algebra-lemma-reverse-long-exact-seq-ext}
$R$ を環、$N$ を $R$-加群とする。
$0 \to M' \to M \to M'' \to 0$ を短完全列とする。
このとき長完全列
$$
\begin{matrix}
0
\to \Hom_R(M'', N)
\to \Hom_R(M, N)
\to \Hom_R(M', N)
\\
\phantom{0\ }
\to \Ext^1_R(M'', N)
\to \Ext^1_R(M, N)
\to \Ext^1_R(M', N)
\to \ldots
\end{matrix}
$$
を得る。
\end{lemma}

\begin{proof}
$M'$ と $M''$ の生成元の集合
$\{m'_{i'}\}_{i' \in I'}$ および
$\{m''_{i''}\}_{i'' \in I''}$ を選ぶ。
各 $i'' \in I''$ に対し、元 $m''_{i''} \in M''$ の
持ち上げ $\tilde m''_{i''} \in M$ を選ぶ。
$F' = \bigoplus_{i' \in I'} R$、
$F'' = \bigoplus_{i'' \in I''} R$、
$F = F' \oplus F''$ とおく。
これらの自由加群の生成元を、対応する選んだ生成元に写すことにより、
全射な $R$-加群写像 $F' \to M'$、$F'' \to M''$、
$F \to M$ を得る。
短完全列の写像
$$
\begin{matrix}
0 & \to & M' & \to & M & \to & M'' & \to & 0 \\
& & \uparrow & & \uparrow & & \uparrow \\
0 & \to & F' & \to & F & \to & F'' & \to & 0 \\
\end{matrix}
$$
を得る。
蛇の補題により、核の列
$0 \to K' \to K \to K'' \to 0$ は
$R$-加群の短完全列である。
したがって、この手続きを無限に続けられる。
言い換えると、次の図式に入る分解の短完全列を得る。
$$
\begin{matrix}
0 & \to & M' & \to & M & \to & M'' & \to & 0 \\
& & \uparrow & & \uparrow & & \uparrow \\
0 & \to & F_\bullet' & \to & F_\bullet & \to & F_\bullet'' & \to & 0 \\
\end{matrix}
$$
各列 $0 \to F'_n \to F_n \to F''_n \to 0$ は
構成により分裂完全なので、関手 $\Hom_R(-, N)$ を適用すると、
複体の短完全列
$$
0 \to
\Hom_R(F''_{\bullet}, N) \to
\Hom_R(F_{\bullet}, N) \to
\Hom_R(F'_{\bullet}, N) \to
0
$$
を得る。
これに蛇の補題を適用して長完全列を得る。
\end{proof}

\begin{lemma}
\label{algebra-lemma-annihilate-ext}
$R$ を環とし、$M$, $N$ を $R$-加群とする。
$xN = 0$ または $xM = 0$ を満たす任意の $x\in R$ は、
各加群 $\Ext^i_R(M, N)$ を零化する。
\end{lemma}

\begin{proof}
$M$ の自由分解 $F_{\bullet}$ を選ぶ。
$\Ext^i_R(M, N)$ は複体 $\Hom_R(F_{\bullet}, N)$ の
コホモロジーとして定義されるので、$xN = 0$ の場合は明らかである。
$xM = 0$ ならば、$F_{\bullet}$ 上の $x$ による乗法は
$M$ 上の零写像を持ち上げる。
したがって補題 \ref{algebra-lemma-ext-welldefined} により、
Ext 群上に零写像と同じ写像を誘導する。
\end{proof}

\begin{lemma}
\label{algebra-lemma-ext-noetherian}
$R$ をネーター環とし、$M$, $N$ を有限 $R$-加群とする。
このとき、すべての $i$ に対して
$\Ext^i_R(M, N)$ は有限 $R$-加群である。
\end{lemma}

\begin{proof}
$\Ext^i_R(M, N)$ は、各 $F_n$ が有限自由 $R$-加群である
複体 $\Hom_R(F_\bullet, N)$ のコホモロジー群として計算されるからである。
補題 \ref{algebra-lemma-resolution-by-finite-free} を参照せよ。
\end{proof}

\section{深さ}
\label{algebra-section-depth}

\noindent
まず定義を与える。

\begin{definition}
\label{algebra-definition-depth}
$R$ を環、$I \subset R$ をイデアルとし、$M$ を有限 $R$-加群とする。
$M$ の {\it $I$-深さ} を $\text{depth}_I(M)$ と書き、次のように定義する。
\begin{enumerate}
\item $IM \not = M$ ならば、$\text{depth}_I(M)$ は、
$I$ に含まれる $M$-正則列の長さ全体の
$\{0, 1, 2, \ldots, \infty\}$ における上限である。
\item $IM = M$ ならば $\text{depth}_I(M) = \infty$ とおく。
\end{enumerate}
$(R, \mathfrak m)$ が局所環ならば、$\text{depth}_{\mathfrak m}(M)$ を単に
$M$ の {\it 深さ} という。
\end{definition}

\noindent
説明。定義 \ref{algebra-definition-regular-sequence} によれば、空列は零加群上の
正則列ではないが、実用上は $0$ 加群の深さを $+\infty$ とおくのが便利である。
$I = R$ ならば、すべての有限 $R$-加群 $M$ に対して
$\text{depth}_I(M) = \infty$ であることに注意する。
$I$ が $R$ の Jacobson 根基に含まれるならば
（例えば、$R$ が局所環で $I \subset \mathfrak m_R$ ならば）、
中山の補題により $M \not = 0 \Rightarrow IM \not = M$ である。
加群 $M$ の $I$-深さが $0$ であるための必要十分条件は、
$M$ が零でなく、$I$ が $M$ 上の非零因子を含まないことである。

\medskip\noindent
例 \ref{algebra-example-global-regular} は、環がネーターであっても深さが
よく振る舞わないことを示し、例 \ref{algebra-example-local-regular} は、
環が局所的でも非ネーターならばよく振る舞わないことを示す。
環がネーター局所環ならば、深さがよく振る舞うことを後で見る。

\begin{lemma}
\label{algebra-lemma-depth-weak-sequence}
$R$ を環、$I \subset R$ をイデアル、$M$ を有限 $R$-加群とする。
このとき $\text{depth}_I(M)$ は、各 $f_i$ が
$M/(f_1, \ldots, f_{i - 1})M$ 上の非零因子であるような列
$f_1, \ldots, f_r \in I$ の長さ全体の上限に等しい。
\end{lemma}

\begin{proof}
$IM = M$ と仮定する。このとき補題 \ref{algebra-lemma-NAK} により、
$f : M \to M$ が $\text{id}_M$ となる $f \in I$ が存在する。
したがって $f, 0, 0, 0, \ldots$ は逐次的な非零因子の無限列であり、
この場合に両者が一致することが分かる。
$IM \not =  M$ ならば、補題に現れる列は $M$-正則列であるから、
やはり両者は一致する。
\end{proof}

\begin{lemma}
\label{algebra-lemma-bound-depth}
$(R, \mathfrak m)$ をネーター局所環とし、
$M$ を零でない有限 $R$-加群とする。このとき
$\dim(\text{Supp}(M)) \geq \text{depth}(M)$ である。
\end{lemma}

\begin{proof}
$\dim(\text{Supp}(M))$ に関する帰納法で証明する。
$\dim(\text{Supp}(M)) = 0$ ならば
$\text{Supp}(M) = \{\mathfrak m\}$ であり、したがって
$\text{Ass}(M) = \{\mathfrak m\}$ である
（補題 \ref{algebra-lemma-ass-support} と \ref{algebra-lemma-ass-zero} による）。
ゆえに、例えば補題 \ref{algebra-lemma-ideal-nonzerodivisor} により、
$M$ の深さは零である。
帰納段階では $\dim(\text{Supp}(M)) > 0$ と仮定する。
$f_i$ が $M/(f_1, \ldots, f_{i - 1})M$ 上の非零因子となるような
$\mathfrak m$ の元の列 $f_1, \ldots, f_d$ をとる。
補題 \ref{algebra-lemma-depth-weak-sequence} により、
$\dim(\text{Supp}(M)) \geq d$ を示せば十分である。
列の長さが零ならば補題は成り立つので、$d > 0$ と仮定してよい。
補題 \ref{algebra-lemma-one-equation-module} により
$\dim(\text{Supp}(M/f_1M)) = \dim(\text{Supp}(M)) - 1$ である。
帰納法により $\dim(\text{Supp}(M/f_1M)) \geq d - 1$ を得て、結論が従う。
\end{proof}

\begin{lemma}
\label{algebra-lemma-depth-finite-noetherian}
$R$ をネーター環、$I \subset R$ をイデアル、$M$ を
$IM \not = M$ を満たす零でない有限 $R$-加群とする。このとき
$\text{depth}_I(M) < \infty$ である。
\end{lemma}

\begin{proof}
$M/IM$ は零でないので、補題 \ref{algebra-lemma-support-zero} により
$\mathfrak p \in \text{Supp}(M/IM)$ を選べる。
すると $(M/IM)_\mathfrak p \not = 0$ であり、これは
$I \subset \mathfrak p$ を含意し、さらに局所化が完全なので
$M_\mathfrak p \not = IM_\mathfrak p$ を含意する。
$f_1, \ldots, f_r \in I$ を $M$-正則列とする。
$(f_1, \ldots, f_r) \subset I$ なので、
$M_\mathfrak p/(f_1, \ldots, f_r)M_\mathfrak p$ は零でない。
局所化は平坦であるから、$f_1, \ldots, f_r$ の像は
$I_\mathfrak p$ 内の $M_\mathfrak p$-正則列をなす。
これは $I$ 内のすべての $M$-正則列に対して成り立つので
$\text{depth}_I(M) \leq \text{depth}_{I_\mathfrak p}(M_\mathfrak p)$
である。右辺は $\leq \text{depth}(M_\mathfrak p)$ であり、
補題 \ref{algebra-lemma-bound-depth} によりこれは $< \infty$ である。
\end{proof}

\begin{lemma}
\label{algebra-lemma-depth-ext}
$R$ を極大イデアル $\mathfrak m$ をもつネーター局所環とし、
$M$ を零でない有限 $R$-加群とする。このとき $\text{depth}(M)$ は、
$\Ext^i_R(R/\mathfrak m, M)$ が零でないような最小の整数 $i$ に等しい。
\end{lemma}

\begin{proof}
$M$ の深さを $\delta(M)$ と書き、$\Ext^i_R(R/\mathfrak m, M)$ が
零でないような最小の整数 $i$ を $i(M)$ と書く。
$i(M) < \infty$ であることはすぐ後で分かる。
補題 \ref{algebra-lemma-ideal-nonzerodivisor} により
$\delta(M) = 0$ であることと $i(M) = 0$ であることは同値である。
実際、$\mathfrak m \in \text{Ass}(M)$ であることは、ちょうど
$i(M) = 0$ であることを意味する。
したがって $\delta(M)$ または $i(M)$ が $> 0$ ならば、
(a) $x$ は $M$ 上の非零因子であり、かつ
(b) $\text{depth}(M/xM) = \delta(M) - 1$ となる
$x \in \mathfrak m$ を選べる。
補題 \ref{algebra-lemma-long-exact-seq-ext} により、短完全列
$0 \to M \to M \to M/xM \to 0$ に付随する Ext 群の長完全列を考える。
$$
\begin{matrix}
0
\to \Hom_R(\kappa, M)
\to \Hom_R(\kappa, M)
\to \Hom_R(\kappa, M/xM)
\\
\phantom{0\ }
\to \Ext^1_R(\kappa, M)
\to \Ext^1_R(\kappa, M)
\to \Ext^1_R(\kappa, M/xM)
\to \ldots
\end{matrix}
$$
$x \in \mathfrak m$ なので、すべての写像
$\Ext^i_R(\kappa, M) \to \Ext^i_R(\kappa, M)$ は零写像である。
補題 \ref{algebra-lemma-annihilate-ext} を参照せよ。
したがって $i(M/xM) = i(M) - 1$ であることは明らかである。
$\delta(M)$ に関する帰納法で証明が完了する。
\end{proof}

\begin{lemma}
\label{algebra-lemma-depth-in-ses}
$R$ をネーター局所環とする。零でない有限 $R$-加群の短完全列
$0 \to N' \to N \to N'' \to 0$ をとる。
\begin{enumerate}
\item
$\text{depth}(N) \geq \min\{\text{depth}(N'), \text{depth}(N'')\}$
\item
$\text{depth}(N'') \geq \min\{\text{depth}(N), \text{depth}(N') - 1\}$
\item
$\text{depth}(N') \geq \min\{\text{depth}(N), \text{depth}(N'') + 1\}$
\end{enumerate}
\end{lemma}

\begin{proof}
補題 \ref{algebra-lemma-depth-ext} の Ext 群 $\Ext^i(\kappa, N)$ による
深さの特徴づけと、補題 \ref{algebra-lemma-long-exact-seq-ext} の長完全コホモロジー列
$$
\begin{matrix}
0
\to \Hom_R(\kappa, N')
\to \Hom_R(\kappa, N)
\to \Hom_R(\kappa, N'')
\\
\phantom{0\ }
\to \Ext^1_R(\kappa, N')
\to \Ext^1_R(\kappa, N)
\to \Ext^1_R(\kappa, N'')
\to \ldots
\end{matrix}
$$
を用いる。
\end{proof}

\begin{lemma}
\label{algebra-lemma-depth-drops-by-one}
$R$ をネーター局所環、$M$ を零でない有限 $R$-加群とする。
\begin{enumerate}
\item $x \in \mathfrak m$ が $M$ 上の非零因子ならば
$\text{depth}(M/xM) = \text{depth}(M) - 1$ である。
\item 任意の $M$-正則列 $x_1, \ldots, x_r$ は、長さ
$\text{depth}(M)$ の $M$-正則列に延長できる。
\end{enumerate}
\end{lemma}

\begin{proof}
(2) は (1) の形式的な帰結である。$x \in R$ を (1) のようにとる。
短完全列 $0 \to M \to M \to M/xM \to 0$ と
補題 \ref{algebra-lemma-depth-in-ses} により、深さは高々 1 だけ下がる。
一方、$x_1, \ldots, x_r \in \mathfrak m$ が $M/xM$ の正則列ならば、
$x, x_1, \ldots, x_r$ は $M$ の正則列である。
したがって深さは少なくとも 1 だけ下がる。
\end{proof}

\begin{lemma}
\label{algebra-lemma-inherit-minimal-primes}
$(R, \mathfrak m)$ をネーター局所環、$M$ を有限 $R$-加群とする。
$x \in \mathfrak m$、$\mathfrak p \in \text{Ass}(M)$ とし、
$\mathfrak q$ は $\mathfrak p + (x)$ 上極小であるとする。
このとき、ある $n \geq 1$ に対して
$\mathfrak q \in \text{Ass}(M/x^nM)$ である。
\end{lemma}

\begin{proof}
$N \cong R/\mathfrak p$ となる部分加群 $N \subset M$ を選ぶ。
Artin--Rees の補題（補題 \ref{algebra-lemma-Artin-Rees}）により、
$N \cap x^nM \subset xN$ となる $n > 0$ を選べる。
$N \to M \to M/x^nM$ の像を $\overline{N} \subset M/x^nM$ とする。
補題 \ref{algebra-lemma-ass} により、$\mathfrak q \in \text{Ass}(\overline{N})$
を示せば十分である。
% P01-E0431: frozen source quotient-parentheses candidate; formula preserved.
$n$ の選び方により全射 $\overline{N} \to N/xN = R/\mathfrak p + (x)$
があり、したがって $\mathfrak q$ は $\overline{N}$ の台に属する。
$\overline{N}$ は $x^n$ と $\mathfrak p$ により零化されるので、
$\mathfrak q$ は $\overline{N}$ の台に属する素イデアルのうち極小である。
ゆえに補題 \ref{algebra-lemma-ass-minimal-prime-support} により、
$\mathfrak q$ は $\overline{N}$ の随伴素イデアルである。
\end{proof}

\begin{lemma}
\label{algebra-lemma-depth-dim-associated-primes}
$(R, \mathfrak m)$ をネーター局所環、$M$ を有限 $R$-加群とする。
$\mathfrak p \in \text{Ass}(M)$ に対して
$\dim(R/\mathfrak p) \geq \text{depth}(M)$ である。
\end{lemma}

\begin{proof}
$\mathfrak m \in \text{Ass}(M)$ ならば、$\mathfrak m$ のすべての元で
零化される零でない元 $x \in M$ が存在する。
したがって $\text{depth}(M) = 0$ である。
特に、この場合には補題が成り立つ。

\medskip\noindent
$\text{depth}(M) = 1$ ならば、最初の段落により
$\mathfrak m \not \in \text{Ass}(M)$ である。
したがって、すべての $\mathfrak p \in \text{Ass}(M)$ に対して
$\dim(R/\mathfrak p) \geq 1$ であり、この場合にも補題は成り立つ。

\medskip\noindent
一般の場合は $\text{depth}(M)$ に関する帰納法で証明する。
ここでその値は $> 1$ と仮定してよい。
$M$ 上の非零因子 $x \in \mathfrak m$ を選ぶ。
補題 \ref{algebra-lemma-ass-zero-divisors} により $x \not \in \mathfrak p$ である。
% P01-E0432: frozen source quotient-parentheses candidate; formula preserved.
補題 \ref{algebra-lemma-one-equation} により
$\dim(R/\mathfrak p + (x)) = \dim(R/\mathfrak p) - 1$ である。
したがって $\mathfrak p + (x)$ 上極小な素イデアル $\mathfrak q$ で
$\dim(R/\mathfrak q) = \dim(R/\mathfrak p) - 1$ となるものが存在する
（短い議論は省略する。ヒント：ネーター局所環 $A$ の次元は、
$A$ の極小素イデアル $\mathfrak r$ にわたる $A/\mathfrak r$ の
次元の最大値である）。
補題 \ref{algebra-lemma-inherit-minimal-primes} のように $n$ を選び、
$\mathfrak q$ が $M/x^nM$ の随伴素イデアルとなるようにする。
補題 \ref{algebra-lemma-depth-drops-by-one} により
$\text{depth}(M/x^nM) = \text{depth}(M) - 1$ なので、
$M/x^nM$ と $\mathfrak q$ に帰納法の仮定を適用できる。
$\dim(R/\mathfrak q) \geq \text{depth}(M/x^nM)$ を得て、結論が従う。
\end{proof}

\begin{lemma}
\label{algebra-lemma-depth-localization}
$R$ をネーター局所環、$M$ を有限 $R$-加群とする。
素イデアル $\mathfrak p \subset R$ に対して
$\text{depth}(M_\mathfrak p) + \dim(R/\mathfrak p) \geq \text{depth}(M)$
である。
\end{lemma}

\begin{proof}
$M_\mathfrak p = 0$ ならば $\text{depth}(M_\mathfrak p) = \infty$ なので、
補題は成り立つ。
$\text{depth}(M) \leq \dim(R/\mathfrak p)$ ならば、補題は明らかである。
$\text{depth}(M) > \dim(R/\mathfrak p)$ ならば、
補題 \ref{algebra-lemma-depth-dim-associated-primes} により、$\mathfrak p$ は
$M$ のいかなる随伴素イデアル $\mathfrak q$ にも含まれない。
したがって補題 \ref{algebra-lemma-silly} と補題 \ref{algebra-lemma-finite-ass} により、
$M$ のどの随伴素イデアルにも含まれない $x \in \mathfrak p$ を選べる。
補題 \ref{algebra-lemma-ass-zero-divisors} により $x$ は $M$ 上の非零因子である。
ゆえに $\text{depth}(M/xM) = \text{depth}(M) - 1$ であり、
$M_\mathfrak p$ が零でなければ
$\text{depth}(M_\mathfrak p / x M_\mathfrak p) =
\text{depth}(M_\mathfrak p) - 1$ である。
補題 \ref{algebra-lemma-depth-drops-by-one} を参照せよ。
したがって $\text{depth}(M)$ に関する帰納法で結論を得る。
\end{proof}

\begin{lemma}
\label{algebra-lemma-depth-goes-down-finite}
% P01-E0433: frozen proof omits its immediate zero-module branch; preserved.
$(R, \mathfrak m)$ をネーター局所環とし、$R \to S$ を有限環写像とする。
$\mathfrak m_1, \ldots, \mathfrak m_n$ を $S$ の極大イデアルとし、
$N$ を有限 $S$-加群とする。このとき
$$
\min\nolimits_{i = 1, \ldots, n} \text{depth}(N_{\mathfrak m_i}) =
\text{depth}_\mathfrak m(N)
$$
である。
\end{lemma}

\begin{proof}
補題 \ref{algebra-lemma-integral-no-inclusion}、\ref{algebra-lemma-integral-going-up}、
および補題 \ref{algebra-lemma-finite-finite-fibres} により、$S$ の極大イデアルは
ちょうど $\mathfrak m$ の上にある $S$ の素イデアルであり、有限個である。
したがって補題の主張には意味がある。
$k = \min\nolimits_{i = 1, \ldots, n} \text{depth}(N_{\mathfrak m_i})$
に関する帰納法で証明する。
$k = 0$ ならば、ある $i$ に対して
$\text{depth}(N_{\mathfrak m_i}) = 0$ である。
補題 \ref{algebra-lemma-depth-ext} により、これは
$\mathfrak m_i S_{\mathfrak m_i}$ が $N_{\mathfrak m_i}$ の
随伴素イデアルであることを意味し、したがって
$\mathfrak m_i$ は $N$ の随伴素イデアルである
（補題 \ref{algebra-lemma-localize-ass}）。
補題 \ref{algebra-lemma-ass-functorial-Noetherian} により、$R$-加群としての
$N$ の随伴素イデアルに $\mathfrak m$ が属する。
ゆえに $\text{depth}_\mathfrak m(N) = 0$ であり、基底の場合が示された。
$k > 0$ ならば $\mathfrak m_i \not \in \text{Ass}_S(N)$ である。
したがって再び補題 \ref{algebra-lemma-ass-functorial-Noetherian} により
$\mathfrak m \not \in \text{Ass}_R(N)$ である。
ゆえに補題 \ref{algebra-lemma-ideal-nonzerodivisor} により、$N$ 上の非零因子
$f \in \mathfrak m$ を選べる。補題 \ref{algebra-lemma-depth-drops-by-one} により、
$N$ から $N/fN$ に移るとすべての深さがちょうど $1$ ずつ下がるので、
残りは帰納法の仮定から従う。
\end{proof}

\section{Ext の関手性}
\label{algebra-section-functoriality-ext}

\noindent
この節では、環の変更などに関する $\Ext$ の関手性を簡単に論じる。
確認すべき事項を列挙する。
\begin{enumerate}
\item $R \to R'$、$R$-加群 $M$、および $R'$-加群 $N'$ が与えられたとき、
$R$-加群 $\Ext^i_R(M, N')$ は自然な $R'$-加群構造をもつ。
さらに標準的な $R'$-線形写像 $\Ext^i_{R'}(M \otimes_R R', N') \to
\Ext^i_R(M, N')$ がある。
\item $R \to R'$ と $R$-加群 $M$, $N$ が与えられたとき、自然な
$R$-加群写像
% P01-E0434: frozen source uses the text operator in two Ext expressions; preserved.
$\Ext^i_R(M, N) \to \text{Ext}^i_R(M, N \otimes_R R')$
がある。
\end{enumerate}

\begin{lemma}
\label{algebra-lemma-flat-base-change-ext}
平坦環写像 $R \to R'$、$R$-加群 $M$、および $R'$-加群 $N'$ が
与えられたとき、自然な写像
$$
\Ext^i_{R'}(M \otimes_R R', N') \to \text{Ext}^i_R(M, N')
$$
は $i \geq 0$ に対して同型である。
\end{lemma}

\begin{proof}
$M$ の自由分解 $F_\bullet$ を選ぶ。
$R \to R'$ は平坦なので、$F_\bullet \otimes_R R'$ は
$R'$ 上の $M \otimes_R R'$ の自由分解である。
主張は、写像
$$
\Hom_{R'}(F_\bullet \otimes_R R', N') \to
\Hom_R(F_\bullet, N')
$$
がホモロジー群上の同型を誘導するということである。
% P01-E0435: frozen source says homology for the cohomological Hom complex; preserved.
これは補題 \ref{algebra-lemma-adjoint-tensor-restrict} により複体の同型なので成り立つ。
\end{proof}






\section{Ext 群の一つの応用}
\label{algebra-section-ext-application}

\noindent
次がその応用である。

\begin{lemma}
\label{algebra-lemma-split-injection-after-completion}
$R$ をネーター環とし、$I \subset R$ を $R$ の Jacobson 根基に含まれる
イデアルとする。$N \to M$ を有限 $R$-加群の準同型とする。
任意に大きい $n$ に対して $N/I^nN \to M/I^nM$ が
分裂単射となるものが存在すると仮定する。
このとき $N \to M$ は分裂単射である。
\end{lemma}

\begin{proof}
$\varphi : N \to M$ が補題の仮定を満たすとする。
これは、任意に大きい $n$ に対して $\Ker(\varphi) \subset I^nN$
であることを含意する。したがって補題
\ref{algebra-lemma-intersection-powers-ideal-module} により、$\varphi$ は単射である。
$Q = M/N$ とおけば、短完全列
$$
0 \to N \to M \to Q \to 0.
$$
を得る。$Q$ の有限自由分解
$$
F_2 \xrightarrow{d_2} F_1 \xrightarrow{d_1} F_0 \to Q \to 0
$$
をとる。写像 $F_0 \to Q$ を持ち上げる写像
$\alpha : F_0 \to M$ を選べる。これは
$\beta \circ d_2 = 0$ を満たす写像 $\beta : F_1 \to N$ を誘導する。
上の拡大が分裂するための必要十分条件は、
$\beta = \gamma \circ d_1$ を満たす写像 $\gamma : F_0 \to N$ が
存在することである。言い換えれば、$\Ext^1_R(Q, N)$ における
$\beta$ の類が、上の短完全列が分裂することへの障害である。

\medskip\noindent
$N/I^nN \to M/I^nM$ が分裂単射となる十分大きい整数 $n$ をとる。
これは
$$
0 \to N/I^nN \to M/I^nM \to Q/I^nQ \to 0.
$$
が依然として短完全列であることを含意する。また列
$$
F_1/I^nF_1 \xrightarrow{d_1} F_0/I^nF_0 \to Q/I^nQ \to 0
$$
も依然として完全である。上と同様に論じると、$\beta$ が誘導する写像
$\overline{\beta} : F_1/I^nF_1 \to N/I^nN$ は、ある写像
$\overline{\gamma_n} : F_0/I^nF_0 \to N/I^nN$ に対して
$\overline{\gamma_n} \circ d_1$ に等しい。
$F_0$ は自由なので、$\overline{\gamma_n}$ を写像
$\gamma_n : F_0 \to N$ に持ち上げることができる。
すると $\beta - \gamma_n \circ d_1$ は $F_1$ から $I^nN$ への写像である。
言い換えれば、この $n$ に対して
$$
\beta \in
\Im\Big(\Hom_R(F_0, N) \to \Hom_R(F_1, N)\Big) + I^n\Hom_R(F_1, N).
$$
を得る。

\medskip\noindent
仮定により任意に大きい $n$ に対してこの性質が成り立つので、
補題 \ref{algebra-lemma-intersection-powers-ideal-module} により、
$\Hom_R(F_0, N)\allowbreak \to \Hom_R(F_1, N)$ の余核における
$\beta$ の像は零である。したがって望むとおり、$\beta$ は写像
$\Hom_R(F_0, N)\allowbreak \to \Hom_R(F_1, N)$ の像に属する。
\end{proof}

\section{Tor 群と平坦性}
\label{algebra-section-tor}

\noindent
この節では、前節で展開したホモロジー代数の一部を用いて
Tor 群とは何かを説明する。すなわち、$R$ を環、$M$, $N$ を
二つの $R$-加群とする。$M$ の自由 $R$-加群による分解
$F_\bullet$ を選ぶ。補題 \ref{algebra-lemma-resolution-by-finite-free} を参照せよ。
ホモロジー複体
$$
F_\bullet \otimes_R N
:
\ldots
\to F_2 \otimes_R N
\to F_1 \otimes_R N
\to F_0 \otimes_R N
$$
を考える。$\text{Tor}^R_i(M, N)$ をこの複体の第 $i$ ホモロジー群と定義する。
次の補題は、これがどの意味で矛盾なく定義されているかを説明する。

\begin{lemma}
\label{algebra-lemma-tor-welldefined}
$R$ を環、$M_1, M_2, N$ を $R$-加群とする。
$F_\bullet$ は加群 $M_1$ の自由分解、$G_\bullet$ は加群 $M_2$ の
自由分解であると仮定する。$\varphi : M_1 \to M_2$ を加群写像とし、
$M_1 = \Coker(d_{F, 1}) \to M_2 = \Coker(d_{G, 1})$ 上で
$\varphi$ を誘導する複体の写像 $\alpha : F_\bullet \to G_\bullet$ をとる。
補題 \ref{algebra-lemma-compare-resolutions} を参照せよ。このとき誘導写像
$$
H_i(\alpha) :
H_i(F_\bullet \otimes_R N)
\longrightarrow
H_i(G_\bullet \otimes_R N)
$$
は $\alpha$ の選び方に依存しない。$\varphi$ が同型ならば、
すべての写像 $H_i(\alpha)$ も同型である。
$M_1 = M_2$、$F_\bullet = G_\bullet$ であり、$\varphi$ が恒等写像ならば、
すべての写像 $H_i(\alpha)$ も恒等写像である。
\end{lemma}

\begin{proof}
この補題の証明は補題 \ref{algebra-lemma-ext-welldefined} の証明と同一である。
\end{proof}

\noindent
この補題は Tor 加群が矛盾なく定義されていることを含意するだけでなく、
構成 $(M, N) \mapsto \text{Tor}_i^R(M, N)$ の第一変数に関する
関手性も与える。第二変数に関する関手性はもちろん明らかである。
各 $\text{Tor}_i^R$ が実際に関手
$$
\text{Mod}_R \times \text{Mod}_R \to \text{Mod}_R.
$$
であることの確認は読者に委ねる。ここで $\text{Mod}_R$ は
$R$-加群の圏を表す。積圏の定義については
Categories, Definition \ref{categories-definition-product-category} を参照せよ。
すなわち、$R$-加群の射 $M_1 \to M_2$ と $N_1 \to N_2$ が
与えられたとき、可換図式
$$
\xymatrix{
\text{Tor}_i^R(M_1, N_1) \ar[r] \ar[d] &
\text{Tor}_i^R(M_1, N_2) \ar[d] \\
\text{Tor}_i^R(M_2, N_1) \ar[r] &
\text{Tor}_i^R(M_2, N_2) \\
}
$$
を得る。

\begin{lemma}
\label{algebra-lemma-long-exact-sequence-tor}
$R$ を環、$M$ を $R$-加群とする。
$0 \to N' \to N \to N'' \to 0$ が $R$-加群の短完全列であると仮定する。
このとき長完全列
$$
\text{Tor}_1^R(M, N')
\to \text{Tor}_1^R(M, N)
\to \text{Tor}_1^R(M, N'')
\to
M \otimes_R N'
\to M \otimes_R N
\to M \otimes_R N''
\to 0
$$
が存在する。
\end{lemma}

\begin{proof}
この証明は補題 \ref{algebra-lemma-long-exact-seq-ext} の証明と同じである。
\end{proof}

\noindent
$R$-加群のホモロジー二重複体
$$
\xymatrix{
\ldots \ar[r]^d &
A_{2, 0} \ar[r]^d &
A_{1, 0} \ar[r]^d &
A_{0, 0} \\
\ldots \ar[r]^d &
A_{2, 1} \ar[r]^d \ar[u]^\delta &
A_{1, 1} \ar[r]^d \ar[u]^\delta &
A_{0, 1} \ar[u]^\delta \\
\ldots \ar[r]^d &
A_{2, 2} \ar[r]^d \ar[u]^\delta &
A_{1, 2} \ar[r]^d \ar[u]^\delta &
A_{0, 2} \ar[u]^\delta \\
&
\ldots \ar[u]^\delta &
\ldots \ar[u]^\delta &
\ldots \ar[u]^\delta \\
}
$$
を考える。これは $d_{i, j} : A_{i, j} \to A_{i-1, j}$ と
$\delta_{i, j} : A_{i, j} \to A_{i, j-1}$ が次の性質をもつことを意味する。
\begin{enumerate}
\item 二つの $d_{i, j}$ の任意の合成は零である。言い換えれば、
二重複体の各行は複体である。
\item 二つの $\delta_{i, j}$ の任意の合成は零である。言い換えれば、
二重複体の各列は複体である。
\item 任意の組 $(i, j)$ に対して $\delta_{i-1, j} \circ d_{i, j}
= d_{i, j-1} \circ \delta_{i, j}$ である。言い換えれば、
すべての正方形は可換である。
\end{enumerate}
本来は、このような二重複体にはスペクトル系列を付随させるべきである。
しかし今のところは、もう少し易しいことを行えば十分である。

\medskip\noindent
すなわち、この節だけでは、二重複体 $(A_{\bullet, \bullet}, d, \delta)$
に対して $R(A)_j = \Coker(A_{1, j} \to A_{0, j})$ および
$U(A)_i = \Coker(A_{i, 1} \to A_{i, 0})$ とおく。
（文字 $R$ と $U$ は Right と Up を示唆している。）
写像 $\delta$ を用いて $R(A)_\bullet$ に複体の構造を入れる。
同様に写像 $d$ を用いて $U(A)_\bullet$ に複体の構造を入れる。
言い換えれば、次の巨大な可換図式を得る。
$$
\xymatrix{
\ldots \ar[r]^d &
U(A)_2 \ar[r]^d &
U(A)_1 \ar[r]^d &
U(A)_0 &
\\
\ldots \ar[r]^d &
A_{2, 0} \ar[r]^d \ar[u] &
A_{1, 0} \ar[r]^d \ar[u] &
A_{0, 0} \ar[r] \ar[u] &
R(A)_0 \\
\ldots \ar[r]^d &
A_{2, 1} \ar[r]^d \ar[u]^\delta &
A_{1, 1} \ar[r]^d \ar[u]^\delta &
A_{0, 1} \ar[r] \ar[u]^\delta &
R(A)_1 \ar[u]^\delta \\
\ldots \ar[r]^d &
A_{2, 2} \ar[r]^d \ar[u]^\delta &
A_{1, 2} \ar[r]^d \ar[u]^\delta &
A_{0, 2} \ar[r] \ar[u]^\delta &
R(A)_2 \ar[u]^\delta \\
&
\ldots \ar[u]^\delta &
\ldots \ar[u]^\delta &
\ldots \ar[u]^\delta &
\ldots \ar[u]^\delta \\
}
$$
（もちろん、これはもはや二重複体ではない。）
二重複体の射 $\Phi : (A_{\bullet, \bullet}, d, \delta)
\to (B_{\bullet, \bullet}, d, \delta)$ が何を意味するかは明らかであり、
これが複体の射 $R(\Phi) : R(A)_\bullet \to R(B)_\bullet$ および
$U(\Phi) : U(A)_\bullet \to U(B)_\bullet$ を誘導することも明らかである。

\begin{lemma}
\label{algebra-lemma-no-spectral-sequence}
$(A_{\bullet, \bullet}, d, \delta)$ を次を満たす二重複体とする。
\begin{enumerate}
\item 各行 $A_{\bullet, j}$ は $R(A)_j$ の分解である。
\item 各列 $A_{i, \bullet}$ は $U(A)_i$ の分解である。
\end{enumerate}
このとき標準同型
$$
H_i(R(A)_\bullet)
\cong
H_i(U(A)_\bullet).
$$
が存在する。この同型は、上の性質をもつ二重複体の射に関して関手的である。
\end{lemma}

\begin{proof}
% P01-E0436: frozen source has an unmatched extra closing parenthesis; preserved.
$H_i(R(A)_\bullet))$ と $H_i(U(A)_\bullet)$ が、第三の群と標準的に
同型であることを示す。すなわち
$$
\mathbf{H}_i(A) :=
\frac{
\{
(a_{i, 0}, a_{i-1, 1}, \ldots, a_{0, i})
\mid
d(a_{i, 0}) = \delta(a_{i-1, 1}), \ldots,
d(a_{1, i-1}) = \delta(a_{0, i})
\}}
{
\{
d(a_{i + 1, 0}) + \delta(a_{i, 1}),
d(a_{i, 1}) + \delta(a_{i-1, 2}),
\ldots,
d(a_{1, i}) + \delta(a_{0, i + 1})
\}
}
$$
とする。ここでは $a_{i, j}$ は $A_{i, j}$ の元を表すという記法を用いる。
言い換えれば、$\mathbf{H}_i$ の元はジグザグによって表される。
$i = 2$ の場合には次のようになる。
$$
\xymatrix{
a_{2, 0} \ar@{|->}[r] & d(a_{2, 0}) = \delta(a_{1, 1}) & \\
& a_{1, 1} \ar@{|->}[u] \ar@{|->}[r] & d(a_{1, 1}) = \delta(a_{0, 2}) \\
& & a_{0, 2} \ar@{|->}[u] \\
}
$$
もちろん「自明な」ジグザグ、すなわち
% P01-E0437: frozen source uses the invalid t-i indices; preserved.
$(0, \ldots, 0, -\delta(a_{t + 1, t-i}),
d(a_{t + 1, t-i}), 0, \ldots, 0)$ の形の元が生成する部分加群で割る。
標準準同型
$$
\mathbf{H}_i(A) \to H_i(R(A)_\bullet), \quad
(a_{i, 0}, a_{i-1, 1}, \ldots, a_{0, i}) \mapsto
\text{class of image of }a_{0, i}
$$
および
$$
\mathbf{H}_i(A) \to H_i(U(A)_\bullet), \quad
(a_{i, 0}, a_{i-1, 1}, \ldots, a_{0, i}) \mapsto
\text{class of image of }a_{i, 0}
$$
があることに注意する。

\medskip\noindent
まず、これらの写像が全射であることを示す。
$\overline{r} \in H_i(R(A)_\bullet)$ と仮定する。
% P01-E0438: frozen source calls the homological representative a cocycle; preserved.
$r \in R(A)_i$ を $\overline{r}$ の類を表すコサイクルとする。
$r$ に写る元 $a_{0, i} \in A_{0, i}$ をとる。
$\delta(r) = 0$ なので、$\delta(a_{0, i})$ は $d$ の像に属する。
したがって $d(a_{1, i-1}) = \delta(a_{0, i})$ を満たす元
$a_{1, i-1} \in A_{1, i-1}$ が存在する。
これはさらに $\delta(a_{1, i-1})$ が $d$ の核に属することを含意する
（なぜなら $d(\delta(a_{1, i-1})) = \delta(d(a_{1, i-1}))
= \delta(\delta(a_{0, i})) = 0$ だからである）。
行の完全性により、$d(a_{2, i-2}) = \delta(a_{1, i-1})$ を満たす元
$a_{2, i-2}$ が見つかる。完全なジグザグが得られるまで同様に続ける。
もちろん $\mathbf{H}_i \to H_i(U(A))$ の全射性も同様に示される。

\medskip\noindent
単射性を示すには、まったく同じように論じる。
すなわち、$H_i(R(A)_\bullet)$ で零に写るジグザグ
$(a_{i, 0}, a_{i-1, 1}, \ldots, a_{0, i})$ が与えられたとする。
これは $a_{0, i}$ が $\Coker(A_{i, 1} \to A_{i, 0})$ の元であって、
% P01-E0439: frozen source transposes the R(A) cokernel indices; preserved.
写像
$\delta : \Coker(A_{i + 1, 1} \to A_{i + 1, 0}) \to
\Coker(A_{i, 1} \to A_{i, 0})$ の像に属する元へ写ることを意味する。
言い換えれば、$a_{0, i}$ は
$\delta \oplus d : A_{0, i + 1} \oplus A_{1, i} \to A_{0, i}$ の像に属する。
自明なジグザグの定義から、与えられたジグザグを自明なもので
変更して $a_{0, i} = 0$ と仮定できる。
これは直ちに $d(a_{1, i-1}) = 0$ を含意する。
各行は完全なので、$a_{1, i-1}$ は
$d : A_{2, i-1} \to A_{1, i-1}$ の像に属する。
したがって再び自明なジグザグで変更し、与えられたジグザグが
$(a_{i, 0}, a_{i-1, 1}, \ldots, a_{2, i-2}, 0, 0)$ の形であると
仮定できる。同様に続ければ、所望の単射性を得る。

\medskip\noindent
$\Phi : (A_{\bullet, \bullet}, d, \delta)
\to (B_{\bullet, \bullet}, d, \delta)$ が、ともに補題の条件を満たす
二重複体の射ならば、明らかに可換図式
$$
\xymatrix{
H_i(U(A)_\bullet) \ar[d] &
\mathbf{H}_i(A) \ar[r] \ar[l] \ar[d] &
H_i(R(A)_\bullet) \ar[d] \\
H_i(U(B)_\bullet) &
\mathbf{H}_i(B) \ar[r] \ar[l] &
H_i(R(B)_\bullet) \\
}
$$
を得る。これで関手性が示された。
\end{proof}

\begin{remark}
\label{algebra-remark-signs-double-complex}
上で構成した同型が「正しい」のは符号を除いてである。
ホモロジー代数のかなりの部分は、さまざまな写像の符号を選び、
適切な符号の介入のもとで図式の可換性を示すことに関わる。
今のところは上の証明で与えた同型をそのまま用い、符号は後で考える。
\end{remark}

\begin{lemma}
\label{algebra-lemma-tor-left-right}
$R$ を環とする。任意の $i \geq 0$ に対して、関手
$\text{Mod}_R \times \text{Mod}_R \to \text{Mod}_R$、
$(M, N) \mapsto \text{Tor}_i^R(M, N)$ と
$(M, N) \mapsto \text{Tor}_i^R(N, M)$ は標準的に同型である。
\end{lemma}

\begin{proof}
$F_\bullet$ を加群 $M$ の自由分解、$G_\bullet$ を加群 $N$ の自由分解とする。
次のように定義される二重複体 $(A_{i, j}, d, \delta)$ を考える。
\begin{enumerate}
\item $A_{i, j} = F_i \otimes_R G_j$ とおく。
% P01-E0440: frozen source omits the base-ring subscript in both codomains; preserved.
\item $d_{i, j} : F_i \otimes_R G_j \to F_{i-1} \otimes G_j$ を
$d_{F, i} \otimes \text{id}$ に等しいものとおく。
\item $\delta_{i, j} : F_i \otimes_R G_j \to F_i \otimes G_{j-1}$ を
$\text{id} \otimes d_{G, j}$ に等しいものとおく。
\end{enumerate}
この二重複体は通常、単に $F_\bullet \otimes_R G_\bullet$ と書かれる。

\medskip\noindent
各 $G_j$ は自由、したがって平坦なので、二重複体の各行は
ホモロジー次数 $0$ を除いて完全である。
各 $F_i$ も自由、したがって平坦なので、二重複体の各列は
ホモロジー次数 $0$ を除いて完全である。
ゆえにこの二重複体は補題 \ref{algebra-lemma-no-spectral-sequence} の条件を満たす。

\medskip\noindent
補題の内容を見るために $R(A)_\bullet$ と $U(A)_\bullet$ を計算する。
すなわち
\begin{eqnarray*}
R(A)_i & = & \Coker(A_{1, i} \to A_{0, i}) \\
& = & \Coker(F_1 \otimes_R G_i \to F_0 \otimes_R G_i) \\
& = & \Coker(F_1 \to F_0) \otimes_R G_i \\
& = & M \otimes_R G_i
\end{eqnarray*}
である。実際、これらの同型は微分 $\delta$ と両立し、ホモロジー複体として
$R(A)_\bullet = M \otimes_R G_\bullet$ であることが分かる。
まったく同様に $U(A)_\bullet = F_\bullet \otimes_R N$ である。したがって
\begin{eqnarray*}
\text{Tor}_i^R(M, N)
& = & H_i(F_\bullet \otimes_R N) \\
& = & H_i(U(A)_\bullet) \\
& = & H_i(R(A)_\bullet) \\
& = & H_i(M \otimes_R G_\bullet) \\
& = & H_i(G_\bullet \otimes_R M) \\
& = & \text{Tor}_i^R(N, M)
\end{eqnarray*}
を得る。ここで第三の等号は補題 \ref{algebra-lemma-no-spectral-sequence} であり、
第五の等号ではテンソル積の同型 $V \otimes W = W \otimes V$ を用いた。

\medskip\noindent
関手性。$R$-加群 $M_\nu$, $N_\nu$、$\nu = 1, 2$ が与えられたとする。
$\varphi : M_1 \to M_2$ と $\psi : N_1 \to N_2$ を $R$-加群の射とする。
$M_\nu$ の自由分解 $F_{\nu, \bullet}$ と、$N_\nu$ の自由分解
$G_{\nu, \bullet}$ があると仮定する。
補題 \ref{algebra-lemma-compare-resolutions} により、$\varphi$ および $\psi$ と両立する
複体の写像 $\alpha : F_{1, \bullet} \to F_{2, \bullet}$ と
$\beta : G_{1, \bullet} \to G_{2, \bullet}$ を選べる。
組 $(\alpha, \beta)$ は二重複体の射
$$
\alpha \otimes \beta :
F_{1, \bullet} \otimes_R G_{1, \bullet}
\longrightarrow
F_{2, \bullet} \otimes_R G_{2, \bullet}
$$
を誘導すると主張する。これは、
$F_{1, i} \otimes_R G_{1, j} \to F_{2, i} \otimes_R G_{2, j}$ が
$\alpha_i \otimes \beta_j$ で与えられることを用いる、きわめて直接的な確認である。
ここで $\alpha_i$、それぞれ $\beta_j$ は、$\alpha$、それぞれ $\beta$ の
次数 $i$、それぞれ $j$ の成分である。また、誘導写像
$R(F_{1, \bullet} \otimes_R G_{1, \bullet})_\bullet
\to R(F_{2, \bullet} \otimes_R G_{2, \bullet})_\bullet$ は、
$\varphi \otimes \beta$ が誘導する写像 $M_1 \otimes_R G_{1, \bullet}
\to M_2 \otimes_R G_{2, \bullet}$ と一致することも容易に確認できる。
$U(-)_\bullet$ 複体に誘導される写像についても同様である。
したがって関手性に関する主張は、補題 \ref{algebra-lemma-no-spectral-sequence} の
関手性に関する主張から従う。
\end{proof}

\begin{remark}
\label{algebra-remark-curiosity-signs-swap}
上で $M = N$ とした場合には興味深いことが起こる。
この場合、標準写像 $\text{Tor}_i^R(M, M)
\to \text{Tor}_i^R(M, M)$ を得る。この写像は恒等写像ではない。
$i = 0$ の場合でさえ恒等写像ではないことに注意する。
例えば $V$ が体上の次元 $n$ のベクトル空間ならば、入れ替え写像
$V \otimes_k V \to V \otimes_k V$ は、固有値 $+1$ を $(n^2 + n)/2$ 個、
固有値 $-1$ を $(n^2-n)/2$ 個もつ。標数 $2$ では対角化可能ですらない。
写像の符号を変えても、この現象は消えないことに注意する。
\end{remark}

\begin{lemma}
\label{algebra-lemma-tor-noetherian}
$R$ をネーター環、$M$, $N$ を有限 $R$-加群とする。
このとき、すべての $p$ に対して $\text{Tor}_p^R(M, N)$ は
有限 $R$-加群である。
\end{lemma}

\begin{proof}
% P01-E0441: frozen source says cohomology for this homological complex; preserved.
$\text{Tor}_p^R(M, N)$ は、各 $F_n$ が有限自由 $R$-加群である複体
$F_\bullet \otimes_R N$ のコホモロジー群として計算されるからである。
補題 \ref{algebra-lemma-resolution-by-finite-free} を参照せよ。
\end{proof}

\begin{lemma}
\label{algebra-lemma-characterize-flat}
$R$ を環、$M$ を $R$-加群とする。次は同値である。
\begin{enumerate}
\item 加群 $M$ は $R$ 上平坦である。
\item すべての $i > 0$ に対して関手 $\text{Tor}_i^R(M, -)$ は零である。
\item 関手 $\text{Tor}_1^R(M, -)$ は零である。
\item すべてのイデアル $I \subset R$ に対して
$\text{Tor}_1^R(M, R/I) = 0$ である。
\item すべての有限生成イデアル $I \subset R$ に対して
$\text{Tor}_1^R(M, R/I) = 0$ である。
\end{enumerate}
\end{lemma}

\begin{proof}
$M$ が平坦であると仮定し、$N$ を $R$-加群とする。
$F_\bullet$ を $N$ の自由分解とする。
$M$ の平坦性により、$F_\bullet \otimes_R M$ は
$N \otimes_R M$ の分解である。したがってすべての高次 Tor 群は消滅する。

\medskip\noindent
残るのは、最後の条件が $M$ の平坦性を含意することを示すことだけである。
$I \subset R$ をイデアルとする。短完全列
$0 \to I \to R \to R/I \to 0$ を考え、
補題 \ref{algebra-lemma-long-exact-sequence-tor} を適用する。完全列
$$
\text{Tor}_1^R(M, R/I) \to
M \otimes_R I \to
M \otimes_R R \to
M \otimes_R R/I \to
0
$$
を得る。明らかに $M \otimes_R R = M$ なので、最後の仮定は、
すべての有限生成イデアル $I$ に対して $M \otimes_R I \to M$ が
単射であることを含意する。したがって補題 \ref{algebra-lemma-flat} により
$M$ は平坦である。
\end{proof}

\begin{remark}
\label{algebra-remark-Tor-ring-mod-ideal}
補題 \ref{algebra-lemma-characterize-flat} の証明は、実際には
$$
\text{Tor}_1^R(M, R/I)
=
\Ker(I \otimes_R M \to M).
$$
を示している。
\end{remark}

\section{Tor の関手性}
\label{algebra-section-functoriality-tor}

\noindent
この節では、環の変更などに関する $\text{Tor}$ の関手性を簡単に論じる。
確認すべき事項を列挙する。
\begin{enumerate}
\item 環写像 $R \to R'$、$R$-加群 $M$、および $R'$-加群 $N'$ が
与えられたとき、$R$-加群 $\text{Tor}_i^R(M, N')$ は
自然な $R'$-加群構造をもつ。
\item 環写像 $R \to R'$ と $R$-加群 $M$, $N$ が与えられたとき、
自然な $R$-加群写像
$\text{Tor}_i^R(M, N) \to \text{Tor}_i^{R'}(M \otimes_R R', N \otimes_R R')$
がある。
\item 環写像 $R \to R'$、$R$-加群 $M$、および $R'$-加群 $N'$ が
与えられたとき、自然な $R'$-加群写像
$\text{Tor}_i^R(M, N') \to \text{Tor}_i^{R'}(M \otimes_R R', N')$
が存在する。
\end{enumerate}

\begin{lemma}
\label{algebra-lemma-flat-base-change-tor}
平坦環写像 $R \to R'$ と $R$-加群 $M$, $N$ が与えられたとき、
自然な $R$-加群写像
$\text{Tor}_i^R(M, N)\otimes_R R'
\to \text{Tor}_i^{R'}(M \otimes_R R', N \otimes_R R')$
はすべての $i$ に対して同型である。
\end{lemma}

\begin{proof}
省略する。$R$ 上の $M$ の自由分解 $F_\bullet$ は、$R$ 上で $R'$ と
テンソル積を取っても完全なままであり、したがって
$(F_\bullet \otimes_R N)\otimes_R R'$ は $R'$ 上の Tor 群を
計算するので、主張は成り立つ。
\end{proof}

\noindent
次の補題はほかに適切な置き場所がないように思われる。

\begin{lemma}
\label{algebra-lemma-tor-commutes-filtered-colimits}
$R$ を環とする。$M = \colim M_i$ を $R$-加群のフィルター余極限とし、
$N$ を $R$-加群とする。このとき、すべての $n$ に対して
$\text{Tor}_n^R(M, N) = \colim \text{Tor}_n^R(M_i, N)$ である。
\end{lemma}

\begin{proof}
$N$ の自由分解 $F_\bullet$ を選ぶ。
補題 \ref{algebra-lemma-tensor-products-commute-with-limits} により、複体として
$F_\bullet \otimes_R M = \colim F_\bullet \otimes_R M_i$ である。
したがって補題 \ref{algebra-lemma-directed-colimit-exact} により結論を得る。
\end{proof}








\section{射影加群}
\label{algebra-section-projective}

\noindent
射影加群に関するいくつかの補題を与える。

\begin{definition}
\label{algebra-definition-projective}
$R$ を環とする。$R$-加群 $P$ が {\it 射影的} であるとは、関手
$\Hom_R(P, -) : \text{Mod}_R \to \text{Mod}_R$ が完全関手であることをいう。
\end{definition}

\noindent
補題 \ref{algebra-lemma-hom-exact} により、任意の $R$-加群 $M$ に対して
関手 $\Hom_R(M, - )$ は左完全である。
したがって $P$ が射影的であるための条件が実際に意味するのは、
$R$-加群の全射 $N \to N'$ が与えられたとき、写像
$\Hom_R(P, N) \to \Hom_R(P, N')$ が全射であるということである。

\begin{lemma}
\label{algebra-lemma-characterize-projective}
$R$ を環、$P$ を $R$-加群とする。次は同値である。
\begin{enumerate}
\item $P$ は射影的である。
\item $P$ は自由 $R$-加群の直和因子である。
\item すべての $R$-加群 $M$ に対して $\Ext^1_R(P, M) = 0$ である。
\end{enumerate}
\end{lemma}

\begin{proof}
$P$ が射影的であると仮定する。$F$ が自由 $R$-加群であるような全射
$\pi : F \to P$ を選ぶ。$P$ は射影的なので、
$\pi \circ i = \text{id}_P$ を満たす $i \in \Hom_R(P, F)$ が存在する。
言い換えれば $F \cong \Ker(\pi) \oplus i(P)$ であり、
$P$ は $F$ の直和因子である。

\medskip\noindent
逆に、$P \oplus Q = F$ が自由 $R$-加群であると仮定する。
自由加群 $F = \bigoplus_{i \in I} R$ は射影的であることに注意する。
実際、$\Hom_R(F, M) = \prod_{i \in I} M$ であり、関手
$M \mapsto \prod_{i \in I} M$ は完全である。
関手として $\Hom_R(F, -) = \Hom_R(P, -) \times \Hom_R(Q, -)$
なので、$P$ と $Q$ はともに射影的である。

\medskip\noindent
$P \oplus Q = F$ が自由 $R$-加群であると仮定する。このとき、
% P01-E0442: frozen source omits the transition arrow after the ellipsis; preserved.
$$
\ldots F \xrightarrow{a} F \xrightarrow{b} F \to P \to 0
$$
の形の自由分解 $F_\bullet$ がある。ここで写像 $a, b$ は交互に現れ、
$P$ および $Q$ への射影に等しい。
したがって複体 $\Hom_R(F_\bullet, M)$ は次数 $\geq 1$ で分裂完全であり、
(3) の消滅が従う。

\medskip\noindent
すべての $R$-加群 $M$ に対して $\Ext^1_R(P, M) = 0$ と仮定する。
自由分解 $F_\bullet \to P$ を選び、
$M = \Im(F_1 \to F_0) = \Ker(F_0 \to P)$ とおく。
商写像 $\pi : F_1 \to M$ の類が与える元
$\xi \in \Ext^1_R(P, M)$ を考える。$\xi$ は零なので、
$\pi = s \circ (F_1 \to F_0)$ を満たす写像 $s : F_0 \to M$ が存在する。
これは明らかに
$$
F_0 = \Ker(s) \oplus \Ker(F_0 \to P) =
P \oplus \Ker(F_0 \to P)
$$
を意味し、結論を得る。
\end{proof}

\begin{lemma}
\label{algebra-lemma-characterize-finite-projective-noetherian}
$R$ をネーター環、$P$ を有限 $R$-加群とする。
すべての有限 $R$-加群 $M$ に対して $\Ext^1_R(P, M) = 0$ ならば、
$P$ は射影的である。
\end{lemma}

\noindent
この補題は強められる。$R$ がネーターであると仮定しなければ、
有限表示 $R$-加群に対する版がある。ネーターの場合には、$M$ が
すべての有限長加群を走る版がある。

\begin{proof}
核が $M$ である全射 $R^{\oplus n} \to P$ を選ぶ。
$\Ext^1_R(P, M) = 0$ なので、この全射は分裂する。
補題 \ref{algebra-lemma-characterize-projective} により結論を得る。
\end{proof}

\begin{lemma}
\label{algebra-lemma-direct-sum-projective}
射影加群の直和は射影的である。
\end{lemma}

\begin{proof}
これは、補題 \ref{algebra-lemma-characterize-projective} における
自由加群の直和因子としての射影加群の特徴づけから従う。
\end{proof}

\begin{lemma}
\label{algebra-lemma-lift-projective-module}
$R$ を環、$I \subset R$ を冪零イデアルとする。
$\overline{P}$ を射影 $R/I$-加群とする。このとき、
$P/IP \cong \overline{P}$ を満たす射影 $R$-加群 $P$ が存在する。
\end{lemma}

\begin{proof}
補題 \ref{algebra-lemma-characterize-projective} により、ある集合 $A$ と
ある $R/I$-加群 $\overline{K}$ に対して直和分解
$\bigoplus_{\alpha \in A} R/I = \overline{P} \oplus \overline{K}$ を選べる。
$A$ 上の自由 $R$-加群を $F = \bigoplus_{\alpha \in A} R$ と書く。
$\bigoplus_{\alpha \in A} R/I$ の直和因子 $\overline{P}$ に付随する
射影子 $\overline{p}$ の持ち上げ $p : F \to F$ を選ぶ。
$p^2 - p \in \text{End}_R(F)$ は $F$ の冪零自己準同型であることに注意する
（$I$ は冪零で、$p^2 - p$ の行列要素は $I$ に属するからである。
より正確には、$I^n = 0$ ならば $(p^2 - p)^n = 0$ である）。
したがって補題 \ref{algebra-lemma-lift-idempotents-noncommutative} により、
$p$ の選び方を修正して $p$ が射影子であると仮定できる。
$P = \Im(p)$ とおく。
\end{proof}

\begin{lemma}
\label{algebra-lemma-lift-finite-projective-module}
$R$ を環、$I \subset R$ を局所冪零イデアルとする。
$\overline{P}$ を有限射影 $R/I$-加群とする。このとき、
$P/IP \cong \overline{P}$ を満たす有限射影 $R$-加群 $P$ が存在する。
\end{lemma}

\begin{proof}
補題 \ref{algebra-lemma-characterize-projective} により、$\overline{P}$ は自由
$R/I$-加群 $\bigoplus_{\alpha \in A} R/I$ の直和因子であることを思い出す。
$\overline{P}$ は有限なので、ある有限部分集合 $A' \subset A$ に対して
$\overline{P}$ は $\bigoplus_{\alpha \in A'} R/I$ に含まれる。
したがって、ある $n$ とある $R/I$-加群 $\overline{K}$ に対して
直和分解 $(R/I)^{\oplus n} = \overline{P} \oplus \overline{K}$ が
あると仮定してよい。$(R/I)^{\oplus n}$ の直和因子 $\overline{P}$ に
付随する射影子 $\overline{p}$ の持ち上げ
$p \in \text{Mat}(n \times n, R)$ を選ぶ。
$p^2 - p \in \text{Mat}(n \times n, R)$ は冪零である。
実際、$I$ は局所冪零であり、$p^2 - p$ の行列要素 $c_{ij}$ は $I$ に
属するので、ある $t > 0$ に対して $c_{ij}^t = 0$ であり、
行列要素を見れば $(p^2 - p)^{tn^2} = 0$ である。
したがって補題 \ref{algebra-lemma-lift-idempotents-noncommutative} により、
$p$ の選び方を修正して $p$ が射影子であると仮定できる。
$P = \Im(p)$ とおく。
\end{proof}

\begin{lemma}
\label{algebra-lemma-lift-projective}
$R$ を環、$I \subset R$ をイデアル、$M$ を $R$-加群とし、次を仮定する。
\begin{enumerate}
\item $I$ は冪零である。
\item $M/IM$ は射影 $R/I$-加群である。
\item $M$ は平坦 $R$-加群である。
\end{enumerate}
このとき $M$ は射影 $R$-加群である。
\end{lemma}

\begin{proof}
補題 \ref{algebra-lemma-lift-projective-module} により、射影 $R$-加群 $P$ と
同型 $P/IP \to M/IM$ を見つけられる。
$M$ が $P$ と同型であることを示せば証明は終わる。
$P$ は射影的なので、写像 $P \to P/IP \to M/IM$ を、$I$ を法として
同型となる $R$-加群写像 $P \to M$ に持ち上げられる。
ある $n$ に対して $I^n = 0$ なので、フィルトレーション
\begin{align*}
0 = I^nM \subset I^{n - 1}M \subset \ldots \subset IM \subset M \\
0 = I^nP \subset I^{n - 1}P \subset \ldots \subset IP \subset P
\end{align*}
を用いれば、誘導写像
$I^aP/I^{a + 1}P \to I^aM/I^{a + 1}M$ が全単射であることを示せば
十分だと分かる。$P$ と $M$ はともに平坦 $R$-加群なので、これは
$P \to M$ が誘導する写像
$$
I^a/I^{a + 1} \otimes_{R/I} P/IP
\longrightarrow
I^a/I^{a + 1} \otimes_{R/I} M/IM
$$
と同一視できる。$P/IP \to M/IM$ が同型となるように $P \to M$ を
選んだので、結論を得る。
\end{proof}

\begin{lemma}
\label{algebra-lemma-projective-modulo-two-ideals}
$R$ を環、$I, J \subset R$ を $I \cap J = 0$ を満たすイデアルとする。
$P/IP$ が射影 $R/I$-加群で、$P/JP$ が射影 $R/J$-加群となる
$R$-加群 $P$ をとる。このとき $P$ は射影 $R$-加群である。
\end{lemma}

\begin{proof}
$F$ が自由 $R$-加群であるような全射 $p : F \to P$ を選ぶ。
$P/IP$ は射影 $R/I$-加群なので、$p \bmod I$ の右逆となる写像
$f : P/IP \to F/IF$ を選べる。商写像 $F/JF \to F/(I + J)F$ と
$p$ が誘導する写像
$q : F/JF \to F/(I + J)F \times_{P/(I + J)P} P/JP$ を考える。
$q$ は $R/J$-加群の全射であることに注意する（細部は省略する）。
$f$ と恒等写像が誘導する写像
$f' : P/JP \to F/(I + J)F \times_{P/(I + J)P} P/JP$ を考える。
$P/JP$ は射影 $R/J$-加群で $q$ は全射なので、
$q \circ g = f'$ を満たす写像 $g : P/JP \to F/JF$ を選べる。
すると $f \bmod I + J = g \bmod I + J$ である。
$I \cap J = 0$ により
$F = F/JF \times_{F/(I + J)F} F/IF$ なので、明確に定まる
$R$-加群準同型 $h = (f, g) : P \to F$ を得る。
写像 $a = p \circ h : P \to P$ は $I$ を法としても $J$ を法としても
恒等写像に簡約される。このとき $a$ は部分加群 $IP$ 上にも
恒等写像を誘導する。実際、$i_\alpha \in I$、$x_\alpha \in P$ を用いて
$x = \sum i_\alpha x_\alpha$ と書けば、
$a(x_\alpha) = x_\alpha + j_\alpha$ となる $j_\alpha \in J$ があり、
$i_\alpha j_\alpha = 0$ なので
$a(x) = \sum i_\alpha a(x_\alpha) = \sum i_\alpha x_\alpha = x$ である。
したがって $a : P \to P$ は $IP$ 上でも $P/IP$ 上でも恒等写像として
作用し、$a$ は $P$ の自己同型であると結論する（細部は省略する）。
ゆえに $P$ は $F$ の直和因子であり、したがって射影的である。
\end{proof}

\section{有限射影加群}
\label{algebra-section-finite-projective-modules}

\begin{definition}
\label{algebra-definition-locally-free}
$R$ を環、$M$ を $R$-加群とする。
\begin{enumerate}
\item $\Spec(R)$ を標準開集合 $D(f_i)$、$i \in I$ で被覆でき、
各 $i \in I$ に対して $M_{f_i}$ が自由 $R_{f_i}$-加群となるとき、
$M$ は {\it 局所自由} であるという。
\item 各 $M_{f_i}$ が有限自由となるように被覆を選べるとき、
$M$ は {\it 有限局所自由} であるという。
\item 各 $M_{f_i}$ が $R_{f_i}^{\oplus r}$ と同型になるように被覆を
選べるとき、$M$ は {\it 階数 $r$ の有限局所自由} であるという。
\end{enumerate}
\end{definition}

\noindent
有限局所自由 $R$-加群は補題 \ref{algebra-lemma-cover} により自動的に有限表示である。
さらに、$M$ が環 $R$ 上の階数 $r$ の有限局所自由加群で、$R$ が零環で
なければ、補題 \ref{algebra-lemma-rank} により $r$ は一意に定まる
（局所化のうち少なくとも一つ $	R_{f_i}$ は零環でないからである）。
% [P01-E1311] The frozen source contains a literal tab inside the preceding math span.

\begin{lemma}
\label{algebra-lemma-finite-projective}
$R$ を環、$M$ を $R$-加群とする。次は同値である。
\begin{enumerate}
\item $M$ は有限表示かつ $R$ 上平坦である。
\item $M$ は有限射影的である。
\item $M$ は有限自由 $R$-加群の直和因子である。
\item $M$ は有限表示であり、すべての
$\mathfrak p \in \Spec(R)$ に対して局所化 $M_{\mathfrak p}$ は自由である。
\item $M$ は有限表示であり、すべての極大イデアル
$\mathfrak m \subset R$ に対して局所化 $M_{\mathfrak m}$ は自由である。
\item $M$ は有限かつ局所自由である。
\item $M$ は有限局所自由である。
\item $M$ は有限であり、すべての素イデアル $\mathfrak p$ に対して
$M_{\mathfrak p}$ は自由で、関数
$$
\rho_M : \Spec(R) \to \mathbf{Z}, \quad
\mathfrak p
\longmapsto
\dim_{\kappa(\mathfrak p)} M \otimes_R \kappa(\mathfrak p)
$$
は Zariski 位相について局所定数である。
\end{enumerate}
\end{lemma}

\begin{proof}
まず $M$ が有限射影的、すなわち (2) が成り立つと仮定する。
全射 $R^n \to M$ をとり、その核を $K$ とする。
$M$ は射影的なので、$0 \to K \to R^n \to M \to 0$ は分裂する。
したがって (2) $\Rightarrow$ (3) である。
(3) $\Rightarrow$ (2) は、射影加群の直和因子が射影的であることから従う。
補題 \ref{algebra-lemma-characterize-projective} を参照せよ。

\medskip\noindent
(3) を仮定すると、$K \oplus M \cong R^{\oplus n}$ と書ける。
ゆえに $K$ は $R^n$ の直和因子であり、したがって有限生成である。
これは $M = R^{\oplus n}/K$ が有限表示であることを示す。
すなわち (3) $\Rightarrow$ (1) である。

\medskip\noindent
$M$ が有限表示かつ平坦、すなわち (1) が成り立つと仮定する。
(7) が成り立つことを示す。任意の素イデアル $\mathfrak p$ と、
$M \otimes_R \kappa(\mathfrak p)$ の基底に写る
$x_1, \ldots, x_r \in M$ を選ぶ。
Nakayama の補題（補題 \ref{algebra-lemma-NAK-localization} の形）により、
ある $g \in R$、$g \not \in \mathfrak p$ に対して、これらの元は $M_g$ を生成する。
対応する全射 $\varphi : R_g^{\oplus r} \to M_g$ は次の二つの性質をもつ。
(a) $\Ker(\varphi)$ は有限 $R_g$-加群であり（補題 \ref{algebra-lemma-extension} を参照）、
(b) $M_g$ の $R_g$ 上の平坦性により
$\Ker(\varphi) \otimes \kappa(\mathfrak p) = 0$ である
（補題 \ref{algebra-lemma-flat-tor-zero} を参照）。
したがって再び Nakayama の補題により、
$\Ker(\varphi)_{g'} = 0$ となる $g' \in R_g$ が存在する。
言い換えれば、$M_{gg'}$ は自由である。
% [P01-E0443] Frozen localization notation retained; see the part-local errata ledger.

\medskip\noindent
有限局所自由加群は補題 \ref{algebra-lemma-cover} により有限加群なので、
(7) $\Rightarrow$ (6) である。
(6) $\Rightarrow$ (7) および (7) $\Rightarrow$ (8) は明らかである。

\medskip\noindent
有限局所自由加群は補題 \ref{algebra-lemma-cover} により有限表示加群なので、
(7) $\Rightarrow$ (4) である。
もちろん (4) は (5) を含意する。
平坦性は局所的に確認できるので（補題 \ref{algebra-lemma-flat-localization} を参照）、
(5) は (1) を含意すると結論する。
ここまでで
$$
\xymatrix{
(2) \ar@{<=>}[r] & (3) \ar@{=>}[r] & (1) \ar@{=>}[r] &
(7)  \ar@{<=>}[r] \ar@{=>}[rd] \ar@{=>}[d] & (6) \\
& & (5) \ar@{=>}[u] & (4) \ar@{=>}[l] & (8)
}
$$
を得た。

\medskip\noindent
$M$ が (1), (4), (5), (6), (7) を満たすと仮定する。
(3) が成り立つことを示す。$M$ が射影的であることを示せば十分である。
$\Hom_R(M, -)$ が完全であることを示さなければならない。
$R$-加群の短完全列 $0 \to N'' \to N \to N'\to 0$ をとる。
次が完全であることを示せばよい。
$0 \to \Hom_R(M, N'') \to \Hom_R(M, N) \to
\Hom_R(M, N') \to 0$。
$M$ は有限局所自由なので、$M_{f_i}$ が有限自由となる被覆
$\Spec(R) = \bigcup D(f_i)$ が存在する。
補題 \ref{algebra-lemma-hom-from-finitely-presented} により、
$$
0 \to \Hom_R(M, N'')_{f_i} \to \Hom_R(M, N)_{f_i} \to
\Hom_R(M, N')_{f_i} \to 0
$$
は
$0 \to \Hom_{R_{f_i}}(M_{f_i}, N''_{f_i}) \to
\Hom_{R_{f_i}}(M_{f_i}, N_{f_i}) \to
\Hom_{R_{f_i}}(M_{f_i}, N'_{f_i}) \to 0$
に等しい。これは、$M_{f_i}$ が自由で、局所化
$0 \to N''_{f_i} \to N_{f_i} \to N'_{f_i} \to 0$ が完全
（局所化は完全だから）なので完全である。したがって補題 \ref{algebra-lemma-cover} により、
$0 \to \Hom_R(M, N'') \to \Hom_R(M, N) \to
\Hom_R(M, N') \to 0$ は完全である。

\medskip\noindent
最後に (8) を仮定する。極大イデアル $\mathfrak m \subset R$ を選ぶ。
$M \otimes_R \kappa(\mathfrak m) = M/\mathfrak mM$ の
$\kappa(\mathfrak m)$-基底に写る $x_1, \ldots, x_r \in M$ を選ぶ。
特に $\rho_M(\mathfrak m) = r$ である。
Nakayama の補題 \ref{algebra-lemma-NAK} により、ある $f \in R$、
$f \not \in \mathfrak m$ に対して $x_1, \ldots, x_r$ は $R_f$ 上 $M_f$ を生成する。
$\rho_M$ が局所定数であるという仮定により、ある $g \in R$、
$g \not \in \mathfrak m$ に対して $\rho_M$ は $D(g)$ 上で定数 $r$ である。
次が同型であると主張する。
$$
\Psi : R_{fg}^{\oplus r} \longrightarrow M_{fg}, \quad
(a_1, \ldots, a_r) \longmapsto \sum a_i x_i
$$
この主張は $M$ が有限局所自由、すなわち (7) が成り立つことを示す。
主張を見るには、補題 \ref{algebra-lemma-characterize-zero-local} により、
すべての $\mathfrak p \in D(fg)$ に対して誘導された局所化写像
$\Psi_{\mathfrak p} : R_{\mathfrak p}^{\oplus r} \to M_{\mathfrak p}$
が同型であることを示せば十分である。
$f$ の選び方により写像 $\Psi_{\mathfrak p}$ は全射である。
仮定 (8) により
$M_{\mathfrak p} \cong R_{\mathfrak p}^{\oplus \rho_M(\mathfrak p)}$
であり、$g$ の選び方により $\rho_M(\mathfrak p) = r$ である。
したがって $\Psi_{\mathfrak p}$ は全射
$R_{\mathfrak p}^{\oplus r} \to
M_{\mathfrak p} \cong R_{\mathfrak p}^{\oplus r}$
を定めるので、補題 \ref{algebra-lemma-fun} により同型である。
（もちろん、この最後の事実は単純な行列の議論からも従う。）
\end{proof}

\begin{lemma}
\label{algebra-lemma-finite-projective-reduced}
$R$ を被約環、$M$ を $R$-加群とする。このとき補題
\ref{algebra-lemma-finite-projective} の同値な条件は、次とも同値である。
\begin{enumerate}
\item[(9)] $M$ は有限であり、関数
$\rho_M : \Spec(R) \to \mathbf{Z}$、$\mathfrak p \mapsto
\dim_{\kappa(\mathfrak p)} M \otimes_R \kappa(\mathfrak p)$
は Zariski 位相について局所定数である。
\end{enumerate}
\end{lemma}

\begin{proof}
極大イデアル $\mathfrak m \subset R$ を選ぶ。
$M \otimes_R \kappa(\mathfrak m) = M/\mathfrak mM$ の
$\kappa(\mathfrak m)$-基底に写る $x_1, \ldots, x_r \in M$ を選ぶ。
特に $\rho_M(\mathfrak m) = r$ である。
Nakayama の補題 \ref{algebra-lemma-NAK} により、ある $f \in R$、
$f \not \in \mathfrak m$ に対して $x_1, \ldots, x_r$ は $R_f$ 上 $M_f$ を生成する。
$\rho_M$ が局所定数であるという仮定により、ある $g \in R$、
$g \not \in \mathfrak m$ に対して $\rho_M$ は $D(g)$ 上で定数 $r$ である。
次が同型であると主張する。
$$
\Psi : R_{fg}^{\oplus r} \longrightarrow M_{fg}, \quad
(a_1, \ldots, a_r) \longmapsto \sum a_i x_i
$$
この主張は $M$ が有限局所自由、すなわち (7) が成り立つことを示す。
$\Psi$ は全射なので、$\Psi$ が単射であることを示せば十分である。
$R_{fg}$ は被約なので、補題 \ref{algebra-lemma-reduced-ring-sub-product-fields} により、
$R_{fg}$ のすべての極小素イデアル $\mathfrak p$ で局所化した後に
$\Psi$ が単射であることを示せば十分である。
しかし補題 \ref{algebra-lemma-minimal-prime-reduced-ring} により
$R_\mathfrak p = \kappa(\mathfrak p)$ であり、
$\rho_M(\mathfrak p) = r$ であることを知っている。したがって
$\Psi_\mathfrak p : R_\mathfrak p^{\oplus r} \to
M \otimes_R \kappa(\mathfrak p)$ は、同じ次元の有限次元ベクトル空間の間の
全射として同型である。
\end{proof}

\begin{remark}
\label{algebra-remark-warning}
有限 $R$-加群が $R$ 上平坦であっても、自動的に射影的になるとは限らない。
反例として、$R = \mathcal{C}^\infty(\mathbf{R})$ を
$\mathbf{R}$ 上の無限回微分可能関数の環とし、
$M = R_{\mathfrak m} = R/I$ とするものがある。ここで
$\mathfrak m = \{f \in R \mid f(0) = 0\}$ および
$I = \{f \in R \mid \exists \epsilon, \epsilon > 0 :
f(x) = 0\ \forall x, |x| < \epsilon\}$ である。
\end{remark}

\begin{lemma}
\label{algebra-lemma-finite-flat-local}
（注意：注記 \ref{algebra-remark-warning} を参照。）
$R$ を局所環、$M$ を有限平坦 $R$-加群とする。このとき $M$ は有限自由である。
\end{lemma}

\begin{proof}
平坦性の方程式的判定法、補題 \ref{algebra-lemma-flat-eq} から従う。
実際、$x_1, \ldots, x_r \in M$ が $M/\mathfrak mM$ の基底に写るとする。
Nakayama の補題 \ref{algebra-lemma-NAK} により、これらの元は $M$ を生成する。
$x_i$ の間に関係がないことを示したい。その代わり、
$x_1, \ldots, x_n \in M$ がベクトル空間 $M/\mathfrak mM$ で線形独立ならば
$R$ 上でも独立であることを、$n$ に関する帰納法で示す。

\medskip\noindent
帰納法の基底は、$x \in M$、$x \not\in \mathfrak mM$ で関係
$fx = 0$ がある場合である。方程式的判定法により、
$x = \sum a_j y_j$ かつすべての $j$ について $fa_j = 0$ となる
$y_j \in M$ と $a_j \in R$ が存在する。
$x \not\in \mathfrak mM$ なので少なくとも一つの $a_j$ は単元であり、
したがって $f = 0 $ である。

\medskip\noindent
$\sum f_i x_i$ を $x_1, \ldots, x_n$ の間の関係とする。
$x_i$ の選び方により $f_i \in \mathfrak m$ である。
平坦性の方程式的判定法により、
$x_i = \sum a_{ij} y_j$ かつ $\sum f_i a_{ij} = 0$ となる
$a_{ij} \in R$ と $y_j \in M$ が存在する。
$x_n \not \in \mathfrak mM$ なので、少なくとも一つの $j$ について
$a_{nj}\not\in \mathfrak m$ である。
$\sum f_i a_{ij} = 0$ であるから
$f_n = \sum_{i = 1}^{n-1} (-a_{ij}/a_{nj}) f_i$ を得る。
関係 $\sum f_i x_i = 0$ は
$\sum_{i = 1}^{n-1} f_i( x_i + (-a_{ij}/a_{nj}) x_n) = 0$ と書き直せる。
元 $x_i + (-a_{ij}/a_{nj}) x_n$ は $M/\mathfrak mM$ の
$n-1$ 個の線形独立な元に写ることに注意する。
帰納法の仮定により、すべての $f_i$、$i \leq n-1$ は零でなければならず、
$f_n = \sum_{i = 1}^{n-1} (-a_{ij}/a_{nj}) f_i$ も零である。
これで帰納段階が示された。
\end{proof}

\begin{lemma}
\label{algebra-lemma-finite-projective-descends}
$R \to S$ を局所環の間の平坦局所準同型とする。
$M$ を有限 $R$-加群とする。このとき $M$ が $R$ 上有限射影的であることと、
$M \otimes_R S$ が $S$ 上有限射影的であることは同値である。
\end{lemma}

\begin{proof}
補題 \ref{algebra-lemma-finite-projective} により、局所環上で有限射影的であることは
有限自由であることと同じである。
$M \otimes_R S$ が有限自由 $S$-加群であると仮定する。
$M/\mathfrak m_RM$ における像が $\kappa(\mathfrak m)$ 上の基底をなす
$x_1, \ldots, x_r \in M$ を選ぶ。
% [P01-E0444] The frozen maximal-ideal subscript mismatch is retained.
すると $x_1 \otimes 1, \ldots, x_r \otimes 1$ は $M \otimes_R S$ の基底である。
これは写像 $R^{\oplus r} \to M, (a_i) \mapsto \sum a_i x_i$ が、
$S$ とテンソル積をとった後に同型になることを含意する。
$R \to S$ の忠実平坦性により（補題 \ref{algebra-lemma-local-flat-ff} を参照）、
この写像は同型である。
\end{proof}

\begin{lemma}
\label{algebra-lemma-locally-free-semi-local-free}
$R$ を半局所環とする。$M$ を有限局所自由加群とする。
$M$ の階数が一定ならば、$M$ は自由である。特に、$R$ のスペクトルが
連結ならば $M$ は自由である。
\end{lemma}

\begin{proof}
省略する。ヒント：まず、$R$ のすべての極大イデアル
$\mathfrak m_1, \ldots, \mathfrak m_n$ に対して、階数が一定であることを用い、
$M/\mathfrak m_iM$ が同じ次元 $d$ をもつことを示す。
次に中国剰余定理により、各 $M/\mathfrak m_iM$ の基底をなす元
$x_1, \ldots, x_d \in M$ が存在することを示す。
最後に $x_1, \ldots, x_d$ が $M$ の基底であることを示す。
\end{proof}

\noindent
次はグルーポイドの章で用いる技術的補題である。

\begin{lemma}
\label{algebra-lemma-semi-local-module-basis-in-submodule}
$R$ を極大イデアル $\mathfrak m$ と無限剰余体をもつ局所環とする。
$R \to S$ を環準同型とする。
$M$ を $S$-加群、$N \subset M$ を $R$-部分加群とする。次を仮定する。
\begin{enumerate}
\item $S$ は半局所で、$\mathfrak mS$ は $S$ の Jacobson 根基に含まれる。
\item $M$ は有限自由 $S$-加群である。
\item $N$ は $S$-加群として $M$ を生成する。
\end{enumerate}
このとき $N$ は $M$ の $S$-基底を含む。
\end{lemma}

\begin{proof}
$M$ が階数 $n$ の自由加群であると仮定する。
$I \subset S$ を Jacobson 根基とする。
Nakayama の補題 \ref{algebra-lemma-NAK} により、元の列
$m_1, \ldots, m_n$ が $M$ の基底であることと、
$\overline{m}_i \in M/IM$ が $M/IM$ を生成することは同値である。
したがって $M$ を $M/IM$ で、$N$ を $N/(N \cap IM)$ で、
$R$ を $R/\mathfrak m$ で、$S$ を $S/IS$ で置き換えてよい。
この場合、$S$ は体の有限直積
$S = k_1 \times \ldots \times k_r$ であり、
$M = k_1^{\oplus n} \times \ldots \times k_r^{\oplus n}$ である。
$N \subset M$ が $S$-加群として $M$ を生成するという事実は、
各因子 $k_i^{\oplus n}$ に非零成分をもつ
$\sum a_j x_j$、ただし $a_j \in S$、となるような $x_j \in N$ が存在することを意味する。
$R = k$ は無限体なので、各因子に非零成分をもつ
$y = \sum c_j x_j$、ただし $c_j \in k$、という線形結合も存在する。
したがって $y \in N$ は自由直和因子 $Sy \subset M$ を生成する。
$n$ に関する帰納法により、結果は $M/Sy$ と部分加群
$\overline{N} = N/(N \cap Sy)$ に対して成り立つ。
言い換えれば、$M/Sy$ を（自由に）生成する
$\overline{y}_2, \ldots, \overline{y}_n$ が $\overline{N}$ に存在する。
すると $y, y_2, \ldots, y_n$ は $M$ を（自由に）生成し、結論を得る。
\end{proof}

\begin{lemma}
\label{algebra-lemma-evaluation-map-iso-finite-projective}
$R$ を環とする。$L$, $M$, $N$ を $R$-加群とする。標準写像
$$
\Hom_R(M, N) \otimes_R L \to \Hom_R(M, N \otimes_R L)
$$
は、$M$ が有限射影的ならば同型である。
\end{lemma}

\begin{proof}
補題 \ref{algebra-lemma-finite-projective} により、$M$ は有限表示であると同時に
有限局所自由である。補題 \ref{algebra-lemma-hom-from-finitely-presented} および
\ref{algebra-lemma-tensor-product-localization} により、矢印の左辺と右辺を作る操作は
局所化と可換である。補題 \ref{algebra-lemma-cover} により、この写像が同型であることは
局所化の後に確認してよい。したがって $M$ が有限自由であると仮定してよい。
この場合、補題は直ちに従う。
\end{proof}

\section{加群写像が定める開集合}
\label{algebra-section-loci-maps}

\noindent
与えられた加群写像が全射または同型となる素イデアルの集合は、
開集合となることがある。有限射影加群の場合には写像の階数を調べられる。

\begin{lemma}
\label{algebra-lemma-map-between-finite}
$R$ を環とする。$\varphi : M \to N$ を $R$-加群の写像とし、
$N$ は有限 $R$-加群であるとする。このとき等式
\begin{align*}
U & = \{\mathfrak p \subset R \mid
\varphi_{\mathfrak p} : M_{\mathfrak p} \to N_{\mathfrak p}
\text{ is surjective}\} \\
& = \{\mathfrak p \subset R \mid
\varphi \otimes \kappa(\mathfrak p) :
M \otimes \kappa(\mathfrak p) \to N \otimes \kappa(\mathfrak p)
\text{ is surjective}\}
\end{align*}
が成り立ち、$U$ は $\Spec(R)$ の開部分集合である。さらに、
$D(f) \subset U$ となる任意の $f \in R$ に対して、写像
$M_f \to N_f$ は全射である。
\end{lemma}

\begin{proof}
表示式の等号は Nakayama の補題から従う。
Nakayama の補題は $U$ が開であることも含意する。
補題 \ref{algebra-lemma-NAK}、特に (3) を参照せよ。
$D(f) \subset U$ ならば、$M_f \to N_f$ は $R_f$ の素イデアルにおける
すべての局所化で全射であり、したがって補題
\ref{algebra-lemma-characterize-zero-local} により全射である。
\end{proof}

\begin{lemma}
\label{algebra-lemma-map-between-finitely-presented}
$R$ を環とする。$\varphi : M \to N$ を $R$-加群の写像とし、
$M$ は有限、$N$ は有限表示であるとする。このとき
$$
U = \{\mathfrak p \subset R \mid
\varphi_{\mathfrak p} : M_{\mathfrak p} \to N_{\mathfrak p}
\text{ is an isomorphism}\}
$$
は $\Spec(R)$ の開部分集合である。
\end{lemma}

\begin{proof}
$\mathfrak p \in U$ とし、表示
$N = R^{\oplus n}/\sum_{j = 1, \ldots, m} R k_j$ を選ぶ。
$R^{\oplus n}$ の第 $i$ 基底ベクトルの $N$ における像を $e_i$ と書く。
各 $i \in \{1, \ldots, n\}$ に対して、ある $f_i \in R$、
$f_i \not \in \mathfrak p$ について $\varphi(m_i) = f_i e_i$ となる元
$m_i \in M_{\mathfrak p}$ を選ぶ。$\varphi_{\mathfrak p}$ は同型なので、
これは可能である。$f = f_1 \ldots f_n$ とおき、
$\psi : R_f^{\oplus n} \to M_f$ を、第 $i$ 基底ベクトルを
$m_i/f_i$ に写す写像とする。
% [P01-E0445] Frozen preimage and denominator declarations retained.
$\varphi_f \circ \psi$ は、与えられた写像 $R^{\oplus n} \to N$ の
$f$ における局所化であることに注意する。
$\varphi_{\mathfrak p}$ は同型なので、$\psi(k_j)$ は $M$ の元であり、
$M_{\mathfrak p}$ で零に写る。したがって、
$g_j \psi(k_j) = 0$ となる $g_j \in R$、$g_j \not \in \mathfrak p$ が存在する。
% [P01-E0446] Frozen codomain assertion retained; see the part-local errata ledger.
$g = g_1 \ldots g_m$ とおくと、$\psi_g$ は $N_{fg}$ を経由して
写像 $\chi : N_{fg} \to M_{fg}$ を与える。
構成により $\chi$ は $\varphi_{fg}$ の右逆である。
したがって $\chi_\mathfrak p$ は同型である。
補題 \ref{algebra-lemma-map-between-finite} により、ある $h \in R$、
$h \not \in \mathfrak p$ に対して
$\chi_h : N_{fgh} \to M_{fgh}$ は全射である。
ゆえに $\varphi_{fgh}$ と $\chi_h$ は互いに逆写像であり、所望の
$D(fgh) \subset U$ が従う。
\end{proof}

\begin{lemma}
\label{algebra-lemma-finitely-presented-localization-free}
$R$ を環、$\mathfrak p \subset R$ を素イデアルとする。
$M$ を有限表示 $R$-加群とする。$M_\mathfrak p$ が自由ならば、
$M_f$ が自由 $R_f$-加群となるような
$f \in R$、$f \not \in \mathfrak p$ が存在する。
\end{lemma}

\begin{proof}
基底 $x_1, \ldots, x_n \in M_\mathfrak p$ を選ぶ。
各 $x_i$ がある $y_i \in M_f$ の像となるような
$f \in R$、$f \not \in \mathfrak p$ を選べる。
$m \gg 0$ に対して $y_i$ を $f^m y_i$ で置き換えれば、
$y_i \in M$ と仮定してよい。実際、これは $x_1, \ldots, x_n$ を
$f^mx_1, \ldots, f^mx_n$ で置き換えるが、$f$ は $R_\mathfrak p$ で単元に
写るので、これも基底である。したがって、$\mathfrak p$ における局所化が
同型となる $R$-加群準同型
$\varphi = (y_1, \ldots, y_n) : R^{\oplus n} \to M$ を得る。
補題 \ref{algebra-lemma-map-between-finitely-presented} により、すべての素イデアル
$\mathfrak q \subset R$、$f \not \in \mathfrak q$ に対して
$\varphi_\mathfrak q$ が同型となるような
$f \in R$、$f \not \in \mathfrak p$ を見つけられる。
すると補題 \ref{algebra-lemma-characterize-zero-local} から $\varphi_f$ は同型となり、
証明が完了する。
\end{proof}

\begin{lemma}
\label{algebra-lemma-cokernel-flat}
$R$ を環とする。$\varphi : P_1 \to P_2$ を有限射影加群の間の写像とする。
このとき次が成り立つ。
\begin{enumerate}
\item $\varphi \otimes \kappa(\mathfrak p)$ が単射となる素イデアル
$\mathfrak p \in \Spec(R)$ の集合 $U$ は開であり、
$D(f) \subset U$ となる任意の $f\in R$ に対して次が成り立つ。
\begin{enumerate}
\item $P_{1, f} \to P_{2, f}$ は単射である。
\item 加群 $\Coker(\varphi)_f$ は $R_f$ 上有限射影的である。
\end{enumerate}
\item $\varphi \otimes \kappa(\mathfrak p)$ が全射となる素イデアル
$\mathfrak p \in \Spec(R)$ の集合 $W$ は開であり、
$D(f) \subset W$ となる任意の $f\in R$ に対して次が成り立つ。
\begin{enumerate}
\item $P_{1, f} \to P_{2, f}$ は全射である。
\item 加群 $\Ker(\varphi)_f$ は $R_f$ 上有限射影的である。
\end{enumerate}
\item $\varphi \otimes \kappa(\mathfrak p)$ が同型となる素イデアル
$\mathfrak p \in \Spec(R)$ の集合 $V$ は開であり、
$D(f) \subset V$ となる任意の $f\in R$ に対して写像
$\varphi : P_{1, f} \to P_{2, f}$ は $R_f$ 上の加群の同型である。
\end{enumerate}
\end{lemma}

\begin{proof}
$U$ が開であることを示すには $\Spec(R)$ 上局所的に議論してよい。
したがって補題 \ref{algebra-lemma-finite-projective} により、$R$ を適当な局所化で
置き換えて $P_1 = R^{n_1}$ および $P_2 = R^{n_2}$ と仮定してよい。
この場合、$\varphi \otimes \kappa(\mathfrak p)$ が単射であることは、
$n_1 \leq n_2$ であり、かつ $\varphi$ の行列のある
$n_1 \times n_1$ 小行列式 $f$ が $\kappa(\mathfrak p)$ で可逆であることと同値である。
したがって $D(f) \subset U$ である。この議論は、
$\mathfrak p \in U$ に対して
$P_{1, \mathfrak p} \to P_{2, \mathfrak p}$ が単射であることも示す。

\medskip\noindent
$D(f) \subset U$ と仮定する。前段落の注意と補題
\ref{algebra-lemma-characterize-zero-local} により
$P_{1, f} \to P_{2, f}$ は単射、すなわち (1)(a) が成り立つ。
補題 \ref{algebra-lemma-finite-projective} により (1)(b) を示すには、
$\Coker(\varphi)$ が $D(f)$ 上局所的に有限射影的であることを示せば十分である。
したがって上で見たように、$P_1 = R^{n_1}$、$P_2 = R^{n_2}$ で、
$\varphi$ の行列のある小行列式が $R$ で可逆であると仮定してよい。
問題の小行列式が $R^{n_2}$ の最初の $n_1$ 個の基底ベクトルに対応するなら、
最後の $n_2 - n_1$ 個の基底ベクトルを用いて写像
$R^{n_2 - n_1} \to R^{n_2} \to \Coker(\varphi)$ を得るが、
これは容易に同型だと分かる。

\medskip\noindent
$W$ の開性と $D(f) \subset W$ に対する (2)(a) は
補題 \ref{algebra-lemma-map-between-finite} から従う。
$P_{2, f}$ は $R_f$ 上射影的なので、
$\varphi_f : P_{1, f} \to P_{2, f}$ は切断をもち、したがって
$\Ker(\varphi)_f$ は $P_{2, f}$ の直和因子である。
% [P01-E0447] Frozen direct-summand ambient module retained.
ゆえに $\Ker(\varphi)_f$ は有限射影的である。したがって (2)(b) も成り立つ。

\medskip\noindent
$V = U \cap W$ が開であることは明らかであり、(3) のもう一つの主張は
(1)(a) と (2)(a) から従う。
\end{proof}

\section{加群の射影性の忠実平坦降下}
\label{algebra-section-ffdescent-projectivity-introduction}

\medskip\noindent
次のいくつかの節では Raynaud と Gruson \cite{GruRay} に従い、
加群の射影性が忠実平坦な環準同型に沿って降下することを証明する。
証明の考え方は、Kaplansky \cite{Kaplansky} 流のデヴィサージュを用いて
可算生成加群の場合に帰着することである。加群 $M$ の適切なフィルトレーションが
与えられたとき、デヴィサージュにより $M$ をフィルトレーション部分加群の
逐次商の直和として表せる（節 \ref{algebra-section-transfinite-devissage} を参照）。
この手法を用いて、射影加群が可算生成加群の直和であることを証明する
（定理 \ref{algebra-theorem-projective-direct-sum}）。可算生成加群について射影性の降下を
証明するため、加群に対する ``Mittag-Leffler'' 条件を導入し、
可算生成加群が射影的であることと、それが平坦かつ Mittag-Leffler であることが
同値であると証明し（定理 \ref{algebra-theorem-projectivity-characterization}）、さらに
Mittag-Leffler 加群であるという性質が降下することを示す
（補題 \ref{algebra-lemma-ffdescent-ML}）。最後に、忠実平坦な環準同型による基底変換が
射影的になる任意の加群 $M$ に対し、逐次商が可算生成射影加群となる部分加群で
$M$ をフィルター付けする。するとデヴィサージュにより $M$ は射影加群の直和、
したがってそれ自身射影的であると結論する
（定理 \ref{algebra-theorem-ffdescent-projectivity}）。

\medskip\noindent
\cite{GruRay} における射影性の忠実平坦降下の証明には誤りがあることを注意する。
そこでは、忠実平坦な環準同型に沿う射影性の降下を、より一般的な種類の
環準同型に沿う射影性の降下から導いている
（\cite[Part II の Example 3.1.4(1)]{GruRay}）。
しかし、このより一般的な種類の写像に沿う降下の証明は正しくない。
\cite{G} で Gruson は何が誤っていたかを説明しているが、ここで問題となる場合の
修正は与えていない。射影性の忠実平坦降下の証明におけるこの穴を埋めることは、
Mittag-Leffler 加群であるという性質が忠実平坦な環準同型に沿って降下することを
証明する問題に帰着する。これは補題 \ref{algebra-lemma-ffdescent-ML} で行う。

\section{平坦性の特徴付け}
\label{algebra-section-characterize-flatness}

\noindent
この節では平坦性の判定法を論じる。この節の主結果は、平坦加群が有限自由加群の
有向系の余極限であると述べる Lazard の定理（下の定理 \ref{algebra-theorem-lazard}）である。
読者に「平坦性の方程式的判定法」、補題 \ref{algebra-lemma-flat-eq} を思い出してもらいたい。
これは一見はるかに強い性質へと変形できることが分かる。

\begin{lemma}
\label{algebra-lemma-flat-factors-free}
$M$ を $R$-加群とする。次は同値である。
\begin{enumerate}
\item $M$ は平坦である。
\item $f: R^n \to M$ が加群写像で $x \in \Ker(f)$ ならば、
$f = g \circ h$ かつ $x \in \Ker(h)$ となる加群写像
$h: R^n \to R^m$ と $g: R^m \to M$ が存在する。
\item $f: R^n \to M$ を加群写像、$N \subset \Ker(f)$ を任意の部分加群、
$h: R^n \to R^{m}$ を $N \subset \Ker(h)$ かつ $f$ が $h$ を経由する写像とする。
このとき任意の $x \in \Ker(f)$ に対して、$N + Rx \subset \Ker(h')$ かつ
$f$ が $h'$ を経由するような写像 $h': R^n \to R^{m'}$ を見つけられる。
\item $f: R^n \to M$ が加群写像で $N \subset \Ker(f)$ が有限生成部分加群ならば、
$f = g \circ h$ かつ $N \subset \Ker(h)$ となる加群写像
$h: R^n \to R^m$ と $g: R^m \to M$ が存在する。
\end{enumerate}
\end{lemma}

\begin{proof}
(1) と (2) の同値性は、平坦性の方程式的判定法の言い換えにすぎない
\footnote{実際、加群写像 $f : R^n \to M$ は $M$ の元
$x_1, x_2, \ldots, x_n$ の選択（すなわち標準基底の元
$e_1, e_2, \ldots, e_n$ の像）に対応する。さらに、元
$x \in \Ker(f)$ はこれらの $x_1, x_2, \ldots, x_n$ の間の関係
（すなわち、$f_i$ を $x$ の座標とした関係 $\sum_i f_i x_i = 0$）に対応する。
加群写像 $h$（$m \times n$ 行列として表される）は補題
\ref{algebra-lemma-flat-eq} の行列 $(a_{ij})$ に対応し、補題
\ref{algebra-lemma-flat-eq} の $y_j$ は $g$ による $R^m$ の標準基底ベクトルの像である。}。
(2) が (3) を含意することを示すため、$f$ が $f = g \circ h$ と分解されるような写像
$g: R^m \to M$ をとる。(2) により、$h''$ が $h(x)$ を零に写し、
$g: R^m \to M$ が $h''$ を経由するような $h'': R^m \to R^{m'}$ を見つける。
このとき $h' = h'' \circ h$ とすればよい。
(3) は、(4) の $N \subset \Ker(f)$ の生成元の個数に関する帰納法により (4) を含意する。
(4) が (2) を含意することは明らかである。
\end{proof}

\begin{lemma}
\label{algebra-lemma-flat-factors-fp}
$M$ を $R$-加群とする。このとき $M$ が平坦であることと次の条件は同値である。
$P$ が有限表示 $R$-加群で $f: P \to M$ が加群写像ならば、
有限自由 $R$-加群 $F$ と加群写像 $h: P \to F$、$g: F \to M$ で
$f = g \circ h$ となるものが存在する。
\end{lemma}

\begin{proof}
これは補題 \ref{algebra-lemma-flat-factors-free} の条件 (4) の言い換えにすぎない。
\end{proof}

\begin{lemma}
\label{algebra-lemma-flat-surjective-hom}
$M$ を $R$-加群とする。このとき $M$ が平坦であることと次の条件は同値である。
すべての有限表示 $R$-加群 $P$ に対し、$N \to M$ が全射 $R$-加群写像ならば、
誘導写像 $\Hom_R(P, N) \to \Hom_R(P, M)$ は全射である。
\end{lemma}

\begin{proof}
まず $M$ が平坦であると仮定する。$P$ が有限表示で写像 $f: P \to M$ が
与えられたとき、それが写像 $N \to M$ を経由することを示さなければならない。
補題 \ref{algebra-lemma-flat-factors-fp} により、写像 $f$ は、$F$ が有限自由であるような
写像 $F \to M$ を経由する。$F$ は自由なので、この写像は $N \to M$ を経由する。
したがって $f$ は $N \to M$ を経由する。

\medskip\noindent
逆に、補題の条件を仮定する。有限表示加群 $P$ からの写像
$f: P \to M$ をとる。$M$ への全射 $N \to M$ をもつ自由加群 $N$ を選ぶ。
すると $f$ は $N \to M$ を経由し、$P$ は有限生成なので、$f$ は $N$ の
有限自由部分加群を経由する。したがって $M$ は補題
\ref{algebra-lemma-flat-factors-fp} の条件を満たし、ゆえに平坦である。
\end{proof}

\begin{theorem}[Lazard の定理]
\label{algebra-theorem-lazard}
$M$ を $R$-加群とする。このとき $M$ が平坦であることと、有限自由
$R$-加群の有向系の余極限であることは同値である。
\end{theorem}

\begin{proof}
有向余極限をとる操作は完全でテンソル積と可換なので、平坦加群の有向系の
余極限は平坦である。したがって $M$ が有限自由加群の有向系の余極限ならば、
$M$ は平坦である。

\medskip\noindent
逆を示すため、まず任意の加群 $M$ が次のように有限表示加群の有向系の余極限として
書けることを思い出す。ある集合 $I$ に対して全射 $f: R^I \to M$ を選び、
その核を $K$ とする。$E$ を順序対 $(J, N)$ の集合とする。ここで $J$ は $I$ の
有限部分集合、$N$ は $R^J \cap K$ の有限生成部分加群である。
$J \subset J'$ かつ $N \subset N'$ であることと $(J, N) \leq (J', N')$ が
同値となるように定めると、$E$ は有向半順序集合となる。
$e = (J, N)$ に対して $M_e = R^J/N$ と定め、$e \leq e'$ に対して
$f_{ee'}: M_e \to M_{e'}$ を自然な写像とする。
すると $(M_e, f_{ee'})$ は有向系であり、自然な写像 $f_e: M_e \to M$ は同型
$\colim_{e
\in E} M_e \xrightarrow{\cong} M$ を誘導する。

\medskip\noindent
ここで $M$ が平坦であると仮定する。$I = M \times \mathbf{Z}$ とし、
$R^{I}$ の標準基底を $(x_i)$ と書く。上の議論における $f: R^I \to M$ として、
$x_i$ を $i$ の $M$ への射影に写す写像をとる。
定理を証明するには、$M_e$ が自由となる $e \in E$ が $E$ の共終部分集合を
なすことを示せば十分である。任意の $e = (J, N) \in E$ をとる。
補題 \ref{algebra-lemma-flat-factors-fp} により、有限自由加群 $F$ と写像
$h: R^J/N \to F$、$g: F \to M$ が存在し、自然な写像
$f_e: R^J/N \to M$ は $R^J/N \xrightarrow{h} F
\xrightarrow{g} M$ と分解する。ある $e' \geq e$ に対して $F$ を $M_{e'}$ として
実現する。

\medskip\noindent
$\{ b_1, \ldots, b_n \}$ を $F$ の有限基底とする。
すべての $\ell$ に対して $i_{\ell} \notin J$ であり、かつ
$f: R^I \to M$ による $x_{i_{\ell}}$ の像が
$g: F \to M$ による $b_{\ell}$ の像に等しくなるような、相異なる
$n$ 個の元 $i_1, \ldots, i_n \in I$ を選ぶ。
$I$ の選び方により $M$ のすべての元は、相異なる無限個の $i \in I$ に対して
$f(x_i)$ と書けるので、これは可能である。
$J' = J \cup \{i_1, \ldots , i_n \}$ とおき、
$i \in J$ に対して $x_i \mapsto h(x_i)$、$\ell = 1, \ldots, n$ に対して
$x_{i_{\ell}} \mapsto b_{\ell}$ として $R^{J'} \to F$ を定める。
$N' = \Ker(R^{J'} \to F)$ とおく。次に注意する。
\begin{enumerate}
\item 図式
$$
\xymatrix{
R^{J'} \ar[r] \ar@{^{(}->}[d] & F \ar[d]^{g} \\
R^{I} \ar[r]_{f} & M
}
$$
は可換である。
%$R^{J'} \to F$ factors $f: R^I \to M$,
したがって $N' \subset K = \Ker(f)$ である。
\item $R^{J'} \to F$ は有限自由加群への全射なので分裂し、したがって
$N'$ は有限生成である。
\item $J \subset J'$ かつ $N \subset N'$ である。
\end{enumerate}
(1) と (2) により $e' = (J', N')$ は $E$ に属し、(3) により
$e' \geq e$ であり、構成により $M_{e'} = R^{J'}/N' \cong F$ は自由である。
\end{proof}

\section{純単射な加群写像}
\label{algebra-section-universally-injective}

\noindent
次に、ある意味で平坦加群と相補的な純単射加群写像を論じる
（補題 \ref{algebra-lemma-flat-universally-injective} を参照）。
Lazard の学位論文 \cite{Autour} に従う。\cite{Lam} も参照せよ。

\begin{definition}
\label{algebra-definition-universally-injective}
$f: M \to N$ を $R$-加群の写像とする。すべての $R$-加群 $Q$ に対して写像
$f \otimes_R \text{id}_Q: M \otimes_R Q \to N \otimes_R Q$ が単射となるとき、
$f$ は {\it 純単射} であるという。$R$-加群の列
$0 \to M_1 \to M_2 \to M_3 \to 0$ が完全であり、
$M_1 \to M_2$ が純単射であるとき、この列は {\it 純完全} であるという。
\end{definition}

\begin{example}
\label{algebra-example-universally-exact}
純完全列の例を挙げる。
\begin{enumerate}
\item テンソル積は直和をとる操作と可換なので、分裂短完全列は純完全である。
\item 純完全列の有向系の余極限は純完全である。これは、有向余極限をとる操作が
完全であり、テンソル積が余極限をとる操作と可換であることから従う。
特に分裂完全列の有向系の余極限は純完全である。逆に、任意の純完全列が
このようにして得られることを後で見る。
\end{enumerate}
\end{example}

\noindent
次に、短完全列が純完全となるための判定条件を列挙する。
これらは上で与えた平坦性の判定条件の類似物である。以下の (3)--(6) はそれぞれ、
補題 \ref{algebra-lemma-flat-eq}、\ref{algebra-lemma-flat-factors-free}、
\ref{algebra-lemma-flat-surjective-hom} および定理 \ref{algebra-theorem-lazard} で与えた
平坦性の判定条件に対応する。

\begin{theorem}
\label{algebra-theorem-universally-exact-criteria}
完全列
$$
0 \to M_1 \xrightarrow{f_1} M_2 \xrightarrow{f_2} M_3 \to 0
$$
を $R$-加群の完全列とする。次は同値である。
\begin{enumerate}
\item 列 $0 \to M_1 \to M_2 \to M_3 \to 0$ は純完全である。
\item すべての有限表示 $R$-加群 $Q$ に対して列
$$
0 \to M_1 \otimes_R Q \to M_2 \otimes_R Q \to
M_3 \otimes_R Q \to 0
$$
は完全である。
\item すべての $i$ について
$$
f_1(x_i) = \sum\nolimits_j a_{ij} y_j,
$$
を満たす元 $x_i \in M_1$ $(i = 1, \ldots, n)$、
$y_j \in M_2$ $(j = 1, \ldots, m)$ および
$a_{ij} \in R$ $(i = 1, \ldots, n, j = 1, \ldots, m)$ が与えられたとき、
すべての $i$ について
$$
x_i = \sum\nolimits_j a_{ij} z_j .
$$
を満たす $z_j \in M_1$ $(j =1, \ldots, m)$ が存在する。
\item $R$-加群写像の可換図式
$$
\xymatrix{
R^n \ar[r] \ar[d] &  R^m \ar[d] \\
M_1 \ar[r]^{f_1}        &  M_2
}
$$
が与えられ、$m$ と $n$ が整数であるとき、上の三角形を可換にする写像
$R^m \to M_1$ が存在する。
\item すべての有限表示 $R$-加群 $P$ に対して、$R$-加群写像
$\Hom_R(P, M_2) \to \Hom_R(P, M_3)$ は全射である。
\item 列 $0 \to M_1 \to M_2 \to M_3 \to 0$ は、
$M_{3, i}$ が有限表示となる形
$$
0 \to M_{1} \to M_{2, i} \to M_{3, i} \to 0
$$
の分裂完全列の有向系の余極限である。
\end{enumerate}
\end{theorem}

\begin{proof}
(1) が (2) を含意することは明らかである。

\medskip\noindent
次に (2) が (3) を含意することを示す。
(3) のような関係 $f_1(x_i) = \sum_j a_{ij} y_j$ をとる。
$R^m$ の基底を $(d_j)$、$R^n$ の基底を $(e_i)$ とし、
$R^m \to R^n$ を $d_j \mapsto \sum_i a_{ij} e_i$ で与えられる写像とする。
$Q$ を $R^m \to R^n$ の余核とする。写像 $f_1: M_1 \to M_2$ により
$R^m \to R^n \to Q \to 0$ とテンソル積をとると、可換図式
$$
\xymatrix{
M_1^{\oplus m} \ar[r] \ar[d] & M_1^{\oplus n} \ar[r] \ar[d] & M_1 \otimes_R Q
\ar[r] \ar[d] & 0 \\
M_2^{\oplus m} \ar[r] & M_2^{\oplus n} \ar[r] & M_2 \otimes_R Q \ar[r] & 0
}
$$
を得る。ここで $M_1^{\oplus m} \to M_1^{\oplus n}$ は
$$
(z_1, \ldots, z_m) \mapsto
(\sum\nolimits_j a_{1j} z_j, \ldots, \sum\nolimits_j a_{nj} z_j),
$$
で与えられ、$M_2^{\oplus m} \to M_2^{\oplus n}$ も同様に与えられる。
$x = (x_1, \ldots, x_n) \in M_1^{\oplus n}$ が
$M_1^{\oplus m} \to M_1^{\oplus n}$ の像に属することを示したい。
(2) により写像 $M_1 \otimes Q \to M_2 \otimes Q$ は単射なので、
上段の完全性から、$x$ が $M_2 \otimes Q$ で $0$ に写ることを示せば十分である。
さらに下段の完全性から、$M_2^{\oplus n}$ における $x$ の像が
$M_2^{\oplus m} \to M_2^{\oplus n}$ の像に属することを示せば十分である。
これは仮定により成り立つ。

\medskip\noindent
条件 (4) は (3) を図式の形に翻訳したものにすぎない。

\medskip\noindent
次に (4) が (5) を含意することを示す。
$\varphi : P \to M_3$ を有限表示 $R$-加群 $P$ からの写像とする。
$\varphi$ が写像 $P \to M_2$ に持ち上がることを示さなければならない。
$P$ の表示
$$
R^n \xrightarrow{g_1} R^m \xrightarrow{g_2} P \to 0.
$$
を選ぶ。$R^n$ と $R^m$ の自由性を用いると、次の図式が可換となるように
$h_2: R^m \to M_2$、次いで $h_1: R^n \to M_1$ を構成できる。
$$
\xymatrix{
         & R^n \ar[r]^{g_1} \ar[d]^{h_1} & R^m \ar[r]^{g_2} \ar[d]^{h_2} & P
\ar[r] \ar[d]^{\varphi} & 0 \\
0 \ar[r] & M_1 \ar[r]^{f_1} & M_2 \ar[r]^{f_2} & M_3 \ar[r] & 0 .
}
$$
(4) により $k_1 \circ g_1 = h_1$ を満たす写像 $k_1: R^m \to M_1$ が存在する。
ここで $h_2' = h_2 - f_1 \circ k_1$ により $h'_2: R^m \to M_2$ を定める。
すると
$$
h'_2 \circ g_1 = h_2 \circ g_1 - f_1 \circ k_1 \circ g_1 =
h_2 \circ g_1 - f_1 \circ h_1 = 0 .
$$
したがって商をとることにより、$h'_2$ は
$\varphi' \circ g_2 = h_2'$ を満たす写像 $\varphi': P \to M_2$ を定める。
図式で書けば
$$
\xymatrix{
R^m \ar[r]^{g_2} \ar[d]_{h'_2} & P \ar[d]^{\varphi} \ar[dl]_{\varphi'} \\
M_2 \ar[r]^{f_2} & M_3.
}
$$
となり、上の三角形は可換である。$\varphi'$ が所望の持ち上げ、すなわち
$f_2 \circ \varphi' = \varphi$ であると主張する。定義から
$$
f_2 \circ \varphi' \circ g_2 = f_2 \circ h'_2 =
f_2 \circ h_2 - f_2 \circ f_1 \circ k_1 = f_2 \circ h_2 =
\varphi \circ g_2.
$$
を得る。$g_2$ は全射なので、これで証明が完了する。

\medskip\noindent
ここで (5) が (6) を含意することを示す。
補題 \ref{algebra-lemma-module-colimit-fp} により、$M_{3}$ を有限表示加群
$M_{3, i}$ の有向系の余極限として書く。$M_{2, i}$ を $M_{3}$ 上の
$M_{3, i}$ と $M_{2}$ のファイバー積とする。定義により、これは
$M_2 \times M_{3, i}$ のうち二つの $M_3$ への射影が等しい元からなる部分加群である。
$M_{1, i}$ を射影 $M_{2, i} \to M_{3, i}$ の核とする。
すると完全列の有向系
$$
0 \to M_{1, i} \to M_{2, i} \to M_{3, i} \to 0,
$$
と、各 $i$ に対して有向系と両立する完全列の写像
$$
\xymatrix{
0  \ar[r] & M_{1, i} \ar[d] \ar[r]  &  M_{2, i}  \ar[r] \ar[d] & M_{3, i} \ar[d]
\ar[r] & 0 \\
0  \ar[r] & M_{1} \ar[r]  &  M_{2}  \ar[r] & M_{3} \ar[r] & 0
}
$$
を得る。ファイバー積 $M_{2, i}$ の定義から、写像
$M_{1, i} \to M_1$ は同型である。(5) により、$M_{3, i} \to M_3$ を持ち上げる
写像 $M_{3, i} \to M_{2}$ が存在し、ファイバー積の普遍性により、これは
$M_{2, i} \to M_{3, i}$ の切断を与える。したがって列
$$
0 \to M_{1, i} \to M_{2, i} \to M_{3, i} \to 0
$$
は分裂する。余極限をとると、行が完全で外側の垂直写像が同型である可換図式
$$
\xymatrix{
0  \ar[r] & \colim M_{1, i} \ar[d]^{\cong} \ar[r]  &  \colim M_{2, i}
 \ar[r]
\ar[d] & \colim M_{3, i} \ar[d]^{\cong} \ar[r] & 0 \\
0  \ar[r] & M_{1} \ar[r]  &  M_{2}  \ar[r] & M_{3} \ar[r] & 0
}
$$
を得る。したがって $\colim
M_{2, i}
\to M_2$ も同型であり、(6) が成り立つ。

\medskip\noindent
条件 (6) は例 \ref{algebra-example-universally-exact} (2) により (1) を含意する。
\end{proof}

\noindent
前の定理は、純完全列が常に分裂短完全列の余極限であることを示す。
純単射写像の余核が有限表示ならば、実際には写像そのものが分裂する。

\begin{lemma}
\label{algebra-lemma-universally-exact-split}
完全列
$$
0 \to M_1 \to M_2 \to M_3 \to 0
$$
を $R$-加群の完全列とし、$M_3$ は有限表示であると仮定する。このとき
$$
0 \to M_1 \to M_2 \to M_3 \to 0
$$
が純完全であることと、それが分裂することは同値である。
\end{lemma}

\begin{proof}
分裂短完全列は常に純完全である。例 \ref{algebra-example-universally-exact} を参照せよ。
逆に列が純完全ならば、定理 \ref{algebra-theorem-universally-exact-criteria} (5) を
$P = M_3$ に適用すると、写像 $M_2 \to M_3$ は切断をもつ。
\end{proof}

\noindent
次の補題は、純単射写像が平坦加群とどのように相補的であるかを示す。

\begin{lemma}
\label{algebra-lemma-flat-universally-injective}
$M$ を $R$-加群とする。このとき $M$ が平坦であることと、任意の
$R$-加群の完全列
$$
0 \to M_1 \to M_2 \to M \to 0
$$
が純完全であることは同値である。
\end{lemma}

\begin{proof}
補題 \ref{algebra-lemma-flat-surjective-hom} と定理
\ref{algebra-theorem-universally-exact-criteria} (5) から従う。
\end{proof}

\begin{example}
\label{algebra-example-universally-exact-non-split-non-flat}
分裂せず、平坦でもない純完全列の例を挙げる。
\begin{enumerate}
\item 補題 \ref{algebra-lemma-universally-exact-split} にもかかわらず、
純完全だが分裂しない $R$-加群の短完全列
$$
0 \to M_1 \to M_2 \to M_3 \to 0
$$
が存在し得る。例えば $R = \mathbf{Z}$ とし、
$M_1 = \bigoplus_{n=1}^{\infty} \mathbf{Z}$、
$M_{2} = \prod_{n = 1}^{\infty} \mathbf{Z}$ とおき、
$M_{3}$ を包含写像 $M_1 \to M_2$ の余核とする。
$M_1, M_2, M_3$ はいずれも捩れがないので平坦である
（More on Algebra, 補題 \ref{more-algebra-lemma-dedekind-torsion-free-flat}）。
したがって補題 \ref{algebra-lemma-flat-universally-injective} により
$$
0 \to M_1 \to M_2 \to M_3 \to 0
$$
は純完全である。しかし切断 $s: M_3 \to M_2$ は存在し得ない。
実際、$x$ を $(2, 2^2, 2^3, \ldots) \in M_2$ の $M_3$ における像とすると、
任意の加群写像 $s: M_3 \to M_2$ は $x$ を零に写さなければならない。
これは、任意の $n \geq 1$ に対して $x \in 2^n M_3$ なので、
すべての $n \geq 1$ に対して $s(x)$ は $2^n$ で割り切れ、したがって
$0$ でなければならないからである。
\item 補題 \ref{algebra-lemma-flat-universally-injective} にもかかわらず、
純完全だが $M_1, M_2, M_3$ がすべて非平坦である $R$-加群の短完全列
$$
0 \to M_1 \to M_2 \to M_3 \to 0
$$
が存在し得る。実際、$M$ が任意の非平坦加群ならば、分裂完全列
$$
0 \to M \to M \oplus M \to M \to 0.
$$
をとればよい。例えば $R = \mathbf{Z}$ 上では、$M$ を任意の捩れ加群とする。
\item (1) のような完全列と (2) のような完全列の直和をとると、
純完全かつ非分裂で、$M_1, M_2, M_3$ がすべて非平坦である
$R$-加群の短完全列
$$
0 \to M_1 \to M_2 \to M_3 \to 0
$$
を得る。
\end{enumerate}
\end{example}

\begin{lemma}
\label{algebra-lemma-ui-flat-domain}
$0 \to M_1 \to M_2 \to M_3 \to 0$ を $R$-加群の純完全列とし、
$M_2$ は平坦であると仮定する。このとき $M_1$ と $M_3$ は平坦である。
\end{lemma}

\begin{proof}
$0 \to N \to N' \to N'' \to 0$ を $R$-加群の短完全列とする。
可換図式
$$
\xymatrix{
M_1 \otimes_R N \ar[r] \ar[d] &
M_2 \otimes_R N \ar[r] \ar[d] &
M_3 \otimes_R N \ar[d] \\
M_1 \otimes_R N' \ar[r] \ar[d] &
M_2 \otimes_R N' \ar[r] \ar[d] &
M_3 \otimes_R N' \ar[d] \\
M_1 \otimes_R N'' \ar[r] &
M_2 \otimes_R N'' \ar[r] &
M_3 \otimes_R N''
}
$$
を考える（境界上の $0$ は省いた）。
仮定により各行は短完全列であり、矢印
$M_2 \otimes N \to M_2 \otimes N'$ は単射である。
これは明らかに $M_1 \otimes N \to M_1 \otimes N'$ が単射であることを含意し、
$M_1$ が平坦だと分かる。特に左列と中央列は短完全列を与える。
図式追跡により矢印 $M_3 \otimes N \to M_3 \otimes N'$ は単射である。
したがって $M_3$ は平坦である。
\end{proof}

\begin{lemma}
\label{algebra-lemma-universally-injective-tensor}
$R$ を環とする。$M \to M'$ を純単射な $R$-加群写像とする。
このとき任意の $R$-加群 $N$ に対して写像
$M \otimes_R N \to M' \otimes_R N$ は純単射である。
\end{lemma}

\begin{proof}
省略する。
\end{proof}

\begin{lemma}
\label{algebra-lemma-composition-universally-injective}
$R$ を環とする。純単射な $R$-加群写像の合成は純単射である。
\end{lemma}

\begin{proof}
省略する。
\end{proof}

\begin{lemma}
\label{algebra-lemma-universally-injective-permanence}
$R$ を環とする。$M \to M'$ と $M' \to M''$ を $R$-加群写像とする。
それらの合成 $M \to M''$ が純単射ならば、$M \to M'$ は純単射である。
\end{lemma}

\begin{proof}
省略する。
\end{proof}

\begin{lemma}
\label{algebra-lemma-universally-injective-check-stalks}
$R \to S$ を環準同型とする。$M \to M'$ を $S$-加群の写像とする。
次は同値である。
\begin{enumerate}
\item $M \to M'$ は $R$-加群の写像として純単射である。
\item $S$ の各素イデアル $\mathfrak q$ に対して、写像
$M_{\mathfrak q} \to M'_{\mathfrak q}$ は $R$-加群の写像として純単射である。
\item $S$ の各極大イデアル $\mathfrak m$ に対して、写像
$M_{\mathfrak m} \to M'_{\mathfrak m}$ は $R$-加群の写像として純単射である。
\item $S$ の各素イデアル $\mathfrak q$ に対して、写像
$M_{\mathfrak q} \to M'_{\mathfrak q}$ は $R_{\mathfrak p}$-加群の写像として
純単射である。ここで $\mathfrak p$ は $R$ における $\mathfrak q$ の逆像である。
\item $S$ の各極大イデアル $\mathfrak m$ に対して、写像
$M_{\mathfrak m} \to M'_{\mathfrak m}$ は $R_{\mathfrak p}$-加群の写像として
純単射である。ここで $\mathfrak p$ は $R$ における $\mathfrak m$ の逆像である。
\end{enumerate}
\end{lemma}

\begin{proof}
$N$ を $R$-加群とする。$\mathfrak q$ を $R$ の素イデアル
$\mathfrak p$ の上にある $S$ の素イデアルとする。このとき
$$
(M \otimes_R N)_{\mathfrak q} =
M_{\mathfrak q} \otimes_R N =
M_{\mathfrak q} \otimes_{R_{\mathfrak p}} N_{\mathfrak p}.
$$
が成り立つ。さらに $M'$ についても同じことが成り立ち、局所化は完全である。
また、$N$ が $R_{\mathfrak p}$-加群ならば $N_{\mathfrak p} = N$ である。
これらを用いれば、各同値性は直接的に証明できる。

\medskip\noindent
例えば (5) が成り立つと仮定する。
$K = \Ker(M \otimes_R N \to M' \otimes_R N)$ とおく。
上の注意により、$S$ の各極大イデアル $\mathfrak m$ に対して
$K_{\mathfrak m} = 0$ である。したがって補題
\ref{algebra-lemma-characterize-zero-local} により $K = 0$ である。
ゆえに (1) が成り立つ。逆に (1) が成り立つと仮定する。
$\mathfrak p \subset R$ の上にある任意の $\mathfrak q \subset S$ をとる。
$R_{\mathfrak p}$ 上の任意の加群 $N$ をとる。
仮定により $\Ker(M \otimes_R N \to M' \otimes_R N) = 0$ である。
したがって上の公式と $N = N_{\mathfrak p}$ であることから
$\Ker(M_{\mathfrak q} \otimes_{R_{\mathfrak p}} N \to
M'_{\mathfrak q} \otimes_{R_{\mathfrak p}} N) = 0$ が分かる。
言い換えれば (4) が成り立つ。もちろん (4) $\Rightarrow$ (5) は直ちに従う。
したがって (1), (4), (5) はすべて同値である。
ほかの同値性の証明は省略する。
\end{proof}

\begin{lemma}
\label{algebra-lemma-universally-injective-localize}
$\varphi : A \to B$ を環準同型とする。$S \subset A$ と $S' \subset B$ を
$\varphi(S) \subset S'$ を満たす乗法的部分集合とする。
$M \to M'$ を $B$-加群の写像とする。
\begin{enumerate}
\item $M \to M'$ が $A$-加群の写像として純単射ならば、
$(S')^{-1}M \to (S')^{-1}M'$ は $A$-加群の写像としても
$S^{-1}A$-加群の写像としても純単射である。
\item $M$ と $M'$ が $(S')^{-1}B$-加群ならば、$M \to M'$ が
$A$-加群の写像として純単射であることと、$S^{-1}A$-加群の写像として
純単射であることは同値である。
\end{enumerate}
\end{lemma}

\begin{proof}
補題 \ref{algebra-lemma-universally-injective-check-stalks} を用いて証明できるが、
次のように直接証明することもできる。
$M \to M'$ が $A$ 上純単射であると仮定する。
$Q$ を $A$-加群とする。このとき $Q \otimes_A M \to Q \otimes_A M'$ は単射である。
局所化は完全なので
$(S')^{-1}(Q \otimes_A M) \to (S')^{-1}(Q \otimes_A M')$ は単射である。
$(S')^{-1}(Q \otimes_A M) = Q \otimes_A (S')^{-1}M$ であり、$M'$ についても
同様なので、$Q \otimes_A (S')^{-1}M \to Q \otimes_A (S')^{-1}M'$ は単射である。
したがって $(S')^{-1}M \to (S')^{-1}M'$ は $A$-加群の写像として純単射である。
これで (1) の第一の部分が示された。
% [P01-E0448] Frozen proof progression retained; the second assertion of (1) is not inserted.
(2) を見るには次の二つの事実を用いればよい。
(a) $Q$ が $S^{-1}A$-加群ならば $Q \otimes_A S^{-1}A = Q$ であり、すなわち
$A$ 上で $Q$ とテンソル積をとることは $S^{-1}A$ 上で $Q$ とテンソル積を
とることと同じである。
(b) $M$ が $S$ の元が可逆に作用する任意の $A$-加群ならば、
$M \otimes_A Q = M \otimes_{S^{-1}A} S^{-1}Q$ である。
(2) はこれから直ちに従う。
\end{proof}

\begin{lemma}
\label{algebra-lemma-check-universally-injective-into-flat}
$R$ を環、$M \to M'$ を $R$-加群の写像とする。
$M'$ が平坦ならば、$M \to M'$ が純単射であることと、$R$ のすべての
有限生成イデアル $I$ に対して $M/IM \to M'/IM'$ が単射であることは同値である。
\end{lemma}

\begin{proof}
定理 \ref{algebra-theorem-universally-exact-criteria} により、すべての有限 $R$-加群 $Q$ に
対して $M \otimes_R Q \to M' \otimes_R Q$ が単射であることを示せば十分である。
補題 \ref{algebra-lemma-trivial-filter-finite-module} により、$Q$ は、逐次商が巡回加群
$R/I_i$ と同型になる部分加群による有限フィルトレーション
$0 = Q_0 \subset Q_1 \subset \ldots \subset Q_n = Q$ をもつ。
$M'$ は平坦なので、逐次商が
$M' \otimes_R R/I_i = M'/I_iM'$ となる部分加群
$M' \otimes_R Q_i$ による $M' \otimes_R Q$ のフィルトレーション
$$
\xymatrix{
M \otimes Q_1 \ar[r] \ar[d] &
M \otimes Q_2 \ar[r] \ar[d] &
\ldots \ar[r] &
M \otimes Q \ar[d] \\
M' \otimes Q_1 \ar@{^{(}->}[r] &
M' \otimes Q_2 \ar@{^{(}->}[r] &
\ldots \ar@{^{(}->}[r] &
M' \otimes Q
}
$$
を得る。単純な帰納法により、$M/I_i M \to M'/I_i M'$ が単射であることを
確認すれば十分だと分かる。有限生成イデアル $I'_i \subset I_i$ の集まりは
有向集合であることに注意する。したがって $M/I_iM = \colim M/I'_iM$ は
フィルター余極限であり、$M'$ についても同様である。
写像 $M/I'_iM \to M'/I'_i M'$ は仮定により単射であり、フィルター余極限は
完全なので（補題 \ref{algebra-lemma-directed-colimit-exact}）、結論を得る。
\end{proof}

\begin{lemma}
\label{algebra-lemma-base-change-along-universally-injective-ring-map}
$R \to S$ を $R$-加群の写像として純単射である環準同型とする。
このとき $R$-加群上の関手 $M \mapsto M \otimes_R S$ は単射、全射および同型を反映する。
\end{lemma}

\begin{proof}
$M \to N$ を $R$-加群の写像とし、その核を $K$、余核を $Q$ とする。
$M \otimes_R S \to N \otimes_R S$ が単射ならば、
$K \otimes_R S \to M \otimes_R S$ の像は零である。
$K \subset K \otimes_R S$ および $M \subset M \otimes_R S$ なので $K = 0$ と結論する。
一方、$M \otimes_R S \to N \otimes_R S$ が全射ならば、テンソル積の右完全性により
$Q \otimes_R S$ は零である。したがって $Q \subset Q \otimes_R S$ も零である。
\end{proof}

\begin{lemma}
\label{algebra-lemma-faithfully-flat-universally-injective}
$R \to S$ を忠実平坦な環準同型とする。このとき $R \to S$ は
$R$-加群の写像として純単射である。特に任意のイデアル
$I \subset R$ に対して $R \cap IS = I$ である。
\end{lemma}

\begin{proof}
$N$ を $R$-加群とする。$N \to N \otimes_R S$ が単射であることを示さなければならない。
$S$ は $R$-加群として忠実平坦なので、$S$ とテンソル積をとった後にこれを証明すれば
十分である。したがって
$N \otimes_R S \to N \otimes_R S \otimes_R S$、
$n \otimes s \mapsto n \otimes 1 \otimes s$ が単射であることを示せば十分である。
これは引き戻し $n \otimes s \otimes s' \mapsto n \otimes ss'$ が存在するので真である。
\end{proof}

\section{有限射影加群の降下}
\label{algebra-section-finite-projective}

\noindent
この節では、{\it 有限} 射影加群であるという性質が忠実平坦な環準同型に沿って
降下するという事実の初等的証明を与える。有限性の条件を外すと、この証明は
適用できない。しかしこの方法は、{\it 可算生成} 射影加群であるという性質の
降下を証明するために用いる方法を示唆している。この節末のコメントを参照せよ。

\begin{lemma}
\label{algebra-lemma-finite-projective-again}
$M$ を $R$-加群とする。このとき $M$ が有限射影的であることと、
$M$ が有限表示かつ平坦であることは同値である。
\end{lemma}

\begin{proof}
これは補題 \ref{algebra-lemma-finite-projective} の一部である。
しかしこの時点では、その補題の含意 (1) $\Rightarrow$ (2) のより洗練された証明を
次のように与えられる。$M$ が有限表示かつ平坦ならば、全射 $R^n \to M$ をとる。
補題 \ref{algebra-lemma-flat-surjective-hom} を $P = M$ に適用すると、写像
$R^n \to M$ は切断をもつ。したがって $M$ は自由加群の直和因子であり、ゆえに射影的である。
\end{proof}

\noindent
次に、降下する加群の性質をいくつか挙げる。

\begin{lemma}
\label{algebra-lemma-descend-properties-modules}
$R \to S$ を忠実平坦な環準同型とする。$M$ を $R$-加群とする。このとき
\begin{enumerate}
\item $S$-加群 $M \otimes_R S$ が有限型ならば、$M$ は有限型である。
\item $S$-加群 $M \otimes_R S$ が有限表示ならば、$M$ は有限表示である。
\item $S$-加群 $M \otimes_R S$ が平坦ならば、$M$ は平坦である。
\item 必要に応じてここにさらに加える。
% [P01-E0449] Frozen publication-facing authoring placeholder retained.
\end{enumerate}
\end{lemma}

\begin{proof}
$M \otimes_R S$ が有限型であると仮定する。
$y_1, \ldots, y_m$ を $S$ 上の $M \otimes_R S$ の生成元とする。
ある $x_1, \ldots, x_n \in M$ を用いて $y_j = \sum x_i \otimes f_i$ と書く。
% [P01-E0450] Frozen missing j-index on f_i retained.
すると写像 $\varphi : R^{\oplus n} \to M$ は、
$\varphi \otimes \text{id}_S : S^{\oplus n} \to M \otimes_R S$ が全射であるという
性質をもつ。$R \to S$ は忠実平坦なので $\varphi$ は全射であり、
$M$ は有限生成である。

\medskip\noindent
$M \otimes_R S$ が有限表示であると仮定する。(1) により $M$ は有限型である。
全射 $R^{\oplus n} \to M$ を選び、その核を $K$ と書く。
$R \to S$ は平坦なので、$K \otimes_R S$ は基底変換
$S^{\oplus n} \to M \otimes_R S$ の核である。
$M \otimes_R S$ は有限表示なので、$K \otimes_R S$ は有限型である。
したがって (1) により $K$ は有限型であり、ゆえに $M$ は有限表示である。

\medskip\noindent
(3) は補題 \ref{algebra-lemma-flatness-descends} である。
\end{proof}

\begin{proposition}
\label{algebra-proposition-ffdescent-finite-projectivity}
$R \to S$ を忠実平坦な環準同型、$M$ を $R$-加群とする。
$S$-加群 $M \otimes_R S$ が有限射影的ならば、$M$ は有限射影的である。
\end{proposition}

\begin{proof}
補題 \ref{algebra-lemma-finite-projective-again} および
\ref{algebra-lemma-descend-properties-modules} から従う。
\end{proof}

\noindent
次のいくつかの節では、デヴィサージュを用いて可算生成の場合に帰着することにより、
有限性の仮定を取り除く。可算生成の場合の方針は、補題
\ref{algebra-lemma-finite-projective-again} に類似した可算生成射影加群の特徴付けを見つけ、
その特徴付けが降下することを直接証明することである。これを行うため、
Mittag-Leffler 加群の概念を導入し、加群 $M$ が可算生成ならば、
それが射影的であることと平坦かつ Mittag-Leffler であることが同値だと証明する
（定理 \ref{algebra-theorem-projectivity-characterization}）。$M$ が有限生成ならば、
この主張は補題 \ref{algebra-lemma-finite-projective-again} に帰着する
（例 \ref{algebra-example-ML} (1) によれば、有限生成加群が Mittag-Leffler であることと
有限表示であることは同値だからである）。

\section{加群の超限デヴィサージュ}
\label{algebra-section-transfinite-devissage}

\noindent
この節では、加群を直和に分解するためのデヴィサージュの手法を導入する。
主結果は、射影加群が可算生成加群の直和であるというものである
（下の定理 \ref{algebra-theorem-projective-direct-sum}）。\cite{Kaplansky} に従う。

\begin{definition}
\label{algebra-definition-devissage}
$M$ を $R$-加群とする。$M$ の {\it 直和デヴィサージュ} とは、順序数 $S$ で
添字付けられ、包含関係について増大する部分加群の族
$(M_{\alpha})_{\alpha \in S}$ で、次を満たすものである。
\begin{enumerate}
\item[(0)] $M_0 = 0$。
\item[(1)] $M = \bigcup_{\alpha} M_{\alpha}$。
\item[(2)] $\alpha \in S$ が極限順序数ならば、
$M_{\alpha} = \bigcup_{\beta < \alpha} M_{\beta}$。
\item[(3)] $\alpha + 1 \in S$ ならば、$M_{\alpha}$ は
$M_{\alpha + 1}$ の直和因子である。
\end{enumerate}
さらに
\begin{enumerate}
\item[(4)] $\alpha + 1 \in S$ に対して
$M_{\alpha + 1}/M_{\alpha}$ は可算生成である。
\end{enumerate}
を満たすとき、$(M_{\alpha})_{\alpha \in S}$ を $M$ の
{\it Kaplansky デヴィサージュ} と呼ぶ。
\end{definition}

\noindent
この用語は次の補題によって正当化される。

\begin{lemma}
\label{algebra-lemma-direct-sum-devissage}
$M$ を $R$-加群とする。$(M_{\alpha})_{\alpha \in S}$ が $M$ の
直和デヴィサージュならば、
$M \cong \bigoplus_{\alpha + 1 \in S} M_{\alpha + 1}/M_{\alpha}$ である。
\end{lemma}

\begin{proof}
直和デヴィサージュの性質 (3) により、各 $\alpha \in S$ に対して包含写像
$M_{\alpha + 1}/M_{\alpha} \to M$ が存在する。これらの包含写像の和で与えられる写像
$$
f : \bigoplus\nolimits_{\alpha + 1\in S} M_{\alpha + 1}/M_{\alpha} \to M
$$
を考える。さらに $\beta\in S$ に対して制限
$$
f_{\beta} :
\bigoplus\nolimits_{\alpha + 1 \leq \beta} M_{\alpha + 1}/M_{\alpha}
\longrightarrow
M
$$
を考える。$S$ 上の超限帰納法により、$f_{\beta}$ の像が $M_{\beta}$ であることを
示せる。$\beta=0$ については $(0)$ により真である。
$\beta+1$ が後続順序数で $\beta$ について真ならば、(3) により
$\beta + 1$ についても真である。また $\beta$ が極限順序数で、
$\alpha < \beta$ に対して真ならば、(2) により $\beta$ についても真である。
したがって (1) により $f$ は全射である。

\medskip\noindent
$S$ 上の超限帰納法は、制限 $f_{\beta}$ が単射であることも示す。
$\beta = 0$ については真である。$\beta+1$ が後続順序数で
$f_{\beta}$ が単射であるとし、$x$ を核の元とする。
その成分 $x_{\alpha + 1} \in M_{\alpha + 1}/M_{\alpha}$ を用いて
$x = (x_{\alpha + 1})_{\alpha + 1 \leq \beta + 1}$ と書く。
性質 (3) と $f_{\beta}$ の像が $M_{\beta}$ であることにより、
$(x_{\alpha + 1})_{\alpha + 1 \leq \beta}$ と $x_{\beta + 1}$ はともに
$0$ に写る。したがって $x_{\beta+1} = 0$ であり、制限 $f_{\beta}$ が単射である
という仮定により、すべての $\alpha + 1 \leq \beta$ に対して
$x_{\alpha + 1} = 0$ でもある。ゆえに $x = 0$ で、$f_{\beta+1}$ は単射である。
$\beta$ が極限順序数ならば、核の元 $x$ を考える。このとき $x$ はすでに
ある $\alpha < \beta$ に対して $f_{\alpha}$ の定義域に含まれる。
したがって $x = 0$ であり、帰納法が完了する。
各 $\beta \in S$ に対して $f_{\beta}$ が単射なので、$f$ は単射だと結論する。
\end{proof}

\begin{lemma}
\label{algebra-lemma-Kaplansky-devissage}
$M$ を $R$-加群とする。このとき $M$ が可算生成 $R$-加群の直和であることと、
Kaplansky デヴィサージュをもつことは同値である。
\end{lemma}

\begin{proof}
補題により「もし」の向きは示される。逆に、各 $N_i$ が可算生成 $R$-加群である
$M = \bigoplus_{i \in I} N_i$ を仮定する。$I$ を整列順序付けし、順序数とみなす。
このとき $M_i = \bigoplus_{j < i} N_j$ とおけば、
$(M_i)_{i \in I}$ は $M$ の Kaplansky デヴィサージュを与える。
% [P01-E0451] Frozen indexing without the terminal ordinal is retained.
\end{proof}

\begin{theorem}
\label{algebra-theorem-kaplansky-direct-sum}
$M$ が可算生成 $R$-加群の直和であると仮定する。$P$ が $M$ の直和因子ならば、
$P$ も可算生成 $R$-加群の直和である。
\end{theorem}

\begin{proof}
$M = P \oplus Q$ と書く。$M$ の Kaplansky デヴィサージュ
$(M_{\alpha})_{\alpha \in S}$ であって、定義の性質 (0)--(4) に加えて
次を満たすものを構成する：
\begin{enumerate}
\item[(5)] 各 $M_{\alpha}$ は $M$ の直和因子である；
\item[(6)] $M_{\alpha} = P_{\alpha} \oplus Q_{\alpha}$ である。ただし $P_{\alpha} =P
\cap M_{\alpha}$ および $Q_\alpha = Q \cap M_{\alpha}$ とする。
\end{enumerate}
（注意：性質 (0)--(2) が成り立つならば、実際、性質 (3) と性質 (5) は同値である。）

\medskip\noindent
これが定理を導くことを見るには、
$(P_{\alpha})_{\alpha \in S}$ が $P$ の Kaplansky デヴィサージュをなすことを示せば十分である。
性質 (0)、(1)、(2) は明らかである。$(M_{\alpha})$ に対する (5) と (6) により、
各 $P_{\alpha}$ は $M$ の直和因子である。$P_{\alpha} \subset P_{\alpha +
1}$ なので、これは $P_{\alpha}$ が $P_{\alpha + 1}$ の直和因子であることを含意する；
したがって $(P_{\alpha})$ に対して (3) が成り立つ。(4) については、
$$
M_{\alpha + 1}/M_{\alpha} \cong P_{\alpha + 1}/P_{\alpha} \oplus
Q_{\alpha + 1}/Q_{\alpha},
$$
であることに注意する。したがって $M_{\alpha + 1}/M_{\alpha}$ が可算生成であることから、
$P_{\alpha + 1}/P_{\alpha}$ も可算生成である。

\medskip\noindent
残るのは $M_{\alpha}$ を構成することである。各 $N_i$ が可算生成 $R$-加群であるように
$M = \bigoplus_{i \in I} N_i$ と書く。$I$ の整列順序を選ぶ。超限再帰により、
各順序数 $\alpha$ に対して $M$ の部分加群 $M_{\alpha}$ を、増大族をなしかつ
$M_{\alpha}$ がいくつかの $N_i$ の直和となるように定める。

\medskip\noindent
$\alpha = 0$ に対しては $M_{0} = 0$ とおく。$\alpha$ が極限順序数で、すべての
$\beta < \alpha$ に対して $M_{\beta}$ が定義されているなら、$M_{\alpha}
= \bigcup_{\beta < \alpha} M_{\beta}$ と定める。各 $\beta <
\alpha$ に対して $M_{\beta}$ はいくつかの $N_i$ の直和なので、
$M_{\alpha}$ についても同じことが成り立つ。$\alpha + 1$ が後続順序数で
$M_{\alpha}$ が定義されているなら、$M_{\alpha + 1}$ を次のように定める。
$M_{\alpha} = M$ なら $M_{\alpha + 1} = M$ とおく。そうでなければ、
$N_j$ が $M_{\alpha}$ に含まれないような最小の $j \in I$ を選ぶ。
次の性質を満たす無限行列 $(x_{mn}), m, n = 1, 2, 3, \ldots$ を構成する：
\begin{enumerate}
\item $N_j$ は、成分 $x_{mn}$ によって生成される $M$ の部分加群に含まれる；
\item 任意の成分 $x_{k\ell}$ をその $P$-成分と $Q$-成分を用いて
$x_{k\ell} = y_{k\ell} + z_{k\ell}$ と書くとき、行列 $(x_{mn})$ は、
$y_{k\ell}$ または $z_{k\ell}$ が非零成分をもつ各 $N_i$ の生成元の集合を含む。
\end{enumerate}
そこで $M_{\alpha + 1}$ を、$M_{\alpha}$ とすべての $x_{mn}$ によって生成される
$M$ の部分加群と定める；行列 $(x_{mn})$ の性質 (2) により、
$M_{\alpha + 1}$ はいくつかの $N_i$ の直和となる。
行列 $(x_{mn})$ を構成するため、$x_{11}, x_{12}, x_{13}, \ldots$ を
$N_j$ の可算な生成元の集合とする。次に、
$x_{11} = y_{11} + z_{11}$ が $P$-成分と $Q$-成分への分解ならば、
$x_{21}, x_{22}, x_{23}, \ldots$ を、$y_{11}$ または $z_{11}$ が非零成分をもつ
$N_i$ たちの和の可算な生成元の集合とする。この過程を $x_{12}$ に対して繰り返し、
行列の第 3 行となる元 $x_{31}, x_{32}, \ldots$ を得る。$x_{21}$ に対して繰り返して
第 4 行を、$x_{13}$ に対して繰り返して第 5 行を得る、というように、
下図に示す連続する反対角線に沿って進める：
$$
\left(
\vcenter{
\xymatrix@R=2mm@C=2mm{
x_{11} & x_{12} \ar[dl] & x_{13} \ar[dl] & x_{14} \ar[dl] & \ldots  \\
x_{21} & x_{22} \ar[dl] & x_{23} \ar[dl] & \ldots  \\
x_{31} & x_{32} \ar[dl] & \ldots \\
x_{41} & \ldots \\
\ldots
}
}
\right).
$$

\medskip\noindent
$I$ 上の超限帰納法により（$N_j$ が $M_{\alpha}$ に含まれないような最小の $j$ に対し、
$M_{\alpha + 1}$ が $N_j$ を含むように構成したことを用いる）、各 $i \in I$ に対して、
$N_i$ はある $M_{\alpha}$ に含まれる。したがって、各 $i \in I$ に対して
$N_i$ がある $\alpha \in S$ に対する $M_{\alpha}$ に含まれるような、十分大きい順序数
$S$ が存在する。これは $(M_{\alpha})_{\alpha \in S}$ が $M$ の Kaplansky デヴィサージュの
性質 (1) を満たすことを意味する。さらに族 $(M_{\alpha})_{\alpha \in S}$ は、
その他の定義上の性質、および上記の (5) と (6) も満たす：
性質 (0)、(2)、(4)、(6) は構成から明らかである；性質 (5) は、各 $M_{\alpha}$ が構成上
いくつかの $N_i$ の直和だから成り立つ；そして (3) は (5) と
$M_{\alpha} \subset M_{\alpha + 1}$ であることから従う。
\end{proof}

\noindent
系として、この節の冒頭で述べた射影加群に関する結果を得る。

\begin{theorem}
\label{algebra-theorem-projective-direct-sum}
\begin{slogan}
任意の射影加群は、可算生成射影加群の直和である。
\end{slogan}
$P$ が射影 $R$-加群ならば、$P$ は可算生成射影 $R$-加群の直和である。
\end{theorem}

\begin{proof}
加群が射影的であることと、それが自由加群の直和因子であることは同値なので、
これは定理 \ref{algebra-theorem-kaplansky-direct-sum} から従う。
\end{proof}



\section{局所環上の射影加群}
\label{algebra-section-projective-local-ring}

\noindent
この節では、たいへん美しい結果を証明する：局所環上の射影加群 $M$ は自由である
（下の定理 \ref{algebra-theorem-projective-free-over-local-ring}）。
$M$ が有限であるという追加の仮定のもとでは、この結果は
補題 \ref{algebra-lemma-finite-flat-local} であることに注意する。一般には次が成り立つ：

\begin{lemma}
\label{algebra-lemma-projective-free}
$R$ を環とする。このとき、すべての射影 $R$-加群が自由であることと、
すべての可算生成射影 $R$-加群が自由であることは同値である。
\end{lemma}

\begin{proof}
定理 \ref{algebra-theorem-projective-direct-sum} から直ちに従う。
\end{proof}

\noindent
以下は、可算生成加群が自由となるための判定条件である。

\begin{lemma}
\label{algebra-lemma-freeness-criteria}
$M$ を、次の性質をもつ可算生成 $R$-加群とする：
$M = N \oplus N'$ かつ $N'$ が有限自由 $R$-加群ならば、
$N$ の任意の元は $N$ のある自由直和因子に含まれる。
このとき $M$ は自由である。
\end{lemma}

\begin{proof}
$x_1, x_2, \ldots$ を $M$ の可算な生成元の集合とする。すべての $n$ に対して、
$F_1 \oplus \ldots \oplus F_n$ が $x_1, \ldots, x_n$ を含む $M$ の直和因子となるような、
$M$ の有限自由直和因子 $F_1, F_2, \ldots$ を帰納的に構成する。
実際、所望の性質をもつ $F_1, \ldots, F_n$ が与えられたとき、
$$
M = F_1 \oplus \ldots \oplus F_n \oplus N
$$
と書き、$x \in N$ を $x_{n + 1}$ の像とする。補題の仮定により、
$x$ を含む自由直和因子 $F_{n + 1} \subset N$ を見つけられる。
もちろん $F_{n + 1}$ を、$F_{n + 1}$ の有限自由直和因子で置き換えることができ、
帰納段階が完了する。このとき $M = \bigoplus_{i = 1}^{\infty} F_i$ は自由である。
\end{proof}

\begin{lemma}
\label{algebra-lemma-projective-freeness-criteria}
$P$ を局所環 $R$ 上の射影加群とする。このとき $P$ の任意の元は、
$P$ のある自由直和因子に含まれる。
\end{lemma}

\begin{proof}
$P$ は射影的なので、ある自由 $R$-加群 $F$ の直和因子であり、$F = P \oplus Q$ と書ける。
$x \in P$ を、$P$ の自由直和因子に含まれることを示したい元とする。
$x$ の表示に必要な基底元の個数が最小となるような $F$ の基底 $B$ をとり、
ある $e_i \in B$ と $a_i \in R$ によって $x
= \sum_{i=1}^n a_i e_i$ と書く。このとき、どの $a_j$ も他の $a_i$ の線形結合として
表すことはできない；実際、ある $b_i \in R$ によって
$a_j = \sum_{i \neq  j} a_i b_i$ と書けるなら、$i \neq j$ に対して $e_i$ を
$e_i + b_ie_j$ で置き換え、$B$ の他の元を変えないことで、$x$ がより短い表示をもつ
$F$ の新しい基底が得られる。

\medskip\noindent
$e_i$ の $P$-成分と $Q$-成分への分解を
$e_i = y_i + z_i, y_i \in P, z_i \in Q$ とする。
$y_i = \sum_{j=1}^{n} b_{ij} e_j + t_i$ と書く。ただし $t_i$ は、
$e_1, \ldots, e_n$ 以外の $B$ の元の線形結合である。証明を終えるには、
行列 $(b_{ij})$ が可逆であることを示せば十分である。実際、そのとき
$i=1, \ldots, n$ に対して $e_i \mapsto y_i$ とし、
$B \setminus \{e_1, \ldots, e_n\}$ を固定する写像 $F \to F$ は同型となる。
したがって $y_1, \ldots, y_n$ と $B \setminus \{e_1, \ldots, e_n\}$ は合わせて
$F$ の基底をなす。このとき $y_1, \ldots, y_n$ が張る部分加群 $N$ は
$P$ の自由部分加群である；$N \subset P$ で、$N$ と $P$ はともに $F$ の直和因子なので、
$N$ は $P$ の直和因子である；さらに $x \in P$ から
$x = \sum_{i=1}^n a_i e_i = \sum_{i=1}^n a_i y_i$ となるので $x \in N$ である。

\medskip\noindent
そこで $(b_{ij})$ が可逆であることを示す。
$y_i = \sum_{j=1}^{n} b_{ij} e_j + t_i$ を
$\sum_{i=1}^n a_i e_i = \sum_{i=1}^n a_i y_i$ に代入し、$e_j$ の係数を比較すると、
$a_j = \sum_{i=1}^n a_i b_{ij}$ を得る。しかし上で述べたように、$B$ の選び方により、
どの $a_j$ も他の $a_i$ の線形結合として書けない。したがって $i \neq
j$ に対して $b_{ij}$ は非単元であり、すべての $i$ に対して $1-b_{ii}$ は非単元である
---ゆえに、とくに $b_{ii}$ は単元である。しかし局所環上で、対角成分が単元、
それ以外の成分が非単元である行列は、その行列式が単元なので可逆である。
\end{proof}

\begin{theorem}
\label{algebra-theorem-projective-free-over-local-ring}
\begin{slogan}
局所環上の射影加群は自由である。
\end{slogan}
$P$ が局所環 $R$ 上の射影加群ならば、$P$ は自由である。
\end{theorem}

\begin{proof}
補題 \ref{algebra-lemma-projective-free}、\ref{algebra-lemma-freeness-criteria}、および
\ref{algebra-lemma-projective-freeness-criteria} から従う。
\end{proof}



\section{Mittag-Leffler 系}
\label{algebra-section-mittag-leffler}

\noindent
この節の目的は、Mittag-Leffler 系を定義し、この概念が有用である理由を説明することである。

\medskip\noindent
以下では、$I$ を有向集合とする。Categories, Definition
\ref{categories-definition-directed-set} を参照せよ。
$I$ で添字付けられた集合または加群の逆系
$(A_i, \varphi_{ji}: A_j \to A_i)$ を考える。Categories, Definition
\ref{categories-definition-directed-system} を参照せよ。$I$ を有向と仮定したので、
これは有向逆系である（Categories, Definition
\ref{categories-definition-directed-system}）。各 $i \in I$ に対し、$j \geq i$ における像
$\varphi_{ji}(A_j) \subset A_i$ は、$A_i$ の部分集合（または部分加群）の減少有向族をなす。
$A'_i = \bigcap_{j \geq i} \varphi_{ji}(A_j)$ とおく。このとき $j \geq i$ に対して
$\varphi_{ji}(A'_j) \subset A'_i$ なので、制限により有向逆系
$(A'_i, \varphi_{ji}|_{A'_j})$ を得る。集合または加群の圏における逆系の極限の構成から、
$\lim A_i = \lim A'_i$ である。$(A_i, \varphi_{ji})$ に対する Mittag-Leffler 条件とは、
$A'_i$ がある $j \geq i$ に対する $\varphi_{ji}(A_j)$ に等しいこと
（したがって、すべての $k \geq j$ に対する $\varphi_{ki}(A_k)$ にも等しいこと）である：

\begin{definition}
\label{algebra-definition-ML-system}
$(A_i, \varphi_{ji})$ を $I$ 上の集合の有向逆系とする。各 $i \in I$ に対して、
$j \geq i$ における族 $\varphi_{ji}(A_j) \subset A_i$ が安定化するとき、
$(A_i, \varphi_{ji})$ は {\it Mittag-Leffler} であるという。明示的には、これは各
$i \in I$ に対してある $j \geq i$ が存在し、$k \geq j$ ならば
$\varphi_{ki}(A_k) = \varphi_{ji}( A_j)$ となることを意味する。
$(A_i, \varphi_{ji})$ が環 $R$ 上の加群の有向逆系であるとき、その台となる集合の逆系が
Mittag-Leffler であるならば、この加群の逆系も Mittag-Leffler であるという。
\end{definition}

\begin{example}
\label{algebra-example-ML-surjective-maps}
$(A_i, \varphi_{ji})$ が集合または加群の有向逆系で、写像 $\varphi_{ji}$ が全射ならば、
この系は明らかに Mittag-Leffler である。逆に、$(A_i, \varphi_{ji})$ が
Mittag-Leffler であると仮定する。$A'_i \subset A_i$ を、$j \geq
i$ における $\varphi_{ji}(A_j)$ の安定した像とする。このとき $j \geq i$ に対して
$\varphi_{ji}|_{A'_j}: A'_j \to A'_i$ は全射であり、$\lim A_i = \lim A'_i$ である。
したがって Mittag-Leffler 系 $(A_i, \varphi_{ji})$ の極限は、全射からなる
$I$ 上の有向逆系の極限としても書ける。
\end{example}

\begin{lemma}
\label{algebra-lemma-ML-limit-nonempty}
$(A_i, \varphi_{ji})$ を $I$ 上の有向逆系とする。$I$ が可算であると仮定する。
$(A_i, \varphi_{ji})$ が Mittag-Leffler で、各 $A_i$ が空でないならば、
$\lim A_i$ は空でない。
\end{lemma}

\begin{proof}
$i_1, i_2, i_3, \ldots$ を $I$ の元の列挙とする。$n = 1, 2, 3, \ldots$ に対する
元の列 $j_n \in I$ を、条件 $j_1 = i_1$、および $j_n \geq i_n$ かつ
$m < n$ に対して $j_n \geq j_m$ によって帰納的に定める。このとき列 $j_n$ は増大し、
$I$ の共終部分集合をなす。したがって $I =\{1, 2, 3, \ldots \}$ と仮定してよい。
よって例 \ref{algebra-example-ML-surjective-maps} により、正の整数で添字付けられた、
空でない集合と全射からなる逆系の極限が空でないことを示せばよい。
これは選択公理から従う。
\end{proof}

\noindent
次の完全性の性質のため、Mittag-Leffler 条件は重要となる。

\begin{lemma}
\label{algebra-lemma-ML-exact-sequence}
$$
0 \to A_i \xrightarrow{f_i} B_i \xrightarrow{g_i} C_i \to 0
$$
を $I$ 上のアーベル群の有向逆系の完全列とする。$I$ が可算であると仮定する。
$(A_i)$ が Mittag-Leffler ならば、
$$
0 \to \lim A_i \to \lim B_i \to \lim C_i\to 0
$$
は完全である。
\end{lemma}

\begin{proof}
有向逆系の極限をとる操作は左完全なので、$\lim B_i \to \lim C_i$ の全射性だけを
証明すればよい。そこで $(c_i) \in \lim
C_i$ とする。各 $i \in I$ に対して $E_i = g_i^{-1}(c_i)$ とおく；
$g_i: B_i \to C_i$ は全射なので、これは空でない。$(B_i)$ に対する写像の系
$\varphi_{ji}: B_j \to B_i$ は写像 $E_j \to E_i$ に制限され、
$(E_i)$ を空でない集合の逆系にする。$(E_i)$ が Mittag-Leffler であることを示せば十分である。
実際、そのとき補題 \ref{algebra-lemma-ML-limit-nonempty} により $\lim E_i$ は空でなく、
$\lim E_i$ の任意の元をとれば、$(c_i)$ に写る $\lim B_i$ の元が得られる。

\medskip\noindent
単射 $f_i: A_i \to B_i$ により、$A_i$ を $B_i$ の部分集合とみなす。
$(A_i)$ は Mittag-Leffler なので、$i \in I$ ならば、$k \geq j$ に対して
$\varphi_{ki}(A_k) = \varphi_{ji}(A_j)$ となるような $j \geq
i$ が存在する。$k \geq j$ に対して $\varphi_{ki}(E_k) = \varphi_{ji}(E_j)$ も成り立つと主張する。
$k \geq j$ に対しては常に $\varphi_{ki}(E_k) \subset \varphi_{ji}(E_j)$ である。
逆の包含を示すため $e_j \in E_j$ をとり、
$\varphi_{ki}(x_k) = \varphi_{ji}(e_j)$ を満たす $x_k \in E_k$ を見つける。
任意の元 $e'_k \in E_k$ をとり、$e'_j = \varphi_{kj}(e'_k)$ とおく。
このとき $g_j(e_j - e'_j) = c_j - c_j = 0$ なので、
$e_j - e'_j = a_j \in A_j$ である。$\varphi_{ki}(A_k) =
\varphi_{ji}(A_j)$ だから、$\varphi_{ki}(a_k) = \varphi_{ji}(a_j)$ を満たす
$a_k \in A_k$ が存在する。したがって
$$
\varphi_{ki}(e'_k + a_k) = \varphi_{ji}(e'_j) + \varphi_{ji}(a_j) =
\varphi_{ji}(e_j),
$$
となるので、$x_k = e'_k + a_k$ ととれる。
\end{proof}




\section{逆系}
\label{algebra-section-inverse-systems}

\noindent
多くの論文では（また、この節でも）、{\it 逆系}という語は、半順序集合
$(\mathbf{N}, \geq)$ 上の逆系を表すために用いられる。
この節ではそのような系を簡単に論じる。この内容は Homology, Section
\ref{homology-section-inverse-systems} でより広く論じる。
環 $R$ と、図示された写像をもつ $R$-加群の列
$$
M_1 \xleftarrow{\varphi_2} M_2 \xleftarrow{\varphi_3} M_3 \leftarrow \ldots
$$
が与えられているとする。連続する写像を合成することにより、$i \geq i'$ のとき写像
$\varphi_{ii'} : M_i \to M_{i'}$ を得て、さらに $i \geq i' \geq i''$ のとき
$\varphi_{ii''} = \varphi_{i'i''} \circ \varphi_{i i'}$ が成り立つ。逆に、写像の系
$\varphi_{ii'}$ が与えられれば、$\varphi_i = \varphi_{i(i-1)}$ とおくことで
上に表示した写像を復元できる。この場合
$$
\lim M_i
=
\{(x_i) \in \prod M_i \mid \varphi_i(x_i) = x_{i - 1}, \ i = 2, 3, \ldots\}
$$
である。Categories, Section \ref{categories-section-limit-sets} と比較せよ。
Homology, Section \ref{homology-section-inverse-systems} で説明するように、
これは実際、Categories, Section \ref{categories-section-limits} で定義される
$R$-加群の圏における極限である。

\begin{lemma}
\label{algebra-lemma-Mittag-Leffler}
$R$ を環とする。$i \geq 1$ に対して、$R$-加群の短完全列
$0 \to K_i \to L_i \to M_i \to 0$ が、短完全列の間の写像
$$
\xymatrix{
0 \ar[r] &
K_i \ar[r] &
L_i \ar[r] &
M_i \ar[r] &
0 \\
0 \ar[r] &
K_{i + 1} \ar[r] \ar[u] &
L_{i + 1} \ar[r] \ar[u] &
M_{i + 1} \ar[r] \ar[u] &
0}
$$
に組み込まれているとする。すべての $i$ に対してある $c = c(i) \geq i$ が存在し、
すべての $j \geq c$ に対して
$\Im(K_c \to K_i) = \Im(K_j \to K_i)$ となるならば、列
$$
0 \to \lim K_i \to \lim L_i \to \lim M_i \to 0
$$
は完全である。
\end{lemma}

\begin{proof}
これは、より一般的な補題 \ref{algebra-lemma-ML-exact-sequence} の特別な場合である。
\end{proof}



\section{Mittag-Leffler 加群}
\label{algebra-section-mittag-leffler-modules}

\noindent
Mittag-Leffler 加群とは、（ごく大まかにいえば）その双対が Mittag-Leffler 系となるような
有向極限として書ける加群である。正確な定義を与えるには、少し準備が必要である。

\begin{definition}
\label{algebra-definition-ML-inductive-system}
$(M_i, f_{ij})$ を $R$-加群の有向系とする。各 $M_i$ が有限表示 $R$-加群であり、
任意の $R$-加群 $N$ に対して逆系
$$
(\Hom_R(M_i, N), \Hom_R(f_{ij}, N))
$$
が Mittag-Leffler であるとき、$(M_i, f_{ij})$ を
{\it 加群の Mittag-Leffler 有向系}という。
\end{definition}

\noindent
加群の Mittag-Leffler 有向系の余極限となる $R$-加群を特徴付ける。

\begin{definition}
\label{algebra-definition-domination}
$f: M \to N$ と $g: M \to M'$ を $R$-加群の写像とする。任意の $R$-加群 $Q$ に対して
$\Ker(f \otimes_R \text{id}_Q) \subset \Ker(g \otimes_R \text{id}_Q)$
であるとき、$g$ は $f$ を {\it 支配する}という。
\end{definition}

\noindent
この条件は、有限表示加群について確かめれば十分である。

\begin{lemma}
\label{algebra-lemma-domination-fp}
$f: M \to N$ と $g: M \to M'$ を $R$-加群の写像とする。$g$ が $f$ を支配するための
必要十分条件は、任意の有限表示 $R$-加群 $Q$ に対して
$\Ker(f \otimes_R \text{id}_Q) \subset \Ker(g \otimes_R
\text{id}_Q)$ となることである。
\end{lemma}

\begin{proof}
すべての有限表示加群 $Q$ に対して
$\Ker(f \otimes_R \text{id}_Q) \subset \Ker(g \otimes_R
\text{id}_Q)$ であると仮定する。$Q$ が任意の加群ならば、有限表示加群 $Q_i$ の
有向系の余極限として $Q = \colim_{i \in I} Q_i$ と書く。このとき、すべての $i$ に対して
$\Ker(f \otimes_R \text{id}_{Q_i}) \subset \Ker(g \otimes_R \text{id}_{Q_i})$ である。
有向余極限をとる操作は完全であり、テンソル積と可換なので、
$\Ker(f \otimes_R \text{id}_Q) \subset \Ker(g
\otimes_R \text{id}_Q)$ が従う。
\end{proof}

\begin{lemma}
\label{algebra-lemma-domination-universally-injective}
$f : M \to N$ と $g : M \to M'$ を $R$-加群の写像とする。$f$ と $g$ の押し出し
$$
\xymatrix{
M  \ar[r]_f \ar[d]_g & N \ar[d]^{g'} \\
M' \ar[r]^{f'} & N'
}
$$
を考える。このとき、$g$ が $f$ を支配することと、$f'$ が純単射であることは同値である。
\end{lemma}

\begin{proof}
$N'$ は、$M' \oplus N$ を、$x \in M$ に対する元 $(g(x), -f(x))$ からなる部分加群で
割ったものであることを思い出そう。$N'$ の構成から短完全列
$$
0 \to \Ker(f) \cap \Ker(g) \to \Ker(f) \to \Ker(f')
\to 0.
$$
を得る。テンソル積をとる操作は押し出しをとる操作と可換なので、任意の $R$-加群 $Q$ に対して
このような短完全列
$$
0 \to \Ker(f \otimes \text{id}_Q ) \cap \Ker(g \otimes
\text{id}_Q) \to \Ker(f \otimes \text{id}_Q)
\to \Ker(f' \otimes \text{id}_Q) \to 0
$$
を得る。したがって、$f'$ が純単射であるための必要十分条件は、任意の $Q$ に対して
$\Ker(f \otimes \text{id}_Q ) \subset \Ker(g \otimes
\text{id}_Q)$ となることであり、これは $g$ が $f$ を支配することと同値である。
\end{proof}

\noindent
上の支配の定義は、次の補題が示すように、写像の支配という通常の概念と関連することがある。

\begin{lemma}
\label{algebra-lemma-domination}
$f: M \to N$ と $g: M \to M'$ を $R$-加群の写像とする。
$\Coker(f)$ が有限表示であると仮定する。このとき $g$ が $f$ を支配するための必要十分条件は、
$g$ が $f$ を経由して分解すること、すなわち $g = h \circ f$ を満たす加群写像
$h: N \to M'$ が存在することである。
\end{lemma}

\begin{proof}
補題 \ref{algebra-lemma-domination-universally-injective} の主張にある $f$ と $g$ の押し出しを考える。
押し出しの構成から $\Coker(f') = \Coker(f)$ が従うので、$\Coker(f')$ は有限表示である。
補題 \ref{algebra-lemma-universally-exact-split} により、$f'$ が純単射であるための必要十分条件は、
$$
0 \to M' \xrightarrow{f'} N' \to \Coker(f') \to 0
$$
が分裂することである。これは $h' \circ f' = \text{id}_{M'}$ を満たす写像
$h' : N' \to M'$ が存在することと同値である。押し出しの普遍性により、このような $h'$ が
存在することは、$g$ が $f$ を経由して分解することと同値である。
\end{proof}


\begin{proposition}
\label{algebra-proposition-ML-characterization}
$M$ を $R$-加群とする。$(M_i, f_{ij})$ を $I$ で添字付けられた有限表示 $R$-加群の
有向系で、$M = \colim M_i$ となるものとする。$f_i:
M_i \to M$ を標準写像とする。次は同値である：
\begin{enumerate}
\item 任意の有限表示 $R$-加群 $P$ と加群写像 $f: P
\to M$ に対して、有限表示 $R$-加群 $Q$ と加群写像 $g: P \to Q$ が存在し、
$g$ と $f$ は互いに支配する。すなわち、任意の $R$-加群 $N$ に対して
$\Ker(f \otimes_R \text{id}_N) = \Ker(g \otimes_R \text{id}_N)$ である。
\item 各 $i \in I$ に対して、$f_{ij}: M_i
\to M_j$ が $f_i: M_i \to M$ を支配するような $j \geq i$ が存在する。
\item 各 $i \in I$ に対して、すべての $k \geq
i$ について $f_{ij}: M_i \to M_j$ が $f_{ik}: M_i \to M_k$ を経由して
分解するような $j \geq i$ が存在する。
\item 任意の $R$-加群 $N$ に対して、逆系
$(\Hom_R(M_i, N), \Hom_R(f_{ij}, N))$ は Mittag-Leffler である。
\item $N = \prod_{s \in I} M_s$ に対して、逆系
$(\Hom_R(M_i, N), \Hom_R(f_{ij}, N))$ は Mittag-Leffler である。
\end{enumerate}
\end{proposition}
\begin{proof}
まず (1) と (2) の同値性を証明する。(1) が成り立つとし、$i
\in I$ とする。写像 $f_i: M_i \to M$ に対応して、(1) のような
$g: M_i \to Q$ を選べる。$M_i$ と $Q$ は有限表示なので、$\Coker(g)$ も有限表示である。
補題 \ref{algebra-lemma-domination} により、$f_i : M_i
\to M$ は $g: M_i \to Q$ を経由して分解する。すなわち、ある $h: Q \to M$ に対して
$f_i = h \circ g$ である。さらに $Q$ は有限表示なので、ある $j \geq i$ に対して
$h$ は $M_j \to M$ を経由して分解する。すなわち、ある $h': Q \to M_j$ に対して
$h = f_j \circ h'$ である。以上より可換図式
$$
\xymatrix{
  &  M  & \\
M_i \ar[dr]_g \ar[ur]^{f_i} \ar[rr]^{f_{ij}} &
& M_j \ar[ul]_{f_j} \\
  & Q  \ar[ur]_{h'} &
}
$$
を得る。したがって $f_{ij}$ は $g$ を支配する。一方、$g$ は $f_i$ を支配するので、
$f_{ij}$ は $f_i$ を支配する。

\medskip\noindent
逆に (2) が成り立つとする。$P$ を有限表示加群とし、$f: P
\to M$ を加群写像とする。このとき、ある $i \in I$ に対して $f$ は
$f_i: M_i \to M$ を経由して分解する。すなわち、ある $g': P \to M_i$ に対して
$f = f_i \circ g'$ である。(2) によって、$f_{ij}$ が $f_i$ を支配するような
$j \geq i$ を選ぶ。可換図式
$$
\xymatrix{
P \ar[d]_{g'} \ar[r]^{f}           & M  \\
M_i \ar[ur]^{f_i} \ar[r]_{f_{ij}} & M_j \ar[u]_{f_j}
}
$$
を得る。この図式と $f_{ij}$ が $f_i$ を支配することから、$f$ と
$f_{ij} \circ g'$ は互いに支配する。したがって $g = f_{ij} \circ g' :
P \to M_j$ ととればよい。

\medskip\noindent
次に、(2) と (3) が同値であることを証明する。$i \in I$ とする。
$f_i$ は $f_{ik}$ を経由して分解するので、$k \geq i$ に対して $f_i$ が
$f_{ik}$ を支配することは常に成り立つ。(2) が成り立つなら、$f_{ij}$ が $f_i$ を
支配するような $j \geq i$ を選ぶ。このとき支配関係の推移性により、$k \geq i$ に対して
$f_{ij}$ は $f_{ik}$ を支配する。すべての $M_i$ は有限表示なので、$k \geq i$ に対して
$\Coker(f_{ik})$ は有限表示である。補題 \ref{algebra-lemma-domination} により、すべての
$k \geq i$ に対して $f_{ij}$ は $f_{ik}$ を経由して分解する。よって (2) は (3) を含意する。
一方、(3) が成り立つなら、任意の $R$-加群 $N$ に対して、$k \geq i$ では
$f_{ij} \otimes_R \text{id}_N$ は $f_{ik}
\otimes_R \text{id}_N$ を経由して分解する。したがって $k
\geq i$ に対して
$\Ker(f_{ik} \otimes_R \text{id}_N) \subset \Ker(f_{ij} \otimes_R \text{id}_N)$
である。しかし $\Ker(f_i \otimes_R \text{id}_N: M_i \otimes_R N
\to M \otimes_R N)$ は、$k \geq i$ における
$\Ker(f_{ik} \otimes_R \text{id}_N)$ の合併である。したがって任意の $R$-加群 $N$ に対して
$\Ker(f_i \otimes_R \text{id}_N) \subset \Ker(f_{ij} \otimes_R \text{id}_N)$ であり、
これは定義により $f_{ij}$ が $f_i$ を支配することを意味する。

\medskip\noindent
(3) が (4) を含意し、(4) が (5) を含意することは自明である。(5) が (3) を含意することを示す。
$N = \prod_{s \in I} M_s$ とおく。(5) が成り立つなら、$i \in I$ が与えられたとき、
すべての $k \geq j$ に対して
$$
\Im( \Hom(M_j, N) \to  \Hom(M_i, N)) =
\Im( \Hom(M_k, N) \to  \Hom(M_i, N))
$$
となるような $j \geq i$ を選ぶ。$s \in I$ にわたる積を $\Hom$ の外へ出し、
積の各成分上の写像を見ると、これは、すべての $k \geq j$ および $s \in I$ に対して
$$
\Im( \Hom(M_j, M_s) \to  \Hom(M_i, M_s)) =
\Im( \Hom(M_k, M_s) \to   \Hom(M_i, M_s))
$$
であることをいう。$s = j$ とすれば、すべての $k \geq j$ に対して
$$
\Im( \Hom(M_j, M_j) \to  \Hom(M_i, M_j)) =
\Im( \Hom(M_k, M_j) \to   \Hom(M_i, M_j))
$$
を得る。$f_{ij}$ は $\Hom(M_j, M_j) \to  \Hom(M_i, M_j)$ による
$\text{id} \in  \Hom(M_j, M_j)$ の像なので、任意の $k \geq j$ に対して、
$f_{ij} = h \circ f_{ik}$ を満たす $h \in \Hom(M_k, M_j)$ が存在することが分かる。
$j \geq k$ なら $h = f_{kj}$ ととれる。したがって (3) が成り立つ。
\end{proof}

\begin{definition}
\label{algebra-definition-mittag-leffler-module}
$M$ を $R$-加群とする。命題 \ref{algebra-proposition-ML-characterization} の同値な条件が
成り立つとき、$M$ は {\it Mittag-Leffler} であるという。
\end{definition}

\noindent
とくに有限表示加群は Mittag-Leffler である。

\begin{remark}
\label{algebra-remark-flat-ML}
$M$ を平坦 $R$-加群とする。Lazard の定理（定理 \ref{algebra-theorem-lazard}）により、
$M_i$ が有限自由 $R$-加群であるような有向系 $(M_i, f_{ij})$ の余極限として
$M = \colim M_i$ と書ける。$M$ が Mittag-Leffler であるためには、双対の逆系
$(\Hom_R(M_i, R), \Hom_R(f_{ij}, R))$ が Mittag-Leffler であれば十分である。
これは、命題 \ref{algebra-proposition-ML-characterization} の判定条件 (4) と、有限自由
$R$-加群 $F$ および任意の $R$-加群 $N$ に対して関手的同型
$\Hom_R(F, R) \otimes_R N \cong \Hom_R(F, N)$ が存在することから従う。
\end{remark}

\begin{lemma}
\label{algebra-lemma-tensor-ML-modules}
$R$ が環で、$M$、$N$ が $R$ 上の Mittag-Leffler 加群ならば、
$M \otimes_R N$ は Mittag-Leffler 加群である。
\end{lemma}

\begin{proof}
$M = \colim_{i \in I} M_i$ および $N = \colim_{j \in J} N_j$ と、
有限表示 $R$-加群の有向余極限として書く。遷移写像を
$f_{ii'} : M_i \to M_{i'}$ および $g_{jj'} : N_j \to N_{j'}$ と記す。このとき
$M_i \otimes_R N_j$ は有限表示 $R$-加群であり（補題 \ref{algebra-lemma-tensor-finiteness} を参照）、
$M \otimes_R N = \colim_{(i, j) \in I \times J} M_i \otimes_R M_j$ である。
% [P01-E0452] Frozen second tensor factor M_j is retained.
$(i, j) \in I \times J$ をとる。Mittag-Leffler 加群の定義により、両方の系に対して
命題 \ref{algebra-proposition-ML-characterization} (3) が成り立つ。言い換えると、
任意に $i'' \geq i$ と $j'' \geq j$ を選ぶたびに
$f_{ii'} = a \circ f_{ii''}$ および $g_{jj'} = b \circ g_{jj''}$ となる写像
$a : M_{i''} \to M_{i'}$ と $b : M_{j''} \to M_{j'}$ が存在するような
$i' \geq i$ と $j' \geq j$ が存在する。
% [P01-E0453] Frozen M-indexed type of b is retained.
このとき、$a \otimes b : M_{i''} \otimes_R N_{j''} \to M_{i'} \otimes_R N_{j'}$ が
系 $(M_i \otimes_R N_j, f_{ii'} \otimes g_{jj'})$ に対して同じ役割を果たすことは明らかである。
したがって命題 \ref{algebra-proposition-ML-characterization} (3) の特徴付けにより、
$M \otimes_R N$ は Mittag-Leffler である。
\end{proof}

\begin{lemma}
\label{algebra-lemma-ML-also}
$R$ を環、$M$ を $R$-加群とする。このとき $M$ が Mittag-Leffler であるための
必要十分条件は、任意の有限自由 $R$-加群 $F$ と加群写像 $f: F \to M$ に対して、
有限表示 $R$-加群 $Q$ と加群写像 $g : F \to Q$ が存在し、$g$ と $f$ が互いに支配すること、
すなわち任意の $R$-加群 $N$ に対して
$\Ker(f \otimes_R \text{id}_N) = \Ker(g \otimes_R \text{id}_N)$ となることである。
\end{lemma}

\begin{proof}
この条件は命題 \ref{algebra-proposition-ML-characterization} の条件 (1) より明らかに弱いので、
Mittag-Leffler 加群はこの条件を満たす。
% [P01-E0454] The frozen source's malformed comparative "clear weaker" is rendered by its evident sense.
逆に、$M$ がこの条件を満たし、$f : P \to M$ が有限表示 $R$-加群 $P$ から $M$ への
$R$-加群写像であると仮定する。$F$ が有限自由 $R$-加群となる全射 $F \to P$ を選ぶ。
仮定により、$Q$ が有限表示 $R$-加群で、$F \to Q$ と $F \to M$ が互いに支配するような
写像 $F \to Q$ を見つけられる。とくに $F \to Q$ の核は $F \to P$ の核を含むので、
$F \to Q$ が合成 $F \to P \to Q$ に等しくなるような $R$-加群写像
$g : P \to Q$ を得る。$N$ を任意の $R$-加群とし、可換図式
$$
\xymatrix{
F \otimes_R N \ar[d] \ar[r] & Q \otimes_R N \\
P \otimes_R N \ar[ru] \ar[r] & M \otimes_R N
}
$$
を考える。仮定により、$F \otimes_R N \to Q \otimes_R N$ と
$F \otimes_R N \to M \otimes_R N$ の核は等しい。したがって
$F \otimes_R N \to P \otimes_R N$ は全射なので、
$P \otimes_R N \to Q \otimes_R N$ と
$P \otimes_R N \to M \otimes_R N$ の核も等しい。
\end{proof}

\begin{lemma}
\label{algebra-lemma-restrict-ML-modules}
$R \to S$ を有限かつ有限表示の環写像とする。$M$ を $S$-加群とする。
$M$ が $S$ 上の Mittag-Leffler 加群ならば、$M$ は $R$ 上の Mittag-Leffler 加群である。
\end{lemma}

\begin{proof}
$M$ が $S$ 上の Mittag-Leffler 加群であると仮定する。有限表示 $S$-加群 $M_i$ の
有向余極限として $M = \colim M_i$ と書く。$M$ は $S$ 上 Mittag-Leffler なので、
各 $i$ に対して、すべての $k \geq j$ について分解
$f_{ij} = h \circ f_{ik}$ が存在するような添字 $j \geq i$ が存在する
（$h$ は $i$、$j$ の選択、および $k$ に依存する）。補題
\ref{algebra-lemma-finite-finitely-presented-extension} により、加群 $M_i$ は $R$-加群としても
有限表示であることに注意する。さらに、写像 $f_{ij}, f_{ik}, h$ はすべて $R$-加群の写像である。
したがって、系 $(M_i, f_{ij})$ を $R$-加群の系とみても同じ条件を満たすことが分かる。
ゆえに $M$ は $R$-加群として Mittag-Leffler である。
\end{proof}

\begin{lemma}
\label{algebra-lemma-mod-ideal-ML-modules}
$R$ を環とする。ある有限生成イデアル $I$ に対して $S = R/I$ とし、$M$ を $S$-加群とする。
このとき $M$ が $R$ 上の Mittag-Leffler 加群であることと、
$M$ が $S$ 上の Mittag-Leffler 加群であることは同値である。
\end{lemma}

\begin{proof}
一方の含意は補題 \ref{algebra-lemma-restrict-ML-modules} から従う。
他方を証明するため、$M$ が $R$-加群として Mittag-Leffler であると仮定する。
有限表示 $S$-加群の有向余極限として $M = \colim M_i$ と書く。$I$ は有限生成なので、
環 $S$ は $R$-代数として有限かつ有限表示であり、したがって加群 $M_i$ は
$R$-加群として有限表示である。補題 \ref{algebra-lemma-finite-finitely-presented-extension} を参照せよ。
次に $N$ を任意の $S$-加群とする。$R \to S$ は全射なので、各 $i$ に対して
$\Hom_R(M_i, N) = \Hom_S(M_i, N)$ であることに注意する。したがって、逆系
$(\Hom_R(M_i, N))_i$ が Mittag-Leffler 条件を満たすことから、系
$(\Hom_S(M_i, N))_i$ も Mittag-Leffler 条件を満たすことが従う。
ゆえに定義により、$M$ は $S$ 上 Mittag-Leffler である。
\end{proof}

\begin{remark}
\label{algebra-remark-go-up-ML-modules}
$R \to S$ を有限かつ有限表示の環写像とする。$M$ を、$R$-加群として Mittag-Leffler である
$S$-加群とする。一般には、$M$ が $S$-加群として Mittag-Leffler になるとは限らない。
たとえば $S$ が $R$ 上の双対数の環、すなわち $S = R \oplus R\epsilon$ かつ
$\epsilon^2 = 0$ であると仮定する。このとき $S$-加群とは、二乗零な $R$-線形自己準同型
$\epsilon : M \to M$ を備えた $R$-加群 $M$ である。
ここで $M_0$ を、Mittag-Leffler でない $R$-加群とする。$F_1$ と $F_0$ が自由
$R$-加群となる表示 $F_1 \xrightarrow{u} F_0 \to M_0 \to 0$ を選ぶ。
$M = F_1 \oplus F_0$ とおき、
$$
\epsilon =
\left(
\begin{matrix}
0 & 0 \\
u & 0
\end{matrix}
\right) : M \longrightarrow M.
$$
とする。このとき $M/\epsilon M \cong F_1 \oplus M_0$ は
$R = S/\epsilon S$ 上 Mittag-Leffler でなく、したがって $S$ 上も Mittag-Leffler でない
（補題 \ref{algebra-lemma-mod-ideal-ML-modules} を参照）。一方、
$M/\epsilon M = M \otimes_S S/\epsilon S$ であり、補題
\ref{algebra-lemma-tensor-ML-modules} により、$M$ が $S$ 上 Mittag-Leffler ならば
これは Mittag-Leffler となるはずである。
\end{remark}


\section{直積とテンソル積の交換}
\label{algebra-section-products-tensor}

\noindent
$M$ を $R$-加群、$(Q_{\alpha})_{\alpha \in A}$ を $R$-加群の族とする。このとき、純テンソル上で
$x \otimes (q_{\alpha}) \mapsto (x \otimes q_{\alpha})$ によって与えられる標準写像
$M \otimes_R \left( \prod_{\alpha
\in A} Q_{\alpha} \right) \to \prod_{\alpha \in A} ( M \otimes_R
Q_{\alpha})$ がある。次の例が示すように、この写像は単射とも全射とも限らない。

\begin{example}
\label{algebra-example-Q-not-ML}
$R = \mathbf{Z}$、$M = \mathbf{Q}$ とし、$n \geq 1$ に対する族
$Q_n = \mathbf{Z}/n$ を考える。このとき $\prod_n (M \otimes Q_n) = 0$ である。
しかし単射 $\mathbf{Z} \to \prod_n Q_n$ を $M$ とテンソルすることで得られる単射
$\mathbf{Q} \to M \otimes (\prod_n Q_n)$ が存在するので、
$M \otimes (\prod_n Q_n)$ は非零である。したがって
$M \otimes (\prod_n Q_n) \to \prod_n (M \otimes Q_n)$ は単射でない。

\medskip\noindent
他方、再び $R = \mathbf{Z}$、$M = \mathbf{Q}$ とし、$n \geq 1$ に対して
$Q_n = \mathbf{Z}$ とする。$M \otimes (\prod_n Q_n)
\to \prod_n (M \otimes Q_n) = \prod_n M$ の像は、$a_n \in \mathbf{Z}$ かつ $m$ が
ある非零整数であるような形 $(a_n/m)_{n \geq 1}$ の列からちょうどなる。
したがってこの写像は全射でない。
\end{example}

\noindent
以下では、$(Q_{\alpha})_{\alpha \in A}$ のすべての選び方に対して写像
$M \otimes_R \left( \prod_{\alpha} Q_{\alpha} \right) \to \prod_{\alpha}
(M \otimes_R Q_{\alpha})$ が全射、全単射、または単射となるために $M$ に必要な条件を
正確に決定する。単射となる加群はちょうど Mittag-Leffler 加群であることが分かるため、
これは重要である（命題 \ref{algebra-proposition-ML-tensor}）。以下、$M$ が $R$-加群で $A$ が集合なら、
$M^A$ で積 $\prod_{\alpha \in A} M$ を表す。

\begin{proposition}
\label{algebra-proposition-fg-tensor}
$M$ を $R$-加群とする。次は同値である：
\begin{enumerate}
\item $M$ は有限生成である。
\item $R$-加群の任意の族 $(Q_{\alpha})_{\alpha \in A}$ に対して、標準写像
$M \otimes_R \left( \prod_{\alpha} Q_{\alpha} \right)
\to \prod_{\alpha} (M \otimes_R Q_{\alpha})$ は全射である。
\item 任意の $R$-加群 $Q$ と任意の集合 $A$ に対して、標準写像
$M \otimes_R Q^{A} \to (M \otimes_R Q)^{A}$ は全射である。
\item 任意の集合 $A$ に対して、標準写像 $M \otimes_R R^{A} \to
M^{A}$ は全射である。
\end{enumerate}
\end{proposition}

\begin{proof}
まず (1) が (2) を含意することを証明する。全射 $R^n \to M$ を選び、可換図式
$$
\xymatrix{
R^n \otimes_R (\prod_{\alpha} Q_{\alpha})  \ar[r]^{\cong} \ar[d] &
\prod_{\alpha} (R^n \otimes_R Q_{\alpha}) \ar[d] \\
M \otimes_R (\prod_{\alpha} Q_{\alpha})  \ar[r] & \prod_{\alpha} ( M
\otimes_R Q_{\alpha}).
}
$$
を考える。上の矢印は同型で、縦の矢印は全射である。したがって下の矢印は全射である。

\medskip\noindent
(2) が (3) を、(3) が (4) を含意することは明らかなので、(4) が (1) を含意することを
証明すればよい。実際、(1) が成り立つためには、$M^M$ の元
$d = (x)_{x \in M}$ が写像 $f: M \otimes_R R^{M} \to M^M$ の像に含まれれば十分である。
この場合、ある $x_i \in M$ と $a_i \in R^M$ に対して
$d = \sum_{i = 1}^{n} f(x_i \otimes a_i)$ である。
$x \in M$ に対して $p_x: M^M \to M$ を $x$ 番目の因子への射影と書けば、
$$
x = p_x(d) = \sum\nolimits_{i = 1}^{n} p_x(f(x_i \otimes a_i)) =
\sum\nolimits_{i=1}^{n} p_x(a_i) x_i.
$$
% [P01-E0455] Frozen use of the M^M projection p_x on a_i in R^M is retained.
したがって $x_1, \ldots, x_n$ は $M$ を生成する。
\end{proof}

\begin{proposition}
\label{algebra-proposition-fp-tensor}
$M$ を $R$-加群とする。次は同値である：
\begin{enumerate}
\item $M$ は有限表示である。
\item $R$-加群の任意の族 $(Q_{\alpha})_{\alpha \in A}$ に対して、標準写像
$M \otimes_R \left( \prod_{\alpha} Q_{\alpha} \right)
\to \prod_{\alpha} (M \otimes_R Q_{\alpha})$ は全単射である。
\item 任意の $R$-加群 $Q$ と任意の集合 $A$ に対して、標準写像
$M \otimes_R Q^{A} \to (M \otimes_R Q)^{A}$ は全単射である。
\item 任意の集合 $A$ に対して、標準写像 $M \otimes_R R^{A} \to
M^{A}$ は全単射である。
\end{enumerate}
\end{proposition}

\begin{proof}
まず (1) が (2) を含意することを証明する。表示 $R^m \to R^n
\to M$ を選び、可換図式
$$
\xymatrix{
R^m \otimes_R (\prod_{\alpha} Q_{\alpha}) \ar[r] \ar[d]^{\cong} & R^n
\otimes_R (\prod_{\alpha} Q_{\alpha}) \ar[r] \ar[d]^{\cong} & M \otimes_R
(\prod_{\alpha} Q_{\alpha}) \ar[r] \ar[d] & 0 \\
\prod_{\alpha} (R^m \otimes_R Q_{\alpha}) \ar[r] & \prod_{\alpha} (R^n
\otimes_R Q_{\alpha}) \ar[r] & \prod_{\alpha} (M \otimes_R Q_{\alpha})
\ar[r] & 0.
}
$$
を考える。最初の二つの縦の矢印は同型で、各行は完全である。これにより写像
$M \otimes_R (\prod_{\alpha} Q_{\alpha}) \to
\prod_{\alpha} ( M \otimes_R Q_{\alpha})$ は全射であり、図式追跡により単射でもある。
したがって (2) が成り立つ。

\medskip\noindent
(2) が (3) を、(3) が (4) を含意することは明らかなので、(4) が (1) を含意することを
証明すればよい。命題 \ref{algebra-proposition-fg-tensor} により、(4) が成り立つなら $M$ はすでに
有限生成であることが分かっている。そこで $F$ が有限自由となる全射 $F \to M$ を選ぶ。
$K$ をその核とする。$K$ が有限生成であることを示さなければならない。任意の集合 $A$ に対して、
可換図式
$$
\xymatrix{
& K \otimes_R R^A \ar[r] \ar[d]_{f_3} & F \otimes_R R^A \ar[r]
\ar[d]_{f_2}^{\cong} & M \otimes_R R^A \ar[r] \ar[d]_{f_1}^{\cong} & 0 \\
0 \ar[r] & K^A \ar[r] & F^A \ar[r] & M^A \ar[r] & 0 .
}
$$
がある。仮定により $f_1$ は同型であり、$F$ は有限自由なので $f_2$ も同型で、各行は完全である。
図式追跡により $f_3$ は全射であることが分かる。したがって命題
\ref{algebra-proposition-fg-tensor} により、$K$ は有限生成である。
\end{proof}

\noindent
次の命題のために以下の補題が必要である。

\begin{lemma}
\label{algebra-lemma-kernel-tensored-fp}
$M$ を $R$-加群、$P$ を有限表示 $R$-加群、$f: P \to M$ を写像とする。
$Q$ を $R$-加群とし、$x \in \Ker(P \otimes Q \to M \otimes Q)$ と仮定する。
このとき、有限表示 $R$-加群 $P'$ と写像 $f': P \to P'$ が存在し、$f$ は $f'$ を
経由して分解し、$x \in \Ker(P \otimes Q \to P' \otimes Q)$ となる。
\end{lemma}

\begin{proof}
$M$ を有限表示加群 $M_i$ の有向系の余極限として $M = \colim_{i \in I} M_i$ と書く。
$P$ は有限表示なので、写像 $f:
P \to M$ は、ある $j \in I$ に対して $M_j \to M$ を経由して分解する。
$Q$ とテンソルすることで可換図式
$$
\xymatrix{
& M_j \otimes Q \ar[dr] & \\
P \otimes Q \ar[ur] \ar[rr] & & M \otimes Q .
}
$$
を得る。$x$ の $M_j \otimes Q$ における像 $y$ は、
$M_j \otimes Q \to M \otimes Q$ の核に含まれる。
$M \otimes Q = \colim_{i \in I} (M_i
\otimes Q)$ なので、これはある $j' \geq j$ に対して $y$ が $M_{j'} \otimes Q$ の
$0$ に写ることを意味する。したがって $P' = M_{j'}$ とし、$f'$ を合成
$P \to M_j \to M_{j'}$ ととればよい。
\end{proof}

\begin{proposition}
\label{algebra-proposition-ML-tensor}
$M$ を $R$-加群とする。次は同値である：
\begin{enumerate}
\item $M$ は Mittag-Leffler である。
\item $R$-加群の任意の族 $(Q_{\alpha})_{\alpha \in A}$ に対して、標準写像
$M \otimes_R \left( \prod_{\alpha} Q_{\alpha} \right)
\to \prod_{\alpha} (M \otimes_R Q_{\alpha})$ は単射である。
\end{enumerate}
\end{proposition}

\begin{proof}
まず (1) が (2) を含意することを証明する。$M$ は Mittag-Leffler であるとし、
$x$ を $M \otimes_R (\prod_{\alpha} Q_{\alpha}) \to
\prod_{\alpha} (M \otimes_R Q_{\alpha})$ の核に属する元とする。
有限表示加群 $M_i$ の有向系の余極限として $M$ を
$M = \colim_{i \in I} M_i$ と書く。このとき
$M \otimes_R (\prod_{\alpha} Q_{\alpha})$ は
$M_i \otimes_R (\prod_{\alpha} Q_{\alpha})$ の余極限である。したがって $x$ はある元
$x_i \in M_i \otimes_R (\prod_{\alpha} Q_{\alpha})$ の像である。
ある $j \geq i$ に対して $x_i$ が $M_j \otimes_R (\prod_{\alpha} Q_{\alpha})$ の
$0$ に写ることを示さなければならない。$M$ は Mittag-Leffler なので、
$M_i \to M_j$ と $M_i \to M$ が互いに支配するような $j \geq i$ を選べる。
そこで可換図式
$$
\xymatrix{
M \otimes_R (\prod_{\alpha} Q_{\alpha}) \ar[r] & \prod_{\alpha} (M
\otimes_R Q_{\alpha}) \\
M_i \otimes_R (\prod_{\alpha} Q_{\alpha}) \ar[r]^{\cong} \ar[d] \ar[u] &
\prod_{\alpha} (M_i \otimes_R Q_{\alpha}) \ar[d] \ar[u] \\
M_j \otimes_R (\prod_{\alpha} Q_{\alpha}) \ar[r]^{\cong} & \prod_{\alpha}
(M_j \otimes_R Q_{\alpha})
}
$$
を考える。命題 \ref{algebra-proposition-fp-tensor} により、下の二つの横写像は同型である。
$x_i$ は $\prod_{\alpha} (M
\otimes_R Q_{\alpha})$ の $0$ に写るので、
$\prod_{\alpha} (M_i \otimes_R Q_{\alpha})$ におけるその像は、写像
$\prod_{\alpha} (M_i \otimes_R Q_{\alpha}) \to \prod_{\alpha} (M \otimes_R Q_{\alpha})$
の核に含まれる。しかし $j$ の選び方により、この核は
$\prod_{\alpha} (M_i \otimes_R Q_{\alpha})
\to \prod_{\alpha} (M_j \otimes_R Q_{\alpha})$ の核に等しい。
したがって $x_i$ は $\prod_{\alpha} (M_j \otimes_R
Q_{\alpha})$ の $0$ に写り、従って $M_j \otimes_R (\prod_{\alpha} Q_{\alpha})$ の
$0$ にも写る。

\medskip\noindent
今度は (2) が成り立つと仮定する。命題 \ref{algebra-proposition-ML-characterization} にある
Mittag-Leffler 性の定式化 (1) を $M$ が満たすことを証明する。
$f: P \to M$ を有限表示加群 $P$ から $M$ への写像とする。
有限表示 $R$-加群の同型類の代表元の集合 $B$ を選ぶ。$A$ を、$Q \in B$ かつ
$x \in \Ker(P \otimes Q \to M \otimes Q)$ であるような対 $(Q, x)$ の集合とする。
$\alpha = (Q, x) \in A$ に対して、$Q_{\alpha}$ で $Q$ を、$x_{\alpha}$ で $x$ を表す。
可換図式
$$
\xymatrix{
M \otimes_R (\prod_{\alpha} Q_{\alpha}) \ar[r] &
\prod_{\alpha} (M \otimes_R Q_{\alpha}) \\
P \otimes_R (\prod_{\alpha} Q_{\alpha}) \ar[r]^{\cong} \ar[u] &
\prod_{\alpha} (P \otimes_R Q_{\alpha}) \ar[u] .
}
$$
を考える。仮定により上の矢印は単射であり、命題 \ref{algebra-proposition-fp-tensor} により
下の矢印は同型である。この同型のもとで
$(x_{\alpha}) \in \prod_{\alpha} (P \otimes_R Q_{\alpha})$ に対応する元を
$x \in P \otimes_R (\prod_{\alpha} Q_{\alpha})$ とする。このとき、図式の上の矢印は
単射なので、$x \in \Ker( P \otimes_R (\prod_{\alpha} Q_{\alpha})
\to M \otimes_R (\prod_{\alpha} Q_{\alpha}))$ である。
補題 \ref{algebra-lemma-kernel-tensored-fp} により、有限表示加群 $P'$ と写像 $f': P \to P'$ を得て、
$f: P \to M$ は $f'$ を経由して分解し、
$x \in \Ker(P \otimes_R
(\prod_{\alpha} Q_{\alpha}) \to P' \otimes_R (\prod_{\alpha}
Q_{\alpha}))$ となる。可換図式
$$
\xymatrix{
P' \otimes_R (\prod_{\alpha} Q_{\alpha}) \ar[r]^{\cong} &
\prod_{\alpha} (P' \otimes_R Q_{\alpha}) \\
P \otimes_R (\prod_{\alpha} Q_{\alpha}) \ar[r]^{\cong} \ar[u] &
\prod_{\alpha} (P \otimes_R Q_{\alpha}) \ar[u] .
}
$$
があり、命題 \ref{algebra-proposition-fp-tensor} により上と下の矢印はともに同型である。
したがって $x$ は左の縦写像の核に含まれるので、$(x_{\alpha})$ は右の縦写像の核に含まれる。
これは、すべての $\alpha \in A$ に対して
$x_{\alpha} \in \Ker(P \otimes_R Q_{\alpha} \to P' \otimes_R
Q_{\alpha})$ となることを意味する。$A$ の定義により、これはすべての有限表示 $Q$ に対して
$\Ker(P \otimes_R Q \to P' \otimes_R Q) \supset \Ker(P \otimes_R
Q \to M \otimes_R Q)$ となることを意味する。さらに $f: P
\to M$ は $f': P \to P'$ を経由して分解するので、実際には等号が成り立つ。
補題 \ref{algebra-lemma-domination-fp} により、$f$ と $f'$ は互いに支配する。
\end{proof}

\begin{lemma}
\label{algebra-lemma-minimal-contains}
$M$ を $R$ 上の平坦 Mittag-Leffler 加群とする。$F$ を $R$-加群とし、
$x \in F \otimes_R M$ とする。このとき、$x \in F' \otimes_R M$ となる最小の部分加群
$F' \subset F$ が存在する。さらに、$F'$ は有限 $R$-加群である。
\end{lemma}

\begin{proof}
$M$ は平坦なので、$F' \subset F$ が部分加群ならば
$F' \otimes_R M \subset F \otimes_R M$ であり、従って主張には意味がある。
$I = \{F' \subset F \mid x \in F' \otimes_R M\}$ とおき、$i \in I$ に対して
対応する部分加群を $F_i \subset F$ と記す。このとき $x$ は写像
$$
F \otimes_R M \longrightarrow \prod (F/F_i \otimes_R M)
$$
によって零に写り、したがって命題 \ref{algebra-proposition-ML-tensor} により、$x$ は写像
$$
F \otimes_R M \longrightarrow \left(\prod F/F_i\right) \otimes_R M
$$
によって零に写る。$M$ は平坦なので、この矢印の核は
$(\bigcap F_i) \otimes_R M$ であり、これは $F' = \bigcap F_i$ であることを証明する。
$F'$ が有限加群であることを見るため、ある $f_j \in F'$ と $m_j \in M$ に対して
$x = \sum_{j = 1, \ldots, m} f_j \otimes m_j$ と仮定する。このとき、
$F'' \subset F'$ を $f_1, \ldots, f_m$ で生成される部分加群とすれば、
$x \in F'' \otimes_R M$ である。もちろん $F'' = F'$ であり、最後の主張が従う。
\end{proof}

\begin{lemma}
\label{algebra-lemma-pure-submodule-ML}
$0 \to M_1 \to M_2 \to M_3 \to 0$ を $R$-加群の純完全列とする。このとき：
\begin{enumerate}
\item $M_2$ が Mittag-Leffler ならば、$M_1$ は Mittag-Leffler である。
\item $M_1$ と $M_3$ が Mittag-Leffler ならば、$M_2$ は Mittag-Leffler である。
\end{enumerate}
\end{lemma}

\begin{proof}
$R$-加群の任意の族 $(Q_{\alpha})_{\alpha \in A}$ に対して、完全な行をもつ可換図式
$$
\xymatrix{
0 \ar[r] & M_1 \otimes_R (\prod_{\alpha} Q_{\alpha}) \ar[r] \ar[d] & M_2
\otimes_R (\prod_{\alpha} Q_{\alpha}) \ar[r] \ar[d] & M_3 \otimes_R
(\prod_{\alpha} Q_{\alpha}) \ar[r] \ar[d] & 0 \\
0 \ar[r] & \prod_{\alpha}(M_1 \otimes Q_{\alpha}) \ar[r] & \prod_{\alpha}(M_2
\otimes Q_{\alpha}) \ar[r] & \prod_{\alpha}(M_3 \otimes Q_{\alpha})\ar[r] & 0
}
$$
がある。従って (1) と (2) は命題 \ref{algebra-proposition-ML-tensor} から従う。
\end{proof}

\begin{lemma}
\label{algebra-lemma-quotient-module-ML}
$M_1 \to M_2 \to M_3 \to 0$ を $R$-加群の完全列とする。
$M_1$ が有限生成で、$M_2$ が Mittag-Leffler ならば、$M_3$ は Mittag-Leffler である。
\end{lemma}

\begin{proof}
$R$-加群の任意の族 $(Q_{\alpha})_{\alpha \in A}$ に対して、テンソル積は右完全なので、
完全な行をもつ可換図式
$$
\xymatrix{
M_1 \otimes_R (\prod_{\alpha} Q_{\alpha}) \ar[r] \ar[d] & M_2
\otimes_R (\prod_{\alpha} Q_{\alpha}) \ar[r] \ar[d] & M_3 \otimes_R
(\prod_{\alpha} Q_{\alpha}) \ar[r] \ar[d] & 0 \\
\prod_{\alpha}(M_1 \otimes Q_{\alpha}) \ar[r] & \prod_{\alpha}(M_2
\otimes Q_{\alpha}) \ar[r] & \prod_{\alpha}(M_3 \otimes Q_{\alpha})\ar[r] & 0
}
$$
がある。命題 \ref{algebra-proposition-fg-tensor} により、左の縦の矢印は全射である。
命題 \ref{algebra-proposition-ML-tensor} により、中央の縦の矢印は単射である。
図式追跡により、右の縦の矢印は単射であることが分かる。従って命題
\ref{algebra-proposition-ML-tensor} により $M_3$ は Mittag-Leffler である。
\end{proof}


\begin{lemma}
\label{algebra-lemma-colimit-universally-injective-ML}
$M = \colim M_i$ が、遷移写像が純単射である Mittag-Leffler $R$-加群 $M_i$ の
有向系の余極限ならば、$M$ は Mittag-Leffler である。
\end{lemma}

\begin{proof}
$(Q_{\alpha})_{\alpha \in A}$ を $R$-加群の族とする。
$M \otimes_R (\prod Q_\alpha) \to \prod M \otimes_R Q_\alpha$ が単射であることを
示さなければならず、各 $i$ に対して
$M_i \otimes_R (\prod Q_\alpha) \to \prod M_i \otimes_R Q_\alpha$ が単射であることは
命題 \ref{algebra-proposition-ML-tensor} から分かっている。$\otimes$ はフィルター余極限と可換なので、
$\prod M_i \otimes_R Q_\alpha \to \prod M \otimes_R Q_\alpha$ が単射であることを示せば
十分である。遷移写像が純単射であるという仮定により、写像
$M_i \otimes_R Q_\alpha \to M \otimes_R Q_\alpha$ はそれぞれ単射なので、これは明らかである。
\end{proof}

\begin{lemma}
\label{algebra-lemma-direct-sum-ML}
$M = \bigoplus_{i \in I} M_i$ が $R$-加群の直和ならば、$M$ が Mittag-Leffler であることと、
各 $M_i$ が Mittag-Leffler であることは同値である。
\end{lemma}

\begin{proof}
「のみならば」の向きは、補題 \ref{algebra-lemma-pure-submodule-ML} (1) と、分裂短完全列が
純完全であることから従う。逆は補題 \ref{algebra-lemma-colimit-universally-injective-ML} から従うが、
次のように直接論じることもできる。まず $I$ が有限ならば、補題
\ref{algebra-lemma-pure-submodule-ML} (2) から従うことに注意する。一般の $I$ に対し、
すべての $M_i$ が Mittag-Leffler ならば、命題 \ref{algebra-proposition-ML-characterization} の
条件 (1) を確かめることで $M$ についても同じことを証明する。
$f: P \to M$ を有限表示加群 $P$ からの写像とする。このとき $I$ のある有限部分集合
$I'$ に対して $f$ は
$P \xrightarrow{f'} \bigoplus_{i' \in I'} M_{i'} \hookrightarrow
\bigoplus_{i \in I} M_i$
と分解する。有限の場合により $\bigoplus_{i' \in I'} M_{i'}$ は Mittag-Leffler であり、
従って有限表示加群 $Q$ と写像 $g: P \to Q$ が存在して、$g$ と $f'$ は互いに支配する。
このとき $g$ と $f$ も互いに支配する。
\end{proof}

\begin{lemma}
\label{algebra-lemma-flat-ML-over-ML-ring}
$R \to S$ を環写像とし、$M$ を $S$-加群とする。$S$ が $R$-加群として Mittag-Leffler で、
$M$ が $S$-加群として平坦かつ Mittag-Leffler ならば、$M$ は $R$-加群として
Mittag-Leffler である。
\end{lemma}

\begin{proof}
これを命題 \ref{algebra-proposition-ML-tensor} の特徴付けから導く。実際、$Q_\alpha$ が
$R$-加群の族であると仮定し、合成
$$
\xymatrix{
M \otimes_R \prod_\alpha Q_\alpha =
M \otimes_S S \otimes_R \prod_\alpha Q_\alpha \ar[d] \\
M \otimes_S \prod_\alpha (S \otimes_R Q_\alpha) \ar[d] \\
\prod_\alpha (M \otimes_S S \otimes_R Q_\alpha) =
\prod_\alpha (M \otimes_R Q_\alpha)
}
$$
を考える。$M$ は $S$ 上平坦で、$S$ は $R$ 上 Mittag-Leffler なので最初の矢印は単射であり、
$M$ は $S$ 上 Mittag-Leffler なので二番目の矢印も単射である。
従って $M$ は $R$ 上 Mittag-Leffler である。
\end{proof}


\section{コヒーレント環}
\label{algebra-section-coherent}

\noindent
$\prod$ と $\otimes$ の交換に関する議論を用いて、どの環に対して平坦加群の積が平坦になるかを
決定する。それらは、いわゆるコヒーレント環であることが分かる。環付き空間上のコヒーレント
$\mathcal{O}_X$-加群という概念のほうが馴染み深いかもしれない。Modules, Section
\ref{modules-section-coherent} を参照せよ。

\begin{definition}
\label{algebra-definition-coherent}
$R$ を環、$M$ を $R$-加群とする。
\begin{enumerate}
\item 有限生成であり、$M$ の任意の有限生成部分加群が $R$ 上有限表示であるような $M$ を
{\it コヒーレント加群}という。
\item 自身上の加群としてコヒーレントである $R$ を
{\it コヒーレント環}という。
\end{enumerate}
\end{definition}

\noindent
したがって、環がコヒーレントであるための必要十分条件は、任意の有限生成イデアルが
加群として有限表示であることである。

\begin{example}
\label{algebra-example-valuation-ring-coherent}
付値環はコヒーレント環である。実際、任意の非零有限生成イデアルは単項であり
（補題 \ref{algebra-lemma-characterize-valuation-ring}）、付値環は整域なので自由であり、
従って有限表示である。
\end{example}

\noindent
コヒーレント加群の圏はアーベル圏である。

\begin{lemma}
\label{algebra-lemma-coherent}
$R$ を環とする。
\begin{enumerate}
\item コヒーレント加群の有限部分加群はコヒーレントである。
\item $\varphi : N \to M$ を有限加群からコヒーレント加群への準同型とする。このとき
$\Ker(\varphi)$ は有限、$\Im(\varphi)$ はコヒーレント、$\Coker(\varphi)$ はコヒーレントである。
\item $\varphi : N \to M$ をコヒーレント加群の準同型とする。このとき
$\Ker(\varphi)$ と $\Coker(\varphi)$ はコヒーレント加群である。
\item $R$-加群の短完全列 $0 \to M_1 \to M_2 \to M_3 \to 0$ が与えられたとき、
三つのうち二つがコヒーレントならば残りもコヒーレントである。
\end{enumerate}
\end{lemma}

\begin{proof}
最初の主張は定義から直ちに従う。

\medskip\noindent
$\varphi : N \to M$ が (2) の仮定を満たすとする。まず $\Im(\varphi)$ は有限であり、
従って (1) によりコヒーレントである。とくに $\Im(\varphi)$ は有限表示なので、完全列
$0 \to \Ker(\varphi) \to N \to \Im(\varphi) \to 0$ に補題
\ref{algebra-lemma-extension} を適用すると、$\Ker(\varphi)$ は有限であることが分かる。
$\Coker(\varphi)$ がコヒーレントであることを証明するため、
$E \subset \Coker(\varphi)$ を有限部分加群とし、$E'$ を $M$ におけるその逆像とする。
完全列 $0 \to \Im(\varphi) \to E' \to E \to 0$ と $\Ker(\varphi)$ が有限であることから、
補題 \ref{algebra-lemma-extension} により $E' \subset M$ は有限であると結論する。
% [P01-E0456] Frozen use of Ker(varphi), rather than the displayed Im(varphi), is retained.
従って $M$ はコヒーレントなので、これは有限表示である。同じ完全列から $E$ は有限表示である
ことも分かり、主張が従う。

\medskip\noindent
(3) は (1) と (2) から直ちに従う。

\medskip\noindent
$0 \to M_1 \xrightarrow{i} M_2 \xrightarrow{p} M_3 \to 0$ を (4) のような
$R$-加群の短完全列とする。残るのは、$M_1$ と $M_3$ がコヒーレントならば $M_2$ も
コヒーレントであることの証明である。補題 \ref{algebra-lemma-extension} により $M_2$ は有限である。
$N_2 \subset M_2$ を有限部分加群とする。
$N_3 = p(N_2) \subset M_3$ および $N_1 = i^{-1}(N_2) \subset M_1$ とおく。
完全列 $0 \to N_1 \to N_2 \to N_3 \to 0$ がある。明らかに $N_3$ は
（$N_2$ の商として）有限であり、従って（$M_3$ の有限部分加群として）有限表示である。
補題 \ref{algebra-lemma-extension} (5) により $N_1$ は有限であり、従って
（$M_1$ の有限部分加群として）有限表示である。補題 \ref{algebra-lemma-extension} (2) により、
$N_2$ は有限表示であると結論する。
\end{proof}

\begin{lemma}
\label{algebra-lemma-coherent-ring}
$R$ を環とする。$R$ がコヒーレントならば、加群がコヒーレントであることと
有限表示であることは同値である。
\end{lemma}

\begin{proof}
コヒーレント加群が有限表示であることは（任意の環上で）明らかである。
逆に $R$ がコヒーレントなら、$R^{\oplus n}$ はコヒーレントであり、
補題 \ref{algebra-lemma-coherent} により任意の写像 $R^{\oplus m} \to R^{\oplus n}$ の余核も
コヒーレントである。
\end{proof}

\begin{lemma}
\label{algebra-lemma-Noetherian-coherent}
ネーター環はコヒーレント環である。
\end{lemma}

\begin{proof}
補題 \ref{algebra-lemma-Noetherian-finite-type-is-finite-presentation} により、任意の有限
$R$-加群は有限表示である。とくに $R$ の任意のイデアルは有限表示である。
\end{proof}

\begin{proposition}
\label{algebra-proposition-characterize-coherent}
\begin{reference}
これは \cite[Theorem 2.1]{Chase} である。
\end{reference}
$R$ を環とする。次は同値である：
\begin{enumerate}
\item $R$ はコヒーレントである、
\item 平坦 $R$-加群の任意の積は平坦である、および
\item 任意の集合 $A$ に対して加群 $R^A$ は平坦である。
\end{enumerate}
\end{proposition}

\begin{proof}
$R$ がコヒーレントであると仮定し、$Q_\alpha$, $\alpha \in A$ を平坦 $R$-加群の集合とする。
補題 \ref{algebra-lemma-flat} により、$R$ の任意の有限生成イデアル $I$ に対して
$I \otimes_R \prod_\alpha Q_\alpha \to \prod Q_\alpha$ が単射であることを示さなければならない。
$R$ はコヒーレントなので、$I$ は有限表示 $R$-加群である。従って命題
\ref{algebra-proposition-fp-tensor} により
$I \otimes_R \prod_\alpha Q_\alpha = \prod I \otimes_R Q_\alpha$ である。
$Q_\alpha$ の平坦性により $I \otimes_R Q_\alpha \to Q_\alpha$ は単射なので、
所望の単射性が従う。

\medskip\noindent
含意 (2) $\Rightarrow$ (3) は自明である。

\medskip\noindent
任意の集合 $A$ に対して $R$-加群 $R^A$ が平坦であると仮定する。
$I$ を $R$ の有限生成イデアルとする。このとき仮定により
$I \otimes_R R^A \to R^A$ は単射である。命題 \ref{algebra-proposition-fg-tensor} と
$I$ の有限性により、その像は $I^A$ に等しい。従って任意の集合 $A$ に対して
$I \otimes_R R^A = I^A$ であり、命題 \ref{algebra-proposition-fp-tensor} により
$I$ は有限表示であると結論する。
\end{proof}







\section{Mittag-Leffler 加群の例と反例}
\label{algebra-section-examples}

\noindent
この節を Mittag-Leffler 加群のいくつかの例と反例で締めくくる。

\begin{example}
\label{algebra-example-ML}
Mittag-Leffler 加群。
\begin{enumerate}
\item 任意の有限表示加群は Mittag-Leffler である。これは、たとえば命題
\ref{algebra-proposition-ML-characterization} (1) から従う。一般に、有限生成加群が
Mittag-Leffler であることと有限表示であることは同値である。これは命題
\ref{algebra-proposition-fg-tensor}、\ref{algebra-proposition-fp-tensor}、および
\ref{algebra-proposition-ML-tensor} から従う。
% [P01-E0457] The frozen source's missing "if" in "if and only" is rendered by its evident sense.
\item 自由加群は命題 \ref{algebra-proposition-ML-characterization} の条件 (1) を満たすので、
Mittag-Leffler である。
\item 前の例と補題 \ref{algebra-lemma-direct-sum-ML} により、射影加群は Mittag-Leffler である。
\end{enumerate}
\end{example}

\noindent
さらに、例の一覧にネーター環 $R$ 上の形式冪級数環を加えたい。
これは次の補題の帰結である。
% [P01-E0458] The frozen source's missing preposition in "consequence the following lemma" is rendered by its evident sense.

\begin{lemma}
\label{algebra-lemma-flat-ML-criterion}
$M$ を平坦 $R$-加群とする。次は同値である：
\begin{enumerate}
\item $M$ は Mittag-Leffler である、および
\item $F$ が有限自由 $R$-加群で $x \in F \otimes_R M$ ならば、
$x \in F' \otimes_R M$ となる $F$ の最小の部分加群 $F'$ が存在する。
\end{enumerate}
\end{lemma}

\begin{proof}
含意 (1) $\Rightarrow$ (2) は補題 \ref{algebra-lemma-minimal-contains} の特別な場合である。
(2) を仮定する。定理 \ref{algebra-theorem-lazard} により、各項が有限自由 $R$-加群であるような
有向系 $(M_i, f_{ij})$ の余極限 $M = \colim_{i \in I} M_i$ として $M$ を書ける。
注意 \ref{algebra-remark-flat-ML} により、逆系
$(\Hom_R(M_i, R), \Hom_R(f_{ij}, R))$ が Mittag-Leffler であることを示せば十分である。
言い換えると、$i \in I$ を固定し、$j \geq i$ に対して $Q_j$ を
$\Hom_R(M_j, R) \to \Hom_R(M_i, R)$ の像とすれば、$Q_j$ が安定化することを
示さなければならない。

\medskip\noindent
$M_i$ は有限自由なので、すべての $j$ に対して同一視
$\Hom_R(M_i, M_j) = \Hom_R(M_i, R) \otimes_R  M_j$ ができる。
$M_j$ が自由であることを用いると、$j \geq i$ に対して $Q_j$ は、
$f_{ij} \in Q_j \otimes_R M_j$ となるような $\Hom_R(M_i, R)$ の最小の部分加群である。
同一視 $\Hom_R(M_i, M) = \Hom_R(M_i, R) \otimes_R M$ のもとで、標準写像
$f_i: M_i \to M$ は $\Hom_R(M_i, R) \otimes_R M$ に属する。
$M$ に関する仮定により、$f_i \in Q \otimes_R M$ となるような
$\Hom_R(M_i, R)$ の最小の部分加群 $Q$ が存在する。$Q_j$ が $Q$ に安定化することを示す。

\medskip\noindent
$j \geq i$ に対して可換図式
$$
\xymatrix{
Q_j \otimes_R M_j \ar[r] \ar[d] & \Hom_R(M_i, R) \otimes_R M_j
\ar[d] \\
Q_j \otimes_R M \ar[r] & \Hom_R(M_i, R) \otimes_R M.
}
$$
がある。$f_{ij} \in Q_j \otimes_R M_j$ は
$f_i \in \Hom_R(M_i, R) \otimes_R M$ に写るので、$f_i \in Q_j \otimes_R M$ が従う。
従って $Q$ の選び方により、すべての $j \geq i$ に対して $Q \subset Q_j$ である。

\medskip\noindent
$Q_j$ は減少し、すべての $j \geq i$ に対して $Q \subset Q_j$ なので、$Q_j$ が $Q$ に
安定化することを示すには、$Q_j \subset Q$ となる $j \geq i$ を見つければ十分である。
$$
\Hom_R(M_i, R) \otimes_R M = \colim_{j \in J}
(\Hom_R(M_i, R) \otimes_R
M_j),
$$
の元として、$f_i$ は $j \geq i$ における $f_{ij}$ の余極限であり、また $f_i$ は部分加群
$$
\colim_{j \in J} (Q \otimes_R M_j) \subset \colim_{j \in J}
(\Hom_R(M_i, R)
\otimes_R M_j).
$$
にも属する。
% [P01-E0459] Frozen J-indexed colimits are retained.
従って、ある $j \geq i$ に対して $f_{ij}$ は $Q \otimes_R M_j$ に属する。
$Q_j$ は $f_{ij} \in Q_j \otimes_R M_j$ となる $\Hom_R(M_i, R)$ の最小の部分加群なので、
$Q_j\subset Q$ と結論する。
\end{proof}

\begin{lemma}
\label{algebra-lemma-product-over-Noetherian-ring}
$R$ をネーター環、$A$ を集合とする。このとき $M = R^A$ は平坦かつ Mittag-Leffler な
$R$-加群である。
\end{lemma}

\begin{proof}
補題 \ref{algebra-lemma-Noetherian-coherent} と命題 \ref{algebra-proposition-characterize-coherent} を合わせると、
$M$ は $R$ 上平坦であることが分かる。$M$ が補題 \ref{algebra-lemma-flat-ML-criterion} の条件を
満たすことを示す。$F$ を有限自由 $R$-加群とする。$F'$ が $F$ の任意の部分加群ならば、
$R$ はネーター的なのでこれは有限表示である。従って命題
\ref{algebra-proposition-fp-tensor} により可換図式
$$
\xymatrix{
F' \otimes_R M \ar[r] \ar[d]^{\cong} & F \otimes_R M \ar[d]^{\cong} \\
(F')^A \ar[r] & F^A
}
$$
があり、これによって写像 $F' \otimes_R M \to F \otimes_R M$ を
$(F')^A \to F^A$ と同一視できる。従って $x \in F \otimes_R M$ が
$(x_\alpha) \in F^A$ に対応するならば、$x_\alpha$ で生成される $F$ の部分加群 $F'$ は、
$x \in F' \otimes_R M$ となる $F$ の最小の部分加群である。
% [P01-E0460] The frozen malformed apposition around F' is rendered by its evident sense.
\end{proof}

\begin{lemma}
\label{algebra-lemma-power-series-ML}
$R$ をネーター環、$n$ を正の整数とする。このとき $R$-加群
$M = R[[t_1, \ldots, t_n]]$ は平坦かつ Mittag-Leffler である。
\end{lemma}

\begin{proof}
$R$-加群として、ある（可算）集合 $A$ に対して $M = R^A$ である。
従って、この補題は補題 \ref{algebra-lemma-product-over-Noetherian-ring} の特別な場合である。
\end{proof}

\begin{example}
\label{algebra-example-not-ML}
Mittag-Leffler でない加群。
\begin{enumerate}
\item 例 \ref{algebra-example-Q-not-ML} と命題 \ref{algebra-proposition-ML-tensor} により、
$\mathbf{Q}$ は Mittag-Leffler $\mathbf{Z}$-加群でない。
\item 後で（定理 \ref{algebra-theorem-projectivity-characterization}）、平坦かつ可算生成な加群について、
射影的であることと Mittag-Leffler であることが同値であることを証明する。
従って、平坦かつ可算生成で射影的でない任意の加群 $M$ は、Mittag-Leffler でない加群の例である。
そのような例については注意 \ref{algebra-remark-warning} を参照せよ。
\item $k$ を体とし、$R = k[[x]]$ とする。$R$-加群
$M = \prod_{n \in \mathbf{N}} R/(x^n)$ は Mittag-Leffler でない。
実際、$\xi_{2^m} = x^{2^{m - 1}}$ かつ、それ以外では $\xi_n = 0$ によって定義される元
$\xi = (\xi_1, \xi_2, \xi_3, \ldots)$ を考える。すると
$$
\xi = (0, x, 0, x^2, 0, 0, 0, x^4, 0, 0, 0, 0, 0, 0, 0, x^8, \ldots)
$$
である。このとき $m \gg 0$ に対して、$M/x^{2^m}M$ における $\xi$ の零化イデアルは
$x^{2^{m - 1}}$ で生成される。
% [P01-E0461] The frozen source's omitted "by" is rendered by its evident sense.
しかし $M$ が Mittag-Leffler ならば、有限 $R$-加群 $Q$ と元 $\xi' \in Q$ が存在し、
すべての $l \geq 1$ に対して $Q/x^l Q$ における $\xi'$ の零化イデアルと
$M/x^l M$ における $\xi$ の零化イデアルが一致するはずである。命題
\ref{algebra-proposition-ML-characterization} (1) を参照せよ。
ここで、$\xi' \in Q$ が捩れ元かどうかに応じて、すべての $l \gg 0$ に対し
$Q/x^l Q$ における $\xi'$ の零化イデアルが $x^a$ または $x^{l - a}$ のいずれかで
生成されるような整数 $a \geq 0$ が存在することを証明できる。以上を合わせると、
すべての $l = 2^m >> 0$ に対して $a = l/2$ または $l - a = l/2$ という等式を得るが、
これは不合理である。
% [P01-E0462] Frozen ordinary double greater-than signs are retained instead of silently changing them to \gg.
\item 同じ議論により、$\bigoplus_{n \in \mathbf{N}} R/(x^n)$ の $(x)$-進完備化は
$R = k[[x]]$ 上 Mittag-Leffler でない（ヒント：$\xi$ は実際、この完備化の元である）。
\item $R = k[a, b]/(a^2, ab, b^2)$ とする。$S$ を表示 $S = R[t]/(at - b)$ をもつ
有限表示 $R$-代数とする。このとき $R$-加群として $S$ は可算生成かつ直既約である
（詳細は省略する）。他方、$R$ はアルティン局所環であり、従って完備局所環、従って
ヘンゼル局所環である。補題 \ref{algebra-lemma-complete-henselian} を参照せよ。
$S$ が $R$-加群として Mittag-Leffler ならば、補題
\ref{algebra-lemma-split-ML-henselian} により、有限 $R$-加群の直和となるはずである。
従って $S$ は $R$-加群として Mittag-Leffler でないと結論する。
\end{enumerate}
\end{example}




\section{可算生成 Mittag-Leffler 加群}
\label{algebra-section-ML-countable}

\noindent
可算生成 Mittag-Leffler 加群は、とりわけ単純な構造をもつことが分かる。

\begin{lemma}
\label{algebra-lemma-ML-countable-colimit}
$M$ を $R$-加群とする。$(M_i, f_{ij})$ が有限表示 $R$-加群の有向系であるように
$M = \colim_{i \in I} M_i$ と書く。$M$ が Mittag-Leffler かつ可算生成ならば、
$M \cong \colim_{i \in I'} M_i$ となる有向な可算部分集合 $I' \subset I$ が存在する。
\end{lemma}

\begin{proof}
$x_1, x_2, \ldots$ を $M$ の可算生成集合とする。各 $x_n$ に対して、$x_n$ が標準写像
$f_i: M_i \to M$ の像に入るような $i \in I$ を選び、これらすべての $i$ の集合を
$I'_{0} \subset I$ とする。いま $M$ は Mittag-Leffler なので、各 $i \in I'_{0}$ に対して、
$j \geq i$ であり、かつすべての $k \geq i$ について $f_{ij}: M_i \to M_j$ が
$f_{ik}: M_i \to M_k$ を経由するような $j \in I$ を選べる（命題
\ref{algebra-proposition-ML-characterization} の条件 (3)）。$I'_1$ を $I'_0$ とこれらすべての
$j$ との合併とする。$I'_1$ は可算なので、これを可算有向集合 $I'_{2} \subset I$ に
拡大できる。次に $I'_{0}$ に対して行ったのと同じ手続きを $I'_{2}$ に適用し、
新たな可算集合 $I'_{3} \subset I$ を得る。さらに $I'_{3}$ を可算有向集合 $I'_{4}$ に
拡大する。このように続ける。すなわち、$\ell$ が奇数なら各 $ i \in I'_{\ell}$ に対して
命題 \ref{algebra-proposition-ML-characterization} (3) にあるような $j$ を加え、$\ell$ が偶数なら
$I'_{\ell}$ を有向集合に拡大する。こうして $\ell \geq 0$ に対する部分集合の列
$I'_{\ell} \subset I$ を得る。合併 $I' = \bigcup I'_{\ell}$ は次を満たす：
% [P01-E0463] The frozen source's malformed "is a countable" is rendered by its evident grammatical sense.
% [P01-E0464] The frozen odd/even stage description is retained, although it reverses the explicitly constructed parity.
\begin{enumerate}
\item $I'$ は可算かつ有向である；
\item 各 $x_n$ は、ある $i \in I'$ に対する $f_i: M_i \to M$ の像に入る；
\item $i \in I'$ ならば、$j \geq i$ であり、すべての $k \in I$ で $k \geq i$ となるものに
ついて $f_{ij}: M_i \to M_j$ が $f_{ik}: M_i \to M_k$ を経由するような $j \in I'$ が
存在する。とくに $k \geq i$ に対して $\Ker(f_{ik}) \subset \Ker(f_{ij})$ である。
\end{enumerate}
標準写像 $\colim_{i \in I'} M_i \to \colim_{i \in I} M_i = M$ は同型であると主張する。
(2) により全射である。単射性について、$x \in \colim_{i \in I'} M_i$ が
$\colim_{i \in I} M_i$ において $0$ に写るとする。ある $i \in I'$ に対する元
$\tilde{x} \in M_i$ によって $x$ を表すと、これはある $k \in I, k \geq i$ に対して
$f_{ik}(\tilde{x}) = 0$ であることを意味する。しかし (3) により、$f_{ij}(\tilde{x}) = 0$
となる $j \in I', j \geq i,$ が存在する。従って $\colim_{i \in I'} M_i$ において
$x = 0$ である。
\end{proof}

\noindent
補題 \ref{algebra-lemma-ML-countable-colimit} により、$R$ 上の可算生成 Mittag-Leffler 加群 $M$ は、
各 $M_n$ が有限表示 $R$-加群であるような系
$$
M_1 \to M_2 \to M_3 \to M_4 \to \ldots
$$
の余極限である。これを見るには、補題 \ref{algebra-lemma-ML-limit-nonempty} の証明と同様に論じて、
可算有向集合が $(\mathbf{N}, \geq)$ と同型な共終部分集合をもつことを確認すればよい。
% [P01-E0465] Frozen (N, >=) order is retained for this displayed forward direct system.
$R = k[x_1, x_2, x_3, \ldots]$ および $M = R/(x_i)$ とする。このとき $M$ は有限生成だが
有限表示でなく、従って Mittag-Leffler でない（例 \ref{algebra-example-ML} の (1) を参照）。
しかし $M_n = R/(x_1, \ldots, x_n)$ とおけば、もちろん $M = \colim_n M_n$ と書ける。
従って、$M$ をこのような極限として書けるという条件から $M$ が Mittag-Leffler であるとは限らない。

\begin{lemma}
\label{algebra-lemma-ML-countable}
$R$ を環とする。
$M$ を $R$-加群とする。
$M$ が Mittag-Leffler かつ可算生成であると仮定する。
$P$ が有限生成であるような任意の $R$-加群写像 $f : P \to M$ に対して、次を満たす
自己準同型 $\alpha : M \to M$ が存在する：
\begin{enumerate}
\item $\alpha : M \to M$ は有限表示 $R$-加群を経由する、および
\item $\alpha \circ f = f$。
\end{enumerate}
\end{lemma}

\begin{proof}
補題 \ref{algebra-lemma-ML-countable-colimit} により、$I$ が可算である有限表示 $R$-加群の有向余極限
$M = \colim_{i \in I} M_i$ と書く。遷移写像を $f_{ij}$ と書き、$M$ への標準写像を
$f_i : M_i \to M$ と書く。$N = \prod_{s \in I} M_s$ とおく。
$$
M_i^* = \Hom_R(M_i, N) = \prod\nolimits_{s \in I} \Hom_R(M_i, M_s)
$$
と書けば、$(M_i^*)$ は $I$ 上の $R$-加群の逆系である。
$\Hom_R(M, N) = \lim M_i^*$ であることに注意する。$M$ は Mittag-Leffler なので、
各 $i \in I$ に対して、次を満たす添字 $k(i) \geq i$ が見つかる：
$$
E_i := \bigcap\nolimits_{i' \geq i} \Im(M_{i'}^* \to M_i^*)
=
\Im(M_{k(i)}^* \to M_i^*)
$$
$\Im(P \to M) \subset \Im(M_j \to M)$ となる $j \in I$ を選んで固定する。
$P$ は有限生成なので、これは可能である。$k = k(j)$ とおく。
$x = (0, \ldots, 0, \text{id}_{M_k}, 0, \ldots, 0) \in M_k^*$ とし、これが
$y = (0, \ldots, 0, f_{jk}, 0, \ldots, 0) \in M_j^*$ に写ることに注意する。
$k$ の選び方により $y \in E_j$ である。例 \ref{algebra-example-ML-surjective-maps} により、
各 $i \geq j$ に対して遷移写像 $E_i \to E_j$ は全射であり、
$\lim E_i = \lim M_i^* = \Hom_R(M, N)$ である。従って補題
\ref{algebra-lemma-ML-limit-nonempty} により、$E_j \subset M_j^*$ において $y$ に写る元
$z \in \Hom_R(M, N)$ が存在する。$z_k$ を $z$ の第 $k$ 成分とする。このとき
$z_k : M \to M_k$ は、図式
$$
\xymatrix{
M \ar[r]_{z_k} & M_k \\
M_j \ar[ru]_{f_{jk}} \ar[u]^{f_j}
}
$$
を可換にする準同型である。$\alpha : M \to M$ を合成
$f_k \circ z_k : M \to M_k \to M$ とする。構成により $\alpha$ は有限表示加群を経由し、
$\alpha \circ f_j = f_j$ である。$f$ の像は $f_j$ の像に含まれるので、これから
$\alpha \circ f = f$ も従う。
\end{proof}

\noindent
後で（補題 \ref{algebra-lemma-split-ML-henselian} を参照）、補題 \ref{algebra-lemma-ML-countable} が、
ヘンゼル局所環上の可算生成 Mittag-Leffler 加群は有限表示加群の直和であることを
意味するのを見よう。




\section{射影加群の特徴づけ}
\label{algebra-section-characterize-projective}

\noindent
この節の目標は、加群が射影的であることと、それが平坦、Mittag-Leffler、かつ可算生成加群の
直和であることが同値であると証明することである（下の定理
\ref{algebra-theorem-projectivity-characterization}）。

\begin{lemma}
\label{algebra-lemma-countgen-projective}
$M$ を $R$-加群とする。$M$ が平坦、Mittag-Leffler、かつ可算生成ならば、$M$ は射影的である。
\end{lemma}

\begin{proof}
Lazard の定理（定理 \ref{algebra-theorem-lazard}）により、集合 $I$ で添字づけられた有限自由
$R$-加群の有向系 $(M_i, f_{ij})$ に対して $M = \colim_{i \in I} M_i$ と書ける。
補題 \ref{algebra-lemma-ML-countable-colimit} により、$I$ は可算と仮定してよい。いま
$$
0 \to N_1 \to N_2 \to N_3 \to 0
$$
を $R$-加群の完全列とする。$\Hom_R(M, -)$ を適用しても完全性が保たれることを示さなければ
ならない。$M_i$ は有限自由なので、各 $i$ に対して
$$
0 \to \Hom_R(M_i, N_1) \to \Hom_R(M_i, N_2) \to
\Hom_R(M_i, N_3) \to 0
$$
は完全である。$M$ は Mittag-Leffler なので、$(\Hom_R(M_i, N_{1}))$ は
Mittag-Leffler 逆系である。従って補題 \ref{algebra-lemma-ML-exact-sequence} により
$$
0 \to \lim_{i \in I} \Hom_R(M_i, N_1) \to
\lim_{i \in I} \Hom_R(M_i, N_2) \to
\lim_{i \in I} \Hom_R(M_i, N_3) \to 0
$$
は完全である。しかし任意の $R$-加群 $N$ に対して関手的同型
$\Hom_R(M, N) \cong \lim_{i \in I} \Hom_R(M_i, N)$ があるので、
$$
0 \to \Hom_R(M, N_1) \to \Hom_R(M, N_2) \to
\Hom_R(M, N_3) \to 0
$$
は完全である。
\end{proof}

\begin{remark}
\label{algebra-remark-characterize-projective}
補題 \ref{algebra-lemma-countgen-projective} は可算生成という仮定なしには成り立たない。たとえば
$\mathbf Z$-加群 $M = \mathbf{Z}[[x]]$ は平坦かつ Mittag-Leffler だが射影的でない。
補題 \ref{algebra-lemma-power-series-ML} により Mittag-Leffler である。自由アーベル群の部分群は
自由なので、射影的な $\mathbf Z$-加群は実際には自由であり、その部分加群も自由である。
従って $M$ が射影的でないことを示すには、自由でない部分加群を与えれば十分である。
素数 $p$ を固定し、すべての整数 $m \geq 1$ に対して、有限個を除くすべての $i$ について
$p^m$ が $a_i$ を割り切るような冪級数 $f(x) = \sum a_i x^i$ からなる部分加群 $N$ を考える。
すると、すべての $a_i \in \mathbf{Z}$ に対して $\sum a_i p^i x^i$ は $N$ に属するので、
$N$ は非可算である。従って $N$ が自由ならばその階数は非可算となり、
$\mathbf{Z}/p$ 上の $N/pN$ の次元も非可算となるはずである。しかしこれは正しくない。
$i \geq 0$ に対する元 $x^i \in N/pN$ が $N/pN$ を生成するからである。
\end{remark}

\begin{theorem}
\label{algebra-theorem-projectivity-characterization}
$M$ を $R$-加群とする。このとき $M$ が射影的であることと次は同値である：
\begin{enumerate}
\item $M$ は平坦である、
\item $M$ は Mittag-Leffler である、
\item $M$ は可算生成 $R$-加群の直和である。
\end{enumerate}
\end{theorem}

\begin{proof}
まず $M$ が射影的であるとする。このとき $M$ は自由加群の直和因子なので、これらの性質が
直和因子に移ることから $M$ は平坦かつ Mittag-Leffler である。Kaplansky の定理
（定理 \ref{algebra-theorem-projective-direct-sum}）により、$M$ は (3) を満たす。

\medskip\noindent
逆に $M$ が (1)--(3) を満たすとする。平坦性および Mittag-Leffler 性は直和因子に移るので、
$M$ は平坦、Mittag-Leffler、可算生成な $R$-加群の直和である。補題
\ref{algebra-lemma-countgen-projective} により、$M$ は射影加群の直和である。従って $M$ は射影的である。
\end{proof}

\begin{lemma}
\label{algebra-lemma-ML-ui-descent}
$f: M \to N$ を $R$-加群の普遍単射な写像とする。$M$ は可算生成 $R$-加群の直和であり、
$N$ は平坦かつ Mittag-Leffler であると仮定する。このとき $M$ は射影的である。
% [P01-E0466] The frozen source's missing article before "universally injective map" is rendered by its evident sense.
\end{lemma}

\begin{proof}
補題 \ref{algebra-lemma-ui-flat-domain} および \ref{algebra-lemma-pure-submodule-ML} により、
$M$ は平坦かつ Mittag-Leffler である。従って定理
\ref{algebra-theorem-projectivity-characterization} から結論が従う。
\end{proof}

\begin{lemma}
\label{algebra-lemma-universally-injective-submodule-powerseries}
$R$ をネーター環、$M$ を $R$-加群とする。$M$ は可算生成 $R$-加群の直和であり、
ある $n$ に対して普遍単射な写像 $M \to R[[t_1, \ldots, t_n]]$ が存在すると仮定する。
このとき $M$ は射影的である。
% [P01-E0467] The frozen source's article error before "R-module" is rendered by its evident sense.
\end{lemma}

\begin{proof}
補題 \ref{algebra-lemma-ML-ui-descent} および \ref{algebra-lemma-power-series-ML} から従う。
\end{proof}




\section{加群の性質の上昇}
\label{algebra-section-ascent}

\noindent
定理 \ref{algebra-theorem-projectivity-characterization} に現れる加群の性質はすべて、任意の環写像に沿って
上昇する：

\begin{lemma}
\label{algebra-lemma-ascend-properties-modules}
$R \to S$ を環写像とし、$M$ を $R$-加群とする。このとき：
\begin{enumerate}
\item $M$ が平坦ならば、$S$-加群 $M \otimes_R S$ は平坦である。
\item $M$ が Mittag-Leffler ならば、$S$-加群 $M \otimes_R S$ は Mittag-Leffler である。
\item $M$ が可算生成 $R$-加群の直和ならば、$S$-加群 $M \otimes_R S$ は可算生成
$S$-加群の直和である。
\item $M$ が射影的ならば、$S$-加群 $M \otimes_R S$ は射影的である。
\end{enumerate}
\end{lemma}

\begin{proof}
(2) 以外はすべて明らかである。(2) については、命題
\ref{algebra-proposition-ML-characterization} における Mittag-Leffler 性の定式化 (3) と、テンソル積が
余極限を取ることと可換であるという事実を用いる。あるいは、命題
\ref{algebra-proposition-ML-tensor} の特徴づけを用いてもよい。
\end{proof}




\section{加群の性質の降下}
\label{algebra-section-descent}

\noindent
定理 \ref{algebra-theorem-projectivity-characterization} において射影性を特徴づける諸性質の
忠実平坦降下を扱う。平坦性を仮定すれば、Mittag-Leffler 加群であるという性質は降下する：

\begin{lemma}
\label{algebra-lemma-ffdescent-ML}
\begin{reference}
Juan Pablo Acosta Lopez からの 2014 年 12 月 20 日付電子メール。
\end{reference}
$R \to S$ を忠実平坦な環写像とし、$M$ を $R$-加群とする。$S$-加群
$M \otimes_R S$ が Mittag-Leffler ならば、$M$ は Mittag-Leffler である。
\end{lemma}

\begin{proof}
$M_i$ が有限表示 $R$-加群であるような有向余極限
$M = \colim_{i\in I} M_i$ として書く。命題
\ref{algebra-proposition-ML-characterization} を用いると、各 $i \in I$ に対して、
$M_i\rightarrow M_j$ が $M_i\rightarrow M$ を支配するような $i \leq j$, $j\in I$ が
存在することを示せばよいと分かる。

\medskip\noindent
$N$ を押し出し
% [P01-E0468] The frozen source's missing "to be" after "Take N" is rendered by its evident sense.
$$
\xymatrix{
M_i  \ar[r] \ar[d] & M_j \ar[d] \\
M \ar[r] & N
}
$$
とする。このとき補題は、$M_j\rightarrow N$ が普遍単射となるような $j$ の存在と同値である。
補題 \ref{algebra-lemma-domination-universally-injective} を参照せよ。
$S$ によるテンソル化で得られる図式
$$
\xymatrix{
M_i\otimes_R S  \ar[r] \ar[d] & M_j\otimes_R S \ar[d] \\
M\otimes_R S \ar[r] & N\otimes_R S
}
$$
が押し出し図式であることに注意する。
% [P01-E0469] The frozen source's sentence-internal capitalized "Is" is rendered normally.
$M \otimes_R S = \colim_{i\in I} M_i \otimes_R S$ は $M\otimes_R S$ を有限表示
$S$-加群の余極限として表し、かつ $M\otimes_R S$ は Mittag-Leffler なので、
$M_j\otimes_R S\rightarrow N\otimes_R S$ が普遍単射となる $j \geq i$ が存在する。
$R\rightarrow S$ は忠実平坦だから、$M_j\rightarrow N$ も普遍単射であると結論する。
\end{proof}

\begin{lemma}
\label{algebra-lemma-ffdescent-countable}
$R \to S$ を忠実平坦な環写像とし、$M$ を $R$-加群とする。$S$-加群
$M \otimes_R S$ が可算生成ならば、$M$ は可算生成である。
\end{lemma}

\begin{proof}
$M \otimes_R S$ が元 $y_i$, $i = 1, 2, 3, \ldots$ で生成されるとする。ある
$n_i \geq 0$, $x_{ij} \in M$, $s_{ij} \in S$ に対して
$y_i = \sum_{j = 1, \ldots, n_i} x_{ij} \otimes s_{ij}$ と書く。可算個の元 $x_{ij}$ で
生成される部分加群を $M' \subset M$ と書く。このとき $M' \otimes_R S \to M \otimes_R S$
の像は生成元 $y_i$ を含むので、この写像は全射である。$S$ は $R$ 上忠実平坦だから、
望みどおり $M' = M$ と結論する。
\end{proof}

\noindent
ここまでで、可算生成射影加群の忠実平坦降下は容易に従う。

\begin{lemma}
\label{algebra-lemma-ffdescent-countable-projectivity}
$R \to S$ を忠実平坦な環写像とし、$M$ を $R$-加群とする。$S$-加群
$M \otimes_R S$ が可算生成かつ射影的ならば、$M$ は可算生成かつ射影的である。
\end{lemma}

\begin{proof}
補題 \ref{algebra-lemma-descend-properties-modules}、\ref{algebra-lemma-ffdescent-ML}、
\ref{algebra-lemma-ffdescent-countable}、および定理
\ref{algebra-theorem-projectivity-characterization} から従う。
\end{proof}

\noindent
残るのは、デヴィサージュを用いて一般の場合の射影性の降下を可算生成の場合へ帰着することだけで
ある。まず二つの簡単な補題を示す。

\begin{lemma}
\label{algebra-lemma-lift-countably-generated-submodule}
$R \to S$ を環写像、$M$ を $R$-加群、$Q$ を $M \otimes_R S$ の可算生成
$S$-部分加群とする。このとき、$\Im(P \otimes_R S \to M \otimes_R S)$ が $Q$ を含むような、
$M$ の可算生成 $R$-部分加群 $P$ が存在する。
\end{lemma}

\begin{proof}
$y_1, y_2, \ldots$ を $Q$ の生成元とし、ある $x_{jk} \in M$ および $s_{jk} \in S$ に対して
$y_j = \sum_k x_{jk} \otimes s_{jk}$ と書く。このとき $P$ を $x_{jk}$ で生成される $M$ の
部分加群とすればよい。
% [P01-E0470] The frozen source's missing "to" in "take P be" is rendered by its evident sense.
\end{proof}

\begin{lemma}
\label{algebra-lemma-adapted-submodule}
$R \to S$ を環写像、$M$ を $R$-加群とする。
$M \otimes_R S = \bigoplus_{i \in I} Q_i$ は可算生成 $S$-加群 $Q_i$ の直和であると仮定する。
$N$ が $M$ の可算生成部分加群ならば、$N' \supset N$ であり、ある部分集合 $I' \subset I$ に
対して
$\Im(N' \otimes_R S \to M \otimes_R S) = \bigoplus_{i \in I'} Q_i$ となるような、
$M$ の可算生成部分加群 $N'$ が存在する。
\end{lemma}

\begin{proof}
$N'_0 = N$ とおく。可算生成部分加群の増大列 $N'_{\ell} \subset M$ を
$\ell = 0, 1, 2, \ldots$ に関して帰納的に構成する。$I'_{\ell}$ を、
$\Im(N'_{\ell} \otimes_R S \to M \otimes_R S)$ の $Q_i$ への射影が非零となる
$i \in I$ の集合とするとき、すべての $i \in I'_{\ell}$ に対して
$\Im(N'_{\ell + 1} \otimes_R S \to M \otimes_R S)$ が $Q_i$ を含むようにする。
$N'_\ell$ から $N'_{\ell + 1}$ を構成するには、$Q$ を $i \in I'_{\ell}$ に対する
（可算個の）$Q_i$ の和とし、補題 \ref{algebra-lemma-lift-countably-generated-submodule} にある $P$ を選び、
$N'_{\ell + 1} = N'_{\ell} + P$ とおく。$N'_{\ell}$ を構成し終えたら、
$N' = \bigcup_{\ell} N'_{\ell}$ および $I' = \bigcup_{\ell} I'_{\ell}$ とすればよい。
\end{proof}

\begin{theorem}
\label{algebra-theorem-ffdescent-projectivity}
$R \to S$ を忠実平坦な環写像とし、$M$ を $R$-加群とする。$S$-加群
$M \otimes_R S$ が射影的ならば、$M$ は射影的である。
\end{theorem}

\begin{proof}
$M$ の Kaplansky デヴィサージュを構成し、それが射影加群の直和、従って射影加群であることを
示す。定理 \ref{algebra-theorem-projective-direct-sum} により、$M \otimes_R S =
\bigoplus_{i \in I} Q_i$ を可算生成 $S$-加群 $Q_i$ の直和として書ける。$M$ に整列順序を
選ぶ。超限再帰を用いて、各順序数 $\alpha$ に対する $M$ の部分加群の増大族
$M_{\alpha}$ を、$M_{\alpha} \otimes_R S$ が $Q_i$ のある部分集合の直和となるように定義する。

\medskip\noindent
$\alpha = 0$ に対して $M_0 = 0$ とおく。$\alpha$ が極限順序数で、すべての
$\beta < \alpha$ に対して $M_{\beta}$ が定義されているならば、
$M_\alpha = \bigcup_{\beta < \alpha} M_{\beta}$ と定義する。各 $\beta < \alpha$ に対して
$M_{\beta} \otimes_R S$ は $Q_i$ のある部分集合の直和なので、
$M_{\alpha} \otimes_R S$ についても同じことが成り立つ。$\alpha + 1$ が後続順序数で
$M_{\alpha}$ が定義されているならば、$M_{\alpha + 1}$ を次のように定義する。
$M_{\alpha} = M$ ならば $M_{\alpha +1} = M$ とおく。そうでなければ、固定した整列順序に
関して $x \notin M_{\alpha}$ となる最小の $x \in M$ を選ぶ。$S$ は $R$ 上平坦なので、
$(M/M_{\alpha}) \otimes_R S = M \otimes_R S/M_{\alpha} \otimes_R S$ である。従って
$M_{\alpha} \otimes_R S$ がある $Q_i$ の直和であることから、
$(M/M_{\alpha}) \otimes_R S$ についても同じことが成り立つ。補題
\ref{algebra-lemma-adapted-submodule} により、$M/M_{\alpha}$ における $x$ の像を含み、かつ
$P \otimes_R S$（$S$ は $R$ 上平坦なので、これは
$\Im(P \otimes_R S \to M \otimes_R S)$ に等しい）がある $Q_i$ の直和となるような、
$M/M_{\alpha}$ の可算生成 $R$-部分加群 $P$ を見つけられる。
% [P01-E0471] The frozen image codomain M tensor S is retained although P lies in M/M_alpha here.
$M \otimes_R S = \bigoplus_{i \in I} Q_i$ は射影的であり、射影性は直和因子に移るので、
$P \otimes_R S$ も射影的である。従って補題
\ref{algebra-lemma-ffdescent-countable-projectivity} により、$P$ は射影的である。最後に
$M_{\alpha + 1}$ を $M$ における $P$ の逆像として定義する。このとき
$M_{\alpha + 1}/M_{\alpha} = P$ は可算生成かつ射影的である。とくに
$M_{\alpha + 1}/M_{\alpha}$ の射影性により列
$0 \to M_{\alpha} \to M_{\alpha + 1} \to M_{\alpha + 1}/M_{\alpha} \to 0$ が分裂するので、
$M_{\alpha}$ は $M_{\alpha + 1}$ の直和因子である。

\medskip\noindent
$M$ に関する超限帰納法（$M_{\alpha}$ に含まれない最小の $x \in M$ を含むように
$M_{\alpha + 1}$ を構成したという事実を用いる）により、各 $x \in M$ はある
$M_{\alpha}$ に含まれる。従って、各 $x \in M$ に対して $x \in M_{\alpha}$ となる
$\alpha \in S$ が存在するという性質を満たす、十分大きな順序数 $S$ が存在する。
% [P01-E0472] Frozen "transfinite induction on M" wording is retained with its recursion context visible.
% [P01-E0473] Frozen reuse of S for both the target ring and an ordinal is retained.
これは $(M_{\alpha})_{\alpha \in S}$ が $M$ の Kaplansky デヴィサージュの性質 (1) を
満たすことを意味する。他の性質は構成から明らかである。従って
$M = \bigoplus_{\alpha + 1 \in S} M_{\alpha + 1}/M_{\alpha}$ と結論する。
各 $M_{\alpha + 1}/M_{\alpha}$ は構成により射影的なので、$M$ は射影的である。
\end{proof}




\section{完備化}
\label{algebra-section-completion}

\noindent
$R$ を環、$I$ をイデアルとする。{\it $I$ に関する $R$ の完備化}を極限
$$
R^\wedge = \lim_n R/I^n.
$$
として定義する。$R^\wedge$ の元は、すべての $n$ に対して
$f_n \equiv f_{n + 1} \bmod I^n$ を満たす元 $f_n \in R/I^n$ の列によって与えられる。
$R^\wedge$ を $R$-代数とみなす。同様に、$M$ が $R$-加群ならば、
{\it $I$ に関する $M$ の完備化}を極限
$$
M^\wedge = \lim_n M/I^nM.
$$
として定義する。$M^\wedge$ の元は、すべての $n$ に対して
$m_n \equiv m_{n + 1} \bmod I^nM$ を満たす元 $m_n \in M/I^nM$ の列によって与えられる。
$M^\wedge$ を $R^\wedge$-加群とみなす。この記述から、標準写像
$$
M \longrightarrow M^\wedge
\quad\text{and}\quad
M \otimes_R R^\wedge \longrightarrow M^\wedge.
$$
が常に存在することは明らかである。さらに、加群の写像 $\varphi : M \to N$ が与えられると、
図式
$$
\xymatrix{
M \ar[r] \ar[d] & N \ar[d] \\
M^\wedge \ar[r] & N^\wedge
}
$$
を可換にする完備化上の誘導写像 $\varphi^\wedge : M^\wedge \to N^\wedge$ を得る。
一般に完備化は完全関手ではない。Examples の節
\ref{examples-section-completion-not-exact} を参照せよ。まずいくつかの肯定的な結果を示す。

\begin{lemma}
\label{algebra-lemma-completion-generalities}
$R$ を環、$I \subset R$ をイデアル、$\varphi : M \to N$ を $R$-加群の写像とする。
\begin{enumerate}
\item $M/IM \to N/IN$ が全射ならば、$M^\wedge \to N^\wedge$ は全射である。
\item $M \to N$ が全射ならば、$M^\wedge \to N^\wedge$ は全射である。
\item $0 \to K \to M \to N \to 0$ が $R$-加群の短完全列で $N$ が平坦ならば、
$0 \to K^\wedge \to M^\wedge \to N^\wedge \to 0$ は短完全列である。
\item 任意の有限 $R$-加群 $M$ に対して、写像 $M \otimes_R R^\wedge \to M^\wedge$ は全射である。
\end{enumerate}
\end{lemma}

\begin{proof}
$M/IM \to N/IN$ が全射であると仮定する。このとき Nakayama の補題により、各
$n \geq 1$ に対して写像 $M/I^nM \to N/I^nN$ は全射である。より正確には、環 $R/I^n$ と
冪零イデアル $I/I^n$ 上の写像 $M/I^nM \to N/I^nN$ に補題 \ref{algebra-lemma-NAK} (11) を
適用すればよい。$K_n = \{x \in M \mid \varphi(x) \in I^nN\}$ とおく。こうして短完全列
$$
0 \to K_n/I^nM \to M/I^nM \to N/I^nN \to 0
$$
を得る。標準写像 $K_{n + 1}/I^{n + 1}M \to K_n/I^nM$ は全射であると主張する。
実際、$x \in K_n$ ならば、$z_j \in I^n$, $n_j \in N$ を用いて
$\varphi(x) = \sum z_j n_j$ と書く。仮定により、$m_j \in M$, $z_{jk} \in I$,
$n_{jk} \in N$ を用いて $n_j = \varphi(m_j) + \sum z_{jk}n_{jk}$ と書ける。従って
$$
\varphi(x - \sum z_j m_j) = \sum z_jz_{jk} n_{jk}.
$$
これは $x' = x - \sum z_j m_j \in K_{n + 1}$ が $x \bmod I^nM$ に写ることを意味し、
主張が証明された。ここで上の短完全列の逆系に補題 \ref{algebra-lemma-Mittag-Leffler} を適用して
(1) を得る。(2) は (1) の特別な場合である。(3) の仮定が成り立つならば、各 $n$ に対して列
$$
0 \to K/I^nK \to M/I^nM \to N/I^nN \to 0
$$
は補題 \ref{algebra-lemma-flat-tor-zero} により短完全である。従って補題
\ref{algebra-lemma-Mittag-Leffler} を直接適用して (3) が成り立つと結論できる。(4) を示すため、
生成元 $x_i \in M$, $i = 1, \ldots, n$ を選ぶ。このとき写像
$R^{\oplus n} \to M$, $(a_1, \ldots, a_n) \mapsto \sum a_ix_i$ は全射である。
従って (2) により、写像
$(R^\wedge)^{\oplus n} \to M^\wedge$, $(a_1, \ldots, a_n) \mapsto \sum a_ix_i$ は
全射である。これから主張 (4) が従う。
\end{proof}

\begin{definition}
\label{algebra-definition-complete}
$R$ を環、$I \subset R$ をイデアル、$M$ を $R$-加群とする。写像
$$
M \longrightarrow M^\wedge = \lim_n M/I^nM
$$
が同型ならば、$M$ は {\it $I$-進完備}であるという\footnote{これには条件
$\bigcap I^nM = 0$ も含まれる。}。$R$ が $R$-加群として {\it $I$-進完備}ならば、
$R$ は {\it $I$-進完備}であるという。
\end{definition}

\noindent
$I$ に関する $R$-加群 $M$ の完備化が $I$-進完備であるとは限らない。例については
Examples の節 \ref{examples-section-noncomplete-completion} を参照せよ。イデアルが有限生成ならば、
その完備化は完備である。

\begin{lemma}
\label{algebra-lemma-hathat-finitely-generated}
\begin{reference}
\cite[Theorem 15]{Matlis}。ここで与える簡潔な証明は、Bjorn Poonen からの
2016 年 11 月 5 日付電子メールによる。
\end{reference}
$R$ を環、$I$ を $R$ の有限生成イデアル、$M$ を $R$-加群とする。このとき
\begin{enumerate}
\item 完備化 $M^\wedge$ は $I$-進完備である、および
\item すべての $n \geq 1$ に対して
$I^nM^\wedge = \Ker(M^\wedge \to M/I^nM) = (I^nM)^\wedge$ である。
\end{enumerate}
とくに $R^\wedge$ は $I$-進完備であり、$I^nR^\wedge = (I^n)^\wedge$ および
$R^\wedge/I^nR^\wedge = R/I^n$ が成り立つ。
\end{lemma}

\begin{proof}
$I$ は有限生成なので $I^n$ も有限生成である。その生成元を $f_1, \ldots, f_r$ とする。
全射 $(f_1, \ldots, f_r) : M^{\oplus r} \to I^n M$ に補題
\ref{algebra-lemma-completion-generalities} (2) を適用すると、全射
$$
(M^\wedge)^{\oplus r} \xrightarrow{(f_1, \ldots, f_r)} (I^n M)^\wedge =
\lim_{m \geq n} I^n M/I^m M = \Ker(M^\wedge \to M/I^n M).
$$
を得る。他方、$(f_1, \ldots, f_r) : (M^\wedge)^{\oplus r} \to M^\wedge$ の像は
$I^n M^\wedge$ である。従って $M^\wedge / I^n M^\wedge \simeq M/I^n M$ である。
逆極限を取ると $(M^\wedge)^\wedge \simeq M^\wedge$ を得る。すなわち、
$M^\wedge$ は $I$-進完備である。
\end{proof}

\begin{lemma}
\label{algebra-lemma-completion-differ-by-torsion}
$R$ を環、$I \subset R$ をイデアルとする。
$0 \to M \to N \to Q \to 0$ を $R$-加群の完全列とし、$Q$ は $I$ のある冪で
零化されるとする。このとき完備化により完全列
$0 \to M^\wedge \to N^\wedge \to Q \to 0$ が得られる。
\end{lemma}

\begin{proof}
$I^c Q = 0$ とする。このとき $n \geq c$ に対して $Q/I^nQ = Q$ である。他方、
$n \geq c$ に対して
$I^nM \subset M \cap I^nN \subset I^{n - c}M$ であることは明らかである。
従って $M^\wedge = \lim M/(M \cap I^n N)$ である。$n \geq c$ に対する完全列の系
$$
0 \to M/(M \cap I^n N) \to N/I^n N \to Q \to 0
$$
に補題 \ref{algebra-lemma-Mittag-Leffler} を適用すれば結論を得る。
\end{proof}

\begin{lemma}
\label{algebra-lemma-hathat}
\begin{reference}
Lenstra と de Smit の未刊ノートによる。
\end{reference}
$R$ を環、$I \subset R$ をイデアル、$M$ を $R$-加群とする。
$K_n = \Ker(M^\wedge \to M/I^nM)$ と書く。このとき $M^\wedge$ が $I$-進完備であることと、
すべての $n \geq 1$ に対して $K_n$ が $I^nM^\wedge$ に等しいことは同値である。
\end{lemma}

\begin{proof}
加群 $I^n M^\wedge$ は $K_n$ に含まれる。従って各 $n \geq 1$ に対して標準完全列
$$
0 \to K_n/I^nM^\wedge \to M^\wedge/I^nM^\wedge \to M/I^nM \to 0.
$$
がある。$I^nM^\wedge$ は $I^nM/I^{n + 1}M$ の上へ写るので、
$K_{n + 1} + I^n M^\wedge = K_n$ である。従って逆系
$\{K_n/I^n M^\wedge\}_{n \geq 1}$ の遷移写像は全射である。補題
\ref{algebra-lemma-Mittag-Leffler} により、短完全列
$$
0 \to
\lim_n K_n/I^n M^\wedge \to
(M^\wedge)^\wedge \to
M^\wedge \to 0
$$
がある。従って $M^\wedge$ が完備であることと、すべての $n \geq 1$ に対して
$K_n/I^n M^\wedge = 0$ であることは同値である。
\end{proof}

\begin{lemma}
\label{algebra-lemma-radical-completion}
$R$ を環、$I \subset R$ をイデアルとし、$R^\wedge = \lim R/I^n$ とおく。
\begin{enumerate}
\item $R^\wedge$ の元で、その $R/I$ における像が単元であるものは単元である、
\item $1 + I$ の任意の元は $R^\wedge$ の可逆元に写る、
\item $1 + IR^\wedge$ の任意の元は $R^\wedge$ において可逆である、および
\item イデアル $IR^\wedge$ と $\Ker(R^\wedge \to R/I)$ は $R^\wedge$ の
Jacobson 根基に含まれる。
\end{enumerate}
\end{lemma}

\begin{proof}
$x \in R^\wedge$ が $R/I$ において単元 $x_1$ に写るとする。補題
\ref{algebra-lemma-locally-nilpotent-unit} により、各 $n$ に対して $x$ は $R/I^n$ の単元 $x_n$ に
写る。従って $y = (x_n^{-1}) \in \lim R/I^n = R^\wedge$ は $x$ の逆元である。
(2) と (3) は (1) から直ちに従う。(4) は (1) と補題
\ref{algebra-lemma-contained-in-radical} から従う。
\end{proof}

\begin{lemma}
\label{algebra-lemma-when-surjective-to-completion}
$A$ を環、$I = (f_1, \ldots, f_r)$ を有限生成イデアルとする。各 $i$ に対して
$M \to \lim M/f_i^nM$ が全射ならば、$M \to \lim M/I^nM$ は全射である。
% [P01-E0474] The frozen source omits the hypothesis introducing M as an A-module.
\end{lemma}

\begin{proof}
$I^n \supset (f_1^n, \ldots, f_r^n) \supset I^{rn}$ なので、
$\lim M/I^nM = \lim M/(f_1^n, \ldots, f_r^n)M$ であることに注意する。
$\lim M/(f_1^n, \ldots, f_r^n)M$ の元 $\xi$ は、記号的に
$$
\xi = \sum\nolimits_{n \geq 0} \sum\nolimits_i f_i^n x_{n, i}
$$
と書ける。ここで $x_{n, i} \in M$ である。$M \to \lim M/f_i^nM$ が全射ならば、
$\lim M/f_i^nM$ において $\sum x_{n, i} f_i^n$ に写る $x_i \in M$ が存在する。
このとき $x = \sum x_i$ は $\lim M/I^nM$ において $\xi$ に写る。
\end{proof}

\begin{lemma}
\label{algebra-lemma-complete-by-sub}
$A$ を環、$I \subset J \subset A$ をイデアルとする。$M$ が $J$-進完備であり、
$I$ が有限生成ならば、$M$ は $I$-進完備である。
% [P01-E0475] The frozen source omits the hypothesis introducing M as an A-module.
\end{lemma}

\begin{proof}
$M$ が $J$-進完備で $I$ が有限生成であると仮定する。
$\bigcap J^nM = 0$ なので $\bigcap I^nM = 0$ である。補題
\ref{algebra-lemma-when-surjective-to-completion} により、$I$ が一つの元で生成される場合に
$M \to \lim M/I^nM$ の全射性を証明すれば十分である。$I = (f)$ とする。
$x_{n + 1} - x_n \in f^nM$ を満たす $x_n \in M$ を取る。すべての $n$ に対して
$x_n - x \in f^nM$ となる $x \in M$ が存在することを示さなければならない。
$x_{n + 1} - x_n \in J^nM$ であり、$M$ は $J$-進完備なので、
$x_n - x \in J^nM$ となる元 $x \in M$ が存在する。$x_n$ を $x_n - x$ に置き換えれば、
$x_n \in J^nM$ と仮定してよい。証明を終えるため、これから $x_n \in I^nM$ が従うことを
示す。実際、$x_n - x_{n + 1} = f^nz_n$ と書く。このとき
$$
x_n = f^n(z_n + fz_{n + 1} + f^2z_{n + 2} + \ldots)
$$
である。$f^c \in J^c$ なので、和 $z_n + fz_{n + 1} + f^2z_{n + 2} + \ldots$ は
$M$ において収束する。和 $f^n(z_n + fz_{n + 1} + f^2z_{n + 2} + \ldots)$ は
$M$ において $x_n$ に収束する。部分和が $x_n - x_{n + c}$ に等しく、かつ
$x_{n + c} \in J^{n + c}M$ だからである。
\end{proof}

\begin{lemma}
\label{algebra-lemma-change-ideal-completion}
$R$ を環とする。$I$, $J$ を $R$ のイデアルとする。
$I^c \subset J$ および $J^d \subset I$ となる整数 $c, d > 0$ が存在すると仮定する。
このとき任意の $R$-加群について、$I$ に関する完備化は $J$ に関する完備化と一致する。
とくに $R$-加群 $M$ が $I$-進完備であることと $J$-進完備であることは同値である。
\end{lemma}

\begin{proof}
写像の系 $M/I^nM \to M/J^{\lfloor n/c \rfloor}M$ と写像の系
$M/J^mM \to M/I^{\lfloor m/d \rfloor}M$ を考えると、完備化の間の互いに逆な写像が得られる。
\end{proof}

\begin{lemma}
\label{algebra-lemma-quotient-complete}
$R$ を環、$I$ を $R$ のイデアルとする。$M$ を $I$-進完備 $R$-加群とし、
$K \subset M$ を $R$-部分加群とする。次は同値である：
\begin{enumerate}
\item $K = \bigcap (K + I^nM)$、および
\item $M/K$ は $I$-進完備である。
\end{enumerate}
\end{lemma}

\begin{proof}
$N = M/K$ とおく。補題 \ref{algebra-lemma-completion-generalities} により、写像
$M = M^\wedge \to N^\wedge$ は全射である。従って $N \to N^\wedge$ は全射である。
$N \to N^\wedge$ の核が加群 $\bigcap (K + I^nM) / K$ であることは容易に分かる。
\end{proof}

\begin{lemma}
\label{algebra-lemma-when-finite-module-complete-over-complete-ring}
$R$ を環、$I$ を $R$ のイデアル、$M$ を $R$-加群とする。
(a) $R$ は $I$-進完備、(b) $M$ は有限 $R$-加群、かつ
(c) $\bigcap I^nM = (0)$ ならば、$M$ は $I$-進完備である。
\end{lemma}

\begin{proof}
補題 \ref{algebra-lemma-completion-generalities} により、写像
$M = M \otimes_R R = M \otimes_R R^\wedge \to M^\wedge$ は全射である。
この写像の核は $\bigcap I^nM$ であり、仮定により零である。従って
$M \cong M^\wedge$ であり、$M$ は完備である。
\end{proof}

\begin{lemma}
\label{algebra-lemma-finite-over-complete-ring}
$R$ を環、$I \subset R$ をイデアル、$M$ を $R$-加群とする。次を仮定する：
\begin{enumerate}
\item $R$ は $I$-進完備である、
\item $\bigcap_{n \geq 1} I^nM = (0)$、および
\item $M/IM$ は有限 $R/I$-加群である。
\end{enumerate}
このとき $M$ は有限 $R$-加群である。
\end{lemma}

\begin{proof}
$x_1, \ldots, x_n \in M$ の $M/IM$ における像が $M/IM$ を $R/I$-加群として生成するように
選ぶ。$x_1, \ldots, x_n$ で生成される $R$-部分加群を $M' \subset M$ と書く。
% [P01-E0476] The frozen source's article error before R/I-module is rendered by its evident sense.
% [P01-E0477] The frozen source's missing \"by\" in the naming construction is rendered by its evident sense.
補題 \ref{algebra-lemma-completion-generalities} により、写像 $(M')^\wedge \to M^\wedge$ は全射である。
$\bigcap I^nM = 0$ なので、とくに $\bigcap I^nM' = (0)$ である。従って補題
\ref{algebra-lemma-when-finite-module-complete-over-complete-ring} により $M'$ は完備であり、
$M' \to M^\wedge$ は全射であると結論する。最後に、写像 $M \to M^\wedge$ の核は
$\bigcap I^nM = (0)$ に等しいので零である。従って
$M \cong M' \cong M^\wedge$ は有限生成であると結論する。
\end{proof}




\section{ネーター環に対する完備化}
\label{algebra-section-completion-noetherian}

\noindent
この節では、ネーター環のイデアルに関する完備化を論じる。

\begin{lemma}
\label{algebra-lemma-completion-tensor}
$I$ をネーター環 $R$ のイデアルとする。$I$ に関する完備化を ${}^\wedge$ と書く。
% [P01-E0478] The frozen source's missing "by" in the notation declaration is rendered by its evident sense.
\begin{enumerate}
\item $K \to N$ が有限 $R$-加群の単射ならば、完備化上の写像
$K^\wedge \to N^\wedge$ は単射である。
\item $0 \to K \to N \to M \to 0$ が有限 $R$-加群の短完全列ならば、
$0 \to K^\wedge \to N^\wedge \to M^\wedge \to 0$ は短完全列である。
\item $M$ が有限 $R$-加群ならば、$M^\wedge = M \otimes_R R^\wedge$ である。
\end{enumerate}
\end{lemma}

\begin{proof}
$M = N/K$ とおけば、(1) が (2) から従うことが分かる。
$0 \to K \to N \to M \to 0$ を (2) にある列とする。各 $n$ に対して短完全列
$$
0 \to K/(I^nN \cap K) \to N/I^nN \to M/I^nM \to 0.
$$
を得る。補題 \ref{algebra-lemma-Mittag-Leffler} により完全列
$$
0 \to \lim K/(I^nN \cap K) \to N^\wedge \to M^\wedge \to 0.
$$
を得る。Artin--Rees の補題 \ref{algebra-lemma-Artin-Rees} により、
$n \geq c$ に対して
$I^nK \subset I^n N \cap K \subset I^{n-c} K$ となるような $c$ を選べる。
従って $K^\wedge = \lim K/I^nK = \lim K/(I^nN \cap K)$ であり、(2) が成り立つと結論する。

\medskip\noindent
$M$ を (3) にある加群とし、$0 \to K \to R^{\oplus t} \to M \to 0$ を $M$ の表示とする。
可換図式
$$
\xymatrix{
&
K \otimes_R R^\wedge \ar[r] \ar[d] &
R^{\oplus t} \otimes_R R^\wedge \ar[r] \ar[d] &
M \otimes_R R^\wedge \ar[r] \ar[d] &
0 \\
0 \ar[r] &
K^\wedge \ar[r] &
(R^{\oplus t})^\wedge \ar[r] &
M^\wedge \ar[r] & 0
}
$$
を得る。上の行は完全である。節 \ref{algebra-section-flat} を参照せよ。下の行は (2) により完全である。
補題 \ref{algebra-lemma-completion-generalities} により縦の矢印は全射である。中央の縦の矢印は同型である。
蛇の補題 \ref{algebra-lemma-snake} により (3) が成り立つと結論する。
\end{proof}

\begin{lemma}
\label{algebra-lemma-completion-flat}
$I$ をネーター環 $R$ のイデアルとする。$I$ に関する完備化を ${}^\wedge$ と書く。
% [P01-E0479] The repeated frozen notation declaration is rendered by its evident sense.
\begin{enumerate}
\item 環写像 $R \to R^\wedge$ は平坦である。
\item 関手 $M \mapsto M^\wedge$ は有限生成 $R$-加群の圏上完全である。
\end{enumerate}
\end{lemma}

\begin{proof}
$J$ を $R$ の任意のイデアルとして、
$J \otimes_R R^\wedge \to R \otimes_R R^\wedge = R^\wedge$ を考える。
補題 \ref{algebra-lemma-completion-tensor} によれば、これは $J^\wedge \to R^\wedge$ と同一視され、
$J^\wedge \to R^\wedge$ は単射である。補題 \ref{algebra-lemma-flat} により (1) が従う。
(2) は補題 \ref{algebra-lemma-completion-tensor} (2) の言い換えである。
\end{proof}

\begin{lemma}
\label{algebra-lemma-completion-faithfully-flat}
$I$ をネーター環 $R$ のイデアルとする。$R^\wedge$ を $I$ に関する $R$ の完備化と書く。
$I$ が $R$ の Jacobson 根基に含まれるならば、環写像 $R \to R^\wedge$ は忠実平坦である。
とくに $(R, \mathfrak m)$ がネーター局所環ならば、完備化
$\lim_n R/\mathfrak m^n$ は忠実平坦である。
\end{lemma}

\begin{proof}
補題 \ref{algebra-lemma-completion-flat} により平坦である。合成 $R \to R^\wedge \to R/I$ を考える。
ここで最後の写像は射影 $R^\wedge \to R/I$ である。これにより、$R$ の任意の極大イデアルは
$\Spec(R^\wedge) \to \Spec(R)$ の像に入ることが分かる。従って補題 \ref{algebra-lemma-ff} により
写像は忠実平坦である。
\end{proof}

\begin{lemma}
\label{algebra-lemma-completion-complete}
$R$ をネーター環、$I$ を $R$ のイデアル、$M$ を $R$-加群とする。
$I$ に関する $M$ の完備化 $M^\wedge$ は $I$-進完備であり、
$I^n M^\wedge = (I^nM)^\wedge$ および $M^\wedge/I^nM^\wedge = M/I^nM$ が成り立つ。
\end{lemma}

\begin{proof}
$I$ は有限生成イデアルなので、これは補題
\ref{algebra-lemma-hathat-finitely-generated} の特別な場合である。
\end{proof}

\begin{lemma}
\label{algebra-lemma-completion-Noetherian}
$I$ を環 $R$ のイデアルとし、次を仮定する：
\begin{enumerate}
\item $R/I$ はネーター環である、
\item $I$ は有限生成である。
\end{enumerate}
このとき $I$ に関する $R$ の完備化 $R^\wedge$ はネーター環であり、
$IR^\wedge$ に関して完備である。
\end{lemma}

\begin{proof}
補題 \ref{algebra-lemma-hathat-finitely-generated} により、$R^\wedge$ は $I$-進完備である。
従って $IR^\wedge$-進にも完備である。$R^\wedge/IR^\wedge = R/I$ はネーター的なので、
$R$ を $R^\wedge$ で置き換えることにより、仮定 (1)、(2) に加えて $R$ 自身も
$I$-進完備であると仮定してよい。

\medskip\noindent
$f_1, \ldots, f_t$ を $I$ の生成元とする。このとき環の全射
$R/I[T_1, \ldots, T_t] \to \bigoplus I^n/I^{n + 1}$ があり、$T_i$ を
$\overline{f}_i \in I/I^2$ に写す。従って $\bigoplus I^n/I^{n + 1}$ はネーター環である。
$J \subset R$ をイデアルとし、イデアル
$$
\bigoplus J \cap I^n/J \cap I^{n + 1} \subset \bigoplus I^n/I^{n + 1}.
$$
を考える。$\overline{g}_1, \ldots, \overline{g}_m$ をこのイデアルの生成元とする。
$\overline{g}_j$ は次数 $d_j$ の斉次元として選べ、さらに
$\overline{g}_j \in J \cap I^{d_j}/J \cap I^{d_j + 1}$ に写る
$g_j \in J \cap I^{d_j}$ を選べる。$g_1, \ldots, g_m$ が $J$ を生成すると主張する。

\medskip\noindent
$x \in J \cap I^n$ とする。このとき、$x - \sum a_j g_j \in J \cap I^{n + 1}$ となる
$a_j \in I^{\max(0, n - d_j)}$ が存在する。その理由は、$g_1, \ldots, g_m$ の選び方により
$J \cap I^n/J \cap I^{n + 1}$ が
$\sum \overline{g}_j I^{n - d_j}/I^{n - d_j + 1}$ に等しいからである。
% [P01-E0480] Frozen unrestricted negative ideal-power exponents are retained.
従って $x \in J$ から始めて、$a_{j, n} \in I^{\max(0, n - d_j)}$ であり
$$
x =
\sum\nolimits_{n = 0, \ldots, N}
\sum\nolimits_{j = 1, \ldots, m} a_{j, n} g_j \bmod I^{N + 1}
$$
を満たすベクトルの列 $(a_{1, n}, \ldots, a_{m, n})_{n \geq 0}$ を見つけられる。
$A_j = \sum_{n \geq 0} a_{j, n}$ とおく。$R$ は完備なので
$x = \sum A_j g_j$ である。従って $J$ は有限生成であり、証明が終わる。
\end{proof}

\begin{lemma}
\label{algebra-lemma-completion-Noetherian-Noetherian}
$R$ をネーター環、$I$ を $R$ のイデアルとする。$I$ に関する $R$ の完備化
$R^\wedge$ はネーター的である。
\end{lemma}

\begin{proof}
これは補題 \ref{algebra-lemma-completion-Noetherian} の帰結である。次のように直接見ることもできる。
$f_1, \ldots, f_n$ を $I$ の生成元として、写像
$$
R[[x_1, \ldots, x_n]] \longrightarrow R^\wedge,
\quad
x_i \longmapsto f_i.
$$
を考える。これは well-defined な全射環写像である（詳細は省略する）。
$R[[x_1, \ldots, x_n]]$ はネーター的なので（補題
\ref{algebra-lemma-Noetherian-power-series} を参照）、証明が終わる。
\end{proof}

\noindent
$R \to S$ を局所環 $(R, \mathfrak m)$ から $(S, \mathfrak n)$ への局所準同型とする。
$S^\wedge$ を $\mathfrak n$ に関する $S$ の完備化とする。一般に $S^\wedge$ は
$S$ の $\mathfrak m$-進完備化ではない。ある $t \geq 1$ に対して
$\mathfrak n^t \subset \mathfrak mS$ ならば、補題
\ref{algebra-lemma-change-ideal-completion} により
$S^\wedge = \lim S/\mathfrak m^nS$ となる。場合によっては、これからさらに
$S^\wedge$ が $R^\wedge$ 上有限であることも従う。

\begin{lemma}
\label{algebra-lemma-finite-after-completion}
$R \to S$ を局所環 $(R, \mathfrak m)$ から $(S, \mathfrak n)$ への局所準同型とする。
$R^\wedge$、resp.\ $S^\wedge$ を、それぞれ $\mathfrak m$、resp.\ $\mathfrak n$ に関する
$R$、resp.\ $S$ の完備化とする。$\mathfrak m$ と $\mathfrak n$ が有限生成であり、
$\dim_{\kappa(\mathfrak m)} S/\mathfrak mS < \infty$ ならば、
\begin{enumerate}
\item $S^\wedge$ は $S$ の $\mathfrak m$-進完備化に等しい、および
\item $S^\wedge$ は有限 $R^\wedge$-加群である。
\end{enumerate}
\end{lemma}

\begin{proof}
$R \to S$ は局所環写像なので $\mathfrak mS \subset \mathfrak n$ である。
仮定 $\dim_{\kappa(\mathfrak m)} S/\mathfrak mS < \infty$ により
$S/\mathfrak mS$ はアルティン環である。補題
\ref{algebra-lemma-finite-dimensional-algebra} を参照せよ。従ってその次元は $0$ である。
% [P01-E0481] The frozen source's subjectless \"Hence has dimension 0\" is rendered by its evident subject.
補題 \ref{algebra-lemma-Noetherian-dimension-0} を参照せよ。従って
$\mathfrak n = \sqrt{\mathfrak mS}$ である。これと $\mathfrak n$ が有限生成であることから、
ある $t \geq 1$ に対して $\mathfrak n^t \subset \mathfrak mS$ が従う。補題
\ref{algebra-lemma-change-ideal-completion} により、$S^\wedge$ を $S$ の $\mathfrak m$-進完備化と
同一視できる。$\mathfrak m$ は有限生成なので、補題
\ref{algebra-lemma-hathat-finitely-generated} により $S^\wedge$ と $R^\wedge$ は
$\mathfrak m$-進完備である。この時点で $S^\wedge$ を $R^\wedge$-加群とみなし、
補題 \ref{algebra-lemma-finite-over-complete-ring} を適用すれば結論を得る。
\end{proof}

\begin{lemma}
\label{algebra-lemma-completion-finite-extension}
$R$ をネーター環、$R \to S$ を有限環写像とする。$\mathfrak p \subset R$ を素イデアルとし、
$\mathfrak q_1, \ldots, \mathfrak q_m$ を $\mathfrak p$ の上にある $S$ の素イデアルとする
（補題 \ref{algebra-lemma-finite-finite-fibres}）。このとき
$$
R_\mathfrak p^\wedge \otimes_R S =
(S_\mathfrak p)^\wedge =
S_{\mathfrak q_1}^\wedge \times \ldots \times S_{\mathfrak q_m}^\wedge
$$
である。ここで $(S_\mathfrak p)^\wedge$ は $\mathfrak p$ に関する完備化であり、局所環
$R_\mathfrak p$ と $S_{\mathfrak q_i}$ はそれぞれの極大イデアルに関して完備化されている。
\end{lemma}

\begin{proof}
$R$ を局所化 $R_\mathfrak p$ で、$S$ を
$S_\mathfrak p = S \otimes_R R_\mathfrak p$ で置き換えてよい。従って $R$ はネーター局所環で、
$\mathfrak p = \mathfrak m$ はその極大イデアルであると仮定してよい。
補題 \ref{algebra-lemma-finite-after-completion} により、$\mathfrak q_iS_{\mathfrak q_i}$-進完備化
$S_{\mathfrak q_i}^\wedge$ は $\mathfrak m$-進完備化に等しい。各 $n \geq 1$ に対して、
$S/\mathfrak m^nS$ の素イデアルは $S$ の極大イデアル
$\mathfrak q_1, \ldots, \mathfrak q_m$ と一対一に対応する（$R$ 上の $S$ に対する
going up による。補題 \ref{algebra-lemma-integral-going-up} を参照）。従って補題
\ref{algebra-lemma-artinian-finite-length} により
$S/\mathfrak m^nS = \prod S_{\mathfrak q_i}/\mathfrak m^nS_{\mathfrak q_i}$ である
（たとえば、$S/\mathfrak m^nS$ がアルティン的であることを見るには命題
\ref{algebra-proposition-dimension-zero-ring} を用いる）。従って $S$ の $\mathfrak m$-進完備化
$S^\wedge$ は $\prod S_{\mathfrak q_i}^\wedge$ に等しい。最後に、補題
\ref{algebra-lemma-completion-tensor} により $R^\wedge \otimes_R S = S^\wedge$ である。
\end{proof}

\begin{lemma}
\label{algebra-lemma-split-completed-sequence}
$R$ を環、$I \subset R$ をイデアルとし、
$0 \to K \to P \to M \to 0$ を $R$-加群の短完全列とする。$M$ が $R$ 上平坦で、
$M/IM$ が射影的 $R/I$-加群ならば、$I$-進完備化の列
$$
0 \to K^\wedge \to P^\wedge \to M^\wedge \to 0
$$
は分裂完全列である。
\end{lemma}

\begin{proof}
$M$ は平坦なので、各列
$$
0 \to K/I^nK \to P/I^nP \to M/I^nM \to 0
$$
は短完全である。補題 \ref{algebra-lemma-flat-tor-zero} を参照せよ。また、列
$0 \to K^\wedge \to P^\wedge \to M^\wedge \to 0$ も短完全である。補題
\ref{algebra-lemma-completion-generalities} を参照せよ。$s_{n + 1} \bmod I^n = s_n$ を満たす分裂
$s_n : M/I^nM \to P/I^nP$ を見つければ十分である。これらの $s_n$ を $n$ に関する帰納法で
構成する。$M/IM$ は射影的 $R/I$-加群なので存在する任意の分裂 $s_1$ を選ぶ。
ある $n > 0$ に対して $s_n$ が与えられたと仮定し、
$P_{n + 1} = \{x \in P \mid x \bmod I^nP \in \Im(s_n)\}$ とおく。写像
$\pi : P_{n + 1}/I^{n + 1}P_{n + 1} \to M/I^{n + 1}M$ は全射である（詳細は省略する）。
補題 \ref{algebra-lemma-lift-projective} により $M/I^{n + 1}M$ は射影的
$R/I^{n + 1}$-加群なので、$\pi$ の切断
$t : M/I^{n + 1}M \to P_{n + 1}/I^{n + 1}P_{n + 1}$ を選べる。
$s_{n + 1}$ を $t$ と標準写像
$P_{n + 1}/I^{n + 1}P_{n + 1} \to P/I^{n + 1}P$ との合成に等しくおけばよい。
\end{proof}

\begin{lemma}
\label{algebra-lemma-complete-modulo-nilpotent}
$A$ をネーター環、$I, J \subset A$ をイデアルとする。$A$ が $I$-進完備で、
$A/I$ が $J$-進完備ならば、$A$ は $J$-進完備である。
\end{lemma}

\begin{proof}
$B$ を $A$ の $(I + J)$-進完備化とする。補題 \ref{algebra-lemma-completion-flat} により
$B/IB$ は $A/I$ の $J$-進完備化であり、従って仮定により $A/I$ と同型である。
さらに補題 \ref{algebra-lemma-complete-by-sub} により $B$ は $I$-進完備である。従って補題
\ref{algebra-lemma-finite-over-complete-ring} により $B$ は有限 $A$-加群である。
Nakayama の補題（補題 \ref{algebra-lemma-NAK}；$I$ が $A$ の Jacobson 根基に
含まれることは補題 \ref{algebra-lemma-radical-completion} による）により、
$A \to B$ は全射である。
% [P01-E0482] The frozen source's missing complementizer in the Nakayama parenthesis is rendered by its evident sense.
補題 \ref{algebra-lemma-completion-flat} により写像 $A \to B$ は平坦である。
$\Spec(B) \to \Spec(A)$ の像は $V(I)$ を含み、$I$ は $A$ の Jacobson 根基に含まれるので、
$A \to B$ は忠実平坦である（補題 \ref{algebra-lemma-ff-rings}）。従って $A \to B$ は単射である。
従って $A$ は $I + J$ に関して完備であり、まして $J$ に関して完備である。
\end{proof}




\section{加群の極限をとること}
\label{algebra-section-limits}

\noindent
この節では、加群の極限をとるときに何が起こるかを論じる。

\begin{lemma}
\label{algebra-lemma-limit-complete-pre}
$I \subset A$ を環の有限生成イデアルとする。
$(M_n)$ を $I^n M_n = 0$ を満たす $A$-加群の逆系とする。
このとき $M = \lim M_n$ は $I$-進完備である。
\end{lemma}

\begin{proof}
$M \to M/I^nM \to M_n$ がある。極限をとると
$M \to M^\wedge \to M$ を得る。従って $M$ は $M^\wedge$ の直和因子である。
$M^\wedge$ は補題 \ref{algebra-lemma-hathat-finitely-generated} により
$I$-進完備なので、$M$ もそうである。
\end{proof}

\begin{lemma}
\label{algebra-lemma-limit-complete}
$I \subset A$ を環の有限生成イデアルとする。
$(M_n)$ を $M_n = M_{n + 1}/I^nM_{n + 1}$ を満たす
$A$-加群の逆系とする。$M = \lim M_n$ とおく。
このとき $M/I^nM = M_n$ であり、$M$ は $I$-進完備である。
\end{lemma}

\begin{proof}
補題 \ref{algebra-lemma-limit-complete-pre} により $M$ は $I$-進完備である。
遷移写像は全射なので、写像 $M \to M_n$ は全射である。
$N_n$ を定める短完全列の逆系
$$
0 \to N_n \to M \to M_n \to 0
$$
を考える。$M_n = M_{n + 1}/I^nM_{n + 1}$ なので、写像
$N_{n + 1} + I^nM \to N_n$ は全射である。従って
$N_{n + 1}/(N_{n + 1} \cap I^{n + 1}M) \to N_n/(N_n \cap I^nM)$
は全射である。短完全列
$$
0 \to N_n/(N_n \cap I^nM) \to M/I^nM \to M_n \to 0
$$
の逆極限をとると、完全列
$$
0 \to \lim N_n/(N_n \cap I^nM) \to M^\wedge \to M
$$
を得る。$M$ は $I$-進完備なので、
$\lim N_n/(N_n \cap I^nM) = 0$ と結論する。従って遷移写像の全射性により、
すべての $n$ に対して $N_n/(N_n \cap I^nM) = 0$ を得る。
従って望んだとおり $M_n = M/I^nM$ である。
\end{proof}

\begin{lemma}
\label{algebra-lemma-finiteness-graded}
$A$ をネーター次数付き環とする。$I \subset A_+$ を斉次イデアルとする。
$(N_n)$ を $N_n = N_{n + 1}/I^n N_{n + 1}$ を満たす有限次数付き
$A$-加群の逆系とする。このとき、すべての $n$ に対して次数付き加群として
$N_n = N/I^nN$ となる有限次数付き $A$-加群 $N$ が存在する。
\end{lemma}

\begin{proof}
$N_1$ を生成する、次数 $d_1, \ldots, d_r$ の斉次元
$x_{1, 1}, \ldots, x_{1, r} \in N_1$ と $r$ を選ぶ。
遷移写像は全射なので、$x_{1, i}$ を持ち上げる斉次元の両立する系
$x_{n, i} \in N_n$ を選べる。次数付き Nakayama の補題
（補題 \ref{algebra-lemma-graded-NAK}）により、$N_n$ は次数
$d_1, \ldots, d_r$ にある元 $x_{n, 1}, \ldots, x_{n, r}$ により生成される。
従って $m \leq n$ に対して、写像 $N_n \to N_n/I^m N_n$ は
次数 $< \min(d_i) + m$ において同型であることが分かる
（それらの次数では $I^mN_n$ が零だからである）。従って次数 $d$ の部分の逆系
$$
\ldots
= N_{2 + d - \min(d_i), d}
= N_{1 + d - \min(d_i), d}
= N_{d - \min(d_i), d} \to N_{-1 + d - \min(d_i), d} \to \ldots
$$
は示されたように安定する。
% [P01-E0483] The frozen stabilization index is retained; its strict-bound and N_0 defect remains sidecar-only.
$N$ を、その第 $d$ 次斉次部分がこの安定値である次数付き $A$-加群とする。
特に $N$ の元 $x_i = \lim x_{n, i}$ を得る。$x_i$ が $N$ を生成すると主張する。
任意の $x \in N_d$ は $x_1, \ldots, x_r$ の線形結合である。
実際、$x_{d - \min(d_i), i}$ が $N_{d - \min(d_i)}$ を生成するので、
これが成り立つ $N_{d - \min(d_i), d}$ で確かめればよい。
最後に読者は、全射 $N/I^nN \to N_n$ が同型であることを、前と同様に
各次数で何が起こるかを確かめることによって検証できる。詳細は省略する。
\end{proof}

\begin{lemma}
\label{algebra-lemma-daniel-litt}
$A$ を次数付き環とする。$I \subset A_+$ を斉次イデアルとする。
$A' = \lim A/I^n$ と書く。$(G_n)$ を、$G_n$ が $I^n$ によって零化される
次数付き $A$-加群の逆系とする。$M$ を次数付き $A$-加群とし、
$\varphi_n : M \to G_n$ を両立する次数付き $A$-加群写像の系とする。
誘導写像
$$
\varphi : M \otimes_A A' \longrightarrow \lim G_n
$$
が同型ならば、すべての $d \in \mathbf{Z}$ に対して
$M_d \to \lim G_{n, d}$ は同型である。
\end{lemma}

\begin{proof}
慣例により次数付き環は次数 $\geq 0$ にあり、次数付き加群は任意の次数に
非零な部分をもってよい。節 \ref{algebra-section-graded} を参照せよ。
$G_n$ は $I^n$ によって零化されるので $\lim G_n$ は
$A'$ 上の加群であり、従って写像 $\varphi$ が存在する。
もう一つ覚えておくと有用なのは
$$
\bigoplus\nolimits_{d \in \mathbf{Z}} \lim G_{n, d} \subset
\lim G_n \subset
\prod\nolimits_{d \in \mathbf{Z}} \lim G_{n, d}
$$
が成り立つことである。ここで下付き添字 ${\ }_d$ は第 $d$ 次斉次部分を表す。

\medskip\noindent
単射性。$x \in M_d$ とする。$x \mapsto 0$ が $\lim G_{n, d}$ において成り立つならば、
$M \otimes_A A'$ において $x \otimes 1 = 0$ である。このとき $x \in M'$ となる
有限生成部分加群 $M' \subset M$ で、$M' \otimes_A A'$ において
$x \otimes 1$ が零になるものを見つけられる。$M'$ は次数
$d_1, \ldots, d_r$ にある斉次元で生成されるとする。$n = d - \min(d_i) + 1$ とおく。
$A'$ から $A/I^n$ への写像があり、$A \to A/I^n$ は
次数 $\leq n - 1$ において同型なので、
$M' \to M' \otimes_A A'$ は次数 $\leq n - 1$ において単射である。
% [P01-E0484] The frozen module-degree bound is retained; its missing generator-degree shift remains sidecar-only.
従って望んだとおり $x = 0$ である。

\medskip\noindent
全射性。$y \in \lim G_{n, d}$ とする。$M \otimes_A A'$ において $y$ に写る有限和
$\sum x_i \otimes f'_i$ を選ぶ。$x_i$ は斉次であると仮定してよく、その次数を $d_i$ とする。
$A'$ は次数付き環ではないが、次数付き環 $A/I^nA$ の極限であり、さらに任意の
固定された次数において遷移写像は最終的に同型になることに注意する（上を参照）。
従って
$$
A = \bigoplus\nolimits_{d \geq 0} A_d \subset A' \subset
\prod\nolimits_{d \geq 0} A_d
$$
を得る。従って
$$
f'_i = \sum\nolimits_{j = 0, \ldots, d - d_i - 1} f_{i, j} + f_i + g'_i
$$
と書ける。ここで $f_{i, j} \in A_j$、$f_i \in A_{d - d_i}$ であり、
$g'_i \in A'$ は $\prod_{j \leq d - d_i} A_j$ において零に写る。
ここで $\varphi_n(\sum_{i, j} f_{i, j}x_i) \in G_n$ を計算すると、
次数 $< d$ の斉次元の和を得る。従って
$\varphi(\sum x_i \otimes f_{i, j})$ は $\lim G_{n, d}$ において零に写る。
同様に計算すると、元 $\varphi(\sum x_i \otimes g'_i)$ は
$\prod_{d' \leq d} \lim G_{n, d'}$ において零に写る。
$\varphi(\sum x_i \otimes f'_i)$ が $y$ であることを知っているので、
望んだとおり $\sum f_ix_i \in M_d$ は $y$ に写ると結論する。
\end{proof}




\section{平坦性の判定法}
\label{algebra-section-criteria-flatness}

\noindent
この節ではネーターの場合にいくつかの重要な技術的補題を証明する。
これらは節 \ref{algebra-section-more-flatness-criteria} で非ネーターの場合へ（部分的に）一般化する。

\begin{lemma}
\label{algebra-lemma-mod-injective}
$R \to S$ を局所環の局所準同型とし、$S$ はネーターであると仮定する。
$\mathfrak m$ を $R$ の極大イデアルと書く。
% [P01-E0485] Three frozen notation declarations omit “by”; Japanese uses the ordinary declaration.
$M$ を平坦 $R$-加群、$N$ を有限 $S$-加群とする。
$u : N \to M$ を $R$-加群の写像とする。
$\overline{u} : N/\mathfrak m N \to M/\mathfrak m M$ が単射ならば、
$u$ は単射である。この場合 $M/u(N)$ は $R$ 上平坦である。
\end{lemma}

\begin{proof}
まず、すべての $n \geq 1$ に対して
$u_n : N/{\mathfrak m}^nN \to M/{\mathfrak m}^nM$ が単射であると主張する。
帰納法で証明する。基底の場合は $\overline{u} = u_1$ が単射であることである。
$M$ が $R$ 上平坦であるという仮定により、短完全列
$0 \to M \otimes_R {\mathfrak m}^n/{\mathfrak m}^{n + 1}
\to M/{\mathfrak m}^{n + 1}M \to M/{\mathfrak m}^n M \to 0$
がある。また
$M \otimes_R {\mathfrak m}^n/{\mathfrak m}^{n + 1}
= M/{\mathfrak m}M \otimes_{R/{\mathfrak m}}
{\mathfrak m}^n/{\mathfrak m}^{n + 1}$ である。
$N$ に対しても同様の完全列
$N \otimes_R {\mathfrak m}^n/{\mathfrak m}^{n + 1}
\to N/{\mathfrak m}^{n + 1}N \to N/{\mathfrak m}^n N \to 0$
があるが、左端の零はない。また
$N \otimes_R {\mathfrak m}^n/{\mathfrak m}^{n + 1}
= N/{\mathfrak m}N \otimes_{R/{\mathfrak m}}
{\mathfrak m}^n/{\mathfrak m}^{n + 1}$ である。従って $u_n$ と写像
$\overline{u} \otimes \text{id}_{{\mathfrak m}^n/{\mathfrak m}^{n + 1}}$
がともに単射なので、写像 $u_{n + 1}$ も単射である。

\medskip\noindent
環 $S$ 上の $N$ とイデアル $\mathfrak mS$ に Krull の交叉定理
（補題 \ref{algebra-lemma-intersect-powers-ideal-module-zero}）を適用すると
$\bigcap \mathfrak m^nN = 0$ を得る。従ってすべての $n$ に対する
$u_n$ の単射性から $u$ は単射である。

\medskip\noindent
$M/u(N)$ が $R$ 上平坦であることを示すには、すべてのイデアル
$I \subset R$ に対して $\text{Tor}_1^R(M/u(N),\allowbreak R/I)\allowbreak = 0$ を示せば十分である。
補題 \ref{algebra-lemma-characterize-flat} を参照せよ。短完全列
$$
0 \to N \xrightarrow{u} M \to M/u(N) \to 0
$$
と $M$ の平坦性から、Tor の完全列
$$
0 \to \text{Tor}_1^R(M/u(N), R/I) \to N/IN \to M/IM
$$
を得る。補題 \ref{algebra-lemma-long-exact-sequence-tor} を参照せよ。
従って $N/IN$ が $M/IM$ に単射的に写ることを示せば十分である。
$R/I \to S/IS$ は $S/IS$ がネーターである局所環の局所準同型であり、
$N/IN \to M/IM$ は $R/I$-加群の写像であり、$N/IN$ は $S/IS$ 上有限であり、
$M/IM$ は $R/I$ 上平坦であり、さらに
$u \bmod I : N/IN \to M/IM$ は $\mathfrak m$ を法として単射である。
従って証明の最初の部分を $u \bmod I$ に適用でき、結論を得る。
\end{proof}

\begin{lemma}
\label{algebra-lemma-grothendieck}
$R \to S$ をネーター局所環の平坦な局所環準同型とする。
$\mathfrak m$ を $R$ の極大イデアルと書く。
$f \in S$ が $S/{\mathfrak m}S$ における非零因子であると仮定する。
このとき $S/fS$ は $R$ 上平坦であり、$f$ は $S$ における非零因子である。
\end{lemma}

\begin{proof}
補題 \ref{algebra-lemma-mod-injective} から直ちに従う。
\end{proof}

\begin{lemma}
\label{algebra-lemma-grothendieck-regular-sequence}
$R \to S$ をネーター局所環の平坦な局所環準同型とする。
$\mathfrak m$ を $R$ の極大イデアルと書く。
$f_1, \ldots, f_c$ を $S$ の元の列とし、その像
$\overline{f}_1, \ldots, \overline{f}_c$ が
$S/{\mathfrak m}S$ において正則列をなすと仮定する。
このとき $f_1, \ldots, f_c$ は $S$ において正則列をなし、各商
$S/(f_1, \ldots, f_i)$ は $R$ 上平坦である。
\end{lemma}

\begin{proof}
帰納法と補題 \ref{algebra-lemma-grothendieck} による。
\end{proof}

\begin{lemma}
\label{algebra-lemma-free-fibre-flat-free}
$R \to S$ をネーター局所環の局所準同型とする。
$\mathfrak m$ を $R$ の極大イデアルとする。$M$ を非零な有限 $S$-加群とする。
(a) $M/\mathfrak mM$ が自由 $S/\mathfrak mS$-加群であり、
(b) $M$ が $R$ 上平坦であると仮定する。
このとき $M$ は自由であり、$S$ は $R$ 上平坦である。
\end{lemma}

\begin{proof}
$\overline{x}_1, \ldots, \overline{x}_n$ を自由加群
$M/\mathfrak mM$ の基底とする。$x_i$ が $\overline{x}_i$ に写るような
$x_1, \ldots, x_n \in M$ を選ぶ。第 $i$ 標準基底ベクトルを $x_i$ に写す写像
$u : S^{\oplus n} \to M$ をとる。補題 \ref{algebra-lemma-mod-injective} により
$u$ は単射である。一方、Nakayama の補題 \ref{algebra-lemma-NAK} によりこの写像は全射である。
補題が従う。
\end{proof}

\begin{lemma}
\label{algebra-lemma-complex-exact-mod}
$R \to S$ をネーター局所環の局所準同型とする。
$\mathfrak m$ を $R$ の極大イデアルとする。
$0 \to F_e \to F_{e-1} \to \ldots \to F_0$ を有限 $S$-加群の有限複体とする。
各 $F_i$ が $R$-平坦であり、複体
$0 \to F_e/\mathfrak m F_e \to F_{e-1}/\mathfrak m F_{e-1}
\to \ldots \to F_0 / \mathfrak m F_0$ が完全であると仮定する。
このとき $0 \to F_e \to F_{e-1} \to \ldots \to F_0$ は完全であり、
さらに加群 $\Coker(F_1 \to F_0)$ は $R$-平坦である。
\end{lemma}

\begin{proof}
$e$ に関する帰納法による。$e = 1$ ならば、これはまさに
補題 \ref{algebra-lemma-mod-injective} である。$e > 1$ ならば、補題
\ref{algebra-lemma-mod-injective} により $F_e \to F_{e-1}$ は単射であり、
$C = \Coker(F_e \to F_{e-1})$ は $R$ 上平坦な有限 $S$-加群である。
従って複体 $0 \to C \to F_{e-2} \to \ldots \to F_0$ に帰納法の仮定を適用できる。
$C \to F_{e-2}$ が単射であると分かり、複体の完全性が従う。
$F_1 \to F_0$ の余核の平坦性も同様に従う。
\end{proof}

\noindent
この節の残りでは、いわゆる ``{\it 平坦性の局所判定法}'' の二つの版を証明する。
以下の興味深い補題 \ref{algebra-lemma-CM-over-regular-flat} にも注意せよ。

\begin{lemma}
\label{algebra-lemma-prepare-local-criterion-flatness}
$R$ を極大イデアル $\mathfrak m$ と剰余体
$\kappa = R/\mathfrak m$ をもつ局所環とする。$M$ を $R$-加群とする。
$\text{Tor}_1^R(\kappa, M) = 0$ ならば、任意の有限長 $R$-加群 $N$ に対して
$\text{Tor}_1^R(N, M) = 0$ が成り立つ。
\end{lemma}

\begin{proof}
$N$ の長さに関する帰納法による。
% [P01-E0486] The frozen “descending induction” misdescription is rendered as the induction actually used.
$N$ の長さが $1$ ならば $N \cong \kappa$ なので結論は成り立つ。
$N$ の長さが $1$ より大きければ、$N$ を短完全列
$0 \to N' \to N \to N'' \to 0$ に組み込むことができ、ここで $N'$、$N''$ は
より小さい長さをもつ有限長 $R$-加群である。
$\text{Tor}_1^R(N, M)$ の消滅は $\text{Tor}_1^R(N', M)$ と
$\text{Tor}_1^R(N'', M)$ の消滅（帰納法の仮定）、および Tor 群の長完全列から従う。
補題 \ref{algebra-lemma-long-exact-sequence-tor} を参照せよ。
\end{proof}

\begin{lemma}[平坦性の局所判定法]
\label{algebra-lemma-local-criterion-flatness}
$R \to S$ をネーター局所環の局所準同型とする。
$\mathfrak m$ を $R$ の極大イデアルとし、$\kappa = R/\mathfrak m$ とおく。
$M$ を有限 $S$-加群とする。$\text{Tor}_1^R(\kappa, M) = 0$ ならば、
$M$ は $R$ 上平坦である。
\end{lemma}

\begin{proof}
$I \subset R$ をイデアルとする。補題 \ref{algebra-lemma-flat} により、
$I \otimes_R M \to M$ が単射であることを示せば十分である。
注意 \ref{algebra-remark-Tor-ring-mod-ideal} により、この核は
$\text{Tor}_1^R(M, R/I)$ に等しい。補題
\ref{algebra-lemma-prepare-local-criterion-flatness} により、有限余長のすべてのイデアルについて
$J \otimes_R M \to M$ が単射である。
% [P01-E0487] The frozen source leaves J syntactically unintroduced; Japanese supplies the evident quantification.

\medskip\noindent
$n >> 0$ を選び、次の短完全列
$$
0
\to I \cap \mathfrak m^n
\to I \oplus \mathfrak m^n
\to I + \mathfrak m^n
\to 0
$$
を考える。これは短完全列 $0 \to R \to R^{\oplus 2} \to R \to 0$ の部分列である。
従って図式
$$
\xymatrix{
(I\cap \mathfrak m^n) \otimes_R M \ar[r] \ar[d] &
I \otimes_R M \oplus \mathfrak m^n \otimes_R M \ar[r] \ar[d] &
(I + \mathfrak m^n) \otimes_R M \ar[d] \\
M \ar[r] &
M \oplus M \ar[r] &
M
}
$$
を得る。$I + \mathfrak m^n$ と $\mathfrak m^n$ は有限余長のイデアルであることに注意する。
従って図式追跡により
$\Ker((I \cap \mathfrak m^n)\otimes_R M \to M)
\to \Ker(I \otimes_R M \to M)$
が全射であることが分かる。特に
$K = \Ker(I \otimes_R M \to M)$ は、$(I \cap \mathfrak m^n) \otimes_R M$ の
$I \otimes_R M$ における像に含まれる。Artin--Rees の補題
\ref{algebra-lemma-Artin-Rees} により、ある $c > 0$ とすべての $n >> 0$ に対して
$K$ は $\mathfrak m^{n-c}(I \otimes_R M)$ に含まれる。
$I \otimes_R M$ は有限 $S$-加群であり（！）、$S$ はネーターなので、
これは $K = 0$ を意味する。実際、上のことから $K$ は
$I \otimes_R M$ の $\mathfrak mS$-進完備化において零に写る。
しかし補題 \ref{algebra-lemma-completion-faithfully-flat} により、$S$ からその
$\mathfrak mS$-進完備化への写像は忠実平坦である。
従って望んだとおり $K = 0$ である。
\end{proof}

\noindent
以下ではしばしば「$M/IM$ は $R/I$ 上平坦であり、
$\text{Tor}_1^R(R/I, M) = 0$ である」という条件に出会う。
次の補題はこれらの条件のいくつかの帰結を与える
（補題 \ref{algebra-lemma-prepare-local-criterion-flatness} の一般化である）。

\begin{lemma}
\label{algebra-lemma-what-does-it-mean}
$R$ を環とする。$I \subset R$ をイデアルとする。$M$ を $R$-加群とする。
$M/IM$ が $R/I$ 上平坦であり、$\text{Tor}_1^R(R/I, M) = 0$ ならば、
\begin{enumerate}
\item すべての $n \geq 1$ に対して $M/I^nM$ は $R/I^n$ 上平坦であり、
\item ある $m \geq 0$ に対して $I^m$ によって零化される任意の加群 $N$ に対して
$\text{Tor}_1^R(N, M) = 0$ が成り立つ。
\end{enumerate}
特に $I$ が冪零ならば、$M$ は $R$ 上平坦である。
\end{lemma}

\begin{proof}
$M/IM$ が $R/I$ 上平坦であり、$\text{Tor}_1^R(R/I, M) = 0$ であると仮定する。
$N$ を $R/I$-加群とする。短完全列
$$
0 \to K \to \bigoplus\nolimits_{i \in I} R/I \to N \to 0
$$
を選ぶ。
% [P01-E0488] The frozen ideal/index-set collision in this free presentation is retained.
$\text{Tor}$ の長完全列と $\text{Tor}_1^R(R/I, M)$ の消滅により
$$
0 \to \text{Tor}_1^R(N, M) \to K \otimes_R M \to
(\bigoplus\nolimits_{i \in I} R/I) \otimes_R M \to N \otimes_R M \to 0
$$
を得る。しかし $K$、$\bigoplus_{i \in I} R/I$、$N$ はすべて
$I$ によって零化されるので、
\begin{align*}
K \otimes_R M & = K \otimes_{R/I} M/IM, \\
(\bigoplus\nolimits_{i \in I} R/I) \otimes_R M & =
(\bigoplus\nolimits_{i \in I} R/I) \otimes_{R/I} M/IM, \\
N \otimes_R M & = N \otimes_{R/I} M/IM.
\end{align*}
$M/IM$ は $R/I$ 上平坦なので
$$
0 \to K \otimes_{R/I} M/IM \to
(\bigoplus\nolimits_{i \in I} R/I) \otimes_{R/I} M/IM \to
N \otimes_{R/} M/IM \to 0
$$
は完全である。
% [P01-E0489] The frozen malformed tensor base R/ is retained; the R/I repair remains sidecar-only.
これと上のことを合わせると、$I$ によって零化される任意の $R$-加群 $N$ に対して
$\text{Tor}_1^R(N, M) = 0$ と結論する。

\medskip\noindent
(2) を $m$ に関する帰納法で証明しよう。$m = 1$ の場合は前の段落で証明した。
$m > 1$ で $N$ が $I^m$ によって零化されるとき、
$N'$ と $N''$ が $I^{m - 1}$ によって零化される完全列
$0 \to N' \to N \to N'' \to 0$ を選べる。
例えば $N' = IN$、$N'' = N/IN$ ととれる。このとき完全列
$$
\text{Tor}_1^R(N', M) \to
\text{Tor}_1^R(N, M) \to
\text{Tor}_1^R(N'', M)
$$
と帰納法により、求める消滅が証明される。

\medskip\noindent
最後に (1) を証明する。$n \geq 1$ が与えられたとき、
$M/I^nM$ が $R/I^n$ 上平坦であることを示さなければならない。
言い換えれば、$I^n$ によって零化される $R$-加群 $N$ の圏上で関手
$N \mapsto N \otimes_{R/I^n} M/I^nM$ が完全であることを示さなければならない。
しかしそのような $N$ に対して
$N \otimes_{R/I^n} M/I^nM = N \otimes_R M$ である。
(2) の $\text{Tor}_1$ の消滅により、ある $I$ の冪によって零化される
$N$ の圏上で関手 $N \mapsto N \otimes_R M$ は完全であり、結論を得る。
\end{proof}

\begin{lemma}
\label{algebra-lemma-what-does-it-mean-again}
$R$ を環とする。$I \subset R$ をイデアルとする。
$M$ を $R$-加群とする。
\begin{enumerate}
\item $M/IM$ が $R/I$ 上平坦であり、
$M \otimes_R I/I^2 \to IM/I^2M$ が単射ならば、
$M/I^2M$ は $R/I^2$ 上平坦である。
\item $M/IM$ が $R/I$ 上平坦であり、
$n = 1, \ldots, k$ に対して
$M \otimes_R I^n/I^{n + 1} \to I^nM/I^{n + 1}M$ が単射ならば、
$M/I^{k + 1}M$ は $R/I^{k + 1}$ 上平坦である。
\end{enumerate}
\end{lemma}

\begin{proof}
第一の主張は、$R$ を $R/I^2$ で、$M$ を $M/I^2M$ で置き換えて
補題 \ref{algebra-lemma-what-does-it-mean} を適用した帰結である。このとき
$$
\text{Tor}_1^{R/I^2}(M/I^2M, R/I) =
\Ker(M \otimes_R I/I^2 \to IM/I^2M),
$$
を用いる。注意 \ref{algebra-remark-Tor-ring-mod-ideal} を参照せよ。
第二の主張も同様に従う。すなわち、$n$ に関する帰納法で、
$n = 1, \ldots, k$ に対して $M/I^{n + 1}M$ が
$R/I^{n + 1}$ 上平坦であることを示す。ここでは、すべての $n$ に対して
$$
\text{Tor}_1^{R/I^{n + 1}}(M/I^{n + 1}M, R/I^n) =
\Ker(M \otimes_R I^n/I^{n + 1} \to I^nM/I^{n + 1}M)
$$
であることを用いる。
\end{proof}

\begin{lemma}[局所判定法の変形]
\label{algebra-lemma-variant-local-criterion-flatness}
$R \to S$ をネーター局所環の局所準同型とする。
$I \not = R$ を $R$ のイデアルとする。$M$ を有限 $S$-加群とする。
$\text{Tor}_1^R(M, R/I) = 0$ であり、$M/IM$ が $R/I$ 上平坦ならば、
$M$ は $R$ 上平坦である。
\end{lemma}

\begin{proof}
第一の証明：補題 \ref{algebra-lemma-what-does-it-mean} により
$\text{Tor}_1^R(\kappa, M)$ は零である。ここで $\kappa$ は
$R$ の剰余体である。従って補題 \ref{algebra-lemma-local-criterion-flatness} により、
$M$ は $R$ 上平坦である。

\medskip\noindent
第二の証明：$\mathfrak m$ を $R$ の極大イデアルとする。
$\mathfrak m \otimes_R M \to M$ が単射であることを示し、その後
補題 \ref{algebra-lemma-local-criterion-flatness} を適用する。
$\sum f_i \otimes x_i \in \mathfrak m \otimes_R M$ であり、
$M$ において $\sum f_i x_i = 0$ であると仮定する。
$R/I$ 上の $M/IM$ に平坦性の等式判定法、補題 \ref{algebra-lemma-flat-eq} を適用すると、
$\overline{a}_{ij} \in R/I$ と $\overline{y}_j \in M/IM$ で
$x_i \bmod IM = \sum_j \overline{a}_{ij} \overline{y}_j $ および
$0 = \sum_i (f_i \bmod I) \overline{a}_{ij}$ を満たすものが存在する。
$a_{ij} \in R$ を $\overline{a}_{ij}$ の持ち上げとし、同様に
$y_j \in M$ を $\overline{y}_j$ の持ち上げとする。このとき
\begin{eqnarray*}
\sum f_i \otimes x_i
& = &
\sum f_i \otimes x_i +
\sum f_ia_{ij} \otimes y_j -
\sum f_i \otimes a_{ij} y_j
\\
& = &
\sum f_i \otimes (x_i - \sum a_{ij} y_j) +
\sum (\sum f_i a_{ij}) \otimes y_j
\end{eqnarray*}
となる。$x_i - \sum a_{ij} y_j \in IM$ かつ
$\sum f_i a_{ij} \in I$ なので、$I \otimes_R M$ のある元が
$\mathfrak m \otimes_R M$ における与えられた元
$\sum f_i \otimes x_i$ に写る。しかし仮定により
$I \otimes_R M \to M$ は単射である（注意
\ref{algebra-remark-Tor-ring-mod-ideal} を参照せよ）。これで結論を得る。
\end{proof}

\noindent
特に、補題 \ref{algebra-lemma-variant-local-criterion-flatness} の状況で、
$I = (x)$ が $R$ の非零因子である一つの元 $x$ によって生成されると仮定する。
このとき $\text{Tor}_1^R(M, R/(x)) = (0)$ であることと、
$x$ が $M$ 上の非零因子であることは同値である。

\begin{lemma}
\label{algebra-lemma-flat-module-powers}
$R \to S$ を環写像とする。$I \subset R$ をイデアルとする。
$M$ を $S$-加群とする。次を仮定する。
\begin{enumerate}
\item $R$ はネーター環である。
\item $S$ はネーター環である。
\item $M$ は有限 $S$-加群である。
\item 各 $n \geq 1$ に対して加群 $M/I^n M$ は $R/I^n$ 上平坦である。
\end{enumerate}
このとき、すべての $\mathfrak q \in V(IS)$ に対して局所化
$M_{\mathfrak q}$ は $R$ 上平坦である。
特に $S$ が局所環であり、$IS$ がその極大イデアルに含まれるならば、
$M$ は $R$ 上平坦である。
\end{lemma}

\begin{proof}
補題 \ref{algebra-lemma-variant-local-criterion-flatness} を用いる。
仮定により $M/IM$ は $R/I$ 上平坦である。従って
$\text{Tor}_1^R(M, R/I)$ が $\mathfrak q$ で局所化すると零になることを確かめれば十分である。
注意 \ref{algebra-remark-Tor-ring-mod-ideal} により、この Tor 群は
$K = \Ker(I \otimes_R M \to M)$ に等しい。
すべての $n \geq 1$ に対して
$I/I^n \otimes_{R/I^n} M/I^nM \to M/I^nM$ の核は零である。
従って $K$ の元は $I/I^n \otimes_{R/I^n} M/I^nM$ において零に写る。
また
$$
I/I^n \otimes_{R/I^n} M/I^nM = I/I^n \otimes_R M =
(I \otimes_R M)/I^{n - 1}(I \otimes_R M)
$$
なので、すべての $n \geq 1$ に対して
$K \subset I^{n - 1}(I \otimes_R M)$ と結論する。
Artin--Rees の補題、より正確には補題
\ref{algebra-lemma-intersection-powers-ideal-module} により、
望んだとおり $K_{\mathfrak q} = 0$ と結論する。
\end{proof}

\begin{lemma}
\label{algebra-lemma-surjective-on-tor-one}
$R \to R' \to R''$ を環写像とする。$M$ を $R$-加群とする。
$M \otimes_R R'$ が $R'$ 上平坦であると仮定する。このとき自然な写像
$\text{Tor}_1^R(M, R') \otimes_{R'} R'' \to
\text{Tor}_1^R(M, R'')$ は全射である。
\end{lemma}

\begin{proof}
$F_\bullet$ を $R$ 上の $M$ の自由分解とする。
複体 $F_2 \otimes_R R' \to F_1\otimes_R R' \to F_0 \otimes_R R'$ は
$\text{Tor}_1^R(M, R')$ を計算する。
複体 $F_2 \otimes_R R'' \to F_1\otimes_R R'' \to F_0 \otimes_R R''$ は
$\text{Tor}_1^R(M, R'')$ を計算する。
$F_i \otimes_R R' \otimes_{R'} R'' = F_i \otimes_R R''$ に注意する。
$K' = \Ker(F_1\otimes_R R' \to F_0 \otimes_R R')$ とおき、同様に
$K'' = \Ker(F_1\otimes_R R'' \to F_0 \otimes_R R'')$ とおく。
従って完全列
$$
0 \to K' \to F_1\otimes_R R' \to F_0 \otimes_R R' \to M \otimes_R R' \to 0.
$$
を得る。$M \otimes_R R'$ が $R'$ 上平坦であるという仮定により、列
$$
K' \otimes_{R'} R'' \to
F_1 \otimes_R R'' \to
F_0 \otimes_R R'' \to
M \otimes_R R'' \to 0
$$
はなお完全である。これは $K' \otimes_{R'} R'' \to K''$ が全射であることを意味する。
$\text{Tor}_1^R(M, R')$ は $K'$ の商であり、
$\text{Tor}_1^R(M, R'')$ は $K''$ の商なので、結論を得る。
\end{proof}

\begin{lemma}
\label{algebra-lemma-surjective-on-tor-one-trivial}
$R \to R'$ を環写像とする。$I \subset R$ をイデアルとし、
$I' = IR'$ とおく。$M$ を $R$-加群とし、$M' = M \otimes_R R'$ とおく。
自然な写像
$\text{Tor}_1^R(R'/I', M) \to \text{Tor}_1^{R'}(R'/I', M')$
は全射である。
\end{lemma}

\begin{proof}
$F_2 \to F_1 \to F_0 \to M \to 0$ を $R$ 上の $M$ の自由分解とする。
$F_i' = F_i \otimes_R R'$ とおく。列
$F_2' \to F_1' \to F_0' \to M' \to 0$ は $F_1'$ においてもはや
完全でないかもしれない。従って $R'$ 上の $M'$ の自由分解は、適当な自由
$R'$-加群 $F_2''$ に対して
$$
F_2' \oplus F_2'' \to F_1' \to F_0' \to M' \to 0
$$
という形をもつ。次に
$F_i \otimes_R R'/I' = F_i' / IF_i' = F_i'/I'F_i'$ に注意する。
従って複体 $F_2'/I'F_2' \to F_1'/I'F_1' \to F_0'/I'F_0'$ は
$\text{Tor}_1^R(M, R'/I')$ を計算する。一方
$F_i' \otimes_{R'} R'/I' = F_i'/I'F_i'$ であり、$F_2''$ についても同様である。
従って複体
$F_2'/I'F_2' \oplus F_2''/I'F_2'' \to F_1'/I'F_1' \to F_0'/I'F_0'$ は
$\text{Tor}_1^{R'}(M', R'/I')$ を計算する。複体間の垂直写像
$$
\xymatrix{
F_2'/I'F_2' \ar[r] \ar[d] &
F_1'/I'F_1' \ar[r] \ar[d] &
F_0'/I'F_0' \ar[d] \\
F_2'/I'F_2' \oplus F_2''/I'F_2'' \ar[r] &
F_1'/I'F_1' \ar[r] &
F_0'/I'F_0'
}
$$
が明らかにホモロジー上の全射を誘導するので、結論を得る。
% [P01-E0490] The frozen “cohomology” wording is rendered by the homology actually computed.
\end{proof}

\begin{lemma}
\label{algebra-lemma-another-variant-local-criterion-flatness}
局所ネーター環の局所準同型からなる可換図式
$$
\xymatrix{
S \ar[r] & S' \\
R \ar[r] \ar[u] & R' \ar[u]
}
$$
が与えられているとする。$I \subset R$ を真のイデアルとする。
$M$ を有限 $S$-加群とする。$I' = IR'$ および
$M' = M \otimes_S S'$ と書く。
% [P01-E0491] The frozen notation declaration omits “by”; Japanese uses the ordinary declaration.
次を仮定する。
\begin{enumerate}
\item $S'$ はテンソル積 $S \otimes_R R'$ の局所化である。
\item $M/IM$ は $R/I$ 上平坦である。
\item $\text{Tor}_1^R(M, R/I) \to \text{Tor}_1^{R'}(M', R'/I')$
は零写像である。
\end{enumerate}
このとき $M'$ は $R'$ 上平坦である。
\end{lemma}

\begin{proof}
$S'$ は $S \otimes_R R'$ の局所化なので、$M'$ は
$M \otimes_R R'$ の局所化である。補題 \ref{algebra-lemma-flat-base-change} により加群
$M/IM \otimes_{R/I} R'/I'
= M \otimes_R R' /I'(M \otimes_R R')$ は $R'/I'$ 上平坦である。
従って平坦加群の局所化は平坦なので、$M'/I'M'$ も $R'/I'$ 上平坦である。
補題 \ref{algebra-lemma-variant-local-criterion-flatness} により、
$\text{Tor}_1^{R'}(M', R'/I')$ が零であることを示せば十分である。
$M'$ は $M \otimes_R R'$ の局所化なので、最後の仮定により、
$\text{Tor}_1^R(M, R/I) \otimes_R R'
\to \text{Tor}_1^{R'}(M \otimes_R R', R'/I')$
が全射であることを示せば十分である。

\medskip\noindent
補題 \ref{algebra-lemma-surjective-on-tor-one-trivial} により
$\text{Tor}_1^R(M, R'/I') \to
\text{Tor}_1^{R'}(M \otimes_R R', R'/I')$ は全射である。
従って今や
$\text{Tor}_1^R(M, R/I) \otimes_R R'
\to \text{Tor}_1^R(M, R'/I')$
が全射であることを示せば十分である。これは環写像
$R \to R/I \to R'/I'$ と加群 $M$ を考えることにより、
補題 \ref{algebra-lemma-surjective-on-tor-one} から従う。
\end{proof}

\noindent
次の補題を、補題 \ref{algebra-lemma-criterion-flatness-fibre-nilpotent}
（冪零イデアルの場合）および補題 \ref{algebra-lemma-criterion-flatness-fibre}
（有限表示代数の場合）と比較されたい。

\begin{lemma}[ファイバーによる平坦性判定法；ネーターの場合]
\label{algebra-lemma-criterion-flatness-fibre-Noetherian}
$R$、$S$、$S'$ をネーター局所環とし、$R \to S \to S'$ を
局所環準同型とする。$\mathfrak m \subset R$ を極大イデアルとする。
$M$ を $S'$-加群とする。次を仮定する。
\begin{enumerate}
\item 加群 $M$ は $S'$ 上有限である。
\item 加群 $M$ は零でない。
\item 加群 $M/\mathfrak m M$ は $S/\mathfrak m S$ 上平坦である。
\item 加群 $M$ は $R$ 上平坦である。
\end{enumerate}
このとき $S$ は $R$ 上平坦であり、$M$ は $S$ 上平坦である。
\end{lemma}

\begin{proof}
$I = \mathfrak mS \subset S$ とおく。このとき (3) により
$M/IM$ は平坦 $S/I$-加群である。
$\mathfrak m \otimes_R S' \to I \otimes_S S'$ は全射なので、
$\mathfrak m \otimes_R M \to I \otimes_S M$ も全射である。
$$
\mathfrak m \otimes_R M \to I \otimes_S M \to M.
$$
を考える。$M$ は $R$ 上平坦なので合成は単射であり、従って両方の矢印が単射である。
特に $\text{Tor}_1^S(S/I, M) = 0$ である。注意
\ref{algebra-remark-Tor-ring-mod-ideal} を参照せよ。
補題 \ref{algebra-lemma-variant-local-criterion-flatness} により
$M$ は $S$ 上平坦である。Nakayama の補題 \ref{algebra-lemma-NAK} により
$M/\mathfrak m_{S'}M$ は零でないので、実際には補題 \ref{algebra-lemma-ff} により
$M$ は $S$ 上忠実平坦であることに注意する
（なぜなら $M/\mathfrak m_SM \not = 0$ を強制するからである）。

\medskip\noindent
完全列 $0 \to \mathfrak m \to R \to \kappa \to 0$ を考える。
これは完全列
$0 \to \text{Tor}_1^R(\kappa, S) \to \mathfrak m \otimes_R S \to I \to 0$
を与える。$M$ は $S$ 上平坦なので、これは完全列
$0 \to \text{Tor}_1^R(\kappa, S)\otimes_S M \to
\mathfrak m \otimes_R M \to I \otimes_S M \to 0$ を与える。
上のことにより、これは $\text{Tor}_1^R(\kappa, S)\otimes_S M = 0$ を意味する。
$M$ は $S$ 上忠実平坦なので、これは $\text{Tor}_1^R(\kappa, S) = 0$ を意味する。
補題 \ref{algebra-lemma-local-criterion-flatness} により、$S$ は $R$ 上平坦であると結論する。
\end{proof}

\begin{lemma}
\label{algebra-lemma-flatness-scallop}
$A$ を環、$M$ を $A$-加群、$f \in A$ とする。次を仮定する。
\begin{enumerate}
\item $f$ は $A$ および $M$ 上の非零因子である。
\item $M_f$ は平坦 $A_f$-加群である。
\item $M/fM$ は平坦 $A/fA$-加群である。
\end{enumerate}
このとき $M$ は平坦 $A$-加群である。「平坦」を「忠実平坦」に置き換えても同じである。
\end{lemma}

\begin{proof}
$0 \to A \to A \to A/fA \to 0$ と $0 \to M \to M \to M/fM \to 0$ は
完全なので、$i = 1$ および $i = 2$ に対して
$\text{Tor}_i^A(M, A/fA) = 0$ を得る。
補題 \ref{algebra-lemma-what-does-it-mean} により、$f$ によって零化される
すべての $A$-加群 $N$ に対して $\text{Tor}_1^A(M, N) = 0$ と結論する
（ここで $A/fA$ 上の $M/fM$ の平坦性を用いる）。
$f$ によって零化される $A$-加群 $N$ が与えられたとき、
$F$ が $A/fA$ のコピーの直和であるような $A$-加群の短完全列
$0 \to N' \to F \to N \to 0$ を選べる。完全列
$$
\text{Tor}_2^A(M, F) \to
\text{Tor}_2^A(M, N) \to
\text{Tor}_1^A(M, N')
$$
から、$f$ によって零化されるすべての $A$-加群 $N$ に対して
$\text{Tor}_2^A(M, N) = 0$ と結論する。
次に $K$ を任意の $A$-加群とする。写像 $f : K \to K$ を二つの短完全列
$$
0 \to K[f] \to K \to K' \to 0
\quad\text{and}\quad
0 \to K' \to K \to K/fK \to 0
$$
に分けられる。ここで $K' = K/K[f] \cong fK$ である。
Tor の完全列を適用して完全列
$$
\text{Tor}_1^A(K[f], M) \to
\text{Tor}_1^A(K, M) \to
\text{Tor}_1^A(K', M)
$$
および
$$
\text{Tor}_2^A(K/fK, M) \to
\text{Tor}_1^A(K', M) \to
\text{Tor}_1^A(K, M)
$$
を得る。$\text{Tor}_1^A(K[f], M)$ と
$\text{Tor}_2^A(K/fK, M)$ の消滅を用いると、
$f : K \to K$ は $\text{Tor}^A_1(M, K)$ 上の単射を誘導すると結論する。
言い換えれば、$\text{Tor}^A_1(M, K)$ が零であることと
$\text{Tor}^A_1(M, K) \otimes_A A_f$ が零であることは同値である。
補題 \ref{algebra-lemma-flat-base-change-tor} により
$$
\text{Tor}^A_1(M, K) \otimes_A A_f = \text{Tor}^{A_f}_1(M_f, K_f)  = 0
$$
である。最後の消滅は $A_f$ 上の $M_f$ の平坦性による
（補題 \ref{algebra-lemma-characterize-flat}）。
% [P01-E0492] Both frozen sentence fragments lacking a finite verb are rendered by their evident predicate.
すべての $A$-加群 $K$ に対して $\text{Tor}_1^A(M, K) = 0$ と結論する。
従って補題 \ref{algebra-lemma-characterize-flat} により $M$ は $A$ 上平坦である。

\medskip\noindent
忠実平坦加群の場合の議論は省略する。
\end{proof}




\begin{lemma}
\label{algebra-lemma-flatness-scallop-pre}
$A$ を環、$M$ を $A$-加群とする。
$I = (f_1, \ldots, f_r)$ を $r \geq 1$ 個の元で生成される $A$ のイデアルとする。
次を仮定する。
\begin{enumerate}
\item $i = 1, \ldots, r$ に対して $M_{f_i}$ は平坦 $A_{f_i}$-加群である。
\item $M/IM$ は平坦 $A/I$-加群である。
\item $i = 1, \ldots, r + 1$ に対して $\text{Tor}_i^A(M, A/I) = 0$ である。
\end{enumerate}
このとき $M$ は平坦 $A$-加群である。「平坦」を「忠実平坦」に置き換えても同じである。
\end{lemma}

\begin{proof}
補題 \ref{algebra-lemma-what-does-it-mean} により、$I$ によって零化される
すべての $A$-加群 $K$ に対して $\text{Tor}_1^A(M, K) = 0$ である。
$I$ によって零化される $A$-加群 $K$ が与えられたとき、
$F$ が $A/I$ のコピーの直和であるような $A$-加群の短完全列
$0 \to N \to F \to K \to 0$ を選べる。完全列
$$
\text{Tor}_i^A(M, F) \to
\text{Tor}_i^A(M, K) \to
\text{Tor}_{i - 1}^A(M, N)
$$
を得る。従って仮定 (3) と $i$ に関する帰納法を用いると、
$I$ によって零化されるすべての $A$-加群 $K$ および
$i = 1, \ldots, r + 1$ に対して $\text{Tor}_i^A(M, K) = 0$ と結論する。

\medskip\noindent
ある $1 \leq j \leq r$ に対して、$f_1, \ldots, f_j$ によって零化される
すべての $A$-加群 $K$ および $i = 1, \ldots, j + 1$ に対して
$\text{Tor}_i^A(M, K) = 0$ を示したと仮定する。
$K$ を $f_1, \ldots, f_{j - 1}$ によって零化される $A$-加群とする。
写像 $f_j : K \to K$ を二つの短完全列
$$
0 \to K[f_j] \to K \to K' \to 0
\quad\text{and}\quad
0 \to K' \to K \to K/f_jK \to 0
$$
に分けられる。ここで $K' = K/K[f_j] \cong f_jK$ である。
$1 \leq i \leq j$ とする。Tor の完全列を適用すると、完全列
$$
\text{Tor}_i^A(K[f_j], M) \to
\text{Tor}_i^A(K, M) \to
\text{Tor}_i^A(K', M)
$$
および
$$
\text{Tor}_{i + 1}^A(K/f_jK, M) \to
\text{Tor}_i^A(K', M) \to
\text{Tor}_i^A(K, M)
$$
を得る。$\text{Tor}_i^A(K[f_j], M)$ と
$\text{Tor}_{i + 1}^A(K/f_jK, M)$ の消滅を用いると、
$f_j : K \to K$ は $\text{Tor}^A_i(M, K)$ 上の単射を誘導すると結論する。
言い換えれば、$\text{Tor}^A_i(M, K)$ が零であることと、
$\text{Tor}^A_i(M, K) \otimes_A A_{f_j}$ が零であることは同値である。
補題 \ref{algebra-lemma-flat-base-change-tor} により
$$
\text{Tor}^A_i(M, K) \otimes_A A_{f_j} =
\text{Tor}^{A_{f_j}}_i(M_{f_j}, K_{f_j})  = 0
$$
である。最後の消滅は $A_{f_j}$ 上の $M_{f_j}$ の平坦性による
（補題 \ref{algebra-lemma-characterize-flat}）。
$f_1, \ldots, f_{j - 1}$ によって零化されるすべての $A$-加群 $K$ および
$i = 1, \ldots, j$ に対して $\text{Tor}_i^A(M, K) = 0$ と結論する。
$j$ に関する降下帰納法により、これは $j = 0$ に対して成り立つ。
すなわち、すべての $A$-加群 $K$ に対して
$\text{Tor}_1^A(M, K) = 0$ である。従って補題
\ref{algebra-lemma-characterize-flat} により $M$ は $A$ 上平坦である。

\medskip\noindent
忠実平坦加群の場合の議論は省略する。
\end{proof}




\section{底変換と平坦性}
\label{algebra-section-base-change-flat}

\noindent
底変換を行うとき平坦性に何が起こるかを扱ういくつかの補題を与える。

\begin{lemma}
\label{algebra-lemma-base-change-flat-up-down}
局所環の局所準同型からなる可換図式
$$
\xymatrix{
S \ar[r] & S' \\
R \ar[r] \ar[u] & R' \ar[u]
}
$$
が与えられているとする。$S'$ はテンソル積 $S \otimes_R R'$ の局所化であると仮定する。
$M$ を $S$-加群とし、$M' = S' \otimes_S M$ とおく。
\begin{enumerate}
\item $M$ が $R$ 上平坦ならば、$M'$ は $R'$ 上平坦である。
\item $M'$ が $R'$ 上平坦であり、$R \to R'$ が平坦ならば、
$M$ は $R$ 上平坦である。
\end{enumerate}
特に次が成り立つ。
\begin{enumerate}
\item[(3)] $S$ が $R$ 上平坦ならば、$S'$ は $R'$ 上平坦である。
\item[(4)] $R' \to S'$ と $R \to R'$ が平坦ならば、$S$ は $R$ 上平坦である。
\end{enumerate}
\end{lemma}

\begin{proof}
(1) の証明。$M$ が $R$ 上平坦ならば、補題 \ref{algebra-lemma-flat-base-change} により
$M \otimes_R R'$ は $R'$ 上平坦である。
$W \subset S \otimes_R R'$ を $W^{-1}(S \otimes_R R') = S'$ となる
乗法的部分集合とすると、$M' = W^{-1}(M \otimes_R R')$ である。
従って平坦加群の局所化として $M'$ は $R'$ 上平坦である。
補題 \ref{algebra-lemma-flat-localization} の (5) を参照せよ。
これで (1) が証明され、特に (3) が成り立つことが分かる。

\medskip\noindent
(2) の証明。$M'$ が $R'$ 上平坦であり、$R \to R'$ が平坦であると仮定する。
北西対角線に関して反転した図式に (3) を適用すると、$S \to S'$ は平坦である。
従って補題 \ref{algebra-lemma-local-flat-ff} により $S \to S'$ は忠実平坦である。
補題 \ref{algebra-lemma-flat} の判定法 (\ref{algebra-item-f-ideal}) を用いて、
$M$ が平坦であることを示す。$I \subset R$ をイデアルとする。
$I \otimes_R M \to M$ に核があるならば、
$(I \otimes_R M) \otimes_S S' \to M \otimes_S S' = M'$ にも核がある。
% [P01-E0493] The frozen conditional merely says “has a kernel”; its required nonzero qualifier remains sidecar-only.
$R \to R'$ は平坦なので $I \otimes_R R' = IR'$ であること、および
$$
(I \otimes_R M) \otimes_S S' =
(I \otimes_R R') \otimes_{R'} (M \otimes_S S') =
IR' \otimes_{R'} M'.
$$
であることに注意する。$R'$ 上の $M'$ の平坦性から、これは $M'$ に単射的に写ると分かる。
これで (2) の証明が終わり、従って (4) も成り立つ。
\end{proof}

\noindent
平坦性の局所判定法のさらにもう一つの応用を与える。

\begin{lemma}
\label{algebra-lemma-yet-another-variant-local-criterion-flatness}
局所環と局所準同型からなる可換図式
$$
\xymatrix{
S \ar[r] & S' \\
R \ar[r] \ar[u] & R' \ar[u]
}
$$
を考える。$M$ を有限 $S$-加群とする。次を仮定する。
\begin{enumerate}
\item 水平の矢印は平坦な環写像である。
\item $M$ は $R$ 上平坦である。
\item $\mathfrak m_R R' = \mathfrak m_{R'}$ である。
\item $R'$ と $S'$ はネーターである。
\end{enumerate}
このとき $M' = M \otimes_S S'$ は $R'$ 上平坦である。
\end{lemma}

\begin{proof}
$\mathfrak m_R \subset R$ かつ $R \to R'$ は平坦なので、仮定 (3) により
$\mathfrak m_R \otimes_R R' = \mathfrak m_R R' = \mathfrak m_{R'}$ を得る。
$M'$ は有限 $S'$-加群であり、補題
\ref{algebra-lemma-flatness-descends-more-general} により $R$ 上平坦であることに注意する。
従って $\mathfrak m_R \otimes_R M' \to M'$ は単射である。また
$$
\mathfrak m_R \otimes_R M' =
\mathfrak m_R \otimes_R R' \otimes_{R'} M' =
\mathfrak m_{R'} \otimes_{R'} M'
$$
を得る。従って $\mathfrak m_{R'} \otimes_{R'} M' \to M'$ は単射である。
これは $\text{Tor}_1^{R'}(\kappa_{R'}, M') = 0$ を示す
（注意 \ref{algebra-remark-Tor-ring-mod-ideal}）。
従って補題 \ref{algebra-lemma-local-criterion-flatness} により
$M'$ は $R'$ 上平坦である。
\end{proof}




\section{アルティン環上の平坦性判定法}
\label{algebra-section-flatness-artinian}

\noindent
アルティン環上の加群に対するいくつかの平坦性判定法を論じる。
アルティン局所環の極大イデアルは冪零なので、次の二つの補題は
アルティン局所環に適用できることに注意する。

\begin{lemma}
\label{algebra-lemma-local-artinian-basis-when-flat}
$(R, \mathfrak m)$ を冪零極大イデアル $\mathfrak m$ をもつ局所環とする。
$M$ を平坦 $R$-加群とする。$A$ が集合であり、
$x_\alpha \in M$、$\alpha \in A$ が $M$ の元の族ならば、次は同値である。
\begin{enumerate}
\item $\{\overline{x}_\alpha\}_{\alpha \in A}$ は
$R/\mathfrak m$ 上のベクトル空間 $M/\mathfrak mM$ の基底をなし、
\item $\{x_\alpha\}_{\alpha \in A}$ は $R$ 上の $M$ の基底をなす。
\end{enumerate}
\end{lemma}

\begin{proof}
(2) $\Rightarrow$ (1) は直ちに従う。(1) を仮定する。
Nakayama の補題 \ref{algebra-lemma-NAK} により、元 $x_\alpha$ は $M$ を生成する。
従って短完全列
$$
0 \to K \to \bigoplus\nolimits_{\alpha \in A} R \to M \to 0
$$
を得る。$R/\mathfrak m$ とテンソルし、補題
\ref{algebra-lemma-flat-tor-zero} を用いると $K/\mathfrak mK = 0$ を得る。
Nakayama の補題 \ref{algebra-lemma-NAK} により $K = 0$ と結論する。
\end{proof}

\begin{lemma}
\label{algebra-lemma-local-artinian-characterize-flat}
$R$ を冪零極大イデアルをもつ局所環とする。$M$ を $R$-加群とする。
次は同値である。
\begin{enumerate}
\item $M$ は $R$ 上平坦である。
\item $M$ は自由 $R$-加群である。
\item $M$ は射影 $R$-加群である。
\end{enumerate}
\end{lemma}

\begin{proof}
任意の射影加群は（自由加群の直和因子として）平坦であり、
任意の自由加群は射影的なので、平坦加群が自由であることを証明すれば十分である。
$M$ を平坦加群とする。$A$ を集合とし、$x_\alpha \in M$、$\alpha \in A$ を、
$\overline{x_\alpha} \in M/\mathfrak m M$ が $R$ の剰余体上の基底をなすような元とする。
補題 \ref{algebra-lemma-local-artinian-basis-when-flat} により
$x_\alpha$ は $R$ 上の $M$ の基底であり、結論を得る。
\end{proof}

\begin{lemma}
\label{algebra-lemma-lift-basis}
$R$ を環とする。$I \subset R$ をイデアルとする。$M$ を $R$-加群とする。
$A$ を集合とし、$x_\alpha \in M$、$\alpha \in A$ を $M$ の元の族とする。
次を仮定する。
\begin{enumerate}
\item $I$ は冪零である。
\item $\{\overline{x}_\alpha\}_{\alpha \in A}$ は
$R/I$ 上の $M/IM$ の基底をなす。
\item $\text{Tor}_1^R(R/I, M) = 0$ である。
\end{enumerate}
このとき $M$ は $R$ 上 $\{x_\alpha\}_{\alpha \in A}$ を基底とする自由加群である。
\end{lemma}

\begin{proof}
$R$、$I$、$M$、$\{x_\alpha\}_{\alpha \in A}$ は補題のとおりで、
仮定 (1)、(2)、(3) を満たすとする。
Nakayama の補題 \ref{algebra-lemma-NAK} により、元 $x_\alpha$ は $R$ 上 $M$ を生成する。
仮定 $\text{Tor}_1^R(R/I, M) = 0$ は短完全列
$$
0 \to I \otimes_R M \to M \to M/IM \to 0.
$$
があることを意味する。$\sum f_\alpha x_\alpha = 0$ を $M$ における関係式とする。
$x_\alpha$ の選び方により $f_\alpha \in I$ である。
従って $I \otimes_R M$ において
$\sum f_\alpha \otimes x_\alpha = 0$ と結論する。
写像 $I \otimes_R M \to I/I^2 \otimes_{R/I} M/IM$ と、
$\{x_\alpha\}_{\alpha \in A}$ が $M/IM$ の基底をなすことから
$f_\alpha \in I^2$ が従う！ 従って $M/I^2M$ における $x_\alpha$ の像の間には
関係がないと結論する。言い換えれば、$M/I^2M$ は $x_\alpha$ の像を基底とする自由加群である。
次に写像 $I \otimes_R M \to I/I^3 \otimes_{R/I^2} M/I^2M$ を用いると
$f_\alpha \in I^3$ と結論する！ 以下同様である。
仮定 (1) によりある $n$ に対して $I^n = 0$ なので、結論を得る。
\end{proof}

\begin{lemma}
\label{algebra-lemma-prepare-lift-flatness}
$\varphi : R \to R'$ を環写像とする。$I \subset R$ をイデアルとする。
$M$ を $R$-加群とする。次を仮定する。
\begin{enumerate}
\item $M/IM$ は $R/I$ 上平坦である。
\item $R' \otimes_R M$ は $R'$ 上平坦である。
\end{enumerate}
$I_2 = \varphi^{-1}(\varphi(I^2)R')$ とおく。
このとき $M/I_2M$ は $R/I_2$ 上平坦である。
\end{lemma}

\begin{proof}
$R$、$M$、$R'$ をそれぞれ $R/I_2$、$M/I_2M$、
$R'/\varphi(I)^2R'$ で置き換えてよい。このとき $I^2 = 0$ であり、
$\varphi$ は単射である。補題 \ref{algebra-lemma-what-does-it-mean} と
$I^2 = 0$ であることにより、
$\text{Tor}^R_1(R/I, M) = K = \Ker(I \otimes_R M \to M)$ が零であることを
証明すれば十分である。$M' = M \otimes_R R'$、$I' = IR'$ とおく。
仮定により写像 $I' \otimes_{R'} M' \to M'$ は単射である。
従って $K$ は
$$
I' \otimes_{R'} M' = I' \otimes_R M = I' \otimes_{R/I} M/IM.
$$
において零に写る。また $I \to I'$ は $R/I$-加群の単射である。
$M/IM$ は $R/I$ 上平坦なので、写像
$$
I \otimes_{R/I} M/IM \longrightarrow I' \otimes_{R/I} M/IM
$$
は単射である。これは望んだとおり、
$I \otimes_R M = I \otimes_{R/I} M/IM$ において $K$ が零であることを意味する。
\end{proof}

\begin{lemma}
\label{algebra-lemma-lift-flatness}
$\varphi : R \to R'$ を環写像とする。$I \subset R$ をイデアルとする。
$M$ を $R$-加群とする。次を仮定する。
\begin{enumerate}
\item $I$ は冪零である。
\item $R \to R'$ は単射である。
\item $M/IM$ は $R/I$ 上平坦である。
\item $R' \otimes_R M$ は $R'$ 上平坦である。
\end{enumerate}
このとき $M$ は $R$ 上平坦である。
\end{lemma}

\begin{proof}
帰納的に $I_1 = I$ および
$n \geq 1$ に対して $I_{n + 1} = \varphi^{-1}(\varphi(I_n)^2R')$ と定義する。
補題 \ref{algebra-lemma-prepare-lift-flatness} により、各 $n \geq 1$ に対して
$M/I_nM$ は $R/I_n$ 上平坦である。
$\varphi(I_{n + 1}) \subset \varphi(I)^{2^n}R'$ であることは明らかである。
$I$ は冪零なので、ある $n$ に対して $\varphi(I_n) = 0$ である。
$\varphi$ は単射なので、ある $n$ に対して $I_n = 0$ と結論し、証明を終える。
\end{proof}

\noindent
平坦性の局所判定法の局所アルティン版を与える。

\begin{lemma}
\label{algebra-lemma-artinian-variant-local-criterion-flatness}
$R$ をアルティン局所環とする。$M$ を $R$-加群とする。
$I \subset R$ を真のイデアルとする。次は同値である。
\begin{enumerate}
\item $M$ は $R$ 上平坦である。
\item $M/IM$ は $R/I$ 上平坦であり、$\text{Tor}_1^R(R/I, M) = 0$ である。
\end{enumerate}
\end{lemma}

\begin{proof}
(1) $\Rightarrow$ (2) は定義から直ちに従う。
$M/IM$ が $R/I$ 上平坦であり、$\text{Tor}_1^R(R/I, M) = 0$ であると仮定する。
補題 \ref{algebra-lemma-local-artinian-characterize-flat} により、
これは $M/IM$ が $R/I$ 上自由であることを意味する。
集合 $A$ と、$M/IM$ における像が基底をなす元 $x_\alpha \in M$ を選ぶ。
補題 \ref{algebra-lemma-lift-basis} により $M$ は自由であり、特に平坦である。
\end{proof}

\noindent
始域がアルティン環である単射準同型に沿って平坦性が降下することが分かる。

\begin{lemma}
\label{algebra-lemma-descent-flatness-injective-map-artinian-rings}
$R \to S$ を環写像とする。$M$ を $R$-加群とする。次を仮定する。
\begin{enumerate}
\item $R$ はアルティン的である。
\item $R \to S$ は単射である。
\item $M \otimes_R S$ は平坦 $S$-加群である。
\end{enumerate}
このとき $M$ は平坦 $R$-加群である。
\end{lemma}

\begin{proof}
第一の証明：$I \subset R$ を $R$ の Jacobson 根基とする。
$I$ は冪零であり、$R/I$ は体の積なので $M/IM$ は $R/I$ 上平坦である。
節 \ref{algebra-section-artinian} を参照せよ。
従って補題 \ref{algebra-lemma-lift-flatness} を適用すると $M$ は平坦である。

\medskip\noindent
第二の証明：補題 \ref{algebra-lemma-artinian-finite-length} により、
$R = \prod R_i$ を局所アルティン環の有限積として書ける。
これは $R$ と $S$ の双方に同様の積分解を誘導する。
% [P01-E0496] The frozen source names R and S here; the likely intended pair S and M remains sidecar-only.
従って $R$ が局所アルティンである場合に帰着する（詳細は省略する）。

\medskip\noindent
$R \to S$ と $M$ が補題のとおり (1)、(2)、(3) を満たし、
さらに $R$ は極大イデアル $\mathfrak m$ をもつ局所環であると仮定する。
$A$ を集合とし、$x_\alpha \in A$ を、
$\overline{x}_\alpha$ が $R/\mathfrak m$ 上の $M/\mathfrak mM$ の基底をなすような元とする。
% [P01-E0494] The frozen x_alpha-in-A typing error is retained; the intended membership in M remains sidecar-only.
Nakayama の補題 \ref{algebra-lemma-NAK} により、元 $x_\alpha$ は $R$-加群として $M$ を生成する。
$N = S \otimes_R M$、$I = \mathfrak mS$ とおく。
このとき $\{1 \otimes x_\alpha\}_{\alpha \in A}$ は $N$ の元の族であり、
$N/IN$ の基底をなす。さらに $N$ は $S$ 上平坦なので
$\text{Tor}_1^S(S/I, N) = 0$ である。
従って補題 \ref{algebra-lemma-lift-basis} により、
$N$ は $\{1 \otimes x_\alpha\}_{\alpha \in A}$ 上自由である。
$R \to S$ の単射性により、係数を $R$ にもつ $x_\alpha$ 間の
非自明な関係は存在しえない。
\end{proof}

\noindent
次の補題を、補題 \ref{algebra-lemma-criterion-flatness-fibre-Noetherian}
（ネーター局所環の場合）、補題 \ref{algebra-lemma-criterion-flatness-fibre}
（有限表示代数の場合）、および
補題 \ref{algebra-lemma-criterion-flatness-fibre-locally-nilpotent}
（局所冪零イデアルの場合）と比較されたい。

\begin{lemma}[ファイバーによる平坦性判定法：冪零の場合]
\label{algebra-lemma-criterion-flatness-fibre-nilpotent}
環の圏における可換図式
$$
\xymatrix{
S \ar[rr] & & S' \\
& R \ar[lu] \ar[ru]
}
$$
が与えられているとする。$I \subset R$ を冪零イデアルとし、
$M$ を $S'$-加群とする。次を仮定する。
\begin{enumerate}
\item 加群 $M/IM$ は $S/IS$ 上平坦である。
\item 加群 $M$ は $R$ 上平坦である。
\end{enumerate}
このとき $M$ は $S$ 上平坦であり、
$M \otimes_S \kappa(\mathfrak q)$ が非零となるすべての
$\mathfrak q \subset S$ に対して $S_{\mathfrak q}$ は $R$ 上平坦である。
\end{lemma}

\begin{proof}
$M$ は $R$ 上平坦なので、短完全列 $0 \to I \to R \to R/I \to 0$ と
テンソルすると短完全列
$$
0 \to I \otimes_R M \to M \to M/IM \to 0.
$$
を得る。$I \otimes_R M \to IS \otimes_S M$ が全射であることに注意する。
これを上のことと合わせると
$$
I \otimes_R M \to IS \otimes_S M \to M
$$
の両方の写像が単射であることを意味する。従って
$\text{Tor}_1^S(IS, M) = 0$ であり（注意
\ref{algebra-remark-Tor-ring-mod-ideal} を参照せよ）、
補題 \ref{algebra-lemma-what-does-it-mean} により $M$ は平坦 $S$-加群である。
% [P01-E0495] The frozen Tor argument remains typed with IS rather than S/IS; its repair is sidecar-only.
最後に、$M \otimes_S \kappa(\mathfrak q)$ が非零となる任意の素イデアル
$\mathfrak q \subset S$ に対して、$S_{\mathfrak q}$ が $R$ 上平坦であることを示す必要がある。
これは補題 \ref{algebra-lemma-ff} と \ref{algebra-lemma-flat-permanence} から従う。
\end{proof}




\section{複体を完全にするものは何か？}
\label{algebra-section-complex-exact}

\noindent
この内容の一部は Buchsbaum と Eisenbud の論文 \cite{WhatExact} に見られる。

\begin{situation}
\label{algebra-situation-complex}
ここで $R$ は環であり、複体
$$
0
\to
R^{n_e}
\xrightarrow{\varphi_e}
R^{n_{e-1}}
\xrightarrow{\varphi_{e-1}}
\ldots
\xrightarrow{\varphi_{i + 1}}
R^{n_i}
\xrightarrow{\varphi_i}
R^{n_{i-1}}
\xrightarrow{\varphi_{i-1}}
\ldots
\xrightarrow{\varphi_1}
R^{n_0}
$$
が与えられている。言い換えれば、
$i = 1, \ldots, e - 1$ に対して
$\varphi_i \circ \varphi_{i + 1} = 0$ を要求する。
\end{situation}

\begin{lemma}
\label{algebra-lemma-add-trivial-complex}
$R$ を環とする。有限自由 $R$-加群の複体
$$
\ldots
\xrightarrow{\varphi_{i + 1}}
R^{n_i}
\xrightarrow{\varphi_i}
R^{n_{i-1}}
\xrightarrow{\varphi_{i-1}}
\ldots
$$
が与えられているとする。ある $i$ に対し、写像 $\varphi_i$ のある行列係数が
可逆であると仮定する。このとき、表示された複体は複体
$$
\ldots \to
R^{n_{i + 2}} \xrightarrow{\varphi_{i + 2}}
R^{n_{i + 1}} \to
R^{n_i - 1} \to
R^{n_{i - 1} - 1} \to
R^{n_{i - 2}} \xrightarrow{\varphi_{i - 2}}
R^{n_{i - 3}} \to
\ldots
$$
と複体 $\ldots \to 0 \to R \to R \to 0 \to \ldots$ の直和に同型である。
ここで写像 $R \to R$ は恒等写像である。
\end{lemma}

\begin{proof}
仮定は、$R^{n_i}$ と $R^{n_{i-1}}$ の基底を変えた後、
$R^{n_i}$ の第一基底ベクトルが $\varphi_i$ によって
$R^{n_{i-1}}$ の第一基底ベクトルに写ることを意味する。
$e_j$ を $R^{n_i}$ の第 $j$ 基底ベクトル、
$f_k$ を $R^{n_{i-1}}$ の第 $k$ 基底ベクトルと書く。
$\varphi_i(e_j) = \sum a_{jk} f_k$ と書く。
このとき $k = 1$ でない限り $a_{1k} = 0$ であり、$a_{11} = 1$ である。
$j > 1$ に対して $e'_j = e_j - a_{j1} e_1$ とおき、
$R^{n_i}$ の基底をもう一度変える。この座標変換の後、
$j > 1$ に対して $a_{j1} = 0$ となる。
$R^{n_{i + 1}} \to R^{n_i}$ の像は、
$j > 1$ に対する $e_j$ が生成する部分加群に含まれることに注意する。
% [P01-E0497] The frozen “subspace” over an arbitrary ring is rendered as the required submodule.
また $R^{n_{i-1}} \to R^{n_{i-2}}$ は、像に含まれる $f_1$ を
零化しなければならない。
これらの条件と $\varphi_i$ の行列 $(a_{jk})$ の形から補題が従う。
\end{proof}

\noindent
状況 \ref{algebra-situation-complex} において、形
$$
0 \to \ldots \to 0 \to R \xrightarrow{1} R \to 0 \to \ldots \to 0
$$
または形
$$
0 \to \ldots \to 0 \to R
$$
の複体を {\it 自明} であるという。
より正確には、$0 \to R^{n_e} \to R^{n_{e-1}} \to \ldots \to R^{n_0}$
について、$e \geq i \geq 1$ で
$n_i = n_{i - 1} = 1$、$\varphi_i = \text{id}_R$、
$j \not \in \{i, i - 1\}$ に対して $n_j = 0$ となるものが存在するか、
または $n_0 = 1$ かつ $i > 0$ に対して $n_i = 0$ ならば、自明であるという。
上の補題は明らかに、局所環上の有限自由加群の任意の有限複体は、
自明な複体との直和を除けば、すべての写像のすべての行列係数が
極大イデアルに入る複体と同じであることを述べている。

\begin{lemma}
\label{algebra-lemma-exact-depth-zero-local}
状況 \ref{algebra-situation-complex} において、$R$ を極大イデアル
$\mathfrak m$ をもつネーター局所環とする。
$\mathfrak m \in \text{Ass}(R)$、言い換えれば $R$ の深さが $0$ であると仮定する。
$0 \to R^{n_e} \to R^{n_{e-1}} \to \ldots \to R^{n_0}$ が
$R^{n_e}, \ldots, R^{n_1}$ において完全であると仮定する。
このとき複体は自明な複体の直和に同型である。
\end{lemma}

\begin{proof}
$\mathfrak m x = 0$ を満たす $x \in R$、$x \not = 0$ を選ぶ。
$n_i > 0$ となる最大の添字を $i$ とする。$i = 0$ ならば主張は成り立つ。
$i > 0$ ならば、$f_1$ を $R^{n_i}$ の第一基底ベクトルと書く。
% [P01-E0499] The frozen notation declaration omits “by”; Japanese uses the ordinary declaration.
複体の完全性により $xf_1$ は零に写らないので、
写像 $R^{n_i} \to R^{n_{i - 1}}$ のある行列係数は
$\mathfrak m$ に入らないと分かる。
すると補題 \ref{algebra-lemma-add-trivial-complex} により
$n_e + \ldots + n_1$ を減らせる。帰納法で証明が終わる。
\end{proof}

\begin{lemma}
\label{algebra-lemma-exact-artinian-local}
状況 \ref{algebra-situation-complex} において、$R$ をアルティン局所環とする。
% [P01-E0498] The frozen English article error has no Japanese analogue.
$0 \to R^{n_e} \to R^{n_{e-1}} \to \ldots \to R^{n_0}$ が
$R^{n_e}, \ldots, R^{n_1}$ において完全であると仮定する。
このとき複体は自明な複体の直和に同型である。
\end{lemma}

\begin{proof}
アルティン局所環の深さは $0$ なので、これは
補題 \ref{algebra-lemma-exact-depth-zero-local} の特別な場合である。
\end{proof}

\noindent
以下で有限自由加群の写像の階数を定義する。
これは階数の一つの可能な定義にすぎない。この節で機能する定義であり、
他の状況ではより便利な別の定義もある。

\begin{definition}
\label{algebra-definition-rank}
$R$ を環とする。$\varphi : R^m \to R^n$ を有限自由加群の写像とする。
\begin{enumerate}
\item $\varphi$ の {\it 階数} とは、
$\wedge^r \varphi : \wedge^r R^m \to \wedge^r R^n$ が非零となる最大の $r$ である。
\item $I(\varphi) \subset R$ を $\varphi$ の行列の $r \times r$ 小行列式が生成する
イデアルとする。ここで $r$ は上で定義した階数である。
\end{enumerate}
\end{definition}

\noindent
$\varphi : R^m \to R^n$ の階数が $0$ であることと
$\varphi = 0$ であることは同値であり、この場合 $I(\varphi) = R$ である。

\begin{lemma}
\label{algebra-lemma-trivial-case-exact}
状況 \ref{algebra-situation-complex} において、複体が自明な複体の直和に同型であると仮定する。
このとき次が成り立つ。
\begin{enumerate}
\item 写像 $\varphi_i$ の階数は
$r_i = n_i - n_{i + 1} + \ldots + (-1)^{e-i-1} n_{e-1} + (-1)^{e-i} n_e$ である。
\item すべての $i$、$1 \leq i \leq e - 1$ に対して
$\text{rank}(\varphi_{i + 1}) + \text{rank}(\varphi_i) = n_i$ である。
\item 各 $I(\varphi_i) = R$ である。
\end{enumerate}
\end{lemma}

\begin{proof}
複体は自明な複体の直和であると仮定してよい。
このとき各 $i$ に対して、$R^{n_i}$ の標準基底ベクトルを、
$R^{n_{i-1}}$ の基底ベクトルに写るものと零に写るもの
（$i > 0$ のとき後者には $R^{n_{i + 1}}$ の基底ベクトルが全射的に写る）
に分けられる。$i = e$ から始める降下帰納法により、
$R^{n_i}$ には零に写る $r_{i + 1}$ 個の基底ベクトルと、
$R^{n_{i-1}}$ の基底ベクトルに写る $r_i$ 個の基底ベクトルがあることを容易に示せる。
% [P01-E0500] The frozen hyphenated count is rendered by its evident cardinal meaning.
これから結果が従う。
\end{proof}

\begin{lemma}
\label{algebra-lemma-div-x-exact-one-less}
状況 \ref{algebra-situation-complex} において、$R$ を極大イデアル
$\mathfrak m$ をもつ局所環とする。
$0 \to R^{n_e} \to R^{n_{e-1}} \to \ldots \to R^{n_0}$ が
$R^{n_e}, \ldots, R^{n_1}$ において完全であると仮定する。
$x \in \mathfrak m$ を非零因子とする。このとき複体
$0 \to (R/xR)^{n_e} \to \ldots \to (R/xR)^{n_1}$ は
$(R/xR)^{n_e}, \ldots, (R/xR)^{n_2}$ において完全である。
\end{lemma}

\begin{proof}
$F_i = R^{n_i}$ を項とし、微分を $\varphi_i$ で与える複体を
$F_\bullet$ と書く。
% [P01-E0499] The frozen “Denote F_\bullet the complex” omission is rendered by the ordinary Japanese declaration.
このとき複体の短完全列
$$
0 \to F_\bullet \xrightarrow{x} F_\bullet \to F_\bullet/xF_\bullet \to 0
$$
がある。蛇の補題を適用すると長完全列
$$
H_i(F_\bullet) \xrightarrow{x} H_i(F_\bullet) \to
H_i(F_\bullet/xF_\bullet) \to H_{i - 1}(F_\bullet)
\xrightarrow{x} H_{i - 1}(F_\bullet)
$$
を得る。補題はこれから従う。
\end{proof}

\begin{lemma}[非輪状性補題]
\label{algebra-lemma-acyclic}
\begin{reference}
\cite[Lemma 1.8]{Peskine-Szpiro}
\end{reference}
$R$ をネーター局所環とする。
$0 \to M_e \to M_{e-1} \to \ldots \to M_0$ を
有限 $R$-加群の複体とする。
$\text{depth}(M_i) \geq i$ と仮定する。
複体が $M_i$ において完全でない最大の添字を $i$ とする。
$i > 0$ ならば
$\Ker(M_i \to M_{i-1})/\Im(M_{i + 1} \to M_i)$
の深さは $\geq 1$ である。
\end{lemma}

\begin{proof}
$H = \Ker(M_i \to M_{i-1})/\Im(M_{i + 1} \to M_i)$ を問題の
同調群とする。
% [P01-E0501] The frozen source calls H a cohomology group; the displayed chain-complex quotient is rendered as homology.
複体を正合列
$0 \to M_e \to M_{e-1} \to K_{e-2} \to 0$、
$i + 2 \leq j \leq e-2 $ に対する
$0 \to K_j \to M_j \to K_{j-1} \to 0$、
$0 \to K_{i + 1} \to M_{i + 1} \to B_i \to 0$、
$0 \to K_i \to M_i \to M_{i-1}$、および
$0 \to B_i \to K_i \to H \to 0$ に分解できる。
% [P01-E0502] The frozen source calls every displayed row short exact, although one row has no terminal zero; Japanese states only exactness.
これらの列を上へたどり、補題 \ref{algebra-lemma-depth-in-ses} を繰り返し用いて、
深さについての主張を証明する。
まず $\text{depth}(M_e) \geq e$ かつ
$\text{depth}(M_{e-1}) \geq e-1$ なので、
$\text{depth}(K_{e-2}) \geq e - 1$ と分かる。
この時点で、$i + 2 \leq j \leq e-2 $ に対する列
$0 \to K_j \to M_j \to K_{j-1} \to 0$ から、同様に
$i + 2 \leq j \leq e-2$ に対して
$\text{depth}(K_{j-1}) \geq j$ が従う。次に列
$0 \to K_{i + 1} \to M_{i + 1} \to B_i \to 0$ から
$\text{depth}(B_i) \geq i + 1$ と分かる。列
$0 \to K_i \to M_i \to M_{i-1}$ からは、
仮定により $M_i$ の深さが $\geq i \geq 1$ なので
$\text{depth}(K_i) \geq 1$ と分かる。そこで列
$0 \to B_i \to K_i \to H \to 0$ から結果が従う。
\end{proof}

\begin{proposition}
\label{algebra-proposition-what-exact}
\begin{reference}
\cite[Corollary 1]{WhatExact}
\end{reference}
状況 \ref{algebra-situation-complex} において、$R$ をネーター局所環とする。
次は同値である。
\begin{enumerate}
\item $0 \to R^{n_e} \to R^{n_{e-1}} \to \ldots \to R^{n_0}$ は
$R^{n_e}, \ldots, R^{n_1}$ において完全である。
\item すべての $i$、$1 \leq i \leq e$ に対して、
次の二条件が満たされる。
\begin{enumerate}
\item $\text{rank}(\varphi_i) = r_i$ である。ここで
$r_i = n_i - n_{i + 1} + \ldots + (-1)^{e-i-1} n_{e-1} + (-1)^{e-i} n_e$ である。
\item $I(\varphi_i) = R$ であるか、または $I(\varphi_i)$ は
長さ $i$ の正則列を含む。
\end{enumerate}
\end{enumerate}
\end{proposition}

\begin{proof}
ある $i$ に対して $\varphi_i$ のある行列係数が
$\mathfrak m$ に入らないならば、補題 \ref{algebra-lemma-add-trivial-complex} を適用する。
ある複体に対する命題と、その複体に自明な複体を加えたものに対する命題が
同値であることは容易に分かる。従って、各 $\varphi_i$ のすべての行列成分が
極大イデアルの元であると仮定してよい。また $e \geq 1$ と仮定してよい。

\medskip\noindent
複体が $R^{n_e}, \ldots, R^{n_1}$ において完全であると仮定する。
$\mathfrak q \in \text{Ass}(R)$ とする。
環 $R_{\mathfrak q}$ の深さは $0$ であり、
複体は $\mathfrak q$ における局所化後も完全であることに注意する。
$R_{\mathfrak q}$ 上の局所化された複体に、補題
\ref{algebra-lemma-exact-depth-zero-local} と \ref{algebra-lemma-trivial-case-exact} を適用する。
すべての $i$ に対して $\varphi_{i, \mathfrak q}$ の階数は $r_i$ であると結論する。
$R \to \bigoplus_{\mathfrak q \in \text{Ass}(R)} R_\mathfrak q$ は単射なので
（補題 \ref{algebra-lemma-zero-at-ass-zero}）、定義 \ref{algebra-definition-rank} で与えた
階数の定義により、$\varphi_i$ の $R$ 上の階数は $r_i$ であると結論する。
従って階数は変わらないから
$I(\varphi_i)_\mathfrak q = I(\varphi_{i, \mathfrak q})$ である。
上で参照した補題により、$e \geq i \geq 1$ に対するすべてのイデアル
$I(\varphi_i)_{\mathfrak q}$ は $R_{\mathfrak q}$ に等しいので、
どのイデアル $I(\varphi_i)$ も $\mathfrak q$ に含まれないと結論する。
これは $I(\varphi_e)I(\varphi_{e-1})\ldots I(\varphi_1)$ が
$R$ のどの随伴素イデアルにも含まれないことを意味する。
補題 \ref{algebra-lemma-silly} により
$x \in I(\varphi_e)I(\varphi_{e - 1})\ldots I(\varphi_1)$ で、
すべての $\mathfrak q \in \text{Ass}(R)$ に対して
$x \not \in \mathfrak q$ となる元を選べる。
$x$ は非零因子であることに注意する（補題 \ref{algebra-lemma-ass-zero-divisors}）。
補題 \ref{algebra-lemma-div-x-exact-one-less} によれば、複体
$0 \to (R/xR)^{n_e} \to \ldots \to (R/xR)^{n_1}$ は
$(R/xR)^{n_e}, \ldots, (R/xR)^{n_2}$ において完全である。
$e$ に関する帰納法により、すべてのイデアル $I(\varphi_i)/xR$ は
長さ $i - 1$ の正則列をもつ。これにより $I(\varphi_i)$ は
長さ $i$ の正則列を含むことが証明された。

\medskip\noindent
(2)(a) と (2)(b) が成り立つと仮定する。
$\dim(R)$ に関する帰納法により (1) を証明する。
$\dim(R) = 0$ ならば、(2)(b) により
$1 \leq i \leq e$ に対して $I(\varphi_i) = R$ でなければならない。
$\varphi_i$ の係数は極大イデアルに含まれるので、
これはすべての $i$ に対して $r_i = 0$ である場合に限って起こりうる。
(2)(a) により $e = 0$ と結論し、(1) が成り立つ。
$\dim(R) > 0$ と仮定する。
任意の素イデアル $\mathfrak p \subset R$ に対して、写像
$\varphi_{i, \mathfrak p}$ をもつ $R_\mathfrak p$ 上の複体
$0 \to R_\mathfrak p^{n_e} \to R_\mathfrak p^{n_{e - 1}} \to \ldots \to
R_\mathfrak p^{n_0}$ について条件 (2)(a) と (2)(b) が成り立つと主張する。
実際、$I(\varphi_i)$ は非零因子を含むので、
$R_\mathfrak p$ における $I(\varphi_i)$ の像は非零である。
これは $\varphi_{i, \mathfrak p}$ の階数が $\varphi_i$ の階数と同じであることを意味する。
すなわち、環 $R$ 上の行列 $\varphi$ について上で定義した階数は、
$R$-代数 $R'$ へ移るとき減少することしかなく、それが起こるのは
$I(\varphi)$ が $R'$ において零に写る場合に限る。従って (2)(a) が成り立つ。
以上から $I(\varphi_{i, \mathfrak p}) = I(\varphi_i)_\mathfrak p$ と分かり、
(2)(b) も局所化で保たれることが分かる。
$R$ の次元に関する帰納法により、$R$ の任意の非極大素イデアル
$\mathfrak p$ において局所化すると複体は完全であると仮定してよい。
従って $\Ker(\varphi_i)/\Im(\varphi_{i + 1})$ の台は
$\{\mathfrak m\}$ に含まれ、ゆえに非零ならば深さは $0$ である。
証明の第一段落で述べたことにより、すべての $i$ に対して
$I(\varphi_i) \subset \mathfrak m$ であるから、(2)(b) は
$\text{depth}(R) \geq e$ を意味する。
% [P01-E0503] The frozen universal containment I(phi_i) subset m is preserved; any qualification remains sidecar-only.
補題 \ref{algebra-lemma-acyclic} により、複体は
$R^{n_e}, \ldots, R^{n_1}$ において完全であると分かり、証明が終わる。
\end{proof}

\begin{remark}
\label{algebra-remark-what-exact}
命題 \ref{algebra-proposition-what-exact} の同値な条件 (1) と (2) が満たされるならば、
$I(\varphi_i) = R$ であることと $i \geq j$ であることが同値となるような
$j$ が存在する。命題の証明と同様に、すべての行列の係数が
$R$ の極大イデアル $\mathfrak m$ に入る場合を見れば十分である。
この場合、$I(\varphi_j) = R$ であることと $\varphi_j = 0$ であることは同値である。
しかし $\varphi_j = 0$ ならば、任意に長い完全複体
$0 \to R^{n_e} \to R^{n_{e-1}} \to \ldots
\to R^{n_j} \to 0 \to 0 \to \ldots \to 0$
が得られる。従って命題により、$i > j$ に対する $I(\varphi_i)$ は
$R$ でなければならない（さもなくば、ネーター局所環の真のイデアルが
任意に長い正則列を含むことになり、これは不可能である）。
\end{remark}
\section{Cohen--Macaulay 加群}
\label{algebra-section-CM}

\noindent
ここでは Cohen--Macaulay 加群が良い性質をもつことを示す。
Ext 群を用いて双対性などとの関係を確立することは後に回す。

\begin{definition}
\label{algebra-definition-CM}
$R$ をネーター局所環とする。$M$ を有限 $R$-加群とする。
$\dim(\text{Supp}(M)) = \text{depth}(M)$ ならば、
$M$ は {\it Cohen--Macaulay} であるという。
\end{definition}

\noindent
最初の目標は命題 \ref{algebra-proposition-CM-module} を確立することである。
これは、深さに関する先の結果をほとんど用いない
（おそらく標準的ではない）一連の初等的補題によって行う。
いくつか記法を導入しよう。

\medskip\noindent
$R$ をネーター局所環とする。$M$ を Cohen--Macaulay 加群とし、
$f_1, \ldots, f_d$ を $d = \dim(\text{Supp}(M))$ を満たす
$M$-正則列とする。すべての $i = 0, 1, \ldots, d-1$ に対して
$\dim (\text{Supp}(M) \cap V(g, f_1, \ldots, f_i))
= d - i - 1$
が成り立つとき、$g \in \mathfrak m$ は
{\it $(M, f_1, \ldots, f_d)$ に関して良い} という。
これは、$i = 0, 1, \ldots, d - 1$ に対して
$\dim(\text{Supp}(M/(f_1, \ldots, f_i)M) \cap V(g)) =
d - i - 1$ であるという条件と同値である。

\begin{lemma}
\label{algebra-lemma-good-element}
記法と仮定は上のとおりとする。$g$ が
$(M, f_1, \ldots, f_d)$ に関して良いならば、
(a) $g$ は $M$ 上の非零因子であり、(b) $M/gM$ は
$f_1, \ldots, f_{d - 1}$ を極大正則列としてもつ
Cohen--Macaulay 加群である。
\end{lemma}

\begin{proof}
$d$ に関する帰納法で補題を証明する。
$d = 0$ ならば $M$ は有限であり、補題を適用する場合は存在しない。
% [P01-E0504] The frozen vacuous d=0 case is preserved; its false nonapplicability claim remains sidecar-only.
$d = 1$ ならば、$g : M \to M$ が単射であることを示せばよい。
仮定により $\dim \text{Supp}(M) \cap V(g) = 0$ なので、
核 $K$ の台は $\{\mathfrak m\}$ である。従って $K$ の長さは有限である。
従って $f_1 : K \to K$ が単射であることから、その像の長さは
$K$ の長さに等しく、ゆえに $f_1 K = K$ である。
Nakayama の補題 \ref{algebra-lemma-NAK} により、これは $K = 0$ を意味する。
また $\dim \text{Supp}(M/gM) = 0$ なので、
$M/gM$ は深さ $0$ の Cohen--Macaulay 加群である。

\medskip\noindent
$d > 1$ と仮定する。定義から容易に分かるように、
$g$ は $(M/f_1M, f_2, \ldots, f_d)$ に関して良い。
帰納法により、(a) $g$ は $M/f_1M$ 上の非零因子であり、
(b) $M/(g, f_1)M$ は $f_2, \ldots, f_{d - 1}$ を
極大正則列としてもつ Cohen--Macaulay 加群である。
補題 \ref{algebra-lemma-permute-xi} により、$g, f_1$ は $M$-正則列である。
従って $g$ は $M$ 上の非零因子であり、
$f_1, \ldots, f_{d - 1}$ は $M/gM$-正則列である。
\end{proof}

\begin{lemma}
\label{algebra-lemma-CM-one-g}
$R$ をネーター局所環とする。$M$ を $R$ 上の
Cohen--Macaulay 加群とする。
$g \in \mathfrak m$ が
$\dim(\text{Supp}(M) \cap V(g))
= \dim(\text{Supp}(M)) - 1$ を満たすと仮定する。
このとき (a) $g$ は $M$ 上の非零因子であり、
(b) $M/gM$ は深さが一つ小さい Cohen--Macaulay 加群である。
\end{lemma}

\begin{proof}
$d = \dim(\text{Supp}(M))$ を満たす $M$-正則列
$f_1, \ldots, f_d$ を選ぶ。
% [P01-E0506] The frozen wrong English article before M-regular has no Japanese analogue.
$g$ が $(M, f_1, \ldots, f_d)$ に関して良ければ、
補題 \ref{algebra-lemma-good-element} により結論を得る。
特に $d = 1$ ならば補題は成り立つ（$d = 0$ の場合は起こらない）。
$d > 1$ と仮定する。(\romannumeral1) $h$ は
$(M, f_1, \ldots, f_d)$ に関して良く、かつ
(\romannumeral2) $\dim(\text{Supp}(M) \cap V(h, g)) = d - 2$
となる元 $h \in R$ を選ぶ。
$h$ の存在を見るため、閉集合
$\text{Supp}(M)$、
$i = 1, \ldots, d - 1$ に対する
$\text{Supp}(M)\cap V(f_1, \ldots, f_i)$、および
$\text{Supp}(M) \cap V(g)$ の極小素イデアルからなる
（有限）集合を $\{\mathfrak q_j\}$ とする。
これらの $\mathfrak q_j$ はいずれも $\mathfrak m$ に等しくないので、
補題 \ref{algebra-lemma-silly} により
$h \in \mathfrak m$、$h \not \in \mathfrak q_j$ を満たす元を見つけられる。
$h$ が (\romannumeral1) と (\romannumeral2) を満たすことは明らかである。
補題 \ref{algebra-lemma-good-element} から $M/hM$ は Cohen--Macaulay である。
(\romannumeral2) により、組 $(M/hM, g)$ は帰納法の仮定を満たす。
従って $M/(h, g)M$ は Cohen--Macaulay であり、
$g : M/hM \to M/hM$ は単射である。
補題 \ref{algebra-lemma-permute-xi} により、
$g : M \to M$ と $h : M/gM \to M/gM$ は単射である。
$M/(g, h)M$ が Cohen--Macaulay であることと合わせて、証明が終わる。
\end{proof}

\begin{proposition}
\label{algebra-proposition-CM-module}
$R$ を極大イデアル $\mathfrak m$ をもつネーター局所環とする。
$M$ を、台の次元が $d$ である $R$ 上の Cohen--Macaulay 加群とする。
$g_1, \ldots, g_c$ が $\mathfrak m$ の元で、
$\dim(\text{Supp}(M/(g_1, \ldots, g_c)M))
= d - c$ を満たすと仮定する。
このとき $g_1, \ldots, g_c$ は $M$-正則列であり、
極大 $M$-正則列へ拡張できる。
\end{proposition}

\begin{proof}
$Z = \text{Supp}(M) \subset \Spec(R)$ とおく。
鎖
$Z \supset Z \cap V(g_1) \supset \ldots \supset Z \cap V(g_1, \ldots, g_c)$
では、補題 \ref{algebra-lemma-one-equation} により各段階で次元は高々 $1$ だけ減少する。
従って仮定により、各段階で次元は毎回ちょうど $1$ だけ減少する。
そこで加群 $M/(g_1, \ldots, g_i)$ と元 $g_{i + 1}$ に
補題 \ref{algebra-lemma-CM-one-g} を逐次適用できる。
% [P01-E0505] The frozen quotient missing its terminal module factor M is preserved exactly.

\medskip\noindent
$c < d$ のとき $g_1, \ldots, g_c$ を一つの元で拡張するには、
補題 \ref{algebra-lemma-silly} を用いて、
$Z \cap V(g_1, \ldots, g_c)$ の有限個の極小素イデアルの
いずれにも入らない元 $g_{c + 1} \in \mathfrak m$ を選べばよい。
\end{proof}

\noindent
命題 \ref{algebra-proposition-CM-module} を証明したので、標準理論の展開を続ける。

\begin{lemma}
\label{algebra-lemma-nonzerodivisor-on-CM}
$R$ を極大イデアル $\mathfrak m$ をもつネーター局所環とする。
$M$ を有限 $R$-加群とする。$x \in \mathfrak m$ を
$M$ 上の非零因子とする。このとき $M$ が Cohen--Macaulay であることと
$M/xM$ が Cohen--Macaulay であることは同値である。
\end{lemma}

\begin{proof}
補題 \ref{algebra-lemma-depth-drops-by-one} により\hfill\break
$\text{depth}(M/xM)\allowbreak = \text{depth}(M)-1$ である。
補題 \ref{algebra-lemma-one-equation-module} により\hfill\break
$\dim(\text{Supp}(M/xM))\allowbreak = \dim(\text{Supp}(M)) - 1$ である。
\end{proof}

\begin{lemma}
\label{algebra-lemma-CM-over-quotient}
$R \to S$ をネーター局所環の全射準同型とする。
$N$ を有限 $S$-加群とする。このとき $N$ が $S$-加群として
Cohen--Macaulay であることと、$N$ が $R$-加群として
Cohen--Macaulay であることは同値である。
\end{lemma}

\begin{proof}
省略する。
\end{proof}

\begin{lemma}
\label{algebra-lemma-CM-ass-minimal-support}
\begin{reference}
\cite[Chapter 0, Proposition 16.5.4]{EGA}
\end{reference}
$R$ をネーター局所環とする。$M$ を有限 Cohen--Macaulay
$R$-加群とする。$\mathfrak p \in \text{Ass}(M)$ ならば、
$\dim(R/\mathfrak p) = \dim(\text{Supp}(M))$ であり、
$\mathfrak p$ は $M$ の台における極小素イデアルである。
特に $M$ は埋め込まれた随伴素イデアルをもたない。
\end{lemma}

\begin{proof}
補題 \ref{algebra-lemma-depth-dim-associated-primes} により
$\text{depth}(M) \leq \dim(R/\mathfrak p)$ である。
$\mathfrak p \in \text{Supp}(M)$ なので（補題 \ref{algebra-lemma-ass-support}）、
もちろん $\dim(R/\mathfrak p) \leq \dim(\text{Supp}(M))$ である。
$M$ は Cohen--Macaulay だから、両方の不等式で等号が成り立つ。
すると $\mathfrak p$ は $\text{Supp}(M)$ において極小でなければならない。
さもなくば $\dim(R/\mathfrak p) < \dim(\text{Supp}(M))$ となるからである。
最後に、$M$ の台における極小素イデアルは
$\text{Ass}(M)$ の極小元に等しいので
（命題 \ref{algebra-proposition-minimal-primes-associated-primes}）、
$M$ は埋め込まれた随伴素イデアルをもたない
（定義 \ref{algebra-definition-embedded-primes}）。
\end{proof}

\begin{definition}
\label{algebra-definition-maximal-CM}
$R$ をネーター局所環とする。有限 $R$-加群 $M$ が
$\text{depth}(M) = \dim(R)$ を満たすとき、
これを {\it 極大 Cohen--Macaulay} 加群という。
\end{definition}

\noindent
言い換えれば、ネーター局所環上の極大 Cohen--Macaulay 加群とは、
その環上で可能な最大の深さをもつ有限加群である。
同値な言い方として、ネーター局所環 $R$ 上の極大 Cohen--Macaulay 加群とは、
次元が環の次元に等しい Cohen--Macaulay 加群である。
特に $M$ が $\Spec(R) = \text{Supp}(M)$ を満たす
Cohen--Macaulay $R$-加群ならば、$M$ は極大 Cohen--Macaulay である。
従って次の二つの補題は極大 Cohen--Macaulay 加群に関するものである。

\begin{lemma}
\label{algebra-lemma-maximal-chain-maximal-CM}
\begin{slogan}
局所 Cohen--Macaulay 環では、素イデアルの任意の極大鎖の長さは
次元に等しい。
\end{slogan}
$R$ をネーター局所環とする。
$\Spec(R) = \text{Supp}(M)$ を満たす Cohen--Macaulay 加群
$M$ が存在すると仮定する。このとき素イデアルの任意の極大鎖
$\mathfrak p_0 \subset
\mathfrak p_1 \subset \ldots \subset \mathfrak p_n$
の長さは $n = \dim(R)$ である。
\end{lemma}

\begin{proof}
$\dim(R)$ に関する帰納法で証明する。$\dim(R) = 0$ ならば主張は明らかである。
$\dim(R) > 0$ と仮定する。このとき $n > 0$ である。
$x \in \mathfrak p_1$ を満たし、$x$ が $R$ のどの極小素イデアルにも入らず、
特に $x \not \in \mathfrak p_0$ である元を選ぶ
（補題 \ref{algebra-lemma-silly} を参照せよ）。
補題 \ref{algebra-lemma-one-equation} により
$\dim(R/xR) = \dim(R) - 1$ である。
命題 \ref{algebra-proposition-CM-module} と補題 \ref{algebra-lemma-CM-over-quotient} により、
加群 $M/xM$ は $R/xR$ 上 Cohen--Macaulay である。
補題 \ref{algebra-lemma-support-quotient} により、
$M/xM$ の台は $\Spec(R/xR)$ である。
ある $n$ に対して $x$ を $x^n$ で置き換えた後、
$\mathfrak p_1$ が $M/xM$ の随伴素イデアルであると仮定してよい
（補題 \ref{algebra-lemma-inherit-minimal-primes} を参照せよ）。
補題 \ref{algebra-lemma-CM-ass-minimal-support} により、
$\mathfrak p_1/(x)$ は $R/xR$ の極小素イデアルであると結論する。
従って鎖
$\mathfrak p_1/(x) \subset \ldots \subset \mathfrak p_n/(x)$
は $R/xR$ における素イデアルの極大鎖である。
帰納法により、この鎖の長さは
$\dim(R/xR) = \dim(R) - 1$ であり、望む結論を得る。
\end{proof}

\begin{lemma}
\label{algebra-lemma-dim-formula-maximal-CM}
$R$ をネーター局所環とする。
$\Spec(R) = \text{Supp}(M)$ を満たす Cohen--Macaulay 加群
$M$ が存在すると仮定する。このとき素イデアル
$\mathfrak p \subset R$ に対して
$$
\dim(R) = \dim(R_{\mathfrak p}) + \dim(R/\mathfrak p).
$$
が成り立つ。
\end{lemma}

\begin{proof}
補題 \ref{algebra-lemma-maximal-chain-maximal-CM} から直ちに従う。
\end{proof}

\begin{lemma}
\label{algebra-lemma-localize-CM-module}
$R$ をネーター局所環とする。$M$ を $R$ 上の
Cohen--Macaulay 加群とする。任意の素イデアル
$\mathfrak p \subset R$ に対して、加群 $M_{\mathfrak p}$ は
$R_\mathfrak p$ 上 Cohen--Macaulay である。
\end{lemma}

\begin{proof}
$\mathfrak p \not = \mathfrak m$ かつ $M$ は非零であると仮定してよいし、
実際そう仮定する。
% [P01-E0507] The frozen missing English finite verb in “M not zero” has no Japanese analogue.
素イデアルの極大鎖
$\mathfrak p = \mathfrak p_c \subset
\mathfrak p_{c - 1} \subset \ldots \subset \mathfrak p_1 \subset \mathfrak m$
を選ぶ。$R_{\mathfrak p_1}$ 上の $M_{\mathfrak p_1}$ に対して
結果を証明すれば、$c$ に関する帰納法で補題が従う。
従って $\mathfrak p$ と $\mathfrak m$ の間に素イデアルが
真に存在しないと仮定してよい。
$M_\mathfrak p$ の台における素イデアルの任意の鎖は、
$M$ の台でさらに一つの素イデアル（すなわち $\mathfrak m$）を
付け加えて延長できるので、
$\dim(\text{Supp}(M_\mathfrak p)) \leq \dim(\text{Supp}(M)) - 1$
である。一方、補題 \ref{algebra-lemma-depth-localization} と
$\mathfrak p$ の選び方により
$\text{depth}(M_\mathfrak p) \geq \text{depth}(M) - \dim(R/\mathfrak p) =
\text{depth}(M) - 1$ である。
従って望むとおり
$\text{depth}(M_\mathfrak p) \geq \dim(\text{Supp}(M_\mathfrak p))$
である（逆向きの不等式は補題 \ref{algebra-lemma-bound-depth} である）。
\end{proof}

\begin{definition}
\label{algebra-definition-module-CM}
$R$ をネーター環とする。$M$ を有限 $R$-加群とする。
$R$ のすべての素イデアル $\mathfrak p$ に対して
$M_\mathfrak p$ が $R_\mathfrak p$ 上の Cohen--Macaulay 加群であるとき、
$M$ は {\it Cohen--Macaulay} であるという。
\end{definition}

\noindent
補題 \ref{algebra-lemma-localize-CM-module} により、
$R$ の極大イデアルについて調べるだけで十分である。

\begin{lemma}
\label{algebra-lemma-maximal-CM-polynomial-algebra}
$R$ をネーター環とする。$M$ を $R$ 上の Cohen--Macaulay 加群とする。
このとき $M \otimes_R R[x_1, \ldots, x_n]$ は
$R[x_1, \ldots, x_n]$ 上の Cohen--Macaulay 加群である。
\end{lemma}

\begin{proof}
変数の個数に関する帰納法により、$R[x]$ 上の
$M[x] = M \otimes_R R[x]$ について証明すれば十分である。
$\mathfrak m \subset R[x]$ を極大イデアルとし、
$\mathfrak p = R \cap \mathfrak m$ とおく。
$R_{\mathfrak p}$ の極大イデアルに含まれる、長さ
$d = \dim(\text{Supp}(M_{\mathfrak p}))$ の
$M_\mathfrak p$-正則列 $f_1, \ldots, f_d$ をとる。
% [P01-E0506] The second frozen wrong English article before M-regular has no Japanese analogue.
$R[x]$ は $R$ 上平坦なので、局所化
$R[x]_{\mathfrak m}$ は $R_{\mathfrak p}$ 上平坦であることに注意する。
従って補題 \ref{algebra-lemma-flat-increases-depth} により、
列 $f_1, \ldots, f_d$ は $R[x]_{\mathfrak m}$ における長さ $d$ の
$M[x]_{\mathfrak m}$-正則列である。商
$$
Q = M[x]_{\mathfrak m}/(f_1, \ldots, f_d)M[x]_{\mathfrak m} =
M_{\mathfrak p}/(f_1, \ldots, f_d)M_{\mathfrak p}
\otimes_{R_\mathfrak p} R[x]_{\mathfrak m}
$$
の台は $\mathfrak p$ の上にある素イデアル全体に等しい。
実際、$R_\mathfrak p \to R[x]_\mathfrak m$ は平坦であり、
$M_{\mathfrak p}/(f_1, \ldots,\allowbreak f_d)M_{\mathfrak p}$ の台は
$\{\mathfrak p\}$ に等しいからである
（詳細は省略する。ヒント：補題 \ref{algebra-lemma-annihilator-flat-base-change} と
\ref{algebra-lemma-support-closed} から従う）。従ってその次元は $1$ である。
証明を終えるには、$Q$ 上の非零因子である
$f \in \mathfrak m$ を見つければ十分である。
$\mathfrak m$ は極大イデアルなので、体拡大
$\kappa(\mathfrak m)/\kappa(\mathfrak p)$ は有限である
（定理 \ref{algebra-theorem-nullstellensatz}）。
従って、$x$ の多項式と見たとき最高次係数が
$\mathfrak p$ に入らない $f \in \mathfrak m$ を見つけられる。
そのような $f$ は
$$
M_{\mathfrak p}/(f_1, \ldots, f_d)M_{\mathfrak p} \otimes_R R[x] =
\bigoplus\nolimits_{n \geq 0}
M_{\mathfrak p}/(f_1, \ldots, f_d)M_{\mathfrak p} \cdot x^n
$$
上で非零因子として作用し、従って $Q$ 上でも非零因子として作用する。
\end{proof}

\section{Cohen--Macaulay 環}
\label{algebra-section-CM-ring}

\noindent
この節の結果の大部分は、節 \ref{algebra-section-CM} の結果の特別な場合である。

\begin{definition}
\label{algebra-definition-local-ring-CM}
ネーター局所環 $R$ が自身の上の加群として Cohen--Macaulay であるとき、
これを {\it Cohen--Macaulay} という。
\end{definition}

\noindent
これは、極大イデアルの $R$-正則列 $x_1, \ldots, x_d$ で、
$R/(x_1, \ldots, x_d)$ の次元が $0$ となるものの存在を
要求することと同値である。
% [P01-E0508] The frozen wrong English article before R-regular has no Japanese analogue.
通常は単に「正則列」といい、「$R$-正則列」とはいわない。

\begin{lemma}
\label{algebra-lemma-reformulate-CM}
\begin{slogan}
Cohen--Macaulay 局所環における正則列は、正しい次元のものを
切り出すことによって特徴づけられる。
\end{slogan}
$R$ を極大イデアル $\mathfrak m $ をもつネーター局所
Cohen--Macaulay 環とする。$x_1, \ldots, x_c \in \mathfrak m$ を元とする。
このとき
$$
x_1, \ldots, x_c
\text{ is a regular sequence }
\Leftrightarrow
\dim(R/(x_1, \ldots, x_c)) = \dim(R) - c
$$
である。これらが成り立つならば、$x_1, \ldots, x_c$ は
長さ $\dim(R)$ の正則列に拡張でき、各商
$R/(x_1, \ldots, x_i)$ は次元 $\dim(R) - i$ の
Cohen--Macaulay 環である。
\end{lemma}

\begin{proof}
命題 \ref{algebra-proposition-CM-module} の特別な場合である。
\end{proof}

\begin{lemma}
\label{algebra-lemma-maximal-chain-CM}
$R$ をネーター局所環とする。
% [P01-E0509] The frozen incomplete English ring hypothesis is rendered by its evident Japanese sense.
$R$ は次元 $d$ の Cohen--Macaulay 環であると仮定する。
理想の任意の極大鎖
$\mathfrak p_0 \subset
\mathfrak p_1 \subset \ldots \subset \mathfrak p_n$
の長さは $n = d$ である。
% [P01-E0510] The frozen statement says ideals rather than prime ideals; Japanese preserves that theorem wording.
\end{lemma}

\begin{proof}
補題 \ref{algebra-lemma-maximal-chain-maximal-CM} の特別な場合である。
\end{proof}

\begin{lemma}
\label{algebra-lemma-CM-dim-formula}
$R$ を次元 $d$ のネーター局所 Cohen--Macaulay 環とする。
任意の素イデアル $\mathfrak p \subset R$ に対して
$$
\dim(R) = \dim(R_{\mathfrak p}) + \dim(R/\mathfrak p).
$$
が成り立つ。
\end{lemma}

\begin{proof}
補題 \ref{algebra-lemma-maximal-chain-CM} から直ちに従う
（これは補題 \ref{algebra-lemma-dim-formula-maximal-CM} の特別な場合でもある）。
\end{proof}

\begin{lemma}
\label{algebra-lemma-localize-CM}
$R$ を Cohen--Macaulay 局所環とする。
任意の素イデアル $\mathfrak p \subset R$ に対して、
環 $R_{\mathfrak p}$ も Cohen--Macaulay である。
\end{lemma}

\begin{proof}
補題 \ref{algebra-lemma-localize-CM-module} の特別な場合である。
\end{proof}

\begin{definition}
\label{algebra-definition-ring-CM}
ネーター環 $R$ のすべての局所環が Cohen--Macaulay であるとき、
これを {\it Cohen--Macaulay} という。
\end{definition}

\begin{lemma}
\label{algebra-lemma-CM-polynomial-algebra}
$R$ をネーター Cohen--Macaulay 環とする。
$R$ 上の任意の多項式環は Cohen--Macaulay である。
\end{lemma}

\begin{proof}
補題 \ref{algebra-lemma-maximal-CM-polynomial-algebra} の特別な場合である。
\end{proof}

\begin{lemma}
\label{algebra-lemma-dimension-shift}
$R$ を次元 $d$ のネーター局所 Cohen--Macaulay 環とする。
$0 \to K \to R^{\oplus n} \to M \to 0$ を $R$-加群の完全列とする。
このとき $M = 0$ であるか、$\text{depth}(K) > \text{depth}(M)$ であるか、
または $\text{depth}(K) = \text{depth}(M) = d$ である。
\end{lemma}

\begin{proof}
補題 \ref{algebra-lemma-depth-in-ses} の特別な場合である。
\end{proof}

\begin{lemma}
\label{algebra-lemma-mcm-resolution}
$R$ を次元 $d$ のネーター局所 Cohen--Macaulay 環とする。
$M$ を深さ $e$ の有限 $R$-加群とする。
各 $F_i$ が有限自由で $K$ が極大 Cohen--Macaulay である完全複体
$$
0 \to K \to F_{d-e-1} \to \ldots \to F_0 \to M \to 0
$$
が存在する。
% [P01-E0511] The frozen d=e endpoint and F_{-1} indexing remain unchanged; the zero-step case is sidecar-only.
\end{lemma}

\begin{proof}
定義と補題 \ref{algebra-lemma-dimension-shift} から直ちに従う。
\end{proof}

\begin{lemma}
\label{algebra-lemma-find-sequence-image-regular}
$\varphi : A \to B$ を局所環の写像とする。
$B$ はネーターかつ Cohen--Macaulay であり、
$\mathfrak m_B = \sqrt{\varphi(\mathfrak m_A) B}$ であると仮定する。
このとき $A$ の元の列 $f_1, \ldots, f_{\dim(B)}$ で、
$\varphi(f_1), \ldots, \varphi(f_{\dim(B)})$ が
$B$ における正則列となるものが存在する。
\end{lemma}

\begin{proof}
$\dim(B)$ に関する帰納法により、次を証明すれば十分である。
$\dim(B) \geq 1$ ならば、$B$ における非零因子に写る
$A$ の元 $f$ を見つけられる。
補題 \ref{algebra-lemma-reformulate-CM} により、
$B$ の有限個の極小素イデアル
$\mathfrak q_1, \ldots, \mathfrak q_r$ のいずれにも
$B$ における像が含まれないような $f \in A$ を見つければ十分である。
仮定 $\mathfrak m_B = \sqrt{\varphi(\mathfrak m_A) B}$ により
$\mathfrak m_A \not \subset \varphi^{-1}(\mathfrak q_i)$ と分かる。
従って補題 \ref{algebra-lemma-silly} により $f$ を見つけられる。
\end{proof}

\section{鎖状環}
\label{algebra-section-catenary}

\noindent
Topology, Section \ref{topology-section-catenary-spaces} と比較せよ。

\begin{definition}
\label{algebra-definition-catenary}
環 $R$ が {\it 鎖状} であるとは、素イデアルの任意の組
$\mathfrak p \subset \mathfrak q$ に対して、素イデアルのすべての有限鎖
$\mathfrak p = \mathfrak p_0 \subset \mathfrak p_1 \subset \ldots \subset
\mathfrak p_e = \mathfrak q$ の長さを抑える整数が存在し、
かつそのようなすべての極大鎖が同じ長さをもつことをいう。
\end{definition}

\begin{lemma}
\label{algebra-lemma-catenary}
環 $R$ が鎖状であることと、位相空間 $\Spec(R)$ が鎖状であることは同値である
（Topology, Definition \ref{topology-definition-catenary} を参照せよ）。
\end{lemma}

\begin{proof}
定義と、補題 \ref{algebra-lemma-irreducible} における既約閉部分集合の特徴づけから
直ちに従う。
\end{proof}

\noindent
一般に、$R$ が鎖状でも有限生成 $R$-代数が鎖状であるとは限らない。
そこで次の定義を置く。

\begin{definition}
\label{algebra-definition-universally-catenary}
ネーター環 $R$ の任意の有限型 $R$-代数が鎖状であるとき、
この環は {\it 普遍鎖状} であるという。
\end{definition}

\noindent
非ネーター環に対してこの定義が適切か明らかでないため、
ネーター環に制限する。補題 \ref{algebra-lemma-quotient-catenary} により、
ネーター環 $R$ が普遍鎖状であることを調べるには、
各多項式環 $R[x_1, \ldots, x_n]$ が鎖状であることを調べれば十分である。

\begin{lemma}
\label{algebra-lemma-localization-catenary}
鎖状環の任意の局所化は鎖状である。
ネーター普遍鎖状環の任意の局所化は普遍鎖状である。
\end{lemma}

\begin{proof}
$A$ を環とし、$S \subset A$ を乗法的部分集合とする。
補題 \ref{algebra-lemma-spec-localization} における $\Spec(S^{-1}A)$ の記述から、
$A$ が鎖状ならば $S^{-1}A$ も鎖状であると分かる。
$S^{-1}A \to C$ が有限型ならば、ある有限型環写像
$A \to B$ に対して $C = S^{-1}B$ である。
従って $A$ がネーターかつ普遍鎖状ならば $B$ は鎖状であり、
$C$ も鎖状であることが分かる。これで補題が証明された。
\end{proof}

\begin{lemma}
\label{algebra-lemma-universally-catenary}
$A$ をネーター普遍鎖状環とする。
$A$ 上本質的に有限型である任意の $A$-代数は普遍鎖状である。
\end{lemma}

\begin{proof}
$B$ が有限型 $A$-代数ならば、補題 \ref{algebra-lemma-Noetherian-permanence} により
$B$ はネーターである。任意の有限型 $B$-代数は有限型 $A$-代数であり、
$A$ が普遍鎖状であるという仮定により鎖状である。
従って $B$ は普遍鎖状である。補題 \ref{algebra-lemma-localization-catenary} により
$B$ の任意の局所化も普遍鎖状であり、証明が終わる。
\end{proof}

\begin{lemma}
\label{algebra-lemma-catenary-check-local}
$R$ を環とする。次は同値である。
\begin{enumerate}
\item $R$ は鎖状である。
\item すべての素イデアル $\mathfrak p$ に対して
$R_\mathfrak p$ は鎖状である。
\item すべての極大イデアル $\mathfrak m$ に対して
$R_\mathfrak m$ は鎖状である。
\end{enumerate}
$R$ がネーターであると仮定する。次は同値である。
\begin{enumerate}
\item $R$ は普遍鎖状である。
\item すべての素イデアル $\mathfrak p$ に対して
$R_\mathfrak p$ は普遍鎖状である。
\item すべての極大イデアル $\mathfrak m$ に対して
$R_\mathfrak m$ は普遍鎖状である。
\end{enumerate}
\end{lemma}

\begin{proof}
いずれの場合も (1) $\Rightarrow$ (2) は
補題 \ref{algebra-lemma-localization-catenary} から従い、
(2) $\Rightarrow$ (3) は直ちに従う。
$R$ のすべての極大イデアル $\mathfrak m$ に対して
$R_\mathfrak m$ が鎖状であると仮定する。
$\mathfrak p \subset \mathfrak q$ が $R$ の素イデアルならば、
$\mathfrak q \subset \mathfrak m$ となる極大イデアルを選ぶ。
$\mathfrak p$ と $\mathfrak q$ の間の素イデアル鎖は、
$\mathfrak pR_\mathfrak m$ と $\mathfrak qR_\mathfrak m$ の間の
素イデアル鎖と一対一に対応するので、$R$ は鎖状であると分かる。
% [P01-E0512] The frozen “primes ideals” grammar is rendered as the intended prime-ideal chains.
$R$ はネーターであり、$R$ のすべての極大イデアル $\mathfrak m$ に対して
$R_\mathfrak m$ が普遍鎖状であると仮定する。
$R \to S$ を有限型環写像とする。$\mathfrak q$ を、
素イデアル $\mathfrak p \subset R$ の上にある $S$ の素イデアルとする。
$R$ において $\mathfrak p \subset \mathfrak m$ となる極大イデアルを選ぶ。
すると $R_\mathfrak p$ は $R_\mathfrak m$ の局所化なので、
補題 \ref{algebra-lemma-localization-catenary} により普遍鎖状である。
従って $S_\mathfrak p$ は $R_\mathfrak p$ 上の有限型環として鎖状である。
ゆえにその局所化 $S_\mathfrak q$ も鎖状である。
先に扱った第一の場合により $S$ は鎖状である。
\end{proof}

\begin{lemma}
\label{algebra-lemma-quotient-catenary}
鎖状環の任意の商は鎖状である。
ネーター普遍鎖状環の任意の商は普遍鎖状である。
\end{lemma}

\begin{proof}
$A$ を環とし、$I \subset A$ をイデアルとする。
補題 \ref{algebra-lemma-spec-closed} における $\Spec(A/I)$ の記述から、
$A$ が鎖状ならば $A/I$ も鎖状であると分かる。
第二の主張は補題 \ref{algebra-lemma-universally-catenary} の特別な場合である。
\end{proof}

\begin{lemma}
\label{algebra-lemma-catenary-check-irreducible}
$R$ をネーター環とする。
\begin{enumerate}
\item $R$ が鎖状であることと、すべての極小素イデアル
$\mathfrak p$ に対して $R/\mathfrak p$ が鎖状であることは同値である。
\item $R$ が普遍鎖状であることと、すべての極小素イデアル
$\mathfrak p$ に対して $R/\mathfrak p$ が普遍鎖状であることは同値である。
\end{enumerate}
\end{lemma}

\begin{proof}
$\mathfrak a \subset \mathfrak b$ が $R$ の素イデアルの包含ならば、
$\mathfrak p \subset \mathfrak a$ となる極小素イデアルを見つけられるので、
第一の主張は明らかである。第二の証明は省略する。
\end{proof}

\begin{lemma}
\label{algebra-lemma-CM-ring-catenary}
ネーター Cohen--Macaulay 環は普遍鎖状である。
より一般に、$R$ がネーター環で、$M$ が
$\text{Supp}(M) = \Spec(R)$ を満たす Cohen--Macaulay
$R$-加群ならば、$R$ は普遍鎖状である。
\end{lemma}

\begin{proof}
補題 \ref{algebra-lemma-CM-polynomial-algebra} により
$R$ 上の多項式環は Cohen--Macaulay なので、
Cohen--Macaulay 環が鎖状であることを示せば十分である。
$R$ を Cohen--Macaulay とし、
$\mathfrak p \subset \mathfrak q$ を $R$ の素イデアルとする。
定義により $R_{\mathfrak q}$ と $R_{\mathfrak p}$ は
Cohen--Macaulay である。素イデアルの極大鎖
$\mathfrak p = \mathfrak p_0 \subset \mathfrak p_1 \subset
\ldots \subset \mathfrak p_n = \mathfrak q$
をとる。次に素イデアルの極大鎖
$\mathfrak q_0 \subset \mathfrak q_1 \subset \ldots \subset
\mathfrak q_m = \mathfrak p$
を選ぶ。補題 \ref{algebra-lemma-maximal-chain-CM} により
$n + m = \dim(R_{\mathfrak q})$ である。
同じ補題により $m = \dim(R_{\mathfrak p})$ である。
従って $n = \dim(R_{\mathfrak q}) - \dim(R_{\mathfrak p})$ は
選び方に依存しない。

\medskip\noindent
より一般の主張を証明するには、補題
\ref{algebra-lemma-maximal-CM-polynomial-algebra} と
\ref{algebra-lemma-maximal-chain-maximal-CM} を用いて、上とまったく同様に論じる。
\end{proof}

\begin{lemma}
\label{algebra-lemma-catenary-Noetherian-local}
$(A, \mathfrak m)$ をネーター局所環とする。次は同値である。
\begin{enumerate}
\item $A$ は鎖状である。
\item $\mathfrak p \mapsto \dim(A/\mathfrak p)$ は
$\Spec(A)$ 上の次元関数である。
\end{enumerate}
\end{lemma}

\begin{proof}
$A$ が鎖状ならば、Topology, Lemma
\ref{topology-lemma-locally-dimension-function} と
補題 \ref{algebra-lemma-catenary} により、
$\Spec(A)$ は次元関数 $\delta$ をもつ。
$\delta(\mathfrak m) = 0$ と仮定してよい。このとき
Topology, Lemma \ref{topology-lemma-dimension-function-catenary} により
$$
\delta(\mathfrak p) = \text{codim}(V(\mathfrak m), V(\mathfrak p)) =
\dim(A/\mathfrak p)
$$
であると分かる。このようにして (1) が (2) を含意することが分かる。
逆の含意も Topology, Lemma
\ref{topology-lemma-dimension-function-catenary} から従う。
\end{proof}

\section{正則局所環}
\label{algebra-section-regular}

\noindent
正則局所環は定義 \ref{algebra-definition-regular-local} で定義されている。
正則局所環のすべての素イデアルにおける局所化が正則であることを
示すのは、それほど容易ではない。実際、ここまで展開してきた内容の
かなりの部分は、この事実の証明を目指したものである。
命題 \ref{algebra-proposition-finite-gl-dim-regular} を参照し、
そこから参照を遡られたい。

\begin{lemma}
\label{algebra-lemma-regular-graded}
$(R, \mathfrak m, \kappa)$ を次元 $d$ の正則局所環とする。
次数付き環 $\bigoplus \mathfrak m^n / \mathfrak m^{n + 1}$ は
次数付き多項式環 $\kappa[X_1, \ldots, X_d]$ に同型である。
\end{lemma}

\begin{proof}
$x_1, \ldots, x_d$ を極大イデアル $\mathfrak m$ の極小生成系とする
（定義 \ref{algebra-definition-regular-local} を参照せよ）。
$X_i$ を $\mathfrak m/\mathfrak m^2$ における $x_i$ の類へ写す全射
$\kappa[X_1, \ldots, X_d]
\to \bigoplus \mathfrak m^n/\mathfrak m^{n + 1}$
がある。命題 \ref{algebra-proposition-dimension} により $d(R) = d$ なので、
数値多項式
$n \mapsto \dim_\kappa \mathfrak m^n/\mathfrak m^{n + 1}$
の次数は $d - 1$ である。補題 \ref{algebra-lemma-quotient-smaller-d} により、
全射 $\kappa[X_1, \ldots, X_d]
\to \bigoplus \mathfrak m^n/\mathfrak m^{n + 1}$ は同型であると結論する。
\end{proof}

\begin{lemma}
\label{algebra-lemma-regular-domain}
任意の正則局所環は整域である。
\end{lemma}

\begin{proof}
補題 \ref{algebra-lemma-intersect-powers-ideal-module-zero} による
$\bigcap \mathfrak m^n = 0$ を用いる。
$fg = 0$ を満たす $f, g \in R$ をとる。
$a, b$ は極大であり、
$f \in \mathfrak m^a$ かつ $g \in \mathfrak m^b$ であると仮定する。
% [P01-E0513] The frozen proof chooses maximal adic orders without first excluding zero factors; no corrective hypothesis is rendered.
$fg = 0 \in \mathfrak m^{a + b + 1}$ なので、
補題 \ref{algebra-lemma-regular-graded} の結果から
$f \in \mathfrak m^{a + 1}$ または
$g \in \mathfrak m^{b + 1}$ と分かる。これは矛盾である。
\end{proof}

\begin{lemma}
\label{algebra-lemma-regular-ring-CM}
$R$ を正則局所環とし、$x_1, \ldots, x_d$ を
極大イデアル $\mathfrak m$ の極小生成系とする。
このとき $x_1, \ldots, x_d$ は正則列であり、
各 $R/(x_1, \ldots, x_c)$ は次元 $d - c$ の正則局所環である。
特に $R$ は Cohen--Macaulay である。
\end{lemma}

\begin{proof}
補題 \ref{algebra-lemma-one-equation} により $R/x_1R$ は
次元 $\geq d - 1$ のネーター局所環であり、
$x_2, \ldots, x_d$ はその極大イデアルを生成する。
従って定義により正則局所環である。
補題 \ref{algebra-lemma-regular-domain} により $R$ は整域なので、
$x_1$ は非零因子である。
\end{proof}

\begin{lemma}
\label{algebra-lemma-regular-quotient-regular}
$R$ を正則局所環とする。$I \subset R$ を、
$R/I$ も正則局所環となるイデアルとする。
このとき $R$ の極大イデアル $\mathfrak m$ の極小生成系
$x_1, \ldots, x_d$ で、ある $0 \leq c \leq d$ に対して
$I = (x_1, \ldots, x_c)$ となるものが存在する。
\end{lemma}

\begin{proof}
$\dim(R) = d$ および $\dim(R/I) = d - c$ とする。
$\overline{\mathfrak m} = \mathfrak m/I$ を $R/I$ の極大イデアルと書く。
$\kappa = R/\mathfrak m$ とおく。正則局所環の定義により
$$
\dim_\kappa((I + \mathfrak m^2)/\mathfrak m^2) =
\dim_\kappa(\mathfrak m/\mathfrak m^2)
- \dim(\overline{\mathfrak m}/\overline{\mathfrak m}^2) = d - (d - c) = c
$$
である。
% [P01-E0514] The frozen third dimension operator remains without its kappa subscript.
従って、$\mathfrak m/\mathfrak m^2$ における像が一次独立となる
$x_1, \ldots, x_c \in I$ を選び、$x_{c + 1}, \ldots, x_d$ を補って
$\mathfrak m$ の極小生成系を得ることができる。
誘導写像 $R/(x_1, \ldots, x_c) \to R/I$ は、
同じ次元をもつ正則局所環の間の全射である
（補題 \ref{algebra-lemma-regular-ring-CM}）。
従ってその核は零、すなわち $I = (x_1, \ldots, x_c)$ である。
実際、そうでなければ補題 \ref{algebra-lemma-regular-domain} と
\ref{algebra-lemma-one-equation} により
$\dim(R/I) < \dim(R/(x_1, \ldots, x_c))$ となる。
\end{proof}

\begin{lemma}
\label{algebra-lemma-free-mod-x}
$R$ をネーター局所環とする。$x \in \mathfrak m$ とする。
$M$ を有限 $R$-加群とし、$x$ は $M$ 上の非零因子であり、
$M/xM$ は $R/xR$ 上自由であると仮定する。
このとき $M$ は $R$ 上自由である。
\end{lemma}

\begin{proof}
$M/xM$ の $R/xR$-基底に写る $M$ の元を
$m_1, \ldots, m_r$ とする。Nakayama の補題 \ref{algebra-lemma-NAK} により、
$m_1, \ldots, m_r$ は $M$ を生成する。
$\sum a_i m_i = 0$ が関係ならば、すべての $i$ に対して
$a_i \in xR$ である。従ってある $b_i \in R$ に対して
$a_i = b_i x$ である。ゆえに $R^r \to M$ の核 $K$ は
$xK = K$ を満たし、Nakayama の補題により零である。
\end{proof}

\begin{lemma}
\label{algebra-lemma-regular-mcm-free}
$R$ を正則局所環とする。$R$ 上の任意の極大
Cohen--Macaulay 加群は自由である。
\end{lemma}

\begin{proof}
$M$ を $R$ 上の極大 Cohen--Macaulay 加群とする。
$x \in \mathfrak m$ を、$\mathfrak m$ を生成する正則列の一部とする。
命題 \ref{algebra-proposition-CM-module} により $x$ は $M$ 上の非零因子であり、
$M/xM$ は $R/xR$ 上の極大 Cohen--Macaulay 加群である。
$\dim(R)$ に関する帰納法により $M/xM$ は自由である。
補題 \ref{algebra-lemma-free-mod-x} により結論を得る。
\end{proof}

\begin{lemma}
\label{algebra-lemma-regular-mod-x}
$R$ をネーター局所環とする。
$x \in \mathfrak m$ を、$R/xR$ が正則局所環となる非零因子とする。
このとき $R$ は正則局所環である。
より一般に、$x_1, \ldots, x_r$ が $R$ における正則列で、
$R/(x_1, \ldots, x_r)$ が正則局所環ならば、
$R$ は正則局所環である。
\end{lemma}

\begin{proof}
$x$ と、$R/xR$ の極大イデアルの極小生成系の持ち上げを合わせると、
$\mathfrak m$ の $\dim(R)$ 個の生成元が得られるので、これは成り立つ。
補題 \ref{algebra-lemma-one-equation} を用いよ。
最後の主張は第一の主張と帰納法から従う。
\end{proof}

\begin{lemma}
\label{algebra-lemma-colimit-regular}
$(R_i, \varphi_{ii'})$ を、遷移写像が局所環写像である
局所環の有向系とする。各 $R_i$ が正則局所環であり、
$R = \colim R_i$ がネーターならば、$R$ は正則局所環である。
\end{lemma}

\begin{proof}
$\mathfrak m \subset R$ を極大イデアルとする。これは極大イデアル
$\mathfrak m_i \subset R_i$ の余極限である。
$d = \dim \mathfrak m/\mathfrak m^2$ に関する帰納法で補題を証明する。
% [P01-E0516] The frozen induction parameter remains ungrouped and lacks the residue-field subscript.
$d = 0$ ならば $R = R/\mathfrak m$ は体であり、$R$ は正則局所環である。
$d > 0$ ならば $x \in \mathfrak m$、$x \not \in \mathfrak m^2$ を選ぶ。
ある $i$ に対して、$x$ に写る $x_i \in \mathfrak m_i$ を見つけられる。
$R/xR = \colim_{i' \geq i} R_{i'}/x_iR_{i'}$ は
ネーター局所環であることに注意する。
補題 \ref{algebra-lemma-regular-ring-CM} により
$R_{i'}/x_iR_{i'}$ は正則局所環である。
従って帰納法により $R/xR$ は正則局所環である。
各 $R_i$ は整域なので（補題 \ref{algebra-lemma-regular-graded}）、
$R$ は整域であると分かる。
% [P01-E0515] The frozen reference to lemma-regular-graded is preserved; the direct domain lemma remains sidecar-only.
従って $x$ は非零因子であり、補題 \ref{algebra-lemma-regular-mod-x} により
$R$ は正則局所環であると結論する。
\end{proof}

\section{環のエピ射}
\label{algebra-section-epimorphism}

\noindent
任意の圏には {\it エピ射} の概念がある。
この内容の一部は \cite{Autour} と \cite{Mazet} から採られている。

\begin{lemma}
\label{algebra-lemma-epimorphism}
$R \to S$ を環写像とする。次は同値である。
\begin{enumerate}
\item $R \to S$ はエピ射である。
\item 二つの環写像 $S \to S \otimes_R S$ は等しい。
\item いずれの環写像 $S \to S \otimes_R S$ も同型である。
\item 環写像 $S \otimes_R S \to S$ は同型である。
\end{enumerate}
\end{lemma}

\begin{proof}
省略する。
\end{proof}

\begin{lemma}
\label{algebra-lemma-composition-epimorphism}
環の二つのエピ射の合成はエピ射である。
\end{lemma}

\begin{proof}
省略する。ヒント：これは任意の圏で成り立つ。
\end{proof}

\begin{lemma}
\label{algebra-lemma-base-change-epimorphism}
$R \to S$ が環のエピ射で、$R \to R'$ が任意の環写像ならば、
$R' \to R' \otimes_R S$ はエピ射である。
\end{lemma}

\begin{proof}
省略する。ヒント：これは押し出しをもつ任意の圏で成り立つ。
\end{proof}

\begin{lemma}
\label{algebra-lemma-permanence-epimorphism}
$A \to B \to C$ が環写像で、$A \to C$ がエピ射ならば、
$B \to C$ もエピ射である。
\end{lemma}

\begin{proof}
省略する。ヒント：これは任意の圏で成り立つ。
\end{proof}

\noindent
これは特に、$R \to S$ が像 $\overline{R} \subset S$ をもつエピ射ならば、
$\overline{R} \to S$ がエピ射であることを意味する。
% [P01-E0517] The frozen misplaced English comma has no Japanese analogue.
従ってエピ射に関する結果を証明するとき、写像が単射であると
仮定できることが多い。次の補題は特に、すべての局所化が
エピ射であることを意味する。

\begin{lemma}
\label{algebra-lemma-epimorphism-local}
$R \to S$ を環写像とする。次は同値である。
\begin{enumerate}
\item $R \to S$ はエピ射である。
\item $R$ の各素イデアル $\mathfrak p$ に対して
$R_{\mathfrak p} \to S_{\mathfrak p}$ はエピ射である。
\end{enumerate}
\end{lemma}

\begin{proof}
$S_{\mathfrak p} = R_{\mathfrak p} \otimes_R S$ なので
（補題 \ref{algebra-lemma-tensor-localization} を参照せよ）、
補題 \ref{algebra-lemma-base-change-epimorphism} により (1) は (2) を含意する。
逆に (2) が成り立つと仮定する。
$a, b : S \to A$ を、$R \to S$ との合成が等しい、
$S$ から環 $A$ への二つの環写像とする。
% [P01-E0518] The frozen malformed “equalizing the map” phrase is rendered by equality of the two composites.
仮定により、$R$ のすべての素イデアル $\mathfrak p$ に対して
誘導写像 $a_{\mathfrak p}, b_{\mathfrak p} : S_{\mathfrak p} \to A_{\mathfrak p}$
は等しい。従って $A \subset \prod_{\mathfrak p} A_{\mathfrak p}$ なので
$a = b$ である（補題 \ref{algebra-lemma-characterize-zero-local} を参照せよ）。
\end{proof}

\begin{lemma}
\label{algebra-lemma-finite-epimorphism-surjective}
\begin{slogan}
環写像が全射であることと、それが有限エピ射であることは同値である。
\end{slogan}
$R \to S$ を環写像とする。次は同値である。
\begin{enumerate}
\item $R \to S$ は有限なエピ射である。
\item $R \to S$ は全射である。
\end{enumerate}
\end{lemma}

\begin{proof}
（この補題は文献で何度も再証明されてきたようであり、
多くの異なる証明がある。）
全射環写像がエピ射であることは明らかである。
$R \to S$ を有限環写像とし、$S \otimes_R S \to S$ が同型であると仮定する。
目標は $R \to S$ が全射であることを示すことである。
$S/R$ が零でないと仮定する。完全列 $R \to S \to S/R \to 0$ から
完全列
$$
R \otimes_R S \to S \otimes_R S \to S/R \otimes_R S \to 0.
$$
が得られる。仮定により第一の矢印は同型なので、
$S/R \otimes_R S = 0$ と結論する。
従って $S/R \otimes_R S/R = 0$ でもある。
補題 \ref{algebra-lemma-trivial-filter-finite-module} により、ある真のイデアル
$I \subset R$ に対して $R$-加群の全射 $S/R \to R/I$ が存在する。
従って全射
$S/R \otimes_R S/R \to R/I \otimes_R R/I = R/I \not = 0$
が存在し、矛盾である。
\end{proof}

\begin{lemma}
\label{algebra-lemma-faithfully-flat-epimorphism}
忠実平坦なエピ射は同型である。
\end{lemma}

\begin{proof}
写像 $S \to S \otimes_R S$ は写像 $R \to S$ と $S$ との
テンソル積なので、補題 \ref{algebra-lemma-epimorphism} の (3) から明らかである。
\end{proof}

\begin{lemma}
\label{algebra-lemma-epimorphism-over-field}
$k \to S$ がエピ射で $k$ が体ならば、$S = k$ または $S = 0$ である。
\end{lemma}

\begin{proof}
$k$ 上の任意の非零代数は忠実平坦なので、
補題 \ref{algebra-lemma-faithfully-flat-epimorphism} の結果から明らかである。
あるいは、$\dim_k(R) \leq 1$ でない限り
$R \to R \otimes_k R$ は全射になりえないことを直接論じてもよい。
% [P01-E0519] The frozen proof's undefined R is preserved; replacing it by S remains sidecar-only.
\end{proof}

\begin{lemma}
\label{algebra-lemma-epimorphism-injective-spec}
$R \to S$ を環のエピ射とする。このとき
\begin{enumerate}
\item $\Spec(S) \to \Spec(R)$ は単射である。
\item $\mathfrak p \subset R$ の上にある $\mathfrak q \subset S$ に対して
$\kappa(\mathfrak p) = \kappa(\mathfrak q)$ である。
\end{enumerate}
\end{lemma}

\begin{proof}
$\mathfrak p$ を $R$ の素イデアルとする。この写像のファイバーは
ファイバー環 $S \otimes_R \kappa(\mathfrak p)$ のスペクトルである。
補題 \ref{algebra-lemma-base-change-epimorphism} により写像
$\kappa(\mathfrak p) \to S \otimes_R \kappa(\mathfrak p)$ はエピ射である。
従って補題 \ref{algebra-lemma-epimorphism-over-field} により
$S \otimes_R \kappa(\mathfrak p) = 0$ または
$S \otimes_R \kappa(\mathfrak p) = \kappa(\mathfrak p)$ であり、
(1) と (2) が証明される。
\end{proof}

\begin{lemma}
\label{algebra-lemma-relations}
$R$ を環とする。$M$, $N$ を $R$-加群とする。
$\{x_i\}_{i \in I}$ を $M$ の生成元集合とし、
$\{y_j\}_{j \in J}$ を $N$ の生成元集合とする。
$\{m_j\}_{j \in J}$ を $M$ の元の族とし、
有限個を除くすべての $j$ に対して $m_j = 0$ とする。このとき
$$
\sum\nolimits_{j \in J} m_j \otimes y_j = 0 \text{ in } M \otimes_R N
$$
であることと、次のことは同値である。
有限個を除くすべての組 $(i, j)$ に対して $a_{i, j} = 0$ となる
$a_{i, j} \in R$ で、
\begin{align*}
m_j & = \sum\nolimits_{i \in I} a_{i, j} x_i \quad\text{for all } j \in J, \\
0 & = \sum\nolimits_{j \in J} a_{i, j} y_j \quad\text{for all } i \in I.
\end{align*}
を満たすものが存在する。
\end{lemma}

\begin{proof}
十分性は直ちに分かる。
$\sum_{j \in J} m_j \otimes y_j = 0$ と仮定する。
$\bigoplus\nolimits_{j \in J} R$ の第 $j$ 基底ベクトルが
$y_j$ に写る短完全列
$$
0 \to K \to \bigoplus\nolimits_{j \in J} R \to N \to 0
$$
を考える。これと $M$ とのテンソル積をとると、完全列
$$
K \otimes_R M \to \bigoplus\nolimits_{j \in J} M \to N \otimes_R M \to 0.
$$
を得る。仮定により、$\sum k_i \otimes x_i$ が中央の項の元
$(m_j)_{j \in J}$ に写るような元 $k_i \in K$ が存在する。
$k_i = (a_{i, j})_{j \in J}$ と書けば、望むものを得る。
% [P01-E0520] The frozen participial-clause grammar is rendered naturally.
\end{proof}

\begin{lemma}
\label{algebra-lemma-kernel-difference-projections}
$\varphi : R \to S$ を環写像とする。$g \in S$ とする。
次は同値である。
\begin{enumerate}
\item $S \otimes_R S$ において $g \otimes 1 = 1 \otimes g$ である。
\item ある $n \geq 0$ と元 $y_i, z_j \in S$ および
$1 \leq i, j \leq n$ に対する $x_{i, j} \in R$ が存在して、
\begin{enumerate}
\item $g = \sum_{i, j \leq n} x_{i, j} y_i z_j$ である。
\item 各 $j$ に対して $\sum x_{i, j}y_i \in \varphi(R)$ である。
\item 各 $i$ に対して $\sum x_{i, j}z_j \in \varphi(R)$ である。
\end{enumerate}
\end{enumerate}
\end{lemma}

\begin{proof}
(2) が (1) を含意することは明らかである。
逆に $g \otimes 1 = 1 \otimes g$ と仮定する。
$0, 1 \in I$、$s_0 = 1$、$s_1 = g$ となる
$R$-加群としての $S$ の生成元 $\{s_i\}_{i \in I}$ を選ぶ。
関係 $g \otimes s_0 + (-1) \otimes s_1 = 0$ に
補題 \ref{algebra-lemma-relations} を適用する。
すると $g = \sum_i a_{i, 0} s_i$、
$-1 = \sum_i a_{i, 1} s_i$、$j \not = 0, 1$ に対して
$0 = \sum_i a_{i, j} s_i$ を満たし、さらにすべての $i$ に対して
$\sum_j a_{i, j}s_j = 0$ を満たす $a_{i, j} \in R$ が存在する。
そこで
$$
\sum\nolimits_{i, j \not = 0} a_{i, j} s_i s_j = -g + a_{0, 0}
$$
であり、各 $j \not = 0$ に対して
$\sum_{i \not = 0} a_{i, j}s_i \in R$ である。
% [P01-E0521] The frozen membership in R is preserved; the image phi(R) remains sidecar-only.
これは $-g + a_{0, 0}$ が (2) の形に書けることを証明する。
従って $g$ は (2) の形に書ける。詳細は省略する。
ヒント：(2) の形の表示をもつ $S$ の元全体が
$S$ の $R$-部分代数をなすことを示せ。
\end{proof}

\begin{remark}
\label{algebra-remark-matrices-associated-to-elements-epicenter}
$R \to S$ を環写像とする。$g \otimes 1 = 1 \otimes g$ を満たす
元 $g \in S$ の集合を $S$ の {\it エピセンター} と呼ぶことがある。
これは $R$-代数である。補題
\ref{algebra-lemma-kernel-difference-projections} の構成により、
エピセンターの各 $g$ に対して行列分解
$$
(g) = Y X Z
$$
で、$X \in \text{Mat}(n \times n, R)$、
$Y \in \text{Mat}(1 \times n, S)$、および
$Z \in \text{Mat}(n \times 1, S)$ を満たすものを得る。
実際、$x_{i, j}, y_i, z_j$ を補題の (2) のとおりとする。
$X = (x_{i, j})$ とおき、$y$ を成分が $y_i$ である行ベクトル、
$z$ を成分が $z_j$ である列ベクトルとする。
% [P01-E0522] The frozen lowercase y,z declarations and uppercase Y,Z products are preserved.
この記法のもとで、補題 \ref{algebra-lemma-kernel-difference-projections} の
条件 (b) と (c) は、ちょうど
$Y X \in \text{Mat}(1 \times n, R)$、
$X Z \in \text{Mat}(n \times 1, R)$
を意味する。行列の三つ組 $(X, YX, XZ)$ を考えることは
非常に便利である。$n \in \mathbf{N}$ と三つ組 $(P, U, V)$ が
与えられたとき、上のような行列分解で
$P = X$、$U = YX$、$V = XZ$ を満たすものが存在するならば、
$(P, U, V)$ は {\it $g$ に随伴する $n$-三つ組} であるという。
\end{remark}

\begin{lemma}
\label{algebra-lemma-epimorphism-cardinality}
$R \to S$ を環のエピ射とする。このとき $S$ の濃度は
$R$ の濃度以下である。式で書けば $|S| \leq |R|$ である。
\end{lemma}

\begin{proof}
$R \to S$ がエピ射であるという条件は、各 $g \in S$ が
$g \otimes 1 = 1 \otimes g$ を満たすことを意味する
（補題 \ref{algebra-lemma-epimorphism} を参照せよ）。
注意 \ref{algebra-remark-matrices-associated-to-elements-epicenter} で
導入した記法を用いる。
$g, g' \in S$ とし、$(P, U, V)$ が $g$ と $g'$ の双方に
随伴する $n$-三つ組であると仮定する。このとき $g = g'$ と主張する。
実際、$g$ の行列分解 $(g) = YXZ$ に対して
$(P, U, V) = (X, YX, XZ)$ と書き、$g'$ の行列分解
$(g') = Y'X'Z'$ に対して
$(P, U, V) = (X', Y'X', X'Z')$ と書く。このとき
$$
(g) = YXZ = UZ = Y'X'Z = Y'PZ = Y'XZ = Y'V = Y'X'Z' = (g')
$$
であるから $g = g'$ である。
従って $S$ の濃度は可能な三つ組の個数で抑えられ、その濃度は高々
$\sup_{n \in \mathbf{N}} |R|^n$ である。
$R$ が無限ならば、これは高々 $|R|$ である
（\cite[Ch. I, 10.13]{Kunen} を参照せよ）。

\medskip\noindent
$R$ が有限環ならば、上の議論から分かるのは $S$ が高々可算であることだけである。
実際この場合 $R$ はアルティン的であり、写像 $R \to S$ は全射である。
この場合の証明は省略する。
\end{proof}

\begin{lemma}
\label{algebra-lemma-epimorphism-modules}
環準同型 $R \to S$ に対して、次は同値である。
\begin{enumerate}
\item $R \to S$ は環のエピ射である。
\item 任意の $S$-加群 $N_1, N_2$ に対して
$\Hom_S(N_1, N_2) = \Hom_R(N_1, N_2)$ である。
\item 制限関手 $\text{Mod}_S \to \text{Mod}_R$ は充満忠実である。
\end{enumerate}
\end{lemma}

\begin{proof}
(2) と (3) が定義により同値であることに注意する。

\medskip\noindent
(1) を仮定し、$N_1, N_2$ を $S$-加群、
$\varphi : N_1 \to N_2$ を $R$-線形写像とする。
任意の $x \in N_1$ に対し、規則
$g \otimes g' \mapsto g\varphi(g'x)$ で定義される写像
$S \otimes_R S \to N_2$ を考える。
二つの写像 $S \to S \otimes_R S$ はともに同型なので
（補題 \ref{algebra-lemma-epimorphism}）、
$g \varphi(g'x) = gg'\varphi(x) = \varphi(gg' x)$ と結論する。
従って $\varphi$ は $S$-線形である。

\medskip\noindent
(2) を仮定する。$N_1 = S \otimes_R S$ を $S$-加群とみなし、
その左 $S$-加群作用を用いる。また $N_2 = S \otimes_R S$ には
右 $S$-加群作用を入れる。$R$-加群としては $N_1 = N_2$ なので、
(2) により $S \otimes_R S$ 上の二つの $S$-加群構造は一致する。
補題 \ref{algebra-lemma-epimorphism} により、これは (1) を含意する。
\end{proof}

\section{純イデアル}
\label{algebra-section-pure-ideals}

\noindent
この節の内容は多くの論文で論じられている。例えば
\cite{Lazard}、\cite{Bkouche}、\cite{DeMarco} を参照せよ。

\begin{definition}
\label{algebra-definition-pure-ideal}
$R$ を環とする。商環 $R/I$ が $R$ 上平坦であるとき、
$I \subset R$ は {\it 純} であるという。
\end{definition}

\begin{lemma}
\label{algebra-lemma-pure}
$R$ を環とし、$I \subset R$ をイデアルとする。次は同値である。
\begin{enumerate}
\item $I$ は純である。
\item すべてのイデアル $J \subset R$ に対して $J \cap I = IJ$ である。
\item すべての有限生成イデアル $J \subset R$ に対して
$J \cap I = JI$ である。
\item すべての $x \in R$ に対して $(x) \cap I = xI$ である。
\item すべての $x \in I$ に対して、ある $y \in I$ が存在して $x = yx$ である。
\item すべての $x_1, \ldots, x_n \in I$ に対して、
すべての $i = 1, \ldots, n$ について $x_i = yx_i$ を満たす
$y \in I$ が存在する。
\item $R$ のすべての素イデアル $\mathfrak p$ に対して
$IR_{\mathfrak p} = 0$ または $IR_{\mathfrak p} = R_{\mathfrak p}$ である。
\item $\text{Supp}(I) = \Spec(R) \setminus V(I)$ である。
\item $I$ は写像 $R \to (1 + I)^{-1}R$ の核である。
\item $R$ のある乗法的部分集合 $S$ に対して、$R$-代数として
$R/I \cong S^{-1}R$ である。
\item $R$-代数として $R/I \cong (1 + I)^{-1}R$ である。
\end{enumerate}
\end{lemma}

\begin{proof}
$R$ の任意のイデアル $J$ に対して短完全列
$0 \to J \to R \to R/J \to 0$ がある。これと $R/I$ との
テンソル積をとると完全列
$J \otimes_R R/I \to R/I \to R/I + J \to 0$ を得て、
$J \otimes_R R/I = J/JI$ である。
% [P01-E0523] The frozen third term R/I + J is preserved; R/(I + J) remains sidecar-only.
従って (1)、(2)、(3) の同値性は補題 \ref{algebra-lemma-flat} から従う。
さらに、これらは (4) を含意する。

\medskip\noindent
(4) $\Rightarrow$ (5) は自明である。
(5) を仮定し、$x_1, \ldots, x_n \in I$ とする。
$x_i = y_ix_i$ を満たす $y_i \in I$ を選ぶ。
$1 - y = \prod_{i = 1, \ldots, n} (1 - y_i)$ を満たす元
$y \in I$ をとる。このとき、すべての $i = 1, \ldots, n$ に対して
$x_i = yx_i$ である。従って (6) が成り立ち、
(5) $\Leftrightarrow$ (6) である。

\medskip\noindent
(5) を仮定する。$x \in I$ とすると、ある $y \in I$ に対して
$x = yx$ である。従って $x(1 - y) = 0$ であり、これは $x$ が
$(1 + I)^{-1}R$ で零に写ることを示す。
もちろん写像 $R \to (1 + I)^{-1}R$ の核は常に $I$ に含まれる。
従って (5) は (9) を含意する。(9) を仮定すると、任意の
$x \in I$ に対して、ある $y \in I$ が存在して $x(1 - y) = 0$ である。
言い換えれば $x = yx$ である。従って (5) と (9) は同値である。

\medskip\noindent
(5) を仮定する。$\mathfrak p$ を $R$ の素イデアルとする。
$\mathfrak p \not \in V(I)$ ならば $IR_{\mathfrak p} = R_{\mathfrak p}$ である。
$\mathfrak p \in V(I)$、すなわち $I \subset \mathfrak p$ ならば、
$x \in I$ に対して、ある $y \in I$ が存在して $x(1 - y) = 0$ である。
従って $x$ は $R_{\mathfrak p}$ で零に写り、すなわち
$IR_{\mathfrak p} = 0$ である。
% [P01-E0524] The frozen repeated-implies grammar is rendered as the unambiguous inference.
従って (7) が成り立つ。

\medskip\noindent
(7) を仮定する。このとき、$R$ の任意の素イデアル $\mathfrak p$ に対して
$(R/I)_{\mathfrak p}$ は $0$ または $R_{\mathfrak p}$ である。
従って補題 \ref{algebra-lemma-flat-localization} により (1) が成り立つ。
ここまでで (1) -- (7) と (9) はすべて同値である。

\medskip\noindent
$IR_{\mathfrak p} = I_{\mathfrak p}$ なので、(7) は (8) を含意する。
最後に (8) が成り立つならば、これは $I_{\mathfrak p}$ が零加群であることと
$\mathfrak p \in V(I)$ であることが同値であるということにほかならず、
明らかに (7) を意味する。これで (1) -- (9) は同値である。

\medskip\noindent
(1) -- (9) が成り立つと仮定する。(9) により
$R/I \subset (1 + I)^{-1}R$ であり、(7) が与える素イデアルでの
局所化の記述により、写像 $R/I \to (1 + I)^{-1}R$ は全射でもある。
従って (11) が成り立つ。

\medskip\noindent
(11) $\Rightarrow$ (10) は自明である。また、$R$ の局所化は
$R$ 上平坦なので、(10) は (1) を含意する。補題
\ref{algebra-lemma-flat-localization} を参照せよ。
\end{proof}

\begin{lemma}
\label{algebra-lemma-pure-ideal-determined-by-zero-set}
\begin{slogan}
純イデアルはその零点集合によって決まる。
\end{slogan}
$R$ を環とする。$I, J \subset R$ が純イデアルならば、
$V(I) = V(J)$ は $I = J$ を含意する。
\end{lemma}

\begin{proof}
例えば補題 \ref{algebra-lemma-pure} の性質 (7) により、
$I = \Ker(R \to \prod_{\mathfrak p \in V(I)} R_{\mathfrak p})$
だから、このイデアルはそれに付随する閉部分集合から復元できる。
\end{proof}

\begin{lemma}
\label{algebra-lemma-pure-open-closed-specializations}
$R$ を環とする。対応 $I \mapsto V(I)$ は全単射
$$
\{I \subset R \text{ pure}\}
\leftrightarrow
\{Z \subset \Spec(R)\text{ closed and closed under generalizations}\}
$$
を定める。
\end{lemma}

\begin{proof}
$I$ を純イデアルとする。$R \to R/I$ は平坦なので、下降性により
一般化は写像 $\Spec(R/I) \to \Spec(R)$ に沿って持ち上がる。
従って $V(I)$ は一般化について閉じている。これにより、この写像は
良定義である。補題 \ref{algebra-lemma-pure-ideal-determined-by-zero-set} により
この写像は単射である。
$Z \subset \Spec(R)$ が閉じており、かつ一般化について閉じているとする。
$Z = V(J)$ を満たす根基イデアル $J \subset R$ をとる。
$I = \{x \in R : x \in xJ\}$ とおく。$I$ はイデアルであることに注意する。
実際、$x, y \in I$ ならば、$x = xf$ と $y = yg$ を満たす
$f, g \in J$ が存在する。このとき
$$
x + y = (x + y)(f + g - fg)
$$
である。確認は読者に委ねる。
$I$ は純であり $V(I) = V(J)$ であると主張する。
この主張が正しければ、補題の写像は全射であり、補題が従う。

\medskip\noindent
$I \subset J$ なので $V(J) \subset V(I)$ である。
$I \subset \mathfrak p$ を満たす素イデアルをとる。
乗法的部分集合 $S = (R \setminus \mathfrak p)(1 + J)$ を考える。
$I$ の定義と $I \subset \mathfrak p$ により $0 \not \in S$ である。
従って $S$ と交わらない $R$ の素イデアル $\mathfrak q$ をとれる。
補題 \ref{algebra-lemma-localization-zero} と \ref{algebra-lemma-spec-localization} を参照せよ。
従って $\mathfrak q \subset \mathfrak p$ かつ
$\mathfrak q \cap (1 + J) = \emptyset$ である。
これは $\mathfrak q + J$ が $R$ の真のイデアルであることを含意する。
$\mathfrak q + J$ を含む極大イデアル $\mathfrak m$ をとる。このとき
$\mathfrak m \in V(J)$ であり、$Z$ は一般化について閉じていると
仮定したので $\mathfrak q \in V(J) = Z$ である。
さらに $\mathfrak q \subset \mathfrak p$ なので、これは
$\mathfrak p \in V(J)$ を含意する。従って $V(I) = V(J)$ である。

\medskip\noindent
最後に、$V(I) = V(J)$ であり $J$ は根基なので $J = \sqrt{I}$ である。
$x \in I$ をとる。定義により、ある $y \in J$ に対して $x = xy$ である。
このとき $x = xy = xy^2 = \ldots = xy^n$ である。
ある $n > 0$ に対して $y^n \in I$ なので、補題
\ref{algebra-lemma-pure} の性質 (5) が成り立ち、$I$ は実際に純である。
\end{proof}

\begin{lemma}
\label{algebra-lemma-finitely-generated-pure-ideal}
$R$ を環とし、$I \subset R$ をイデアルとする。次は同値である。
\begin{enumerate}
\item $I$ は純かつ有限生成である。
\item $I$ は冪等元によって生成される。
\item $I$ は純で、$V(I)$ は開である。
\item $R/I$ は射影的 $R$-加群である。
\end{enumerate}
\end{lemma}

\begin{proof}
(1) が成り立つならば、補題 \ref{algebra-lemma-pure} により
$I = I \cap I = I^2$ である。従って補題
\ref{algebra-lemma-ideal-is-squared-union-connected} により、
$I$ は冪等元で生成される。従って (1) $\Rightarrow$ (2) である。
(2) が成り立つならば $I = (e)$ であり、$R$-加群として
$R = (1 - e) \oplus (e)$ である。従って $R/I$ は平坦で、
$I$ は純であり、$V(I) = D(1 - e)$ は開である。
従って (2) $\Rightarrow$ (1) $+$ (3) である。
最後に (3) を仮定する。このとき $V(I)$ は開かつ閉であるから、
$R$ のある冪等元 $e$ に対して $V(I) = D(1 - e)$ である。
補題 \ref{algebra-lemma-disjoint-decomposition} を参照せよ。
イデアル $J = (e)$ は $V(J) = V(I)$ を満たす純イデアルなので、
補題 \ref{algebra-lemma-pure-ideal-determined-by-zero-set} により $I = J$ である。
従って (3) $\Rightarrow$ (2) である。補題
\ref{algebra-lemma-finite-projective} により、(4) は $I$ が純かつ
$R/I$ が有限表示であるという主張と同値である。
さらに、補題 \ref{algebra-lemma-extension} により、$R/I$ が有限表示であることと
$I$ が有限生成であることは同値である。従って (4) は (1) と同値である。
\end{proof}

\noindent
上の結果を用いて、すべての有限平坦加群が有限表示となる環を
特徴づけることができる。

\begin{lemma}
\label{algebra-lemma-finite-flat-module-finitely-presented}
$R$ を環とする。次は同値である。
\begin{enumerate}
\item 閉じており、かつ一般化について閉じている任意の
$Z \subset \Spec(R)$ は開でもある。
\item 任意の有限平坦 $R$-加群は有限局所自由である。
\end{enumerate}
\end{lemma}

\begin{proof}
任意の有限平坦 $R$-加群が有限局所自由ならば、純イデアル $I$ に対する
$R/I$ の台は開である。従って (2) $\Rightarrow$ (1) は補題
\ref{algebra-lemma-pure-ideal-determined-by-zero-set} から従う。
% [P01-E0525] The frozen zero-set uniqueness reference is preserved; the correspondence lemma remains sidecar-only.

\medskip\noindent
逆に、$R$ が (1) を満たすと仮定する。
$M$ を有限平坦 $R$-加群とする。$M$ の台
$Z = \text{Supp}(M)$ は閉である。補題 \ref{algebra-lemma-support-closed} を参照せよ。
一方、$\mathfrak p \subset \mathfrak p'$ ならば、補題
\ref{algebra-lemma-finite-flat-local} により加群 $M_{\mathfrak p'}$ は自由であり、
$M_{\mathfrak p} = M_{\mathfrak p'} \otimes_{R_{\mathfrak p'}} R_{\mathfrak p}$ である。
% [P01-E0526] Japanese supplies the missing sentence-final punctuation after the frozen equality.
従って
$\mathfrak p' \in \text{Supp}(M) \Rightarrow \mathfrak p \in \text{Supp}(M)$
であり、言い換えれば台は一般化について閉じている。
$R$ は (1) を満たすので、$M$ の台は開かつ閉である。
$M$ が $r$ 個の元 $m_1, \ldots, m_r$ で生成されるとする。
加群 $\wedge^i(M)$（$i = 1, \ldots, r$）も有限平坦 $R$-加群である。
実際、$\wedge^i(M)_{\mathfrak p} = \wedge^i(M_{\mathfrak p})$ は
$R_{\mathfrak p}$ 上自由である。また
$\text{Supp}(\wedge^{i + 1}(M)) \subset \text{Supp}(\wedge^i(M))$ である。
従って開かつ閉な部分集合による分解
$$
\Spec(R) = U_0 \amalg U_1 \amalg \ldots \amalg U_r
$$
で、すべての $i = 0, \ldots, r$ に対して $\wedge^i(M)$ の台が
$U_r \cup \ldots \cup U_i$ となるものが存在する。
$\mathfrak p$ を $R$ の素イデアルとし、$\mathfrak p \in U_i$ とする。
等式
$\wedge^i(M) \otimes_R \kappa(\mathfrak p) =
\wedge^i(M \otimes_R \kappa(\mathfrak p))$
に注意する。必要なら $m_1, \ldots, m_r$ の番号を付け替えて、
$m_1, \ldots, m_i$ が $M \otimes_R \kappa(\mathfrak p)$ を
生成すると仮定できる。中山の補題 \ref{algebra-lemma-NAK} により、ある
$f \in R$、$f \not \in \mathfrak p$ に対して全射
$$
R_f^{\oplus i} \longrightarrow M_f, \quad
(a_1, \ldots, a_i) \longmapsto \sum a_im_i
$$
を得る。さらに $D(f) \subset U_i$ と仮定してよい。
これは $\wedge^i(M_f) = \wedge^i(M)_f$ が、その台が
全体 $\Spec(R_f)$ である平坦 $R_f$-加群だということを意味する。
上の結果により、これは一つの元、すなわち
$m_1 \wedge \ldots \wedge m_i$ で生成される。
従って、$V(J) = \Spec(R_f)$ を満たすある純イデアル
$J \subset R_f$ に対して $\wedge^i(M)_f \cong R_f/J$ である。
明らかにこれは $J = (0)$ を意味する。補題
\ref{algebra-lemma-pure-ideal-determined-by-zero-set} を参照せよ。
従って $m_1 \wedge \ldots \wedge m_i$ は $\wedge^i(M_f)$ の基底であり、
上の写像は全射であると同時に単射でもある。
これにより、望みどおり $M$ が有限局所自由であることが証明された。
\end{proof}

\section{有限大域次元の環}
\label{algebra-section-ring-finite-gl-dim}

\noindent
次の補題は、与えられた加群の異なる射影分解を比較するために
しばしば用いられる。

\begin{lemma}[Schanuel の補題]
\label{algebra-lemma-Schanuel}
$R$ を環、$M$ を $R$-加群とする。
$$
0 \to K \xrightarrow{c_1} P_1 \xrightarrow{p_1} M \to 0
\quad\text{and}\quad
0 \to L \xrightarrow{c_2} P_2 \xrightarrow{p_2} M \to 0
$$
が二つの短完全列で、$P_i$ は射影的であると仮定する。
このとき $K \oplus P_2 \cong L \oplus P_1$ である。
より正確には、可換図式
$$
\xymatrix{
0 \ar[r] &
K \oplus P_2 \ar[r]_{(c_1, \text{id})} \ar[d] &
P_1 \oplus P_2 \ar[r]_{(p_1, 0)} \ar[d] &
M \ar[r] \ar@{=}[d] &
0 \\
0 \ar[r] &
P_1 \oplus L \ar[r]^{(\text{id}, c_2)} &
P_1 \oplus P_2 \ar[r]^{(0, p_2)} &
M \ar[r] &
0
}
$$
で、その縦の矢印が同型となるものが存在する。
\end{lemma}

\begin{proof}
短完全列 $0 \to N \to P_1 \oplus P_2 \to M \to 0$ によって
定義される加群 $N$ を考える。ただし最後の写像は二つの写像
$P_i \to M$ の和である。射影 $N \to P_1$ は核 $L$ をもつ全射であり、
$N \to P_2$ は核 $K$ をもつ全射であることが容易に分かる。
$P_i$ は射影的なので $N \cong K \oplus P_2
\cong L \oplus P_1$ である。これで最初の主張が証明された。

\medskip\noindent
第二の主張を証明するために（同時に第一の主張も再証明するために）、
$p_1 = p_2 \circ a$ および $p_2 = p_1 \circ b$ を満たす
$a : P_1 \to P_2$ と $b : P_2 \to P_1$ を選ぶ。
$P_1$ と $P_2$ が射影的なので、これは可能である。このとき可換図式
$$
\xymatrix{
0 \ar[r] &
K \oplus P_2 \ar[r]_{(c_1, \text{id})} &
P_1 \oplus P_2 \ar[r]_{(p_1, 0)} &
M \ar[r] &
0 \\
0 \ar[r] &
N \ar[r] \ar[d] \ar[u] &
P_1 \oplus P_2 \ar[r]_{(p_1, p_2)}
\ar[d]_S \ar[u]^T &
M \ar[r] \ar@{=}[d] \ar@{=}[u] &
0 \\
0 \ar[r] &
P_1 \oplus L \ar[r]^{(\text{id}, c_2)} &
P_1 \oplus P_2 \ar[r]^{(0, p_2)} &
M \ar[r] &
0
}
$$
を得る。ここで $T$ と $S$ は行列
$$
S = \left(
\begin{matrix}
\text{id} & 0 \\
a & \text{id}
\end{matrix}
\right)
\quad\text{and}\quad
T = \left(
\begin{matrix}
\text{id} & b \\
0 & \text{id}
\end{matrix}
\right)
$$
で与えられる。このとき $S$、$T$、写像 $N \to P_1 \oplus L$、
および写像 $N \to K \oplus P_2$ はすべて所望の同型である。
\end{proof}

\begin{definition}
\label{algebra-definition-finite-proj-dim}
$R$ を環、$M$ を $R$-加群とする。射影的 $R$-加群による
有限長の分解をもつとき、$M$ は {\it 有限射影次元} をもつという。
そのような分解の最小の長さを $M$ の {\it 射影次元} と呼ぶ。
\end{definition}

\noindent
$M$ の射影次元が $0$ であることと、$M$ が射影加群であることは
明らかに同値である。次の補題は、射影次元が射影分解の選択に
どの程度依存しないかを説明する。

\begin{lemma}
\label{algebra-lemma-independent-resolution}
$R$ を環とし、$M$ を射影次元 $d$ の $R$-加群とする。
$F_i$ が射影的で、$e \geq d - 1$ である完全列
$F_e \to F_{e-1} \to \ldots \to F_0 \to M \to 0$ があると仮定する。
このとき $F_e \to F_{e-1}$ の核は射影的である
（$e = 0$ の場合には $F_0 \to M$ の核が射影的である）。
\end{lemma}

\begin{proof}
$d$ に関する帰納法で証明する。$d = 0$ ならば $M$ は射影的である。
この場合、分裂 $F_0 = \Ker(F_0 \to M) \oplus M$ があるので、
$\Ker(F_0 \to M)$ は射影的である。$e = 0$ ならばこれで証明は終わる。
$e > 0$ ならば、$M$ を $\Ker(F_0 \to M)$ で置き換えると $e$ が減少する。

\medskip\noindent
次に $d > 0$ と仮定する。$P_i$ が射影的である最小長の有限分解
$0 \to P_d \to P_{d-1} \to \ldots \to P_0 \to M \to 0$ をとる。
Schanuel の補題 \ref{algebra-lemma-Schanuel} により
$P_0 \oplus \Ker(F_0 \to M) \cong F_0 \oplus \Ker(P_0 \to M)$ である。
$d = 1$、$e = 0$ の場合、右辺は射影的な
$F_0 \oplus P_1$ なので、これで証明される。
従って以後 $e > 0$ と仮定してよい。加群
$F_0 \oplus \Ker(P_0 \to M)$ は長さ $d - 1$ の有限射影分解
$$
0 \to P_d \to P_{d-1} \to \ldots \to
P_2 \to P_1 \oplus F_0 \to \Ker(P_0 \to M) \oplus F_0 \to 0
$$
をもつ。長さ $e - 1$ の完全列
$$
F_e \to F_{e-1} \to \ldots \to F_2 \to P_0 \oplus F_1 \to
P_0 \oplus \Ker(F_0 \to M) \to 0
$$
に帰納法を適用すると、$e \geq 2$ ならば
$\Ker(F_e \to F_{e - 1})$ が射影的であると結論する。
または $\Ker(F_1 \oplus P_0 \to F_0 \oplus P_0)$ が射影的である。
これは補題を含意する。
\end{proof}

\begin{lemma}
\label{algebra-lemma-what-kind-of-resolutions}
$R$ を環、$M$ を $R$-加群とし、$d \geq 0$ とする。
次は同値である。
\begin{enumerate}
\item $M$ の射影次元は $\leq d$ である。
\item $P_i$ が射影的である分解
$0 \to P_d \to P_{d - 1} \to \ldots \to P_0 \to M \to 0$ が存在する。
\item $P_i$ が射影的であるある分解
$\ldots \to P_2 \to P_1 \to P_0 \to M \to 0$ について、
$d \geq 2$ ならば $\Ker(P_{d - 1} \to P_{d - 2})$ は射影的であり、
$d = 1$ ならば $\Ker(P_0 \to M)$ は射影的であり、
$d = 0$ ならば $M$ は射影的である。
\item $P_i$ が射影的である任意の分解
$\ldots \to P_2 \to P_1 \to P_0 \to M \to 0$ について、
$d \geq 2$ ならば $\Ker(P_{d - 1} \to P_{d - 2})$ は射影的であり、
$d = 1$ ならば $\Ker(P_0 \to M)$ は射影的であり、
$d = 0$ ならば $M$ は射影的である。
\end{enumerate}
\end{lemma}

\begin{proof}
(1) と (2) の同値性は射影次元の定義である。定義
\ref{algebra-definition-finite-proj-dim} を参照せよ。
補題 \ref{algebra-lemma-independent-resolution} により (2) $\Rightarrow$ (4) である。
(4) $\Rightarrow$ (3) と (3) $\Rightarrow$ (2) は直ちに分かる。
\end{proof}

\begin{lemma}
\label{algebra-lemma-what-kind-of-resolutions-local}
$R$ を局所環、$M$ を $R$-加群とし、$d \geq 0$ とする。
補題 \ref{algebra-lemma-what-kind-of-resolutions} の同値な条件 (1) -- (4) は、
さらに次と同値である。
\begin{enumerate}
\item[(5)] $P_i$ が自由である分解
$0 \to P_d \to P_{d - 1} \to \ldots \to P_0 \to M \to 0$ が存在する。
\end{enumerate}
\end{lemma}

\begin{proof}
補題 \ref{algebra-lemma-what-kind-of-resolutions} と定理
\ref{algebra-theorem-projective-free-over-local-ring} から従う。
\end{proof}

\begin{lemma}
\label{algebra-lemma-what-kind-of-resolutions-Noetherian}
$R$ をネーター環、$M$ を有限 $R$-加群とし、$d \geq 0$ とする。
補題 \ref{algebra-lemma-what-kind-of-resolutions} の同値な条件 (1) -- (4) は、
さらに次と同値である。
\begin{enumerate}
\item[(6)] $P_i$ が有限射影的である分解
$0 \to P_d \to P_{d - 1} \to \ldots \to P_0 \to M \to 0$ が存在する。
\end{enumerate}
\end{lemma}

\begin{proof}
$F_i$ が有限自由である分解
$\ldots \to F_2 \to F_1 \to F_0 \to M \to 0$ を選ぶ
（補題 \ref{algebra-lemma-resolution-by-finite-free}）。補題
\ref{algebra-lemma-what-kind-of-resolutions} により、少なくとも $d \geq 2$ ならば
$P_d = \Ker(F_{d - 1} \to F_{d - 2})$ は射影的である。
さらに $P_d$ は有限 $R$-加群である。実際、$R$ はネーター的であり、
$P_d \subset F_{d - 1}$ の右辺は有限自由である。
従って $0 \to P_d \to F_{d - 1} \to \ldots \to F_1 \to F_0 \to M \to 0$
が所望の分解である。
% [P01-E0527] The frozen proof omits the d = 0 and d = 1 cases; no mathematical repair is inserted here.
\end{proof}

\begin{lemma}
\label{algebra-lemma-what-kind-of-resolutions-Noetherian-local}
$R$ を局所ネーター環、$M$ を有限 $R$-加群とし、$d \geq 0$ とする。
補題 \ref{algebra-lemma-what-kind-of-resolutions} の同値な条件 (1) -- (4)、
補題 \ref{algebra-lemma-what-kind-of-resolutions-local} の条件 (5)、および
補題 \ref{algebra-lemma-what-kind-of-resolutions-Noetherian} の条件 (6) は、
さらに次と同値である。
\begin{enumerate}
\item[(7)] $F_i$ が有限自由である分解
$0 \to F_d \to F_{d - 1} \to \ldots \to F_0 \to M \to 0$ が存在する。
\end{enumerate}
\end{lemma}

\begin{proof}
これは補題 \ref{algebra-lemma-what-kind-of-resolutions}、
\ref{algebra-lemma-what-kind-of-resolutions-local}、
\ref{algebra-lemma-what-kind-of-resolutions-Noetherian} と、局所環上の有限射影加群が
有限自由であることから従う。補題 \ref{algebra-lemma-finite-projective} を参照せよ。
\end{proof}

\begin{lemma}
\label{algebra-lemma-projective-dimension-ext}
$R$ を環、$M$ を $R$-加群とし、$n \geq 0$ とする。次は同値である。
\begin{enumerate}
\item $M$ の射影次元は $\leq n$ である。
\item すべての $R$-加群 $N$ とすべての $i \geq n + 1$ に対して
$\Ext^i_R(M, N) = 0$ である。
\item すべての $R$-加群 $N$ に対して $\Ext^{n + 1}_R(M, N) = 0$ である。
\end{enumerate}
\end{lemma}

\begin{proof}
(1) を仮定する。$M$ の自由分解 $F_\bullet \to M$ を選び、
$d_e : F_e \to F_{e - 1}$ と書く。補題
\ref{algebra-lemma-independent-resolution} により、$e \geq n - 1$ に対して
$P_e = \Ker(d_e)$ は射影的である。これは $e \geq n$ に対して
$F_e \cong P_e \oplus P_{e - 1}$ であることを含意する。
ここで $d_e$ は直和因子 $P_{e - 1}$ を $F_{e - 1}$ 内の
$P_{e - 1}$ へ同型に写す。
% [P01-E0528] The frozen decomposition uses P_{-1} when n = 0; the indexing is preserved without repair.
従って、任意の $R$-加群 $N$ に対して複体
$\Hom_R(F_\bullet, N)$ は次数 $\geq n + 1$ で分裂完全である。
従って (2) が成り立つ。(2) $\Rightarrow$ (3) は自明である。

\medskip\noindent
(3) が成り立つと仮定する。$n = 0$ ならば補題
\ref{algebra-lemma-characterize-projective} により $M$ は射影的なので、
(1) が成り立つ。$n > 0$ ならば、自由 $R$-加群 $F$ と、核 $K$ をもつ
全射 $F \to M$ を選ぶ。補題 \ref{algebra-lemma-reverse-long-exact-seq-ext} と、
部分 (1) によるすべての $i > 0$ についての
$\Ext_R^i(F, N)$ の消滅から、すべての $R$-加群 $N$ に対して
$\Ext_R^n(K, N) = 0$ であることが分かる。
従って帰納法により $K$ の射影次元は $\leq n - 1$ である。
$K$ の任意の有限射影分解の先頭に $F$ を加えれば、$M$ の一つ長い
射影分解が得られるので、$M$ の射影次元は $\leq n$ である。
\end{proof}

\begin{lemma}
\label{algebra-lemma-exact-sequence-projective-dimension}
$R$ を環とし、$0 \to M' \to M \to M'' \to 0$ を
$R$-加群の短完全列とする。
\begin{enumerate}
\item $M$ の射影次元が $\leq n$ で、$M''$ の射影次元が
$\leq n + 1$ ならば、$M'$ の射影次元は $\leq n$ である。
\item $M'$ と $M''$ の射影次元がともに $\leq n$ ならば、
$M$ の射影次元は $\leq n$ である。
\item $M'$ の射影次元が $\leq n$ で、$M$ の射影次元が
$\leq n + 1$ ならば、$M''$ の射影次元は $\leq n + 1$ である。
\end{enumerate}
\end{lemma}

\begin{proof}
補題 \ref{algebra-lemma-projective-dimension-ext} の射影次元の特徴づけと、
補題 \ref{algebra-lemma-reverse-long-exact-seq-ext} の Ext 群の長完全列を組み合わせる。
\end{proof}

\begin{definition}
\label{algebra-definition-finite-gl-dim}
$R$ を環とする。すべての $R$-加群が長さ高々 $n$ の射影的
$R$-加群による分解をもつような整数 $n$ が存在するとき、
環 $R$ は {\it 有限大域次元} をもつという。そのような最小の $n$ を
$R$ の {\it 大域次元} と呼ぶ。
\end{definition}

\noindent
次の補題の証明における議論は Auslander の論文
\cite{Auslander} に見られる。

\begin{lemma}
\label{algebra-lemma-colimit-projective-dimension}
$R$ を環とする。部分加群 $M_e$ が包含によって整列されている
$M = \bigcup_{e \in E} M_e$ があると仮定する。
商 $M_e/\bigcup\nolimits_{e' < e} M_{e'}$ の射影次元が
$\leq n$ であると仮定する。このとき $M$ の射影次元は $\leq n$ である。
\end{lemma}

\begin{proof}
$n$ に関する帰納法で証明する。

\medskip\noindent
基底の場合：$n = 0$ とする。このとき
$P_e = M_e/\bigcup_{e' < e} M_{e'}$ は射影的である。
従って射影 $M_e \to P_e$ の切断 $P_e \to M_e$ を選べる。
誘導写像 $\psi : \bigoplus_{e \in E} P_e \to M$ は同型であると主張する。
実際、$x = \sum x_e \in \bigoplus P_e$ が零でないならば、
$x_{e_{max}}$ が零でないような最大の $e_{max}$ をとる。
このとき $y = \psi(x) = \psi(\sum x_e)$ は零でない。実際、
$y \in M_{e_{max}}$ であり、$P_{e_{max}}$ における像
$x_{e_{max}}$ は零でない。
一方、$y \in M$ とする。このとき、ある $e$ に対して $y \in M_e$ である。
$e$ に関する超限帰納法で $y \in \Im(\psi)$ を示す。
$x_e \in P_e$ を $y$ の像とする。このとき
$y - \psi(x_e) \in \bigcup_{e' < e} M_{e'}$ である。
帰納法の仮定により $y - \psi(x_e) \in \Im(\psi)$ なので、
$y \in \Im(\psi)$ である。従って主張は正しく、$\psi$ は同型である。
射影加群の直和として $M$ は射影的である。補題
\ref{algebra-lemma-direct-sum-projective} を参照せよ。

\medskip\noindent
$n > 0$ の場合、$e \in E$ に対して $F_e$ を $M_e$ の元の集合上の
自由 $R$-加群とする。このとき整列集合 $E$ 上の短完全列系
$$
0 \to K_e \to F_e \to M_e \to 0
$$
がある。遷移写像 $F_{e'} \to F_e$ と $K_{e'} \to K_e$ も
単射であることに注意する。$F = \bigcup F_e$ および
$K = \bigcup K_e$ とおく。このとき
$$
0 \to
K_e/\bigcup\nolimits_{e' < e} K_{e'} \to
F_e/\bigcup\nolimits_{e' < e} F_{e'} \to
M_e/\bigcup\nolimits_{e' < e} M_{e'} \to 0
$$
も $R$-加群の短完全列である。また
$F_e/\bigcup_{e' < e} F_{e'}$ は、
$\bigcup_{e' < e} M_{e'}$ に含まれない $M_e$ の元の集合上の
自由 $R$-加群である。従って補題
\ref{algebra-lemma-exact-sequence-projective-dimension} により
$K_e/\bigcup_{e' < e} K_{e'}$ の射影次元は高々 $n - 1$ である。
帰納法により $K$ の射影次元は高々 $n - 1$ である。
従って $M$ の射影次元は高々 $n$ であり、証明が終わる。
\end{proof}

\begin{lemma}
\label{algebra-lemma-finite-gl-dim}
$R$ を環とする。次は同値である。
\begin{enumerate}
\item $R$ の大域次元は有限で $\leq n$ である。
\item すべての有限 $R$-加群の射影次元は $\leq n$ である。
\item すべての巡回 $R$-加群 $R/I$ の射影次元は $\leq n$ である。
\end{enumerate}
\end{lemma}

\begin{proof}
(1) $\Rightarrow$ (2) と (2) $\Rightarrow$ (3) は明らかである。
(3) を仮定する。$M$ の生成元集合 $E \subset M$ を選び、
$E$ 上の整列順序を選ぶ。$e \in E$ に対し、$e' \leq e$ を満たす
元 $e' \in E$ によって生成される $M$ の部分加群を $M_e$ と書く。
このとき $M = \bigcup_{e \in E} M_e$ である。各 $e \in E$ に対して商
$$
M_e/\bigcup\nolimits_{e' < e} M_{e'}
$$
は零であるか一つの元で生成されるので、(3) によりその射影次元は
$\leq n$ である。補題 \ref{algebra-lemma-colimit-projective-dimension} により、
これは $M$ の射影次元が $\leq n$ であることを意味する。
\end{proof}

\begin{lemma}
\label{algebra-lemma-localize-finite-gl-dim}
$R$ を環、$M$ を $R$-加群とし、$S \subset R$ を乗法的部分集合とする。
\begin{enumerate}
\item $M$ の射影次元が $\leq n$ ならば、$S^{-1}M$ の
$S^{-1}R$ 上の射影次元は $\leq n$ である。
\item $R$ の大域次元が有限で $\leq n$ ならば、
$S^{-1}R$ の大域次元も有限で $\leq n$ である。
\end{enumerate}
\end{lemma}

\begin{proof}
$0 \to P_n \to P_{n - 1} \to \ldots \to P_0 \to M \to 0$ を
射影分解とする。局所化は完全であり（命題
\ref{algebra-proposition-localization-exact}）、各 $S^{-1}P_i$ は射影的
$S^{-1}R$-加群なので（補題 \ref{algebra-lemma-ascend-properties-modules}）、
$0 \to S^{-1}P_n \to \ldots \to S^{-1}P_0 \to S^{-1}M \to 0$ は
$S^{-1}M$ の射影分解である。これで (1) が証明された。
$M'$ を $S^{-1}R$-加群とする。$M' = S^{-1}M'$ に注意すれば、
(2) は (1) から従う。
\end{proof}

\section{正則環と大域次元}
\label{algebra-section-regular-finite-gl-dim}

\noindent
有限大域次元の環についての結果を用いて、正則局所環の別の
特徴づけを与えることができる。

\begin{proposition}
\label{algebra-proposition-regular-finite-gl-dim}
$R$ を次元 $d$ の正則局所環とする。深さ $e$ の任意の有限
$R$-加群 $M$ は有限自由分解
$$
0 \to F_{d-e} \to \ldots \to F_0 \to M \to 0.
$$
をもつ。特に、正則局所環の大域次元は $\leq d$ である。
\end{proposition}

\begin{proof}
第一の部分は補題 \ref{algebra-lemma-regular-mcm-free} と補題
\ref{algebra-lemma-mcm-resolution} から従う。最後の部分はこれと補題
\ref{algebra-lemma-finite-gl-dim} から従う。
\end{proof}

\begin{lemma}
\label{algebra-lemma-finite-gl-dim-primes}
$R$ をネーター環、$n \geq 0$ を整数とする。
$R$ の大域次元が有限で $\leq n$ であることと、$R$ のすべての
極大イデアル $\mathfrak m$ に対して環 $R_{\mathfrak m}$ の
大域次元が $\leq n$ であることは同値である。
\end{lemma}

\begin{proof}
補題 \ref{algebra-lemma-localize-finite-gl-dim} で、$R$ の大域次元が有限で
$n$ ならば、すべての局所化 $R_{\mathfrak m}$ の大域次元が
高々 $n$ であることを見た。
% [P01-E0529] The frozen malformed citation clause is rendered naturally.
逆に、すべての $R_{\mathfrak m}$ の大域次元が $\leq n$ であると仮定する。
$M$ を有限 $R$-加群とする。$F_i$ が有限自由である分解
$0 \to K_n \to F_{n-1} \to \ldots \to F_0 \to M \to 0$
をとる。このとき $K_n$ は有限 $R$-加群である。
% [P01-E0530] The frozen n = 0 negative-index display is preserved without a syzygy convention.
補題 \ref{algebra-lemma-independent-resolution} と仮定により、すべての加群
$K_n \otimes_R R_{\mathfrak m}$ は射影的である。従って補題
\ref{algebra-lemma-finite-projective} により、加群 $K_n$ は有限射影的である。
\end{proof}

\begin{lemma}
\label{algebra-lemma-length-resolution-residue-field}
$R$ を極大イデアル $\mathfrak m$ と剰余体 $\kappa$ をもつ
ネーター局所環とする。このとき $\kappa$ の射影次元は
$\geq \dim_\kappa \mathfrak m / \mathfrak m^2$ である。
\end{lemma}

\begin{proof}
$\mathfrak m / \mathfrak m^2$ における像が基底をなす
$\mathfrak m$ の元 $x_1 , \ldots, x_n$ をとる。
$x_1, \ldots, x_n$ 上の {\it Koszul 複体} を考える。これは写像
$$
0 \to \wedge^n R^n \to \wedge^{n-1} R^n \to \wedge^{n-2} R^n \to
\ldots \to \wedge^i R^n \to \ldots \to R^n \to R
$$
からなる複体であり、その写像は
$$
e_{j_1} \wedge \ldots \wedge e_{j_i}
\longmapsto
\sum_{a = 1}^i (-1)^{a + 1} x_{j_a} e_{j_1} \wedge \ldots
\wedge \hat e_{j_a} \wedge \ldots \wedge e_{j_i}
$$
で与えられる。これが複体 $K_{\bullet}(R, x_{\bullet})$ をなすことは
容易に分かる。補題 \ref{algebra-lemma-NAK} の部分 (8) により、
$K_{\bullet}(R, x_{\bullet})$ の最後の写像の余核は $\kappa$ である。

\medskip\noindent
$\kappa$ の射影次元が有限で $d$ ならば、有限自由 $R$-加群による
長さ $d$ の分解 $F_{\bullet} \to \kappa$ をとれる
（補題 \ref{algebra-lemma-what-kind-of-resolutions-Noetherian-local}）。
補題 \ref{algebra-lemma-add-trivial-complex} により、複体 $F_{\bullet}$ の
すべての写像が $\Im(F_i \to F_{i-1})
\subset \mathfrak m F_{i-1}$ を満たすと仮定してよい。
実際、分解から自明な直和因子を除いても、分解は高々短くなるだけである。
補題 \ref{algebra-lemma-compare-resolutions} により、$\kappa$ 上で恒等写像を誘導する
複体の写像 $\alpha : K_{\bullet}(R, x_{\bullet}) \to F_{\bullet}$ をとれる。
写像 $\alpha_i : \wedge^i R^n = K_i(R, x_{\bullet}) \to F_i$ が、
$\alpha_i \otimes \kappa : \wedge^i \kappa^n \to F_i \otimes \kappa$
を単射にすることを帰納法で証明する。これにより $F_n \not = 0$、
従って所望の $d \geq n$ が分かる。

\medskip\noindent
$R \xrightarrow{\alpha_0} F_0 \to \kappa$ の合成は零でないので、
$i = 0$ の場合は明らかである。
$F_0$ の階数は $1$ でなければならないことに注意する。
さもなければ、余核が $\kappa$ 一個である写像 $F_1 \to F_0$ の像を
$\mathfrak m F_0$ に含めることはできない。

\medskip\noindent
次に $i = 1$ の場合を確認する。この議論は示唆的だと思われるからである。
帰納段階では $i = 0$ の場合から $i = 1$ の場合が従うので、
読者はここを飛ばしてもよい。上で $F_0 = R$ であり、
$F_1 \to F_0 = R$ の像が $\mathfrak m$ であることを見た。可換図式
$$
\begin{matrix}
R^n & = & K_1(R, x_{\bullet}) & \to & K_0(R, x_{\bullet}) & = & R \\
& & \downarrow & & \downarrow & & \downarrow \\
& & F_1 & \to & F_0 & = & R
\end{matrix}
$$
があり、最右の縦の矢印は単元による乗法で与えられる。
従って合成 $R^n \to F_1 \to F_0 = R$ の像も $\mathfrak m$ に等しい。
ゆえに $\dim_\kappa (\mathfrak m / \mathfrak m^2) = n$ なので、写像
$R^n \otimes \kappa \to F_1 \otimes \kappa$ は単射でなければならない。

\medskip\noindent
$i \geq 1$ とし、すべての $j \leq i - 1$ に対して
$\alpha_j \otimes \kappa$ の単射性が証明済みであると仮定する。
可換図式
$$
\begin{matrix}
\wedge^i R^n & = & K_i(R, x_{\bullet}) & \to & K_{i-1}(R, x_{\bullet})
& = & \wedge^{i-1} R^n \\
& & \downarrow & & \downarrow & & \\
& & F_i & \to & F_{i-1} & &
\end{matrix}
$$
を考える。写像 $\wedge^{i-1} \kappa^n \to F_{i-1} \otimes \kappa$ が
単射であることは分かっている。従って
$\wedge^{i-1} \kappa^n \otimes_{\kappa} \mathfrak m/\mathfrak m^2
\to F_{i-1} \otimes \mathfrak m/\mathfrak m^2$ も単射である。
また、複体の選び方により $F_i$ は $\mathfrak mF_{i-1}$ に写り、
Koszul 複体についても同様である。従って可換図式
$$
\begin{matrix}
\wedge^i \kappa^n & \to &
\wedge^{i-1} \kappa^n \otimes \mathfrak m/\mathfrak m^2 \\
\downarrow & & \downarrow \\
F_i \otimes \kappa & \to & F_{i-1} \otimes \mathfrak m/\mathfrak m^2
\end{matrix}
$$
を得る。ここで、写像
$\wedge^i \kappa^n \to
\wedge^{i-1} \kappa^n \otimes \mathfrak m/\mathfrak m^2$ が
単射であることを確認すれば十分であり、これは直接確かめられる。
\end{proof}

\begin{lemma}
\label{algebra-lemma-dim-gl-dim}
$R$ をネーター局所環とする。剰余体 $\kappa$ の $R$ 上の
射影次元が有限で $n$ であると仮定する。このとき $\dim(R) \geq n$ である。
\end{lemma}

\begin{proof}
$F_{\bullet}$ を有限自由 $R$-加群による $\kappa$ の有限分解とする
（補題 \ref{algebra-lemma-what-kind-of-resolutions-Noetherian-local}）。
補題 \ref{algebra-lemma-add-trivial-complex} により、複体 $F_{\bullet}$ の
すべての写像が性質 $\Im(F_i \to F_{i-1})
\subset \mathfrak m F_{i-1}$ をもつと仮定してよい。
% [P01-E0531] The frozen “have to property” typo is rendered by its evident sense.
実際、分解から自明な直和因子を除いても、分解は高々短くなるだけである。
$F_n \not = 0$ かつ $i > n$ に対して $F_i = 0$ とする。
すると $\kappa$ の射影次元は $n$ である。
命題 \ref{algebra-proposition-what-exact} により、複体の選び方から
$I(\varphi_n)$ は $R$ に等しくなりえないので、
$\text{depth}_{I(\varphi_n)}(R) \geq n$ である。
従って補題 \ref{algebra-lemma-bound-depth} により $\dim(R) \geq n$ でもある。
\end{proof}

\begin{proposition}
\label{algebra-proposition-finite-gl-dim-regular}
$(R, \mathfrak m, \kappa)$ をネーター局所環とする。次は同値である。
\begin{enumerate}
\item $\kappa$ の $R$-加群としての射影次元は有限である。
\item $R$ の大域次元は有限である。
\item $R$ は正則局所環である。
\end{enumerate}
さらにこの場合、$R$ の大域次元は
$\dim(R) = \dim_\kappa(\mathfrak m/\mathfrak m^2)$ に等しい。
\end{proposition}

\begin{proof}
命題 \ref{algebra-proposition-regular-finite-gl-dim} により (3) $\Rightarrow$ (2) である。
(2) $\Rightarrow$ (1) は自明である。(1) を仮定する。補題
\ref{algebra-lemma-length-resolution-residue-field} と \ref{algebra-lemma-dim-gl-dim} により
$\dim(R) \geq \dim_\kappa(\mathfrak m /\mathfrak m^2)$ である。
従って $R$ は正則である。定義 \ref{algebra-definition-regular-local} と
その直前の議論を参照せよ。同値な条件 (1) -- (3) が成り立つと仮定する。
命題 \ref{algebra-proposition-regular-finite-gl-dim} により $R$ の大域次元は
高々 $\dim(R)$ であり、補題 \ref{algebra-lemma-length-resolution-residue-field} により
少なくとも $\dim_\kappa(\mathfrak m/\mathfrak m^2)$ である。
従って主張された等式が成り立つ。
\end{proof}

\begin{lemma}
\label{algebra-lemma-localization-of-regular-local-is-regular}
ネーター局所環 $R$ が正則局所環であることと、有限大域次元をもつことは
同値である。この場合、すべての素イデアル $\mathfrak p$ に対して
$R_{\mathfrak p}$ は正則局所環である。
\end{lemma}

\begin{proof}
命題 \ref{algebra-proposition-finite-gl-dim-regular} と
\ref{algebra-proposition-regular-finite-gl-dim} により、ネーター局所環が
正則局所環であることと有限大域次元をもつことは同値である。
さらに、任意の局所化 $R_{\mathfrak p}$ は有限大域次元をもつ
（補題 \ref{algebra-lemma-localize-finite-gl-dim} を参照せよ）ので、正則局所環である。
\end{proof}

\noindent
補題 \ref{algebra-lemma-localization-of-regular-local-is-regular} により、
次の定義は正則局所環についての先の定義と矛盾しないので妥当である。

\begin{definition}
\label{algebra-definition-regular}
ネーター環 $R$ が {\it 正則} であるとは、すべての素イデアルでの局所化
$R_{\mathfrak p}$ が正則局所環であることをいう。
\end{definition}

\noindent
極大イデアルにおける局所環だけが正則であることを要求すれば十分である。
$R$ をネーター的と仮定しても、これは $R$ が有限大域次元をもつことと
同じではないことに注意する。実際、次元が有限でないために有限大域次元を
もたない正則ネーター環の例がある。

\begin{lemma}
\label{algebra-lemma-finite-gl-dim-finite-dim-regular}
$R$ をネーター環とする。次は同値である。
\begin{enumerate}
\item $R$ の大域次元は有限で $n$ である。
\item $R$ は次元 $n$ の正則環である。
\item すべての極大イデアルでの局所化 $R_{\mathfrak m}$ が
次元 $\leq n$ の正則環であり、少なくとも一つの $\mathfrak m$ に対して
等号が成り立つような整数 $n$ が存在する。
\item すべての素イデアルでの局所化 $R_{\mathfrak p}$ が
次元 $\leq n$ の正則環であり、少なくとも一つの $\mathfrak p$ に対して
等号が成り立つような整数 $n$ が存在する。
\end{enumerate}
\end{lemma}

\begin{proof}
上の議論から従う。より正確には、定義 \ref{algebra-definition-regular}、補題
\ref{algebra-lemma-finite-gl-dim-primes}、および命題
\ref{algebra-proposition-finite-gl-dim-regular} を組み合わせればよい。
\end{proof}

\begin{lemma}
\label{algebra-lemma-flat-under-regular}
$R \to S$ を局所ネーター環の局所準同型とする。
$R \to S$ は平坦であり、$S$ は正則であると仮定する。
このとき $R$ は正則である。
\end{lemma}

\begin{proof}
$\mathfrak m \subset R$ を極大イデアルとし、
$\kappa = R/\mathfrak m$ を剰余体とする。$d = \dim S$ とおく。
各 $F_i$ が有限自由 $R$-加群である任意の分解
$F_\bullet \to \kappa$ を選び、
$K_d = \Ker(F_{d - 1} \to F_{d - 2})$ とおく。
% [P01-E0532] The frozen d = 0 negative-index syzygy is preserved without an augmentation convention.
$R \to S$ の平坦性により、複体
$0 \to K_d \otimes_R S \to F_{d - 1} \otimes_R S \to \ldots
\to F_0 \otimes_R S \to \kappa \otimes_R S \to 0$
は依然として完全である。$S$ の大域次元は $d$ なので
（命題 \ref{algebra-proposition-regular-finite-gl-dim} を参照せよ）、
$K_d \otimes_R S$ は有限自由 $S$-加群である
（補題 \ref{algebra-lemma-independent-resolution} も参照せよ）。
% [P01-E0533] The frozen equality claim and cited proposition are preserved; the precise alternative remains sidecar-only.
補題 \ref{algebra-lemma-finite-projective-descends} により、$K_d$ は有限自由
$R$-加群である。従って $\kappa$ は有限射影次元をもち、命題
\ref{algebra-proposition-finite-gl-dim-regular} により $R$ は正則である。
\end{proof}

\section{Auslander--Buchsbaum}
\label{algebra-section-Auslander-Buchsbaum}

\noindent
次の結果は \cite{Auslander-Buchsbaum} に見られる。

\begin{proposition}
\label{algebra-proposition-Auslander-Buchsbaum}
$R$ をネーター局所環とする。$M$ を有限射影次元
$\text{pd}_R(M)$ をもつ非零有限 $R$-加群とする。このとき
$$
\text{depth}(R) = \text{pd}_R(M) + \text{depth}(M)
$$
である。
\end{proposition}

\begin{proof}
$\text{depth}(M)$ に関する帰納法で証明する。最も興味深いのは
$\text{depth}(M) = 0$ の場合である。この場合、
$$
0 \to R^{n_e} \to R^{n_{e-1}} \to \ldots \to R^{n_0} \to M \to 0
$$
を極小有限自由分解とする。従って $e = \text{pd}_R(M)$ である。
補題 \ref{algebra-lemma-add-trivial-complex} により、複体の写像のすべての
行列係数が $R$ の極大イデアルに含まれると仮定してよい。
一方では、命題 \ref{algebra-proposition-what-exact} により
$\text{depth}(R) \geq e$ である。他方では、この長完全列を短完全列
\begin{align*}
0 \to R^{n_e} \to R^{n_{e - 1}} \to K_{e - 2} \to 0,\\
0 \to K_{e - 2} \to R^{n_{e - 2}} \to K_{e - 3} \to 0,\\
\ldots,\\
0 \to K_0 \to R^{n_0} \to M \to 0
\end{align*}
に分解し、補題 \ref{algebra-lemma-depth-in-ses} を用いると、
\begin{align*}
\text{depth}(K_{e - 2}) \geq \text{depth}(R) - 1,\\
\text{depth}(K_{e - 3}) \geq \text{depth}(R) - 2,\\
\ldots,\\
\text{depth}(K_0) \geq \text{depth}(R) - (e - 1),\\
\text{depth}(M) \geq \text{depth}(R) - e
\end{align*}
が分かる。
% [P01-E0534] The frozen short-exact decomposition is undefined for e = 0 or e = 1; no low-index repair is inserted.
$\text{depth}(M) = 0$ なので $\text{depth}(R) \leq e$ と結論する。
これで $\text{depth}(M) = 0$ の場合の証明が終わる。

\medskip\noindent
帰納段階。$\text{depth}(M) > 0$ ならば、$M$ と $R$ の双方上の
非零因子である $x \in \mathfrak m$ を選ぶ。これは可能である。
実際、$\text{pd}_R(M) > 0$ ならば前述の命題
\ref{algebra-proposition-what-exact} により $\text{depth}(R) > 0$ である。
または $\text{pd}_R(M) = 0$ ならば $M$ は有限自由であるから、
$\text{depth}(R) = \text{depth}(M) > 0$ でもある。
従って例えば補題 \ref{algebra-lemma-depth-in-ses} により
$\text{depth}(R \oplus M) > 0$ であり、$R$ と $M$ の双方上の
非零因子である $x \in \mathfrak m$ をとれる。
% [P01-E0535] Both frozen choices of x are retained; the duplicated choice is not silently removed.
上と同様の極小分解
$$
0 \to R^{n_e} \to R^{n_{e-1}} \to \ldots \to R^{n_0} \to M \to 0
$$
をとる。蛇の補題を適用すると、
$$
0 \to (R/xR)^{n_e} \to (R/xR)^{n_{e-1}} \to \ldots \to (R/xR)^{n_0} \to
M/xM \to 0
$$
も極小分解であることが分かる。従って
$\text{pd}_R(M) = \text{pd}_{R/xR}(M/xM)$ である。
補題 \ref{algebra-lemma-depth-drops-by-one} により
$\text{depth}(R/xR) = \text{depth}(R) - 1$ かつ
$\text{depth}(M/xM) = \text{depth}(M) - 1$ である。
ここまでは深さをすべて $R$-加群としての深さとしてきたが、
$\text{depth}_R(M/xM) = \text{depth}_{R/xR}(M/xM)$ であり、
$R/xR$ についても同様であることに注意する。
帰納法の仮定により、$R/xR$ 上の $M/xM$ に対して
Auslander--Buchsbaum の公式が成り立つ。$R/xR$ と $M/xM$ の深さは
ともに一つ減少し、射影次元は変わっていないので、結論が従う。
\end{proof}

\section{準同型と次元}
\label{algebra-section-homomorphism-dimension}

\noindent
この節では、環の間に写像があるときの次元を関係づける
いくつかの容易な結果を集める。

\begin{lemma}
\label{algebra-lemma-dimension-going-up}
$R \to S$ を、上昇性（定義 \ref{algebra-definition-going-up-down}）または
下降性（定義 \ref{algebra-definition-going-up-down}）を満たす環写像とする。
さらに $\Spec(S) \to \Spec(R)$ が全射であると仮定する。
このとき $\dim(R) \leq \dim(S)$ である。
\end{lemma}

\begin{proof}
上昇性を仮定する。$R$ の素イデアルの任意の鎖
$\mathfrak p_0 \subset \mathfrak p_1 \subset \ldots
\subset \mathfrak p_e$ をとる。全射性により、$\mathfrak p_0$ に写る
素イデアル $\mathfrak q_0$ を選べる。上昇性により、これを
$\mathfrak p_i$ の上にある素イデアル $\mathfrak q_i$ の長さ $e$ の
鎖へ延長できる。従って $\dim(S) \geq \dim(R)$ である。
下降性の場合もまったく同様である。純粋に位相的な版については、
Topology の補題 \ref{topology-lemma-dimension-specializations-lift} も参照せよ。
\end{proof}

\begin{lemma}
\label{algebra-lemma-going-up-maximal-on-top}
$R \to S$ を上昇性をもつ環写像とする。定義
\ref{algebra-definition-going-up-down} を参照せよ。
$\mathfrak q \subset S$ が極大イデアルならば、$\mathfrak q$ の $R$ における
逆像も極大イデアルである。
% [P01-E0536] The frozen if-clause fragment is rendered as one natural conditional.
\end{lemma}

\begin{proof}
自明である。
\end{proof}

\begin{lemma}
\label{algebra-lemma-integral-dim-up}
$R \to S$ を、$S$ が $R$ 上整となる環写像とする。
このとき $\dim (R) \geq \dim(S)$ であり、$\Spec(S)$ のすべての閉点は
$\Spec(R)$ の閉点に写る。
\end{lemma}

\begin{proof}
補題 \ref{algebra-lemma-integral-no-inclusion}、
\ref{algebra-lemma-going-up-maximal-on-top}、および定義から直ちに従う。
\end{proof}

\begin{lemma}
\label{algebra-lemma-integral-sub-dim-equal}
$R \subset S$ であり、$S$ が $R$ 上整であるとする。
このとき $\dim(R) = \dim(S)$ である。
\end{lemma}

\begin{proof}
補題 \ref{algebra-lemma-integral-going-up}、
\ref{algebra-lemma-integral-overring-surjective}、
\ref{algebra-lemma-dimension-going-up}、および
\ref{algebra-lemma-integral-dim-up} を組み合わせればよい。
\end{proof}

\begin{definition}
\label{algebra-definition-fibre}
$R \to S$ を環写像とする。$\mathfrak q \subset S$ を、$R$ の
素イデアル $\mathfrak p$ の上にある素イデアルとする。
{\it $\mathfrak q$ におけるファイバーの局所環} とは局所環
$$
S_{\mathfrak q}/\mathfrak pS_{\mathfrak q}
=
(S/\mathfrak pS)_{\mathfrak q}
=
(S \otimes_R \kappa(\mathfrak p))_{\mathfrak q}
$$
のことである。
\end{definition}

\begin{lemma}
\label{algebra-lemma-dimension-base-fibre-total}
$R \to S$ をネーター環の準同型とする。$\mathfrak q \subset S$ を
素イデアル $\mathfrak p$ の上にある素イデアルとする。このとき
$$
\dim(S_{\mathfrak q})
\leq
\dim(R_{\mathfrak p})
+
\dim(S_{\mathfrak q}/\mathfrak pS_{\mathfrak q}).
$$
である。
\end{lemma}

\begin{proof}
命題 \ref{algebra-proposition-dimension} の次元の特徴づけを用いる。
$x_1, \ldots, x_d$ を、$d = \dim(R_{\mathfrak p})$ であり、
$R_{\mathfrak p}$ の定義イデアルを生成する $\mathfrak p$ の元とする。
$y_1, \ldots, y_e$ を、$e = \dim(S_{\mathfrak q}/\mathfrak pS_{\mathfrak q})$
であり、$S_{\mathfrak q}/\mathfrak pS_{\mathfrak q}$ の定義イデアルを
生成する $\mathfrak q$ の元とする。
$S_{\mathfrak q}/(x_1, \ldots, x_d, y_1, \ldots, y_e)$ の
極大イデアルが冪零であることは明らかである。従って
$x_1, \ldots, x_d, y_1, \ldots, y_e$ は $S_{\mathfrak q}$ の
定義イデアルを生成する。
\end{proof}

\begin{lemma}
\label{algebra-lemma-dimension-base-fibre-equals-total}
$R \to S$ をネーター環の準同型とする。$\mathfrak q \subset S$ を
素イデアル $\mathfrak p$ の上にある素イデアルとする。
$R \to S$ が下降性を満たすと仮定する
（例えば $R \to S$ が平坦ならば、補題 \ref{algebra-lemma-flat-going-down} を参照せよ）。
このとき
$$
\dim(S_{\mathfrak q})
=
\dim(R_{\mathfrak p})
+
\dim(S_{\mathfrak q}/\mathfrak pS_{\mathfrak q}).
$$
である。
\end{lemma}

\begin{proof}
補題 \ref{algebra-lemma-dimension-base-fibre-total} により不等式
$\dim(S_{\mathfrak q}) \leq
\dim(R_{\mathfrak p}) + \dim(S_{\mathfrak q}/\mathfrak pS_{\mathfrak q})$
が成り立つ。等号を得るために、
$d = \dim(S_{\mathfrak q}/\mathfrak pS_{\mathfrak q})$ である素イデアルの鎖
$\mathfrak pS \subset \mathfrak q_0 \subset \mathfrak q_1 \subset \ldots
\subset \mathfrak q_d = \mathfrak q$ を選ぶ。一方、
$e = \dim(R_{\mathfrak p})$ である素イデアルの鎖
$\mathfrak p_0 \subset \mathfrak p_1 \subset \ldots \subset \mathfrak p_e
= \mathfrak p$ を選ぶ。下降定理により、$\mathfrak p_{e-1}$ の上にある
$\mathfrak q_{-1} \subset \mathfrak q_0$ を選べる。次に
$\mathfrak p_{e-2}$ の上にある
$\mathfrak q_{-2} \subset \mathfrak q_{-1}$ を選べる。
帰納的に続けると、長さ $e + d$ の鎖
$\mathfrak q_{-e} \subset \ldots \subset \mathfrak q_d$ を得る。
\end{proof}

\begin{lemma}
\label{algebra-lemma-flat-over-regular-with-regular-fibre}
$R \to S$ を局所ネーター環の局所準同型とする。次を仮定する。
\begin{enumerate}
\item $R$ は正則である。
\item $S/\mathfrak m_RS$ は正則である。
\item $R \to S$ は平坦である。
\end{enumerate}
このとき $S$ は正則である。
\end{lemma}

\begin{proof}
補題 \ref{algebra-lemma-dimension-base-fibre-equals-total} により
$\dim(S) = \dim(R) + \dim(S/\mathfrak m_RS)$ である。
$d = \dim(R)$ である生成元 $x_1, \ldots, x_d \in \mathfrak m_R$ を選び、
$e = \dim(S/\mathfrak m_RS)$ であり、$S/\mathfrak m_RS$ の極大イデアルを
生成する $y_1, \ldots, y_e \in \mathfrak m_S$ を選ぶ。
このとき $x_1, \ldots, x_d, y_1, \ldots, y_e$ は $S$ の極大イデアルを
生成する元であり、$e + d = \dim(S)$ である。
\end{proof}

\noindent
次の補題は後で、体上有限型の環が Cohen--Macaulay であることと、
多項式環上準有限平坦であることが同値であると示すために用いられる。
これは補題 \ref{algebra-lemma-CM-over-regular-flat} の部分的な逆である。

\begin{lemma}
\label{algebra-lemma-finite-flat-over-regular-CM}
$R \to S$ をネーター局所環の局所準同型とする。
$R$ は Cohen--Macaulay であると仮定する。
$S$ が $R$ 上有限平坦であるか、または $S$ が $R$ 上平坦で
$\dim(S) \leq \dim(R)$ ならば、$S$ は Cohen--Macaulay であり、
$\dim(R) = \dim(S)$ である。
\end{lemma}

\begin{proof}
$x_1, \ldots, x_d \in \mathfrak m_R$ を、長さ
$d = \dim(R)$ の正則列とする。補題 \ref{algebra-lemma-flat-increases-depth} により、
これは $S$ の正則列に写る。従って $\dim(S) \leq d$ ならば
$S$ は Cohen--Macaulay である。$S$ が $R$ 上有限平坦ならば、補題
\ref{algebra-lemma-integral-sub-dim-equal} によりこれは成り立つ。
第二の場合には仮定したとおりである。
\end{proof}

\section{次元公式}
\label{algebra-section-dimension-formula}

\noindent
鎖状（定義 \ref{algebra-definition-catenary}）および普遍鎖状
（定義 \ref{algebra-definition-universally-catenary}）の定義を思い出そう。

\begin{lemma}
\label{algebra-lemma-dimension-formula}
$R \to S$ を環写像とする。$\mathfrak q$ を、$R$ の素イデアル
$\mathfrak p$ の上にある $S$ の素イデアルとする。次を仮定する。
\begin{enumerate}
\item $R$ はネーター的である。
\item $R \to S$ は有限型である。
\item $R$, $S$ は整域である。
\item $R \subset S$ である。
\end{enumerate}
このとき
$$
\text{height}(\mathfrak q)
\leq
\text{height}(\mathfrak p) + \text{trdeg}_R(S)
- \text{trdeg}_{\kappa(\mathfrak p)} \kappa(\mathfrak q)
$$
であり、$R$ が普遍鎖状ならば等号が成り立つ。
\end{lemma}

\begin{proof}
$R \subset S' \subset S$ とし、$S'$ は $S$ の有限生成
$R$-部分代数であるとする。この場合
$\mathfrak q' = S' \cap \mathfrak q$ とおく。体の塔における超越次数の
加法性（Fields、補題 \ref{fields-lemma-transcendence-degree-tower}）により、
環写像 $R \to S'$ と $S' \to S$ に対する補題から $R \to S$ に対する
補題が従う。従って $R$ 上の $S$ の生成元の個数に関する帰納法を用いて、
$S$ が $R$ 上一つの元で生成される場合に帰着できる。

\medskip\noindent
場合 I：$S = R[x]$ は $R$ 上の多項式代数である。
この場合 $\text{trdeg}_R(S) = 1$ である。また $R \to S$ は平坦なので
$$
\dim(S_{\mathfrak q}) =
\dim(R_{\mathfrak p}) + \dim(S_{\mathfrak q}/\mathfrak pS_{\mathfrak q})
$$
である。補題 \ref{algebra-lemma-dimension-base-fibre-equals-total} を参照せよ。
$\mathfrak r = \mathfrak pS$ とおく。このとき
$\text{trdeg}_{\kappa(\mathfrak p)} \kappa(\mathfrak q) = 1$ であることは
$\mathfrak q = \mathfrak r$ と同値であり、
$\dim(S_{\mathfrak q}/\mathfrak pS_{\mathfrak q}) = 0$ を含意する。
同様に、$\text{trdeg}_{\kappa(\mathfrak p)} \kappa(\mathfrak q) = 0$
であることは真の包含 $\mathfrak r \subset \mathfrak q$ があることと同値で、
$\dim(S_{\mathfrak q}/\mathfrak pS_{\mathfrak q}) = 1$ を含意する。
従って場合 I はすべて等号をもって証明された。

\medskip\noindent
場合 II：$S = R[x]/\mathfrak n$ かつ $\mathfrak n \not = 0$ とする。
この場合 $\text{trdeg}_R(S) = 0$ である。
$\mathfrak q$ に対応する素イデアルを $\mathfrak q' \subset R[x]$ と書く。
% [P01-E0537] The frozen missing-preposition designation is rendered naturally.
このとき
$$
S_{\mathfrak q} = (R[x])_{\mathfrak q'}/\mathfrak n(R[x])_{\mathfrak q'}
$$
である。前の場合により
$\dim((R[x])_{\mathfrak q'}) =
\dim(R_{\mathfrak p}) + 1
- \text{trdeg}_{\kappa(\mathfrak p)} \kappa(\mathfrak q)$ である。
$\mathfrak n \not = 0$ なので、$S_{\mathfrak q}$ の次元は少なくとも一つ減少する。
補題 \ref{algebra-lemma-one-equation} を参照せよ。これで補題の不等式が証明された。
$R$ が普遍鎖状の場合に等号が成り立つことを見るため、
$\mathfrak n \subset R[x]$ は高さ一の素イデアルであることに注意する。
% [P01-E0538] The frozen malformed conditional phrase is rendered naturally.
実際、これは $R$ の分数体 $F$ に対して $F[x]$ の非零素イデアルに対応する。
従って $S_\mathfrak q = R[x]_{\mathfrak q'}/\mathfrak nR[x]_{\mathfrak q'}$
における任意の素イデアルの極大鎖は、$R[x]$ の $(0)$ と
$\mathfrak q'$ の間にある、長さが一つ大きい素イデアルの極大鎖に対応する。
% [P01-E0539] The frozen modifier order is rendered by the unambiguous chain-length comparison.
$R$ が普遍鎖状ならば、これらはすべて $\mathfrak q'$ の高さに等しい
同じ長さをもつ。従って
$\dim(S_\mathfrak q) = \dim(R[x]_{\mathfrak q'}) - 1$ であり、
所望の等号が従う。
\end{proof}

\noindent
次の補題は、生成的有限写像が余次元 $1$ で準有限となる傾向をもつことを述べる。

\begin{lemma}
\label{algebra-lemma-finite-in-codim-1}
$A \to B$ を環写像とする。次を仮定する。
\begin{enumerate}
\item $A \subset B$ は整域の拡大である。
\item 誘導される分数体の拡大は有限である。
\item $A$ はネーター的である。
\item $A \to B$ は有限型である。
\end{enumerate}
$\mathfrak p \subset A$ を高さ $1$ の素イデアルとする。
このとき $\mathfrak p$ の上にある $B$ の素イデアルは有限個以下であり、
それらはすべて高さ $1$ である。
\end{lemma}

\begin{proof}
次元公式（補題 \ref{algebra-lemma-dimension-formula}）により、
$\mathfrak p$ の上にある任意の素イデアル $\mathfrak q$ に対して
$$
\dim(B_{\mathfrak q}) \leq
\dim(A_{\mathfrak p}) - \text{trdeg}_{\kappa(\mathfrak p)} \kappa(\mathfrak q).
$$
である。整域 $B_\mathfrak q$ は少なくとも $2$ 個の素イデアルをもつので、
$\dim(B_{\mathfrak q}) \geq 1$ である。従って
$\dim(B_{\mathfrak q}) = 1$ であり、拡大
$\kappa(\mathfrak p) \subset \kappa(\mathfrak q)$ は代数的である。
ゆえに $\mathfrak q$ はそのファイバー
$\Spec(B \otimes_A \kappa(\mathfrak p))$ の閉点を定める。
補題 \ref{algebra-lemma-finite-residue-extension-closed} を参照せよ。
$B \otimes_A \kappa(\mathfrak p)$ はネーター環なので、ファイバー
$\Spec(B \otimes_A \kappa(\mathfrak p))$ はネーター位相空間である。
補題 \ref{algebra-lemma-Noetherian-topology} を参照せよ。
閉点のみからなるネーター位相空間は有限である。例えば Topology、補題
\ref{topology-lemma-Noetherian} を参照せよ。
\end{proof}

\section{体上有限型代数の次元}
\label{algebra-section-dimension-finite-type-algebras}

\noindent
この節では、体上の多項式環の次元を計算する。また、体上有限型の整域の
次元が、極大イデアルにおけるその局所環の次元であることを証明する。
基礎体上の超越次数との関係は節
\ref{algebra-section-dimension-finite-type-algebras-reprise} で確立する。

\begin{lemma}
\label{algebra-lemma-dim-affine-space}
$\mathfrak m$ を $k[x_1, \ldots, x_n]$ の極大イデアルとする。
イデアル $\mathfrak m$ は $n$ 個の元で生成される。
$k[x_1, \ldots, x_n]_{\mathfrak m}$ の次元は $n$ である。
従って $k[x_1, \ldots, x_n]_{\mathfrak m}$ は次元 $n$ の正則局所環である。
\end{lemma}

\begin{proof}
Hilbert の零点定理（定理 \ref{algebra-theorem-nullstellensatz}）により、剰余体
$\kappa = \kappa(\mathfrak m)$ は $k$ の有限拡大である。
$x_i$ の像を $\alpha_i \in \kappa$ と書く。
% [P01-E0540] The frozen missing-preposition designation is rendered naturally.
$\kappa_i = k(\alpha_1, \ldots, \alpha_i)
\subset \kappa$（$i = 1, \ldots, n$）および $\kappa_0 = k$ とおく。
% [P01-E0541] The frozen malformed definitional construction is rendered naturally.
体論により $\kappa_i = k[\alpha_1, \ldots, \alpha_i]$ であることに注意する。
元 $f_i \in \mathfrak m \cap k[x_1, \ldots, x_i]$ を次のように帰納的に定める。
$P_i(T) \in \kappa_{i-1}[T]$ を、$\kappa_{i-1}$ 上の $\alpha_i $ の
モニック最小多項式とする。$Q_i(T) \in k[x_1, \ldots, x_{i-1}][T]$ を
$P_i(T)$ の同じ次数をもつモニックな持ち上げとし、$f_i = Q_i(x_i)$ とおく。
$d_i = \deg_T(P_i) = \deg_T(Q_i) = \deg_{x_i}(f_i)$ ならば、Fields の補題
\ref{fields-lemma-multiplicativity-degrees} と
\ref{fields-lemma-degree-minimal-polynomial} により
$d_1d_2\ldots d_i = [\kappa_i : k]$ であることに注意する。

\medskip\noindent
すべての $i = 0, 1, \ldots, n$ に対して同型
$$
\psi_i : k[x_1, \ldots, x_i] /(f_1, \ldots, f_i) \cong \kappa_i.
$$
があると主張する。構成により、合成
$k[x_1, \ldots, x_i] \to k[x_1, \ldots, x_n] \to \kappa$ は
$\kappa_i$ の上への全射であり、$f_1, \ldots, f_i$ はその核に属する。
従って全射準同型が得られる。$\psi_i$ が単射であることを帰納法で示す。
$i = 0$ の場合は明らかである。$i$ に対する主張から
$i + 1$ に対する主張を示す。環拡大
$k[x_1, \ldots, x_i]/(f_1, \ldots, f_i) \to
k[x_1, \ldots, x_{i + 1}]/(f_1, \ldots, f_{i + 1})$ は、
体上の $1$ 個の元と一つの既約方程式によって生成される。
初等的な体論により
$k[x_1, \ldots, x_{i + 1}]/(f_1, \ldots, f_{i + 1})$ は体であり、
従って $\psi_i$ は単射である。
% [P01-E0542] The frozen concluding psi_i index is preserved; psi_{i+1} remains sidecar-only.

\medskip\noindent
これは $\mathfrak m = (f_1, \ldots, f_n)$ を含意する。さらに
$$
k[x_1, \ldots, x_n]/(f_1, \ldots, f_i)
\cong
\kappa_i[x_{i + 1}, \ldots, x_n].
$$
とも結論する。従って $(f_1, \ldots, f_i)$ は素イデアルである。ゆえに
$$
(0) \subset (f_1) \subset (f_1, f_2) \subset \ldots \subset
(f_1, \ldots, f_n) = \mathfrak m
$$
は長さ $n$ の素イデアルの鎖である。補題が従う。
\end{proof}

\begin{proposition}
\label{algebra-proposition-finite-gl-dim-polynomial-ring}
体上 $n$ 変数の多項式代数は正則環である。その大域次元は $n$ である。
極大イデアルにおけるすべての局所化は次元 $n$ の正則局所環である。
\end{proposition}

\begin{proof}
補題 \ref{algebra-lemma-dim-affine-space} により、極大イデアルにおける
すべての局所化 $k[x_1, \ldots, x_n]_{\mathfrak m}$ は次元 $n$ の
正則局所環である。従って補題
\ref{algebra-lemma-finite-gl-dim-finite-dim-regular} により結論を得る。
\end{proof}

\begin{lemma}
\label{algebra-lemma-dimension-height-polynomial-ring}
$k$ を体とし、$\mathfrak p \subset \mathfrak q \subset k[x_1, \ldots, x_n]$
を素イデアルの対とする。$\mathfrak p$ と $\mathfrak q$ の間の
任意の素イデアルの極大鎖の長さは
$\text{height}(\mathfrak q) - \text{height}(\mathfrak p)$ である。
\end{lemma}

\begin{proof}
命題 \ref{algebra-proposition-finite-gl-dim-polynomial-ring} により、
$k[x_1, \ldots, x_n]$ の任意の局所環は正則である。
従って補題 \ref{algebra-lemma-regular-ring-CM} により、すべての局所環は
Cohen--Macaulay である。極大イデアルにおける局所環の次元は $n$ なので、
補題 \ref{algebra-lemma-maximal-chain-CM} により、$k[x_1, \ldots, x_n]$ の
素イデアルの任意の極大鎖の長さは $n$ である。
従って、例えば補題 \ref{algebra-lemma-CM-dim-formula} により、$(0)$ と
$\mathfrak p$ の間の素イデアルの任意の極大鎖の長さは
$\text{height}(\mathfrak p)$ である。これらを合わせれば補題の主張を得る。
\end{proof}

\begin{lemma}
\label{algebra-lemma-dimension-spell-it-out}
$k$ を体とし、$S$ を整域である有限型 $k$-代数とする。
このとき $S$ の任意の極大イデアル $\mathfrak m$ に対して
$\dim(S) = \dim(S_{\mathfrak m})$ である。言い換えれば、
素イデアルの任意の極大鎖の長さは $S$ の次元に等しい。
\end{lemma}

\begin{proof}
$S = k[x_1, \ldots, x_n]/\mathfrak p$ と書く。命題
\ref{algebra-proposition-finite-gl-dim-polynomial-ring} と補題
\ref{algebra-lemma-dimension-height-polynomial-ring} により、$S$ における
素イデアルのすべての極大鎖（必ず極大イデアルで終わる）の長さは
$n - \text{height}(\mathfrak p)$ である。従ってこの数が $S$ の次元であり、
$S$ の任意の極大イデアル $\mathfrak m$ に対する
$S_{\mathfrak m}$ の次元でもある。
\end{proof}

\noindent
位相空間 $X$ の点 $x$ における次元 $\dim_x(X)$ を Topology の定義
\ref{topology-definition-Krull} で定めたことを思い出そう。

\begin{lemma}
\label{algebra-lemma-dimension-at-a-point-finite-type-over-field}
$k$ を体、$S$ を有限型 $k$-代数とする。
$X = \Spec(S)$ とおく。$\mathfrak p \subset S$ を素イデアル、
$x \in X$ を対応する点とする。次の数は等しい。
\begin{enumerate}
\item $\dim_x(X)$。
\item $x$ を通る $X$ の既約成分 $Z$ 全体についての $\max \dim(Z)$。
\item $\mathfrak p \subset \mathfrak m$ を満たす極大イデアル
$\mathfrak m$ 全体についての $\min \dim(S_{\mathfrak m})$。
\end{enumerate}
\end{lemma}

\begin{proof}
$X = \bigcup_{i \in I} Z_i$ を $X$ の既約成分への分解とする。
それらは有限個である（補題 \ref{algebra-lemma-obvious-Noetherian} と
\ref{algebra-lemma-Noetherian-topology} を参照せよ）。
$I' = \{i \mid x \in Z_i\}$ とし、$T = \bigcup_{i \not \in I'} Z_i$ とおく。
このとき $U = X \setminus T$ は点 $x$ を含む $X$ の開部分集合である。
数 (2) は $\max_{i \in I'} \dim(Z_i)$ である。
$x \in W$ を満たす任意の開集合 $W \subset U$ に対して、$W$ の既約成分は
$i \in I'$ に対する既約集合 $W_i = Z_i \cap W$ であり、$x$ は
それらすべてに含まれる。各 $W_i$（$i \in I'$）は閉点を含むことに注意する。
実際 $X$ は Jacobson である。節 \ref{algebra-section-ring-jacobson} を参照せよ。
$W_i \subset Z_i$ なので $\dim(W_i) \leq \dim(Z_i)$ である。
閉点の存在と補題 \ref{algebra-lemma-dimension-spell-it-out} により、開集合 $W_i$ には
長さ $\dim(Z_i)$ の既約閉部分集合の鎖がある。
従って任意の $i \in I'$ に対して $\dim(W_i) = \dim(Z_i)$ である。
ゆえに $\dim(W)$ は数 (2) に等しい。これで (1) $ = $ (2) が証明された。

\medskip\noindent
$\mathfrak p$ を含む任意の極大イデアル
$\mathfrak m \supset \mathfrak p$ をとり、$x_0 \in X$ を対応する点とする。
まず $x_0$ は、$i \in I'$ に対するすべての既約成分 $Z_i$ に含まれる。
$\mathfrak q_i$ を既約成分 $Z_i$ に対応する $S$ の極小素イデアルとする。
$x_0 \in Z_i$ を満たす各 $i$ に対して
（これは $\mathfrak m \supset \mathfrak q_i$ と同値である）、全射
$$
S_{\mathfrak m} \longrightarrow
S_\mathfrak m/\mathfrak q_i S_\mathfrak m =(S/\mathfrak q_i)_{\mathfrak m}
$$
がある。さらに、このように得られる素イデアル
$\mathfrak q_i S_\mathfrak m$ はネーター局所環 $S_{\mathfrak m}$ の
極小素イデアルを尽くす。補題
\ref{algebra-lemma-irreducible-components-containing-x} を参照せよ。
補題 \ref{algebra-lemma-dimension-spell-it-out} を用いると、$S_{\mathfrak m}$ の
次元は、$x_0$ を通る $Z_i$ の次元の最大値であると結論する。
補題の証明を終えるには、
$x_0 \in Z_i \Rightarrow i \in I'$ となるように $x_0$ を選べることを
示せば十分である。$S$ は Jacobson なので（上で見たとおり）、
$V(\mathfrak p) \setminus T$（$T$ は上のもの）が空でないことを示せばよい。
これは点 $x$、すなわち $\mathfrak p$ を含むので明らかである。
\end{proof}

\begin{lemma}
\label{algebra-lemma-dimension-closed-point-finite-type-field}
$k$ を体、$S$ を有限型 $k$-代数とし、$X = \Spec(S)$ とおく。
$\mathfrak m \subset S$ を極大イデアル、$x \in X$ を付随する閉点とする。
このとき $\dim_x(X) = \dim(S_{\mathfrak m})$ である。
\end{lemma}

\begin{proof}
補題 \ref{algebra-lemma-dimension-at-a-point-finite-type-over-field} の特別な場合である。
\end{proof}

\begin{lemma}
\label{algebra-lemma-disjoint-decomposition-CM-algebra}
$k$ を体、$S$ を有限型 $k$ 代数とする。
% [P01-E0543] The frozen missing k-algebra hyphen has no Japanese grammatical analogue.
$S$ は Cohen--Macaulay であると仮定する。このとき
$\Spec(S) = \coprod T_d$ は開かつ閉な部分集合 $T_d$ の有限非交和であり、
$T_d$ は次元 $d$ の等次元である（Topology、定義
\ref{topology-definition-equidimensional} を参照せよ）。同値な言い方をすれば、
$S$ は環 $S_d$（$d = 0, \ldots, \dim(S)$）の積であり、$S_d$ の
すべての極大イデアル $\mathfrak m$ の高さは $d$ である。
\end{lemma}

\begin{proof}
二つの主張の同値性は補題 \ref{algebra-lemma-disjoint-implies-product} から従う。
$\mathfrak m \subset S$ を極大イデアルとする。補題
\ref{algebra-lemma-maximal-chain-CM} により、$S_{\mathfrak m}$ の素イデアルの
任意の極大鎖は $\dim(S_{\mathfrak m})$ に等しい同じ長さをもつ。
従って、補題 \ref{algebra-lemma-dimension-spell-it-out} により、$\mathfrak m$ に
対応する点を通る既約成分はすべて $\dim(S_{\mathfrak m})$ に等しい
次元をもつ。
% [P01-E0544] The frozen singular/plural mismatch is rendered with the intended plural subject.
$\Spec(S)$ は Jacobson 位相空間なので、その任意の二つの既約成分の
交わりが空でなければ閉点を含む。補題
\ref{algebra-lemma-finite-type-field-Jacobson} と \ref{algebra-lemma-jacobson} を参照せよ。
従って、交わる任意の二つの既約成分は同じ次元をもつことが示された。
補題はこれと、$\Spec(S)$ が有限個の既約成分をもつことから容易に従う
（補題 \ref{algebra-lemma-obvious-Noetherian} と
\ref{algebra-lemma-Noetherian-topology} を参照せよ）。
\end{proof}

\section{ネーター正規化}
\label{algebra-section-Noether-normalization}

\noindent
この節では、ネーター正規化補題のいくつかの変種を証明する。
用いる鍵となる材料は、次の二つの補題に含まれている。

\begin{lemma}
\label{algebra-lemma-helper}
$n \in \mathbf{N}$ とする。
$N$ を多重指数 $\nu = (\nu_1, \ldots, \nu_n)$ からなる有限非空集合とする。
$e = (e_1, \ldots, e_n)$ が与えられたとき、
$e \cdot \nu = \sum e_i\nu_i$ とおく。このとき
$e_1 \gg e_2 \gg \ldots \gg e_{n-1} \gg e_n$ ならば、次が成り立つ：
$\nu, \nu' \in N$ ならば
$$
(e \cdot \nu = e \cdot \nu')
\Leftrightarrow
(\nu = \nu')
$$
である。
\end{lemma}

\begin{proof}
$N = \{\nu_j\}$、$\nu_j = (\nu_{j1}, \ldots, \nu_{jn})$ と書く。
$A_i = \max_j \nu_{ji} - \min_j \nu_{ji}$ とする。
% [P01-E0545] The frozen unqualified index range, including its undefined e_0 case, is preserved.
各 $i$ について
$e_{i - 1} > A_ie_i + A_{i + 1}e_{i + 1} + \ldots + A_ne_n$
であれば、補題は成り立つ。実際、
$e \cdot (\nu - \nu') = 0$ と仮定する。このとき $n \ge 2$ ならば、
$$
e_1(\nu_1 - \nu'_1) = \sum\nolimits_{i = 2}^n e_i(\nu'_i - \nu_i).
$$
$(\nu_1 - \nu'_1) \ge 0$ と仮定してよい。
$(\nu_1 - \nu'_1) > 0$ ならば、
$$
e_1(\nu_1 - \nu'_1) \ge e_1 >
A_2e_2 + \ldots + A_ne_n \ge
\sum\nolimits_{i = 2}^n e_i|\nu'_i - \nu_i| \ge
\sum\nolimits_{i = 2}^n e_i(\nu'_i - \nu_i).
$$
これは矛盾なので $\nu'_1 = \nu_1$ である。帰納法により、
$2 \le i \le n$ に対して $\nu'_i = \nu_i$ である。
\end{proof}

\begin{lemma}
\label{algebra-lemma-helper-polynomial}
$R$ を環とする。$g \in R[x_1, \ldots, x_n]$ を非定数元、すなわち
$g \not \in R$ とする。
$e_1 \gg e_2 \gg \ldots \gg e_{n-1} \gg e_n = 1$ に対し、多項式
$$
g(x_1 + x_n^{e_1}, x_2 + x_n^{e_2}, \ldots, x_{n - 1} + x_n^{e_{n - 1}}, x_n)
=
ax_n^d + \text{lower order terms in }x_n
$$
と書ける。ここで $d > 0$ であり、$a \in R$ は $g$ の非零係数の
一つである。
\end{lemma}

\begin{proof}
% [P01-E0546] The frozen awkward coefficient predicate is rendered naturally.
$g = \sum_{\nu \in N} a_\nu x^\nu$ と書く。ただし $a_\nu \in R$ は
零でない。ここで $N$ は補題 \ref{algebra-lemma-helper} におけるような
多重指数の有限集合であり、
$x^\nu = x_1^{\nu_1} \ldots x_n^{\nu_n}$ である。
$$
(x_1 + x_n^{e_1})^{\nu_1} \ldots (x_{n-1} + x_n^{e_{n-1}})^{\nu_{n-1}}
x_n^{\nu_n}
\quad\text{is}\quad
x_n^{e_1\nu_1 + \ldots + e_{n-1}\nu_{n-1} + \nu_n}.
$$
における最高次項に注意する。従って、補題 \ref{algebra-lemma-helper} から
主張が従う。実際、この補題により、$g$ の非零項 $a_\nu x^\nu$ で
$g(x_1 + x_n^{e_1}, x_2 + x_n^{e_2}, \ldots, x_{n - 1} + x_n^{e_{n - 1}},
x_n)$ の最高次項を生じるものはちょうど一つである。すなわち、
$e \cdot \nu$ が最大となる一意な $\nu \in N$ に対して
$a = a_\nu$ である。
\end{proof}

\begin{lemma}
\label{algebra-lemma-one-relation}
$k$ を体とする。ある真のイデアル $I$ に対し
$S = k[x_1, \ldots, x_n]/I$ とする。$I \not = 0$ ならば、
$S$ が $k[y_1, \ldots, y_{n-1}]$ 上有限となるような
$y_1, \ldots, y_{n-1} \in k[x_1, \ldots, x_n]$ が存在する。
さらに、$y_i$ は $x_1, \ldots, x_n$ によって生成される
$k[x_1, \ldots, x_n]$ の $\mathbf{Z}$-部分代数に属するように
選べる。
\end{lemma}

\begin{proof}
$f \in I$、$f\not = 0$ を選ぶ。$S$ は
$k[x_1, \ldots, x_n]/(f)$ の商なので、この環について補題を示せば
十分である。適当な整数 $e_i$ に対し
$y_i = x_i - x_n^{e_i}$、$i = 1, \ldots, n-1$ ととる。
これはいつうまくいくだろうか。$\overline{x_n} \in
k[x_1, \ldots, x_n]/(f)$ が環 $k[y_1, \ldots, y_{n-1}]$ 上整である
ことを示せば十分である。この環上の $\overline{x_n}$ の方程式は
$$
f(y_1 + x_n^{e_1}, \ldots, y_{n-1} + x_n^{e_{n-1}}, x_n) = 0.
$$
である。
% [P01-E0547] The frozen subject-verb disagreement is rendered naturally.
従って、上の方程式を $x_n$ の多項式と見たときの最高次係数が
$k$ の非零元となるような整数 $e_i$ が存在することを示せばよい。
これは、例えば $e_1 \gg e_2 \gg \ldots \gg e_{n-1}$ と選ぶことで
達成できる。補題 \ref{algebra-lemma-helper-polynomial} を参照せよ。
\end{proof}

\begin{lemma}
\label{algebra-lemma-Noether-normalization}
\begin{slogan}
ネーター正規化
\end{slogan}
$k$ を体とする。あるイデアル $I$ に対して
$S = k[x_1, \ldots, x_n]/I$ とする。$I \neq (1)$ ならば、
$r\geq 0$ と $y_1, \ldots, y_r \in k[x_1, \ldots, x_n]$ が存在して、
次を満たす。(a) 始域を $y_1, \ldots, y_r$ 上の多項式環とした写像
$k[y_1, \ldots, y_r] \to S$ は単射であり、(b) 写像
$k[y_1, \ldots, y_r] \to S$ は有限である。このとき整数 $r$ は
$S$ の次元である。さらに、$y_i$ は $x_1, \ldots, x_n$ により
生成される $k[x_1, \ldots, x_n]$ の $\mathbf{Z}$-部分代数に属する
ように選べる。
\end{lemma}

\begin{proof}
$n$ に関する帰納法による。$n = 0$ の場合は自明である。
$I = 0$ ならば、$r = n$ および $y_i = x_i$ ととる。
$I \not = 0$ ならば、補題 \ref{algebra-lemma-one-relation} のように
$y_1, \ldots, y_{n-1}$ を選ぶ。$S' \subset S$ を $y_i$ の像で
生成される部分環とする。帰納法により、(a), (b) が
$k[z_1, \ldots, z_r] \to S'$ に対して成り立つような $r$ と
$z_1, \ldots, z_r \in k[y_1, \ldots, y_{n-1}]$ を選べる。
$S' \to S$ は単射かつ有限なので、(a), (b) は
$k[z_1, \ldots, z_r] \to S$ に対しても成り立つ。
% [P01-E0548] The frozen “last assertion” attribution remains literal; the corrected target is sidecar-only.
最後の主張は補題 \ref{algebra-lemma-integral-sub-dim-equal} から従う。
\end{proof}

\begin{lemma}
\label{algebra-lemma-Noether-normalization-at-point}
$k$ を体とする。
% [P01-E0549] The frozen missing k-algebra hyphen has no Japanese grammatical analogue.
$S$ を有限型 $k$ 代数とし、$X = \Spec(S)$ とおく。
$\mathfrak q$ を $S$ の素イデアルとし、$x \in X$ を対応する点とする。
$\dim(S_g) = \dim_x(X) =: d$ であり、かつ有限単射
$k[y_1, \ldots, y_d] \to S_g$ が存在するような
$g \in S$、$g \not \in \mathfrak q$ が存在する。
\end{lemma}

\begin{proof}
定義により、$\dim_x(X)$ は $g \in S$、$g \not \in \mathfrak q$ に
対する $S_g$ の次元の最小値であり、すなわち最小値は達成される。
従って補題は補題 \ref{algebra-lemma-Noether-normalization} から従う。
\end{proof}

\begin{lemma}
\label{algebra-lemma-refined-Noether-normalization}
$k$ を体とする。$\mathfrak q \subset k[x_1, \ldots, x_n]$ を
素イデアルとする。$r = \text{trdeg}_k\ \kappa(\mathfrak q)$ とおく。
このとき、$\varphi^{-1}(\mathfrak q) = (y_{r + 1}, \ldots, y_n)$ を
満たす有限環準同型
$\varphi : k[y_1, \ldots, y_n] \to k[x_1, \ldots, x_n]$ が存在する。
\end{lemma}

\begin{proof}
$n$ に関する帰納法による。$n = 0$ の場合は明らかである。
$n > 0$ と仮定する。$r = n$ ならば $\mathfrak q = (0)$ であり、
結果は明らかである。非零元 $f \in \mathfrak q$ を選ぶ。
もちろん $f$ は非定数である。次の形の自己同型
$$
k[x_1, \ldots, x_n] \longrightarrow k[x_1, \ldots, x_n],
\quad
x_n \mapsto x_n,
\quad
x_i \mapsto x_i + x_n^{e_i}\ (i < n)
$$
を施した後、補題 \ref{algebra-lemma-helper-polynomial} により、
% [P01-E0550] The frozen full coefficient ring is preserved; the intended smaller coefficient ring remains sidecar-only.
$f$ は $k[x_1, \ldots, x_n]$ 上 $x_n$ に関してモニックであると
仮定してよい。従って環準同型
$$
k[y_1, \ldots, y_n] \longrightarrow k[x_1, \ldots, x_n],
\quad
y_n \mapsto f,
\quad
y_i \mapsto x_i\ (i < n)
$$
は有限である。さらに構成により
$y_n \in \mathfrak q \cap k[y_1, \ldots, y_n]$ である。従って
$\mathfrak q \cap k[y_1, \ldots, y_n] = \mathfrak pk[y_1, \ldots, y_n] + (y_n)$
である。ここで $\mathfrak p \subset k[y_1, \ldots, y_{n - 1}]$ は
素イデアルである。
$\kappa(\mathfrak p) \subset \kappa(\mathfrak q)$ は有限であるから、
$r = \text{trdeg}_k\ \kappa(\mathfrak p)$ である。
組 $(k[y_1, \ldots, y_{n - 1}], \mathfrak p)$ に帰納法の仮定を適用すると、
$\mathfrak p \cap k[z_1, \ldots, z_{n - 1}] =
(z_{r + 1}, \ldots, z_{n - 1})$ を満たす有限環準同型
$k[z_1, \ldots, z_{n - 1}] \to k[y_1, \ldots, y_{n - 1}]$ を得る。
環準同型
$k[z_1, \ldots, z_{n - 1}] \to k[y_1, \ldots, y_{n - 1}]$ を、
$z_n$ を $y_n$ に送ることにより環準同型
$k[z_1, \ldots, z_n] \to k[y_1, \ldots, y_n]$ に拡張する。
環準同型の合成
$$
k[z_1, \ldots, z_n] \to k[y_1, \ldots, y_n] \to k[x_1, \ldots, x_n]
$$
が問題を解く。
\end{proof}

\begin{lemma}
\label{algebra-lemma-Noether-normalization-over-a-domain}
$R \to S$ を単射な有限型環準同型とする。$R$ は整域であると仮定する。
このとき整数 $d$ と、単射による分解
$$
R \to R[y_1, \ldots, y_d] \to S' \to S
$$
が存在し、$S'$ は $R[y_1, \ldots, y_d]$ 上有限であり、ある非零元
$f \in R$ に対して $S'_f \cong S_f$ である。
\end{lemma}

\begin{proof}
$S$ を $R$ 上生成する $x_1, \ldots, x_n \in S$ を選ぶ。
$K$ を $R$ の分数体とし、$S_K = S \otimes_R K$ とする。
補題 \ref{algebra-lemma-Noether-normalization} により、
$K[y_1, \ldots, y_d] \to S_K$ が有限単射となるような
$y_1, \ldots, y_d \in S$ を見つけられる。
% [P01-E0551] The frozen i/j index switch is preserved.
$y_j$ を $x_1, \ldots, x_n$ で生成される $\mathbf{Z}$-代数内に
選べるので、$y_i \in S$ であることに注意する。
有限環準同型は整であるから（補題 \ref{algebra-lemma-finite-is-integral} を
参照せよ）、$P_i(x_i) = 0$ が $S_K$ で成り立つようなモニック多項式
$P_i \in K[y_1, \ldots, y_d][T]$ を見つけられる。
すべての $i$ に対して $fP_i \in R[y_1, \ldots, y_d][T]$ となる
非零元 $f \in R$ をとる。このとき $fP_i(x_i)$ は $S_K$ で零に
写る。従って、$f$ を $R$ の別の非零元で置き換えることにより、
$fP_i(x_i)$ が $S$ で零であるとも仮定してよい。
$x_i' = fx_i$ とおき、$S' \subset S$ を
$y_1, \ldots, y_d$ および $x'_1, \ldots, x'_n$ で生成される
$R$-部分代数とする。$Q_i(x_i') = 0$ であり、ここで
$Q_i = f^{\deg_T(P_i)}P_i(T/f)$ は、$f$ の選び方により
$R[y_1, \ldots, y_d]$ に係数をもつ $T$ のモニック多項式なので、
$x'_i$ は $R[y_1, \ldots, y_d]$ 上整である。
従って、補題 \ref{algebra-lemma-characterize-finite-in-terms-of-integral} により
$R[y_1, \ldots, y_d] \subset S'$ は有限である。
$S' \subset S$ なので $S'_f \subset S_f$ である
（局所化は完全である）。一方、$S'_f$ の元 $x_i = x'_i/f$ は
$S_f$ を $R_f$ 上生成するので、$S'_f \to S_f$ は全射である。
従って $S'_f \cong S_f$ となり、証明が完了する。
\end{proof}

\section{体上有限型代数の次元（再論）}
\label{algebra-section-dimension-finite-type-algebras-reprise}

\noindent
この節は節 \ref{algebra-section-dimension-finite-type-algebras} の続きである。
ここでは、体上の有限型整域について、その次元と基礎体上の超越次数との
関係を確立する。

\begin{lemma}
\label{algebra-lemma-dimension-prime-polynomial-ring}
$k$ を体とする。
% [P01-E0552] The frozen missing k-algebra hyphen has no Japanese grammatical analogue.
$S$ を整域である有限型 $k$ 代数とする。
$K$ を $S$ の分数体とする。
$r = \text{trdeg}(K/k)$ を $k$ 上の $K$ の超越次数とする。
このとき $\dim(S) = r$ である。さらに、$S$ を任意の極大イデアルで
局所化して得られる局所環の次元は $r$ である。
\end{lemma}

\begin{proof}
$S = k[x_1, \ldots, x_n]/\mathfrak p$ と書ける。
補題 \ref{algebra-lemma-dimension-height-polynomial-ring} により、$S$ を
極大イデアルで局所化して得られる局所環はすべて同じ次元をもつ。
補題 \ref{algebra-lemma-Noether-normalization} を適用すると、
$d = \dim(S)$ である有限単射な環準同型
$$
k[y_1, \ldots, y_d] \to S
$$
を得る。明らかに $k(y_1, \ldots, y_d) \subset K$ は有限拡大であり、
主張が従う。
\end{proof}

\begin{lemma}
\label{algebra-lemma-tr-deg-specialization}
$k$ を体とする。$S$ を有限型 $k$-代数とする。
$\mathfrak q \subset \mathfrak q' \subset S$ を相異なる素イデアルとする。
このとき
$\text{trdeg}_k\ \kappa(\mathfrak q') < \text{trdeg}_k\ \kappa(\mathfrak q)$
である。
\end{lemma}

\begin{proof}
補題 \ref{algebra-lemma-dimension-prime-polynomial-ring} により、
$\dim V(\mathfrak q) = \text{trdeg}_k\ \kappa(\mathfrak q)$ であり、
$\mathfrak q'$ についても同様である。従って、真の包含
$V(\mathfrak q') \subset V(\mathfrak q)$ が次元の真の不等式を
含意することから結果が従う。
\end{proof}

\noindent
次の補題は補題 \ref{algebra-lemma-dimension-closed-point-finite-type-field} を
一般化する。

\begin{lemma}
\label{algebra-lemma-dimension-at-a-point-finite-type-field}
$k$ を体とする。
% [P01-E0553] The frozen missing k-algebra hyphen has no Japanese grammatical analogue.
$S$ を有限型 $k$ 代数とする。$X = \Spec(S)$ とする。
$\mathfrak p \subset S$ を素イデアルとし、$x \in X$ を対応する点とする。
このとき
$$
\dim_x(X) = \dim(S_{\mathfrak p}) + \text{trdeg}_k\ \kappa(\mathfrak p).
$$
が成り立つ。
\end{lemma}

\begin{proof}
補題 \ref{algebra-lemma-dimension-prime-polynomial-ring} により、
$r = \text{trdeg}_k\ \kappa(\mathfrak p)$ は $V(\mathfrak p)$ の
次元に等しい。$S$ において $\mathfrak p$ から始まる素イデアルの
任意の極大鎖
$\mathfrak p \subset \mathfrak p_1 \subset \ldots \subset \mathfrak p_r$
を選ぶ。補題 \ref{algebra-lemma-dimension-spell-it-out} により、これは長さ
$r$ をもつ。$\mathfrak q_j$、$j \in J$ を、$\mathfrak p$ に含まれる
$S$ の極小素イデアルとする。これらは規則
$\mathfrak q_j \mapsto \mathfrak q_jS_{\mathfrak p}$ により
$S_{\mathfrak p}$ の極小素イデアルと $1-1$ に対応する。
補題 \ref{algebra-lemma-dimension-at-a-point-finite-type-over-field} により、
$\dim_x(X)$ は環 $S/\mathfrak q_j$ の次元の最大値に等しい。
各 $j$ に対して素イデアルの極大鎖
$\mathfrak q_j \subset \mathfrak p'_1 \subset \ldots \subset \mathfrak p'_{s(j)}
= \mathfrak p$
を選ぶ。このとき $\dim(S_{\mathfrak p}) = \max_{j \in J} s(j)$ である。
さて、各鎖
$$
\mathfrak q_j \subset \mathfrak p'_1 \subset \ldots \subset
\mathfrak p'_{s(j)} = \mathfrak p \subset
\mathfrak p_1 \subset \ldots \subset \mathfrak p_r
$$
は $S/\mathfrak q_j$ における極大鎖であり、上で述べたことから
$\dim_x(X) = \max_{j \in J} r + s(j)$ である。従って補題が従う。
\end{proof}

\noindent
次の補題は、ある有限型スペクトルの別の有限型スペクトルにおける
余次元が高さの差であることを述べる。

\begin{lemma}
\label{algebra-lemma-codimension}
$k$ を体とする。
% [P01-E0554] The frozen missing k-algebra hyphen has no Japanese grammatical analogue.
$S' \to S$ を有限型 $k$ 代数の全射とする。
$\mathfrak p \subset S$ を素イデアルとし、$\mathfrak p'$ を
$S'$ の対応する素イデアルとする。$X = \Spec(S)$、resp.\
$X' = \Spec(S')$ とし、$x \in X$、resp. $x'\in X'$ を
$\mathfrak p$、resp.\ $\mathfrak p'$ に対応する点とする。このとき
$$
\dim_{x'} X' - \dim_x X =
\text{height}(\mathfrak p') - \text{height}(\mathfrak p).
$$
である。
\end{lemma}

\begin{proof}
補題 \ref{algebra-lemma-dimension-at-a-point-finite-type-field} から直ちに従う。
\end{proof}

\begin{lemma}
\label{algebra-lemma-dimension-preserved-field-extension}
$k$ を体とする。$S$ を有限型 $k$-代数とし、$K/k$ を体拡大とする。
このとき $\dim(S) = \dim(K \otimes_k S)$ である。
\end{lemma}

\begin{proof}
補題 \ref{algebra-lemma-Noether-normalization} により、
$d = \dim(S)$ を満たす有限単射
$k[y_1, \ldots, y_d] \to S$ が存在する。$K$ は $k$ 上平坦なので、
有限単射 $K[y_1, \ldots, y_d] \to K \otimes_k S$ も得られる。
結果は補題 \ref{algebra-lemma-integral-sub-dim-equal} から従う。
\end{proof}

\begin{lemma}
\label{algebra-lemma-dimension-at-a-point-preserved-field-extension}
$k$ を体とする。$S$ を有限型 $k$-代数とし、$X = \Spec(S)$ とおく。
$K/k$ を体拡大とし、$S_K = K \otimes_k S$ および
$X_K = \Spec(S_K)$ とおく。$\mathfrak q \subset S$ を
$x \in X$ に対応する素イデアルとし、$\mathfrak q_K \subset S_K$ を
$\mathfrak q$ の上にあり $x_K \in X_K$ に対応する素イデアルとする。
このとき $\dim_x X = \dim_{x_K} X_K$ である。
\end{lemma}

\begin{proof}
表示 $S = k[x_1, \ldots, x_n]/I$ を選ぶ。これにより表示
$K \otimes_k S = K[x_1, \ldots, x_n]/(K \otimes_k I)$ を得る。
$\mathfrak q_K' \subset K[x_1, \ldots, x_n]$、resp.\
$\mathfrak q' \subset k[x_1, \ldots, x_n]$ を対応する素イデアルとする。
次のネーター局所環の可換図式を考える。
$$
\xymatrix{
K[x_1, \ldots, x_n]_{\mathfrak q_K'} \ar[r] &
(K \otimes_k S)_{\mathfrak q_K} \\
k[x_1, \ldots, x_n]_{\mathfrak q'} \ar[r] \ar[u] &
S_{\mathfrak q} \ar[u]
}
$$
二つの垂直矢印は、それぞれ平坦環準同型 $S \to S_K$ および
$k[x_1, \ldots, x_n] \to K[x_1, \ldots, x_n]$ の局所化なので平坦である。
さらに、垂直矢印は同じファイバー環をもつ。従って補題
\ref{algebra-lemma-dimension-base-fibre-equals-total} から
$\text{height}(\mathfrak q') - \text{height}(\mathfrak q)
= \text{height}(\mathfrak q_K') - \text{height}(\mathfrak q_K)$
を得る。
% [P01-E0555] The frozen missing designation preposition is rendered naturally.
$x' \in X' = \Spec(k[x_1, \ldots, x_n])$ および
$x'_K \in X'_K = \Spec(K[x_1, \ldots, x_n])$ を、それぞれ
$\mathfrak q'$ および $\mathfrak q_K'$ に対応する点とする。
補題 \ref{algebra-lemma-codimension} と上で示したことにより、
\begin{eqnarray*}
n - \dim_x X & = & \dim_{x'} X' - \dim_x X \\
& = & \text{height}(\mathfrak q') - \text{height}(\mathfrak q) \\
& = & \text{height}(\mathfrak q_K') - \text{height}(\mathfrak q_K) \\
& = & \dim_{x'_K} X'_K - \dim_{x_K} X_K \\
& = & n - \dim_{x_K} X_K
\end{eqnarray*}
であり、補題が従う。
\end{proof}

\begin{lemma}
\label{algebra-lemma-inequalities-under-field-extension}
$k$ を体とする。$S$ を有限型 $k$-代数とし、$K/k$ を体拡大とする。
$S_K = K \otimes_k S$ とおく。$\mathfrak q \subset S$ を素イデアルとし、
$\mathfrak q_K \subset S_K$ を $\mathfrak q$ の上にある
素イデアルとする。このとき
$$
\dim (S_K \otimes_S \kappa(\mathfrak q))_{\mathfrak q_K} =
\dim (S_K)_{\mathfrak q_K} - \dim S_\mathfrak q =
\text{trdeg}_k \kappa(\mathfrak q) - \text{trdeg}_K \kappa(\mathfrak q_K)
$$
である。さらに、与えられた $\mathfrak q$ に対し、上の数が零に
なるように $\mathfrak q_K$ を常に選べる。
\end{lemma}

\begin{proof}
$S_\mathfrak q \to (S_K)_{\mathfrak q_K}$ は、特殊ファイバー
$(S_K \otimes_S \kappa(\mathfrak q))_{\mathfrak q_K}$ をもつ
ネーター局所環の平坦局所準同型である。
% [P01-E0556] The frozen verbless inference is rendered naturally.
従って最初の等式は補題 \ref{algebra-lemma-dimension-base-fibre-equals-total}
から従う。二番目の等式は、補題
\ref{algebra-lemma-dimension-at-a-point-preserved-field-extension} の記法で
$\dim_x X = \dim_{x_K} X_K$ が成り立つこと、ならびに補題
\ref{algebra-lemma-dimension-at-a-point-finite-type-field} により
$\dim_x X = \dim S_\mathfrak q + \text{trdeg}_k \kappa(\mathfrak q)$
が成り立ち、$\dim_{x_K} X_K$ についても同様であることから従う。
$\mathfrak q_K$ を $\mathfrak q S_K$ 上極小に選べば、
ファイバー環の次元は零となる。
\end{proof}







\section{体上の次数付き代数の次元}
\label{algebra-section-dimension-graded}

\noindent
次が基本的な結果である。

\begin{lemma}
\label{algebra-lemma-dimension-graded}
$k$ を体とする。$S$ を、$k$ 上有限個の次数 $1$ の元で生成される
次数付き $k$-代数とする。$S_0 = k$ と仮定する。すべての
$d \gg 0$ に対し $\dim(S_d) = P(d)$ となる多項式
$P(T) \in \mathbf{Q}[T]$ をとる。命題
\ref{algebra-proposition-graded-hilbert-polynomial} を参照せよ。このとき
\begin{enumerate}
\item 無関係イデアル $S_{+}$ は極大イデアル $\mathfrak m$ である。
\item $S$ の任意の極小素イデアルは斉次イデアルであり、
$S_{+} = \mathfrak m$ に含まれる。
\item $\dim(S) = \deg(P) + 1 = \dim_x\Spec(S)$ である
（$\deg(0) = -1$ という規約を用いる）。ここで $x$ は極大イデアル
$S_{+} = \mathfrak m$ に対応する点である。
\item 局所環 $R = S_{\mathfrak m}$ の Hilbert 関数は $S$ の
Hilbert 関数に等しい。
\end{enumerate}
\end{lemma}

\begin{proof}
最初の主張は明らかである。二番目は補題
\ref{algebra-lemma-graded-ring-minimal-prime} から従う。(2) により、すべての
既約成分は $x$ を通る。従って、補題
\ref{algebra-lemma-dimension-at-a-point-finite-type-over-field} により
$\dim(S) = \dim_x\Spec(S) = \dim(S_\mathfrak m)$ である。
$\mathfrak m^d/\mathfrak m^{d + 1} \cong
\mathfrak m^dS_\mathfrak m/\mathfrak m^{d + 1}S_\mathfrak m$
なので、局所環 $S_\mathfrak m$ の Hilbert 関数は $S$ の Hilbert
関数に等しく、これが (4) である。命題
\ref{algebra-proposition-dimension} により (3) の最後の等式を得る。
\end{proof}







\section{一般平坦性}
\label{algebra-section-generic-flatness}

\noindent
基本的には、これは整域上の有限型代数が、その整域の一つの元を
逆にすると平坦になることを述べている。この結果には、強さが増す
順にいくつかの版がある。

\begin{lemma}
\label{algebra-lemma-generic-flatness-Noetherian}
$R \to S$ を環準同型とし、$M$ を $S$-加群とする。次を仮定する。
\begin{enumerate}
\item $R$ はネーターである。
\item $R$ は整域である。
\item $R \to S$ は有限型である。
\item $M$ は有限型 $S$-加群である。
\end{enumerate}
このとき、$M_f$ が自由 $R_f$-加群となるような非零元 $f \in R$ が
存在する。
\end{lemma}

\begin{proof}
$K$ を $R$ の分数体とし、$S_K = K \otimes_R S$ とおく。これは
$K$ 上有限型の代数である。$d = \dim(S_K)$ に関する帰納法で示す
（これは例えばネーター正規化により有限である。節
\ref{algebra-section-Noether-normalization} を参照せよ）。$d \geq 0$ を固定する。
$\dim(S_K) < d$ であるすべての場合に補題が成り立つと仮定する。

\medskip\noindent
補題のような $R \to S$ と $M$ が与えられ、
$\dim(S_K) = d$ であるとする。補題 \ref{algebra-lemma-filter-Noetherian-module}
により、$M_i/M_{i - 1}$ がある $S$ の素イデアル $\mathfrak q$ に対し
$S/\mathfrak q$ と同型となるようなフィルトレーション
$0 \subset M_1 \subset M_2 \subset \ldots \subset M_n = M$
が存在する。$\dim((S/\mathfrak q)_K) \leq \dim(S_K)$ に注意する。
また、自由加群の拡大は自由であることにも注意する（基本事項
\ref{algebra-item-extension-free} を参照せよ）。従って $M = S$ であり、
$S$ は $R$ 上有限型の整域であると仮定してよい。

\medskip\noindent
$R \to S$ が非自明な核をもつなら、その核の非零元 $f \in R$ をとる。
この場合 $S_f = 0$ なので補題は成り立つ。（これは実際には
$d = -1$ の場合であり、帰納法の出発点である。）従って、
$R \to S$ はネーター整域の有限型拡大であると仮定してよい。

\medskip\noindent
補題 \ref{algebra-lemma-Noether-normalization-over-a-domain} を適用し、
$R$ を $R_f$（ここで $f$ はその補題の元）で置き換えると、分解
$$
R \subset R[y_1, \ldots, y_d] \subset S
$$
を得る。第二の拡大は有限である。$S$ の分数体の
$R[y_1, \ldots, y_d]$ の分数体上の基底をなす
$z_1, \ldots, z_r \in S$ を選ぶ。これにより短完全列
$$
0 \to
R[y_1, \ldots, y_d]^{\oplus r} \xrightarrow{(z_1, \ldots, z_r)}
S \to N \to 0
$$
を得る。構成により、$N$ は有限 $R[y_1, \ldots, y_d]$-加群であり、
その台は $\Spec(R[y_1, \ldots, y_d])$ の一般点 $(0)$ を含まない。
補題 \ref{algebra-lemma-support-closed} により、$N$ を零化する非零元
$g \in R[y_1, \ldots, y_d]$ が存在する。この $g$ により $N$ を
$S' = R[y_1, \ldots, y_d]/(g)$ 上の有限加群と見なせる。
$\dim(S'_K) < d$ なので、帰納法により $N_f$ が自由 $R_f$-加群と
なるような非零元 $f \in R$ が存在する。
% [P01-E0557] The frozen fused inference is rendered grammatically.
$(R[y_1, \ldots, y_d])_f \cong R_f[y_1, \ldots, y_d]$ も自由なので、
自由加群の拡大は自由であるという既述の事実から結論が従う。
\end{proof}

\begin{lemma}
\label{algebra-lemma-generic-flatness-finitely-presented}
\begin{slogan}
一般自由性。
\end{slogan}
$R \to S$ を環準同型とし、$M$ を $S$-加群とする。次を仮定する。
\begin{enumerate}
\item $R$ は整域である。
\item $R \to S$ は有限表示である。
\item $M$ は有限表示 $S$-加群である。
\end{enumerate}
このとき、$M_f$ が自由 $R_f$-加群となるような非零元 $f \in R$ が
存在する。
\end{lemma}

\begin{proof}
$S = R[x_1, \ldots, x_n]/(g_1, \ldots, g_m)$ と書く。
% [P01-E0558] The frozen missing designation preposition is rendered naturally.
$g \in R[x_1, \ldots, x_n]$ に対し、その $S$ における像を
$\overline{g}$ と書く。ある $n_i \in S^{\oplus t}$ に対し
$M = S^{\oplus t}/\sum Sn_i$ と書ける。ある
$g_{ij} \in R[x_1, \ldots, x_n]$ に対し
$n_i = (\overline{g}_{i1}, \ldots, \overline{g}_{it})$ と書く。
$R_0 \subset R$ を、すべての元
$g_i, g_{ij} \in R[x_1, \ldots, x_n]$ のすべての係数で生成される
部分環とする。
$S_0 = R_0[x_1, \ldots, x_n]/(g_1, \ldots, g_m)$ と定め、
$M_0 = S_0^{\oplus t}/\sum S_0n_i$ と定める。このとき $R_0$ は
$\mathbf{Z}$ 上有限型の整域であり、従ってネーターである
（補題 \ref{algebra-lemma-Noetherian-permanence} を参照せよ）。
さらに、単射 $R_0 \to R$ を通じて
$S \cong R \otimes_{R_0} S_0$ および
$M \cong R \otimes_{R_0} M_0$ が成り立つ。
補題 \ref{algebra-lemma-generic-flatness-Noetherian} を適用すると、
$(M_0)_f$ が自由 $(R_0)_f$-加群となるような非零元
$f \in R_0$ を得る。従って
$M_f = R_f \otimes_{(R_0)_f} (M_0)_f$ は自由 $R_f$-加群である。
\end{proof}

\begin{lemma}
\label{algebra-lemma-generic-flatness}
$R \to S$ を環準同型とし、$M$ を $S$-加群とする。次を仮定する。
\begin{enumerate}
\item $R$ は整域である。
\item $R \to S$ は有限型である。
\item $M$ は有限型 $S$-加群である。
\end{enumerate}
このとき、次を満たす非零元 $f \in R$ が存在する。
\begin{enumerate}
\item[(a)] $M_f$ および $S_f$ は自由 $R_f$-加群である。
\item[(b)] $S_f$ は有限表示 $R_f$-代数であり、$M_f$ は
有限表示 $S_f$-加群である。
\end{enumerate}
\end{lemma}

\begin{proof}
まず $S = R[x_1, \ldots, x_n]$ の場合に補題を証明し、次に一般の
結果を導く。

\medskip\noindent
$S = R[x_1, \ldots, x_n]$ と仮定する。$M$ を生成する元
$m_1, \ldots, m_t$ を選ぶ。これにより短完全列
$$
0 \to N \to S^{\oplus t} \xrightarrow{(m_1, \ldots, m_t)} M \to 0.
$$
を得る。
% [P01-E0559] The frozen missing designation prepositions are rendered naturally.
$K$ を $R$ の分数体とし、
$S_K = K \otimes_R S = K[x_1, \ldots, x_n]$ とおく。同様に
$N_K = K \otimes_R N$、$M_K = K \otimes_R M$ とおく。
$R \to K$ は平坦なので、$K$ とテンソルをとった後もこの列は完全である。
$S_K = K[x_1, \ldots, x_n]$ はネーター環なので（補題
\ref{algebra-lemma-Noetherian-permanence} を参照せよ）、これを生成する有限個の
元 $n'_1, \ldots, n'_s \in N_K$ を見つけられる。ある
$a_{ij} \in K$ に対して $n'_i = \sum a_{ij}n_j$ となるような
$n_1, \ldots, n_r \in N$ を選ぶ。
$$
M' = S^{\oplus t}/\sum\nolimits_{i = 1, \ldots, r} Sn_i
$$
とおく。構成により $M'$ は有限表示 $S$-加群であり、
$M'_K \cong M_K$ を誘導する全射 $M' \to M$ が存在する。
$R \to S$ と $M'$ に補題
\ref{algebra-lemma-generic-flatness-finitely-presented} を適用すると、
$M'_f$ が自由 $R_f$-加群となるような $f \in R$ を得る。
従って $M'_f \to M_f$ は整域 $R_f$ 上の加群の全射であり、その始域は
自由加群で、$K$ とテンソルをとると同型になる。
$M'_f$ は分数体 $K$ をもつ整域 $R_f$ 上自由であるから
$M'_f \subset M'_K$ であり、従ってこの写像は単射でもある。
ゆえに $M'_f \to M_f$ は同型であり、結果が証明された。

\medskip\noindent
一般の場合、全射 $R[x_1, \ldots, x_n] \to S$ を選ぶ。$S$ と $M$ の
双方を $R[x_1, \ldots, x_n]$ 上の有限加群と考える。上で証明した
特別な場合により、$S_f$ と $M_f$ が自由 $R_f$-加群であり、かつ
有限表示 $R_f[x_1, \ldots, x_n]$-加群となるような非零元
$f \in R$ が存在する。これは明らかに、$S_f$ が有限表示
$R_f$-代数であり、$M_f$ が有限表示 $S_f$-加群であることを含意する。
\end{proof}

\noindent
$R \to S$ を環準同型とし、$M$ を $S$-加群とする。元 $f \in R$ に
関する次の条件を考える。
\begin{equation}
\label{algebra-equation-flat-and-finitely-presented}
\left\{
\begin{matrix}
S_f & \text{is of finite presentation over }R_f\\
M_f & \text{is of finite presentation as }S_f\text{-module}\\
S_f, M_f & \text{are free as }R_f\text{-modules}
\end{matrix}
\right.
\end{equation}
次で定まる開部分集合を定義する。
\begin{equation}
\label{algebra-equation-good-locus}
U(R \to S, M)
=
\bigcup\nolimits_{f \in
R\text{ with }(\ref{algebra-equation-flat-and-finitely-presented})}
D(f)
\end{equation}
これは $\Spec(R)$ の開部分集合である。

\begin{lemma}
\label{algebra-lemma-generic-flatness-locus-extension}
$R \to S$ を環準同型とする。$0 \to M_1 \to M_2 \to M_3 \to 0$ を
$S$-加群の短完全列とする。このとき
$$
U(R \to S, M_1) \cap U(R \to S, M_3) \subset U(R \to S, M_2).
$$
である。
\end{lemma}

\begin{proof}
$u \in U(R \to S, M_1) \cap U(R \to S, M_3)$ とする。
$u \in D(f_1)$、$u \in D(f_3)$ であり、かつ
(\ref{algebra-equation-flat-and-finitely-presented}) が $f_1$ と $M_1$、
および $f_3$ と $M_3$ に対して成り立つような
$f_1, f_3 \in R$ を選ぶ。$f = f_1f_3$ とおく。このとき
$u \in D(f)$ であり、(\ref{algebra-equation-flat-and-finitely-presented}) は
$f$ と $M_1$、$M_3$ の双方に対して成り立つ。自由加群の拡大は
自由であり、有限表示加群の拡大は有限表示である
（補題 \ref{algebra-lemma-extension}）。従って
(\ref{algebra-equation-flat-and-finitely-presented}) は $f$ と $M_2$ に
対して成り立つ。ゆえに $u \in U(R \to S, M_2)$ であり、主張を得る。
\end{proof}

\begin{lemma}
\label{algebra-lemma-generic-flatness-locus-localize}
$R \to S$ を環準同型とし、$M$ を $S$-加群とし、$f \in R$ とする。
同一視 $\Spec(R_f) = D(f)$ のもとで
$U(R_f \to S_f, M_f) = D(f) \cap U(R \to S, M)$ が成り立つ。
\end{lemma}

\begin{proof}
$u \in U(R_f \to S_f, M_f)$ と仮定する。このとき
$u \in D(g)$ であり、組 $((R_f)_g \to (S_f)_g, (M_f)_g)$ に対して
(\ref{algebra-equation-flat-and-finitely-presented}) が成り立つような
$g \in R_f$ が存在する。ある $a \in R$ に対して $g = a/f^n$ と
書き、$h = af$ とおく。このとき $R_h = (R_f)_g$、
$S_h = (S_f)_g$、$M_h = (M_f)_g$ である。さらに $u \in D(h)$ なので、
$u \in U(R \to S, M)$ である。
逆に $u \in D(f) \cap U(R \to S, M)$ と仮定する。このとき
$u \in D(g)$ であり、組 $(R_g \to S_g, M_g)$ に対して
(\ref{algebra-equation-flat-and-finitely-presented}) が成り立つような
$g \in R$ が存在する。すると
(\ref{algebra-equation-flat-and-finitely-presented}) は明らかに組
$(R_{fg} \to S_{fg}, M_{fg}) = ((R_f)_g \to (S_f)_g, (M_f)_g)$ に
対しても成り立つ。従って $u \in U(R_f \to S_f, M_f)$ となり、
主張を得る。
\end{proof}

\begin{lemma}
\label{algebra-lemma-generic-flatness-locus-reduce}
$R \to S$ を環準同型とし、$M$ を $S$-加群とする。
$U \subset \Spec(R)$ を稠密開集合とする。各 $i \in I$ に対し
$U(R_{f_i} \to S_{f_i}, M_{f_i})$ が $D(f_i)$ で稠密であるような
開集合による被覆 $U = \bigcup_{i \in I} D(f_i)$ が存在すると仮定する。
このとき $U(R \to S, M)$ は $\Spec(R)$ で稠密である。
\end{lemma}

\begin{proof}
補題 \ref{algebra-lemma-generic-flatness-locus-localize} によれば、これは純粋に
位相的な主張である。実際、その補題により、各 $i \in I$ に対して
$U(R \to S, M) \cap D(f_i)$ は $D(f_i)$ で稠密である。
Topology、補題 \ref{topology-lemma-nowhere-dense-local} により、
$U(R \to S, M) \cap U$ は $U$ で稠密である。$U$ は
$\Spec(R)$ で稠密なので、$U(R \to S, M)$ は $\Spec(R)$ で稠密である。
\end{proof}

\begin{lemma}
\label{algebra-lemma-generic-flatness-reduced}
$R \to S$ を環準同型とし、$M$ を $S$-加群とする。次を仮定する。
\begin{enumerate}
\item $R \to S$ は有限型である。
\item $M$ は有限 $S$-加群である。
\item $R$ は被約である。
\end{enumerate}
このとき、次を満たす部分集合 $U \subset \Spec(R)$ が存在する。
\begin{enumerate}
\item $U$ は $\Spec(R)$ の開かつ稠密な部分集合である。
\item すべての $u \in U$ に対し、$u \in D(f) \subset U$ であり、
さらに次を満たす $f \in R$ が存在する。
\begin{enumerate}
\item $M_f$ および $S_f$ は $R_f$ 上自由である。
\item $S_f$ は有限表示 $R_f$-代数である。
\item $M_f$ は有限表示 $S_f$-加群である。
\end{enumerate}
\end{enumerate}
\end{lemma}

\begin{proof}
この補題は、式 (\ref{algebra-equation-good-locus}) の開集合
$U(R \to S, M)$ が $\Spec(R)$ で稠密であるという主張と同値である。
まず $S = R[x_1, \ldots, x_n]$ の場合に補題を証明し、次に一般の
結果を導く。

\medskip\noindent
$S = R[x_1, \ldots, x_n]$ であり、$M$ が $S$ 上の任意の有限加群で
ある場合の証明。この場合 $S_f = R_f[x_1, \ldots, x_n]$ は
$R_f$ 上自由かつ有限表示なので、$S$ に関する条件を考える必要はなく、
$M$ に関する条件だけを考えればよい。$n$ に関する帰納法を用いる。

\medskip\noindent
あるイデアル $J_i \subset S$ に対して
$M_i/M_{i - 1} \cong S/J_i$ となる有限フィルトレーション
$$
0 \subset M_1 \subset M_2 \subset \ldots \subset M_t = M
$$
が存在する。補題 \ref{algebra-lemma-trivial-filter-finite-module} を参照せよ。
有限個の稠密開集合の共通部分は稠密開集合なので、補題
\ref{algebra-lemma-generic-flatness-locus-extension} により、
% [P01-E0560] The frozen R/J_i module is preserved; S/J_i remains sidecar-only.
各加群 $R/J_i$ について補題を証明すれば十分である。
従って、$S = R[x_1, \ldots, x_n]$ のあるイデアル $J$ に対し
$M = S/J$ であると仮定してよい。

\medskip\noindent
$I \subset R$ を $J$ の元の係数で生成されるイデアルとする。
$U_1 = \Spec(R) \setminus V(I)$ とし、
$$
U_2 = \Spec(R) \setminus \overline{U_1}.
$$
とおく。このとき $U = U_1 \cup U_2$ が $\Spec(R)$ で稠密であることは
明らかである。(a) $D(f) \subset U_1$ または (b)
$D(f) \subset U_2$ を満たす元 $f \in R$ をとる。このような任意の
$f$ に対して、組 $(R_f \to R_f[x_1, \ldots, x_n], M_f)$ について
補題が成り立つならば、補題 \ref{algebra-lemma-generic-flatness-locus-reduce}
により $U(R \to S, M)$ は $\Spec(R)$ で稠密である。
従って (a) $I = R$ または (b) $V(I) = \Spec(R)$ のいずれかを
仮定してよい。

\medskip\noindent
場合 (b) では、$R$ が被約なので実際 $I = 0$ である。
従って $J = 0$、$M = S$ であり、この場合補題は成り立つ。

\medskip\noindent
場合 (a) ではもう少し議論が必要である。$J$ の元の係数全体は
イデアルをなすので（詳細は省略する）、$I$ の各元は実際に
$J$ のある元の単項式の係数である。従って、
$$
g = \sum\nolimits_{K \in E} a_K x^K \in J
$$
となる元が存在する。ここで $E$ は多重指数
$K = (k_1, \ldots, k_n)$ の有限集合であり、少なくとも一つの係数
$a_{K_0}$ は $R$ の単元である。場合 (a) では $1 \in I$ なので、
実際、係数の一つが $1$ であるものを見つけられる。
$m = \#\{K \in E \mid a_K \text{ is not a unit}\}$ とおく。
$0 \leq m \leq \# E - 1$ に注意する。$m$ に関する帰納法で示す。

\medskip\noindent
$m = 0$ の場合。この場合、$g$ のすべての係数 $a_K$、$K \in E$ は
単元であり、$E \not = \emptyset$ である。$E = \{K_0\}$ が一元集合で
$K_0 = (0, \ldots, 0)$ ならば、$g$ は単元であり $J = S$ なので、
結果は確かに成り立つ。（特に $n = 0$ のときこれが起こり、
$n$ に関する帰納法の基底を与える。）
$E = \{(0, \ldots, 0)\}$ でなければ、少なくとも一つの $K$ は
$(0, \ldots, 0)$ と等しくない。すなわち $g \not \in R$ である。
ここでネーター正規化の通常の技巧を用いる。すなわち、
$$
G(y_1, \ldots, y_n) =
g(y_1 + y_n^{e_1}, y_2 + y_n^{e_2}, \ldots, y_{n - 1} + y_n^{e_{n - 1}}, y_n)
$$
を $0 \ll e_{n -1} \ll e_{n - 2} \ll \ldots \ll e_1$ とともに考える。
補題 \ref{algebra-lemma-helper-polynomial} により、
$G(y_1, \ldots, y_n)$ は $y_n$ の多項式として
$$
a_K y_n^{k_n + \sum_{i = 1, \ldots, n - 1} e_i k_i} +
\text{lower order terms in }y_n
$$
という形である。$a_K$ は単元なので、
$M = R[x_1, \ldots, x_n]/J$ は $R[y_1, \ldots, y_{n - 1}]$ 上有限である。
従って
$U(R \to R[x_1, \ldots, x_n], M) = U(R \to R[y_1, \ldots, y_{n - 1}], M)$
であり、$n$ に関する帰納法で主張を得る。

\medskip\noindent
$m > 0$ の場合。$a_K$ が単元でないような多重指数 $K \in E$ を選ぶ。
前と同様に $U_1 = \Spec(R_{a_K}) = \Spec(R) \setminus V(a_K)$ とおき、
$$
U_2 = \Spec(R) \setminus \overline{U_1}.
$$
とおく。このとき $U = U_1 \cup U_2$ が $\Spec(R)$ で稠密であることは
明らかである。(a) $D(f) \subset U_1$ または (b)
$D(f) \subset U_2$ を満たす元 $f \in R$ をとる。このような任意の
$f$ に対して、組 $(R_f \to R_f[x_1, \ldots, x_n], M_f)$ について
補題が成り立つならば、補題 \ref{algebra-lemma-generic-flatness-locus-reduce}
により $U(R \to S, M)$ は $\Spec(R)$ で稠密である。
従って (a) $a_KR = R$ または (b) $V(a_K) = \Spec(R)$ のいずれかを
仮定してよい。場合 (a) では $a_K$ が単元になったため、数 $m$ は
減少する。場合 (b) では $R$ が被約なので $a_K = 0$ と分かる。
従って集合 $E$ が小さくなり、数 $m$ も減少する。どちらの場合も
$m$ に関する帰納法で主張を得る。

\medskip\noindent
これで $S = R[x_1, \ldots, x_n]$ の場合に補題を証明した。
補題の条件を満たす任意の組 $(R \to S, M)$ を考える。
全射 $R[x_1, \ldots, x_n] \to S$ を選ぶ。
(\ref{algebra-equation-good-locus}) で導入した記法のもとで、
$$
U(R \to S, M) =
U(R \to R[x_1, \ldots, x_n], S)
\cap
U(R \to R[x_1, \ldots, x_n], M)
$$
が成り立つ。
% [P01-E0561] The frozen malformed inference is rendered naturally.
右辺の二つの開集合が稠密であることをいま証明したので、左辺の
開集合も稠密である。
\end{proof}







\section{Krull--Akizuki 周辺}
\label{algebra-section-krull-akizuki}

\noindent
Krull--Akizuki の一つの応用は、離散付値環が豊富に存在することを
示すことである。より一般に、この節ではネーター局所環を支配する
離散付値環の構成法を示す。

\medskip\noindent
まず、極大イデアルをブローアップすることにより、ネーター局所整域が
$1$ 次元ネーター局所整域によって支配されることを示す。

\begin{lemma}
\label{algebra-lemma-dominate-by-dimension-1}
$R$ を分数体 $K$ をもつネーター局所整域とする。
$R$ は体でないと仮定する。このとき $R \subset R' \subset K$ で
次を満たすものが存在する。
\begin{enumerate}
\item $R'$ は次元 $1$ のネーター局所環である。
\item $R \to R'$ は局所環準同型、すなわち $R'$ は $R$ を支配する。
\item $R \to R'$ は本質的有限型である。
\end{enumerate}
\end{lemma}

\begin{proof}
$R$ を支配する任意の付値環 $A \subset K$ を選ぶ
（補題 \ref{algebra-lemma-dominate} により存在する）。
% [P01-E0562] The frozen missing designation preposition is rendered naturally.
対応する付値を $v$ と書く。$x_1, \ldots, x_r$ を $R$ の極大イデアル
$\mathfrak m$ の極小生成系とする。
$v(x_r) = \min\{v(x_1), \ldots, v(x_r)\}$ と仮定してよいし、そうする。
環
$$
S = R[x_1/x_r, x_2/x_r, \ldots, x_{r - 1}/x_r] \subset K.
$$
を考える。$\mathfrak mS = x_rS$ は単項イデアルである。また、
$S \subset A$ かつ $v(x_r) > 0$ なので $x_rS \not = S$ である。
$x_rS$ 上の極小素イデアル $\mathfrak q$ を選ぶ。このとき補題
\ref{algebra-lemma-minimal-over-1} により $\text{height}(\mathfrak q) = 1$ であり、
$\mathfrak q$ は $\mathfrak m$ の上にある。従って
$R' = S_{\mathfrak q}$ が解である。
\end{proof}

\begin{lemma}[Koll\'ar]
\label{algebra-lemma-hart-serre-loc-thm}
\begin{reference}
これは J\'anos Koll\'ar の近刊論文
『Variants of normality for Noetherian schemes』から採られている。
\end{reference}
$(R, \mathfrak m)$ をネーター局所環とする。このとき、次のちょうど
一つが成り立つ。
\begin{enumerate}
\item $(R, \mathfrak m)$ はアルティン環である。
\item $(R, \mathfrak m)$ は次元 $1$ の正則環である。
\item $\text{depth}(R) \geq 2$ である。
\item 同型でない有限環準同型 $R \to R'$ で、その核と余核が
$\mathfrak m$ のある冪によって零化され、$\mathfrak m$ が $R'$ の
随伴素イデアルでなく、かつ $R' \not = 0$ であるものが存在する。
\end{enumerate}
\end{lemma}

\begin{proof}
$(R, \mathfrak m)$ がアルティン環でないことと
$V(\mathfrak m) \subset \Spec(R)$ が疎なことは同値である。
命題 \ref{algebra-proposition-dimension-zero-ring} を参照せよ。以下これを仮定する。

\medskip\noindent
$J \subset R$ を $\mathfrak m$ のある冪によって零化される最大の
イデアルとする。$J \not = 0$ ならば、$R \to R/J$ により
$(R, \mathfrak m)$ は (4) の場合である。

\medskip\noindent
そうでなければ $J = 0$ である。特に $\mathfrak m$ は $R$ の
随伴素イデアルでなく、補題 \ref{algebra-lemma-ideal-nonzerodivisor} により
非零因子 $x \in \mathfrak m$ が存在する。$\mathfrak m$ が
$R/xR$ の随伴素イデアルでなければ、同じ補題により
$\text{depth}(R) \geq 2$ である。従って残るのは、
$y \not \in xR$ かつ $y \mathfrak m \subset xR$ を満たす
$y \in R$ が存在する場合である。

\medskip\noindent
$y \mathfrak m \subset x \mathfrak m$ ならば、写像
$\varphi : \mathfrak m \to \mathfrak m$、$f \mapsto yf/x$ を考えられる
（$x$ は非零因子なのでこれは良定義である）。補題
\ref{algebra-lemma-charpoly-module} の行列式技巧により、係数が $R$ にある
モニック多項式 $P$ で $P(\varphi) = 0$ を満たすものが存在する。
従って $R_x$ において $P(y/x) = 0$ である。
$R' \subset R_x$ を $R$ と $y/x$ で生成される環とする。このとき
$R \subset R'$ であり、$R'/R$ は $\mathfrak m$ のある冪によって
零化される有限 $R$-加群である。従って $R$ は (4) の場合である。

\medskip\noindent
そうでなければ、ある $R$ の単元 $u$ に対して $y t = u x$ となる
$t \in \mathfrak m$ が存在する。$t$ を $u^{-1}t$ で置き換えると
$yt = x$ を得る。特に $y$ は非零因子である。任意の
$t' \in \mathfrak m$ に対し、ある $s \in R$ について
$y t' = x s$ である。従って $y (t' - s t ) = x s - x s = 0$ である。
$y$ は零因子でないから $t' = ts$ となり、$\mathfrak m = (t)$ である。
従って $(R, \mathfrak m)$ は次元 1 の正則環である。

\medskip\noindent
ここまでの議論により、すべての $R$ は四つの場合のいずれかに入る。
証明を終えるには、性質の六つの組合せのそれぞれについて、二つの性質が
同時には成り立たないことを示さなければならない。ここでは (2) と (3) が
それぞれ (4) を排除することだけを示し、他の場合は読者に委ねる。

\medskip\noindent
$R$ が次元 $1$ の正則環であり、$R \to R'$ が有限環準同型で、その核と
余核が $\mathfrak m$ のある冪によって零化され、$\mathfrak m$ が
$R'$ の随伴素イデアルでないと仮定する。このとき $R'$ は深さが少なくとも
$1$ の有限 $R$-加群であり、補題 \ref{algebra-lemma-regular-mcm-free} により
自由である。
% [P01-E0563] The frozen wrong locative preposition is rendered naturally.
$R \to R'$ は一般点で同型なので、$R'$ は $R$-加群として階数 $1$ の
自由加群でなければならない。すると
$R' \to \text{End}_R(R') \cong R$ は写像 $R \to R'$ の逆写像である。
従って (4) は成り立たない。

\medskip\noindent
$R$ の深さが $\geq 2$ であり、$R \to R'$ が有限環準同型で、その核と
余核が $\mathfrak m$ のある冪によって零化され、$\mathfrak m$ が
$R'$ の随伴素イデアルでないと仮定する。このとき $R \to R'$ は
必ず単射である（さもなければ核は深さ $0$ と深さ $\geq 1$ を同時に
もつ）。また余核の深さは少なくとも $1$ である（補題
\ref{algebra-lemma-depth-in-ses} と、仮定により $R'$ の深さが $\geq 1$ で
あることによる）。従って、余核が $\{\mathfrak m\}$ 上に台をもつ
ならば、それも $0$ でなければならない。ゆえに (4) は成り立たない。
\end{proof}

\begin{lemma}
\label{algebra-lemma-nonregular-dimension-one}
$R$ を極大イデアル $\mathfrak m$ をもつ局所環とする。
$R$ はネーターで次元 $1$ であり、
$\dim(\mathfrak m/\mathfrak m^2) > 1$ であると仮定する。
このとき、次を満たす環準同型 $R \to R'$ が存在する。
\begin{enumerate}
\item $R \to R'$ は有限である。
\item $R \to R'$ は同型でない。
\item $R \to R'$ の核と余核は $\mathfrak m$ のある冪で零化される。
\item $\mathfrak m$ は $R'$ の随伴素イデアルでない。
\end{enumerate}
\end{lemma}

\begin{proof}
補題 \ref{algebra-lemma-hart-serre-loc-thm} と、$R$ がアルティンでなく、
正則でなく、深さが $\geq 2$ でもないという事実から従う
（最後の点は、補題 \ref{algebra-lemma-bound-depth} により深さが次元を
超えないためである）。
\end{proof}

\begin{example}
\label{algebra-example-nonreduced}
ネーター局所環
$$
R = k[[x, y]]/(y^2)
$$
を考える。これは次元 1 で Cohen--Macaulay である。
補題 \ref{algebra-lemma-nonregular-dimension-one} のような拡大の例は
$$
k[[x, y]]/(y^2) \subset k[[x, z]]/(z^2), \ \ y \mapsto xz
$$
である。言い換えれば、これは $R$ に $y/x$ を添加して得られる。
この構成を $n > 1$ 回反復すると、元 $y/x^n$ を添加することになる。
\end{example}

\begin{example}
\label{algebra-example-bad-dvr-char-p}
$k$ を標数 $p > 0$ の体で、$p$ 乗全体の部分体 $k^p$ 上の $k$ の
次数が無限であるものとする。例えば
$k = \mathbf{F}_p(t_1, t_2, t_3, \ldots)$ である。環
$$
A =
\left\{
\sum a_i x^i \in k[[x]] \text{ such that }
[k^p(a_0, a_1, a_2, \ldots) : k^p] < \infty
\right\}
$$
を考える。このとき $A$ は離散付値環であり、その完備化は
$A^\wedge = k[[x]]$ である。$A \subset k[[x]]$ が誘導する分数体の
拡大は無限次純非分離拡大であることに注意する。
任意の $f \in k[[x]]$、$f \not \in A$ を選び、
$R = A[f] \subset k[[x]]$ とおく。このとき $R$ は次元 $1$ の
ネーター局所整域であり、その完備化 $R^\wedge$ は被約でない
（考えてみよ！）。
\end{example}

\begin{remark}
\label{algebra-remark-resolution-dim-1}
$R$ を $1$ 次元半局所ネーター整域とする。$R_{\mathfrak m}$ が正則で
ないような極大イデアル $\mathfrak m \subset R$ が存在するなら、
補題 \ref{algebra-lemma-nonregular-dimension-one} を $(R, \mathfrak m)$ に
適用して有限環拡大 $R \subset R_1$ を得られる。
（例えば、$\Spec(R_1) \to \Spec(R)$ が $\Spec(R)$ のイデアル
$\mathfrak m$ に沿うブローアップとなるようにできる。）
もちろん $R_1$ は $R$ と同じ分数体をもつ $1$ 次元半局所ネーター
整域である。$R_1$ が正則半局所環でなければ、構成を反復して
$R_1 \subset R_2$ を得られる。従って有限環拡大の列
$$
R \subset R_1 \subset R_2 \subset R_3 \subset \ldots
$$
を得る。この列は、ある $n$ に対して $R_n$ が正則ならば停止しうる。
特異点解消とは、最終的に $R_n$ が実際に正則になるという主張で
あろう。しかし実際にはそうではない。すなわち、完備化が被約でない
次元 $1$、標数 $0$ のネーター局所整域 $A$ が存在する。
\cite[Proposition 3.1]{Ferrand-Raynaud} または Examples、節
\ref{examples-section-local-completion-nonreduced} を参照せよ。
標数 $p > 0$ の例については例 \ref{algebra-example-bad-dvr-char-p} を参照せよ。
ブローアップの構成は完備化と可換なので、この列が決して安定しない
ことは容易に分かる。（主として正標数における）議論については
\cite{Bennett} を参照せよ。一方、$R$ のすべての極大イデアルにおける
完備化が被約ならば、この手続きは停止する
% [P01-E0564] The frozen unresolved editorial placeholder is preserved visibly.
（将来の参照をここに挿入）。
\end{remark}

\begin{lemma}
\label{algebra-lemma-characterize-dvr}
$A$ を環とする。次は同値である。
\begin{enumerate}
\item 環 $A$ は離散付値環である。
\item 環 $A$ はネーターで体でない付値環である。
\item 環 $A$ は次元 $1$ の正則局所環である。
\item 環 $A$ は、極大イデアル $\mathfrak m$ が一つの非零元で
生成されるネーター局所整域である。
\item 環 $A$ は次元 $1$ のネーター正規局所整域である。
\end{enumerate}
この場合、$\pi$ が $A$ の極大イデアルの生成元ならば、$A$ の各非零元は
$u\pi^n$ と一意的に書ける。ここで $u \in A$ は単元である。
\end{lemma}

\begin{proof}
(1) と (2) の同値性は補題
\ref{algebra-lemma-valuation-ring-Noetherian-discrete} である。さらに、補題
\ref{algebra-lemma-valuation-ring-Noetherian-discrete} の証明において、
$A$ が離散付値環ならば $A$ は PID であることを見たので、(3) が従う。
正則局所環は整域であることに注意する（補題
\ref{algebra-lemma-regular-domain} を参照せよ）。これを用いると (3) と (4) の
同値性は次元論から従う。節 \ref{algebra-section-dimension} を参照せよ。

\medskip\noindent
(3) を仮定し、$\pi$ を極大イデアル $\mathfrak m$ の生成元とする。
すべての $n \geq 0$ に対して
$\dim_{A/\mathfrak m} \mathfrak m^n/\mathfrak m^{n + 1} = 1$ である。
これは $\pi^n$ で生成されるからである（零にはなりえない）。
特に $\mathfrak m^n = (\pi^n)$ であり、次数付き環
$\bigoplus \mathfrak m^n/\mathfrak m^{n + 1}$ は
% [P01-E0565] The frozen ambiguous coefficient-ring formula is preserved.
多項式環 $A/\mathfrak m[T]$ と同型である。
$x \in A \setminus \{0\}$ に対して
$v(x) = \max\{n \mid x \in \mathfrak m^n\}$ と定める。
言い換えれば、$u \in A^*$ に対して $x = u \pi^{v(x)}$ である。
上の注意により、すべての $x, y \in A \setminus \{0\}$ に対して
$v(xy) = v(x) + v(y)$ が成り立つ。これを $A$ の分数体 $K$ に、
$v(a/b) = v(a) - v(b)$ とおくことで拡張する
（上で示した乗法性により良定義である）。すると $A$ が、付値が
$\geq 0$ である $K$ の元全体であることは明らかである。従って補題
\ref{algebra-lemma-valuation-valuation-ring} により $A$ は付値環である。

\medskip\noindent
補題 \ref{algebra-lemma-valuation-ring-normal} により付値環は正規整域である。
従って同値な条件 (1) -- (3) は (5) を含意する。(5) を仮定する。
矛盾を導くため、$\mathfrak m$ は $1$ 元で生成できないと仮定する。
補題 \ref{algebra-lemma-nonregular-dimension-one} により、任意の
$\mathfrak m$ の非零元を逆にすると同型になるが同型ではない有限環準同型
$A \to A'$ が存在する。特に $A'$ を $A$ の分数体の部分集合と
同一視できる。$A \to A'$ は有限なので整である（補題
\ref{algebra-lemma-finite-is-integral} を参照せよ）。$A$ は正規なので
$A = A'$ となり、矛盾である。
\end{proof}

\begin{definition}
\label{algebra-definition-uniformizer}
$A$ を離散付値環とする。$A$ の極大イデアルを生成する元
$\pi \in A$ を{\it 一意化元}という。
\end{definition}

\noindent
補題 \ref{algebra-lemma-characterize-dvr} により、離散付値環の任意の二つの
一意化元は同伴である。

\begin{lemma}
\label{algebra-lemma-finite-length}
$R$ を分数体 $K$ をもつ整域とし、$M$ を $K^{\oplus r}$ の
$R$-部分加群とする。$R$ は次元 $1$ のネーター局所環であると仮定する。
任意の非零元 $x \in R$ に対して $\text{length}_R(R/xR) < \infty$ であり、
$$
\text{length}_R(M/xM) \leq r \cdot \text{length}_R(R/xR).
$$
である。
\end{lemma}

\begin{proof}
$x$ が単元ならば結果は真である。従って $x \in \mathfrak m$ と仮定して
よい。ここで右辺のイデアルは $R$ の極大イデアルである。$x$ は零でなく
$R$ は整域なので $\dim(R/xR) = 0$ であり、従って $R/xR$ は有限長をもつ。
補題のような $M \subset K^{\oplus r}$ を考える。必要なら
$K^{\oplus r}$ をより小さい部分空間で置き換えることにより、$M$ の元が
$K$-ベクトル空間として $K^{\oplus r}$ を生成すると仮定してよい。

\medskip\noindent
まず $M$ が有限 $R$-加群であると仮定する。この場合、分母を払って
$M \subset R^{\oplus r}$ と仮定できる。
% [P01-E0566] The frozen number-agreement typo is rendered naturally.
$M$ はベクトル空間として $K^{\oplus r}$ を生成するので、
$R^{\oplus r}/M$ は有限長をもつ。特に、
$x^cR^{\oplus r} \subset M$ となる整数 $c \geq 0$ が存在する。
$M \supset xM \supset x^2M \supset \ldots$ は、各逐次商が
$M/xM$ と同型である加群列である。従って
$$
n \text{length}_R(M/xM) = \text{length}_R(M/x^nM).
$$
である。$M = R^{\oplus r}$ に対する同じ議論により、
$$
n \text{length}_R(R^{\oplus r}/xR^{\oplus r}) =
\text{length}_R(R^{\oplus r}/x^nR^{\oplus r}).
$$
である。上の $c$ の選び方により、$x^nM$ は
$x^n R^{\oplus r}$ と $x^{n + c}R^{\oplus r}$ の間に挟まれる。
これから容易に
$$
r(n + c) \text{length}_R(R/xR)
\geq
n \text{length}_R(M/xM)
\geq
r (n - c) \text{length}_R(R/xR)
$$
を得る。従って有限の場合には、実際、補題の結果が等号つきで得られる。

\medskip\noindent
次に $M$ が有限でないと仮定する。ある自然数 $k$ に対して
$M/xM$ の長さが $\geq k$ であると仮定する。このとき
% [P01-E0567] The frozen missing inclusion relation before N_k is preserved.
$$
0 \subset N_0 \subset N_1 \subset N_2 \subset \ldots N_k \subset M/xM
$$
で、$i = 0, \ldots k - 1$ に対して
$N_i \not = N_{i + 1}$ となるものを選べる。
$i = 1, \ldots, k$ に対し、$xM$ を法とする剰余類が $N_i$ に属するが
$N_{i - 1}$ に属さない元 $m_i \in M$ を選ぶ。有限 $R$-加群
$M' = Rm_1 + \ldots + Rm_k \subset M$ を考える。
$N'_i \subset M'/xM'$ を $N_i$ の逆像とする。$m_i$ の選び方から
$N'_i \not =N'_{i + 1}$ であることは明らかである。
従って $\text{length}_R(M'/xM') \geq k$ である。有限の場合により
$k \leq r\text{length}_R(R/xR)$ を得て、望む結果が従う。
\end{proof}

\noindent
次が最初の応用である。

\begin{lemma}
\label{algebra-lemma-finite-extension-residue-fields-dimension-1}
$R \to S$ を、分数体の単射 $K \subset L$ を誘導する整域の準同型とする。
$R$ が次元 $1$ のネーター局所環であり、$[L : K] < \infty$ ならば、
\begin{enumerate}
\item $R$ の極大イデアル $\mathfrak m$ の上にある $S$ の各素イデアル
$\mathfrak n_i$ は極大である。
\item そのようなものは有限個である。
\item 各 $i$ に対して
$[\kappa(\mathfrak n_i) : \kappa(\mathfrak m)] < \infty$ である。
\end{enumerate}
\end{lemma}

\begin{proof}
非零元 $x \in \mathfrak m$ を選ぶ。補題 \ref{algebra-lemma-finite-length} を
$S \subset L \cong K^{\oplus n}$ に適用する。ここで
$n = [L : K]$ である。従って環 $S/xS$ は $R$ 上有限長をもつ。
それゆえ $S/\mathfrak m S$ は $\kappa(\mathfrak m)$ 上有限長をもつ。
言い換えれば、$\dim_{\kappa(\mathfrak m)} S/\mathfrak m S$ は有限である
（補題 \ref{algebra-lemma-dimension-is-length}）。従って
$S/\mathfrak mS$ はアルティン環である（補題
\ref{algebra-lemma-finite-dimensional-algebra}）。
% [P01-E0568] The frozen subject-verb disagreement is rendered naturally.
アルティン環に関する構造結果は (1) と (2) を含意する。例えば補題
\ref{algebra-lemma-artinian-finite-length} を参照せよ。(3) は上で確立した
有限性から従う。
\end{proof}

\begin{lemma}
\label{algebra-lemma-finite-length-global}
$R$ を分数体 $K$ をもつ整域とし、$M$ を $K^{\oplus r}$ の
$R$-部分加群とする。$R$ は次元 $1$ のネーター環であると仮定する。
任意の非零元 $x \in R$ に対して
$\text{length}_R(M/xM) < \infty$ である。
\end{lemma}

\begin{proof}
$R$ は次元 $1$ なので、$x$ は有限個の素イデアル
$\mathfrak m_i$、$i = 1, \ldots, n$ に含まれ、各々は極大である。
$R$ はネーターなので、$R/xR$ はアルティン環であり、命題
\ref{algebra-proposition-dimension-zero-ring} と補題
\ref{algebra-lemma-artinian-finite-length} により
$R/xR = \prod_{i = 1, \ldots, n} (R/xR)_{\mathfrak m_i}$ である。
従って $M/xM$ も同様に、その $\mathfrak m_i$ における局所化の積
$M/xM = \prod (M/xM)_{\mathfrak m_i}$ に分解する。
$R_{\mathfrak m_i}$ 上の $M_{\mathfrak m_i}$ に補題
\ref{algebra-lemma-finite-length} を適用すると、各
$M_{\mathfrak m_i}/xM_{\mathfrak m_i} = (M/xM)_{\mathfrak m_i}$ は
$R_{\mathfrak m_i}$ 上有限長をもつ。従って $M/xM$ は $R$ 上有限長を
もつ。実際、上のことから、$M/xM$ は逐次商が剰余体
$\kappa(\mathfrak m_i)$ と同型であるような $R$-部分加群の
有限フィルトレーションをもつ。
\end{proof}

\begin{lemma}[Krull-Akizuki]
\label{algebra-lemma-krull-akizuki}
$R$ を分数体 $K$ をもつ整域とする。$L/K$ を有限体拡大とする。
$R$ はネーターで $\dim(R) = 1$ であると仮定する。このとき
$R \subset A \subset L$ を満たす任意の環 $A$ はネーターである。
\end{lemma}

\begin{proof}
$I \subset A$ を非零イデアルとする。補題
\ref{algebra-lemma-intersection-not-zero} により、非零元
$x \in I \cap R$ を見つけられる。このとき
$I/xA \subset A/xA$ を得る。補題 \ref{algebra-lemma-finite-length-global}
により、$R$-加群 $A/xA$ は $R$-加群として有限長をもつ。
従って $I/xA$ は $R$-加群として有限長をもつ。ゆえに $I$ は
$A$ のイデアルとして有限生成である。
\end{proof}

\begin{lemma}
\label{algebra-lemma-exists-dvr}
$R$ を分数体 $K$ をもつネーター局所整域とし、$R$ は体でないと
仮定する。$L/K$ を有限生成体拡大とする。
% [P01-E0569] The frozen missing indefinite article has no Japanese analogue.
このとき、分数体が $L$ で $R$ を支配する離散付値環 $A$ が存在する。
\end{lemma}

\begin{proof}
$L$ が $K$ 上有限でなければ、$K$ 上の $L$ の超越基底
$x_1, \ldots, x_r$ を選び、$R$ を、$\mathfrak m_R$ と
$x_1, \ldots, x_r$ で生成される極大イデアルにおける
$R[x_1, \ldots, x_r]$ の局所化で置き換える。従って
$K \subset L$ は有限であると仮定してよい。

\medskip\noindent
補題 \ref{algebra-lemma-dominate-by-dimension-1} により
$\dim(R) = 1$ と仮定してよい。

\medskip\noindent
$A \subset L$ を $L$ における $R$ の整閉包とする。補題
\ref{algebra-lemma-krull-akizuki} によりこれはネーターである。補題
\ref{algebra-lemma-integral-overring-surjective} により、$R$ の極大イデアルの
上にある素イデアル $\mathfrak q \subset A$ が存在する。補題
\ref{algebra-lemma-characterize-dvr} により、環 $A_{\mathfrak q}$ は望むように
$R$ を支配する離散付値環である。
\end{proof}

\section{因数分解}
\label{algebra-section-factoring}

\noindent
ここでは、学部初年次の代数学の講義で通常扱われるいくつかの概念と、
それらの間の関係を述べる。

\begin{definition}
\label{algebra-definition-irreducible-prime-element}
$R$ を整域とする。
\begin{enumerate}
\item $x, y \in R$ が{\it 同伴}であるとは、ある単元 $u \in R^*$ が存在して
$x = uy$ となることをいう。
\item 元 $x \in R$ が{\it 既約}であるとは、それが非零かつ単元でなく、
$x = yz$、$y, z \in R$ ならば、$y$ が単元であるか $x$ と同伴であることをいう。
\item 元 $x \in R$ が{\it 素元}であるとは、$x$ で生成されるイデアルが
素イデアルであることをいう。
\end{enumerate}
\end{definition}

\begin{lemma}
\label{algebra-lemma-easy-divisibility}
$R$ を整域とし、$x, y \in R$ とする。
このとき、$x$ と $y$ が同伴であることと $(x) = (y)$ であることは同値である。
\end{lemma}

\begin{proof}
ある単元 $u \in R$ に対して $x = uy$ ならば $(x) \subset (y)$ であり、
$y = u^{-1}x$ なので $(y) \subset (x)$ でもある。逆に $(x) = (y)$ と仮定する。
% [P01-E0570] The frozen undefined coefficient ring A is preserved.
このとき、ある $f, g \in A$ に対して $x = fy$ および $y = gx$ である。
従って $x = fg x$ であり、$R$ は整域なので $fg = 1$ である。
ゆえに $x$ と $y$ は同伴である。
\end{proof}

\begin{lemma}
\label{algebra-lemma-factorization-exists}
$R$ を整域とする。次の条件を考える。
\begin{enumerate}
\item 環 $R$ は単項イデアルについて昇鎖条件を満たす。
\item 任意の非零非単元 $a \in R$ は、各 $b_i$ が $R$ の既約元である
因数分解 $a = b_1 \ldots b_k$ をもつ。
\end{enumerate}
このとき (1) は (2) を含意する。
\end{lemma}

\begin{proof}
非零かつ単元でなく、既約元への因数分解をもたない元 $x$ をとる。
$x_1 = x$ とおく。
% [P01-E0571] The frozen stronger claim about both factors is preserved.
$y$ と $z$ のいずれも既約でも単元でもないように $x = yz$ と書ける。
このとき、$y$ が既約元への因数分解をもたなければ $x_2 = y$ とおき、
$z$ が既約元への因数分解をもたなければ $x_2 = z$ とおく。
このように続けると、$x_n/x_{n + 1}$ が単元でない $R$ の元の列
$$
\ldots | x_3 | x_2 | x_1
$$
を得る。これは単項イデアルの狭義昇鎖
$(x_1) \subset (x_2) \subset (x_3) \subset \ldots$ を与え、証明が完了する。
\end{proof}

\begin{definition}
\label{algebra-definition-UFD}
{\it 一意分解整域}（略して {\it UFD}）とは、次を満たす整域 $R$ のことである。
$x \in R$ が非零非単元ならば $x$ は既約元への因数分解をもち、
$$
x = a_1 \ldots a_m = b_1 \ldots b_n
$$
が既約元への因数分解ならば $n = m$ であり、ある置換
$\sigma : \{1, \ldots, n\} \to \{1, \ldots, n\}$ が存在して、
$a_i$ と $b_{\sigma(i)}$ は同伴である。
\end{definition}

\begin{lemma}
\label{algebra-lemma-characterize-UFD}
$R$ を整域とし、すべての非零非単元が既約元へ因数分解されると仮定する。
このとき、$R$ が UFD であることと、すべての既約元が素元であることは同値である。
\end{lemma}

\begin{proof}
$R$ が UFD であると仮定し、$x \in R$ を既約元とする。
$ab \in (x)$、すなわち $ab = cx$ とする。因数分解
$a = a_1 \ldots a_n$、$b = b_1 \ldots b_m$、$c = c_1 \ldots c_r$ を選ぶ。
因数分解
$$
a_1 \ldots a_n b_1 \ldots b_m = c_1 \ldots c_r x
$$
の一意性により、$x$ は元
% [P01-E0572] The frozen malformed factor range is preserved.
$a_1, \ldots, b_m$ の一つと同伴である。
言い換えれば、$a \in (x)$ または $b \in (x)$ であり、従って $x$ は素元である。

\medskip\noindent
すべての既約元が素元であると仮定する。既約元への因数分解が、置換と
同伴な元への置換を除いて一意であることを示せばよい。
$a_i$ と $b_j$ が既約であり $a_1 \ldots a_m = b_1 \ldots b_n$ とする。
$a_1$ は素元なので、ある $j$ に対して $b_j \in (a_1)$ である。
番号を付け直して $b_1 \in (a_1)$ と仮定してよい。このとき
$b_1 = a_1 u$ であり、$b_1$ は既約なので $u$ は単元である。
従って $a_1$ と $b_1$ は同伴であり、
% [P01-E0573] The frozen swapped product endpoints are preserved.
$a_2 \ldots a_n = ub_2\ldots b_m$ である。$n + m$ に関する帰納法により、
$n = m$ であり、望むように $i = 2, \ldots, n$ に対して
$a_i$ は $b_{\sigma(i)}$ と同伴である。
\end{proof}

\begin{lemma}
\label{algebra-lemma-characterize-UFD-height-1}
$R$ をネーター整域とする。このとき、$R$ が UFD であることと、
高さ $1$ のすべての素イデアルが単項であることは同値である。
\end{lemma}

\begin{proof}
$R$ が UFD であると仮定し、$\mathfrak p$ を高さ 1 の素イデアルとする。
非零元 $x \in \mathfrak p$ をとり、$x = a_1 \ldots a_n$ を既約元への
因数分解とする。$\mathfrak p$ は素イデアルなので、ある $i$ に対して
$a_i \in \mathfrak p$ である。補題 \ref{algebra-lemma-characterize-UFD} により
イデアル $(a_i)$ は素イデアルである。$\mathfrak p$ の高さは $1$ なので、
$(a_i) = \mathfrak p$ と結論できる。

\medskip\noindent
高さ $1$ のすべての素イデアルが単項であると仮定する。$R$ はネーターなので、
補題 \ref{algebra-lemma-factorization-exists} により任意の非零非単元 $x$ は
既約元への因数分解をもつ。補題 \ref{algebra-lemma-characterize-UFD} により、
既約元 $x$ が素元であることを示せば十分である。
$(x)$ を含む極小素イデアル $(x) \subset \mathfrak p$ をとる。
補題 \ref{algebra-lemma-minimal-over-1} により $\mathfrak p$ の高さは $1$ である。
仮定により $\mathfrak p = (y)$ である。従って $x = yz$ であり、
$x$ は既約なので $z$ は単元である。ゆえに $(x) = (y)$ であり、
$x$ は素元である。
\end{proof}

\begin{lemma}[Nagata の一意分解性判定法]
\label{algebra-lemma-invert-prime-elements}
\begin{reference}
\cite[Lemma 2]{Nagata-UFD}
\end{reference}
$A$ を整域とし、$S \subset A$ を素元で生成される乗法的部分集合とする。
$x \in A$ を既約元とする。このとき
\begin{enumerate}
\item $S^{-1}A$ における $x$ の像は既約元または単元であり、
\item $x$ が素元であることと、$S^{-1}A$ における $x$ の像が
$S^{-1}A$ の素元または単元であることは同値である。
\end{enumerate}
さらにこのとき、$A$ が UFD であることと、$A$ の任意の非零非単元が
既約元への因数分解をもち、かつ $S^{-1}A$ が UFD であることは同値である。
\end{lemma}

\begin{proof}
$\alpha, \beta \in S^{-1}A$ に対して $x = \alpha \beta$ とする。
このとき、ある $a, b \in A$、$s, s' \in S$ に対して
$\alpha = a/s$ および $\beta = b/s'$ である。従って $ss'x = ab$ を得る。
仮定により、ある素元 $p_i$ によって $ss' = p_1 \ldots p_r$ と書ける。
各 $i$ について、元 $p_i$ は $a$ または $b$ のいずれかを割る。
割り算を行うと、ある $s'', s''' \in S$ に対して
$x = a' b'$、$a = s'' a'$、$b = s''' b'$ となる因数分解を得る。
$x$ は既約なので $a'$ または $b'$ のいずれかは単元である。
元に戻ってたどれば、$\alpha$ または $\beta$ のいずれかが単元である。
これで (1) が示された。

\medskip\noindent
$x$ が素元であると仮定する。このとき $A/(x)$ は整域である。従って
$S^{-1}A/xS^{-1}A = S^{-1}(A/(x))$ は整域または零環である。
ゆえに $x$ は素元または単元に写る。

\medskip\noindent
$S^{-1}A$ における $x$ の像が単元であると仮定する。このとき、ある
$s \in S$ と $y \in A$ に対して $y x = s$ である。仮定により、
$p_i$ を素元として $s = p_1 \ldots p_r$ と書ける。各 $i$ について、
$p_i$ は $y$ を割るか、または $p_i$ は $x$ を割る。後者の場合には
$p_i$ と $x$ は同伴であり（$x$ は既約である）、証明は終わる。
しかしすべての $i = 1, \ldots, r$ について前者の場合が起これば、
$x$ は単元となり矛盾する。

\medskip\noindent
$S^{-1}A$ における $x$ の像が素元であると仮定する。
$a, b \in A$ かつ $ab \in (x)$ と仮定する。このとき、ある
$s \in S$ と $y \in A$ に対して $sa = xy$ または $sb = xy$ である。
前者が起こるとする。仮定により、$p_i$ を素元として
$s = p_1 \ldots p_r$ と書ける。各 $i$ について、$p_i$ は $y$ を割るか、
または $p_i$ は $x$ を割る。後者の場合には $p_i$ と $x$ は同伴であり
（$x$ は既約である）、証明は終わる。すべての $i = 1, \ldots, r$ について
前者の場合が起これば、望むように $a \in (x)$ である。これで (2) の証明が完了する。

\medskip\noindent
補題の最後の主張は (1)、(2) および補題 \ref{algebra-lemma-characterize-UFD} から従う。
\end{proof}

\begin{lemma}
\label{algebra-lemma-UFD-ascending-chain-condition-principal-ideals}
UFD は単項イデアルについて昇鎖条件を満たす。
\end{lemma}

\begin{proof}
$R$ の単項イデアルの昇鎖
$(a_1) \subset (a_2) \subset (a_3) \subset \ldots$ を考える。
$p_i$ を素元として $a_1 = p_1^{e_1} \ldots p_r^{e_r}$ と書く。
このとき、ある $0 \leq c_i \leq e_i$ に対して $a_n$ は
$p_1^{c_1} \ldots p_r^{c_r}$ と同伴である。可能性は有限個しかないので結論を得る。
\end{proof}

\begin{lemma}
\label{algebra-lemma-factoring-in-polynomial}
$R$ を整域とし、$R$ が単項イデアルについて昇鎖条件を満たすと仮定する。
このとき、$R$ 上の多項式環も同じ性質を満たす。
\end{lemma}

\begin{proof}
$R[x]$ の単項イデアルの昇鎖
$(f_1) \subset (f_2) \subset (f_3) \subset \ldots$ を考える。
$f_{n + 1}$ は $f_n$ を割るので、この列では次数が減少する。
従って、すべての $n \gg 0$ に対して $f_n$ の次数は一定の $d \geq 0$ である。
$a_n$ を $f_n$ の最高次係数とする。条件 $f_n \in (f_{n + 1})$ は、
すべての $n$ に対して $a_{n + 1}$ が $a_n$ を割ることを含意する。
$R$ に関する仮定と補題 \ref{algebra-lemma-easy-divisibility} により、十分大きいすべての
$n$ について $a_{n + 1}$ と $a_n$ は同伴である。従って大きい $n$ について、
$u \in R$ で $f_n = u f_{n + 1}$ となり、この係数は次数の理由からもとの環に属し、
$a_n$ と $a_{n + 1}$ が同伴であることから単元である。
\end{proof}

\begin{lemma}
\label{algebra-lemma-polynomial-ring-UFD}
UFD 上の多項式環は UFD である。特に、$k$ が体ならば
$k[x_1, \ldots, x_n]$ は UFD である。
\end{lemma}

\begin{proof}
$R$ を UFD とする。このとき $R$ は単項イデアルについて昇鎖条件を満たす
（補題 \ref{algebra-lemma-UFD-ascending-chain-condition-principal-ideals}）。従って
$R[x]$ も単項イデアルについて昇鎖条件を満たし
（補題 \ref{algebra-lemma-factoring-in-polynomial}）、ゆえに $R[x]$ のすべての元は
既約元への因数分解をもつ（補題 \ref{algebra-lemma-factorization-exists}）。
$S \subset R$ を素元で生成される乗法的部分集合とする。
$R$ のすべての非単元は素元の積なので、$K = S^{-1}R$ は $R$ の分数体である。
$R$ の各素元は $R[x]$ の素元に写り、
$S^{-1}(R[x]) = S^{-1}R[x] = K[x]$ は UFD（さらに PID）であることに注意する。
従って補題 \ref{algebra-lemma-invert-prime-elements} を適用して結論を得る。
\end{proof}

\begin{lemma}
\label{algebra-lemma-UFD-normal}
一意分解整域は正規である。
\end{lemma}

\begin{proof}
$R$ を UFD とする。$x$ を $R$ の分数体の元で $R$ 上整なものとする。
ある $a_i \in R$ に対して
$x^d - a_1 x^{d - 1} - \ldots - a_d = 0$ とする。
$u$ を単元、$e_i \in \mathbf{Z}$、$p_1, \ldots, p_r$ を互いに同伴でない
既約元として $x = u p_1^{e_1} \ldots p_r^{e_r}$ と書ける。
補題を示すには $e_i \geq 0$ を示せばよい。そうでなく、例えば
$e_1 < 0$ とする。このとき $N \gg 0$ に対して
$$
u^d p_2^{de_2 + N} \ldots p_r^{de_r + N} =
p_1^{-de_1}p_2^N \ldots p_r^N(
\sum\nolimits_{i = 1, \ldots, d} a_i x^{d - i} ) \in (p_1)
$$
を得るが、これは $R$ における因数分解の一意性に反する。
\end{proof}

\begin{definition}
\label{algebra-definition-PID}
{\it 単項イデアル整域}（略して {\it PID}）とは、すべてのイデアルが
単項イデアルであるような整域 $R$ のことである。
\end{definition}

\begin{lemma}
\label{algebra-lemma-PID-UFD}
単項イデアル整域は一意分解整域である。
\end{lemma}

\begin{proof}
PID はネーターなので、補題 \ref{algebra-lemma-characterize-UFD-height-1} から従う。
\end{proof}

\begin{definition}
\label{algebra-definition-dedekind-domain}
{\it Dedekind 整域}とは、任意の非零イデアル $I \subset R$ が、非零素イデアルの積
$$
I = \mathfrak p_1 \ldots \mathfrak p_r
$$
として、$\mathfrak p_i$ の置換を除いて一意的に書けるような整域 $R$ のことである。
\end{definition}

\begin{lemma}
\label{algebra-lemma-PID-dedekind}
PID は Dedekind 整域である。
\end{lemma}

\begin{proof}
$R$ を PID とする。$R$ のすべての非零イデアルは単項であり、
$R$ は UFD である（補題 \ref{algebra-lemma-PID-UFD}）。さらに $R$ のすべての既約元は
素元なので（補題 \ref{algebra-lemma-characterize-UFD}）、元の因数分解は素イデアルへの
因数分解を与え、主張が従う。
\end{proof}

\begin{lemma}
\label{algebra-lemma-product-ideals-principal}
\begin{slogan}
イデアルの積が可逆加群であることと、両因子が可逆加群であることは同値である。
\end{slogan}
$A$ を環とする。$I$ と $J$ を $A$ の非零イデアルとし、ある非零因子
$f \in A$ に対して $IJ = (f)$ であると仮定する。このとき $I$ と $J$ は
有限生成イデアルであり、$A$-加群として階数 $1$ の有限局所自由加群である。
\end{lemma}

\begin{proof}
補題 \ref{algebra-lemma-finite-projective} により、$I$ と $J$ が階数 $1$ の有限局所自由
$A$-加群であることを示せば十分である。そのために、$x_i \in I$、$y_i \in J$ として
$f = \sum_{i = 1, \ldots, n} x_i y_i$ と書く。また、ある $a_i \in A$ に対して
$x_i y_i = a_i f$ とも書ける。$f$ は非零因子なので $\sum a_i = 1$ である。
従って、各 $I_{a_i}$ および $J_{a_i}$ が $A_{a_i}$ 上階数 $1$ の自由加群であることを
示せば十分である。$A$ を $A_{a_i}$ で置き換えれば、ある $x \in I$ と $y \in J$ に
対して $f = xy$ と結論できる。$x$ と $y$ はともに非零因子であることに注意する。
$I = (x)$ および $J = (y)$ であると主張する。これで証明が完了する。
実際、$x' \in I$ ならば、ある $a \in A$ に対して $x'y = af = axy$ である。
従って $x' = ax$ となり、主張が従う。
\end{proof}

\begin{lemma}
\label{algebra-lemma-characterize-Dedekind}
% [P01-E0574] The frozen missing terminal punctuation is rendered naturally.
$R$ を環とする。次は同値である。
\begin{enumerate}
\item $R$ は Dedekind 整域である。
\item $R$ はネーター整域であり、すべての非零極大イデアル $\mathfrak m$ に対して
局所環 $R_{\mathfrak m}$ は離散付値環である。
\item $R$ はネーター正規整域であり、$\dim(R) \leq 1$ である。
\end{enumerate}
\end{lemma}

\begin{proof}
(1) を仮定する。Dedekind 整域の定義では $R$ がネーターであることを
仮定しなかったので、議論は自明ではない。$\mathfrak p \subset R$ を素イデアルとする。
定義における因数分解の一意性により $\mathfrak p \not = \mathfrak p^2$ である。
$x \not \in \mathfrak p^2$ を満たす $x \in \mathfrak p$ を選ぶ。
$y \in \mathfrak p$ をもう一つの元とする（例えば $y = 0$）。
$(x, y) = \mathfrak p_1 \ldots \mathfrak p_r$ と書く。
$(x, y) \subset \mathfrak p$ なので、素イデアル $\mathfrak p_i$ の少なくとも一つは
$\mathfrak p$ に含まれる。しかし $x \not \in \mathfrak p^2$ なので、そのようなものは
高々一つである。従って $\mathfrak p_1, \ldots, \mathfrak p_r$ のちょうど一つが
$\mathfrak p$ に含まれ、$\mathfrak p_1 \subset \mathfrak p$ としてよい。
任意の $y$ に対して $(x, y)R_\mathfrak p = \mathfrak p_1R_\mathfrak p$ は素イデアルである。
$(x)R_\mathfrak p = \mathfrak pR_\mathfrak p$ と主張する。実際、
$y \in \mathfrak p$ を選ぶ。上の議論を $y^2$ に適用すると
$(x, y^2)R_\mathfrak p$ は素イデアルである。従って
$y \in (x, y^2)R_\mathfrak p$、すなわち $R_\mathfrak p$ において
$y = ax + by^2$ である。ゆえに $(1 - by)y = ax \in (x)R_\mathfrak p$、
すなわち望むように $y \in (x)R_\mathfrak p$ である。

\medskip\noindent
改めて $(x) = \mathfrak p_1 \ldots \mathfrak p_r$ と書き、
$\mathfrak p_1 \subset \mathfrak p$ とすると、
$\mathfrak p_1 R_\mathfrak p = \mathfrak p R_\mathfrak p$、すなわち
$\mathfrak p_1 = \mathfrak p$ と結論できる。さらに補題
\ref{algebra-lemma-product-ideals-principal} により、
$\mathfrak p_1 = \mathfrak p$ は $R$ の有限生成イデアルである。
補題 \ref{algebra-lemma-cohen} により $R$ はネーターである。
% [P01-E0575] The frozen m/p variable mismatch and quantifier are preserved.
さらに、すべての素イデアル $\mathfrak p$ に対して $R_\mathfrak m$ が
離散付値環であることが従う。補題 \ref{algebra-lemma-characterize-dvr} を参照せよ。

\medskip\noindent
(2) と (3) の同値性は補題 \ref{algebra-lemma-normality-is-local} および
\ref{algebra-lemma-characterize-dvr} から従う。(2) と (3) が成り立つと仮定する。
$I \subset R$ をイデアルとする。$I$ の因数分解を構成する。
$I$ が素イデアルならば示すことはない。そうでなければ、
$I \subset \mathfrak p$ かつ $\mathfrak p \subset R$ が極大となるように選ぶ。
$J = \{x \in R \mid x \mathfrak p \subset I\}$ とおく。
$J \mathfrak p = I$ と主張する。$R$ の極大イデアル $\mathfrak m$ における局所化後に
これを確かめれば十分である（$J$ の形成は局所化と可換であり、補題
\ref{algebra-lemma-characterize-zero-local} を用いる）。このとき、
$\mathfrak p R_\mathfrak m = R_\mathfrak m$ なら結果は明らかであり、そうでなければ
$\mathfrak p R_\mathfrak m = \mathfrak m R_\mathfrak m$ である。
後者の場合 $\mathfrak p R_\mathfrak m = (\pi)$ であり、$\mathfrak p$ が単項である場合は
直ちに分かる。ネーター帰納法によりイデアル $J$ は因数分解をもち、
$I$ の望む因数分解を得る。因数分解の一意性の証明は省略する。
\end{proof}

\noindent
次は Krull--Akizuki の補題の変形である。

\begin{lemma}
\label{algebra-lemma-integral-closure-Dedekind}
$A$ を分数体 $K$ をもつ次元 $1$ のネーター整域とする。
$L/K$ を有限拡大とし、$B$ を $L$ における $A$ の整閉包とする。
このとき $B$ は Dedekind 整域であり、
$\Spec(B) \to \Spec(A)$ は全射で有限なファイバーをもち、有限次剰余体拡大を誘導する。
\end{lemma}

\begin{proof}
Krull--Akizuki（補題 \ref{algebra-lemma-krull-akizuki}）により環 $B$ はネーターである。
補題 \ref{algebra-lemma-integral-sub-dim-equal} により $\dim(B) = 1$ である。
従って補題 \ref{algebra-lemma-characterize-Dedekind} により $B$ は Dedekind 整域である。
スペクトル上の写像の全射性は補題
\ref{algebra-lemma-integral-overring-surjective} から従う。最後の二つの主張は補題
\ref{algebra-lemma-finite-extension-residue-fields-dimension-1} から従う。
\end{proof}

\section{消滅位数}
\label{algebra-section-orders-of-vanishing}

\begin{lemma}
\label{algebra-lemma-ord-additive}
$R$ を次元 $1$ の半局所ネーター環とする。$a, b \in R$ が非零因子ならば
$$
\text{length}_R(R/(ab)) =
\text{length}_R(R/(a)) +
\text{length}_R(R/(b))
$$
であり、これらの長さは有限である。
\end{lemma}

\begin{proof}
有限性は補題 \ref{algebra-lemma-finite-length-global} で見た。
$0 \to R/(a) \to R/(ab) \to R/(b) \to 0$ という短完全列があり、
最初の写像は $b$ 倍で与えられるので、加法性も成り立つ。
% [P01-E0576] The frozen omitted complementizer is rendered naturally.
（長さの加法性、補題 \ref{algebra-lemma-length-additive} を用いよ。）
\end{proof}

\begin{definition}
\label{algebra-definition-ord}
$K$ を体とし、$R \subset K$ を局所\footnote{$R$ が半局所である場合にも
定義できるが、実際に望まれることはおそらくない！}
な次元 $1$ のネーター部分環で、その分数体が $K$ であると仮定する。
この場合、{\it $R$ に沿う消滅位数}
$$
\text{ord}_R : K^* \longrightarrow \mathbf{Z}
$$
を、$x \in R$ に対して
$$
\text{ord}_R(x) = \text{length}_R(R/(x))
$$
とし、非零な $x, y \in R$ に対して
$\text{ord}_R(x/y) = \text{ord}_R(x) - \text{ord}_R(y)$ とおくことにより定義する。
\end{definition}

\noindent
消滅位数を用いて、ベクトル空間内の格子を比較できる。まず定義を与える。

\begin{definition}
\label{algebra-definition-lattice}
$R$ を分数体 $K$ をもつ次元 $1$ のネーター局所整域とし、
$V$ を有限次元 $K$-ベクトル空間とする。
{\it $V$ 内の格子}とは、$V = K \otimes_R M$ を満たす有限 $R$-部分加群
$M \subset V$ のことである。
\end{definition}

\noindent
条件 $V = K \otimes_R M$ は、$M$ がベクトル空間 $V$ の基底を含むことを意味する。
文献の多くでは、格子の概念は環 $R$ が離散付値環である場合にのみ定義されることがある。
$R$ が離散付値環ならば任意の格子は自由 $R$-加群であるが、一般にはそうとは限らない。

\begin{lemma}
\label{algebra-lemma-compare-lattices}
$R$ を分数体 $K$ をもつ次元 $1$ のネーター局所整域とし、
$V$ を有限次元 $K$-ベクトル空間とする。
\begin{enumerate}
\item $M$ が $V$ 内の格子であり、$M \subset M' \subset V$ が
$M$ を含む $V$ の $R$-部分加群ならば、次は同値である。
\begin{enumerate}
\item $M'$ は格子である。
\item $\text{length}_R(M'/M)$ は有限である。
\item $M'$ は有限生成である。
\end{enumerate}
\item $M$ が $V$ 内の格子で、$M' \subset M$ が $M$ の $R$-部分加群ならば、
$M'$ が格子であることと $\text{length}_R(M/M')$ が有限であることは同値である。
\item $M$, $M'$ が $V$ 内の格子ならば、$M \cap M'$ および $M + M'$ も格子である。
\item $M \subset M' \subset M'' \subset V$ が $V$ 内の格子ならば
$$
\text{length}_R(M''/M) =
\text{length}_R(M'/M) +
\text{length}_R(M''/M').
$$
\item $M$, $M'$, $N$, $N'$ が $V$ 内の格子で、
$N \subset M \cap M'$、$M + M' \subset N'$ ならば
\begin{eqnarray*}
& & \text{length}_R(M/M \cap M') - \text{length}_R(M'/M \cap M')\\
& = &
\text{length}_R(M/N) - \text{length}_R(M'/N) \\
& = &
\text{length}_R(M + M' / M') - \text{length}_R(M + M'/M) \\
& = &
\text{length}_R(N' / M') - \text{length}_R(N'/M)
\end{eqnarray*}
\end{enumerate}
\end{lemma}

\begin{proof}
(1) の証明。(1)(a) を仮定する。$y_1, \ldots, y_m$ が $M'$ を生成するとする。
このとき各 $y_i = x_i/f_i$ であり、ここで $x_i \in M$、
非零な $f_i \in R$ である。従って
$f_1 \ldots f_m M' \subset M$ である。
$R$ は次元 $1$ のネーター局所環なので、ある $n$ に対して
$\mathfrak m^n \subset (f_1 \ldots f_m)$ である
（例えば、補題 \ref{algebra-lemma-one-equation} と命題
\ref{algebra-proposition-dimension-zero-ring} を組み合わせるか、補題
\ref{algebra-lemma-finite-length} と \ref{algebra-lemma-length-infinite} を組み合わせよ）。
% [P01-E0577] The frozen missing sentence boundary is rendered naturally.
言い換えれば、ある $n$ に対して $\mathfrak m^nM' \subset M$ である。
従って補題 \ref{algebra-lemma-length-finite} により
$\text{length}(M'/M) < \infty$、すなわち (1)(b) が成り立つ。
(1)(b) を仮定する。このとき $M'/M$ は有限 $R$-加群である
（補題 \ref{algebra-lemma-finite-length-finite} 参照）。
従って $M'$ は有限 $R$-加群の拡大として有限 $R$-加群である。
ゆえに (1)(c) が成り立つ。(1)(c) $\Rightarrow$ (1)(a) は、定義
\ref{algebra-definition-lattice} の直後の注意から従う。

\medskip\noindent
(2) の証明。$M$ が $V$ 内の格子であり、$M' \subset M$ が
$R$-部分加群であると仮定する。(1) で、$M'$ が格子ならば
$\text{length}_R(M/M') < \infty$ であることを見た。逆に
$\text{length}_R(M/M') < \infty$ と仮定する。$R$ はネーターであり、
ある $n$ に対して $\mathfrak m^n M \subset M'$ なので
（補題 \ref{algebra-lemma-length-infinite}）、$M'$ は有限生成である。
従って $M'$ は $V$ の基底を含み、$M'$ は格子である。

\medskip\noindent
(3) の証明。$M$, $M'$ が $V$ 内の格子であると仮定する。
$R$ はネーターなので、$M$ の部分加群 $M \cap M'$ は有限である。
$M$ は格子なので、$V$ の $K$-基底をなす
$x_1, \ldots, x_n \in M$ を見つけられる。$M'$ は格子なので、
$y_i \in M'$、$f_i \in R$ として $x_i = y_i/f_i$ と書ける。
従って $f_ix_i \in M \cap M'$ である。ゆえに $M \cap M'$ も格子である。
$M + M'$ が格子であることは (1) から従う。

\medskip\noindent
(4) は長さの加法性（補題 \ref{algebra-lemma-length-additive}）と完全列
$$
0 \to M'/M \to M''/M \to M''/M' \to 0
$$
から従う。
% [P01-E0578] The frozen missing punctuation after the display is rendered naturally.
(5) は (4) を繰り返し適用すれば従う。
\end{proof}

\begin{definition}
\label{algebra-definition-distance}
$R$ を分数体 $K$ をもつ次元 $1$ のネーター局所整域とし、
$V$ を有限次元 $K$-ベクトル空間とする。
$M$, $M'$ を $V$ 内の二つの格子とする。{\it $M$ と $M'$ の距離}とは、
補題 \ref{algebra-lemma-compare-lattices} の (5) に現れる整数
$$
d(M, M') = \text{length}_R(M/M \cap M') - \text{length}_R(M'/M \cap M')
$$
のことである。
\end{definition}

\noindent
特に、$M' \subset M$ ならば
$d(M, M') = \text{length}_R(M/M')$ である。

\begin{lemma}
\label{algebra-lemma-properties-distance-function}
$R$ を分数体 $K$ をもつ次元 $1$ のネーター局所整域とし、
$V$ を有限次元 $K$-ベクトル空間とする。
この距離関数は、$V$ の任意の三つの格子 $M$, $M'$, $M''$ に対して
$$
d(M, M'') = d(M, M') + d(M', M'')
$$
を満たす。特に $d(M, M') = - d(M', M)$ である。
\end{lemma}

\begin{proof}
省略する。
\end{proof}

\begin{lemma}
\label{algebra-lemma-order-vanishing-determinant}
$R$ を分数体 $K$ をもつ次元 $1$ のネーター局所整域とし、
$V$ を有限次元 $K$-ベクトル空間とする。
$\varphi : V \to V$ を $K$-線形同型とする。
任意の格子 $M \subset V$ に対して
$$
d(M, \varphi(M)) = \text{ord}_R(\det(\varphi))
$$
である。
\end{lemma}

\begin{proof}
整数 $d(M, \varphi(M))$ が格子 $M$ に依存しないことを次のように示せる。
$M'$ をもう一つの格子とする。このとき
\begin{eqnarray*}
d(M, \varphi(M)) & = & d(M, M') + d(M', \varphi(M)) \\
& = & d(M, M') + d(\varphi(M'), \varphi(M)) + d(M', \varphi(M'))
\end{eqnarray*}
である。$\varphi$ は同型なので
$d(\varphi(M'), \varphi(M)) = d(M', M) = -d(M, M')$ であり、従って
$d(M, \varphi(M)) = d(M', \varphi(M'))$ である。
さらに、（補題の）等式の両辺は $\varphi$ について加法的である。すなわち
$$
\text{ord}_R(\det(\varphi \circ \psi))
=
\text{ord}_R(\det(\varphi))
+
\text{ord}_R(\det(\psi))
$$
であり、上で示した独立性により
% [P01-E0579] The frozen missing closing parenthesis is preserved.
\begin{eqnarray*}
d(M, \varphi(\psi((M))) & = &
d(M, \psi(M)) + d(\psi(M), \varphi(\psi(M))) \\
& = & d(M, \psi(M)) + d(M, \varphi(M))
\end{eqnarray*}
でもある。従って $\text{GL}(V)$ の生成元について補題を示せば十分である。
同型 $K^{\oplus n} \cong V$ を選ぶ。このとき
$\text{GL}(V) = \text{GL}_n(K)$ は基本行列 $E$ によって生成される。
$E$ が単位行列である場合、結果は明らかである。
$i \not = j$、$\lambda \in K$、$\lambda \not = 0$ に対して
$E = E_{ij}(\lambda)$、例えば
$$
E_{12}(\lambda) =
\left(
\begin{matrix}
1 & \lambda & \ldots \\
0 & 1 & \ldots \\
\ldots & \ldots & \ldots
\end{matrix}
\right)
$$
ならば、別の基底に関して $E_{12}(1)$ を得る。
$R^{\oplus n} \subset K^{\oplus n}$ を格子としてとれば、
$E = E_{12}(1)$ の場合の結果は明らかである。最後に $a \in K^*$ に対して
$E = E_i(a)$、例えば
$$
E_1(a) =
\left(
\begin{matrix}
a & 0 & \ldots \\
0 & 1 & \ldots \\
\ldots & \ldots & \ldots
\end{matrix}
\right)
$$
ならば、
% [P01-E0580] The frozen undefined rank symbol b is preserved.
$E_1(a)(R^{\oplus b}) = aR \oplus R^{\oplus n - 1}$ であり、
$d(R^{\oplus n}, aR \oplus R^{\oplus n - 1})
= \text{ord}_R(a)$ であることは明らかである。
\end{proof}

\begin{lemma}
\label{algebra-lemma-finite-extension-dim-1}
$A \to B$ を環準同型とし、次を仮定する。
\begin{enumerate}
\item $A$ は次元 $1$ のネーター局所整域である。
\item $A \subset B$ は整域の有限拡大である。
\end{enumerate}
$L/K$ を対応する分数体の有限拡大とする。
$y \in L^*$ とし、$x = \text{Nm}_{L/K}(y)$ とおく。
このとき $B$ は半局所である。
$\mathfrak m_i$、$i = 1, \ldots, n$ を $B$ の極大イデアルとすると
$$
\text{ord}_A(x) =
\sum\nolimits_i
[\kappa(\mathfrak m_i) : \kappa(\mathfrak m_A)]
\text{ord}_{B_{\mathfrak m_i}}(y)
$$
である。ここで $\text{ord}$ は定義 \ref{algebra-definition-ord} のように定める。
\end{lemma}

\begin{proof}
補題 \ref{algebra-lemma-finite-in-codim-1} により環 $B$ は半局所である。
ある $b, b' \in B$ によって $y = b/b'$ と書く。
$\text{ord}$ の加法性と $\text{Nm}$ の乗法性により、$y = b$ または
$y = b'$ について補題を示せば十分である。言い換えれば $y \in B$ と仮定してよい。
この場合、公式の右辺は
$$
\sum [\kappa(\mathfrak m_i) : \kappa(\mathfrak m_A)]
\text{length}_{B_{\mathfrak m_i}}((B/yB)_{\mathfrak m_i})
$$
である。補題 \ref{algebra-lemma-pushdown-module} により、これは
$\text{length}_A(B/yB)$ に等しい。補題 \ref{algebra-lemma-order-vanishing-determinant} により
$$
\text{length}_A(B/yB) = d(B, yB) =
\text{ord}_A(\det\nolimits_K(L \xrightarrow{y} L)).
$$
定義により $x = \text{Nm}_{L/K}(y) = \det\nolimits_K(L \xrightarrow{y} L)$ なので、
補題が示された。
\end{proof}

\section{準有限写像}
\label{algebra-section-quasi-finite}

\noindent
有限型の環準同型 $R \to S$ を考える。
写像 $\Spec(S) \to \Spec(R)$ がある点で準有限であるとは、その点が
そのファイバー内で孤立していることをいう。これはファイバーがその点で
零次元であることを意味する。この節では、この重要だが技術的な概念の
基本的性質を調べる。より進んだ内容は次節で扱う。

\begin{lemma}
\label{algebra-lemma-isolated-point}
$k$ を体とし、$S$ を有限型 $k$-代数とする。
$\mathfrak q$ を $S$ の素イデアルとする。次は同値である。
\begin{enumerate}
\item $\mathfrak q$ は $\Spec(S)$ の孤立点である。
\item $S_{\mathfrak q}$ は $k$ 上有限である。
\item $g \in S$、$g \not\in \mathfrak q$ かつ
$D(g) = \{ \mathfrak q \}$ を満たす元が存在する。
\item $\dim_{\mathfrak q} \Spec(S) = 0$ である。
\item $\mathfrak q$ は $\Spec(S)$ の閉点であり、
$\dim(S_{\mathfrak q}) = 0$ である。
\item 体拡大 $\kappa(\mathfrak q)/k$ は有限であり、
$\dim(S_{\mathfrak q}) = 0$ である。
\end{enumerate}
この場合、ある有限型 $k$-代数 $S'$ に対して
$S = S_{\mathfrak q} \times S'$ である。また、(3) の元 $g$ は
$S_{\mathfrak q} = S_g$ を満たす。
\end{lemma}

\begin{proof}
$\mathfrak q$ が $\Spec(S)$ の孤立点、すなわち
$\{\mathfrak q\}$ が $\Spec(S)$ で開であると仮定する。
$\Spec(S)$ は Jacobson 空間なので（補題
\ref{algebra-lemma-finite-type-field-Jacobson} および
\ref{algebra-lemma-jacobson} 参照）、$\mathfrak q$ は閉点である。
従って $\{\mathfrak q\}$ は $\Spec(S)$ で開かつ閉である。
補題 \ref{algebra-lemma-disjoint-decomposition} および
\ref{algebra-lemma-disjoint-implies-product} により、
$\mathfrak q$ が唯一の点 $\Spec(S_1)$ に対応するように
$S = S_1 \times S_2$ と書ける。従って
$S_1 = S_{\mathfrak q}$ は $k$ 上有限型の零次元環である。
ゆえに、例えば補題 \ref{algebra-lemma-Noether-normalization} により、これは $k$ 上有限である。
これで (1) が (2) を含意することが示された。

\medskip\noindent
$S_{\mathfrak q}$ が $k$ 上有限であると仮定する。
補題 \ref{algebra-lemma-finite-dimensional-algebra} により
$S_{\mathfrak q}$ はアルティン局所環である。従って補題
\ref{algebra-lemma-artinian-finite-length} により
$\Spec(S_{\mathfrak q}) = \{\mathfrak qS_{\mathfrak q}\}$ である。
完全列 $0 \to K \to S \to S_{\mathfrak q}
\to Q \to 0$ を考える。明らかに $K_{\mathfrak q} = Q_{\mathfrak q} = 0$ である。
$S$ はネーターなので $K$ は有限 $S$-加群であり、
$S_{\mathfrak q}$ は $k$ 上有限なので $Q$ も有限 $S$-加群である。
従って $g \in S$、$g \not \in \mathfrak q$ で
$K_g = Q_g = 0$ を満たす元が存在する。ゆえに
$S_{\mathfrak q} = S_g$ かつ $D(g) = \{ \mathfrak q \}$ である。
これで (2) が (3) を含意することが示された。

\medskip\noindent
$D(g) =  \{ \mathfrak q \}$ と仮定する。$\Spec(S)$ の位相の構成により
$D(g)$ は開なので、$\mathfrak q$ は $\Spec(S)$ の孤立点である。
これで (3) が (1) を含意することが示された。
言い換えれば、(1)、(2)、(3) は同値である。

\medskip\noindent
$\dim_{\mathfrak q} \Spec(S) = 0$ と仮定する。これは
$\Spec(S)$ における $\mathfrak q$ のある開近傍が零次元であることを意味する。
従って零次元の $D(g)$ という形の開近傍が存在する。
$S_g$ はネーターなので、$S_g$ はアルティンであり、
$D(g) = \Spec(S_g)$ は有限離散集合であると結論できる
（命題 \ref{algebra-proposition-dimension-zero-ring} 参照）。
従って $\mathfrak q$ は $D(g)$ の孤立点であり、上で (1) と (2) の同値性を
$\mathfrak qS_g \subset S_g$ に適用すると、
$S_{\mathfrak q} = (S_g)_{\mathfrak qS_g}$ は $k$ 上有限である。
ゆえに (4) は (2) を含意する。(1) が (4) を含意することは明らかである。
従って (1) -- (4) はすべて同値である。

\medskip\noindent
補題 \ref{algebra-lemma-dimension-closed-point-finite-type-field} は
(5) $\Rightarrow$ (4) を与える。
補題 \ref{algebra-lemma-dimension-at-a-point-finite-type-field} から
(4) $\Rightarrow$ (6) が従い、補題
\ref{algebra-lemma-finite-residue-extension-closed} から
(6) $\Rightarrow$ (5) が従う。これで (1) -- (6) が同値であることが分かった。

\medskip\noindent
% [P01-E0581] The frozen malformed object-verb order is rendered naturally.
補題末尾の二つの主張は、上の (1)、(2)、(3) の同値性の証明の途中で示した。
\end{proof}

\begin{lemma}
\label{algebra-lemma-isolated-point-fibre}
\begin{slogan}
ファイバー内の孤立点に関する同値条件
\end{slogan}
$R \to S$ を有限型の環準同型とする。
$\mathfrak q \subset S$ を $\mathfrak p \subset R$ の上にある素イデアルとする。
$F = \Spec(S \otimes_R \kappa(\mathfrak p))$ を
$\Spec(S) \to \Spec(R)$ のファイバーとする。注意
\ref{algebra-remark-fundamental-diagram} を参照せよ。
% [P01-E0582] The frozen omitted designation preposition has no Japanese analogue.
$\mathfrak q$ に対応する点を $\overline{\mathfrak q} \in F$ と表す。
% [P01-E0583] The frozen missing list punctuation is rendered naturally.
次は同値である。
\begin{enumerate}
\item $\overline{\mathfrak q}$ は $F$ の孤立点である。
\item $S_{\mathfrak q}/\mathfrak pS_{\mathfrak q}$ は
$\kappa(\mathfrak p)$ 上有限である。
\item $D(g)$ の素イデアルのうち $\mathfrak p$ に写るものが
$\mathfrak q$ だけとなるような $g \in S$、$g \not \in \mathfrak q$ が存在する。
\item $\dim_{\overline{\mathfrak q}}(F) = 0$ である。
\item $\overline{\mathfrak q}$ は $F$ の閉点であり、
$\dim(S_{\mathfrak q}/\mathfrak pS_{\mathfrak q}) = 0$ である。
\item 体拡大 $\kappa(\mathfrak q)/\kappa(\mathfrak p)$ は有限であり、
$\dim(S_{\mathfrak q}/\mathfrak pS_{\mathfrak q}) = 0$ である。
\end{enumerate}
\end{lemma}

\begin{proof}
$S_{\mathfrak q}/\mathfrak pS_{\mathfrak q} =
(S \otimes_R \kappa(\mathfrak p))_{\overline{\mathfrak q}}$ であることに注意する。
さらに $S \otimes_R \kappa(\mathfrak p)$ は
$\kappa(\mathfrak p)$ 上有限型である。
各条件は $\kappa(\mathfrak p)$-代数
$S \otimes_R \kappa(\mathfrak p)$ と素イデアル
$\overline{\mathfrak q}$ に関する補題
\ref{algebra-lemma-isolated-point} の条件に正確に対応するので、同値である。
\end{proof}

\begin{definition}
\label{algebra-definition-quasi-finite}
$R \to S$ を有限型の環準同型とし、$\mathfrak q \subset S$ を素イデアルとする。
\begin{enumerate}
\item 補題 \ref{algebra-lemma-isolated-point-fibre} の同値な条件が満たされるとき、
$R \to S$ は{\it $\mathfrak q$ で準有限}であるという。
\item 環準同型 $A \to B$ が有限型であり、$B$ のすべての素イデアルで
準有限であるとき、これを{\it 準有限}という。
\end{enumerate}
\end{definition}

\begin{lemma}
\label{algebra-lemma-quasi-finite-above-p}
$R \to S$ を有限型の環準同型とし、$\mathfrak p$ を $R$ の素イデアルとする。
このとき次は同値である。
\begin{enumerate}
\item $R \to S$ は、$\mathfrak p$ の上にある $S$ のすべての素イデアルで準有限である。
\item $S \otimes_R \kappa(\mathfrak p)$ は有限
$\kappa(\mathfrak p)$-代数である。
\item $\Spec(S \otimes_R \kappa(\mathfrak p))$ は有限集合である。
\end{enumerate}
\end{lemma}

\begin{proof}
条件 (1) は、$\Spec(S \otimes_R \kappa(\mathfrak p))$ の各点が開であり、
その位相が離散であることをいう。従って (1)、(2)、(3) の同値性は補題
\ref{algebra-lemma-finite-type-algebra-finite-nr-primes} から従う。
\end{proof}

\begin{lemma}
\label{algebra-lemma-quasi-finite}
$R \to S$ を有限型の環準同型とする。このとき $R \to S$ が準有限であることと、
すべての素イデアル $\mathfrak p \subset R$ に対して環
$S \otimes_R \kappa(\mathfrak p)$ が $\kappa(\mathfrak p)$ 上有限であることは同値である。
\end{lemma}

\begin{proof}
より一般的な補題 \ref{algebra-lemma-quasi-finite-above-p} から直ちに従う。
\end{proof}

\begin{lemma}
\label{algebra-lemma-quasi-finite-local}
$R \to S$ を有限型の環準同型とし、$\mathfrak q \subset S$ を
$\mathfrak p \subset R$ の上にある素イデアルとする。
$f \in R$、$f \not \in \mathfrak p$ および
$g \in S$、$g \not \in \mathfrak q$ とする。
このとき $R \to S$ が $\mathfrak q$ で準有限であることと、
$R_f \to S_{fg}$ が $\mathfrak qS_{fg}$ で準有限であることは同値である。
\end{lemma}

\begin{proof}
$\Spec(S_{fg}) \to \Spec(R_f)$ のファイバーは、
$\Spec(S) \to \Spec(R)$ のファイバーの開部分集合と同相である。
従って補題は、補題 \ref{algebra-lemma-isolated-point-fibre} の同値条件 (1) から従う。
\end{proof}

\begin{lemma}
\label{algebra-lemma-four-rings}
環と表示された素イデアルの可換図式
$$
\xymatrix{
S \ar[r] & S' & &
\mathfrak q \ar@{-}[r] & \mathfrak q' \\
R \ar[u] \ar[r] &  R' \ar[u] & &
\mathfrak p \ar@{-}[r] \ar@{-}[u] & \mathfrak p' \ar@{-}[u]
}
$$
が与えられたとする。
% [P01-E0584] The frozen missing finite verb is rendered naturally.
$R \to S$ は有限型であり、$S \otimes_R R' \to S'$ は全射であると仮定する。
$R \to S$ が $\mathfrak q$ で準有限ならば、
$R' \to S'$ は $\mathfrak q'$ で準有限である。
\end{lemma}

\begin{proof}
補題 \ref{algebra-lemma-isolated-point} のように、
$S_1$ が $\kappa(\mathfrak p)$ 上有限で、$\mathfrak q$ が $S_1$ の点に
対応するように
$S \otimes_R \kappa(\mathfrak p) = S_1 \times S_2$ と書く。
この積分解は任意の $S \otimes_R \kappa(\mathfrak p)$-代数に対応する
積分解を誘導する。特に
$S' \otimes_{R'} \kappa(\mathfrak p') = S'_1 \times S'_2$ を得る。
$S \otimes_R R' \to S'$ は全射なので、標準写像
$(S \otimes_R \kappa(\mathfrak p)) \otimes_{\kappa(\mathfrak p)}
\kappa(\mathfrak p') \to S' \otimes_{R'} \kappa(\mathfrak p')$
は全射であり、従って
$S_i \otimes_{\kappa(\mathfrak p)} \kappa(\mathfrak p') \to S'_i$
も全射である。このことから $S'_1$ は $\kappa(\mathfrak p')$ 上有限である。
写像 $S' \otimes_{R'} \kappa(\mathfrak p') \to
\kappa(\mathfrak q')$ は $S_1'$ を経由する
（すなわち因子 $S_2'$ を零化する）。これは、写像
$S \otimes_R \kappa(\mathfrak p) \to
\kappa(\mathfrak q)$ が $S_1$ を経由する
（すなわち因子 $S_2$ を零化する）からである。
従ってファイバーの非交和分解
$\Spec(S' \otimes_{R'} \kappa(\mathfrak p'))
= \Spec(S_1') \amalg \Spec(S_2')$
において、$\mathfrak q'$ は $\Spec(S_1')$ の点に対応する。
補題 \ref{algebra-lemma-spec-product} を参照せよ。
% [P01-E0585] The frozen missing article has no Japanese analogue.
$S_1'$ は体上有限なのでアルティン環であり、従って
$\Spec(S_1')$ は有限離散集合である
（命題 \ref{algebra-proposition-dimension-zero-ring} 参照）。
ゆえに望むように $\mathfrak q'$ はそのファイバー内で孤立している。
\end{proof}

\begin{lemma}
\label{algebra-lemma-quasi-finite-composition}
準有限な環準同型の合成は準有限である。
\end{lemma}

\begin{proof}
$A \to B$ と $B \to C$ が準有限な環準同型であると仮定する。
補題 \ref{algebra-lemma-compose-finite-type} により $A \to C$ は有限型である。
$\mathfrak r \subset C$ を、$\mathfrak q \subset B$ および
$\mathfrak p \subset A$ の上にある $C$ の素イデアルとする。
% [P01-E0586] The frozen redundant conditional particle is rendered naturally.
$A \to B$ と $B \to C$ はそれぞれ $\mathfrak q$ と $\mathfrak r$ で
準有限なので、$D(b)$ の素イデアルのうち $\mathfrak p$ に写るものが
$\mathfrak q$ だけとなり、同様に $D(c)$ の素イデアルのうち
$\mathfrak q$ に写るものが $\mathfrak r$ だけとなるような
$b \in B$ と $c \in C$ が存在する。
$c' \in C$ が $b \in B$ の像ならば、$D(cc')$ の素イデアルのうち
$\mathfrak p$ に写るものは $\mathfrak r$ だけである。
従って $A \to C$ は $\mathfrak r$ で準有限である。
\end{proof}

\begin{lemma}
\label{algebra-lemma-quasi-finite-base-change}
$R \to S$ を有限型の環準同型とし、$R \to R'$ を任意の環準同型とする。
$S' = R' \otimes_R S$ とおく。
\begin{enumerate}
\item 集合
$\{\mathfrak q' \mid R' \to S' \text{ quasi-finite at }\mathfrak q'\}$
は、標準写像 $\Spec(S') \to \Spec(S)$ による $\Spec(S)$ の対応する集合の逆像である。
\item $\Spec(R') \to \Spec(R)$ が全射ならば、$R \to S$ が準有限であることと
$R' \to S'$ が準有限であることは同値である。
\item 準有限な環準同型の任意の基底変換は準有限である。
\end{enumerate}
\end{lemma}

\begin{proof}
$\mathfrak p' \subset R'$ を $\mathfrak p \subset R$ の上にある素イデアルとする。
このときファイバー環 $S' \otimes_{R'} \kappa(\mathfrak p')$ は、
ファイバー環 $S \otimes_R \kappa(\mathfrak p)$ を体拡大
$\kappa(\mathfrak p) \to \kappa(\mathfrak p')$ により基底変換したものである。
従って最初の主張は、体拡大による次元の不変性
（補題 \ref{algebra-lemma-dimension-at-a-point-preserved-field-extension}）と補題
\ref{algebra-lemma-isolated-point} から従う。
準有限写像の基底変換に関する安定性は、このことと有限型性の基底変換に関する
安定性から従う。第二の主張は、仮定により任意の素イデアル
$\mathfrak q \subset S$ の上にある素イデアル
$\mathfrak q' \subset S'$ を見つけられることから従う。
\end{proof}

\begin{lemma}
\label{algebra-lemma-quasi-finite-permanence}
$A \to B$ および $B \to C$ を、$A \to C$ が有限型となる環準同型とする。
$\mathfrak r$ を、$\mathfrak q \subset B$ および
$\mathfrak p \subset A$ の上にある $C$ の素イデアルとする。
$A \to C$ が $\mathfrak r$ で準有限ならば、
$B \to C$ は $\mathfrak r$ で準有限である。
\end{lemma}

\begin{proof}
$B \to C$ は有限型なので（補題 \ref{algebra-lemma-compose-finite-type}）、
この主張には意味がある。
補題 \ref{algebra-lemma-isolated-point-fibre} の特徴付け (3) を用いる。
$A \to C$ が $\mathfrak r$ で準有限ならば、ある $c \in C$ が存在して
$$
\{\mathfrak r' \subset C \text{ lying over }\mathfrak p\} \cap D(c) =
\{\mathfrak{r}\}.
$$
となる。$\mathfrak q$ の上にある素イデアル $\mathfrak r' \subset C$ は、
$\mathfrak p$ の上にある素イデアル $\mathfrak r' \subset C$ の部分集合をなすので、
$B \to C$ は $\mathfrak r$ で準有限である。
\end{proof}

\begin{lemma}
\label{algebra-lemma-generically-finite}
$R \to S$ を有限型の環準同型とし、
$\mathfrak p \subset R$ を極小素イデアルとする。
$\mathfrak p$ の上にある $S$ の素イデアルが高々有限個であると仮定する。
このとき $g \in R$、$g \not \in \mathfrak p$ であって、
環準同型 $R_g \to S_g$ が有限となるような元が存在する。
\end{lemma}

\begin{proof}
$x_1, \ldots, x_n$ を $R$ 上の $S$ の生成元とする。
補題 \ref{algebra-lemma-quasi-finite-above-p} により、仮定は
$S \otimes_R \kappa(\mathfrak p) = S_{\mathfrak p}/\mathfrak pS_{\mathfrak p}$
が有限 $\kappa(\mathfrak p)$-代数であることを意味する。
従って、$P_i(x_i)$ が
$S_{\mathfrak p}/\mathfrak pS_{\mathfrak p}$ で零に写るようなモニック多項式
$P_i \in R_{\mathfrak p}[X]$ を見つけられる。
$\mathfrak p$ は極小素イデアルなので、
$\mathfrak pR_{\mathfrak p}$ は局所冪零イデアルである
（補題 \ref{algebra-lemma-minimal-prime-reduced-ring} 参照）。
従って $\mathfrak pS_{\mathfrak p}$ も局所冪零イデアルである
（補題 \ref{algebra-lemma-locally-nilpotent} 参照）。
% [P01-E0587] Both frozen occurrences of the undefined unindexed P are preserved.
ゆえに、$S_{\mathfrak p}$ において $P(x_i)^{e_i} = 0$ となる
$e_i \geq 1$ が存在する。すべての $i$ に対して $P_i$ の係数が
$R[1/g_1]$ に属するような $g_1 \in R$、$g_1 \not \in \mathfrak p$ をとる。
次に、$S_{g_1g_2}$ において $P(x_i)^{e_i} = 0$ となるような
$g_2 \in R$、$g_2 \not \in \mathfrak p$ をとる。
$g = g_1g_2$ とおけばよい。
\end{proof}

\section{Zariski の主定理}
\label{algebra-section-Zariski}

\noindent
この節の目的は Zariski の主定理の代数版を証明することである。
この定理は、Stacks project の後の諸章におけるスキームとスキームの射の理論の
多くの展開の基礎となる。

\medskip\noindent
$R \to S$ を有限型の環準同型とする。この節の目標は、写像が準有限である
$\Spec(S)$ の点の集合が{\it 開}であることを示すことである
（定理 \ref{algebra-theorem-main-theorem}）。
実際、ある有限環準同型 $R \to S'$ が存在し、ある意味で $S/R$ の準有限軌跡が
$\Spec(S')$ の開部分集合となることが分かる
（ただし、この代数の章ではスキームの言葉を展開しないため、それは証明しない。
$R \to S$ が準有限の場合については補題
\ref{algebra-lemma-quasi-finite-open-integral-closure} を参照せよ）。
これらの主張の証明はいくぶん技巧的であり、整拡大と有限環拡大に関する
長い補題列を用いる。この内容は \cite{Henselian} および \cite{Peskine} にある。
Thierry Coquand のノートも有用であった。

\begin{lemma}
\label{algebra-lemma-make-integral-trivial}
$\varphi : R \to S$ を環準同型とする。$t \in S$ が関係式
$\varphi(a_0) + \varphi(a_1)t + \ldots + \varphi(a_n) t^n = 0$
を満たすと仮定する。このとき $\varphi(a_n)t$ は $R$ 上整である。
\end{lemma}

\begin{proof}
実際、等式\hfill\break
$\varphi(a_0)\allowbreak + \varphi(a_1)t\allowbreak + \ldots\allowbreak +
\varphi(a_n) t^n\allowbreak = 0$
に $\varphi(a_n)^{n-1}$ を掛け、
$\varphi(a_0 a_n^{n-1})\allowbreak +
\varphi(a_1 a_n^{n-2}) (\varphi(a_n)t)\allowbreak +
\ldots +
(\varphi(a_n) t)^n = 0$
と書けばよい。
\end{proof}

\noindent
次の補題は、ある意味でこの節の鍵となる補題である。

\begin{lemma}
\label{algebra-lemma-make-integral-trick}
$R$ を環とし、$\varphi : R[x] \to S$ を環準同型とする。
$t \in S$ とする。(a) $t$ は $R[x]$ 上整であり、(b) モニックな
$p \in R[x]$ が存在して $t \varphi(p) \in \Im(\varphi)$ であると仮定する。
このとき、$t - \varphi(q)$ が $R$ 上整となるような $q \in R[x]$ が存在する。
\end{lemma}

\begin{proof}
ある $r \in R[x]$ に対して $t \varphi(p) = \varphi(r)$ と書く。
ユークリッド除法により、$q, r' \in R[x]$、$\deg(r') < \deg(p)$ として
$r = qp + r'$ と書く。$t$ を $t - \varphi(q)$ で置き換えても
$R[x]$ 上整であり、$t \varphi(p) = \varphi(r')$ を得る。
環 $S_t$ では、これを $\varphi(p) - (1/t) \varphi(r') = 0$ と書ける。
これは、$\varphi(x)$ が局所化 $S_t$ の元で、部分環
$\varphi(R)[1/t] \subset S_t$ 上整であることを含意する。
一方、$t$ は部分環 $\varphi(R)[\varphi(x)] \subset S$ 上整である。
これらを合わせると、補題 \ref{algebra-lemma-integral-transitive} により、
$t$ は部分環 $\varphi(R)[1/t] \subset S_t$ 上整であると結論できる。
言い換えれば、$r_{i, j} \in R$ を用いて $S_t$ における形
$$
t^d + \sum\nolimits_{i < d}
\left(\sum\nolimits_{j = 0, \ldots, n_i} \varphi(r_{i, j})/t^j\right) t^i = 0
$$
の等式が存在する。これは、十分大きいある $N$ に対して $S$ で
$t^{d + N}\allowbreak +
\sum_{i < d} \sum_{j = 0, \ldots,\allowbreak n_i} \varphi(r_{i, j}) t^{i + N - j}\allowbreak = 0$
であることを意味する。言い換えれば $t$ は $R$ 上整である。
\end{proof}

\begin{lemma}
\label{algebra-lemma-combine-lemmas}
$R$ を環とし、$\varphi : R[x] \to S$ を環準同型とする。
$t \in S$ とし、$t$ は $R[x]$ 上整であると仮定する。
$p \in R[x]$、$p = a_0 + a_1x + \ldots +
a_k x^k$ が $t \varphi(p) \in \Im(\varphi)$ を満たすとする。
このとき、$\varphi(a_k)^n t - \varphi(q)$ が $R$ 上整となる
$q \in R[x]$ と $n \geq 0$ が存在する。
\end{lemma}

\begin{proof}
$R'$ と $S'$ を、それぞれ元 $a_k$ における $R$ と $S$ の局所化とする。
$\varphi' : R'[x] \to S'$ を $\varphi$ の局所化とし、
$t' \in S'$ を $t$ の像とする。$p' = p/a_k \in R'[x]$ とおく。
$t \varphi(p) \in \Im(\varphi)$ なので
$t' \varphi'(p') \in \Im(\varphi')$ である。
$p'$ はモニックなので、補題 \ref{algebra-lemma-make-integral-trick} により
$t' - \varphi'(q')$ が $R'$ 上整となる $q' \in R'[x]$ が存在する。
$a_k^n q'$ が $q$ の像となるような $n \geq 0$ と $q \in R[x]$ を選べる。
このとき $\varphi(a_k)^n t - \varphi(q)$ は $S$ の元で、その $S'$ における像は
$R'$ 上整である。補題 \ref{algebra-lemma-integral-closure-localize} により、
$\varphi(a_k)^m(\varphi(a_k)^n t - \varphi(q))$ が $R$ 上整となる
$m \geq 0$ が存在する。従って望むように
$\varphi(a_k)^{m + n}t - \varphi(a_k^m q)$ は $R$ 上整である。
\end{proof}

\begin{situation}
\label{algebra-situation-one-transcendental-element}
$R$ を環とし、$\varphi : R[x] \to S$ を有限とする。
$$
J = \{ g \in S \mid gS \subset \Im(\varphi)\}
$$
を $\varphi$ の「導手イデアル」とする。
% [P01-E0588] The frozen missing finite verb is rendered naturally.
$\varphi(R) \subset S$ は $S$ の中で整閉であると仮定する。
\end{situation}

\begin{lemma}
\label{algebra-lemma-leading-coefficient-in-J}
% [P01-E0589] The frozen standalone situation fragment is joined naturally.
状況 \ref{algebra-situation-one-transcendental-element} のもとで、
$u \in S$、$a_0, \ldots, a_k \in R$ とし、
$u \varphi(a_0 + a_1x + \ldots + a_k x^k) \in J$ と仮定する。
このとき $u \varphi(a_k)^m \in J$ となる $m \geq 0$ が存在する。
\end{lemma}

\begin{proof}
$S$ が $R[x]$-加群として $t_1, \ldots, t_n$ で生成されると仮定する。
この場合
$J = \{ g \in S \mid gt_i \in \Im(\varphi)\text{ for all }i\}$ である。
各元 $u t_i$ は $R[x]$ 上整であることに注意する
（補題 \ref{algebra-lemma-finite-is-integral} 参照）。
$\varphi(a_0 + a_1x + \ldots + a_k x^k) u t_i \in
\Im(\varphi)$ である。補題 \ref{algebra-lemma-combine-lemmas} により、各 $i$ について
$\varphi(a_k^{n_i}) u t_i - \varphi(q_i)$ が $R$ 上整となるような整数 $n_i$ と
元 $q_i \in R[x]$ が存在する。仮定によりこの元は $\varphi(R)$ に属し、
従って $\varphi(a_k^{n_i}) u t_i \in \Im(\varphi)$ である。
ゆえに $m = \max\{n_1, \ldots, n_n\}$ とすればよい。
\end{proof}

\begin{lemma}
\label{algebra-lemma-all-coefficients-in-J}
状況 \ref{algebra-situation-one-transcendental-element} のもとで、
$u \in S$、$a_0, \ldots, a_k \in R$ とし、
$u \varphi(a_0 + a_1x + \ldots + a_k x^k) \in \sqrt{J}$ と仮定する。
このとき、すべての $i$ に対して $u \varphi(a_i) \in \sqrt{J}$ である。
\end{lemma}

\begin{proof}
補題の仮定のもとで、ある $n \geq 1$ に対して
$u^n \varphi(a_0 + a_1x + \ldots + a_k x^k)^n \in J$ である。
補題 \ref{algebra-lemma-leading-coefficient-in-J} により、ある $m \geq 1$ に対して
$u^n \varphi(a_k^{nm}) \in J$ であると分かる。
従って $u \varphi(a_k) \in \sqrt{J}$ であり、
$u \varphi(a_0 + a_1x + \ldots + a_k x^k) - u \varphi(a_k x^k) =
u \varphi(a_0 + a_1x + \ldots + a_{k-1} x^{k-1}) \in \sqrt{J}$
である。$k$ に関する帰納法で証明が終わる。
\end{proof}

\noindent
この補題は次の定義を示唆する。

\begin{definition}
\label{algebra-definition-strongly-transcendental}
環の包含 $R \subset S$ と元 $x \in S$ が与えられたとする。
$u \in S$ と $a_i \in R$ に対して
$u(a_0 + a_1 x + \ldots + a_k x^k) = 0$ ならば、すべての $i$ について
$ua_i = 0$ となるとき、$x$ は{\it $R$ 上強超越的}であるという。
\end{definition}

\noindent
$S$ が整域ならば、これは $S$ の分数体の元としての $x$ が
$R$ の分数体上超越的であることと同じである。

\begin{lemma}
\label{algebra-lemma-reduced-strongly-transcendental-minimal-prime}
$R \subset S$ を被約環の包含とし、$x \in S$ は $R$ 上強超越的であると仮定する。
$\mathfrak q \subset S$ を極小素イデアルとし、
$\mathfrak p = R \cap \mathfrak q$ とおく。
このとき $S/\mathfrak q$ における $x$ の像は、部分環
$R/\mathfrak p$ 上強超越的である。
\end{lemma}

\begin{proof}
$u(a_0 + a_1x + \ldots + a_k x^k) \in \mathfrak q$ と仮定する。
補題 \ref{algebra-lemma-minimal-prime-reduced-ring} により局所環
$S_{\mathfrak q}$ は体であり、従って
$u(a_0 + a_1x + \ldots + a_k x^k) $ は $S_{\mathfrak q}$ で零である。
ゆえに、ある $u' \in S$、$u' \not\in \mathfrak q$ に対して
$uu'(a_0 + a_1x + \ldots + a_k x^k) = 0$ である。
$x$ は $R$ 上強超越的なので、すべての $i$ に対して $uu'a_i = 0$ を得る。
これは $ua_i \in \mathfrak q$ を含意する。
\end{proof}

\begin{lemma}
\label{algebra-lemma-domains-transcendental-not-quasi-finite}
$R\subset S$ を整域の包含とし、$x \in S$ とする。
$x$ が $R$ 上（強）超越的であり、$S$ が $R[x]$ 上有限であると仮定する。
このとき $R \to S$ は $S$ のどの素イデアルでも準有限でない。
\end{lemma}

\begin{proof}
まず $R$ が正規である場合を考える。定義
\ref{algebra-definition-ring-normal} を参照せよ。
補題 \ref{algebra-lemma-polynomial-ring-normal} により $R[x]$ は正規である。
素イデアル $\mathfrak q \subset S$ をとり、
$\mathfrak p = R \cap \mathfrak q$ とおく。
拡大 $\kappa(\mathfrak p)
\subset \kappa(\mathfrak q)$ が有限であると仮定する。
$R \to S$ が $\mathfrak q$ で準有限ならばそうなる。
$\mathfrak r = R[x] \cap \mathfrak q$ とおく。
$\kappa(\mathfrak p)
\subset \kappa(\mathfrak r) \subset \kappa(\mathfrak q)$ なので、
拡大 $\kappa(\mathfrak p)
\subset \kappa(\mathfrak r)$ も有限である。
従って包含 $\mathfrak r \supset \mathfrak p R[x]$ は狭義である。
$R[x] \subset S$ に対する下降定理
（命題 \ref{algebra-proposition-going-down-normal-integral} 参照）により、
素イデアル $\mathfrak pR[x]$ の上にある
$\mathfrak q' \subset \mathfrak q$ を見つけられる。
従ってファイバー $\Spec(S \otimes_R \kappa(\mathfrak p))$ は
$\mathfrak q$ と異なる点、すなわち $\mathfrak q'$ を含み、
その閉包は $\mathfrak q$ を含む。ゆえに $\mathfrak q$ はファイバー内で孤立していない。

\medskip\noindent
$R$ が正規でない場合、$R$ の分数体 $K$ における整閉包を $R'$ とし、
$R \subset R' \subset K$ とする。$S$ の分数体 $L$ の部分環で、
$R'$ と $S$ により生成されるものを $S'$ とし、
$S \subset S' \subset L$ とする。構成により写像
$S \otimes_R R' \to S'$ は全射である。これは
$R'[x] \subset S'$ が有限であることを含意する。
また、写像 $S \subset S'$ は $\Spec$ 上の全射を誘導する
（補題 \ref{algebra-lemma-integral-overring-surjective} 参照）。
補題 \ref{algebra-lemma-four-rings} と上で論じた正規の場合から結論を得る。
\end{proof}

\begin{lemma}
\label{algebra-lemma-reduced-strongly-transcendental-not-quasi-finite}
$R \subset S$ を被約環の包含とする。
% [P01-E0590] The frozen ungrammatical finite verb is rendered naturally.
$x \in S$ は $R$ 上強超越的であり、$S$ は $R[x]$ 上有限であると仮定する。
このとき $R \to S$ は $S$ のどの素イデアルでも準有限でない。
\end{lemma}

\begin{proof}
$\mathfrak q \subset S$ を任意の素イデアルとし、
極小素イデアル $\mathfrak q' \subset \mathfrak q$ を選ぶ。
補題 \ref{algebra-lemma-reduced-strongly-transcendental-minimal-prime} および
\ref{algebra-lemma-domains-transcendental-not-quasi-finite} により、拡大
$R/(R \cap \mathfrak q') \subset
S/\mathfrak q'$ は $\mathfrak q$ に対応する素イデアルで準有限でない。
補題 \ref{algebra-lemma-four-rings} により、拡大 $R \to S$ は
$\mathfrak q$ で準有限でない。
\end{proof}

\begin{lemma}
\label{algebra-lemma-quasi-finite-monogenic}
$R$ を環とし、$S = R[x]/I$ とする。
$\mathfrak q \subset S$ を素イデアルとし、
$R \to S$ は $\mathfrak q$ で準有限であると仮定する。
$S' \subset S$ を $S$ における $R$ の整閉包とする。
このとき $g \in S'$、$g \not\in \mathfrak q$ かつ
$S'_g \cong S_g$ を満たす元が存在する。
\end{lemma}

\begin{proof}
$\mathfrak p$ を $\Spec(R)$ における $\mathfrak q$ の像とする。
$f \in I$、$f = a_nx^n + \ldots + a_0$ で、ある $i$ に対して
$a_i \not \in \mathfrak p$ となるものが存在する。
実際、そうでなければファイバー環
$S \otimes_R \kappa(\mathfrak p)$ は $\kappa(\mathfrak p)[x]$ となり、
写像は $\mathfrak p$ の上にあるどの素イデアルでも準有限でない。
従って、$b_j \in S'$、$j = 0, \ldots, m$ で、ある $j$ に対して
$b_j \not \in \mathfrak q \cap S'$ となる関係式
$b_m x^m + \ldots + b_0 = 0$ が存在する。
$m$ に関する帰納法で補題を証明する。
% [P01-E0591] The frozen duplicated copula is rendered naturally.
基底の場合 $m = 0$ は空虚である
（$b_0 = 0$ と $b_0 \not \in \mathfrak q$ は矛盾するため）。

\medskip\noindent
$b_m \not \in \mathfrak q$ の場合。この場合 $x$ は $S'_{b_m}$ 上整であり、
実際、補題 \ref{algebra-lemma-make-integral-trivial} により $b_mx \in S'$ である。
従って単射 $S'_{b_m} \to S_{b_m}$ は全射でもあり、望む同型である。

\medskip\noindent
$b_m \in \mathfrak q$ の場合。この場合も補題
\ref{algebra-lemma-make-integral-trivial} により $b_mx \in S'$ である。
$b'_{m - 1} = b_mx + b_{m - 1}$ とおく。このとき
$$
b'_{m - 1}x^{m - 1} + b_{m - 2}x^{m - 2} + \ldots + b_0 = 0
$$
である。
% [P01-E0592] The frozen missing punctuation after the display is rendered naturally.
$b'_{m - 1}$ は $S' \cap \mathfrak q$ を法として $b_{m - 1}$ と合同なので、
$b'_{m - 1}, b_{m - 2}, \ldots, b_0$ の少なくとも一つが
$S' \cap \mathfrak q$ に属さないという性質は保たれる。
従って $m$ に関する帰納法で結論を得る。
\end{proof}

\begin{theorem}[Zariski の主定理]
\label{algebra-theorem-main-theorem}
$R$ を環とする。
% [P01-E0593] The frozen grammatical classification of the map as an R-algebra is preserved.
$R \to S$ を有限型 $R$-代数とする。
$S' \subset S$ を $S$ における $R$ の整閉包とし、
$\mathfrak q \subset S$ を $S$ の素イデアルとする。
$R \to S$ が $\mathfrak q$ で準有限ならば、
$g \in S'$、$g \not \in \mathfrak q$ かつ
$S'_g \cong S_g$ を満たす元が存在する。
\end{theorem}

\begin{proof}
$S$ が $x_1, \ldots, x_n$ で生成される $R$-部分代数上有限となるような
有限個の元 $x_1, \ldots, x_n \in S$ が存在する
（例えば $R$ 上の $S$ の生成元をとればよい）。
% [P01-E0594] The split English compound is rendered by the standard Japanese term.
% [P01-E0595] The frozen result-class wording is preserved.
そのような数の最小値 $n$ に関する帰納法で命題を証明する。

\medskip\noindent
$n = 0$ の場合は自明である。この場合 $S' = S$ だからである
（補題 \ref{algebra-lemma-finite-is-integral} 参照）。

\medskip\noindent
$n = 1$ の場合。$R$ を $S$ におけるその整閉包で置き換えてよい
（補題 \ref{algebra-lemma-quasi-finite-permanence} により、$R \to S$ は依然として
$\mathfrak q$ で準有限である）。従って $R \subset S$ は $S$ の中で整閉、
すなわち $R = S'$ であると仮定してよい。
写像 $\varphi : R[x] \to S$、$x \mapsto x_1$ を考える
（後で $\varphi$ は単射でないことが分かる）。
仮定により $\varphi$ は有限である。従って状況
\ref{algebra-situation-one-transcendental-element} にある。
$J \subset S$ を状況 \ref{algebra-situation-one-transcendental-element} で定めた
「導手イデアル」とする。図式
$$
\xymatrix{
R[x] \ar[r] & S \ar[r] & S/\sqrt{J} & R/(R \cap \sqrt{J})[x] \ar[l]
\\
& R \ar[lu] \ar[r] \ar[u] & R/(R \cap \sqrt{J}) \ar[u] \ar[ru] &
}
$$
を考える。補題 \ref{algebra-lemma-all-coefficients-in-J} により、
商 $S/\sqrt{J}$ における $x$ の像は
$R/ (R \cap \sqrt{J})$ 上強超越的である。
従って補題 \ref{algebra-lemma-reduced-strongly-transcendental-not-quasi-finite} により、
環準同型 $R/ (R \cap \sqrt{J}) \to S/\sqrt{J}$ は
$S/\sqrt{J}$ のどの素イデアルでも準有限でない。
補題 \ref{algebra-lemma-four-rings} により、
$\mathfrak q$ は $V(J) \subset \Spec(S)$ に属さないと分かる。
従って $s \in J$、$s \not\in \mathfrak q$ となる元が存在する。
$J$ の定義により、ある多項式 $f \in R[x]$ に対して
$s = \varphi(f)$ と書ける。
$I = \Ker(\varphi : R[x] \to S)$ とおく。
$\varphi(f) \in J$ なので $(R[x]/I)_f \cong S_{\varphi(f)}$ である。
また $s \not \in \mathfrak q$ は
$f \not \in \varphi^{-1}(\mathfrak q)$ を意味する。
従って $\varphi^{-1}(\mathfrak q)/I$ は $R[x]/I$ の素イデアルであり、
そこで $R \to R[x]/I$ は準有限である
（補題 \ref{algebra-lemma-quasi-finite-local} 参照）。
$R$ は $S$ の中で整閉なので、$R$ は $R[x]/I$ の中でも整閉である。
補題 \ref{algebra-lemma-quasi-finite-monogenic} により、
$h \in R$、$h \not \in R \cap \mathfrak q$ で
$R_h \cong (R[x]/I)_h$ を満たす元が存在する。
従って $(R[x]/I)_{fh} = S_{\varphi(fh)}$ は、ある
$h' \in R$、$h' \not \in \mathfrak q$ に対する $R$ の単項局所化
$R_{h'}$ と同型である。

\medskip\noindent
$n > 1$ の場合。$R[x_1, \ldots, x_{n-1}]$ の $S$ における整閉包である
部分環 $R' \subset S$ を考える。補題
\ref{algebra-lemma-quasi-finite-permanence} により、拡大 $S/R'$ は
$\mathfrak q$ で準有限である。また $S$ は $R'[x_n]$ 上有限であることに注意する。
上の $n = 1$ の場合により、$g' \in R'$、$g' \not \in \mathfrak q$ かつ
$(R')_{g'} \cong S_{g'}$ となる元が存在する。
この時点では $R'$ は $R$ 上有限型とは限らないため、
$R \to R'$ に帰納法を適用できない。
$S$ は $R$ 上有限生成なので、特に $(R')_{g'}$ は $R$ 上有限生成である。
$y_i \in R'$ とし、元
$g'$ および $y_1/(g')^{n_1}, \ldots, y_N/(g')^{n_N}$ が
$R$ 上 $(R')_{g'}$ を生成するとする。
$R''$ を、$x_1, \ldots, x_{n-1}, y_1, \ldots, y_N, g'$ で生成される
$R'$ の $R$-部分代数とする。これは
$(R'')_{g'} \cong S_{g'}$ という性質をもつ。
% [P01-E0596] The frozen sentence fragments are rendered as complete causal clauses.
全射性は $y_i$ の選び方から従い、単射性は
$R'' \subset R'$ であり局所化が完全であることから従う。
$R'$ の選び方により $R''$ は
$R[x_1, \ldots, x_{n-1}]$ 上有限であることに注意する
（補題 \ref{algebra-lemma-characterize-integral} 参照）。
$\mathfrak q'' = R'' \cap \mathfrak q$ とおく。
$(R'')_{\mathfrak q''} = S_{\mathfrak q}$ なので、
$R \to R''$ は $\mathfrak q''$ で準有限である
（補題 \ref{algebra-lemma-isolated-point-fibre} 参照）。
$R \to R''$、$\mathfrak q''$ および
$x_1, \ldots, x_{n-1} \in R''$ に帰納法の仮定を適用すると、
$R$ 上整な部分環 $R''' \subset R''$ と、
$g'' \in R'''$、$g'' \not \in \mathfrak q''$ かつ
$(R''')_{g''} \cong (R'')_{g''}$ となる元を得る。
$(R'')_{g''}$ における $g'$ の像を、ある $g''' \in  R'''$ により
$g'''/(g'')^n$ と書く。$g = g''g''' \in R'''$ とおく。
このとき $g \not\in
\mathfrak q$ かつ $(R''')_g \cong S_g$ であることは明らかである。
構成により $R''' \subset S'$ なので、望むように $S'_g \cong S_g$ でもある。
\end{proof}

\begin{lemma}
\label{algebra-lemma-quasi-finite-open}
$R \to S$ を有限型の環準同型とする。
$S/R$ が準有限となる $\Spec(S)$ の点 $\mathfrak q$ の集合は
$\Spec(S)$ で開である。
\end{lemma}

\begin{proof}
$\mathfrak q \subset S$ を環準同型が準有限となる点とする。
定理 \ref{algebra-theorem-main-theorem} により、整環準同型
$R \to S'$、$S' \subset S$ と、$g \in S'$、$g\not \in \mathfrak q$ かつ
$S'_g \cong S_g$ を満たす元が存在する。
$S$、従って $S_g$ は $R$ 上有限型なので、$g, y_1, \ldots, y_N$ で生成される
部分 $R$-代数を $S'' \subset S'$ とすると
$S''_g \cong S_g$ となるような有限個の元
$y_1, \ldots, y_N$ を $S'$ から選べる。
$S''$ は $R$ 上有限なので（補題 \ref{algebra-lemma-characterize-integral} 参照）、
$S''$ は $R$ 上準有限である（補題 \ref{algebra-lemma-quasi-finite} 参照）。
これは $S''_g$ が $R$ 上準有限であることを容易に含意する。
例えば、ある素イデアルで準有限であるという性質はその素イデアルでの
局所環だけに依存するからである。従って $S_g$ は $R$ 上準有限である。
同じ理由により、$R \to S$ は $D(g)$ に属する $S$ のすべての素イデアルで
準有限である。
\end{proof}

\begin{lemma}
\label{algebra-lemma-quasi-finite-open-integral-closure}
$R \to S$ を有限型の環準同型とし、$S$ は $R$ 上準有限であると仮定する。
$S' \subset S$ を $S$ における $R$ の整閉包とする。このとき
\begin{enumerate}
\item $\Spec(S) \to \Spec(S')$ はある開部分集合への同相である。
\item $g \in S'$ で $D(g)$ が写像の像に含まれるならば
$S'_g \cong S_g$ である。
\item 有限 $R$-代数 $S'' \subset S'$ で、環準同型
$S'' \to S$ に対して (1) と (2) が成り立つものが存在する。
\end{enumerate}
\end{lemma}

\begin{proof}
$S/R$ は準有限なので、$\Spec(S)$ の各点 $\mathfrak q$ に定理
\ref{algebra-theorem-main-theorem} を適用できる。
$\Spec(S)$ は準コンパクトなので（補題 \ref{algebra-lemma-quasi-compact} 参照）、
$S'_{g_i} = S_{g_i}$ であり、かつ
$g_1, \ldots, g_n$ が $S$ の単位イデアルを生成するような有限個の
$g_i \in S'$、$i = 1, \ldots, n$ を選べる
（言い換えれば、$g_1, \ldots, g_n$ に対応する $\Spec(S)$ の標準開集合が
$\Spec(S)$ 全体を覆う）。

\medskip\noindent
$D(g) \subset \Spec(S')$ が像に含まれると仮定する。
このとき $D(g) \subset \bigcup D(g_i)$ である。
言い換えれば、$g_1, \ldots, g_n$ は $S'_g$ の単位イデアルを生成する。
$g_i$ の選び方により $S'_{gg_i} \cong S_{gg_i}$ であることに注意する。
従って補題 \ref{algebra-lemma-cover} により $S'_g \cong S_g$ である。

\medskip\noindent
(3) のような有限代数 $S'' \subset S'$ を構成する。
各 $S'_{g_i} \cong S_{g_i}$ は有限型 $R$-代数である。
各 $i$ について、$S'_{g_i}$ が $R$-代数として $1/g_i$ と元 $y_{ij}$ で
生成されるような元 $y_{ij} \in S'$ を選ぶ。$S''$ を、すべての $g_i$ と
すべての $y_{ij}$ で生成される $S'$ の部分 $R$-代数とする。詳細は省略する。
\end{proof}

\section{Zariski の主定理の応用}
\label{algebra-section-apply-main-theorem}

\noindent
以下は直接の応用である。$1$ 次元半局所環の準有限写像のうち有限なものは、
補題 \ref{algebra-lemma-finite-extension-dim-1} の公式で常に等号が成立するものとして
特徴づけられる。

\begin{lemma}
\label{algebra-lemma-quasi-finite-extension-dim-1}
$A \subset B$ を整域の拡大とする。次を仮定する。
\begin{enumerate}
\item $A$ は次元 $1$ の Noether 局所環である。
\item $A \to B$ は有限型である。
\item 誘導される分数体の拡大 $L/K$ は有限である。
\end{enumerate}
このとき $B$ は半局所環である。
$x \in \mathfrak m_A$、$x \not = 0$ とする。
$\mathfrak m_i$、$i = 1, \ldots, n$ を
$B$ の極大イデアルとする。このとき
$$
[L : K]\text{ord}_A(x)
\geq
\sum\nolimits_i
[\kappa(\mathfrak m_i) : \kappa(\mathfrak m_A)]
\text{ord}_{B_{\mathfrak m_i}}(x)
$$
である。ここで $\text{ord}$ は定義 \ref{algebra-definition-ord} のように定義される。
等号が成立するための必要十分条件は $A \to B$ が有限であることである。
\end{lemma}

\begin{proof}
補題 \ref{algebra-lemma-finite-in-codim-1} により環 $B$ は半局所環である。
$B'$ を $B$ における $A$ の整閉包とする。
補題 \ref{algebra-lemma-quasi-finite-open-integral-closure} により、
有限 $A$-部分代数 $C \subset B'$ であって、
$\mathfrak n_i = C \cap \mathfrak m_i$ とおけば
$C_{\mathfrak n_i} \cong B_{\mathfrak m_i}$ となり、かつ素イデアル
$\mathfrak n_1, \ldots, \mathfrak n_n$ が互いに異なるものを選べる。
補題 \ref{algebra-lemma-finite-in-codim-1} により環 $C$ は半局所環である。
$\mathfrak p_j$、$j = 1, \ldots, m$ を $C$ の他の極大イデアル
（「欠落点」）とする。補題 \ref{algebra-lemma-finite-extension-dim-1} により
$$
\text{ord}_A(x^{[L : K]}) =
\sum\nolimits_i
[\kappa(\mathfrak n_i) : \kappa(\mathfrak m_A)]
\text{ord}_{C_{\mathfrak n_i}}(x)
+
\sum\nolimits_j
[\kappa(\mathfrak p_j) : \kappa(\mathfrak m_A)]
\text{ord}_{C_{\mathfrak p_j}}(x)
$$
であり、不等式が従う。等号の場合には $m = 0$
（「欠落点」がない）と結論できる。従って $C \subset B$ は半局所環の包含であり、
極大イデアル上の全単射と、すべての極大イデアルにおける局所化上の同型を誘導する。
そこで $b \in B$ ならば、
$I = \{x \in C \mid xb \in C\}$ は $C$ のどの極大イデアルにも含まれない
$C$ のイデアルである。従って $I = C$、ゆえに $b \in C$ である。
従って $B = C$ であり、$B$ は $A$ 上有限である。
\end{proof}

\noindent
次は、局所環の局所準同型の構造に対する Zariski の主定理の、
より標準的な応用である。

\begin{lemma}
\label{algebra-lemma-essentially-finite-type-fibre-dim-zero}
$(R, \mathfrak m_R) \to (S, \mathfrak m_S)$ を局所環の局所準同型とする。
次を仮定する。
\begin{enumerate}
\item $R \to S$ は本質的有限型である。
\item $\kappa(\mathfrak m_R) \subset \kappa(\mathfrak m_S)$ は有限である。
\item $\dim(S/\mathfrak m_RS) = 0$ である。
\end{enumerate}
このとき $S$ は有限 $R$-代数の局所化である。
\end{lemma}

\begin{proof}
$S'$ を有限型 $R$-代数で、$S'$ のある素イデアル $\mathfrak q'$ に対して
$S = S'_{\mathfrak q'}$ となるものとする。
定義 \ref{algebra-definition-quasi-finite} により、
$R \to S'$ は $\mathfrak q'$ において準有限である。
ある $g' \in S'$、$g' \not \in \mathfrak q'$ に対して
$S'$ を $S'_{g'}$ で置き換えることにより、
$R \to S'$ が準有限であると仮定してよい
（補題 \ref{algebra-lemma-quasi-finite-open} 参照）。
すると補題 \ref{algebra-lemma-quasi-finite-open-integral-closure} により、
有限 $R$-代数 $S''$ と元
$g' \in S'$、$g' \not \in \mathfrak q'$ および $g'' \in S''$ が存在して、
$S'_{g'} \cong S''_{g''}$ が $R$-代数として成り立つ。
これで補題が証明された。
\end{proof}

\begin{lemma}
\label{algebra-lemma-completion-at-quasi-finite-prime}
$R \to S$ を環準同型とし、$\mathfrak q$ を $R$ の $\mathfrak p$ の上にある
$S$ の素イデアルとする。次を仮定する。
\begin{enumerate}
\item $R$ は Noether 環である。
\item $R \to S$ は有限型である。
\item $R \to S$ は $\mathfrak q$ において準有限である。
\end{enumerate}
このとき、ある $R_\mathfrak p^\wedge$-代数 $B$ に対して
$R_\mathfrak p^\wedge \otimes_R S = S_\mathfrak q^\wedge \times B$
である。
\end{lemma}

\begin{proof}
有限 $R$-代数 $S' \subset S$ と元
$g \in S'$、$g \not \in \mathfrak q' = S' \cap \mathfrak q$ であって、
$S'_g = S_g$、従って特に
$S'_{\mathfrak q'} = S_\mathfrak q$ となるものが存在する
（補題 \ref{algebra-lemma-quasi-finite-open-integral-closure} 参照）。
補題 \ref{algebra-lemma-completion-finite-extension} により
$$
R_\mathfrak p^\wedge \otimes_R S' = (S'_{\mathfrak q'})^\wedge \times B'
$$
である。この積分解のもとで $g$ は組 $(u, b')$ に写り、
$g \not \in \mathfrak q'$ なので
$u \in (S'_{\mathfrak q'})^\wedge$ は単元であることに注意する。
$R_\mathfrak p^\wedge \otimes_R S'$ の積分解は次の積分解を誘導する。
% [P01-E0597] The frozen source omits punctuation after this displayed product decomposition.
$$
R_\mathfrak p^\wedge \otimes_R S = A \times B
$$
$S'_g = S_g$ なので
$(R_\mathfrak p^\wedge \otimes_R S')_g = (R_\mathfrak p^\wedge \otimes_R S)_g$
でもあり、$g \mapsto (u, b')$ で $u$ が単元なので
$(S'_{\mathfrak q'})^\wedge = A$ である。
同型 $S'_{\mathfrak q'} = S_\mathfrak q$ は完備化上の同型を定めるから、
$A = S_\mathfrak q^\wedge$ も従う。
これで証明は終わるが、誘導された写像
$\phi : R_\mathfrak p^\wedge \otimes_R S \to A = S_\mathfrak q^\wedge$
が自然なものであることを確認しておくべきである。
$x \otimes 1$、$x \in R_\mathfrak p^\wedge$ の形の元については、
自然な写像 $R_\mathfrak p^\wedge \to S_\mathfrak q^\wedge$ が
$(S'_{\mathfrak q'})^\wedge$ を経由するので明らかである。
$1 \otimes y$、$y \in S$ の形の元については次のように論じられる。
ある $n \geq 1$ に対して、元 $g^ny$ はある $y' \in S'$ の像である。
% [P01-E0598] The frozen source says “image ... by the composition” rather than “under”.
従って $\phi(1 \otimes g^ny)$ は、$y'$ の合成写像
$S' \to (S'_{\mathfrak q'})^\wedge \to S_\mathfrak q^\wedge$ による像であり、
$g^ny$ の写像 $S \to S_\mathfrak q^\wedge$ による像に等しい。
$g$ は単元に写るので、$\phi(1 \otimes y)$ も正しい値、すなわち
$y$ の $S \to S_\mathfrak q^\wedge$ による像をもつ。
\end{proof}

\section{ファイバーの次元}
\label{algebra-section-dimension-fibres}

\noindent
Zariski の主定理を用いて、ファイバーの次元の振る舞いを調べる。
位相空間 $X$ の点 $x$ における次元 $\dim_x(X)$ は、
Topology、定義 \ref{topology-definition-Krull} で定義したことを想起する。

\begin{definition}
\label{algebra-definition-relative-dimension}
$R \to S$ が有限型であり、$\mathfrak q \subset S$ を
$R$ の素イデアル $\mathfrak p$ の上にある素イデアルとする。
$S/R$ の $\mathfrak q$ における{\it 相対次元}を
$\dim_{\mathfrak q}(S/R)$ と書き、
$\mathfrak q$ に対応する点における
$\Spec(S \otimes_R \kappa(\mathfrak p))$ の次元として定義する。
$\dim(S/R)$ を、すべての $\mathfrak q$ にわたる
$\dim_{\mathfrak q}(S/R)$ の上限とする。
これは $S/R$ の{\it 相対次元}と呼ばれる。
\end{definition}

\noindent
特に、$R \to S$ が $\mathfrak q$ において準有限であるための必要十分条件は
$\dim_{\mathfrak q}(S/R) = 0$ である。次の補題は多かれ少なかれ
Zariski の主定理の言い換えである。

\begin{lemma}
\label{algebra-lemma-quasi-finite-over-polynomial-algebra}
$R \to S$ を有限型環準同型とする。
$\mathfrak q \subset S$ を素イデアルとする。
$\dim_{\mathfrak q}(S/R) = n$ と仮定する。
このとき $g \in S$、$g \not\in \mathfrak q$ であって、
$S_g$ が多項式代数 $R[t_1, \ldots, t_n]$ 上準有限となるものが存在する。
\end{lemma}

\begin{proof}
環 $\overline{S} = S \otimes_R \kappa(\mathfrak p)$ は
$\kappa(\mathfrak p)$ 上有限型である。
$\overline{\mathfrak q}$ を $\mathfrak q$ に対応する
$\overline{S}$ の素イデアルとする。
一点における位相空間の次元の定義により、
開集合 $U \subset \Spec(\overline{S})$ であって
% [P01-E0599] The frozen source drops mathfrak from the barred prime symbol.
$\overline{q} \in U$ かつ $\dim(U) = n$ となるものが存在する。
$\Spec(\overline{S})$ の位相は $\Spec(S)$ の位相から誘導されるので
（注意 \ref{algebra-remark-fundamental-diagram} 参照）、
$g \in S$、$g \not \in \mathfrak q$ で、その像
$\overline{g} \in \overline{S}$ が
$D(\overline{g}) \subset U$ を満たすものを選べる。
従って $S$ を $S_g$ で置き換えると
$\dim(\overline{S}) = n$ である。

\medskip\noindent
次に、$S$ の $R$-代数としての生成元 $x_1, \ldots, x_N$ を選ぶ。
補題 \ref{algebra-lemma-Noether-normalization} により、
$x_1, \ldots, x_N$ で生成される $S$ の $\mathbf{Z}$-部分代数の元
$y_1, \ldots, y_n$ であって、写像
$R[t_1, \ldots, t_n] \to S$、$t_i \mapsto y_i$ が、
% [P01-E0600] The frozen polynomial algebra omits the comma after t_1.
$\kappa(\mathfrak p)[t_1\ldots, t_n] \to \overline{S}$ を
有限にするという性質をもつものが存在する。
特に、$S$ は $R[t_1, \ldots, t_n]$ 上
$\mathfrak q$ において準有限である。従って補題
\ref{algebra-lemma-quasi-finite-open} により、ある
$g\in S$、$g \not \in \mathfrak q$ に対して $S$ を $S_g$ で置き換え、
$R[t_1, \ldots, t_n] \to S$ が準有限であると仮定してよい。
\end{proof}

\begin{lemma}
\label{algebra-lemma-refined-quasi-finite-over-polynomial-algebra}
$R \to S$ を環準同型とする。$\mathfrak q \subset S$ を
$R$ の素イデアル $\mathfrak p$ の上にある素イデアルとする。
次を仮定する。
\begin{enumerate}
\item $R \to S$ は有限型である。
\item $\dim_{\mathfrak q}(S/R) = n$ である。
\item $\text{trdeg}_{\kappa(\mathfrak p)}\kappa(\mathfrak q) = r$ である。
\end{enumerate}
このとき $f \in R$、$f \not \in \mathfrak p$、
$g \in S$、$g \not\in \mathfrak q$ と、準有限環準同型
$$
\varphi : R_f[x_1, \ldots, x_n] \longrightarrow S_g
$$
であって
% [P01-E0601] The frozen lemma statement omits punctuation after this equality.
$\varphi^{-1}(\mathfrak qS_g) =
(\mathfrak p, x_{r + 1}, \ldots, x_n)R_f[x_1, \ldots, x_n]$
となるものが存在する。
\end{lemma}

\begin{proof}
$S$ を主局所化で置き換えた後、準有限環準同型
$\varphi : R[t_1, \ldots, t_n] \to S$ が存在すると仮定してよい
（補題 \ref{algebra-lemma-quasi-finite-over-polynomial-algebra} 参照）。
$\mathfrak q' = \varphi^{-1}(\mathfrak q)$ とおく。
$\overline{\mathfrak q}' \subset \kappa(\mathfrak p)[t_1, \ldots, t_n]$
を $\mathfrak q'$ に対応する素イデアルとする。
補題 \ref{algebra-lemma-refined-Noether-normalization} により、有限環準同型
$\kappa(\mathfrak p)[x_1, \ldots, x_n] \to
\kappa(\mathfrak p)[t_1, \ldots, t_n]$
であって、$\overline{\mathfrak q}'$ の逆像が
$(x_{r + 1}, \ldots, x_n)$ となるものが存在する。
$\overline{h}_i \in \kappa(\mathfrak p)[t_1, \ldots, t_n]$
を $x_i$ の像とする。
$f \in R$、$f \not \in \mathfrak p$ と
$h_i \in R_f[t_1, \ldots, t_n]$ であって、
$\kappa(\mathfrak p)[t_1, \ldots, t_n]$ において
$\overline{h}_i$ に写るものを選べる。このとき環準同型
$$
R_f[x_1, \ldots, x_n] \longrightarrow R_f[t_1, \ldots, t_n]
$$
は $\kappa(\mathfrak p)$ とテンソルを取ると有限になる。
特に $R_f[t_1, \ldots, t_n]$ は
$R_f[x_1, \ldots, x_n]$ 上、素イデアル
$\mathfrak q'R_f[t_1, \ldots, t_n]$ において準有限である。
従って補題 \ref{algebra-lemma-quasi-finite-open} により、
$g \in R_f[t_1, \ldots, t_n]$、
$g \not \in \mathfrak q'R_f[t_1, \ldots, t_n]$ であって、
$R_f[x_1, \ldots, x_n] \to R_f[t_1, \ldots, t_n, 1/g]$
が準有限となるものが存在する。従って合成
$$
R_f[x_1, \ldots, x_n] \longrightarrow
R_f[t_1, \ldots, t_n, 1/g] \longrightarrow S_{\varphi(g)}
$$
は準有限であり、証明が完了する。
\end{proof}

\begin{lemma}
\label{algebra-lemma-dimension-inequality-quasi-finite}
$R \to S$ を有限型環準同型とする。
$\mathfrak q \subset S$ を $\mathfrak p \subset R$ の上にある素イデアルとする。
$R \to S$ が $\mathfrak q$ において準有限ならば、
$\dim(S_{\mathfrak q}) \leq \dim(R_{\mathfrak p})$ である。
\end{lemma}

\begin{proof}
$R_{\mathfrak p}$ が Noether 環ならば
（このとき $S_{\mathfrak q}$ は $R_{\mathfrak p}$ 上本質的有限型なので、
それも Noether 環である）、
これは補題 \ref{algebra-lemma-dimension-base-fibre-total} と定義から直ちに従う。
一般の場合には、$S'$ を $S_\mathfrak p$ における
$R_\mathfrak p$ の整閉包とする。
Zariski の主定理 \ref{algebra-theorem-main-theorem} により、
% [P01-E0602] The frozen source reverses the direction by saying q' lies over q.
$\mathfrak q$ の上にあるある $\mathfrak q' \subset S'$ に対して
$S_{\mathfrak q} = S'_{\mathfrak q'}$ である。
補題 \ref{algebra-lemma-integral-dim-up} により
$\dim(S') \leq \dim(R_\mathfrak p)$ であり、従ってなおさら
$\dim(S_\mathfrak q) = \dim(S'_{\mathfrak q'}) \leq \dim(R_\mathfrak p)$
である。
\end{proof}

\begin{lemma}
\label{algebra-lemma-dimension-quasi-finite-over-polynomial-algebra}
\begin{slogan}
アフィン n 空間の準有限被覆の次元は高々 n である。
\end{slogan}
$k$ を体とし、$S$ を有限型 $k$-代数とする。
準有限な $k$-代数準同型
$k[t_1, \ldots, t_n] \subset S$ が存在すると仮定する。
このとき $\dim(S) \leq n$ である。
\end{lemma}

\begin{proof}
補題 \ref{algebra-lemma-dim-affine-space} により、
$k[t_1, \ldots, t_n]$ の任意の局所環の次元は高々 $n$ である。
従って補題 \ref{algebra-lemma-dimension-inequality-quasi-finite} から結論が従う。
\end{proof}

\begin{lemma}
\label{algebra-lemma-dimension-fibres-bounded-open-upstairs}
$R \to S$ を有限型環準同型とする。
$\mathfrak q \subset S$ を素イデアルとし、
$\dim_{\mathfrak q}(S/R) = n$ と仮定する。
このとき $\Spec(S)$ における $\mathfrak q$ の開近傍 $V$ であって、
すべての $\mathfrak q' \in V$ に対して
$\dim_{\mathfrak q'}(S/R) \leq n$ となるものが存在する。
\end{lemma}

\begin{proof}
補題 \ref{algebra-lemma-quasi-finite-over-polynomial-algebra} により、
$S$ が多項式代数 $R[t_1, \ldots, t_n]$ 上準有限であると
仮定してよい。ファイバーを考えることにより、
補題 \ref{algebra-lemma-dimension-quasi-finite-over-polynomial-algebra} に帰着する。
\end{proof}

\noindent
言い換えると、この補題は、ファイバーの次元が $\leq n$ となる点の集合が
$\Spec(S)$ で開であることを述べている。次の補題は、この開集合の形成が
基底変換と可換であることを述べる。環準同型が有限表示ならば、
この集合は準コンパクト開である（以下を参照）。

\begin{lemma}
\label{algebra-lemma-dimension-fibres-bounded-open-upstairs-base-change}
$R \to S$ を有限型環準同型とする。
$R \to R'$ を任意の環準同型とする。
$S' = R' \otimes_R S$ とおき、スペクトル上の対応する写像を
$f : \Spec(S') \to \Spec(S)$ と書く。
$n \geq 0$ とする。逆像
$f^{-1}(\{\mathfrak q \in \Spec(S) \mid
\dim_{\mathfrak q}(S/R) \leq n\})$
は
$\{\mathfrak q' \in \Spec(S') \mid
\dim_{\mathfrak q'}(S'/R') \leq n\}$
に等しい。
\end{lemma}

\begin{proof}
この条件は、体上有限型であるファイバー環の次元によって定式化されている。
これと補題 \ref{algebra-lemma-dimension-at-a-point-preserved-field-extension} を
組み合わせれば補題が得られる。
\end{proof}

\begin{lemma}
\label{algebra-lemma-dimension-fibres-bounded-quasi-compact-open-upstairs}
$R \to S$ を有限表示の環準同型とする。
$n \geq 0$ とする。集合
$$
V_n = \{\mathfrak q \in \Spec(S) \mid \dim_{\mathfrak q}(S/R) \leq n\}
$$
は $\Spec(S)$ の準コンパクト開部分集合である。
\end{lemma}

\begin{proof}
補題 \ref{algebra-lemma-dimension-fibres-bounded-open-upstairs} によりこれは開である。
% [P01-E0603] The frozen proof reuses n for both the fibre bound and presentation variables.
$S = R[x_1, \ldots, x_n]/(f_1, \ldots, f_m)$ を $S$ の表示とする。
$R_0$ を、多項式 $f_i$ の係数で生成される $R$ の
$\mathbf{Z}$-部分代数とする。
$S_0 = R_0[x_1, \ldots, x_n]/(f_1, \ldots, f_m)$ とおく。
このとき $S = R \otimes_{R_0} S_0$ である。
補題 \ref{algebra-lemma-dimension-fibres-bounded-open-upstairs-base-change} により、
$V_n$ は、準コンパクト連続写像
$\Spec(S) \to \Spec(S_0)$ による開集合 $V_{0, n}$ の逆像である。
$S_0$ は Noether 環なので $V_{0, n}$ は準コンパクトである。
\end{proof}

\begin{lemma}
\label{algebra-lemma-finite-type-domain-over-valuation-ring-dim-fibres}
$R$ を剰余体 $k$、分数体 $K$ をもつ付値環とする。
$S$ を $R$ を含む整域で、$S$ が $R$ 上有限型なものとする。
$S \otimes_R k$ が零環でなければ
$$
\dim(S \otimes_R k) = \dim(S \otimes_R K)
$$
% [P01-E0604] The frozen source omits punctuation after the displayed equality.
である。実際、$\Spec(S \otimes_R k)$ は等次元である。
\end{lemma}

\begin{proof}
$S$ の $\mathfrak m_RS$ を含む任意の素イデアル $\mathfrak q$ に対して
$\dim_{\mathfrak q}(S/k)$ が $\dim(S \otimes_R K)$ に等しいことを
示せば十分である。そのような素イデアルを一つ選ぶ。
補題 \ref{algebra-lemma-dimension-fibres-bounded-open-upstairs} により、不等式
$\dim_{\mathfrak q}(S/k) \geq \dim(S \otimes_R K)$ が成り立つ。
$n = \dim_{\mathfrak q}(S/k)$ とおく。
補題 \ref{algebra-lemma-quasi-finite-over-polynomial-algebra} により、
ある $g \in S$、$g \not \in \mathfrak q$ に対して $S$ を $S_g$ で
置き換えた後、準有限環準同型
$R[t_1, \ldots, t_n] \to S$ が存在する。
$\dim(S \otimes_R K) < n$ ならば、
$K[t_1, \ldots, t_n] \to S \otimes_R K$ は非零の核をもつ。
$f = \sum a_I t_1^{i_1}\ldots t_n^{i_n}$ と書く。
$f$ を、最小の付値をもつ $f$ の非零係数で割ることにより、
$f\in R[t_1, \ldots, t_n]$ であり、ある $a_I$ の $k$ における像が
零でないと仮定してよい。従って環準同型
$k[t_1, \ldots, t_n] \to S \otimes_R k$ は非零の核をもち、
$\dim(S \otimes_R k) < n$ を含意する。これは矛盾である。
\end{proof}

\section{有限表示の代数と加群}
\label{algebra-section-finite-presentation}

\noindent
この節では、有限型または有限表示という仮定が重要となる標準的な結果を
いくつか論じる。

\begin{lemma}
\label{algebra-lemma-finite-type-descends}
$R \to S$ を環準同型とする。
$R \to R'$ を忠実平坦環準同型とする。
$S' = R'\otimes_R S$ とおく。
$R \to S$ が有限型であるための必要十分条件は、
$R' \to S'$ が有限型であることである。
\end{lemma}

\begin{proof}
$R \to S$ が有限型ならば $R' \to S'$ が有限型であることは明らかである。
$R' \to S'$ が有限型であると仮定する。
$y_1, \ldots, y_m$ が $R'$ 上 $S'$ を生成するとする。
$y_j = \sum_i a_{ij} \otimes x_{ji}$ と書く。ただし
$a_{ij} \in R'$、$x_{ji} \in S$ である。
% [P01-E0605] The frozen proof switches the indices from x_{ji} to x_{ij}.
$A \subset S$ を $x_{ij}$ で生成される $R$-部分代数とする。
平坦性により $A' := R' \otimes_R A \subset S'$ であり、
構成から $y_j \in A'$ である。従って $A' = S'$ である。
忠実平坦性により $A = S$ である。
\end{proof}

\begin{lemma}
\label{algebra-lemma-finite-presentation-descends}
$R \to S$ を環準同型とする。
$R \to R'$ を忠実平坦環準同型とする。
$S' = R'\otimes_R S$ とおく。
$R \to S$ が有限表示であるための必要十分条件は、
$R' \to S'$ が有限表示であることである。
\end{lemma}

\begin{proof}
$R \to S$ が有限表示ならば $R' \to S'$ が有限表示であることは明らかである。
$R' \to S'$ が有限表示であると仮定する。
補題 \ref{algebra-lemma-finite-type-descends} により $R \to S$ は有限型である。
$S = R[x_1, \ldots, x_n]/I$ と書く。
平坦性により $S' = R'[x_1, \ldots, x_n]/R'\otimes I$ である。
$g_1, \ldots, g_m$ が $R'[x_1, \ldots, x_n]$ 上
$R'\otimes I$ を生成するとする。
$g_j = \sum_i a_{ij} \otimes f_{ji}$ と書く。ただし
$a_{ij} \in R'$、$f_{ji} \in I$ である。
% [P01-E0606] The frozen proof switches the indices from f_{ji} to f_{ij}.
$J \subset I$ を $f_{ij}$ で生成されるイデアルとする。
平坦性により $R' \otimes_R J \subset R'\otimes_R I$ であり、
両者は $R'[x_1, \ldots, x_n]$ 上のイデアルである。
構成から $g_j \in R' \otimes_R J$ である。従って
$R' \otimes_R J = R'\otimes_R I$ である。
忠実平坦性により $J = I$ である。
\end{proof}

\begin{lemma}
\label{algebra-lemma-construct-fp-module}
$R$ を環とする。
$I \subset R$ をイデアルとする。
$S \subset R$ を乗法的部分集合とする。
$R' = S^{-1}(R/I) = S^{-1}R/S^{-1}I$ とおく。
\begin{enumerate}
\item 任意の有限 $R'$-加群 $M'$ に対して、
$S^{-1}(M/IM) \cong M'$ となる有限 $R$-加群 $M$ が存在する。
\item 任意の有限表示 $R'$-加群 $M'$ に対して、
$S^{-1}(M/IM) \cong M'$ となる有限表示 $R$-加群 $M$ が存在する。
\end{enumerate}
\end{lemma}

\begin{proof}
(1) の証明。短完全列
$0 \to K' \to (R')^{\oplus n} \to M' \to 0$
を選ぶ。$K \subset R^{\oplus n}$ を、写像
$R^{\oplus n} \to (R')^{\oplus n}$ による $K'$ の逆像とする。
このとき $M = R^{\oplus n}/K$ が求めるものである。

\medskip\noindent
(2) の証明。
表示 $(R')^{\oplus m} \to (R')^{\oplus n} \to M' \to 0$ を選ぶ。
最初の写像が行列 $A' = (a'_{ij})$ で与えられ、
二番目の写像が生成元 $x'_i \in M'$、$i = 1, \ldots, n$ により
決定されるとする。$R' = S^{-1}(R/I)$ なので、
$s \in S$ と、$R$ に係数をもつ行列 $A = (a_{ij})$ であって
$a'_{ij} = a_{ij} / s \bmod S^{-1}I$ となるものを選べる。
$M$ を、表示
$R^{\oplus m} \to R^{\oplus n} \to M \to 0$
をもつ有限表示 $R$-加群とする。ここで最初の写像は行列 $A$ で与えられ、
二番目の写像は生成元 $x_i \in M$、$i = 1, \ldots, n$ により決定される。
このとき写像 $M \to M'$、$x_i \mapsto x'_i$ は同型
$S^{-1}(M/IM) \cong M'$ を誘導する。
\end{proof}

\begin{lemma}
\label{algebra-lemma-construct-fp-module-from-localization}
$R$ を環とする。
$S \subset R$ を乗法的部分集合とする。
$M$ を $R$-加群とする。
\begin{enumerate}
\item $S^{-1}M$ が有限 $S^{-1}R$-加群ならば、有限 $R$-加群 $M'$ と、
同型 $S^{-1}M' \to S^{-1}M$ を誘導する写像 $M' \to M$ が存在する。
\item $S^{-1}M$ が有限表示 $S^{-1}R$-加群ならば、
有限表示 $R$-加群 $M'$ と、同型
$S^{-1}M' \to S^{-1}M$ を誘導する写像 $M' \to M$ が存在する。
\end{enumerate}
\end{lemma}

\begin{proof}
(1) の証明。$x_1, \ldots, x_n \in M$ を、
$S^{-1}R$-加群として $S^{-1}M$ を生成する元とする。
$M'$ を $x_1, \ldots, x_n$ で生成される $M$ の $R$-部分加群とする。

\medskip\noindent
(2) の証明。$x_1, \ldots, x_n \in M$ を、
$S^{-1}R$-加群として $S^{-1}M$ を生成する元とする。
$K = \Ker(R^{\oplus n} \to M)$ とおく。ここで写像は
$(a_1, \ldots, a_n) \mapsto \sum a_i x_i$ で与えられる。
補題 \ref{algebra-lemma-extension} により
$S^{-1}K$ は有限 $S^{-1}R$-加群である。
(1) により、$S^{-1}K' = S^{-1}K$ を満たす有限部分加群
$K' \subset K$ を選べる。
$M' = \Coker(K' \to R^{\oplus n})$ とおく。
\end{proof}

\begin{lemma}
\label{algebra-lemma-construct-fp-module-from-stalk}
$R$ を環とする。
$\mathfrak p \subset R$ を素イデアルとする。
$M$ を $R$-加群とする。
\begin{enumerate}
\item $M_{\mathfrak p}$ が有限 $R_{\mathfrak p}$-加群ならば、
有限 $R$-加群 $M'$ と、同型
$M'_{\mathfrak p} \to M_{\mathfrak p}$ を誘導する写像
$M' \to M$ が存在する。
\item $M_{\mathfrak p}$ が有限表示 $R_{\mathfrak p}$-加群ならば、
有限表示 $R$-加群 $M'$ と、同型
$M'_{\mathfrak p} \to M_{\mathfrak p}$ を誘導する写像
$M' \to M$ が存在する。
\end{enumerate}
\end{lemma}

\begin{proof}
% [P01-E0607] The frozen proof omits terminal punctuation.
これは補題 \ref{algebra-lemma-construct-fp-module-from-localization} の特別な場合である。
\end{proof}

\begin{lemma}
\label{algebra-lemma-local-isomorphism}
$\varphi : R \to S$ を環準同型とする。
$\mathfrak q \subset S$ を $\mathfrak p \subset R$ の上にある素イデアルとする。
次を仮定する。
\begin{enumerate}
\item $S$ は $R$ 上有限表示である。
\item $\varphi$ は同型 $R_\mathfrak p \cong S_\mathfrak q$ を誘導する。
\end{enumerate}
このとき $f \in R$、$f \not \in \mathfrak p$ と
$R_f$-代数 $C$ であって、
$S_f \cong R_f \times C$ が $R_f$-代数として成り立つものが存在する。
\end{lemma}

\begin{proof}
$S = R[x_1, \ldots, x_n]/(g_1, \ldots, g_m)$ と書く。
$a_i \in R_\mathfrak p$ を、$S_\mathfrak q$ における $x_i$ の像に
写る元とする。ある $f \in R$、$f \not \in \mathfrak p$ に対して
$a_i = b_i/f$ と書く。$R$ を $R_f$ で、$x_i$ を $x_i - a_i$ で
置き換えることにより、
$S = R[x_1, \ldots, x_n]/(g_1, \ldots, g_m)$ であり、
$x_i$ が $S_\mathfrak q$ において零に写ると仮定してよい。
$c_j$ を $g_j$ の定数項とすると、$c_j$ は
$R_\mathfrak p$ において零に写る。さらにもう一度 $R$ を置き換えると、
$g_j$ の定数係数 $c_j$ が零であると仮定してよい。
従って、核がイデアル $(x_1, \ldots, x_n)$ である
$R$-代数準同型 $S \to R$、$x_i \mapsto 0$ を得る。

\medskip\noindent
同型 $R_\mathfrak p \to S_\mathfrak q \to R_\mathfrak p$ があり、
$S \to R$ は $x_i$ を零に写す。従って
$S_\mathfrak q = R_\mathfrak p[x_1, \ldots, x_n]/(x_1, \ldots, x_n)$
でなければならず、なおさら
$S_\mathfrak q = S_\mathfrak p/(x_1, \ldots, x_n)S_\mathfrak p$ である。
これは、有限生成イデアル $(x_1, \ldots, x_n)S_\mathfrak p$ が
$S_\mathfrak p$ において純であることを意味する
（定義 \ref{algebra-definition-pure-ideal} 参照）。
従って補題 \ref{algebra-lemma-finitely-generated-pure-ideal} により、
$(x_1, \ldots, x_n)S_\mathfrak p$ は
$S_\mathfrak p$ の冪等元 $e$ で生成される。
ある $f \in R$、$f \not \in \mathfrak p$ に対して
$R \to S$ を $R_f \to S_f$ で置き換えると、
$e$ に写る冪等元 $e' \in S$ を選べる。
すると $e'S$ と $(x_1, \ldots, x_n)S$ は、
$S_\mathfrak p$ において等しくなる有限生成イデアルである。
従って再び、ある $f \in R$、$f \not \in \mathfrak p$ に対して
$R \to S$ を $R_f \to S_f$ で置き換えることにより、
$e'S = (x_1, \ldots, x_n)S$ と仮定してよい。
$C = e'S$ とおけば証明が完了する。
\end{proof}

\begin{lemma}
\label{algebra-lemma-isomorphic-local-rings}
$R$ を環とする。
$S$、$S'$ を $R$ 上有限表示とする。
$\mathfrak q \subset S$ と $\mathfrak q' \subset S'$ を素イデアルとする。
$S_{\mathfrak q} \cong S'_{\mathfrak q'}$ が $R$-代数として成り立つならば、
$g \in S$、$g \not \in \mathfrak q$ と
$g' \in S'$、$g' \not \in \mathfrak q'$ であって、
$S_g \cong S'_{g'}$ が $R$-代数として成り立つものが存在する。
\end{lemma}

\begin{proof}
$\psi : S_{\mathfrak q} \to S'_{\mathfrak q'}$ を補題の仮定の同型とする。
$S = R[x_1, \ldots, x_n]/(f_1, \ldots, f_r)$ および
$S' = R[y_1, \ldots, y_m]/J$ と書く。
各 $i = 1, \ldots, n$ について分数
$h_i/g_i$ を選ぶ。ただし $h_i, g_i \in R[y_1, \ldots, y_m]$ であり、
$g_i \bmod J$ は $\mathfrak q'$ に属さず、この分数は
$\psi$ による $x_i$ の像を表す。
% [P01-E0608] The frozen coordinated replacement omits the second replacement operation.
$S'$ を $S'_{g_1 \ldots g_n}$ で置き換え、
$R[y_1, \ldots, y_m, y_{m + 1}]$
（$y_{m + 1}$ を $1/(g_1\ldots g_n)$ に写す）を用いることにより、
$\psi(x_i)$ がある $h_i \in R[y_1, \ldots, y_m]$ の像であると
仮定してよい。元
$f_j(h_1, \ldots, h_n) \in R[y_1, \ldots, y_m]$ を考える。
$\psi$ は各 $f_j$ を消すから、
$g \in R[y_1, \ldots, y_m]$、$g \bmod J \not \in \mathfrak q'$ であって、
各 $j = 1, \ldots, r$ に対して
$g f_j(h_1, \ldots, h_n) \in J$ となるものが存在する。
$S'$ を $S'_g$ で置き換え、前と同様に
$R[y_1, \ldots, y_m, y_{m + 1}]$ を用いることにより、
$f_j(h_1, \ldots, h_n) \in J$ と仮定してよい。
従って、局所環上で $\psi$ を誘導する環準同型
$S \to S'$、$x_i \mapsto h_i$ を得る。
補題 \ref{algebra-lemma-compose-finite-type} により写像
$S \to S'$ は有限表示である。
補題 \ref{algebra-lemma-local-isomorphism} により
$S' = S \times C$ と仮定してよい。
% [P01-E0609] The frozen source selects the C factor although the conclusion needs S.
従って、因子 $C$ に対応する冪等元において $S'$ を局所化すれば
結論を得る。
\end{proof}

\begin{lemma}
\label{algebra-lemma-finite-type-mod-nilpotent}
$R$ を環とし、$I \subset R$ を冪零イデアルとする。
$S$ を $R$-代数であって、$R/I \to S/IS$ が有限型となるものとする。
このとき $R \to S$ は有限型である。
\end{lemma}

\begin{proof}
$s_1, \ldots, s_n \in S$ を、その $S/IS$ における像が
$R/I$ 上の代数として $S/IS$ を生成するように選ぶ。
補題 \ref{algebra-lemma-NAK} の (11) により、$R$-代数準同型
$R[x_1, \ldots, x_n] \to S$、$x_i \mapsto s_i$ は全射であり、
結論が従う。
\end{proof}

\begin{lemma}
\label{algebra-lemma-surjective-mod-locally-nilpotent}
$R$ を環とし、$I \subset R$ を局所冪零イデアルとする。
$S \to S'$ を $R$-代数準同型で、$S \to S'/IS'$ が全射であり、
$S'$ が $R$ 上有限型となるものとする。このとき $S \to S'$ は全射である。
\end{lemma}

\begin{proof}
あるイデアル $K$ に対して $S' = R[x_1, \ldots, x_m]/K$ と書く。
% [P01-E0610] The frozen proof switches from m variables to undefined n three times.
仮定により
$g_j = x_j + \sum \delta_{j, J} x^J \in R[x_1, \ldots, x_n]$
で、$\delta_{j, J} \in I$ かつ
$g_j \bmod K \in \Im(S \to S')$ となるものが存在する。
従って、$g_1, \ldots, g_m$ が
$R[x_1, \ldots, x_n]$ を生成することを示せば十分である。
$R_0 \subset R$ を、少なくとも $\delta_{j, J}$ を含む
$R$ の有限生成 $\mathbf{Z}$-部分代数とする。
このとき $R_0 \cap I$ は冪零イデアルである
（補題 \ref{algebra-lemma-Noetherian-power} による）。
従って $R_0[x_1, \ldots, x_n]$ は $g_1, \ldots, g_m$ で生成される
（$x_j \mapsto g_j$ が $R_0[x_1, \ldots, x_m]$ の自己同型を定めるからである。
詳細は省略する）。$R$ は部分環 $R_0$ の合併なので証明が完了する。
\end{proof}

\begin{lemma}
\label{algebra-lemma-isomorphism-modulo-ideal}
$R$ を環、$I \subset R$ をイデアルとする。
$S \to S'$ を $R$-代数準同型とする。
$IS \subset \mathfrak q \subset S$ を素イデアルとする。次を仮定する。
\begin{enumerate}
\item $S \to S'$ は全射である。
\item $S_\mathfrak q/IS_\mathfrak q \to S'_\mathfrak q/IS'_\mathfrak q$
は同型である。
\item $S$ は $R$ 上有限型である。
% [P01-E0611] The frozen fourth item omits the finite verb.
\item $S'$ は $R$ 上有限表示である。
\item $S'_\mathfrak q$ は $R$ 上平坦である。
\end{enumerate}
このとき、ある $g \in S$、$g \not \in \mathfrak q$ に対して
$S_g \to S'_g$ は同型である。
\end{lemma}

\begin{proof}
$J = \Ker(S \to S')$ とおく。
補題 \ref{algebra-lemma-compose-finite-type} により
$J$ は有限生成イデアルである。$S'_\mathfrak q$ は
$R$ 上平坦なので
$J_\mathfrak q/IJ_\mathfrak q \subset S_\mathfrak q/IS_{\mathfrak q}$
である（$0 \to J \to S \to S' \to 0$ に
補題 \ref{algebra-lemma-flat-tor-zero} を適用する）。
仮定 (2) により $J_\mathfrak q/IJ_\mathfrak q$ は零である。
Nakayama の補題（補題 \ref{algebra-lemma-NAK}）により、
$g \in S$、$g \not \in \mathfrak q$ であって
$J_g = 0$ となるものが存在する。従って望むように
$S_g \cong S'_g$ である。
\end{proof}

\begin{lemma}
\label{algebra-lemma-isomorphism-modulo-locally-nilpotent}
$R$ を環、$I \subset R$ をイデアルとする。
$S \to S'$ を $R$-代数準同型とする。次を仮定する。
\begin{enumerate}
\item $I$ は局所冪零である。
\item $S/IS \to S'/IS'$ は同型である。
\item $S$ は $R$ 上有限型である。
\item $S'$ は $R$ 上有限表示である。
\item $S'$ は $R$ 上平坦である。
\end{enumerate}
このとき $S \to S'$ は同型である。
\end{lemma}

\begin{proof}
補題 \ref{algebra-lemma-surjective-mod-locally-nilpotent} により写像
$S \to S'$ は全射である。$I$ が局所冪零なので、イデアル
$IS$ と $IS'$ も局所冪零である
（補題 \ref{algebra-lemma-locally-nilpotent}）。
従って $S$ の任意の素イデアル $\mathfrak q$ は $IS$ を含み、
（自明に）
$S_\mathfrak q/IS_\mathfrak q \cong S'_\mathfrak q/IS'_\mathfrak q$
である。従って補題 \ref{algebra-lemma-isomorphism-modulo-ideal} が適用でき、
任意の素イデアル $\mathfrak q \subset S$ に対して
$S_\mathfrak q \to S'_\mathfrak q$ が同型であることが分かる。
% [P01-E0612] The frozen conclusion has the awkward order “injective for example by”.
従って、例えば補題 \ref{algebra-lemma-characterize-zero-local} により
$S \to S'$ は単射である。
\end{proof}

\section{余極限と有限表示の写像}
\label{algebra-section-colimits-flat}

\noindent
% [P01-E0613] The frozen introduction says “prove result” without an article.
この節では、絶対 Noether 化帰着を用いて結果を証明する際に
後で役立ついくつかの準備的補題を証明する。
Categories、節 \ref{categories-section-directed-colimits} では
フィルター余極限を一般に論じた。ここではこの非常に一般的な概念の例を挙げる。

\begin{lemma}
\label{algebra-lemma-ring-colimit-fp-category}
$R \to A$ を環準同型とする。
$R$-代数準同型 $A' \to A$ の図式で、$A'$ が
$R$ 上有限表示であるもの全体の圏 $\mathcal{I}$ を考える。
このとき $\mathcal{I}$ はフィルター圏であり、
$\mathcal{I}$ 上の $A'$ の余極限は $A$ と同型である。
\end{lemma}

\begin{proof}
圏\footnote{集合論的な困難を避けるため、$A'$ が
$R[x_1, x_2, x_3, \ldots]$ の商であるような $A' \to A$ だけを考える。}
$\mathcal{I}$ は空でない。実際 $R \to R$ はその対象である。
$\mathcal{I}$ の二つの対象 $A' \to A$、$A'' \to A$ を考える。
このとき $A' \otimes_R A'' \to A$ は
$\mathcal{I}$ に属する（補題 \ref{algebra-lemma-compose-finite-type} と
\ref{algebra-lemma-base-change-finiteness} を用いる）。
環準同型 $A' \to A' \otimes_R A''$ と
$A'' \to A' \otimes_R A''$ は $\mathcal{I}$ の射を定め、
フィルター圏の二番目の定義条件を証明する
（Categories、定義 \ref{categories-definition-directed} 参照）。
最後に、$\mathcal{I}$ の二つの射
$\sigma, \tau : A' \to A''$ があるとする。
$x_1, \ldots, x_r \in A'$ が $R$-代数として $A'$ を生成するならば、
$A''' = A''/(\sigma(x_i) - \tau(x_i))$ を考えられる。
これは有限表示 $R$-代数であり、与えられた $R$-代数準同型
$A'' \to A$ は全射 $\nu : A'' \to A'''$ を経由する。
従って $\nu$ は $\mathcal{I}$ において $\sigma$ と $\tau$ を
等化する射であり、望む条件が得られる。

\medskip\noindent
% [P01-E0614] The frozen proof calls the filtered index category cofiltered.
添字圏が余フィルター圏であることから、集合の圏で
$B = \colim_{A' \to A} A'$ の値を計算できる
（詳細の一部は省略する。Categories、節
\ref{categories-section-directed-colimits} の議論と比較せよ）。
$B \to A$ が全射であることを見るには、各 $a \in A$ に対して
$R[x] \to A$、$x \mapsto a$ を用いれば、
$a$ が $B \to A$ の像に属することが分かる。
逆に $b \in B$ が $A$ の零に写るならば、
$\mathcal{I}$ の $A' \to A$ と、$b$ に写る $a' \in A'$ を選べる。
すると $A'/(a') \to A$ も $\mathcal{I}$ に属し、写像
$A' \to B$ は $A' \to A'/(a') \to B$ と分解する。
従って望むように $b = 0$ である。
\end{proof}

\noindent
余極限を前順序集合上で考える方が容易なことが多い。
% [P01-E0615] The frozen sentence omits the verb in the preordered-set declaration.
$(\Lambda, \geq)$ を前順序集合とする。
$\Lambda$ 上の環の系は、各 $\lambda \in \Lambda$ に対する環
$R_\lambda$ と、$\lambda \leq \mu$ のときの射
$R_\lambda \to R_\mu$ により与えられる。
これらの射は、すべての $\lambda \leq \mu \leq \nu$ に対して
$R_\lambda \to R_\mu \to R_\nu$ が写像
$R_\lambda \to R_\nu$ に等しいという規則を満たさなければならない。
Categories、節 \ref{categories-section-posets-limits} を参照せよ。
% [P01-E0616] The frozen text switches from Lambda to an undefined I.
しばしば $(I, \leq)$ が{\it 有向}であると仮定する。
これは $\Lambda$ が空でなく、任意の $\lambda, \mu \in \Lambda$ に対して
$\lambda \leq \nu$ かつ $\mu \leq \nu$ となる
$\nu \in \Lambda$ が存在することを意味する。
この場合、余極限 $\colim_\lambda R_\lambda$ は
「直極限」と呼ばれることもあるが、ここではこの用語を用いない。

\medskip\noindent
Categories、補題 \ref{categories-lemma-directed-category-system} は、
フィルター添字圏上の余極限と有向集合上の余極限が同じものであることを
述べている。

\begin{lemma}
\label{algebra-lemma-ring-colimit-fp}
$R \to A$ を環準同型とする。
有限表示 $R$-代数の有向系 $A_\lambda$ であって
$A = \colim_\lambda A_\lambda$ となるものが存在する。
$A$ が $R$ 上有限型ならば、$A_\lambda$ の系のすべての遷移写像が
全射となるように選べる。
\end{lemma}

\begin{proof}
第一の証明は、補題 \ref{algebra-lemma-ring-colimit-fp-category} と
Categories、補題 \ref{categories-lemma-directed-category-system} から従う。

\medskip\noindent
第二の証明。補題 \ref{algebra-lemma-module-colimit-fp} の証明と比較せよ。
任意の有限部分集合 $S \subset A$ と、$S$ の元の間の多項式関係の
任意の有限集合 $E$ を考える。
各 $s \in S$ は $x_s \in A$ に対応し、各 $e \in E$ は、
$f_e(x_s) = 0$ を満たす多項式
$f_e \in R[X_s; s\in S]$ からなる。
$A_{S, E} = R[X_s; s\in S]/(f_e; e\in E)$ とおく。
これは有限表示 $R$-代数である。
標準写像 $A_{S, E} \to A$ がある。
$S \subset S'$ であり、$E$ の元が写像
$R[X_s; s \in S] \to R[X_s; s\in S']$ により
$E'$ の部分集合に対応するならば、
$A$ への写像と可換する明らかな写像
$A_{S, E} \to A_{S', E'}$ がある。
% [P01-E0617] The frozen phrase omits “to” after “setting Lambda equal”.
従って、$\Lambda$ を、上の包含関係で順序づけた対
$(S, E)$ の集合とすれば、有向半順序集合を得る。
この有向系の余極限が $A$ であることは明らかである。

\medskip\noindent
最後の主張について、$A = R[x_1, \ldots, x_n]/I$ とする。
この場合、$S = \{x_1, \ldots, x_n\}$ となる上の系
$(S, E)$ からなる部分集合 $\Lambda' \subset \Lambda$ を考える。
依然として $A = \colim_{\lambda' \in \Lambda'} A_{\lambda'}$
であることは容易に分かる。さらに遷移写像は明らかに全射である。
\end{proof}

\noindent
有限表示の環準同型は次のように特徴づけられる。
これはある意味で、有限表示の代数が $R$-代数の圏における
「コンパクト」対象であることを述べている。

\begin{lemma}
\label{algebra-lemma-characterize-finite-presentation}
$\varphi : R \to S$ を環準同型とする。次は同値である。
\begin{enumerate}
\item $\varphi$ は有限表示である。
\item 任意の $R$-代数の有向系 $A_\lambda$ に対して、写像
$$
\colim_\lambda \Hom_R(S, A_\lambda) \longrightarrow
\Hom_R(S, \colim_\lambda A_\lambda)
$$
は全単射である。
\item 任意の $R$-代数の有向系 $A_\lambda$ に対して、写像
$$
\colim_\lambda \Hom_R(S, A_\lambda) \longrightarrow
\Hom_R(S, \colim_\lambda A_\lambda)
$$
は全射である。
\end{enumerate}
\end{lemma}

\begin{proof}
(1) を仮定し、$S = R[x_1, \ldots, x_n] / (f_1, \ldots, f_m)$
と書く。$A = \colim A_\lambda$ とおく。
$R$-代数準同型 $S \to A$ または $S \to A_\lambda$ は
$x_1, \ldots, x_n$ の像によって決定されることに注意する。
従って
$\colim_\lambda \Hom_R(S, A_\lambda) \to \Hom_R(S, A)$
が単射であることは明らかである。
全射性を見るため、$\chi : S \to A$ を $R$-代数準同型とする。
このとき各 $x_i$ は、ある $A_{\lambda_i}$ の像に属する元に写る。
$\mu \geq \lambda_i$、$i = 1, \ldots, n$ を選び、
$i = 1, \ldots, n$ に対して $\chi(x_i)$ が
$y_i \in A_\mu$ の像であると仮定してよい。
$z_j = f_j(y_1, \ldots, y_n) \in A_\mu$ を考える。
$\chi$ は準同型なので、$z_j$ の
$A = \colim_\lambda A_\lambda$ における像は零である。
従って $z_j$ が $A_{\mu_j}$ において零に写るような
$\mu_j \geq \mu$ が存在する。
$\nu \geq \mu_j$、$j = 1, \ldots, m$ を選ぶ。
このとき $z_1, \ldots, z_m$ の像は $A_\nu$ において零である。
これはまさに、$y_i$ が関係
$f_j(y'_1, \ldots, y'_n) = 0$ を満たす元
$y'_i \in A_\nu$ に写ることを意味する。
従って環準同型 $S \to A_\nu$ を得る。ゆえに (1) は (2) を含意する。

\medskip\noindent
(2) が (3) を含意することは明らかである。(3) を仮定する。
補題 \ref{algebra-lemma-ring-colimit-fp} により、
$S_\lambda$ が $R$ 上有限表示であるように
$S = \colim_\lambda S_\lambda$ と書ける。
このとき、ある $\lambda$ に対して恒等写像は
$$
S \to S_\lambda \to S
$$
と分解する。補題 \ref{algebra-lemma-compose-finite-type} の (4) を
$S \to S_\lambda \to S$ に適用すれば、$S$ は $S_\lambda$ 上
有限表示であることが分かる。同じ補題の (2) を
$R \to S_\lambda \to S$ に適用すれば、
$S$ は $R$ 上有限表示であると結論できる。
\end{proof}

\noindent
% [P01-E0618] The frozen introduction lacks an article before “given type”.
上の基本事項を用いると、代数 $A$ が指定された型の代数の
フィルター余極限となるための次の判定条件を与えられる。

\begin{lemma}
\label{algebra-lemma-when-colimit}
$R \to \Lambda$ を環準同型とする。
$\mathcal{E}$ を $R$-代数の集合で、各 $A \in \mathcal{E}$ が
$R$ 上有限表示となるものとする。次の二つは同値である。
\begin{enumerate}
\item $\Lambda$ は $\mathcal{E}$ の元のフィルター余極限である。
% [P01-E0619] The frozen source writes “R algebra map” without the hyphen.
\item $R$ 代数準同型 $A \to \Lambda$ で $A$ が $R$ 上有限表示である
任意のものに対して、$B \in \mathcal{E}$ を満たす分解
$A \to B \to \Lambda$ を選べる。
\end{enumerate}
\end{lemma}

\begin{proof}
$\mathcal{I} \to \mathcal{E}$、$i \mapsto A_i$ を、
$\Lambda = \colim_i A_i$ となるフィルター図式とする。
$A \to \Lambda$ を、$A$ が $R$ 上有限表示である
$R$-代数準同型とする。補題 \ref{algebra-lemma-characterize-finite-presentation}
を適用すると分解 $A \to A_i \to \Lambda$ を得る。
従って (1) は (2) を含意する。

\medskip\noindent
補題 \ref{algebra-lemma-ring-colimit-fp-category} の圏 $\mathcal{I}$ を考える。
Categories、補題 \ref{categories-lemma-cofinal-in-filtered} により、
$A \in \mathcal{E}$ であるような $A \to \Lambda$ からなる
充満部分圏 $\mathcal{J}$ は $\mathcal{I}$ で共終であり、
フィルター圏である。すると Categories、補題
\ref{categories-lemma-cofinal} により、
$\Lambda$ は $\mathcal{J}$ 上の余極限でもある。
\end{proof}

\noindent
% [P01-E0620] The frozen prose says “limit” although the next lemma proves a colimit.
さらに強いことが成り立つ。すなわち
$R = \colim_\lambda R_\lambda$ が与えられたとき、
有限表示 $R$-加群の圏は有限表示 $R_\lambda$-加群の圏の極限と同値である。
有限表示 $R$-代数の圏についても同様である。

\begin{lemma}
\label{algebra-lemma-module-map-property-in-colimit}
$A$ を環とし、$M, N$ を $A$-加群とする。
$R = \colim_{i \in I} R_i$ を $A$-代数の有向余極限とする。
\begin{enumerate}
\item $M$ が有限 $A$-加群で、$u, u' : M \to N$ が
$u \otimes 1 = u' \otimes 1 : M \otimes_A R \to N \otimes_A R$
を満たす $A$-加群準同型ならば、ある $i$ に対して
$u \otimes 1 = u' \otimes 1 : M \otimes_A R_i \to N \otimes_A R_i$
である。
\item $N$ が有限 $A$-加群で、$u : M \to N$ が
$u \otimes 1 : M \otimes_A R \to N \otimes_A R$ を全射にする
$A$-加群準同型ならば、ある $i$ に対して写像
$u \otimes 1 : M \otimes_A R_i \to N \otimes_A R_i$ は全射である。
\item $N$ が有限表示 $A$-加群で、
$v : N \otimes_A R \to M \otimes_A R$ が $R$-加群準同型ならば、
ある $i$ と $R_i$-加群準同型
$v_i : N \otimes_A R_i \to M \otimes_A R_i$ であって
$v = v_i \otimes 1$ となるものが存在する。
\item $M$ が有限 $A$-加群、$N$ が有限表示 $A$-加群で、
$u : M \to N$ が
$u \otimes 1 : M \otimes_A R \to N \otimes_A R$ を同型にする
$A$-加群準同型ならば、ある $i$ に対して写像
$u \otimes 1 : M \otimes_A R_i \to N \otimes_A R_i$ は同型である。
\end{enumerate}
\end{lemma}

\begin{proof}
(1) を示すため、$u$ が (1) のようなものとし、
$x_1, \ldots, x_m \in M$ を生成元とする。
$N \otimes_A R = \colim_i N \otimes_A R_i$ なので、
% [P01-E0621] The frozen equality is placed in M tensor R_i although the images lie in N.
$M \otimes_A R_i$ において
$u(x_j) \otimes 1 = u'(x_j) \otimes 1$、$j = 1, \ldots, m$
となる $i \in I$ を選べる。このような $i$ に対して
$u \otimes 1 = u' \otimes 1 : M \otimes_A R_i \to N \otimes_A R_i$
である。

\medskip\noindent
(2) を示すため、$u \otimes 1$ が全射であると仮定し、
$y_1, \ldots, y_m \in N$ を生成元とする。
$N \otimes_A R = \colim_i N \otimes_A R_i$ なので、
% [P01-E0622] The frozen source omits the map u tensor 1 on the chosen preimages.
$N \otimes_A R$ における像が $y_j \otimes 1$ に等しくなる
$i \in I$ と $z_j \in M \otimes_A R_i$、$j = 1, \ldots, m$ を選べる。
このような $i$ に対して写像
$u \otimes 1 : M \otimes_A R_i \to N \otimes_A R_i$ は全射である。

\medskip\noindent
(3) を示すため、$y_1, \ldots, y_m \in N$ を生成元とする。
% [P01-E0623] The frozen presentation rule uses undefined x_j rather than y_j.
$K = \Ker(A^{\oplus m} \to N)$ とおく。ここで写像は
$(a_1, \ldots, a_m) \mapsto \sum a_j x_j$ で与えられる。
$k_1, \ldots, k_t$ を $K$ の生成元とし、
$k_s = (k_{s1}, \ldots, k_{sm})$ と書く。
$M \otimes_A R = \colim_i M \otimes_A R_i$ なので、
$M \otimes_A R$ における像が $v(y_j \otimes 1)$ に等しくなる
$i \in I$ と $z_j \in M \otimes_A R_i$、$j = 1, \ldots, m$ を選べる。
$z_j$ を用いて写像
$v_i : N \otimes_A R_i \to M \otimes_A R_i$ を定義したい。
$K \otimes_A R_i \to R_i^{\oplus m} \to N \otimes_A R_i \to 0$
は表示なので、各 $s = 1, \ldots, t$ に対して
$\xi_s = \sum_j k_{sj}z_j$ が
$M \otimes_A R_i$ において零であることを確認すれば十分である。
これは現時点では成り立たないかもしれないが、
$\xi_s$ の $M \otimes_A R$ における像は零なので、
$i$ を少し大きくすれば成り立つ。

\medskip\noindent
(4) を示すため、$u \otimes 1$ が同型であり、
$M$ が有限、$N$ が有限表示であると仮定する。
$v : N \otimes_A R \to M \otimes_A R$ を
$u \otimes 1$ の逆写像とする。(3) を適用して、ある $i$ に対する写像
$v_i : N \otimes_A R_i \to M \otimes_A R_i$ を得る。
% [P01-E0624] The frozen inverse identities switch the tensor base from A to R.
(1) を適用すれば、$i$ を大きくした後
$v_i \circ (u \otimes 1) = \text{id}_{M \otimes_R R_i}$ かつ
$(u \otimes 1) \circ v_i = \text{id}_{N \otimes_R R_i}$
となることが分かる。
\end{proof}

\begin{lemma}
\label{algebra-lemma-colimit-category-fp-modules}
$R = \colim_{\lambda \in \Lambda} R_\lambda$ を環の有向余極限とする。
このとき、有限表示 $R$-加群の圏は、有限表示
$R_\lambda$-加群の圏の余極限である。より正確には次が成り立つ。
\begin{enumerate}
\item 有限表示 $R$-加群 $M$ が与えられると、ある
$\lambda \in \Lambda$ と有限表示 $R_\lambda$-加群
$M_\lambda$ が存在して、$M \cong M_\lambda \otimes_{R_\lambda} R$ となる。
\item $\lambda \in \Lambda$、有限表示 $R_\lambda$-加群
$M_\lambda, N_\lambda$、および $R$-加群準同型
$\varphi : M_\lambda \otimes_{R_\lambda} R \to N_\lambda \otimes_{R_\lambda} R$
が与えられると、ある $\mu \geq \lambda$ と $R_\mu$-加群準同型
$\varphi_\mu : M_\lambda \otimes_{R_\lambda} R_\mu \to
N_\lambda \otimes_{R_\lambda} R_\mu$
が存在して、$\varphi = \varphi_\mu \otimes 1_R$ となる。
\item $\lambda \in \Lambda$、有限表示 $R_\lambda$-加群
$M_\lambda, N_\lambda$、および $R_\lambda$-加群準同型
$\varphi, \psi : M_\lambda \to N_\lambda$
が与えられ、$\varphi \otimes 1_R = \psi \otimes 1_R$ ならば、ある
$\mu \geq \lambda$ に対して
$\varphi \otimes 1_{R_\mu} = \psi \otimes 1_{R_\mu}$ となる。
\end{enumerate}
\end{lemma}

\begin{proof}
(1) を示すため、表示
$R^{\oplus m} \to R^{\oplus n} \to M \to 0$
を選ぶ。最初の写像が行列 $A = (a_{ij})$ で与えられるとする。
$R_\lambda$ に係数をもつ行列
$A_\lambda = (a_{\lambda, ij})$ で、$R$ において $A$ に写るものと、
$\lambda \in \Lambda$ を選べる。
そこで $M_\lambda$ を、最初の矢印が $A_\lambda$ で与えられる表示
$R_\lambda^{\oplus m} \to R_\lambda^{\oplus n} \to M_\lambda \to 0$
をもつ $R_\lambda$-加群とすればよい。

\medskip\noindent
(2) と (3) は補題 \ref{algebra-lemma-module-map-property-in-colimit} から従う。
\end{proof}

\begin{lemma}
\label{algebra-lemma-algebra-map-property-in-colimit}
$A$ を環とし、$B, C$ を $A$-代数とする。
$R = \colim_{i \in I} R_i$ を $A$-代数の有向余極限とする。
\begin{enumerate}
\item $B$ が有限型 $A$-代数で、$u, u' : B \to C$ が
$u \otimes 1 = u' \otimes 1 : B \otimes_A R \to C \otimes_A R$
を満たす $A$-代数準同型ならば、ある $i$ に対して
$u \otimes 1 = u' \otimes 1 : B \otimes_A R_i \to C \otimes_A R_i$
となる。
\item $C$ が有限型 $A$-代数で、$u : B \to C$ が
$u \otimes 1 : B \otimes_A R \to C \otimes_A R$ を全射にする
$A$-代数準同型ならば、ある $i$ に対して写像
$u \otimes 1 : B \otimes_A R_i \to C \otimes_A R_i$ は全射である。
\item $C$ が $A$ 上有限表示で、
$v : C \otimes_A R \to B \otimes_A R$ が $R$-代数準同型ならば、
ある $i$ と $R_i$-代数準同型
$v_i : C \otimes_A R_i \to B \otimes_A R_i$ が存在して
$v = v_i \otimes 1$ となる。
\item $B$ が有限型 $A$-代数、$C$ が有限表示 $A$-代数で、
$u \otimes 1 : B \otimes_A R \to C \otimes_A R$ が同型ならば、
ある $i$ に対して写像
$u \otimes 1 : B \otimes_A R_i \to C \otimes_A R_i$ は同型である。
\end{enumerate}
\end{lemma}

\begin{proof}
(1) を示すため、$u$ が (1) のとおりであると仮定し、
$x_1, \ldots, x_m \in B$ を生成元とする。
% [P01-E0625] Frozen part (1) takes a colimit of B and locates equal C-images in B tensor R_i.
$B \otimes_A R = \colim_i B \otimes_A R_i$ なので、
$B \otimes_A R_i$ において
$u(x_j) \otimes 1 = u'(x_j) \otimes 1$、$j = 1, \ldots, m$
となる $i \in I$ を選べる。このような $i$ に対して
$u \otimes 1 = u' \otimes 1 : B \otimes_A R_i \to C \otimes_A R_i$
である。

\medskip\noindent
(2) を示すため、$u \otimes 1$ が全射であると仮定し、
$y_1, \ldots, y_m \in C$ を生成元とする。
$B \otimes_A R = \colim_i B \otimes_A R_i$ なので、
% [P01-E0626] The frozen choice says the images of z_j in C tensor R without applying u tensor 1.
$C \otimes_A R$ における像が $y_j \otimes 1$ に等しくなる
$i \in I$ と $z_j \in B \otimes_A R_i$、$j = 1, \ldots, m$ を選べる。
このような $i$ に対して写像
$u \otimes 1 : B \otimes_A R_i \to C \otimes_A R_i$ は全射である。

\medskip\noindent
(3) を示すため、$c_1, \ldots, c_m \in C$ を生成元とする。
% [P01-E0627] The frozen kernel presentation has undefined codomain N.
$K = \Ker(A[x_1, \ldots, x_m] \to N)$ とおく。ここで写像は
% [P01-E0628] The frozen presentation sends each x_j to the sum of all c_j.
$x_j \mapsto \sum c_j$ で与えられる。$f_1, \ldots, f_t$ を
$A[x_1, \ldots, x_m]$ のイデアルとしての $K$ の生成元とする。
$f_j = f_j(x_1, \ldots, x_m)$ を多項式とみなす。
$B \otimes_A R = \colim_i B \otimes_A R_i$ なので、
$B \otimes_A R$ における像が $v(c_j \otimes 1)$ に等しくなる
$i \in I$ と $z_j \in B \otimes_A R_i$、$j = 1, \ldots, m$ を選べる。
$z_j$ を用いて写像
$v_i : C \otimes_A R_i \to B \otimes_A R_i$ を定義したい。
$K \otimes_A R_i \to R_i[x_1, \ldots, x_m] \to C \otimes_A R_i \to 0$
は表示なので、各 $s = 1, \ldots, t$ に対して
% [P01-E0629] The frozen relation uses f_j although the indexed element is xi_s.
$\xi_s = f_j(z_1, \ldots, z_m)$ が
$B \otimes_A R_i$ において零であることを確認すれば十分である。
これは現時点では成り立たないかもしれないが、$\xi_s$ の
$B \otimes_A R$ における像は零なので、$i$ を少し大きくすれば成り立つ。

\medskip\noindent
(4) を示すため、$u \otimes 1$ が同型であり、$B$ が有限型
$A$-代数、$C$ が有限表示 $A$-代数であると仮定する。
% [P01-E0630] The frozen alleged inverse v has the same direction as u tensor 1.
$v : B \otimes_A R \to C \otimes_A R$ を $u \otimes 1$ の逆写像とする。
(3) のような
$v_i : C \otimes_A R_i \to B \otimes_A R_i$ を選ぶ。
(1) を適用すれば、$i$ を大きくした後
% [P01-E0631] The frozen inverse identities switch the tensor base from A to R.
$v_i \circ (u \otimes 1) = \text{id}_{B \otimes_R R_i}$ かつ
$(u \otimes 1) \circ v_i = \text{id}_{C \otimes_R R_i}$
となることが分かる。
\end{proof}

\begin{lemma}
\label{algebra-lemma-colimit-category-fp-algebras}
$R = \colim_{\lambda \in \Lambda} R_\lambda$ を環の有向余極限とする。
このとき、有限表示 $R$-代数の圏は、有限表示
$R_\lambda$-代数の圏の余極限である。より正確には次が成り立つ。
\begin{enumerate}
\item 有限表示 $R$-代数 $A$ が与えられると、ある
$\lambda \in \Lambda$ と有限表示 $R_\lambda$-代数
$A_\lambda$ が存在して、$A \cong A_\lambda \otimes_{R_\lambda} R$ となる。
\item $\lambda \in \Lambda$、有限表示 $R_\lambda$-代数
$A_\lambda, B_\lambda$、および $R$-代数準同型
$\varphi : A_\lambda \otimes_{R_\lambda} R \to B_\lambda \otimes_{R_\lambda} R$
が与えられると、ある $\mu \geq \lambda$ と $R_\mu$-代数準同型
$\varphi_\mu : A_\lambda \otimes_{R_\lambda} R_\mu \to
B_\lambda \otimes_{R_\lambda} R_\mu$
が存在して、$\varphi = \varphi_\mu \otimes 1_R$ となる。
\item $\lambda \in \Lambda$、有限表示 $R_\lambda$-代数
$A_\lambda, B_\lambda$、および $R_\lambda$-代数準同型
$\varphi_\lambda, \psi_\lambda : A_\lambda \to B_\lambda$
が与えられ、
% [P01-E0632] The frozen equalities drop the lambda subscripts from both named maps.
$\varphi \otimes 1_R = \psi \otimes 1_R$ ならば、ある
$\mu \geq \lambda$ に対して
$\varphi \otimes 1_{R_\mu} = \psi \otimes 1_{R_\mu}$ となる。
\end{enumerate}
\end{lemma}

\begin{proof}
(1) を示すため、表示
$A = R[x_1, \ldots, x_n]/(f_1, \ldots, f_m)$
を選ぶ。$f_j \in R[x_1, \ldots, x_n]$ に写る元
$f_{\lambda, j} \in R_\lambda[x_1, \ldots, x_n]$ と
$\lambda \in \Lambda$ を選べる。
そこで単に
$A_\lambda =
R_\lambda[x_1, \ldots, x_n]/(f_{\lambda, 1}, \ldots, f_{\lambda, m})$
とおけばよい。

\medskip\noindent
(2) と (3) は補題 \ref{algebra-lemma-algebra-map-property-in-colimit} から従う。
\end{proof}

\begin{lemma}
\label{algebra-lemma-limit-no-condition-local}
$R \to S$ を局所環の局所準同型とする。
ある有向集合 $(\Lambda, \leq)$ と、局所環の局所準同型の系
$R_\lambda \to S_\lambda$ が存在して、次が成り立つ。
\begin{enumerate}
\item 系 $R_\lambda \to S_\lambda$ の余極限は $R \to S$ に等しい。
\item 各 $R_\lambda$ は $\mathbf{Z}$ 上本質的有限型である。
\item 各 $S_\lambda$ は $R_\lambda$ 上本質的有限型である。
\end{enumerate}
\end{lemma}

\begin{proof}
% [P01-E0633] The frozen designation omits “by”; Japanese renders the evident designation.
環準同型を $\varphi : R  \to S$ と表す。
$\mathfrak m \subset R$ を $R$ の極大イデアル、
$\mathfrak n \subset S$ を $S$ の極大イデアルとする。次をおく。
$$
\Lambda = \{
(A, B)
\mid
A \subset R, B \subset S, \# A < \infty, \# B < \infty, \varphi(A) \subset B
\}.
$$
半順序として包含関係をとる。各
$\lambda = (A, B) \in \Lambda$ に対し、$R'_\lambda$ を
$a \in A$ で生成される $\mathbf{Z}$-部分代数とし、
$S'_\lambda$ を $b$、$b \in B$ で生成される
$\mathbf{Z}$-部分代数とする。$R_\lambda$ を
素イデアル $R'_\lambda \cap \mathfrak m$ における
$R'_\lambda$ の局所化とし、$S_\lambda$ を
素イデアル $S'_\lambda \cap \mathfrak n$ における
$S'_\lambda$ の局所化とする。図示すれば
$$
\xymatrix{
B \ar[r] &
S'_\lambda \ar[r] &
S_\lambda \ar[r] &
S \\
A \ar[r] \ar[u] &
R'_\lambda \ar[r] \ar[u] &
R_\lambda \ar[r] \ar[u] &
R \ar[u]
}.
$$
遷移写像は明らかである。残りの主張の証明は読者に委ねる。
\end{proof}

\begin{lemma}
\label{algebra-lemma-limit-essentially-finite-type}
$R \to S$ を局所環の局所準同型とする。
$S$ が $R$ 上本質的有限型であると仮定する。
このとき、ある有向集合 $(\Lambda, \leq)$ と、局所環の局所準同型の系
$R_\lambda \to S_\lambda$ が存在して、次が成り立つ。
\begin{enumerate}
\item 系 $R_\lambda \to S_\lambda$ の余極限は $R \to S$ に等しい。
\item 各 $R_\lambda$ は $\mathbf{Z}$ 上本質的有限型である。
\item 各 $S_\lambda$ は $R_\lambda$ 上本質的有限型である。
\item 各 $\lambda \leq \mu$ に対し、写像
$S_\lambda \otimes_{R_\lambda} R_\mu \to S_\mu$ は、
$S_\mu$ を $S_\lambda \otimes_{R_\lambda} R_\mu$ の商の
局所化として表示する。
\end{enumerate}
\end{lemma}

\begin{proof}
環準同型を $\varphi : R  \to S$ と表す。
$\mathfrak m \subset R$ を $R$ の極大イデアル、
$\mathfrak n \subset S$ を $S$ の極大イデアルとする。
$x_1, \ldots, x_n \in S$ を、$S$ が
$x_1, \ldots, x_n$ で生成される $S$ の $R$-部分代数の局所化と
なるような元とする。言い換えれば、$S$ は多項式環
$R[x_1, \ldots, x_n]$ の局所化の商である。

\medskip\noindent
$\Lambda = \{ A \subset R \mid \# A < \infty\}$ を
$R$ の有限部分集合全体の集合とする。
半順序として包含関係をとる。各
$\lambda = A \in \Lambda$ に対し、$R'_\lambda$ を
$a \in A$ で生成される $\mathbf{Z}$-部分代数とし、
$S'_\lambda$ を $\varphi(a)$、$a \in A$ および元
$x_1, \ldots, x_n$ で生成される $\mathbf{Z}$-部分代数とする。
$R_\lambda$ を素イデアル $R'_\lambda \cap \mathfrak m$ における
$R'_\lambda$ の局所化とし、$S_\lambda$ を素イデアル
$S'_\lambda \cap \mathfrak n$ における $S'_\lambda$ の局所化とする。
図示すれば
$$
\xymatrix{
\varphi(A) \amalg \{x_i\} \ar[r] &
S'_\lambda \ar[r] &
S_\lambda \ar[r] &
S \\
A \ar[r] \ar[u] &
R'_\lambda \ar[r] \ar[u] &
R_\lambda \ar[r] \ar[u] &
R \ar[u]
}
$$
$A \subset B$ が $\Lambda$ における
$\lambda \leq \mu$ に対応するならば、標準写像
$R_\lambda \to R_\mu$ と $S_\lambda \to S_\mu$ が存在し、
有向集合 $\Lambda$ 上の系が得られることは明らかである。

\medskip\noindent
すべての写像 $R_\lambda \to R$ は単射で、$R$ の任意の元は
やがて像に入るので、$R = \colim R_\lambda$ であることは明らかである。
同じ議論が $S = \colim S_\lambda$ にも適用できる。
主張 (2)、(3) は構成から成り立つ。最後の主張は、写像
$S'_\lambda \otimes_{R'_\lambda} R'_\mu
\to S'_\mu$ が明らかに全射であることから従う。
\end{proof}

\begin{lemma}
\label{algebra-lemma-limit-essentially-finite-presentation}
$R \to S$ を局所環の局所準同型とする。
$S$ が $R$ 上本質的有限表示であると仮定する。
このとき、ある有向集合 $(\Lambda, \leq)$ と、局所環の局所準同型の系
% [P01-E0634] The frozen English has singular “homomorphism” after “a system of”.
$R_\lambda \to S_\lambda$ が存在して、次が成り立つ。
\begin{enumerate}
\item 系 $R_\lambda \to S_\lambda$ の余極限は $R \to S$ に等しい。
\item 各 $R_\lambda$ は $\mathbf{Z}$ 上本質的有限型である。
\item 各 $S_\lambda$ は $R_\lambda$ 上本質的有限型である。
\item 各 $\lambda \leq \mu$ に対し、写像
$S_\lambda \otimes_{R_\lambda} R_\mu \to S_\mu$ は、
$S_\mu$ を素イデアルにおける
$S_\lambda \otimes_{R_\lambda} R_\mu$ の局所化として表示する。
\end{enumerate}
\end{lemma}

\begin{proof}
仮定により同型
$\Phi : (R[x_1, \ldots, x_n]/I)_{\mathfrak q} \to S$
を選べる。ここで $I \subset R[x_1, \ldots, x_n]$ は有限生成イデアル、
$\mathfrak q \subset R[x_1, \ldots, x_n]/I$ は素イデアルである。
（$R \cap \mathfrak q$ は $R$ の極大イデアル
$\mathfrak m$ に等しいことに注意せよ。）
さらに、イデアル $I$ の生成元 $f_1, \ldots, f_m \in I$ を選ぶ。
$R$ を、各 $R_\lambda$ が局所的かつ $\mathbf{Z}$ 上本質的有限型となる
有向集合 $(\Lambda, \leq )$ 上の余極限
$R = \colim R_\lambda$ と任意の仕方で書く。
ある $\lambda_0 \in \Lambda$ が存在して、$f_j$ はある
$f_{j, \lambda_0} \in R_{\lambda_0}[x_1, \ldots, x_n]$ の像となる。
すべての $\lambda \geq \lambda_0$ に対し、
$f_{j, \lambda} \in R_{\lambda}[x_1, \ldots, x_n]$ で
$f_{j, \lambda_0}$ の像を表す。こうして環準同型の系
$$
R_\lambda[x_1, \ldots, x_n]/(f_{1, \lambda}, \ldots, f_{m, \lambda})
\to
R[x_1, \ldots, x_n]/(f_1, \ldots, f_m) \to S
$$
を得る。
% [P01-E0635] The frozen designation of q_lambda omits “to be”.
$\mathfrak q_\lambda$ を $\mathfrak q$ の逆像とする。
$S_\lambda = (R_\lambda[x_1, \ldots, x_n]/
(f_{1, \lambda}, \ldots, f_{m, \lambda}))_{\mathfrak q_\lambda}$ とおく。
これでよいことの確認は読者に委ねる。
\end{proof}

\begin{remark}
\label{algebra-remark-suitable-systems-limits}
$R \to S$ を、本質的有限表示である局所環の局所準同型とする。
補題 \ref{algebra-lemma-limit-essentially-finite-type} に列挙した性質をもつ
任意の系 $(\Lambda, \leq)$、$R_\lambda \to S_\lambda$ をとる。
これが「誤った」系、すなわち補題
\ref{algebra-lemma-limit-essentially-finite-presentation} の性質 (4) を
満たさないことがありうる。以下がその例である。
$k = \mathbf{F}_2$ とする。環
$$
R = \text{localization of }
k[z, y_1, y_2, \ldots]/(y_i^2 - zy_{i + 1})
\text{ at }(z, y_1, y_2, \ldots)
$$
を考える。$S = R/zR$ とおく。系として $\Lambda = \mathbf{N}$ と
$$
R_n = \text{localization of }
k[z, y_1, \ldots, y_n]/(\{y_i^2 - zy_{i + 1}\}_{i \leq n-1})
\text{ at }(z, y_1, \ldots, y_n)
$$
をとり、$S_n = R_n/(z, y_n^2)$ とする。
% [P01-E0636] The frozen “all ... not” construction has ambiguous quantifier scope.
写像 $S_n \otimes_{R_n} R_{n + 1} \to S_{n + 1}$ はいずれも
局所化（この場合は同型）ではない。実際、
$1 \otimes y_{n + 1}^2$ は零に写るからである。
代わりに $S_n' = R_n/zR_n$ とすれば、写像
$S'_n \otimes_{R_n} R_{n + 1} \to S'_{n + 1}$ は同型である。
この注意の教訓は、系の選択には多少慎重でなければならないということである。
\end{remark}

\begin{lemma}
\label{algebra-lemma-limit-module-essentially-finite-presentation}
$R \to S$ を局所環の局所準同型とする。
$S$ が $R$ 上本質的有限表示であると仮定する。
$M$ を有限表示 $S$-加群とする。このとき、ある有向集合
$(\Lambda, \leq)$ と局所環の局所準同型の系
$R_\lambda \to S_\lambda$、および $S_\lambda$-加群
$M_\lambda$ が存在して、次が成り立つ。
\begin{enumerate}
\item 系 $R_\lambda \to S_\lambda$ の余極限は $R \to S$ に等しい。
系 $M_\lambda$ の余極限は $M$ である。
\item 各 $R_\lambda$ は $\mathbf{Z}$ 上本質的有限型である。
\item 各 $S_\lambda$ は $R_\lambda$ 上本質的有限型である。
\item 各 $M_\lambda$ は $S_\lambda$ 上有限である。
\item 各 $\lambda \leq \mu$ に対し、写像
$S_\lambda \otimes_{R_\lambda} R_\mu \to S_\mu$ は、
$S_\mu$ を素イデアルにおける
$S_\lambda \otimes_{R_\lambda} R_\mu$ の局所化として表示する。
\item 各 $\lambda \leq \mu$ に対し、写像
$M_\lambda \otimes_{S_\lambda} S_\mu \to M_\mu$ は同型である。
\end{enumerate}
\end{lemma}

\begin{proof}
補題 \ref{algebra-lemma-limit-essentially-finite-presentation} の証明と同様に、
まず $R = \colim R_\lambda$ を、本質的有限型な局所
$\mathbf{Z}$-代数の有向余極限として書ける。次に、ある
$\lambda_1 \in \Lambda$ に対して
$f_{j, \lambda_1} \in R_{\lambda_1}[x_1, \ldots, x_n]$ が存在して
$$
S =
\colim_{\lambda \geq \lambda_1} S_\lambda, \text{ with }
S_\lambda =
(R_\lambda[x_1, \ldots, x_n]/
(f_{1, \lambda}, \ldots, f_{m, \lambda}))_{\mathfrak q_\lambda}
$$
となると仮定してよい。$S$ 上の $M$ の表示
$$
S^{\oplus s} \to S^{\oplus t} \to M \to 0
$$
を選ぶ。$A \in \text{Mat}(t \times s, S)$ をこの表示の行列とする。
ある $\lambda_2 \in \Lambda$、$\lambda_2 \geq \lambda_1$ に対し、
$A$ に写る行列
$A_{\lambda_2} \in \text{Mat}(t \times s, S_{\lambda_2})$ を選べる。
すべての $\lambda \geq \lambda_2$ に対して
$M_\lambda = \Coker(S_\lambda^{\oplus s} \xrightarrow{A_\lambda}
S_\lambda^{\oplus t})$ とおく。これでよいことの確認は読者に委ねる。
\end{proof}

\begin{lemma}
\label{algebra-lemma-limit-no-condition}
$R \to S$ を環準同型とする。このとき、ある有向集合
$(\Lambda, \leq)$ と環準同型の系 $R_\lambda \to S_\lambda$ が存在して、
次が成り立つ。
\begin{enumerate}
\item 系 $R_\lambda \to S_\lambda$ の余極限は $R \to S$ に等しい。
\item 各 $R_\lambda$ は $\mathbf{Z}$ 上有限型である。
\item 各 $S_\lambda$ は $R_\lambda$ 上有限型である。
\end{enumerate}
\end{lemma}

\begin{proof}
これは補題 \ref{algebra-lemma-limit-no-condition-local} の非局所版である。
証明も同様なので読者に委ねる。
\end{proof}

\begin{lemma}
\label{algebra-lemma-limit-integral}
$R \to S$ を環準同型とする。$S$ が $R$ 上整であると仮定する。
このとき、ある有向集合 $(\Lambda, \leq)$ と環準同型の系
$R_\lambda \to S_\lambda$ が存在して、次が成り立つ。
\begin{enumerate}
\item 系 $R_\lambda \to S_\lambda$ の余極限は $R \to S$ に等しい。
\item 各 $R_\lambda$ は $\mathbf{Z}$ 上有限型である。
\item 各 $S_\lambda$ は $R_\lambda$ 上有限である。
\end{enumerate}
\end{lemma}

\begin{proof}
$E \subset R$ が有限部分集合、$F \subset S$ が有限部分集合で、
各元 $f \in F$ が、係数が $E$ に属するモニック多項式
$P(X) \in R[X]$ の根であるような対 $(E, F)$ の集合を
$\Lambda$ とする。$E \subset E'$ かつ $F \subset F'$ のとき
$(E, F) \leq (E', F')$ と定める。
$\lambda = (E, F) \in \Lambda$ が与えられたとき、
$R_\lambda \subset R$ を $E$ で生成される $R$ の
$\mathbf{Z}$-部分代数、$S_\lambda \subset S$ を $F$ と
$S$ における $E$ の像で生成される $\mathbf{Z}$-部分代数とする。
$R = \colim R_\lambda$ であることは明らかである。
% [P01-E0637] The frozen proof says every element of S is integral over S, not over R.
$S$ のすべての元が $S$ 上整であるから、
$S = \colim S_\lambda$ である。環準同型
$R_\lambda \to S_\lambda$ が有限であることは、補題
\ref{algebra-lemma-characterize-finite-in-terms-of-integral} と、対
$(E, F)$ に課した条件により $F$ の元が $R_\lambda$ 上整であり、
$S_\lambda$ がそれらによって $R_\lambda$ 上生成されることから従う。
これで補題が従う。
\end{proof}

\begin{lemma}
\label{algebra-lemma-limit-finite-type}
$R \to S$ を環準同型とする。$S$ が $R$ 上有限型であると仮定する。
このとき、ある有向集合 $(\Lambda, \leq)$ と環準同型の系
$R_\lambda \to S_\lambda$ が存在して、次が成り立つ。
\begin{enumerate}
\item 系 $R_\lambda \to S_\lambda$ の余極限は $R \to S$ に等しい。
\item 各 $R_\lambda$ は $\mathbf{Z}$ 上有限型である。
\item 各 $S_\lambda$ は $R_\lambda$ 上有限型である。
\item 各 $\lambda \leq \mu$ に対し、写像
$S_\lambda \otimes_{R_\lambda} R_\mu \to S_\mu$ は、
$S_\mu$ を $S_\lambda \otimes_{R_\lambda} R_\mu$ の商として表示する。
\end{enumerate}
\end{lemma}

\begin{proof}
これは補題 \ref{algebra-lemma-limit-essentially-finite-type} の非局所版である。
証明も同様なので読者に委ねる。
\end{proof}

\begin{lemma}
\label{algebra-lemma-limit-finite-presentation}
$R \to S$ を環準同型とする。$S$ が $R$ 上有限表示であると仮定する。
このとき、ある有向集合 $(\Lambda, \leq)$ と環準同型の系
$R_\lambda \to S_\lambda$ が存在して、次が成り立つ。
\begin{enumerate}
\item 系 $R_\lambda \to S_\lambda$ の余極限は $R \to S$ に等しい。
\item 各 $R_\lambda$ は $\mathbf{Z}$ 上有限型である。
\item 各 $S_\lambda$ は $R_\lambda$ 上有限型である。
\item 各 $\lambda \leq \mu$ に対し、写像
$S_\lambda \otimes_{R_\lambda} R_\mu \to S_\mu$ は同型である。
\end{enumerate}
\end{lemma}

\begin{proof}
これは補題 \ref{algebra-lemma-limit-essentially-finite-presentation} の非局所版である。
証明も同様なので読者に委ねる。
\end{proof}

\begin{lemma}
\label{algebra-lemma-limit-module-finite-presentation}
$R \to S$ を環準同型とする。$S$ が $R$ 上有限表示であると仮定する。
$M$ を有限表示 $S$-加群とする。このとき、ある有向集合
$(\Lambda, \leq)$ と環準同型の系 $R_\lambda \to S_\lambda$、および
$S_\lambda$-加群 $M_\lambda$ が存在して、次が成り立つ。
\begin{enumerate}
\item 系 $R_\lambda \to S_\lambda$ の余極限は $R \to S$ に等しい。
系 $M_\lambda$ の余極限は $M$ である。
\item 各 $R_\lambda$ は $\mathbf{Z}$ 上有限型である。
\item 各 $S_\lambda$ は $R_\lambda$ 上有限型である。
\item 各 $M_\lambda$ は $S_\lambda$ 上有限である。
\item 各 $\lambda \leq \mu$ に対し、写像
$S_\lambda \otimes_{R_\lambda} R_\mu \to S_\mu$ は同型である。
\item 各 $\lambda \leq \mu$ に対し、写像
$M_\lambda \otimes_{S_\lambda} S_\mu \to M_\mu$ は同型である。
\end{enumerate}
特に、すべての $\lambda \in \Lambda$ に対して
$$
M = M_\lambda \otimes_{S_\lambda} S
= M_\lambda \otimes_{R_\lambda} R.
$$
\end{lemma}

\begin{proof}
これは補題 \ref{algebra-lemma-limit-module-essentially-finite-presentation} の
非局所版である。証明も同様なので読者に委ねる。
\end{proof}














\section{さらなる平坦性判定法}
\label{algebra-section-more-flatness-criteria}

\noindent
次の補題は、正規曲面から非特異曲面への有限射が平坦であることを示すために、
代数幾何でしばしば用いられる。単射な有限局所環準同型は明らかに条件 (3) を
満たすので、これは補題 \ref{algebra-lemma-finite-flat-over-regular-CM} の部分的な逆である。

\begin{lemma}
\label{algebra-lemma-CM-over-regular-flat}
\begin{slogan}
奇跡的平坦性
\end{slogan}
$R \to S$ をネーター局所環の局所準同型とする。次を仮定する。
\begin{enumerate}
\item $R$ は正則である。
% [P01-E0638] The frozen second assumption omits the finite verb “is”.
\item $S$ は Cohen--Macaulay である。
\item $\dim(S) = \dim(R) + \dim(S/\mathfrak m_R S)$.
\end{enumerate}
このとき $R \to S$ は平坦である。
\end{lemma}

\begin{proof}
$\dim(R)$ に関する帰納法を用いる。$\dim(R) = 0$ の場合は、
$R$ が体なので自明である。$\dim(R) > 0$ と仮定する。
(3) により $\dim(S) > 0$ でもある。
$\mathfrak q_1, \ldots, \mathfrak q_r$ を $S$ の極小素イデアルとする。
次に注意する。
$\mathfrak q_i \not \supset \mathfrak m_R S$ である。実際、
$$
\dim(S/\mathfrak q_i) = \dim(S) > \dim(S/\mathfrak m_R S)
$$
であり、最初の等号は補題 \ref{algebra-lemma-maximal-chain-CM} により、
不等号は (3) による。
% [P01-E0639] The frozen English lacks punctuation after the displayed inequality.
従って $\mathfrak p_i = R \cap \mathfrak q_i$ は $\mathfrak m_R$ に等しくない。
補題 \ref{algebra-lemma-silly} により、$x \in \mathfrak m_R$、
$x \not \in \mathfrak m_R^2$、かつ $x \not \in \mathfrak p_i$ となる
元を選ぶ。従って $x$ は $S$ のどの極小素イデアルにも属さない。
従って (2) と補題 \ref{algebra-lemma-reformulate-CM} により、$x$ は $S$ 上の
非零因子であり、$S/xS$ は
$\dim(S/xS) = \dim(S) - 1$ を満たす Cohen--Macaulay 環である。
(1) と補題 \ref{algebra-lemma-regular-ring-CM} により、環 $R/xR$ は
$\dim(R/xR) = \dim(R) - 1$ を満たす正則環である。
帰納法により $R/xR \to S/xS$ は平坦である。従って補題
\ref{algebra-lemma-variant-local-criterion-flatness} とその直後の注意から結論を得る。
\end{proof}

\begin{lemma}
\label{algebra-lemma-flat-over-regular}
$R \to S$ をネーター局所環の準同型とする。
$R$ が正則局所環であり、正則パラメータ系が $S$ の正則列に写ると仮定する。
このとき $R \to S$ は平坦である。
\end{lemma}

\begin{proof}
$x_1, \ldots, x_d$ を、$S$ の正則列に写る $R$ のパラメータ系とする。
後者は体なので、$S/(x_1, \ldots, x_d)S$ は
$R/(x_1, \ldots, x_d)$ 上平坦であることに注意する。
次に $x_d$ は $S/(x_1, \ldots, x_{d - 1})S$ における非零因子なので、
平坦性の局所判定法（補題 \ref{algebra-lemma-variant-local-criterion-flatness} と
その後の注意を参照）により、$S/(x_1, \ldots, x_{d - 1})S$ は
$R/(x_1, \ldots, x_{d - 1})$ 上平坦である。
さらに $x_{d - 1}$ は $S/(x_1, \ldots, x_{d - 2})S$ における
非零因子なので、平坦性の局所判定法（補題
\ref{algebra-lemma-variant-local-criterion-flatness} とその後の注意を参照）により、
$S/(x_1, \ldots, x_{d - 2})S$ は
$R/(x_1, \ldots, x_{d - 2})$ 上平坦である。
この操作を続ければ、$S$ が $R$ 上平坦であるという結論に達する。
\end{proof}

\noindent
次の補題は、基礎環上有限表示な環上の有限表示加群に関する結果を、
ネーターの場合の有限加群に関する対応する結果から導くための鍵である。

\begin{lemma}
\label{algebra-lemma-colimit-eventually-flat}
$R \to S$、$M$、$\Lambda$、$R_\lambda \to S_\lambda$、$M_\lambda$ を
補題 \ref{algebra-lemma-limit-module-essentially-finite-presentation} のとおりとする。
$M$ が $R$ 上平坦であると仮定する。このとき、ある
$\lambda \in \Lambda$ に対して加群 $M_\lambda$ は
$R_\lambda$ 上平坦である。
\end{lemma}

\begin{proof}
ある $\lambda \in \Lambda$ を選び、次を考える。
$$
\text{Tor}_1^{R_\lambda}(M_\lambda, R_\lambda/\mathfrak m_\lambda)
=
\Ker(\mathfrak m_\lambda \otimes_{R_\lambda} M_\lambda
\to M_\lambda).
$$
注意 \ref{algebra-remark-Tor-ring-mod-ideal} を参照せよ。右辺から、これは有限生成
$S_\lambda$-加群であることが分かる（$S_\lambda$ がネーターで、
問題の加群が有限だからである）。$\xi_1, \ldots, \xi_n$ を生成元とする。
$M$ は $R$ 上平坦なので、
$0 = \Ker(\mathfrak m_\lambda R \otimes_R M \to M)$ である。
$\otimes$ は余極限と可換するから、各 $\xi_i$ が
$\mathfrak m_{\lambda}R_{\lambda'} \otimes_{R_{\lambda'}} M_{\lambda'}$
において零に写るような $\lambda' \geq \lambda$ が存在する。
従って写像
$$
\text{Tor}_1^{R_\lambda}(M_\lambda, R_\lambda/\mathfrak m_\lambda)
\longrightarrow
\text{Tor}_1^{R_{\lambda'}}(M_{\lambda'},
R_{\lambda'}/\mathfrak m_{\lambda}R_{\lambda'})
$$
は零である。$M_\lambda \otimes_{R_\lambda} R_\lambda/\mathfrak m_\lambda$
は体である環 $R_\lambda/\mathfrak m_\lambda$ 上平坦であることに注意する。
従って補題 \ref{algebra-lemma-another-variant-local-criterion-flatness} を適用して、
$M_{\lambda'}$ が $R_{\lambda'}$ 上平坦であることを得る。
\end{proof}

\noindent
上の補題を用いて、節 \ref{algebra-section-criteria-flatness} の結果を
非ネーターの場合に再証明し始めることができる。

\begin{lemma}
\label{algebra-lemma-mod-injective-general}
$R \to S$ を局所環の局所準同型とする。
% [P01-E0640] One of three frozen nearby designation constructions omits “by”.
$\mathfrak m$ を $R$ の極大イデアルとする。
$u : M \to N$ を $S$-加群の準同型とする。次を仮定する。
\begin{enumerate}
\item $S$ は $R$ 上本質的有限表示である。
\item $M$、$N$ は $S$ 上有限表示である。
\item $N$ は $R$ 上平坦である。
\item $\overline{u} : M/\mathfrak mM \to N/\mathfrak mN$ は単射である。
\end{enumerate}
このとき $u$ は単射であり、$N/u(M)$ は $R$ 上平坦である。
\end{lemma}

\begin{proof}
補題 \ref{algebra-lemma-limit-module-essentially-finite-presentation} とその証明により、
$S_\lambda$-加群の準同型
$u_\lambda : M_\lambda \to N_\lambda$ を伴う局所環準同型の系
$R_\lambda \to S_\lambda$ を、$N$ と $M$ の双方について同補題の
% [P01-E0641] Japanese makes the frozen modifier attachment explicit.
結論 (1) -- (6) を満たし、写像 $u_\lambda$ の余極限が $u$ となるように選べる。
補題 \ref{algebra-lemma-colimit-eventually-flat} により、十分大きいすべての
$\lambda$ に対して $N_\lambda$ が $R_\lambda$ 上平坦であると
仮定してよい。$\mathfrak m_\lambda \subset R_\lambda$ を極大イデアルとし、
$\kappa_\lambda = R_\lambda / \mathfrak m_\lambda$、および
$\kappa = R/\mathfrak m$ をそれぞれ剰余体とする。

\medskip\noindent
写像
$$
\Psi_\lambda :
M_\lambda/\mathfrak m_\lambda M_\lambda \otimes_{\kappa_\lambda} \kappa
\longrightarrow
M/\mathfrak m M.
$$
を考える。$S_\lambda/\mathfrak m_\lambda S_\lambda$ は体
$\kappa_\lambda$ 上本質的有限型なので、テンソル積
$S_\lambda/\mathfrak m_\lambda S_\lambda \otimes_{\kappa_\lambda} \kappa$
は $\kappa$ 上本質的有限型である。従ってこれはネーター環であり、
$\Psi_\lambda$ の核は有限生成であると結論できる。
$M/\mathfrak m M$ は系 $M_\lambda/\mathfrak m_\lambda M_\lambda$ の余極限で、
$\kappa$ は体 $\kappa_\lambda$ の余極限なので、
% [P01-E0642] The frozen strict relation is stronger than directedness supplies.
$\lambda' > \lambda$ が存在して、$\Psi_\lambda$ の核は次の写像の核で
生成される。
$$
\Psi_{\lambda, \lambda'} :
M_\lambda/\mathfrak m_\lambda M_\lambda
\otimes_{\kappa_\lambda}
\kappa_{\lambda'}
\longrightarrow
M_{\lambda'}/\mathfrak m_{\lambda'} M_{\lambda'}.
$$
構成により、乗法的部分集合
$W \subset S_\lambda \otimes_{R_\lambda} R_{\lambda'}$ が存在して、
$S_{\lambda'} = W^{-1}(S_\lambda \otimes_{R_\lambda} R_{\lambda'})$ かつ
$$
W^{-1}(M_\lambda/\mathfrak m_\lambda M_\lambda
\otimes_{\kappa_\lambda}
\kappa_{\lambda'})
=
M_{\lambda'}/\mathfrak m_{\lambda'} M_{\lambda'}.
$$
ここで $x$ を写像
$$
\Psi_{\lambda'} :
M_{\lambda'}/\mathfrak m_{\lambda'} M_{\lambda'}
\otimes_{\kappa_{\lambda'}} \kappa
\longrightarrow
M/\mathfrak m M.
$$
の核の元とする。するとある $w \in W$ に対して
% [P01-E0643] The frozen membership omits the tensor base and representative/image qualification.
$wx \in M_\lambda/\mathfrak m_\lambda M_\lambda \otimes \kappa$ となる。
従って $wx \in \Ker(\Psi_\lambda)$ である。ゆえに $wx$ は
$\Psi_{\lambda, \lambda'}$ の核の元の線形結合である。
従って $M_{\lambda'}/\mathfrak m_{\lambda'} M_{\lambda'}
\otimes_{\kappa_{\lambda'}} \kappa$ において $wx = 0$ であり、
$w$ は $S_{\lambda'}$ で可逆なので $x = 0$ である。
十分大きいすべての $\lambda'$ に対して $\Psi_{\lambda'}$ の核は
零であると結論する。

\medskip\noindent
前段落の結果により、十分大きいすべての $\lambda$ に対して
$\Psi_\lambda$ の核が零であると仮定してよい。このことから写像
$M_\lambda/\mathfrak m_\lambda M_\lambda \to M/\mathfrak m M$ は単射である。
これと $\overline{u}$ の単射性を合わせると、形式的に
$\overline{u_\lambda} : M_\lambda/\mathfrak m_\lambda M_\lambda
\to N_\lambda/\mathfrak m_\lambda N_\lambda$ も単射であることが従う。
補題 \ref{algebra-lemma-mod-injective} により、十分大きいすべての
$\lambda$ に対して写像 $u_\lambda$ は単射であり、
$N_\lambda/u_\lambda(M_\lambda)$ は $R_\lambda$ 上平坦である。
これで補題が従う。
\end{proof}

\begin{lemma}
\label{algebra-lemma-grothendieck-general}
$R \to S$ を局所環の局所環準同型とする。
% [P01-E0640] Second frozen nearby designation construction omits “by”.
$\mathfrak m$ を $R$ の極大イデアルとする。次を仮定する。
\begin{enumerate}
\item $S$ は $R$ 上本質的有限表示である。
\item $S$ は $R$ 上平坦である。
\item $f \in S$ は $S/{\mathfrak m}S$ における非零因子である。
\end{enumerate}
このとき $S/fS$ は $R$ 上平坦であり、$f$ は $S$ における非零因子である。
\end{lemma}

\begin{proof}
% [P01-E0645] The frozen proof is a subjectless sentence fragment.
これは補題 \ref{algebra-lemma-mod-injective-general} から直ちに従う。
\end{proof}

\begin{lemma}
\label{algebra-lemma-grothendieck-regular-sequence-general}
$R \to S$ を局所環の局所環準同型とする。
% [P01-E0640] Third frozen nearby designation construction omits “by”.
$\mathfrak m$ を $R$ の極大イデアルとする。次を仮定する。
\begin{enumerate}
\item $R \to S$ は本質的有限表示である。
\item $R \to S$ は平坦である。
\item $f_1, \ldots, f_c$ は $S$ の元の列であり、その像
$\overline{f}_1, \ldots, \overline{f}_c$ は
$S/{\mathfrak m}S$ において正則列をなす。
\end{enumerate}
このとき $f_1, \ldots, f_c$ は $S$ の正則列であり、各商
$S/(f_1, \ldots, f_i)$ は $R$ 上平坦である。
\end{lemma}

\begin{proof}
% [P01-E0646] The frozen proof consists only of a noun phrase.
帰納法と補題 \ref{algebra-lemma-grothendieck-general} を用いる。
\end{proof}

\noindent
局所有限表示な局所環準同型の場合に対する平坦性の局所判定法は次のとおりである。

\begin{lemma}
\label{algebra-lemma-variant-local-criterion-flatness-general}
$R \to S$ を局所環の局所準同型とする。
$I \not = R$ を $R$ のイデアルとし、$M$ を $S$-加群とする。次を仮定する。
\begin{enumerate}
\item $S$ は $R$ 上本質的有限表示である。
\item $M$ は $S$ 上有限表示である。
\item $\text{Tor}_1^R(M, R/I) = 0$ である。
\item $M/IM$ は $R/I$ 上平坦である。
\end{enumerate}
このとき $M$ は $R$ 上平坦である。
\end{lemma}

\begin{proof}
$\Lambda$、$R_\lambda \to S_\lambda$、$M_\lambda$ を補題
\ref{algebra-lemma-limit-module-essentially-finite-presentation} のとおりとする。
% [P01-E0647] The frozen designation of the inverse-image ideal omits “by”.
$I_\lambda \subset R_\lambda$ を $I$ の逆像とする。この場合、系
% [P01-E0648] The frozen quotient system starts from R_lambda rather than its quotient.
$R/I \to S/IS$、$M/IM$、$R_\lambda \to S_\lambda/I_\lambda S_\lambda$、
および $M_\lambda/I_\lambda M_\lambda$ も補題
\ref{algebra-lemma-limit-module-essentially-finite-presentation} の結論を満たす。
従って補題 \ref{algebra-lemma-colimit-eventually-flat} により、添字集合
$\Lambda$ を縮小した後、
% [P01-E0649] The frozen sentence leaves the flatness base unspecified.
すべての $\lambda$ に対して $M_\lambda/I_\lambda M_\lambda$ が
平坦であると仮定してよい。ある $\lambda$ を選び、次を考える。
$$
\text{Tor}_1^{R_\lambda}(M_\lambda, R_\lambda/I_\lambda)
=
\Ker(I_\lambda \otimes_{R_\lambda} M_\lambda
\to M_\lambda).
$$
注意 \ref{algebra-remark-Tor-ring-mod-ideal} を参照せよ。右辺から、これは有限生成
$S_\lambda$-加群であることが分かる（$S_\lambda$ がネーターで、
問題の加群が有限だからである）。$\xi_1, \ldots, \xi_n$ を生成元とする。
$\text{Tor}_1^R(M, R/I) = 0$ であり、$\otimes$ は余極限と可換するので、
各 $\xi_i$ が
$\text{Tor}_1^{R_{\lambda'}}(M_{\lambda'}, R_{\lambda'}/I_{\lambda'})$
において零に写るような $\lambda' \geq \lambda$ が存在する。
写像の合成
$$
\xymatrix{
R_{\lambda'} \otimes_{R_\lambda}
\text{Tor}_1^{R_\lambda}(M_\lambda, R_\lambda/I_\lambda)
\ar[d]^{\text{surjective by Lemma \ref{algebra-lemma-surjective-on-tor-one}}} \\
\text{Tor}_1^{R_\lambda}(M_\lambda, R_{\lambda'}/I_\lambda R_{\lambda'})
\ar[d]^{\text{surjective up to localization by
Lemma \ref{algebra-lemma-surjective-on-tor-one-trivial}}} \\
\text{Tor}_1^{R_{\lambda'}}(M_{\lambda'}, R_{\lambda'}/I_\lambda R_{\lambda'})
\ar[d]^{\text{surjective by Lemma \ref{algebra-lemma-surjective-on-tor-one}}} \\
\text{Tor}_1^{R_{\lambda'}}(M_{\lambda'}, R_{\lambda'}/I_{\lambda'}).
}
$$
は、示した理由により局所化を除いて全射である。
$M_{\lambda'}$ は $M_\lambda \otimes_{R_\lambda} R_{\lambda'}$ に
等しくないので、局所化が必要である。実際、これは
$M_\lambda \otimes_{S_\lambda} S_{\lambda'}$ に等しく、
$S_{\lambda'}$ は $S_{\lambda} \otimes_{R_\lambda} R_{\lambda'}$ の
局所化である。従って、
% [P01-E0650] Japanese supplies a complete consequence for the frozen fragment.
局所化を除いた（または $S_{\lambda'}$ をテンソルした）主張が得られる。
補題 \ref{algebra-lemma-surjective-on-tor-one} が最初と三番目の矢印に適用できるのは、
$M_\lambda/I_\lambda M_\lambda$ が $R_\lambda/I_\lambda$ 上平坦であり、
$M_{\lambda'}/I_\lambda M_{\lambda'}$ が
$R_{\lambda'}/I_\lambda R_{\lambda'}$ 上平坦だからである。後者は平坦加群
$M_\lambda/I_\lambda M_\lambda$ の基底変換だからである。
上で説明したように、この合成は生成元 $\xi_i$ を零に写す。
最後に
$\text{Tor}_1^{R_{\lambda'}}(M_{\lambda'}, R_{\lambda'}/I_{\lambda'})$
は零であると結論する。補題 \ref{algebra-lemma-variant-local-criterion-flatness} により、
これは $M_{\lambda'}$ が $R_{\lambda'}$ 上平坦であることを含意する。
\end{proof}

\noindent
次の補題を、補題 \ref{algebra-lemma-criterion-flatness-fibre-Noetherian}
（ネーター局所環の場合）および補題
\ref{algebra-lemma-criterion-flatness-fibre-nilpotent}
（基礎環の冪零イデアルの場合）と比較せよ。

\begin{lemma}[ファイバーによる平坦性判定法]
\label{algebra-lemma-criterion-flatness-fibre}
$R$、$S$、$S'$ を局所環とし、$R \to S \to S'$ を局所環準同型とする。
$M$ を $S'$-加群とする。$\mathfrak m \subset R$ を極大イデアルとする。
次を仮定する。
\begin{enumerate}
\item 環準同型 $R \to S$ と $R \to S'$ は本質的有限表示である。
\item 加群 $M$ は $S'$ 上有限表示である。
\item 加群 $M$ は零でない。
\item 加群 $M/\mathfrak mM$ は平坦 $S/\mathfrak mS$-加群である。
\item 加群 $M$ は平坦 $R$-加群である。
\end{enumerate}
このとき $S$ は $R$ 上平坦であり、$M$ は平坦 $S$-加群である。
\end{lemma}

\begin{proof}
補題 \ref{algebra-lemma-limit-essentially-finite-presentation} の証明と同様に、
まず $R = \colim R_\lambda$ を、本質的有限型な局所
$\mathbf{Z}$-代数の有向余極限として書ける。
% [P01-E0651] The frozen designation of the maximal ideal omits “by”.
$\mathfrak p_\lambda$ を $R_\lambda$ の極大イデアルとする。
次に、ある $\lambda_1 \in \Lambda$ に対して
$f_{j, \lambda_1} \in R_{\lambda_1}[x_1, \ldots, x_n]$ が存在して
$$
S =
\colim_{\lambda \geq \lambda_1} S_\lambda, \text{ with }
S_\lambda =
(R_\lambda[x_1, \ldots, x_n]/
(f_{1, \lambda}, \ldots, f_{u, \lambda}))_{\mathfrak q_\lambda}
$$
となると仮定してよい。ある $\lambda_2 \in \Lambda$、
$\lambda_2 \geq \lambda_1$ に対して
$g_{j, \lambda_2} \in R_{\lambda_2}[x_1, \ldots, x_n, y_1, \ldots, y_m]$
で、その像
$\overline{g}_{j, \lambda_2} \in S_{\lambda_2}[y_1, \ldots, y_m]$ が
$$
S' =
\colim_{\lambda \geq \lambda_2} S'_\lambda, \text{ with }
S'_\lambda =
(S_\lambda[y_1, \ldots, y_m]/
(\overline{g}_{1, \lambda}, \ldots,
\overline{g}_{v, \lambda}))_{\overline{\mathfrak q}'_\lambda}
$$
を満たすものが存在する。これはさらに次を含意することに注意する。
% [P01-E0652] The frozen displayed presentation omits the f-relation family.
$$
S'_\lambda =
(R_\lambda[x_1, \ldots, x_n, y_1, \ldots, y_m]/
(g_{1, \lambda}, \ldots, g_{v, \lambda}))_{\mathfrak q'_\lambda}
$$
$S'$ 上の $M$ の表示
$$
(S')^{\oplus s} \to (S')^{\oplus t} \to M \to 0
$$
を選ぶ。$A \in \text{Mat}(t \times s, S')$ をこの表示の行列とする。
ある $\lambda_3 \in \Lambda$、$\lambda_3 \geq \lambda_2$ に対して、
% [P01-E0653] The frozen approximating matrix is placed over S_lambda rather than S'_lambda.
$A$ に写る行列 $A_{\lambda_3} \in \text{Mat}(t \times s, S_{\lambda_3})$
を選べる。すべての $\lambda \geq \lambda_3$ に対して
$M_\lambda = \Coker((S'_\lambda)^{\oplus s} \xrightarrow{A_\lambda}
(S'_\lambda)^{\oplus t})$ とおく。

\medskip\noindent
これらの選択のもと、各 $\lambda_3 \leq \lambda \leq \mu$ に対して
$S_\lambda \otimes_{R_{\lambda}} R_\mu \to S_\mu$ は局所化、
$S'_\lambda \otimes_{S_{\lambda}} S_\mu \to S'_\mu$ は局所化であり、写像
$M_\lambda \otimes_{S'_\lambda} S'_\mu \to M_\mu$ は同型である。
これはさらに
$S'_\lambda \otimes_{R_{\lambda}} R_\mu \to S'_\mu$ が局所化で
あることを含意する。従って $M$ は $R$ 上平坦なので、補題
\ref{algebra-lemma-colimit-eventually-flat} により、十分大きいすべての
$\lambda$ に対して加群 $M_\lambda$ は $R_\lambda$ 上平坦である。
さらに、次に注意する。
$
\mathfrak m = \colim \mathfrak p_\lambda
$,
$
S/\mathfrak mS = \colim S_\lambda/\mathfrak p_\lambda S_\lambda
$,
$
S'/\mathfrak mS' = \colim S'_\lambda/\mathfrak p_\lambda S'_\lambda
$,
および
$
M/\mathfrak mM = \colim M_\lambda/\mathfrak p_\lambda M_\lambda
$。また、各 $\lambda_3 \leq \lambda \leq \mu$ に対して、上に列挙した
性質から
$$
S'_\lambda/\mathfrak p_\lambda S'_\lambda
\otimes_{S_{\lambda}/\mathfrak p_\lambda S_\lambda}
S_\mu/\mathfrak p_\mu S_\mu
\longrightarrow
S'_\mu/\mathfrak p_\mu S'_\mu
$$
は局所化であり、写像
$$
M_\lambda / \mathfrak p_\lambda M_\lambda
\otimes_{S'_\lambda/\mathfrak p_\lambda S'_\lambda}
S'_\mu /\mathfrak p_\mu S'_\mu
\longrightarrow
M_\mu/\mathfrak p_\mu M_\mu
$$
は同型であることが分かる。従って系
$(S_\lambda/\mathfrak p_\lambda S_\lambda \to
S'_\lambda/\mathfrak p_\lambda S'_\lambda,
M_\lambda/\mathfrak p_\lambda M_\lambda)$
も補題 \ref{algebra-lemma-limit-module-essentially-finite-presentation} のような系である。
$M/\mathfrak m M$ は $S/\mathfrak mS$ 上平坦と仮定されているので、
補題 \ref{algebra-lemma-colimit-eventually-flat} を再び適用でき、十分大きいすべての
$\lambda$ に対して $M_\lambda/\mathfrak p_\lambda M_\lambda$ は
$S_\lambda/\mathfrak p_\lambda S_\lambda$ 上平坦であることが分かる。
従って十分大きい $\lambda$ に対し、データ
$R_\lambda \to S_\lambda \to S'_\lambda, M_\lambda$ は補題
\ref{algebra-lemma-criterion-flatness-fibre-Noetherian} の仮定を満たす。
そのような $\lambda$ を選ぶ。このとき、
% [P01-E0654] The frozen proof identifies S with an unlocalized tensor product.
$S = S_\lambda \otimes_{R_\lambda} R$ は $R$ 上平坦であり、
$M = M_\lambda \otimes_{S_\lambda} S$ は $S$ 上平坦である
（平坦加群の基底変換は平坦だからである）。
\end{proof}

\noindent
次は「ファイバーによる平坦性判定法」補題
\ref{algebra-lemma-criterion-flatness-fibre} の容易な帰結である。
この種のさらなる結果については More on Flatness、節
\ref{flat-section-introduction} を参照せよ。

\begin{lemma}
\label{algebra-lemma-criterion-flatness-fibre-fp-over-ft}
$R$、$S$、$S'$ を局所環とし、$R \to S \to S'$ を局所環準同型とする。
$M$ を $S'$-加群とする。$\mathfrak m \subset R$ を極大イデアルとする。
次を仮定する。
\begin{enumerate}
\item $R \to S'$ は本質的有限表示である。
\item $R \to S$ は本質的有限型である。
\item $M$ は $S'$ 上有限表示である。
\item $M$ は零でない。
\item $M/\mathfrak mM$ は平坦 $S/\mathfrak mS$-加群である。
\item $M$ は平坦 $R$-加群である。
\end{enumerate}
このとき $S$ は $R$ 上本質的有限表示かつ平坦であり、
$M$ は平坦 $S$-加群である。
\end{lemma}

\begin{proof}
% [P01-E0655] The frozen proof assumes finite presentation where the hypothesis has finite type.
$S$ は $R$ 上本質的有限表示なので、ある有限型 $R$-代数 $C$ に対して
$S = C_{\overline{\mathfrak q}}$ と書ける。
$C = R[x_1, \ldots, x_n]/I$ と書く。
% [P01-E0656] The frozen sentence mixes “denote” with “be”.
$\mathfrak q \subset R[x_1, \ldots, x_n]$ を
$\overline{\mathfrak q}$ に対応する素イデアルとする。
すると $S = B/J$ である。ここで
$B = R[x_1, \ldots, x_n]_{\mathfrak q}$ は $R$ 上本質的有限表示で、
$J = IB$ である。$f_1, \ldots, f_k \in J$ を、それらの像
$\overline{f}_i \in B/\mathfrak mB$ がネーター環
$B/\mathfrak mB$ における $J$ の像 $\overline{J}$ を生成するように選べる。
従って、
% [P01-E0657] The frozen plural existential is paired with one ideal symbol.
$J' \subset J$ を満たす有限生成イデアルであって、$B/J' \to B/J$ が同型
$$
(B/J') \otimes_R R/\mathfrak m \longrightarrow
B/J \otimes_R R/\mathfrak m = S/\mathfrak mS.
$$
を誘導するものが存在する。上のような任意の $J'$ に対して、補題
\ref{algebra-lemma-criterion-flatness-fibre} を環準同型
$$
R \longrightarrow B/J' \longrightarrow S'
$$
および加群 $M$ に適用できる。従って、上のような任意の選択 $J'$ に対して
$B/J'$ は $R$ 上平坦である。ここで、上のような二つの有限生成イデアル
% [P01-E0658] The frozen inclusion chain repeats J' instead of introducing J''.
$J' \subset J' \subset J$ があるとすると、
$B/J' \to B/J''$ は、本質的有限表示で平坦な $R$-代数の間の全射であり、
% [P01-E0659] Japanese separates the two frozen conflicting relative clauses.
さらに $\mathfrak m$ を法として同型であると結論できる。
従って補題 \ref{algebra-lemma-mod-injective-general} により
$B/J' = B/J''$、すなわち $J' = J''$ である。
これは明らかに $J$ が有限生成、すなわち $S$ が $R$ 上本質的有限表示で
あることを意味する。従って補題 \ref{algebra-lemma-criterion-flatness-fibre} を
$R \to S \to S'$ に適用でき、証明が終わる。
\end{proof}

\begin{lemma}[ファイバーによる平坦性判定法：局所冪零の場合]
\label{algebra-lemma-criterion-flatness-fibre-locally-nilpotent}
可換図式
$$
\xymatrix{
S \ar[rr] & & S' \\
& R \ar[lu] \ar[ru]
}
$$
を環の圏で考える。$I \subset R$ を局所冪零イデアルとし、
$M$ を $S'$-加群とする。次を仮定する。
\begin{enumerate}
\item $R \to S$ は有限型である。
\item $R \to S'$ は有限表示である。
\item $M$ は有限表示 $S'$-加群である。
\item $M/IM$ は $S/IS$-加群として平坦である。
\item $M$ は $R$-加群として平坦である。
\end{enumerate}
このとき $M$ は平坦 $S$-加群であり、
$M \otimes_S \kappa(\mathfrak q)$ が零でないようなすべての
$\mathfrak q \subset S$ に対して、$S_\mathfrak q$ は $R$ 上平坦かつ
本質的有限表示である。
\end{lemma}

\begin{proof}
$M \otimes_S \kappa(\mathfrak q)$ が零でなければ、
$S' \otimes_S \kappa(\mathfrak q)$ も零でない。従って補題
\ref{algebra-lemma-in-image} により、$\mathfrak q$ の上にある素イデアル
$\mathfrak q' \subset S'$ が存在する。
$\mathfrak p \subset R$ を $\Spec(R)$ における $\mathfrak q$ の像とする。
$I$ は局所冪零なので $I \subset \mathfrak p$ であり、従って
$M/\mathfrak p M$ は $S/\mathfrak pS$ 上平坦である。
従って補題 \ref{algebra-lemma-criterion-flatness-fibre-fp-over-ft} を
$R_\mathfrak p \to S_\mathfrak q \to S'_{\mathfrak q'}$ と
$M_{\mathfrak q'}$ に適用できる。$M_{\mathfrak q'}$ は $S$ 上平坦であり、
$S_\mathfrak q$ は $R$ 上平坦かつ本質的有限表示であると結論する。
$\mathfrak q'$ は $S'$ の任意の素イデアルだったので、補題
\ref{algebra-lemma-flat-localization} により $M$ も $S$ 上平坦であることが分かる。
\end{proof}


\section{平坦軌跡の開性}
\label{algebra-section-open-flat}

\noindent
補題 \ref{algebra-lemma-colimit-eventually-flat} を用いてネーターの場合に帰着する。
ネーターの場合は、節 \ref{algebra-section-complex-exact} で与えた完全複体の
特徴づけを用いて扱う。

\begin{lemma}
\label{algebra-lemma-CM-dim-finite-type}
$k$ を体とし、$S$ を有限型 $k$-代数とする。
$f_1, \ldots, f_i$ を $S$ の元とする。
$S$ が次元 $d$ の Cohen--Macaulay かつ等次元であり、
$\dim V(f_1, \ldots, f_i) \leq d - i$ であると仮定する。
このとき等号が成り立ち、
% [P01-E0660] The frozen compound subject takes a singular verb.
$f_1, \ldots, f_i$ は、
% [P01-E0661] Japanese naturally renders the intended membership relation.
$V(f_1, \ldots, f_i)$ のすべての素イデアル $\mathfrak q$ に対して
$S_{\mathfrak q}$ の正則列をなす。
\end{lemma}

\begin{proof}
$S$ が次元 $d$ の Cohen--Macaulay かつ等次元ならば、$S$ のすべての
極大イデアルに対して $\dim(S_{\mathfrak m}) = d$ である。
補題 \ref{algebra-lemma-disjoint-decomposition-CM-algebra} を参照せよ。
命題 \ref{algebra-proposition-CM-module} により、
$\mathfrak m \in V(f_1, \ldots, f_i)$ を満たすすべての極大イデアル
$\mathfrak m$ に対して、この列は $S_{\mathfrak m}$ の正則列であり、
局所環 $S_{\mathfrak m}/(f_1, \ldots, f_i)$ は次元 $d - i$ の
Cohen--Macaulay 環である。このことは実際、
$S/(f_1, \ldots, f_i)$ が次元 $d - i$ の Cohen--Macaulay かつ
等次元な環であることを意味する。
\end{proof}

\begin{lemma}
\label{algebra-lemma-open-regular-sequence}
$R \to S$ を有限型環準同型とする。すべてのファイバー
$S \otimes_R \kappa(\mathfrak p)$ が次元 $d$ の Cohen--Macaulay かつ
等次元な環となる整数 $d$ をとる。$f_1, \ldots, f_i$ を $S$ の元とする。
集合
$$
\{ \mathfrak q \in V(f_1, \ldots, f_i)
\mid f_1, \ldots, f_i
\text{ are a regular sequence in }
S_{\mathfrak q}/\mathfrak p S_{\mathfrak q}
\text{ where }\mathfrak p = R \cap \mathfrak q
\}
$$
は $V(f_1, \ldots, f_i)$ の開集合である。
\end{lemma}

\begin{proof}
$\overline{S} = S/(f_1, \ldots, f_i)$ と書く。
$\mathfrak q$ を補題で定義した集合の元とし、$\mathfrak p$ を対応する
$R$ の素イデアルとする。定義 \ref{algebra-definition-relative-dimension} で
定義した相対次元を用いる。まず、補題
\ref{algebra-lemma-dimension-at-a-point-finite-type-field} により
$d = \dim_{\mathfrak q}(S/R) =
\dim(S_{\mathfrak q}/\mathfrak pS_{\mathfrak q}) +
\text{trdeg}_{\kappa(\mathfrak p)}\ \kappa(\mathfrak q)$ であることに注意する。
$f_1, \ldots, f_i$ はネーター局所環
$S_{\mathfrak q}/\mathfrak pS_{\mathfrak q}$ の正則列をなすので、補題
\ref{algebra-lemma-one-equation} から
$\dim(\overline{S}_{\mathfrak q}/\mathfrak p\overline{S}_{\mathfrak q})
= \dim(S_{\mathfrak q}/\mathfrak pS_{\mathfrak q}) - i$ を得る。
従って補題 \ref{algebra-lemma-dimension-at-a-point-finite-type-field} により
$\dim_{\mathfrak q}(\overline{S}/R)
= \dim(\overline{S}_{\mathfrak q}/\mathfrak p\overline{S}_{\mathfrak q}) +
\text{trdeg}_{\kappa(\mathfrak p)}\ \kappa(\mathfrak q)
= d - i$ である。補題 \ref{algebra-lemma-dimension-fibres-bounded-open-upstairs}
により、$\mathfrak q$ のある近傍内のすべての
$\mathfrak q' \in V(f_1, \ldots, f_i) = \Spec(\overline{S})$ に対して
$\dim_{\mathfrak q'}(\overline{S}/R) \leq d - i$ である。
従って、ある $g \in S$、$g \not \in \mathfrak q$ により $S$ を $S_g$ で
置き換えた後、この不等式がすべての $\mathfrak q'$ に対して成り立つと
仮定してよい。補題 \ref{algebra-lemma-CM-dim-finite-type} から結論が従う。
\end{proof}

\begin{lemma}
\label{algebra-lemma-exact-on-fibres-open}
$R \to S$ を環準同型とする。有限自由 $S$-加群からなる有限ホモロジー複体
$$
F_{\bullet} :
0
\to
S^{n_e}
\xrightarrow{\varphi_e}
S^{n_{e-1}}
\xrightarrow{\varphi_{e-1}}
\ldots
\xrightarrow{\varphi_{i + 1}}
S^{n_i}
\xrightarrow{\varphi_i}
S^{n_{i-1}}
\xrightarrow{\varphi_{i-1}}
\ldots
\xrightarrow{\varphi_1}
S^{n_0}
$$
を考える。$S$ の各素イデアル $\mathfrak q$ に対し、
$\mathfrak p$ を $R$ における $\mathfrak q$ の逆像として、複体
$\overline{F}_{\bullet, \mathfrak q} =
F_{\bullet, \mathfrak q} \otimes_R \kappa(\mathfrak p)$ を考える。
% [P01-E0662] Japanese supplies the missing coordinated finite predicates.
$R$ はネーターであり、$R \to S$ は有限型かつ平坦で、そのファイバー
$S \otimes_R \kappa(\mathfrak p)$ が次元 $d$ の Cohen--Macaulay 環となる
整数 $d$ が存在すると仮定する。集合
$$
\{\mathfrak q \in \Spec(S) \mid
\overline{F}_{\bullet, \mathfrak q}\text{ is exact}\}
$$
は $\Spec(S)$ の開集合である。
\end{lemma}

\begin{proof}
$\mathfrak q$ を補題で定義した集合の元とする。命題
\ref{algebra-proposition-what-exact} を用いて、$D(g)$ が補題で定義した集合に
含まれるような $g \in S$、$g \not \in \mathfrak q$ が存在することを示す。
言い換えれば、$S$ を $S_g$ で置き換えた後、補題の集合が
$\Spec(S)$ 全体になることを示す。そのため、証明中に有限回、$S$ を
このような局所化で置き換える。

命題 \ref{algebra-proposition-what-exact} は、環が局所環の場合、写像
$\varphi_i$ の階数とイデアル $I(\varphi_i)$ によって複体の完全性を
特徴づけることを思い出そう。まず階数条件を扱う。
$r_i = n_i - n_{i + 1} + \ldots + (-1)^{e - i} n_e$ とおく。
$r_i + r_{i + 1} = n_i$ であり、$r_i$ は完全な場合に期待される
$\varphi_i$ の階数であることに注意する。

\medskip\noindent
補題 \ref{algebra-lemma-complex-exact-mod} により、
$\overline{F}_{\bullet, \mathfrak q}$ が完全ならば、局所化
$F_{\bullet, \mathfrak q}$ は完全である。特に、複体 $F_\bullet$ は
ある $g \in S$、$g \not \in \mathfrak q$ による局所化の後に完全になる。
この場合、$S$ のすべての素イデアルにおける局所化へ命題
\ref{algebra-proposition-what-exact} を適用すると、$\varphi_i$ のすべての
$(r_i + 1) \times (r_i + 1)$-小行列式は零である。従って
$\varphi_i$ の階数は高々 $r_i$ である。

\medskip\noindent
$I_i \subset S$ を、$\varphi_i$ の行列の $r_i \times r_i$-小行列式で
生成されるイデアルとする。命題 \ref{algebra-proposition-what-exact} により、複体
$\overline{F}_{\bullet, \mathfrak q}$ が完全であるための必要十分条件は、
すべての $1 \leq i \leq e$ に対して、
$(I_i)_{\mathfrak q} = S_{\mathfrak q}$ であるか、または
$(I_i)_{\mathfrak q}$ が長さ $i$ の
$S_{\mathfrak q}/\mathfrak p S_{\mathfrak q}$-正則列を含むことである。
実際、上での $r_i$ の選び方と $\varphi_i$ の階数に対する上界により、
命題 \ref{algebra-proposition-what-exact} の条件を満たす方法はこれしかない。

\medskip\noindent
$(I_i)_{\mathfrak q} = S_{\mathfrak q}$ ならば、ある
$g \not\in \mathfrak q$ で $S$ を局所化した後、$I_i = S$ と仮定してよい。
これは明らかに開条件である。

\medskip\noindent
$(I_i)_{\mathfrak q} \not = S_{\mathfrak q}$ ならば、
% [P01-E0663] The frozen proof uses localized elements globally without clearing denominators.
$S_{\mathfrak q}/\mathfrak pS_{\mathfrak q}$ の正則列をなす
$f_1, \ldots, f_i \in (I_i)_{\mathfrak q}$ が存在する。
$(f_1, \ldots, f_i) \not \subset \mathfrak q'$ を満たす任意の素イデアル
$\mathfrak q' \subset S$ に対して $(I_i)_{\mathfrak q'} = S_{\mathfrak q'}$
であることに注意する。従って補題 \ref{algebra-lemma-open-regular-sequence} から
結論が従う。
\end{proof}


\begin{theorem}
\label{algebra-theorem-openness-flatness}
$R$ を環とし、$R \to S$ を有限表示環準同型とする。
$M$ を有限表示 $S$-加群とする。集合
$$
\{ \mathfrak q \in \Spec(S) \mid
M_{\mathfrak q}\text{ is flat over }R\}
$$
は $\Spec(S)$ の開集合である。
\end{theorem}

\begin{proof}
$\mathfrak q \in \Spec(S)$ を素イデアルとする。
$\mathfrak p \subset R$ を $R$ における $\mathfrak q$ の逆像とする。
$M_{\mathfrak q}$ が $R$ 上平坦であることと $R_{\mathfrak p}$ 上平坦で
あることは同値である。$M_{\mathfrak q}$ が $R$ 上平坦であると仮定する。
ある $g \in S$、$g \not \in \mathfrak q$ が存在して、$M_g$ が $R$ 上
平坦になると主張する。

\medskip\noindent
まず、$R$ と $S$ が $\mathbf{Z}$ 上有限型である場合に帰着する。
補題 \ref{algebra-lemma-limit-module-finite-presentation} のような有向集合
$\Lambda$ と系 $(R_\lambda \to S_\lambda, M_\lambda)$ を選ぶ。
$\mathfrak p_\lambda$ を $R_\lambda$ における $\mathfrak p$ の逆像とし、
$\mathfrak q_\lambda$ を $S_\lambda$ における $\mathfrak q$ の逆像とする。
すると系
$$
((R_\lambda)_{\mathfrak p_\lambda},
(S_\lambda)_{\mathfrak q_\lambda},
(M_\lambda)_{\mathfrak q_{\lambda}})
$$
は補題 \ref{algebra-lemma-limit-module-essentially-finite-presentation} のような系である。
従って補題 \ref{algebra-lemma-colimit-eventually-flat} により、ある $\lambda$ に対して
加群 $M_\lambda$ は素イデアル $\mathfrak q_{\lambda}$ において
$R_\lambda$ 上平坦である。系
$(R_\lambda \to S_\lambda, M_\lambda, \mathfrak q_{\lambda})$ に対して
主張を証明できると仮定する。すなわち、ある $g \in S_\lambda$、
$g \not\in \mathfrak q_\lambda$ が存在して $(M_\lambda)_g$ が
$R_\lambda$ 上平坦であると仮定する。補題
\ref{algebra-lemma-limit-module-finite-presentation} により
$M = M_\lambda \otimes_{R_\lambda} R$ であり、従って
$M_g = (M_\lambda)_g \otimes_{R_\lambda} R$ でもある。
ゆえに補題 \ref{algebra-lemma-flat-base-change} により、系
$(R \to S, M, \mathfrak q)$ に対する主張が従う。

\medskip\noindent
ここで $R$ と $S$ は $\mathbf{Z}$ 上有限型であると仮定してよい。
$S$ を多項式環 $R[x_1, \ldots, x_n]$ の商として書ける。
もちろん $S$ を $R[x_1, \ldots, x_n]$ で置き換え、$S$ が $R$ 上の
多項式環であると仮定してよい。特に $R \to S$ は平坦であり、
% [P01-E0664] The frozen English uses the plural noun attributively.
すべてのファイバー環 $S \otimes_R \kappa(\mathfrak p)$ の大域次元は
$n$ である。

\medskip\noindent
各 $F_i$ が有限自由となる $S$ 上の $M$ の分解 $F_\bullet$ を選ぶ。
補題 \ref{algebra-lemma-resolution-by-finite-free} を参照せよ。
$K_n = \Ker(F_{n-1} \to F_{n-2})$ とおく。各 $F_i$ は $R$ 上平坦であり、
$M$ に関する仮定もあるので、補題 \ref{algebra-lemma-flat-ses} により
$(K_n)_{\mathfrak q}$ は $R$ 上平坦である。さらに列
$$
0 \to
K_n/\mathfrak p K_n \to
F_{n-1}/ \mathfrak p F_{n-1} \to
\ldots \to
F_0 / \mathfrak p F_0 \to
M/\mathfrak p M \to
0
$$
は、$\text{Tor}_i^{R_\mathfrak p}(\kappa(\mathfrak p), M_{\mathfrak q})$ が
消えるため、$\mathfrak q$ で局所化すると完全である。
% [P01-E0665] The frozen proof claims exact global dimension n rather than at most n.
$S_\mathfrak q/\mathfrak p S_{\mathfrak q}$ の大域次元は $n$ なので、
$K_n / \mathfrak p K_n$ を $\mathfrak q$ で局所化したものは
$S_\mathfrak q/\mathfrak p S_{\mathfrak q}$ 上有限自由であると結論する。
補題 \ref{algebra-lemma-free-fibre-flat-free} により $(K_n)_{\mathfrak q}$ は
$S_{\mathfrak q}$ 上自由である。特に、ある $g \in S$、
$g \not \in \mathfrak q$ が存在して、$(K_n)_g$ は $S_g$ 上有限自由である。

\medskip\noindent
補題 \ref{algebra-lemma-exact-on-fibres-open} により、さらに $S_g$ を局所化して、
複体
$$
0 \to K_n \to F_{n-1} \to \ldots \to F_0
$$
が $R \to S$ の {\it すべてのファイバー} 上で完全になるようにできる。
補題 \ref{algebra-lemma-complex-exact-mod} により、これは $F_1 \to F_0$ の余核が
平坦であることを含意する。これでネーターの場合の定理が証明された。
\end{proof}

\section{Cohen--Macaulay 軌跡の開性}
\label{algebra-section-CM-open}

\noindent
この節では、有限型代数の Cohen--Macaulay 性を平坦性によって特徴づける。
次いで、これを用いて、そのような代数が Cohen--Macaulay となる点の集合が
開であることを証明する。
% [P01-E0666] Japanese makes the frozen nested object clause explicit.

\begin{lemma}
\label{algebra-lemma-where-CM}
$S$ を体 $k$ 上の有限型代数とする。
$\varphi : k[y_1, \ldots, y_d] \to S$ を準有限環準同型とする。
$\Spec(S)$ の部分集合として次の等式が成り立つ。
$$
\{ \mathfrak q \mid
S_{\mathfrak q} \text{ flat over }k[y_1, \ldots, y_d]\}
=
\{ \mathfrak q \mid
S_{\mathfrak q} \text{ CM and }\dim_{\mathfrak q}(S/k) = d\}
$$
% [P01-E0667, P01-E0668] The two frozen predicates inside the display omit “is”.
記法については定義 \ref{algebra-definition-relative-dimension} を参照せよ。
\end{lemma}

\begin{proof}
$\mathfrak q \subset S$ を素イデアルとする。
% [P01-E0669] Japanese supplies a complete designation construction.
$\mathfrak p = k[y_1, \ldots, y_d] \cap \mathfrak q$ とおく。
例えば補題 \ref{algebra-lemma-dimension-inequality-quasi-finite} により、常に
$\dim(S_{\mathfrak q}) \leq \dim(k[y_1, \ldots, y_d]_{\mathfrak p})$
であることに注意する。また、体の拡大
$\kappa(\mathfrak q)/\kappa(\mathfrak p)$ は有限であり、従って
$\text{trdeg}_k(\kappa(\mathfrak p)) = \text{trdeg}_k(\kappa(\mathfrak q))$
である。

\medskip\noindent
$\mathfrak q$ を左辺の元とする。すると補題
\ref{algebra-lemma-finite-flat-over-regular-CM} を適用でき、
$S_{\mathfrak q}$ は Cohen--Macaulay であり、
$\dim(S_{\mathfrak q}) = \dim(k[y_1, \ldots, y_d]_{\mathfrak p})$
であると結論する。これを上の超越次数の等式および補題
\ref{algebra-lemma-dimension-at-a-point-finite-type-field} と合わせると、
$\dim_{\mathfrak q}(S/k) = d$ が従う。従って $\mathfrak q$ は右辺の元である。

\medskip\noindent
$\mathfrak q$ を右辺の元とする。上の超越次数の等式、仮定
$\dim_{\mathfrak q}(S/k) = d$、および補題
\ref{algebra-lemma-dimension-at-a-point-finite-type-field} により、
$\dim(S_{\mathfrak q}) = \dim(k[y_1, \ldots, y_d]_{\mathfrak p})$
と結論する。従って補題 \ref{algebra-lemma-CM-over-regular-flat} を適用でき、
$\mathfrak q$ は左辺の元であることが分かる。
\end{proof}

\begin{lemma}
\label{algebra-lemma-finite-type-over-field-CM-open}
$S$ を体 $k$ 上の有限型代数とする。$S_{\mathfrak q}$ が
Cohen--Macaulay となる素イデアル $\mathfrak q$ の集合は $S$ において開である。
% [P01-E0670] The frozen statement names the ring S rather than Spec(S).
\end{lemma}

\noindent
この補題は下の補題 \ref{algebra-lemma-finite-presentation-flat-CM-locus-open} の
特別な場合であるから、望むなら同補題の証明まで読み飛ばしてよい。

\begin{proof}
$\mathfrak q \subset S$ を、$S_{\mathfrak q}$ が Cohen--Macaulay となる
素イデアルとする。ある $g \in S$、$g \not \in \mathfrak q$ が存在して
環 $S_g$ が Cohen--Macaulay になることを示さねばならない。
任意の $g \in S$、$g \not \in \mathfrak q$ に対して、$S$ を $S_g$ で、
$\mathfrak q$ を $\mathfrak qS_g$ で置き換えてよい。これと補題
\ref{algebra-lemma-Noether-normalization-at-point} および
\ref{algebra-lemma-dimension-at-a-point-finite-type-field} を合わせると、
$d = \dim(S_{\mathfrak q}) + \text{trdeg}_k(\kappa(\mathfrak q))$
を満たす有限単射環準同型 $k[y_1, \ldots, y_d] \to S$ が存在すると
仮定してよい。$\mathfrak p = k[y_1, \ldots, y_d] \cap \mathfrak q$ とおく。
構成により、$\mathfrak q$ は補題 \ref{algebra-lemma-where-CM} の表示等式の
右辺の元であり、従って左辺の元でもある。

\medskip\noindent
定理 \ref{algebra-theorem-openness-flatness} により、ある $g \in S$、
$g \not \in \mathfrak q$ に対して環 $S_g$ は
$k[y_1, \ldots, y_d]$ 上平坦である。従って補題
\ref{algebra-lemma-where-CM} の等式を再び用いると、$S_g$ のすべての局所環は
望みどおり Cohen--Macaulay である。
\end{proof}

\begin{lemma}
\label{algebra-lemma-generic-CM}
$k$ を体とし、$S$ を有限型 $k$ 代数とする。
% [P01-E0671] The frozen English compound lacks standard hyphenation.
Cohen--Macaulay な素イデアルの集合は稠密な開集合
$U \subset \Spec(S)$ をなす。
\end{lemma}

\begin{proof}
補題 \ref{algebra-lemma-finite-type-over-field-CM-open} により、この集合は開である。
また、すべての極小素イデアル $\mathfrak q \subset S$ を含む。
実際、極小素イデアルにおける局所環 $S_{\mathfrak q}$ は次元零であり、
従って Cohen--Macaulay である。
\end{proof}

\begin{lemma}
\label{algebra-lemma-finite-presentation-flat-CM-locus-open}
$R$ を環とする。$R \to S$ は有限表示かつ平坦であるとする。
任意の $d \geq 0$ に対して集合
$$
\left\{
\begin{matrix}
\mathfrak q \in \Spec(S)
\text{ such that setting }\mathfrak p = R \cap \mathfrak q
\text{ the fibre ring}\\
S_{\mathfrak q}/\mathfrak pS_{\mathfrak q}
\text{ is Cohen-Macaulay}
\text{ and } \dim_{\mathfrak q}(S/R) = d
\end{matrix}
\right\}
$$
は $\Spec(S)$ の開集合である。
\end{lemma}

\begin{proof}
$\mathfrak q$ を表示した集合の元とし、$\mathfrak p$ を対応する
$R$ の素イデアルとする。ある $g \in S$、$g \not \in \mathfrak q$ が
存在して、$R \to S_g$ のすべてのファイバー環が Cohen--Macaulay に
なることを示さねばならない。
% [P01-E0672] The frozen local goal omits the fixed relative-dimension-d condition.
証明中、有限回に限り、ある $g \in S$、$g \not \in \mathfrak q$ に対して
$S$ を $S_g$ で置き換えてよい。従って補題
\ref{algebra-lemma-quasi-finite-over-polynomial-algebra} により、
$d = \dim_{\mathfrak q}(S/R)$ を満たす準有限環準同型
$R[t_1, \ldots, t_d] \to S$ が存在すると仮定してよい。
$\mathfrak q' = R[t_1, \ldots, t_d] \cap \mathfrak q$ とおく。
補題 \ref{algebra-lemma-where-CM} により、環準同型
$$
R[t_1, \ldots, t_d]_{\mathfrak q'} /
\mathfrak p R[t_1, \ldots, t_d]_{\mathfrak q'}
\longrightarrow
S_{\mathfrak q}/\mathfrak p S_{\mathfrak q}
$$
は平坦である。従ってファイバーによる平坦性判定法である補題
\ref{algebra-lemma-criterion-flatness-fibre} により、
$R[t_1, \ldots, t_d]_{\mathfrak q'} \to S_{\mathfrak q}$ は平坦である。
従って定理 \ref{algebra-theorem-openness-flatness} により、ある $g \in S$、
% [P01-E0673] Japanese punctuates the frozen interrupting condition naturally.
$g \not \in \mathfrak q$ に対して環準同型
$R[t_1, \ldots, t_d] \to S_g$ は平坦である。
$S$ を $S_g$ で置き換えると、すべての素イデアル
$\mathfrak r \subset S$ に対して、
$\mathfrak r' = R[t_1, \ldots, t_d] \cap \mathfrak r$ および
$\mathfrak p' = R \cap \mathfrak r$ とおけば、局所環準同型
$R[t_1, \ldots, t_d]_{\mathfrak r'} \to S_{\mathfrak r}$ は平坦である。
従ってその基底変換
$$
R[t_1, \ldots, t_d]_{\mathfrak r'} /
\mathfrak p' R[t_1, \ldots, t_d]_{\mathfrak r'}
\longrightarrow
S_{\mathfrak r}/\mathfrak p' S_{\mathfrak r}
$$
も平坦である。従って $k = \kappa(\mathfrak p')$ として補題
\ref{algebra-lemma-where-CM} を適用すると、望みどおり
$\mathfrak r$ は補題の集合に属することが分かる。
\end{proof}

\begin{lemma}
\label{algebra-lemma-generic-CM-flat-finite-presentation}
$R$ を環とし、$R \to S$ を有限表示な平坦環準同型とする。
$\mathfrak p = R \cap \mathfrak q$ としたときにファイバー環
$S_{\mathfrak q} \otimes_R \kappa(\mathfrak p)$ が Cohen--Macaulay となる
素イデアル $\mathfrak q$ の集合は、$\Spec(S) \to \Spec(R)$ の
各ファイバー内で開かつ稠密である。
\end{lemma}

\begin{proof}
この集合を $W$ と呼ぶ。補題
\ref{algebra-lemma-finite-presentation-flat-CM-locus-open} により開である。
また、$W$ と一つのファイバーとの共通部分は、そのファイバーに対して補題
\ref{algebra-lemma-generic-CM} を適用する集合そのものなので、各ファイバー内で稠密である。
\end{proof}

\begin{lemma}
\label{algebra-lemma-extend-field-CM-locus}
$k$ を体とし、$S$ を有限型 $k$-代数とする。
$K/k$ を体の拡大とし、$S_K = K \otimes_k S$ とおく。
$\mathfrak q \subset S$ を $S$ の素イデアルとし、
$\mathfrak q_K \subset S_K$ を $\mathfrak q$ の上にある
$S_K$ の素イデアルとする。このとき $S_{\mathfrak q}$ が
Cohen--Macaulay であることと $(S_K)_{\mathfrak q_K}$ が
Cohen--Macaulay であることは同値である。
\end{lemma}

\begin{proof}
証明中、有限回に限り、任意の $g \in S$、$g \not \in \mathfrak q$ に対して
$S$ を $S_g$ で置き換えてよい。従って補題
\ref{algebra-lemma-Noether-normalization-at-point} を用いると、
$\dim(S) = \dim_{\mathfrak q}(S/k) =: d$ と仮定し、有限単射準同型
$k[x_1, \ldots, x_d] \to S$ を選んでよい。
これは基底変換により有限単射準同型
$K[x_1, \ldots, x_d] \to S_K$ も誘導することに注意する。
補題 \ref{algebra-lemma-dimension-at-a-point-preserved-field-extension} により
$\dim_{\mathfrak q_K}(S_K/K) = d$ である。
$\mathfrak p = k[x_1, \ldots, x_d] \cap \mathfrak q$ および
$\mathfrak p_K = K[x_1, \ldots, x_d] \cap \mathfrak q_K$ とおく。
ネーター局所環の次の可換図式を考える。
$$
\xymatrix{
S_{\mathfrak q} \ar[r] &
(S_K)_{\mathfrak q_K} \\
k[x_1, \ldots, x_d]_{\mathfrak p} \ar[r] \ar[u] &
K[x_1, \ldots, x_d]_{\mathfrak p_K} \ar[u]
}
$$
補題 \ref{algebra-lemma-where-CM} により、左の垂直矢印が平坦であることと
右の垂直矢印が平坦であることが同値だと示せばよい。
下の矢印が平坦なので、この同値は補題
\ref{algebra-lemma-base-change-flat-up-down} により成り立つ。
\end{proof}

\begin{lemma}
\label{algebra-lemma-CM-locus-commutes-base-change}
$R$ を環とし、$R \to S$ を有限型環準同型とする。
$R \to R'$ を任意の環準同型とし、$S' = R' \otimes_R S$ とおく。
% [P01-E0674] Japanese supplies a complete designation construction.
$S \to S'$ に付随する写像を
$f : \Spec(S') \to \Spec(S)$ と書く。
$\mathfrak p = R \cap \mathfrak q$ としたときにファイバー環
$S_{\mathfrak q} \otimes_R \kappa(\mathfrak p)$ が Cohen--Macaulay となる
素イデアル $\mathfrak q$ の集合を $W$ とし、$S'/R'$ に対する同様の集合を
$W'$ とする。このとき $W' = f^{-1}(W)$ である。
\end{lemma}

\begin{proof}
% [P01-E0675] Japanese supplies a finite verb for the frozen fragment.
これは補題 \ref{algebra-lemma-extend-field-CM-locus} と定義から直ちに従う。
\end{proof}

\begin{lemma}
\label{algebra-lemma-relative-dimension-CM}
$R$ を環とする。$R \to S$ を、(a) 平坦、(b) 有限表示、
(c) Cohen--Macaulay なファイバーをもつ環準同型とする。
% [P01-E0676] Japanese presents the frozen three predicates in parallel.
このとき $S = S_0 \times \ldots \times S_n$ を $R$-代数 $S_d$ の積として
書くことができ、各 $S_d$ は (a)、(b)、(c) を満たし、そのすべての
ファイバーは次元 $d$ で等次元である。
\end{lemma}

\begin{proof}
% [P01-E0677] Japanese supplies a complete designation construction.
各整数 $d$ に対し、補題
\ref{algebra-lemma-finite-presentation-flat-CM-locus-open} で定義した集合を
$W_d \subset \Spec(S)$ と書く。明らかに
$\Spec(S) = \coprod W_d$ であり、引用した補題により各 $W_d$ は開である。
従って補題 \ref{algebra-lemma-disjoint-implies-product} から結論が従う。
\end{proof}
\section{微分}
\label{algebra-section-differentials}

\noindent
本節では環準同型の微分加群を定義する。

\begin{definition}
\label{algebra-definition-derivation}
$\varphi : R \to S$ を環準同型とし、$M$ を $S$-加群とする。
$M$ に値をとる {\it 導分}、より正確には
{\it $R$-導分}とは、加法的であり、$\varphi(R)$ の元を零に写し、
{\it Leibniz 則} $D(ab) = aD(b) + bD(a)$ を満たす写像
$D : S \to M$ のことである。
\end{definition}

\noindent
$r \in R$ および $a \in S$ ならば $D(ra) = rD(a)$ であることに注意する。
同値な定義として、$R$-導分とは Leibniz 則を満たす
$R$-線形写像 $D : S \to M$ のことである。
すべての $R$-導分の集合は $S$-加群をなす。実際、二つの
$R$-導分 $D, D'$ が与えられたとき、その和
$D + D' : S \to M$, $a \mapsto D(a)+D'(a)$ は $R$-導分であり、
$R$-導分 $D$ と元 $c\in S$ が与えられたとき、スカラー倍
$cD : S \to M$, $a \mapsto cD(a)$ も $R$-導分である。
この $S$-加群を
$$
\text{Der}_R(S, M).
$$
と書く。また、$\alpha : M \to N$ が $S$-加群準同型ならば、
合成 $\alpha \circ D$ は $N$ に値をとる $R$-導分である。
このようにして、対応 $M \mapsto \text{Der}_R(S, M)$ は共変関手となる。

\medskip\noindent
次の自由 $S$-加群の準同型を考える。
$$
\bigoplus\nolimits_{(a, b)\in S^2} S[(a, b)] \oplus
\bigoplus\nolimits_{(f, g)\in S^2} S[(f, g)] \oplus
\bigoplus\nolimits_{r\in R} S[r]
\longrightarrow
\bigoplus\nolimits_{a\in S} S[a]
$$
これは、自然な記法のもとで
$$
[(a, b)] \longmapsto [a + b] - [a] - [b],\quad
[(f, g)] \longmapsto [fg] -f[g] - g[f],\quad
[r]      \longmapsto [\varphi(r)]
$$
という規則により定義される。この準同型の余核を $\Omega_{S/R}$ とする。
$a$ を余核における $[a]$ の類 $\text{d}a$ に写す写像
$\text{d} : S \to \Omega_{S/R}$ がある。$\Omega_{S/R}$ に課した関係により、
これは $R$-導分である。すなわち、
$$
\text{d}(a + b) = \text{d}a + \text{d}b, \quad
\text{d}(fg) = f\text{d}g + g\text{d}f, \quad
\text{d}\varphi(r) = 0
$$
が $a,b,f,g \in S$ および $r \in R$ に対して成り立つ。

\begin{definition}
\label{algebra-definition-differentials}
対 $(\Omega_{S/R}, \text{d})$ を、$R$ 上の $S$ の
{\it K\"ahler 微分加群}、または単に {\it 微分加群}という。
\end{definition}

\begin{lemma}
\label{algebra-lemma-universal-omega}
\begin{slogan}
微分加群からの写像は導分と同じものである。
\end{slogan}
$R$ 上の $S$ の微分加群は次の普遍性をもつ。写像
$$
\Hom_S(\Omega_{S/R}, M)
\longrightarrow
\text{Der}_R(S, M), \quad
\alpha
\longmapsto
\alpha \circ \text{d}
$$
は関手の同型である。
\end{lemma}

\begin{proof}
定義により、$R$-導分とは、各 $a \in S$ に元 $D(a) \in M$ を
対応させる規則である。従って $D$ から写像
$[D] : \bigoplus S[a] \to M$ が得られる。一方、$R$-導分であるための
条件は、ちょうど $[D]$ が上で与えた $\Omega_{S/R}$ の表示に現れる
写像の像を零に写すことを意味する。
\end{proof}

\begin{lemma}
\label{algebra-lemma-trivial-differential-surjective}
$R \to S$ が全射であるとする。このとき $\Omega_{S/R} = 0$ である。
\end{lemma}

\begin{proof}
これは、すべての $R$-導分が明らかに零でなければならないことからも、
$\Omega_{S/R}$ の表示に現れる写像が全射であることからも分かる。
\end{proof}

\noindent
\begin{equation}
\label{algebra-equation-functorial-omega}
\vcenter{
\xymatrix{
S \ar[r]_\varphi
&
S'
\\
R \ar[r]^\psi \ar[u]^\alpha
&
R' \ar[u]_\beta
}
}
\end{equation}
が環の可換図式であるとする。このとき、次の可換図式に収まる
微分加群の自然な準同型が存在する。
$$
\xymatrix{
\Omega_{S/R} \ar[r] &
\Omega_{S'/R'}
\\
S \ar[u]^{\text{d}} \ar[r]^{\varphi}
&
S' \ar[u]_{\text{d}}
}
$$
この準同型を構成するには、$\Omega_{S/R}$ と $\Omega_{S'/R'}$ の
表示の間の自然な写像を用いればよい。すなわち、
\begin{equation}
\label{algebra-equation-map-presentations}
\vcenter{
\xymatrix{
\bigoplus S'[(a', b')]
\oplus
\bigoplus S'[(f', g')]
\oplus
\bigoplus S'[r'] \ar[r]
&
\bigoplus S' [a'] \\
 \\
\bigoplus S[(a, b)]
\oplus
\bigoplus S[(f, g)]
\oplus
\bigoplus S[r] \ar[r]
\ar[uu]^{
\begin{matrix}
[(a, b)] \mapsto [(\varphi(a), \varphi(b))] \\
[(f, g)] \mapsto [(\varphi(f), \varphi(g))] \\
[r]\mapsto [\psi(r)]
\end{matrix}
} &
\bigoplus S[a] \ar[uu]_{[a] \mapsto [\varphi(a)]}
}
}
\end{equation}
である。その結果は単に、$f\text{d}g \in \Omega_{S/R}$ が
$\varphi(f)\text{d}\varphi(g)$ に写るということである。


\begin{lemma}
\label{algebra-lemma-colimit-differentials}
$I$ を有向集合とする。
$(R_i \to S_i, \varphi_{ii'})$ を $I$ 上の環準同型の系とする
（Categories, Section \ref{categories-section-posets-limits} を参照）。
このとき
% [P01-E0678] The frozen display punctuation is retained before the continuing clause.
$$
\Omega_{S/R} =
\colim_i \Omega_{S_i/R_i}.
$$
が成り立つ。ここで $R \to S = \colim (R_i \to S_i)$ である。
\end{lemma}

\begin{proof}
% [P01-E0679] Japanese renders the intended antecedent explicitly.
これは $\Omega_{S/R}$ の定義表示と、上で述べたその関手性から明らかである。
\end{proof}

\begin{lemma}
\label{algebra-lemma-differential-surjective}
図式 (\ref{algebra-equation-functorial-omega}) において、
$S \to S'$ が全射で、その核が $I \subset S$ であるとする。
このとき $\Omega_{S/R} \to \Omega_{S'/R'}$ は全射であり、その核は
$\varphi(a) \in \beta(R')$ を満たす $a \in S$ に対する元
$\text{d}a$ によって $S$-加群として生成される。
（特に、元 $\text{d}(i)$, $i \in I$ もこれに含まれる。）
\end{lemma}

\begin{proof}[第一の証明]
表示の写像 (\ref{algebra-equation-map-presentations}) を考える。
自由加群の右側の垂直写像は明らかに全射である。従ってこの写像は全射である。
$\Omega_{S/R}$ の元 $\eta$ が $\Omega_{S'/R'}$ で零に写るとする。
ある $s_i, a_i \in S$ に対して、$\eta$ を $\sum s_i[a_i]$ の像として書く。
すると $\sum \varphi(s_i)[\varphi(a_i)]$ は、図式の左上隅にある元
$$
\theta = \sum s_j'[a_j', b_j'] + \sum s_k'[f_k', g_k'] + \sum s_l'[r_l']
$$
の像であることが分かる。$\varphi$ は全射なので、項
$s_j'[a_j', b_j']$ および $s_k'[f_k', g_k']$ は
% [P01-E0680] Japanese names the geometrically correct lower-left source corner.
左下隅の元の像に入る。従って、これらの元の像を引いて
$\eta$ と $\theta$ を取り替えることにより、
$\theta = \sum s_l'[r_l']$ と仮定してよい。
言い換えると、$\sum \varphi(s_i)[\varphi(a_i)]$ は
$\sum s'_l [\beta(r'_l)]$ という形である。
% [P01-E0681] The reduction to a common direct-sum component is retained as stated.
次に、すべての $i$ に対して $a' = \varphi(a_i)$、すべての $l$ に対して
$a' = \beta(r_l')$ を満たすある $a' \in S'$ が存在すると仮定してよい。
これは図式の右上隅の直和分解から明らかである。
$\varphi(a) = a'$ を満たす $a \in S$ を選ぶ。
すると、ある $x_i \in I$ に対して $a_i = a + x_i$ と書ける。
従って、許されている関係を用いれば、すべての $a_i$ が $a$ に等しいと
仮定してよい。しかしこのとき、考えている元は $s[a]$ という形だと
仮定してよい。なお
% [P01-E0682] Frozen source formula is retained verbatim despite the ill-typed symbol.
$\varphi(s)[a'] = \sum \varphi(s_l')[\beta(r_l')]$ である。
従って、$\varphi(s) = 0$ ならば証明は終わり、そうでなければ
$a' = \varphi(a)$ は $\beta$ の像に入るので、この場合も証明は終わる。
\end{proof}

\begin{proof}[第二の証明]
以下では、補題 \ref{algebra-lemma-universal-omega} で与えた微分加群の普遍性を
その都度断らずに用いる。

\medskip\noindent
% [P01-E0683] Japanese supplies the missing article naturally.
(\ref{algebra-equation-functorial-omega}) において
$R'' = S \times_{S'} R'$ とおく。このとき次の図式がある。
$$
\xymatrix{
S \ar[r] & S \ar[r] & S' \\
R \ar[r] \ar[u] & R'' \ar[r] \ar[u] & R' \ar[u]
}
$$
$M$ を $S$-加群とする。定義から直ちに、$R$-導分
$D : S \to M$ が $R''$-導分であることと、
$R'' \to S$ の像に含まれる元を零に写すことは同値である。
% [P01-E0684] The frozen source's generator wording is retained; no silent repair.
普遍性により、これは自然な写像
$\Omega_{S/R} \to \Omega_{S/R''}$ が全射であり、その核が
$R''$ の像によって $S$-加群として生成されるという主張に言い換えられる。

\medskip\noindent
前段落から、$S \to S'$ が全射で
$R = S \times_{S'} R'$ であるときに
$\Omega_{S/R} \to \Omega_{S'/R'}$ が同型であることを示せば十分である。
$M'$ を $S'$-加群とする。任意の $R'$-導分 $D' : S' \to M'$ は、
$S \to S'$ を前合成することにより $R$-導分を与える。
逆に、$M$ を $S$-加群とし、$D : S \to M$ を $R$-導分とする。
$i \in I$ ならば、$R = S \times_{S'} R'$ であるため、
$\alpha(a) = i$ を満たす $a \in R$ が存在する。
従って $D(i) = 0$ であり、すべての $s \in S$ に対して
$0 = D(is) = iD(s)$ となる。それゆえ $D$ の像は、$I$ によって零化される
元からなる部分加群 $M' \subset M$ に含まれる。さらに誘導写像
$S \to M'$ は $R'$-導分 $S' \to M'$ を経由する。
これが $\Omega_{S/R} \to \Omega_{S'/R'}$ が同型であることを意味するのを
普遍性から示すのは演習とし、詳細は省略する。
\end{proof}

\begin{lemma}
\label{algebra-lemma-exact-sequence-differentials}
$A \to B \to C$ を環準同型とする。このとき $C$-加群の標準的な完全列
$$
C \otimes_B \Omega_{B/A} \to
\Omega_{C/A} \to
\Omega_{C/B} \to 0
$$
が存在する。
\end{lemma}

\begin{proof}
$R = A$, $S = C$, $R' = B$, $S' = C$ とおくことにより、
図式 (\ref{algebra-equation-functorial-omega}) が得られる。
補題 \ref{algebra-lemma-differential-surjective} により
$\Omega_{C/A} \to \Omega_{C/B}$ は全射であり、その核は
$B \to C$ の像に入る $c \in C$ に対する元 $\text{d}(c)$ によって生成される。
従って補題が従う。
\end{proof}


\begin{lemma}
\label{algebra-lemma-differentials-localize}
$\varphi : A \to B$ を環準同型とする。
\begin{enumerate}
\item $S \subset A$ が、$B$ の可逆元に写る乗法的部分集合ならば、
$\Omega_{B/A} = \Omega_{B/S^{-1}A}$ である。
\item $S \subset B$ が乗法的部分集合ならば、
$S^{-1}\Omega_{B/A} = \Omega_{S^{-1}B/A}$ である。
\end{enumerate}
\end{lemma}

\begin{proof}
(1) の等式を示すには、任意の $A$-導分 $D : B \to M$ が
元 $\varphi(s)^{-1}$ を零に写すことを示せば十分である。
これは $1 = \varphi(s) \varphi(s)^{-1}$ に Leibniz 則を適用すれば明らかである。
% [P01-E0697] Japanese punctuation supplies the introductory boundary.
% [P01-E0685] Japanese supplies the natural article and designation.
(2) を示すため、自然な写像
$S^{-1}\Omega_{B/A} \to \Omega_{S^{-1}B/A}$ があることに注意する。
これが同型であることを示すには、
$S^{-1}\Omega_{B/A}$ に値をとる $S^{-1}B$ の $A$-導分
$\text{d}'$ が存在することを示せば十分である。
その定義は単に
$\text{d}'(b/s) = (1/s)\text{d}b - (1/s^2)b\text{d}s$
とすればよい。詳細は省略する。
\end{proof}

\begin{lemma}
\label{algebra-lemma-differential-seq}
図式 (\ref{algebra-equation-functorial-omega}) において、
$S \to S'$ が全射で、その核が $I \subset S$ であるとし、
$R' = R$ と仮定する。このとき $S'$-加群の標準的な完全列
% [P01-E0692] Frozen display punctuation is retained before the following sentence.
$$
I/I^2
\longrightarrow
\Omega_{S/R} \otimes_S S'
\longrightarrow
\Omega_{S'/R}
\longrightarrow
0
$$
が存在する。最左の写像は、$f \in I$ を
$\text{d}f \otimes 1$ に写すという規則で特徴づけられる。
\end{lemma}

\begin{proof}
中央の項は $\Omega_{S/R} \otimes_S S/I$ である。
% [P01-E0686] Japanese supplies the designation particle explicitly.
$f \in I$ に対し、$I/I^2$ における $f$ の像を $\overline{f}$ と書く。
写像 $\overline{f} \mapsto \text{d}f \otimes 1$ が well-defined であることを
示すには、$f_1, f_2 \in I$ ならば
$\text{d} f_1f_2 \otimes 1 = 0$ であることを確かめればよい。
これは Leibniz 則
$\text{d} f_1f_2 \otimes 1 = (f_1 \text{d}f_2 + f_2 \text{d} f_1 )\otimes 1 =
\text{d}f_2 \otimes f_1 + \text{d}f_1 \otimes f_2 = 0$
から明らかである。
% [P01-E0687] Japanese uses the grammatically correct predicate.
同様の計算により、この写像が $S' = S/I$-線形であることも分かる。

\medskip\noindent
写像 $\Omega_{S/R} \otimes_S S' \to \Omega_{S'/R}$ は、
$S$-線形写像 $\Omega_{S/R} \to \Omega_{S'/R}$ に付随する
標準的な $S'$-線形写像である。
補題 \ref{algebra-lemma-differential-surjective} により
$\Omega_{S/R} \to \Omega_{S'/R}$ が全射なので、この写像も全射である。

\medskip\noindent
二つの写像の合成は零である。実際、$f \in I$ に対して
$\text{d}f$ は $\Omega_{S'/R}$ で零に写る。
完全性が意味するのは、$\Omega_{S/R} \to \Omega_{S'/R}$ の核が
部分加群 $I\Omega_{S/R}$ と、$f \in I$ に対する元 $\text{d}f$ とによって
$S$-部分加群として生成されるということにほかならない。
補題 \ref{algebra-lemma-differential-surjective} により、この核は
ある $r \in R$ に対して $\varphi(a) = \beta(r)$ を満たす
元 $\text{d}(a)$ によって生成されることが分かっている。
しかし $a = \alpha(r) + a - \alpha(r)$ なので、
$\text{d}(a) = \text{d}(a - \alpha(r))$ である。また
$\varphi(a - \alpha(r)) =
\varphi(a) - \varphi(\alpha(r)) = \beta(r) - \beta(r) = 0$ だから、
$a - \alpha(r) \in I$ である。
% [P01-E0696] Japanese supplies the missing complementizer.
従って、$f \in I$ に対する元 $\text{d}f$ だけで
核が $S$-加群として生成されることが分かり、望む結論を得る。
\end{proof}

\begin{lemma}
\label{algebra-lemma-differential-seq-split}
図式 (\ref{algebra-equation-functorial-omega}) において、
$S \to S'$ が全射で、その核が $I \subset S$ であるとし、
$R' = R$ と仮定する。さらに、$S \to S'$ の右逆である
$R$-代数準同型 $S' \to S$ が存在すると仮定する。
このとき、補題 \ref{algebra-lemma-differential-seq} の $S'$-加群の完全列は短完全列
$$
0 \longrightarrow
I/I^2
\longrightarrow
\Omega_{S/R} \otimes_S S'
\longrightarrow
\Omega_{S'/R}
\longrightarrow
0
$$
となり、しかもこれは分裂短完全列である。
\end{lemma}

\begin{proof}
$\beta : S' \to S$ を、全射 $\alpha : S \to S'$ の右逆とする。
写像
% [P01-E0693] Frozen display punctuation is retained before the following sentence.
$$
D : S \longrightarrow I/I^2,
\quad
x \longmapsto x - \beta(\alpha(x))
$$
を考える。$D$ が $R$-導分であることは容易に示せる（省略する）。
さらに、$x \in I, s \in S$ ならば $x D(s) = 0$ である。
従って普遍性により、$D$ は写像
$\tau : \Omega_{S/R} \otimes_S S' \to I/I^2$ を誘導する。
これが $\text{d} : I/I^2  \to \Omega_{S/R} \otimes_S S'$ の
左逆であることの確認は省略する。従って結論を得る。
\end{proof}

\begin{lemma}
\label{algebra-lemma-differential-mod-power-ideal}
$R \to S$ を環準同型とし、$I \subset S$ をイデアルとする。
$n \geq 1$ を整数とし、$S' = S/I^{n + 1}$ とおく。
写像 $\Omega_{S/R} \to \Omega_{S'/R}$ は同型
$$
\Omega_{S/R} \otimes_S S/I^n
\longrightarrow
\Omega_{S'/R} \otimes_{S'} S/I^n.
$$
を誘導する。
\end{lemma}

\begin{proof}
これは補題 \ref{algebra-lemma-differential-seq} と、
$\text{d}$ に対する Leibniz 則により
$\text{d}(I^{n + 1}) \subset I^n\Omega_{S/R}$ となることから従う。
\end{proof}

\begin{lemma}
\label{algebra-lemma-differentials-base-change}
環準同型 $R \to R'$ および $R \to S$ が与えられているとする。
$S' = S \otimes_R R'$ とおけば、図式
(\ref{algebra-equation-functorial-omega}) が得られる。このとき、上で定義した
標準的な写像は同型 $\Omega_{S/R} \otimes_R R' = \Omega_{S'/R'}$ を誘導する。
\end{lemma}

\begin{proof}
$\text{d}' : S' = S \otimes_R R' \to \Omega_{S/R} \otimes_R R'$ を、
$\text{d}'( \sum a_i \otimes x_i ) = \sum \text{d}(a_i) \otimes x_i$ で
与えられる写像とする。写像
$S \times R' \to \Omega_{S/R} \otimes_R R'$,
$(a, x)\mapsto \text{d}a \otimes_R x$ が $R$-双線形なので、これは存在する。
簡単な計算により、これが $R'$-導分であることを確かめられる。
$(\Omega_{S/R} \otimes_R R', \text{d}')$ が普遍性を満たすことを示そう。
$D : S' \to M'$ を、ある $S'$-加群に値をとる
$R'$-導分とする。合成 $S \to S' \to M'$ は $R$-導分なので、
$S$-線形写像 $\varphi_D : \Omega_{S/R} \to M'$ が得られる。
これを $R'$ とテンソルして写像 $\varphi'_D :
\Omega_{S/R} \otimes_R R' \to M'$,
$\varphi'_D(\eta \otimes x) = x\varphi_D(\eta)$ を得る。
$D = \varphi'_D \circ \text{d}'$ であることは明らかである。
\end{proof}

\noindent
乗法写像 $S \otimes_R S \to S$ とは、
$a \otimes b$ を $S$ の $ab$ に写す $R$-代数準同型である。
二つの写像 $S \to S \otimes_R S$ のいずれかを介して
$S \otimes_R S$ を $S$-代数とみなせば、これは $S$-代数準同型でもある。

\begin{lemma}
\label{algebra-lemma-differentials-diagonal}
$R \to S$ を環準同型とする。$J = \Ker(S \otimes_R S \to S)$ を
乗法写像の核とする。$S$-加群の標準的な同型
$\Omega_{S/R} \to J/J^2$,
$a \text{d} b \mapsto a \otimes b - ab \otimes 1$ が存在する。
\end{lemma}

\begin{proof}[第一の証明]
補題 \ref{algebra-lemma-differential-seq-split} を可換図式
$$
\xymatrix{
S \otimes_R S \ar[r] & S \\
S \ar[r] \ar[u] & S \ar[u]
}
$$
に適用する。ここで左の垂直矢印は $a \mapsto a \otimes 1$ である。
完全列
$0 \to J/J^2 \to
\Omega_{S \otimes_R S/S} \otimes_{S \otimes_R S} S \to \Omega_{S/S} \to 0$
が得られる。補題 \ref{algebra-lemma-trivial-differential-surjective} により
項 $\Omega_{S/S}$ は $0$ なので、残りの二項の間の同型を得る。
$S \to S \otimes_R S$ は $R \to S$ の基底変換だから、
補題 \ref{algebra-lemma-differentials-base-change} により
$\Omega_{S \otimes_R S/S} = \Omega_{S/R} \otimes_S (S \otimes_R S)$
であり、従って
% [P01-E0694] Frozen display punctuation is retained before the following sentence.
$$
\Omega_{S \otimes_R S/S} \otimes_{S \otimes_R S} S =
\Omega_{S/R} \otimes_S (S \otimes_R S) \otimes_{S \otimes_R S} S =
\Omega_{S/R}
$$
となる。この写像が補題の規則で与えられることの確認は省略する。
\end{proof}

\begin{proof}[第二の証明]
まず、規則 $a \text{d} b \mapsto a \otimes b - ab \otimes 1$ が
well-defined であることを示す。そのためには、
% [P01-E0688] Frozen source math token d(ab) is retained verbatim.
$\text{d}r$ および $a\text{d}b + b \text{d}a - d(ab)$ が
零に写ることを示さなければならない。
% [P01-E0689] Japanese supplies complete causal sentences for the frozen fragments.
前者については、テンソル積の定義により
$r \otimes 1 - 1 \otimes r = 0$ だからである。後者については、
$$
(a \otimes b - ab \otimes 1) +
(b \otimes a - ba \otimes 1) -
(1 \otimes ab - ab \otimes 1)
=
(a \otimes 1 - 1\otimes a)(1\otimes b - b \otimes 1)
$$
が $J^2$ に属するからである。

\medskip\noindent
逆向きの写像を構成する。
$S \to S \otimes_R S$, $a \mapsto a \otimes 1$ を
$R \to S$ の基底変換とみなすことができる。
従って補題 \ref{algebra-lemma-differentials-base-change} により
$\Omega_{S \otimes_R S/S} = \Omega_{S/R} \otimes_S (S \otimes_R S)$
である。ここで、補題 \ref{algebra-lemma-differential-seq} の列は写像
$$
J/J^2  \to \Omega_{S \otimes_R S/ S} \otimes_{S \otimes_R S} S
= (\Omega_{S/R} \otimes_S (S \otimes_R S))\otimes_{S \otimes_R S} S
= \Omega_{S/R}.
$$
を与える。
% [P01-E0690] Japanese supplies the missing complementizer.
これが上の写像の逆写像であることの確認は読者に任せる。
\end{proof}




\begin{lemma}
\label{algebra-lemma-differentials-polynomial-ring}
$S = R[x_1, \ldots, x_n]$ ならば、$\Omega_{S/R}$ は
基底 $\text{d}x_1, \ldots, \text{d}x_n$ をもつ有限自由 $S$-加群である。
\end{lemma}

\begin{proof}
まず、$\text{d}x_1, \ldots, \text{d}x_n$ が $\Omega_{S/R}$ を
$S$-加群として生成することを示す。そのため、任意の
$g \in R[x_1, \ldots, x_n]$ に対して $\text{d}g$ が和
$\sum g_i \text{d}x_i$ として表せることを示す。
$g$ の（全）次数に関する帰納法を用いる。
$g$ の次数が $0$ ならば $\text{d}g = 0$ だから明らかである。
$g$ の次数が $> 0$ ならば、$c\in R$ および
$\deg(g_i) < \deg(g)$ を満たす形、すなわち $g$ を
$c + \sum g_i x_i$ に書ける。
Leibniz 則により
$\text{d}g = \sum g_i \text{d} x_i + \sum x_i \text{d}g_i$ だから、
帰納法によって結論を得る。

\medskip\noindent
$R$-導分 $\partial / \partial x_i :
R[x_1, \ldots, x_n] \to R[x_1, \ldots, x_n]$ を考える。
（その定義は読者に任せる。定義的性質は
$\partial / \partial x_i (x_j) = \delta_{ij}$ である。）
普遍性により、これは $\text{d}x_i$ を $1$ に写し、
$j \not = i$ ならば $\text{d}x_j$ を $0$ に写す
$S$-加群準同型 $l_i :
\Omega_{S/R} \to R[x_1, \ldots, x_n]$ に対応する。
従って、元 $\text{d}x_1, \ldots, \text{d}x_n$ の間に
$S$-線形関係がないことは明らかである。
\end{proof}

\begin{lemma}
\label{algebra-lemma-differentials-finitely-presented}
$R \to S$ が有限表示ならば、$\Omega_{S/R}$ は有限表示
$S$-加群である。
\end{lemma}

\begin{proof}
$S = R[x_1, \ldots, x_n]/(f_1, \ldots, f_m)$ と書き、
$I = (f_1, \ldots, f_m)$ と書く。
補題 \ref{algebra-lemma-differential-seq} によれば、$S$-加群の完全列
% [P01-E0695] Frozen display punctuation is retained before the following sentence.
$$
I/I^2
\to
\Omega_{R[x_1, \ldots, x_n]/R} \otimes_{R[x_1, \ldots, x_n]} S
\to
\Omega_{S/R}
\to
0
$$
がある。$I/I^2$ が（$f_i$ の像によって生成される）有限
$S$-加群であること、および補題
\ref{algebra-lemma-differentials-polynomial-ring} により中央の項が
有限自由であることから、結論が従う。
\end{proof}

\begin{lemma}
\label{algebra-lemma-differentials-finitely-generated}
$R \to S$ が有限型ならば、$\Omega_{S/R}$ は有限生成
% [P01-E0691] Japanese supplies the missing article naturally.
$S$-加群である。
\end{lemma}

\begin{proof}
これは補題 \ref{algebra-lemma-differentials-finitely-presented} の証明と
ほぼ同じだが、それより容易である。
\end{proof}







\section{de Rham 複体}
\label{algebra-section-de-rham-complex}

\noindent
$A \to B$ を環準同型とする。Section \ref{algebra-section-differentials} で構成した
% [P01-E0698] Japanese identifies the frozen type-confused designation as the universal derivation.
微分加群への普遍 $A$-導分を
$\text{d} : B \to \Omega_{B/A}$ と書く。
Section \ref{algebra-section-tensor-algebra} と同様に、$i \geq 0$ に対して
$\Omega_{B/A}^i = \wedge^i_B(\Omega_{B/A})$ を第 $i$ 外冪とする。
{\it $A$ 上の $B$ の de Rham 複体}とは、以下で構成し記述する
$A$-線形微分を備えた複体
$$
\Omega_{B/A}^0 \to \Omega_{B/A}^1 \to \Omega_{B/A}^2 \to \ldots
$$
のことである。

\medskip\noindent
写像 $\text{d} : \Omega^0_{B/A} \to \Omega^1_{B/A}$ は
普遍導分 $\text{d} : B \to \Omega_{B/A}$ である。
これは実際に $A$-線形であることに注意する。

\medskip\noindent
$p \geq 1$ に対し、一意な $A$-線形写像
$\text{d} : \Omega_{B/A}^p \to \Omega_{B/A}^{p + 1}$ が存在し、
次を満たすと主張する。
\begin{equation}
\label{algebra-equation-rule}
\text{d}\left(b_0\text{d}b_1 \wedge \ldots \wedge \text{d}b_p\right) =
\text{d}b_0 \wedge \text{d}b_1 \wedge \ldots \wedge \text{d}b_p
\end{equation}
% [P01-E0699] Japanese supplies the sentence boundary after the frozen equation.
$\Omega_{B/A}$ は元 $\text{d}b$ によって $B$-加群として生成されることを
思い出そう。従って $\Omega^p_{B/A}$ は元
$b_0\text{d}b_1 \wedge \ldots \wedge \text{d}b_p$ によって
$A$-加群として生成されるので、写像
$\text{d} : \Omega^p_{B/A} \to \Omega^{p + 1}_{B/A}$ が存在するならば
一意である。

\medskip\noindent
$\text{d} : \Omega_{B/A}^1 \to \Omega_{B/A}^2$ の構成。
Definition \ref{algebra-definition-differentials} により、元 $\text{d}b$ は
関係式 $\text{d}a = 0$（$a \in A$）および
$\text{d}(b' + b'') = \text{d}b' + \text{d}b''$、
$\text{d}(b'b'') = b'\text{d}b'' + b''\text{d}b'$（$b', b'' \in B$）の
もとで、$\Omega_{B/A}$ を $B$-加群として自由に生成する。
従って、規則
$$
\sum b'_i \text{d}b_i \longmapsto \sum \text{d}b'_i \wedge \text{d}b_i
$$
が well-defined であることを示すには、$a \in A$ および
$b, b', b'' \in B$ に対する元
$$
b\text{d}a,
\quad\text{and}\quad
b\text{d}(b' + b'') - b\text{d}b' - b\text{d}b''
\quad\text{and}\quad
b\text{d}(b'b'') - bb'\text{d}b'' - bb''\text{d}b'
$$
が零に写ることを示さなければならない。
これは $\text{d}$ に対する Leibniz 則を用いた直接計算から明らかである。

\medskip\noindent
合成 $\Omega^0_{B/A} \to \Omega^1_{B/A} \to \Omega^2_{B/A}$ は零である。
実際、
$\text{d}(\text{d}(b)) = \text{d}(1 \text{d}b) =
\text{d}(1) \wedge \text{d}(b) = 0 \wedge \text{d}b = 0$ である。
ここで、$1 \in B$ は $A \to B$ の像に入るので
$\text{d}(1) = 0$ である。以下ではこのことを用いる。

\medskip\noindent
$p \geq 2$ に対する
$\text{d} : \Omega_{B/A}^p \to \Omega_{B/A}^{p + 1}$ の構成。
$A$-線形写像
$$
\gamma : \Omega^1_{B/A} \otimes_A \ldots \otimes_A \Omega^1_{B/A}
\longrightarrow
\Omega_{B/A}^{p + 1}
$$
を公式
$$
\omega_1 \otimes \ldots \otimes \omega_p
\longmapsto
\sum (-1)^{i + 1}
\omega_1 \wedge \ldots \wedge \text{d}(\omega_i) \wedge \ldots \wedge \omega_p
$$
% [P01-E0700] The frozen unbounded summation is retained verbatim.
で定義する。これが自然な全射
$\Omega^1_{B/A} \otimes_A \ldots \otimes_A \Omega^1_{B/A} \to \Omega^p_{B/A}$
を経由し、求める写像
$\text{d} : \Omega^p_{B/A} \to \Omega^{p + 1}_{B/A}$ を与えることを示す。
補題 \ref{algebra-lemma-present-wedge} によれば、
$\Omega^1_{B/A} \otimes_A \ldots \otimes_A \Omega^1_{B/A} \to \Omega^p_{B/A}$
の核は、ある $i \not = j$ に対して $\omega_i = \omega_j$ である元
$\omega_1 \otimes \ldots \otimes \omega_p$、およびある $f \in B$ に対する元
$\omega_1 \otimes \ldots \otimes f\omega_i \otimes \ldots \otimes \omega_p -
\omega_1 \otimes \ldots \otimes f\omega_j \otimes \ldots \otimes \omega_p$
によって $A$-加群として生成される。
直接計算により、第一の型の元は $\gamma$ によって $0$ に写る。
言い換えると、$\gamma$ は交代的である。
証明を終えるには、
$$
\gamma(
\omega_1 \otimes \ldots \otimes f\omega_i \otimes \ldots \otimes \omega_p)
=
\gamma(
\omega_1 \otimes \ldots \otimes f\omega_j \otimes \ldots \otimes \omega_p)
$$
が $f \in B$ に対して成り立つことを示さなければならない。
$A$-線形性と交代性により、$p = 2$, $i = 1$, $j = 2$,
$\omega_1 = b \text{d}b'$ および $\omega_2 = c \text{d} c'$、
ここで $b, b', c, c' \in B$、の場合を示せば十分である。
従って、示すべき等式は
\begin{align*}
& \text{d}(fb) \wedge \text{d}b' \wedge c \text{d}c'
- fb \text{d}b' \wedge \text{d}c \wedge \text{d}c' \\
& =
\text{d}b \wedge \text{d}b' \wedge fc\text{d}c'
- b \text{d}b' \wedge \text{d}(fc) \wedge \text{d}c'
\end{align*}
% [P01-E0701] Japanese supplies the continuing reformulation boundary.
である。言い換えると、
$$
(c \text{d}(fb) + fb \text{d}c - fc \text{d}b - b \text{d}(fc))
\wedge \text{d}b' \wedge \text{d}c' = 0.
$$
である。これは Leibniz 則から従う。
元
$b_0\text{d}b_1 \otimes \text{d}b_2 \otimes \ldots \otimes \text{d}b_p$
における $\gamma$ の値は
$\text{d}b_0 \wedge \text{d}b_1 \wedge \ldots \wedge \text{d}b_p$
であることに注意する。従って、このように得られた写像
$\text{d} : \Omega^p_{B/A} \to \Omega^{p + 1}_{B/A}$ は
(\ref{algebra-equation-rule}) を満たす。

\medskip\noindent
最後に、$\Omega^p_{B/A}$ は元
$b_0\text{d}b_1 \wedge \ldots \wedge \text{d}b_p$ によって加法的に生成され、
$\text{d}(b_0\text{d}b_1 \wedge \ldots \wedge \text{d}b_p) =
\text{d}b_0 \wedge \ldots \wedge \text{d}b_p$ であるから、まったく同様に、
$p \geq 1$ に対して合成
$
\Omega^p_{B/A} \to
\Omega^{p + 1}_{B/A} \to
\Omega^{p + 2}_{B/A}
$
が零であることが分かる。従って de Rham 複体は実際に複体である。

\medskip\noindent
環 $R$ だけが与えられた場合には $\Omega_R = \Omega_{R/\mathbf{Z}}$ とおく。
これは $R$ の絶対微分加群と呼ばれることもある。この呼称は妥当である。
実際、$\Omega_R$ を、$\text{d}1$ の消滅を仮定せず Leibniz 則だけを
仮定して得られる微分加群とすると、Leibniz 則から
$\text{d}1 = \text{d}(1 \cdot 1) =
1 \text{d}1 + 1 \text{d}1 = 2 \text{d}1$、従って
$\Omega_R$ において $\text{d}1 = 0$ を得る。
この場合、{\it $R$ の絶対 de Rham 複体}とは対応する複体
% [P01-E0702] Japanese supplies the continuing clause boundary.
$$
\Omega_R^0 \to \Omega_R^1 \to \Omega_R^2 \to \ldots
$$
のことである。ここで $\Omega^i_R = \Omega^i_{R/\mathbf{Z}}$ などとおく。

\medskip\noindent
環の可換図式
% [P01-E0703] Japanese supplies a sentence boundary after the frozen diagram.
$$
\xymatrix{
B \ar[r] & B' \\
A \ar[r] \ar[u] & A' \ar[u]
}
$$
が与えられているとする。de Rham 複体の自然な写像
% [P01-E0704] Japanese supplies a sentence boundary after the frozen display.
$$
\Omega^\bullet_{B/A} \longrightarrow \Omega^\bullet_{B'/A'}
$$
が存在する。具体的には、次数 $0$ では写像 $B \to B'$、
次数 $1$ では Section \ref{algebra-section-differentials} で構成した写像
$\Omega_{B/A} \to \Omega_{B'/A'}$ であり、$p \geq 2$ に対しては
誘導写像
$\Omega^p_{B/A} = \wedge^p_B(\Omega_{B/A}) \to
\wedge^p_{B'}(\Omega_{B'/A'}) = \Omega^p_{B'/A'}$ である。
微分との両立性は、公式 (\ref{algebra-equation-rule}) による微分の特徴づけから従う。

\begin{lemma}
\label{algebra-lemma-de-rham-complex-base-change}
環準同型 $A \to A'$ および $A \to B$ が与えられているとする。
$B' = B \otimes_A A'$ とおけば、上のような図式が得られる。
このとき、上で定義した標準的な写像は複体の同型
$\Omega^\bullet_{B/A} \otimes_A A' = \Omega^\bullet_{B'/A'}$
を誘導する。
\end{lemma}

\begin{proof}
これは補題 \ref{algebra-lemma-differentials-base-change} と、
外冪をとる操作が基底変換と可換であることから従う。
\end{proof}

\begin{lemma}
\label{algebra-lemma-de-rham-complex}
$A \to B$ を環準同型とする。$\pi : \Omega_{B/A} \to \Omega$ を
全射 $B$-加群準同型とする。
% [P01-E0705] Japanese uses a complete designation construction.
$\pi$ と普遍導分 $\text{d}_{B/A} : B \to \Omega_{B/A}$ の合成を
$\text{d} : B \to \Omega$ と書く。$\Omega^i = \wedge_B^i(\Omega)$ とおく。
$\pi$ の核は、$\text{d}_{B/A}(\omega) \in \Omega_{B/A}^2$ が
$\Omega^2$ において零に写るような元
$\omega \in \Omega_{B/A}$ によって $B$-加群として生成されると仮定する。
このとき de Rham 複体
$$
\Omega^0 \to \Omega^1 \to \Omega^2 \to \ldots
$$
が存在し、その微分は規則
$$
\text{d} : \Omega^p \to \Omega^{p + 1},\quad
\text{d}\left(f_0\text{d}f_1 \wedge \ldots \wedge \text{d}f_p\right) =
\text{d}f_0 \wedge \text{d}f_1 \wedge \ldots \wedge \text{d}f_p
$$
% [P01-E0706] Japanese supplies terminal punctuation after the frozen display.
によって定義される。
\end{lemma}

\begin{proof}
可換図式
$$
\xymatrix{
\Omega_{B/A}^0 \ar[d] \ar[r]_{\text{d}_{B/A}} &
\Omega_{B/A}^1 \ar[d]_\pi \ar[r]_{\text{d}_{B/A}} &
\Omega_{B/A}^2 \ar[d]_{\wedge^2\pi} \ar[r]_{\text{d}_{B/A}} &
\ldots \\
\Omega^0 \ar[r]^{\text{d}} &
\Omega^1 \ar[r]^{\text{d}} &
\Omega^2 \ar[r]^{\text{d}} &
\ldots
}
$$
% [P01-E0707] Japanese separates the two clauses after the frozen diagram.
が存在することを示す。写像 $\text{d}$ の記述は、上で与えた
$A$ 上の $B$ の de Rham 複体の微分
$\text{d}_{B/A} : \Omega^p_{B/A} \to \Omega^{p + 1}_{B/A}$ の
構成から従う。
% [P01-E0708] Japanese renders the standard closed adjective as 最左.
最左の垂直矢印は同型なので、第一の正方形が得られる。
$\pi$ は全射だから、第二の正方形を得るには
$\text{d}_{B/A}$ が $\pi$ の核を $\wedge^2\pi$ の核に写すことを
示せば十分である。$\pi$ の核の任意の元は
$\pi(\omega_i) = 0$ および
$\wedge^2\pi(\text{d}_{B/A}(\omega_i)) = 0$ を満たす形
$\sum b_i\omega_i$ であることが仮定されている。
$\text{d}_{B/A}$ に対する Leibniz 則により
$\text{d}_{B/A}(\sum b_i\omega_i) = \sum b_i \text{d}_{B/A}(\omega_i) +
\sum \text{d}_{B/A}(b_i) \wedge \omega_i$ である。
従って、これは $\wedge^2\pi$ によって零に写る。

\medskip\noindent
$i > 1$ に対し、$\wedge^i \pi$ は全射であり、その核は
$\Ker(\pi) \wedge \Omega^{i - 1}_{B/A}
\to \Omega_{B/A}^i$ の像であることに注意する。
$\omega_1 \in \Ker(\pi)$ および
$\omega_2 \in \Omega^{i - 1}_{B/A}$ に対して
% [P01-E0709] Japanese supplies the clause boundary after the frozen equation.
$$
\text{d}_{B/A}(\omega_1 \wedge \omega_2) =
\text{d}_{B/A}(\omega_1) \wedge \omega_2 -
\omega_1 \wedge \text{d}_{B/A}(\omega_2)
$$
であり、上で示したことによりこれは $\wedge^{i + 1}\pi$ の核に入る。
従って上の図式の第 $(i + 1)$ 正方形が得られる。
これで証明が完了する。
\end{proof}

\section{有限階微分作用素}
\label{algebra-section-differential-operators}

\noindent
本節では有限階の微分作用素を導入する。

\begin{definition}
\label{algebra-definition-differential-operators}
$R \to S$ を環準同型とし、$M$, $N$ を $S$-加群とする。
$k \geq 0$ を整数とする。{\it 階数 $k$ の微分作用素 $D : M \to N$} を、
すべての $g \in S$ に対して写像
$m \mapsto D(gm) - gD(m)$ が階数 $k - 1$ の微分作用素となるような
$R$-線形写像として帰納的に定義する。基底の場合 $k = 0$ には、
階数 $0$ の微分作用素を $S$-線形写像と定義する。
\end{definition}

\noindent
$D : M \to N$ が階数 $k$ の微分作用素ならば、すべての $g \in S$ に対して
写像 $gD$ も階数 $k$ の微分作用素である。二つの階数 $k$ の
微分作用素の和も同じ階数の微分作用素である。従って、これらすべての集合
$$
\text{Diff}^k(M, N) = \text{Diff}^k_{S/R}(M, N)
$$
は $S$-加群である。また
% [P01-E0710] Japanese supplies terminal punctuation after the frozen display.
$$
\text{Diff}^0(M, N) \subset
\text{Diff}^1(M, N) \subset
\text{Diff}^2(M, N) \subset \ldots
$$
が成り立つ。

\begin{lemma}
\label{algebra-lemma-composition-differential-operators}
$R \to S$ を環準同型とし、$L, M, N$ を $S$-加群とする。
$D : L \to M$ と $D' : M \to N$ がそれぞれ階数 $k$ と $k'$ の
微分作用素ならば、$D' \circ D$ は階数 $k + k'$ の微分作用素である。
\end{lemma}

\begin{proof}
$g \in S$ とする。このとき $x \in L$ を
$$
D'(D(gx)) - gD'(D(x)) = D'(D(gx)) - D'(gD(x)) + D'(gD(x)) - gD'(D(x))
$$
に写す写像は、より低い階数の微分作用素二つの合成の和である。
従って $k + k'$ に関する帰納法により補題が従う。
\end{proof}

\begin{lemma}
\label{algebra-lemma-module-principal-parts}
$R \to S$ を環準同型とし、$M$ を $S$-加群とする。
$k \geq 0$ とする。このとき $S$-加群 $P^k_{S/R}(M)$ と、
$S$-加群 $N$ に関して関手的な標準的同型
$$
\text{Diff}^k_{S/R}(M, N) = \Hom_S(P^k_{S/R}(M), N)
$$
が存在する。
\end{lemma}

\noindent
Yoneda の補題により、これは階数 $k$ の「普遍」微分作用素
$D_{univ} : M \to P^k_{S/R}(M)$ が存在し、任意の階数 $k$ の
微分作用素 $D : M \to N$ に対して、一意な $S$-線形写像
$\alpha : P^k_{S/R}(M) \to N$ で
% [P01-E0711] Frozen ill-typed composition order is retained verbatim.
$D = D_{univ} \circ \alpha$ を満たすものが存在することを示している。
対 $(P^k_{S/R}(M), D_{univ})$ は一意な同型を除いて一意である。

\begin{proof}
% [P01-E0712] The published editorial placeholder is preserved in Japanese.
$P^k_{S/R}(M)$ の存在は一般的な圏論的議論
（ここに将来の参照を挿入する）から従うが、構成も与える。
$F = \bigoplus_{m \in M} S[m]$ とおく。ここで $[m]$ は、
$m$ に対応する直和因子の基底元を表す記号である。
任意の微分作用素 $D : M \to N$ が与えられると、$[m]$ を
$D(m)$ に写す $S$-線形写像 $L_D : F \to N$ が得られる。
$D$ の階数が $0$ ならば、$L_D$ は元
$$
[m + m'] - [m] - [m'],\quad
g_0[m] - [g_0m]
$$
を零に写す。ここで $g_0 \in S$ および $m, m' \in M$ である。
$D$ の階数が $1$ ならば、$L_D$ は元
$$
[m + m'] - [m] - [m'],\quad
f[m] - [fm], \quad
g_0g_1[m] - g_0[g_1m] - g_1[g_0m] + [g_1g_0m]
$$
を零に写す。ここで
$f \in R$, $g_0, g_1 \in S$, $m \in M$ である。
$D$ の階数が $k$ ならば、$L_D$ は元
$[m + m'] - [m] - [m']$, $f[m] - [fm]$ および
% [P01-E0713] Japanese supplies terminal punctuation after the frozen display.
$$
g_0g_1\ldots g_k[m] - \sum g_0 \ldots \hat g_i \ldots g_k[g_im] + \ldots
+(-1)^{k + 1}[g_0\ldots g_km]
$$
を零に写す。逆に、$L : F \to N$ が前文に列挙したすべての元を
零に写す $S$-線形写像ならば、$m \mapsto L([m])$ は階数 $k$ の
微分作用素である。従って $P^k_{S/R}(M)$ は、これらの元が生成する
部分加群による $F$ の商であることが分かる。
\end{proof}

\begin{definition}
\label{algebra-definition-module-principal-parts}
$R \to S$ を環準同型とし、$M$ を $S$-加群とする。
補題 \ref{algebra-lemma-module-principal-parts} で構成した加群
$P^k_{S/R}(M)$ を $M$ の{\it 階数 $k$ の主部加群}という。
\end{definition}

\noindent
包含
$$
\text{Diff}^0(M, N) \subset
\text{Diff}^1(M, N) \subset
\text{Diff}^2(M, N) \subset \ldots
$$
は、Yoneda の補題（Categories, Lemma \ref{categories-lemma-yoneda}）を介して
全射
% [P01-E0714] Japanese supplies terminal punctuation after the frozen display.
$$
\ldots \to P^2_{S/R}(M) \to P^1_{S/R}(M) \to P^0_{S/R}(M) = M
$$
に対応する。

\begin{example}
\label{algebra-example-derivations-and-differential-operators}
$R \to S$ を環準同型とし、$N$ を $S$-加群とする。
$\text{Diff}^1(S, N) = \text{Der}_R(S, N) \oplus N$ であることに注意する。
実際、$D : S \to N$ が階数 $1$ の微分作用素ならば、
$\sigma_D(g) := D(g) - gD(1)$ によって定義される
$\sigma_D : S \to N$ は $R$-導分であり、
$D = \sigma_D + \lambda_{D(1)}$ である。ここで
$\lambda_x : S \to N$ は $g$ を $gx$ に写す線形写像である。
$\Omega_{S/R}$ の普遍性により
$P^1_{S/R} = \Omega_{S/R} \oplus S$ が従う。
\end{example}

\begin{lemma}
\label{algebra-lemma-sequence-of-principal-parts}
$R \to S$ を環準同型とし、$M$ を $S$-加群とする。
$M$ に関して関手的な標準的短完全列
$$
0 \to \Omega_{S/R} \otimes_S M \to P^1_{S/R}(M) \to M \to 0
$$
が存在し、これを{\it 主部列}という。
\end{lemma}

\begin{proof}
写像 $P^1_{S/R}(M) \to M$ は上で与えた。
$N$ を $S$-加群とし、$D : M \to N$ を階数 $1$ の微分作用素とする。
$m \in M$ に対し、写像
$$
g \longmapsto D(gm) - gD(m)
$$
は、階数 $1$ の微分作用素の公理により $R$-導分 $S \to N$ である。
従って、これは規則 $a\text{d}b \mapsto aD(bm) - abD(m)$ によって
決まる線形写像 $D_m : \Omega_{S/R} \to N$ に対応する
（補題 \ref{algebra-lemma-universal-omega} を参照）。写像
$$
\Omega_{S/R} \times M \longrightarrow N,\quad
(\eta, m) \longmapsto D_m(\eta)
$$
は $S$-双線形であり（詳細は省略する）、従って $S$-線形写像
% [P01-E0715] Japanese supplies terminal punctuation after the frozen display.
$$
\sigma_D : \Omega_{S/R} \otimes_S M \to N
$$
を定める。このようにして、$N$ に関して関手的な写像
$\text{Diff}^1(M, N) \to \Hom_S(\Omega_{S/R} \otimes_S M, N)$,
$D \mapsto \sigma_D$ を得る。
% [P01-E0716] Japanese uses the complete Yoneda correspondence construction.
Yoneda の補題により、これは写像
$\Omega_{S/R} \otimes_S M \to P^1_{S/R}(M)$ に対応する。
この写像が $M$ に関して関手的であることは構成から直ちに分かる。
列
$$
\Omega_{S/R} \otimes_S M \to P^1_{S/R}(M) \to M \to 0
$$
は完全である。実際、任意の加群 $N$ に対して列
$$
0 \to \Hom_S(M, N) \to
\text{Diff}^1(M, N) \to
\Hom_S(\Omega_{S/R} \otimes_S M, N)
$$
が直接の確認により完全だからである。

\medskip\noindent
$\Omega_{S/R} \otimes_S M \to P^1_{S/R}(M)$ が単射であることを
次のように示す。$F$ が自由 $S$-加群である完全列
$$
0 \to M' \to F \to M \to 0
$$
を選ぶ。これは、すべての $N$ に対して完全列
$$
0 \to \text{Diff}^1(M, N) \to \text{Diff}^1(F, N) \to \text{Diff}^1(M', N)
$$
を誘導する。従って、可換図式
$$
\xymatrix{
0 \ar[r] &
\Omega_{S/R} \otimes_S M' \ar[r] \ar[d] &
P^1_{S/R}(M') \ar[r] \ar[d] &
M' \ar[r] \ar[d] &
0 \\
0 \ar[r] &
\Omega_{S/R} \otimes_S F \ar[r] \ar[d] &
P^1_{S/R}(F) \ar[r] \ar[d] &
F \ar[r] \ar[d] &
0 \\
0 \ar[r] &
\Omega_{S/R} \otimes_S M \ar[r] \ar[d] &
P^1_{S/R}(M) \ar[r] \ar[d] &
M \ar[r] \ar[d] &
0 \\
& 0 & 0 & 0
}
$$
の中央の列は完全である。左の列は
$\Omega_{S/R} \otimes_S -$ の右完全性により完全である。
蛇の補題（Section \ref{algebra-section-snake} を参照）により、
自由加群 $F$ に対して左側の完全性を証明すれば十分である。
$P^1_{S/R}(-)$ が直和と可換であることを用いて、$M = S$ の場合に帰着する。
この場合は Example \ref{algebra-example-derivations-and-differential-operators} の
議論から従う。
\end{proof}

\begin{remark}
\label{algebra-remark-functoriality-principal-parts}
環の可換図式
% [P01-E0717] Japanese supplies the coordinated-list boundary after the diagram.
$$
\xymatrix{
B \ar[r] & B' \\
A \ar[u] \ar[r] & A' \ar[u]
}
$$
と、$B$-加群 $M$、$B'$-加群 $M'$、および $B$-線形写像
$M \to M'$ が与えられているとする。このとき、両立する加群準同型の系
% [P01-E0718] Japanese supplies terminal punctuation after the frozen diagram.
$$
\xymatrix{
\ldots \ar[r] &
P^2_{B'/A'}(M') \ar[r] &
P^1_{B'/A'}(M') \ar[r] &
P^0_{B'/A'}(M') \\
\ldots \ar[r] &
P^2_{B/A}(M) \ar[r] \ar[u] &
P^1_{B/A}(M) \ar[r] \ar[u] &
P^0_{B/A}(M) \ar[u]
}
$$
が得られる。これらの写像は、この型の写像をさらに合成することと両立する。
これを見る最も容易な方法は、補題
\ref{algebra-lemma-module-principal-parts} の証明における加群
$P^k_{B/A}(M)$ の生成元と関係式による記述を用いることである。
% [P01-E0719] Japanese separates the independent clauses.
一方、これらの加群の普遍性から直接見ることもできる。
さらに、これらの写像は補題
\ref{algebra-lemma-sequence-of-principal-parts} の短完全列と両立する。
\end{remark}

\begin{lemma}
\label{algebra-lemma-principal-parts-base-change}
環準同型 $A \to A'$ および $A \to B$ と、$B$-加群 $M$ が
与えられているとする。$B' = B \otimes_A A'$ とおき、
$M' = M \otimes_A A'$ を $B'$-加群とみなす。
Remark \ref{algebra-remark-functoriality-principal-parts} の写像は同型
$P^k_{B/A}(M) \otimes_A A' = P^k_{B'/A'}(M')$ を誘導する。
\end{lemma}

\begin{proof}
$D_{univ} : M \to P^k_{B/A}(M)$ を普遍微分作用素とする。
誘導される $A'$-線形写像
$D_{univ} \otimes 1 : M' = M \otimes_A A' \to P^k_{B/A}(M) \otimes_A A'$
が $M'/B'/A'$ に対する階数 $k$ の微分作用素であることは容易に確かめられる。
$D_{univ} \otimes 1$ が普遍性を満たすことを示す。
$D' : M' \to N'$ を、$B'$-加群 $N'$ に値をとる階数 $k$ の
微分作用素とする。このとき、合成 $D : M \to M' \to N'$ は、
$N'$ を $B$-加群とみなしたものに値をとる微分作用素である。
従って $D = \gamma \circ D_{univ}$ を満たす $B$-線形写像
$\gamma : P^k_{B/A}(M) \to N'$ が存在する。
誘導される $B'$-線形写像
$\gamma' : P^k_{B/A}(M) \otimes_A A' \to N'$ が
$\gamma' \circ (D_{univ} \otimes 1) = D'$ を満たすことは
読者が容易に確認できる。$\gamma'$ の一意性の証明は省略する。
\end{proof}

\begin{lemma}
\label{algebra-lemma-principal-parts-diagonal}
$R \to S$ を環準同型とし、$M$ を $S$-加群とする。
$J = \Ker(S \otimes_R S \to S)$ を乗法写像の核とする。
$S$-加群の標準的同型
$$
P^k_{S/R}(M) \longrightarrow (S \otimes_R M)/J^{k + 1}(S \otimes_R M)
$$
が存在する。ここで $s \in S$ は、$s \otimes 1$ による乗法を介して
右辺に作用する。
% [P01-E0720] Japanese supplies the missing singular correspondence predicate.
さらに、階数 $k$ の普遍微分作用素は、
$m \mapsto \text{class of }1 \otimes m$ で与えられる写像に対応する。
\end{lemma}

\begin{proof}
% [P01-E0721] Frozen tensor product without the base-ring subscript is retained verbatim.
% [P01-E0722] Japanese supplies terminal punctuation after the frozen inline map.
写像 $T : M \to S \otimes M$, $m \mapsto 1 \otimes m$ を考える。
$T$ は $R$-線形であり、また
\begin{align*}
g_0g_1\ldots g_k T(m)
- \sum g_0 \ldots \hat g_i \ldots g_k T(g_im) + \ldots
+(-1)^{k + 1}T(g_0\ldots g_km) \\
=
\prod_{i = 0,\ldots,k} (g_i \otimes 1 - 1 \otimes g_i) 1 \otimes m
\end{align*}
が $m \in M$ および $g_0, \ldots, g_k \in S$ に対して成り立つ。
従って、補題の主張にある規則
$m \mapsto \text{class of }1 \otimes m$ は階数 $k$ の微分作用素である
（補題 \ref{algebra-lemma-module-principal-parts} の証明を参照）。
$P^k_{S/R}(M)$ の普遍性により、補題の主張にある矢印を得る。
他方、$D : M \to N$ が階数 $k$ の微分作用素ならば、
$S$-線形写像 $L : S \otimes_R M \to N$,
$g \otimes m \mapsto gD(m)$ を考えることができる。
上と同様の計算と、$J$ が形
$g \otimes 1 - 1 \otimes g$ の元によって生成されることを用いると、
$L$ が $J^{k + 1}(S \otimes_R M)$ を零に写すことを読者は確認できる。
特に、これを普遍微分作用素
$D_{univ} : M \to P^k_{S/R}(M)$ に適用すれば、$S$-線形写像
$(S \otimes_R M)/J^{k + 1}(S \otimes_R M) \to P^k_{S/R}(M)$ が得られる。
これが補題の主張の写像の逆であることの確認は読者に任せる。
\end{proof}

\begin{lemma}
\label{algebra-lemma-differentials-de-rham-complex-order-1}
$A \to B$ を環準同型とする。微分
$\text{d} : \Omega^i_{B/A}  \to \Omega^{i + 1}_{B/A}$ は
階数 $1$ の微分作用素である。
\end{lemma}

\begin{proof}
$b \in B$ が与えられたとき、
$\text{d} \circ b - b \circ \text{d}$ が線形作用素であることを
示さなければならない。従って、
% [P01-E0723] Japanese supplies terminal punctuation after the frozen display.
$$
\text{d} \circ b \circ b' - b \circ \text{d} \circ b' -
b' \circ \text{d} \circ b + b' \circ b \circ \text{d} = 0
$$
を示せばよい。これを確かめるには $\Omega^i_{B/A}$ の加法的生成元上で
確認すれば十分である。従って、
$$
\text{d}(bb'b_0\text{d}b_1 \wedge \ldots \wedge \text{d}b_i) -
b\text{d}(b'b_0\text{d}b_1 \wedge \ldots \wedge \text{d}b_i) -
b'\text{d}(bb_0\text{d}b_1 \wedge \ldots \wedge \text{d}b_i) +
bb'\text{d}(b_0\text{d}b_1 \wedge \ldots \wedge \text{d}b_i)
$$
が零であることを示せばよい。これは Leibniz 則を用いる心地よい計算であり、
読者に任せる。
\end{proof}

\begin{lemma}
\label{algebra-lemma-check-differential-operators}
$A \to B$ を環準同型とし、$g_i \in B$, $i \in I$ を
$A$-代数としての $B$ の生成元の集合とする。$M, N$ を $B$-加群とし、
$D : M \to N$ を $A$-線形写像とする。
% [P01-E0724] Frozen scope without k >= 1 is retained verbatim.
$D$ が階数 $k$ の微分作用素であることを示すには、$i \in I$ に対して
$D \circ g_i - g_i \circ D$ が階数 $k - 1$ の微分作用素であることを
示せば十分である。
\end{lemma}

\begin{proof}
実際、$D \circ g - g \circ D$ が階数 $k - 1$ の微分作用素となる
元 $g \in B$ の集合は $B$ の $A$-部分代数であると主張する。
これは関係式
$$
D \circ (g + g') - (g + g') \circ D =
(D \circ g - g \circ D) + (D \circ g' - g' \circ D)
$$
および
% [P01-E0725] Japanese supplies terminal punctuation after the frozen display.
$$
D \circ gg' - gg' \circ D =
(D \circ g - g \circ D) \circ g' + g \circ (D \circ g' - g' \circ D)
$$
から従う。厳密にいえば、積について結論するためには補題
\ref{algebra-lemma-composition-differential-operators} も用いる。
\end{proof}

\begin{lemma}
\label{algebra-lemma-invert-system-differential-operators}
$A \to B$ を環準同型とし、$M, N$ を $B$-加群とする。
$S \subset B$ を乗法的部分集合とする。
任意の階数 $k$ の微分作用素 $D : M \to N$ は、一意に階数 $k$ の
微分作用素 $E : S^{-1}M \to S^{-1}N$ に拡張される。
\end{lemma}

\begin{proof}
$k$ に関する帰納法による。$k = 0$ ならば $D$ は $B$-線形であるから、
局所化の関手性により拡張が得られる。$b \in B$ が与えられたとき、
作用素 $L_b : m \mapsto D(bm) - bD(m)$ の階数は $k - 1$ である。
従って帰納法により、これは一意な階数 $k - 1$ の微分作用素
$E_b : S^{-1}M \to S^{-1}N$ に拡張される。
さらに、計算により $L_{b'b} = L_{b'} \circ b + b' \circ L_b$ が分かる。
% [P01-E0728] Japanese supplies the inference boundary.
従って一意性から
$E_{b'b} = E_{b'} \circ b + b' \circ E_b$ を得る。
同様に
$E_{b'} \circ b - b \circ E_{b'} =
E_b \circ b' - b' \circ E_b$ を得る。
ここで $m \in M$ および $g \in S$ に対して
% [P01-E0729] Japanese supplies terminal punctuation after the frozen display.
$$
E(m/g) = (1/g)(D(m) - E_g(m/g))
$$
とおく。これが well-defined であることを示すには、$g' \in S$ に対して
代表元 $g'm/g'g$ を用いても同じ結果を得ることを示せば十分である。計算すると
\begin{align*}
(1/g'g)(D(g'm) - E_{g'g}(g'm/gg'))
& =
(1/gg')(g'D(m) + E_{g'}(m) - E_{g'g}(g'm/gg')) \\
& =
(1/g'g)(g'D(m) - g' E_g(m/g))
\end{align*}
% [P01-E0730] Japanese supplies the clause boundary after the frozen display.
となり、これは先ほどと同じである。
% [P01-E0726] Frozen undefined base-ring symbol R is retained verbatim.
$D$ と $E_g$ は $R$-線形なので、$E$ が $R$-線形であることは明らかである。
$g = 1$ とし、$E_1 = 0$ を用いると、$E$ が $D$ を拡張することが分かる。
補題 \ref{algebra-lemma-check-differential-operators} により、
$E$ が階数 $k$ の微分作用素であることを示すには、
$b \in B$ に対する $E \circ b - b \circ E$ と、
$g' \in S$ に対する $E \circ 1/g' - 1/g' \circ E$ が
階数 $k - 1$ の微分作用素であることを示せば十分である。
前者について、$S^{-1}M$ の元 $m/g$ を選び、
\begin{align*}
E(b m/g) - bE(m/g)
& =
(1/g)(D(bm) - bD(m) - E_g(bm/g) + bE_g(m/g)) \\
& =
(1/g)(L_b(m) - E_b(m) + gE_b(m/g)) \\
& =
E_b(m/g)
\end{align*}
% [P01-E0731] Japanese supplies the clause boundary after the frozen display.
であることに注意する。これは階数 $k - 1$ の微分作用素である。
最後に、
\begin{align*}
E(m/g'g) - (1/g')E(m/g)
& =
(1/g'g)(D(m) - E_{g'g}(m/g'g)) -
(1/g'g)(D(m) - E_g(m/g)) \\
& =
-(1/g')E_{g'}(m/g'g)
\end{align*}
% [P01-E0732] Japanese supplies the clause boundary after the frozen display.
である。これも、線形写像（$1/g'$ による乗法および符号）と
$E_{g'}$ の合成として、階数 $k - 1$ の微分作用素である。
一意性の証明は省略する。
\end{proof}

\begin{lemma}
\label{algebra-lemma-base-change-differential-operators}
$R \to A$ と $R \to B$ を環準同型とし、$M$ と $M'$ を $A$-加群とする。
$D : M \to M'$ を $R \to A$ に関する階数 $k$ の微分作用素とする。
$N$ を任意の $B$-加群とする。このとき写像
$$
D \otimes \text{id}_N : M \otimes_R N \to M' \otimes_R N
$$
は $B \to A \otimes_R B$ に関する階数 $k$ の微分作用素である。
\end{lemma}

\begin{proof}
$D' = D \otimes \text{id}_N$ が $B$-線形であることは明らかである。
% [P01-E0727] Frozen induction omits the k = 0 base case; no silent repair.
補題 \ref{algebra-lemma-check-differential-operators} により、
$$
D' \circ a \otimes 1 -  a \otimes 1 \circ D' =
(D \circ a - a \circ D) \otimes \text{id}_N
$$
が階数 $k - 1$ の微分作用素であることを示せば十分であり、
これは $k$ に関する帰納法から従う。
\end{proof}




\section{素朴な余接複体}
\label{algebra-section-netherlander}

\noindent
$R \to S$ を環準同型とする。元 $s \in S$ を変数とする多項式環を
% [P01-E0733] Japanese uses a complete designation construction.
$R[S]$ と書く。$s \in S$ に対応する変数を $[s] \in R[S]$ と書こう。
従って $R[S]$ は、$S$ の元からなるすべての非順序列
$s_1, \ldots, s_n$ を走らせたときの基底元
% [P01-E0734] Japanese has no number-agreement defect.
$[s_1] \ldots [s_n]$ 上の自由 $R$-加群である。
標準的全射
\begin{equation}
\label{algebra-equation-canonical-presentation}
R[S] \longrightarrow S,\quad [s] \longmapsto s
\end{equation}
があり、その核を $I \subset R[S]$ と書く。$I$ が元
$[s + s'] - [s] - [s']$, $[s][s'] - [ss']$, $[r] - r$
によって生成されることは簡単に分かる。
補題 \ref{algebra-lemma-differential-seq} により、標準的写像
\begin{equation}
\label{algebra-equation-naive-cotangent-complex}
I/I^2 \longrightarrow \Omega_{R[S]/R} \otimes_{R[S]} S
\end{equation}
が存在し、その余核は標準的に $\Omega_{S/R}$ と同型である。$S$-加群
$\Omega_{R[S]/R} \otimes_{R[S]} S$ は生成元 $\text{d}[s]$ 上自由である。

\begin{definition}
\label{algebra-definition-naive-cotangent-complex}
$R \to S$ を環準同型とする。{\it 素朴な余接複体}
$\NL_{S/R}$ とは、鎖複体 (\ref{algebra-equation-naive-cotangent-complex})
$$
\NL_{S/R} = \left(I/I^2 \longrightarrow \Omega_{R[S]/R} \otimes_{R[S]} S\right)
$$
であって、$I/I^2$ を（ホモロジー）次数 $1$ に置き、
$\Omega_{R[S]/R} \otimes_{R[S]} S$ を次数 $0$ に置いたものである。
% [P01-E0735] Japanese states the notation relation explicitly.
次数 $1$ のホモロジーを
$H_1(L_{S/R}) = H_1(\NL_{S/R})$\footnote{文献ではこの加群を
$\Gamma_{S/R}$ と書くこともある。}と書く。
\end{definition}

\noindent
先へ進む前に、実際の余接複体について少し述べよう
（Cotangent, Section \ref{cotangent-section-cotangent-ring-map}）。
環準同型 $R \to S$ が与えられると、各項が多項式代数であり、
標準的ホモトピー同値
% [P01-E0736] Japanese supplies terminal punctuation after the frozen display.
$$
P_\bullet \longrightarrow S
$$
を備えた標準的単体的 $R$-代数 $P_\bullet$ が存在する。
$R$ 上の $S$ の余接複体 $L_{S/R}$ は、単体的加群
% [P01-E0737] Japanese supplies terminal punctuation after the frozen display.
$$
\Omega_{P_\bullet/R} \otimes_{P_\bullet} S
$$
に付随する鎖複体として定義される。
上で定義した素朴な余接複体は切断
$\tau_{\leq 1}L_{S/R}$ と標準的に同型である
（Homology, Section \ref{homology-section-truncations} および
Cotangent, Section \ref{cotangent-section-surjections} を参照）。
特に、実際に $H_1(\NL_{S/R}) = H_1(L_{S/R})$ であり、
我々の定義は余接複体を用いる定義と両立する。
さらに、上で見たように
$H_0(L_{S/R}) = H_0(\NL_{S/R}) = \Omega_{S/R}$ である。

\medskip\noindent
$R \to S$ を環準同型とする。{\it $R$ 上の $S$ の表示}とは、
$P$ が（ある変数集合上の）多項式代数であるような
$R$-代数の全射 $\alpha : P \to S$ のことである。
$S$ が $R$ 上有限型である場合、しばしば
「$R[x_1, \ldots, x_n] \to S$ を $S/R$ の表示とする」と述べる。
表示の核が $I$ であることも示したいときには、
「$0 \to I \to R[x_1, \ldots, x_n] \to S \to 0$ を
$S/R$ の表示とする」と述べる。
素朴な余接複体の定義に用いた写像 $R[S] \to S$ は表示の一例である。

\medskip\noindent
任意の表示 $\alpha$ から、$S$-加群の二項鎖複体
% [P01-E0738] Japanese uses the established 二項 compound.
$$
\NL(\alpha) : I/I^2 \longrightarrow \Omega_{P/R} \otimes_P S.
$$
が得られることに注意する。
% [P01-E0739] Frozen prose formulas without tensor bases are retained verbatim.
ここで項 $I/I^2$ は次数 $1$ に、項
$\Omega_{P/R} \otimes S$ は次数 $0$ に置く。
$I/I^2$ における $f \in I$ の類は
$\Omega_{P/R} \otimes S$ における $\text{d}f \otimes 1$ に写る。
この複体の余核は標準的に $\Omega_{S/R}$ である
（補題 \ref{algebra-lemma-differential-seq} を参照）。
複体 $\NL(\alpha)$ を、$S/R$ の表示
$\alpha : P \to S$ に付随する{\it 素朴な余接複体}という。
$P = R[S]$ であり、$S$ への標準的全射を備えているならば、
$\NL_{S/R}$ が回復される。
% [P01-E0740] Japanese supplies the omitted subject.
$P = R[x_1, \ldots, x_n]$ ならば、この複体を
$I/I^2 \to \bigoplus_{i = 1, \ldots, n} S\text{d}x_i$
と書くこともある。

\medskip\noindent
環の可換図式
\begin{equation}
\label{algebra-equation-functoriality-NL}
\vcenter{
\xymatrix{
S \ar[r]_{\phi} & S' \\
R \ar[r] \ar[u] & R' \ar[u]
}
}
\end{equation}
が与えられているとする。$\alpha : P \to S$ を $R$ 上の $S$ の表示、
$\alpha' : P' \to S'$ を $R'$ 上の $S'$ の表示とする。
{\it 表示 $\alpha : P \to S$ から表示
$\alpha' : P' \to S'$ への射}とは、等式
$\phi \circ \alpha = \alpha' \circ \varphi$ を満たす
$R$-代数準同型
$$
\varphi : P \to P'
$$
のことである。このとき $\varphi(I) \subset I'$ であることに注意する。
ここで $I = \Ker(\alpha)$ および $I' = \Ker(\alpha')$ である。
従って $\varphi$ は $S$-加群の写像 $I/I^2 \to I'/(I')^2$ を誘導し、
微分の関手性により $S$-加群の写像
% [P01-E0741] Frozen formulas without tensor bases are retained verbatim.
$\Omega_{P/R} \otimes S \to \Omega_{P'/R'} \otimes S'$ も誘導する。
これらの写像は $\NL(\alpha)$ および $\NL(\alpha')$ の微分と両立し、
素朴な余接複体の写像
$$
\NL(\alpha) \longrightarrow \NL(\alpha').
$$
を得る。誘導写像
$\NL(\alpha) \otimes_S S' \to \NL(\alpha')$ を考えるのが便利なことも多い。

\medskip\noindent
特別な場合 $P = R[S]$ および $P' = R'[S']$ には、
写像 $\phi : S \to S'$ は規則 $[s] \mapsto [\phi(s)]$ により
標準的環準同型 $\varphi : P \to P'$ を誘導する。
従って上の構成は、図式 (\ref{algebra-equation-functoriality-NL}) に付随する
鎖複体の標準的（！）写像
$$
\NL_{S/R} \longrightarrow \NL_{S'/R'},\quad\text{and}\quad
\NL_{S/R} \otimes_S S' \longrightarrow \NL_{S'/R'}
$$
を定める。この構成は合成と両立することに注意する。すなわち、可換図式
% [P01-E0742] Japanese supplies the continuing clause boundary.
$$
\xymatrix{
S \ar[r]_{\phi} & S' \ar[r]_{\phi'} & S'' \\
R \ar[r] \ar[u] & R' \ar[u] \ar[r] & R'' \ar[u]
}
$$
が与えられるとき、写像
$$
\NL_{S/R} \longrightarrow \NL_{S'/R'} \longrightarrow \NL_{S''/R''}
$$
の合成は、外側の正方形が与える写像
$\NL_{S/R} \to \NL_{S''/R''}$ である。

\medskip\noindent
実は $\NL(\alpha)$ は $\NL_{S/R}$ とホモトピー同値であり、
上で構成した写像はホモトピーを除いて well-defined である
（複体の写像のホモトピーは
Homology, Section \ref{homology-section-complexes} で論じられているが、
次の補題の主張の正確な意味もその証明中に詳述する）。

\begin{lemma}
\label{algebra-lemma-NL-homotopy}
図式 (\ref{algebra-equation-functoriality-NL}) が与えられているとする。
$\alpha : P \to S$ および $\alpha' : P' \to S'$ を表示とする。
\begin{enumerate}
\item $\alpha$ から $\alpha'$ への表示の射が存在する。
\item 任意の二つの表示の射は、ホモトピーな複体の射
$\NL(\alpha) \to \NL(\alpha')$ を誘導する。
\item この構成は表示の射の合成と両立する
（正確な主張については証明を参照）。
\item $R \to R'$ および $S \to S'$ が同型ならば、
$\alpha$ から $\alpha'$ への任意の表示の写像 $\varphi$ が誘導する写像
$\NL(\alpha) \to \NL(\alpha')$ はホモトピー同値かつ擬同型である。
\end{enumerate}
特に、$\alpha$ と標準表示 (\ref{algebra-equation-canonical-presentation}) を
比較すると、ホモトピーを除いて well-defined であり、すべての関手性と
（ホモトピーを除いて）両立する擬同型
$\NL(\alpha) \to \NL_{S/R}$ が存在することが分かる。
\end{lemma}

\begin{proof}
$P$ は $R$ 上の多項式代数なので、ある集合 $A$ に対して
$P = R[x_a, a \in A]$ と書ける。
$\alpha'$ は全射だから、すべての $a \in A$ に対して
$\alpha'(f_a) = \phi(\alpha(x_a))$ を満たす元 $f_a \in P'$ を選べる。
$\varphi(x_a) = f_a$ を満たす一意な $R$-代数準同型
$\varphi : P = R[x_a, a \in A] \to P'$ をとる。
これが (1) の射を与える。

\medskip\noindent
% [P01-E0743] Japanese supplies the omitted imperative copula.
$\varphi$ と $\varphi'$ を $\alpha$ から $\alpha'$ への表示の射とする。
$I = \Ker(\alpha)$ および $I' = \Ker(\alpha')$ とおく。
図式
$$
\xymatrix{
I/I^2 \ar[r]^-{\text{d}}
\ar@<1ex>[d]^{\varphi'_1} \ar@<-1ex>[d]_{\varphi_1}
&
\Omega_{P/R} \otimes_P S
\ar@<1ex>[d]^{\varphi'_0} \ar@<-1ex>[d]_{\varphi_0}
\ar[ld]_h
\\
I'/(I')^2 \ar[r]^-{\text{d}}
&
\Omega_{P'/R'} \otimes_{P'} S'
}
$$
の対角写像 $h$ を構成しなければならない。
% [P01-E0744] Japanese makes h the subject of both homotopy identities.
垂直写像は $\varphi$, $\varphi'$ によって誘導され、この対角写像は
$$
\varphi_1 - \varphi'_1 = h \circ \text{d}
\quad\text{and}\quad
\varphi_0 - \varphi'_0 = \text{d} \circ h
$$
を満たすものとする。写像 $\varphi - \varphi' : P \to P'$ を考える。
$\varphi$ と $\varphi'$ はともに $\alpha$ および $\alpha'$ と
両立するため、$\varphi - \varphi' : P \to I'$ が得られる。
これから、$\varphi, \varphi' : P \to P'$ が $I'/(I')^2$ に誘導する
$P$-加群構造は同じである。実際、
$\varphi(p)i' - \varphi'(p)i' = (\varphi - \varphi')(p)i' \in (I')^2$
だからである。また $\varphi - \varphi'$ は $R$-線形であり、
$$
(\varphi - \varphi')(fg) =
\varphi(f)(\varphi - \varphi')(g) + (\varphi - \varphi')(f)\varphi'(g)
$$
を満たす。
% [P01-E0745] Japanese is unaffected by the frozen English article error.
従って誘導写像 $D : P \to I'/(I')^2$ は $R$-導分である。
よって $D = h \circ \text{d}$ を満たす標準的写像
$h : \Omega_{P/R} \otimes_P S \to I'/(I')^2$ を得る。
計算（省略する）により、$h$ が求めるホモトピーであることが分かる。

\medskip\noindent
可換図式
$$
\xymatrix{
S \ar[r]_{\phi} & S' \ar[r]_{\phi'} & S'' \\
R \ar[r] \ar[u] & R' \ar[u] \ar[r] & R'' \ar[u]
}
$$
が与えられ、
\begin{enumerate}
\item $\alpha : P \to S$,
\item $\alpha' : P' \to S'$, and
\item $\alpha'' : P'' \to S''$
\end{enumerate}
が表示であるとする。また、
\begin{enumerate}
\item $\varphi : P \to P'$ は $\alpha$ から $\alpha'$ への表示の射であり、
\item $\varphi' : P' \to P''$
は $\alpha'$ から $\alpha''$ への表示の射である
\end{enumerate}
とする。このとき、$\varphi' \circ \varphi : P \to P''$ が
$\alpha$ から $\alpha''$ への表示の射であること、および、
それが誘導する素朴な余接複体の写像
$\NL(\alpha) \to \NL(\alpha'')$ が、$\varphi$ と $\varphi'$ の誘導する
写像 $\NL(\alpha) \to \NL(\alpha')$ および
$\NL(\alpha') \to \NL(\alpha'')$ の合成であることは直ちに分かる。

\medskip\noindent
二項複体という単純な場合、擬同型とは、項の間の写像の余核と核の
両方に同型を誘導する写像にほかならない。
% [P01-E0746] Japanese uses the established 二項 compound.
上で説明したように、二項複体のホモトピーな写像は核と余核に同じ写像を定める。
従って、$\varphi$ が $R$ 上の $S$ の表示 $\alpha$ からそれ自身への写像ならば、
誘導写像 $\NL(\alpha) \to \NL(\alpha)$ は、(2) により恒等写像と
% [P01-E0747] Japanese states the logical implication explicitly.
ホモトピーであるため、擬同型である。
(4) を一般に証明するため、(1) により存在する
$\alpha'$ から $\alpha$ への射 $\varphi'$ を考える。
% [P01-E0748] Frozen attribution to part (3) is retained without adding the missing premise.
合成 $\NL(\alpha) \to \NL(\alpha') \to \NL(\alpha)$ および
$\NL(\alpha') \to \NL(\alpha) \to \NL(\alpha')$ は (3) により
恒等写像とホモトピーである。従って定義により、これらの写像は
ホモトピー同値である。すると形式的に、両方の写像
$\NL(\alpha) \to \NL(\alpha')$ および
$\NL(\alpha') \to \NL(\alpha)$ が擬同型であることが従う。
いくつかの詳細は省略する。
\end{proof}
\begin{lemma}
\label{algebra-lemma-NL-polynomial-algebra}
$A \to B$ が多項式代数ならば、$\NL_{B/A}$ は
$\Omega_{B/A}$ を次数 $0$ に置いた鎖複体
$(0 \to \Omega_{B/A})$ とホモトピー同値である。
\end{lemma}

\begin{proof}
% [P01-E0749] Japanese supplies a complete proof sentence.
これは補題 \ref{algebra-lemma-NL-homotopy} と、
$\text{id}_B : B \to B$ が核零の $A$ 上の $B$ の表示であることから従う。
\end{proof}

\noindent
次の補題は、素朴な余接複体を導入する動機の一部である。
余接複体を用いれば、これを真のコホモロジー長完全列へ拡張できる。
$B \to C$ が局所完全交叉ならば、列の左端に零を付け加えられる
（More on Algebra, Lemma
\ref{more-algebra-lemma-transitive-lci-at-end} を参照）。

\begin{lemma}[Jacobi--Zariski 列]
\label{algebra-lemma-exact-sequence-NL}
$A \to B \to C$ を環準同型とする。核を $I$ とする表示
$\alpha : A[x_s, s \in S] \to B$ を選び、核を $J$ とする表示
$\beta : B[y_t, t \in T] \to C$ を選ぶ。
$\gamma : A[x_s, y_t] \to C$ を、$C$ の誘導された表示とし、その核を
$K$ とする。このとき標準的可換図式
% [P01-E0750] Japanese supplies the continuing exact-row boundary.
$$
\xymatrix{
0 \ar[r] &
\Omega_{A[x_s]/A} \otimes C \ar[r] &
\Omega_{A[x_s, y_t]/A} \otimes C \ar[r] &
\Omega_{B[y_t]/B} \otimes C \ar[r] &
0 \\
&
I/I^2 \otimes C \ar[r] \ar[u] &
K/K^2 \ar[r] \ar[u] &
J/J^2 \ar[r] \ar[u] &
0
}
$$
が得られ、その各行は完全である。さらに、補題
\ref{algebra-lemma-exact-sequence-differentials} の列を拡張する
$C$-加群のホモロジー群の完全列
% [P01-E0751] Japanese places the C-module qualifier before the frozen display.
$$
H_1(\NL_{B/A} \otimes_B C) \to
H_1(L_{C/A}) \to
H_1(L_{C/B}) \to
C \otimes_B \Omega_{B/A} \to
\Omega_{C/A} \to
\Omega_{C/B} \to 0
$$
が得られる。もし $\text{Tor}_1^B(\Omega_{B/A}, C) = 0$ および
$\text{Tor}_2^B(\Omega_{B/A}, C) = 0$ ならば、
$H_1(\NL_{B/A} \otimes_B C) = H_1(L_{B/A}) \otimes_B C$ である。
\end{lemma}

\begin{proof}
写像の正確な定義は省略する。上の行の完全性は、
$\text{d}x_s$, $\text{d}y_t$ が中央の加群の基底をなすことから従う。
写像 $\gamma$ は
$$
A[x_s, y_t] \to B[y_t] \to C
$$
と分解する。第一の矢印は全射であり、第二の矢印は $\beta$ に等しい。
従って $K \to J$ は全射である。さらに、最初に表示した矢印の核は
$IA[x_s, y_t]$ である。それゆえ $I/I^2 \otimes C$ は
$K/K^2 \to J/J^2$ の核へ全射する。最後に補題
\ref{algebra-lemma-NL-homotopy} を用いて、各項を素朴な余接複体の
ホモロジー群と同一視できる。

\medskip\noindent
最後の主張はホモロジー代数の命題である。
$\NL_{B/A} = (N^{-1} \to N^0)$ は、$N^0$ が自由であり、
コホモロジー加群が $H^0 = \Omega_{B/A}$ および
$H^{-1} = H_1(L_{B/A})$ である二項 $B$-加群複体だったことを思い出す。
微分の像を $M \subset N^0$ と書く。
$\text{Tor}_1^B(H^0, C) = 0$ ならば、完全列
$$
0 \to M \otimes_B C \to N^0 \otimes_B C \to H^0 \otimes_B C \to 0
$$
がある。$N^0$ は自由なので
$\text{Tor}_2^B(H^0, C) =
\text{Tor}_1^B(M, C)$ でもある。従って
$\text{Tor}_2^B(H^0, C) = 0$ ならば、完全列
$$
0 \to H^{-1} \otimes_B C \to N^{-1} \otimes_B C \to M \otimes_B C \to 0
$$
も得られる。以上をまとめると、
$\text{Tor}_1^B(H^0, C) = 0$ および $\text{Tor}_2^B(H^0, C) = 0$
ならば、望んだとおり $H^{-1} \otimes_B C$ は
$N^{-1} \otimes_B C \to N^0 \otimes_B C$ の核である。
\end{proof}

\begin{remark}
\label{algebra-remark-composition-homotopy-equivalent-to-zero}
$A \to B$ および $\phi : B \to C$ を環準同型とする。
このとき合成 $\NL_{B/A} \to \NL_{C/A} \to \NL_{C/B}$ は
零とホモトピー同値である。実際、この合成は正方形
% [P01-E0752] Japanese supplies terminal punctuation after the frozen diagram.
$$
\xymatrix{
B \ar[r]_\phi & C \\
A \ar[r] \ar[u] & B \ar[u]
}
$$
に対する素朴な余接複体の関手性である。
$J = \Ker(B[C] \to C)$ と書く。明示的なホモトピーは、
基底元 $\text{d}[b]$ を $J/J^2$ における
$[\phi(b)] - b$ の類に写す写像
$\Omega_{A[B]/A} \otimes_{A[B]} B \to J/J^2$ によって与えられる。
\end{remark}

\begin{lemma}
\label{algebra-lemma-NL-surjection}
$A \to B$ を核 $I$ の全射環準同型とする。
このとき $\NL_{B/A}$ は、$I/I^2$ を次数 $1$ に置いた鎖複体
$(I/I^2 \to 0)$ とホモトピー同値である。
特に $H_1(L_{B/A}) = I/I^2$ である。
\end{lemma}

\begin{proof}
% [P01-E0753] Japanese supplies a complete proof sentence.
これは補題 \ref{algebra-lemma-NL-homotopy} と、
$A \to B$ が $A$ 上の $B$ の表示であることから従う。
\end{proof}

\begin{lemma}
\label{algebra-lemma-application-NL}
$A \to B \to C$ を環準同型とし、$A \to C$ が全射であると仮定する
（従って $B \to C$ も全射である）。
$I = \Ker(A \to C)$ および $J = \Ker(B \to C)$ と書く。
このとき列
$$
I/I^2 \to J/J^2 \to \Omega_{B/A} \otimes_B B/J \to 0
$$
は完全である。
\end{lemma}

\begin{proof}
% [P01-E0754] Japanese supplies a complete proof sentence.
これは補題 \ref{algebra-lemma-exact-sequence-NL} と、補題
\ref{algebra-lemma-NL-surjection} における素朴な余接複体
$\NL_{C/B}$ および $\NL_{C/A}$ の記述から従う。
\end{proof}

\begin{lemma}[平坦基底変換]
\label{algebra-lemma-change-base-NL}
$R \to S$ を環準同型とし、$\alpha : P \to S$ を表示とする。
$R \to R'$ を平坦環準同型とする。
$\alpha' : P \otimes_R R' \to S' = S \otimes_R R'$ を誘導された表示とする。
このとき
$\NL(\alpha) \otimes_R R' = \NL(\alpha) \otimes_S S' = \NL(\alpha')$ である。
特に、$R \to R'$ が平坦ならば、標準的写像
$$
\NL_{S/R} \otimes_S S' \longrightarrow \NL_{S \otimes_R R'/R'}
$$
はホモトピー同値である。
\end{lemma}

\begin{proof}
これは、$R \to R'$ が平坦であるため
$\Ker(\alpha') = R' \otimes_R \Ker(\alpha)$ となることから分かる。
\end{proof}

\begin{lemma}
\label{algebra-lemma-colimits-NL}
$R_\lambda \to S_\lambda$ を、有向集合 $\Lambda$ 上の環準同型の系とする。
$R = \colim R_\lambda$ および $S = \colim S_\lambda$ とおく。
このとき $\NL_{S/R} = \colim \NL_{S_\lambda/R_\lambda}$ である。
\end{lemma}

\begin{proof}
$\NL_{S/R}$ は複体
$I/I^2 \to \bigoplus_{s \in S} S\text{d}[s]$ であったことを思い出す。
ここで $I \subset R[S]$ は標準表示 $R[S] \to S$ の核である。
いま $R[S] = \colim R_\lambda[S_\lambda]$ であり、同様に
$I = \colim I_\lambda$ であることは明らかである。ここで
$I_\lambda = \Ker(R_\lambda[S_\lambda] \to S_\lambda)$ である。
従って補題は明らかである。
\end{proof}

\begin{lemma}
\label{algebra-lemma-NL-of-localization}
$S \subset A$ が $A$ の乗法的部分集合ならば、
$\NL_{S^{-1}A/A}$ は零複体とホモトピー同値である。
\end{lemma}

\begin{proof}
$A \to S^{-1}A$ は平坦なので、平坦基底変換
（補題 \ref{algebra-lemma-change-base-NL}）により
$\NL_{S^{-1}A/A} \otimes_A S^{-1}A \to \NL_{S^{-1}A/S^{-1}A}$
はホモトピー同値である。矢印の始域は $\NL_{S^{-1}A/A}$ と同型であり、
終域は零（補題 \ref{algebra-lemma-NL-surjection} による）なので、結論を得る。
\end{proof}

\begin{lemma}
\label{algebra-lemma-NL-localize-bottom}
% [P01-E0755] Japanese renders the introduction naturally.
$S \subset A$ を $A$ の乗法的部分集合とし、
$S^{-1}A \to B$ を環準同型とする。
このとき $\NL_{B/A} \to \NL_{B/S^{-1}A}$ はホモトピー同値である。
\end{lemma}

\begin{proof}
$A$ 上の $B$ の表示 $\alpha : P \to B$ を選ぶ。
すると $\beta : S^{-1}P \to B$ は $S^{-1}A$ 上の $B$ の表示である。
直接計算により $\NL(\alpha) = \NL(\beta)$ が分かる。
% [P01-E0756] Japanese states the causal step explicitly.
補題 \ref{algebra-lemma-NL-homotopy} により素朴な余接複体はホモトピーを除いて
well-defined なので、これで補題が証明される。
\end{proof}

\begin{lemma}
\label{algebra-lemma-principal-localization-NL}
\begin{slogan}
素朴な余接複体の形成は一元での局所化と可換である。
\end{slogan}
$A \to B$ を環準同型とし、$g \in B$ とする。
$\alpha : P \to B$ が核 $I$ の表示であると仮定する。
このとき、$A$ 上の $B_g$ の表示は、$\alpha$ を拡張し
$x$ を $1/g$ に写す写像
$$
\beta : P[x] \longrightarrow B_g
$$
である。$\beta$ の核 $J$ は、$I$ と元 $f x - 1$ によって生成される。
ここで $f \in P$ は $\alpha$ によって $g \in B$ に写る元である。
この状況では
\begin{enumerate}
\item $J/J^2 = (I/I^2)_g \oplus B_g (f x - 1)$,
\item $\Omega_{P[x]/A} \otimes_{P[x]} B_g =
\Omega_{P/A} \otimes_P B_g \oplus B_g \text{d}x$,
\item $\NL(\beta) \cong
\NL(\alpha) \otimes_B B_g \oplus (B_g \xrightarrow{g} B_g)$
\end{enumerate}
% [P01-E0757] Japanese supplies terminal punctuation after the frozen enumeration.
従って標準的写像 $\NL_{B/A} \otimes_B B_g \to \NL_{B_g/A}$ は
ホモトピー同値である。
\end{lemma}

\begin{proof}
$P[x]/(I, fx - 1) = B[x]/(gx - 1) = B_g$ なので、
$I$ と $fx - 1$ が $J$ を生成するという主張を得る。
可換図式
% [P01-E0758] Frozen formulas without tensor bases are retained verbatim.
$$
\xymatrix{
0 \ar[r] &
\Omega_{P/A} \otimes B_g \ar[r] &
\Omega_{P[x]/A} \otimes B_g \ar[r] &
\Omega_{B[x]/B} \otimes B_g \ar[r] &
0 \\
&
(I/I^2)_g \ar[r] \ar[u] &
J/J^2 \ar[r] \ar[u] &
(gx - 1)/(gx - 1)^2 \ar[r] \ar[u] &
0
}
$$
を考える。補題 \ref{algebra-lemma-exact-sequence-NL} により各行は完全である。
$B_g$-加群 $\Omega_{B[x]/B} \otimes B_g$ は
$\text{d}x$ 上階数 $1$ の自由加群である。
$B_g$-加群 $\Omega_{P[x]/A} \otimes B_g$ の元 $\text{d}x$ は
上の行の分裂を与える。元 $gx - 1 \in (gx - 1)/(gx - 1)^2$ は
$\Omega_{B[x]/B} \otimes B_g$ において $g\text{d}x$ に写る。
従って $(gx - 1)/(gx - 1)^2$ は $B_g$ 上階数 $1$ の自由加群である。
（これは次のように見ることもできる。$gx - 1$ は可逆な定数項をもつ
多項式なので $B[x]$ の非零因子であり、任意の非零因子は補題
\ref{algebra-lemma-regular-quasi-regular} により長さ $1$ の擬正則列を与える。）

\medskip\noindent
% [P01-E0759] Japanese supplies the omitted copula.
$(I/I^2)_g \to J/J^2$ が単射であることを証明しよう。
$x$ を $1/f$ に写す $P$-代数準同型
$$
\pi : P[x] \to (P/I^2)_f = P_f/I_f^2
$$
を考える。$J$ は $I$ と $fx - 1$ によって生成されるので、
$\pi(J) \subset (I/I^2)_f = (I/I^2)_g$ である。
これは平方零イデアルだから $\pi(J^2) = 0$ である。
% [P01-E0760] Japanese states the intended membership unambiguously.
$a \in I$ の $J$ における像が $J^2$ に属するならば、
$\pi(a) = 0$ であり、従って $a$ は $I_f/I_f^2$ で零に写る。
これで望む単射性が証明される。

\medskip\noindent
従って二項複体の短完全列
% [P01-E0761] Japanese uses the established 二項 compound.
$$
0 \to \NL(\alpha) \otimes_B B_g \to \NL(\beta)
\to (B_g \xrightarrow{g} B_g) \to 0
$$
がある。
% [P01-E0762] Frozen false blanket splitting claim is retained verbatim in sense.
このような短完全列は、複体の圏において常に分裂できる。
この場合には具体的な分裂として
% [P01-E0763] Japanese supplies terminal punctuation after the frozen display.
$$
J/J^2 = (I/I^2)_g \oplus B_g (fx - 1)\quad\text{and}\quad
\Omega_{P[x]/A} \otimes B_g = \Omega_{P/A} \otimes B_g \oplus
B_g (g^{-2}\text{d}f + \text{d}x)
$$
をとることができる。これは
$\text{d}(fx - 1) = x\text{d}f + f \text{d}x =
g(g^{-2}\text{d}f + \text{d}x)$
が $\Omega_{P[x]/A} \otimes B_g$ において成り立つからである。
\end{proof}
\begin{lemma}
\label{algebra-lemma-localize-NL}
$A \to B$ を環準同型とし、$S \subset B$ を乗法的部分集合とする。
標準的写像 $\NL_{B/A} \otimes_B S^{-1}B \to \NL_{S^{-1}B/A}$ は
擬同型である。
\end{lemma}

\begin{proof}
$S$ を（整除関係で順序づけた）有向集合とみなせば、
$S^{-1}B = \colim_{g \in S} B_g$ である
（補題 \ref{algebra-lemma-localization-colimit} を参照）。
% [P01-E0764] Japanese has no subject--verb agreement defect.
補題 \ref{algebra-lemma-principal-localization-NL} により、各写像
$\NL_{B/A} \otimes_B B_g \to \NL_{B_g/A}$ は擬同型である。
補題 \ref{algebra-lemma-colimits-NL} から補題が従う。
\end{proof}

\begin{lemma}
\label{algebra-lemma-sum-two-terms}
$R$ を環とする。
% [P01-E0765] Japanese renders the coordination and 二項 compound naturally.
$A_1 \to A_0$ および $B_1 \to B_0$ を二項複体とする。
複体の射 $\varphi : A_\bullet \to B_\bullet$ および
$\psi : B_\bullet \to A_\bullet$ が存在し、
$\varphi \circ \psi$ と $\psi \circ \varphi$ は恒等写像とホモトピーであるとする。
このとき $R$-加群として
$A_1 \oplus B_0 \cong B_1 \oplus A_0$ である。
\end{lemma}

\begin{proof}
写像 $h : A_0 \to A_1$ で
$$
\text{id}_{A_1} - \psi_1 \circ \varphi_1 = h \circ d_A
\text{ and }
\text{id}_{A_0} - \psi_0 \circ \varphi_0 = d_A \circ h.
$$
を満たすものを選ぶ。同様に、写像 $h' : B_0 \to B_1$ で
$$
\text{id}_{B_1} - \varphi_1 \circ \psi_1 = h' \circ d_B
\text{ and }
\text{id}_{B_0} - \varphi_0 \circ \psi_0 = d_B \circ h'.
$$
を満たすものを選ぶ。自明な計算により
% [P01-E0766] Japanese supplies terminal punctuation after the frozen display.
$$
\left(
\begin{matrix}
\text{id}_{A_1} & -h' \circ \psi_1 + h \circ \psi_0 \\
0 & \text{id}_{B_0}
\end{matrix}
\right)
=
\left(
\begin{matrix}
\psi_1 & h \\
-d_B & \varphi_0
\end{matrix}
\right)
\left(
\begin{matrix}
\varphi_1 & - h' \\
d_A & \psi_0
\end{matrix}
\right)
$$
が分かる。
% [P01-E0767] Frozen inference from one product to both factors is retained.
これにより、右辺の両方の行列が可逆であることが分かり、補題が証明される。
\end{proof}

\begin{lemma}
\label{algebra-lemma-conormal-module}
$R \to S$ を有限型環準同型とする。
% [P01-E0768] Japanese renders the coordinated presentations naturally.
任意の表示 $\alpha : R[x_1, \ldots, x_n] \to S$ および
$\beta : R[y_1, \ldots, y_m] \to S$ に対して
$$
I/I^2 \oplus S^{\oplus m} \cong J/J^2 \oplus S^{\oplus n}
$$
が $S$-加群として成り立つ。
% [P01-E0769] Japanese separates the supplementary definitions.
ここで $I = \Ker(\alpha)$ および $J = \Ker(\beta)$ である。
\end{lemma}

\begin{proof}
補題 \ref{algebra-lemma-NL-homotopy} および
\ref{algebra-lemma-sum-two-terms} を参照。
\end{proof}

\begin{lemma}
\label{algebra-lemma-conormal-module-localize}
$R \to S$ を有限型環準同型とし、$g \in S$ とする。
% [P01-E0770] Japanese renders the coordinated presentations naturally.
任意の表示 $\alpha : R[x_1, \ldots, x_n] \to S$ および
$\beta : R[y_1, \ldots, y_m] \to S_g$ に対して
$$
(I/I^2)_g \oplus S^{\oplus m}_g \cong J/J^2 \oplus S_g^{\oplus n}
$$
が $S_g$-加群として成り立つ。
% [P01-E0771] Japanese separates the supplementary definitions.
ここで $I = \Ker(\alpha)$ および $J = \Ker(\beta)$ である。
\end{lemma}

\begin{proof}
$\beta' : R[x_1, \ldots, x_n, x] \to S_g$ を、
$\alpha$ から出発して補題 \ref{algebra-lemma-principal-localization-NL} で
構成した表示とする。このとき
$\NL(\alpha) \otimes_S S_g$ は $\NL(\beta')$ とホモトピー同値である。
補題 \ref{algebra-lemma-NL-homotopy} により
$\NL(\beta)$ と $\NL(\beta')$ はホモトピー同値である。
従って $\NL(\alpha) \otimes_S S_g$ は $\NL(\beta)$ と
ホモトピー同値である。最後に補題 \ref{algebra-lemma-conormal-module} を適用する。
\end{proof}

\section{局所完全交叉}
\label{algebra-section-lci}

\noindent
局所完全交叉であるという性質は、ネーター局所環の内在的性質である。
これについては Divided Power Algebra、第 \ref{dpa-section-lci} 節で論じる。
しかし当面は、この性質を体上の有限型代数に対してのみ定義する。

\begin{definition}
\label{algebra-definition-lci-field}
$k$ を体とし、$S$ を有限型 $k$-代数とする。
\begin{enumerate}
\item $\dim(S) = n - c$ を満たす表示
$S = k[x_1, \ldots, x_n]/(f_1, \ldots, f_c)$ が存在するとき、
$S$ は {\it $k$ 上大域的完全交叉}であるという。
\item 各環 $S_{g_i}$ が $k$ 上大域的完全交叉となる被覆
$\Spec(S) = \bigcup D(g_i)$ が存在するとき、$S$ は
{\it $k$ 上局所完全交叉}であるという。
\end{enumerate}
零環も $k$ 上大域的完全交叉であるという規約を用いる。
\end{definition}

\noindent
定義 \ref{algebra-definition-lci-field} のように、$S$ が大域的完全交叉
$S = k[x_1, \ldots, x_n]/(f_1, \ldots, f_c)$ であると仮定する。
極大イデアル $\mathfrak m \subset k[x_1, \ldots, x_n]$ に対して
$\dim(k[x_1, \ldots, x_n]_\mathfrak m) = n$ である
（補題 \ref{algebra-lemma-dim-affine-space}）。
$(f_1, \ldots, f_c) \subset \mathfrak m$ ならば、補題
\ref{algebra-lemma-one-equation} により $\dim(S_\mathfrak m) \geq n - c$ である。
% P01-E0772: the frozen paragraph indexes S by a maximal ideal of the polynomial ring without first replacing it by the corresponding maximal ideal of S.
定義 \ref{algebra-definition-lci-field} より $\dim(S) = n - c$ なので、
$S$ のすべての極大イデアルに対して
$\dim(S_\mathfrak m) = n - c$ であり、$\Spec(S)$ は次元 $n - c$ の
等次元空間であると分かる
（Topology、定義 \ref{topology-definition-equidimensional} および
補題 \ref{algebra-lemma-dimension-at-a-point-finite-type-over-field} を参照せよ）。
以下ではこのことを断りなく用いることが多い。

\begin{lemma}
\label{algebra-lemma-localize-lci}
$k$ を体、$S$ を有限型 $k$-代数、$g \in S$ とする。
\begin{enumerate}
\item $S$ が大域的完全交叉ならば、$S_g$ も大域的完全交叉である。
\item $S$ が局所完全交叉ならば、$S_g$ も局所完全交叉である。
\end{enumerate}
\end{lemma}

\begin{proof}
第二の主張は第一の主張から直ちに従う。
% P01-E0773: the frozen “Proof of the first statement.” is a sentence fragment.
第一の主張を示す。$S_g$ が零環ならば主張は成り立つので、
$S_g$ は零環でないと仮定する。定義 \ref{algebra-definition-lci-field} のように
$n - c = \dim(S)$ を満たす
$S = k[x_1, \ldots, x_n]/(f_1, \ldots, f_c)$ と書く。
定義直後の注意により $\dim(S_g) = n - c$ である。
$g' \in k[x_1, \ldots, x_n]$ を、その剰余類が $g$ に対応する元とする。このとき
$S_g =  k[x_1, \ldots, x_n, x_{n + 1}]/(f_1, \ldots, f_c, x_{n + 1}g' - 1)$
となり、主張を得る。
\end{proof}

\begin{lemma}
\label{algebra-lemma-lci-CM}
$k$ を体、$S$ を有限型 $k$-代数とする。
$S$ が局所完全交叉ならば、$S$ は Cohen--Macaulay 環である。
\end{lemma}

\begin{proof}
% P01-E0774: the frozen proof says “maximal prime” where the local criterion uses a maximal ideal.
$S$ の極大素イデアル $\mathfrak m$ を一つとる。
$S_\mathfrak m$ が Cohen--Macaulay であることを示せばよい。
仮定により、$\dim(S) = n - c$ を満たす
$S = k[x_1, \ldots, x_n]/(f_1, \ldots, f_c)$ としてよい。
$\mathfrak m' \subset k[x_1, \ldots, x_n]$ を $\mathfrak m$ に対応する
極大イデアルとする。命題 \ref{algebra-proposition-finite-gl-dim-polynomial-ring} によれば、
局所環 $k[x_1, \ldots, x_n]_{\mathfrak m'}$ は次元 $n$ の正則局所環である。
従って特に、補題 \ref{algebra-lemma-regular-ring-CM} により Cohen--Macaulay である。
補題 \ref{algebra-lemma-one-equation} を $c$ 回適用すると、局所環
$S_{\mathfrak m} = k[x_1, \ldots, x_n]_{\mathfrak m'}/(f_1, \ldots, f_c)$
は次元 $\geq n - c$ をもつ。一方、仮定より
$\dim(S_{\mathfrak m}) \leq n - c$ であるから等号が成り立つ。
従って $f_1, \ldots, f_c$ は
$k[x_1, \ldots, x_n]_{\mathfrak m'}$ の正則列であり、
$S_{\mathfrak m}$ は Cohen--Macaulay である
（命題 \ref{algebra-proposition-CM-module} を参照せよ）。
\end{proof}

\noindent
次の補題は、この節の残りの内容に対する技術的な鍵である。
重要なのは、環 $S$ の任意の表示を選べる一方、条件 (1) は
その選択に依存しないという点である。

\begin{lemma}
\label{algebra-lemma-lci}
$k$ を体、$S$ を有限型 $k$-代数、$\mathfrak q$ を $S$ の
素イデアルとする。任意の表示 $S = k[x_1, \ldots, x_n]/I$ を選び、
$\mathfrak q'$ を $\mathfrak q$ に対応する
$k[x_1, \ldots, x_n]$ の素イデアルとする。
$c = \text{height}(\mathfrak q') - \text{height}(\mathfrak q)$、
すなわち $\dim_{\mathfrak q}(S) = n - c$ とおく
（補題 \ref{algebra-lemma-codimension} を参照せよ）。次の条件は同値である。
\begin{enumerate}
\item $g \in S$, $g \not \in \mathfrak q$ であって、$S_g$ が
$k$ 上大域的完全交叉となるものが存在する。
\item イデアル $I_{\mathfrak q'} \subset k[x_1, \ldots, x_n]_{\mathfrak q'}$
は $c$ 個の元で生成できる。
\item 余法加群 $(I/I^2)_{\mathfrak q}$ は $S_{\mathfrak q}$ 上
$c$ 個の元で生成できる。
\item 余法加群 $(I/I^2)_{\mathfrak q}$ は階数 $c$ の自由
$S_{\mathfrak q}$-加群である。
\item イデアル $I_{\mathfrak q'}$ は正則局所環
$k[x_1, \ldots, x_n]_{\mathfrak q'}$ の正則列で生成できる。
\end{enumerate}
この場合、$I_{\mathfrak q'}/\mathfrak q'I_{\mathfrak q'}$ を生成する
$I_{\mathfrak q'}$ の任意の $c$ 個の元は、局所環
$k[x_1, \ldots, x_n]_{\mathfrak q'}$ の正則列をなす。
\end{lemma}

\begin{proof}
$R = k[x_1, \ldots, x_n]_{\mathfrak q'}$ とおく。これは、例えば
補題 \ref{algebra-lemma-lci-CM} により、次元
$\text{height}(\mathfrak q')$ の Cohen--Macaulay 局所環である。
さらに
$\overline{R} = R/IR = R/I_{\mathfrak q'} = S_{\mathfrak q}$ は
次元 $\text{height}(\mathfrak q)$ の商である。
$(I/I^2)_{\mathfrak q}$ を生成する元
$f_1, \ldots, f_c \in I_{\mathfrak q'}$ をとる。
補題 \ref{algebra-lemma-NAK} により、$f_1, \ldots, f_c$ は
$I_{\mathfrak q'}$ を生成する。次元がちょうど合うので、命題
\ref{algebra-proposition-CM-module} により
$f_1, \ldots, f_c$ は $R$ の正則列である。
補題 \ref{algebra-lemma-regular-quasi-regular} により
$(I/I^2)_{\mathfrak q}$ は自由である。
以上から (2), (3), (4) は同値であり、これらは補題の最後の主張を含意し、
従って (5) も含意する。

\medskip\noindent
(5) が成り立ち、$I_{\mathfrak q'}$ が長さ $e$ の正則列で
生成されるとする。このとき次元論により
$\text{height}(\mathfrak q) = \dim(S_{\mathfrak q}) =
\dim(k[x_1, \ldots, x_n]_{\mathfrak q'}) - e =
\text{height}(\mathfrak q') - e$ である
（第 \ref{algebra-section-dimension} 節を参照せよ）。従って $e = c$ であり、
(5) は (2) を含意する。

\medskip\noindent
第一段落で導入した記法を引き続き用いる。各 $f_i$ に対して、
$d_i \not \in \mathfrak q'$ かつ
$f_i' = d_i f_i \in k[x_1, \ldots, x_n]$ となる
$d_i \in k[x_1, \ldots, x_n]$ をとれる。
このとき依然として $I_{\mathfrak q'} = (f_1', \ldots, f_c')R$ である。
従って $g' \not \in \mathfrak q'$ かつ
$I_{g'} = (f_1', \ldots, f_c')$ となる
$g' \in k[x_1, \ldots, x_n]$ が存在する。
さらに $g'' \not \in \mathfrak q'$ かつ
$\dim(S_{g''}) = \dim_{\mathfrak q} \Spec(S)$ となる
$g'' \in k[x_1, \ldots, x_n]$ をとる。
補題 \ref{algebra-lemma-codimension} により、この次元は $n - c$ に等しい。
最後に $g$ を $S$ における $g'g''$ の像とする。このとき
$$
S_g \cong k[x_1, \ldots, x_n, x_{n + 1}]
/
(f_1', \ldots, f_c', x_{n + 1}g'g'' - 1)
$$
であり、$g''$ の選び方からこの環の次元は $n - c$ である。
従ってこれは大域的完全交叉である。
ゆえに (2), (3), (4) のそれぞれは (1) を含意する。

\medskip\noindent
(1) を仮定する。$S_g$ の大域的完全交叉としての表示を
$S_g \cong k[y_1, \ldots, y_m]/(f_1, \ldots, f_t)$ とする。
$J = (f_1, \ldots, f_t)$ と書き、$\mathfrak q'' \subset k[y_1, \ldots, y_m]$
を $\mathfrak qS_g$ に対応する素イデアルとする。このとき
$t = m - \dim(S_g) =
\text{height}(\mathfrak q'') - \text{height}(\mathfrak q)$
であることに注意する。最後の等号については補題
\ref{algebra-lemma-codimension} を参照せよ。
補題 \ref{algebra-lemma-lci-CM} の証明（および上の議論）で見たように、
$f_1, \ldots, f_t$ は局所環
$k[y_1, \ldots, y_m]_{\mathfrak q''}$ の正則列をなす。
補題 \ref{algebra-lemma-regular-quasi-regular} により
$(J/J^2)_{\mathfrak q}$ は階数 $t$ の自由加群である。
補題 \ref{algebra-lemma-conormal-module-localize} により
$$
J/J^2 \oplus S_g^n \cong (I/I^2)_g \oplus S_g^m
$$
である。
% P01-E0775: the frozen free-module summands use ordinary powers S_g^n and S_g^m instead of direct-sum notation.
従って $(I/I^2)_{\mathfrak q}$ は階数
$t + n - m = m - \dim(S_g) + n - m = n - \dim(S_g) =
\text{height}(\mathfrak q') - \text{height}(\mathfrak q) = c$
の自由加群である。従って (4) を得る。
\end{proof}

\noindent
補題 \ref{algebra-lemma-lci} の結果は次の定義を示唆する。

\begin{definition}
\label{algebra-definition-lci-local-ring}
$k$ を体とし、$S$ を $k$ 上本質的有限型な局所 $k$-代数とする。
局所 $k$-代数 $R$ と元 $f_1, \ldots, f_c \in \mathfrak m_R$ が存在して
次を満たすとき、$S$ は {\it （$k$ 上の）完全交叉}であるという。
\begin{enumerate}
\item $R$ は $k$ 上本質的有限型である。
\item $R$ は正則局所環である。
\item $f_1, \ldots, f_c$ は $R$ の正則列をなす。
\item $k$-代数として $S \cong R/(f_1, \ldots, f_c)$ である。
\end{enumerate}
\end{definition}

\noindent
Cohen 構造定理（定理 \ref{algebra-theorem-cohen-structure-theorem} を参照）により、
任意の完備ネーター局所環は、ある正則完備局所環の商として書ける。
従って上の定義を用いて、任意の完備ネーター局所環に対する
完全交叉環の概念を定義できる。これについては
Divided Power Algebra、第 \ref{dpa-section-lci} 節で論じる。
当面は、次の補題がそのような定義に意味があることを示す。

\begin{lemma}
\label{algebra-lemma-ci-well-defined}
$A \to B \to C$ を全射な局所環準同型とする。
$A$ と $B$ は正則局所環であると仮定する。次の条件は同値である。
% P01-E0776: the frozen sentence introducing the equivalence list lacks a colon.
\begin{enumerate}
\item $\Ker(A \to C)$ は正則列で生成される。
\item $\Ker(A \to C)$ は $\dim(A) - \dim(C)$ 個の元で生成される。
\item $\Ker(B \to C)$ は正則列で生成される。
\item $\Ker(B \to C)$ は $\dim(B) - \dim(C)$ 個の元で生成される。
\end{enumerate}
\end{lemma}

\begin{proof}
正則局所環は Cohen--Macaulay である
（補題 \ref{algebra-lemma-regular-ring-CM} を参照せよ）。
従って (1) $\Leftrightarrow$ (2) および
(3) $\Leftrightarrow$ (4) が成り立つ
（命題 \ref{algebra-proposition-CM-module} を参照せよ）。
% P01-E0777: the frozen “Hence the equivalences ...” clause has no finite verb.
補題 \ref{algebra-lemma-regular-quotient-regular} により、イデアル
$\Ker(A \to B)$ は $\dim(A) - \dim(B)$ 個の元で生成できる。
従って (4) は (2) を含意する。

\medskip\noindent
残るのは (1) が (4) を含意することの証明である。
$\dim(A) - \dim(B)$ に関する帰納法でこれを示す。
$\dim(A) - \dim(B) = 0$ の場合は自明である。
$\dim(A) > \dim (B)$ と仮定する。
$I = \Ker(A \to C)$ および $J = \Ker(A \to B)$ と書く。
$J \subset I$ であることに注意する。仮定により、$I$ の最小生成元数は
$\dim(A) - \dim(C)$ である。$\mathfrak m \subset A$ を極大イデアルとし、
写像
$$
J/ \mathfrak m J \to  I / \mathfrak m I \to \mathfrak m /\mathfrak m^2
$$
を考える。
% P01-E0778: the frozen display lacks terminal punctuation before a new sentence.
補題 \ref{algebra-lemma-regular-quotient-regular} とその証明により、この合成は単射である。
$J /\mathfrak mJ$ において零でない任意の元 $x \in J$ をとる。
上の議論と中山の補題により、$x$ は $I$ のある最小生成系の元である。
従って $A$ を $A/xA$ で、$I$ を $I/xA$ で置き換えられ、
$\dim(A)$ と $I$ の最小生成元数はいずれも $1$ だけ減少する。これで示された。
\end{proof}

\begin{lemma}
\label{algebra-lemma-lci-local}
$k$ を体とし、$S$ を $k$ 上本質的有限型な局所 $k$-代数とする。
次の条件は同値である。
\begin{enumerate}
\item $S$ は $k$ 上完全交叉である。
\item $R$ が $k$ 上本質的有限表示な正則局所環である任意の全射
$R \to S$ に対して、イデアル $\Ker(R \to S)$ は正則列で生成できる。
\item $R$ が $k$ 上本質的有限表示な正則局所環である全射
$R \to S$ が存在して、イデアル $\Ker(R \to S)$ は
$\dim(R) - \dim(S)$ 個の元で生成できる。
\item $k$ 上大域的完全交叉な $A$ と $A$ の素イデアル
$\mathfrak a$ が存在して、$S \cong A_{\mathfrak a}$ である。
\item $k$ 上局所完全交叉な $A$ と $A$ の素イデアル
$\mathfrak a$ が存在して、$S \cong A_{\mathfrak a}$ である。
\end{enumerate}
\end{lemma}

\begin{proof}
(2) が (1) を、(1) が (3) を含意することは明らかである。
(4) が (5) を含意することも明らかである。(3) が (4) を含意することを示そう。
従って、$R$ は $k$ 上本質的有限表示な正則局所環であり、イデアル
$\Ker(R \to S)$ は $\dim(R) - \dim(S)$ 個の元で生成できるような全射
$R \to S$ が存在すると仮定する。
ある $J \subset k[x_1, \ldots, x_n]$ および
$J \subset \mathfrak q$ を満たす素イデアル
$\mathfrak q \subset k[x_1, \ldots, x_n]$ によって
$R = (k[x_1, \ldots, x_n]/J)_{\mathfrak q}$ と書ける。
$I \subset k[x_1, \ldots, x_n]$ を写像
$k[x_1, \ldots, x_n] \to S$ の核とすると、
$S \cong (k[x_1, \ldots, x_n]/I)_{\mathfrak q}$ である。
仮定により $(I/J)_{\mathfrak q}$ は
$\dim(R) - \dim(S)$ 個の元で生成される。
補題 \ref{algebra-lemma-ci-well-defined} により、
$I_{\mathfrak q}$ は
$\dim(k[x_1, \ldots, x_n]_{\mathfrak q}) - \dim(S)$ 個の元で生成できる。
補題 \ref{algebra-lemma-lci} により、ある
$g \in k[x_1, \ldots, x_n]$, $g \not \in \mathfrak q$ に対して、代数
$(k[x_1, \ldots, x_n]/I)_g$ は大域的完全交叉であり、
$S$ はその局所環に同型である。

\medskip\noindent
補題の証明を終えるには、(5) が (2) を含意することを示せばよい。
(5) を仮定し、$\pi : R \to S$ を、$R$ が $k$ 上本質的有限型な
正則局所 $k$-代数である全射とする。
仮定により、$k$ 上のある局所完全交叉 $A$ に対して
$S = A_{\mathfrak a}$ である。
$J \subset \mathfrak q \subset k[y_1, \ldots, y_m]$ を満たす表示
$R = (k[y_1, \ldots, y_m]/J)_{\mathfrak q}$ を選ぶ。
$J$ は写像 $k[y_1, \ldots, y_m] \to R$ の核であるとしてよく、そうする。
$I \subset k[y_1, \ldots, y_m]$ を写像
$k[y_1, \ldots, y_m] \to S = A_{\mathfrak a}$ の核とする。
このとき $J \subset I$ であり、$(I/J)_{\mathfrak q}$ は全射
$\pi : R \to S$ の核である。従って
$S = (k[y_1, \ldots, y_m]/I)_{\mathfrak q}$ である。

\medskip\noindent
補題 \ref{algebra-lemma-isomorphic-local-rings} により、
$g \in A$, $g \not \in \mathfrak a$ および
$g' \in k[y_1, \ldots, y_m]$, $g' \not \in \mathfrak q$ であって
$A_g \cong (k[y_1, \ldots, y_m]/I)_{g'}$ となるものが存在する。
$A$ を $A_g$ で、$k[y_1, \ldots, y_m]$ を
$k[y_1, \ldots, y_{m + 1}]$ で置き換えた後、
$A \cong k[y_1, \ldots, y_m]/I$ と仮定してよい。
局所環の全射
$$
k[y_1, \ldots, y_m]_{\mathfrak q} \to R \to S.
$$
を考える。$R \to S$ の核が正則列で生成されることを示さねばならない。
補題 \ref{algebra-lemma-lci} により、
$k[y_1, \ldots, y_m]_{\mathfrak q} \to A_{\mathfrak a} = S$
の核は正則列で生成される（$A$ は $k$ 上局所完全交叉だからである）。
補題 \ref{algebra-lemma-ci-well-defined} により結論を得る。
\end{proof}

\begin{lemma}
\label{algebra-lemma-lci-at-prime}
$k$ を体、$S$ を有限型 $k$-代数、$\mathfrak q$ を $S$ の
素イデアルとする。次の条件は同値である。
\begin{enumerate}
\item 局所環 $S_{\mathfrak q}$ は完全交叉環である
（定義 \ref{algebra-definition-lci-local-ring}）。
\item $g \in S$, $g \not \in \mathfrak q$ であって、$S_g$ が
$k$ 上局所完全交叉となるものが存在する。
\item $g \in S$, $g \not \in \mathfrak q$ であって、$S_g$ が
$k$ 上大域的完全交叉となるものが存在する。
\item $S = k[x_1, \ldots, x_n]/I$ の任意の表示に対し、
$\mathfrak q' \subset k[x_1, \ldots, x_n]$ が $\mathfrak q$ に対応するとき、
補題 \ref{algebra-lemma-lci} の同値な条件 (1) -- (5) のいずれかが成り立つ。
% P01-E0779: the frozen plural phrase “any of the equivalent conditions ... hold” has number disagreement and is logically clearer as “any one ... holds”.
\end{enumerate}
\end{lemma}

\begin{proof}
これは補題 \ref{algebra-lemma-lci}、補題 \ref{algebra-lemma-lci-local}、および
各定義を組み合わせたものである。
\end{proof}

\begin{lemma}
\label{algebra-lemma-lci-global}
$k$ を体、$S$ を有限型 $k$-代数とする。次の条件は同値である。
\begin{enumerate}
\item 環 $S$ は $k$ 上局所完全交叉である。
\item $S$ のすべての局所環は $k$ 上完全交叉環である。
\item $S$ の極大イデアルにおけるすべての局所化は
$k$ 上完全交叉環である。
\end{enumerate}
\end{lemma}

\begin{proof}
これは補題 \ref{algebra-lemma-lci-at-prime}、$\Spec(S)$ の準コンパクト性、
および各定義から従う。
\end{proof}

\noindent
次の補題は、完全交叉であるという性質が基礎体の変更で
（強い意味において）保たれることを述べる。

\begin{lemma}
\label{algebra-lemma-lci-field-change-local}
$K/k$ を体拡大とする。$S$ を $k$ 上有限型な代数とする。
$\mathfrak q_K$ を $S_K = K \otimes_k S$ の素イデアルとし、
$\mathfrak q$ をそれに対応する $S$ の素イデアルとする。
このとき $S_{\mathfrak q}$ が $k$ 上完全交叉であること
（定義 \ref{algebra-definition-lci-local-ring}）と、
$(S_K)_{\mathfrak q_K}$ が $K$ 上完全交叉であることは同値である。
\end{lemma}

\begin{proof}
表示 $S = k[x_1, \ldots, x_n]/I$ を選ぶ。
これから表示 $S_K = K[x_1, \ldots, x_n]/I_K$ が得られる。
ここで $I_K = K \otimes_k I$ である。
$\mathfrak q_K' \subset K[x_1, \ldots, x_n]$ および
$\mathfrak q' \subset k[x_1, \ldots, x_n]$ を、それぞれ対応する
素イデアルとする。補題 \ref{algebra-lemma-lci} の同値な条件が
組 $(S = k[x_1, \ldots, x_n]/I, \mathfrak q)$ に対して成り立つことと、
組 $(S_K = K[x_1, \ldots, x_n]/I_K, \mathfrak q_K)$ に対して
成り立つことが同値であると示す。そうすれば補題が従う
（補題 \ref{algebra-lemma-lci-at-prime} を参照せよ）。

\medskip\noindent
補題 \ref{algebra-lemma-dimension-at-a-point-preserved-field-extension} により
$\dim_{\mathfrak q} S = \dim_{\mathfrak q_K} S_K$ である。
従って補題 \ref{algebra-lemma-lci} に現れる整数 $c$ は、組
$(S = k[x_1, \ldots, x_n]/I, \mathfrak q)$ と組
$(S_K = K[x_1, \ldots, x_n]/I_K, \mathfrak q_K)$ に対して同じである。
他方、
\begin{eqnarray*}
I \otimes_{k[x_1, \ldots, x_n]} \kappa(\mathfrak q')
\otimes_{\kappa(\mathfrak q')} \kappa(\mathfrak q_K')
& = &
I \otimes_{k[x_1, \ldots, x_n]} \kappa(\mathfrak q_K') \\
& = &
I \otimes_{k[x_1, \ldots, x_n]} K[x_1, \ldots, x_n]
\otimes_{K[x_1, \ldots, x_n]} \kappa(\mathfrak q_K') \\
& = &
(K \otimes_k I) \otimes_{K[x_1, \ldots, x_n]} \kappa(\mathfrak q_K') \\
& = &
I_K \otimes_{K[x_1, \ldots, x_n]} \kappa(\mathfrak q'_K).
\end{eqnarray*}
% P01-E0780: the frozen display alternates between q_K' and q'_K for the same extended prime.
従って
$\dim_{\kappa(\mathfrak q')}
I \otimes_{k[x_1, \ldots, x_n]} \kappa(\mathfrak q')
=
\dim_{\kappa(\mathfrak q'_K)}
I_K \otimes_{K[x_1, \ldots, x_n]} \kappa(\mathfrak q_K')$
である。従って中山の補題 \ref{algebra-lemma-NAK} により、
$I_{\mathfrak q'}$ の最小生成元数と
$(I_K)_{\mathfrak q'_K}$ の最小生成元数は等しい。
ゆえに補題 \ref{algebra-lemma-lci} の特徴づけ (2) から結論が従う。
\end{proof}

\begin{lemma}
\label{algebra-lemma-lci-field-change}
$k \to K$ を体拡大とし、$S$ を有限型 $k$-代数とする。
このとき $S$ が $k$ 上局所完全交叉であることと、
$S \otimes_k K$ が $K$ 上局所完全交叉であることは同値である。
\end{lemma}

\begin{proof}
これは補題 \ref{algebra-lemma-lci-global} と補題
\ref{algebra-lemma-lci-field-change-local} を組み合わせれば従う。
ここでは（同じ原理に基づく）別の証明も与える。

\medskip\noindent
$S' = S \otimes_k K$ とおく。$\alpha : k[x_1, \ldots, x_n] \to S$
を核が $I$ である表示とし、
$\alpha' : K[x_1, \ldots, x_n] \to S'$ を核が $I'$ である
誘導された表示とする。

\medskip\noindent
$S$ が局所完全交叉であると仮定する。
素イデアル $\mathfrak q \subset S'$ を一つとる。
$\mathfrak q'$ を $K[x_1, \ldots, x_n]$ の対応する素イデアル、
$\mathfrak p$ を $S$ の対応する素イデアル、
$\mathfrak p'$ を $k[x_1, \ldots, x_n]$ の対応する素イデアルとする。
% P01-E0781: the frozen designation sentence omits “by/as” after “Denote”.
ネーター局所環の次の図を考える。
$$
\xymatrix{
S'_{\mathfrak q} &  K[x_1, \ldots, x_n]_{\mathfrak q'} \ar[l] \\
S_{\mathfrak p}\ar[u] &  k[x_1, \ldots, x_n]_{\mathfrak p'} \ar[u] \ar[l]
}
$$
% P01-E0782: the frozen diagram lacks terminal punctuation before the next sentence.
補題 \ref{algebra-lemma-lci} により、$S_{\mathfrak p}$ は
$k[x_1, \ldots, x_n]_{\mathfrak p'}$ において、ある正則列
$f_1, \ldots, f_c$ で切り出される。
右側の縦射は平坦なので、$f_1, \ldots, f_c$ の像は
$K[x_1, \ldots, x_n]_{\mathfrak q'}$ の正則列をなす。
$k$ 上 $K$ をテンソルする関手は完全なので
$S'_{\mathfrak q} = K[x_1, \ldots, x_n]_{\mathfrak q'}/(f_1, \ldots, f_c)$
である。従って再び補題 \ref{algebra-lemma-lci} により、$S'$ は
$\mathfrak q$ の近傍で局所完全交叉である。
$\mathfrak q$ は任意だったので、$S'$ は $K$ 上局所完全交叉である。

\medskip\noindent
$S'$ が局所完全交叉であると仮定する。
$S$ の極大イデアル $\mathfrak m$ を一つとり、
$\mathfrak m'$ を $k[x_1, \ldots, x_n]$ の対応する極大イデアルとする。
剰余体を $\kappa = \kappa(\mathfrak m)$ と書く。
% P01-E0783: the frozen residue-field designation omits “by/as”.
注意 \ref{algebra-remark-fundamental-diagram} により、$\mathfrak m$ 上にある
$S'$ の素イデアルは $K \otimes_k \kappa$ の素イデアルに対応する。
Hilbert の零点定理 \ref{algebra-theorem-nullstellensatz} により
$[\kappa : k] < \infty$ である。
従って $K \otimes_k \kappa$ は $K$ 上有限かつ零でない。
ゆえに $K \otimes_k \kappa$ は有限個、しかも $> 0$ 個の素イデアルをもち、
それらはすべて極大で、それぞれの剰余体は $K$ 上有限である
（第 \ref{algebra-section-artinian} 節を参照せよ）。
% P01-E0784: the frozen prose twice says “a finite number > 0” of primes.
従って $\mathfrak m$ 上にある素イデアル
$\mathfrak n \subset S'$ は有限個、しかも $> 0$ 個あり、それぞれ極大で、その剰余体は
$K$ 上有限である。その一つ $\mathfrak n \subset S'$ をとり、
$\mathfrak n' \subset K[x_1, \ldots, x_n]$ を
$K[x_1, \ldots, x_n]$ の対応する素イデアルとする。
$V(\mathfrak mS')$ は有限なので、$\mathfrak n$ はその孤立閉点である。
従って $\mathfrak mS'_{\mathfrak n}$ は
$S'_{\mathfrak n}$ の定義イデアルである。
このことから、例えば補題 \ref{algebra-lemma-dimension-base-fibre-equals-total} により
$\dim(S_{\mathfrak m}) = \dim(S'_{\mathfrak n})$ が従う。
（これは補題
\ref{algebra-lemma-dimension-at-a-point-preserved-field-extension} からも分かる。）
対応するネーター局所環の図
$$
\xymatrix{
S'_{\mathfrak n} &  K[x_1, \ldots, x_n]_{\mathfrak n'} \ar[l] \\
S_{\mathfrak m}\ar[u] &  k[x_1, \ldots, x_n]_{\mathfrak m'} \ar[u] \ar[l]
}
$$
を考える。
% P01-E0785: the frozen diagram lacks terminal punctuation before the next sentence.
補題 \ref{algebra-lemma-change-base-NL} によれば
$\NL(\alpha) \otimes_S S' = \NL(\alpha')$ であり、特に
$I'/(I')^2 = I/I^2 \otimes_S S'$ である。従って
$(I/I^2)_{\mathfrak m} \otimes_{S_{\mathfrak m}} \kappa$ と
$(I'/(I')^2)_{\mathfrak n} \otimes_{S'_{\mathfrak n}} \kappa(\mathfrak n)$
の次元は等しい。$(I'/(I')^2)_{\mathfrak n}$ は階数
$n - \dim S'_{\mathfrak n}$ の自由加群なので、
$(I/I^2)_{\mathfrak m}$ は
$n - \dim S'_{\mathfrak n} = n - \dim S_{\mathfrak m}$ 個の元で生成できる。
補題 \ref{algebra-lemma-lci} により、$S$ は $\mathfrak m$ の近傍で
局所完全交叉である。$\mathfrak m$ は任意の極大イデアルだったので、
$S$ は局所完全交叉である。
\end{proof}

\noindent
最後に、後で次のことを証明するために用いる補題を示す。
環写像 $T \to A \to B$ が与えられ、$B$ が $T$ 上 syntomic であり、
かつ $B$ が $A$ 上 syntomic ならば、$A$ は $T$ 上 syntomic である。

\begin{lemma}
\label{algebra-lemma-lci-permanence-initial}
局所環の可換図
$$
\xymatrix{
B & S \ar[l] \\
A \ar[u] & R \ar[l] \ar[u]
}
$$
を考える。
% P01-E0786: the frozen diagram lacks terminal punctuation before “be a commutative square”.
次を仮定する。
\begin{enumerate}
\item $R$ と $\overline{S} = S/\mathfrak m_R S$ は正則局所環である。
\item あるイデアル $I$, $J$ に対して $A = R/I$ および $B = S/J$ である。
\item $J \subset S$ および
$\overline{J} = J/\mathfrak m_R \cap J \subset \overline{S}$ は
正則列で生成される。
% P01-E0787: the frozen quotient omits parentheses and likely the extension of m_R to S.
\item $A \to B$ と $R \to S$ は平坦である。
\end{enumerate}
このとき $I$ は正則列で生成される。
\end{lemma}

\begin{proof}
$\overline{B} = B/\mathfrak m_RB = B/\mathfrak m_AB$ とおく。このとき
$\overline{B} = \overline{S}/\overline{J}$ である。
$\overline{f}_1, \ldots, \overline{f}_{\overline{c}} \in \overline{J}$ が
$\overline{J}$ を生成する正則列をなすような元
$f_1, \ldots, f_{\overline{c}} \in J$ をとる。
補題 \ref{algebra-lemma-ci-well-defined} により
$\overline{c} = \dim(\overline{S}) - \dim(\overline{B})$ であることに注意する。
補題 \ref{algebra-lemma-grothendieck-regular-sequence} により、環
$S/(f_1, \ldots, f_{\overline{c}})$ は $R$ 上平坦である。
従って $S/(f_1, \ldots, f_{\overline{c}}) + IS$ は $A$ 上平坦である。
% P01-E0788: the frozen quotient by the sum lacks parentheses in both occurrences.
従って写像 $S/(f_1, \ldots, f_{\overline{c}}) + IS \to B$ は、
$A$ 上平坦な有限 $S/IS$-加群の全射であり、
$\mathfrak m_A$ を法として同型であるから、
補題 \ref{algebra-lemma-mod-injective} により同型である。言い換えれば、
$J = (f_1, \ldots, f_{\overline{c}}) + IS$ である。

\medskip\noindent
再び補題 \ref{algebra-lemma-ci-well-defined} により、イデアル $J$ は
$c = \dim(S) - \dim(B)$ 個の元からなる正則列で生成される。
従って $J/\mathfrak m_SJ$ は次元 $c$ のベクトル空間である。
上の $J$ の記述により、
$g_1, \ldots, g_{c - \overline{c}} \in I$ であって、
$J$ が
$f_1, \ldots, f_{\overline{c}}, g_1, \ldots, g_{c - \overline{c}}$
で生成されるものが存在する
（中山の補題 \ref{algebra-lemma-NAK} を用いよ）。
環 $A' = R/(g_1, \ldots, g_{c - \overline{c}})$ と全射
$A' \to A$ を考える。上の議論から
$B = S/(f_1, \ldots, f_{\overline{c}}, g_1, \ldots, g_{c - \overline{c}})$
は $A'$ 上平坦である
（$S/(f_1, \ldots, f_{\overline{c}})$ が $R$ 上平坦だからである）。
従って $A' \to B$ は単射である
（これは忠実平坦だからである。補題 \ref{algebra-lemma-local-flat-ff} を参照せよ）。
この写像は $A$ を経由して因子分解するので $A' = A$ を得る。
補題 \ref{algebra-lemma-dimension-base-fibre-equals-total} により
$\dim(B) = \dim(A) + \dim(\overline{B})$ および
$\dim(S) = \dim(R) + \dim(\overline{S})$ であることに注意する。
従って初等的な代数計算により
$c - \overline{c} = \dim(R) -\dim(A)$ である。
% P01-E0789: the frozen dimension difference omits spacing before the second dim command.
ゆえに補題 \ref{algebra-lemma-ci-well-defined} により
$I = (g_1, \ldots, g_{c - \overline{c}})$ は正則列で生成される。
\end{proof}

\section{syntomic 射}
\label{algebra-section-syntomic}

\noindent
syntomic 環写像とは、平坦かつ有限表示で、すべてのファイバーが
局所完全交叉である環写像のことである。一般の局所完全交叉環写像については
More on Algebra、第 \ref{more-algebra-section-lci} 節で論じる。

\begin{definition}
\label{algebra-definition-lci}
環写像 $R \to S$ が平坦かつ有限表示で、そのすべてのファイバー環
$S \otimes_R \kappa(\mathfrak p)$ が局所完全交叉であるとき、
この写像は {\it syntomic} である、または $S$ は
{\it $R$ 上平坦局所完全交叉}であるという
（定義 \ref{algebra-definition-lci-field} を参照せよ）。
\end{definition}

\noindent
体上の代数がその体上 syntomic であることと、局所完全交叉であることは
明らかに同値である。この定義には次の好ましい特徴がある。

\begin{lemma}
\label{algebra-lemma-syntomic-descends}
\begin{slogan}
syntomic 性は基底上 fpqc 局所的である。
\end{slogan}
$R \to S$ を環写像、$R \to R'$ を忠実平坦な環写像とし、
$S' = R'\otimes_R S$ とおく。
このとき $R \to S$ が syntomic であることと
$R' \to S'$ が syntomic であることは同値である。
\end{lemma}

\begin{proof}
補題 \ref{algebra-lemma-finite-presentation-descends} および補題
\ref{algebra-lemma-flatness-descends} により、平坦であるという性質と
有限表示であるという性質については主張が成り立つ。
補題 \ref{algebra-lemma-ff-rings} により写像
$\Spec(R') \to \Spec(R)$ は全射である。
従って $\mathfrak p \subset R$ 上にある素イデアル
$\mathfrak p' \subset R'$ が与えられたとき、
$S \otimes_R \kappa(\mathfrak p)$ が局所完全交叉であることと
$S' \otimes_{R'} \kappa(\mathfrak p')$ が局所完全交叉であることの
同値性を示せば十分である。
% P01-E0790: the frozen “it suffices to show given primes ... that” construction is grammatically malformed.
ここで
$S' \otimes_{R'} \kappa(\mathfrak p') =
S \otimes_R \kappa(\mathfrak p)
\otimes_{\kappa(\mathfrak p)} \kappa(\mathfrak p')$
であることに注意する。従って補題 \ref{algebra-lemma-lci-field-change} を適用できる。
\end{proof}

\begin{lemma}
\label{algebra-lemma-base-change-syntomic}
syntomic 写像の任意の基底変換は syntomic である。
\end{lemma}

\begin{proof}
平坦性、有限表示性、およびファイバーが局所完全交叉であることの
いずれについても、補題 \ref{algebra-lemma-flat-base-change}、
\ref{algebra-lemma-compose-finite-type}、\ref{algebra-lemma-lci-field-change} により成り立つ。
\end{proof}

\begin{lemma}
\label{algebra-lemma-local-syntomic}
$R \to S$ を環写像とする。単位イデアルを生成する
$g_1, \ldots g_m \in S$ があり、各 $R \to S_{g_i}$ が
syntomic であると仮定する。このとき $R \to S$ は syntomic である。
% P01-E0791: the frozen mathematical list omits the comma after \ldots.
\end{lemma}

\begin{proof}
平坦性と有限表示性については、補題 \ref{algebra-lemma-flat-localization} および
\ref{algebra-lemma-cover-upstairs} により成り立つ。
ファイバー環が局所完全交叉であるという性質は、その定義から
$S$ 上局所的である（定義 \ref{algebra-definition-lci-field} を参照せよ）。
\end{proof}

\begin{definition}
\label{algebra-definition-relative-global-complete-intersection}
$R \to S$ を環写像とする。表示
$S = R[x_1, \ldots, x_n]/(f_1, \ldots, f_c)$ が存在し、
$\Spec(S) \to \Spec(R)$ の空でない各ファイバーの次元が $n - c$ であるとき、
$R \to S$ は {\it 相対的大域的完全交叉}であるという。
この状況を示すため、「$S = R[x_1, \ldots, x_n]/(f_1, \ldots, f_c)$ を
相対的大域的完全交叉とする」ともいう。
\end{definition}

\noindent
大域的表示を見つける際には、次の補題が役立つことがある。

\begin{lemma}
\label{algebra-lemma-huber}
$S$ を有限表示 $R$-代数とし、$I/I^2$ が $S$ 上自由であるような表示
$S = R[x_1, \ldots, x_n]/I$ をもつとする。このとき $S$ は、
$(f_1, \ldots, f_c)/(f_1, \ldots, f_c)^2$ が
$f_1, \ldots, f_c$ の類を基底とする自由加群となるような表示
$S = R[y_1, \ldots, y_m]/(f_1, \ldots, f_c)$ をもつ。
\end{lemma}

\begin{proof}
補題 \ref{algebra-lemma-finite-presentation-independent} により、
$I$ は有限生成イデアルである。$I/I^2$ の基底に写る元
$f_1, \ldots, f_c \in I$ をとる。
中山の補題（補題 \ref{algebra-lemma-NAK}）により、ある $g \in 1 + I$ が存在して
$$
g \cdot I \subset (f_1, \ldots, f_c)
$$
となり、$I_g \cong (f_1, \ldots, f_c)_g$ である。従って
$$
S \cong R[x_1, \ldots, x_n]/(f_1, \ldots, f_c)[1/g]
\cong R[x_1, \ldots, x_n, x_{n + 1}]/(f_1, \ldots, f_c, gx_{n + 1} - 1)
$$
となり、望む表示を得る。例えば補題
\ref{algebra-lemma-principal-localization-NL} を適用すれば、
$f_1, \ldots,\allowbreak f_c,gx_{n + 1} - 1$ は
$(f_1, \ldots,\allowbreak f_c, gx_{n + 1} - 1)\allowbreak/(f_1, \ldots,\allowbreak f_c, gx_{n + 1} - 1)^2$
の基底をなすことが従う。
% P01-E0792: the frozen source omits a space after the comma before gx_{n+1}-1.
\end{proof}



\begin{example}
\label{algebra-example-factor-polynomials}
$n , m \geq 1$ を整数とする。環写像
\begin{eqnarray*}
R = \mathbf{Z}[a_1, \ldots, a_{n + m}]
& \longrightarrow &
S = \mathbf{Z}[b_1, \ldots, b_n, c_1, \ldots, c_m] \\
a_1 & \longmapsto & b_1 + c_1 \\
a_2 & \longmapsto & b_2 + b_1 c_1 + c_2 \\
\ldots & \ldots & \ldots \\
a_{n + m} & \longmapsto & b_n c_m
\end{eqnarray*}
を考える。言い換えれば、これは多項式環のただ一つの環写像であって、
上に示した通りに定まり、多項式の因数分解
$$
x^{n + m} + a_1 x^{n + m - 1} + \ldots + a_{n + m}
=
(x^n + b_1 x^{n - 1} + \ldots + b_n)
(x^m + c_1 x^{m - 1} + \ldots + c_m)
$$
を満たすものである。$S$ は $R$ 上 $n + m$ 個の元
（すなわち $b_i, c_j$）で生成され、$n + m$ 個の関係式
（すなわち $a_k = a_k(b_i, c_j)$）があることに注意する。
$S$ が $R$ 上相対的大域的完全交叉であることを示すには、
すべてのファイバーの次元が $0$ であることを証明すれば十分である。

\medskip\noindent
これを示すため、$R \to k$ を体 $k$ への環写像とする。
$a_i$ が $\alpha_i \in k$ に写るとする。
ファイバー環 $S_k = k \otimes_R S$ を考え、$k \to K$ を体拡大とする。
$S_k \to K$ の $k$-代数写像を与えることは、
$\beta_1, \ldots, \beta_n, \gamma_1, \ldots, \gamma_m \in K$ であって
% P01-E0793: the frozen prose says “a map of S_k to K” instead of “from S_k to K”.
$$
x^{n + m} + \alpha_1 x^{n + m - 1} + \ldots + \alpha_{n + m}
=
(x^n + \beta_1 x^{n - 1} + \ldots + \beta_n)
(x^m + \gamma_1 x^{m - 1} + \ldots + \gamma_m).
$$
を満たすものを見つけることと同じである。従って $K$ における
そのような $n + m$-組の選び方は高々有限個である。
% P01-E0794: the frozen “at most finitely many” is redundant.
これにより、すべてのファイバーは有限個の閉点しかもたない
（例えば、それらがすべて $\overline{k}$ における解に対応することを見るには
Hilbert の零点定理を用いよ）。従って $R \to S$ は相対的大域的完全交叉である。

\medskip\noindent
別の議論として、写像
$\mathbf{Z}[a_1, \ldots, a_{n + m}] \to
\mathbf{Z}[b_1, \ldots, b_n, c_1, \ldots, c_m]$
が実は {\it 有限}な環写像でもあることを示せばよい。
実際、補題 \ref{algebra-lemma-polynomials-divide} により各 $b_i, c_j$ は
$R$ 上整であるから、補題 \ref{algebra-lemma-characterize-integral} により
$R \to S$ は有限である。
\end{example}

\begin{example}
\label{algebra-example-roots-universal-polynomial}
環写像
\begin{eqnarray*}
R = \mathbf{Z}[a_1, \ldots, a_n]
& \longrightarrow &
S = \mathbf{Z}[\alpha_1, \ldots, \alpha_n] \\
a_1 & \longmapsto &
\alpha_1 + \ldots + \alpha_n \\
\ldots & \ldots & \ldots \\
a_n & \longmapsto & \alpha_1 \ldots \alpha_n
\end{eqnarray*}
を考える。言い換えれば、これは上に示した通りに定まる多項式環の
ただ一つの環写像であって、
$$
x^n + a_1 x^{n - 1} + \ldots + a_n
=
\prod\nolimits_{i = 1}^n (x + \alpha_i)
$$
を $\mathbf{Z}[\alpha_i, x]$ において満たすものである。
別の言い方をすれば、$a_i$ は
$\alpha_1, \ldots, \alpha_n$ の $i$ 番目の基本対称式に写る。
% P01-E0795: the frozen English typesets “i” and “th” without a hyphen.
基本対称多項式の通常の理論（詳細は省略する）により、環 $S$ は
$R$ 上有限自由であり、その基底は
$0 \leq e_i \leq n - i$ を満たす元
$\alpha_1^{e_1} \alpha_2^{e_2} \ldots \alpha_n^{e_n}$ からなる。
従ってなおさら $R \to S$ のファイバー環は有限であり、次元 $0$ である。
他方、$S$ は $R$ 上 $n$ 個の元で生成され、$n$ 個の関係式をもつ。
従って $S$ は $R$ 上相対的大域的完全交叉である。
$S$ の $R$ 上の階数は正なので、$S$ は $R$ 上忠実平坦でもある。
\end{example}

\begin{lemma}
\label{algebra-lemma-base-change-relative-global-complete-intersection}
$S = R[x_1, \ldots, x_n]/(f_1, \ldots, f_c)$ を相対的大域的完全交叉とする
（定義 \ref{algebra-definition-relative-global-complete-intersection}）。
% P01-E0796: the frozen statement lacks punctuation before the enumeration.
\begin{enumerate}
\item 任意の $R \to R'$ に対して、基底変換
$R' \otimes_R S = R'[x_1, \ldots, x_n]/(f_1, \ldots, f_c)$ は
相対的大域的完全交叉である。
\item $h \in R[x_1, \ldots, x_n]$ の像である任意の $g \in S$ に対して、環
$S_g\allowbreak = R[x_1, \ldots,\allowbreak x_n, x_{n + 1}]\allowbreak/(f_1, \ldots,\allowbreak f_c, hx_{n + 1} - 1)$
は相対的大域的完全交叉である。
\item ある $f \in R$ に対して $R \to S$ が
$R \to R_f \to S$ と因子分解すると仮定する。
% P01-E0797: the frozen source incorrectly separates the if-clause and “Then” with a period.
このとき環 $S = R_f[x_1, \ldots, x_n]/(f_1, \ldots, f_c)$ は
$R_f$ 上相対的大域的完全交叉である。
\end{enumerate}
\end{lemma}

\begin{proof}
補題 \ref{algebra-lemma-dimension-preserved-field-extension} により、
基底変換のファイバーは元の写像のファイバーと同じ次元をもつ。
さらに
$R' \otimes_R R[x_1, \ldots, x_n]/(f_1, \ldots, f_c)
= R'[x_1, \ldots, x_n]/(f_1, \ldots, f_c)$ である。従って (1) が従う。
(2) は、一元における局所化が
$S_g \cong S[x_{n + 1}]/(gx_{n + 1} - 1)$ と書けることから従う。
(3) の仮定の下では $R_f \otimes_R S \cong S$ なので、
(3) は (1) から従う。
\end{proof}

\begin{lemma}
\label{algebra-lemma-localize-relative-complete-intersection}
$R$ を環とし、$S = R[x_1, \ldots, x_n]/(f_1, \ldots, f_c)$ とする。
次の各場合に、$g \in S$ に写る $h \in R[x_1, \ldots, x_n]$ であって、
$$
S_g = R[x_1, \ldots, x_n, x_{n + 1}]/(f_1, \ldots, f_c, hx_{n + 1} - 1)
$$
が定義 \ref{algebra-definition-relative-global-complete-intersection} の形の表示をもつ
相対的大域的完全交叉となるものを見つけることができる。
\begin{enumerate}
\item $I \subset R$ をイデアルとする。
$\Spec(S/IS) \to \Spec(R/I)$ のファイバーの次元が $n - c$ ならば、
上のような $(h, g)$ で、$g$ が $1 \in S/IS$ に
写るものをとれる。
\item $\mathfrak p \subset R$ を素イデアルとする。
$\dim(S \otimes_R \kappa(\mathfrak p)) = n - c$ ならば、
上のような $(h, g)$ で、$g$ が
$S \otimes_R \kappa(\mathfrak p)$ の単元に写るものをとれる。
\item $\mathfrak q \subset S$ を $\mathfrak p \subset R$ 上にある
素イデアルとする。$\dim_{\mathfrak q}(S/R) = n - c$ ならば、
上のような $(h, g)$ で $g \not \in \mathfrak q$ となるものをとれる。
\end{enumerate}
\end{lemma}

\begin{proof}
(1) について。補題 \ref{algebra-lemma-dimension-fibres-bounded-open-upstairs} により、
$V(IS)$ を含む開部分集合 $W \subset \Spec(S)$ であって、
$W \to \Spec(R)$ のすべてのファイバーの次元が $\leq n - c$ となるものが存在する。
$W = \Spec(S) \setminus V(J)$ と書く。このとき
$V(J) \cap V(IS) = \emptyset$ だから、
$1 \in S/IS$ に写る $g \in J$ をとれる。
$h \in R[x_1, \ldots, x_n]$ を $g$ の任意の逆像とする。

\medskip\noindent
(2) について。補題 \ref{algebra-lemma-dimension-fibres-bounded-open-upstairs} により、
$\Spec(S \otimes_R \kappa(\mathfrak p))$ を含む開部分集合
$W \subset \Spec(S)$ であって、$W \to \Spec(R)$ のすべてのファイバーの
次元が $\leq n - c$ となるものが存在する。
$W = \Spec(S) \setminus V(J)$ と書く。このとき
$V(J \cdot S \otimes_R \kappa(\mathfrak p)) = \emptyset$ である。
従って $S \otimes_R \kappa(\mathfrak p)$ において単元に写る
$g \in J$ をとれる（詳細は省略する）。
$h \in R[x_1, \ldots, x_n]$ を $g$ の任意の逆像とする。

\medskip\noindent
(3) について。補題 \ref{algebra-lemma-dimension-fibres-bounded-open-upstairs} により、
$g \in S$, $g \not \in \mathfrak q$ であって、
$R \to S_g$ の空でないすべてのファイバーの次元が
$\leq n - c$ となるものが存在する。
$h \in R[x_1, \ldots, x_n]$ を $g$ に写る任意の元とする。
\end{proof}

\noindent
次の補題は、相対的大域的完全交叉に対して
絶対ネーター近似を行えることを述べる。

\begin{lemma}
\label{algebra-lemma-relative-global-complete-intersection-Noetherian}
$R$ を環とし、$S = R[x_1, \ldots, x_n]/(f_1, \ldots, f_c)$ を
相対的大域的完全交叉とする
（定義 \ref{algebra-definition-relative-global-complete-intersection}）。
$f_i \in R_0[x_1, \ldots, x_n]$ であり、かつ
$$
S_0 = R_0[x_1, \ldots, x_n]/(f_1, \ldots, f_c)
$$
が相対的大域的完全交叉となるような有限型
$\mathbf{Z}$-部分代数 $R_0 \subset R$ が存在する。
% P01-E0798: the frozen “There exist a ... subalgebra” has subject-verb disagreement.
\end{lemma}

\begin{proof}
$R_0 \subset R$ を、多項式 $f_1, \ldots, f_c$ のすべての係数で生成される
$R$ の $\mathbf{Z}$-代数とする。
$S_0 = R_0[x_1, \ldots, x_n]/(f_1, \ldots, f_c)$ とおく。
明らかに $S = R \otimes_{R_0} S_0$ である。
素イデアル $\mathfrak q \subset S$ をとり、その下にある
$\mathfrak p \subset R$, $\mathfrak q_0 \subset S_0$,
$\mathfrak p_0 \subset R_0$ をそれぞれ対応する素イデアルとする。
% P01-E0799: the frozen “denote ... the primes” construction omits by/as and compresses a chain of contractions into “it lies over”.
$\dim (S \otimes_R \kappa(\mathfrak p) ) = n - c$ なので、
補題 \ref{algebra-lemma-dimension-preserved-field-extension} により
$\dim (S_0 \otimes_{R_0} \kappa(\mathfrak p_0)) = n - c$ でもある。
補題 \ref{algebra-lemma-dimension-fibres-bounded-open-upstairs} により、
$g \in S_0$, $g \not \in \mathfrak q_0$ であって、
$R_0 \to (S_0)_g$ の空でないすべてのファイバーの次元が
$\leq n - c$ となるものが存在する。
$\mathfrak q$ は任意であり $\Spec(S)$ は準コンパクトなので、
有限個の $g_1, \ldots, g_m \in S_0$ をとって、(a) 各
$j = 1, \ldots, m$ に対して
$R_0 \to (S_0)_{g_j}$ の空でないファイバーの次元が
$\leq n - c$ であり、(b) $\Spec(S) \to \Spec(S_0)$ の像が
$D(g_1) \cup \ldots \cup D(g_m)$ に含まれるようにできる。
言い換えれば、$S = R \otimes_{R_0} S_0$ における
$g_1, \ldots, g_m$ の像は単位イデアルを生成する。
$R_0$ を大きくした後、$g_1, \ldots, g_m$ が $S_0$ において
単位イデアルを生成すると仮定してよい。(a) により、
$R_0 \to S_0$ の空でないすべてのファイバーの次元は
$\leq n - c$ であり、結論を得る。
\end{proof}

\begin{lemma}
\label{algebra-lemma-relative-global-complete-intersection-conormal}
$R$ を環とし、$S = R[x_1, \ldots, x_n]/(f_1, \ldots, f_c)$ を
相対的大域的完全交叉とする
（定義 \ref{algebra-definition-relative-global-complete-intersection}）。
$S$ の各素イデアル $\mathfrak q$ に対し、
$\mathfrak q'$ を $R[x_1, \ldots, x_n]$ の対応する素イデアルとする。
このとき次が成り立つ。
\begin{enumerate}
\item $f_1, \ldots, f_c$ は局所環
$R[x_1, \ldots, x_n]_{\mathfrak q'}$ の正則列をなす。
% P01-E0800: the frozen English uses singular “is” with the list f_1,...,f_c.
\item 各環
$R[x_1, \ldots, x_n]_{\mathfrak q'}/(f_1, \ldots, f_i)$ は $R$ 上平坦である。
\item $S$-加群 $(f_1, \ldots, f_c)/(f_1, \ldots, f_c)^2$ は自由であり、
元 $f_i \bmod (f_1, \ldots, f_c)^2$ を基底とする。
\end{enumerate}
\end{lemma}

\begin{proof}
$R$ はネーターであると仮定する。
$\mathfrak p = R \cap \mathfrak q'$ とおく。
例えば補題 \ref{algebra-lemma-lci} により、
$f_1, \ldots, f_c$ は局所環
$R[x_1, \ldots, x_n]_{\mathfrak q'} \otimes_R \kappa(\mathfrak p)$ の
正則列をなす。さらに局所環
$R[x_1, \ldots, x_n]_{\mathfrak q'}$ は $R_{\mathfrak p}$ 上平坦である。
$R$、従って $R[x_1, \ldots, x_n]_{\mathfrak q'}$ はネーターなので、
補題 \ref{algebra-lemma-grothendieck-regular-sequence} から (1) と (2) が従う。

\medskip\noindent
$R$ を一般の環とする。これを有限型 $\mathbf{Z}$-部分代数の
フィルター付き余極限として
$R = \colim_{\lambda \in \Lambda} R_\lambda$ と書く
（第 \ref{algebra-section-colimits-flat} 節と比較せよ）。
すべての $\lambda$ に対して
$f_1, \ldots, f_c \in R_\lambda[x_1, \ldots, x_n]$ と仮定してよい。
$R_0 \subset R$ を補題
\ref{algebra-lemma-relative-global-complete-intersection-Noetherian} のものとする。
すべての $\lambda$ に対して $R_0 \subset R_\lambda$ と仮定してよい。
このとき $S_\lambda = R_\lambda[x_1, \ldots, x_n]/(f_1, \ldots, f_c)$ は
相対的大域的完全交叉である
（$R_0 \to R_\lambda$ による $S_0$ の基底変換だからである。補題
\ref{algebra-lemma-base-change-relative-global-complete-intersection} を参照せよ）。
$\mathfrak p_\lambda$, $\mathfrak q_\lambda$, $\mathfrak q'_\lambda$ を、
$\mathfrak p$, $\mathfrak q$, $\mathfrak q'$ から誘導される
$R_\lambda$, $S_\lambda$, $R_\lambda[x_1, \ldots, x_n]$ の
素イデアルとする。
% P01-E0801: the frozen designation omits by/as and uses singular “the prime” for three primes.
この記法の下で、各 $\lambda$ に対して (1) と (2) が成り立つ。
また
$$
R[x_1, \ldots, x_n]_{\mathfrak q'}/(f_1, \ldots, f_i)
=
\colim R_\lambda[x_1, \ldots, x_n]_{\mathfrak q_\lambda'}/(f_1, \ldots, f_i)
$$
なので、補題 \ref{algebra-lemma-colimit-rings-flat} により (2) の $R$ 上の
平坦性を得る。さらに
\begin{align*}
R[x_1, \ldots, x_n]_{\mathfrak q'}/(f_1, \ldots, f_i)
\xrightarrow{f_{i + 1}}
R[x_1, \ldots, x_n]_{\mathfrak q'}/(f_1, \ldots, f_i) \\
=
\colim
\left(
R_\lambda[x_1, \ldots, x_n]_{\mathfrak q_\lambda'}/(f_1, \ldots, f_i)
\xrightarrow{f_{i + 1}}
R_\lambda[x_1, \ldots, x_n]_{\mathfrak q_\lambda'}/(f_1, \ldots, f_i)
\right)
\end{align*}
であり、フィルター付き余極限は完全なので
（補題 \ref{algebra-lemma-directed-colimit-exact}）、(1) が従う。

\medskip\noindent
(3) を示す。$N$ を $S$-加群
$(f_1, \ldots, f_c)/(f_1, \ldots, f_c)^2$ とし、
$e_i \in N$ を $f_i$ の像とする。
% P01-E0802: the frozen “Denote N the S-module” construction omits by/as.
補題 \ref{algebra-lemma-regular-quasi-regular} と (1) により、
$S$ のすべての素イデアル $\mathfrak q$ に対して
$e_1, \ldots, e_c$ は $N_\mathfrak q$ の基底をなす。
% P01-E0803: the frozen subject e_1,...,e_c takes singular “is a basis”; “form a basis” is required.
補題 \ref{algebra-lemma-characterize-zero-local} により (3) が従う。
\end{proof}

\begin{lemma}
\label{algebra-lemma-relative-global-complete-intersection}
相対的大域的完全交叉は syntomic、すなわち平坦である。
% P01-E0804: the frozen “syntomic, i.e., flat” can misleadingly identify syntomicity with flatness outside the already-constrained relative-global-CI context.
\end{lemma}

\begin{proof}
$R \to S$ を相対的大域的完全交叉とする。
ファイバーは大域的完全交叉であり、$S$ は $R$ 上有限表示である。
従って示すべきことは $R \to S$ の平坦性だけである。
これは補題 \ref{algebra-lemma-relative-global-complete-intersection-conormal} の
(2) により成り立つ。
\end{proof}

\begin{lemma}
\label{algebra-lemma-adjoin-roots}
$A$ を環とし、
$P(x) = x^n + b_1 x^{n-1} + \ldots + b_n \in A[x]$ を
$A$ 上のモニック多項式とする。このとき syntomic、有限自由、忠実平坦な
環拡大 $A \subset A'$ であって、ある $\beta_i \in A'$ に対して
$P(x) = \prod_{i = 1, \ldots, n} (x - \beta_i)$ となるものが存在する。
\end{lemma}

\begin{proof}
$R$ と $S$ を例 \ref{algebra-example-roots-universal-polynomial} のものとし、
$R \to A$ が $a_i$ を $b_i$ に写すとき、
$A' = A \otimes_R S$ とおき、$\beta_i = -1 \otimes \alpha_i$ とする。
$R \to S$ は忠実平坦かつ有限自由であり、補題
\ref{algebra-lemma-relative-global-complete-intersection} により syntomic である。
これらの性質は基底変換 $A \to A'$ に継承される。細部は省略する。
\end{proof}

\begin{lemma}
\label{algebra-lemma-syntomic}
$R \to S$ を環写像とし、$\mathfrak q \subset S$ を
$R$ の素イデアル $\mathfrak p$ 上にある素イデアルとする。
次の条件は同値である。
\begin{enumerate}
\item $g \in S$, $g \not \in \mathfrak q$ であって
$R \to S_g$ が syntomic となる元が存在する。
\item $g \in S$, $g \not \in \mathfrak q$ であって
$S_g$ が $R$ 上相対的大域的完全交叉となる元が存在する。
\item $g \in S$, $g \not \in \mathfrak q$ であって、
$R \to S_g$ が有限表示で、局所環写像
$R_{\mathfrak p} \to S_{\mathfrak q}$ が平坦であり、局所環
$S_{\mathfrak q}/\mathfrak pS_{\mathfrak q}$ が
$\kappa(\mathfrak p)$ 上完全交叉環となる元が存在する
（定義 \ref{algebra-definition-lci-local-ring} を参照せよ）。
\end{enumerate}
\end{lemma}

\begin{proof}
(1) $\Rightarrow$ (3) は補題 \ref{algebra-lemma-lci-at-prime} であり、
(2) $\Rightarrow$ (1) は補題
\ref{algebra-lemma-relative-global-complete-intersection} である。
残るのは (3) が (2) を含意することの証明である。

\medskip\noindent
(3) を仮定する。ある $g \in S$, $g\not\in \mathfrak q$ に対して
$S$ を $S_g$ で置き換えた後、$S$ は $R$ 上有限表示であると仮定してよい。
% P01-E0805: the frozen sentence lacks a comma after the localization condition.
表示 $S = R[x_1, \ldots, x_n]/I$ を選び、
$\mathfrak q' \subset R[x_1, \ldots, x_n]$ を
$\mathfrak q$ に対応する素イデアルとする。
$\kappa(\mathfrak p) = k$ と書く。
$S \otimes_R k = k[x_1, \ldots, x_n]/\overline{I}$ であることに注意する。
ここで $\overline{I} \subset k[x_1, \ldots, x_n]$ は
$I$ の像で生成されるイデアルである。
$\overline{\mathfrak q}' \subset k[x_1, \ldots, x_n]$ を
$\mathfrak q'$ の像で生成される素イデアルとする。
補題 \ref{algebra-lemma-lci-at-prime} により、補題 \ref{algebra-lemma-lci} の同値な条件は
$\overline{I}$ と $\overline{\mathfrak q}'$ に対して成り立つ。
$\overline{I}_{\overline{\mathfrak q}'}/
\overline{\mathfrak q}'\overline{I}_{\overline{\mathfrak q}'}$ の
$\kappa(\overline{\mathfrak q}')$ 上の次元を $c$ とする。
このベクトル空間の基底に写る元
$f_1, \ldots, f_c \in I$ をとる。
$\overline{f}_j \in \overline{I}$ の像は
$\overline{I}_{\overline{\mathfrak q}'}$ を生成する
（補題 \ref{algebra-lemma-lci} による）。
$S' = R[x_1, \ldots, x_n]/(f_1, \ldots, f_c)$ とおき、
$J$ を全射 $S' \to S$ の核とする。$S$ は有限表示なので、
$J$ は有限生成イデアルである
（補題 \ref{algebra-lemma-compose-finite-type}）。短完全列
$$
0 \to J \to S' \to S \to 0
$$
を考える。
% P01-E0806: the frozen short exact display lacks terminal punctuation before the next sentence.
$S_\mathfrak q$ は $R$ 上平坦なので、写像
$J_{\mathfrak q'} \otimes_R k \to S'_{\mathfrak q'} \otimes_R k$
は単射である（補題 \ref{algebra-lemma-flat-tor-zero}）。
一方、構成により $S'_{\mathfrak q'} \otimes_R k$ は
$S_\mathfrak q \otimes_R k$ に同型に写る。従って
$J_{\mathfrak q'} \otimes_R k =
J_{\mathfrak q'}/\mathfrak pJ_{\mathfrak q'} = 0$ である。
中山の補題（補題 \ref{algebra-lemma-NAK}）により、
$g \in R[x_1, \ldots, x_n]$, $g \not \in \mathfrak q'$ であって
$J_g = 0$ となるものが存在する。言い換えれば
$S'_g \cong S_g$ である。さらに局所化すれば、補題
\ref{algebra-lemma-localize-relative-complete-intersection} により
$S'$（従って $S$）は相対的大域的完全交叉となる。これが望む結論である。
\end{proof}

\begin{lemma}
\label{algebra-lemma-syntomic-presentation-ideal-mod-squares}
$R$ を環とし、ある有限生成イデアル $I$ に対して
$S = R[x_1, \ldots, x_n]/I$ とする。
$g \in S$ が、$S_g$ が $R$ 上 syntomic となる元ならば、
$(I/I^2)_g$ は有限射影 $S_g$-加群である。
\end{lemma}

\begin{proof}
補題 \ref{algebra-lemma-syntomic} により、$S_g$ において単位イデアルを生成する
有限個の元 $g_1, \ldots, g_m \in S$ であって、各
$S_{gg_j}$ が $R$ 上相対的大域的完全交叉となるものが存在する。
補題 \ref{algebra-lemma-finite-projective} により
$(I/I^2)_{gg_j}$ が有限射影であることを示せば十分なので、
$S_g$ は相対的大域的完全交叉であると仮定してよい。
この場合、補題 \ref{algebra-lemma-conormal-module-localize} および
\ref{algebra-lemma-relative-global-complete-intersection-conormal} から結論が従う。
\end{proof}

\begin{lemma}
\label{algebra-lemma-composition-syntomic}
$R \to S$, $S \to S'$ を環写像とする。
\begin{enumerate}
\item $R \to S$ と $S \to S'$ が syntomic ならば、
$R \to S'$ も syntomic である。
\item $R \to S$ と $S \to S'$ が相対的大域的完全交叉ならば、
$R \to S'$ も相対的大域的完全交叉である。
\end{enumerate}
\end{lemma}

\begin{proof}
(2) を示す。
% P01-E0807: the frozen “Proof of (2).” is a sentence fragment.
$R \to S$ と $S \to S'$ が相対的大域的完全交叉であり、
定義 \ref{algebra-definition-relative-global-complete-intersection} のような表示
$S =  R[x_1, \ldots, x_n]/(f_1, \ldots, f_c)$ および
$S' = S[y_1, \ldots, y_m]/(h_1, \ldots, h_d)$ があるとする。このとき
$$
S' \cong
R[x_1, \ldots, x_n, y_1, \ldots, y_m]/(f_1, \ldots, f_c, h'_1, \ldots, h'_d)
$$
となる。ここで $h_j' \in R[x_1, \ldots, x_n, y_1, \ldots, y_m]$ は
$h_j$ の持ち上げである。
% P01-E0808: the frozen display and prose place the prime on h in inconsistent source order (h'_j versus h_j').
従ってファイバー環の次元を上から評価すれば十分であり、
$R = k$ は体であると仮定してよい。
この場合、$k$ 上有限型で次元 $n - c$ の等次元環 $S$ と、
空でないすべてのファイバー環が次元 $m - d$ の等次元環となる
有限型環写像 $S \to S'$ がある。従って、例えば補題
\ref{algebra-lemma-dimension-base-fibre-total} を $S'$ の極大イデアルにおける
局所化へ適用すれば、望み通り
$\dim(S') \leq n - c + m - d$ を得る。

\medskip\noindent
(1) を (2) に帰着する。$R \to S$ と $S \to S'$ が syntomic であると仮定する。
$\mathfrak q' \subset S$ を $\mathfrak q \subset S$ 上にある素イデアルとする。
% P01-E0809: the frozen source places both q' and q in S; q' must be a prime of S' for the subsequent localization and containment statements.
補題 \ref{algebra-lemma-syntomic} により、
$g' \in S'$, $g' \not \in \mathfrak q'$ であって
$S \to S'_{g'}$ が相対的大域的完全交叉となるものが存在する。
同様に、$g \in S$, $g \not \in \mathfrak q$ であって
$R \to S_g$ が相対的大域的完全交叉となるものをとる。
補題 \ref{algebra-lemma-base-change-relative-global-complete-intersection} により
環写像 $S_g \to S_{gg'}$ は相対的大域的完全交叉である。
% P01-E0810: the frozen base-change target is S_{gg'}, but g' belongs to S'; the target should be S'_{gg'}.
(2) により、$R \to S_{gg'}$ は相対的大域的完全交叉であり、
$gg' \not \in \mathfrak q'$ である。
% P01-E0811: the frozen composite repeats S_{gg'} where the target must be the localization of S'.
$\mathfrak q'$ は任意だったので、補題 \ref{algebra-lemma-syntomic} と
\ref{algebra-lemma-local-syntomic} を組み合わせれば、$R \to S'$ は syntomic である
（ここでは $S'$ のスペクトルが準コンパクトであることも用いる。
補題 \ref{algebra-lemma-quasi-compact} を参照せよ）。
\end{proof}

\noindent
次の補題は後で改良する
（Smoothing Ring Maps、命題 \ref{smoothing-proposition-lift-smooth} を参照せよ）。

\begin{lemma}
\label{algebra-lemma-lift-syntomic}
$R$ を環、$I \subset R$ をイデアルとする。
$R/I \to \overline{S}$ を syntomic 写像とする。
$\overline{S}$ の単位イデアルを生成する元
$\overline{g}_i \in \overline{S}$ であって、各
$\overline{S}_{g_i} \cong S_i/IS_i$ が $R$ 上のある
相対的大域的完全交叉 $S_i$ によって書けるものが存在する。
% P01-E0812: the frozen “Then there exists elements” has subject-verb disagreement.
% P01-E0813: the frozen localization uses g_i although only \overline{g}_i is introduced; the same mismatch recurs in the proof.
\end{lemma}

\begin{proof}
補題 \ref{algebra-lemma-syntomic} により、$\overline{S}$ の単位イデアルを生成する
元の族 $\overline{g}_i \in \overline{S}$ であって、
各 $\overline{S}_{g_i}$ が $R/I$ 上相対的大域的完全交叉となるものが存在する。
従って $\overline{S}$ 自身が相対的大域的完全交叉であると仮定してよい。
定義 \ref{algebra-definition-relative-global-complete-intersection} のように
$\overline{S} =
(R/I)[x_1, \ldots, x_n]/(\overline{f}_1, \ldots, \overline{f}_c)$
と書く。$\overline{f}_1, \ldots, \overline{f}_c$ を持ち上げる
$f_1, \ldots, f_c \in R[x_1, \ldots, x_n]$ を選び、
$S = R[x_1, \ldots, x_n]/(f_1, \ldots, f_c)$ とおく。
$S/IS \cong \overline{S}$ であることに注意する。
補題 \ref{algebra-lemma-localize-relative-complete-intersection} により、
$\overline{S}$ において $1$ に写る $g \in S$ であって、
$S_g$ が $R$ 上相対的大域的完全交叉となるものをとれる。
$\overline{S} \cong S_g/IS_g$ なので証明は終わる。
\end{proof}

\section{滑らかな環写像}
\label{algebra-section-smooth}

\noindent
滑らかな環写像の定義を一つの例によって動機づけよう。
$R$ を環とし、ある零でない $f \in R[x, y]$ に対して
$S = R[x, y]/(f)$ とする。この場合、完全列
$$
S \to
S\text{d}x \oplus S\text{d}y \to
\Omega_{S/R} \to 0
$$
があり、最初の矢印は $1$ を
$\frac{\partial f}{\partial x} \text{d}x +
\frac{\partial f}{\partial y} \text{d}y$ に写す。第
% P01-E0814: the frozen formula is followed by “see Section” without punctuation.
\ref{algebra-section-netherlander} 節を参照せよ。
$f$ の偏導関数が $S$ の単位イデアルを生成するならば、
$\Omega_{S/R}$ は階数 $1$ の局所自由加群であることが分かる。
この場合、$S$ は $R$ 上相対次元 $1$ で滑らかである。
しかし、$f$ の二つの偏導関数がともに零ならば、
$\Omega_{S/R}$ が階数 $2$ の局所自由加群になることもある。
例えば、素数 $p$ に対して $R$ で $p = 0$ であり、
$f = x^p + y^p$ ならば、このことが起こる。
ここで $R \to S$ は相対次元 $1$ の相対的大域的完全交叉だが、
滑らかではない。従って、環写像が滑らかであることを確かめるには、
微分加群が自由であるかどうかを確かめるだけでは十分でない。
正しい条件は次のものである。

\begin{definition}
\label{algebra-definition-smooth}
環写像 $R \to S$ が有限表示であり、その素朴な余接複体
$\NL_{S/R}$ が次数 $0$ に置かれた有限射影 $S$-加群に擬同型であるとき、
その環写像は {\it 滑らか}であるという。これは
$H_1(\NL_{S/R}) = 0$ であり、かつ $\Omega_{S/R}$ が有限射影
$S$-加群であることを意味する。
\end{definition}

\noindent
$R \to S$ を環写像とする。補題 \ref{algebra-lemma-NL-homotopy} により、
任意の表示 $\alpha : P \to S$ に対して
$H_1(\NL(\alpha)) = H_1(\NL_{S/R})$ および
$H_0(\NL(\alpha)) = \Omega_{S/R}$ である。従って、
$R \to S$ が滑らかならば、核を $I$ とする任意の全射
$\alpha : R[x_1, \ldots, x_n] \to S$ に対して、写像
$$
I/I^2
\longrightarrow
\Omega_{R[x_1, \ldots, x_n]/R} \otimes_{R[x_1, \ldots, x_n]} S
$$
は単射であり、その余核は有限射影 $S$-加群である。
従って、表示された矢印は分裂単射である。言い換えると、
$S$-加群として
$\bigoplus_{i = 1}^n S \text{d}x_i \cong I/I^2 \oplus \Omega_{S/R}$
であり、$I/I^2$ も有限射影 $S$-加群である。逆に、
$R \to S$ が有限表示であり、核を $I$ とする全射
$\alpha : R[x_1, \ldots, x_n] \to S$ で、表示された矢印が
分裂単射となるものがあれば、$R \to S$ は滑らかである。

\begin{lemma}
\label{algebra-lemma-localize-smooth}
$R \to S$ を滑らかな環写像とする。
任意の局所化 $S_g$ は $R$ 上滑らかである。
$f \in R$ の像が $S$ で可逆ならば、$R_f \to S$ は滑らかである。
\end{lemma}

\begin{proof}
補題 \ref{algebra-lemma-localize-NL} により、$S_g$ の $R$ 上の素朴な余接複体は、
$S$ の $R$ 上の素朴な余接複体の基底変換である。仮定により、
$S/R$ の素朴な余接複体は $\Omega_{S/R}$ であり、これは有限射影
% P01-E0816: the frozen proof says the complex “is” Omega although the definition only gives a quasi-isomorphism.
$S$-加群である。従って、その基底変換も同様である。
ゆえに $S_g$ は $R$ 上滑らかである。

\medskip\noindent
第二の主張も、補題 \ref{algebra-lemma-NL-localize-bottom} から同じように従う。
\end{proof}

\begin{lemma}
\label{algebra-lemma-base-change-smooth}
\begin{slogan}
滑らかさは基底変換で保たれる
% P01-E0815: the frozen slogan lacks terminal punctuation.
\end{slogan}
$R \to S$ を滑らかな環写像とする。
$R \to R'$ を任意の環写像とする。
このとき、基底変換 $R' \to S' = R' \otimes_R S$ は滑らかである。
\end{lemma}

\begin{proof}
核を $I$ とする表示
$\alpha : R[x_1, \ldots, x_n] \to S$ をとる。
$\alpha' : R'[x_1, \ldots, x_n] \to R' \otimes_R S$ を
誘導された表示とし、$I' = \Ker(\alpha')$ とおく。
$0 \to I \to R[x_1, \ldots, x_n] \to S \to 0$ は完全だから、列
$R' \otimes_R I \to R'[x_1, \ldots, x_n] \to R' \otimes_R S \to 0$
は完全である。従って $R' \otimes_R I \to I'$ は全射である。
定義 \ref{algebra-definition-smooth} により短完全列
$$
0 \to I/I^2 \to
\Omega_{R[x_1, \ldots, x_n]/R} \otimes_{R[x_1, \ldots, x_n]} S \to
\Omega_{S/R} \to
0
$$
があり、$S$-加群 $\Omega_{S/R}$ は有限射影である。
特に $I/I^2$ は
$\Omega_{R[x_1, \ldots, x_n]/R} \otimes_{R[x_1, \ldots, x_n]} S$
の直和因子である。可換図式
$$
\xymatrix{
R' \otimes_R (I/I^2) \ar[r] \ar[d] &
R' \otimes_R (\Omega_{R[x_1, \ldots, x_n]/R} \otimes_{R[x_1, \ldots, x_n]} S)
\ar[d] \\
I'/(I')^2 \ar[r] &
\Omega_{R'[x_1, \ldots, x_n]/R'}
\otimes_{R'[x_1, \ldots, x_n]} (R' \otimes_R S)
}
$$
を考える。右の垂直写像は同型なので、上で述べたことから、
左の垂直写像は単射かつ全射であることが分かる。
従って $\NL(\alpha')$ は次数 $0$ に置かれた
$\Omega_{S'/R'} \cong S' \otimes_S \Omega_{S/R}$ に擬同型である。
この加群は有限射影加群の基底変換だから有限射影である。
\end{proof}

\begin{lemma}
\label{algebra-lemma-smooth-over-field}
$k$ を体とする。
$S$ を滑らかな $k$-代数とする。
このとき $S$ は局所完全交叉である。
\end{lemma}

\begin{proof}
補題 \ref{algebra-lemma-base-change-smooth} および
\ref{algebra-lemma-lci-field-change} により、$k$ が代数閉である場合を
証明すれば十分である。核を $I$ とする表示
$\alpha : k[x_1, \ldots, x_n] \to S$ を選ぶ。
$\mathfrak m$ を $S$ の極大イデアルとし、
$\mathfrak m' \supset I$ を $k[x_1, \ldots, x_n]$ の対応する
極大イデアルとする。補題 \ref{algebra-lemma-lci} の条件 (5) が
（$\mathfrak q$ の代わりに $\mathfrak m$ を用いて）成り立つことを示す。
$k$ は代数閉なので、定理 \ref{algebra-theorem-nullstellensatz} により、
ある $a_i \in k$ に対して
$\mathfrak m' = (x_1 - a_1, \ldots, x_n - a_n)$ と書ける。
$k \to S$ が滑らかであるという仮定により、$S$-加群写像
$\text{d} : I/I^2 \to \bigoplus_{i = 1}^n S \text{d}x_i$
は分裂単射である。従って、対応する写像
$I/\mathfrak m' I \to \bigoplus \kappa(\mathfrak m') \text{d}x_i$
は単射である。
$\dim_{\kappa(\mathfrak m')}(I/\mathfrak m' I) = c$ とおき、
$I/\mathfrak m' I$ の $\kappa(\mathfrak m')$-基底に写る
$f_1, \ldots, f_c \in I$ を選ぶ。中山の補題 \ref{algebra-lemma-NAK} により、
$f_1, \ldots, f_c$ は
$k[x_1, \ldots, x_n]_{\mathfrak m'}$ 上で
$I_{\mathfrak m'}$ を生成する。可換図式
$$
\xymatrix{
I \ar[r] \ar[d] & I/I^2 \ar[rr] \ar[d] & &
I/\mathfrak m'I \ar[d] \\
\Omega_{k[x_1, \ldots, x_n]/k} \ar[r] &
\bigoplus S\text{d}x_i \ar[rr]^{\text{d}x_i \mapsto x_i - a_i} & &
\mathfrak m'/(\mathfrak m')^2
}
$$
を考える（可換性の証明は省略する）。中央の垂直写像は、
$\alpha$ の素朴な余接複体を定める写像である。右下の水平矢印は同型
$\bigoplus \kappa(\mathfrak m') \text{d}x_i \to
\mathfrak m'/(\mathfrak m')^2$ を誘導することに注意する。
従って $I_{\mathfrak m'}$ の生成元 $f_1, \ldots, f_c$ は、
$k[x_1, \ldots, x_n]_{\mathfrak m'}$ の元の族であって、その
$\mathfrak m'/(\mathfrak m')^2$ における類が
$\kappa(\mathfrak m')$ 上線形独立であるような族に写る。
従って、補題 \ref{algebra-lemma-regular-ring-CM} により、それらは環
$k[x_1, \ldots, x_n]_{\mathfrak m'}$ の正則列をなす。
これは補題 \ref{algebra-lemma-lci} の条件 (5) を確かめるので、
% P01-E0817: the frozen citation runs directly into “hence” without punctuation.
ある $g \in S$, $g \not \in \mathfrak m$ に対して $S_g$ は
$k$ 上の大域的完全交叉である。これは $S$ のどの極大イデアルに
対しても成り立つから、$S$ は $k$ 上の局所完全交叉である。
\end{proof}

\begin{definition}
\label{algebra-definition-standard-smooth}
$R$ を環とする。整数 $n \geq c \geq 0$ および
$f_1, \ldots, f_c \in R[x_1, \ldots, x_n]$ が与えられたとき、
% P01-E0818: the frozen fronted data lack a comma before the main clause.
$$
R[x_1, \ldots, x_n]/(f_1, \ldots, f_c)
$$
が $R$ 上の {\it 標準滑らかな代数}であるとは、多項式
$$
g =
\det
\left(
\begin{matrix}
\partial f_1/\partial x_1 &
\partial f_2/\partial x_1 &
\ldots &
\partial f_c/\partial x_1 \\
\partial f_1/\partial x_2 &
\partial f_2/\partial x_2 &
\ldots &
\partial f_c/\partial x_2 \\
\ldots & \ldots & \ldots & \ldots \\
\partial f_1/\partial x_c &
\partial f_2/\partial x_c &
\ldots &
\partial f_c/\partial x_c
\end{matrix}
\right)
$$
の $R[x_1, \ldots, x_n]/(f_1, \ldots, f_c)$ における像が
可逆であることをいう。ある $n \geq c \geq 0$ および
$f_1, \ldots, f_c \in R[x_1, \ldots, x_n]$ が存在して、
$R[x_1, \ldots, x_n]/(f_1, \ldots, f_c)$ が $R$ 上の
標準滑らかな代数であり、かつ $S$ が $R$-代数として
$R[x_1, \ldots, x_n]/(f_1, \ldots, f_c)$ と同型であるとき、
$R$-代数 $S$ は {\it 標準滑らか}である、または環写像
$R \to S$ は {\it 標準滑らか}であるという。
\end{definition}

\begin{lemma}
\label{algebra-lemma-standard-smooth}
$S = R[x_1, \ldots, x_n]/(f_1, \ldots, f_c) = R[x_1, \ldots, x_n]/I$
を標準滑らかな代数とする。このとき
\begin{enumerate}
\item 環写像 $R \to S$ は滑らかである。
\item $S$-加群 $\Omega_{S/R}$ は
$\text{d}x_{c + 1}, \ldots, \text{d}x_n$ を基底とする自由加群である。
\item $S$-加群 $I/I^2$ は $f_1, \ldots, f_c$ の類を基底とする
自由加群である。
\item 任意の $g \in S$ に対して、環写像 $R \to S_g$ は標準滑らかである。
\item 任意の環写像 $R \to R'$ に対して、基底変換
$R' \to R'\otimes_R S$ は標準滑らかである。
\item $f \in R$ の像が $S$ で可逆ならば、$R_f \to S$ は
標準滑らかである。
\item 環 $S$ は $R$ 上の相対的大域的完全交叉である。
\end{enumerate}
\end{lemma}

\begin{proof}
与えられた表示の素朴な余接複体
$$
(f_1, \ldots, f_c)/(f_1, \ldots, f_c)^2
\longrightarrow
\bigoplus\nolimits_{i = 1}^n S \text{d}x_i
$$
を考える。
% P01-E0819: the frozen display lacks terminal punctuation before the next sentence.
この写像と、直和の最初の $c$ 個の直和因子への射影とを合成する。
標準滑らかな代数の定義により、類
$f_i \bmod (f_1, \ldots, f_c)^2$ は
$\bigoplus_{i = 1}^c S\text{d}x_i$ の基底に写る。
従って $(f_1, \ldots, f_c)/(f_1, \ldots, f_c)^2$ は階数 $c$ の
自由加群で、その基底は元 $f_i \bmod (f_1, \ldots, f_c)^2$ により
与えられる。また、素朴な余接複体の次数 $0$ のホモロジー、すなわち
$\Omega_{S/R}$ は、$\text{d}x_{c + j}$ の像
（$j = 1, \ldots, n - c$）を基底とする自由 $S$-加群である。
特に、これで $R \to S$ が滑らかであることが示された。

\medskip\noindent
(4) と (6) の証明は省略する。ただし、下の例および補題
\ref{algebra-lemma-base-change-relative-global-complete-intersection} の証明を参照せよ。

\medskip\noindent
$\varphi : R \to R'$ を任意の環写像とする。
$S' = R'[x_1, \ldots, x_n]/(f_1^\varphi, \ldots, f_c^\varphi)$ と書く。
% P01-E0820: the frozen designation “Denote S' = ... where ...” omits by/as.
ここで $f^\varphi$ は、$f \in R[x_1, \ldots, x_n]$ のすべての係数に
$\varphi$ を適用して得られる多項式である。このとき
$S' \cong R' \otimes_R S$ である。さらに、定義
\ref{algebra-definition-standard-smooth} に現れる $S'/R'$ の行列式は
$g^\varphi$ に等しい。従ってその $S'$ における像は、
$R[x_1, \ldots, x_n] \to S \to S'$ による $g$ の像であり、ゆえに
可逆である。これで (5) が示された。

\medskip\noindent
(7) を証明するには、任意の素イデアル $\mathfrak p \subset R$ に対して
$S \otimes_R \kappa(\mathfrak p)$ の次元が $n - c$ であることを
示せば十分である。(5) により、体 $k$ 上の任意の標準滑らかな代数
$k[x_1, \ldots, x_n]/(f_1, \ldots, f_c)$ の次元が $n - c$ であることを
示せば十分である。補題 \ref{algebra-lemma-smooth-over-field} により、
$k[x_1, \ldots, x_n]/(f_1, \ldots, f_c)$ は局所完全交叉であることが
すでに分かっている。従って、$I/I^2$ は階数 $c$ の自由加群だから、
例えば補題 \ref{algebra-lemma-lci} により、
$k[x_1, \ldots, x_n]/(f_1, \ldots, f_c)$ の次元は $n - c$ である。
\end{proof}

\begin{example}
\label{algebra-example-make-standard-smooth}
$R$ を環とする。
$f_1, \ldots, f_c \in R[x_1, \ldots, x_n]$ とする。
$$
h =
\det
\left(
\begin{matrix}
\partial f_1/\partial x_1 &
\partial f_2/\partial x_1 &
\ldots &
\partial f_c/\partial x_1 \\
\partial f_1/\partial x_2 &
\partial f_2/\partial x_2 &
\ldots &
\partial f_c/\partial x_2 \\
\ldots & \ldots & \ldots & \ldots \\
\partial f_1/\partial x_c &
\partial f_2/\partial x_c &
\ldots &
\partial f_c/\partial x_c
\end{matrix}
\right).
$$
とおく。
$S = R[x_1, \ldots, x_{n + 1}]/(f_1, \ldots, f_c, x_{n + 1}h - 1)$
とおく。これは標準滑らかな代数の例である。ただし、表示が誤っていて、
変数は次の順序でなければならない：
$x_1, \ldots, x_c, x_{n + 1}, x_{c + 1}, \ldots, x_n$。
\end{example}

\begin{lemma}
\label{algebra-lemma-compose-standard-smooth}
標準滑らかな環写像の合成は標準滑らかである。
\end{lemma}

\begin{proof}
$R \to S$ と $S \to S'$ が標準滑らかであると仮定する。表示
$S =  R[x_1, \ldots, x_n]/(f_1, \ldots, f_c)$
および
$S' = S[y_1, \ldots, y_m]/(g_1, \ldots, g_d)$
を選ぶ。$g_j$ に写る元
$g_j' \in R[x_1, \ldots, x_n, y_1, \ldots, y_m]$ を選ぶ。
こうして
$S' = R[x_1, \ldots, x_n, y_1, \ldots, y_m]/
(f_1, \ldots, f_c, g'_1, \ldots, g'_d)$
であることが分かる。
% P01-E0821: the frozen source writes the same lifted polynomial family as g_j' and g'_j.
$S'$ が標準滑らかであることを示すには、行列式
$$
\det
\left(
\begin{matrix}
\partial f_1/\partial x_1 &
\ldots &
\partial f_c/\partial x_1 &
\partial g_1/\partial x_1 &
\ldots &
\partial g_d/\partial x_1 \\
\ldots &
\ldots &
\ldots &
\ldots &
\ldots &
\ldots  \\
\partial f_1/\partial x_c &
\ldots &
\partial f_c/\partial x_c &
\partial g_1/\partial x_c &
\ldots &
\partial g_d/\partial x_c \\
0 &
\ldots &
0 &
\partial g_1/\partial y_1 &
\ldots &
\partial g_d/\partial y_1 \\
\ldots &
\ldots &
\ldots &
\ldots &
\ldots &
\ldots  \\
0 &
\ldots &
0 &
\partial g_1/\partial y_d &
\ldots &
\partial g_d/\partial y_d
\end{matrix}
\right)
$$
が $S'$ で可逆であることを確かめれば十分である。
これは、仮定により可逆な二つの行列式の積だから明らかである。
\end{proof}

\begin{lemma}
\label{algebra-lemma-smooth-syntomic}
$R \to S$ を滑らかな環写像とする。
$\Spec(S)$ には標準開集合 $D(g)$ による開被覆で、各 $S_g$ が
$R$ 上標準滑らかとなるものが存在する。特に $R \to S$ は syntomic である。
\end{lemma}

\begin{proof}
核を $I = (f_1, \ldots, f_m)$ とする表示
$\alpha : R[x_1, \ldots, x_n] \to S$ を選ぶ。各部分集合
$E \subset \{1, \ldots, m\}$ に対して、類 $f_e, e\in E$ が有限射影
$S$-加群 $I/I^2$ を自由に生成する開集合 $U_E$ を考える。
補題 \ref{algebra-lemma-cokernel-flat} を参照せよ。
$\Spec(S)$ を、それぞれがいずれかの $U_E$ に全体として含まれる
標準開集合 $D(g)$ で被覆できる。そのような $g$ に対して、表示
$$
\beta : R[x_1, \ldots, x_n, x_{n + 1}] \longrightarrow S_g
$$
で $x_{n + 1}$ を $1/g$ に写すものを考える。
$J = \Ker(\beta)$ とおき、補題 \ref{algebra-lemma-principal-localization-NL}
を用いると、$J/J^2 \cong (I/I^2)_g \oplus S_g$ は自由である。
$S$ を $S_g$ に置き換えてよいし、そうする。すると補題
\ref{algebra-lemma-huber} により、核が $I = (f_1, \ldots, f_c)$ で、
$I/I^2$ が $f_1, \ldots, f_c$ の類を自由基底とするような表示
$\alpha : R[x_1, \ldots, x_n] \to S$ があると仮定してよい。

\medskip\noindent
前段の終わりに得られた表示 $\alpha$ を用いて、
$\Omega_{R[x_1, \ldots, x_n]/R}$ の基底元
$\text{d}x_1, \ldots, \text{d}x_n$ に関して、ほぼ同じ議論を繰り返す。
すなわち、濃度が $c$ の任意の部分集合
$E \subset \{1, \ldots, n\}$ に対し、$\NL(\alpha)$ の微分と射影との合成
$$
S^{\oplus c} \cong I/I^2
\longrightarrow
\Omega_{R[x_1, \ldots, x_n]/R} \otimes_{R[x_1, \ldots, x_n]} S
\longrightarrow
\bigoplus\nolimits_{i \in E} S\text{d}x_i
$$
が同型となる $\Spec(S)$ の開集合 $U_E$ を考える。
再び、$\Spec(S)$ を（有限個の）標準開集合 $D(g)$ で、各 $D(g)$ が
いずれかの $U_E$ に全体として含まれるように被覆できる。
番号を付け直して、$E = \{1, \ldots, c\}$ と仮定してよい。
$D(g) \subset U_E$ を満たす $g$ に対して、表示
$$
\beta : R[x_1, \ldots, x_n, x_{n + 1}] \to S_g
$$
で $x_{n + 1}$ を $1/g$ に写すものを考える。
$J = \Ker(\beta)$ とおく。補題 \ref{algebra-lemma-principal-localization-NL}
から、$\alpha(f) = g$ となる元に対して
$J = (f_1, \ldots, f_c, fx_{n + 1} - 1)$ であり、合成
$$
J/J^2 \longrightarrow
\Omega_{R[x_1, \ldots, x_{n + 1}]/R} \otimes_{R[x_1, \ldots, x_{n + 1}]} S_g
\longrightarrow
\bigoplus\nolimits_{i = 1}^c S_g\text{d}x_i \oplus S_g \text{d}x_{n + 1}
$$
が同型であることが分かる。座標を
$x_1, \ldots, x_c, x_{n + 1}, x_{c + 1}, \ldots, x_n$
の順に並べ直すと、望み通り $S_g$ は $R$ 上標準滑らかである。

\medskip\noindent
標準滑らかな代数は syntomic であり
（補題 \ref{algebra-lemma-standard-smooth} および
\ref{algebra-lemma-relative-global-complete-intersection}）、また $R$ 上 syntomic
であることは $S$ 上局所的だから（補題 \ref{algebra-lemma-local-syntomic}）、
これで証明が完了する。
\end{proof}

\begin{definition}
\label{algebra-definition-smooth-at-prime}
$R \to S$ を環写像とする。
$\mathfrak q$ を $S$ の素イデアルとする。
$g \in S$, $g \not \in \mathfrak q$ で $R \to S_g$ が滑らかとなるものが
存在するとき、$R \to S$ は {\it $\mathfrak q$ において滑らか}であるという。
\end{definition}

\noindent
有限表示の環写像については、これは次のように特徴づけられる。

\begin{lemma}
\label{algebra-lemma-smooth-at-point}
$R \to S$ を有限表示とし、$\mathfrak q$ を $S$ の素イデアルとする。
次は同値である：
% P01-E0822: the frozen statement lacks a colon before the enumeration.
\begin{enumerate}
\item $R \to S$ は $\mathfrak q$ において滑らかである。
\item $H_1(L_{S/R})_\mathfrak q = 0$ であり、
$\Omega_{S/R, \mathfrak q}$ は有限自由 $S_\mathfrak q$-加群である。
\item $H_1(L_{S/R})_\mathfrak q = 0$ であり、
$\Omega_{S/R, \mathfrak q}$ は射影 $S_\mathfrak q$-加群である。
\item $H_1(L_{S/R})_\mathfrak q = 0$ であり、
$\Omega_{S/R, \mathfrak q}$ は平坦 $S_\mathfrak q$-加群である。
% P01-E0823: the frozen criterion switches from the section's NL notation to unexplained L notation.
\end{enumerate}
\end{lemma}

\begin{proof}
素朴な余接複体の形成が局所化と可換することを、以下では断りなく用いる。
第 \ref{algebra-section-netherlander} 節、特に補題 \ref{algebra-lemma-localize-NL} を参照せよ。
$\Omega_{S/R}$ は有限表示 $S$-加群であることに注意する。
補題 \ref{algebra-lemma-differentials-finitely-presented} を参照せよ。
従って補題 \ref{algebra-lemma-finite-projective} により (2)、(3)、(4) は同値である。
(1) がこれら同値な条件 (2)、(3)、(4) を含意することは明らかである。
(2) が成り立つと仮定する。$S_\mathfrak q$ を主局所化の余極限として書き、
補題 \ref{algebra-lemma-colimit-category-fp-modules} を用いると、
$(\Omega_{S/R})_g$ が有限自由となるような
$g \in S$, $g \not \in \mathfrak q$ を見つけられる。
核を $I$ とする表示 $\alpha : R[x_1, \ldots, x_n] \to S$ を選ぶ。
補題 \ref{algebra-lemma-NL-homotopy} により、$\NL_{S/R}$ の代わりに
$\NL(\alpha)$ を用いてよい。全射
$$
\Omega_{R[x_1, \ldots, x_n]/R} \otimes_{R[x_1, \ldots, x_n]} S
\to \Omega_{S/R} \to 0
$$
は、$(\Omega_{S/R})_g$ が射影だから、$g$ を可逆にした後で右逆をもつ。
従って
$\text{d} : (I/I^2)_g\allowbreak \to
\Omega_{R[x_1, \ldots,\allowbreak x_n]/R}\allowbreak \otimes_{R[x_1, \ldots,\allowbreak x_n]} S_g$
の像は直和因子であり、この写像も右逆をもつ。
% P01-E0824: the frozen proof gives d itself a right inverse; only the surjection onto Im(d) is justified to split.
従って $H_1(L_{S/R})_g$ は $(I/I^2)_g$ の商である。
特に $H_1(L_{S/R})_g$ は有限 $S_g$-加群である。
ゆえに $H_1(L_{S/R})_{\mathfrak q}$ が消えるならば、ある
$g' \in S$, $g' \not \in \mathfrak q$ に対して
$H_1(L_{S/R})_{gg'}$ が消える。このとき定義により
$R \to S_{gg'}$ は滑らかである。
\end{proof}

\begin{lemma}
\label{algebra-lemma-locally-smooth}
\begin{slogan}
環写像が滑らかであるための必要十分条件は、対象のすべての素イデアルにおいて滑らかなことである
% P01-E0825: the frozen slogan lacks terminal punctuation.
\end{slogan}
$R \to S$ を環写像とする。このとき $R \to S$ が滑らかであるための
必要十分条件は、$R \to S$ が $S$ のすべての素イデアル $\mathfrak q$ において
滑らかなことである。
\end{lemma}

\begin{proof}
順方向は自明である。$R \to S$ が $S$ のすべての素イデアル
$\mathfrak q$ において滑らかであると仮定する。補題
\ref{algebra-lemma-quasi-compact} により $\Spec(S)$ は準コンパクトだから、各
$S_{g_i}$ が滑らかとなる有限被覆
$\Spec(S) = \bigcup D(g_i)$ が存在する。補題
\ref{algebra-lemma-cover-upstairs} により、これは $S$ が $R$ 上有限表示であることを
含意する。補題 \ref{algebra-lemma-localize-NL} によれば、
$\NL_{S/R} \otimes_S S_{g_i}$ は次数 $0$ に置かれた有限射影
$S_{g_i}$-加群に擬同型である。補題 \ref{algebra-lemma-finite-projective} により、
これは $\NL_{S/R}$ が次数 $0$ に置かれた有限射影 $S$-加群に
擬同型であることを含意する。
\end{proof}

\begin{lemma}
\label{algebra-lemma-compose-smooth}
滑らかな環写像の合成は滑らかである。
\end{lemma}

\begin{proof}
これは多くの異なる方法で証明できる。一つは蛇の補題
（補題 \ref{algebra-lemma-snake}）と Jacobi--Zariski 完全列
（補題 \ref{algebra-lemma-exact-sequence-NL}）を、射影加群は自由加群の
直和因子であるという特徴づけ（補題 \ref{algebra-lemma-characterize-projective}）と
組み合わせる方法である。別の証明は、補題
\ref{algebra-lemma-smooth-syntomic}、\ref{algebra-lemma-compose-standard-smooth}、
\ref{algebra-lemma-locally-smooth} を組み合わせることで得られる。
\end{proof}

\begin{lemma}
\label{algebra-lemma-product-smooth}
$R$ を環とする。$S = S' \times S''$ を $R$-代数の積とする。
このとき $S$ が $R$ 上滑らかであるための必要十分条件は、
$S'$ と $S''$ がともに $R$ 上滑らかなことである。
\end{lemma}

\begin{proof}
省略する。ヒント：補題 \ref{algebra-lemma-locally-smooth} により、滑らかさは
素イデアルごとに確かめられる。補題 \ref{algebra-lemma-spec-product} により
$\Spec(S)$ は $\Spec(S')$ と $\Spec(S'')$ の非交和なので、
$R \to S$ が $\mathfrak q$ において滑らかであることは、
$R \to S'$ が対応する素イデアルにおいて滑らかであること、または
$R \to S''$ が対応する素イデアルにおいて滑らかであることに対応する。
\end{proof}

\begin{lemma}
\label{algebra-lemma-relative-global-complete-intersection-smooth}
$R$ を環とする。
$S = R[x_1, \ldots, x_n]/(f_1, \ldots, f_c)$ を相対的大域的完全交叉とする。
$\mathfrak q \subset S$ を素イデアルとする。このとき $R \to S$ が
$\mathfrak q$ において滑らかであるための必要十分条件は、濃度が $c$ の
部分集合 $I \subset \{1, \ldots, n\}$ で、多項式
$$
g_I = \det (\partial f_j/\partial x_i)_{j = 1, \ldots, c, \ i \in I}.
$$
が $\mathfrak q$ の元に写らないものが存在することである。
% P01-E0826: the frozen display already ends with a period, leaving the following predicate as a fragment.
\end{lemma}

\begin{proof}
補題 \ref{algebra-lemma-relative-global-complete-intersection-conormal} により、
$S$ の与えられた表示に付随する素朴な余接複体は
$$
\bigoplus\nolimits_{j = 1}^c S \cdot f_j
\longrightarrow
\bigoplus\nolimits_{i = 1}^n S \cdot \text{d}x_i, \quad
f_j \longmapsto \sum \frac{\partial f_j}{\partial x_i} \text{d}x_i.
$$
である。
% P01-E0827: the frozen summation has no index range although the target direct sum runs from i=1 to n.
この写像を与える行列の極大小行列式は、ちょうど多項式 $g_I$ である。

\medskip\noindent
$g_I$ が $g \in S$, $g \not \in \mathfrak q$ に写ると仮定する。
このとき代数 $S_g$ は $R$ 上滑らかである。実際、補題
\ref{algebra-lemma-localize-NL} により、その素朴な余接複体は、上の複体を
$g$ で局所化したものに擬同型である。また構成により、それは次数 $0$
に置かれた自由な階数 $n - c$ の加群に擬同型である。
% P01-E0828: the frozen phrase “a free rank n-c module” is nonidiomatic.

\medskip\noindent
逆に、すべての $g_I$ が $\mathfrak q$ に入ると仮定する。
この場合、上の複体を $\kappa(\mathfrak q)$ とテンソルしたものは
極大階数をもたない。従って、$g \in S$, $g \not \in \mathfrak q$ で、
それによる局所化の後にこの写像が分裂単射となるものは存在しない。
再び補題 \ref{algebra-lemma-localize-NL} により、$R$ 上滑らかとなるそのような
局所化は存在しない。
\end{proof}

\begin{lemma}
\label{algebra-lemma-flat-fibre-smooth}
$R \to S$ を環写像とする。
$\mathfrak q \subset S$ を、$R$ の素イデアル $\mathfrak p$ の上にある
素イデアルとする。次を仮定する。
\begin{enumerate}
\item $g \in S$, $g \not\in \mathfrak q$ で、$R \to S_g$ が有限表示となる
ものが存在する。
\item 局所環準同型 $R_{\mathfrak p} \to S_{\mathfrak q}$ は平坦である。
\item ファイバー $S \otimes_R \kappa(\mathfrak p)$ は、
$\mathfrak q$ に対応する素イデアルにおいて
$\kappa(\mathfrak p)$ 上滑らかである。
\end{enumerate}
このとき $R \to S$ は $\mathfrak q$ において滑らかである。
\end{lemma}

\begin{proof}
補題 \ref{algebra-lemma-syntomic} および \ref{algebra-lemma-smooth-over-field} により、
$S_g$ が相対的大域的完全交叉となるような $g \in S$ が存在する。
% P01-E0829: the frozen application omits the required condition g notin q before replacing S by S_g.
$S$ を $S_g$ で置き換え、
$S = R[x_1, \ldots, x_n]/(f_1, \ldots, f_c)$ が相対的大域的完全交叉で
あると仮定してよい。濃度が $c$ の任意の部分集合
$I \subset \{1, \ldots, n\}$ に対し、補題
\ref{algebra-lemma-relative-global-complete-intersection-smooth} の多項式
$g_I = \det (\partial f_j/\partial x_i)_{j = 1, \ldots, c, i \in I}$
を考える。$g_I$ の多項式環
$\kappa(\mathfrak p)[x_1, \ldots, x_n]$ における像
$\overline{g}_I$ は、$f_j$ の像 $\overline{f}_j$ の環
$\kappa(\mathfrak p)[x_1, \ldots, x_n]$ における偏導関数の行列式である。
従って、補題 \ref{algebra-lemma-relative-global-complete-intersection-smooth} を
$R \to S$ と
$\kappa(\mathfrak p) \to S \otimes_R \kappa(\mathfrak p)$ の両方に適用すれば、
本補題が従う。
\end{proof}

\noindent
次の補題に現れる集合 $U, V$ は定義により開であることに注意する。

\begin{lemma}
\label{algebra-lemma-flat-base-change-locus-smooth}
$R \to S$ を有限表示の環写像とする。
$R \to R'$ を平坦な環写像とする。
$S' = R' \otimes_R S$ を基底変換と記す。
% P01-E0830: the frozen “Denote S'=... the base change” construction omits by/as.
$U \subset \Spec(S)$ を $R \to S$ が滑らかである素イデアルの集合とする。
$V \subset \Spec(S')$ を $R' \to S'$ が滑らかである素イデアルの集合とする。
% P01-E0831: the frozen “Let V ... the set” sentence omits the verb “be”.
このとき、写像 $f : \Spec(S') \to \Spec(S)$ のもとで、$V$ は $U$ の逆像である。
\end{lemma}

\begin{proof}
補題 \ref{algebra-lemma-change-base-NL} により、
$\NL_{S/R} \otimes_S S'$ は $\NL_{S'/R'}$ とホモトピー同値である。
これは既に $f^{-1}(U) \subset V$ を含意する。

\medskip\noindent
$\mathfrak q' \subset S'$ を $\mathfrak q \subset S$ の上にある
素イデアルとし、$\mathfrak q' \in V$ と仮定する。
$\mathfrak q \in U$ を示さなければならない。
$S \to S'$ は平坦だから、補題 \ref{algebra-lemma-local-flat-ff} により
$S_{\mathfrak q} \to S'_{\mathfrak q'}$ は忠実平坦である。
従って $H_1(L_{S'/R'})_{\mathfrak q'}$ の消滅は
$H_1(L_{S/R})_{\mathfrak q}$ の消滅を含意する。
% P01-E0832: the frozen proof again switches from NL to the unexplained L notation.
補題 \ref{algebra-lemma-finite-projective-descends} を $S_{\mathfrak q}$-加群
$(\Omega_{S/R})_{\mathfrak q}$ と写像
$S_{\mathfrak q} \to S'_{\mathfrak q'}$ に適用すると、
$(\Omega_{S/R})_{\mathfrak q}$ は射影であることが分かる。
従って補題 \ref{algebra-lemma-smooth-at-point} により、$R \to S$ は
$\mathfrak q$ において滑らかである。
\end{proof}

\begin{lemma}
\label{algebra-lemma-smooth-field-change-local}
$K/k$ を体拡大とする。
$S$ を $k$ 上有限型の代数とする。
$\mathfrak q_K$ を $S_K = K \otimes_k S$ の素イデアルとし、
$\mathfrak q$ を $S$ の対応する素イデアルとする。
このとき $S$ が $k$ 上 $\mathfrak q$ において滑らかであるための
必要十分条件は、$S_K$ が $K$ 上 $\mathfrak q_K$ において
滑らかなことである。
\end{lemma}

\begin{proof}
これは補題 \ref{algebra-lemma-flat-base-change-locus-smooth} の特別な場合である。
\end{proof}

\begin{lemma}
\label{algebra-lemma-lift-smooth}
$R$ を環、$I \subset R$ をイデアルとする。
$R/I \to \overline{S}$ を滑らかな環写像とする。
$\overline{S}$ の単位イデアルを生成する元
$\overline{g}_i \in \overline{S}$ であって、各
$\overline{S}_{g_i} \cong S_i/IS_i$ が $R$ 上のある（標準）滑らかな環
$S_i$ によって書けるものが存在する。
% P01-E0833: the frozen “there exists elements” has subject-verb disagreement.
% P01-E0834: the frozen localization drops the bar from the only introduced localizer, both here and in the proof.
\end{lemma}

\begin{proof}
補題 \ref{algebra-lemma-smooth-syntomic} により、$\overline{S}$ の単位イデアルを
生成する元の族 $\overline{g}_i \in \overline{S}$ であって、各
$\overline{S}_{g_i}$ が $R/I$ 上標準滑らかとなるものを得る。
従って $\overline{S}$ は $R/I$ 上標準滑らかであると仮定してよい。
定義 \ref{algebra-definition-standard-smooth} のように
$\overline{S} =
(R/I)[x_1, \ldots, x_n]/(\overline{f}_1, \ldots, \overline{f}_c)$
と書く。$\overline{f}_1, \ldots, \overline{f}_c$ の持ち上げ
$f_1, \ldots, f_c \in R[x_1, \ldots, x_n]$ を選ぶ。
$S = R[x_1, \ldots, x_n, x_{n + 1}]/(f_1, \ldots, f_c, x_{n + 1}\Delta - 1)$
とおく。ここで、例 \ref{algebra-example-make-standard-smooth} と同様に
$\Delta = \det(\frac{\partial f_j}{\partial x_i})_{i, j = 1, \ldots, c}$
である。これで補題が証明された。
\end{proof}

\section{形式的に滑らかな写像}
\label{algebra-section-formally-smooth}

\noindent
この節では、形式的に滑らかな環写像を定義する。有限表示の環写像が
形式的に滑らかであるための必要十分条件は滑らかであることであると後に分かる。
命題 \ref{algebra-proposition-smooth-formally-smooth} を参照せよ。

\begin{definition}
\label{algebra-definition-formally-smooth}
$R \to S$ を環写像とする。
任意の、実線部分が可換な図式
$$
\xymatrix{
S \ar[r] \ar@{-->}[rd] & A/I \\
R \ar[r] \ar[u] & A \ar[u]
}
$$
であって、$I \subset A$ が平方零イデアルであるものに対し、
図式を可換にする点線矢印が存在するとき、$S$ は
\textit{$R$ 上形式的に滑らか}であるという。
\end{definition}

\begin{lemma}
\label{algebra-lemma-base-change-fs}
$R \to S$ を形式的に滑らかな環写像とする。
$R \to R'$ を任意の環写像とする。
このとき、基底変換 $S' = R' \otimes_R S$ は $R'$ 上形式的に滑らかである。
\end{lemma}

\begin{proof}
定義 \ref{algebra-definition-formally-smooth} にあるような、実線部分が可換な図式
$$
\xymatrix{
S \ar[r] \ar@{-->}[rrd] & R' \otimes_R S \ar[r] \ar@{-->}[rd] & A/I \\
R  \ar[u] \ar[r] & R' \ar[r] \ar[u] & A \ar[u]
}
$$
が与えられたとする。仮定により長い方の点線矢印が存在する。
テンソル積の普遍性から短い方の点線矢印を得る。
\end{proof}

\begin{lemma}
\label{algebra-lemma-compose-formally-smooth}
形式的に滑らかな環写像の合成は形式的に滑らかである。
\end{lemma}

\begin{proof}
省略する。（ヒント：これは完全に形式的であり、適切な図式を考えれば従う。）
\end{proof}

\begin{lemma}
\label{algebra-lemma-polynomial-ring-formally-smooth}
$R$ 上の多項式環は $R$ 上形式的に滑らかである。
\end{lemma}

\begin{proof}
$S = R[x_j; j \in J]$ として、定義 \ref{algebra-definition-formally-smooth} にあるような
図式が与えられたとする。このとき、上側の水平矢印による $x_j$ の像である
$A/I$ の元を元 $a_j \in A$ に持ち上げるだけで、点線矢印が得られる。
% P01-E0835: the frozen relative clause redundantly ends with a second “to”.
\end{proof}

\begin{lemma}
\label{algebra-lemma-characterize-formally-smooth}
$R \to S$ を環写像とする。$P \to S$ を、多項式環 $P$ から $S$ への
全射な $R$-代数写像とする。その核を $J \subset P$ とする。
% P01-E0836: the frozen designation omits “by/as”.
このとき、$R \to S$ が形式的に滑らかであるための必要十分条件は、
全射 $P/J^2 \to S$ の右逆となる $R$-代数写像
$\sigma : S \to P/J^2$ が存在することである。
\end{lemma}

\begin{proof}
$R \to S$ が形式的に滑らかであると仮定する。
可換図式
$$
\xymatrix{
S \ar[r] \ar@{-->}[rd] & P/J \\
R \ar[r] \ar[u] &  P/J^2\ar[u]
}
$$
を考える。仮定により点線矢印が存在する。これは $\sigma$ の存在を示す。

\medskip\noindent
逆に、補題にあるような $\sigma$ が与えられているとする。
定義 \ref{algebra-definition-formally-smooth} にあるような、実線部分が可換な図式
$$
\xymatrix{
S \ar[r] \ar@{-->}[rd] & A/I \\
R \ar[r] \ar[u] & A \ar[u]
}
$$
が与えられたとする。
補題 \ref{algebra-lemma-polynomial-ring-formally-smooth} により $P$ は形式的に滑らかだから、
写像 $P \to S \to A/I$ を持ち上げる $R$-代数準同型
$\psi : P \to A$ が存在する。
明らかに $\psi(J) \subset I$ であり、$I^2 = 0$ だから
$\psi(J^2) = 0$ を得る。従って $\psi$ は
$\overline{\psi} : P/J^2 \to A$ を経由する。求める点線矢印は合成
$\overline{\psi} \circ \sigma : S \to A$ である。
\end{proof}

\begin{remark}
\label{algebra-remark-lemma-characterize-formally-smooth}
補題 \ref{algebra-lemma-characterize-formally-smooth} は、より一般に
$P$ が $R$ 上形式的に滑らかである場合にも成り立つ。
\end{remark}

\begin{lemma}
\label{algebra-lemma-characterize-formally-smooth-again}
$R \to S$ を環写像とする。
$P \to S$ を、多項式環 $P$ から $S$ への全射な
$R$-代数写像とする。その核を $J \subset P$ とする。
% P01-E0837: the frozen repeated designation omits “by/as”.
このとき、$R \to S$ が形式的に滑らかであるための必要十分条件は、
補題 \ref{algebra-lemma-differential-seq} の列
$$
0 \to J/J^2 \to \Omega_{P/R} \otimes_P S \to \Omega_{S/R} \to 0
$$
が分裂完全列であることである。
\end{lemma}

\begin{proof}
$S$ が $R$ 上形式的に滑らかであると仮定する。
補題 \ref{algebra-lemma-characterize-formally-smooth} によれば、これは標準写像
$P/J^2 \to S$ の右逆である $R$-代数写像
$S \to P/J^2$ が存在することを意味する。
補題 \ref{algebra-lemma-differential-mod-power-ideal} により
$\Omega_{P/R} \otimes_P S = \Omega_{(P/J^2)/R} \otimes_{P/J^2} S$
である。補題 \ref{algebra-lemma-differential-seq-split} により、この列は分裂する。

\medskip\noindent
補題の完全列が分裂完全であると仮定する。
分裂 $\sigma : \Omega_{S/R} \to \Omega_{P/R} \otimes_P S$ を選ぶ。
各 $\lambda \in S$ に対し、$\lambda$ に写る $x_\lambda \in P$ を選ぶ。
次に、各 $\lambda \in S$ に対し
$$
\text{d}f_\lambda = \text{d}x_\lambda - \sigma(\text{d}\lambda)
$$
が完全列の中間項で成り立つような $f_\lambda \in J$ を選ぶ。
写像 $s : \lambda \mapsto x_\lambda - f_\lambda \mod J^2$ は
$R$-代数準同型 $s : S \to P/J^2$ であると主張する。
これを示すため、$h \in J$ かつ
$\text{d}h = 0$ が $\Omega_{P/R} \otimes_R S$ で成り立つならば
$h \in J^2$ であることを繰り返し用いる。
% P01-E0838: the frozen middle term tensors over R rather than P.
$\lambda, \mu \in S$ とする。このとき
$\sigma(\text{d}\lambda + \text{d}\mu - \text{d}(\lambda + \mu)) = 0$ である。
これは
$$
\text{d}(x_\lambda + x_\mu - x_{\lambda + \mu}
- f_\lambda - f_\mu + f_{\lambda + \mu}) = 0
$$
を含意する。従って
$x_\lambda + x_\mu - x_{\lambda + \mu}
- f_\lambda - f_\mu + f_{\lambda + \mu} \in J^2$ であり、さらに
$s(\lambda) + s(\mu) = s(\lambda + \mu)$ となる。
同様に
$\sigma(\lambda \text{d}\mu + \mu \text{d}\lambda - \text{d}\lambda \mu) = 0$
である。
% P01-E0839: the frozen last differential lacks parentheses around the product.
これは完全列の中間項において
$$
\mu(\text{d}x_\lambda - \text{d}f_\lambda) +
\lambda(\text{d}x_\mu - \text{d}f_\mu) -
\text{d}x_{\lambda\mu} + \text{d}f_{\lambda\mu} = 0
$$
であることを含意する。また、やはり中間項において
$$
\text{d}(x_\lambda x_\mu) =
x_\lambda \text{d}x_\mu + x_\mu \text{d}x_\lambda =
\lambda \text{d}x_\mu + \mu \text{d} x_\lambda
$$
である。これらを合わせると
$x_\lambda x_\mu - x_{\lambda\mu}
- \mu f_\lambda - \lambda f_\mu + f_{\lambda\mu} \in J^2$ を得る。
$f_\lambda f_\mu \in J^2$ でもあるから、
$(x_\lambda - f_\lambda)(x_\mu - f_\mu) -
(x_{\lambda\mu} - f_{\lambda\mu}) \in J^2$ であり、これは
$s(\lambda)s(\mu) = s(\lambda\mu)$ を意味する。
$\lambda \in R$ ならば $\text{d}\lambda = 0$ であり、
$\text{d}f_\lambda = \text{d}x_\lambda$ であることが分かる。従って
$\lambda - x_\lambda + f_\lambda \in J^2$、ゆえに所望の
$s(\lambda) = \lambda$ を得る。ここで
補題 \ref{algebra-lemma-characterize-formally-smooth} を適用すれば、
$S/R$ が形式的に滑らかであると結論できる。
\end{proof}

\begin{proposition}
\label{algebra-proposition-characterize-formally-smooth}
$R \to S$ を環写像とする。形式的に滑らかな $R$-代数 $P$ と、
核が $J$ である全射 $P \to S$ を考える。次の条件は同値である。
\begin{enumerate}
\item $S$ は $R$ 上形式的に滑らかである、
\item 上のようなある $P \to S$ について $P/J^2 \to S$ の切断が存在する、
\item 上のようなすべての $P \to S$ について $P/J^2 \to S$ の切断が存在する、
\item 上のようなある $P \to S$ について、列
$0 \to J/J^2 \to \Omega_{P/R} \otimes S \to \Omega_{S/R} \to 0$ は分裂完全である、
\item 上のようなすべての $P \to S$ について、列
$0 \to J/J^2 \to \Omega_{P/R} \otimes S \to \Omega_{S/R} \to 0$ は分裂完全である、
% P01-E0840: the frozen two conormal sequences omit the tensor base P.
および
\item 素朴な余接複体 $\NL_{S/R}$ は、次数 $0$ に置かれた射影
$S$-加群と擬同型である。すなわち、$H_1(\NL_{S/R}) = 0$ であり、
$\Omega_{S/R}$ は射影 $S$-加群である。
\end{enumerate}
\end{proposition}

\begin{proof}
（1）が（3）を、（3）が（2）を含意することは明らかである。
補題 \ref{algebra-lemma-characterize-formally-smooth} の証明の前半を参照せよ。
また、（3）は（5）を、（5）は（4）を含意し、（2）も（4）を含意する。
補題 \ref{algebra-lemma-characterize-formally-smooth-again} の証明の前半を参照せよ。
最後に、補題 \ref{algebra-lemma-characterize-formally-smooth-again} を
標準全射 $R[S] \to S$（\ref{algebra-equation-canonical-presentation}）に適用すると、
（1）が（6）を含意することが分かる。

\medskip\noindent
（4）を仮定して（6）を示そう。環写像 $R \to P \to S$ に付随する、
補題 \ref{algebra-lemma-exact-sequence-NL} の列を考える。
上で示した（1）$\Rightarrow$（6）により、
$\NL_{P/R} \otimes_R S$ は次数 $0$ に置かれた
$\Omega_{P/R} \otimes_P S$ と擬同型である。
% P01-E0841: the frozen first base-change complex tensors over R, then over P.
従って $H_1(\NL_{P/R} \otimes_P S) = 0$ である。
$P \to S$ は全射だから、$\NL_{S/P}$ は次数 $1$ に置かれた
$J/J^2$ とホモトピー同値である（補題 \ref{algebra-lemma-NL-surjection}）。
従って完全列
$0 \to H_1(L_{S/R}) \to J/J^2 \to \Omega_{P/R} \otimes_P S \to
\Omega_{S/R} \to 0$ を得る。
仮定から $H_1(L_{S/R}) = 0$ であり、$\Omega_{S/R}$ は射影
$S$-加群であることが分かる。従って（6）が従う。
% P01-E0842: the frozen proof switches twice from NL to unexplained L notation.

\medskip\noindent
最後に、（6）が（1）を含意することを示そう。仮定は、$P = R[S]$ で
$P \to S$ が標準全射（\ref{algebra-equation-canonical-presentation}）であるとき、
複体 $J/J^2 \to \Omega_{P/R} \otimes S$ が次数 $0$ に置かれた
射影 $S$-加群と擬同型であることを意味する。
% P01-E0843: the frozen conormal complex again omits the tensor base P.
従って補題 \ref{algebra-lemma-characterize-formally-smooth-again} により、
$S$ は $R$ 上形式的に滑らかである。
\end{proof}

\begin{lemma}
\label{algebra-lemma-ses-formally-smooth}
$A \to B \to C$ を環写像とする。$B \to C$ は形式的に滑らかであると仮定する。
このとき、補題 \ref{algebra-lemma-exact-sequence-differentials} の列
$$
0 \to \Omega_{B/A} \otimes_B C \to \Omega_{C/A} \to \Omega_{C/B} \to 0
$$
は分裂短完全列である。
\end{lemma}

\begin{proof}
命題 \ref{algebra-proposition-characterize-formally-smooth} と
補題 \ref{algebra-lemma-exact-sequence-NL} から従う。
\end{proof}

\begin{lemma}
\label{algebra-lemma-differential-seq-formally-smooth}
$A \to B \to C$ を環写像とし、$A \to C$ は形式的に滑らか、
$B \to C$ は全射で、その核は $J \subset B$ であるとする。
このとき、補題 \ref{algebra-lemma-differential-seq} の完全列
$$
0 \to J/J^2 \to \Omega_{B/A} \otimes_B C \to \Omega_{C/A} \to 0
$$
は分裂完全である。
\end{lemma}

\begin{proof}
命題 \ref{algebra-proposition-characterize-formally-smooth}、
補題 \ref{algebra-lemma-exact-sequence-NL}、および
補題 \ref{algebra-lemma-differential-seq} から従う。
\end{proof}

\begin{lemma}
\label{algebra-lemma-application-NL-formally-smooth}
$A \to B \to C$ を環写像とする。$A \to C$ は全射（従って
$B \to C$ も全射）であり、$A \to B$ は形式的に滑らかであると仮定する。
$I = \Ker(A \to C)$ および $J = \Ker(B \to C)$ とおく。
% P01-E0844: the frozen designation omits “by/as”.
このとき、補題 \ref{algebra-lemma-application-NL} の列
$$
0 \to I/I^2 \to J/J^2 \to \Omega_{B/A} \otimes_B B/J \to 0
$$
は分裂完全である。
\end{lemma}

\begin{proof}
$A \to B$ は形式的に滑らかだから、$A \to B$ との合成が商写像
$A \to A/I^2$ に等しい環写像 $\sigma : B \to A/I^2$ が存在する。
このとき $\sigma$ は写像 $J/J^2 \to I/I^2$ を誘導し、これは写像
$I/I^2 \to J/J^2$ の逆である。
\end{proof}

\begin{lemma}
\label{algebra-lemma-lift-formal-smoothness}
$R \to S$ を環写像とする。$I \subset R$ をイデアルとする。次を仮定する。
\begin{enumerate}
\item $I^2 = 0$、
\item $R \to S$ は平坦である、および
\item $R/I \to S/IS$ は形式的に滑らかである。
\end{enumerate}
このとき $R \to S$ は形式的に滑らかである。
\end{lemma}

\begin{proof}
（1）、（2）、（3）を仮定する。
核が $J$ である $R$-代数の全射
$P = R[\{x_t\}_{t \in T}] \to S$ をとる。すると
$0 \to J \to P \to S \to 0$ は平坦 $R$-加群の短完全列である。
これは $I \otimes_R S = IS$、$I \otimes_R P = IP$、
$I \otimes_R J = IJ$、および $J \cap IP = IJ$ を含意する。
証明を通して
$$
\Omega_{(S/IS)/(R/I)} = \Omega_{S/R} \otimes_S (S/IS)
= \Omega_{S/R} \otimes_R R/I = \Omega_{S/R} / I\Omega_{S/R}
$$
を用い、$P$ についても同様の等式を用いる
（補題 \ref{algebra-lemma-differentials-base-change} を参照せよ）。
補題 \ref{algebra-lemma-characterize-formally-smooth-again} により、列
\begin{equation}
\label{algebra-equation-split}
0 \to J/(IJ + J^2) \to
\Omega_{P/R} \otimes_P S/IS \to
\Omega_{S/R} \otimes_S S/IS \to 0
\end{equation}
は分裂完全である。もちろん中間項は
$\bigoplus_{t \in T} S/IS \text{d}x_t$ である。分裂
$\sigma : \Omega_{P/R} \otimes_P S/IS \to J/(IJ + J^2)$ を選ぶ。
各 $t \in T$ に対し、$J/(IJ + J^2)$ において
$\sigma(\text{d}x_t)$ に写る元 $f_t \in J$ を選ぶ。これにより、一意な
$S$-加群写像
$$
\tilde \sigma : \Omega_{P/R} \otimes_R S
= \bigoplus S\text{d}x_t \longrightarrow J/J^2
$$
であって、$\tilde\sigma(\text{d}x_t) = f_t$ を満たすものが定まる。
% P01-E0845: the frozen source tensors over R despite the displayed direct-sum identification.
$\sigma$ は $\text{d}$ の切断だから、差
$$
\Delta = \text{id}_{J/J^2} - \tilde \sigma \circ \text{d}
$$
は自己写像 $J/J^2 \to J/J^2$ であり、その像は
$(IJ + J^2)/J^2$ に含まれる。特に $I^2 = 0$ だから
$\Delta((IJ + J^2)/J^2) = 0$ である。これは $\Delta$ が
$$
J/J^2 \to J/(IJ + J^2) \xrightarrow{\overline{\Delta}}
(IJ + J^2)/J^2 \to J/J^2
$$
と経由することを意味する。ここで $\overline{\Delta}$ は
$S/IS$-加群写像である。
列（\ref{algebra-equation-split}）が分裂することを再び用いると、
$\overline{\delta} \circ d$ が $\overline{\Delta}$ に等しくなるような
$S/IS$-加群写像
$\overline{\delta} : \Omega_{P/R} \otimes_P S/IS \to (IJ + J^2)/J^2$
を見つけられる。
% P01-E0846: the frozen differential here is a bare italic d unlike adjacent upright notation.
上と同じ方法で、写像 $\overline{\delta}$ は $S$-加群写像
$\delta : \Omega_{P/R} \otimes_P S \to J/J^2$ を定める。
$\tilde \sigma$ を $\tilde \sigma + \delta$ で置き換えると、
簡単な計算から $\Delta = 0$ を得る。言い換えれば、$\tilde \sigma$ は
$J/J^2 \to \Omega_{P/R} \otimes_P S$ の切断である。
補題 \ref{algebra-lemma-characterize-formally-smooth-again} により、
$R \to S$ は形式的に滑らかである。
\end{proof}

\begin{proposition}
\label{algebra-proposition-smooth-formally-smooth}
$R \to S$ を環写像とする。次の条件は同値である。
% P01-E0847: the frozen introductory clause lacks a colon before the enumeration.
\begin{enumerate}
\item $R \to S$ は有限表示かつ形式的に滑らかである、
\item $R \to S$ は滑らかである。
\end{enumerate}
\end{proposition}

\begin{proof}
命題 \ref{algebra-proposition-characterize-formally-smooth} と
定義 \ref{algebra-definition-smooth} から従う。
（$R \to S$ が有限表示ならば、$\Omega_{S/R}$ は有限表示
$S$-加群であることに注意せよ。補題
\ref{algebra-lemma-differentials-finitely-presented} を参照せよ。）
\end{proof}

\begin{lemma}
\label{algebra-lemma-finite-presentation-fs-Noetherian}
$R \to S$ を滑らかな環写像とする。このとき、$\mathbf{Z}$ 上有限型の
部分環 $R_0 \subset R$ と滑らかな環写像 $R_0 \to S_0$ であって、
$S \cong R \otimes_{R_0} S_0$ を満たすものが存在する。
\end{lemma}

\begin{proof}
滑らかであることが有限表示かつ形式的に滑らかであることと同値であることを
用いる。命題 \ref{algebra-proposition-smooth-formally-smooth} を参照せよ。
$S = R[x_1, \ldots, x_n]/(f_1, \ldots, f_m)$ と書き、
$I = (f_1, \ldots, f_m)$ とおく。
補題 \ref{algebra-lemma-characterize-formally-smooth} にあるように、$S$ への射影の
右逆 $\sigma : S \to R[x_1, \ldots, x_n]/I^2$ を選ぶ。
$\sigma(x_i \bmod I) = h_i \bmod I^2$ を満たす
$h_i \in R[x_1, \ldots, x_n]$ を選ぶ。$x_i - h_i \in I$ だから、
$$
x_i - h_i = \sum\nolimits_j b_{ij} f_j
$$
を満たす $b_{ij} \in R[x_1, \ldots, x_n]$ が存在する。
$\sigma$ が $R$-代数準同型
$R[x_1, \ldots, x_n]/I \to R[x_1, \ldots, x_n]/I^2$ であることは、
ある $a_{kl} \in R[x_1, \ldots, x_n]$ に対して
$$
f_j(h_1, \ldots, h_n) = \sum\nolimits_{j_1 j_2} a_{j_1 j_2} f_{j_1} f_{j_2}
$$
が成り立つことと同値である。
$R_0 \subset R$ を、多項式 $f_j, h_i, a_{kl}, b_{ij}$ のすべての係数により
$\mathbf{Z}$ 上生成される部分環とする。
$S_0 = R_0[x_1, \ldots, x_n]/(f_1, \ldots, f_m)$、
$I_0 = (f_1, \ldots, f_m)$ とおく。二つ目の表示式は
$R_0[x_1, \ldots, x_n]$ で成り立つから、規則
$x_i \mapsto h_i \bmod I_0^2$ により定まる $R_0$-代数写像を
$\sigma_0 : S_0 \to R_0[x_1, \ldots, x_n]/I_0^2$ とすることができる。
一つ目の表示式は $R_0[x_1, \ldots, x_n]$ で成り立つから、
$\sigma_0$ は射影
$R_0[x_1, \ldots, x_n] / I_0^2 \to R_0[x_1, \ldots, x_n] / I_0 = S_0$
の右逆である。従って補題 \ref{algebra-lemma-characterize-formally-smooth} により、
環 $S_0$ は $R_0$ 上形式的に滑らかである。
\end{proof}

\begin{lemma}
\label{algebra-lemma-smooth-descends-through-colimit}
$A = \colim A_i$ を環のフィルター余極限とする。
$A \to B$ を滑らかな環写像とする。このとき、ある $i$ と滑らかな環写像
$A_i \to B_i$ であって、$B = B_i \otimes_{A_i} A$ を満たすものが存在する。
\end{lemma}

\begin{proof}
補題 \ref{algebra-lemma-finite-presentation-fs-Noetherian} から従う。実際、
補題 \ref{algebra-lemma-characterize-finite-presentation} により、$R_0 \to A$ は
ある $i$ について $A_i$ を経由する。
\end{proof}

\begin{lemma}
\label{algebra-lemma-descent-formally-smooth}
$R \to S$ を環写像とし、$R \to R'$ を忠実平坦な環写像とする。
$S' = S \otimes_R R'$ とおく。このとき、$R \to S$ が形式的に滑らかであるための
必要十分条件は $R' \to S'$ が形式的に滑らかであることである。
\end{lemma}

\begin{proof}
$R \to S$ が形式的に滑らかならば、補題 \ref{algebra-lemma-base-change-fs} により
$R' \to S'$ は形式的に滑らかである。
逆を示すため、$R' \to S'$ が形式的に滑らかであると仮定する。
任意の $S$-加群 $N$ に対して
$N \otimes_R R' = N \otimes_S S'$ であることに注意する。
特に $S \to S'$ も忠実平坦である。
多項式環 $P = R[\{x_i\}_{i \in I}]$ と、核が $J$ である
$R$-代数の全射 $P \to S$ を選ぶ。
$P' = P \otimes_R R'$ は $R'$ 上の多項式代数である。
$R \to R'$ は平坦だから、全射 $P' \to S'$ の核 $J'$ は
$J \otimes_R R'$ である。従って、分裂完全列
（補題 \ref{algebra-lemma-characterize-formally-smooth-again} を参照せよ）
$$
0 \to J'/(J')^2 \to \Omega_{P'/R'} \otimes_{P'} S' \to \Omega_{S'/R'} \to 0
$$
は、$S \to S'$ による対応する列
$$
J/J^2 \to \Omega_{P/R} \otimes_P S \to \Omega_{S/R} \to 0
$$
の基底変換である。補題 \ref{algebra-lemma-differential-seq} を参照せよ。
$S \to S'$ は忠実平坦だから、次の二つが従う。（1）この列
（${}'$ の付かないもの）も完全であり、（2）$\Omega_{S/R}$ は射影
$S$-加群である。実際、$\Omega_{S'/R'}$ は自由加群
$\Omega_{P'/R'} \otimes_{P'} S'$ の直和因子として射影であり、上で述べたことから
$\Omega_{S/R} \otimes_S {S'} = \Omega_{S'/R'}$ である。
従って（2）は忠実平坦な環写像による射影性の降下から従う。
定理 \ref{algebra-theorem-ffdescent-projectivity} を参照せよ。
ゆえに列
$0 \to J/J^2 \to \Omega_{P/R} \otimes_P S \to \Omega_{S/R} \to 0$
も完全であり、補題 \ref{algebra-lemma-characterize-formally-smooth-again} をもう一度
適用すれば証明が完了する。
\end{proof}

\noindent
滑らかな環写像は次の強い持ち上げ性を満たすことが分かる。

\begin{lemma}
\label{algebra-lemma-smooth-strong-lift}
$R \to S$ を滑らかな環写像とする。実線部分が可換な図式
$$
\xymatrix{
S \ar[r] \ar@{-->}[rd] & A/I \\
R \ar[r] \ar[u] & A \ar[u]
}
$$
が与えられているとする。ここで $I \subset A$ は局所冪零イデアルとする。
このとき、図式を可換にする点線矢印が存在する。
\end{lemma}

\begin{proof}
補題 \ref{algebra-lemma-finite-presentation-fs-Noetherian} により、図式を可換図式
$$
\xymatrix{
S_0 \ar[r] & S \ar[r] \ar@{-->}[rd] & A/I \\
R_0 \ar[r] \ar[u] & R \ar[r] \ar[u] & A \ar[u]
}
$$
へ拡張できる。ここで $R_0 \to S_0$ は滑らかで、$R_0$ は
$\mathbf{Z}$ 上有限型であり、$S = S_0 \otimes_{R_0} R$ である。
$x_1, \ldots, x_n \in S_0$ を $R_0$ 上の $S_0$ の生成元とする。
$a_1, \ldots, a_n$ を、$A/I$ において $x_1, \ldots, x_n$ と同じ元へ
写る $A$ の元とする。$A_0 \subset A$ を、$R_0$ の像と元
$a_1, \ldots, a_n$ によって生成される部分環とする。
% P01-E0848: the frozen designation of A_0 omits “by/as”.
$I_0 = A_0 \cap I$ とおく。このとき $A_0/I_0 \subset A/I$ であり、
$S_0 \to A/I$ は $A_0/I_0$ の中に写る。従って、図式
$$
\xymatrix{
S_0 \ar[r] \ar@{-->}[rd] & A_0/I_0 \\
R_0 \ar[r] \ar[u] & A_0 \ar[u]
}
$$
における点線矢印を見つければ十分である。
構成により環 $A_0$ は $\mathbf{Z}$ 上有限型である。
従って $A_0$ はネーター環であり、よって $I_0$ は冪零である。
補題 \ref{algebra-lemma-Noetherian-power} を参照せよ。
$I_0^n = 0$ としよう。命題 \ref{algebra-proposition-smooth-formally-smooth} により、
$R_0$-代数写像 $S_0 \to A_0/I_0$ を順次
$S_0 \to A_0/I_0^2$、$S_0 \to A_0/I_0^3$、$\ldots$、
そして最後に $S_0 \to A_0/I_0^n = A_0$ へ持ち上げることができる。
\end{proof}

\section{滑らかさと微分}
\label{algebra-section-smooth-differential}

\noindent
微分と滑らかな環写像に関するいくつかの結果を示す。

\begin{lemma}
\label{algebra-lemma-triangle-differentials-smooth}
$A \to B \to C$ を環写像とし、$B \to C$ は滑らかであるとする。
% P01-E0849: the frozen sentence combines “Given” with a redundant “then”.
このとき、補題 \ref{algebra-lemma-exact-sequence-differentials} の列
$$
0 \to C \otimes_B \Omega_{B/A} \to \Omega_{C/A} \to \Omega_{C/B} \to 0
$$
は完全である。
\end{lemma}

\begin{proof}
滑らかな環写像は形式的に滑らかだから、これはより一般的な
補題 \ref{algebra-lemma-ses-formally-smooth} から従う。
命題 \ref{algebra-proposition-smooth-formally-smooth} を参照せよ。
一方、滑らかさの定義には $H_1(L_{C/B}) = 0$ が含まれるから、
補題 \ref{algebra-lemma-exact-sequence-NL} から直接にも従う。
% P01-E0850: the frozen proof switches from NL to unexplained L notation.
ここで滑らかさは $B \to C$ の滑らかさをいう。
\end{proof}

\begin{lemma}
\label{algebra-lemma-differential-seq-smooth}
$A \to B \to C$ を環写像とし、$A \to C$ は滑らか、
$B \to C$ は全射で、その核は $J \subset B$ であるとする。
このとき、補題 \ref{algebra-lemma-differential-seq} の完全列
$$
0 \to J/J^2 \to \Omega_{B/A} \otimes_B C \to \Omega_{C/A} \to 0
$$
は分裂完全である。
\end{lemma}

\begin{proof}
滑らかな環写像は形式的に滑らかだから、これはより一般的な
補題 \ref{algebra-lemma-differential-seq-formally-smooth} から従う。
命題 \ref{algebra-proposition-smooth-formally-smooth} を参照せよ。
\end{proof}

\begin{lemma}
\label{algebra-lemma-application-NL-smooth}
$A \to B \to C$ を環写像とする。$A \to C$ は全射（従って
$B \to C$ も全射）であり、$A \to B$ は滑らかであると仮定する。
$I = \Ker(A \to C)$ および $J = \Ker(B \to C)$ とおく。
% P01-E0851: the frozen designation omits “by/as”.
このとき、補題 \ref{algebra-lemma-application-NL} の列
$$
0 \to I/I^2 \to J/J^2 \to \Omega_{B/A} \otimes_B B/J \to 0
$$
は完全である。
\end{lemma}

\begin{proof}
滑らかな環写像は形式的に滑らかだから、これはより一般的な
補題 \ref{algebra-lemma-application-NL-formally-smooth} から従う。
命題 \ref{algebra-proposition-smooth-formally-smooth} を参照せよ。
\end{proof}

\begin{lemma}
\label{algebra-lemma-section-smooth}
\begin{slogan}
$R$ が $S$ の直和因子であり、$S$ が $R$ 上滑らかならば、
$S$ の $I$-進完備化はしばしば $R$ 上の形式冪級数環となる。
ここで $I$ は $S$ から $R$ への射影の核である。
% P01-E0852: the frozen slogan says “a power series” instead of “a power series ring”.
% P01-E0853: the frozen slogan lacks terminal punctuation.
\end{slogan}
$\varphi : R \to S$ を滑らかな環写像とする。
$\sigma : S \to R$ を $\varphi$ の左逆とする。
$I = \Ker(\sigma)$ とおく。このとき、
\begin{enumerate}
\item $I/I^2$ は有限局所自由 $R$-加群である、および
\item $I/I^2$ が自由ならば、$S^\wedge$ は $S$ の $I$-進完備化として、
$R$-代数の同型 $S^\wedge \cong R[[t_1, \ldots, t_d]]$ が成り立つ。
\end{enumerate}
\end{lemma}

\begin{proof}
補題 \ref{algebra-lemma-differential-seq-split} を
$R \to S \to R$ に適用すると、
$I/I^2 = \Omega_{S/R} \otimes_{S, \sigma} R$ であることが分かる。
滑らかな射の定義により加群 $\Omega_{S/R}$ は $S$ 上有限局所自由だから、
（1）が従う。
$I/I^2$ が自由ならば、$I/I^2$ における像が $R$-基底をなす
$f_1, \ldots, f_d \in I$ を選ぶ。次で定義される $R$-代数写像を考える。
$$
\Psi : R[[x_1, \ldots, x_d]] \longrightarrow S^\wedge, \quad
x_i \longmapsto f_i.
$$
$P = R[[x_1, \ldots, x_d]]$ および $J = (x_1, \ldots, x_d) \subset P$ とおく。
% P01-E0854: the frozen designation omits “by/as”.
誘導される商環の写像を $\Psi_n : P/J^n \to S/I^n$ と書く。
$S/I^2 = \varphi(R) \oplus I/I^2$ であることに注意する。
従って $\Psi_2$ は同型である。$\Psi_2$ の逆を
$\sigma_2 : S/I^2 \to P/J^2$ とする。
% P01-E0855: the frozen designation of the inverse omits “by/as”.
すべての $n > 2$ について、$\Psi_n$ の逆
$\sigma_n : S/I^n \to P/J^n$ が存在することを $n$ に関する帰納法で示す。
実際、$S$ は $R$ 上形式的に滑らかだから
（命題 \ref{algebra-proposition-smooth-formally-smooth}）、$R$-代数の実線図式
$$
\xymatrix{
S \ar@{..>}[r] \ar[rd]_{\sigma_{n - 1}} & P/J^n \ar[d] \\
 & P/J^{n - 1}
}
$$
において、図式を可換にするある $R$-代数写像
$\tau : S \to P/J^n$ で点線矢印を埋められる。
これは $J^{n - 1}$ を法として $\sigma_{n - 1}$ に等しい
$R$-代数写像 $\overline{\tau} : S/I^n \to P/J^n$ を誘導する。
構成により写像 $\Psi_n$ は全射である。
いま $\overline{\tau} \circ \Psi_n$ は $P/J^n$ の $R$-代数自己準同型であり、
$x_i$ を $x_i + \delta_{i, n}$ に写す。ここで
$\delta_{i, n} \in J^{n - 1}/J^n$ である。
従って $\Psi_n$ は同型であり、ゆえに逆 $\sigma_n$ をもつ。
これで補題が証明された。
\end{proof}

\section{体上の滑らかな代数}
\label{algebra-section-smooth-over-field}

\noindent
注意：次の二つの補題は、一般の非完全体上では成り立たない。

\begin{lemma}
\label{algebra-lemma-rank-omega}
$k$ を代数閉体とする。
$S$ を有限型 $k$-代数とする。
$\mathfrak m \subset S$ を極大イデアルとする。このとき
$$
\dim_{\kappa(\mathfrak m)} \Omega_{S/k} \otimes_S \kappa(\mathfrak m)
=
\dim_{\kappa(\mathfrak m)} \mathfrak m/\mathfrak m^2.
$$
\end{lemma}

\begin{proof}
補題 \ref{algebra-lemma-differential-seq} の完全列
$$
\mathfrak m/\mathfrak m^2 \to
\Omega_{S/k} \otimes_S \kappa(\mathfrak m) \to
\Omega_{\kappa(\mathfrak m)/k} \to 0
$$
を考える。最初の写像が同型であることを示したい。
$k$ は代数閉体だから、定理 \ref{algebra-theorem-nullstellensatz} により合成
$k \to \kappa(\mathfrak m)$ は同型である。
従って全射 $S \to \kappa(\mathfrak m)$ は $k$-代数写像として分裂し、
補題 \ref{algebra-lemma-differential-seq-split} により上の列は左でも完全である。
$\Omega_{\kappa(\mathfrak m)/k} = 0$ だから結論を得る。
\end{proof}

\begin{lemma}
\label{algebra-lemma-characterize-smooth-kbar}
$k$ を代数閉体とする。
$S$ を有限型 $k$-代数とする。
$\mathfrak m \subset S$ を極大イデアルとする。次の条件は同値である。
\begin{enumerate}
\item 環 $S_{\mathfrak m}$ は正則局所環である。
\item
$\dim_{\kappa(\mathfrak m)} \Omega_{S/k} \otimes_S \kappa(\mathfrak m)
\leq \dim(S_{\mathfrak m})$ である。
\item
$\dim_{\kappa(\mathfrak m)} \Omega_{S/k} \otimes_S \kappa(\mathfrak m)
= \dim(S_{\mathfrak m})$ である。
\item $g \in S$、$g \not \in \mathfrak m$ を満たし、$S_g$ が $k$ 上
滑らかとなるものが存在する。言い換えれば、$S/k$ は $\mathfrak m$ において
滑らかである。
\end{enumerate}
\end{lemma}

\begin{proof}
補題 \ref{algebra-lemma-rank-omega} と定義 \ref{algebra-definition-regular} により、
（1）、（2）、（3）が同値であることに注意する。

\medskip\noindent
$S$ は $\mathfrak m$ において滑らかであると仮定する。
補題 \ref{algebra-lemma-smooth-syntomic} により、適切な
$g \in S$、$g \not \in \mathfrak m$ に対して $S_g$ は $k$ 上
標準滑らかである。従って補題 \ref{algebra-lemma-standard-smooth} により、
$\Omega_{S_g/k}$ は階数 $\dim(S_g)$ の自由加群である。
従って補題 \ref{algebra-lemma-rank-omega} により
$\dim(S_{\mathfrak m}) = \dim (\mathfrak m/\mathfrak m^2)$ であり、
言い換えれば $S_\mathfrak m$ は正則である。

\medskip\noindent
逆に、$S_{\mathfrak m}$ は正則であると仮定する。
$d = \dim(S_{\mathfrak m}) = \dim \mathfrak m/\mathfrak m^2$ とおく。
各 $i$ について $x_i$ が $\mathfrak m$ の元に写るような表示
$S = k[x_1, \ldots, x_n]/I$ を選ぶ。言い換えれば、
$\mathfrak m'' = (x_1, \ldots, x_n)$ は $k[x_1, \ldots, x_n]$ の
対応する極大イデアルである。列
$$
I/\mathfrak m''I \to \mathfrak m''/(\mathfrak m'')^2
\to \mathfrak m/(\mathfrak m)^2 \to 0
$$
が完全であることに注意する。
% P01-E0856: the frozen prose incorrectly calls this sequence short exact.
$\mathfrak m/\mathfrak m^2$ への写像の核を
$\mathfrak m''/(\mathfrak m'')^2$ において張る像をもつ
$c = n - d$ 個の元 $f_1, \ldots, f_c \in I$ を選ぶ。
これは明らかに可能である。$J = (f_1, \ldots, f_c)$ とおく。
従って $J \subset I$ である。
$S' = k[x_1, \ldots, x_n]/J$ とおくと、全射 $S' \to S$ がある。
$S'$ の対応する極大イデアルを $\mathfrak m' = \mathfrak m''S'$ とおく。
% P01-E0857: the frozen three consecutive designation clauses omit “by/as”.
従って
$$
\xymatrix{
k[x_1, \ldots, x_n] \ar[r] & S' \ar[r] & S \\
\mathfrak m'' \ar[u] \ar[r] & \mathfrak m' \ar[r] \ar[u] &
\mathfrak m \ar[u]
}
$$
を得る。$J$ の選び方により、完全列
$$
J/\mathfrak m''J \to \mathfrak m''/(\mathfrak m'')^2
\to \mathfrak m'/(\mathfrak m')^2 \to 0
$$
から $\dim( \mathfrak m'/(\mathfrak m')^2 ) = d$ が分かる。
$S'_{\mathfrak m'}$ は $S_{\mathfrak m}$ へ全射だから、
$\dim(S_{\mathfrak m'}) \geq d$ である。
% P01-E0858: the frozen dimension localizes S at the maximal ideal of S' rather than S'.
従って定義 \ref{algebra-definition-regular-local} の前の議論により、
$S'_{\mathfrak m'}$ も次元 $d$ の正則環である。
$S'$ は $c = n - d$ 本の方程式で切り出されていたから、
$g' \in S'$、$g' \not \in \mathfrak m'$ であって、$S'_{g'}$ が
$k$ 上大域的完全交叉となるものが存在する。
補題 \ref{algebra-lemma-lci} を参照せよ。
また、写像 $S'_{\mathfrak m'} \to S_{\mathfrak m}$ は同じ次元の
ネーター局所整域の間の全射であり、従って同型である。
ゆえに $S' \to S$ は有限生成の核をもつ全射であり、
$\mathfrak m'$ における局所化の後に同型となる。
従って $g' \in S'$、$g \not \in \mathfrak m'$ であって、
$S'_{g'} \to S_{g'}$ が同型となるものを見つけられる。
% P01-E0859: the frozen sentence switches from g' to undefined g and localizes S at g'.
以上より、$S$ をある単項局所化で置き換えれば、$S$ は
大域的完全交叉であると仮定してよい。

\medskip\noindent
ここで $S = k[x_1, \ldots, x_n]/(f_1, \ldots, f_c)$、
$\dim S = n - c$ と書ける。この代数の素朴な余接複体は
$$
\bigoplus S \cdot f_j
\to
\bigoplus S \cdot \text{d}x_i
$$
で与えられることを思い出そう。
補題 \ref{algebra-lemma-relative-global-complete-intersection-conormal} を参照せよ。
補題 \ref{algebra-lemma-relative-global-complete-intersection-smooth} により、
$S$ が $\mathfrak m$ において滑らかであることを示すには、上の写像を与える
行列「$A$」の $c \times c$ 小行列式 $g_I$ の一つが
$\mathfrak m$ で消えないことを示せばよい。
補題 \ref{algebra-lemma-rank-omega} により、行列 $A \bmod \mathfrak m$ の階数は
$c$ である。従って結論を得る。
\end{proof}

\begin{lemma}
\label{algebra-lemma-characterize-smooth-over-field}
$k$ を任意の体とする。
$S$ を有限型 $k$-代数とする。
$X = \Spec(S)$ とおく。
$\mathfrak q \subset S$ を、$x \in X$ に対応する素イデアルとする。
次の条件は同値である。
\begin{enumerate}
\item $k$-代数 $S$ は $k$ 上 $\mathfrak q$ において滑らかである。
\item
$\dim_{\kappa(\mathfrak q)} \Omega_{S/k} \otimes_S \kappa(\mathfrak q)
\leq \dim_x X$ である。
\item
$\dim_{\kappa(\mathfrak q)} \Omega_{S/k} \otimes_S \kappa(\mathfrak q)
= \dim_x X$ である。
\end{enumerate}
さらに、この場合、局所環 $S_{\mathfrak q}$ は正則である。
\end{lemma}

\begin{proof}
$S$ が $k$ 上 $\mathfrak q$ において滑らかなら、補題
\ref{algebra-lemma-smooth-syntomic} により、$g \in S$、
$g \not \in \mathfrak q$ であって $S_g$ が $k$ 上標準滑らかとなるものが
存在する。$k$ 上の標準滑らかな代数の微分加群は、その次元に等しい階数の
自由加群である。補題 \ref{algebra-lemma-standard-smooth} を参照せよ
（体上の相対大域的完全交叉の次元が、変数の個数から方程式の個数を
引いたものに等しいことを用いる）。従って（1）は（3）を含意する。
補題の証明を終えるには、（2）が（1）を含意し、さらに
$S_{\mathfrak q}$ が正則であることを含意することを示せば十分である。

\medskip\noindent
（2）を仮定する。ナカヤマの補題 \ref{algebra-lemma-NAK} により、
$\Omega_{S/k, \mathfrak q}$ は $\leq \dim_x X$ 個の元で生成できる。
$g \in S$、$g \not \in \mathfrak q$ であって、
$\Omega_{S/k}$ が高々 $\dim_x X$ 個の元で生成されるものを選び、
$S$ を $S_g$ で置き換えてよい。
$K/k$ を、$k$-代数写像
$\psi : \kappa(\mathfrak q) \to K$ が存在するような代数閉拡大体とする。
$S_K = K \otimes_k S$ を考える。$\mathfrak m \subset S_K$ を、全射
$$
\xymatrix{
S_K = K \otimes_k S \ar[r] &
K \otimes_k \kappa(\mathfrak q)
\ar[r]^-{\text{id}_K \otimes \psi} &
K.
}
$$
に対応する極大イデアルとする。
$\mathfrak m \cap S = \mathfrak q$、言い換えれば
$\mathfrak m$ は $\mathfrak q$ の上にあることに注意する。
補題 \ref{algebra-lemma-dimension-at-a-point-preserved-field-extension} により、
$\mathfrak m$ に対応する点における
$X_K = \Spec(S_K)$ の次元は $\dim_x X$ である。
補題 \ref{algebra-lemma-dimension-closed-point-finite-type-field} により、これは
$\dim((S_K)_{\mathfrak m})$ に等しい。
補題 \ref{algebra-lemma-differentials-base-change} により、$S_K$ の $K$ 上の
微分加群は $\Omega_{S/k}$ の基底変換であり、従ってこれも高々
$\dim_x X = \dim((S_K)_{\mathfrak m})$ 個の元で生成される。
補題 \ref{algebra-lemma-characterize-smooth-kbar} により、$S_K$ は $K$ 上
$\mathfrak m$ において滑らかである。
補題 \ref{algebra-lemma-flat-base-change-locus-smooth} により、これは
$S$ が $k$ 上 $\mathfrak q$ において滑らかであることを含意する。
これで（1）が示された。さらに、補題
\ref{algebra-lemma-characterize-smooth-kbar} により局所環
$(S_K)_{\mathfrak m}$ は正則である。
$S_{\mathfrak q} \to (S_K)_{\mathfrak m}$ は平坦だから、補題
\ref{algebra-lemma-flat-under-regular} より $S_{\mathfrak q}$ は正則である。
\end{proof}

\noindent
次の補題は（いくつかの異なる仕方で）大幅に一般化できる。

\begin{lemma}
\label{algebra-lemma-computation-differential}
$k$ を体とする。
$R$ を $k$ を含むネーター局所環とする。
剰余体 $\kappa = R/\mathfrak m$ が $k$ の有限生成分離拡大であると
仮定する。このとき写像
$$
\text{d} :
\mathfrak m/\mathfrak m^2
\longrightarrow
\Omega_{R/k} \otimes_R \kappa(\mathfrak m)
$$
は単射である。
\end{lemma}

\begin{proof}
$R$ を $R/\mathfrak m^2$ で置き換えてよい。従って
$\mathfrak m^2 = 0$ と仮定してよい。仮定により
$\kappa = k(\overline{x}_1, \ldots, \overline{x}_r, \overline{y})$
と書ける。ただし $\overline{x}_1, \ldots, \overline{x}_r$ は
$\kappa$ の $k$ 上の超越基底であり、$\overline{y}$ は
$k(\overline{x}_1, \ldots, \overline{x}_r)$ 上分離代数的である。
$P(T)$ を $k(\overline{x}_1, \ldots, \overline{x}_r)[T]$ に属する、
$\overline{y}$ の $k(\overline{x}_1, \ldots, \overline{x}_r)$ 上の
最小多項式とする。
すると $P(\overline{y}) = 0$ だが $P'(\overline{y}) \not = 0$ である。
各 $\overline{x}_i \in \kappa$ の任意の持ち上げ $x_i \in R$ を選ぶ。
これにより $k$-代数の可換図式
$$
\xymatrix{
R \ar[r] & \kappa \\
& k(\overline{x}_1, \ldots, \overline{x}_r) \ar[lu]^\varphi \ar[u]
}
$$
を得る。左上向きの矢印 $\varphi$ を、$\kappa$ から $R$ への
$k$-代数写像に拡張したい。そのために、$\overline{y}$ の任意の持ち上げ
$y \in R$ を選ぶ。それが
$\kappa \cong k(\overline{x}_1, \ldots, \overline{x}_r)[T]/(P)$
上で定義された $k$-代数写像を与えることを見るには、
$P^\varphi(y) = 0$ となるように $y$ を選べることを示せばよい。
そうでなければ、$\delta \in \mathfrak m$ に対して
$$
P^\varphi(y + \delta) = P^\varphi(y) + (P')^\varphi(y)\delta
$$
と計算される。これは $\mathfrak m^2 = 0$ による。
$(P')^\varphi(y)$ は $R$ の単元だから、望むように選択を調整できる。
これにより $k$-代数として
$R \cong \kappa \oplus \mathfrak m$ であることが分かる。
ここで $\Omega_{\kappa \oplus \mathfrak m/k}$ を直接計算するか、補題
\ref{algebra-lemma-differential-seq-split} を適用すれば証明が終わる。
\end{proof}

\begin{lemma}
\label{algebra-lemma-separable-smooth}
$k$ を体とする。
$S$ を有限型 $k$-代数とする。
$\mathfrak q \subset S$ を素イデアルとする。
$\kappa(\mathfrak q)$ が $k$ 上分離的であると仮定する。
次の条件は同値である。
\begin{enumerate}
\item 代数 $S$ は $k$ 上 $\mathfrak q$ において滑らかである。
\item 環 $S_{\mathfrak q}$ は正則である。
\end{enumerate}
\end{lemma}

\begin{proof}
% P01-E0860: the frozen designation omits “ideal” after “maximal”.
$R = S_{\mathfrak q}$ とおき、その極大イデアルを $\mathfrak m$、
剰余体を $\kappa$ とおく。
% P01-E0861: the frozen coordinated citation lacks “Lemma” before its second reference.
補題 \ref{algebra-lemma-computation-differential} および補題
\ref{algebra-lemma-differential-seq} により、短完全列
$$
0 \to \mathfrak m/\mathfrak m^2 \to
\Omega_{R/k} \otimes_R \kappa \to
\Omega_{\kappa/k} \to 0
$$
を得る。補題 \ref{algebra-lemma-differentials-localize} により
$\Omega_{R/k} = \Omega_{S/k, \mathfrak q}$ であることに注意する。
さらに、$\kappa$ は $k$ 上分離的だから
$\dim_{\kappa} \Omega_{\kappa/k} = \text{trdeg}_k(\kappa)$ である。
従って
$$
\dim_{\kappa} \Omega_{R/k} \otimes_R \kappa
=
\dim_\kappa \mathfrak m/\mathfrak m^2 + \text{trdeg}_k (\kappa)
\geq
\dim R + \text{trdeg}_k (\kappa)
=
\dim_{\mathfrak q} S
$$
を得る（最後の等号については補題
\ref{algebra-lemma-dimension-at-a-point-finite-type-field} を参照せよ）。
% P01-E0862: the frozen prose lacks punctuation before the equality criterion.
ここで等号が成り立つことと $R$ が正則であることは同値である。
従って補題 \ref{algebra-lemma-characterize-smooth-over-field} を適用すればよい。
\end{proof}

\begin{lemma}
\label{algebra-lemma-characteristic-zero}
$R \to S$ を $\mathbf{Q}$-代数の写像とする。
$f \in S$ とし、ある $S$-部分加群 $C$ に対して
$\Omega_{S/R} = S \text{d}f \oplus C$ であると仮定する。このとき
\begin{enumerate}
\item $f$ は冪零でない。
\item $S$ がネーター局所環なら、$f$ は $S$ の非零因子である。
\end{enumerate}
\end{lemma}

\begin{proof}
$a \in S$ に対し、ある $\theta(a) \in S$ と $c(a) \in C$ を用いて
$\text{d}(a) = \theta(a)\text{d}f + c(a)$ と書く。
$R$-導分 $S \to S$、$a \mapsto \theta(a)$ を考える。
$\theta(f) = 1$ であることに注意する。

\medskip\noindent
$f^n = 0$ で $n > 1$ が最小なら、
$0 = \theta(f^n) = n f^{n - 1}$ となり、$n$ の最小性に矛盾する。
従って $f$ は冪零でない。

\medskip\noindent
$fa = 0$ と仮定する。$f$ が単元なら $a = 0$ であり、結論を得る。
$f$ は単元でないと仮定する。このときライプニッツ則により
$0 = \theta(fa) = f\theta(a) + a$ であり、従って $a \in (f)$ である。
帰納的に $fa = 0 \Rightarrow a \in (f^n)$ を示したと仮定する。
$a = f^nb$ と書くと
$0 = \theta(f^{n + 1}b) = (n + 1)f^nb + f^{n + 1}\theta(b)$ を得る。
従って
$a = f^n b = -f^{n + 1}\theta(b)/(n + 1) \in (f^{n + 1})$ である。
ネーター局所環 $S$ では $\bigcap (f^n) = 0$ であるから、結論を得る。
補題 \ref{algebra-lemma-intersect-powers-ideal-module-zero} を参照せよ。
\end{proof}

\noindent
次の結果は、おそらく応用上ほとんど役に立たない。

\begin{lemma}
\label{algebra-lemma-characteristic-zero-local-smooth}
$k$ を標数 $0$ の体とする。
$S$ を有限型 $k$-代数とする。
$\mathfrak q \subset S$ を素イデアルとする。
次の条件は同値である。
\begin{enumerate}
\item 代数 $S$ は $k$ 上 $\mathfrak q$ において滑らかである。
\item $S_{\mathfrak q}$-加群 $\Omega_{S/k, \mathfrak q}$ は
（有限）自由である。
\item 環 $S_{\mathfrak q}$ は正則である。
\end{enumerate}
\end{lemma}

\begin{proof}
標数零では任意の体の拡大が分離的だから、（1）と（3）の同値性は
補題 \ref{algebra-lemma-separable-smooth} から従う。また（1）は、滑らかな代数の
定義により（2）を含意する。
$\Omega_{S/k, \mathfrak q}$ が $S_{\mathfrak q}$ 上自由であると仮定する。
補題 \ref{algebra-lemma-separable-smooth} の証明で導入した記法と観察を用いる。
従って $R = S_{\mathfrak q}$ とし、その極大イデアルを $\mathfrak m$、
剰余体を $\kappa$ とする。目標は $R$ が正則であることを示すことである。

\medskip\noindent
$\mathfrak m/\mathfrak m^2 = 0$ なら $\mathfrak m = 0$ であり、
$R \cong \kappa$ である。従って $R$ は正則であり、結論を得る。

\medskip\noindent
$\mathfrak m/ \mathfrak m^2 \not = 0$ なら、
$\mathfrak m/ \mathfrak m^2$ における像が零でない任意の
$f \in \mathfrak m$ を選ぶ。補題 \ref{algebra-lemma-computation-differential} により、
$\text{d}f$ は $\Omega_{R/k}/\mathfrak m\Omega_{R/k}$ において
零でない像をもつ。仮定により
$\Omega_{R/k} = \Omega_{S/k, \mathfrak q}$ は有限自由であるから、
ナカヤマの補題 \ref{algebra-lemma-NAK} により $\text{d}f$ は直和因子を生成する。
補題 \ref{algebra-lemma-characteristic-zero} を適用すると、$f$ は $R$ の
非零因子である。さらに補題 \ref{algebra-lemma-differential-seq} により完全列
$$
(f)/(f^2) \to \Omega_{R/k} \otimes_R R/fR \to \Omega_{(R/fR)/k} \to 0
$$
を得る。これは $\Omega_{(R/fR)/k}$ も有限自由であることを含意する。
従って帰納法により $R/fR$ は正則局所環である。
$f \in \mathfrak m$ は非零因子だったから、補題
\ref{algebra-lemma-regular-mod-x} により $R$ は正則である。
\end{proof}

\begin{example}
\label{algebra-example-characteristic-p}
補題 \ref{algebra-lemma-characteristic-zero-local-smooth} は標数
$p > 0$ では成り立たない。標準的な例は環写像
$$
\mathbf{F}_p \longrightarrow \mathbf{F}_p[x]/(x^p)
$$
である。その微分加群は自由だが、明らかに滑らかでない。また
（$p > 2$ のとき）環写像
% P01-E0863: the frozen example does not define or constrain alpha.
$$
\mathbf{F}_p(t) \to \mathbf{F}_p(t)[x, y]/(x^p + y^2 + \alpha)
$$
は素イデアル $\mathfrak q = (y, x^p + \alpha)$ において滑らかでないが、
正則である。
\end{example}

\noindent
上の結果を用いると、一般点における滑らかさを体の拡大によって
特徴づけられる。

\begin{lemma}
\label{algebra-lemma-smooth-at-generic-point}
$R \to S$ を有限型の単射な環写像とし、$R$ と $S$ は整域であるとする。
このとき $R \to S$ が $\mathfrak q = (0)$ において滑らかであることと、
誘導される分数体の拡大 $L/K$ が分離的であることは同値である。
\end{lemma}

\begin{proof}
$R \to S$ が $(0)$ において滑らかであると仮定する。
ある非零元 $g \in S$ に対して $S$ を $S_g$ で置き換え、
$R \to S$ が滑らかであると仮定してよい。このとき
$K \to S \otimes_R K$ は滑らかである
（補題 \ref{algebra-lemma-base-change-smooth}）。さらに、任意の体の拡大
$K'/K$ に対して環写像 $K' \to S \otimes_R K'$ も滑らかである。
従って補題 \ref{algebra-lemma-characterize-smooth-over-field} により
$S \otimes_R K'$ は正則環であり、特に被約である。
% P01-E0864: the frozen English predicate has a spurious indefinite article.
従って $S \otimes_R K$ は $K$ 上幾何学的被約である。
補題 \ref{algebra-lemma-geometrically-reduced-permanence} により、
$L$ は $K$ 上幾何学的被約である。
従って補題 \ref{algebra-lemma-characterize-separable-field-extensions} により、
$L/K$ は分離的である。

\medskip\noindent
逆に、$L/K$ が分離的であると仮定する。補題
\ref{algebra-lemma-generic-finite-presentation} により、$R \to S$ は
有限表示であると仮定してよい。補題
\ref{algebra-lemma-flat-base-change-locus-smooth} により、
$K \to S \otimes_R K$ が $(0)$ において滑らかであることを示せば十分である。
これは補題 \ref{algebra-lemma-separable-smooth}、体が正則環であるという事実、
および $L/K$ が分離的であるという仮定から従う。
\end{proof}

\section{ネーターの場合の滑らかな環写像}
\label{algebra-section-smooth-Noetherian}

\begin{definition}
\label{algebra-definition-small-extension}
$\varphi : B' \to B$ を環写像とする。
$B'$ と $B$ がアルティン局所環であり、$\varphi$ が全射で、かつ
$I = \Ker(\varphi)$ が $B'$-加群として長さ $1$ をもつとき、
$\varphi$ を {\it 小拡大} と呼ぶ。
\end{definition}

\noindent
これは明らかに、$I^2 = 0$ であり、さらにある $x \in B'$ に対して
$I = (x)$ かつ $\mathfrak m' x = 0$ となることを意味する。
ここで $\mathfrak m' \subset B'$ は極大イデアルである。

\begin{lemma}
\label{algebra-lemma-smooth-test-artinian}
$R \to S$ を環写像とする。$\mathfrak q$ を、
$\mathfrak p \subset R$ の上にある $S$ の素イデアルとする。
% P01-E0865: the frozen second coordinated predicate omits “is”.
$R$ はネーターであり、$R \to S$ は有限型であると仮定する。
次の条件は同値である。
\begin{enumerate}
\item $R \to S$ は $\mathfrak q$ において滑らかである。
\item 局所 $R$-代数の任意の全射
$(B', \mathfrak m') \to (B, \mathfrak m)$ で、
$\Ker(B' \to B)$ の平方が零であるものと、実線からなる任意の可換図式
$$
\xymatrix{
S \ar[r] \ar@{-->}[rd] & B \\
R \ar[r] \ar[u] & B' \ar[u]
}
$$
で $\mathfrak q = S \cap \mathfrak m$ を満たすものに対し、図式を
可換にする点線の矢印が存在する。
% P01-E0866: the frozen intersection formula should use inverse image because S -> B need not be injective.
\item （2）と同じであるが、$B' \to B$ は小拡大のみを走る。
\item （2）と同じであるが、$B' \to B$ は、さらに $S \to B$ が同型
$\kappa(\mathfrak q) \cong \kappa(\mathfrak m)$ を誘導する小拡大のみを走る。
\end{enumerate}
\end{lemma}

\begin{proof}
（1）を仮定する。これは、ある $g \in S$、$g \not \in \mathfrak q$ に対して
$R \to S_g$ が滑らかであることを意味する。命題
\ref{algebra-proposition-smooth-formally-smooth} により、$R \to S_g$ は
形式的に滑らかである。（2）のような任意の図式が与えられると、写像
$S \to B$ は自動的に $S_{\mathfrak q}$ を経由し、従ってなおさら
$S_g$ を経由することに注意する。$S_g$ の $R$ 上の形式的滑らかさにより、
左上の頂点を $S_g$ とした同様の図式に入る射 $S_g \to B'$ を得る。
これと $S \to S_g$ を合成すれば、求める矢印が得られる。
言い換えれば、（1）が（2）を含意することを示した。

\medskip\noindent
（2）が（3）を含意し、（3）が（4）を含意することは明らかである。

\medskip\noindent
（4）を仮定する。（1）が成り立つことを示せば、補題の証明が終わる。
表示 $S\allowbreak = R[x_1, \ldots,\allowbreak x_n]\allowbreak/(f_1, \ldots,\allowbreak f_m)$ を選ぶ。
$S$ は $R$ 上有限型であり、従って有限表示だから、これは可能である
（補題 \ref{algebra-lemma-Noetherian-finite-type-is-finite-presentation} を参照せよ）。
$I = (f_1, \ldots, f_m)$ とおく。この表示の素朴な余接複体
$$
\text{d} : I/I^2
\longrightarrow
\bigoplus\nolimits_{j = 1}^n S\text{d}x_j
$$
を考える（節 \ref{algebra-section-netherlander} を参照せよ）。
この複体を $\mathfrak q$ で局所化したとき、その写像が分裂単射になることを
示せば十分である。補題 \ref{algebra-lemma-smooth-at-point} を参照せよ。
% P01-E0867: the frozen designation of S' omits “by/as”.
$S' = R[x_1, \ldots, x_n]/I^2$ とおく。
補題 \ref{algebra-lemma-differential-mod-power-ideal} により
$$
S \otimes_{S'} \Omega_{S'/R} =
S \otimes_{R[x_1, \ldots, x_n]} \Omega_{R[x_1, \ldots, x_n]/R} =
\bigoplus\nolimits_{j = 1}^m S\text{d}x_j.
$$
% P01-E0868: the frozen final direct sum runs to m although the polynomial ring has n variables.
従って写像
$$
\text{d} :
I/I^2
\longrightarrow
S \otimes_{S'} \Omega_{S'/R}
$$
は上の素朴な余接複体の写像と同じである。特に、示そうとしている主張の
真偽は三つの環 $R \to S' \to S$ のみに依存する。
$\mathfrak q' \subset R[x_1, \ldots, x_n]$ を、$\mathfrak q$ に対応する
素イデアルとする。局所化は微分加群を取る操作と可換だから
（補題 \ref{algebra-lemma-differentials-localize}）、写像
\begin{equation}
\label{algebra-equation-target-map}
\text{d} :
I_{\mathfrak q'}/I_{\mathfrak q'}^2
\longrightarrow
S_{\mathfrak q} \otimes_{S'_{\mathfrak q'}} \Omega_{S'_{\mathfrak q'}/R}
\end{equation}
が分裂単射であることを示せば十分である。この写像は
$R \to S'_{\mathfrak q'} \to S_{\mathfrak q}$ から得られる。

\medskip\noindent
$N \in \mathbf{N}$ を整数とする。環
$$
B'_N = S'_{\mathfrak q'} / (\mathfrak q')^N S'_{\mathfrak q'}
= (S'/(\mathfrak q')^N S')_{\mathfrak q'}
$$
と、その商 $B_N = B'_N/IB'_N$ を考える。
$B_N \cong S_{\mathfrak q}/\mathfrak q^NS_{\mathfrak q}$ であることに
注意する。$B'_N$ は局所ネーター環をその極大イデアルの冪で割った商だから、
アルティン局所環である。$B'_N \to B_N$ の核 $I_N$ の
$B'_N$-部分加群によるフィルトレーション
$$
0 \subset J_{N, 1} \subset J_{N, 2} \subset \ldots \subset J_{N, n(N)} = I_N
$$
で、各逐次商 $J_{N, i}/J_{N, i - 1}$ が長さ $1$ をもつものを考える。
（$B'_N$ はアルティンだから、このようなフィルトレーションが存在する。）
これにより小拡大の列
$$
B'_N  \to B'_N/J_{N, 1} \to B'_N/J_{N, 2} \to \ldots \to
B'_N/J_{N, n(N)} = B'_N/I_N
= B_N = S_{\mathfrak q}/\mathfrak q^NS_{\mathfrak q}
$$
を得る。写像 $S \to B_N$ から始めて、これらの小拡大に条件（4）を
順次適用すると、可換図式
% P01-E0869: the frozen English content clause omits “that”.
$$
\xymatrix{
S \ar[r] \ar[rd] & B_N  \\
R \ar[r] \ar[u] & B'_N \ar[u]
}
$$
が存在することが分かる。環写像 $S \to B'_N$ は明らかに
$S \to S_{\mathfrak q} \to B'_N$ と分解し、ここで
$S_{\mathfrak q} \to B'_N$ は局所環の局所準同型である。
さらに、$B'_N$ の極大イデアルの $N$ 乗は零だから、
$S_{\mathfrak q} \to B'_N$ は
$S_{\mathfrak q}/(\mathfrak q)^NS_{\mathfrak q} = B_N$ を経由する。
言い換えれば、すべての $N \in \mathbf{N}$ に対し、$R$-代数の全射
$B'_N \to B_N$ は分裂をもつことを示した。

\medskip\noindent
核が $I_N$ である全射 $B'_N \to B_N$ から得られる表示
$$
I_N \to B_N \otimes_{B'_N} \Omega_{B'_N/R} \to \Omega_{B_N/R} \to 0
$$
を考える（補題 \ref{algebra-lemma-differential-seq} を参照せよ）。
上で示したように、$R$-代数写像 $B'_N \to B_N$ は右逆をもつ。
従って補題 \ref{algebra-lemma-differential-seq-split} により、上の列は
分裂完全である。従ってすべての $N$ に対して写像
$$
I_N \longrightarrow B_N \otimes_{B'_N} \Omega_{B'_N/R}
$$
は分裂単射である。証明の残りは、これが正確に何を意味するかを
展開することで得られる。ここで
$$
I_N = I_{\mathfrak q'}/
(I_{\mathfrak q'}^2 + (\mathfrak q')^N \cap I_{\mathfrak q'})
$$
である。アルティン--リースの補題（補題 \ref{algebra-lemma-Artin-Rees}）により、
ある $c \geq 0$ が存在して
$$
S_{\mathfrak q}/\mathfrak q^{N - c}S_{\mathfrak q}
\otimes_{S_{\mathfrak q}} I_N =
S_{\mathfrak q}/\mathfrak q^{N - c}S_{\mathfrak q}
\otimes_{S_{\mathfrak q}}
I_{\mathfrak q'}/I_{\mathfrak q'}^2
$$
がすべての $N \geq c$ に対して成り立つ
（これらのテンソル積は、$\mathfrak q^{N - c}$ で割ることを
凝った仕方で表したにすぎない）。
% P01-E0870: the frozen English uses singular “tensor product” after “these”.
もちろん $c \geq 1$ と仮定してよい。
補題 \ref{algebra-lemma-differential-mod-power-ideal} により
$$
S'_{\mathfrak q'}/(\mathfrak q')^{N - c}S'_{\mathfrak q'}
\otimes_{S'_{\mathfrak q'}} \Omega_{B'_N/R} =
S'_{\mathfrak q'}/(\mathfrak q')^{N - c}S'_{\mathfrak q'}
\otimes_{S'_{\mathfrak q'}} \Omega_{S'_{\mathfrak q'}/R}
$$
であることが分かる。
% P01-E0871: the frozen prose lacks a sentence boundary around this display.
さらにこれを $B_N = S_{\mathfrak q}/\mathfrak q^N$ でテンソルすると
$$
S_{\mathfrak q}/\mathfrak q^{N - c}S_{\mathfrak q}
\otimes_{S'_{\mathfrak q'}} \Omega_{B'_N/R} =
S_{\mathfrak q}/\mathfrak q^{N - c}S_{\mathfrak q}
\otimes_{S'_{\mathfrak q'}} \Omega_{S'_{\mathfrak q'}/R}.
$$
を得る。分裂単射は任意のものとのテンソルを取っても分裂単射のままだから、
$$
S_{\mathfrak q}/\mathfrak q^{N - c}S_{\mathfrak q}
\otimes_{S_{\mathfrak q}}
(\ref{algebra-equation-target-map}) =
S_{\mathfrak q}/\mathfrak q^{N - c}S_{\mathfrak q}
\otimes_{S_{\mathfrak q}/\mathfrak q^N S_{\mathfrak q}}
(I_N \longrightarrow B_N \otimes_{B'_N} \Omega_{B'_N/R})
$$
はすべての $N \geq c$ に対して分裂単射である。補題
\ref{algebra-lemma-split-injection-after-completion} により、
(\ref{algebra-equation-target-map}) は分裂単射である。これで証明が終わる。
\end{proof}

\section{滑らかな環写像に関する結果の概観}
\label{algebra-section-smooth-overview}

\noindent
前節までに証明した滑らかな環写像に関する結果を列挙する。
より正確な主張と定義については、挙げた参照先を見よ。
\begin{enumerate}
\item 環写像 $R \to S$ が有限表示であり、$S/R$ の素朴な余接複体が
次数 $0$ に置かれた有限射影的 $S$-加群と擬同型なら、この写像は滑らかである。
定義 \ref{algebra-definition-smooth} を参照せよ。
\item $S$ が $R$ 上滑らかなら、$\Omega_{S/R}$ は有限射影的
$S$-加群である。定義 \ref{algebra-definition-smooth} の後の議論を参照せよ。
\item 滑らかであるという性質は $S$ 上局所的である。
補題 \ref{algebra-lemma-locally-smooth} を参照せよ。
\item 滑らかであるという性質は基底変換で保たれる。
補題 \ref{algebra-lemma-base-change-smooth} を参照せよ。
\item 滑らかであるという性質は合成で保たれる。
補題 \ref{algebra-lemma-compose-smooth} を参照せよ。
\item 滑らかな環写像はシントミックであり、特に平坦である。
補題 \ref{algebra-lemma-smooth-syntomic} を参照せよ。
\item 有限表示かつ平坦な環写像のファイバー環が滑らかなら、その環写像は
滑らかである。補題 \ref{algebra-lemma-flat-fibre-smooth} を参照せよ。
\item 有限表示な環写像 $R \to S$ が滑らかであることと、形式的に
滑らかであることは同値である。命題
\ref{algebra-proposition-smooth-formally-smooth} を参照せよ。
\item $R \to S$ が有限型の環写像で $R$ がネーターなら、
$R \to S$ が滑らかであることを調べるには、アルティン局所環の小拡大に
沿う形式的滑らかさの持ち上げ性質を調べれば十分である。
補題 \ref{algebra-lemma-smooth-test-artinian} を参照せよ。
\item 滑らかな環写像 $R \to S$ は、ある滑らかな環写像
$R_0 \to S_0$ で $R_0$ が $\mathbf{Z}$ 上有限型であるものの
基底変換である。補題 \ref{algebra-lemma-finite-presentation-fs-Noetherian} を参照せよ。
\item 環写像が滑らかとなる点の集合を作る操作は、平坦基底変換と可換である。
補題 \ref{algebra-lemma-flat-base-change-locus-smooth} を参照せよ。
\item $S$ が代数閉体 $k$ 上有限型であり、
$\mathfrak m \subset S$ が極大イデアルなら、次は同値である。
% P01-E0872: the frozen premise omits “is” and the list introduction lacks a colon.
\begin{enumerate}
\item $S$ は $\mathfrak m$ のある近傍で $k$ 上滑らかである。
\item $S_{\mathfrak m}$ は正則局所環である。
\item $\dim(S_{\mathfrak m}) =
\dim_{\kappa(m)} \Omega_{S/k} \otimes_S \kappa(\mathfrak m)$ である。
% P01-E0873: the frozen dimension subscript uses kappa(m) although only mathfrak m is defined.
\end{enumerate}
補題 \ref{algebra-lemma-characterize-smooth-kbar} を参照せよ。
\item $S$ が体 $k$ 上有限型であり、$\mathfrak q \subset S$ が
素イデアルなら、次は同値である。
% P01-E0874: three frozen prime-ideal premises omit “is”; this one also lacks a colon.
\begin{enumerate}
\item $S$ は $\mathfrak q$ のある近傍で $k$ 上滑らかである。
\item $\dim_{\mathfrak q}(S/k) =
\dim_{\kappa(\mathfrak q)} \Omega_{S/k} \otimes_S \kappa(\mathfrak q)$ である。
\end{enumerate}
補題 \ref{algebra-lemma-characterize-smooth-over-field} を参照せよ。
\item $S$ が体上滑らかなら、そのすべての局所環は正則である。
補題 \ref{algebra-lemma-characterize-smooth-over-field} を参照せよ。
\item $S$ が体 $k$ 上有限型で、$\mathfrak q \subset S$ が素イデアルであり、
体の拡大 $\kappa(\mathfrak q)/k$ が分離的で、さらに
$S_{\mathfrak q}$ が正則なら、$S$ は $\mathfrak q$ において $k$ 上
滑らかである。補題 \ref{algebra-lemma-separable-smooth} を参照せよ。
\item $S$ が体 $k$ 上有限型で、$k$ の標数が $0$ であり、
$\mathfrak q \subset S$ が素イデアルで、さらに
$\Omega_{S/k, \mathfrak q}$ が自由なら、$S$ は $\mathfrak q$ において
$k$ 上滑らかである。補題
\ref{algebra-lemma-characteristic-zero-local-smooth} を参照せよ。
\end{enumerate}
これらの結果の一部は、標準滑らかな環写像という概念を用いて証明された。
定義 \ref{algebra-definition-standard-smooth} を参照せよ。これは、シントミック射の場合に
相対大域的完全交叉写像が果たす役割に対応する。また、例を作る最も簡単な
方法でもある。

\section{エタール環写像}
\label{algebra-section-etale}

\noindent
エタール環写像とは、相対次元が零に等しい滑らかな環写像のことである。
これは、次のもう少し直接的な定義と同じである。

\begin{definition}
\label{algebra-definition-etale}
$R \to S$ を環写像とする。有限表示であり、素朴な余接複体
$\NL_{S/R}$ が零と擬同型であるとき、すなわち
$H_1(\NL_{S/R}) = 0$ かつ $\Omega_{S/R} = 0$ であるとき、
$R \to S$ は {\it エタール} であるという。
$S$ の素イデアル $\mathfrak q$ が与えられたとき、ある
$g \in S$、$g \not \in \mathfrak q$ に対して $R \to S_g$ が
エタールなら、$R \to S$ は {\it $\mathfrak q$ においてエタール} であるという。
\end{definition}

\noindent
特に、$S$ が $R$ 上エタールなら $\Omega_{S/R} = 0$ である。
$R \to S$ が滑らかなとき、$R \to S$ がエタールであることと
$\Omega_{S/R} = 0$ であることは同値である。
滑らかな環写像についての結果から、エタール写像についても非常に多くの
結果が自動的に得られる。これらを下の補題 \ref{algebra-lemma-etale} にまとめる。
その前に、{\it 任意の} エタール環写像が標準滑らかであることを証明する。

\begin{lemma}
\label{algebra-lemma-etale-standard-smooth}
任意のエタール環写像は標準滑らかである。より正確には、
$R \to S$ がエタールなら、表示
$S = R[x_1, \ldots, x_n]/(f_1, \ldots, f_n)$ で、
$\det(\partial f_j/\partial x_i)$ の像が $S$ で可逆となるものが存在する。
\end{lemma}

\begin{proof}
$R \to S$ をエタールとする。表示
$S = R[x_1, \ldots, x_n]/I$ を選ぶ。$R \to S$ はエタールだから、
$$
\text{d} :
I/I^2
\longrightarrow
\bigoplus\nolimits_{i = 1, \ldots, n} S\text{d}x_i
$$
は同型であり、特に $I/I^2$ は自由 $S$-加群である。
従って補題 \ref{algebra-lemma-huber} により、必要なら表示を変えることで、
$I = (f_1, \ldots, f_c)$ かつ類 $f_i \bmod I^2$ が $I/I^2$ の
基底をなすと仮定してよい。上の表示写像が同型であることから直ちに
$c = n$ であり、$\det(\partial f_j/\partial x_i)$ は $S$ で可逆だと分かる。
\end{proof}

\begin{lemma}
\label{algebra-lemma-etale}
エタール環写像について、次の結果が成り立つ。
\begin{enumerate}
\item 任意の環 $R$ と任意の $f \in R$ に対して、環写像
$R \to R_f$ はエタールである。
\item エタール環写像の合成はエタールである。
\item エタール環写像の基底変換はエタールである。
\item エタールであるという性質は局所的である。環写像 $R \to S$ と、
単位イデアルを生成する元 $g_1, \ldots, g_m \in S$ が与えられ、
$j = 1, \ldots, m$ に対して $R \to S_{g_j}$ がエタールなら、
$R \to S$ はエタールである。
\item 有限表示な $R \to S$ と平坦な環写像 $R \to R'$ が与えられたとし、
$S' = R' \otimes_R S$ とおく。$R' \to S'$ がエタールとなる素イデアルの
集合は、$R \to S$ がエタールとなる素イデアルの集合の、
$\Spec(S') \to \Spec(S)$ による逆像である。
\item エタール環写像はシントミックであり、特に平坦である。
\item $S$ が体 $k$ 上有限型なら、$S$ が $k$ 上エタールであることと
$\Omega_{S/k} = 0$ であることは同値である。
\item 任意のエタール環写像 $R \to S$ は、あるエタール環写像
$R_0 \to S_0$ で $R_0$ が $\mathbf{Z}$ 上有限型であるものの基底変換である。
\item $A = \colim A_i$ を環のフィルター余極限とする。
$A \to B$ をエタール環写像とする。このとき、ある $i$ とエタール環写像
$A_i \to B_i$ が存在して、$B \cong A \otimes_{A_i} B_i$ となる。
\item $A$ を環とし、$S$ を $A$ の積閉部分集合とする。
$S^{-1}A \to B'$ をエタールとする。このときエタール環写像
$A \to B$ で $B' \cong S^{-1}B$ となるものが存在する。
\item $A$ を環とし、$B = B' \times B''$ を $A$-代数の直積とする。
$B$ が $A$ 上エタールであることと、$B'$ と $B''$ の双方が $A$ 上
エタールであることは同値である。
\end{enumerate}
\end{lemma}

\begin{proof}
各場合に、対応する滑らかな環写像についての結果を用い、
$\Omega_{S/R}$ が零であることを示す短い議論を加える。

\medskip\noindent
（1）の証明。環写像 $R \to R_f$ は滑らかであり、
$\Omega_{R_f/R} = 0$ である。

\medskip\noindent
（2）の証明。滑らかな写像 $A \to B$ と $B \to C$ の合成
$A \to C$ は滑らかである。補題 \ref{algebra-lemma-compose-smooth} を参照せよ。
補題 \ref{algebra-lemma-exact-sequence-differentials} により、
$\Omega_{C/B}$ と $\Omega_{B/A}$ はともに零だから
$\Omega_{C/A}$ も零である。

\medskip\noindent
（3）の証明。$R \to S$ をエタールとし、$R \to R'$ を任意とする。
補題 \ref{algebra-lemma-base-change-smooth} により
$R' \to S' = R' \otimes_R S$ は滑らかである。補題
\ref{algebra-lemma-differentials-base-change} により
$\Omega_{S'/R'} = S' \otimes_S \Omega_{S/R}$ だから、
$\Omega_{S'/R'} = 0$ である。従って $R' \to S'$ はエタールである。

\medskip\noindent
（4）の証明。（4）の仮定を置く。補題 \ref{algebra-lemma-locally-smooth} により
$R \to S$ は滑らかである。また、すべての $i$ について
$\Omega_{S_{g_i}/R} = (\Omega_{S/R})_{g_i} = 0$ であることが
与えられている。すると補題 \ref{algebra-lemma-cover} により
$\Omega_{S/R} = 0$ である。

\medskip\noindent
（5）の証明。滑らかな写像についての結果は補題
\ref{algebra-lemma-flat-base-change-locus-smooth} である。その補題の証明では、
$\NL_{S/R} \otimes_S S'$ が $\NL_{S'/R'}$ とホモトピー同値であることを
用いた。従って、有限表示 $S$-加群 $M$ に対し、
$(M \otimes_S S')_{\mathfrak q'} = 0$ となる $S'$ の素イデアル
$\mathfrak q'$ の集合が、$M_{\mathfrak q} = 0$ となる $S$ の素イデアル
$\mathfrak q$ の集合の逆像であることを示せばよい。
これは補題 \ref{algebra-lemma-support-base-change} から従う。

\medskip\noindent
（6）の証明。滑らかな環写像についての対応する結果
（補題 \ref{algebra-lemma-smooth-syntomic}）から直ちに従う。

\medskip\noindent
（7）の証明。補題 \ref{algebra-lemma-characterize-smooth-over-field} と定義から従う。

\medskip\noindent
（8）の証明。補題 \ref{algebra-lemma-finite-presentation-fs-Noetherian} は
滑らかな環写像について結果を与える。得られた滑らかな環写像
$R_0 \to S_0$ は補題 \ref{algebra-lemma-relative-dimension-CM} の仮定を満たす。
従って $S_0$ を $R_0$ 上相対次元 $0$ の因子で置き換えてよい。

\medskip\noindent
（9）の証明。補題 \ref{algebra-lemma-characterize-finite-presentation} により、
ある $i$ に対して $R_0 \to A$ は $A_i$ を経由するから、（8）より従う。

\medskip\noindent
（10）の証明。$S^{-1}A$ は $A$ の単項局所化のフィルター余極限だから、
（9）、（1）、（2）より従う。

\medskip\noindent
（11）の証明。補題 \ref{algebra-lemma-product-smooth} を用いて滑らかさについての
結果を得た後、$\Omega_{B/A}$ が零であることと
$\Omega_{B'/A}$ と $\Omega_{B''/A}$ がともに零であることが同値だと用いる。
\end{proof}

\noindent
次に、体上エタールであることの意味を、より詳しく調べる。

\begin{lemma}
\label{algebra-lemma-etale-over-field}
$k$ を体とする。環写像 $k \to S$ がエタールであることと、
$S$ が $k$-代数として $k$ の有限分離拡大体の有限直積と同型であることは
同値である。
\end{lemma}

\begin{proof}
以下では次の事実を断りなく用いる。$S = S_1 \times \ldots \times S_n$ が
$k$-代数の有限直積なら、$S$ が $k$ 上エタールであることと、各
$S_i$ が $k$ 上エタールであることは同値である。
補題 \ref{algebra-lemma-etale} の（11）を参照せよ。

\medskip\noindent
$k'/k$ が有限分離拡大なら、
$k' = k(\alpha) \cong k[x]/(f)$ と書ける。ここで $f$ は元
$\alpha$ の最小多項式である。$k'$ は $k$ 上分離的だから
$\gcd(f, f') = 1$ である。これは
$\text{d} : k'\cdot f \to k' \cdot \text{d}x$ が同型であることを
含意する。従って $k \to k'$ はエタールである。ゆえに $S$ が
$k$ の有限分離拡大体の有限直積なら、$S$ は $k$ 上エタールである。

\medskip\noindent
逆に、$k \to S$ がエタールであると仮定する。このとき $S$ は
$k$ 上滑らかであり、$\Omega_{S/k} = 0$ である。補題
\ref{algebra-lemma-characterize-smooth-over-field} により、$S$ の各極大イデアル
$\mathfrak m$ に対して $\dim_\mathfrak m \Spec(S) = 0$ である。
従って $\dim(S) = 0$ である。命題
\ref{algebra-proposition-dimension-zero-ring} により、$S$ はアルティン局所環の
有限直積である。先にも用いた補題
\ref{algebra-lemma-characterize-smooth-over-field} により、これらの局所環は体である。
従って $S = k'$ が体であると仮定してよい。ヒルベルトの零点定理
（定理 \ref{algebra-theorem-nullstellensatz}）により、拡大 $k'/k$ は有限である。
$k \to k'$ の滑らかさから、補題 \ref{algebra-lemma-smooth-at-generic-point} により
$k'/k$ は分離拡大である。これで証明が完了する。
\end{proof}

\begin{lemma}
\label{algebra-lemma-etale-at-prime}
$R \to S$ を環写像とする。$\mathfrak q \subset S$ を、$R$ の
$\mathfrak p$ の上にある素イデアルとする。
$S/R$ が $\mathfrak q$ においてエタールなら、
\begin{enumerate}
\item $\mathfrak p S_{\mathfrak q} = \mathfrak qS_{\mathfrak q}$ であり、
これは局所環 $S_{\mathfrak q}$ の極大イデアルである。
\item 体の拡大 $\kappa(\mathfrak q)/\kappa(\mathfrak p)$ は有限分離的である。
\end{enumerate}
\end{lemma}

\begin{proof}
まず、ある $g \in S$、$g \not \in \mathfrak q$ に対して $S$ を $S_g$ で
置き換え、$R \to S$ がエタールであると仮定してよい。
すると $S \otimes_R \kappa(\mathfrak p)$ が
$\kappa(\mathfrak p)$ 上エタールであるという事実を展開することで、
補題 \ref{algebra-lemma-etale-over-field} から結論が従う。
\end{proof}

\begin{lemma}
\label{algebra-lemma-etale-quasi-finite}
エタール環写像は準有限である。
\end{lemma}

\begin{proof}
$R \to S$ をエタール環写像とする。定義により $R \to S$ は有限型である。
任意の素イデアル $\mathfrak p \subset R$ に対し、ファイバー環
$S \otimes_R \kappa(\mathfrak p)$ は $\kappa(\mathfrak p)$ 上エタールであり、
従って $\kappa(\mathfrak p)$ 上有限分離的な体の有限直積である。
% P01-E0875: the frozen English has “a finite products of fields finite separable”.
特に、これは $\kappa(\mathfrak p)$ 上有限である。
従って補題 \ref{algebra-lemma-quasi-finite} により $R \to S$ は準有限である。
\end{proof}

\begin{lemma}
\label{algebra-lemma-characterize-etale}
$R \to S$ を環写像とする。$\mathfrak q$ を、$R$ の素イデアル
$\mathfrak p$ の上にある $S$ の素イデアルとする。次を仮定する。
\begin{enumerate}
\item $R \to S$ は有限表示である。
\item $R_{\mathfrak p} \to S_{\mathfrak q}$ は平坦である。
\item $\mathfrak p S_{\mathfrak q}$ は局所環 $S_{\mathfrak q}$ の
極大イデアルである。
\item 体の拡大 $\kappa(\mathfrak q)/\kappa(\mathfrak p)$ は
有限分離的である。
\end{enumerate}
このとき $R \to S$ は $\mathfrak q$ においてエタールである。
\end{lemma}

\begin{proof}
補題 \ref{algebra-lemma-isolated-point-fibre} を適用し、
$g \in S$、$g \not \in \mathfrak q$ であって、$\mathfrak q$ が
$S_g$ における $\mathfrak p$ の上の唯一の素イデアルとなるものを得る。
$S$ を $S_g$ で置き換える。すると
$S \otimes_R \kappa(\mathfrak p)$ は素イデアルを一つしかもたず、従って
局所環であり、ゆえに
$S_{\mathfrak q}/\mathfrak pS_{\mathfrak q}
\cong \kappa(\mathfrak q)$ に等しい。
補題 \ref{algebra-lemma-flat-fibre-smooth} により、ある
$g \in S$、$g \not \in \mathfrak q$ に対して $R \to S_g$ は滑らかである。
再び $S$ を $S_g$ で置き換えれば、$R \to S$ は滑らかだと仮定してよい。
% P01-E0876: the frozen replacement sentence is a run-on construction.
補題 \ref{algebra-lemma-smooth-syntomic} により、さらに $R \to S$ は標準滑らか、
例えば $S = R[x_1, \ldots, x_n]/(f_1, \ldots, f_c)$ であると仮定してよい。
$S \otimes_R \kappa(\mathfrak p) = \kappa(\mathfrak q)$ の次元は
$0$ だから $n = c$ である。すなわち $R \to S$ はエタールである。
\end{proof}

\noindent
ここで、まったく新しい現象が現れる。

\begin{lemma}
\label{algebra-lemma-map-between-etale}
$R \to S$ と $R \to S'$ がエタールであるとする。このとき任意の
$R$-代数写像 $S' \to S$ はエタールである。
\end{lemma}

\begin{proof}
まず、補題 \ref{algebra-lemma-compose-finite-type} により $S' \to S$ は
有限表示であることに注意する。$\mathfrak q \subset S$ を、
素イデアル $\mathfrak q' \subset S'$ と $\mathfrak p \subset R$ の上にある
素イデアルとする。補題 \ref{algebra-lemma-etale-at-prime} により、環写像
$S'_{\mathfrak q'}/\mathfrak p S'_{\mathfrak q'} \to
S_{\mathfrak q}/\mathfrak p S_{\mathfrak q}$ は
$\kappa(\mathfrak p)$ の有限分離拡大体の間の写像である。
特に平坦である。従って補題 \ref{algebra-lemma-criterion-flatness-fibre} により
$S'_{\mathfrak q'} \to S_{\mathfrak q}$ は平坦である。
従って $S' \to S$ は平坦である。さらに、上の議論から
$\mathfrak q'S_{\mathfrak q}$ は $S_{\mathfrak q}$ の極大イデアルであり、
$S'_{\mathfrak q'} \to S_{\mathfrak q}$ の剰余体拡大は有限分離的である。
従って補題 \ref{algebra-lemma-characterize-etale} から、$S' \to S$ は
$\mathfrak q$ においてエタールであると分かる。エタールであるという性質は
局所的だから（補題 \ref{algebra-lemma-etale} を参照せよ）、結論を得る。
\end{proof}

\begin{lemma}
\label{algebra-lemma-surjective-flat-finitely-presented}
$\varphi : R \to S$ を環写像とする。$R \to S$ が全射、平坦かつ
有限表示なら、冪等元 $e \in R$ で $S = R_e$ となるものが存在する。
% P01-E0877: the frozen existential verb disagrees with the singular idempotent.
\end{lemma}

\begin{proof}[第一の証明]
$I$ を $\varphi$ の核とする。$\varphi$ は有限表示だから、補題
\ref{algebra-lemma-finite-presentation-independent} により $I$ は有限生成である。
さらに $S$ は $R$ 上平坦だから、完全列
$0 \to I \to R \to S \to 0$ と $S$ の $R$ 上のテンソル積を取ると
$I/I^2 = 0$ を得る。ここで補題
\ref{algebra-lemma-ideal-is-squared-union-connected} を適用すればよい。
\end{proof}

\begin{proof}[第二の証明]
$\Spec(S) \to \Spec(R)$ は閉部分集合の上への同相写像であり
（補題 \ref{algebra-lemma-spec-closed} を参照せよ）、かつ開写像だから
（命題 \ref{algebra-proposition-fppf-open} を参照せよ）、その像はある冪等元
$e \in R$ に対する $D(e)$ である（補題
\ref{algebra-lemma-disjoint-decomposition} を参照せよ）。従って
$R_e \to S$ はスペクトル上の全単射を誘導する。この写像はすべての局所環上で
同型を誘導する。例えば補題 \ref{algebra-lemma-finite-flat-local} および補題
\ref{algebra-lemma-NAK} を用いよ。すると $R_e \to S$ は単射でもある。
例えば補題 \ref{algebra-lemma-characterize-zero-local} を参照せよ。
\end{proof}

\begin{lemma}
\label{algebra-lemma-lift-etale}
\begin{slogan}
エタール環写像は環の全射に沿って持ち上がる
\end{slogan}
$R$ を環、$I \subset R$ をイデアルとする。
$R/I \to \overline{S}$ をエタール環写像とする。このときエタール環写像
$R \to S$ で、$R/I$-代数として
$\overline{S} \cong S/IS$ となるものが存在する。
\end{lemma}

\begin{proof}
補題 \ref{algebra-lemma-etale-standard-smooth} により、定義
\ref{algebra-definition-standard-smooth} のように
$\overline{S}\allowbreak =
(R/I)[x_1, \ldots,\allowbreak x_n]\allowbreak/(\overline{f}_1, \ldots,\allowbreak \overline{f}_n)$
と書け、$\overline{\Delta} =
\det(\frac{\partial \overline{f}_i}{\partial x_j})_{i, j = 1, \ldots,\allowbreak n}$
は $\overline{S}$ で可逆である。$f_i$ の持ち上げをいくつか選び、
$S\allowbreak = R[x_1, \ldots,\allowbreak x_n, x_{n+1}]\allowbreak/(f_1, \ldots,\allowbreak f_n, x_{n + 1}\Delta - 1)$
とおけばよい。ここで
$\Delta\allowbreak = \det(\frac{\partial f_i}{\partial x_j})_{i, j = 1, \ldots,\allowbreak n}$
である。例 \ref{algebra-example-make-standard-smooth} を参照せよ。
これで補題が示された。
\end{proof}

\begin{lemma}
\label{algebra-lemma-lift-etale-infinitesimal}
可換図式
$$
\xymatrix{
0 \ar[r] &
J \ar[r] &
B' \ar[r] &
B \ar[r] & 0 \\
0 \ar[r] &
I \ar[r] \ar[u] &
A' \ar[r] \ar[u] &
A \ar[r] \ar[u] & 0
}
$$
を考える。各行は完全であり、$B' \to B$ と $A' \to A$ は、核の平方が
零イデアルとなる全射な環写像であるとする。$A \to B$ がエタールで、
$J = I \otimes_A B$ なら、$A' \to B'$ はエタールである。
\end{lemma}

\begin{proof}
補題 \ref{algebra-lemma-lift-etale} により、エタール環写像 $A' \to C$ で
$C/IC = B$ となるものが存在する。すると $A' \to C$ は形式的に滑らか
だから（命題 \ref{algebra-proposition-smooth-formally-smooth}）、
$A'$-代数写像 $\varphi : C \to B'$ を得る。
$A' \to C$ は平坦だから
$I \otimes_A B = I \otimes_A C/IC = IC$ である。従って
$J = I \otimes_A B$ という仮定により、$\varphi$ は同型
$IC \to J$ と同型 $C/IC \to B'/IB'$ を誘導し、ゆえに
$\varphi$ は同型である。
\end{proof}

\begin{example}
\label{algebra-example-factor-polynomials-etale}
$n , m \geq 1$ を整数とする。例 \ref{algebra-example-factor-polynomials} の環写像
% P01-E0878: the frozen mathematical span has anomalous spacing before its comma.
\begin{eqnarray*}
R = \mathbf{Z}[a_1, \ldots, a_{n + m}]
& \longrightarrow &
S = \mathbf{Z}[b_1, \ldots, b_n, c_1, \ldots, c_m] \\
a_1 & \longmapsto & b_1 + c_1 \\
a_2 & \longmapsto & b_2 + b_1 c_1 + c_2 \\
\ldots & \ldots & \ldots \\
a_{n + m} & \longmapsto & b_n c_m
\end{eqnarray*}
を考える。記号的に
$$
S = R[b_1, \ldots, c_m]/(\{a_k(b_i, c_j) - a_k\}_{k = 1, \ldots, n + m})
$$
と書く。例えば $a_1(b_i, c_j) = b_1 + c_1$ である。
偏微分の行列は
$$
\left(
\begin{matrix}
1 & c_1 & \ldots & c_m & 0 & \ldots & \ldots & 0 \\
0 & 1 & c_1 & \ldots & c_m & 0 & \ldots & 0 \\
\ldots & \ldots & \ldots & \ldots & \ldots & \ldots & \ldots & \ldots \\
0 & \ldots & 0 & 1 & c_1 & c_2 & \ldots & c_m \\
1 & b_1 & \ldots & b_{n - 1} & b_n & 0 & \ldots & 0 \\
0 & 1 & b_1 & \ldots & b_{n - 1} & b_n & \ldots & 0 \\
\ldots & \ldots & \ldots & \ldots & \ldots & \ldots & \ldots & \ldots \\
0 & \ldots & \ldots & 0 & 1 & b_1 & \ldots & b_n
\end{matrix}
\right)
$$
である。この行列の行列式 $\Delta$ は、多項式
$g = x^n + b_1 x^{n - 1} + \ldots + b_n$ と
$h = x^m + c_1 x^{m - 1} + \ldots + c_m$ の {\it 終結式} として
よく知られている。また上の行列は、$g, h$ に付随する
{\it シルベスター行列} と呼ばれる。式で書けば
$\Delta = \text{Res}_x(g, h)$ である。シルベスター行列は線形写像
\begin{eqnarray*}
S[x]_{< m} \oplus S[x]_{< n} & \longrightarrow & S[x]_{< n + m} \\
a \oplus b & \longmapsto & ag + bh
\end{eqnarray*}
の行列の転置である。$\mathfrak q \subset S$ を任意の素イデアルとする。
上の議論により、次の条件は同値である。
\begin{enumerate}
\item $R \to S$ は $\mathfrak q$ においてエタールである。
\item $\Delta = \text{Res}_x(g, h) \not \in \mathfrak q$ である。
\item 多項式 $g, h$ の像
$\overline{g}, \overline{h} \in \kappa(\mathfrak q)[x]$ は
$\kappa(\mathfrak q)[x]$ において互いに素である。
\end{enumerate}
（2）と（3）の同値性は、シルベスター行列の
$\text{Mat}(n + m, \kappa(\mathfrak q))$ における像が核をもつことと、
多項式 $\overline{g}, \overline{h}$ が共通因子をもつことが同値であるために
成り立つ。従って環写像
$$
R \longrightarrow S[\frac{1}{\Delta}] = S[\frac{1}{\text{Res}_x(g, h)}]
$$
はエタールである。
\end{example}


\begin{lemma}
\label{algebra-lemma-factor-mod-lift-etale}
$R$ を環とする。$f \in R[x]$ をモニック多項式とする。
$\mathfrak p$ を $R$ の素イデアルとする。
$f \bmod \mathfrak p = \overline{g} \overline{h}$ を、$f$ の
$\kappa(\mathfrak p)[x]$ における像の因数分解とする。
$\gcd(\overline{g}, \overline{h}) = 1$ なら、次が存在する。
\begin{enumerate}
\item エタール環写像 $R \to R'$。
\item $\mathfrak p$ の上にある素イデアル $\mathfrak p' \subset R'$。
\item $R'[x]$ における因数分解 $f = g h$。
\end{enumerate}
さらに、これらは次を満たす。
\begin{enumerate}
\item $\kappa(\mathfrak p) = \kappa(\mathfrak p')$。
\item $\overline{g} = g \bmod \mathfrak p'$ および
$\overline{h} = h \bmod \mathfrak p'$。
\item 多項式 $g, h$ は $R'[x]$ で単位イデアルを生成する。
\end{enumerate}
\end{lemma}

\begin{proof}
$\overline{g} = \overline{b}_0 x^n + \overline{b}_1 x^{n - 1} + \ldots
+ \overline{b}_n$ および
$\overline{h} = \overline{c}_0 x^m + \overline{c}_1 x^{m - 1} + \ldots
+ \overline{c}_m$ とし、$\overline{b}_0, \overline{c}_0 \in \kappa(\mathfrak p)$
は零でないとする。$R$ を $\mathfrak p$ に含まれない $R$ のある元で
局所化した後、$\overline{b}_0$ は可逆元 $b_0 \in R$ の像であると
仮定してよい。$\overline{g}$ を $\overline{g}/b_0$ で、
$\overline{h}$ を $b_0\overline{h}$ で置き換えれば、
$\overline{g}$ と $\overline{h}$ がモニックである場合に帰着する
（確認は省略する）。
$\overline{g} = x^n + \overline{b}_1 x^{n - 1} + \ldots + \overline{b}_n$
および
$\overline{h} = x^m + \overline{c}_1 x^{m - 1} + \ldots + \overline{c}_m$
とする。また
$f = x^{n + m} + a_1 x^{n + m - 1} + \ldots + a_{n + m}$ と書く。
ファイバー積
$$
R' = R \otimes_{\mathbf{Z}[a_1, \ldots, a_{n + m}]}
\mathbf{Z}[b_1, \ldots, b_n, c_1, \ldots, c_m]
$$
を考える。ここで写像 $\mathbf{Z}[a_k] \to \mathbf{Z}[b_i, c_j]$ は、
例 \ref{algebra-example-factor-polynomials} および
\ref{algebra-example-factor-polynomials-etale} のものとする。構成により
$R$-代数写像
$$
R' = R \otimes_{\mathbf{Z}[a_1, \ldots, a_{n + m}]}
\mathbf{Z}[b_1, \ldots, b_n, c_1, \ldots, c_m]
\longrightarrow
\kappa(\mathfrak p)
$$
で $b_i$ を $\overline{b}_i$ に、$c_j$ を $\overline{c}_j$ に写すものがある。
この写像の核を $\mathfrak p' \subset R'$ とおく。
% P01-E0879: the frozen designation of p' omits “by/as”.
仮定により多項式 $\overline{g}, \overline{h}$ は互いに素だから、元
$\Delta = \text{Res}_x(g, h) \in \mathbf{Z}[b_i, c_j]$
（例 \ref{algebra-example-factor-polynomials-etale} を参照せよ）は、表示した写像の
もとで $\kappa(\mathfrak p)$ の零に写らない。従って $R \to R'$ は
$\mathfrak p'$ においてエタールである。実際、補題で与えた問題の解は
環写像 $R \to R'[1/\Delta]$ と素イデアル
$\mathfrak p' R'[1/\Delta]$ である。
% P01-E0880: localization is at Delta = Res_x(g,h), but the frozen proof next invokes Res_x(f,g).
$\text{Res}_x(f, g)$ はこの環で可逆だから、シルベスター行列は
$R'[1/\Delta]$ 上可逆であり、従ってある $a, b \in R'[1/\Delta][x]$ に対して
$1 = a g +  b h$ である。例 \ref{algebra-example-factor-polynomials-etale} を参照せよ。
\end{proof}

\section{エタール環写像の局所構造}
\label{algebra-section-etale-local-structure}

\noindent
補題 \ref{algebra-lemma-etale-standard-smooth} が示すように、標準エタール射を
相対次元 $0$ の標準滑らかな射として定義しても、あまり意味はない。
エタール射のモデルとして、体の有限分離拡大 $k'/k$ が与える例を採る。
すなわち、ある元 $\alpha \in k'$ で $k' = k(\alpha)$ となり、
$\alpha$ の最小多項式 $f(x) \in k[x]$ の導多項式 $f'$ が $f$ と
互いに素となるものを、常に見つけられる。

\begin{definition}
\label{algebra-definition-standard-etale}
$R$ を環とし、$g , f  \in R[x]$ とする。
% P01-E0881: the frozen mathematical list has anomalous spacing around its comma.
$f$ はモニックであり、導多項式 $f'$ は局所化 $R[x]_g/(f)$ で
可逆であると仮定する。この場合、環写像 $R \to R[x]_g/(f)$ を
{\it 標準エタール} という。
\end{definition}

\noindent
命題 \ref{algebra-proposition-etale-locally-standard} で、任意のエタール環写像が
局所的に標準エタールであることを示す。

\begin{lemma}
\label{algebra-lemma-standard-etale}
$R \to R[x]_g/(f)$ を標準エタールとする。
\begin{enumerate}
\item 環写像 $R \to R[x]_g/(f)$ はエタールである。
\item 任意の環写像 $R \to R'$ に対し、標準エタール環写像
$R \to R[x]_g/(f)$ の基底変換 $R' \to R'[x]_g/(f)$ は
標準エタールである。
\item $R[x]_g/(f)$ の任意の単項局所化は $R$ 上標準エタールである。
\item 標準エタール写像の合成は一般には {\bf 標準エタールでない}。
\end{enumerate}
\end{lemma}

\begin{proof}
省略する。（4）の例を挙げる。
環写像 $\mathbf{F}_2 \to \mathbf{F}_{2^2}$ は標準エタールである。
環写像
$\mathbf{F}_{2^2} \to \mathbf{F}_{2^2} \times \mathbf{F}_{2^2}
\times \mathbf{F}_{2^2} \times \mathbf{F}_{2^2}$ は標準エタールである。
しかし環写像
$\mathbf{F}_2 \to \mathbf{F}_{2^2} \times \mathbf{F}_{2^2}
\times \mathbf{F}_{2^2} \times \mathbf{F}_{2^2}$ は標準エタールでない。
\end{proof}

\noindent
標準エタール射は、エタール写像を作る便利な方法である。例を挙げよう。

\begin{lemma}
\label{algebra-lemma-make-etale-map-prescribed-residue-field}
$R$ を環とする。$\mathfrak p$ を $R$ の素イデアルとする。
$L/\kappa(\mathfrak p)$ を有限分離拡大体とする。このとき、エタール環写像
$R \to R'$ と $\mathfrak p$ の上にある素イデアル $\mathfrak p'$ で、
体の拡大 $\kappa(\mathfrak p')/\kappa(\mathfrak p)$ が
$\kappa(\mathfrak p) \subset L$ と同型となるものが存在する。
\end{lemma}

\begin{proof}
原始元定理により、$L = \kappa(\mathfrak p)[\alpha]$ と書ける。
$\overline{f} \in \kappa(\mathfrak p)[x]$ を $\alpha$ の最小多項式とする
（特にこれはモニックである）。ある $c \in R$、
$c\not \in \mathfrak p$ に対して $\alpha$ を $c\alpha$ で置き換えれば、
$\overline{f}$ の係数はすべて $R \to \kappa(\mathfrak p)$ の像に属すると
仮定してよい（確認は省略する）。従って、
$\kappa(\mathfrak p)[x]$ で $\overline{f}$ に写るモニック多項式
$f \in R[x]$ を見つけられる。$\kappa(\mathfrak p) \subset L$ は
分離的だから $\gcd(\overline{f}, \overline{f}') = 1$ である。
従って、ある元 $\gamma \in L$ が存在して
$\overline{f}'(\alpha) \gamma = 1$ となる。ゆえに $R$-代数写像
\begin{eqnarray*}
R[x, 1/f']/(f) & \longrightarrow & L \\
x & \longmapsto & \alpha \\
1/f' & \longmapsto & \gamma
\end{eqnarray*}
を得る。左辺は $R$ 上の標準エタール代数 $R'$ であり、環写像の核が
求める素イデアルを与える。
\end{proof}

\begin{proposition}
\label{algebra-proposition-etale-locally-standard}
$R \to S$ を環写像とする。$\mathfrak q \subset S$ を素イデアルとする。
$R \to S$ が $\mathfrak q$ においてエタールなら、ある
$g \in S$、$g \not \in \mathfrak q$ が存在して、$R \to S_g$ は
標準エタールである。
\end{proposition}

\begin{proof}
次の証明はいささか回りくどく、短くする方法があるかもしれない。

\medskip\noindent
ステップ 1。定義 \ref{algebra-definition-etale} により、ある
$g \in S$、$g \not \in \mathfrak q$ に対して $R \to S_g$ は
エタールである。従って $S$ は $R$ 上エタールであると仮定してよい。

\medskip\noindent
ステップ 2。補題 \ref{algebra-lemma-etale} により、$\mathbf{Z}$ 上有限型の
$R_0$ を始域とするエタール環写像 $R_0 \to S_0$ と環写像
$R_0 \to R$ が存在して、$R = R \otimes_{R_0} S_0$ となる。
% P01-E0882: the frozen base-change identity has R on both the left and first tensor factor.
$\mathfrak q$ に対応する $S_0$ の素イデアルを $\mathfrak q_0$ とおく。
% P01-E0883: the frozen designation of q_0 omits “by/as”.
$(R_0 \to S_0, \mathfrak q_0)$ について結果を示せば、基底変換により
$(R \to S, \mathfrak q)$ についても従う。従って $R$ はネーターであると
仮定してよい。

\medskip\noindent
ステップ 3。補題 \ref{algebra-lemma-etale-quasi-finite} により
$R \to S$ は準有限である。補題
\ref{algebra-lemma-quasi-finite-open-integral-closure} により、有限環写像
$R \to S'$、$R$-代数写像 $S' \to S$、および元
$g' \in S'$、$g' \not \in \mathfrak q$ で、$S' \to S$ が同型
$S'_{g'} \cong S_{g'}$ を誘導するものが存在する。
% P01-E0884: the frozen sentence compares g' in S' directly with q in S and repeats “such that”.
（もちろん、一般に $S'$ は $R$ 上エタールではない。）
従って、（a）$R$ はネーター、（b）$R \to S$ は有限、（c）
$R \to S$ は $\mathfrak q$ においてエタールであると仮定してよい
（ただし、もはやすべての素イデアルにおいてエタールとは限らない）。

\medskip\noindent
ステップ 4。$\mathfrak p \subset R$ を $\mathfrak q$ に対応する
素イデアルとする。ファイバー環
$S \otimes_R \kappa(\mathfrak p)$ を考える。これは
$\kappa(\mathfrak p)$ 上の有限代数である。従ってアルティン環であり
（補題 \ref{algebra-lemma-finite-dimensional-algebra} を参照せよ）、局所環の有限積
$$
S \otimes_R \kappa(\mathfrak p) = \prod\nolimits_{i = 1}^n A_i
$$
である（命題 \ref{algebra-proposition-dimension-zero-ring} を参照せよ）。因子の一つ、
例えば $A_1$ は局所環
$S_{\mathfrak q}/\mathfrak pS_{\mathfrak q}$ であり、
$\kappa(\mathfrak q)$ と同型である（補題
\ref{algebra-lemma-etale-at-prime} を参照せよ）。他の因子は、
$\mathfrak p$ の上にある $S$ の他の素イデアル、例えば
$\mathfrak q_2, \ldots, \mathfrak q_n$ に対応する。

\medskip\noindent
ステップ 5。有限分離体拡大
$\kappa(\mathfrak q)/\kappa(\mathfrak p)$ を生成する非零元
$\alpha \in \kappa(\mathfrak q)$ を選べる（従って体拡大が自明な場合でも
$\alpha = 0$ は許さない）。任意の
$\lambda \in \kappa(\mathfrak p)^*$ に対して、元
$\lambda \alpha$ も $\kappa(\mathfrak p)$ 上
$\kappa(\mathfrak q)$ を生成することに注意せよ。元
$$
\overline{t} =
(\alpha, 0, \ldots, 0) \in
\prod\nolimits_{i = 1}^n A_i =
S \otimes_R \kappa(\mathfrak p).
$$
を考える。必要なら上のように $\alpha$ を $\lambda \alpha$ で置き換えることで、
$\overline{t}$ はある $t \in S$ の像であると仮定してよい。
$x$ を $t$ に写す $R$-代数写像 $R[x] \to S$ の核を
$I \subset R[x]$ とする。$S' = R[x]/I$ とおけば、$S' \subset S$ である。
図式は次のとおりである。
$$
\xymatrix{
R[x] \ar[r] & S' \ar[r] & S \\
R \ar[u] \ar[ru] \ar[rru] & &
}
$$
構成により、$S$ の素イデアル $\mathfrak q_j$（$j \geq 2$）はすべて
$R[x]$ の素イデアル $(\mathfrak p, x)$ の上にある。一方、
$\alpha \not = 0$ だから、$\mathfrak q$ は $R[x]$ の別の素イデアルの上にある。

\medskip\noindent
ステップ 6。$\mathfrak q$ に対応する $S'$ の素イデアルを
$\mathfrak q' \subset S'$ と書く。
% P01-E0886: the frozen designation of q' omits “by/as”.
上で見たことから、$\mathfrak q$ は $\mathfrak q'$ の上にある $S$ の唯一の
素イデアルである。従って
$S_{\mathfrak q} = S_{\mathfrak q'}$ である。補題
\ref{algebra-lemma-unique-prime-over-localize-below} を参照せよ（$R \to S$ が
有限なので $S' \to S$ は有限であり、補題
\ref{algebra-lemma-integral-going-up} により $S' \to S$ に対して上昇定理が成り立つ）。
ゆえに $S'_{\mathfrak q'} \to S_{\mathfrak q}$ は、有限単射な環写像
$S' \to S$ の局所化として有限かつ単射である。局所環の写像
$$
R_{\mathfrak p} \to S'_{\mathfrak q'} \to S_{\mathfrak q}
$$
を考える。第 2 の写像は有限かつ単射である。また
$S_{\mathfrak q}/\mathfrak pS_{\mathfrak q} = \kappa(\mathfrak q)$ である
（補題 \ref{algebra-lemma-etale-at-prime} を参照せよ）。従ってなおさら
$S_{\mathfrak q}/\mathfrak q'S_{\mathfrak q} = \kappa(\mathfrak q)$ である。
さらに
$$
\kappa(\mathfrak p) \subset \kappa(\mathfrak q') \subset \kappa(\mathfrak q)
$$
であり、$\alpha$ は $\kappa(\mathfrak q')$ から
$\kappa(\mathfrak q)$ への像に属するので、
$\kappa(\mathfrak q') = \kappa(\mathfrak q)$ と結論できる。
従って $S'_{\mathfrak q'}$-加群写像
$S'_{\mathfrak q'} \to S_{\mathfrak q}$ に中山の補題
\ref{algebra-lemma-NAK} を適用すると、写像
$S'_{\mathfrak q'} \to S_{\mathfrak q}$ は全射である。言い換えれば
$S'_{\mathfrak q'} \cong S_{\mathfrak q}$ である。

\medskip\noindent
ステップ 7。補題 \ref{algebra-lemma-isomorphic-local-rings} により、ある
$g \in S$、$g \not \in \mathfrak q$ および
$g' \in S'$、$g' \not \in \mathfrak q'$ が存在して
$S'_{g'} \cong S_g$ となる。$R$ はネーター環であり、$S'$ は有限
$R$-加群 $S$ の $R$-部分加群だから、この環は $R$ 上有限である。
従って $S$ を $S'$ で置き換えることで、（a）$R$ はネーター、
（b）$S$ は $R$ 上有限、（c）$S$ は $\mathfrak q$ において
$R$ 上エタール、（d）$S = R[x]/I$ と仮定してよい。
% P01-E0888: the frozen condition (b) omits “is”.

\medskip\noindent
ステップ 8。環
$S \otimes_R \kappa(\mathfrak p) = \kappa(\mathfrak p)[x]/\overline{I}$
を考える。ここで
$\overline{I} = I \cdot \kappa(\mathfrak p)[x]$ は
$\kappa(\mathfrak p)[x]$ において $I$ が生成するイデアルである。
$\kappa(\mathfrak p)[x]$ は単項イデアル整域なので、あるモニック多項式
$\overline{h} \in \kappa(\mathfrak p)[x]$ に対して
$\overline{I} = (\overline{h})$ である。ある
$\lambda \in \kappa(\mathfrak p)$ によって $\overline{h}$ を
$\lambda \cdot \overline{h}$ で置き換えることで、$\overline{h}$ はある
$h \in I \subset R[x]$ の像であると仮定してよい（問題は、$h$ を
モニックに選べるかどうかが分からないことである）。またステップ 4 と同様に、
$S \otimes_R \kappa(\mathfrak p) = A_1 \times \ldots \times A_n$ であり、
$A_1 = \kappa(\mathfrak q)$ は $\kappa(\mathfrak p)$ の有限分離拡大、
$A_2, \ldots, A_n$ は局所環である。これより
$$
\overline{h} = \overline{h}_1 \overline{h}_2^{e_2} \ldots \overline{h}_n^{e_n}
$$
と書ける。ここで $\overline{h}_i \in \kappa(\mathfrak p)[x]$ は
相異なる互いに素な既約モニック多項式であり、
$e_2, \ldots, e_n \geq 1$ である。番号は
% P01-E0889: the frozen formula uses e_i also for i=1 although only e_2,...,e_n were introduced.
$A_i = \kappa(\mathfrak p)[x]/(\overline{h}_i^{e_i})$ が
$\kappa(\mathfrak p)[x]$-代数として成り立つように選んだ。
$\overline{h}_1$ は $\alpha \in \kappa(\mathfrak q)$ の最小多項式であり、
従って分離多項式である（その導多項式は自身と互いに素である）。

\medskip\noindent
ステップ 9。$m \in I$ をモニックな元とする。このような元は、環拡大
$R \to R[x]/I$ が有限、従って整だから存在する。
$\kappa(\mathfrak p)[x]$ における像を $\overline{m}$ と書く。
% P01-E0890: the frozen designation of the image omits “by/as”.
次のように因数分解できる。
$$
\overline{m} = \overline{k}
\overline{h}_1^{d_1} \overline{h}_2^{d_2} \ldots \overline{h}_n^{d_n}
$$
ここで $d_1 \geq 1$、$d_j \geq e_j$（$j = 2, \ldots, n$）であり、
$\overline{k} \in \kappa(\mathfrak p)[x]$ はすべての
$\overline{h}_i$ と互いに素である。
$l \deg(m) > \deg(h)$ かつ $l \geq 2$ となるように
$f = m^l + h$ とおく。このとき $f$ は $R$ 上のモニック多項式である。
また、$\kappa(\mathfrak p)[x]$ における $f$ の像 $\overline{f}$ は
$$
\overline{f} =
\overline{h}_1 \overline{h}_2^{e_2} \ldots \overline{h}_n^{e_n}
+
\overline{k}^l \overline{h}_1^{ld_1} \overline{h}_2^{ld_2}
\ldots \overline{h}_n^{ld_n}
=
\overline{h}_1(\overline{h}_2^{e_2} \ldots \overline{h}_n^{e_n}
+
\overline{k}^l
\overline{h}_1^{ld_1 - 1} \overline{h}_2^{ld_2} \ldots \overline{h}_n^{ld_n})
= \overline{h}_1 \overline{w}
$$
と因数分解される。ここで $\overline{w}$ は $\overline{h}_1$ と互いに素な
多項式である。$g = f'$（$x$ に関する導多項式）とおく。

\medskip\noindent
ステップ 10。環写像 $R[x] \to S = R[x]/I$ は次の性質をもつ。
（1）$f$ を零に写し、（2）$g$ を $S \setminus \mathfrak q$ の元に写す。
第 1 の主張は $f$ が $I$ の元なので明らかである。第 2 の主張については、
$g$ が
$\kappa(\mathfrak q) = \kappa(\mathfrak p)[x]/(\overline{h}_1)$
で零に写らないことを示せばよい。$\kappa(\mathfrak p)[x]$ における
$g$ の像は $\overline{f}$ の導多項式である。従って（2）は
$$
\overline{g} =
\frac{\text{d}\overline{f}}{\text{d}x} =
\overline{w}\frac{\text{d}\overline{h}_1}{\text{d}x} +
\overline{h}_1\frac{\text{d}\overline{w}}{\text{d}x},
$$
$\overline{w}$ は $\overline{h}_1$ と互いに素であり、
$\overline{h}_1$ は分離的であることから明らかである。

\medskip\noindent
ステップ 11。$\varphi : R[x]/(f) \to S$ は全射な環写像であり、
$R[x]_g/(f)$ は $R$ 上エタールであり（標準エタールだからである。補題
\ref{algebra-lemma-standard-etale} を参照せよ）、かつ
$\varphi(g) \not \in \mathfrak q$ であると結論できる。
さらに $\varphi(g') \not \in \mathfrak q$ であり、
$S_{\varphi(g')}$ が $R$ 上エタールとなる元
$g' \in R[x]/(f)$ を選ぶ（$S$ は $\mathfrak q$ において
$R$ 上エタールなので、このような元は存在する）。すると環写像
$R[x]_{gg'}/(f) \to S_{\varphi(gg')}$ は $R$ 上のエタール代数の間の
全射である。従って補題 \ref{algebra-lemma-map-between-etale} によりエタールである。
ゆえに補題 \ref{algebra-lemma-surjective-flat-finitely-presented} により局所化である。
従って、$\mathfrak q$ に属さない元による $S$ のある局所化は、
$R$ 上の標準エタール代数のある局所化と同型である。これは示すべきことであった。
\end{proof}

\noindent
次の二つの補題は、エタール位相がザリスキー被覆と有限平坦射によって生成される
位相より粗いことを述べている。初読時には飛ばしてよい。

\begin{lemma}
\label{algebra-lemma-standard-etale-finite-flat-Zariski}
$R \to S$ を標準エタール射とする。次の性質をもつ環写像
$R \to S'$ が存在する。
% P01-E0891: the frozen list introduction omits terminal punctuation.
\begin{enumerate}
\item $R \to S'$ は有限、有限表示、かつ平坦である
（言い換えれば $S'$ は有限射影的 $R$-加群である）。
\item $\Spec(S') \to \Spec(R)$ は全射である。
\item $\mathfrak p \subset R$ の上にある任意の素イデアル
$\mathfrak q \subset S$ と、$\mathfrak p$ の上にある任意の素イデアル
$\mathfrak q' \subset S'$ に対して、ある
$g' \in S'$、$g' \not \in \mathfrak q'$ が存在し、環写像
$R \to S'_{g'}$ は写像 $\varphi : S \to S'_{g'}$ を介して因数分解し、
$\varphi^{-1}(\mathfrak q'S'_{g'}) = \mathfrak q$ となる。
\end{enumerate}
\end{lemma}

\begin{proof}
定義 \ref{algebra-definition-standard-etale} のような $S$ の表示を
$S = R[x]_g/(f)$ とする。
$a_i \in R$ を用いて
$f = x^n + a_1 x^{n - 1} + \ldots + a_n$ と書く。
補題 \ref{algebra-lemma-adjoin-roots} により、ある有限自由かつ忠実平坦な環写像
$R \to S'$ が存在し、ある $\alpha_i \in S'$ に対して
$f = \prod (x - \alpha_i)$ となる。従って $R \to S'$ は条件（1）、（2）を満たす。
$\mathfrak q \subset R[x]/(f)$ を、$g \not \in \mathfrak q$ を満たす
素イデアルとする（すなわち $S$ の素イデアルに対応する）。
$\mathfrak p = R \cap \mathfrak q$ とおき、
$\mathfrak q' \subset S'$ を $\mathfrak p$ の上にある素イデアルとする。
$R$-代数の写像が $n$ 個あり、それらは
\begin{eqnarray*}
\varphi_i : R[x]/(f) & \longrightarrow & S' \\
x & \longmapsto & \alpha_i
\end{eqnarray*}
で与えられることに注意せよ。証明を終えるには、ある $i$ に対して
（a）$\varphi_i(g)$ の $\kappa(\mathfrak q')$ における像が零でなく、
（b）$\varphi_i^{-1}(\mathfrak q') = \mathfrak q$ であることを示せばよい。
実際、その $i$ に対して単に $g' = \varphi_i(g)$、$\varphi = \varphi_i$ ととれる。

\medskip\noindent
$\kappa(\mathfrak p)[x]$ における $f$ の像を $\overline{f}$ と書く。
$\Spec(\kappa(\mathfrak p)[x]/(\overline{f}))$ の点として、
素イデアル $\mathfrak q$ は $\overline{f}$ の既約因子 $f_1$ に対応する。
さらに $g \not \in \mathfrak q$ とは、$f_1$ が
$\kappa(\mathfrak p)[x]$ における $g$ の像 $\overline{g}$ を割らないことを
意味する。$\kappa(\mathfrak q')$ における
$\alpha_1, \ldots, \alpha_n$ の像を
$\overline{\alpha}_1, \ldots, \overline{\alpha}_n$ と書く。
% P01-E0892: the frozen designation of the root images omits “by/as”.
多項式 $\overline{f}$ は $\kappa(\mathfrak q')[x]$ において完全分解し、実際
$$
\overline{f} = \prod\nolimits_i (x - \overline{\alpha}_i)
$$
である。また $\varphi_i(g)$ は $\overline{g}(\overline{\alpha}_i)$ に帰着する。
従って $f_1(\overline{\alpha}_i) = 0$ かつ
$\overline{g}(\overline{\alpha}_i) \not = 0$ となる $i$ を選べる。
この $i$ に対して性質（a）、（b）が成り立つ。詳細の一部は省略する。
\end{proof}

\begin{lemma}
\label{algebra-lemma-etale-finite-flat-zariski}
$R \to S$ を環写像とする。次を仮定する。
\begin{enumerate}
\item $R \to S$ はエタールである。
\item $\Spec(S) \to \Spec(R)$ は全射である。
\end{enumerate}
このとき次の性質をもつ環写像 $R \to S'$ が存在する。
\begin{enumerate}
\item $R \to S'$ は有限、有限表示、かつ平坦である
（言い換えれば、有限射影的 $R$-加群である）。
\item $\Spec(S') \to \Spec(R)$ は全射である。
\item 任意の素イデアル $\mathfrak q' \subset S'$ に対して、ある
$g' \in S'$、$g' \not \in \mathfrak q'$ が存在し、環写像
$R \to S'_{g'}$ は $R \to S \to S'_{g'}$ と因数分解する。
\end{enumerate}
\end{lemma}

\begin{proof}
命題 \ref{algebra-proposition-etale-locally-standard} と
$\Spec(S)$ の準コンパクト性（補題 \ref{algebra-lemma-quasi-compact} を参照せよ）により、
$S$ の単位イデアルを生成する $g_1, \ldots, g_n \in S$ を、各
$R \to S_{g_i}$ が標準エタールとなるように選べる。環写像
$R \to \prod_{i = 1, \ldots, n} S_{g_i}$ に対して補題を示せば、環写像
$R \to S$ に対する補題が従う。従って
$S = \prod_{i = 1, \ldots, n} S_i$ は標準エタール射の有限積であると
仮定してよい。

\medskip\noindent
各 $i$ に対し、標準エタール射 $R \to S_i$ に合わせて補題
\ref{algebra-lemma-standard-etale-finite-flat-Zariski} のような環写像
$R \to S_i'$ を選ぶ。
$S' = S_1' \otimes_R \ldots \otimes_R S_n'$ とおく。以下では
$R$-代数写像 $S_i' \to S'$ を断りなく用いる。これでよいことを示す。
性質（1）、（2）は直ちに従う。性質（3）のため、
$\mathfrak q' \subset S'$ を素イデアルとする。
$\Spec(R)$ におけるその像を $\mathfrak p$ と書く。
% P01-E0893: the frozen designation of the base prime omits “by/as”.
$\mathfrak p$ が $\Spec(S_i) \to \Spec(R)$ の像に属するような
$i \in \{1, \ldots, n\}$ を選ぶ。これは仮定により可能である。
$S_i'$ のスペクトルにおける $\mathfrak q'$ の像を
$\mathfrak q_i' \subset S_i'$ とおく。
% P01-E0894: the frozen designation of the contracted prime omits “by/as”.
$S'_i$ の構成により、ある $g'_i \in S_i'$ が存在し、
% P01-E0895: the frozen proof omits the required condition g_i' notin q_i'.
$R \to (S_i')_{g_i'}$ は
$R \to S_i \to (S_i')_{g_i'}$ と因数分解する。従って
$R \to S'_{g_i'}$ も
$$
R \to S_i \to (S_i')_{g_i'} \to S'_{g_i'}
$$
と因数分解する。これは求める性質である。
\end{proof}

\section{準有限環写像のエタール局所構造}
\label{algebra-section-etale-local-quasi-finite}

\noindent
次の補題群は、大まかに言えば、準有限環写像がエタール拡大の後に
有限となることを述べる。結果の解釈を助けるため、有限型環写像が
準有限となる軌跡は開であり（補題 \ref{algebra-lemma-quasi-finite-open} を参照せよ）、
この軌跡の形成は任意の基底変換と可換であること（補題
\ref{algebra-lemma-quasi-finite-base-change} を参照せよ）を思い出そう。

\begin{lemma}
\label{algebra-lemma-produce-finite}
$R \to S' \to S$ を環写像とする。
$\mathfrak p \subset R$ を素イデアルとし、$g \in S'$ を元とする。
次を仮定する。
\begin{enumerate}
\item $R \to S'$ は整である。
\item $R \to S$ は有限型である。
\item $S'_g \cong S_g$ である。
\item $g$ は $S' \otimes_R \kappa(\mathfrak p)$ で可逆である。
% P01-E0896: the frozen predicate omits the finite verb “is”.
\end{enumerate}
このとき、ある $f \in R$、$f \not \in \mathfrak p$ が存在して
$R_f \to S_f$ は有限となる。
% P01-E0897: the frozen existential uses an ungrammatical article before f.
\end{lemma}

\begin{proof}
仮定により、$V(g) \subset \Spec(S')$ の射
$\Spec(S') \to \Spec(R)$ による像 $T$ は $\mathfrak p$ を含まない。
節 \ref{algebra-section-going-up}、特に補題 \ref{algebra-lemma-going-up-closed} により
$T$ は閉である。
% P01-E0898: the frozen citation phrase leaves “especially” unattached.
$T \cap D(f) = \emptyset$ となる $f \in R$、$f \not \in \mathfrak p$ を選ぶ。
すると $g$ は $S'_f$ で可逆になる。従って
$S'_f \cong S_f$ である。ゆえに $S_f$ は $R_f$ 上有限型かつ整であり、
従って有限である。
\end{proof}

\begin{lemma}
\label{algebra-lemma-etale-makes-quasi-finite-finite-one-prime}
$R \to S$ を環写像とする。$\mathfrak q \subset S$ を
$\mathfrak p \subset R$ の上にある素イデアルとする。
$R \to S$ は有限型であり、$\mathfrak q$ において準有限であると仮定する。
% P01-E0899: the frozen finite-type predicate omits “is”.
このとき次が存在する。
\begin{enumerate}
\item エタール環写像 $R \to R'$。
\item $\mathfrak p$ の上にある素イデアル $\mathfrak p' \subset R'$。
\item 積分解
$$
R' \otimes_R S = A \times B
$$
\end{enumerate}
これらは次の性質をもつ。
% P01-E0900: the frozen list introduction lacks terminal punctuation.
\begin{enumerate}
\item $\kappa(\mathfrak p) = \kappa(\mathfrak p')$ である。
\item $R' \to A$ は有限である。
\item $A$ は $\mathfrak p'$ の上にある素イデアル
$\mathfrak r$ をちょうど一つもつ。
\item $\mathfrak r$ は $\mathfrak q$ の上にある。
\item $B$ は $\mathfrak q$ と $\mathfrak p'$ の上にある素イデアルをもたない。
\end{enumerate}
\end{lemma}
 
\begin{proof}
$S' \subset S$ を $S$ における $R$ の整閉包とする。
$\mathfrak q' = S' \cap \mathfrak q$ とおく。ザリスキーの主定理
\ref{algebra-theorem-main-theorem} により、ある $g \in S'$、
$g \not \in \mathfrak q'$ が存在して $S'_g \cong S_g$ となる。
ファイバー環 $F = S \otimes_R \kappa(\mathfrak p)$ および
$F' = S' \otimes_R \kappa(\mathfrak p)$ を考える。
$\mathfrak q'$ に対応する $F'$ の素イデアルを
$\overline{\mathfrak q}'$ と書く。
% P01-E0901: the frozen designation of the fibre prime omits “by/as”.
$F'$ は $\kappa(\mathfrak p)$ 上整だから、
$\overline{\mathfrak q}'$ は $\Spec(F')$ の閉点である（補題
\ref{algebra-lemma-integral-over-field} を参照せよ）。また $\mathfrak q$ は
$\Spec(F)$ の孤立閉点 $\overline{\mathfrak q}$ を定めることに注意せよ
（定義 \ref{algebra-definition-quasi-finite} を参照せよ）。
$S'_g \cong S_g$ であるから $F'_g \cong F_g$ であり、従って
$\overline{\mathfrak q}$ と $\overline{\mathfrak q}'$ はそれぞれ
$\Spec(F)$ と $\Spec(F')$ において同型な開近傍をもつ。ゆえに集合
$\{\overline{\mathfrak q}'\} \subset \Spec(F')$ は開である。
上で示した $\mathfrak q'$ の閉性と合わせると、
$\overline{\mathfrak q}'$ も $\Spec(F')$ の孤立閉点を定める。
% P01-E0902: the frozen proof cites q' rather than qbar' as the closed point of Spec(F').

\medskip\noindent
さらに小さな注意として、写像 $\Spec(F) \to \Spec(F')$ のもとで
$\overline{\mathfrak q}'$ に写る点は $\overline{\mathfrak q}$ だけである。
これは上の議論から従う。

\medskip\noindent
補題 \ref{algebra-lemma-disjoint-implies-product} により、
$F' = F'_1 \times F'_2$ かつ
$\Spec(F'_1) = \{\overline{\mathfrak q}'\}$ となるように書ける。
$F' = S' \otimes_R \kappa(\mathfrak p)$ だから、ある $s' \in S'$ が存在し、
ある $r \in R$、$r \not \in \mathfrak p$ に対する元
$(r, 0) \in F'_1 \times F'_2 = F'$ に写る。実際、以下で $s'$ について
用いるのは、それが $S'$ の元であって $\mathfrak q'$ に含まれず、
$\mathfrak p$ の上にある他の任意の素イデアルに含まれることだけである。

\medskip\noindent
$f(x) \in R[x]$ を $f(s') = 0$ となるモニック多項式とする。
その像を $\overline{f} \in \kappa(\mathfrak p)[x]$ と書く。
% P01-E0903: the frozen designation of the polynomial image omits “by/as”.
$\overline{h}(0) \not = 0$ となるように
$\overline{f} = x^e \overline{h}$ と因数分解できる。必要なら $f$ を
$x f$ で置き換えることで、$e \geq 1$ と仮定してよい。
補題 \ref{algebra-lemma-factor-mod-lift-etale} により、エタール環拡大
$R \to R'$、$\mathfrak p$ の上にある素イデアル $\mathfrak p'$、および
$R'[x]$ における因数分解 $f = h i$ で、
$\kappa(\mathfrak p) = \kappa(\mathfrak p')$、
$\overline{h} = h \bmod \mathfrak p'$、
$x^e = i \bmod \mathfrak p'$ を満たし、さらに（適当な $a, b$ に対して）
$R'[x]$ で $a h + b i = 1$ と書けるものを見つけられる。

\medskip\noindent
$h(s'), i(s') \in R' \otimes_R S'$ を考える。構成により
$h(s')i(s') = f(s') = 0$ である。一方、
$a(s')h(s') + b(s')i(s') = 1$ だから、これらは単位イデアルを生成する。
従って $R' \otimes_R S'$ は、これらの元における局所化の積
$$
R' \otimes_R S'
=
(R' \otimes_R S')_{i(s')}
\times
(R' \otimes_R S')_{h(s')}
=
S'_1 \times S'_2
$$
である。さらに、この積分解は先に得たファイバー環 $F'$ の積分解と両立する。
これは、$s', i, h$ を、$\overline{\mathfrak q}'$ が $F'$ の素イデアルのうち
$F'$ における $i(s')$ の像を含まない唯一のものとなるように選んだからである。
ここで、$\mathfrak p'$ における $R'$ 上の $R'\otimes_R S'$ の
ファイバー環は、$\kappa(\mathfrak p) = \kappa(\mathfrak p')$ により
$F'$ と同じであることを用いた。従って $S'_1$ は
% P01-E0904: the frozen source doubles the space before “has exactly”.
$\mathfrak p'$ の上にある素イデアルをちょうど一つもち、それを
$\mathfrak r'$ とすれば、この素イデアルは $\mathfrak q'$ の上にある。
ゆえに元 $g \in S'$ は $S'_1$ の $\mathfrak r'$ に含まれない元に写る。

\medskip\noindent
基底変換 $R'\otimes_R S$ も同様の積分解を受け継ぐ。
$$
R' \otimes_R S
=
(R' \otimes_R S)_{i(s')}
\times
(R' \otimes_R S)_{h(s')}
=
S_1 \times S_2
$$
上の議論から、$S_1$ は $\mathfrak p'$ の上にある素イデアルをちょうど一つ
もつ（上と同様にファイバー環を考えよ）。これを $\mathfrak r$ とすれば、
この素イデアルは $\mathfrak q$ の上にある。

\medskip\noindent
ここで補題 \ref{algebra-lemma-produce-finite} を環写像
$R' \to S'_1 \to S_1$、素イデアル $\mathfrak p'$、および元 $g$ に適用する。
すると、$R'$ を単項局所化で置き換えた後、$S_1$ は $R'$ 上有限であると
仮定できる。これは求める性質である。
\end{proof}

\begin{lemma}
\label{algebra-lemma-etale-makes-quasi-finite-finite}
$R \to S$ を環写像とし、$\mathfrak p \subset R$ を素イデアルとする。
$R \to S$ は有限型であると仮定する。このとき次が存在する。
% P01-E0905: the frozen existential list introduction lacks punctuation.
\begin{enumerate}
\item エタール環写像 $R \to R'$。
\item $\mathfrak p$ の上にある素イデアル $\mathfrak p' \subset R'$。
\item 積分解
$$
R' \otimes_R S = A_1 \times \ldots \times A_n \times B
$$
\end{enumerate}
これらは次の性質をもつ。
% P01-E0906: the frozen properties introduction lacks terminal punctuation.
\begin{enumerate}
\item $\kappa(\mathfrak p) = \kappa(\mathfrak p')$ である。
\item 各 $A_i$ は $R'$ 上有限である。
\item 各 $A_i$ は $\mathfrak p'$ の上にある素イデアル
$\mathfrak r_i$ をちょうど一つもつ。
\item $R' \to B$ は $\mathfrak p'$ の上にあるどの素イデアルにおいても
準有限でない。
% P01-E0907: the frozen quasi-finite predicate omits “is”.
\end{enumerate}
\end{lemma}

\begin{proof}
$F = S \otimes_R \kappa(\mathfrak p)$ を、素イデアル $\mathfrak p$ における
$S/R$ のファイバー環と書く。
% P01-E0908: the frozen designation of F omits “by/as”.
$F$ は $\kappa(\mathfrak p)$ 上有限型だからネーターであり、従って
$\Spec(F)$ は孤立閉点を有限個しかもたない。孤立閉点がない場合、すなわち
$\mathfrak p$ の上にあって $S/R$ が $\mathfrak q$ において準有限となる
$S$ の素イデアル $\mathfrak q$ がない場合、補題は成り立つ。
そのような素イデアル $\mathfrak q$ が少なくとも一つ存在するなら、補題
\ref{algebra-lemma-etale-makes-quasi-finite-finite-one-prime} を適用できる。
これにより、同補題のような図式
$$
\xymatrix{
S \ar[r] & R'\otimes_R S \ar@{=}[r] & A_1 \times B' \\
R \ar[r] \ar[u] & R' \ar[u] \ar[ru]
}
$$
を得る。$\mathfrak p$ と $\mathfrak p'$ における剰余体は同じなので、
$S/R$ と $(A_1 \times B')/R'$ のファイバー環は同じである。従って
ファイバーの孤立閉点の個数に関する帰納法により、補題が
$R' \to B'$ と $\mathfrak p'$ に対して成り立つと仮定してよい。
こうしてエタール環写像 $R' \to R''$、素イデアル
$\mathfrak p'' \subset R''$、および分解
$$
R'' \otimes_{R'} B' = A_2 \times \ldots \times A_n \times B
$$
を得る。環写像 $R \to R''$、素イデアル $\mathfrak p''$、および得られた分解
$$
R'' \otimes_R S = (R'' \otimes_{R'} A_1) \times
A_2 \times \ldots \times A_n \times B
$$
が補題で課した問題の解であることの確認は省略する。
% P01-E0909: the frozen display completing the prior sentence has no terminal punctuation.
\end{proof}

\begin{lemma}
\label{algebra-lemma-etale-makes-quasi-finite-finite-variant}
$R \to S$ を環写像とし、$\mathfrak p \subset R$ を素イデアルとする。
$R \to S$ は有限型であると仮定する。
% P01-E0910: the frozen finite-type predicate omits “is”.
このとき次が存在する。
% P01-E0911: the frozen existential list introduction lacks punctuation.
\begin{enumerate}
\item エタール環写像 $R \to R'$。
\item $\mathfrak p$ の上にある素イデアル $\mathfrak p' \subset R'$。
\item 積分解
$$
R' \otimes_R S = A_1 \times \ldots \times A_n \times B
$$
\end{enumerate}
これらは次の性質をもつ。
% P01-E0912: the frozen properties introduction lacks terminal punctuation.
\begin{enumerate}
\item 各 $A_i$ は $R'$ 上有限である。
\item 各 $A_i$ は $\mathfrak p'$ の上にある素イデアル
$\mathfrak r_i$ をちょうど一つもつ。
\item 有限体拡大 $\kappa(\mathfrak r_i)/\kappa(\mathfrak p')$ は
純非分離的である。
\item $R' \to B$ は $\mathfrak p'$ の上にあるどの素イデアルにおいても
準有限でない。
% P01-E0913: the frozen quasi-finite predicate omits “is”.
\end{enumerate}
\end{lemma}

\begin{proof}
証明の方針は二つのエタール環拡大を作ることである。まず剰余体を制御し、
次に補題 \ref{algebra-lemma-etale-makes-quasi-finite-finite} を適用する。

\medskip\noindent
$F = S \otimes_R \kappa(\mathfrak p)$ を、素イデアル $\mathfrak p$ における
$S/R$ のファイバー環と書く。
% P01-E0914: the frozen designation of F omits “by/as”.
補題 \ref{algebra-lemma-etale-makes-quasi-finite-finite} の証明と同様に、
環写像 $R \to S$ が準有限となる、$R$ の上にある $S$ の素イデアルは
有限個である。これらを $\mathfrak q_1, \ldots, \mathfrak q_n$ とする。
% P01-E0915: the frozen source has “finitely may” rather than “finitely many”.
% P01-E0916: the frozen source says the primes of S lie over R rather than p.
$\kappa(\mathfrak p) \subset L_i \subset \kappa(\mathfrak q_i)$ を、
$\kappa(\mathfrak p) \subset L_i$ が分離的で、体拡大
$\kappa(\mathfrak q_i)/L_i$ が純非分離的となる部分体とする。
$L/\kappa(\mathfrak p)$ を、すべての $L_i$（$i = 1, \ldots, n$）が
埋め込まれる有限ガロア拡大とする。補題
\ref{algebra-lemma-make-etale-map-prescribed-residue-field} により、エタール環拡大
$R \to R'$ と $\mathfrak p$ の上にある素イデアル $\mathfrak p'$ で、
体拡大 $\kappa(\mathfrak p')/\kappa(\mathfrak p)$ が
$\kappa(\mathfrak p) \subset L$ と同型となるものを見つけられる。
従って $\mathfrak p'$ における $R' \otimes_R S$ のファイバー環は
$F \otimes_{\kappa(\mathfrak p)} L$ と同型である。
$\mathfrak q_i$ の上にある素イデアルは
$\kappa(\mathfrak q_i) \otimes_{\kappa(\mathfrak p)} L$ の素イデアルに
対応する。これは $L$ の選択と初等的な体論により、$L$ 上純非分離な体の
積である。また補題 \ref{algebra-lemma-quasi-finite-base-change} により、これらは
$R' \to R' \otimes_R S$ が準有限となる $\mathfrak p'$ 上の素イデアルの
すべてである。従って、$R$ を $R'$ で、$\mathfrak p$ を
$\mathfrak p'$ で、$S$ を $R' \otimes_R S$ で置き換えれば、
$\mathfrak p$ の上にあって $S/R$ が準有限となるすべての素イデアル
$\mathfrak q$ に対し、体拡大
$\kappa(\mathfrak q)/\kappa(\mathfrak p)$ は純非分離的であると
仮定してよい。

\medskip\noindent
次に補題 \ref{algebra-lemma-etale-makes-quasi-finite-finite} を適用する。
このエタール環拡大のもとで体拡大は変わらないから、得られる結果が
求めるものである。
\end{proof}

\section{局所準同型}
\label{algebra-section-local-homomorphisms}

\noindent
自然に収まる節のない補題をいくつか述べる。最初の補題は、大まかにいえば、
局所環のエタール写像が、非単元の単項イデアルのすべての冪を法として
同型になることを述べる。

\begin{lemma}
\label{algebra-lemma-lindel}
\begin{reference}
\cite[Lemma on page 321]{Lindel}, \cite[Lemma 4.1.5]{KC}
\end{reference}
$(R, \mathfrak m_R) \to (S, \mathfrak m_S)$ を局所環の局所準同型とする。
$S$ は $R$ のエタール環拡大の局所化であり、
$\kappa(\mathfrak m_R) \to \kappa(\mathfrak m_S)$ は同型であると仮定する。
このとき、ある $t \in \mathfrak m_R$ が存在して、すべての $n \geq 1$ に対し
$R/t^nR \to S/t^nS$ は同型となる。
% P01-E0917: the frozen source uses the wrong indefinite article before t.
% P01-E0918: the frozen source misspells “isomorphism”.
\end{lemma}

\begin{proof}
あるエタール $R$-代数 $T$ と $\mathfrak m_R$ の上にある素イデアル
$\mathfrak q \subset T$ に対して $S = T_{\mathfrak q}$ と書く。
命題 \ref{algebra-proposition-etale-locally-standard} により、
$R \to T$ は標準エタールであると仮定してよい。定義
\ref{algebra-definition-standard-etale} のように $T = R[x]_g/(f)$ と書く。
剰余体に関する仮定により、$x$ と $a$ が
$\kappa(\mathfrak q) = \kappa(\mathfrak m_S) = \kappa(\mathfrak m_R)$
で同じ像をもつような $a \in R$ を選べる。そこで $x$ を $x - a$ で
置き換えれば、$\mathfrak q$ は $T$ において $x$ と
$\mathfrak m_R$ で生成されると仮定してよい。特に
$t = f(0) \in \mathfrak m_R$ である。$t = f(0)$ が求める元であることを示す。

\medskip\noindent
$f = x^d + \sum_{i = 1, \ldots, d - 1} a_i x^i + t$ と書く。
$R \to T$ は標準エタールなので、$a_1$ は $R$ の単元である。
$f$ の導多項式は $T$ で可逆であり、特に $\mathfrak q$ に含まれないからである。
% P01-E0919: the frozen source runs these two finite clauses together.
$h = a_1 + a_2 x + \ldots + a_{d - 1} x^{d - 2} + x^{d - 1} \in R[x]$
とおけば、$R[x]$ において $f = t + xh$ である。
$h \not \in \mathfrak q$ だから、$T$ を $R[x]_{hg}/(f)$ で置き換えてよい。
この置換の後
$$
T/tT = (R/tR)[x]_{hg}/(f) = (R/tR)[x]_{hg}/(xh) = (R/tR)[x]_{hg}/(x)
$$
であり、これは $R/tR$ の商である。補題
\ref{algebra-lemma-surjective-mod-locally-nilpotent} により、すべての $n \geq 1$ に対し
$R/t^nR \to T/t^nT$ は全射である。一方、平坦な局所環写像
$R/t^nR \to S/t^nS$ は、すべての $n$ に対して
$R/t^nR \to T/t^nT$ を介して因数分解することが分かっている。従ってこれらの
写像は単射でもある（局所環の平坦な局所準同型は忠実平坦であり、従って単射で
ある。補題 \ref{algebra-lemma-local-flat-ff} および
\ref{algebra-lemma-faithfully-flat-universally-injective} を参照せよ）。
$S$ は $T$ の局所化だから、$S/t^nS$ は
$T/t^nT = R/t^nR$ を極大イデアルの上にある素イデアルで局所化したものである。
しかしこの環はすでに局所環であるから、証明は完了する。
\end{proof}

\begin{lemma}
\label{algebra-lemma-etale-under-finite-flat}
$(R, \mathfrak m_R) \to (S, \mathfrak m_S)$ を局所環の局所準同型とする。
$S$ は $R$ のエタール環拡大の局所化であると仮定する。このとき、有限、
有限表示、かつ忠実平坦な環写像 $R \to S'$ で、$S'$ の任意の極大イデアル
$\mathfrak m'$ に対し因数分解
$$
R \to S \to S'_{\mathfrak m'}.
$$
が環写像 $R \to S'_{\mathfrak m'}$ に対して存在するものがある。
\end{lemma}

\begin{proof}
あるエタール $R$-代数 $T$ に対して $S = T_{\mathfrak q}$ と書く。命題
\ref{algebra-proposition-etale-locally-standard} により、$T$ は標準エタールであると
仮定してよい。環写像 $R \to T$ に補題
\ref{algebra-lemma-standard-etale-finite-flat-Zariski} を適用して $R \to S'$ を得る。
すると特に、$S'$ の任意の極大イデアル $\mathfrak m'$ に対して、ある
$g' \not \in \mathfrak m'$ に対する因数分解
$\varphi : T \to S'_{g'}$ で、
$\mathfrak q = \varphi^{-1}(\mathfrak m'S'_{g'})$ を満たすものが得られる。
従って $\varphi$ は求める局所環写像
$S \to S'_{\mathfrak m'}$ を誘導する。
\end{proof}

\section{整閉包と滑らかな基底変換}
\label{algebra-section-integral-closure-smooth-base-change}

\begin{lemma}
\label{algebra-lemma-trick}
$R$ を環とし、$f \in R[x]$ をモニック多項式とする。
$R \to B$ を環写像とする。$h \in B[x]/(f)$ が $R$ 上整ならば、元
$f' h$ は $f'h = \sum_i b_i x^i$ と書ける。ここで $b_i \in B$ は
$R$ 上整である。
\end{lemma}

\begin{proof}
環 $B[x]/(f)$ において、ある $r_i \in R$ に対して
$h^e + r_1 h^{e - 1} + \ldots + r_e = 0$ であるとする。
ある有限自由環拡大 $B \subset B'$ と、ある $\alpha_i \in B'$ が存在して
$f = (x - \alpha_1) \ldots (x - \alpha_d)$ となる（補題
\ref{algebra-lemma-adjoin-roots} を参照せよ）。各 $\alpha_i$ は $R$ 上整である。
$h_i \in B$ を用いて
$h = h_0 + h_1 x + \ldots + h_{d - 1} x^{d - 1}$ と表示できる。
このとき普遍的事実として
$$
f' h =
\sum\nolimits_{i = 1, \ldots, d}
h(\alpha_i)
(x - \alpha_1) \ldots \widehat{(x - \alpha_i)} \ldots (x - \alpha_d)
$$
が $B'[x]/(f)$ の元として成り立つ。これは環
$B_{univ} = \mathbf{Z}[\alpha_i, h_j]$ 上で両辺を点 $\alpha_i$ に代入して
証明する（詳細の一部は省略する）。仮定により $h$ は環 $B[x]/(f)$ で
$h^e + r_1 h^{e - 1} + \ldots + r_e = 0$ を満たすから
$$
h(\alpha_i)^e + r_1 h(\alpha_i)^{e - 1} + \ldots + r_e = 0
$$
が $B'$ で成り立つ。従って $h(\alpha_i)$ は $R$ 上整である。
上の式から、$B'[x]/(f)$ において
$f'h \equiv \sum_{j = 0, \ldots, d - 1} b'_j x^j$ と書ける。ここで
$b'_j \in B'$ は $R$ 上整である。しかし $f' h \in B[x]/(f)$ であり、
$1, x, \ldots, x^{d - 1}$ は $B'[x]/(f)$ の $B'$-基底なので、求めるように
$b'_j \in B$ である。
\end{proof}

\begin{lemma}
\label{algebra-lemma-integral-closure-commutes-etale}
$R \to S$ をエタール環写像とし、$R \to B$ を任意の環写像とする。
$A \subset B$ を $B$ における $R$ の整閉包とし、
$A' \subset S \otimes_R B$ を $S \otimes_R B$ における $S$ の整閉包とする。
このとき、標準写像 $S \otimes_R A \to A'$ は同型である。
\end{lemma}

\begin{proof}
$A \subset B$ であり $R \to S$ は平坦なので、写像
$S \otimes_R A \to A'$ は単射である。整閉包を取る操作が局所化と可換である
こと（補題 \ref{algebra-lemma-integral-closure-localize}）を以下で繰り返し用いる。
従って、補題 \ref{algebra-lemma-cover}（$S$-加群の写像が同型であるかを判定する基準）
により、$S$ 上で局所化してよい。よって
$S = R[x]_g/(f) = (R[x]/(f))_g$ が $R$ 上標準エタールであると仮定してよい（命題
\ref{algebra-proposition-etale-locally-standard} を参照せよ）。もう一度局所化を
適用すると、$A'$ は $(A'')_g$ である。ここで $A''$ は
$B[x]/(f)$ における $R[x]/(f)$ の整閉包である。$a \in A''$ とする。
$a$ が $S \otimes_R A$ に属することを示せば十分である。補題
\ref{algebra-lemma-trick} により $f' a = \sum a_i x^i$ と書け、$a_i \in A$ である。
$f'$ は $S$ で可逆である（標準エタール環写像の定義による）から、求めるように
$a \in S \otimes_R A$ である。
\end{proof}

\begin{example}
\label{algebra-example-fourier}
$p$ を素数とする。環拡大
$$
R = \mathbf{Z}[1/p] \subset
R' = \mathbf{Z}[1/p][x]/(x^{p - 1} + \ldots + x + 1)
$$
は次の性質をもつ。$d < p$ ならば、ある元
$\alpha_0, \ldots, \alpha_{d - 1} \in R'$ が存在して
$$
\prod\nolimits_{0 \leq i < j < d} (\alpha_i - \alpha_j)
$$
は $R'$ の単元となる。実際、$i = 0, \ldots, p - 1$ に対し
$\alpha_i$ を $R'$ における $x^i$ の類とすればよい。このとき
$$
T^p - 1 = \prod\nolimits_{i = 0, \ldots, p - 1} (T - \alpha_i)
$$
が $R'[T]$ で成り立つ。実際、円分多項式
$x^{p - 1} + \ldots + x + 1$ は $\mathbf{Q}$ 上既約なので、環
$\mathbf{Q}[x]/(x^{p - 1} + \ldots + x + 1)$ は体であり、
$\alpha_i$ は $T^p - 1$ の相異なる根である。従ってこの等式を得る。
両辺を微分して $T = \alpha_i$ を代入すると
% P01-E0920: the frozen differentiated product repeats alpha_1 at both endpoints.
$$
p \alpha_i^{p - 1}
=
(\alpha_i - \alpha_1) \ldots
\widehat{(\alpha_i - \alpha_i)} \ldots
(\alpha_i - \alpha_1)
$$
を得る。従ってこれは $R'$ で可逆である。
\end{example}

\begin{lemma}
\label{algebra-lemma-integral-closure-commutes-smooth}
$R \to S$ を滑らかな環写像とし、$R \to B$ を任意の環写像とする。
$A \subset B$ を $B$ における $R$ の整閉包とし、
$A' \subset S \otimes_R B$ を $S \otimes_R B$ における $S$ の整閉包とする。
このとき、標準写像 $S \otimes_R A \to A'$ は同型である。
\end{lemma}

\begin{proof}
補題 \ref{algebra-lemma-integral-closure-commutes-etale} の証明と同様に論じれば、
$S$ 上で局所化してよい。従って $R \to S$ は標準滑らかな環写像であると
仮定してよい（補題 \ref{algebra-lemma-smooth-syntomic} を参照せよ）。標準滑らかな
環写像の定義により、$S$ は多項式環 $R[x_1, \ldots, x_n]$ 上エタールである。
エタール環拡大の場合はすでに示したので（補題
\ref{algebra-lemma-integral-closure-commutes-etale}）、$S = R[x]$ の場合に帰着する。
従って示すべきことは
$$
f = \sum b_i x^i
\text{ integral over }R[x]
\Leftrightarrow
\text{each }b_i\text{ integral over }R.
$$
である。右から左への含意は、$B[x]$ において $R[x]$ 上整な元の集合が環であり
（補題 \ref{algebra-lemma-integral-closure-is-ring}）、かつ $x$ を含むことから従う。

\medskip\noindent
$f \in B[x]$ が $R[x]$ 上整であるとし、次数 $< d$ の
$f = \sum_{i < d} b_i x^i$ と書く。整閉包と局所化は可換なので、各 $b_i$ が
$R[1/p]$ 上でも $R[1/q]$ 上でも整となる相異なる素数 $p, q$ が存在することを
示せば十分である。従って、有限自由環拡大 $R \subset R'$ で、$R'$ が
% P01-E0921: the frozen transition does not explicitly pass to R[1/p] or R[1/q].
$\prod_{i < j} (\alpha_i - \alpha_j)$ が $R'$ の単元となるような
$\alpha_1, \ldots, \alpha_d$ を含むものを見つけられる（例
\ref{algebra-example-fourier} を参照せよ）。この場合、普遍的等式
$$
f
=
\sum_i
f(\alpha_i)
\frac{(x - \alpha_1) \ldots \widehat{(x - \alpha_i)} \ldots (x - \alpha_d)}
{(\alpha_i - \alpha_1) \ldots \widehat{(\alpha_i - \alpha_i)} \ldots
(\alpha_i - \alpha_d)}.
$$
が成り立つ。
% P01-E0922: the frozen source begins the next sentence with the editorial artifact “OK,”.
元 $f(\alpha_i)$ は $R'$ 上整である。実際、
$(R' \otimes_R B)[x] \to R' \otimes_R B$, $h \mapsto h(\alpha_i)$ は環写像である。
従って、$(R' \otimes_R B)[x]$ における $f$ の係数は $R'$ 上整である。
$R'$ は $R$ 上有限（従って $R$ 上整）なので、それらは $R$ 上も整であり、これで
証明が完了する。
\end{proof}

\begin{lemma}
\label{algebra-lemma-integral-closure-commutes-colim-smooth}
$R \to S$ および $R \to B$ を環写像とする。
$A \subset B$ を $B$ における $R$ の整閉包とし、
$A' \subset S \otimes_R B$ を $S \otimes_R B$ における $S$ の整閉包とする。
$S$ が滑らかな $R$-代数のフィルター余極限ならば、標準写像
$S \otimes_R A \to A'$ は同型である。
\end{lemma}

\begin{proof}
テンソル積を取る操作と整閉包を取る操作がフィルター余極限と可換であるという
直接的な事実、および補題
% P01-E0923: the frozen English compound subject has singular “commutes”.
\ref{algebra-lemma-integral-closure-commutes-smooth} から従う。
\end{proof}

\section{形式的に不分岐な写像}
\label{algebra-section-formally-unramified}

\noindent
形式的エタール写像の概念を導入する前に、形式的に不分岐な写像の概念を
定義する方が論理的に効率がよいことが分かる。

\begin{definition}
\label{algebra-definition-formally-unramified}
$R \to S$ を環写像とする。任意の可換な実線図式
$$
\xymatrix{
S \ar[r] \ar@{-->}[rd] & A/I \\
R \ar[r] \ar[u] & A \ar[u]
}
$$
において、$I \subset A$ が平方零イデアルであり、図式を可換にする破線の
矢印が高々一つしか存在しないとき、$S$ は {\it $R$ 上形式的に不分岐である}
という。
\end{definition}

\begin{lemma}
\label{algebra-lemma-base-change-formally-unramified}
$R \to S$ を形式的に不分岐な写像とし、$R \to R'$ を任意の環写像とする。
このとき、基底変換 $S' = R' \otimes_R S$ は $R'$ 上形式的に不分岐である。
\end{lemma}

\begin{proof}
定義 \ref{algebra-definition-formally-unramified} に現れる形の実線図式
$$
\xymatrix{
S \ar[r] \ar@{-->}[rrd] & R' \otimes_R S \ar[r] \ar@{-->}[rd] & A/I \\
R  \ar[u] \ar[r] & R' \ar[r] \ar[u] & A \ar[u]
}
$$
が与えられたとする。仮定により、長い方の破線の矢印は高々一つしか存在しない。
テンソル積の普遍性から、短い方の破線の矢印も高々一つしか存在しない。
\end{proof}

\begin{lemma}
\label{algebra-lemma-characterize-formally-unramified}
$R \to S$ を環写像とする。次は同値である。
\begin{enumerate}
\item $R \to S$ は形式的に不分岐である。
\item 微分加群 $\Omega_{S/R}$ は零である。
\end{enumerate}
\end{lemma}

\begin{proof}
$J = \Ker(S \otimes_R S \to S)$ を乗法写像の核とし、
$A_{univ} = S \otimes_R S/J^2$ とおく。
$I_{univ} = J/J^2$ は $\Omega_{S/R}$ と同型であった（補題
\ref{algebra-lemma-differentials-diagonal} を参照せよ）。さらに、二つの
$R$-代数写像
$\sigma_1, \sigma_2 : S \to A_{univ}$、すなわち
$\sigma_1(s) = s \otimes 1 \bmod J^2$ と
$\sigma_2(s) = 1 \otimes s \bmod J^2$ の差は、普遍導分
$\text{d} : S \to \Omega_{S/R} = I_{univ}$ である。

\medskip\noindent
$R \to S$ が形式的に不分岐であると仮定する。このとき
$\sigma_1 = \sigma_2$ である。従って、すべての $s \in S$ に対して
$\text{d}(s) = 0$ である。よって $\Omega_{S/R} = 0$ である。

\medskip\noindent
$\Omega_{S/R} = 0$ と仮定する。$A, I, R \to A, S \to A/I$ を、
定義 \ref{algebra-definition-formally-unramified} に現れる実線図式とする。
$\tau_1, \tau_2 : S \to A$ を、図式を可換にする二つの破線の矢印とする。
規則 $s_1 \otimes s_2 \mapsto \tau_1(s_1)\tau_2(s_2)$ で定まる
$R$-代数写像 $A_{univ} \to A$ を考える。これが良定義であることの確認は省略する。
$I_{univ} = \Omega_{S/R} = 0$ なので $A_{univ} \cong S$ であり、
$\tau_1 = \tau_2$ を得る。
\end{proof}

\begin{lemma}
\label{algebra-lemma-formally-unramified-local}
$R \to S$ を環写像とする。次は同値である。
\begin{enumerate}
\item $R \to S$ は形式的に不分岐である。
\item $S$ のすべての素イデアル $\mathfrak q$ に対し、
$R \to S_{\mathfrak q}$ は形式的に不分岐である。
\item $S$ のすべての素イデアル $\mathfrak q$ に対し、
$\mathfrak p = R \cap \mathfrak q$ とおけば
$R_{\mathfrak p} \to S_{\mathfrak q}$ は形式的に不分岐である。
\end{enumerate}
\end{lemma}

\begin{proof}
補題 \ref{algebra-lemma-characterize-formally-unramified} で、(1) は
$\Omega_{S/R} = 0$ と同値であることを見た。同様に、補題
\ref{algebra-lemma-differentials-localize} により、(2) と (3) はすべての
$\mathfrak q$ に対して $(\Omega_{S/R})_{\mathfrak q} = 0$ であることと
同値である。従って、同値性は補題 \ref{algebra-lemma-characterize-zero-local}
から従う。
\end{proof}

\begin{lemma}
\label{algebra-lemma-formally-unramified-localize}
$A \to B$ を形式的に不分岐な環写像とする。
\begin{enumerate}
\item $S \subset A$ を乗法的部分集合とすると、
$S^{-1}A \to S^{-1}B$ は形式的に不分岐である。
\item $S \subset B$ を乗法的部分集合とすると、
$A \to S^{-1}B$ は形式的に不分岐である。
\end{enumerate}
\end{lemma}

\begin{proof}
% P01-E0924: the frozen English proof begins with a subjectless sentence fragment.
両方の主張は補題 \ref{algebra-lemma-formally-unramified-local} から従う。
（補題 \ref{algebra-lemma-characterize-formally-unramified} と補題
\ref{algebra-lemma-differentials-localize} を組み合わせても導ける。）
\end{proof}

\begin{lemma}
\label{algebra-lemma-colimit-formally-unramified}
$R$ を環とし、$I$ を有向集合とする。
$(S_i, \varphi_{ii'})$ を $I$ 上の $R$-代数の系とする。
各 $R \to S_i$ が形式的に不分岐ならば、
$S = \colim_{i \in I} S_i$ は $R$ 上形式的に不分岐である。
% P01-E0925: the frozen English lemma statement lacks terminal punctuation.
\end{lemma}

\begin{proof}
定義 \ref{algebra-definition-formally-unramified} に現れる形の図式を考える。
仮定により、合成 $S_i \to S \to A/I$ を持ち上げる
$R$-代数写像 $S_i \to A$ は高々一つしか存在しない。$S$ の各元は
写像 $S_i \to S$ のいずれかの像に属するから、この図式に収まる写像
$S \to A$ も高々一つしか存在しない。
\end{proof}

\section{余法加群と普遍的肥厚化}
\label{algebra-section-conormal}

\noindent
第一無限小近傍は、スキームの閉埋め込みに対してだけでなく、任意の形式的に
不分岐な射に対してすでに定義できることが分かる。その基礎となるのが次の
代数的事実である。

\begin{lemma}
\label{algebra-lemma-universal-thickening}
$R \to S$ を形式的に不分岐な環写像とする。核が平方零イデアルである
$R$-代数の全射 $S' \to S$ で、次の普遍性をもつものが存在する。
$I \subset A$ が平方零イデアルである任意の可換図式
$$
\xymatrix{
S \ar[r]_a & A/I \\
R \ar[r]^b \ar[u] & A \ar[u]
}
$$
が与えられたとき、$S' \to A \to A/I$ が $S' \to S \to A/I$ に等しくなる
一意な $R$-代数写像 $a' : S' \to A$ が存在する。
\end{lemma}

\begin{proof}
$S$ の $R$-代数としての生成元の集合 $z_i \in S$, $i \in I$ を選ぶ。
% P01-E0926: the frozen proof reuses I for both the generator index set and the square-zero ideal.
$P = R[\{x_i\}_{i \in I}]$ を、生成元 $x_i$, $i \in I$ をもつ多項式環とする。
$x_i$ を $z_i$ に写す $R$-代数写像 $P \to S$ を考え、
$J = \Ker(P \to S)$ とおく。写像
$$
\text{d} : J/J^2 \longrightarrow \Omega_{P/R} \otimes_P S
$$
を考える（補題 \ref{algebra-lemma-differential-seq} を参照せよ）。仮定により
$\Omega_{S/R} = 0$ なので、これは全射である（補題
\ref{algebra-lemma-characterize-formally-unramified} を参照せよ）。
$\Omega_{P/R}$ は $\text{d}x_i$ を基底とする自由加群であり、従って
$\Omega_{P/R} \otimes_P S$ は $S$ 上自由であることに注意する。
そこで上の全射の分裂を一つ選び
% P01-E0927: the frozen direct-sum display lacks terminal punctuation.
$$
J/J^2 = K \oplus \Omega_{P/R} \otimes_P S
$$
と書く。$J'/J^2$ が上の分解の第二の直和因子となるような $P$ のイデアル
$J^2 \subset J' \subset J$ を取る。$S' = P/J'$ とおく。短完全列
$$
0 \to J/J' \to S' \to S \to 0
$$
を得て、$J/J' \cong K$ は $S'$ の平方零イデアルであることが分かる。
従って
$$
\xymatrix{
S \ar[r]_1 & S \\
R \ar[r] \ar[u] & S' \ar[u]
}
$$
は上の形の図式である。実際、これは図式の圏における始対象であると主張する。
すなわち、$(I \subset A, a, b)$ を任意の図式とする。可換図式
$$
\xymatrix{
S \ar[r]_1 & S \ar[r]_a & A/I \\
R \ar[r] \ar@/_/[rr]_b \ar[u] & P \ar[u] \ar[r]^\beta & A \ar[u]
}
$$
が得られるような $R$-代数写像 $\beta : P \to A$ を選べる。
ここで $\beta(J') = 0$ とは限らない。言い換えれば、$\beta$ が
$S' = P/J'$ を経由するとは限らない。しかし $\beta(J') \subset I$ であり、
$\beta(J) \subset I$ かつ $I^2 = 0$ なので $\beta(J^2) = 0$ であることは
明らかである。従って、$(J/J' \subset S', 1, R \to S')$ から
$(I \subset A, a, b)$ への射を見つけるための「障害」は、対応する
$S$-線形写像 $\overline{\beta} : J'/J^2 \to I$ である。
$\beta$ を選ぶ自由度は $\beta(x_i)$ の選択にある。$\beta$ の別の選択
$\beta'$ は、$\delta_i \in I$ を用いて
$\beta'(x_i) = \beta(x_i) + \delta_i$ とすることで得られる。
この場合、$g \in J'$ に対して
% P01-E0928: the frozen Taylor display lacks terminal punctuation.
$$
\beta'(g) =
\beta(g) + \sum\nolimits_i \delta_i \beta(\frac{\partial g}{\partial x_i})
$$
を得る。構成により
% P01-E0929: the frozen differential formula has a free i and omits the sum.
$\text{d}|_{J'/J^2} : J'/J^2 \to \Omega_{P/R} \otimes_P S$ は
$g \mapsto \frac{\partial g}{\partial x_i}\text{d}x_i \otimes 1$ で与えられる
同型なので、すべての $g \in J'$ に対して $\beta'(g) = 0$ となる
$\delta_i \in I$ の選択が一意に存在する。（すなわち、$\delta_i$ は
$-\overline{\beta}(g)$ である。ここで $g \in J'/J^2$ は
$\text{d}|_{J'/J^2}$ により $\text{d}x_i \otimes 1$ に写る一意な元である。）
% P01-E0930: the frozen parenthetical leaves the i-dependence of g implicit.
解の一意性から、補題で要求される一意性が従う。
\end{proof}

\noindent
補題 \ref{algebra-lemma-universal-thickening} の状況では、
$R$-代数写像 $S' \to S$ は一意な同型を除いて一意である。

\begin{definition}
\label{algebra-definition-universal-thickening}
$R \to S$ を形式的に不分岐な環写像とする。
\begin{enumerate}
\item $R$ 上の $S$ の {\it 普遍的一次肥厚化} とは、補題
\ref{algebra-lemma-universal-thickening} の $R$-代数の全射 $S' \to S$ のことである。
\item $R \to S$ の {\it 余法加群} とは、普遍的一次肥厚化
$S' \to S$ の核 $I$ を $S$-加群とみなしたものである。
\end{enumerate}
この状況で余法加群をしばしば {\it $C_{S/R}$} と記す。
\end{definition}

\begin{lemma}
\label{algebra-lemma-universal-thickening-quotient}
$I \subset R$ を環のイデアルとする。$R$ 上の $R/I$ の普遍的一次肥厚化は
全射 $R/I^2 \to R/I$ である。$R$ 上の $R/I$ の余法加群は
$C_{(R/I)/R} = I/I^2$ である。
\end{lemma}

\begin{proof}
省略する。
\end{proof}

\begin{lemma}
\label{algebra-lemma-universal-thickening-localize}
$A \to B$ を形式的に不分岐な環写像とする。
$\varphi : B' \to B$ を $A$ 上の $B$ の普遍的一次肥厚化とする。
\begin{enumerate}
\item $S \subset A$ を乗法的部分集合とする。このとき
$S^{-1}B' \to S^{-1}B$ は $S^{-1}A$ 上の $S^{-1}B$ の普遍的一次肥厚化である。
特に $S^{-1}C_{B/A} = C_{S^{-1}B/S^{-1}A}$ である。
\item $S \subset B$ を乗法的部分集合とする。このとき
$S' = \varphi^{-1}(S)$ は $B'$ の乗法的部分集合であり、
$(S')^{-1}B' \to S^{-1}B$ は $A$ 上の $S^{-1}B$ の普遍的一次肥厚化である。
特に $S^{-1}C_{B/A} = C_{S^{-1}B/A}$ である。
\end{enumerate}
この補題が意味をもつことは補題
\ref{algebra-lemma-formally-unramified-localize} から分かる。
\end{lemma}

\begin{proof}
% P01-E0931: both frozen proof paragraphs begin with sentence fragments.
(1) の記法と仮定を用いる。$(S^{-1}B)' \to S^{-1}B$ を
$S^{-1}A$ 上の $S^{-1}B$ の普遍的一次肥厚化とする。
$S^{-1}B' \to S^{-1}B$ は、核が平方零である $S^{-1}A$-代数の全射である。
従って定義により、$S^{-1}B$ への写像と両立する写像
$(S^{-1}B)' \to S^{-1}B'$ を得る。任意の可換図式
$$
\xymatrix{
B \ar[r] & S^{-1}B \ar[r] & D/I \\
A \ar[r] \ar[u] & S^{-1}A \ar[r] \ar[u] & D \ar[u]
}
$$
を考える。ここで $I \subset D$ は平方零イデアルである。
$B'$ は $A$ 上の $B$ の普遍的一次肥厚化なので、$A$-代数写像
$B' \to D$ を得る。しかし $D$ における $S$ の像が $D$ の可逆元に
写ることは明らかであり、従って両立する写像 $S^{-1}B' \to D$ を得る。
これを $D = (S^{-1}B)'$ に適用すると、写像
$S^{-1}B' \to (S^{-1}B)'$ を得る。この写像が上で述べた写像の逆であることの
確認は省略する。

\medskip\noindent
(2) の記法と仮定を用いる。$(S^{-1}B)' \to S^{-1}B$ を
$A$ 上の $S^{-1}B$ の普遍的一次肥厚化とする。
$(S')^{-1}B' \to S^{-1}B$ は、核が平方零である $A$-代数の全射である。
従って定義により、$S^{-1}B$ への写像と両立する写像
$(S^{-1}B)' \to (S')^{-1}B'$ を得る。任意の可換図式
$$
\xymatrix{
B \ar[r] & S^{-1}B \ar[r] & D/I \\
A \ar[r] \ar[u] & A \ar[r] \ar[u] & D \ar[u]
}
$$
を考える。ここで $I \subset D$ は平方零イデアルである。
$B'$ は $A$ 上の $B$ の普遍的一次肥厚化なので、$A$-代数写像
$B' \to D$ を得る。しかし $D$ における $S'$ の像が $D$ の可逆元に
写ることは明らかであり、従って両立する写像 $(S')^{-1}B' \to D$ を得る。
これを $D = (S^{-1}B)'$ に適用すると、写像
$(S')^{-1}B' \to (S^{-1}B)'$ を得る。この写像が上で述べた写像の逆である
ことの確認は省略する。
\end{proof}

\begin{lemma}
\label{algebra-lemma-differentials-universal-thickening}
$R \to A  \to B$ を環写像とし、$A \to B$ は形式的に不分岐であると仮定する。
$B' \to B$ を $A$ 上の $B$ の普遍的一次肥厚化とする。このとき $B'$ は
$A$ 上形式的に不分岐であり、標準写像
$\Omega_{A/R} \otimes_A B \to \Omega_{B'/R} \otimes_{B'} B$ は同型である。
\end{lemma}

\begin{proof}
原理的には定義から形式的にこれらの結果を導けるはずだが、補題
\ref{algebra-lemma-universal-thickening} の証明における $B'$ の構成を用いる。
すなわち、$P = A[x_i]$ が $A$ 上の多項式環となる表示 $B = P/J$ を選ぶ。
次に、$\Omega_{P/A} \otimes_P B$ において
$\text{d}f_i = \text{d}x_i \otimes 1$ となる元 $f_i \in J$ を選ぶ。
これらを選ぶと、$J' = (f_i) + J^2$ として $B' = P/J'$ である（補題
\ref{algebra-lemma-universal-thickening} の証明を参照せよ）。

\medskip\noindent
標準完全列
$$
J'/(J')^2 \to \Omega_{P/A} \otimes_P B' \to \Omega_{B'/A} \to 0
$$
を考える（補題 \ref{algebra-lemma-differential-seq} を参照せよ）。構成により、
$f_i \in J'$ の類は、構成により $J'/J^2$ を法として加群
$\Omega_{P/A} \otimes_P B'$ を生成する元に写る。
% P01-E0932: the frozen Nakayama step uses J'/J^2 where the nilpotent kernel is J/J'.
$J'/J^2$ は冪零イデアルなので、これらの元は加群全体を生成する
（中山の補題 \ref{algebra-lemma-NAK} による）。従って $\Omega_{B'/A} = 0$ であり、
$B'$ は $A$ 上形式的に不分岐である（補題
\ref{algebra-lemma-characterize-formally-unramified} を参照せよ）。

\medskip\noindent
$P$ は $A$ 上の多項式環なので
$\Omega_{P/R} = \Omega_{A/R} \otimes_A P \oplus \bigoplus P\text{d}x_i$
である。この分解を用いる。完全列
$$
J'/(J')^2 \to
\Omega_{P/R} \otimes_P B' \to
\Omega_{B'/R} \to 0
$$
を考える（補題 \ref{algebra-lemma-differential-seq} を参照せよ）。これに $B$ を
テンソルして完全列
% P01-E0933: the frozen tensored exact-sequence display lacks terminal punctuation.
$$
J'/(J')^2 \otimes_{B'} B \to
\Omega_{P/R} \otimes_P B \to
\Omega_{B'/R} \otimes_{B'} B \to 0
$$
を得る。$J' = (f_i) + J^2$ を思い出すと、最初の矢印は部分加群
$J^2/(J')^2$ を零化することが分かる。上で与えた直和分解
$\Omega_{P/R} \otimes_P B =
\Omega_{A/R} \otimes_A B \oplus \bigoplus B\text{d}x_i $ に関して、
% P01-E0934: the frozen prose omits “above,” after “given”.
部分加群 $(f_i)/(J')^2 \otimes_{B'} B$ は直和因子
$\bigoplus B\text{d}x_i$ の上へ同型に写ることが分かる。従ってこの完全列に
残るものは求める同型
$\Omega_{A/R} \otimes_A B \to \Omega_{B'/R} \otimes_{B'} B$
である。
\end{proof}

\section{形式的エタール写像}
\label{algebra-section-formally-etale}

\begin{definition}
\label{algebra-definition-formally-etale}
$R \to S$ を環写像とする。$I \subset A$ が平方零イデアルである任意の
可換な実線図式
$$
\xymatrix{
S \ar[r] \ar@{-->}[rd] & A/I \\
R \ar[r] \ar[u] & A \ar[u]
}
$$
に対して、図式を可換にする破線の矢印が一意に存在するとき、
$S$ は {\it $R$ 上形式的エタールである} という。
\end{definition}

\noindent
環写像が形式的エタールであるための必要十分条件は、それが形式的に滑らかで、
かつ形式的に不分岐であることである。

\begin{lemma}
\label{algebra-lemma-base-change-formally-etale}
$R \to S$ を形式的エタール写像とし、$R \to R'$ を任意の環写像とする。
このとき、基底変換 $S' = R' \otimes_R S$ は $R'$ 上形式的エタールである。
\end{lemma}

\begin{proof}
補題 \ref{algebra-lemma-base-change-fs} と
\ref{algebra-lemma-base-change-formally-unramified} を組み合わせればよい。
\end{proof}

\begin{lemma}
\label{algebra-lemma-formally-etale-etale}
$R \to S$ を有限表示の環写像とする。次は同値である。
\begin{enumerate}
\item $R \to S$ は形式的エタールである。
\item $R \to S$ はエタールである。
\end{enumerate}
\end{lemma}

\begin{proof}
$R \to S$ が形式的エタールであると仮定する。命題
\ref{algebra-proposition-smooth-formally-smooth} により $R \to S$ は滑らかである。
補題 \ref{algebra-lemma-characterize-formally-unramified} により
$\Omega_{S/R} = 0$ である。従って定義により $R \to S$ はエタールである。

\medskip\noindent
$R \to S$ がエタールであると仮定する。命題
\ref{algebra-proposition-smooth-formally-smooth} により $R \to S$ は形式的に滑らかである。
補題 \ref{algebra-lemma-characterize-formally-unramified} により、それは形式的に
不分岐である。従って $R \to S$ は形式的エタールである。
\end{proof}

\begin{lemma}
\label{algebra-lemma-colimit-formally-etale}
$R$ を環とし、$I$ を有向集合とする。
$(S_i, \varphi_{ii'})$ を $I$ 上の $R$-代数の系とする。各
$R \to S_i$ が形式的エタールならば、
$S = \colim_{i \in I} S_i$ は $R$ 上形式的エタールである。
% P01-E0935: the frozen English lemma conclusion lacks terminal punctuation.
\end{lemma}

\begin{proof}
定義 \ref{algebra-definition-formally-etale} に現れる形の図式を考える。
% P01-E0936: the frozen English proof obscures the per-index uniqueness.
仮定により、各 $S_i \to S \to A/I$ の合成を持ち上げる一意な
$R$-代数写像 $S_i \to A$ を得る。従ってこれらは遷移写像
$\varphi_{ii'}$ と両立し、持ち上げ $S \to A$ を定める。これで存在が示された。
一意性は各 $S_i$ への制限を取れば明らかである。
\end{proof}

\begin{lemma}
\label{algebra-lemma-localization-formally-etale}
$R$ を環とし、$S \subset R$ を任意の乗法的部分集合とする。
このとき環写像 $R \to S^{-1}R$ は形式的エタールである。
\end{lemma}

\begin{proof}
$I \subset A$ を平方零イデアルとする。ここで述べているのは、すべての
$f \in S$ に対して $\varphi(f) \mod I$ が可逆となる環写像
$\varphi : R \to A$ が与えられれば、すべての $f \in S$ に対して
$\varphi(f)$ も $A$ で可逆になるということである。これは、標準写像
$A \to A/I$ による $(A/I)^*$ の逆像が $A^*$ であることから従う。
\end{proof}

\begin{lemma}
\label{algebra-lemma-formally-etale-lift-infinitesimal}
$R \to S$ を環写像とする。$J \subset S$ を $R \to S/J$ が全射となる
イデアルとし、$I \subset R$ をその核とする。$R \to S$ が形式的エタール
ならば、すべての $n$ に対して $R/I^n \to S/J^n$ は同型であり、
$\bigoplus I^n/I^{n + 1} \to \bigoplus J^n/J^{n + 1}$ は次数付き環の
同型である。
\end{lemma}

\begin{proof}
持ち上げ性を帰納的に用いて、次の破線の矢印を得る。
% P01-E0937: the frozen three-diagram display lacks terminal punctuation.
$$
\xymatrix{
S \ar[r] \ar@{-->}[rd] & S/J = R/I \\
R \ar[r] \ar[u] & R/I^2 \ar[u]
}
\quad
\xymatrix{
S \ar[r] \ar@{-->}[rd] & R/I^2 \\
R \ar[r] \ar[u] & R/I^3 \ar[u]
}
\quad
\xymatrix{
S \ar[r] \ar@{-->}[rd] & R/I^3 \\
R \ar[r] \ar[u] & R/I^4 \ar[u]
}
$$
対応する写像 $S/J^n \to R/I^n$ は同型である。実際、合成
$S/J^n \to R/I^n \to S/J^n$ は、形式的エタール環写像の持ち上げ性に
おける一意性により（帰納的に）恒等写像となる。
\end{proof}

\begin{lemma}
\label{algebra-lemma-formally-etale-omega}
$R \to S \to S'$ を環写像とする。$J$ と $J'$ をそれぞれ乗法写像
$S \otimes_R S \to S$ と $S' \otimes_{R'} S' \to S'$ の核とする。
% P01-E0938: the frozen statement uses the undefined base ring R' in the second tensor product.
$S \to S'$ が形式的エタールならば、すべての $k \geq 0$ に対して写像
$$
S' \otimes_S \left((S \otimes_R S)/J^{k + 1}\right)
\longrightarrow
(S' \otimes_R S')/(J')^{k + 1}
$$
は同型である。特に、写像
$S' \otimes_S \Omega_{S/R} \to \Omega_{S'/R}$ は同型である。
\end{lemma}

\begin{proof}
$S' \otimes_S (S \otimes_R S) = S' \otimes_R S$ であり、イデアル $J$ が
$S' \otimes_R S$ に生成するイデアルは全射
$S' \otimes_R S \to S'$ の核 $I$ であることに注意する。
従って左辺は $S' \otimes_R S/I^{k + 1}$ に等しい。写像
$S' \otimes_R S \to S' \otimes_R S'$ は補題
\ref{algebra-lemma-base-change-formally-etale} により形式的エタールである。
そこで補題 \ref{algebra-lemma-formally-etale-lift-infinitesimal} により、補題の主張で
表示された矢印は同型である。最後の主張はこのことと、
$\Omega_{S/R} = J/J^2$ および $\Omega_{S'/R} = J'/(J')^2$ であること
（補題 \ref{algebra-lemma-differentials-diagonal}）から従う。
\end{proof}

\begin{lemma}
\label{algebra-lemma-formally-etale-principal-parts}
$R \to S \to S'$ を環写像とし、$S \to S'$ は形式的エタール
（例えばエタール）であるとする。$M$ を $S$-加群とし、
$M' = S' \otimes_R M$ とおく。
% P01-E0939: the frozen statement tensors M over R although extension along S requires tensoring over S.
このとき
% P01-E0940: the frozen displayed identity lacks terminal punctuation.
$$
S' \otimes_S P^k_{S/R}(M) = P^k_{S'/R}(M')
$$
が成り立つ。従って、任意の $S$-加群 $N$ と任意の有限階微分作用素
$D : M \to N$ に対して、$D$ の（一意な）拡張
$D' : M' \to S' \otimes_S N$ で、同じかそれ以下の階数をもつ微分作用素で
あるものが存在する。
\end{lemma}

\begin{proof}
$J$ と $J'$ を補題 \ref{algebra-lemma-formally-etale-omega} の主張におけるものとする。
このとき
\begin{align*}
S' \otimes_S P^k_{S/R}(M)
& =
S' \otimes_S \left((S \otimes_R M)/J^{k + 1}(S \otimes_R M)\right) \\
& =
S' \otimes_S \left((S \otimes_R S)/J^{k + 1}\right) \otimes_S M \\
& =
\left(S' \otimes_R S'/(J')^{k + 1}\right) \otimes_S M \\
& =
\left(S' \otimes_R S'/(J')^{k + 1}\right) \otimes_{S'} M' \\
& =
S' \otimes_R M'/(J')^{k + 1}(S' \otimes_R M') \\
& =
P^k_{S'/R'}(M')
\end{align*}
% P01-E0941: the frozen final equality uses undefined R' instead of R.
最初と最後の等式は補題 \ref{algebra-lemma-principal-parts-diagonal} による。
第三の等式は補題 \ref{algebra-lemma-formally-etale-omega} である。
最後の主張は次のように従う。$D$ が線形写像
$\gamma : P^k_{S/R}(M) \to N$ に対応するならば、
$D' : M' \to S' \otimes_R N$ を、線形写像
$1 \otimes \gamma : S' \otimes_R P^k_{S/R}(M) \to S' \otimes_R N$ に
対応させればよい。
% P01-E0942: the frozen proof tensors over R where the lemma requires extension over S.
\end{proof}

\begin{remark}
\label{algebra-remark-formally-etale-differential-operators}
$R \to S \to S'$ を環写像とし、$S \to S'$ は形式的エタール
（例えばエタール）であるとする。$M_i$, $i = 1, 2, 3$ を $S$-加群とし、
$D_i : M_i \to M_{i + 1}$, $i = 1, 2$ を有限階微分作用素とする。補題
\ref{algebra-lemma-formally-etale-principal-parts} のように、$D_i$ の
$M'_i = S' \otimes_S M_i$ への拡張を
$D'_i : M'_i \to M'_{i + 1}$, $i = 1, 2$ とすれば、
$D'_2 \circ D'_1$ は $D_2 \circ D_1$ の拡張である。特に、$M$ が
$S$-加群ならば、$M' = S' \otimes_S M$ は $S$-代数
$\text{Diff}_{S/R}(M, M)$ 上の加群である。
\end{remark}
\section{不分岐環写像}
\label{algebra-section-unramified}

\noindent
ここで用いる不分岐環写像の定義は \cite{Henselian} による。
ここで G-不分岐環写像と呼ぶものは、EGA が不分岐写像と呼ぶものである。
% P01-E0943: the frozen source has the incorrect article “an G-unramified ring map”.

\begin{definition}
\label{algebra-definition-unramified}
$R \to S$ を環写像とする。
\begin{enumerate}
\item $R \to S$ が有限型であり、かつ $\Omega_{S/R} = 0$ であるとき、
$R \to S$ は {\it 不分岐} であるという。
\item $R \to S$ が有限表示であり、かつ $\Omega_{S/R} = 0$ であるとき、
$R \to S$ は {\it G-不分岐} であるという。
\item $S$ の素イデアル $\mathfrak q$ が与えられたとする。
$g \in S$, $g \not \in \mathfrak q$ であって $R \to S_g$ が不分岐となる
ものが存在するとき、$S$ は {\it $\mathfrak q$ において不分岐} であるという。
\item $S$ の素イデアル $\mathfrak q$ が与えられたとする。
$g \in S$, $g \not \in \mathfrak q$ であって $R \to S_g$ が G-不分岐となる
ものが存在するとき、$S$ は {\it $\mathfrak q$ において G-不分岐} であるという。
\end{enumerate}
\end{definition}

\noindent
もちろん、G-不分岐写像は不分岐である。

\begin{lemma}
\label{algebra-lemma-formally-unramified-unramified}
$R \to S$ を環写像とする。次は同値である。
\begin{enumerate}
\item $R \to S$ は形式的に不分岐かつ有限型である。
\item $R \to S$ は不分岐である。
\end{enumerate}
さらに、次も同値である。
\begin{enumerate}
\item $R \to S$ は形式的に不分岐かつ有限表示である。
\item $R \to S$ は G-不分岐である。
\end{enumerate}
\end{lemma}

\begin{proof}
補題 \ref{algebra-lemma-characterize-formally-unramified} と定義から従う。
\end{proof}

\begin{lemma}
\label{algebra-lemma-unramified}
不分岐環写像および G-不分岐環写像には次の性質がある。
\begin{enumerate}
\item 不分岐環写像の基底変換は不分岐である。
G-不分岐環写像の基底変換は G-不分岐である。
\item 不分岐環写像の合成は不分岐である。
G-不分岐環写像の合成は G-不分岐である。
\item 任意の主局所化 $R \to R_f$ は G-不分岐かつ不分岐である。
\item $I \subset R$ がイデアルならば、$R \to R/I$ は不分岐である。
$I \subset R$ が有限生成イデアルならば、$R \to R/I$ は G-不分岐である。
\item エタール環写像は G-不分岐かつ不分岐である。
\item $R \to S$ が有限型（それぞれ有限表示）、
$\mathfrak q \subset S$ が素イデアル、かつ
$(\Omega_{S/R})_{\mathfrak q} = 0$ ならば、$R \to S$ は
$\mathfrak q$ において不分岐（それぞれ G-不分岐）である。
\item $R \to S$ が有限型（それぞれ有限表示）、
$\mathfrak q \subset S$ が素イデアル、かつ
$\Omega_{S/R} \otimes_S \kappa(\mathfrak q) = 0$ ならば、
$R \to S$ は $\mathfrak q$ において不分岐（それぞれ G-不分岐）である。
\item $R \to S$ が有限型（それぞれ有限表示）、
$\mathfrak q \subset S$ が $\mathfrak p \subset R$ の上にある素イデアル、かつ
$(\Omega_{S \otimes_R \kappa(\mathfrak p)/\kappa(\mathfrak p)})_{\mathfrak q}
= 0$ ならば、$R \to S$ は $\mathfrak q$ において不分岐
（それぞれ G-不分岐）である。
\item $R \to S$ が有限型（それぞれ有限表示）、
% P01-E0944: the frozen source says “resp. presentation”, omitting “finite”.
$\mathfrak q \subset S$ が $\mathfrak p \subset R$ の上にある素イデアル、かつ
$(\Omega_{S \otimes_R \kappa(\mathfrak p)/\kappa(\mathfrak p)})
\otimes_{S \otimes_R \kappa(\mathfrak p)} \kappa(\mathfrak q) = 0$
ならば、$R \to S$ は $\mathfrak q$ において不分岐
（それぞれ G-不分岐）である。
\item $R \to S$ を環写像とし、$g_1, \ldots, g_m \in S$ が単位イデアルを
生成するとする。$j = 1, \ldots, m$ に対して $R \to S_{g_j}$ が不分岐
（それぞれ G-不分岐）ならば、$R \to S$ は不分岐
（それぞれ G-不分岐）である。
\item 環写像 $R \to S$ が $S$ のすべての素イデアルにおいて不分岐
（それぞれ G-不分岐）ならば、$R \to S$ は不分岐
（それぞれ G-不分岐）である。
\item $R \to S$ が G-不分岐ならば、有限型 $\mathbf{Z}$-代数 $R_0$、
G-不分岐環写像 $R_0 \to S_0$、および環写像 $R_0 \to R$ であって、
$S = R \otimes_{R_0} S_0$ となるものが存在する。
\item $R \to S$ が不分岐ならば、有限型 $\mathbf{Z}$-代数 $R_0$、
不分岐環写像 $R_0 \to S_0$、および環写像 $R_0 \to R$ であって、
$S$ が $R \otimes_{R_0} S_0$ の商となるものが存在する。
\end{enumerate}
\end{lemma}

\begin{proof}
各項を順に証明する。

\medskip\noindent
(1) について。補題 \ref{algebra-lemma-differentials-base-change} および
\ref{algebra-lemma-base-change-finiteness} から従う。

\medskip\noindent
(2) について。補題 \ref{algebra-lemma-exact-sequence-differentials} および
\ref{algebra-lemma-base-change-finiteness} から従う。
% P01-E0945: the frozen proof of composition cites the base-change finiteness lemma; finite-type/finite-presentation composition is lemma-compose-finite-type.

\medskip\noindent
(3) について。省略する $\Omega_{R_f/R}$ の直接計算から従う。

\medskip\noindent
(4) について。補題 \ref{algebra-lemma-trivial-differential-surjective} を参照すれば
$\Omega_{(R/I)/R} = 0$ であり、環写像 $R \to R/I$ は有限型である。
$I$ が有限生成イデアルならば、$R \to R/I$ は有限表示である。

\medskip\noindent
(5) について。定義 \ref{algebra-definition-etale} に続く議論を参照せよ。

\medskip\noindent
(6) について。この場合、$\Omega_{S/R}$ は有限 $S$-加群である
（補題 \ref{algebra-lemma-differentials-finitely-generated} を参照せよ）。従って
$g \in S$, $g \not \in \mathfrak q$ であって
$(\Omega_{S/R})_g = 0$ となるものが存在する。補題
\ref{algebra-lemma-differentials-localize} により、これは
$\Omega_{S_g/R} = 0$ を意味し、従って求める通り $R \to S_g$ は
不分岐である。

\medskip\noindent
(7) について。中山の補題（補題 \ref{algebra-lemma-NAK}）を用いれば、
この条件が (6) の条件と同値であることが分かる。

\medskip\noindent
(8) および (9) について。これらは (6) と (7) が同値であるのと同じ仕方で
同値である。さらに補題 \ref{algebra-lemma-differentials-base-change} により
$\Omega_{S \otimes_R \kappa(\mathfrak p)/\kappa(\mathfrak p)} =
\Omega_{S/R} \otimes_S (S \otimes_R \kappa(\mathfrak p))$ である。
従って、両者に現れる $\kappa(\mathfrak q)$-ベクトル空間が標準的に同型なので、
(9) は (7) と同値である。

\medskip\noindent
(10) について。補題 \ref{algebra-lemma-cover} および
\ref{algebra-lemma-differentials-localize} から従う。

\medskip\noindent
(11) について。(6)、(7)、および $S$ のスペクトルが準コンパクトであることから
従う。
% P01-E0946: the frozen proof cites (6) and (7), which assume global finite type/presentation; the local-at-every-prime hypothesis plus quasi-compactness instead yields a finite principal cover and invokes (10).

\medskip\noindent
(12) について。$S = R[x_1, \ldots, x_n]/(g_1, \ldots, g_m)$ と書く。
$\Omega_{S/R} = 0$ なので、ある
$h_{ij}, a_{ijk} \in R[x_1, \ldots, x_n]$ によって
$\Omega_{R[x_1, \ldots, x_n]/R}$ の中で
$$
\text{d}x_i = \sum h_{ij}\text{d}g_j + \sum a_{ijk}g_j\text{d}x_k
$$
と書ける。多項式 $g_i, h_{ij}, a_{ijk}$ のすべての係数を含む有限生成
$\mathbf{Z}$-部分代数 $R_0 \subset R$ を選ぶ。
$S_0 = R_0[x_1, \ldots, x_n]/(g_1, \ldots, g_m)$ とおけばよい。

\medskip\noindent
(13) について。$S = R[x_1, \ldots, x_n]/I$ と書く。
$\Omega_{S/R} = 0$ なので、ある
$h_{ij} \in R[x_1, \ldots, x_n]$ および $g_{ij}, g'_{ik} \in I$ によって
$\Omega_{R[x_1, \ldots, x_n]/R}$ の中で
$$
\text{d}x_i = \sum h_{ij}\text{d}g_{ij} + \sum g'_{ik}\text{d}x_k
$$
と書ける。多項式 $g_{ij}, h_{ij}, g'_{ik}$ のすべての係数を含む有限生成
$\mathbf{Z}$-部分代数 $R_0 \subset R$ を選ぶ。
$S_0 = R_0[x_1, \ldots, x_n]/(g_{ij}, g'_{ik})$ とおけばよい。
\end{proof}

\begin{lemma}
\label{algebra-lemma-diagonal-unramified}
$R \to S$ を環写像とする。$R \to S$ が不分岐ならば、冪等元
$e \in S \otimes_R S$ であって、$S \otimes_R S \to S$ が
$S \otimes_R S \to (S \otimes_R S)_e$ と同型になるものが存在する。
\end{lemma}

\begin{proof}
$J = \Ker(S \otimes_R S \to S)$ とおく。仮定と補題
\ref{algebra-lemma-differentials-diagonal} により $J/J^2 = 0$ である。
$S$ は $R$ 上有限型なので $J$ は有限生成である。より具体的には、
$x_i$ が $R$ 上 $S$ を生成するとき、イデアルの生成元として
$x_i \otimes 1 - 1 \otimes x_i$ を取れる。
補題 \ref{algebra-lemma-ideal-is-squared-union-connected} を適用すればよい。
\end{proof}

\begin{lemma}
\label{algebra-lemma-unramified-at-prime}
$R \to S$ を環写像とする。$\mathfrak q \subset S$ を、$R$ の
$\mathfrak p$ の上にある素イデアルとする。$S/R$ が $\mathfrak q$ において
不分岐ならば、
\begin{enumerate}
\item $\mathfrak p S_{\mathfrak q} = \mathfrak qS_{\mathfrak q}$ であり、
これは局所環 $S_{\mathfrak q}$ の極大イデアルである。
% P01-E0947: the frozen source reads “we have pS_q = qS_q is the maximal ideal”, joining two predicates ungrammatically.
\item 体拡大 $\kappa(\mathfrak q)/\kappa(\mathfrak p)$ は有限分離拡大である。
\end{enumerate}
\end{lemma}

\begin{proof}
まずある $g \in S$, $g \not \in \mathfrak q$ によって
$S$ を $S_g$ で置き換え、$R \to S$ が不分岐であると仮定してよい。
補題 \ref{algebra-lemma-unramified} により、基底変換
$S \otimes_R \kappa(\mathfrak p)$ は $\kappa(\mathfrak p)$ 上不分岐である。
補題 \ref{algebra-lemma-characterize-smooth-over-field} により、それは
$\kappa(\mathfrak p)$ 上滑らかであり、従ってエタールである。
それゆえ、補題 \ref{algebra-lemma-etale-over-field} により
$F = S \otimes_R \kappa(\mathfrak p)$ は $\kappa(\mathfrak p)$ の
有限分離拡大の有限直積である。注記 \ref{algebra-remark-local-ring-fibre} の記法と
結果を用いると、
$F_{\overline{\mathfrak q}} = S_\mathfrak q/\mathfrak pS_\mathfrak q$
は $\kappa(\mathfrak q)$ に等しい。これは補題を導く。
\end{proof}

\begin{lemma}
\label{algebra-lemma-unramified-quasi-finite}
$R \to S$ を有限型環写像とし、$\mathfrak q$ を $S$ の素イデアルとする。
$R \to S$ が $\mathfrak q$ において不分岐ならば、$R \to S$ は
$\mathfrak q$ において準有限である。特に、不分岐環写像は準有限である。
\end{lemma}

\begin{proof}
不分岐環写像は有限型である。従って第二の主張が第一の主張から従うことは
明らかである。第一の主張を見るには、補題
\ref{algebra-lemma-isolated-point-fibre} の (2) の特徴づけを、補題
\ref{algebra-lemma-unramified-at-prime} を用いて適用すればよい。
\end{proof}

\begin{lemma}
\label{algebra-lemma-characterize-unramified}
$R \to S$ を環写像とする。$\mathfrak q$ を $S$ の素イデアルであって、
$R$ の素イデアル $\mathfrak p$ の上にあるものとする。次を仮定する。
\begin{enumerate}
\item $R \to S$ は有限型である。
\item $\mathfrak p S_{\mathfrak q}$ は局所環 $S_{\mathfrak q}$ の
極大イデアルである。
\item 体拡大 $\kappa(\mathfrak q)/\kappa(\mathfrak p)$ は有限分離拡大である。
\end{enumerate}
このとき $R \to S$ は $\mathfrak q$ において不分岐である。
\end{lemma}

\begin{proof}
補題 \ref{algebra-lemma-unramified} (8) により、
$\Omega_{S \otimes_R \kappa(\mathfrak p) / \kappa(\mathfrak p)}$ が
$\mathfrak q$ において局所化したとき零であることを示せば十分である。
従って $S$ を $S \otimes_R \kappa(\mathfrak p)$ で、$R$ を
$\kappa(\mathfrak p)$ で置き換えてよい。言い換えれば、$R = k$ が体であり、
$S$ が有限型 $k$-代数であると仮定してよい。この場合、仮定は
$S_{\mathfrak q} \cong \kappa(\mathfrak q)$ を含意する。従って求める通り
$(\Omega_{S/k})_{\mathfrak q} = \Omega_{S_\mathfrak q/k} =
\Omega_{\kappa(\mathfrak q)/k}$ は零である（第一の等式は補題
\ref{algebra-lemma-differentials-localize} による）。
\end{proof}

\begin{lemma}
\label{algebra-lemma-etale-flat-unramified-finite-presentation}
$R \to S$ を環写像とする。次は同値である。
\begin{enumerate}
\item $R \to S$ はエタールである。
\item $R \to S$ は平坦かつ G-不分岐である。
\item $R \to S$ は平坦、不分岐、かつ有限表示である。
\end{enumerate}
\end{lemma}

\begin{proof}
(2) と (3) は定義により同値である。(1) $\Rightarrow$ (3) は、エタール環写像が
有限表示であること、補題 \ref{algebra-lemma-etale}（エタール写像の平坦性）、および
補題 \ref{algebra-lemma-unramified}（エタール写像は不分岐である）から従う。
逆に、補題 \ref{algebra-lemma-characterize-etale} におけるエタール環写像の特徴づけと、
補題 \ref{algebra-lemma-unramified-at-prime} における不分岐環写像の構造から、
(3) は (1) を含意する。
% P01-E0948: the frozen compound subject “the characterization ... and the structure ... shows” has singular agreement.
（ここでは、すべての素イデアル $\mathfrak q \subset S$ において
$R \to S$ がエタールならば $R \to S$ はエタールであることを用いている。
補題 \ref{algebra-lemma-etale} を参照せよ。）
\end{proof}

\begin{lemma}
\label{algebra-lemma-characterize-etale-over-polynomial-ring}
$k$ を体とする。
$$
\varphi : k[x_1, \ldots, x_n] \to A, \quad x_i \longmapsto a_i
$$
を有限型環写像とする。このとき $\varphi$ がエタールであるための必要十分条件は
次の二条件である。(a) $A$ の極大イデアルにおける局所環の次元が $n$ であり、
(b) 元 $\text{d}(a_1), \ldots, \text{d}(a_n)$ が $A$-加群として
$\Omega_{A/k}$ を生成する。
\end{lemma}

\begin{proof}
(a) と (b) を仮定する。条件 (b) は
$\Omega_{A/k[x_1, \ldots, x_n]} = 0$ を含意するので、$\varphi$ は
不分岐である。従って補題
\ref{algebra-lemma-etale-flat-unramified-finite-presentation} により、
$\varphi$ が平坦であることを証明すれば十分である。
$\mathfrak m \subset A$ を極大イデアルとする。$X = \Spec(A)$ とおき、
$\mathfrak m$ に対応する閉点を $x \in X$ と書く。このとき補題
\ref{algebra-lemma-dimension-closed-point-finite-type-field} により
$\dim(A_\mathfrak m)$ は $\dim_x X$ である。従って補題
\ref{algebra-lemma-characterize-smooth-over-field} により、(a) と (b) が成り立てば、
各極大イデアル $\mathfrak m$ に対して $A_\mathfrak m$ は正則局所環である。
すると補題 \ref{algebra-lemma-CM-over-regular-flat} により
$k[x_1, \ldots, x_n]_{\varphi^{-1}(\mathfrak m)} \to A_\mathfrak m$
は平坦である（正則局所環が CM であることも用いる。補題
\ref{algebra-lemma-regular-ring-CM} を参照せよ）。ゆえに補題
\ref{algebra-lemma-flat-localization} により $\varphi$ は平坦である。

\medskip\noindent
$\varphi$ がエタールであると仮定する。このとき
$\Omega_{A/k[x_1, \ldots, x_n]} = 0$ なので (b) が成り立つ。一方、
エタール環写像は平坦（補題 \ref{algebra-lemma-etale}）かつ準有限
（補題 \ref{algebra-lemma-etale-quasi-finite}）である。従って $A$ の各極大イデアル
$\mathfrak m$ に対し、補題 \ref{algebra-lemma-dimension-base-fibre-equals-total} を
$k[x_1, \ldots, x_n]_{\varphi^{-1}(\mathfrak m)} \to A_\mathfrak m$
に適用して $\dim(A_\mathfrak m) = n$ を得る。従って (a) が成り立つ。
% P01-E0949: the frozen source has the typo “we my apply” instead of “we may apply”.
\end{proof}

\section{不分岐環写像の局所構造}
\label{algebra-section-local-structure-unramified}

\noindent
不分岐射は局所的には（適切な意味で）閉埋め込みとエタール射との合成である。
本節では、この事実の代数的基礎を論じる。

\begin{proposition}
\label{algebra-proposition-unramified-locally-standard}
$R \to S$ を環写像とし、$\mathfrak q \subset S$ を素イデアルとする。
$R \to S$ が $\mathfrak q$ において不分岐ならば、次が存在する。
\begin{enumerate}
\item $g \in S$, $g \not \in \mathfrak q$ なる元。
\item 標準エタール環写像 $R \to S'$。
\item 全射な $R$-代数写像 $S' \to S_g$。
\end{enumerate}
\end{proposition}

\begin{proof}
この証明は命題 \ref{algebra-proposition-etale-locally-standard} の証明と「同じ」である。
やや迂回的な証明であり、短くする方法があるかもしれない。

\medskip\noindent
ステップ 1。定義 \ref{algebra-definition-unramified} により、
$g \in S$, $g \not \in \mathfrak q$ であって $R \to S_g$ が不分岐となる
ものが存在する。従って $S$ は $R$ 上不分岐であると仮定してよい。

\medskip\noindent
ステップ 2。補題 \ref{algebra-lemma-unramified} により、$R_0$ が
$\mathbf{Z}$ 上有限型となる不分岐環写像 $R_0 \to S_0$ と、
$S$ が $R \otimes_{R_0} S_0$ の商となる環写像 $R_0 \to R$ が存在する。
$\mathfrak q$ に対応する $S_0$ の素イデアルを $\mathfrak q_0$ と書く。
$(R_0 \to S_0, \mathfrak q_0)$ に対して結果を示せば、基底変換により
$(R \to S, \mathfrak q)$ に対する結果が従う。従って $R$ はネーター環であると
仮定してよい。

\medskip\noindent
ステップ 3。補題 \ref{algebra-lemma-unramified-quasi-finite} により、
$R \to S$ は準有限である。補題
\ref{algebra-lemma-quasi-finite-open-integral-closure} により、有限環写像
$R \to S'$、$R$-代数写像 $S' \to S$、および
$g' \in S'$, $g' \not \in \mathfrak q$ なる元であって、$S' \to S$ が同型
$S'_{g'} \cong S_{g'}$ を誘導するものが存在する。
% P01-E0950: the frozen source compares g' in S' directly with q in S and repeats “such that”; it should refer to the image of g' (or the contracted prime) once.
（$S'$ は $R$ 上不分岐とは限らないことに注意する。）従って、
(a) $R$ はネーター環、(b) $R \to S$ は有限、(c) $R \to S$ は
$\mathfrak q$ において不分岐であると仮定してよい
（ただし、もはやすべての素イデアルにおいて不分岐とは限らない）。

\medskip\noindent
ステップ 4。$\mathfrak p \subset R$ を $\mathfrak q$ に対応する素イデアルとする。
ファイバー環 $S \otimes_R \kappa(\mathfrak p)$ を考える。これは
$\kappa(\mathfrak p)$ 上有限な代数であるからアルティン環であり
（補題 \ref{algebra-lemma-finite-dimensional-algebra} を参照せよ）、従って局所環の有限直積
$$
S \otimes_R \kappa(\mathfrak p) = \prod\nolimits_{i = 1}^n A_i
$$
である。命題 \ref{algebra-proposition-dimension-zero-ring} を参照せよ。一つの因子、例えば
$A_1$ は局所環 $S_{\mathfrak q}/\mathfrak pS_{\mathfrak q}$ であり、
補題 \ref{algebra-lemma-unramified-at-prime} により $\kappa(\mathfrak q)$ と同型である。
他の因子は、$\mathfrak p$ の上にある $S$ の他の素イデアル
$\mathfrak q_2, \ldots, \mathfrak q_n$ に対応する。

\medskip\noindent
ステップ 5。有限分離体拡大
$\kappa(\mathfrak q)/\kappa(\mathfrak p)$ を生成する非零元
$\alpha \in \kappa(\mathfrak q)$ を選べる（従って、体拡大が自明な場合でも
$\alpha = 0$ は許さない）。任意の $\lambda \in \kappa(\mathfrak p)^*$ に対し、
元 $\lambda \alpha$ も $\kappa(\mathfrak p)$ 上 $\kappa(\mathfrak q)$ を
生成することに注意する。元
$$
\overline{t} =
(\alpha, 0, \ldots, 0) \in
\prod\nolimits_{i = 1}^n A_i =
S \otimes_R \kappa(\mathfrak p).
$$
を考える。必要なら上のように $\alpha$ を $\lambda \alpha$ で置き換えることにより、
$\overline{t}$ はある $t \in S$ の像であると仮定してよい。
$x$ を $t$ に写す $R$-代数写像 $R[x] \to S$ の核を $I \subset R[x]$ とする。
$S' = R[x]/I$ とおけば $S' \subset S$ である。図式
$$
\xymatrix{
R[x] \ar[r] & S' \ar[r] & S \\
R \ar[u] \ar[ru] \ar[rru] & &
}
$$
を得る。構成により、$S$ の素イデアル $\mathfrak q_j$, $j \geq 2$ はすべて
$R[x]$ の素イデアル $(\mathfrak p, x)$ の上にある。一方、
$\alpha \not = 0$ なので、$\mathfrak q$ は $R[x]$ の別の素イデアルの上にある。

\medskip\noindent
ステップ 6。$\mathfrak q$ に対応する $S'$ の素イデアルを
$\mathfrak q' \subset S'$ と書く。上の議論により、$\mathfrak q$ は
$\mathfrak q'$ の上にある $S$ の唯一の素イデアルである。従って補題
\ref{algebra-lemma-unique-prime-over-localize-below} により
$S_{\mathfrak q} = S_{\mathfrak q'}$ である（$R \to S$ が有限なので
$S' \to S$ も有限であり、補題 \ref{algebra-lemma-integral-going-up} により
$S' \to S$ に対して上昇性が成り立つ）。従って
$S'_{\mathfrak q'} \to S_{\mathfrak q}$ は、有限単射環写像 $S' \to S$ の
局所化として有限かつ単射である。局所環の写像
$$
R_{\mathfrak p} \to S'_{\mathfrak q'} \to S_{\mathfrak q}
$$
を考える。第二の写像は有限かつ単射である。補題
\ref{algebra-lemma-unramified-at-prime} により
$S_{\mathfrak q}/\mathfrak pS_{\mathfrak q} = \kappa(\mathfrak q)$ である。
従ってなおさら
$S_{\mathfrak q}/\mathfrak q'S_{\mathfrak q} = \kappa(\mathfrak q)$ である。
また
$$
\kappa(\mathfrak p) \subset \kappa(\mathfrak q') \subset \kappa(\mathfrak q)
$$
であり、$\alpha$ は $\kappa(\mathfrak q')$ の $\kappa(\mathfrak q)$ における像に
含まれるので、$\kappa(\mathfrak q') = \kappa(\mathfrak q)$ と結論できる。
そこで中山の補題 \ref{algebra-lemma-NAK} を $S'_{\mathfrak q'}$-加群写像
$S'_{\mathfrak q'} \to S_{\mathfrak q}$ に適用すると、写像
$S'_{\mathfrak q'} \to S_{\mathfrak q}$ は全射である。
言い換えれば、$S'_{\mathfrak q'} \cong S_{\mathfrak q}$ である。

\medskip\noindent
ステップ 7。補題 \ref{algebra-lemma-isomorphic-local-rings} により、
$g \in S$, $g \not \in \mathfrak q$ および
$g' \in S'$, $g' \not \in \mathfrak q'$ であって
$S'_{g'} \cong S_g$ となるものが存在する。$R$ はネーター環なので、
$S'$ は有限 $R$-加群 $S$ の $R$-部分加群として $R$ 上有限である。
従って $S$ を $S'$ で置き換えた後、(a) $R$ はネーター環、
(b) $S$ は $R$ 上有限、
% P01-E0951: the frozen source omits “is” in “(b) S finite over R”.
(c) $S$ は $R$ 上 $\mathfrak q$ において不分岐、かつ
(d) $S = R[x]/I$ と仮定してよい。

\medskip\noindent
ステップ 8。環
$S \otimes_R \kappa(\mathfrak p) = \kappa(\mathfrak p)[x]/\overline{I}$
を考える。ここで $\overline{I} = I \cdot \kappa(\mathfrak p)[x]$ は
$I$ が $\kappa(\mathfrak p)[x]$ の中で生成するイデアルである。
$\kappa(\mathfrak p)[x]$ は主イデアル整域なので、あるモニックな
$\overline{h} \in \kappa(\mathfrak p)$ に対して
$\overline{I} = (\overline{h})$ である。
% P01-E0952: the frozen source places the monic polynomial h-bar in kappa(p) instead of kappa(p)[x].
ある $\lambda \in \kappa(\mathfrak p)$ に対し $\overline{h}$ を
$\lambda \cdot \overline{h}$ で置き換えた後、$\overline{h}$ はある
$h \in R[x]$ の像であると仮定してよい。
% P01-E0953: the scaling factor must be nonzero to preserve the principal generator, but the frozen source omits the star on kappa(p).
（問題は、$h$ をモニックに選べるかどうか分からないことである。）
また、ステップ 4 と同様に
$S \otimes_R \kappa(\mathfrak p) = A_1 \times \ldots \times A_n$ であり、
$A_1 = \kappa(\mathfrak q)$ は $\kappa(\mathfrak p)$ の有限分離拡大、
$A_2, \ldots, A_n$ は局所環であることが分かっている。これは
$$
\overline{h} = \overline{h}_1 \overline{h}_2^{e_2} \ldots \overline{h}_n^{e_n}
$$
を含意する。
% P01-E0954: after the frozen source replaces the monic generator by an arbitrary nonzero scalar multiple, monicity is no longer assured; this factorization into monic factors needs the resulting leading unit as a coefficient.
ここで $\overline{h}_i \in \kappa(\mathfrak p)[x]$ は互いに素な不可約モニック
多項式であり、$e_2, \ldots, e_n \geq 1$ である。番号は
$A_i = \kappa(\mathfrak p)[x]/(\overline{h}_i^{e_i})$ が
$\kappa(\mathfrak p)[x]$-代数として成り立つように選んでいる。
% P01-E0955: the frozen formula uses e_i also for i = 1, although only e_2,...,e_n are defined; it should set e_1 = 1 or state the i = 1 case separately.
$\overline{h}_1$ は $\alpha \in \kappa(\mathfrak q)$ の最小多項式であり、
従って可分多項式である（その導関数は自身と互いに素である）ことに注意する。

\medskip\noindent
ステップ 9。$m \in I$ をモニックな元とする。そのような元は、環拡大
$R \to R[x]/I$ が有限であり、従って整であることから存在する。
$\kappa(\mathfrak p)[x]$ における像を $\overline{m}$ と書く。ある
$d_1 \geq 1$, $d_j \geq e_j$, $j = 2, \ldots, n$ および、すべての
$\overline{h}_i$ と互いに素な
$\overline{k} \in \kappa(\mathfrak p)[x]$ によって
$$
\overline{m} = \overline{k}
\overline{h}_1^{d_1} \overline{h}_2^{d_2} \ldots \overline{h}_n^{d_n}
$$
と分解できる。$l \deg(m) > \deg(h)$ かつ $l \geq 2$ となるようにして
$f = m^l + h$ とおく。このとき $f$ は $R$ 上のモニック多項式である。
また、$f$ の $\kappa(\mathfrak p)[x]$ における像 $\overline{f}$ は
$$
\overline{f} =
\overline{h}_1 \overline{h}_2^{e_2} \ldots \overline{h}_n^{e_n}
+
\overline{k}^l \overline{h}_1^{ld_1} \overline{h}_2^{ld_2}
\ldots \overline{h}_n^{ld_n}
=
\overline{h}_1(\overline{h}_2^{e_2} \ldots \overline{h}_n^{e_n}
+
\overline{k}^l
\overline{h}_1^{ld_1 - 1} \overline{h}_2^{ld_2} \ldots \overline{h}_n^{ld_n})
= \overline{h}_1 \overline{w}
$$
と分解する。ここで $\overline{w}$ は $\overline{h}_1$ と互いに素な多項式である。
$g = f'$（$x$ に関する導関数）とおく。

\medskip\noindent
ステップ 10。環写像 $R[x] \to S = R[x]/I$ は次の性質をもつ。
(1) $f$ を零に写す。
(2) $g$ を $S \setminus \mathfrak q$ の元に写す。
第一の主張は $f$ が $I$ の元なので明らかである。
% P01-E0956: the frozen source only says h lies in R[x], but Step 10 needs h in I to conclude f = m^l + h lies in I.
第二の主張については、$g$ が
$\kappa(\mathfrak q) = \kappa(\mathfrak p)[x]/(\overline{h}_1)$ で
零に写らないことを示せばよい。$\kappa(\mathfrak p)[x]$ における $g$ の像は
$\overline{f}$ の導関数である。従って (2) は、
$$
\overline{g} =
\frac{\text{d}\overline{f}}{\text{d}x} =
\overline{w}\frac{\text{d}\overline{h}_1}{\text{d}x} +
\overline{h}_1\frac{\text{d}\overline{w}}{\text{d}x},
$$
$\overline{w}$ が $\overline{h}_1$ と互いに素であり、
$\overline{h}_1$ が可分であることから明らかである。

\medskip\noindent
ステップ 11。以上により、$\varphi : R[x]/(f) \to S$ は全射な環写像であり、
$R[x]_g/(f)$ は $R$ 上エタールであり（標準エタールだからである。補題
\ref{algebra-lemma-standard-etale} を参照せよ）、かつ
$\varphi(g) \not \in \mathfrak q$ である。従って写像
$(R[x]/(f))_g \to S_{\varphi(g)}$ が求める全射である。
\end{proof}



\begin{lemma}
\label{algebra-lemma-etale-makes-unramified-closed-at-prime}
$R \to S$ を環写像とする。$\mathfrak q$ を、$\mathfrak p \subset R$ の上にある
$S$ の素イデアルとする。$R \to S$ が有限型で、$\mathfrak q$ において
不分岐であると仮定する。このとき次が存在する。
\begin{enumerate}
\item エタール環写像 $R \to R'$。
\item $\mathfrak p$ の上にある素イデアル $\mathfrak p' \subset R'$。
\item 直積分解
$$
R' \otimes_R S = A \times B
$$
\end{enumerate}
これらは次の性質をもつ。
\begin{enumerate}
\item $R' \to A$ は全射である。
\item $\mathfrak p'A$ は $\mathfrak p'$ と $\mathfrak q$ の上にある
$A$ の素イデアルである。
\end{enumerate}
\end{lemma}

\begin{proof}
$(R \to S, \mathfrak p, \mathfrak q)$ を、エタール環写像
$R \to R'$ による任意の基底変換
$(R' \to R'\otimes_R S, \mathfrak p', \mathfrak q')$ で置き換えてよい。
ここで $\mathfrak p'$ は $\mathfrak p$ の上にある素イデアルであり、
$\mathfrak q'$ は $\mathfrak q$ と $\mathfrak p'$ の両方の上にあるように
選んでいる。$R \to R'$ と $\mathfrak p'$ が与えられたとき、適切な
$\mathfrak q'$ は常に見つかることにも注意する。

\medskip\noindent
$R \to S$ が有限型であるという仮定から、補題
\ref{algebra-lemma-etale-makes-quasi-finite-finite-variant} を適用できる。従って
$S = A_1 \times \ldots \times A_n \times B$ であり、各 $R \to A_i$ は有限で、
$\mathfrak p$ の上にある素イデアル $\mathfrak r_i$ をちょうど一つもち、
$\kappa(\mathfrak p) \subset \kappa(\mathfrak r_i)$ は純非分離であり、さらに
$R \to B$ は $\mathfrak p$ の上にあるどの素イデアルにおいても準有限でないと
仮定してよい。不分岐射は準有限なので（補題
\ref{algebra-lemma-unramified-quasi-finite} を参照せよ）、明らかにある $i$ に対して
$\mathfrak q = \mathfrak r_i$ である。$\mathfrak q = \mathfrak r_1$ としよう。
補題 \ref{algebra-lemma-unramified-at-prime} により
$\kappa(\mathfrak r_1)/\kappa(\mathfrak p)$ は分離拡大であり、従って自明な
体拡大である。また、$\mathfrak p(A_1)_{\mathfrak r_1}$ は極大イデアルである。
さらに、補題 \ref{algebra-lemma-unique-prime-over-localize-below} により
（有限環写像は補題 \ref{algebra-lemma-integral-going-up} により上昇性を満たすので、
これは $R \to A_1$ に適用できる）、
$(A_1)_{\mathfrak r_1} = (A_1)_{\mathfrak p}$ である。
中山の補題 \ref{algebra-lemma-NAK} から、局所環の写像
$R_{\mathfrak p} \to (A_1)_{\mathfrak p} = (A_1)_{\mathfrak r_1}$ は全射である。
$A_1$ は $R$ 上有限なので、$f \in R$, $f \not \in \mathfrak p$ であって
$R_f \to (A_1)_f$ が全射となるものが存在する。
$R$ を $R_f$ で置き換えれば結論を得る。
\end{proof}

\begin{lemma}
\label{algebra-lemma-etale-makes-unramified-closed}
\begin{slogan}
不分岐環写像では、エタール近傍へ移ることによってファイバー内の点を分離できる。
\end{slogan}
$R \to S$ を環写像とし、$\mathfrak p$ を $R$ の素イデアルとする。
$R \to S$ が不分岐ならば、次が存在する。
\begin{enumerate}
\item エタール環写像 $R \to R'$。
\item $\mathfrak p$ の上にある素イデアル $\mathfrak p' \subset R'$。
\item 直積分解
$$
R' \otimes_R S = A_1 \times \ldots \times A_n \times B
$$
\end{enumerate}
これらは次の性質をもつ。
\begin{enumerate}
\item $R' \to A_i$ は全射である。
\item $\mathfrak p'A_i$ は $\mathfrak p'$ の上にある $A_i$ の素イデアルである。
\item $B$ には $\mathfrak p'$ の上にある素イデアルが存在しない。
\end{enumerate}
\end{lemma}

\begin{proof}
補題 \ref{algebra-lemma-etale-makes-quasi-finite-finite-variant} を適用できる。
従って、エタール基底変換の後、
$S = A_1 \times \ldots \times A_n \times B$ であり、各 $R \to A_i$ は有限で、
$\mathfrak p$ の上にある素イデアル $\mathfrak r_i$ をちょうど一つもち、
$\kappa(\mathfrak p) \subset \kappa(\mathfrak r_i)$ は純非分離であり、さらに
$R \to B$ は $\mathfrak p$ の上にあるどの素イデアルにおいても準有限でないと
仮定してよい。$R \to S$ は準有限なので（補題
\ref{algebra-lemma-unramified-quasi-finite} を参照せよ）、$B$ には $\mathfrak p$ の上にある
素イデアルが存在しない。補題 \ref{algebra-lemma-unramified-at-prime} により
$\kappa(\mathfrak r_i)/\kappa(\mathfrak p)$ は分離拡大であり、従って自明な
体拡大である。また、$\mathfrak p(A_i)_{\mathfrak r_i}$ は極大イデアルである。
さらに、補題 \ref{algebra-lemma-unique-prime-over-localize-below} により
（有限環写像は補題 \ref{algebra-lemma-integral-going-up} により上昇性を満たすので、
これは $R \to A_i$ に適用できる）、
$(A_i)_{\mathfrak r_i} = (A_i)_{\mathfrak p}$ である。
中山の補題 \ref{algebra-lemma-NAK} から、局所環の写像
$R_{\mathfrak p} \to (A_i)_{\mathfrak p} = (A_i)_{\mathfrak r_i}$ は全射である。
$A_i$ は $R$ 上有限なので、$f \in R$, $f \not \in \mathfrak p$ であって
$R_f \to (A_i)_f$ が全射となるものが存在する。
$R$ を $R_f$ で置き換えれば結論を得る。
\end{proof}

\section{ヘンゼル局所環}
\label{algebra-section-henselian}

\noindent
この節では、ヘンゼル局所環の概念を少し論じる。
$(R, \mathfrak m, \kappa)$ を局所環とする。
$a \in R$ に対して、$\kappa$ における $a$ の像を $\overline{a}$ と書く。
多項式 $f \in R[T]$ に対しては、$\kappa[T]$ における $f$ の像を
$\overline{f}$ と書くことが多い。
多項式 $f \in R[T]$ が与えられたとき、$T$ に関する $f$ の導関数を
$f'$ と書く。$\overline{f}' = \overline{f'}$ であることに注意せよ。

\begin{definition}
\label{algebra-definition-henselian}
$(R, \mathfrak m, \kappa)$ を局所環とする。
\begin{enumerate}
\item 任意のモニック多項式 $f \in R[T]$ と、$\overline{f}$ の任意の根
$a_0 \in \kappa$ で $\overline{f'}(a_0) \not = 0$ を満たすものに対し、
$f(a) = 0$ かつ $a_0 = \overline{a}$ を満たす $a \in R$ が存在するとき、
$R$ は {\it ヘンゼル的} であるという。
\item $R$ がヘンゼル的であり、その剰余体が分離代数閉であるとき、
$R$ は {\it 強ヘンゼル的} であるという。
\end{enumerate}
\end{definition}

\noindent
条件 $\overline{f'}(a_0) \not = 0$ は、$a_0$ が多項式 $\overline{f}$ の
単純根であるという条件と同値であることに注意せよ。実際、この条件は、
持ち上げ $a \in R$ が存在するならば一意であることも意味する。

\begin{lemma}
\label{algebra-lemma-uniqueness}
$(R, \mathfrak m, \kappa)$ を局所環とする。
$f \in R[T]$ とする。$a, b \in R$ が $f(a) = f(b) = 0$、
$a = b \bmod \mathfrak m$、および $f'(a) \not \in \mathfrak m$ を
満たすとする。このとき $a = b$ である。
\end{lemma}

\begin{proof}
$R[x, y]$ において
$f(x + y) - f(x) = f'(x)y + g(x, y) y^2$ と書く
（これは $f(x + y)$ を展開すれば分かるので、詳細は省略する）。
すると、ある $c \in R$ に対して
$0 = f(b) - f(a) = f(a + (b - a)) - f(a) =
f'(a)(b - a) + c (b - a)^2$ である。仮定により
$f'(a)$ は $R$ の単元である。従って
$(b - a)(1 + f'(a)^{-1}c(b - a)) = 0$ である。
仮定より $b - a \in \mathfrak m$ なので、
$1 + f'(a)^{-1}c(b - a)$ は $R$ の単元である。
ゆえに $R$ において $b - a = 0$ である。
\end{proof}

\noindent
以下がヘンゼル局所環の特徴づけである。

\begin{lemma}
\label{algebra-lemma-characterize-henselian}
\begin{slogan}
ヘンゼル局所環の特徴づけ
\end{slogan}
$(R, \mathfrak m, \kappa)$ を局所環とする。
次の条件は同値である。
\begin{enumerate}
\item $R$ はヘンゼル的である。
\item 任意の $f \in R[T]$ と、$\overline{f}$ の任意の根
$a_0 \in \kappa$ で $\overline{f'}(a_0) \not = 0$ を満たすものに対し、
$f(a) = 0$ かつ $a_0 = \overline{a}$ を満たす $a \in R$ が存在する。
\item 任意のモニック多項式 $f \in R[T]$ と、
$\gcd(g_0, h_0) = 1$ を満たす任意の分解
$\overline{f} = g_0 h_0$ に対し、$g_0 = \overline{g}$ および
$h_0 = \overline{h}$ を満たす分解 $f = gh$ が $R[T]$ に存在する。
\item 任意のモニック多項式 $f \in R[T]$ と、
$\gcd(g_0, h_0) = 1$ を満たす任意の分解
$\overline{f} = g_0 h_0$ に対し、$g_0 = \overline{g}$、
$h_0 = \overline{h}$、さらに $\deg_T(g) = \deg_T(g_0)$ を満たす
分解 $f = gh$ が $R[T]$ に存在する。
\item 任意の $f \in R[T]$ と、$\gcd(g_0, h_0) = 1$ を満たす
任意の分解 $\overline{f} = g_0 h_0$ に対し、
$g_0 = \overline{g}$ および $h_0 = \overline{h}$ を満たす
分解 $f = gh$ が $R[T]$ に存在する。
\item 任意の $f \in R[T]$ と、$\gcd(g_0, h_0) = 1$ を満たす
任意の分解 $\overline{f} = g_0 h_0$ に対し、
$g_0 = \overline{g}$、$h_0 = \overline{h}$、さらに
$\deg_T(g) = \deg_T(g_0)$ を満たす分解 $f = gh$ が $R[T]$ に存在する。
\item 任意のエタール環写像 $R \to S$ と、$\mathfrak m$ の上にあり
$\kappa = \kappa(\mathfrak q)$ を満たす $S$ の素イデアル
$\mathfrak q$ に対して、$R \to S$ のレトラクション
$\tau : S \to R$ が存在する。
\item 任意のエタール環写像 $R \to S$ と、$\mathfrak m$ の上にあり
$\kappa = \kappa(\mathfrak q)$ を満たす $S$ の素イデアル
$\mathfrak q$ に対して、$\mathfrak q = \tau^{-1}(\mathfrak m)$ を満たす
$R \to S$ の一意なレトラクション $\tau : S \to R$ が存在する。
\item 任意の有限 $R$-代数は局所環の積である。
\item 任意の有限 $R$-代数は局所環の有限積である。
\item 任意の有限型 $R$-代数 $S$ は、$R \to A$ が有限であり、
$R \to B$ が $\mathfrak m$ の上にあるどの素イデアルにおいても
準有限でないような $A \times B$ として書ける。
\item 任意の有限型 $R$-代数 $S$ は、$R \to A$ が有限であり、
$\Spec(B \otimes_R \kappa)$ の各既約成分の次元が $\geq 1$ となるような
$A \times B$ として書ける。
\item 任意の準有限 $R$-代数 $S$ は、$R \to A$ が有限で
$B \otimes_R \kappa = 0$ となるような $S = A \times B$ として書ける。
\end{enumerate}
\end{lemma}

\begin{proof}
まず、容易な含意を列挙する。
\begin{enumerate}
\item 2$\Rightarrow$1。(2) ではすべての多項式を考え、
(1) ではモニックなものだけを考えるからである。
\item 5$\Rightarrow$3。(5) ではすべての多項式を考え、
(3) ではモニックなものだけを考えるからである。
\item 6$\Rightarrow$4。(6) ではすべての多項式を考え、
(4) ではモニックなものだけを考えるからである。
\item 4$\Rightarrow$3 は明らかである。
\item 6$\Rightarrow$5 は明らかである。
\item 8$\Rightarrow$7 は明らかである。
\item 10$\Rightarrow$9 は明らかである。
\item 11$\Leftrightarrow$12 は、素イデアルにおいて準有限であることの定義による。
\item 11$\Rightarrow$13 は準有限であることの定義による。
\end{enumerate}

\noindent
1$\Rightarrow$8 の証明。(1) を仮定する。
$R \to S$ をエタールとし、$\mathfrak q \subset S$ を
$\kappa(\mathfrak q) \cong \kappa$ を満たす素イデアルとする。
命題 \ref{algebra-proposition-etale-locally-standard} により、
$g \in S$, $g \not \in \mathfrak q$ で $R \to S_g$ が標準エタールと
なるものが存在する。$S$ を $S_g$ で置き換えることにより、
$S = R[t]_g/(f)$ が標準エタールであると仮定してよい（詳細は省略する）。
素イデアル $\mathfrak q$ の剰余体は $\kappa$ なので、これは
$\overline{f}$ の根 $a_0$ で $\overline{g}$ の根でないものに対応する。
標準エタール代数の定義により、これは
$\overline{f'}(a_0) \not = 0$ であることも意味する。
また標準エタール代数の定義により $f$ はモニックなので、
$R$ がヘンゼル的であることから、$a_0 = \overline{a}$ かつ
$f(a) = 0$ を満たす $a \in R$ が存在する。
これより $g(a)$ は $R$ の単元であり、規則 $t \mapsto a$ によって
求める写像 $\tau : S = R[t]_g/(f) \to R$ を得る。
構成により $\tau^{-1}(\mathfrak m) = \mathfrak q$ である。
補題 \ref{algebra-lemma-uniqueness} により、写像 $\tau$ は一意である。
従って (8) が成り立つ。

\medskip\noindent
7$\Rightarrow$8 の証明。（これは実際には重要でなく、省略すべきである。）
(7) を仮定し、$R \to S$ がエタールであるとする。
$\mathfrak q_1, \ldots, \mathfrak q_r$ を、$\mathfrak m$ の上にある
$S$ の他の素イデアルとする。このとき、$g \in S$ で
$g \not \in \mathfrak q$ かつ $i = 1, \ldots, r$ に対して
$g \in \mathfrak q_i$ となるものが存在する。実際、
$\bigcap_{i=1}^{r} \mathfrak{q}_{i} \not\subset \mathfrak{q}$ と論じられる。
そうでなければ、ある $i$ に対して
$\mathfrak{q}_{i} \subset \mathfrak{q}$ となるが、エタール射のファイバーは
離散的なので、これは起こり得ない（例えば補題
\ref{algebra-lemma-etale-over-field} を用いよ）。
エタール環写像 $R \to S_g$ と素イデアル $\mathfrak qS_g$ に (7) を適用する。
すると、合成 $\tau : S \to S_g \to R$ が
$\tau^{-1}(\mathfrak m) = \mathfrak q$ を満たすようなレトラクション
$\tau_g : S_g \to R$ を得る。詳細は省略する。

\medskip\noindent
8$\Rightarrow$11 の証明。(8) を仮定し、$R \to S$ を有限型環写像とする。
補題 \ref{algebra-lemma-etale-makes-quasi-finite-finite} を適用する。
すると、エタール環写像 $R \to R'$ と、$\mathfrak m$ の上にある
素イデアル $\mathfrak m' \subset R'$ で $\kappa = \kappa(\mathfrak m')$ を
満たすものが存在し、$R' \otimes_R S = A' \times B'$、
$A'$ は $R'$ 上有限、かつ $B'$ は $\mathfrak m'$ の上にあるどの素イデアルでも
$R'$ 上準有限でない。 (8) を適用して、
$\mathfrak m = \tau^{-1}(\mathfrak m')$ を満たすレトラクション
$\tau : R' \to R$ を得る。
% P01-E0957: the frozen inverse-image equality is ill-typed and reversed; for tau : R' -> R it should be m' = tau^{-1}(m).
ここで
$$
S = (S \otimes_R R') \otimes_{R', \tau} R
= (A' \times B') \otimes_{R', \tau} R
= (A' \otimes_{R', \tau} R)  \times  (B' \otimes_{R', \tau} R)
$$
を用いれば、(11) のような分解を得る。

\medskip\noindent
8$\Rightarrow$10 の証明。(8) を仮定し、$R \to S$ を有限環写像とする。
補題 \ref{algebra-lemma-etale-makes-quasi-finite-finite} を適用する。
すると、エタール環写像 $R \to R'$ と、$\mathfrak m$ の上にある
素イデアル $\mathfrak m' \subset R'$ で $\kappa = \kappa(\mathfrak m')$ を
満たすものが存在し、
$R' \otimes_R S = A'_1 \times \ldots \times A'_n \times B'$、各 $A'_i$ は
$R'$ 上有限で $\mathfrak m'$ の上に素イデアルをちょうど一つもち、
$B'$ は $\mathfrak m'$ の上にあるどの素イデアルでも $R'$ 上準有限でない。
(8) を適用して、$\mathfrak m' = \tau^{-1}(\mathfrak m)$ を満たす
レトラクション $\tau : R' \to R$ を得る。すると
\begin{align*}
S & = (S \otimes_R R') \otimes_{R', \tau} R \\
& = (A'_1 \times \ldots \times A'_n \times B') \otimes_{R', \tau} R \\
& = (A'_1 \otimes_{R', \tau} R)  \times
\ldots \times (A'_1 \otimes_{R', \tau} R) \times
(B' \otimes_{R', \tau} R) \\
& = A_1 \times \ldots \times A_n \times B
\end{align*}
% P01-E0958: the frozen third line repeats A'_1 in the final A-factor; it should end with A'_n.
因子 $B$ は $R$ 上有限だが、$R \to B$ は $\mathfrak m$ の上にある
どの素イデアルにおいても準有限でない。従って $B = 0$ である。
因子 $A_i$ は有限 $R$-代数で、$\mathfrak m$ の上にある素イデアルを
ちょうど一つもつので、局所環である。これにより、$S$ が局所環の有限積で
あることが示された。

\medskip\noindent
9$\Rightarrow$10 の証明。$S$ が局所環 $R$ 上有限ならば、極大イデアルは
高々有限個しかないので、この含意が成り立つ。実際、$R \to S$ に対する
上昇性により $S$ のすべての極大イデアルは $\mathfrak m$ の上にあり、
$S/\mathfrak mS$ はアルティン環なので素イデアルを有限個しかもたない。

\medskip\noindent
10$\Rightarrow$1 の証明。(10) を仮定する。$f \in R[T]$ をモニック多項式とし、
$a_0 \in \kappa$ を $\overline{f}$ の単純根とする。
このとき $S = R[T]/(f)$ は有限 $R$-代数である。(10) を適用すると、
$S = A_1 \times \ldots \times A_r$ は局所 $R$-代数の有限積となる。
特に、$S/\mathfrak mS = \prod A_i/\mathfrak mA_i$ は
$\kappa[T]/(\overline{f})$ の局所環の積への分解である。
従って、因子の一つ、例えば $A_1/\mathfrak mA_1$ は、商写像
$\kappa[T]/(\overline{f}) \to \kappa[T]/(T - a_0)$ である。
$A_1$ は有限自由 $R$-加群 $S$ の直和因子なので、それ自身有限自由
$R$-加群である。$A_1/\mathfrak mA_1$ は次元 1 の $\kappa$-ベクトル空間なので、
$R$-加群として $A_1 \cong R$ である。これは明らかに
$R \to A_1$ が同型であることを意味する。写像
$R[T] \to S \to A_1 \to R$ による $T$ の像を $a \in R$ とする。
すると、望む通り $f(a) = 0$ かつ $\overline{a} = a_0$ である。

\medskip\noindent
13$\Rightarrow$1 の証明。(13) を仮定する。$f \in R[T]$ をモニック多項式とし、
$a_0 \in \kappa$ を $\overline{f}$ の単純根とする。
このとき $S_1 = R[T]/(f)$ は有限 $R$-代数である。
$\overline{g} = \overline{f}/(T - a_0)$ を満たす任意の元
$g \in R[T]$ をとる。このとき $S = (S_1)_g$ は準有限 $R$-代数であり、
$S \otimes_R \kappa \cong \kappa[T]_{\overline{g}}/(\overline{f})
\cong \kappa[T]/(T - a_0) \cong \kappa$ を満たす。
(13) を $S$ に適用すると、$A$ が $R$ 上有限かつ
$B \otimes_R \kappa = 0$ となる $S = A \times B$ を得る。
特に $\kappa \cong S/\mathfrak mS = A/\mathfrak mA$ である。
$A$ は平坦 $R$-代数 $S$ の直和因子なので有限平坦、従って $R$ 上自由である。
$A/\mathfrak mA$ は次元 1 の $\kappa$-ベクトル空間なので、
$R$-加群として $A \cong R$ である。これは明らかに $R \to A$ が
同型であることを意味する。写像 $R[T] \to S \to A \to R$ による
$T$ の像を $a \in R$ とする。すると、望む通り
$f(a) = 0$ かつ $\overline{a} = a_0$ である。

\medskip\noindent
8$\Rightarrow$2 の証明。(8) を仮定する。$f \in R[T]$ を任意の多項式とし、
$a_0 \in \kappa$ を単純根とする。
% P01-E0959: the frozen proof calls a_0 a simple root without specifying that it is a root of f-bar.
このとき代数 $S = R[T]_{f'}/(f)$ は $R$ 上エタールである。
$b \in R$ を $\overline{b} = a_0$ を満たす任意の元とし、
$\mathfrak q \subset S$ を $\mathfrak m$ と $T - b$ で生成される素イデアルとする。
(8) を $S$ と $\mathfrak q$ に適用して $\tau : S \to R$ を得る。
このとき像 $\tau(T) = a \in R$ が (2) の条件を満たす。

\medskip\noindent
ここまでで、(1), (2), (7), (8), (9), (10), (11), (12), (13) が
すべて同値であることが分かった。(3), (4), (5), (6) のうち最も弱い主張は
(3) で、最も強いものは (6) である。従って、なお (3) が (1) を含意し、
(1) が (6) を含意することを示せばよい。

\medskip\noindent
3$\Rightarrow$1 の証明。(3) を仮定する。$f \in R[T]$ をモニックとし、
$a_0 \in \kappa$ を $\overline{f}$ の単純根とする。すると
$h_0(a_0) \not = 0$ を満たす分解
$\overline{f} = (T - a_0)h_0$ が得られるので、
$\gcd(T - a_0, h_0) = 1$ である。(3) を適用して、
$\overline{g} = T - a_0$ および $\overline{h} = h_0$ を満たす分解
$f = gh$ を得る。有限自由 $R$-代数 $S = R[T]/(f)$ とおく。
$g$, $h$ の $S$ における像もそれぞれ $g$, $h$ と書く。
法 $\mathfrak m$ で等式が成り立つので、中山の補題
\ref{algebra-lemma-NAK} により $gS + hS = S$ である。
$S$ において $gh = f = 0$ なので、これより $gS \cap hS = 0$ も従う。
従って中国剰余定理により $S = S/(g) \times S/(h)$ を得る。
ゆえに $A = S/(g)$ は有限自由 $R$-加群の直和因子であり、従って有限自由である。
さらに $A/\mathfrak mA = \kappa[T]/(T - a_0)$ なので、$A$ の階数は $1$ である。
従って写像 $R \to A$ は同型である。写像
$R[T] \to S \to A \to R$ による $T$ の像と等しい $R$ の元を
$a \in R$ とおけば、
$f(a) = 0$ かつ $\overline{a} = a_0$ を満たす元が得られる。

\medskip\noindent
1$\Rightarrow$6 の証明。(1)、または同値な (1), (2), (7), (8), (9),
(10), (11), (12), (13) のすべてを仮定する。
$f \in R[T]$ を多項式とする。
$\gcd(g_0, h_0) = 1$ を満たす分解
$\overline{f} = g_0h_0$ があるとする。$g_0$ はモニックであると仮定してよく、
実際そうする。$S = R[T]/(f)$ を考える。この分解があることから、
$f$ の係数は $R$ の単位イデアルを生成する。
% P01-E0960: the frozen assertion fails for the valid coprime factorization g_0 = 1 and h_0 = 0; the trivial case must be handled separately.
これより $S$ の $R$ 上のファイバーは有限であり、従って $R$ 上準有限である。
また補題 \ref{algebra-lemma-grothendieck-general} により $S$ は $R$ 上平坦である。
(13) と (10) を組み合わせると、
$S = A_1 \times \ldots \times A_n \times B$ と書ける。ここで各 $A_i$ は
局所環で $R$ 上有限であり、$B \otimes_R \kappa = 0$ である。
因子 $A_1, \ldots, A_n$ を並べ替えることにより、$\kappa[T]$ の商として
$$
\kappa[T]/(g_0) =
A_1/\mathfrak m A_1 \times \ldots \times A_r/\mathfrak mA_r,
\ \kappa[T]/(h_0) =
A_{r + 1}/\mathfrak mA_{r + 1} \times \ldots \times A_n/\mathfrak mA_n
$$
と仮定してよい。有限平坦 $R$-代数
$A = A_1 \times \ldots \times A_r$ は $R$-加群として自由である。補題
\ref{algebra-lemma-finite-flat-local} を参照せよ。その階数は $\deg_T(g_0)$ である。
$g \in R[T]$ を $R$-線形作用素 $T : A \to A$ の特性多項式とする。
すると $g$ は次数 $\deg_T(g) = \deg_T(g_0)$ のモニック多項式であり、
さらに $\overline{g} = g_0$ である。ケイリー--ハミルトンの定理
（補題 \ref{algebra-lemma-charpoly}）により、$T_A$ を $A$ における $T$ の像とすると
$g(T_A) = 0$ である。従って、全射な写像 $R[T]/(g) \to A$ を得るが、
これは中山の補題 \ref{algebra-lemma-NAK} により同型である。
構成により写像 $R[T] \to A$ は $R[T]/(f)$ を経由するので、
ある $h$ に対して $f = gh$ と書ける。これで証明が完了する。
\end{proof}

\begin{lemma}
\label{algebra-lemma-finite-over-henselian}
$(R, \mathfrak m, \kappa)$ をヘンゼル局所環とする。
\begin{enumerate}
\item $R \to S$ が有限環写像ならば、$S$ は、それぞれ $R$ 上有限な
ヘンゼル局所環の有限積である。
\item $R \to S$ が有限環写像で $S$ が局所環ならば、$S$ は
ヘンゼル局所環であり、$R \to S$ は（有限）局所環写像である。
\item $R \to S$ が有限型環写像で、$\mathfrak q$ が $\mathfrak m$ の上にある
$S$ の素イデアルであり、その点で $R \to S$ が準有限ならば、
$S_{\mathfrak q}$ はヘンゼル的で $R$ 上有限である。
\item $R \to S$ が準有限ならば、$\mathfrak m$ の上にある任意の素イデアル
$\mathfrak q$ に対して $S_{\mathfrak q}$ はヘンゼル的で $R$ 上有限である。
\end{enumerate}
\end{lemma}

\begin{proof}
(2) から (1) が従う。実際、(1) の $S$ は補題
\ref{algebra-lemma-characterize-henselian} の (10) により、$\mathfrak m$ の上にある
素イデアルでの局所化の有限積である。
また (2) も補題 \ref{algebra-lemma-characterize-henselian} の (10) から従う。
任意の有限 $S$-代数は有限 $R$-代数でもあるからである
（もちろん、局所環の間の任意の有限環写像は局所的である）。

\medskip\noindent
$R \to S$ と $\mathfrak q$ を (3) の通りとする。
補題 \ref{algebra-lemma-characterize-henselian} の (11) により、$A$ は $R$ 上有限で、
$B$ は $\mathfrak m$ の上にあるどの素イデアルにおいても $R$ 上準有限でない
ように $S = A \times B$ と書く。このとき $S_\mathfrak q$ は
$A$ のある極大イデアルでの局所化なので、(1) から (3) が従う。
(4) は (3) から従う。
\end{proof}

\begin{lemma}
\label{algebra-lemma-mop-up}
$(R, \mathfrak m, \kappa)$ をヘンゼル局所環とする。
任意の有限型 $R$-代数 $S$ は
$S = A_1 \times \ldots \times A_n \times B$ と書ける。ここで $A_i$ は
局所環かつ $R$ 上有限であり、$R \to B$ は $\mathfrak m$ の上にある
$B$ のどの素イデアルにおいても準有限でない。
\end{lemma}

\begin{proof}
これは補題 \ref{algebra-lemma-characterize-henselian} の (11) と (10) を
組み合わせたものである。
\end{proof}

\begin{lemma}
\label{algebra-lemma-mop-up-strictly-henselian}
$(R, \mathfrak m, \kappa)$ を強ヘンゼル局所環とする。
任意の有限型 $R$-代数 $S$ は
$S = A_1 \times \ldots \times A_n \times B$ と書ける。ここで $A_i$ は
局所環かつ $R$ 上有限であり、$\kappa \subset \kappa(\mathfrak m_{A_i})$ は
有限純非分離拡大で、$R \to B$ は $\mathfrak m$ の上にある
$B$ のどの素イデアルにおいても準有限でない。
\end{lemma}

\begin{proof}
まず、補題 \ref{algebra-lemma-mop-up} のように
$S = A_1 \times \ldots \times A_n \times B$ と書く。
体拡大 $\kappa(\mathfrak m_{A_i})/\kappa$ は有限であり、$\kappa$ は
分離代数閉なので、この拡大は有限純非分離である。
\end{proof}

\begin{lemma}
\label{algebra-lemma-henselian-cat-finite-etale}
$(R, \mathfrak m, \kappa)$ をヘンゼル局所環とする。
有限エタール環拡大 $R \to S$ の圏は、関手
$S \mapsto S/\mathfrak mS$ により、有限エタール代数
$\kappa \to \overline{S}$ の圏と同値である。
\end{lemma}

\begin{proof}
主張に現れる圏の間の関手を $\mathcal{C} \to \mathcal{D}$ と書く。
$R \to S$ が有限エタールであるとする。このとき
$$
S = A_1 \times \ldots \times A_n
$$
と書ける。ここで $A_i$ は局所環で $S$ 上有限エタールである。
補題 \ref{algebra-lemma-mop-up}、または補題
\ref{algebra-lemma-characterize-henselian} の (10) を用いよ。
特に $A_i/\mathfrak mA_i$ は $\kappa$ の有限分離体拡大である。補題
\ref{algebra-lemma-etale-at-prime} を参照せよ。
従って、$\mathcal{C}$ と $\mathcal{D}$ の各対象は標準的に既約な部分へ
分解し、それらは与えられた関手によって対応する。
次に、$S_1$, $S_2$ を $R$ 上有限エタールで、
$\kappa_1 = S_1/\mathfrak mS_1$ と $\kappa_2 = S_2/\mathfrak mS_2$ が
体（$\kappa$ 上有限分離）であるとする。$S_1 \otimes_R S_2$ は
$R$ 上有限エタールであり、先と同様に
$$
S_1 \otimes_R S_2 = A_1 \times \ldots \times A_n
$$
と書ける。このとき $\Hom_R(S_1, S_2)$ は、$S_2 \to A_i$ が同型となる
添字 $i \in \{1, \ldots, n\}$ の集合と同一視される。
これを見るには、任意の $R$-代数写像 $\varphi : S_1 \to S_2$ に対して写像
$\varphi \times 1 : S_1 \otimes_R S_2 \to S_2$ が全射であり、従って
因子 $A_i$ の一つへの射影に等しいことを用いればよい。
しかし全く同様に、$\Hom_\kappa(\kappa_1, \kappa_2)$ は、
$\kappa_2 \to A_i/\mathfrak mA_i$ が同型となる添字
$i \in \{1, \ldots, n\}$ の集合と同一視される。
上の議論によりこれらの添字集合は一致するので、この関手は充満忠実である。
最後に、$\kappa'/\kappa$ を有限分離体拡大とする。補題
\ref{algebra-lemma-make-etale-map-prescribed-residue-field} により、エタール環写像
$R \to S$ と、$\mathfrak m$ の上にある $S$ の素イデアル $\mathfrak q$ で、
$\kappa \subset \kappa(\mathfrak q)$ が与えられた拡大と同型になるものが存在する。
% P01-E0961: the frozen proof cites an undefined part (1); the required decomposition is lemma-mop-up.
(1) により $S = A_1 \times \ldots \times A_n \times B$ と書ける。
$R \to S$ は準有限なので、$B$ には $\mathfrak m$ の上にある素イデアルが
存在しない。従って、ある $i$ に対して $S_{\mathfrak q}$ は $A_i$ に等しい。
ゆえに $R \to A_i$ は有限エタールであり、与えられた剰余体拡大を生じる。
従って、この関手は本質的全射であり、結論を得る。
\end{proof}

\begin{lemma}
\label{algebra-lemma-unramified-over-strictly-henselian}
$(R, \mathfrak m, \kappa)$ を強ヘンゼル局所環とする。
$R \to S$ を不分岐環写像とする。このとき
$$
S = A_1 \times \ldots \times A_n \times B
$$
と書ける。ここで各 $R \to A_i$ は全射であり、$B$ には
$\mathfrak m$ の上にある素イデアルが存在しない。
\end{lemma}

\begin{proof}
まず、補題 \ref{algebra-lemma-mop-up} のように
$S = A_1 \times \ldots \times A_n \times B$ と書く。
このとき $R \to A_i$ は有限不分岐であり、$A_i$ は局所環である。
従って $A_i$ の極大イデアルは $\mathfrak mA_i$ であり、その剰余体
$A_i / \mathfrak m A_i$ は $\kappa$ の有限分離拡大である。補題
\ref{algebra-lemma-unramified-at-prime} を参照せよ。
しかし $R$ が強ヘンゼル的であるという条件は $\kappa$ が分離代数閉であることを
意味するので、$\kappa = A_i / \mathfrak m A_i$ である。
中山の補題 \ref{algebra-lemma-NAK} により、望む通り $R \to A_i$ は全射である。
\end{proof}

\begin{lemma}
\label{algebra-lemma-complete-henselian}
\begin{slogan}
完備局所環はニュートン法によりヘンゼル的である
\end{slogan}
$(R, \mathfrak m, \kappa)$ を完備局所環とする。定義
\ref{algebra-definition-complete-local-ring} を参照せよ。
このとき $R$ はヘンゼル的である。
\end{lemma}

\begin{proof}
$f \in R[T]$ をモニックとする。
その像を $f_n \in R/\mathfrak m^{n + 1}[T]$ と書く。
$T$ に関する $f_n$ の導関数を $f'_n$ と書く。
$a_0 \in \kappa$ を $f_0$ の単純根とする。以下のように帰納的に、これを
$R$ 上の $f$ の解へ持ち上げる。$a_n \bmod \mathfrak m = a_0$ かつ
$f_n(a_n) = 0$ を満たす $a_n \in R/\mathfrak m^{n + 1}$ が与えられたとする。
$a_n = b \bmod \mathfrak m^{n + 1}$ を満たす任意の元
$b \in R/\mathfrak m^{n + 2}$ を選ぶ。このとき
$f_{n + 1}(b) \in \mathfrak m^{n + 1}/\mathfrak m^{n + 2}$ である。
ここで
$$
a_{n + 1} = b - f_{n + 1}(b)/f'_{n + 1}(b)
$$
とおく（ニュートン法）。$a_0$ に関する条件により
$f'_{n + 1}(b) \in R/\mathfrak m^{n + 2}$ は可逆なので、これは意味をもつ。
すると $R/\mathfrak m^{n + 2}$ において
$f_{n + 1}(a_{n + 1}) = f_{n + 1}(b) - f_{n + 1}(b) = 0$ と計算できる。
このように構成した元の系 $a_n \in R/\mathfrak m^{n + 1}$ は両立するので、
元 $a \in \lim R/\mathfrak m^n = R$ を得る
（ここで $R$ の完備性を用いた）。さらに、これは各
$R/\mathfrak m^n$ で零に写るので $f(a) = 0$ である。
最後に $\overline{a} = a_0$ であり、結論を得る。
\end{proof}

\begin{lemma}
\label{algebra-lemma-local-dimension-zero-henselian}
\begin{slogan}
次元零の局所環はヘンゼル的である。
\end{slogan}
$(R, \mathfrak m)$ を次元 $0$ の局所環とする。
このとき $R$ はヘンゼル的である。
\end{lemma}

\begin{proof}
$R \to S$ を有限環写像とする。補題
\ref{algebra-lemma-characterize-henselian} により、$S$ が局所環の積であることを
示せば十分である。補題 \ref{algebra-lemma-finite-finite-fibres} により、
$S$ は有限個の素イデアル $\mathfrak m_1, \ldots, \mathfrak m_r$ をもち、
それらはすべて $\mathfrak m$ の上にある。補題
\ref{algebra-lemma-integral-no-inclusion} により、これらの素イデアルの間には
包含関係がないので、すべて極大である。
補題 \ref{algebra-lemma-Zariski-topology} により
$\mathfrak m_1 \cap \ldots \cap \mathfrak m_r$ の各元は冪零である。
従って補題 \ref{algebra-lemma-product-local} により、$S$ は素イデアル
$\mathfrak m_i$ での $S$ の局所化の積である。
\end{proof}

\noindent
次の補題は、ヘンゼル化および強ヘンゼル化の一意性と関手性にとって
鍵となる。

\begin{lemma}
\label{algebra-lemma-map-into-henselian}
$R \to S$ を環写像とし、$S$ はヘンゼル局所環であるとする。次を与える。
\begin{enumerate}
\item エタール環写像 $R \to A$。
\item $\mathfrak p = R \cap \mathfrak m_S$ の上にある $A$ の素イデアル
$\mathfrak q$。
% P01-E0962: for a possibly noninjective map R -> S, the frozen expression R intersect m_S should be the inverse image of m_S in R.
\item $\kappa(\mathfrak p)$-代数写像
$\tau : \kappa(\mathfrak q) \to S/\mathfrak m_S$。
\end{enumerate}
このとき、$\mathfrak q = f^{-1}(\mathfrak m_S)$ を満たし、$f$ が剰余体上で
写像 $\tau$ を誘導する $R$-代数準同型 $f : A \to S$ が一意に存在する。
\end{lemma}

\begin{proof}
$A \otimes_R S$ を考える。これは $S$ 上のエタール代数である。補題
\ref{algebra-lemma-etale} を参照せよ。さらに、核
$$
\mathfrak q' = \Ker(A \otimes_R S \to
\kappa(\mathfrak q) \otimes_{\kappa(\mathfrak p)} \kappa(\mathfrak m_S)
\xrightarrow{\tau \otimes 1}
\kappa(\mathfrak m_S))
$$
は $\mathfrak m_S$ の上にある素イデアルであり、その剰余体は $S$ の
剰余体に等しい。従って補題 \ref{algebra-lemma-characterize-henselian} により、
$\sigma^{-1}(\mathfrak m_S) = \mathfrak q'$ を満たす一意なレトラクション
$\sigma : A \otimes_R S \to S$ が存在する。
$f$ を合成 $A \to A \otimes_R S \to S$ とする。
$f$ の性質の確認は省略する。$f$ の一意性は $\sigma$ の一意性から従う
（詳細は省略する）。
\end{proof}

\begin{lemma}
\label{algebra-lemma-strictly-henselian-solutions}
$\varphi : R \to S$ を強ヘンゼル局所環の局所準同型とする。
$P_1, \ldots, P_n \in R[x_1, \ldots, x_n]$ を、
$R[x_1, \ldots, x_n]/(P_1, \ldots, P_n)$ が $R$ 上エタールとなる
多項式とする。このとき写像
$$
R^n \longrightarrow S^n, \quad
(h_1, \ldots, h_n) \longmapsto (\varphi(h_1), \ldots, \varphi(h_n))
$$
は、
$$
\{
(r_1, \ldots, r_n) \in R^n
\mid
P_i(r_1, \ldots, r_n) = 0, \ i = 1, \ldots, n
\}
$$
と
$$
\{
(s_1, \ldots, s_n) \in S^n
\mid
P^\varphi_i(s_1, \ldots, s_n) = 0, \ i = 1, \ldots, n
\}
$$
との間の全単射を誘導する。ここで
$P^\varphi_i \in S[x_1, \ldots, x_n]$ は $\varphi$ による $P_i$ の像である。
\end{lemma}

\begin{proof}
第一の解集合は、集合
$$
\Hom_R(R[x_1, \ldots, x_n]/(P_1, \ldots, P_n), R).
$$
と標準的に同型である。$R$ はヘンゼル的なので、写像
$R \to R/\mathfrak m_R$ は、この集合と剰余体 $R/\mathfrak m_R$ における
解集合との間の全単射を誘導する。補題
\ref{algebra-lemma-characterize-henselian} を参照せよ。$S$ についても同様である。
ここで $R[x_1, \ldots, x_n]/(P_1, \ldots, P_n)$ は $R$ 上エタールであり、
$R/\mathfrak m_R$ は分離代数閉なので、
$R/\mathfrak m_R[x_1, \ldots, x_n]/(\overline{P}_1, \ldots, \overline{P}_n)$
は $R/\mathfrak m_R$ のコピーの有限積である。ここで $\overline{P}_i$ は
$R/\mathfrak m_R[x_1, \ldots, x_n]$ における $P_i$ の像である。従ってテンソル積
$$
R/\mathfrak m_R[x_1, \ldots, x_n]/(\overline{P}_1, \ldots, \overline{P}_n)
\otimes_{R/\mathfrak m_R} S/\mathfrak m_S
=
S/\mathfrak m_S[x_1, \ldots, x_n]/
(\overline{P}^\varphi_1, \ldots, \overline{P}^\varphi_n)
$$
% P01-E0963: the frozen residue-field polynomial-ring expressions omit parentheses around R/m_R and S/m_S, making the quotients formally ill-typed or ambiguous.
もまた、同じ添字集合をもつ $S/\mathfrak m_S$ のコピーの有限積である。
これで補題が示された。
\end{proof}

\begin{lemma}
\label{algebra-lemma-split-ML-henselian}
$R$ をヘンゼル局所環とする。$R$ 上の任意の可算生成 Mittag-Leffler 加群は、
有限表示 $R$-加群の直和である。
\end{lemma}

\begin{proof}
$M$ を可算生成 Mittag-Leffler $R$-加群とする。
任意の元 $x \in M$ に対し、$x \in N$、$N$ は有限表示、$K$ は
Mittag-Leffler となる直和分解 $M = N \oplus K$ が存在すると主張する。

\medskip\noindent
この主張が正しいとする。$M$ の生成元 $x_1, x_2, x_3, \ldots$ を選ぶ。
この主張により、帰納的に直和分解
$$
M = N_1 \oplus N_2 \oplus \ldots \oplus N_n \oplus K_n
$$
を見いだせる。ここで $N_i$ は有限表示であり、
$x_1, \ldots, x_n \in N_1 \oplus \ldots \oplus N_n$、また $K_n$ は
Mittag-Leffler である。これを無限に繰り返すと $M = \bigoplus N_i$ を得る。

\medskip\noindent
主張を証明しなければならない。$x \in M$ とする。補題
\ref{algebra-lemma-ML-countable} により、$\alpha$ が有限表示加群を経由し、かつ
$\alpha (x) = x$ を満たす自己準同型 $\alpha : M \to M$ が存在する。
$\alpha$ が
$$
\xymatrix{
M \ar[r]^\pi & P \ar[r]^i & M
}
$$
と分解するとする。$a = \pi \circ \alpha \circ i : P \to P$ とおくと、
$i \circ a \circ \pi = \alpha^3$ である。補題
\ref{algebra-lemma-charpoly-module} により、$P(a) = 0$ を満たすモニック多項式
$P \in R[T]$ が存在する。
% P01-E0964: the frozen proof reuses P for both the finitely presented module and the annihilating monic polynomial, making later references ambiguous.
これは形式的に $\alpha^2 P(\alpha) = 0$ を意味する。
従って $M$ を $R[T]/(T^2P)$ 上の加群とみなしてよい。
$x \not = 0$ と仮定する。$\alpha(x) = x$ より
$0 = \alpha^2P(\alpha)x = P(1)x$ であり、従って
$I = \{r \in R \mid rx = 0\}$ を $x$ の零化イデアルとすると、
$R/I$ において $P(1) = 0$ である。
$x \not = 0$ なので $I \subset \mathfrak m_R$ であり、従って
$1$ は $\overline{P} = P \bmod \mathfrak m_R \in R/\mathfrak m_R[T]$ の根である。
$R$ はヘンゼル的なので、補題 \ref{algebra-lemma-characterize-henselian} により、
ある $Q_1, Q_2 \in R[T]$ に対して分解
$$
T^2P = (T^2 Q_1) Q_2
$$
を見いだせる。ここで
$Q_2 = (T - 1)^e \bmod \mathfrak m_R R[T]$ かつ
$Q_1(1) \not = 0 \bmod \mathfrak m_R$ である。
$N = \Im(\alpha^2Q_1(\alpha) : M \to M)$ および
$K = \Im(Q_2(\alpha) : M \to M)$ とおく。
$T^2Q_1$ と $Q_2$ は $R[T]$ の単位イデアルを生成するので、
直和分解 $M = N \oplus K$ を得る。さらに、$Q_2$ は $N$ 上零として作用し、
$T^2Q_1$ は $K$ 上零として作用する。$N$ は $P$ の商なので有限生成である。
また、$\alpha^2Q_1(\alpha)x = Q_1(1)x$ かつ $Q_1(1)$ は $R$ の単元なので、
$x \in N$ である。補題 \ref{algebra-lemma-direct-sum-ML} により、加群 $N$ と $K$ は
Mittag-Leffler である。最後に、有限生成 Mittag-Leffler 加群は有限表示なので、
有限生成加群 $N$ は有限表示である。例 \ref{algebra-example-ML} の (1) を参照せよ。
\end{proof}

\section{エタール環写像のフィルター余極限}
\label{algebra-section-ind-etale}

\noindent
この節は ind-エタール環写像を扱う節
（Pro-エタールコホモロジー、第 \ref{proetale-section-ind-etale} 節）の準備である。
この内容は、局所環のヘンゼル化および強ヘンゼル化の一意性を証明する際にも
有用となる。

\begin{lemma}
\label{algebra-lemma-base-change-colimit-etale}
$R \to A$ および $R \to R'$ を環写像とする。$A$ がエタール
$R$-代数のフィルター余極限ならば、$R' \otimes_R A$ はエタール
$R'$-代数のフィルター余極限である。
% P01-E0965: the frozen English statement duplicates its predicate: then so is ... is a filtered colimit.
\end{lemma}

\begin{proof}
余極限はテンソル積と可換し、エタール環写像は基底変換によって保たれる
（補題 \ref{algebra-lemma-etale}）ので、これは正しい。
\end{proof}

\begin{lemma}
\label{algebra-lemma-composition-colimit-etale}
$A \to B \to C$ を環写像とする。$B$ がエタール $A$-代数の
フィルター余極限で、$C$ がエタール $B$-代数のフィルター余極限ならば、
$C$ はエタール $A$-代数のフィルター余極限である。
\end{lemma}

\begin{proof}
補題 \ref{algebra-lemma-when-colimit} の判定法を用いる。
$A \to C$ を、$P$ が $A$ 上有限表示となるように
$A \to P \to C$ と分解する。$I$ を有向集合、$B_i$ をエタール
$A$-代数として $B = \colim_{i \in I} B_i$ と書く。
$J$ を有向集合、$C_j$ をエタール $B$-代数として
$C = \colim_{j \in J} C_j$ と書く。
補題 \ref{algebra-lemma-characterize-finite-presentation} により、ある $j$ に対して
$P \to C$ を $P \to C_j \to C$ と分解できる。
補題 \ref{algebra-lemma-etale} により、$i \in I$ とエタール環写像
$B_i \to C'_j$ で $C_j = B \otimes_{B_i} C'_j$ を満たすものが存在する。
このとき $C_j = \colim_{i' \geq i} B_{i'} \otimes_{B_i} C'_j$ であり、
再び $P \to C_j$ が
$P \to B_{i'} \otimes_{B_i} C'_j \to C$ と分解することが分かる。
合成およびテンソル積はエタール環写像を保つので、
$A \to C' = B_{i'} \otimes_{B_i} C'_j$ はエタールである。
% P01-E0966: the frozen English begins this sentence with As and then starts a second causal as, leaving a sentence fragment.
従って $P \to C$ を $P \to C' \to C$ と分解でき、$C'$ は $A$ 上
エタールであるから、補題 \ref{algebra-lemma-when-colimit} の判定法が適用できる。
\end{proof}

\begin{lemma}
\label{algebra-lemma-colimit-colimit-etale}
$R$ を環とする。$A = \colim A_i$ を $R$-代数のフィルター余極限とし、
各 $A_i$ はエタール $R$-代数のフィルター余極限であるとする。
このとき $A$ はエタール $R$-代数のフィルター余極限である。
\end{lemma}

\begin{proof}
$J_i$ を有向集合、$A_j$ をエタール $R$-代数として
$A_i = \colim_{j \in J_i} A_j$ と書く。
各 $i \leq i'$ と $j \in J_i$ に対して、図式
$$
\xymatrix{
A_i \ar[r] & A_{i'} \\
A_j \ar[u] \ar[r]^{\varphi_{jj'}} & A_{j'} \ar[u]
}
$$
を可換にする $j' \in J_{i'}$ と $R$-代数写像
$\varphi_{jj'} : A_j \to A_{j'}$ が存在する。
これは $R \to A_j$ が有限表示であり、従って補題
\ref{algebra-lemma-characterize-finite-presentation} が適用できるからである。
$\mathcal{J}$ を、対象が $\coprod_{i \in I} J_i$ で、射が上のような三つ組
$(j, j', \varphi_{jj'})$ である圏とする（合成則は明らかなものとする）。
このとき $\mathcal{J}$ はフィルター圏であり、
$A = \colim_\mathcal{J} A_j$ である。詳細は省略する。
\end{proof}

\begin{lemma}
\label{algebra-lemma-colimit-colimit-etale-better}
$I$ を有向集合とする。$i \mapsto (R_i \to A_i)$ を $I$ 上の環の射の系とする。
$R = \colim R_i$ および $A = \colim A_i$ とおく。各 $A_i$ がエタール
$R_i$-代数のフィルター余極限ならば、$A$ はエタール $R$-代数の
フィルター余極限である。
\end{lemma}

\begin{proof}
$A = A \otimes_R R = \colim A_i \otimes_{R_i} R$ なので、補題
\ref{algebra-lemma-colimit-colimit-etale} を適用できる。実際、補題
\ref{algebra-lemma-base-change-colimit-etale} により $R \to A_i \otimes_{R_i} R$ は
エタール環写像のフィルター余極限である。
\end{proof}

\begin{lemma}
\label{algebra-lemma-colimits-of-etale}
$R$ を環とし、$A \to B$ を $R$-代数準同型とする。$A$ と $B$ が
エタール $R$-代数のフィルター余極限ならば、$B$ はエタール
$A$-代数のフィルター余極限である。
\end{lemma}

\begin{proof}
$A_i$ と $B_j$ が $R$ 上エタールとなるフィルター余極限として
$A = \colim A_i$ および $B = \colim B_j$ と書く。
各 $i$ に対し、$A_i \to B$ が $B_j$ を経由するような $j$ をとれる。
補題 \ref{algebra-lemma-characterize-finite-presentation} を参照せよ。
分解 $A_i \to B_j$ は補題 \ref{algebra-lemma-map-between-etale} によりエタールである。
$A \to A \otimes_{A_i} B_j$ はエタールなので（補題
\ref{algebra-lemma-etale}）、$A_i \to B$ の分解 $A_i \to B_j \to B$ と組
$(i, j)$ 全体にわたる余極限として
$B = \colim A \otimes_{A_i} B_j$ を示せば十分である
（これは有向系である。詳細は省略する）。これは余極限がテンソル積と可換し、
従って $\colim A \otimes_{A_i} B_j = A \otimes_A B = B$ となることから明らかである。
\end{proof}

\begin{lemma}
\label{algebra-lemma-map-into-henselian-colimit}
$R \to S$ を環写像とし、$S$ はヘンゼル局所環であるとする。次を与える。
\begin{enumerate}
\item エタール $R$-代数のフィルター余極限である $R$-代数 $A$。
\item $\mathfrak p = R \cap \mathfrak m_S$ の上にある $A$ の素イデアル
$\mathfrak q$。
% P01-E0967: as in P01-E0962, the frozen source uses an intersection for contraction along a possibly noninjective ring map.
\item $\kappa(\mathfrak p)$-代数写像
$\tau : \kappa(\mathfrak q) \to S/\mathfrak m_S$。
\end{enumerate}
このとき、$\mathfrak q = f^{-1}(\mathfrak m_S)$ および
$f \bmod \mathfrak q = \tau$ を満たす一意な $R$-代数準同型
$f : A \to S$ が存在する。
\end{lemma}

\begin{proof}
エタール $R$-代数のフィルター余極限として $A = \colim A_i$ と書く。
$\mathfrak q_i = A_i \cap \mathfrak q$ とおく。
% P01-E0968: the frozen source again writes an intersection although A_i -> A need not be injective; contraction is the inverse image.
補題 \ref{algebra-lemma-map-into-henselian} を適用して $f_i : A_i \to S$ を得る。
$f = \colim f_i$ とおく。
\end{proof}

\begin{lemma}
\label{algebra-lemma-uniqueness-henselian}
$R$ を環とする。環写像の可換図式
$$
\xymatrix{
S \ar[r] & K \\
R \ar[u] \ar[r] & S' \ar[u]
}
$$
が与えられたとする。ここで $S$, $S'$ はヘンゼル局所環で、
$S$, $S'$ はエタール $R$-代数のフィルター余極限、$K$ は体であり、
射 $S \to K$ と $S' \to K$ は $K$ を $S$ と $S'$ の双方の剰余体と
同一視する。このとき、$K$ への写像と両立する $R$-代数同型
$S \to S'$ が一意に存在する。
\end{lemma}

\begin{proof}
補題 \ref{algebra-lemma-map-into-henselian-colimit} から直ちに従う。
\end{proof}


\noindent
次の補題は、厳密にはエタール環写像の余極限に関するものではない。

\begin{lemma}
\label{algebra-lemma-colimit-henselian}
局所準同型に沿う（強）ヘンゼル局所環のフィルター余極限は
（強）ヘンゼル的である。
\end{lemma}

\begin{proof}
圏論、補題 \ref{categories-lemma-directed-category-system} によれば、これは実際には
有向集合上の（強）ヘンゼル局所環の余極限についての問題にすぎない。
各 $\varphi_{ii'}$ が局所的であるような系 $(R_i, \varphi_{ii'})$ をとる。
このとき $R = \colim_i R_i$ は局所環で、その剰余体 $\kappa$ は
$\colim \kappa_i$ である（議論は省略する）。各 $\kappa_i$ が分離代数閉ならば
$\colim \kappa_i$ も分離代数閉であることは容易に分かる。従って、各 $R_i$ が
ヘンゼル的ならば $R$ もヘンゼル的であることを示せば十分である。
$f \in R[T]$ がモニックで、$a_0 \in \kappa$ が $\overline{f}$ の単純根で
あるとする。このとき十分大きなある $i$ に対して、$f$ に写る
$f_i \in R_i[T]$ と、$a_0$ に写る $a_{0, i} \in \kappa_i$ が存在する。
$\overline{f_i}(a_{0, i}) \in \kappa_i$ および
$\overline{f_i'}(a_{0, i}) \in \kappa_i$ はそれぞれ
$0 = \overline{f}(a_0) \in \kappa$ および
$0 \not = \overline{f'}(a_0) \in \kappa$ に写るので、
$a_{0, i}$ は $\overline{f_i}$ の単純根であると結論できる。
$R_i$ はヘンゼル的なので、$f_i(a_i) = 0$ および
$a_{0, i} = \overline{a_i}$ を満たす $a_i \in R_i$ が存在する。
$a_i$ の像 $a \in R$ が求める解である。従って $R$ はヘンゼル的である。
\end{proof}

\section{ヘンゼル化と強ヘンゼル化}
\label{algebra-section-henselization}

\noindent
この節ではヘンゼル化を構成する。この内容を読む際には、補題
\ref{algebra-lemma-uniqueness-henselian} ですでに証明した一意性と、補題
\ref{algebra-lemma-map-into-henselian-colimit} で指摘した関手的な振る舞いを
念頭に置くことを勧める。

\begin{lemma}
\label{algebra-lemma-henselization}
$(R, \mathfrak m, \kappa)$ を局所環とする。次の性質をもつ局所環写像
$R \to R^h$ が存在する。
\begin{enumerate}
\item $R^h$ はヘンゼル的である。
\item $R^h$ はエタール $R$-代数のフィルター余極限である。
\item $\mathfrak m R^h$ は $R^h$ の極大イデアルである。
\item $\kappa = R^h/\mathfrak m R^h$ である。
\end{enumerate}
\end{lemma}

\begin{proof}
$R \to S$ がエタール環写像で、$\mathfrak q$ が $\mathfrak m$ の上にある
$S$ の素イデアルで $\kappa = \kappa(\mathfrak q)$ を満たすような組
$(S, \mathfrak q)$ の圏を考える。組の射
$(S, \mathfrak q) \to (S', \mathfrak q')$ は、
$\varphi^{-1}(\mathfrak q') = \mathfrak q$ を満たす $R$-代数写像
$\varphi : S \to S'$ によって与えられる。次のようにおく。
$$
R^h = \colim_{(S, \mathfrak q)} S.
$$
組の圏がフィルター的であることを示そう。圏論、定義
\ref{categories-definition-directed} を参照せよ。
この圏は組 $(R, \mathfrak m)$ を含むので空でなく、これにより圏論、定義
\ref{categories-definition-directed} の (1) が示される。
任意の組 $(S, \mathfrak q)$ に対して、合成
$\kappa \to S/\mathfrak q \to \kappa(\mathfrak q)$ は同型なので、
素イデアル $\mathfrak q$ は極大で、その剰余体は $\kappa$ である。
$(S, \mathfrak q)$ と $(S', \mathfrak q')$ を二つの対象とする。
$S'' = S \otimes_R S'$ および
$\mathfrak q'' = \mathfrak qS'' + \mathfrak q'S''$ とおく。
上で述べたことにより
$S''/\mathfrak q'' = S/\mathfrak q \otimes_R S'/\mathfrak q' = \kappa$ である。
さらに補題 \ref{algebra-lemma-etale} により $R \to S''$ はエタールである。
これにより圏論、定義 \ref{categories-definition-directed} の (2) が示された。
次に、$\varphi, \psi : (S, \mathfrak q) \to (S', \mathfrak q')$ を
組の二つの射とする。このとき補題 \ref{algebra-lemma-map-between-etale} により、
$\varphi$, $\psi$ および $S' \otimes_R S' \to S'$ はエタール環写像である。
$$
S'' = (S' \otimes_{\varphi, S, \psi} S')
\otimes_{S' \otimes_R S'} S'
$$
を、素イデアル
$$
\mathfrak q'' =
(\mathfrak q' \otimes S' + S' \otimes \mathfrak q') \otimes S'
+
(S' \otimes_{\varphi, S, \psi} S') \otimes \mathfrak q'
$$
とともに考える。上と同様の議論（エタール写像の基底変換はエタールであり、
エタール写像の合成はエタールである）により、$S''$ は $R$ 上エタールである。
さらに、標準写像 $S' \to S''$（例えば最も右の因子を用いる）は
$\varphi$ と $\psi$ を等化する。これにより圏論、定義
\ref{categories-definition-directed} の (3) が示された。
従って、$R^h$ の元は $f \in S$ を伴う三つ組
$(S, \mathfrak q, f)$ からなり、二つの三つ組
$(S, \mathfrak q, f)$、$(S', \mathfrak q', f')$ が $R^h$ の同じ元を
定めるための必要十分条件は、組 $(S'', \mathfrak q'')$ と組の射
$\varphi : (S, \mathfrak q) \to (S'', \mathfrak q'')$ および
$\varphi' : (S', \mathfrak q') \to (S'', \mathfrak q'')$ で
$\varphi(f) = \varphi'(f')$ を満たすものが存在することである。

\medskip\noindent
$x \in R^h$ とする。$x$ を三つ組 $(S, \mathfrak q, f)$ で表す。
$\mathfrak q_1, \ldots, \mathfrak q_r$ を、$\mathfrak m$ の上にある
$S$ の他の素イデアルとする。上で $\mathfrak q$ が極大であることを見たので、
$\mathfrak q \not \subset \mathfrak q_i$ である。従って
$\mathfrak q$ は素イデアルなので、$g \in S$ で
$g \not \in \mathfrak q$ かつ $i = 1, \ldots, r$ に対して
$g \in \mathfrak q_i$ となるものをとれる。組の射
$(S, \mathfrak q) \to (S_g, \mathfrak qS_g)$ を考える。
このようにして、$x$ は、$\mathfrak q$ が $\mathfrak m$ の上にある
$S$ の唯一の素イデアル、すなわち
$\sqrt{\mathfrak mS} = \mathfrak q$ となる三つ組
$(S, \mathfrak q, f)$ で常に与えられると仮定してよい。
しかし $R \to S$ はエタールなので、補題
\ref{algebra-lemma-etale-at-prime} により
$\mathfrak mS_{\mathfrak q} = \mathfrak qS_{\mathfrak q}$ である。
従って実際に $\mathfrak mS = \mathfrak q$ を得る。

\medskip\noindent
$x \not \in \mathfrak mR^h$ とする。$x$ を
$\mathfrak mS = \mathfrak q$ となる三つ組 $(S, \mathfrak q, f)$ で表す。
このとき $f \not \in \mathfrak mS$、すなわち $f \not \in \mathfrak q$ である。
従って $(S, \mathfrak q) \to (S_f, \mathfrak qS_f)$ は、$f$ の像が可逆となる
組の射である。ゆえに $x$ は可逆で、その逆元は三つ組
$(S_f, \mathfrak qS_f, 1/f)$ で表される。従って $R^h$ は極大イデアル
$\mathfrak mR^h$ をもつ局所環である。剰余体は $\kappa$ である。
実際、三つ組 $(S, \mathfrak q, f)$ を $\mathfrak q$ を法とする $f$ の
剰余類へ写すことにより $R^h/\mathfrak mR^h \to \kappa$ を定義できる。

\medskip\noindent
$R^h$ がヘンゼル的であることを示さなければならない。
$P \in R^h[T]$ をモニック多項式とし、$a_0 \in \kappa$ を
還元 $\overline{P} \in \kappa[T]$ の単純根とする。
このとき、$P$ があるモニック多項式 $Q \in S[T]$ の像となるような
組 $(S, \mathfrak q)$ が存在する。$S' = S[T]/(Q)$ とおき、
$a' \in S$ を $a_0$ の任意の持ち上げとして、$\mathfrak q' \subset S'$ を
極大イデアル $\mathfrak q' = \mathfrak qS' + (T - a')S'$ とする。
構成により $S \to S'$ は $\mathfrak q'$ においてエタールで、
$\kappa = \kappa(\mathfrak q')$ である。
$g \in S'$, $g \not \in \mathfrak q'$ で $S'' = S'_g$ が $S$ 上
エタールとなるものを選ぶ。このとき
$(S, \mathfrak q) \to (S'', \mathfrak q'S'')$ は組の射である。
ここで三つ組 $(S'', \mathfrak q'S'', \text{class of }T)$ は、
$P(a) = 0$ および $\overline{a} = a_0$ を満たす元 $a \in R^h$ を定める。
% P01-E0969: the frozen English says Now that triple, substituting that for the article the.
これは求める性質である。
\end{proof}

\begin{lemma}
\label{algebra-lemma-strict-henselization}
$(R, \mathfrak m, \kappa)$ を局所環とする。
$\kappa \subset \kappa^{sep}$ を分離代数閉包とする。次の可換図式が存在する。
$$
\xymatrix{
\kappa \ar[r] & \kappa \ar[r] & \kappa^{sep} \\
R \ar[r] \ar[u] & R^h \ar[r] \ar[u] & R^{sh} \ar[u]
}
$$
これは次の性質をもつ。
\begin{enumerate}
\item 写像 $R^h \to R^{sh}$ は局所的である。
\item $R^{sh}$ は強ヘンゼル的である。
\item $R^{sh}$ はエタール $R$-代数のフィルター余極限である。
\item $\mathfrak m R^{sh}$ は $R^{sh}$ の極大イデアルである。
\item $\kappa^{sep} = R^{sh}/\mathfrak m R^{sh}$ である。
\end{enumerate}
\end{lemma}

\begin{proof}
補題 \ref{algebra-lemma-henselization} と全く同じ証明で示される。
唯一の違いは、組の代わりに三つ組 $(S, \mathfrak q, \alpha)$ を用いることである。
ここで $R \to S$ はエタール、$\mathfrak q$ は $\mathfrak m$ の上にある
$S$ の素イデアル、$\alpha : \kappa(\mathfrak q) \to \kappa^{sep}$ は
$\kappa$ の拡大の埋め込みである。
\end{proof}

\begin{definition}
\label{algebra-definition-henselization}
$(R, \mathfrak m, \kappa)$ を局所環とする。
\begin{enumerate}
\item 補題 \ref{algebra-lemma-henselization} で構成した局所環写像
$R \to R^h$ を $R$ の {\it ヘンゼル化} という。
\item 分離代数閉包 $\kappa \subset \kappa^{sep}$ が与えられたとき、補題
\ref{algebra-lemma-strict-henselization} で構成した局所環写像 $R \to R^{sh}$ を
{\it $\kappa \subset \kappa^{sep}$ に関する $R$ の強ヘンゼル化} という。
\item 局所環写像 $R \to R^{sh}$ が、補題
\ref{algebra-lemma-strict-henselization} で構成した局所環写像の一つと同型であるとき、
これを $R$ の {\it 強ヘンゼル化} という。
% P01-E0970: the frozen third definition item lacks terminal punctuation.
\end{enumerate}
\end{definition}

\noindent
写像 $R \to R^h \to R^{sh}$ は平坦な局所環準同型である。
補題 \ref{algebra-lemma-uniqueness-henselian} により、$R$-代数 $R^h$ と $R^{sh}$ は、
それぞれ剰余体 $\kappa$ と $\kappa^{sep}$ をもち、エタール
$R$-代数のフィルター余極限であるヘンゼル局所環という条件によって、
一意な同型を除き定まる。この節の残りでは、主に（強）ヘンゼル化の
関手性を論じる。$R$ とそのヘンゼル化との関係についてのより精緻な結果は、
代数学詳論、第 \ref{more-algebra-section-permanence-henselization} 節で論じる。

\begin{remark}
\label{algebra-remark-construct-sh-from-h}
$R^{sh}$ を $R^h$ から構成することもできる。実際、任意の有限分離中間拡大
$\kappa^{sep}/\kappa'/\kappa$ に対して、剰余体拡大が与えられた拡大を再現する
有限エタール局所環拡大 $R^h \subset R^h(\kappa')$ が、一意な同型を除き
一意に存在する。補題 \ref{algebra-lemma-henselian-cat-finite-etale} を参照せよ。
従って
$$
R^{sh} =
\bigcup\nolimits_{\kappa \subset \kappa' \subset \kappa^{sep}}
R^h(\kappa')
$$
とおける。剰余体の水準の射と両立するこの系の射は、補題
\ref{algebra-lemma-henselian-cat-finite-etale} により存在する。
補題 \ref{algebra-lemma-colimit-henselian} により、これはヘンゼル局所環を与える。
実際、各環 $R^h(\kappa')$ は補題 \ref{algebra-lemma-finite-over-henselian} により
ヘンゼル的である。構成により、$R^h \to R^{sh}$ が誘導する
剰余体拡大は体拡大 $\kappa^{sep}/\kappa$ である。
従って、このように構成した $R^{sh}$ は強ヘンゼル的である。
補題 \ref{algebra-lemma-composition-colimit-etale} により、$R$-代数 $R^{sh}$ は
エタール $R$-代数の余極限である。従って補題
\ref{algebra-lemma-uniqueness-henselian} の一意性から、$R^{sh}$ は強ヘンゼル化である。
\end{remark}

\begin{lemma}
\label{algebra-lemma-henselian-functorial-prepare}
$R \to S$ を局所環の局所写像とする。$S \to S^h$ をヘンゼル化とする。
$R \to A$ をエタール環写像とし、$\mathfrak q$ を $\mathfrak m_R$ の上にある
$A$ の素イデアルで $R/\mathfrak m_R \cong \kappa(\mathfrak q)$ を
満たすものとする。このとき、可換図式
$$
\xymatrix{
A \ar[r]_f & S^h \\
R \ar[u] \ar[r] & S \ar[u]
}
$$
に入る環の射 $f : A \to S^h$ で
$f^{-1}(\mathfrak m_{S^h}) = \mathfrak q$ を満たすものが一意に存在する。
\end{lemma}

\begin{proof}
これは補題 \ref{algebra-lemma-map-into-henselian} の特別な場合である。
\end{proof}

\begin{lemma}
\label{algebra-lemma-henselian-functorial}
$R \to S$ を局所環の局所写像とする。
$R \to R^h$ と $S \to S^h$ をそれぞれのヘンゼル化とする。
可換図式
$$
\xymatrix{
R^h \ar[r]_f & S^h \\
R \ar[u] \ar[r] & S \ar[u]
}
$$
に入る局所環写像 $R^h \to S^h$ が一意に存在する。
\end{lemma}

\begin{proof}
補題 \ref{algebra-lemma-map-into-henselian-colimit} から直ちに従う。
\end{proof}

\noindent
以下は、ヘンゼル化の少し異なる構成である。

\begin{lemma}
\label{algebra-lemma-henselization-different}
$R$ を環とし、$\mathfrak p \subset R$ を素イデアルとする。
$R \to S$ がエタールで、$\mathfrak q$ が $\mathfrak p$ の上にある素イデアルで
$\kappa(\mathfrak p) = \kappa(\mathfrak q)$ を満たすような組
$(S, \mathfrak q)$ の圏を考える。この圏はフィルター的であり、標準的に
$$
(R_{\mathfrak p})^h = \colim_{(S, \mathfrak q)} S
= \colim_{(S, \mathfrak q)} S_{\mathfrak q}
$$
である。
\end{lemma}

\begin{proof}
組の射 $(S, \mathfrak q) \to (S', \mathfrak q')$ は、
$\varphi^{-1}(\mathfrak q') = \mathfrak q$ を満たす $R$-代数写像
$\varphi : S \to S'$ によって与えられる。組の圏がフィルター的であることを
示そう。圏論、定義 \ref{categories-definition-directed} を参照せよ。
この圏は組 $(R, \mathfrak p)$ を含むので空でなく、これにより圏論、定義
\ref{categories-definition-directed} の (1) が示される。
$(S, \mathfrak q)$ と $(S', \mathfrak q')$ を二つの組とする。
$\mathfrak q$ と $\mathfrak q'$ はそれぞれ、剰余体が
$\kappa(\mathfrak p)$ であるファイバー環
$S \otimes \kappa(\mathfrak p)$ と $S' \otimes \kappa(\mathfrak p)$ の
素イデアルに対応し、従ってそれぞれ
$S \otimes \kappa(\mathfrak p)$ と $S' \otimes \kappa(\mathfrak p)$ の
極大イデアルに対応することに注意せよ。
$S'' = S \otimes_R S'$ とおく。上のことから、$\mathfrak q$ と
$\mathfrak q'$ の上にあり、剰余体が $\kappa(\mathfrak p)$ である
一意な素イデアル $\mathfrak q'' \subset S''$ が存在する。
補題 \ref{algebra-lemma-etale} により環写像 $R \to S''$ はエタールである。
これにより圏論、定義 \ref{categories-definition-directed} の (2) が示された。
次に、$\varphi, \psi : (S, \mathfrak q) \to (S', \mathfrak q')$ を
組の二つの射とする。このとき補題 \ref{algebra-lemma-map-between-etale} により、
$\varphi$, $\psi$ および $S' \otimes_R S' \to S'$ はエタール環写像である。
次を考える。
$$
S'' = (S' \otimes_{\varphi, S, \psi} S')
\otimes_{S' \otimes_R S'} S'
$$
上と同様の議論（エタール写像の基底変換はエタールであり、エタール写像の
合成はエタールである）により、$S''$ は $R$ 上エタールである。
$\mathfrak p$ 上の $S''$ のファイバー環は
$$
F'' = (F' \otimes_{\varphi, F, \psi} F')
\otimes_{F' \otimes_{\kappa(\mathfrak p)} F'} F'
$$
である。ここで $F', F$ はそれぞれ $S'$ と $S$ のファイバー環である。
$\varphi$ と $\psi$ は組の射なので、$\mathfrak p'$ に対応する写像
$F' \to \kappa(\mathfrak p)$ は写像 $F'' \to \kappa(\mathfrak p)$ へ延長し、
% P01-E0971: the frozen source refers to an undefined prime p' here; the pair carries q'.
これはさらに、剰余体が $\kappa(\mathfrak p)$ である素イデアル
$\mathfrak q'' \subset S''$ に対応する。標準写像 $S' \to S''$
（例えば最も右の因子を用いる）は、$\varphi$ と $\psi$ を等化する組の射
$(S', \mathfrak q') \to (S'', \mathfrak q'')$ である。
これにより圏論、定義 \ref{categories-definition-directed} の (3) が示された。
従ってこの圏はフィルター的である。

\medskip\noindent
補題 \ref{algebra-lemma-henselization} の証明では、$R_{\mathfrak p}$ とその極大イデアル
$\mathfrak pR_{\mathfrak p}$ から始め、対応する余極限として
$(R_{\mathfrak p})^h$ を構成したことを思い出そう。
ここで $(R, \mathfrak p)$ に対する任意の組 $(S, \mathfrak q)$ から、
$(R_{\mathfrak p}, \mathfrak pR_{\mathfrak p})$ に対する組
$(S_{\mathfrak p}, \mathfrak qS_{\mathfrak p})$ が得られる。
さらに、この状況では
$$
S_{\mathfrak p} = \colim_{f \in R, f \not \in \mathfrak p} S_f.
$$
従って補題の等式を示すには、
$(R_{\mathfrak p}, \mathfrak pR_{\mathfrak p})$ に対する任意の組
$(S_{loc}, \mathfrak q_{loc})$ が、$(R, \mathfrak p)$ 上のある組
$(S, \mathfrak q)$ に対する $(S_{\mathfrak p}, \mathfrak qS_{\mathfrak p})$ の形で
あることを示せば十分である（いくつかの詳細は省略する）。
これは補題 \ref{algebra-lemma-etale} から従う。
\end{proof}

\begin{lemma}
\label{algebra-lemma-henselian-functorial-improve}
$R \to S$ を環写像とする。$\mathfrak q \subset S$ を
$\mathfrak p \subset R$ の上にある素イデアルとする。
$R \to R^h$ と $S \to S^h$ を、それぞれ $R_\mathfrak p$ と
$S_\mathfrak q$ のヘンゼル化とする。補題
\ref{algebra-lemma-henselian-functorial} の局所環写像 $R^h \to S^h$ は、$S^h$ を
$R^h \otimes_R S$ の、$\mathfrak m^h$ と $\mathfrak q$ の上にある一意な
素イデアルにおけるヘンゼル化と同一視する。
\end{lemma}

\begin{proof}
補題 \ref{algebra-lemma-henselization-different} により、$R^h$ と $S^h$ はそれぞれ
エタール $R$-代数とエタール $S$-代数のフィルター余極限である。
従って $R^h \otimes_R S$ はエタール $S$-代数 $A_i$ のフィルター余極限である
（補題 \ref{algebra-lemma-etale}）。補題 \ref{algebra-lemma-colimits-of-etale} により、
$S^h$ はエタール $R^h \otimes_R S$-代数のフィルター余極限である。
さらに $S^h$ は剰余体が $\kappa(\mathfrak q)$ に等しいヘンゼル局所環なので、
主張は補題 \ref{algebra-lemma-uniqueness-henselian} の一意性から従う。
\end{proof}

\begin{lemma}
\label{algebra-lemma-strictly-henselian-functorial-prepare}
$\varphi : R \to S$ を局所環の局所写像とする。
$S/\mathfrak m_S \subset \kappa^{sep}$ を分離代数閉包とする。
$S \to S^{sh}$ を $S/\mathfrak m_S \subset \kappa^{sep}$ に関する $S$ の
強ヘンゼル化とする。$R \to A$ をエタール環写像とし、$\mathfrak q$ を
$\mathfrak m_R$ の上にある $A$ の素イデアルとする。任意の可換図式
$$
\xymatrix{
\kappa(\mathfrak q) \ar[r]_{\phi} & \kappa^{sep} \\
R/\mathfrak m_R \ar[r]^{\varphi} \ar[u] & S/\mathfrak m_S \ar[u]
}
$$
が与えられると、可換図式
$$
\xymatrix{
A \ar[r]_f & S^{sh} \\
R \ar[u] \ar[r]^{\varphi} & S \ar[u]
}
$$
に入る環の射 $f : A \to S^{sh}$ で、
$f^{-1}(\mathfrak m_{S^h}) = \mathfrak q$ を満たし、誘導写像
$\kappa(\mathfrak q) \to \kappa^{sep}$ が与えられた写像となるものが一意に存在する。
% P01-E0972: the frozen statement takes the inverse image of m_{S^h}; the codomain is S^{sh}, so this should be m_{S^{sh}}.
\end{lemma}

\begin{proof}
これは補題 \ref{algebra-lemma-map-into-henselian} の特別な場合である。
\end{proof}

\begin{lemma}
\label{algebra-lemma-strictly-henselian-functorial}
$R \to S$ を局所環の局所写像とする。分離代数閉包
$R/\mathfrak m_R \subset \kappa_1^{sep}$ および
$S/\mathfrak m_S \subset \kappa_2^{sep}$ を選ぶ。
$R \to R^{sh}$ と $S \to S^{sh}$ を対応する強ヘンゼル化とする。
任意の可換図式
$$
\xymatrix{
\kappa_1^{sep} \ar[r]_{\phi} & \kappa_2^{sep} \\
R/\mathfrak m_R \ar[r]^{\varphi} \ar[u] & S/\mathfrak m_S \ar[u]
}
$$
が与えられると、可換図式
% P01-E0973: the frozen introductory clause Given any commutative diagram is followed by capitalized There exists without completion or punctuation.
$$
\xymatrix{
R^{sh} \ar[r]_f & S^{sh} \\
R \ar[u] \ar[r] & S \ar[u]
}
$$
に入る局所環写像 $R^{sh} \to S^{sh}$ で、$R^{sh}$ と $S^{sh}$ の
剰余体上で $\phi$ を誘導するものが一意に存在する。
\end{lemma}

\begin{proof}
補題 \ref{algebra-lemma-map-into-henselian-colimit} から直ちに従う。
\end{proof}

\begin{lemma}
\label{algebra-lemma-strict-henselization-different}
$R$ を環とし、$\mathfrak p \subset R$ を素イデアルとする。
$\kappa(\mathfrak p) \subset \kappa^{sep}$ を分離代数閉包とする。
$R \to S$ がエタール、$\mathfrak q$ が $\mathfrak p$ の上にある素イデアル、
$\phi : \kappa(\mathfrak q) \to \kappa^{sep}$ が
$\kappa(\mathfrak p)$-代数写像となる三つ組 $(S, \mathfrak q, \phi)$ の圏を
考える。この圏はフィルター的であり、標準的に
$$
(R_{\mathfrak p})^{sh} =
\colim_{(S, \mathfrak q, \phi)} S =
\colim_{(S, \mathfrak q, \phi)} S_{\mathfrak q}
$$
である。
\end{lemma}

\begin{proof}
三つ組の射
$(S, \mathfrak q, \phi) \to (S', \mathfrak q', \phi')$ は、
$\varphi^{-1}(\mathfrak q') = \mathfrak q$ および
$\phi' \circ \varphi = \phi$ を満たす $R$-代数写像
$\varphi : S \to S'$ によって与えられる。組の圏がフィルター的であることを
示そう。圏論、定義 \ref{categories-definition-directed} を参照せよ。
% P01-E0974: the frozen proof says category of pairs although the objects and morphisms here are triples.
この圏は三つ組
$(R, \mathfrak p, \kappa(\mathfrak p) \subset \kappa^{sep})$ を含むので空でなく、
これにより圏論、定義 \ref{categories-definition-directed} の (1) が示される。
$(S, \mathfrak q, \phi)$ と $(S', \mathfrak q', \phi')$ を二つの三つ組とする。
$\mathfrak q$ と $\mathfrak q'$ はそれぞれ、剰余体が
$\kappa(\mathfrak p)$ 上有限分離であるファイバー環
$S \otimes \kappa(\mathfrak p)$ と $S' \otimes \kappa(\mathfrak p)$ の
素イデアルに対応し、$\phi$ と $\phi'$ はそれぞれ $\kappa^{sep}$ への写像に
対応することに注意せよ。従って、このデータは $\kappa(\mathfrak p)$-代数写像
$$
\phi : S \otimes_R \kappa(\mathfrak p) \longrightarrow \kappa^{sep},
\quad
\phi' : S' \otimes_R \kappa(\mathfrak p) \longrightarrow \kappa^{sep}.
$$
に対応する。$S'' = S \otimes_R S'$ とおく。上の写像を組み合わせると、
一意な $\kappa(\mathfrak p)$-代数写像
% P01-E0975: the frozen source says Combining the maps the above, with the article and modifier reversed.
$$
\phi'' = \phi \otimes \phi' :
S'' \otimes_R \kappa(\mathfrak p)
\longrightarrow
\kappa^{sep}
$$
を得る。その核は $\mathfrak q$ と $\mathfrak q'$ の上にある素イデアル
$\mathfrak q'' \subset S''$ に対応し、その剰余体は $\phi''$ によって、
$\kappa^{sep}$ における $\phi(\kappa(\mathfrak q))$ と
$\phi'(\kappa(\mathfrak q'))$ の合成体へ写る。
補題 \ref{algebra-lemma-etale} により環写像 $R \to S''$ はエタールである。
従って $(S'', \mathfrak q'', \phi'')$ は
$(S, \mathfrak q, \phi)$ と $(S', \mathfrak q', \phi')$ の双方を支配する
三つ組である。これにより圏論、定義
\ref{categories-definition-directed} の (2) が示された。
次に、$\varphi, \psi : (S, \mathfrak q, \phi) \to
(S', \mathfrak q', \phi')$ を組の二つの射とする。
% P01-E0976: the frozen source again says morphisms of pairs although these are morphisms of triples.
補題 \ref{algebra-lemma-map-between-etale} により、$\varphi$, $\psi$ および
$S' \otimes_R S' \to S'$ はエタール環写像である。次を考える。
$$
S'' = (S' \otimes_{\varphi, S, \psi} S')
\otimes_{S' \otimes_R S'} S'
$$
上と同様の議論（エタール写像の基底変換はエタールであり、エタール写像の
合成はエタールである）により、$S''$ は $R$ 上エタールである。
$\mathfrak p$ 上の $S''$ のファイバー環は
$$
F'' = (F' \otimes_{\varphi, F, \psi} F')
\otimes_{F' \otimes_{\kappa(\mathfrak p)} F'} F'
$$
である。ここで $F', F$ はそれぞれ $S'$ と $S$ のファイバー環である。
$\varphi$ と $\psi$ は三つ組の射なので、写像
$\phi' : F' \to \kappa^{sep}$ は写像
$\phi'' : F'' \to \kappa^{sep}$ へ延長し、これはさらに素イデアル
$\mathfrak q'' \subset S''$ に対応する。標準写像 $S' \to S''$
（例えば最も右の因子を用いる）は、$\varphi$ と $\psi$ を等化する三つ組の射
$(S', \mathfrak q', \phi') \to (S'', \mathfrak q'', \phi'')$ である。
これにより圏論、定義 \ref{categories-definition-directed} の (3) が示された。
従ってこの圏はフィルター的である。

\medskip\noindent
この系の余極限 $R_{colim}$ が、$\kappa^{sep}$ に関する
$R_{\mathfrak p}$ の強ヘンゼル化に等しいことをなお示さなければならない。
これを見るため、三つ組 $(S, \mathfrak q, \phi)$ の系は、補題
\ref{algebra-lemma-henselization-different} の組 $(S, \mathfrak q)$ を部分系として
含むことに注意する。従って、その補題により $R_{colim}$ は
$R_{\mathfrak p}^h$ を含む。さらに、
$R_{\mathfrak p}^h \subset R_{colim}$ がエタール環拡大の有向余極限であることは
明らかである。従って補題 \ref{algebra-lemma-finite-over-henselian} および
\ref{algebra-lemma-colimit-henselian} により $R_{colim}$ はヘンゼル的である。
最後に、補題 \ref{algebra-lemma-make-etale-map-prescribed-residue-field} により、
$R_{colim}$ の剰余体は $\kappa^{sep}$ に等しい。
従って $R_{colim}$ は強ヘンゼル的であり、求める通り
$R_{\mathfrak p}$ の強ヘンゼル化に等しい。いくつかの詳細は省略する。
\end{proof}

\begin{lemma}
\label{algebra-lemma-strictly-henselian-functorial-improve}
$R \to S$ を環写像とする。$\mathfrak q \subset S$ を
$\mathfrak p \subset R$ の上にある素イデアルとする。分離代数閉包
$\kappa(\mathfrak p) \subset \kappa_1^{sep}$ および
$\kappa(\mathfrak q) \subset \kappa_2^{sep}$ を選ぶ。
$R^{sh}$ と $S^{sh}$ を、それぞれ $R_\mathfrak p$ と $S_\mathfrak q$ の
対応する強ヘンゼル化とする。任意の可換図式
$$
\xymatrix{
\kappa_1^{sep} \ar[r]_{\phi} & \kappa_2^{sep} \\
\kappa(\mathfrak p) \ar[r]^{\varphi} \ar[u] & \kappa(\mathfrak q) \ar[u]
}
$$
が与えられると、補題 \ref{algebra-lemma-strictly-henselian-functorial} の局所環写像
$R^{sh} \to S^{sh}$ は、$S^{sh}$ を、$\mathfrak q$ と極大イデアル
$\mathfrak m^{sh} \subset R^{sh}$ の上にある素イデアルにおける
$R^{sh} \otimes_R S$ の強ヘンゼル化と同一視する。
\end{lemma}

\begin{proof}
証明は補題 \ref{algebra-lemma-henselian-functorial-improve} の証明と同一であるが、
補題 \ref{algebra-lemma-strict-henselization-different} を補題
\ref{algebra-lemma-henselization-different} の代わりに用いる。
\end{proof}

\begin{lemma}
\label{algebra-lemma-sh-from-h-map}
$R \to S$ を環写像とする。$\mathfrak q \subset S$ を
$\mathfrak p \subset R$ の上にある素イデアルとし、
$\kappa(\mathfrak p) \to \kappa(\mathfrak q)$ は同型であるとする。
$\kappa(\mathfrak p) = \kappa(\mathfrak q)$ の分離代数閉包
$\kappa^{sep}$ を選ぶ。このとき
$$
(S_\mathfrak q)^{sh} =
(S_\mathfrak q)^h \otimes_{(R_\mathfrak p)^h} (R_\mathfrak p)^{sh}
$$
である。
\end{lemma}

\begin{proof}
これは、注意 \ref{algebra-remark-construct-sh-from-h} における局所環の強ヘンゼル化の
別の構成と、剰余体が等しいという事実から従う。いくつかの詳細は省略する。
\end{proof}

\section{ヘンゼル化と準有限環写像}
\label{algebra-section-henselization-quasi-finite}

\noindent
この節では、準有限環写像に対する（強）ヘンゼル化の間の関手的写像に
関するいくつかの結果を証明する。

\begin{lemma}
\label{algebra-lemma-quasi-finite-henselization}
$R \to S$ を環写像とする。
$\mathfrak q$ を $R$ の $\mathfrak p$ の上にある $S$ の素イデアルとする。
$R \to S$ は $\mathfrak q$ において準有限であると仮定する。補題
\ref{algebra-lemma-henselian-functorial} の可換図式
$$
\xymatrix{
R_{\mathfrak p}^h \ar[r] & S_{\mathfrak q}^h \\
R_{\mathfrak p} \ar[u] \ar[r] & S_{\mathfrak q} \ar[u]
}
$$
は、$S_{\mathfrak q}^h$ を
$R_{\mathfrak p}^h \otimes_{R_{\mathfrak p}} S_{\mathfrak q}$ の
$\mathfrak q$ により生成される素イデアルにおける局所化と同一視する。
さらに、環写像 $R_{\mathfrak p}^h \to S_{\mathfrak q}^h$ は有限である。
\end{lemma}

\begin{proof}
$R_{\mathfrak p}^h \otimes_R S$ は、$\mathfrak q$ に対応する素イデアルにおいて
$R_{\mathfrak p}^h$ 上準有限であることに注意せよ。補題
\ref{algebra-lemma-four-rings} を参照せよ。従って、
$R_{\mathfrak p}^h \otimes_{R_{\mathfrak p}} S_{\mathfrak q}$ の局所化
$S'$ はヘンゼル的で、$R_{\mathfrak p}^h$ 上有限である。補題
\ref{algebra-lemma-finite-over-henselian} を参照せよ。局所化として、$S'$ は
$R_{\mathfrak p}^h \otimes_{R_{\mathfrak p}} S_{\mathfrak q}$ 上の
エタール代数のフィルター余極限である。補題
\ref{algebra-lemma-henselian-functorial-improve} により、$S_\mathfrak q^h$ は
$R_{\mathfrak p}^h \otimes_{R_{\mathfrak p}} S_{\mathfrak q}$ の
ヘンゼル化である。従って補題 \ref{algebra-lemma-uniqueness-henselian} の
一意性の結果により $S' = S_\mathfrak q^h$ である。
\end{proof}

\begin{lemma}
\label{algebra-lemma-quotient-henselization}
\begin{slogan}
ヘンゼル化は商と両立する。
\end{slogan}
$R$ をヘンゼル化 $R^h$ をもつ局所環とする。
$I \subset \mathfrak m_R$ とする。
このとき $R^h/IR^h$ は $R/I$ のヘンゼル化である。
\end{lemma}

\begin{proof}
これは補題 \ref{algebra-lemma-quasi-finite-henselization} の特別な場合である。
\end{proof}

\begin{lemma}
\label{algebra-lemma-quasi-finite-strict-henselization}
$R \to S$ を環写像とする。
$\mathfrak q$ を $R$ の $\mathfrak p$ の上にある $S$ の素イデアルとする。
$R \to S$ は $\mathfrak q$ において準有限であると仮定する。
$\kappa_2^{sep}/\kappa(\mathfrak q)$ を可分代数閉包とし、
$\kappa(\mathfrak p)$ 上可分代数的な元からなる部分体を
$\kappa_1^{sep} \subset \kappa_2^{sep}$ と記す
（体、補題 \ref{fields-lemma-separable-first}）。補題
\ref{algebra-lemma-strictly-henselian-functorial} の可換図式
$$
\xymatrix{
R_{\mathfrak p}^{sh} \ar[r] & S_{\mathfrak q}^{sh} \\
R_{\mathfrak p} \ar[u] \ar[r] & S_{\mathfrak q} \ar[u]
}
$$
は、$S_{\mathfrak q}^{sh}$ を
$R_{\mathfrak p}^{sh} \otimes_{R_{\mathfrak p}} S_{\mathfrak q}$ の、写像
$$
R_{\mathfrak p}^{sh} \otimes_{R_{\mathfrak p}} S_{\mathfrak q}
\longrightarrow
\kappa_1^{sep} \otimes_{\kappa(\mathfrak p)} \kappa(\mathfrak q)
\longrightarrow
\kappa_2^{sep}
$$
の核である素イデアルにおける局所化と同一視する。さらに、環写像
$R_{\mathfrak p}^{sh} \to S_{\mathfrak q}^{sh}$ は局所環の有限局所準同型であり、
その剰余体拡大は、有限かつ純非分離である拡大
$\kappa_2^{sep}/\kappa_1^{sep}$ である。
\end{lemma}

\begin{proof}
$R \to S$ は $\mathfrak q$ において準有限なので、拡大
$\kappa(\mathfrak q)/\kappa(\mathfrak p)$ は有限である。定義
\ref{algebra-definition-quasi-finite} および補題
\ref{algebra-lemma-isolated-point-fibre} を参照せよ。従って
$\kappa_1^{sep}$ は $\kappa(\mathfrak p)$ の可分代数閉包である
（小さな詳細は省略する）。特に補題
\ref{algebra-lemma-strictly-henselian-functorial} は実際に適用できる。
次に、$\kappa_2^{sep}$ における $\kappa(\mathfrak q)$ と
$\kappa_1^{sep}$ の合成体は可分代数閉であり、従って
$\kappa_2^{sep}$ に等しい。これから
$\kappa_2^{sep}/\kappa_1^{sep}$ が有限であると結論する。構成により、拡大
$\kappa_2^{sep}/\kappa_1^{sep}$ は純非分離である。環写像
$R_{\mathfrak p}^{sh} \to S_{\mathfrak q}^{sh}$ は実際に局所的であり、
実際に有限純非分離である剰余体拡大
$\kappa_2^{sep}/\kappa_1^{sep}$ を誘導する。

\medskip\noindent
$R_{\mathfrak p}^{sh} \otimes_R S$ は、補題の主張で与えた素イデアル
$\mathfrak q'$ において $R_{\mathfrak p}^{sh}$ 上準有限であることに注意せよ。
補題 \ref{algebra-lemma-four-rings} を参照せよ。従って、
$R_{\mathfrak p}^{sh} \otimes_{R_{\mathfrak p}} S_{\mathfrak q}$ の
$\mathfrak q'$ における局所化 $S'$ はヘンゼル的で、
$R_{\mathfrak p}^{sh}$ 上有限である。補題
\ref{algebra-lemma-finite-over-henselian} を参照せよ。前段落の議論により写像
$\kappa_1^{sep} \otimes_{\kappa(\mathfrak p)} \kappa(\mathfrak q)
\to \kappa_2^{sep}$ は全射なので、$S'$ の剰余体は
$\kappa_2^{sep}$ であることに注意せよ。さらに、局所化として $S'$ は
$R_{\mathfrak p}^{sh} \otimes_{R_{\mathfrak p}} S_{\mathfrak q}$ 上の
エタール代数のフィルター余極限である。補題
\ref{algebra-lemma-strictly-henselian-functorial-improve} により、
$S_{\mathfrak q}^{sh}$ は
$R_{\mathfrak p}^{sh} \otimes_{R_{\mathfrak p}} S_{\mathfrak q}$ の
$\mathfrak q'$ における強ヘンゼル化である。従って補題
\ref{algebra-lemma-uniqueness-henselian} の一意性の結果により
$S' = S_\mathfrak q^{sh}$ である。
\end{proof}

\begin{lemma}
\label{algebra-lemma-quotient-strict-henselization}
$R$ を強ヘンゼル化 $R^{sh}$ をもつ局所環とする。
$I \subset \mathfrak m_R$ とする。
このとき $R^{sh}/IR^{sh}$ は $R/I$ の強ヘンゼル化である。
\end{lemma}

\begin{proof}
これは補題 \ref{algebra-lemma-quasi-finite-strict-henselization} の特別な場合である。
\end{proof}

\begin{lemma}
\label{algebra-lemma-local-tensor-with-integral}
$A \to B$ と $A \to C$ を局所環の局所準同型とする。
$A \to C$ が整であり、
$\kappa(\mathfrak m_C)/\kappa(\mathfrak m_A)$ または
$\kappa(\mathfrak m_B)/\kappa(\mathfrak m_A)$ のいずれかが純非分離ならば、
$D = B \otimes_A C$ は局所環であり、$B \to D$ と $C \to D$ は局所的である。
\end{lemma}

\begin{proof}
$D$ の任意の極大イデアルは、整環写像 $B \to D$ に対する上昇性により
$B$ の極大イデアルの上にある（補題 \ref{algebra-lemma-integral-going-up}）。いま
$D/\mathfrak m_B D = \kappa(\mathfrak m_B) \otimes_A C =
\kappa(\mathfrak m_B) \otimes_{\kappa(\mathfrak m_A)} C/\mathfrak m_A C$ である。
$C/\mathfrak m_A C$ のスペクトルは一点、すなわち $\mathfrak m_C$ からなる。
従って $D/\mathfrak m_B D$ のスペクトルは
$\kappa(\mathfrak m_B) \otimes_{\kappa(\mathfrak m_A)} \kappa(\mathfrak m_C)$
のスペクトルと同じであり、仮定により
$\kappa(\mathfrak m_C)/\kappa(\mathfrak m_A)$ または
$\kappa(\mathfrak m_B)/\kappa(\mathfrak m_A)$ のいずれかが純非分離なので、
これは一点である。これにより $D$ は局所的であり、環写像
$B \to D$ と $C \to D$ は局所的であることが証明された。
\end{proof}

\begin{lemma}
\label{algebra-lemma-base-change-strict-henselization-quasi-finite}
$A \to B$ と $A \to C$ を環写像とする。$\kappa$ を可分代数閉体とし、
$B \otimes_A C \to \kappa$ を環準同型とする。対応する強ヘンゼル化の写像を
（証明を参照）
$$
\xymatrix{
B^{sh} \ar[r] & (B \otimes_A C)^{sh} \\
A^{sh} \ar[u] \ar[r] & C^{sh} \ar[u]
}
$$
と記す。次のいずれかを仮定する。
\begin{enumerate}
\item $A \to B$ は素イデアル
$\mathfrak p_B = \Ker(B \to \kappa)$ において準有限である、または
\item $B$ は準有限 $A$-代数のフィルター余極限である、または
\item $B_{\mathfrak p_B}$ は $A_{\mathfrak p_A}$ 上の準有限代数の
フィルター余極限である、または
\item $B$ は $A$ 上整である。
\end{enumerate}
このとき $B^{sh} \otimes_{A^{sh}} C^{sh} \to (B \otimes_A C)^{sh}$ は同型である。
\end{lemma}

\begin{proof}
$D = B \otimes_A C$ と書く。
$\mathfrak p_A = \Ker(A \to \kappa)$ と記し、$\mathfrak p_B$、
$\mathfrak p_C$、$\mathfrak p_D$ も同様に記す。
$\kappa_A \subset \kappa$ を、$\kappa$ における
$\kappa(\mathfrak p_A)$ の可分代数閉包と記し、$\kappa_B$、$\kappa_C$、
$\kappa_D$ も同様に記す。$A^{sh}$ を、可分代数閉包
$\kappa_A/\kappa(\mathfrak p_A)$ を用いて構成した
$A_{\mathfrak p_A}$ の強ヘンゼル化と記す。$B^{sh}$、$C^{sh}$、$D^{sh}$
についても同様とする。補題 \ref{algebra-lemma-strictly-henselian-functorial} の関手性から、
補題の可換図式を得る。

\medskip\noindent
可換図式から得られる写像
$$
c : B^{sh} \otimes_{A^{sh}} C^{sh} \to D^{sh} = (B \otimes_A C)^{sh}
$$
を考える。$A \to B$ が
$\mathfrak p_B = \Ker(B \to \kappa)$ において準有限ならば、補題
\ref{algebra-lemma-four-rings} により環写像 $C \to D$ は
$\mathfrak p_D$ において準有限である。従って補題
\ref{algebra-lemma-quasi-finite-strict-henselization}
（および補題 \ref{algebra-lemma-base-change-integral}）により、環写像 $c$ は有限
$C^{sh}$-代数の準同型であり、ある素イデアル $\mathfrak q$ と
$\mathfrak r$ に対して
$$
B^{sh} = (B \otimes_A A^{sh})_{\mathfrak q}
\quad\text{and}\quad
D^{sh} = (D \otimes_C C^{sh})_{\mathfrak r} =
(B \otimes_A C^{sh})_{\mathfrak r}
$$
である。また
$$
B^{sh} \otimes_{A^{sh}} C^{sh} =
(B \otimes_A A^{sh})_{\mathfrak q} \otimes_{A^{sh}} C^{sh} =
\text{a localization of }
B \otimes_A C^{sh}
$$
なので、$c$ の始域と終域はいずれも $B \otimes_A C^{sh}$ の局所化である
（写像と両立する）。従って、
$B^{sh} \otimes_{A^{sh}} C^{sh}$ が局所的であることを示せば十分である
（小さな詳細は省略する）。これは補題
\ref{algebra-lemma-local-tensor-with-integral} と、すでに用いた補題
\ref{algebra-lemma-quasi-finite-strict-henselization} により
$A^{sh} \to B^{sh}$ が純非分離剰余体拡大をもつ有限写像であるという
事実から従う。これで補題の (1) が証明された。

\medskip\noindent
(2) の場合、$B = \colim B_i$ を準有限 $A$-代数のフィルター余極限として書く。
対応して $D_i = B_i \otimes_A C$ として $D = \colim D_i$ を得る。
$B^{sh} = \colim B_i^{sh}$ であることに注意せよ。実際、補題
\ref{algebra-lemma-colimit-henselian} により、環 $\colim B_i^{sh}$ は強ヘンゼル局所環である。
また、補題 \ref{algebra-lemma-colimit-colimit-etale-better} により
$\colim B_i^{sh}$ はエタール $B$-代数のフィルター余極限である。
最後に、$\colim B_i^{sh}$ の剰余体は $\kappa(\mathfrak p_B)$ の
可分代数閉包である（詳細は省略する）。従って
$B^{sh} = \colim B_i^{sh}$ であると結論する。定義
\ref{algebra-definition-henselization} に続く議論を参照せよ。同様に
$D^{sh} = \colim D_i^{sh}$ を得る。このとき (1) の場合から
$$
D^{sh} = \colim D_i^{sh} = \colim B_i^{sh} \otimes_{A^{sh}} C^{sh} =
B^{sh} \otimes_{A^{sh}} C^{sh}
$$
と結論する。これはフィルター余極限がテンソル積と可換するからである。
% P01-E0977: the frozen English has singular filtered colimit with the plural verb commute.

\medskip\noindent
(3) の場合。$A$、$B$、$C$ をそれぞれ $\mathfrak p_A$、
$\mathfrak p_B$、$\mathfrak p_C$ における局所化で置き換えてよい。
従って (3) は (2) から従う。

\medskip\noindent
整環写像は有限環写像のフィルター余極限なので、(4) も (2) から従う。
\end{proof}

\section{Serre の正規性判定法}
\label{algebra-section-serre-criterion}

\noindent
ネーター環に対して次の性質を導入する。

\begin{definition}
\label{algebra-definition-conditions}
$R$ をネーター環とする。
$k \geq 0$ を整数とする。
\begin{enumerate}
\item 高さが $\leq k$ であるすべての素イデアル $\mathfrak p$ に対して
局所環 $R_{\mathfrak p}$ が正則であるとき、$R$ は性質 {\it $(R_k)$} を
もつという。また、$R$ は {\it 余次元 $\leq k$ で正則} であるともいう。
\item すべての素イデアル $\mathfrak p$ に対して、局所環
$R_{\mathfrak p}$ の深さが
$\min\{k, \dim(R_{\mathfrak p})\}$ 以上であるとき、$R$ は性質
{\it $(S_k)$} をもつという。
\item $M$ を有限 $R$-加群とする。すべての素イデアル $\mathfrak p$ に対して、
加群 $M_{\mathfrak p}$ の深さが
$\min\{k, \dim(\text{Supp}(M_{\mathfrak p}))\}$ 以上であるとき、
$M$ は性質 $(S_k)$ をもつという。
\end{enumerate}
\end{definition}

\noindent
任意のネーター環は性質 $(S_0)$ をもち、その上の任意の有限加群も同様である。
零加群の深さを $\infty$ とする規約（第 \ref{algebra-section-depth} 節を参照）および
空集合の次元を $-\infty$ とする規約（位相、第
\ref{topology-section-krull-dimension} 節を参照）により、零加群はすべての
$k$ に対して性質 $(S_k)$ をもつ。

\begin{lemma}
\label{algebra-lemma-criterion-no-embedded-primes}
$R$ をネーター環とする。
$M$ を有限 $R$-加群とする。次は同値である。
\begin{enumerate}
\item $M$ は埋め込まれた随伴素イデアルをもたない。
\item $M$ は性質 $(S_1)$ をもつ。
\end{enumerate}
\end{lemma}

\begin{proof}
$\mathfrak p$ を $M$ の埋め込まれた随伴素イデアルとする。このとき、
$\mathfrak p \supset \mathfrak q$ を満たす別の $M$ の随伴素イデアル
$\mathfrak q$ が存在する。特に、$\mathfrak q$ も台に属するので、これは
$\dim(\text{Supp}(M_{\mathfrak p})) \geq 1$ を意味する。一方、
$\mathfrak pR_{\mathfrak p}$ は $M_{\mathfrak p}$ に随伴する
（補題 \ref{algebra-lemma-associated-primes-localize}）ので、
$\text{depth}(M_{\mathfrak p}) = 0$ である
（補題 \ref{algebra-lemma-ideal-nonzerodivisor} を参照）。言い換えれば、
$(S_1)$ は成り立たない。逆に、$(S_1)$ が成り立たないならば、
$\dim(\text{Supp}(M_{\mathfrak p})) \geq 1$ かつ
$\text{depth}(M_{\mathfrak p}) = 0$ を満たす素イデアル $\mathfrak p$ が存在する。
$\text{depth}(M_{\mathfrak p}) = 0$ なので、二つの補題
\ref{algebra-lemma-associated-primes-localize} および
\ref{algebra-lemma-ideal-nonzerodivisor} により
$\mathfrak p \in \text{Ass}(M)$ である。
$\dim(\text{Supp}(M_{\mathfrak p})) \geq 1$ なので、
$\mathfrak q \subset \mathfrak p$、$\mathfrak q \not = \mathfrak p$ を満たす
素イデアル $\mathfrak q \in \text{Supp}(M)$ が存在する。このような
$\mathfrak q$ を $\text{Supp}(M)$ において極小にとれる。このとき命題
\ref{algebra-proposition-minimal-primes-associated-primes} により
$\mathfrak q \in \text{Ass}(M)$ であり、従って $\mathfrak p$ は
埋め込まれた随伴素イデアルである。
\end{proof}

\begin{lemma}
\label{algebra-lemma-criterion-reduced}
\begin{slogan}
被約であることは R0 と S1 に等しい。
\end{slogan}
$R$ をネーター環とする。次は同値である。
\begin{enumerate}
\item $R$ は被約である。
\item $R$ は性質 $(R_0)$ と $(S_1)$ をもつ。
\end{enumerate}
\end{lemma}

\begin{proof}
$R$ は被約であると仮定する。補題
\ref{algebra-lemma-minimal-prime-reduced-ring} により、$R$ の任意の極小素イデアル
$\mathfrak p$ に対して $R_{\mathfrak p}$ は体である。従って $(R_0)$ が成り立つ。
$\mathfrak p$ を高さ $\geq 1$ の素イデアルとする。このとき
$A = R_{\mathfrak p}$ は次元 $\geq 1$ の被約局所環である。従って、その極大
イデアル $\mathfrak m$ は随伴素イデアルではない。実際、もしそうなら、零化イデアルが
$\mathfrak m$ である $x \in \mathfrak m$ が存在して $x^2 = 0$ となるからである。
従って補題 \ref{algebra-lemma-ass-zero-divisors} により
$A = R_{\mathfrak p}$ の深さは少なくとも一である。これは $(S_1)$ が
成り立つことを示す。

\medskip\noindent
逆に、$R$ が $(R_0)$ と $(S_1)$ を満たすと仮定する。
$\mathfrak p$ が $R$ の極小素イデアルならば、$(R_0)$ により
$R_{\mathfrak p}$ は体であり、従って被約である。
$\mathfrak p$ が極小でなければ、$(S_1)$ により $R_{\mathfrak p}$ の深さが
$\geq 1$ であることが分かり、
$R_{\mathfrak p} \to R_{\mathfrak p}[1/t]$ が単射となるような
$t \in \mathfrak pR_{\mathfrak p}$ が存在すると結論する。
補題 \ref{algebra-lemma-characterize-zero-local} により、$R_\mathfrak p[1/t]$ は
その素イデアルにおける局所化の積に含まれる。これは、$R_{\mathfrak p}$ が、
$\mathfrak p \supset \mathfrak q$ かつ $t \not \in \mathfrak q$ を満たす
素イデアルにおける $R$ の局所化の積の部分環であることを意味する。
これらの素イデアルはより小さい高さをもつので、高さに関する帰納法により
$R$ は被約であると結論する。
\end{proof}

\begin{lemma}[Serre の正規性判定法]
\label{algebra-lemma-criterion-normal}
\begin{reference}
\cite[IV, Theorem 5.8.6]{EGA}
\end{reference}
\begin{slogan}
正規であることは R1 と S2 に等しい。
\end{slogan}
$R$ をネーター環とする。次は同値である。
\begin{enumerate}
\item $R$ は正規環である。
\item $R$ は性質 $(R_1)$ と $(S_2)$ をもつ。
\end{enumerate}
\end{lemma}

\begin{proof}
(1) $\Rightarrow$ (2) の証明。$R$ は正規、すなわち素イデアルにおける
すべての局所化 $R_{\mathfrak p}$ は正規整域であると仮定する。特に補題
\ref{algebra-lemma-criterion-reduced} により、$R$ は $(R_0)$ と $(S_1)$ をもつ。
従って、次元 $d$ の局所ネーター正規整域 $R$ が深さ
$\geq \min(2, d)$ をもち、$d = 1$ ならば正則であることを示せば十分である。
$d = 1$ の場合の主張は補題 \ref{algebra-lemma-characterize-dvr} から従う。

\medskip\noindent
$R$ を、極大イデアル $\mathfrak m$ と次元 $d \geq 2$ をもつ
局所ネーター正規整域とする。補題 \ref{algebra-lemma-hart-serre-loc-thm} を
$R$ に適用する。$R$ がその補題の (1) または (2) の場合に該当しないことは
明らかである。その補題の (4) にあるような $R \to R'$ をとる。
$R$ は整域なので $R \subset R'$ である。$\mathfrak m$ は $R'$ の
随伴素イデアルではないので、$R'$ 上の非零因子である
$x \in \mathfrak m$ が存在する。このとき $R_x = R'_x$ なので、
$R$ と $R'$ は同じ分数体をもつ整域である。しかし $R \subset R'$ の
有限性により、$R'$ のすべての元は $R$ 上整である
（補題 \ref{algebra-lemma-finite-is-integral}）。$R$ は正規なので、
$R = R'$ と結論する。これは (4) が起こらないことを意味する。
従って残る可能性 (3)、すなわち望む通り
$\text{depth}(R) \geq 2$ を得る。

\medskip\noindent
(2) $\Rightarrow$ (1) の証明。$R$ は $(R_1)$ と $(S_2)$ を満たすと仮定する。
補題 \ref{algebra-lemma-criterion-reduced} により $R$ は被約である。従って、
$R$ が次元 $d$ の被約局所ネーター環で $(S_2)$ と $(R_1)$ を満たすならば、
$R$ が正規整域であることを示せば十分である。$d = 0$ なら結果は明らかである。
$d = 1$ なら結果は補題 \ref{algebra-lemma-characterize-dvr} から従う。

\medskip\noindent
$R$ を、極大イデアル $\mathfrak m$ と次元 $d \geq 2$ をもち、
$(R_1)$ と $(S_2)$ を満たす被約局所ネーター環とする。補題
\ref{algebra-lemma-characterize-reduced-ring-normal} により、$R$ がその全商環
$Q(R)$ において整閉であることを示せば十分である。$R$ 上整である
$x \in Q(R)$ をとる。このとき $R' = R[x]$ は $R$ の有限環拡大である
（補題 \ref{algebra-lemma-characterize-finite-in-terms-of-integral}）。
すべての非極大素イデアル $\mathfrak p \subset R$ に対して
$\dim(R_\mathfrak p) < d$ なので、帰納法により
$R_\mathfrak p = R'_\mathfrak p$ である。従って $R'/R$ の台は
$\{\mathfrak m\}$ である。これから $R'/R$ は $\mathfrak m$ のある冪で
零化される（補題 \ref{algebra-lemma-Noetherian-power-ideal-kills-module}）。補題
\ref{algebra-lemma-hart-serre-loc-thm} により、これは $R$ の深さが
$\geq 2 = \min(2, d)$ であるという仮定に反する。これで証明は完了する。
\end{proof}

\begin{lemma}
\label{algebra-lemma-regular-normal}
正則環は正規である。
\end{lemma}

\begin{proof}
$R$ を正則環とする。補題 \ref{algebra-lemma-criterion-normal} により、
$R$ が $(R_1)$ と $(S_2)$ をもつことを証明すれば十分である。
正則局所環は Cohen--Macaulay なので（補題
\ref{algebra-lemma-regular-ring-CM} を参照）、$R$ が $(S_2)$ をもつことは明らかである。
性質 $(R_1)$ は直ちに従う。
\end{proof}

\begin{lemma}
\label{algebra-lemma-normal-domain-intersection-localizations-height-1}
$R$ を分数体 $K$ をもつネーター正規整域とする。このとき
\begin{enumerate}
\item 任意の非零元 $a \in R$ に対して、商 $R/aR$ は埋め込まれた
素イデアルをもたず、そのすべての随伴素イデアルは高さ $1$ をもつ
% P01-E0978: the frozen first item lacks terminal punctuation before the next item.
\item
$$
R = \bigcap\nolimits_{\text{height}(\mathfrak p) = 1} R_{\mathfrak p}
$$
\item 任意の非零元 $x \in K$ に対して、商 $R/(R \cap xR)$ は埋め込まれた
素イデアルをもたず、そのすべての随伴素イデアルは高さ $1$ をもつ。
% P01-E0979: the frozen third item says associates primes rather than associated primes.
\end{enumerate}
\end{lemma}

\begin{proof}
補題 \ref{algebra-lemma-criterion-normal} により、$R$ は $(S_2)$ をもつ。従って、
任意の非零元 $a \in R$ に対して $R/aR$ は $(S_1)$ をもつ
（例えば補題 \ref{algebra-lemma-depth-in-ses} を用いよ）。
% P01-E0980: the frozen parenthetical citation has no terminal punctuation before Hence.
従って $R/aR$ は埋め込まれた素イデアルをもたない
（補題 \ref{algebra-lemma-criterion-no-embedded-primes}）。ゆえに $R/aR$ の随伴
素イデアルは、$(a)$ の上の極小素イデアル $\mathfrak p$ にちょうど等しく、
$a$ は零でないので、それらは高さ $1$ をもつ
（補題 \ref{algebra-lemma-minimal-over-1}）。これで (1) が証明された。

\medskip\noindent
従って、$b \in R$ が与えられたとき、$b \in aR$ であることと、
$(a)$ の上のすべての極小素イデアル $\mathfrak p$ に対して
$b \in aR_{\mathfrak p}$ であることとは同値である
（補題 \ref{algebra-lemma-zero-at-ass-zero} を参照）。上で見た通り、これらの
素イデアルはすべて高さ $1$ をもつので、$b/a \in R$ であることと、
高さ 1 のすべての素イデアルに対して $b/a \in R_{\mathfrak p}$ であることとは
同値である。従って (2) が成り立つ。

\medskip\noindent
(3) について $x = a/b$ と書く。$\mathfrak p_1, \ldots, \mathfrak p_r$ を
$(ab)$ の上の極小素イデアルとする。上により、これらはすべて高さ 1 をもつ。
このとき、補題の (2) により
$R \cap xR = \bigcap_{i = 1, \ldots, r} (R \cap xR_{\mathfrak p_i})$ である。
従って $R/(R \cap xR)$ は
$\bigoplus R/(R \cap xR_{\mathfrak p_i})$ の部分加群である。
$R_{\mathfrak p_i}$ は離散付値環なので（ネーター正規整域 $R$ に対する
性質 $(R_1)$ による。補題 \ref{algebra-lemma-criterion-normal} を参照）、ある
$e_i \in \mathbf{Z}$ に対して
$xR_{\mathfrak p_i} = \mathfrak p_i^{e_i}R_{\mathfrak p_i}$ である。
従って定義 \ref{algebra-definition-symbolic-power} により、直和は
$\bigoplus_{e_i > 0} R/\mathfrak p_i^{(e_i)}$ に等しい。補題
\ref{algebra-lemma-symbolic-power-associated} により、加群
$R/\mathfrak p^{(n)}$ の唯一の随伴素イデアルは $\mathfrak p$ である。
従って $R/(R \cap xR)$ の随伴素イデアルの集合は
$\{\mathfrak p_i\}$ の部分集合であり、それらの間に包含関係はない。
% P01-E0981: the frozen proof says associate primes rather than associated primes.
これで (3) が証明された。
\end{proof}

\section{体の形式的滑らかさ}
\label{algebra-section-p-bases}

\noindent
この節では、体拡大が形式的に滑らかであることと可分であることが同値であると示す。
ただし、まず有限生成体拡大が可分代数的であることと形式的に不分岐であることが
同値であると証明する。

\begin{lemma}
\label{algebra-lemma-characterize-separable-algebraic-field-extensions}
$K/k$ を有限生成体拡大とする。次は同値である。
\begin{enumerate}
\item $K$ は $k$ の有限可分体拡大である。
\item $\Omega_{K/k} = 0$ である。
\item $K$ は $k$ 上形式的に不分岐である。
\item $K$ は $k$ 上不分岐である。
\item $K$ は $k$ 上形式的エタールである。
\item $K$ は $k$ 上エタールである。
\end{enumerate}
\end{lemma}

\begin{proof}
(2) と (3) の同値性は補題
\ref{algebra-lemma-characterize-formally-unramified} である。補題
\ref{algebra-lemma-etale-over-field} により、(1) と (6) は同値である。
性質 (6) は (5) と (4) を含意し、これらはいずれも (3) を含意する
（補題 \ref{algebra-lemma-formally-etale-etale}、\ref{algebra-lemma-unramified}、および
\ref{algebra-lemma-formally-unramified-unramified}）。従って (2) が (1) を含意することを
示せば十分である。$K$ が整域 $A$ の分数体となるような有限生成
$k$-部分代数 $A \subset K$ を選ぶ。$S = A \setminus \{0\}$ とおく。
$0 = \Omega_{K/k} = S^{-1}\Omega_{A/k}$ であり
（補題 \ref{algebra-lemma-differentials-localize}）、かつ
$\Omega_{A/k}$ は有限生成なので
（補題 \ref{algebra-lemma-differentials-finitely-generated}）、$A$ を局所化
$A_f$ で置き換えることにより、$\Omega_{A/k} = 0$ の場合に帰着できる
（詳細は省略する）。このとき $A$ は $k$ 上不分岐であり、従って例えば
$\mathfrak q = (0)$ として補題 \ref{algebra-lemma-unramified-at-prime} を適用すれば、
$K/k$ は有限可分である。
\end{proof}

\begin{lemma}
\label{algebra-lemma-derivative-zero-pth-power}
$k$ を標数 $p > 0$ の完全体とする。$K/k$ を拡大とする。
$a \in K$ とする。このとき、$\Omega_{K/k}$ において
$\text{d}a = 0$ であることと、$a$ が $p$ 乗であることとは同値である。
\end{lemma}

\begin{proof}
補題 \ref{algebra-lemma-colimit-differentials} により、$L/k$ が有限生成体拡大であり、
$\text{d}a$ が $\Omega_{L/k}$ において零となるような部分体
$k \subset L \subset K$ が存在する。従って、$K$ は $k$ の有限生成体拡大であると
仮定してよい。

\medskip\noindent
$K$ が $k(x_1, \ldots, x_r)$ 上有限可分となるような超越基底
$x_1, \ldots, x_r \in K$ を選ぶ。これは定義により可能である。定義
\ref{algebra-definition-perfect} および
\ref{algebra-definition-separable-field-extension} を参照せよ。この結果は純超越部分体
$k(x_1, \ldots, x_r) \subset K$ に対して成り立つことに注意する。実際、
$$
\Omega_{k(x_1, \ldots, x_r)/k} =
\bigoplus\nolimits_{i = 1}^r k(x_1, \ldots, x_r) \text{d}x_i
$$
であり、すべての偏微分が零である任意の有理関数は $p$ 乗である。さらに、
$k(x_1, \ldots, x_r) \subset K$ は有限可分なので（計算は省略する）、
$$
\Omega_{K/k} =
\bigoplus\nolimits_{i = 1}^r K\text{d}x_i
$$
でもある。微分加群において $\text{d}a = 0$ である元 $a \in K$ をとる。
$x_i$ の選び方により、$a$ の最小多項式
$P(T) \in k(x_1, \ldots, x_r)[T]$ は可分である。
$$
P(T) = T^d + \sum\nolimits_{i = 1}^d a_i T^{d - i}
$$
と書けば、$\Omega_{K/k}$ において
$$
0 = \text{d}P(a) = \sum\nolimits_{i = 1}^d a^{d - i}\text{d}a_i
$$
である。上の $\Omega_{K/k}$ の記述と、$P$ が $a$ の最小多項式であることから、
これは $\text{d}a_i = 0$ を含意する。従って各 $i$ に対し
$a_i = b_i^p$ である。それゆえ、体、補題
\ref{fields-lemma-pth-root} により、$a$ は $p$ 乗である。
\end{proof}

\begin{lemma}
\label{algebra-lemma-size-extension-pth-roots}
$k$ を標数 $p > 0$ の体とする。
$a_1, \ldots, a_n \in k$ を、$\Omega_{k/\mathbf{F}_p}$ において
$\text{d}a_1, \ldots, \text{d}a_n$ が線形独立となるような元とする。
このとき、体拡大 $k(a_1^{1/p}, \ldots, a_n^{1/p})$ の $k$ 上の次数は
$p^n$ である。
\end{lemma}

\begin{proof}
$n$ に関する帰納法による。$n = 1$ なら、結果は補題
\ref{algebra-lemma-derivative-zero-pth-power} である。帰納段階として、
$k(a_1^{1/p}, \ldots, a_{n - 1}^{1/p})$ の $k$ 上の次数が
$p^{n - 1}$ であると仮定する。$a_n$ が
$k(a_1^{1/p}, \ldots, a_{n - 1}^{1/p})$ において $p$ 乗に写らないことを
示さなければならない。もし写るなら、
\begin{align*}
a_n & =
\left(\sum\nolimits_{I = (i_1, \ldots, i_{n - 1}),\ 0 \leq i_j \leq p - 1}
\lambda_I a_1^{i_1/p} \ldots a_{n - 1}^{i_{n - 1}/p}\right)^p \\
& = \sum\nolimits_{I = (i_1, \ldots, i_{n - 1}),\ 0 \leq i_j \leq p - 1}
\lambda_I^p a_1^{i_1} \ldots a_{n - 1}^{i_{n - 1}}
\end{align*}
と書ける。$\text{d}$ を適用すると、$\text{d}a_n$ は
$\text{d}a_i$、$i < n$ に線形従属する。これは矛盾である。
\end{proof}

\begin{lemma}
\label{algebra-lemma-separable-differentials}
$k$ を標数 $p > 0$ の体とする。次は同値である。
\begin{enumerate}
\item 体拡大 $K/k$ は可分である
（定義 \ref{algebra-definition-separable-field-extension} を参照）。
\item 写像
$K \otimes_k \Omega_{k/\mathbf{F}_p} \to \Omega_{K/\mathbf{F}_p}$
は単射である。
\end{enumerate}
\end{lemma}

\begin{proof}
$K$ を有限生成体拡大 $K_i/k$ の有向余極限 $K = \colim_i K_i$ として書く。
定義により、$K$ が可分であることと各 $K_i$ が $k$ 上可分であることとは
同値であり、補題 \ref{algebra-lemma-colimit-differentials} により、
$K \otimes_k \Omega_{k/\mathbf{F}_p} \to \Omega_{K/\mathbf{F}_p}$ が
単射であることと、各
$K_i \otimes_k \Omega_{k/\mathbf{F}_p} \to \Omega_{K_i/\mathbf{F}_p}$ が
単射であることとは同値である。従って $K/k$ は有限生成体拡大であると仮定してよい。

\medskip\noindent
$K/k$ は可分な有限生成体拡大であると仮定する。補題
\ref{algebra-lemma-generating-finitely-generated-separable-field-extensions} にあるように
$x_1, \ldots, x_{r + 1} \in K$ を選ぶ。この場合、
$G(x_1, \ldots, x_{r + 1}) = 0$ かつ
$\partial G/\partial X_{r + 1}$ が恒等的には零でないような既約多項式
$G(X_1, \ldots, X_{r + 1}) \in k[X_1, \ldots, X_{r + 1}]$ が存在する。
さらに $K$ はこの整域の分数体である。
$S = K[X_1, \ldots, X_{r + 1}]/(G)$ とおく。
% P01-E0982: the frozen source leaves the domain unnamed and uses K, rather than k, in the displayed presentation of S.
$$
G = \sum a_I X^I, \quad X^I = X_1^{i_1}\ldots X_{r + 1}^{i_{r + 1}}.
$$
と書く。上の $S$ の表示を用いると、
$$
\Omega_{S/\mathbf{F}_p}
=
\frac{
S \otimes_k \Omega_k \oplus
\bigoplus\nolimits_{i = 1, \ldots, r + 1} S\text{d}X_i
}{
\langle
\sum X^I \text{d}a_I + \sum \partial G/\partial X_i \text{d}X_i
\rangle
}
$$
を得る。
% P01-E0983: the frozen differential presentations use undefined Omega_k instead of the relative module over F_p.
$\Omega_{K/\mathbf{F}_p}$ は $S$-加群 $\Omega_{S/\mathbf{F}_p}$ の局所化なので
（補題 \ref{algebra-lemma-differentials-localize} を参照）、
$$
\Omega_{K/\mathbf{F}_p}
=
\frac{
K \otimes_k \Omega_k \oplus
\bigoplus\nolimits_{i = 1, \ldots, r + 1} K\text{d}X_i
}{
\langle
\sum X^I \text{d}a_I + \sum \partial G/\partial X_i \text{d}X_i
\rangle
}
$$
と結論する。いま、多項式 $\partial G/\partial X_{r + 1}$ は恒等的には零でないので、
写像 $K \otimes_k \Omega_{k/\mathbf{F}_p} \to \Omega_{S/\mathbf{F}_p}$ は
望む通り単射である。
% P01-E0984: the frozen conclusion uses Omega_S although the lemma and localized presentation require Omega_K.

\medskip\noindent
$K/k$ は有限生成体拡大で、
$K \otimes_k \Omega_{k/\mathbf{F}_p} \to \Omega_{K/\mathbf{F}_p}$ が
単射であると仮定する。
（証明のこの部分は、補題
\ref{algebra-lemma-characterize-separable-field-extensions} を証明する議論と同じである。）
$K$ の $k$ 上の超越基底 $x_1, \ldots, x_r$ を、有限拡大
$k(x_1, \ldots, x_r) \subset K$ の非分離次数が最小となるように選ぶ。
$K$ が $k(x_1, \ldots, x_r)$ 上可分ならば、結論を得る。
そうでないと仮定して矛盾を導く。このとき、
$k(x_1, \ldots, x_r)$ 上可分でない元 $\alpha \in K$ が存在する。その最小多項式を
$P(T) \in k(x_1, \ldots, x_r)[T]$ とする。$\alpha$ は可分でないので、実際
$P$ は $T^p$ の多項式である。分母を払えば、
$$
G(X_1, \ldots, X_r, T) = \sum a_{I, i} X^I T^i \in k[X_1, \ldots, X_r, T]
$$
という既約多項式を得て、$K$ において
$G(x_1, \ldots, x_r, \alpha) = 0$ となる。これは
$k[X_1, \ldots, X_r, T]/(G) \subset K$ を意味することに注意せよ。
ある組 $(I_0, i_0)$ に対して係数 $a_{I_0, i_0} = 1$ であると仮定してよい。
少なくとも一つの $i$ に対し $\text{d}G/\text{d}X_i$ が恒等的には零でないと
主張する。実際、そうでなければ $G$ は
$X_1^p, \ldots, X_r^p, T^p$ の多項式である。これは
$$
\sum\nolimits_{(I, i) \not = (I_0, i_0)} x^I\alpha^i \text{d}a_{I, i}
$$
が $\Omega_{K/\mathbf{F}_p}$ において零であることを意味する。$K$ の元の集合
$$
\{x^I\alpha^i \mid a_{I, i} \not = 0 \text{ and } (I, i) \not = (I_0, i_0)\}
$$
の間には $k$-線形関係がないことに注意せよ。従って、
$K \otimes_k \Omega_{k/\mathbf{F}_p} \to \Omega_{K/\mathbf{F}_p}$ が単射であるという
仮定は、すべての $(I, i)$ に対し $\Omega_{k/\mathbf{F}_p}$ において
$\text{d}a_{I, i} = 0$ であることを含意する。補題
\ref{algebra-lemma-derivative-zero-pth-power} により、各 $a_{I, i}$ は $p$ 乗である。
これは $G$ が $p$ 乗であることを意味し、$G$ の既約性に反する。
従って、番号を付け直して、$\text{d}G/\text{d}X_1$ は零でないと仮定してよい。
すると $x_1$ は $k(x_2, \ldots, x_r, \alpha)$ 上可分代数的であり、
$x_2, \ldots, x_r, \alpha$ は $K$ の $k$ 上の超越基底であることが分かる。
これは、有限拡大 $k(x_2, \ldots, x_r, \alpha) \subset K$ の非分離次数が、
有限拡大 $k(x_1, \ldots, x_r) \subset K$ の非分離次数より小さいことを意味し、
矛盾である。
\end{proof}

\begin{lemma}
\label{algebra-lemma-formally-smooth-implies-separable}
$K/k$ を体の拡大とする。$K$ が $k$ 上形式的に滑らかならば、
$K$ は $k$ の可分拡大である。
\end{lemma}

\begin{proof}
$K$ は $k$ 上形式的に滑らかであると仮定する。補題
\ref{algebra-lemma-ses-formally-smooth} により、
$K \otimes_k \Omega_{k/\mathbf{Z}} \to \Omega_{K/\mathbf{Z}}$ は単射である。
従って補題 \ref{algebra-lemma-separable-differentials} により、$K$ は $k$ 上可分である。
% P01-E0985: the frozen proof invokes a positive-characteristic criterion without separately disposing of characteristic zero.
\end{proof}

\begin{lemma}
\label{algebra-lemma-characterize-formally-smooth-field-extension}
$K/k$ を体の拡大とする。このとき、$K$ が $k$ 上形式的に滑らかであることと
$H_1(L_{K/k}) = 0$ であることとは同値である。
\end{lemma}

\begin{proof}
これは命題 \ref{algebra-proposition-characterize-formally-smooth} と、ベクトル空間が
自由（従って射影的）であるという事実から従う。
% P01-E0986: the frozen English says a vector spaces is free.
\end{proof}

\begin{lemma}
\label{algebra-lemma-formally-smooth-extensions-easy}
$K/k$ を体の拡大とする。
\begin{enumerate}
\item $K$ が $k$ 上純超越ならば、$K$ は $k$ 上形式的に滑らかである。
\item $K$ が $k$ 上可分代数的ならば、$K$ は $k$ 上形式的に滑らかである。
\item $K$ が $k$ 上可分ならば、$K$ は $k$ 上形式的に滑らかである。
\end{enumerate}
\end{lemma}

\begin{proof}
(1) について $K = k(x_j; j \in J)$ と書く。$A$ を $k$-代数とし、
$I \subset A$ を平方零イデアルとする。
$\varphi : K \to A/I$ を $k$-代数写像とする。
$a_j \mod I = \varphi(x_j)$ となる元 $a_j \in A$ をとる。
このとき、$x_j$ を $a_j$ に写し、$I$ を法として $\varphi$ に帰着する一意な
$k$-代数写像 $K \to A$ が存在することは容易に分かる。従って
$k \subset K$ は形式的に滑らかである。

\medskip\noindent
(2) の場合、$k \subset K$ はエタール環拡大の余極限であることが分かる。
エタール環写像は形式的エタールである
（補題 \ref{algebra-lemma-formally-etale-etale}）。従って、この場合は補題
\ref{algebra-lemma-colimit-formally-etale} と、形式的エタールな環写像は形式的に
滑らかであるという自明な観察から従う。

\medskip\noindent
(3) の場合、$K = \colim K_i$ をその有限生成 $k$-部分拡大のフィルター余極限として
書く。定義 \ref{algebra-definition-separable-field-extension} により、各 $K_i$ は
$k$ の純超越拡大上可分代数的である。従って (1)、(2) および補題
\ref{algebra-lemma-compose-formally-smooth} により、$K_i/k$ は形式的に滑らかである。
% P01-E0987: the frozen English has the malformed compound finitely generated sub k-extensions.
従って補題 \ref{algebra-lemma-characterize-formally-smooth-field-extension} により
$H_1(L_{K_i/k}) = 0$ である。ゆえに補題 \ref{algebra-lemma-colimits-NL} により
$H_1(L_{K/k}) = 0$ である。従って再び補題
\ref{algebra-lemma-characterize-formally-smooth-field-extension} により、$K/k$ は
形式的に滑らかである。
\end{proof}

\begin{lemma}
\label{algebra-lemma-fields-are-formally-smooth}
\begin{slogan}
体拡大では、形式的に滑らかであることは可分であることに等しい。
\end{slogan}
$k$ を体とする。
\begin{enumerate}
\item $k$ の標数が零ならば、$k$ の任意の拡大体は $k$ 上形式的に滑らかである。
\item $k$ の標数が $p > 0$ ならば、$K/k$ が形式的に滑らかであることと、
可分体拡大であることとは同値である。
\end{enumerate}
\end{lemma}

\begin{proof}
補題 \ref{algebra-lemma-formally-smooth-implies-separable} と
\ref{algebra-lemma-formally-smooth-extensions-easy} を組み合わせよ。
\end{proof}

\noindent
ここで、可分体拡大のさまざまな特徴付けをすべてまとめる。

\begin{proposition}
\label{algebra-proposition-characterize-separable-field-extensions}
$K/k$ を体拡大とする。$k$ の標数が零ならば
% P01-E0988: the frozen characteristic-zero clause introduces the list without saying the conditions are equivalent.
\begin{enumerate}
\item $K$ は $k$ 上可分である。
\item $K$ は $k$ 上幾何学的に被約である。
\item $K$ は $k$ 上形式的に滑らかである。
\item $H_1(L_{K/k}) = 0$ である。
\item 写像 $K \otimes_k \Omega_{k/\mathbf{Z}} \to \Omega_{K/\mathbf{Z}}$ は
単射である。
\end{enumerate}
$k$ の標数が $p > 0$ ならば、次は同値である。
\begin{enumerate}
\item $K$ は $k$ 上可分である。
\item 環 $K \otimes_k k^{1/p}$ は被約である。
\item $K$ は $k$ 上幾何学的に被約である。
\item 写像 $K \otimes_k \Omega_{k/\mathbf{F}_p} \to \Omega_{K/\mathbf{F}_p}$ は
単射である。
\item $H_1(L_{K/k}) = 0$ である。
\item $K$ は $k$ 上形式的に滑らかである。
\end{enumerate}
\end{proposition}

\begin{proof}
これは補題 \ref{algebra-lemma-characterize-separable-field-extensions}、
\ref{algebra-lemma-fields-are-formally-smooth}
\ref{algebra-lemma-formally-smooth-implies-separable}、および
\ref{algebra-lemma-separable-differentials} を組み合わせたものである。
% P01-E0989: the frozen citation list lacks punctuation between two adjacent lemma references.
\end{proof}

\noindent
有限生成可分体拡大のさらに別の特徴付けを与える。

\begin{lemma}
\label{algebra-lemma-localization-smooth-separable}
$K/k$ を有限生成体拡大とする。このとき、$K$ が $k$ 上可分であることと、
$K$ が滑らかな $k$-代数の局所化であることとは同値である。
\end{lemma}

\begin{proof}
分数体が $K$ である整域となる有限型 $k$-代数 $R$ を選ぶ。補題
\ref{algebra-lemma-smooth-at-generic-point} によれば、$k \to R$ が
$(0)$ において滑らかであることと $K/k$ が可分であることとは同値である。
これで補題が証明された。
\end{proof}

\begin{lemma}
\label{algebra-lemma-colimit-syntomic}
$K/k$ を体拡大とする。このとき $K$ は $k$ 上の大域的完全交差代数の
フィルター余極限である。$K/k$ が可分ならば、$K$ は $k$ 上の滑らかな代数の
フィルター余極限である。
\end{lemma}

\begin{proof}
$E \subset K$ を有限部分集合とする。$E$ を含み、$k$ 上大域的完全交差
（それぞれ滑らか）である $k$-部分代数 $A \subset K$ が存在することを
示せば十分である。可分／滑らかな場合は補題
\ref{algebra-lemma-localization-smooth-separable} から従う。一般の場合、$L \subset K$ を
$E$ により生成される部分体とする。$k$ 上の超越基底
$x_1, \ldots, x_d \in L$ を選ぶ。拡大 $L/k(x_1, \ldots, x_d)$ は有限である。
$L = k(x_1, \ldots, x_d)[y_1, \ldots, y_r]$ とする。多項式
$P_i \in k(x_1, \ldots, x_d)[Y_1, \ldots, Y_r]$ を、
$P_i = P_i(Y_1, \ldots, Y_i)$ が
$k(x_1, \ldots, x_d)[Y_1, \ldots, Y_{i - 1}]$ 上 $Y_i$ に関してモニックで、
$k(x_1, \ldots, x_d)[y_1, \ldots, y_{i - 1}][Y_i]$ において
$y_i$ の最小多項式に写るよう帰納的に選ぶ。
% P01-E0990: the frozen English says minimum polynomial instead of minimal polynomial.
このとき、$P_1, \ldots, P_r$ は
$k(x_1, \ldots, x_r)[Y_1, \ldots, Y_r]$ における正則列であり、
$L = k(x_1, \ldots, x_r)[Y_1, \ldots, Y_r]/(P_1, \ldots, P_r)$ であることは
明らかである。
% P01-E0991: the frozen two formulas use x_r although the transcendence basis ends at x_d.
$h \in k[x_1, \ldots, x_d]$ を、
$P_i \in k[x_1, \ldots, x_d, 1/h, Y_1, \ldots, Y_r]$ となるような多項式とする。
このとき、$P_1, \ldots, P_r$ は
$k[x_1, \ldots, x_d, 1/h, Y_1, \ldots, Y_r]$ における正則列であり、
$A = k[x_1, \ldots, x_d, 1/h, Y_1, \ldots, Y_r]/(P_1, \ldots, P_r)$ は
大域的完全交差である。$h$ の選び方を調整すれば $E \subset A$ と仮定でき、
結論を得る。
\end{proof}

\section{平坦環写像の構成}
\label{algebra-section-constructing-flat}

\noindent
次の補題はときどき有用である。

\begin{lemma}
\label{algebra-lemma-flat-local-given-residue-field}
$(R, \mathfrak m, k)$ を局所環とする。$K/k$ を体拡大とする。
局所環 $(R', \mathfrak m', k')$ と、$\mathfrak m' = \mathfrak mR'$ を満たす
平坦な局所環写像 $R \to R'$ であって、$k'$ が $k$ の拡大として $K$ と
同型になるものが存在する。
\end{lemma}

\begin{proof}
$k' = k(\alpha)$ が $k$ の単生成拡大であると仮定する。このとき $k'$ は、
補題にあるような平坦局所拡大 $R \subset R'$ の剰余体である。実際、
$\alpha$ が $k$ 上超越的ならば、$R'$ を素イデアル $\mathfrak mR[x]$ における
$R[x]$ の局所化とする。$\alpha$ が最小多項式
$T^d + \sum \overline{\lambda}_iT^{d - i}$ をもつ代数的元ならば、
$R' = R[T]/(T^d + \sum \lambda_i T^{d - i})$ とする。

\medskip\noindent
三つ組 $(k', R \to R', \phi)$ 全体を考える。ここで
$k \subset k' \subset K$ は部分体、$R \to R'$ は補題にあるような局所環写像、
$\phi : R' \to k'$ は $k$-拡大の同型
$R'/\mathfrak mR' \cong k'$ を誘導するものとする。これらは「大きな」圏
$\mathcal{C}$ をなし、射
$(k_1, R_1, \phi_1) \to (k_2, R_2, \phi_2)$ は、図式
$$
\xymatrix{
R_1 \ar[d]_\psi \ar[r]_{\phi_1} & k_1 \ar[r] & K \ar@{=}[d] \\
R_2 \ar[r]^{\phi_2} & k_2 \ar[r] & K
}
$$
が可換となるような環写像 $\psi : R_1 \to R_2$ によって与えられる。
これは $k_1 \subset k_2$ を含意する。

\medskip\noindent
$I$ を有向集合とし、$((R_i, k_i, \phi_i), \psi_{ii'})$ を $I$ 上の系とする。
% P01-E0992: the frozen system and upper-bound tuples reverse the component order used for the category objects.
圏論、第 \ref{categories-section-posets-limits} 節を参照せよ。この場合、
$$
R' = \colim_{i \in I} R_i
$$
を考えられる。これは極大イデアル $\mathfrak mR'$ と剰余体
$k' = \bigcup_{i \in I} k_i$ をもつ局所環である。さらに、環写像
$R \to R'$ は平坦写像の余極限なので平坦である（また、テンソル積は有向余極限と
可換する）。従って $(R', k', \phi')$ はこの系の「上界」であることが分かる。

\medskip\noindent
$\mathcal{C}$ が集合ならば、Zorn の補題のほとんど自明な適用で証明は完了するが、
実際には集合ではない。
% P01-E0993: the frozen counterfactual says was a set rather than were a set.
（実際、現れる局所環の濃度に妥当な上界を見つければ、この方法を機能させられる。）
この問題を避けるため、$K$ 上に整列順序を選ぶ。$x \in K$ に対して、
$K(x)$ を、$\leq x$ である $K$ のすべての元によって生成される $K$ の部分体とする。
$x \in K$ に関する超限再帰により、剰余体拡大 $K(x)/k$ をもつ、補題にあるような
環写像 $R \subset R(x)$ を構成する。さらに構成により、すべての
$y \leq x$ に対して $R(x)$ は $R(y)$ を含む。実際、$x$ が直前元 $x'$ をもつなら、
$K(x) = K(x')[x]$ であり、$R(x') \subset R(x)$ を証明の第一段落で構成した
局所環拡大とすればよい。$x$ が直前元をもたないなら、まず
$R'(x) = \colim_{x' < x} R(x')$ と、証明の第三段落にあるようにおく。
$R'(x)$ の剰余体は $K'(x) = \bigcup_{x' < x} K(x')$ である。
$K(x) = K'(x)[x]$ なので、証明の第一段落の構成を用いて
$R'(x) \subset R(x)$ を作れる。これで補題の証明が完了する。
% P01-E0994: the frozen two field-generation formulas use square brackets although x may be transcendental.
\end{proof}

\begin{lemma}
\label{algebra-lemma-colimit-finite-etale-given-residue-field}
$(R, \mathfrak m, k)$ を局所環とする。$k \subset K$ が可分代数拡大ならば、
有向集合 $I$ と、局所環の有限エタール拡大 $R \subset R_i$、$i \in I$ の系で、
$R' = \colim R_i$ の剰余体が $K$（$k$ の拡大として）となるものが存在する。
\end{lemma}

\begin{proof}
$R \subset R'$ を補題 \ref{algebra-lemma-flat-local-given-residue-field} の証明で構成した
拡大とする。構成により、$A$ を整列集合として
$R' = \colim_{\alpha \in A} R_\alpha$ であり、遷移写像
$R_\alpha \to R_{\alpha + 1}$ は有限エタールで、$\alpha$ が後続元でなければ
$R_\alpha = \colim_{\beta < \alpha} R_\beta$ である。超限帰納法により結果を証明する。

\medskip\noindent
$R_\alpha$ に対して結果が成り立つ、すなわち $R_\alpha = \colim R_i$ であり
$R_i$ は $R$ 上有限エタールであると仮定する。
$R_\alpha \to R_{\alpha + 1}$ は有限エタールなので、ある $i$ と有限エタール拡大
$R_i \to R_{i, 1}$ が存在して、
$R_{\alpha + 1} = R_\alpha \otimes_{R_i} R_{i, 1}$ となる。従って
$R_{\alpha + 1} = \colim_{i' \geq i} R_{i'} \otimes_{R_i} R_{i, 1}$ であり、
結果は $\alpha + 1$ に対して成り立つ。$\alpha$ は後続元でなく、すべての
$\beta < \alpha$ に対して結果が $R_\beta$ について成り立つと仮定する。
任意の有限部分集合 $E \subset R_\alpha$ はある $\beta < \alpha$ に対して
$R_\beta$ に含まれるので、仮定により $E$ は有限エタール部分拡大に含まれることが
分かる。従って結果は $R_\alpha$ に対して成り立つ。
% P01-E0995: the frozen English inserts and after the subordinate Since clause.
\end{proof}

\begin{lemma}
\label{algebra-lemma-finite-free-given-residue-field-extension}
$R$ を環とする。$\mathfrak p \subset R$ を素イデアルとし、
$L/\kappa(\mathfrak p)$ を有限体拡大とする。このとき、
$\mathfrak q = \mathfrak pS$ が素イデアルで、
$\kappa(\mathfrak q)/\kappa(\mathfrak p)$ が与えられた拡大
$L/\kappa(\mathfrak p)$ と同型になるような有限自由環写像 $R \to S$ が存在する。
\end{lemma}

\begin{proof}
$\kappa(\mathfrak p) \subset L$ の次数に関する帰納法による。
% P01-E0996: the frozen English says induction of the degree rather than induction on the degree.
次数が $1$ なら $R = S$ とする。一般に、部分拡大
$\kappa(\mathfrak p) \subset L' \subset L$ が存在するなら、次数に関する帰納法で
結論を得る（まず $L'/\kappa(\mathfrak p)$ に対応する $R \subset S'$ を構成し、
次に $L/L'$ に対応する $S' \subset S$ を構成する）。
% P01-E0997: the frozen condition does not require a proper intermediate field and says then construction.
従って、$L \supset \kappa(\mathfrak p)$ は単一の元 $\alpha \in L$ により
生成されると仮定してよい。$\alpha$ の $\kappa(\mathfrak p)$ 上の最小多項式を
$X^d + \sum_{i < d} a_iX^i$ とする。従って $a_i \in \kappa(\mathfrak p)$ である。
ある $f_i, g \in R$ と $g \not \in \mathfrak p$ に対して、$a_i$ を
$f_i/g$ の像として書ける。$\alpha$ を $g\alpha$ で置き換え（それに応じて
$a_i$ を $g^{d - i}a_i$ で置き換え）れば、$a_i$ はある $f_i \in R$ の像であると
仮定してよい。このとき単に $S = R[x]/(x^d + \sum f_ix^i)$ とすればよい。
\end{proof}

\begin{lemma}
\label{algebra-lemma-cofinal-system-flat}
$A$ を環とする。$\kappa = \max(|A|, \aleph_0)$ とおく。このとき、任意の平坦
$A$-代数 $B$ は、濃度 $|B'| \leq \kappa$ の平坦な $A$-部分代数
$B' \subset B$ のフィルター余極限である。（$B$ が $A$ 上忠実平坦なら、
$B'$ も $A$ 上忠実平坦であることに注意せよ。）
\end{lemma}

\begin{proof}
$B$ の濃度が $\leq \kappa$ なら、これは正しい。$E \subset B$ を
$|E| \leq \kappa$ である $A$-部分代数とする。$E$ が
$|B'| \leq \kappa$ を満たす平坦な $A$-部分代数 $B'$ に含まれることを示す。
次の理由で補題が従う。(a) $B$ の任意の有限部分集合は、濃度が高々
$\kappa$ である $A$-部分代数に含まれる。(b) 濃度が高々 $\kappa$ である
$B$ の任意の二つの $A$-部分代数は、濃度が高々 $\kappa$ である
$A$-部分代数に含まれる。詳細は省略する。

\medskip\noindent
$A$-部分代数の列
$$
E = E_0 \subset E_1 \subset E_2 \subset \ldots
$$
を帰納的に構成する。各々の濃度は $\leq \kappa$ である。証明を完了するため、
$B' = \bigcup E_k$ が $A$ 上平坦であることを示す。

\medskip\noindent
構成は次の通りである。$E_0 = E$ とおく。$k \geq 0$ に対して $E_k$ が
与えられたとき、$A$ の係数をもつ $E_k$ の元の間の関係全体の集合 $S_k$ を考える。
従って、元 $s \in S_k$ は、整数 $n \geq 1$、
$a_1, \ldots, a_n \in A$ および $e_1, \ldots, e_n \in E_k$ で、
$E_k$ において $\sum a_i e_i = 0$ となるものにより与えられる。
$A \to B$ の平坦性から、補題 \ref{algebra-lemma-flat-eq} により、各
$s = (n, a_1, \ldots, a_n, e_1, \ldots, e_n) \in S_k$ に対して
$$
(m_s, b_{s, 1}, \ldots, b_{s, m_s}, a_{s, 11}, \ldots, a_{s, nm_s})
$$
を選べる。ここで $m_s \geq 0$ は整数、$b_{s, j} \in B$、
$a_{s, ij} \in A$ であり、
$$
e_i = \sum\nolimits_j a_{s, ij} b_{s, j}, \forall i,
\quad\text{and}\quad
0 = \sum\nolimits_i a_i a_{s, ij}, \forall j.
$$
これらを選んだとき、$E_{k + 1} \subset B$ を、次の元により生成される
$A$-部分代数とする。
% P01-E0998: the frozen English misspells choices as choicse.
\begin{enumerate}
\item $E_k$。
\item 各 $s \in S_k$ に対する元 $b_{s, 1}, \ldots, b_{s, m_s}$。
\end{enumerate}
いくらかの集合論（省略する）により、$E_{k + 1}$ の濃度は高々
$\kappa$ である（これは、帰納的に $|E_k| \leq \kappa$ が分かっており、
従って $S_k$ の濃度も高々 $\kappa$ であることを用いる）。
% P01-E0999: the frozen English has the nonidiomatic word order has at most cardinality kappa.

\medskip\noindent
$B' = \bigcup E_k$ が $A$ 上平坦であることを示すため、$A$ に係数をもつ
$B'$ における関係 $\sum_{i = 1, \ldots, n} a_i b'_i = 0$ を考える。
$i = 1, \ldots, n$ に対して $b'_i \in E_k$ となるよう十分大きな $k$ を選ぶ。
すると $(n, a_1, \ldots, a_n, b'_1, \ldots, b'_n) \in S_k$ であり、
従ってこの関係は $E_{k + 1}$ において自明であり、まして $B'$ においても
自明である。従って補題 \ref{algebra-lemma-flat-eq} により $A \to B'$ は平坦である。
\end{proof}

\section{Cohen の構造定理}
\label{algebra-section-cohen-structure-theorem}

\noindent
ここでは、可換代数における基本的な概念を導入する。

\begin{definition}
\label{algebra-definition-complete-local-ring}
$(R, \mathfrak m)$ を局所環とする。$R$ の $\mathfrak m$ に関する
完備化への標準写像
$$
R \longrightarrow \lim_n R/\mathfrak m^n
$$
が同型であるとき、$R$ を {\it 完備局所環}という\footnote{ここには
$\bigcap \mathfrak m^n = (0)$ という条件も含まれる。文献によっては、
これを $R$ が完備かつ分離的であると表現することがある。注意すべきことに、
局所環の完備化 $\lim_n R/\mathfrak m^n$ が完備でないこともありうる。
例、補題 \ref{examples-lemma-noncomplete-completion} を参照せよ。
$\mathfrak m$ が有限生成ならばこの現象は起こらない。補題
\ref{algebra-lemma-hathat-finitely-generated} を参照せよ。この場合、完備化は
ネーター環である。補題 \ref{algebra-lemma-completion-Noetherian} を参照せよ。}。
\end{definition}

\noindent
アルティン局所環 $R$ は完備局所環であることに注意しよう。実際、ある
$n > 0$ に対して $\mathfrak m_R^n = 0$ となる。本節では主として
ネーター完備局所環を扱う。

\begin{lemma}
\label{algebra-lemma-quotient-complete-local}
$R$ をネーター完備局所環とする。$R$ の任意の商環もネーター完備局所環
である。有限環写像 $R \to S$ が与えられると、$S$ はネーター完備局所環
の直積である。
% P01-E1000: frozen English uses the doubled construction “Given ..., then ...”.
\end{lemma}

\begin{proof}
補題 \ref{algebra-lemma-Noetherian-permanence} により、環 $S$ はネーター環である。
補題 \ref{algebra-lemma-completion-tensor} により、$R$-加群としての $S$ は完備である。
従って補題 \ref{algebra-lemma-completion-finite-extension} により、$S$ はその各極大
イデアルにおける完備化の直積である。
\end{proof}

\begin{lemma}
\label{algebra-lemma-complete-local-ring-Noetherian}
$(R, \mathfrak m)$ を完備局所環とする。$\mathfrak m$ が有限生成イデアル
ならば、$R$ はネーター環である。
\end{lemma}

\begin{proof}
補題 \ref{algebra-lemma-completion-Noetherian} を参照せよ。
\end{proof}

\begin{definition}
\label{algebra-definition-coefficient-ring}
$(R, \mathfrak m)$ を完備局所環とする。部分環
$\Lambda \subset R$ が次の条件を満たすとき、これを
{\it 係数環}と呼ぶ。
\begin{enumerate}
\item $\Lambda$ は極大イデアル $\Lambda \cap \mathfrak m$ をもつ
完備局所環である。
\item $\Lambda$ の剰余体から $R$ の剰余体への写像は同型である。
\item $\Lambda \cap \mathfrak m = p\Lambda$ である。ここで $p$ は
$R$ の剰余体の標数である。
\end{enumerate}
\end{definition}

\noindent
この定義についていくつか注意を述べよう。議論を次の各場合に分ける。
\begin{enumerate}
\item 局所環 $R$ が体を含む場合。これは
$\mathbf{Q} \subset R$ であるか、または $p$ を
$R/\mathfrak m$ の標数として $pR = 0$ である場合に起こる。
この場合、係数環 $\Lambda$ は $R$ に含まれる体であり、
$R/\mathfrak m$ へ同型に写る。
\item $R/\mathfrak m$ の標数が $p > 0$ であるが、$R$ において
$p$ のどの冪も零でない場合。この場合、$\Lambda$ は $p$ を一意化元とし、
$R/\mathfrak m$ を剰余体とする完備離散付値環である。
\item $R/\mathfrak m$ の標数が $p > 0$ であり、ある $n > 1$ に対して
$R$ において $p^{n - 1} \not = 0$、$p^n = 0$ となる場合。
この場合、$\Lambda$ は極大イデアルが $p$ で生成され、剰余体が
$R/\mathfrak m$ であるアルティン局所環である。
\end{enumerate}
上に現れた、$p$ を一意化元とする完備離散付値環は特別な役割を果たすので、
次のように命名する。

\begin{definition}
\label{algebra-definition-cohen-ring}
{\it Cohen 環}とは、素数 $p$ を一意化元にもつ完備離散付値環である。
\end{definition}

\begin{lemma}
\label{algebra-lemma-cohen-rings-exist}
$p$ を素数とする。$k$ を標数 $p$ の体とする。
$\Lambda/p\Lambda \cong k$ を満たす Cohen 環 $\Lambda$ が存在する。
\end{lemma}

\begin{proof}
まず、$p$ 進整数 $\mathbf{Z}_p$ は $\mathbf{F}_p$ に対する
Cohen 環をなす。$k$ を標数 $p$ の任意の体とする。
$\mathfrak m_R = pR$ かつ $R/pR = k$ を満たす平坦局所環写像
$\mathbf{Z}_p \to R$ をとる。補題
\ref{algebra-lemma-flat-local-given-residue-field} を参照せよ。
補題 \ref{algebra-lemma-completion-Noetherian} により、完備化
$\Lambda = R^\wedge$ はネーター環である。
$\Lambda/p\Lambda = R/pR$ は体なので、$\Lambda$ は極大イデアル
$(p)$ をもつ完備ネーター局所環である
（等式と完備性については補題 \ref{algebra-lemma-hathat-finitely-generated}）。
任意の $n$ に対し、写像
$\mathbf{Z}/p^n\mathbf{Z} \to \Lambda/p^n\Lambda = R/p^nR$
は平坦である（等式には再び補題
\ref{algebra-lemma-hathat-finitely-generated} を用いる）。
従って、例えば補題 \ref{algebra-lemma-flat-module-powers} により
$\mathbf{Z}_p \to \Lambda$ は平坦である。ゆえに $p$ は
$\Lambda$ の非零因子である。従って $\Lambda$ の次元は $\geq 1$
であり（補題 \ref{algebra-lemma-one-equation}）、$\Lambda$ は次元 $1$ の
正則環、すなわち補題 \ref{algebra-lemma-characterize-dvr} により
離散付値環である。従って $\Lambda$ は $k$ に対する Cohen 環である。
\end{proof}

\begin{lemma}
\label{algebra-lemma-cohen-ring-formally-smooth}
$p > 0$ を素数とする。$\Lambda$ を、剰余体の標数が $p$ である
Cohen 環とする。任意の $n \geq 1$ に対して、環写像
$$
\mathbf{Z}/p^n\mathbf{Z} \to \Lambda/p^n\Lambda
$$
は形式的に滑らかである。
\end{lemma}

\begin{proof}
$n = 1$ ならば、これは命題
\ref{algebra-proposition-characterize-separable-field-extensions} から従う。
一般の $n$ については $n$ に関する帰納法を用いる。
すなわち、$\mathbf{Z}/p^n\mathbf{Z} \to \Lambda/p^n\Lambda$ が
形式的に滑らかであると仮定すると、環写像
$\mathbf{Z}/p^{n + 1}\mathbf{Z} \to \Lambda/p^{n + 1}\Lambda$
とイデアル $I = (p^n) \subset \mathbf{Z}/p^{n + 1}\mathbf{Z}$ に
補題 \ref{algebra-lemma-lift-formal-smoothness} を適用できる。
\end{proof}

\begin{theorem}[Cohen の構造定理]
\label{algebra-theorem-cohen-structure-theorem}
$(R, \mathfrak m)$ を完備局所環とする。
\begin{enumerate}
\item $R$ は係数環をもつ（定義
\ref{algebra-definition-coefficient-ring} を参照せよ）。
\item $\mathfrak m$ が有限生成イデアルならば、$R$ は
$$
\Lambda[[x_1, \ldots, x_n]]/I
$$
という形の商環に同型である。ここで $\Lambda$ は体または Cohen 環である。
\end{enumerate}
\end{theorem}

\begin{proof}
まず係数環の存在を証明する。初めに、剰余体 $\kappa$ の標数が零である場合を
扱う。この場合、$n > 0$ に関する帰納法により、標準写像
$R/\mathfrak m^n \to \kappa = R/\mathfrak m$ の切断
$$
\varphi_n : \kappa \longrightarrow R/\mathfrak m^n
$$
が存在することを示す。$n = 1$ の場合は明らかである。$n > 1$ とし、
$\varphi_{n - 1}$ が与えられているとする。体拡大
$\kappa/\mathbf{Q}$ は命題
\ref{algebra-proposition-characterize-separable-field-extensions} により
形式的に滑らかである。従って、次の図式の点線矢印を構成できる。
$$
\xymatrix{
R/\mathfrak m^{n - 1} &
R/\mathfrak m^n \ar[l] \\
\kappa \ar[u]^{\varphi_{n - 1}} \ar@{..>}[ru] & \mathbf{Q} \ar[l] \ar[u]
}
$$
これで帰納段階が証明された。これらの写像を合わせると、
$$
\lim_n\ \varphi_n : \kappa \longrightarrow
R = \lim_n\ R/\mathfrak m^n
$$
が得られ、その像が求める係数環である。

\medskip\noindent
次に、剰余体 $\kappa$ の標数が $p > 0$ である場合に係数環の存在を
証明する。$\kappa = \Lambda/p\Lambda$ を満たす Cohen 環
$\Lambda$ を選ぶ。補題 \ref{algebra-lemma-cohen-rings-exist} を参照せよ。
この場合、$n > 0$ に関する帰納法により、写像
$$
\varphi_n :
\Lambda/p^n\Lambda
\longrightarrow
R/\mathfrak m^n
$$
であって、還元写像 $R/\mathfrak m^n \to \kappa$ との合成が
与えられた同型 $\Lambda/p\Lambda = \kappa$ を与えるものが存在することを
示す。$n = 1$ の場合は明らかである。$n > 1$ とし、
$\varphi_{n - 1}$ が与えられているとする。環写像
$\mathbf{Z}/p^n\mathbf{Z} \to \Lambda/p^n\Lambda$ は補題
\ref{algebra-lemma-cohen-ring-formally-smooth} により形式的に滑らかである。
従って、次の図式の点線矢印を構成できる。
$$
\xymatrix{
R/\mathfrak m^{n - 1} &
R/\mathfrak m^n \ar[l] \\
\Lambda/p^n\Lambda \ar[u]^{\varphi_{n - 1}} \ar@{..>}[ru] &
\mathbf{Z}/p^n\mathbf{Z} \ar[l] \ar[u]
}
$$
これで帰納段階が証明された。これらの写像を合わせると、
$$
\lim_n\ \varphi_n :
\Lambda = \lim_n\ \Lambda/p^n\Lambda
\longrightarrow
R = \lim_n\ R/\mathfrak m^n
$$
が得られ、その像が求める係数環である。

\medskip\noindent
定理の最後の主張も直ちに従う。実際、$y_1, \ldots, y_n$ をイデアル
$\mathfrak m$ の生成元とすれば、今構成した写像 $\Lambda \to R$ を用いて
写像
$$
\Lambda[[x_1, \ldots, x_n]] \longrightarrow R,
\quad x_i \longmapsto y_i.
$$
を得る。両辺は $(x_1, \ldots, x_n)$ 進完備であり、構成によりこの写像は
$(x_1, \ldots, x_n)$ を法として全射であるから、補題
\ref{algebra-lemma-completion-generalities} により全射である。
\end{proof}

\begin{remark}
\label{algebra-remark-Noetherian-complete-local-ring-universally-catenary}
$k$ が体ならば、冪級数環 $k[[X_1, \ldots, X_d]]$ は次元 $d$ の
ネーター完備正則局所環である。$\Lambda$ が Cohen 環ならば、
$\Lambda[[X_1, \ldots, X_d]]$ は次元 $d + 1$ の完備局所ネーター正則環
である。従って Cohen の構造定理から、任意のネーター完備局所環は正則局所環
の商環であることが従う。特に、ネーター完備局所環は普遍鎖状である。
補題 \ref{algebra-lemma-CM-ring-catenary} および補題
\ref{algebra-lemma-regular-ring-CM} を参照せよ。
\end{remark}

\begin{lemma}
\label{algebra-lemma-regular-complete-containing-coefficient-field}
$(R, \mathfrak m)$ をネーター完備局所環とし、$R$ は正則であると仮定する。
\begin{enumerate}
\item $R$ が $\mathbf{F}_p$ または $\mathbf{Q}$ を含むならば、$R$ は
その剰余体上の冪級数環に同型である。
\item $k$ が体であり、$k \to R$ が同型
$k \to R/\mathfrak m$ を誘導する環写像ならば、$R$ は $k$-代数として
$k$ 上の冪級数環に同型である。
\end{enumerate}
\end{lemma}

\begin{proof}
(1) の場合、Cohen の構造定理
（定理 \ref{algebra-theorem-cohen-structure-theorem}）により係数環が存在する。
それは剰余体へ同型に写る体でなければならない。従って (2) を示せば十分である。
(2) の場合、$\mathfrak m/\mathfrak m^2$ の基へ写る
$f_1, \ldots, f_d \in \mathfrak m$ を選び、$x_i$ を $f_i$ へ写す
連続 $k$-代数写像 $k[[x_1, \ldots, x_d]] \to R$ を考える。
始域と終域はともに $(x_1, \ldots, x_d)$ 進完備なので、補題
\ref{algebra-lemma-completion-generalities} によりこの写像は全射である。
一方、これは単射でなければならない。そうでなければ補題
\ref{algebra-lemma-one-equation} により $R$ の次元は $< d$ となるからである。
\end{proof}

\begin{lemma}
\label{algebra-lemma-complete-local-Noetherian-domain-finite-over-regular}
$(R, \mathfrak m)$ をネーター完備局所整域とする。このとき、次の性質をもつ
$R_0 \subset R$ が存在する。
% P01-E1001: frozen English says “there exists a R_0 subset R”.
\begin{enumerate}
\item $R_0$ は正則完備局所環である。
\item $R_0 \subset R$ は有限であり、剰余体上の同型を誘導する。
\item $R_0$ は、$k$ を体として $k[[X_1, \ldots, X_d]]$ に同型であるか、
または $\Lambda$ を Cohen 環として
$\Lambda[[X_1, \ldots, X_d]]$ に同型である。
\end{enumerate}
\end{lemma}

\begin{proof}
$\Lambda$ を $R$ の係数環とする。$R$ は整域なので、$\Lambda$ は体であるか、
または $\Lambda$ は Cohen 環である。

\medskip\noindent
場合 I：$\Lambda = k$ が体であるとする。$d = \dim(R)$ とおく。
定義イデアル $I \subset R$ を生成する
$x_1, \ldots, x_d \in \mathfrak m$ を選ぶ。
（第 \ref{algebra-section-dimension} 節を参照せよ。）
補題 \ref{algebra-lemma-change-ideal-completion} により、$R$ は $I$ 進についても
完備である。$X_i$ を $x_i$ へ写す写像
$R_0 = k[[X_1, \ldots, X_d]] \to R$ を考える。
$R_0$ はイデアル $I_0 = (X_1, \ldots, X_d)$ に関して完備であり、
$R/I_0R \cong R/IR$ は $k = R_0/I_0$ 上有限であることに注意せよ
（$\dim(R/I) = 0$ であるため。第 \ref{algebra-section-dimension} 節を参照せよ）。
従って補題 \ref{algebra-lemma-finite-over-complete-ring} により
$R_0 \to R$ は有限である。$\dim(R) = \dim(R_0)$ なので、
これは $R_0 \to R$ が単射であることを意味する
（補題 \ref{algebra-lemma-integral-dim-up} を参照せよ）。これで補題が証明された。

\medskip\noindent
場合 II：$\Lambda$ が Cohen 環であるとする。$d + 1 = \dim(R)$ とおく。
$p > 0$ を剰余体 $k$ の標数とする。$R$ は整域なので、$p$ は $R$ の
非零因子である。従って $\dim(R/pR) = d$ である。補題
\ref{algebra-lemma-one-equation} を参照せよ。
$R/pR$ において定義イデアルを生成する
$x_1, \ldots, x_d \in R$ を選ぶ。このとき
$I = (p, x_1, \ldots, x_d)$ は $R$ の定義イデアルである。
補題 \ref{algebra-lemma-change-ideal-completion} により、$R$ は $I$ 進についても
完備である。$X_i$ を $x_i$ へ写す写像
$R_0 = \Lambda[[X_1, \ldots, X_d]] \to R$ を考える。
$R_0$ はイデアル $I_0 = (p, X_1, \ldots, X_d)$ に関して完備であり、
$R/I_0R \cong R/IR$ は $k = R_0/I_0$ 上有限であることに注意せよ
（$\dim(R/I) = 0$ であるため。第 \ref{algebra-section-dimension} 節を参照せよ）。
従って補題 \ref{algebra-lemma-finite-over-complete-ring} により
$R_0 \to R$ は有限である。$\dim(R) = \dim(R_0)$ なので、
これは $R_0 \to R$ が単射であることを意味する
（補題 \ref{algebra-lemma-integral-dim-up} を参照せよ）。これで補題が証明された。
\end{proof}

\section{日本環}
\label{algebra-section-japanese}

\noindent
本節では、整閉包の有限性についての議論を始める。

\begin{definition}
\label{algebra-definition-N}
\begin{reference}
\cite[Chapter 0, Definition 23.1.1]{EGA}
\end{reference}
$R$ を分数体 $K$ をもつ整域とする。
\begin{enumerate}
\item $K$ における $R$ の整閉包が有限 $R$-加群であるとき、
$R$ は {\it N-1} であるという。
\item 任意の有限体拡大 $L/K$ に対し、$L$ における $R$ の整閉包が
$R$ 上有限であるとき、$R$ は {\it N-2}、または{\it 日本環}であるという。
\end{enumerate}
\end{definition}

\noindent
これらの概念が主に重要となるのはネーター環の場合であるが、
ここで非ネーター的な例を一つ挙げる。

\begin{example}
\label{algebra-example-Japanese-not-Noetherian}
$k$ を体とする。整域 $R = k[x_1, x_2, x_3, \ldots]$ は N-2 だが、
ネーター環ではない。その理由は次の通りである。$R \subset L$ であり、
体 $L$ が $R$ の分数体の有限拡大であると仮定する。このときある整数
$n$ が存在し、$L$ は有限拡大 $L_0/k(x_1, \ldots, x_n)$ に
（超越的な）元 $x_{n + 1}, x_{n + 2}$ などを添加して得られる。
$S_0$ を $L_0$ における $k[x_1, \ldots, x_n]$ の整閉包とする。
後の命題 \ref{algebra-proposition-ubiquity-nagata} により、$S_0$ は
$k[x_1, \ldots, x_n]$ 上有限である。さらに、$L$ における $R$ の
整閉包は $S = S_0[x_{n + 1}, x_{n + 2}, \ldots]$ であり
（補題 \ref{algebra-lemma-polynomial-domain-normal} を用いる）、従って
$R$ 上有限である。同じ議論は
$R = \mathbf{Z}[x_1, x_2, x_3, \ldots]$ にも成り立つ。
\end{example}

\begin{lemma}
\label{algebra-lemma-localize-N}
$R$ を整域とする。$R$ が N-1 ならば、$R$ の任意の局所化も N-1 である。
N-2 についても同様である。
\end{lemma}

\begin{proof}
これらの主張は、整閉包をとる操作が局所化と可換することから従う。
補題 \ref{algebra-lemma-integral-closure-localize} を参照せよ。
\end{proof}

\begin{lemma}
\label{algebra-lemma-Japanese-local}
$R$ を整域とする。$f_1, \ldots, f_n \in R$ が単位イデアルを生成するとする。
各整域 $R_{f_i}$ が N-1 ならば、$R$ も N-1 である。
N-2 についても同様である。
\end{lemma}

\begin{proof}
$R_{f_i}$ が N-2（または N-1）であると仮定する。
$L$ を $R$ の分数体の有限拡大とする
（N-1 の場合には分数体そのものとする）。$S$ を $L$ における
$R$ の整閉包とする。補題 \ref{algebra-lemma-integral-closure-localize} により、
$S_{f_i}$ は $L$ における $R_{f_i}$ の整閉包である。従って仮定により、
$S_{f_i}$ は $R_{f_i}$ 上有限である。ゆえに補題 \ref{algebra-lemma-cover} により、
$S$ は $R$ 上有限である。
\end{proof}

\begin{lemma}
\label{algebra-lemma-quasi-finite-over-Noetherian-japanese}
$R$ を整域とする。$R \subset S$ を整域の準有限拡大
（例えば有限拡大）とする。$R$ が N-2 かつネーターであると仮定する。
このとき $S$ は N-2 である。
\end{lemma}

\begin{proof}
$L/K$ を誘導される分数体の拡大とする。これは有限体拡大であることに
注意せよ（例えばファイバー $S \otimes_R K$ に補題
\ref{algebra-lemma-isolated-point-fibre} (2) を適用し、準有限環写像の定義を
用いればよい）。
$S'$ を $S$ における $R$ の整閉包とする。このとき $S'$ は
$L$ における $R$ の整閉包に含まれ、後者は仮定により $R$ 上有限である。
$R$ はネーターなので、これは $S'$ が $R$ 上有限であることを意味する。
補題 \ref{algebra-lemma-quasi-finite-open-integral-closure} により、元
$g_1, \ldots, g_n \in S'$ で、$S'_{g_i} \cong S_{g_i}$ を満たし、
かつ $g_1, \ldots, g_n$ が $S$ の単位イデアルを生成するものが存在する。
従って補題 \ref{algebra-lemma-localize-N} と
\ref{algebra-lemma-Japanese-local} により、$S'$ が N-2 であることを示せば十分で
ある。こうして $S$ が $R$ 上有限である場合に帰着した。

\medskip\noindent
$R \subset S$ が補題の仮定を満たし、さらに $S$ は $R$ 上有限であると
仮定する。$M$ を $S$ の分数体の有限体拡大とする。このとき $M$ は
$K$ の有限体拡大でもあるから、$M$ における $R$ の整閉包 $T$ は
$R$ 上有限である。補題 \ref{algebra-lemma-integral-closure-transitive} により、
$T$ は $M$ における $S$ の整閉包でもあり、補題
\ref{algebra-lemma-integral-permanence} によって結論を得る。
\end{proof}

\begin{lemma}
\label{algebra-lemma-Laurent-ring-N-1}
$R$ をネーター整域とする。$R[z, z^{-1}]$ が N-1 ならば、
$R$ も N-1 である。
\end{lemma}

\begin{proof}
$R'$ をその分数体 $K$ における $R$ の整閉包とする。
$S'$ をその分数体における $R[z, z^{-1}]$ の整閉包とする。
明らかに $R' \subset S'$ である。$K[z, z^{-1}]$ は正規整域なので、
$S' \subset K[z, z^{-1}]$ である。
$f_1, \ldots, f_n \in S'$ が $R[z, z^{-1}]$-加群として $S'$ を
生成すると仮定する。$a_{ij} \in K$ として
$f_i = \sum a_{ij}z^j$（有限和）と書く。任意の $x \in R'$ に対し、
$$
x = \sum h_i f_i
$$
と書ける。ここで $h_i \in R[z, z^{-1}]$ である。従って $R'$ は有限
$R$-部分加群 $\sum Ra_{ij} \subset K$ に含まれる。$R$ はネーターなので、
$R'$ は有限 $R$-加群である。
\end{proof}

\begin{lemma}
\label{algebra-lemma-finite-extension-N-2}
$R$ をネーター整域とし、$R \subset S$ を整域の有限拡大とする。
$S$ が N-1 ならば $R$ も N-1 である。
$S$ が N-2 ならば $R$ も N-2 である。
\end{lemma}

\begin{proof}
省略する。（ヒント：拡大体における $R$ の整閉包は、拡大体における
$S$ の整閉包に含まれる。）
\end{proof}

\begin{lemma}
\label{algebra-lemma-Noetherian-normal-domain-finite-separable-extension}
$R$ を分数体 $K$ をもつネーター正規整域とする。
$L/K$ を有限可分体拡大とする。このとき、$L$ における $R$ の整閉包は
$R$ 上有限である。
\end{lemma}

\begin{proof}
トレース対
（Fields、定義 \ref{fields-definition-trace-pairing}）
$$
L \times L \longrightarrow K,
\quad (x, y) \longmapsto \langle x, y\rangle := \text{Trace}_{L/K}(xy).
$$
を考える。$L/K$ は可分なので、この対は非退化である
（Fields、補題 \ref{fields-lemma-separable-trace-pairing}）。
さらに、$x \in L$ が $R$ 上整ならば、
$\text{Trace}_{L/K}(x)$ は $R$ に属する。実際、$K$ 上の $x$ の
最小多項式の係数は $R$ に属し
（補題 \ref{algebra-lemma-minimal-polynomial-normal-domain}）、
$\text{Trace}_{L/K}(x)$ はそれらの係数の一つの整数倍だからである
（Fields、補題
\ref{fields-lemma-trace-and-norm-from-minimal-polynomial}）。
$R$ 上整であり、かつ $L$ の $K$-基をなす
$x_1, \ldots, x_n \in L$ を選ぶ。このとき、整閉包 $S \subset L$ は
$R$-加群
$$
M = \{y \in L \mid \langle x_i, y\rangle \in R, \ i = 1, \ldots, n\}
$$
に含まれる。線形代数により、$R$-加群として
$M \cong R^{\oplus n}$ である。$R$ はネーターなので、
$S \subset R^{\oplus n}$ は有限生成 $R$-加群である。
\end{proof}

\begin{example}
\label{algebra-example-bad-invariants}
環がネーターでない場合、補題
\ref{algebra-lemma-Noetherian-normal-domain-finite-separable-extension} は成り立たない。
例えば $G = \{+1, -1\}$ が
$A = \mathbf{C}[x_1, x_2, x_3, \ldots]$ に作用し、$-1$ が $x_i$ を
$-x_i$ に写す場合を考える。不変式環 $R = A^G$ は、すべての $x_ix_j$
によって生成される $\mathbf{C}$-代数である。従って $R \subset A$ は
有限でない。しかし $R$ は正規整域であり、その分数体
$K = L^G$ は $A$ の分数体 $L$ における $G$-不変部分である。
また明らかに、$A$ は $L$ における $R$ の整閉包である。
% P01-E1002: frozen English lacks a separator after K = L^G.
\end{example}

\noindent
純非分離拡大の場合、次の補題を補題
\ref{algebra-lemma-Noetherian-normal-domain-finite-separable-extension}
の代わりに用いられることがある。

\begin{lemma}
\label{algebra-lemma-Noetherian-normal-domain-insep-extension}
$R$ を、分数体 $K$ の標数が $p > 0$ であるネーター正規整域とする。
$a \in K$ を、$D(a) \not = 0$ を満たす導分 $D : R \to R$ が存在する
ような元とする。このとき、$L = K[x]/(x^p - a)$ における $R$ の
整閉包は $R$ 上有限である。
\end{lemma}

\begin{proof}
ある $f \in R$ に対して $x$ を $fx$ で、$a$ を $f^pa$ で置き換える
ことにより、$a \in R$ と仮定してよい。従って $D(a) \in R$ でもある。
$i \leq p - 1$ に関する帰納法によって、もし
$$
y = a_0 + a_1x + \ldots + a_i x^i,\quad a_j \in K
$$
が $R$ 上整ならば $D(a)^i a_j \in R$ であることを示す。
すると整閉包は、$D(a)^{-p + 1}x^j$、$j = 0, \ldots, p - 1$ を
基底にもつ有限 $R$-加群に含まれる。$R$ はネーターなので、
これで補題が証明される。

\medskip\noindent
$i = 0$ ならば、$y = a_0$ が $R$ 上整であることと $a_0 \in R$
であることは同値であり、主張は正しい。ある $i < p - 1$ に対して
主張が成り立つと仮定し、
$$
y = a_0 + a_1x + \ldots + a_{i + 1} x^{i + 1},\quad a_j \in K
$$
が $R$ 上整であると仮定する。このとき
$$
y^p = a_0^p + a_1^p a + \ldots + a_{i + 1}^pa^{i + 1}
$$
は $R$ の元である（$K$ に属しかつ $R$ 上整だからである）。
$D$ を適用すると、
$$
(a_1^p + 2a_2^p a + \ldots + (i + 1)a_{i + 1}^p a^i)D(a)
$$
が $R$ に属することが分かる。従って
$$
D(a)a_1 + 2D(a) a_2 x + \ldots + (i + 1)D(a) a_{i + 1} x^i
$$
は $R$ 上整である。帰納法により、$j = 1, \ldots, i + 1$ に対して
$D(a)^{i + 1}a_j \in R$ を得る
（ここで $1, \ldots, i + 1$ が可逆であることを用いる）。
さらに $D(a)^{i + 1}a_0$ も $R$ に属する。実際、これは
$D(a)^{i + 1}y$ と $\sum_{j > 0} D(a)^{i + 1}a_jx^j$ の差であり、
両者は $R$ 上整である
（$a \in R$ なので $x$ は $R$ 上整であることを用いた）。
\end{proof}

\begin{lemma}
\label{algebra-lemma-domain-char-zero-N-1-2}
分数体の標数が零であるネーター整域が N-1 であることと、
N-2（すなわち日本環）であることは同値である。
\end{lemma}

\begin{proof}
標数零ではすべての体拡大が可分なので、これは補題
\ref{algebra-lemma-Noetherian-normal-domain-finite-separable-extension}
から明らかである。
\end{proof}

\begin{lemma}
\label{algebra-lemma-domain-char-p-N-1-2}
$R$ を、分数体 $K$ の標数が $p > 0$ であるネーター整域とする。
このとき $R$ が N-2 であることと、任意の有限純非分離拡大 $L/K$ に対し、
$L$ における $R$ の整閉包が $R$ 上有限であることは同値である。
\end{lemma}

\begin{proof}
$K$ の任意の有限純非分離体拡大における $R$ の整閉包が有限であると
仮定する。$L/K$ を任意の有限拡大とする。$L$ における $R$ の整閉包が
$R$ 上有限であることを示さなければならない。$L$ を含む有限正規体拡大
$M/K$ を選ぶ。$R$ はネーターなので、$M$ における $R$ の整閉包が
$R$ 上有限であることを示せば十分である。Fields、補題
\ref{fields-lemma-normal-case} により、$M_{insep}/K$ が純非分離で
$M/M_{insep}$ が可分であるような部分拡大 $M/M_{insep}/K$ が存在する。
仮定により、$M_{insep}$ における $R$ の整閉包 $R'$ は $R$ 上有限である。
補題 \ref{algebra-lemma-Noetherian-normal-domain-finite-separable-extension} により、
$M$ における $R'$ の整閉包 $R''$ は $R'$ 上有限である。
すると補題 \ref{algebra-lemma-finite-transitive} により $R''$ は $R$ 上有限である。
$R''$ は $M$ における $R$ の整閉包でもあるので
（補題 \ref{algebra-lemma-integral-closure-transitive} を参照せよ）、結論を得る。
\end{proof}

\begin{lemma}
\label{algebra-lemma-polynomial-ring-N-2}
$R$ をネーター整域とする。
$R$ が N-1 ならば $R[x]$ は N-1 である。
$R$ が N-2 ならば $R[x]$ は N-2 である。
\end{lemma}

\begin{proof}
$R$ が N-1 であると仮定する。$R'$ を $R$ の整閉包とすると、
これは $R$ 上有限である。従って $R'[x]$ も $R[x]$ 上有限である。
環 $R'[x]$ は正規であるから
（補題 \ref{algebra-lemma-polynomial-domain-normal} を参照せよ）、N-1 である。
これで最初の主張が証明された。

\medskip\noindent
二番目の主張については、補題 \ref{algebra-lemma-finite-extension-N-2} により、
$R'[x]$ が N-2 であることを示せば十分である。言い換えれば、
$R$ は正規 N-2 整域であると仮定してよく、実際そう仮定する。
標数零の場合は補題 \ref{algebra-lemma-domain-char-zero-N-1-2} により終わる。
標数 $p > 0$ の場合、$K$ を $R$ の分数体として、任意の有限純非分離拡大
$L/K(x)$ における $R[x]$ の整閉包が有限であることを示さなければならない。
% P01-E1003: frozen English says “extension of L/K(x)”.
ある有限純非分離体拡大 $L'/K$ と $q = p^e$ が存在して
$L \subset L'(x^{1/q})$ となる。細部は省略する。
$R[x]$ はネーターなので、$L'(x^{1/q})$ における $R[x]$ の整閉包が
$R[x]$ 上有限であることを示せば十分である。この整閉包は
$R'[x^{1/q}]$ に等しい。ここで $R \subset R' \subset L'$ は
$L'$ における $R$ の整閉包を表す。
% P01-E1004: frozen English omits “where” before the R' apposition.
$R$ は N-2 なので、$R'$ は $R$ 上有限であり、従って
$R'[x^{1/q}]$ は $R[x]$ 上有限である。
\end{proof}

\begin{lemma}
\label{algebra-lemma-openness-normal-locus}
$R$ をネーター整域とする。$R_f$ が正規となるような非零元
$f \in R$ が存在すると仮定する。このとき
$$
U = \{\mathfrak p \in \Spec(R) \mid R_{\mathfrak p} \text{ is normal}\}
$$
は $\Spec(R)$ の開集合である。
\end{lemma}

\begin{proof}
標準開集合 $D(f)$ が $U$ に含まれることは明らかである。
Serre の正規性判定法、補題 \ref{algebra-lemma-criterion-normal} によれば、
$\mathfrak p \not \in U$ であることから、ある
$\mathfrak q \subset \mathfrak p$ に対して次のいずれかが成り立つ。
\begin{enumerate}
\item 場合 I：$\text{depth}(R_{\mathfrak q}) < 2$ かつ
$\dim(R_{\mathfrak q}) \geq 2$ である。
\item 場合 II：$R_{\mathfrak q}$ は正則でなく、
$\dim(R_{\mathfrak q}) = 1$ である。
\end{enumerate}
特にこれは $R_{\mathfrak q}$ が正規でないことも意味し、従って
$f \in \mathfrak q$ である。場合 I では
$\text{depth}(R_{\mathfrak q}) =
\text{depth}(R_{\mathfrak q}/fR_{\mathfrak q}) + 1$
である。このような素イデアル $\mathfrak q$ は、ちょうど
$R/fR$ の埋め込まれた随伴素イデアルである。場合 II では
$\mathfrak q$ は高さ 1 の $R/fR$ の随伴素イデアルである。
従って、そのような素イデアル $\mathfrak q$ の有限集合 $E$ が存在し
（補題 \ref{algebra-lemma-finite-ass} を参照せよ）、
$$
\Spec(R) \setminus U
=
\bigcup\nolimits_{\mathfrak q \in E} V(\mathfrak q)
$$
となる。これが示すべきことであった。
\end{proof}

\begin{lemma}
\label{algebra-lemma-characterize-N-1}
$R$ をネーター整域とする。このとき $R$ が N-1 であることと、
次の二条件が成り立つことは同値である。
\begin{enumerate}
\item $R_f$ が正規となるような非零元 $f \in R$ が存在する。
\item 任意の極大イデアル $\mathfrak m \subset R$ に対し、局所環
$R_{\mathfrak m}$ は N-1 である。
\end{enumerate}
\end{lemma}

\begin{proof}
まず $R$ が N-1 であると仮定する。$R'$ を、その分数体 $K$ における
$R$ の整閉包とする。仮定により、$R$-加群として $R'$ を生成する元
$x_1, \ldots, x_n$ を $R'$ の中にとれる。$R' \subset K$ なので、
$f_i x_i \in R$ を満たす非零元 $f_i \in R$ をとれる。このとき
$f = f_1 \ldots f_n$ とおけば $R_f \cong R'_f$ である。
従って $R_f$ は正規であり、(1) を得る。
(2) は補題 \ref{algebra-lemma-localize-N} から従う。

\medskip\noindent
(1) と (2) を仮定する。$K$ を $R$ の分数体とする。
$R \subset R' \subset K$ が、$K$ に含まれる $R$ の有限拡大であると
仮定する。$R_f$ はすでに正規なので $R_f = R'_f$ であることに注意せよ。
従って補題 \ref{algebra-lemma-openness-normal-locus} により、
$R'_{\mathfrak p'}$ が正規でないような素イデアル
$\mathfrak p' \in \Spec(R')$ の集合は $\Spec(R')$ の閉集合である。
$\Spec(R') \to \Spec(R)$ は閉写像なので、この集合の像は
$\Spec(R)$ の閉集合である。このような環 $R'$ に対し、その像を
$Z_{R'} \subset \Spec(R)$ と書く。
% P01-E1005: frozen English has malformed word order in the Z_{R'} definition.

\medskip\noindent
極大イデアル $\mathfrak m \subset R$ を選ぶ。
$R_{\mathfrak m} \subset R_{\mathfrak m}'$ を、$K$ におけるこの局所環の
整閉包とする。仮定により、これは有限環拡大である。補題
\ref{algebra-lemma-integral-closure-localize} により、$R$ 上整な有限個の元
$x_1, \ldots, x_n \in K$ で、$R_{\mathfrak m}'$ が
$R_{\mathfrak m}$ 上 $x_1, \ldots, x_n$ により生成されるようなものを
とれる。$R' = R[x_1, \ldots, x_n] \subset K$ とおく。この選び方から
$\mathfrak m \not \in Z_{R'}$ は明らかである。

\medskip\noindent
$\Spec(R)$ は準コンパクトなので、上の議論から、
$\bigcap Z_{R'_i} = \emptyset$ を満たす有限個の
$R \subset R'_i \subset K$ をとれる。$R'$ を、これらすべてで生成される
$K$ の部分環とする。これは $R$ 上有限である。また
$Z_{R'} = \emptyset$ である。実際、任意の素イデアル $\mathfrak p'$ は、
ある素イデアル $\mathfrak p'_i$ の上にあり、
$(R'_i)_{\mathfrak p'_i}$ は正規である。これは
$R'_{\mathfrak p'} = (R'_i)_{\mathfrak p'_i}$ も正規であることを意味する。
従って $R'$ は正規である。言い換えれば、$R'$ は $K$ における
$R$ の整閉包である。
\end{proof}

\begin{lemma}[Tate]
\label{algebra-lemma-tate-japanese}
\begin{reference}
\cite[Theorem 23.1.3]{EGA}
\end{reference}
$R$ を環とする。$x \in R$ とする。次を仮定する。
\begin{enumerate}
\item $R$ はネーター正規整域である。
\item $R/xR$ は整域かつ N-2 である。
\item $R \cong \lim_n R/x^nR$ は $x$ に関して完備である。
\end{enumerate}
このとき $R$ は N-2 である。
\end{lemma}

\begin{proof}
そうでなければ補題は明らかなので、$x \not = 0$ と仮定してよい。
$K$ を $R$ の分数体とする。$K$ の標数が零ならば、補題は (1) と
補題 \ref{algebra-lemma-domain-char-zero-N-1-2} から従う。従って
$K$ の標数が $p > 0$ であると仮定してよく、補題
\ref{algebra-lemma-domain-char-p-N-1-2} を適用できる。そこで有限純非分離体拡大
$L/K$ が与えられたとき、$L$ における $R$ の整閉包 $S$ が
$R$ 上有限であることを示せばよい。

\medskip\noindent
$L^q \subset K$ を満たす $p$ の冪 $q$ をとる。必要なら $L$ を拡大して、
$y^q = x$ を満たす元 $y \in L$ が存在すると仮定してよい。
$R \to S$ はスペクトルの同相写像を誘導するので
（補題 \ref{algebra-lemma-p-ring-map} を参照せよ）、素イデアル
$\mathfrak p = xR$ の上にある素イデアル $\mathfrak q \subset S$ は
一意である。$y^q = x$ なので、明らかに
$$
\mathfrak q = \{f \in S \mid f^q \in \mathfrak p\} = yS
$$
である。補題 \ref{algebra-lemma-characterize-dvr} により、
$R_{\mathfrak p}$ は離散付値環である。すると Krull--Akizuki
（補題 \ref{algebra-lemma-krull-akizuki}）により $S_{\mathfrak q}$ は
ネーター環である。そこで再び補題 \ref{algebra-lemma-characterize-dvr} により、
$S_{\mathfrak q}$ は離散付値環であると結論できる。補題
\ref{algebra-lemma-finite-extension-residue-fields-dimension-1} により、
$\kappa(\mathfrak q)/\kappa(\mathfrak p)$ は有限体拡大である。
従って仮定 (2) により、$R/xR$ の整閉包
$S' \subset \kappa(\mathfrak q)$ は $R/xR$ 上有限である。
$S/yS \subset S'$ なので、これは $S/yS$ が $R$ 上有限であることを
意味する。$S/y^nS$ は、その部分商が
$y^iS/y^{i + 1}S \cong S/yS$ である有限フィルトレーションをもつ。
従って各 $S/y^nS$ は $R$ 上有限である。特に $S/xS$ は $R$ 上有限で
ある。また $\bigcap x^nS = (0)$ であることは明らかである。
実際、この共通部分の元の $q$ 乗は
$\bigcap x^nR = (0)$ に含まれる
（補題 \ref{algebra-lemma-intersect-powers-ideal-module-zero}）。
従って補題 \ref{algebra-lemma-finite-over-complete-ring} を適用して、
$S$ は $R$ 上有限であると結論できる。
\end{proof}

\begin{lemma}
\label{algebra-lemma-power-series-over-N-2}
$R$ を環とする。$R$ がネーター整域かつ N-2 ならば、
$R[[x]]$ もそうである。
\end{lemma}

\begin{proof}
補題 \ref{algebra-lemma-Noetherian-power-series} により $R[[x]]$ は
ネーター環である。$R' \supset R$ を、その分数体における $R$ の
整閉包とする。$R$ は N-2 なので、これは $R$ 上有限である。
従って $R'[[x]]$ は $R[[x]]$ 上有限である。補題
\ref{algebra-lemma-power-series-over-Noetherian-normal-domain} により、
$R'[[x]]$ は正規整域である。$x \in R'[[x]]$ に補題
\ref{algebra-lemma-tate-japanese} を適用すると、$R'[[x]]$ が N-2 であることが
分かる。すると補題 \ref{algebra-lemma-finite-extension-N-2} により
$R[[x]]$ は N-2 である。
\end{proof}

\section{永田環}
\label{algebra-section-nagata}

\noindent
まず定義を与える。

\begin{definition}
\label{algebra-definition-nagata}
$R$ を環とする。
\begin{enumerate}
\item 任意の有限型環準同型 $R \to S$ に対し、$S$ が整域ならば
$S$ が N-2（すなわち日本環）であるとき、$R$ を
{\it 普遍的日本環}という。
\item $R$ がネーター環であり、任意の素イデアル $\mathfrak p$ に対して
$R/\mathfrak p$ が N-2 であるとき、$R$ を {\it 永田環}という。
\end{enumerate}
\end{definition}

\noindent
ネーターな普遍的日本環が永田環であることは明らかである。
ここでの目標は、永田環が普遍的日本環であることを示すことである。
これは決して自明ではなく、いくらか準備が必要になる。まず有用な補題を示す。

\begin{lemma}
\label{algebra-lemma-nagata-in-reduced-finite-type-finite-integral-closure}
$R$ を永田環とする。
$R \to S$ が本質的有限型で、$S$ が被約であるとする。
このとき、$S$ における $R$ の整閉包 $A$ は $R$ 上有限である。
\end{lemma}

\begin{proof}
$S$ は $R$ 上本質的有限型なのでネーター環であり、有限個の極小素イデアル
$\mathfrak q_1, \ldots, \mathfrak q_m$ をもつ（補題
\ref{algebra-lemma-Noetherian-irreducible-components} を参照）。
$S$ は被約なので $S \subset \prod S_{\mathfrak q_i}$ であり、各
$S_{\mathfrak q_i} = K_i$ は体である（補題
\ref{algebra-lemma-total-ring-fractions-no-embedded-points} および
\ref{algebra-lemma-minimal-prime-reduced-ring} を参照）。
各 $K_i$ における $R$ の整閉包 $A_i'$ が $R$ 上有限であることを
示せば十分である。実際、$R$ はネーター環であり
$A \subset \prod A_i'$ だからである。
$\mathfrak q_i$ に対応する $R$ の素イデアルを
$\mathfrak p_i \subset R$ とする。
$S$ は $R$ 上本質的有限型なので、
$K_i = S_{\mathfrak q_i} = \kappa(\mathfrak q_i)$ は
$\kappa(\mathfrak p_i)$ の有限生成体拡大である。従って、
$K_i$ における $\kappa(\mathfrak p_i)$ の代数閉包 $L_i$ は
$\kappa(\mathfrak p_i)$ 上有限である（Fields, 補題
\ref{fields-lemma-algebraic-closure-in-finitely-generated} を参照）。
$A_i'$ は $L_i$ における $R/\mathfrak p_i$ の整閉包であることが
明らかなので、永田環の定義から結論が従う。
\end{proof}

\begin{lemma}
\label{algebra-lemma-check-universally-japanese}
$R$ を環とする。
$R$ が普遍的日本環であることを確認するには、次を示せば十分である：
$R \to S$ が有限型で、$S$ が整域ならば、$S$ は N-1 である。
\end{lemma}

\begin{proof}
実際、補題の条件を仮定する。
$R \to S$ を有限型環準同型とし、$S$ を整域とする。
$L$ を $S$ の分数体の有限拡大とする。
このとき有限環拡大 $S \subset S' \subset L$ であって、
$L$ が $S'$ の分数体となるものが存在する。
仮定により $S'$ は N-1 であり、従って $L$ における $S'$ の整閉包
$S''$ は $S'$ 上有限である。ゆえに $S''$ は $S$ 上有限であり
（補題 \ref{algebra-lemma-finite-transitive}）、かつ $S''$ は $L$ における $S$ の
整閉包である（補題 \ref{algebra-lemma-integral-closure-transitive}）。
従って $R$ は普遍的日本環である。
\end{proof}

\begin{lemma}
\label{algebra-lemma-universally-japanese}
$R$ が普遍的日本環ならば、$R$ 上本質的有限型な任意の代数も
普遍的日本環である。
\end{lemma}

\begin{proof}
$R$ 上有限型な代数の場合は定義から直ちに従う。一般の場合は補題
\ref{algebra-lemma-localize-N} を適用すれば従う。
\end{proof}

\begin{lemma}
\label{algebra-lemma-quasi-finite-over-nagata}
$R$ を永田環とする。
$R \to S$ が準有限環準同型（例えば有限環準同型）ならば、
$S$ も永田環である。
\end{lemma}

\begin{proof}
まず、$R$ はネーター環であり、準有限環準同型は有限型なので、
$S$ はネーター環である。
$\mathfrak q \subset S$ を素イデアルとし、
$\mathfrak p = R \cap \mathfrak q$ とおく。このとき
$R/\mathfrak p \subset S/\mathfrak q$ は準有限であるから、補題
\ref{algebra-lemma-quasi-finite-over-Noetherian-japanese} により
$S/\mathfrak q$ は N-2 であり、結論を得る。
\end{proof}

\begin{lemma}
\label{algebra-lemma-nagata-localize}
永田環の局所化は永田環である。
\end{lemma}

\begin{proof}
補題 \ref{algebra-lemma-localize-N} から明らかである。
\end{proof}

\begin{lemma}
\label{algebra-lemma-nagata-local}
$R$ を環とする。$f_1, \ldots, f_n \in R$ が単位イデアルを生成するとする。
\begin{enumerate}
\item 各 $R_{f_i}$ が普遍的日本環ならば、$R$ も普遍的日本環である。
\item 各 $R_{f_i}$ が永田環ならば、$R$ も永田環である。
\end{enumerate}
\end{lemma}

\begin{proof}
$\varphi : R \to S$ を有限型環準同型とし、$S$ を整域とする。
このとき $\varphi(f_1), \ldots, \varphi(f_n)$ は $S$ の単位イデアルを
生成する。従って各 $S_{f_i} = S_{\varphi(f_i)}$ が N-1 ならば、
$S$ も N-1 である（補題 \ref{algebra-lemma-Japanese-local} を参照）。
これで (1) が示された。

\medskip\noindent
各 $R_{f_i}$ が永田環ならば、各 $R_{f_i}$ はネーター環であり、従って
$R$ はネーター環である（補題 \ref{algebra-lemma-cover} を参照）。また、
$\mathfrak p \subset R$ が素イデアルならば、各
$R_{f_i}/\mathfrak pR_{f_i} = (R/\mathfrak p)_{f_i}$ は N-2 である。
従って補題 \ref{algebra-lemma-Japanese-local} により $R/\mathfrak p$ は N-2
である。これで (2) が示された。
\end{proof}

\begin{lemma}
\label{algebra-lemma-Noetherian-complete-local-Nagata}
ネーター完備局所環は永田環である。
\end{lemma}

\begin{proof}
$R$ をネーター完備局所環とする。
$\mathfrak p \subset R$ を素イデアルとする。
このとき $R/\mathfrak p$ もネーター完備局所環である（補題
\ref{algebra-lemma-quotient-complete-local} を参照）。
従って、完備局所ネーター整域 $R$ が N-2 であることを示せば十分である。
補題 \ref{algebra-lemma-quasi-finite-over-Noetherian-japanese} および
\ref{algebra-lemma-complete-local-Noetherian-domain-finite-over-regular} により、
$k$ が体である $R = k[[X_1, \ldots, X_d]]$ の場合、または
$\Lambda$ が Cohen 環である $R = \Lambda[[X_1, \ldots, X_d]]$ の場合に
帰着する。

\medskip\noindent
$k[[X_1, \ldots, X_d]]$ の場合、補題
\ref{algebra-lemma-power-series-over-N-2} により、体が N-2 であるという主張に
帰着し、これは明らかである。
$\Lambda[[X_1, \ldots, X_d]]$ の場合には、Cohen 環 $\Lambda$ が
N-2 であるという主張に帰着する。$x = p \in \Lambda$ として補題
\ref{algebra-lemma-tate-japanese} をもう一度適用すれば、さらに体の場合へ
帰着する。従って結論を得る。
\end{proof}

\begin{definition}
\label{algebra-definition-analytically-unramified}
$(R, \mathfrak m)$ をネーター局所環とする。
その完備化 $R^\wedge = \lim_n R/\mathfrak m^n$ が被約であるとき、
$R$ は {\it 解析的不分岐}であるという。
素イデアル $\mathfrak p \subset R$ について、$R/\mathfrak p$ が
解析的不分岐であるとき、この素イデアルは {\it 解析的不分岐}であるという。
\end{definition}

\noindent
ここまでに、任意のネーター局所環 $R$ について次が成り立つことが
分かっている。写像 $R \to R^\wedge$ は忠実平坦な局所環準同型である
（補題 \ref{algebra-lemma-completion-faithfully-flat}）。完備化 $R^\wedge$ は
ネーター環であり（補題 \ref{algebra-lemma-completion-Noetherian}）、かつ完備である
（補題 \ref{algebra-lemma-completion-complete}）。従って完備化 $R^\wedge$ は
永田環である（補題 \ref{algebra-lemma-Noetherian-complete-local-Nagata}）。
さらに、節 \ref{algebra-section-cohen-structure-theorem} で見たように、
$R^\wedge$ は正則局所環の商であり
（定理 \ref{algebra-theorem-cohen-structure-theorem}）、従って普遍鎖状である
（注意 \ref{algebra-remark-Noetherian-complete-local-ring-universally-catenary}）。

\begin{lemma}
\label{algebra-lemma-analytically-unramified-easy}
$(R, \mathfrak m)$ をネーター局所環とする。
\begin{enumerate}
\item $R$ が解析的不分岐ならば、$R$ は被約である。
\item $R$ が解析的不分岐ならば、$R$ の各極小素イデアルは
解析的不分岐である。
\item $R$ が被約で極小素イデアル $\mathfrak q_1, \ldots, \mathfrak q_t$
をもち、各 $\mathfrak q_i$ が解析的不分岐ならば、$R$ は
解析的不分岐である。
\item $R$ が解析的不分岐ならば、全商環 $Q(R)$ における $R$ の
整閉包は $R$ 上有限である。
\item $R$ が整域かつ解析的不分岐ならば、$R$ は N-1 である。
\end{enumerate}
\end{lemma}

\begin{proof}
この証明では、定義 \ref{algebra-definition-analytically-unramified} の直後の
注意を用いる。$R \to R^\wedge$ は忠実平坦な局所環準同型なので
単射であり、(1) が従う。

\medskip\noindent
$\mathfrak q$ を $R$ の極小素イデアルとし、$R$ が解析的不分岐であると
仮定する。このとき $\mathfrak q$ は $R$ の随伴素イデアルである
（命題 \ref{algebra-proposition-minimal-primes-associated-primes} を参照）。
従って、$\{x \in R \mid fx = 0\} = \mathfrak q$ を満たす
$f \in R$ が存在する。完備化は完全なので（補題
\ref{algebra-lemma-completion-flat}）、
$(R/\mathfrak q)^\wedge = R^\wedge/\mathfrak q^\wedge$ であり、かつ
$\{x \in R^\wedge \mid fx = 0\} = \mathfrak q^\wedge$ であることに
注意する。$x \in R^\wedge$ が $x^2 \in \mathfrak q^\wedge$ を満たせば、
$fx^2 = 0$、従って $(fx)^2 = 0$、従って $fx = 0$、従って
$x \in \mathfrak q^\wedge$ である。ゆえに $\mathfrak q$ は
解析的不分岐であり、(2) が成り立つ。

\medskip\noindent
$R$ が被約で極小素イデアル $\mathfrak q_1, \ldots, \mathfrak q_t$ をもち、
各 $\mathfrak q_i$ が解析的不分岐であると仮定する。このとき
$R \to R/\mathfrak q_1 \times \ldots \times R/\mathfrak q_t$ は単射である。
完備化は完全なので（補題 \ref{algebra-lemma-completion-flat} を参照）、
$R^\wedge \subset (R/\mathfrak q_1)^\wedge \times \ldots \times
(R/\mathfrak q_t)^\wedge$ である。従って (3) は明らかである。

\medskip\noindent
$R$ が解析的不分岐であると仮定する。
$\mathfrak p_1, \ldots, \mathfrak p_s$ を $R^\wedge$ の極小素イデアルとする。
このとき
$$
Q(R^\wedge) =
R^\wedge_{\mathfrak p_1} \times \ldots \times R^\wedge_{\mathfrak p_s}
$$
であり、$R^\wedge$ は被約なので、各 $R^\wedge_{\mathfrak p_i}$ は
体である（補題 \ref{algebra-lemma-total-ring-fractions-no-embedded-points} を参照）。
従って、$Q(R^\wedge)$ における $R^\wedge$ の整閉包 $S$ は
$S = S_1 \times \ldots \times S_s$ に等しい。ここで $S_i$ は
$R^\wedge/\mathfrak p_i$ の分数体におけるその整閉包である。
補題 \ref{algebra-lemma-Noetherian-complete-local-Nagata} により、環 $S_i$ は
$R^\wedge/\mathfrak p_i$ 上有限である。従って $S$ は $R^\wedge$ 上有限である。
% P01-E1006: frozen English omits “by” in the designation of R-prime.
$Q(R)$ における $R$ の整閉包を $R'$ と書く。
$R \to R^\wedge$ は平坦なので、
$R' \otimes_R R^\wedge \subset Q(R) \otimes_R R^\wedge \subset Q(R^\wedge)$
である。さらに $R' \otimes_R R^\wedge$ は $R^\wedge$ 上整である
（補題 \ref{algebra-lemma-base-change-integral}）。
% P01-E1007: frozen English uses the article “a” before R-wedge-submodule.
従って $R' \otimes_R R^\wedge \subset S$ は $R^\wedge$-部分加群である。
$R^\wedge$ はネーター環なので、これは有限 $R^\wedge$-加群である。
従って、$R' \otimes_R R^\wedge$ が $R^\wedge$-加群として元
$f_i \otimes 1$ で生成されるような $f_1, \ldots, f_n \in R'$ をとれる。
忠実平坦性により、$R'$ は $R$-加群として $f_1, \ldots, f_n$ で
生成される。これで (4) が示された。

\medskip\noindent
(5) は (4) の特別な場合である。
\end{proof}

\begin{lemma}
\label{algebra-lemma-codimension-1-analytically-unramified}
$R$ をネーター局所環とする。
$\mathfrak p \subset R$ を素イデアルとする。次を仮定する。
\begin{enumerate}
\item $R_{\mathfrak p}$ は離散付値環であり、
\item $\mathfrak p$ は解析的不分岐である。
\end{enumerate}
このとき、$R^\wedge/\mathfrak pR^\wedge$ の任意の随伴素イデアル
$\mathfrak q$ に対して、局所環 $(R^\wedge)_{\mathfrak q}$ は
離散付値環である。
\end{lemma}

\begin{proof}
仮定 (2) は、$R^\wedge/\mathfrak pR^\wedge$ が被約環であることを意味する。
従って、$R^\wedge/\mathfrak pR^\wedge$ の随伴素イデアル
$\mathfrak q \subset R^\wedge$ は、$\mathfrak pR^\wedge$ 上の
極小素イデアルと同じものである。特に、$(R^\wedge)_{\mathfrak q}$ の
極大イデアルは $\mathfrak p(R^\wedge)_{\mathfrak q}$ である。
$xR_{\mathfrak p} = \mathfrak pR_{\mathfrak p}$ を満たす $x \in R$ をとる。
上で見たことから、$x \in (R^\wedge)_{\mathfrak q}$ は極大イデアルを
生成する。$R \to R^\wedge$ は忠実平坦なので、$x$ は
$(R^\wedge)_{\mathfrak q}$ の非零因子である。従って結論を得る。
\end{proof}

\begin{lemma}
\label{algebra-lemma-criterion-analytically-unramified}
$(R, \mathfrak m)$ をネーター局所整域とする。
$x \in \mathfrak m$ とする。次を仮定する。
\begin{enumerate}
\item $x \not = 0$、
\item $R/xR$ は埋没素イデアルをもたず、
\item $R/xR$ の各随伴素イデアル $\mathfrak p \subset R$ に対して、
\begin{enumerate}
\item 局所環 $R_{\mathfrak p}$ は正則であり、
\item $\mathfrak p$ は解析的不分岐である。
\end{enumerate}
\end{enumerate}
このとき $R$ は解析的不分岐である。
\end{lemma}

\begin{proof}
$\mathfrak p_1, \ldots, \mathfrak p_t$ を $R$-加群 $R/xR$ の
随伴素イデアルとする。$R/xR$ は埋没素イデアルをもたないので、
各 $\mathfrak p_i$ は高さ $1$ であり、$(x)$ 上の極小素イデアルである。
各 $i$ に対し、$\mathfrak q_{i1}, \ldots, \mathfrak q_{is_i}$ を
$R^\wedge$-加群 $R^\wedge/\mathfrak p_iR^\wedge$ の随伴素イデアルとする。
補題 \ref{algebra-lemma-codimension-1-analytically-unramified} により、
$(R^\wedge)_{\mathfrak q_{ij}}$ は正則である。
補題 \ref{algebra-lemma-bourbaki} により、
$$
\text{Ass}_{R^\wedge}(R^\wedge/xR^\wedge)
=
\bigcup\nolimits_{\mathfrak p \in \text{Ass}_R(R/xR)}
\text{Ass}_{R^\wedge}(R^\wedge/\mathfrak pR^\wedge)
=
\{\mathfrak q_{ij}\}.
$$
$y^2 = 0$ を満たす $y \in R^\wedge$ をとる。
$(R^\wedge)_{\mathfrak q_{ij}}$ は正則であり、従って整域なので
（補題 \ref{algebra-lemma-regular-domain}）、$y$ は
$(R^\wedge)_{\mathfrak q_{ij}}$ において零に写る。
従って補題 \ref{algebra-lemma-zero-at-ass-zero} により、$y$ は
$R^\wedge/xR^\wedge$ において零に写る。従って $y = xy'$ である。
$x$ は非零因子なので（$R \to R^\wedge$ が平坦だからである）、
$(y')^2 = 0$ である。ゆえに
$y \in \bigcap x^nR^\wedge = (0)$
である（補題 \ref{algebra-lemma-intersect-powers-ideal-module-zero}）。
\end{proof}

\begin{lemma}
\label{algebra-lemma-local-nagata-domain-analytically-unramified}
$(R, \mathfrak m)$ を局所環とする。
$R$ がネーター整域かつ永田環ならば、$R$ は解析的不分岐である。
\end{lemma}

\begin{proof}
$\dim(R)$ に関する帰納法による。
$\dim(R) = 0$ の場合は自明である。従って $\dim(R) = d$ とし、
次元 $< d$ のすべてのネーター永田整域について補題が成り立つと仮定する。

\medskip\noindent
$R \subset S$ を $R$ の分数体における $R$ の整閉包とする。
仮定により $S$ は有限 $R$-加群である。
補題 \ref{algebra-lemma-quasi-finite-over-nagata} により $S$ は永田環である。
補題 \ref{algebra-lemma-integral-sub-dim-equal} により
$\dim(R) = \dim(S)$ である。
$\mathfrak m_1, \ldots, \mathfrak m_t$ を $S$ の極大イデアルとする。
これらはいずれも $R$ の極大イデアル $\mathfrak m$ の上にある。
さらに $S/\mathfrak mS$ はアルティン環なので、十分大きい $n$ に対し
$$
(\mathfrak m_1 \cap \ldots \cap \mathfrak m_t)^n \subset \mathfrak mS
$$
である。補題 \ref{algebra-lemma-completion-flat} により
$R^\wedge \to S^\wedge$ は単射であり、中国剰余定理
\ref{algebra-lemma-chinese-remainder} と補題 \ref{algebra-lemma-change-ideal-completion} を
合わせると $S^\wedge = \prod S^\wedge_i$ を得る。ここで
$S^\wedge_i$ は極大イデアル $\mathfrak m_i$ に関する $S$ の完備化である。
従って、$S_{\mathfrak m_i}$ が解析的不分岐であることを示せば十分である。
言い換えれば、$R$ がネーター正規永田整域である場合に帰着した。

\medskip\noindent
$R$ をネーター正規局所永田整域と仮定する。
極大イデアルの中から非零元 $x \in \mathfrak m$ をとる。
補題 \ref{algebra-lemma-criterion-analytically-unramified} を適用する。
性質 (1)、(2)、(3)(a)、(3)(b) を確認しなければならない。
性質 (1) は明らかである。
補題 \ref{algebra-lemma-normal-domain-intersection-localizations-height-1} により、
$R/xR$ は埋没素イデアルをもたない。従って性質 (2) が成り立つ。
同じ補題により、$R/xR$ の各随伴素イデアル $\mathfrak p$ は高さ $1$ である。
従って $R_{\mathfrak p}$ は $1$ 次元正規整域であり、ゆえに正則である
（補題 \ref{algebra-lemma-characterize-dvr}）。これで (3)(a) が成り立つ。
最後に、$R/\mathfrak p$ は永田環なので（補題
\ref{algebra-lemma-quasi-finite-over-nagata}、または定義から直接）、帰納法の仮定に
より (3)(b) が成り立つ。従って $R$ は解析的不分岐である。
\end{proof}

\begin{lemma}
\label{algebra-lemma-local-nagata-and-analytically-unramified}
% P01-E1008: frozen English omits the colon introducing this equivalence list.
$(R, \mathfrak m)$ をネーター局所環とする。次は同値である：
\begin{enumerate}
\item $R$ は永田環である。
\item 有限準同型 $R \to S$ で $S$ が整域であり、
$\mathfrak m' \subset S$ が極大である任意の場合に、局所環
$S_{\mathfrak m'}$ は解析的不分岐である。
% P01-E1009: frozen English doubles the conditional in item 3.
\item $(R, \mathfrak m) \to (S, \mathfrak m')$ が有限な局所環準同型で、
$S$ が整域である任意の場合に、$S$ は解析的不分岐である。
\end{enumerate}
\end{lemma}

\begin{proof}
$R$ が永田環であると仮定し、$R \to S$ と $\mathfrak m' \subset S$ を
(2) の通りとする。補題 \ref{algebra-lemma-quasi-finite-over-nagata} により
$S$ は永田環である。従って局所環 $S_{\mathfrak m'}$ は永田環である
（補題 \ref{algebra-lemma-nagata-localize}）。ゆえに補題
\ref{algebra-lemma-local-nagata-domain-analytically-unramified} により、これは
解析的不分岐である。(2) が (3) を含意することは明らかである。

\medskip\noindent
(3) が成り立つと仮定する。$\mathfrak p \subset R$ を素イデアルとし、
$L/\kappa(\mathfrak p)$ を有限体拡大とする。
% P01-E1010: frozen English leaves the integral closure's ambient field implicit.
(1) を示すには、この体拡大における $R/\mathfrak p$ の整閉包が
$R/\mathfrak p$ 上有限であることを示さなければならない。
$\kappa(\mathfrak p)$ 上 $L$ を生成する元
$x_1, \ldots, x_n \in L$ をとる。各 $i$ に対して、
$P_i(T) = T^{d_i} + a_{i, 1} T^{d_i - 1} + \ldots + a_{i, d_i}$ を
$\kappa(\mathfrak p)$ 上の $x_i$ の最小多項式とする。
適当な $f_i \in R$、$f_i \not \in \mathfrak p$ により $x_i$ を
$f_i x_i$ で置き換えることで、$a_{i, j} \in R/\mathfrak p$ と仮定してよい。
実際、さらに $\mathfrak m$ の元を掛けることで、すべての $i, j$ に対し
$a_{i, j} \in \mathfrak m/\mathfrak p \subset R/\mathfrak p$ と仮定できる。
% P01-E1011: frozen formula omits parentheses around the quotient coefficient ring.
こうして $S = R/\mathfrak p[x_1, \ldots, x_n] \subset L$ とおく。
このとき $S$ は $R$ 上有限な整域であり、$S/\mathfrak m S$ は
% P01-E1012: frozen formula omits parentheses around the quotient coefficient ring.
$R/\mathfrak m[T_1, \ldots, T_n]/(T_1^{d_1}, \ldots, T_n^{d_n})$ の
商である。従って $S$ は局所環である。(3) により $S$ は
解析的不分岐であり、補題 \ref{algebra-lemma-analytically-unramified-easy} により、
$L$ におけるその整閉包 $S'$ は $S$ 上有限である。
$S'$ はまた $L$ における $R/\mathfrak p$ の整閉包なので、結論を得る。
\end{proof}

\noindent
次の命題から、特に永田環上の有限型代数は永田環であることが分かる。

\begin{proposition}[Nagata]
\label{algebra-proposition-nagata-universally-japanese}
$R$ を環とする。次は同値である：
\begin{enumerate}
\item $R$ は永田環である。
\item 任意の有限型 $R$-代数は永田環である。
\item $R$ は普遍的日本環かつネーター環である。
\end{enumerate}
\end{proposition}

\begin{proof}
ネーターな普遍的日本環について任意の有限型代数が永田環であること、
すなわち条件 (2) が成り立つことは明らかである。$R$ を永田環とする。
任意の有限生成 $R$-代数 $S$ が永田環であることを示す。
これにより命題が証明される。

\medskip\noindent
ステップ 1. 各 $R_i \to R_{i + 1}$ が一つの元で生成されるような環準同型の列
$R = R_0 \to R_1 \to R_2 \to \ldots \to R_n = S$ が存在する。
従って帰納法により、$S \cong R[x]/I$ の場合に $S$ が永田環であることを
証明すれば十分である。

\medskip\noindent
ステップ 2. $\mathfrak q \subset S$ を $S$ の素イデアルとし、
$\mathfrak p \subset R$ を対応する $R$ の素イデアルとする。
$S/\mathfrak q$ が N-2 であることを示さなければならない。
従って次の主張を証明することに帰着した：
(*) 永田整域 $R$ と整域の単生成拡大 $R \subset S$ が与えられたとき、
$S$ は N-2 である。

\medskip\noindent
ステップ 3：$R$ を永田整域とし、$R \subset S$ を整域の単生成拡大とする。
$R$ と $S$ の分数体の誘導された拡大が純超越であると仮定する。
この場合 $S = R[x]$ である。補題 \ref{algebra-lemma-polynomial-ring-N-2} により
$S$ は N-2 である。従って次の主張を証明することに帰着した：
(**) 永田整域 $R$ と、分数体の有限拡大を誘導する整域の単生成拡大
$R \subset S$ が与えられたとき、$S$ は N-2 である。

\medskip\noindent
ステップ 4. $R$ を正規永田整域とし、$R \subset S$ を分数体の有限拡大
$L/K$ を誘導する整域の単生成拡大とする。
$S$ を $R$-代数として生成する元 $x \in S$ をとる。
$M/L$ を有限体拡大とする。
$M$ における $R$ の整閉包を $R'$ とする。
このとき、$M$ における $S$ の整閉包 $S'$ は、$M$ における
$R'[x]$ の整閉包に等しい。また $R'$ の分数体は $M$ であり、
$R \subset R'$ は有限である（$R$ の永田性による）。
これは $R'$ が永田環であることを意味する
（補題 \ref{algebra-lemma-quasi-finite-over-nagata}）。
$S'$ が $S$ 上有限であることを示すのは、$S'$ が $R'[x]$ 上有限であることを
示すのと同じである。$R$ を $R'$ で、$S$ を $R'[x]$ で置き換えると、
次の主張に帰着する：
(***) 分数体 $K$ をもつ正規永田整域 $R$ と $x \in K$ が与えられたとき、
$R$ と $x$ で生成される環 $S \subset K$ は N-1 である。

\medskip\noindent
ステップ 5. $R$ を分数体 $K$ をもつ正規永田整域とする。
$x = b/a \in K$ とする。$R$ と $x$ で生成される環 $S \subset K$ が
N-1 であることを示さなければならない。$S_a \cong R_a$ は正規である。
従って補題 \ref{algebra-lemma-characterize-N-1} により、$S$ の各極大イデアル
$\mathfrak m$ について $S_{\mathfrak m}$ が N-1 であることを示せば十分である。

\medskip\noindent
前段落と同じ仮定のもと、このような極大イデアルを一つとり、
$\mathfrak n = R \cap \mathfrak m$ とおく。剰余体拡大
$\kappa(\mathfrak m)/\kappa(\mathfrak n)$ は有限であり
（定理 \ref{algebra-theorem-nullstellensatz}）、$x$ の像で生成される。
% P01-E1013: frozen English omits the subject before “may assume”.
ある元 $a \in R$、$a \not \in \mathfrak n$ に対し $R$ を
$R_a$ で、$S$ を $S_a$ で置き換えることにより、次のようなモニック多項式が
存在すると仮定してよい：
% P01-E1014: frozen formula places f(X) in R[x] rather than R[X].
% P01-E1015: frozen summation gives a comma-separated list instead of bounds.
$f(X) = X^d + \sum_{i = 1, \ldots, d} a_iX^{d - i}$ が $R[x]$ に属し、
$f(x) \in \mathfrak m$ である（細部は省略する）。
多項式 $f(X)$ が $K''[X]$ で完全に分解するような有限体拡大
$K''/K$ をとる。$K''$ における $R$ の整閉包を $R'$ とする。
部分環 $S' \subset K''$ を、$R'$ と $x$ で生成されるものとする。
$R$ は永田環なので、$R'$ は $R$ 上有限かつ永田環である
（補題 \ref{algebra-lemma-quasi-finite-over-nagata}）。
さらに $S'$ は $S$ 上有限である。$S'$ のすべての極大イデアル
$\mathfrak m'$ に対して局所環 $S'_{\mathfrak m'}$ が N-1 ならば、補題
\ref{algebra-lemma-characterize-N-1} により $S'_{\mathfrak m}$ は N-1 である。
すると補題 \ref{algebra-lemma-finite-extension-N-2} により、
$S_{\mathfrak m}$ も N-1 である。
$R$ を $R'$ で、$S$ を $S'$ で、$\mathfrak m$ を $\mathfrak m$ 上にある
任意の極大イデアル $\mathfrak m'$ で置き換えると、上の多項式 $f$ が
完全に分解する状況に至る：
% P01-E1313: frozen product uses a comma-separated index list instead of bounds.
$f(X) = \prod_{i = 1, \ldots, d} (X - a_i)$、ここで $a_i \in R$ である。
$f(x) \in \mathfrak m$ なので、ある $i$ に対し
$x - a_i \in \mathfrak m$ である。最後に $x$ を $x - a_i$ で
置き換えることで、$x \in \mathfrak m$ と仮定してよい。

\medskip\noindent
まとめると、$R$ は分数体 $K$ をもつ正規永田整域であり、$x \in K$、
$S$ は $x$ と $R$ で生成される $K$ の部分環であり、最後に
$\mathfrak m \subset S$ は $x \in \mathfrak m$ を満たす極大イデアルである。
$S_{\mathfrak m}$ が N-1 であることを示さなければならない。

\medskip\noindent
局所環 $S_{\mathfrak m}$ と元 $x$ に補題
\ref{algebra-lemma-criterion-analytically-unramified} が適用できることを示す。
これにより $S_{\mathfrak m}$ は解析的不分岐となり、補題
\ref{algebra-lemma-analytically-unramified-easy} により N-1 であることが分かる。

\medskip\noindent
性質 (1)、(2)、(3)(a)、(3)(b) を確認しなければならない。
性質 (1) は自明である。$X \mapsto x$ である
$I = \Ker(R[X] \to S)$ とおく。
$K$ において $ax = b$ を満たす、すべての一次式 $aX - b$ により
$I$ が生成されると主張する。これらの一次式がすべて $I$ に属することは
明らかである。$g = a_d X^d + \ldots a_1 X + a_0 \in I$ ならば、
% P01-E1314: frozen polynomial omits the plus sign before its a_1 X term.
$a_dx$ は $R$ 上整であり（補題 \ref{algebra-lemma-make-integral-trivial}）、
$R$ は正規なので $b := a_dx \in R$ である。
このとき $g - (a_dX - b)X^{d - 1} \in I$ であり、次数に関する帰納法で
結論を得る。その帰結として
$$
S/xS = R[X]/(X, I) = R/J
$$
であり、ここで
$$
J = \{b \in R \mid ax = b \text{ for some }a \in R\} = xR \cap R
$$
である。補題 \ref{algebra-lemma-normal-domain-intersection-localizations-height-1} により、
$S/xS = R/J$ は $R$-加群として埋没素イデアルをもたず、従って
$R/J$-加群としても、$S/xS$-加群としても、$S$-加群としても
埋没素イデアルをもたない。これで性質 (2) が示された。
$\mathfrak q \subset \mathfrak m$ を満たすこのような随伴素イデアル
$\mathfrak q \subset S$ をとる（従ってこれは
$S_{\mathfrak m}/xS_{\mathfrak m}$ の随伴素イデアルであるが、以下の議論には
影響しない）。このとき $\mathfrak q$ は $xS$ 上極小であり、従って高さは
$1$ である。上の等式列から、$\mathfrak p = R \cap \mathfrak q$ は
$R/J$ の随伴素イデアルであり、従って高さ $1$ である
（補題 \ref{algebra-lemma-normal-domain-intersection-localizations-height-1} を参照）。
従って $R_{\mathfrak p}$ は離散付値環であり、それゆえ
$R_{\mathfrak p} \subset S_{\mathfrak q}$ は等号である。
これは $S_{\mathfrak q}$ が正則であることを示す。これで性質 (3)(a) が
示された。最後に、$(S/\mathfrak q)_{\mathfrak m}$ は
$S/\mathfrak q$ の局所化であり、後者は $S/xS = R/J$ の商である。
従って $(S/\mathfrak q)_{\mathfrak m}$ は永田環 $R$ の商の局所化であり、
ゆえに永田環である（補題 \ref{algebra-lemma-quasi-finite-over-nagata} および
\ref{algebra-lemma-nagata-localize}）。従って補題
\ref{algebra-lemma-local-nagata-domain-analytically-unramified} により
解析的不分岐である。これで (3)(b) が成り立ち、証明が完了する。
\end{proof}

\begin{proposition}
\label{algebra-proposition-ubiquity-nagata}
次の種類の環は永田環であり、特に普遍的日本環である：
\begin{enumerate}
\item 体、
\item ネーター完備局所環、
\item $\mathbf{Z}$、
\item 分数体の標数が零である Dedekind 整域、
\item 上のいずれかの有限型環拡大。
\end{enumerate}
\end{proposition}

\begin{proof}
ネーター完備局所環の場合は補題
\ref{algebra-lemma-Noetherian-complete-local-Nagata} である。
他の場合には、環の各素イデアル $\mathfrak p$ に対し
$R/\mathfrak p$ が N-2 であるかを確認するだけでよい。
$R/\mathfrak p$ が体の場合、すなわち $\mathfrak p$ が極大の場合には
これは明らかである。従って Dedekind 環の場合、
$\mathfrak p = (0)$ のときだけ確認すればよい。しかし分数体の標数は零と
仮定しているので、補題 \ref{algebra-lemma-domain-char-zero-N-1-2} が適用できる。
\end{proof}

\begin{example}
\label{algebra-example-nonjapanese-dvr}
離散付値環が永田環であることと N-2 であることは同値である
（極大イデアルによる商が体であり、従って N-2 だからである）。
例 \ref{algebra-example-bad-dvr-char-p} の離散付値環 $A$ は永田環ではない。
すなわち N-2 ではない。実際、有限拡大 $A \subset R = A[f]$ は
N-1 ではない。これを見るため $f = \sum a_i x^i$ と書く。
各 $n \geq 1$ に対し $g_n = \sum_{i < n} a_i x^i \in A$ とおく。
すると $h_n = (f - g_n)/x^n$ は $R$ の分数体の元であり、
$h_n^p \in k^p[[x]] \subset A$ である。従って $R$ の整閉包 $R'$ は
$h_1, h_2, h_3, \ldots$ を含む。さて、$R'$ が $R$ 上、従って $A$ 上
有限であるならば、すべての $n$ に対して
$f  = x^n h_n + g_n$ は部分加群 $A + x^nR'$ に含まれる。
Artin--Rees により、これは $f \in A$ を含意する
（補題 \ref{algebra-lemma-intersect-powers-ideal-module-zero}）。これは矛盾である。
\end{example}

\begin{lemma}
\label{algebra-lemma-nagata-pth-roots}
$(A, \mathfrak m)$ を、永田環であり、分数体の標数が $p$ である
ネーター局所整域とする。$a \in A$ が $A^\wedge$ に $p$ 乗根をもつならば、
$a$ は $A$ に $p$ 乗根をもつ。
\end{lemma}

\begin{proof}
$\alpha \in A^\wedge$ を $a$ の $p$ 乗根とする。矛盾を導くため、
$a$ は $A$ に $p$ 乗根をもたないと仮定する。
このとき $a$ は $A$ の分数体 $K$ にも $p$ 乗根をもたない。
実際、$b, c \in A$、$c \not = 0$ に対して $\beta = b/c$ が
$\beta^p = a$ を満たすならば、$A^\wedge$ において
$\alpha = b/c$ である。$A \to A^\wedge$ の忠実平坦性により、これは
$A$ において $c$ が $b$ を割ることを意味し、従って $\beta \in A$ となり
矛盾する。ゆえに $K[x]/(x^p - a)$ は体なので、
$B = A[x]/(x^p - a)$ は整域である。しかし $\alpha$ の存在という仮定により、
$B$ の完備化は被約ではない。これは先の結果に矛盾する。実際、
$B$ は永田環であり（命題 \ref{algebra-proposition-nagata-universally-japanese}）、
従って補題 \ref{algebra-lemma-local-nagata-domain-analytically-unramified} により
解析的不分岐である。
\end{proof}

\section{性質の上昇}
\label{algebra-section-ascending-properties}

\noindent
この節では、環の性質の「上昇」に関するいくつかの代数的事実を証明し始める。
環の深さについてこれを行うには、加群の深さの上昇に関する次の結果を用いる
（\cite[IV, Proposition 6.3.1]{EGA} を参照）。

\begin{lemma}
\label{algebra-lemma-apply-grothendieck-module}
\begin{reference}
\cite[IV, Proposition 6.3.1]{EGA}
\end{reference}
次が成り立つ：
$$
\text{depth}(M \otimes_R N)
=
\text{depth}(M) + \text{depth}(N/\mathfrak m_RN)
$$
ここで $R \to S$ はネーター局所環の局所準同型、$M$ は有限
$R$-加群、$N$ は $R$ 上平坦な有限 $S$-加群である。
\end{lemma}

\begin{proof}
主張および以下の証明では、$M$ の深さは $R$-加群としての深さ、
$M \otimes_R N$ の深さは $S$-加群としての深さ、
$N/\mathfrak m_RN$ の深さは $S/\mathfrak m_RS$-加群としての深さとする。
% P01-E1016: frozen English omits “by” in the designation of n.
右辺を $n$ と書く。まず $n$ が零であると仮定する。
このとき $\text{depth}(M) = 0$ かつ
$\text{depth}(N/\mathfrak m_RN) = 0$ である。
これは、零化イデアルが $\mathfrak m_R$ である $z \in M$ と、
零化イデアルが $\mathfrak m_S/\mathfrak m_RS$ である
$\overline{y} \in N/\mathfrak m_RN$ が存在することを意味する。
$\overline{y}$ の持ち上げ $y \in N$ をとる。
$N$ は $R$ 上平坦なので、写像 $z : R/\mathfrak m_R \to M$ から
単射 $N/\mathfrak m_RN \to M \otimes_R N$ が得られる。
従って $z \otimes y$ の零化イデアルは $\mathfrak m_S$ である。
ゆえに $\text{depth}(M \otimes_R N) = 0$ でもある。

\medskip\noindent
$n > 0$ と仮定する。$\text{depth}(N/\mathfrak m_RN) > 0$ ならば、
$N/\mathfrak m_RN$ 上の非零因子 $\overline{f} \in S/\mathfrak m_RS$ に
写る $f \in \mathfrak m_S$ をとれる。このとき補題
\ref{algebra-lemma-depth-drops-by-one} により
$\text{depth}(N/\mathfrak m_RN) =
\text{depth}(N/(f, \mathfrak m_R)N) + 1$ である。
補題 \ref{algebra-lemma-mod-injective} により、元 $f \in S$ は $N$ 上の
非零因子であり、$N/fN$ は $R$ 上平坦である。
従って $n$ に関する帰納法により
$$
\text{depth}(M \otimes_R N/fN) =
\text{depth}(M) + \text{depth}(N/(f, \mathfrak m_R)N).
$$
$N/fN$ は $R$ 上平坦なので、最初の写像が $f$ による乗法である列
$$
0 \to M \otimes_R N \to M \otimes_R N \to M \otimes_R N/fN \to 0
$$
は完全である（補題 \ref{algebra-lemma-flat-tor-zero}）。従って補題
\ref{algebra-lemma-depth-drops-by-one} により
$\text{depth}(M \otimes_R N) = \text{depth}(M \otimes_R N/fN) + 1$
となり、補題の式において等号が成り立つと結論できる。

\medskip\noindent
$n > 0$ であるが $\text{depth}(N/\mathfrak m_RN) = 0$ ならば、
$M$ 上の非零因子 $f \in \mathfrak m_R$ をとれる。
$N$ は $R$ 上平坦なので、$f$ は $M \otimes_R N$ 上でも非零因子である。
再び $n$ に関する帰納法により
$$
\text{depth}(M/fM \otimes_R N) =
\text{depth}(M/fM) + \text{depth}(N/\mathfrak m_RN).
$$
この場合、補題 \ref{algebra-lemma-depth-drops-by-one} により
$\text{depth}(M \otimes_R N) = \text{depth}(M/fM \otimes_R N) + 1$
かつ $\text{depth}(M) = \text{depth}(M/fM) + 1$ である。
従って補題の式において等号が成り立つ。
\end{proof}

\begin{lemma}
\label{algebra-lemma-apply-grothendieck}
$R \to S$ をネーター局所環の平坦な局所環準同型とする。このとき
$$
\text{depth}(S) = \text{depth}(R) + \text{depth}(S/\mathfrak m_RS).
$$
\end{lemma}

\begin{proof}
これは補題 \ref{algebra-lemma-apply-grothendieck-module} の特別な場合である。
\end{proof}

\begin{lemma}
\label{algebra-lemma-CM-goes-up}
$R \to S$ をネーター局所環の平坦な局所準同型とする。
% P01-E1017: frozen English omits the colon introducing this equivalence list.
このとき次は同値である：
\begin{enumerate}
\item $S$ は Cohen--Macaulay である。
\item $R$ および $S/\mathfrak m_RS$ は Cohen--Macaulay である。
\end{enumerate}
\end{lemma}

\begin{proof}
定義ならびに補題 \ref{algebra-lemma-apply-grothendieck} および
\ref{algebra-lemma-dimension-base-fibre-equals-total} から従う。
\end{proof}

\begin{lemma}
\label{algebra-lemma-Sk-goes-up}
$\varphi : R \to S$ を環準同型とする。次を仮定する。
\begin{enumerate}
\item $R$ はネーター環である。
\item $S$ はネーター環である。
\item $\varphi$ は平坦である。
\item ファイバー環 $S \otimes_R \kappa(\mathfrak p)$ は $(S_k)$ を満たす。
\item $R$ は性質 $(S_k)$ をもつ。
\end{enumerate}
このとき $S$ は性質 $(S_k)$ をもつ。
\end{lemma}

\begin{proof}
$\mathfrak q$ を $R$ の素イデアル $\mathfrak p$ の上にある
$S$ の素イデアルとする。補題 \ref{algebra-lemma-apply-grothendieck} により
$$
\text{depth}(S_{\mathfrak q}) =
\text{depth}(S_{\mathfrak q}/\mathfrak pS_{\mathfrak q}) +
\text{depth}(R_{\mathfrak p}).
$$
一方、補題 \ref{algebra-lemma-dimension-base-fibre-total} により
$$
\dim(R_{\mathfrak p})
+
\dim(S_{\mathfrak q}/\mathfrak pS_{\mathfrak q})
\geq
\dim(S_{\mathfrak q})
$$
である。（実際には補題
\ref{algebra-lemma-dimension-base-fibre-equals-total} により等号が成り立つが、
厳密にはここでは必要ない。）
最後に、写像のファイバー環は $(S_k)$ を満たすと仮定しているので、
$\text{depth}(S_{\mathfrak q}/\mathfrak pS_{\mathfrak q})
\geq \min(k, \dim(S_{\mathfrak q}/\mathfrak pS_{\mathfrak q}))$ である。
従って次の不等式列から補題が従う：
\begin{eqnarray*}
\text{depth}(S_{\mathfrak q}) & = &
\text{depth}(S_{\mathfrak q}/\mathfrak pS_{\mathfrak q}) +
\text{depth}(R_{\mathfrak p}) \\
& \geq &
\min(k, \dim(S_{\mathfrak q}/\mathfrak pS_{\mathfrak q})) +
\min(k, \dim(R_{\mathfrak p})) \\
& = &
\min(2k, \dim(S_{\mathfrak q}/\mathfrak pS_{\mathfrak q}) + k,
k + \dim(R_\mathfrak p),
\dim(S_{\mathfrak q}/\mathfrak pS_{\mathfrak q}) +
\dim(R_{\mathfrak p})) \\
& \geq &
\min(k, \dim(S_{\mathfrak q}))
\end{eqnarray*}
を得る。
\end{proof}

\begin{lemma}
\label{algebra-lemma-Rk-goes-up}
$\varphi : R \to S$ を環準同型とする。次を仮定する。
\begin{enumerate}
\item $R$ はネーター環である。
% P01-E1018: frozen English omits punctuation after the second assumption.
\item $S$ はネーター環である。
\item $\varphi$ は平坦である。
\item ファイバー環 $S \otimes_R \kappa(\mathfrak p)$ は性質 $(R_k)$ をもつ。
\item $R$ は性質 $(R_k)$ をもつ。
\end{enumerate}
このとき $S$ は性質 $(R_k)$ をもつ。
\end{lemma}

\begin{proof}
$\mathfrak q$ を $R$ の素イデアル $\mathfrak p$ の上にある
$S$ の素イデアルとする。$\dim(S_{\mathfrak q}) \leq k$ と仮定する。
補題 \ref{algebra-lemma-dimension-base-fibre-equals-total} により
$\dim(S_{\mathfrak q}) = \dim(R_{\mathfrak p})
+ \dim(S_{\mathfrak q}/\mathfrak pS_{\mathfrak q})$ なので、
$\dim(R_{\mathfrak p}) \leq k$ かつ
$\dim(S_{\mathfrak q}/\mathfrak pS_{\mathfrak q}) \leq k$ である。
従って仮定により、$R_{\mathfrak p}$ および
$S_{\mathfrak q}/\mathfrak pS_{\mathfrak q}$ は正則である。
補題 \ref{algebra-lemma-flat-over-regular-with-regular-fibre} により、
$S_{\mathfrak q}$ は正則である。
\end{proof}

\begin{lemma}
\label{algebra-lemma-reduced-goes-up-noetherian}
$\varphi : R \to S$ を環準同型とする。次を仮定する。
\begin{enumerate}
\item $R$ はネーター環である。
% P01-E1019: frozen English omits punctuation after the second assumption.
\item $S$ はネーター環である。
\item $\varphi$ は平坦である。
\item ファイバー環 $S \otimes_R \kappa(\mathfrak p)$ は被約である。
\item $R$ は被約である。
\end{enumerate}
このとき $S$ は被約である。
\end{lemma}

\begin{proof}
ネーター環が被約であることは、性質 $(S_1)$ と $(R_0)$ をもつことと
同値である（補題 \ref{algebra-lemma-criterion-reduced} を参照）。
従って $R$ とファイバー環がこれらの性質をもつことが分かっている。
補題 \ref{algebra-lemma-Sk-goes-up} および \ref{algebra-lemma-Rk-goes-up} を適用すると、
$S$ は $(S_1)$ と $(R_0)$ をもち、再び補題
\ref{algebra-lemma-criterion-reduced} により、言い換えれば被約である。
\end{proof}

\begin{lemma}
\label{algebra-lemma-reduced-goes-up}
$\varphi : R \to S$ を環準同型とする。次を仮定する。
\begin{enumerate}
\item $\varphi$ は滑らかである。
\item $R$ は被約である。
\end{enumerate}
このとき $S$ は被約である。
\end{lemma}

\begin{proof}
$R \to S$ は正則ファイバーをもつ平坦写像であることに注意する
（節 \ref{algebra-section-smooth-overview} の滑らかな環準同型に関する結果の一覧を
参照）。特にファイバーは被約である。従って $R$ がネーター環ならば、
$S$ もネーター環であり、補題 \ref{algebra-lemma-reduced-goes-up-noetherian} から
結果を得る。

\medskip\noindent
一般の場合、有限生成 $\mathbf{Z}$-部分代数 $R_0 \subset R$ と、
$S \cong R \otimes_{R_0} S_0$ となる滑らかな環準同型
$R_0 \to S_0$ をとれる（節 \ref{algebra-section-smooth-overview} の注意 (10) を
参照）。さて、$x \in S$ が $x^2 = 0$ を満たすならば、$R_0$ を拡大して、
$x$ が元 $x_0 \in S_0$ から来ると仮定できる。さらに $R_0$ を拡大して、
$S_0$ において $x_0^2 = 0$ と仮定できる。
% P01-E1020: frozen English grammatically predicates reducedness of an inclusion.
しかし $R_0 \subset R$ は被約環の部分環を表すので、その始域は被約である。
従って $S_0$ は被約であり、望む通り $x_0 = 0$ である。
\end{proof}

\begin{lemma}
\label{algebra-lemma-normal-goes-up-noetherian}
$\varphi : R \to S$ を環準同型とする。次を仮定する。
\begin{enumerate}
\item $R$ はネーター環である。
\item $S$ はネーター環である。
\item $\varphi$ は平坦である。
\item ファイバー環 $S \otimes_R \kappa(\mathfrak p)$ は正規である。
\item $R$ は正規である。
\end{enumerate}
このとき $S$ は正規である。
\end{lemma}

\begin{proof}
ネーター環が正規であることは、性質 $(S_2)$ と $(R_1)$ をもつことと
同値である（補題 \ref{algebra-lemma-criterion-normal} を参照）。
従って $R$ とファイバー環がこれらの性質をもつことが分かっている。
補題 \ref{algebra-lemma-Sk-goes-up} および \ref{algebra-lemma-Rk-goes-up} を適用すると、
$S$ は $(S_2)$ と $(R_1)$ をもち、再び補題
\ref{algebra-lemma-criterion-normal} により、言い換えれば正規である。
\end{proof}

\begin{lemma}
\label{algebra-lemma-normal-goes-up}
$\varphi : R \to S$ を環準同型とする。次を仮定する。
\begin{enumerate}
\item $\varphi$ は滑らかである。
\item $R$ は正規である。
\end{enumerate}
このとき $S$ は正規である。
\end{lemma}

\begin{proof}
$R \to S$ は正則ファイバーをもつ平坦写像であることに注意する
（節 \ref{algebra-section-smooth-overview} の滑らかな環準同型に関する結果の一覧を
参照）。特にファイバーは正規である。従って $R$ がネーター環ならば、
$S$ もネーター環であり、補題 \ref{algebra-lemma-normal-goes-up-noetherian} から
結果を得る。

\medskip\noindent
次に一般の場合を扱う。まず $R$ は被約であり、従って補題
\ref{algebra-lemma-reduced-goes-up} により $S$ は被約である。
$\mathfrak q$ を $S$ の素イデアルとし、$\mathfrak p$ を対応する
$R$ の素イデアルとする。$R_{\mathfrak p}$ は正規整域であることに注意する。
$S_{\mathfrak q}$ が正規整域であることを示さなければならない。
そのため $R$ を $R_{\mathfrak p}$ で、$S$ を $S_{\mathfrak p}$ で
置き換えてよい。従って $R$ は正規整域であると仮定してよい。

\medskip\noindent
$R \to S$ が滑らかで、$R$ が正規整域であると仮定する。
有限生成 $\mathbf{Z}$-部分代数 $R_0 \subset R$ と、
$S \cong R \otimes_{R_0} S_0$ となる滑らかな環準同型
$R_0 \to S_0$ をとれる（節 \ref{algebra-section-smooth-overview} の注意 (10) を
参照）。$R_0$ は永田整域なので（命題
\ref{algebra-proposition-ubiquity-nagata} を参照）、その整閉包 $R_0'$ は
$R_0$ 上有限である。さらに $R$ は正規整域なので、
$R_0' \subset R$ は明らかである。従って $R_0$ を $R_0'$ で、
$S_0$ を $R_0' \otimes_{R_0} S_0$ で置き換えることで、$R_0$ は
ネーター正規整域であると仮定できる。証明の第 1 段落から、
$S_0$ は正規環であると結論できる（もちろん整域であるとは限らない）。
このようにして、$R = \bigcup R_\lambda$ はネーター正規整域の合併であり、
それに応じて $S = \colim R_\lambda \otimes_{R_0} S_0$ は正規環の
余極限であることが分かる。これは $S$ が正規環であることを含意する。
細部は省略する。
\end{proof}

\begin{lemma}
\label{algebra-lemma-regular-goes-up}
\begin{slogan}
正則性は滑らかな環準同型に沿って上昇する。
\end{slogan}
$\varphi : R \to S$ を環準同型とする。次を仮定する。
\begin{enumerate}
\item $\varphi$ は滑らかである。
\item $R$ は正則環である。
\end{enumerate}
このとき $S$ は正則である。
\end{lemma}

\begin{proof}
% P01-E1021: frozen English attaches “for all” to the property notation.
すべての $(R_k)$ に対して補題 \ref{algebra-lemma-Rk-goes-up} を適用し、補題
\ref{algebra-lemma-characterize-smooth-over-field} により仮定が満たされることを
確認すれば従う。
\end{proof}

\section{性質の下降}
\label{algebra-section-descending-properties}

\noindent
この節では、環の性質の「下降」に関するいくつかの代数的事実を証明し始める。
性質を上昇させるより下降させる方がしばしば「容易」であることが分かる。
% P01-E1022: frozen English has singular “assumption” with plural “are”.
言い換えれば、環準同型 $R \to S$ に関する仮定は、前節の対応する補題に
おける仮定より弱いことが多い。しかし、対応する上昇も示せない限り、
下降に関する結果は役に立たないことが多いので注意されたい。
この現象を示す典型的な結果は次である。

\begin{lemma}
\label{algebra-lemma-descent-Noetherian}
$R \to S$ を環準同型とする。次を仮定する。
\begin{enumerate}
\item $R \to S$ は忠実平坦である。
\item $S$ はネーター環である。
\end{enumerate}
このとき $R$ はネーター環である。
\end{lemma}

\begin{proof}
$I_0 \subset I_1 \subset I_2 \subset \ldots$ を $R$ のイデアルの増大列とする。
仮定により、ある $n$ に対して
$I_nS = I_{n+1}S = I_{n+2}S = \ldots$ である。
忠実平坦性により、拡大して縮約すると同じイデアルに戻る。すなわち、
$R$ の各イデアル $I$ に対して $I = R \cap IS$ である
（補題 \ref{algebra-lemma-faithfully-flat-universally-injective}）。従って望む通り
$I_n = I_{n+1} = I_{n+2} = \ldots$ である。
\end{proof}

\begin{lemma}
\label{algebra-lemma-descent-reduced}
$R \to S$ を環準同型とする。次を仮定する。
\begin{enumerate}
\item $R \to S$ は忠実平坦である。
\item $S$ は被約である。
\end{enumerate}
このとき $R$ は被約である。
\end{lemma}

\begin{proof}
補題 \ref{algebra-lemma-faithfully-flat-universally-injective} により
$R \to S$ は単射なので明らかである。
\end{proof}

\begin{lemma}
\label{algebra-lemma-descent-normal}
$R \to S$ を環準同型とする。次を仮定する。
\begin{enumerate}
\item $R \to S$ は忠実平坦である。
\item $S$ は正規環である。
\end{enumerate}
このとき $R$ は正規環である。
\end{lemma}

\begin{proof}
$S$ は被約なので $R$ も被約である。$R$ の素イデアル $\mathfrak p$ をとる。
$R_{\mathfrak p}$ が正規整域であることを示さなければならない。
% P01-E1023: frozen English omits “flat” after “faithfully”.
$S_{\mathfrak p}$ は $R_{\mathfrak p}$ 上でも忠実平坦なので、$R$ は
極大イデアル $\mathfrak m$ をもつ局所環であると仮定してよい。
$\mathfrak q$ を $\mathfrak m$ の上にある $S$ の素イデアルとする。
このとき $R \to S_{\mathfrak q}$ は忠実平坦である
（補題 \ref{algebra-lemma-local-flat-ff}）。従って $S$ も局所環であると仮定してよい。
特に $S$ は正規整域である。$R \to S$ は忠実平坦であり、$S$ は
正規整域なので、$R$ は整域である。次に $a, b \in R$ に対して
$a/b$ が $R$ 上整であると仮定する。$S$ は正規なので $a/b \in S$、従って
$a \in bS$ である。これは $a : R \to R/bR$ が $S$ への基底変換後に
零写像となることを意味する。忠実平坦性により $a \in bR$、従って
$a/b \in R$ である。ゆえに $R$ は正規である。
\end{proof}

\begin{lemma}
\label{algebra-lemma-descent-regular}
$R \to S$ を環準同型とする。次を仮定する。
\begin{enumerate}
\item $R \to S$ は忠実平坦である。
\item $S$ は正則環である。
\end{enumerate}
このとき $R$ は正則環である。
\end{lemma}

\begin{proof}
補題 \ref{algebra-lemma-descent-Noetherian} により $R$ はネーター環である。
素イデアル $\mathfrak p \subset R$ をとり、$\mathfrak p$ の上にある
素イデアル $\mathfrak q \subset S$ を選ぶ。このとき補題
\ref{algebra-lemma-flat-under-regular} を $R_\mathfrak p \to S_\mathfrak q$ に適用すると、
$R_\mathfrak p$ は正則であると結論できる。$\mathfrak p$ は任意だったので、
$R$ は正則である。
\end{proof}

\begin{lemma}
\label{algebra-lemma-descent-Sk}
$R \to S$ を環準同型とする。次を仮定する。
\begin{enumerate}
\item $R \to S$ は忠実平坦である。
\item $S$ はネーター環であり、性質 $(S_k)$ をもつ。
\end{enumerate}
このとき $R$ はネーター環であり、性質 $(S_k)$ をもつ。
\end{lemma}

\begin{proof}
(1) と (2) から $R$ がネーター環であることは既に見た
（補題 \ref{algebra-lemma-descent-Noetherian} を参照）。
素イデアル $\mathfrak p \subset R$ をとる。ファイバー環
$S \otimes_R \kappa(\mathfrak p)$ の極小素イデアルに対応する、
$\mathfrak p$ の上にある素イデアル $\mathfrak q \subset S$ を選ぶ。
このとき $A = R_{\mathfrak p} \to S_{\mathfrak q} = B$ はネーター局所環の
平坦な局所環準同型であり、$\mathfrak m_AB$ は $B$ の定義イデアルである。
従って $\dim(A) = \dim(B)$（補題
\ref{algebra-lemma-dimension-base-fibre-equals-total}）かつ
$\text{depth}(A) = \text{depth}(B)$（補題
\ref{algebra-lemma-apply-grothendieck}）である。ゆえに $B$ は $(S_k)$ をもつので、
$A$ も $(S_k)$ をもつ。
\end{proof}

\begin{lemma}
\label{algebra-lemma-descent-Rk}
$R \to S$ を環準同型とする。次を仮定する。
\begin{enumerate}
\item $R \to S$ は忠実平坦である。
\item $S$ はネーター環であり、性質 $(R_k)$ をもつ。
\end{enumerate}
このとき $R$ はネーター環であり、性質 $(R_k)$ をもつ。
\end{lemma}

\begin{proof}
(1) と (2) から $R$ がネーター環であることは既に見た
（補題 \ref{algebra-lemma-descent-Noetherian} を参照）。
素イデアル $\mathfrak p \subset R$ をとり、
$\dim(R_{\mathfrak p}) \leq k$ と仮定する。ファイバー環
$S \otimes_R \kappa(\mathfrak p)$ の極小素イデアルに対応する、
$\mathfrak p$ の上にある素イデアル $\mathfrak q \subset S$ を選ぶ。
このとき $A = R_{\mathfrak p} \to S_{\mathfrak q} = B$ はネーター局所環の
平坦な局所環準同型であり、$\mathfrak m_AB$ は $B$ の定義イデアルである。
従って $\dim(A) = \dim(B)$ である
（補題 \ref{algebra-lemma-dimension-base-fibre-equals-total}）。
$S$ は $(R_k)$ をもつので、$B$ は正則局所環である。
補題 \ref{algebra-lemma-flat-under-regular} により $A$ は正則である。
\end{proof}

\begin{lemma}
\label{algebra-lemma-descent-nagata}
$R \to S$ を環準同型とする。次を仮定する。
\begin{enumerate}
\item $R \to S$ は滑らかで、スペクトル上全射である。
\item $S$ は永田環である。
\end{enumerate}
このとき $R$ は永田環である。
\end{lemma}

\begin{proof}
永田環はネーター普遍的日本環と同じものであることを思い出す
（命題 \ref{algebra-proposition-nagata-universally-japanese}）。
補題 \ref{algebra-lemma-descent-Noetherian} において $R$ がネーター環であることは
既に見た。$R \to A$ を整域への有限型環準同型とする。
補題 \ref{algebra-lemma-check-universally-japanese} により、$A$ が N-1 であることを
確認すれば十分である。$B = A \otimes_R S$ は有限型 $S$-代数であり、
従って永田であることは明らかである
（命題 \ref{algebra-proposition-nagata-universally-japanese}）。
$A \to B$ は滑らかなので（補題 \ref{algebra-lemma-base-change-smooth}）、
$B$ は被約である（補題 \ref{algebra-lemma-reduced-goes-up}）。
$B$ はネーター環なので、極小素イデアルは有限個
$\mathfrak q_1, \ldots, \mathfrak q_t$ しかない
（補題 \ref{algebra-lemma-Noetherian-irreducible-components} を参照）。
% P01-E1024: frozen English omits terminal punctuation here.
$A \to B$ は平坦なので、これらはすべて $(0) \subset A$ の上にある
（下降定理、補題 \ref{algebra-lemma-flat-going-down} を参照）。
全商環 $Q(B)$ は $L_i = \kappa(\mathfrak q_i)$ の積である
（補題 \ref{algebra-lemma-total-ring-fractions-no-embedded-points} および
\ref{algebra-lemma-minimal-prime-reduced-ring}）。さらに、$Q(B)$ における $B$ の
整閉包 $B'$ は、各因子 $L_i$ における $B/\mathfrak q_i$ の整閉包 $B_i'$ の
積である（補題 \ref{algebra-lemma-characterize-reduced-ring-normal} と比較せよ）。
$B$ は普遍的日本環なので、環拡大 $B/\mathfrak q_i \subset B_i'$ は有限であり、
$B' = \prod B_i'$ は $B$ 上有限であると結論できる。
$A \to B$ は平坦なので、$A$ の任意の非零因子は $B$ の非零因子に写る。
対応する写像
$$
Q(A) \otimes_A B = (A \setminus \{0\})^{-1}A \otimes_A B
= (A \setminus \{0\})^{-1}B \to Q(B)
$$
は単射である（補題 \ref{algebra-lemma-tensor-localization} を用いた）。
% P01-E1025: frozen English first uses A-prime without defining it.
この写像を通じて $A'$ は $B'$ に写る。これにより単射
$$
A' \otimes_A B \longrightarrow B'
$$
が誘導される（上記と $A \to B$ の平坦性による）。$B'$ は有限
$B$-加群であり、$B$ はネーター環なので、$A' \otimes_A B$ は有限
$B$-加群である。従って有限個の元 $x_i \in A'$ が存在し、元
$x_i \otimes 1$ は $A' \otimes_A B$ を $B$-加群として生成する。
% P01-E1026: frozen English uses past tense “generated” in this conclusion.
最後に $A \to B$ の忠実平坦性により、$x_i$ は $A'$ を
$A$-加群としても生成すると結論でき、証明が終わる。
\end{proof}

\begin{remark}
\label{algebra-remark-universally-catenary-does-not-descend}
「普遍鎖状」であるという性質は下降しない。エタール環準同型に沿ってさえ
下降しない。Examples の節
\ref{examples-section-non-catenary-Noetherian-local} には、局所ネーター環で
普遍鎖状でない $A$ と、二つの極大イデアル $\mathfrak m$, $\mathfrak n$ を
もち、$B_{\mathfrak m}$ と $B_{\mathfrak n}$ がそれぞれ次元 $2$ と $1$ の
正則環である半局所環 $B$ との有限環準同型 $A \to B$ の構成がある。
% P01-E1027: frozen English mismatches plural “fields” with singular “that”.
また、それぞれの剰余体は $A$ の剰余体と同じである。さらに、
$\mathfrak m_A$ は $B_{\mathfrak m}$ と $B_{\mathfrak n}$ の双方で
極大イデアルを生成する（従って $A \to B$ は有限であるだけでなく
不分岐でもある）。補題 \ref{algebra-lemma-etale-makes-unramified-closed} により、
$B \otimes_A A' = B_1 \times B_2$ と分解し、$A' \to B_i$ が全射となるような
局所エタール環準同型 $A \to A'$ が存在する。これは、$A'$ が二つの
極小素イデアル $\mathfrak q_i$ をもち、$A'/\mathfrak q_i \cong B_i$ となる
ことを示す。$B_i$ は正則局所環なので（$B_{\mathfrak m}$ または
$B_{\mathfrak n}$ のいずれか上エタールだからである）、$A'$ は
普遍鎖状であると結論できる。
\end{remark}

\section{幾何学的正規代数}
\label{algebra-section-geometrically-normal}

\noindent
この節では、環の性質の上昇と下降のいくつかの応用を述べる。

\begin{lemma}
\label{algebra-lemma-geometrically-normal}
$k$ を体、$A$ を $k$-代数とする。$A$ に関する次の性質は同値である。
\begin{enumerate}
\item 任意の体拡大 $k'/k$ に対して $k' \otimes_k A$ は正規環である。
\item 任意の有限生成体拡大 $k'/k$ に対して $k' \otimes_k A$ は正規環である。
\item 任意の有限純非分離拡大 $k'/k$ に対して $k' \otimes_k A$ は正規環である。
\item $k^{perf} \otimes_k A$ は正規環である。
\end{enumerate}
ここで正規環は定義 \ref{algebra-definition-ring-normal} で定義されたものである。
\end{lemma}

\begin{proof}
(1) $\Rightarrow$ (2) $\Rightarrow$ (3) および
(1) $\Rightarrow$ (4) は明らかである。

\medskip\noindent
$k'/k$ が有限純非分離拡大ならば、$k$-拡大の埋め込み
$k' \to k^{perf}$ が存在する。環準同型
$k' \otimes_k A \to k^{perf} \otimes_k A$ は忠実平坦なので、
$k^{perf} \otimes_k A$ が正規ならば、補題 \ref{algebra-lemma-descent-normal} により
$k' \otimes_k A$ は正規である。これにより (4) $\Rightarrow$ (3) が分かる。

\medskip\noindent
(2) を仮定し、$k'/k$ を任意の体拡大とする。この体を有限生成
体拡大の有向余極限 $k' = \colim_i k_i$ と書ける。従って
$k' \otimes_k A = \colim_i k_i \otimes_k A$ は正規環の有向余極限である。
ゆえに補題 \ref{algebra-lemma-colimit-normal-ring} により $k' \otimes_k A$ は
正規環であり、(1) が成り立つ。

\medskip\noindent
(3) を仮定し、$K/k$ を有限生成体拡大とする。補題
\ref{algebra-lemma-make-separable} により、可換図式
$$
\xymatrix{
K \ar[r] & K' \\
k \ar[u] \ar[r] & k' \ar[u]
}
$$
であって、$k'/k$, $K'/K$ が有限純非分離体拡大であり、$K'/k'$ が
分離的となるものをとれる。補題 \ref{algebra-lemma-localization-smooth-separable} により、
滑らかな $k'$-代数 $B$ が存在し、$K'$ は $B$ の分数体である。
ここで次のように論じる。
ステップ 1：仮定 (3) により $k' \otimes_k A$ は正規環である。
ステップ 2：$k' \otimes_k A \to B \otimes_{k'} k' \otimes_k A$ は
滑らかであり（補題 \ref{algebra-lemma-base-change-smooth}）、滑らかな写像に沿って
正規性は上昇するので（補題 \ref{algebra-lemma-normal-goes-up}）、
$B \otimes_{k'} k' \otimes_k A$ は正規環である。
ステップ 3：$K' \otimes_{k'} k' \otimes_k A = K' \otimes_k A$ は
正規環の局所化なので正規環である（補題
\ref{algebra-lemma-localization-normal-ring}）。
ステップ 4：最後に、忠実平坦な環準同型
$K \otimes_k A \to K' \otimes_k A$ に沿う正規性の下降により
（補題 \ref{algebra-lemma-descent-normal}）、$K \otimes_k A$ は正規環である。
これで補題が証明された。
\end{proof}

\begin{definition}
\label{algebra-definition-geometrically-normal}
$k$ を体とする。$k$-代数 $R$ が補題
\ref{algebra-lemma-geometrically-normal} の同値な条件を満たすとき、この代数は $k$ 上
{\it 幾何学的正規}であるという。
\end{definition}

\begin{lemma}
\label{algebra-lemma-localization-geometrically-normal-algebra}
\begin{slogan}
局所化は幾何学的正規性を保つ。
\end{slogan}
$k$ を体とする。幾何学的正規な $k$-代数の局所化は幾何学的正規である。
\end{lemma}

\begin{proof}
環が正規であることは素イデアルにおける局所化で確認できるので明らかである。
\end{proof}

\begin{lemma}
\label{algebra-lemma-separable-field-extension-geometrically-normal}
$k$ を体とし、$K/k$ を分離体拡大とする。このとき $K$ は $k$ 上
幾何学的正規である。
\end{lemma}

\begin{proof}
$k^{perf} \otimes_k K$ が体なので成り立つ。実際、補題
\ref{algebra-lemma-separable-extension-preserves-reducedness} によりこれは被約である。
補題 \ref{algebra-lemma-perfection}（または定義 \ref{algebra-definition-perfection}）により、
体拡大 $k^{perf}/k$ は純非分離である。従って補題
\ref{algebra-lemma-radicial-integral-bijective} により、環 $k^{perf} \otimes_k K$ は
ただ一つの素イデアルをもつ。ただ一つの素イデアルをもつ被約環は体である。
\end{proof}

\begin{lemma}
\label{algebra-lemma-geometrically-normal-tensor-normal}
$k$ を体とし、$A, B$ を $k$-代数とする。$A$ は $k$ 上幾何学的正規であり、
$B$ は正規環であると仮定する。このとき $A \otimes_k B$ は正規環である。
\end{lemma}

\begin{proof}
$\mathfrak r$ を $A \otimes_k B$ の素イデアルとする。
% P01-E1028: frozen English has a malformed respective designation.
$\mathfrak p$, $\mathfrak q$ をそれぞれ $A$, $B$ の対応する素イデアルとする。
このとき $(A \otimes_k B)_{\mathfrak r}$ は
$A_{\mathfrak p} \otimes_k B_{\mathfrak q}$ の局所化である。従って補題
\ref{algebra-lemma-localization-normal-ring} および補題
\ref{algebra-lemma-localization-geometrically-normal-algebra} により、環
$A_{\mathfrak p} \otimes_k B_{\mathfrak q}$ について結果を証明すれば十分である。
ゆえに $A$ と $B$ は整域であると仮定してよい。

\medskip\noindent
% P01-E1029: frozen English twice says “fractions fields”.
$A$ と $B$ を分数体 $K$, $L$ をもつ整域とする。$B$ は有限型正規
% P01-E1030: frozen English twice splits the compound “subalgebra”.
$k$-部分代数のフィルター余極限であることに注意する（$k$ は永田環であるため、
命題 \ref{algebra-proposition-ubiquity-nagata} を参照。また、有限型 $k$-部分代数の
整閉包も命題 \ref{algebra-proposition-nagata-universally-japanese} により有限型
$k$-部分代数である）。補題 \ref{algebra-lemma-colimit-normal-ring} により、
$B$ が $k$ 上有限型である場合に帰着する。

\medskip\noindent
% P01-E1029 and P01-E1031: frozen English repeats “fractions fields” and omits “is”.
$A$ と $B$ を分数体 $K$, $L$ をもつ整域とし、$B$ は $k$ 上有限型であるとする。
この場合、環 $K \otimes_k B$ は $K$ 上有限型であり、従ってネーター環である
（補題 \ref{algebra-lemma-Noetherian-permanence}）。特に $K \otimes_k B$ は有限個の
極小素イデアルしかもたない（補題
\ref{algebra-lemma-Noetherian-irreducible-components}）。$A \to A \otimes_k B$ は
平坦なので、これは $A \otimes_k B$ が有限個の極小素イデアルしか
もたないことを含意する（平坦な環準同型に対する下降定理、補題
\ref{algebra-lemma-flat-going-down} により、これらの素イデアルはすべて
$(0) \subset A$ の上にある）。従って $A \otimes_k B$ がその全商環において
整閉であることを証明すれば十分である
（補題 \ref{algebra-lemma-characterize-reduced-ring-normal}）。

\medskip\noindent
$K \otimes_k B$ と $A \otimes_k L$ はともに正規環であると主張する。
これが成り立てば、$A \otimes_k B$ 上整である $Q(A \otimes_k B)$ の任意の元
$x$ は（補題 \ref{algebra-lemma-normal-ring-integrally-closed} により）
$K \otimes_k B \cap A \otimes_k L = A \otimes_k B$ に含まれ、証明が終わる。
% P01-E1032: frozen source has K rather than k as this tensor base.
仮定により $A \otimes_K L$ は正規環なので、$K \otimes_k B$ が正規であることを
証明すれば十分である。

\medskip\noindent
$A$ は $k$ 上幾何学的正規なので、補題
\ref{algebra-lemma-localization-geometrically-normal-algebra} により $K$ は $k$ 上
幾何学的正規であり、従って $K$ は $k$ 上幾何学的被約である。
ゆえに $K = \bigcup K_i$ は、$k$ の幾何学的被約な有限生成体拡大の合併である
（補題 \ref{algebra-lemma-subalgebra-separable}）。各 $K_i$ は滑らかな $k$-代数の
局所化である（補題 \ref{algebra-lemma-localization-smooth-separable}）。従って
$K_i \otimes_k B$ は滑らかな $B$-代数の局所化であり、ゆえに正規である
（補題 \ref{algebra-lemma-normal-goes-up}）。こうして $K \otimes_k B$ は正規環である
（補題 \ref{algebra-lemma-colimit-normal-ring}）。以上で証明が終わる。
\end{proof}

\begin{lemma}
\label{algebra-lemma-geometrically-normal-over-separable-algebraic}
$k'/k$ を分離代数体拡大とし、$A$ を $k'$ 上の代数とする。このとき、$A$ が
$k$ 上幾何学的正規であることと、$k'$ 上幾何学的正規であることは同値である。
\end{lemma}

\begin{proof}
$L/k$ を有限純非分離体拡大とする。このとき
$L' = k' \otimes_k L$ は体であり（Fields の節
\ref{fields-section-algebraic} の内容を参照）、
$A \otimes_k L = A \otimes_{k'} L'$ である。従って $A$ が $k'$ 上
幾何学的正規ならば、$A$ は $k$ 上幾何学的正規である。

\medskip\noindent
$A$ が $k$ 上幾何学的正規であると仮定し、$K/k'$ を体拡大とする。このとき
$$
K \otimes_{k'} A = (K \otimes_k A) \otimes_{(k' \otimes_k k')} k'
$$
補題 \ref{algebra-lemma-separable-algebraic-diagonal} により
$k' \otimes_k k' \to k'$ は局所化なので、$K \otimes_{k'} A$ は
正規環の局所化であり、従って正規である。
\end{proof}

\section{幾何学的正則代数}
\label{algebra-section-geometrically-regular}

\noindent
$k$ を体とする。$A$ をネーター $k$-代数とし、$K/k$ を有限生成体拡大とする。
このとき環 $K \otimes_k A$ もネーター環である
（補題 \ref{algebra-lemma-Noetherian-field-extension} を参照）。従って次の補題には
意味がある。

\begin{lemma}
\label{algebra-lemma-geometrically-regular}
$k$ を体、$A$ を $k$-代数とし、$A$ はネーターであると仮定する。
$A$ に関する次の性質は同値である。
\begin{enumerate}
\item 任意の有限生成体拡大 $k'/k$ に対して $k' \otimes_k A$ は正則である。
\item 任意の有限純非分離拡大 $k'/k$ に対して $k' \otimes_k A$ は正則である。
\end{enumerate}
ここで正則環は定義 \ref{algebra-definition-regular} の意味で用いる。
\end{lemma}

\begin{proof}
補題直前の注意により、この補題には意味がある。
(1) $\Rightarrow$ (2) は明らかである。

\medskip\noindent
(2) を仮定し、$K/k$ を有限生成体拡大とする。補題
\ref{algebra-lemma-make-separable} により、可換図式
$$
\xymatrix{
K \ar[r] & K' \\
k \ar[u] \ar[r] & k' \ar[u]
}
$$
であって、$k'/k$, $K'/K$ が有限純非分離体拡大であり、$K'/k'$ が
分離的となるものをとれる。補題 \ref{algebra-lemma-localization-smooth-separable} により、
滑らかな $k'$-代数 $B$ が存在し、$K'$ は $B$ の分数体である。
ここで次のように論じる。
ステップ 1：仮定 (2) により $k' \otimes_k A$ は正則環である。
ステップ 2：$k' \otimes_k A \to B \otimes_{k'} k' \otimes_k A$ は
滑らかであり（補題 \ref{algebra-lemma-base-change-smooth}）、滑らかな写像に沿って
正則性は上昇するので（補題 \ref{algebra-lemma-regular-goes-up}）、
$B \otimes_{k'} k' \otimes_k A$ は正則環である。
ステップ 3：$K' \otimes_{k'} k' \otimes_k A = K' \otimes_k A$ は
正則環の局所化なので正則環である（定義から直ちに従う）。
ステップ 4：最後に、忠実平坦な環準同型
$K \otimes_k A \to K' \otimes_k A$ に沿う正則性の下降により
（補題 \ref{algebra-lemma-descent-regular}）、$K \otimes_k A$ は正則環である。
これで補題が証明された。
\end{proof}

\begin{definition}
\label{algebra-definition-geometrically-regular}
$k$ を体とし、$R$ をネーター $k$-代数とする。$k$-代数 $R$ が補題
\ref{algebra-lemma-geometrically-regular} の同値な条件を満たすとき、この代数は $k$ 上
{\it 幾何学的正則}であるという。
\end{definition}

\noindent
定義から、$k$ の任意の有限生成体拡大 $K$ に対して
$K \otimes_k R$ は $K$ 上幾何学的正則な代数であることが明らかである。
後に More on Algebra の命題
\ref{more-algebra-proposition-characterization-geometrically-regular} で、
$k \subset k' \subset k^{1/p}$（有限）であるときに $R \otimes_k k'$ が
正則であることを確認すれば十分であると分かる。

\begin{lemma}
\label{algebra-lemma-geometrically-regular-descent}
\begin{slogan}
幾何学的正則性は代数の忠実平坦写像に沿って下降する。
\end{slogan}
$k$ を体とし、$A \to B$ を忠実平坦な $k$-代数準同型とする。
$B$ が $k$ 上幾何学的正則ならば、$A$ もそうである。
\end{lemma}

\begin{proof}
$B$ は $k$ 上幾何学的正則であると仮定する。$k'/k$ を有限純非分離拡大とする。
このとき $A \to B$ の基底変換である
$A \otimes_k k' \to B \otimes_k k'$ は忠実平坦である
（補題 \ref{algebra-lemma-surjective-spec-radical-ideal} および
\ref{algebra-lemma-flat-base-change} による）。また、$B$ に関する仮定により
$B \otimes_k k'$ は $k$ 上正則である。従って補題
\ref{algebra-lemma-descent-regular} により $A \otimes_k k'$ は正則である。
\end{proof}

\begin{lemma}
\label{algebra-lemma-geometrically-regular-goes-up}
$k$ を体とし、$A \to B$ を $k$-代数の滑らかな環準同型とする。
$A$ が $k$ 上幾何学的正則ならば、$B$ も $k$ 上幾何学的正則である。
\end{lemma}

\begin{proof}
$k'/k$ を有限生成体拡大とする。このとき
$A \otimes_k k' \to B \otimes_k k'$ は滑らかな環準同型であり
（補題 \ref{algebra-lemma-base-change-smooth}）、$A \otimes_k k'$ は正則である。
従って補題 \ref{algebra-lemma-regular-goes-up} により $B \otimes_k k'$ は正則である。
\end{proof}

\begin{lemma}
\label{algebra-lemma-geometrically-regular-over-subfields}
$k$ を体とし、$A$ を $k$ 上の代数とする。$k = \colim k_i$ を部分体の
有向余極限とする。$A$ が各 $k_i$ 上幾何学的正則ならば、$A$ は $k$ 上
幾何学的正則である。
\end{lemma}

\begin{proof}
$k'/k$ を有限純非分離体拡大とする。有限個の変数を $k$ に添加し、有限個の
多項式関係を課すことで $k'$ を得られる。従って、ある $i$ と有限純非分離
体拡大 $k_i'/k_i$ が存在し、
% P01-E1033: frozen source has k_i rather than k-prime on the left.
$k_i = k \otimes_{k_i} k_i'$ となる。従って
$A \otimes_k k' = A \otimes_{k_i} k_i'$ であり、補題は明らかである。
\end{proof}

\begin{lemma}
\label{algebra-lemma-geometrically-regular-over-separable-algebraic}
$k'/k$ を分離代数体拡大とし、$A$ を $k'$ 上の代数とする。このとき、$A$ が
$k$ 上幾何学的正則であることと、$k'$ 上幾何学的正則であることは同値である。
\end{lemma}

\begin{proof}
$L/k$ を有限純非分離体拡大とする。このとき
$L' = k' \otimes_k L$ は体であり（Fields の節
\ref{fields-section-algebraic} の内容を参照）、
$A \otimes_k L = A \otimes_{k'} L'$ である。従って $A$ が $k'$ 上
幾何学的正則ならば、$A$ は $k$ 上幾何学的正則である。

\medskip\noindent
$A$ が $k$ 上幾何学的正則であると仮定する。$k'$ は $k$ の有限拡大の
フィルター余極限なので、補題
\ref{algebra-lemma-geometrically-regular-over-subfields} により、$k'/k$ は
有限分離拡大であると仮定してよい。環準同型
$$
k' \to A \otimes_k k' \to A.
$$
を考える。$A \otimes_k k'$ は $A$ の $k'$ への基底変換として、$k'$ 上
幾何学的正則である。$A \otimes_k k' \to A$ は、写像 $k' \to A$ による
$k' \otimes_k k' \to k'$ の基底変換であることに注意する。$k'/k$ は
エタールな環拡大なので、補題 \ref{algebra-lemma-etale} により
$k' \otimes_k k' \to k'$ はエタールである。従って補題
\ref{algebra-lemma-geometrically-regular-goes-up} により $A$ は $k'$ 上
幾何学的正則である。
\end{proof}

\section{幾何学的 Cohen--Macaulay 代数}
\label{algebra-section-geometrically-CM}

\noindent
この節の名称はいささか不適切である。Cohen--Macaulay 代数は自動的に
幾何学的 Cohen--Macaulay だからである。以下の補題
\ref{algebra-lemma-extend-field-CM-locus} および補題
\ref{algebra-lemma-CM-geometrically-CM} を参照されたい。

\begin{lemma}
\label{algebra-lemma-tensor-fields-CM}
$k$ を体とし、$K/k$ と $L/k$ を二つの体拡大で、その一方が有限型体拡大で
あるものとする。このとき $K \otimes_k L$ はネーター Cohen--Macaulay 環である。
\end{lemma}

\begin{proof}
補題 \ref{algebra-lemma-Noetherian-field-extension} により環 $K \otimes_k L$ は
ネーター環である。$K$ は純超越拡大 $k(t_1, \ldots, t_r)$ の有限拡大であると
する。このとき
$k(t_1, \ldots, t_r) \otimes_k L \to K \otimes_k L$ は有限自由な環準同型である。
補題 \ref{algebra-lemma-finite-flat-over-regular-CM} により、
$k(t_1, \ldots, t_r) \otimes_k L$ が Cohen--Macaulay であることを示せば十分である。
これは多項式環 $L[t_1, \ldots, t_r]$ の局所化なので明らかである
（多項式環が Cohen--Macaulay であることについては、例えば補題
\ref{algebra-lemma-CM-polynomial-algebra} を参照）。
\end{proof}

\begin{lemma}
\label{algebra-lemma-CM-geometrically-CM}
$k$ を体とし、$S$ をネーター $k$-代数とする。$K/k$ を有限生成体拡大とし、
$S_K = K \otimes_k S$ とおく。$\mathfrak q \subset S$ を $S$ の素イデアルとし、
$\mathfrak q_K \subset S_K$ を $\mathfrak q$ の上にある $S_K$ の素イデアルとする。
このとき $S_{\mathfrak q}$ が Cohen--Macaulay であることと
$(S_K)_{\mathfrak q_K}$ が Cohen--Macaulay であることは同値である。
\end{lemma}

\begin{proof}
補題 \ref{algebra-lemma-Noetherian-field-extension} により環 $S_K$ はネーター環である。
従って $S_{\mathfrak q} \to (S_K)_{\mathfrak q_K}$ はネーター局所環の
平坦な局所準同型である。そのファイバー
$$
(S_K)_{\mathfrak q_K} / \mathfrak q (S_K)_{\mathfrak q_K}
\cong (\kappa(\mathfrak q) \otimes_k K)_{\mathfrak q'}
$$
は、Cohen--Macaulay 環（補題 \ref{algebra-lemma-tensor-fields-CM}）
$\kappa(\mathfrak q) \otimes_k K$ の適当な素イデアル $\mathfrak q'$ における
局所化である。従って補題 \ref{algebra-lemma-CM-goes-up} から補題が従う。
\end{proof}

\section{余極限と有限表示の写像 II}
\label{algebra-section-colimits-finite-presentation}

\noindent
この節は節 \ref{algebra-section-colimits-flat} の続きである。

\medskip\noindent
まず平坦性の開性の応用から始める。これは平坦加群を平坦加群で近似できると
いう有用な結果である。

\begin{lemma}
\label{algebra-lemma-flat-finite-presentation-limit-flat}
$R \to S$ を環準同型とし、$M$ を $S$-加群とする。次を仮定する。
\begin{enumerate}
\item $R \to S$ は有限表示である。
\item $M$ は有限表示 $S$-加群である。
\item $M$ は $R$ 上平坦である。
\end{enumerate}
このとき次が成り立つ。
\begin{enumerate}
\item 有限型 $\mathbf{Z}$-代数 $R_0$、有限型環準同型 $R_0 \to S_0$、および
有限 $S_0$-加群 $M_0$ で、$M_0$ が $R_0$ 上平坦であるものが存在する。
% P01-E1034: frozen English says singular “a ring maps”.
さらに環準同型 $R_0 \to R$, $S_0 \to S$ と $S_0$-加群準同型
$M_0 \to M$ が存在して、$S \cong R \otimes_{R_0} S_0$ および
$M = S \otimes_{S_0} M_0$ が成り立つ。
\item $R = \colim_{\lambda \in \Lambda} R_\lambda$ が有向余極限として
書かれているならば、ある $\lambda$、有限表示の環準同型
$R_\lambda \to S_\lambda$、および有限表示 $S_\lambda$-加群 $M_\lambda$ が
存在し、$M_\lambda$ は $R_\lambda$ 上平坦であり、
$S = R \otimes_{R_\lambda} S_\lambda$ および
$M = S \otimes_{S_{\lambda}} M_\lambda$ が成り立つ。
\item
$$
(R \to S, M) =
\colim_{\lambda \in \Lambda}
(R_\lambda \to S_\lambda, M_\lambda)
$$
が次を満たす有向余極限として書かれているとする。
\begin{enumerate}
\item $\mu \geq \lambda$ に対して
$R_\mu \otimes_{R_\lambda} S_\lambda \to S_\mu$ および
$S_\mu \otimes_{S_\lambda} M_\lambda \to M_\mu$ は同型である。
\item $R_\lambda \to S_\lambda$ は有限表示である。
\item $M_\lambda$ は有限表示 $S_\lambda$-加群である。
\end{enumerate}
このとき、十分大きいすべての $\lambda$ に対して加群 $M_\lambda$ は
$R_\lambda$ 上平坦である。
\end{enumerate}
\end{lemma}

\begin{proof}
まず $(R \to S, M)$ を、補題
\ref{algebra-lemma-limit-module-finite-presentation} におけるような系
$(R_\lambda \to S_\lambda, M_\lambda)$ の有向余極限として書く。
素イデアル $\mathfrak q \subset S$ をとる。$\mathfrak p \subset R$、
$\mathfrak q_\lambda \subset S_\lambda$、および
$\mathfrak p_\lambda \subset R_\lambda$ を対応する素イデアルとする。
定理 \ref{algebra-theorem-openness-flatness} の証明で見たように、
$$
((R_\lambda)_{\mathfrak p_\lambda},
(S_\lambda)_{\mathfrak q_\lambda},
(M_\lambda)_{\mathfrak q_{\lambda}})
$$
は補題 \ref{algebra-lemma-limit-module-essentially-finite-presentation} におけるような
系である。従って補題 \ref{algebra-lemma-colimit-eventually-flat} により、ある
$\lambda_{\mathfrak q} \in \Lambda$ が存在し、すべての
$\lambda \geq \lambda_{\mathfrak q}$ に対して、加群 $M_\lambda$ は素イデアル
$\mathfrak q_{\lambda}$ において $R_\lambda$ 上平坦である。

\medskip\noindent
定理 \ref{algebra-theorem-openness-flatness} により、$U_\lambda$ のすべての
素イデアルにおいて $M_\lambda$ が $R_\lambda$ 上平坦となる開部分集合
$U_\lambda \subset \Spec(S_\lambda)$ を得る。$\Spec(S) \to \Spec(S_\lambda)$ に
よる $U_\lambda$ の逆像を $V_\lambda \subset \Spec(S)$ と書く。
上の議論により、各 $\mathfrak q \in \Spec(S)$ に対して、ある
$\lambda_{\mathfrak q}$ が存在し、すべての
$\lambda \geq \lambda_{\mathfrak q}$ について $\mathfrak q \in V_\lambda$ となる。
$\Spec(S)$ は準コンパクトなので、これはただ一つの
$\lambda_0 \in \Lambda$ が存在して $V_{\lambda_0} = \Spec(S)$ となることを
含意する。

\medskip\noindent
補集合 $\Spec(S_{\lambda_0}) \setminus U_{\lambda_0}$ は、あるイデアル
$I \subset S_{\lambda_0}$ に対する $V(I)$ である。
$V_{\lambda_0} = \Spec(S)$ なので $IS = S$ である。
% P01-E1035: frozen source pairs f_1,...,f_r with s_1,...,s_n in one sum.
$f_1, \ldots, f_r \in I$ と $s_1, \ldots, s_n \in S$ を
$\sum f_i s_i = 1$ となるように選ぶ。$\colim S_\lambda = S$ なので、
$\lambda_0$ を大きくすれば、$\sum f_i s_{i, \lambda_0} = 1$ となる
$s_{i, \lambda_0} \in S_{\lambda_0}$ が存在すると仮定できる。
従ってこの $\lambda_0$ に対して $U_{\lambda_0} = \Spec(S_{\lambda_0})$ である。
これで (1) が証明された。

\medskip\noindent
(2) を証明する。$(R_0 \to S_0, M_0)$ を (1) におけるものとし、
$R = \colim R_\lambda$ と仮定する。$R_0$ は有限型 $\mathbf{Z}$-代数なので、
ある $\lambda$ と写像 $R_0 \to R_\lambda$ が存在し、
$R_0 \to R_\lambda \to R$ は与えられた写像 $R_0 \to R$ となる
（補題 \ref{algebra-lemma-characterize-finite-presentation} を参照）。このとき
$S_\lambda = R_\lambda \otimes_{R_0} S_0$ および
$M_\lambda = S_\lambda \otimes_{S_0} M_0$ とおけば (2) が従う。

\medskip\noindent
最後に (3) を証明する。$(R_\lambda \to S_\lambda, M_\lambda)$ を (3) に
おけるものとする。$(R_0 \to S_0, M_0)$ と $R_0 \to R$ を (1) に
おけるものとして選ぶ。(2) の証明と同様に、ある $\lambda_0$ と環準同型
$R_0 \to R_{\lambda_0}$ が存在し、$R_0 \to R_{\lambda_0} \to R$ は
与えられた写像 $R_0 \to R$ となる。$S_0$ は $R_0$ 上有限表示で、かつ
$S = \colim S_\lambda$ なので、ある $\lambda_1 \geq \lambda_0$ に対して
$R_0$-代数準同型 $S_0 \to S_{\lambda_1}$ が得られ、その合成
$S_0 \to S_{\lambda_1} \to S$ は与えられた写像 $S_0 \to S$ となる
（補題 \ref{algebra-lemma-characterize-finite-presentation} を参照）。
すべての $\lambda \geq \lambda_1$ に対して、これにより写像
$$
\Psi_{\lambda} :
R_\lambda \otimes_{R_0} S_0
\longrightarrow
R_\lambda \otimes_{R_{\lambda_1}} S_{\lambda_1}
\cong
S_\lambda
$$
を得る。最後の同型は仮定による。構成から
$\colim_\lambda \Psi_\lambda$ は同型である。従って補題
\ref{algebra-lemma-colimit-category-fp-algebras} により、十分大きいすべての
$\lambda$ に対して $\Psi_\lambda$ は同型である。同様に、ある
$\lambda_2 \geq \lambda_1$ と $S_0$-加群準同型 $M_0 \to M_{\lambda_2}$ が
存在し、$M_0 \to M_{\lambda_2} \to M$ は与えられた写像 $M_0 \to M$ となる
（補題 \ref{algebra-lemma-module-map-property-in-colimit} を参照）。
$\lambda \geq \lambda_2$ に対して誘導される写像
$$
S_\lambda \otimes_{S_0} M_0
\longrightarrow
S_\lambda \otimes_{S_{\lambda_2}} M_{\lambda_2}
\cong
M_\lambda
$$
があり、十分大きい $\lambda$ に対してこの写像は補題
\ref{algebra-lemma-colimit-category-fp-modules} により同型である。
$M_0$ は $R_0$ 上平坦なので、これは (3) を含意する。
\end{proof}

\begin{lemma}
\label{algebra-lemma-descend-faithfully-flat-finite-presentation}
$R \to A \to B$ を環準同型とし、$A \to B$ は有限表示かつ忠実平坦であると
仮定する。このとき可換図式
$$
\xymatrix{
R \ar[r] \ar@{=}[d] &
A_0 \ar[d] \ar[r] &
B_0 \ar[d] \\
R \ar[r] & A \ar[r] & B
}
$$
であって、$R \to A_0$ が有限表示、$A_0 \to B_0$ が有限表示かつ忠実平坦で、
$B = A \otimes_{A_0} B_0$ となるものが存在する。
\end{lemma}

\begin{proof}
% P01-E1036: frozen English omits “by” after “replaced”.
まず $R$ を $\mathbf{Z}$ で置き換えた場合に補題を証明する。
補題 \ref{algebra-lemma-flat-finite-presentation-limit-flat} により図式
$$
\xymatrix{
A_0 \ar[r] \ar[d] & A \ar[d] \\
B_0 \ar[r] & B
}
$$
が存在する。ここで $A_0$ は $\mathbf{Z}$ 上有限型であり、$B_0$ は $A_0$ 上
有限表示かつ平坦で、$B = A \otimes_{A_0} B_0$ となる。
$A_0 \to B_0$ は有限表示かつ平坦なので、
$\Spec(B_0) \to \Spec(A_0)$ の像は開である
（命題 \ref{algebra-proposition-fppf-open} を参照）。従って像の補集合は、ある
イデアル $I_0 \subset A_0$ に対する $V(I_0)$ である。$A \to B$ は
忠実平坦なので、写像 $\Spec(B) \to \Spec(A)$ は全射である
（補題 \ref{algebra-lemma-ff-rings} を参照）。ここで、像の基底変換は基底変換の像で
あることを用いる。従って $I_0A = A$ である。$r_i \in A$, $f_i \in I_0$ に
対して関係式 $\sum f_i r_i = 1$ を選ぶ。$A_0$ を元 $r_i$ を含むように拡大し
（それに応じて $B_0$ も拡大し）た後、$A_0 \to B_0$ もスペクトル上全射、
すなわち忠実平坦であることが分かる。

\medskip\noindent
従って $R = \mathbf{Z}$ の場合に補題は成り立つ。一般の場合、今得た解
$A_0' \to B_0'$ をとり、$A_0 = A_0' \otimes_{\mathbf{Z}} R$、
$B_0 = B_0' \otimes_{\mathbf{Z}} R$ とおく。
\end{proof}

\begin{lemma}
\label{algebra-lemma-colimit-finite}
$A = \colim_{i \in I} A_i$ を環の有向余極限とする。
% P01-E1037: frozen English omits “be” in this and six parallel declarations.
$0 \in I$ とし、$\varphi_0 : B_0 \to C_0$ を $A_0$-代数の準同型とする。
次を仮定する。
\begin{enumerate}
\item $A \otimes_{A_0} B_0 \to A \otimes_{A_0} C_0$ は有限である。
\item $C_0$ は $B_0$ 上有限型である。
\end{enumerate}
このとき、ある $i \geq 0$ に対して写像
$A_i \otimes_{A_0} B_0 \to A_i \otimes_{A_0} C_0$ は有限である。
\end{lemma}

\begin{proof}
$x_1, \ldots, x_m$ を $B_0$ 上の $C_0$ の生成元とする。
$P_j(1 \otimes x_j) = 0$ が $A \otimes_{A_0} C_0$ において成り立つような
モニック多項式 $P_j \in A \otimes_{A_0} B_0[T]$ を選ぶ。ある
$i \geq 0$ に対して、$P_j$ に写る
$P_{j, i} \in A_i \otimes_{A_0} B_0[T]$ をとれる。テンソル積 $\otimes$ は
余極限と可換なので、必要なら $i$ を大きくすることで、
$P_{j, i}(1 \otimes x_j)$ は $A_i \otimes_{A_0} C_0$ において零となる。
この $i$ が求めるものである。
\end{proof}

\begin{lemma}
\label{algebra-lemma-colimit-surjective}
$A = \colim_{i \in I} A_i$ を環の有向余極限とする。$0 \in I$ とし、
$\varphi_0 : B_0 \to C_0$ を $A_0$-代数の準同型とする。次を仮定する。
\begin{enumerate}
\item $A \otimes_{A_0} B_0 \to A \otimes_{A_0} C_0$ は全射である。
\item $C_0$ は $B_0$ 上有限型である。
\end{enumerate}
このとき、ある $i \geq 0$ に対して写像
$A_i \otimes_{A_0} B_0 \to A_i \otimes_{A_0} C_0$ は全射である。
\end{lemma}

\begin{proof}
$x_1, \ldots, x_m$ を $B_0$ 上の $C_0$ の生成元とする。
$A \otimes_{A_0} C_0$ において $1 \otimes x_j$ に写る
$b_j \in A \otimes_{A_0} B_0$ を選ぶ。ある $i \geq 0$ に対して、$b_j$ に写る
$b_{j, i} \in A_i \otimes_{A_0} B_0$ をとれる。$i$ を大きくすれば、すべての
$j = 1, \ldots, m$ に対して $b_{j, i}$ は $A_i \otimes_{A_0} C_0$ における
$1 \otimes x_j$ に写ると仮定できる。この $i$ が求めるものである。
\end{proof}

\begin{lemma}
\label{algebra-lemma-colimit-unramified}
$A = \colim_{i \in I} A_i$ を環の有向余極限とする。$0 \in I$ とし、
$\varphi_0 : B_0 \to C_0$ を $A_0$-代数の準同型とする。次を仮定する。
\begin{enumerate}
\item $A \otimes_{A_0} B_0 \to A \otimes_{A_0} C_0$ は不分岐である。
\item $C_0$ は $B_0$ 上有限型である。
\end{enumerate}
このとき、ある $i \geq 0$ に対して写像
$A_i \otimes_{A_0} B_0 \to A_i \otimes_{A_0} C_0$ は不分岐である。
\end{lemma}

\begin{proof}
$B_i = A_i \otimes_{A_0} B_0$, $C_i = A_i \otimes_{A_0} C_0$,
$B = A \otimes_{A_0} B_0$, $C = A \otimes_{A_0} C_0$ とおく。
$x_1, \ldots, x_m$ を $B_0$ 上の $C_0$ の生成元とする。このとき
$\text{d}x_1, \ldots, \text{d}x_m$ は $C_0$ 上 $\Omega_{C_0/B_0}$ を生成し、
その像は $C_i$ 上 $\Omega_{C_i/B_i}$ を生成する
（補題 \ref{algebra-lemma-differentials-polynomial-ring} および
\ref{algebra-lemma-differential-seq}）。
$0 = \Omega_{C/B} = \colim \Omega_{C_i/B_i}$ であることに注意する
（補題 \ref{algebra-lemma-colimit-differentials}）。従って、ある $i$ に対して
$\text{d}x_1, \ldots, \text{d}x_m$ は零に写り、望む通り
$\Omega_{C_i/B_i} = 0$ となる。
\end{proof}

\begin{lemma}
\label{algebra-lemma-colimit-isomorphism}
$A = \colim_{i \in I} A_i$ を環の有向余極限とする。$0 \in I$ とし、
$\varphi_0 : B_0 \to C_0$ を $A_0$-代数の準同型とする。次を仮定する。
\begin{enumerate}
\item $A \otimes_{A_0} B_0 \to A \otimes_{A_0} C_0$ は同型である。
\item $B_0 \to C_0$ は有限表示である。
\end{enumerate}
このとき、ある $i \geq 0$ に対して写像
$A_i \otimes_{A_0} B_0 \to A_i \otimes_{A_0} C_0$ は同型である。
\end{lemma}

\begin{proof}
補題 \ref{algebra-lemma-colimit-surjective} により、写像
$A_i \otimes_{A_0} B_0 \to A_i \otimes_{A_0} C_0$ が全射となる $i$ が存在する。
この写像は有限表示なので、その核は有限生成イデアルである。
$g_1, \ldots, g_r \in A_i \otimes_{A_0} B_0$ を核の生成元とする。
このとき $g_j$ が $A_{i'} \otimes_{A_0} B_0$ において零に写るような
$i' \geq i$ を選べる。すると
$A_{i'} \otimes_{A_0} B_0 \to A_{i'} \otimes_{A_0} C_0$ は同型である。
\end{proof}

\begin{lemma}
\label{algebra-lemma-colimit-etale}
$A = \colim_{i \in I} A_i$ を環の有向余極限とする。$0 \in I$ とし、
$\varphi_0 : B_0 \to C_0$ を $A_0$-代数の準同型とする。次を仮定する。
\begin{enumerate}
\item $A \otimes_{A_0} B_0 \to A \otimes_{A_0} C_0$ はエタールである。
\item $B_0 \to C_0$ は有限表示である。
\end{enumerate}
このとき、ある $i \geq 0$ に対して写像
$A_i \otimes_{A_0} B_0 \to A_i \otimes_{A_0} C_0$ はエタールである。
\end{lemma}

\begin{proof}
$C_0 = B_0[x_1, \ldots, x_n]/(f_{1, 0}, \ldots, f_{m, 0})$ と書く。
$B_i = A_i \otimes_{A_0} B_0$, $C_i = A_i \otimes_{A_0} C_0$ と書く。
ここで $C_i = B_i[x_1, \ldots, x_n]/(f_{1, i}, \ldots, f_{m, i})$ であり、
$f_{j, i}$ は $B_i$ 上の多項式環における $f_{j, 0}$ の像である。
$B = A \otimes_{A_0} B_0$, $C = A \otimes_{A_0} C_0$ と書く。
ここで $C = B[x_1, \ldots, x_n]/(f_1, \ldots, f_m)$ であり、$f_j$ は
$B$ 上の多項式環における $f_{j, 0}$ の像である。仮定は写像
$$
\text{d} :
(f_1, \ldots, f_m)/(f_1, \ldots, f_m)^2
\longrightarrow
\bigoplus C \text{d}x_k
$$
が同型であるということである。従って十分大きい $i$ に対して、
$$
\xi_{k, i} \in (f_{1, i}, \ldots, f_{m, i})/(f_{1, i}, \ldots, f_{m, i})^2
$$
であって、$\bigoplus C_i \text{d}x_k$ において
$\text{d}\xi_{k, i} = \text{d}x_k$ となる元をとれる。さらに、必要なら
$i$ を大きくすることで、極限において成り立つことから
$\sum (\partial f_{j, i}/\partial x_k) \xi_{k, i} =
f_{j, i} \bmod (f_{1, i}, \ldots, f_{m, i})^2$ となる。
この $i$ が求めるものである。
\end{proof}

\begin{lemma}
\label{algebra-lemma-colimit-smooth}
$A = \colim_{i \in I} A_i$ を環の有向余極限とする。$0 \in I$ とし、
$\varphi_0 : B_0 \to C_0$ を $A_0$-代数の準同型とする。次を仮定する。
\begin{enumerate}
\item $A \otimes_{A_0} B_0 \to A \otimes_{A_0} C_0$ は滑らかである。
\item $B_0 \to C_0$ は有限表示である。
\end{enumerate}
このとき、ある $i \geq 0$ に対して写像
$A_i \otimes_{A_0} B_0 \to A_i \otimes_{A_0} C_0$ は滑らかである。
\end{lemma}

\begin{proof}
$C_0 = B_0[x_1, \ldots, x_n]/(f_{1, 0}, \ldots, f_{m, 0})$ と書く。
$B_i = A_i \otimes_{A_0} B_0$, $C_i = A_i \otimes_{A_0} C_0$ と書く。
ここで $C_i = B_i[x_1, \ldots, x_n]/(f_{1, i}, \ldots, f_{m, i})$ であり、
$f_{j, i}$ は $B_i$ 上の多項式環における $f_{j, 0}$ の像である。
$B = A \otimes_{A_0} B_0$, $C = A \otimes_{A_0} C_0$ と書く。
ここで $C = B[x_1, \ldots, x_n]/(f_1, \ldots, f_m)$ であり、$f_j$ は
$B$ 上の多項式環における $f_{j, 0}$ の像である。仮定は写像
$$
\text{d} :
(f_1, \ldots, f_m)/(f_1, \ldots, f_m)^2
\longrightarrow
\bigoplus C \text{d}x_k
$$
が分裂単射であるということである。
$\xi_k \in (f_1, \ldots, f_m)/(f_1, \ldots, f_m)^2$ を
$\sum (\partial f_j/\partial x_k) \xi_k =
f_j \bmod (f_1, \ldots, f_m)^2$ となる元とする。このとき十分大きい $i$ に
対して、
$$
\xi_{k, i} \in (f_{1, i}, \ldots, f_{m, i})/(f_{1, i}, \ldots, f_{m, i})^2
$$
であって、極限において成り立つことから
$\sum (\partial f_{j, i}/\partial x_k) \xi_{k, i} =
f_{j, i} \bmod (f_{1, i}, \ldots, f_{m, i})^2$ となる元をとれる。
この $i$ が求めるものである。
\end{proof}

\begin{lemma}
\label{algebra-lemma-colimit-lci}
$A = \colim_{i \in I} A_i$ を環の有向余極限とする。$0 \in I$ とし、
$\varphi_0 : B_0 \to C_0$ を $A_0$-代数の準同型とする。次を仮定する。
\begin{enumerate}
\item $A \otimes_{A_0} B_0 \to A \otimes_{A_0} C_0$ はシントミックである
（それぞれ、相対的大域的完全交差である）。
\item $C_0$ は $B_0$ 上有限表示である。
\end{enumerate}
このとき、ある $i \geq 0$ に対して写像
$A_i \otimes_{A_0} B_0 \to A_i \otimes_{A_0} C_0$ はシントミックである
（それぞれ、相対的大域的完全交差である）。
\end{lemma}

\begin{proof}
$A \otimes_{A_0} B_0 \to A \otimes_{A_0} C_0$ は相対的大域的完全交差であると
仮定する。補題 \ref{algebra-lemma-relative-global-complete-intersection-Noetherian} により、
有限型 $\mathbf{Z}$-代数 $R$、環準同型
$R \to A \otimes_{A_0} B_0$、相対的大域的完全交差 $R \to S$、および同型
$$
(A \otimes_{A_0} B_0) \otimes_R S
\longrightarrow
A \otimes_{A_0} C_0
$$
が存在する。$R$ は $\mathbf{Z}$ 上有限型（従って有限表示）なので、ある
$i$ と、写像 $R \to A \otimes_{A_0} B_0$ を持ち上げる写像
$R \to A_i \otimes_{A_0} B_0$ が存在する
（補題 \ref{algebra-lemma-characterize-finite-presentation} を参照）。同じ補題により、
ある $i' \geq i$ が存在して、
$(A_i \otimes_{A_0} B_0) \otimes_R S \to A \otimes_{A_0} C_0$ は写像
$(A_i \otimes_{A_0} B_0) \otimes_R S \to A_{i'} \otimes_{A_0} C_0$ から得られる。
従って $i$ を $i'$ で置き換えた後、表示された写像は
$A_i \otimes_{A_0} B_0$-代数準同型
$$
(A_i \otimes_{A_0} B_0) \otimes_R S
\longrightarrow
A_i \otimes_{A_0} C_0
$$
から得られると仮定してよい。補題 \ref{algebra-lemma-colimit-isomorphism} により、
$i$ を大きくすればこの写像は同型である。相対的大域的完全交差の基底変換は
補題 \ref{algebra-lemma-base-change-relative-global-complete-intersection} により
相対的大域的完全交差なので、この場合の証明は終わる。

\medskip\noindent
$A \otimes_{A_0} B_0 \to A \otimes_{A_0} C_0$ はシントミックであると仮定する。
単位イデアルを生成する元 $g_1, \ldots, g_m$ が
$A \otimes_{A_0} C_0$ に存在し、
$A \otimes_{A_0} B_0 \to (A \otimes_{A_0} C_0)_{g_j}$ は相対的大域的完全交差で
ある（補題 \ref{algebra-lemma-syntomic} を参照）。ある $i$ と、$g_j$ に写る元
$g_{i, j} \in A_i \otimes_{A_0} C_0$ をとれる。$i$ を大きくすれば、
$g_{i, 1}, \ldots, g_{i, m}$ は $A_i \otimes_{A_0} C_0$ の単位イデアルを
生成すると仮定できる。前段落の結果により、さらに $i$ を大きくして、写像
$A_i \otimes_{A_0} B_0 \to (A_i \otimes_{A_0} C_0)_{g_{i, j}}$ は
相対的大域的完全交差であると仮定できる。従って補題
\ref{algebra-lemma-local-syntomic}（および既に用いた補題 \ref{algebra-lemma-syntomic}）により、
$A_i \otimes_{A_0} B_0 \to A_i \otimes_{A_0} C_0$ はシントミックである。
\end{proof}














\noindent
次の補題は上の結果の応用であるが、ほかのどこにも収まりがよくないように思われる。

\begin{lemma}
\label{algebra-lemma-fppf-fpqf}
$R \to S$ を有限表示の忠実平坦環準同型とする。このとき、可換図式
$$
\xymatrix{
S \ar[rr] & & S' \\
& R \ar[lu] \ar[ru]
}
$$
であって、$R \to S'$ が準有限かつ忠実平坦で有限表示となるものが存在する。
\end{lemma}

\begin{proof}
まず、この補題を $R$ が $\mathbf{Z}$ 上有限型である場合に帰着する。
補題 \ref{algebra-lemma-descend-faithfully-flat-finite-presentation} により、図式
$$
\xymatrix{
S_0 \ar[r] & S \\
R_0 \ar[u] \ar[r] & R \ar[u]
}
$$
であって、$R_0$ が $\mathbf{Z}$ 上有限型であり、$S_0$ が $R_0$ 上
有限表示の忠実平坦代数で、$S = R \otimes_{R_0} S_0$ となるものが存在する。
環準同型 $R_0 \to S_0$ に対して補題を証明すれば、$R \to S$ に対する補題は
基底変換によって従う。実際、準有限環準同型の基底変換は準有限である
（補題 \ref{algebra-lemma-quasi-finite-base-change} を参照）。もちろんここでは、
平坦準同型の基底変換が平坦であり、有限表示の準同型の基底変換が
有限表示であることも用いている。

\medskip\noindent
$R \to S$ は有限表示の忠実平坦環準同型であり、$R$ はネーターであると仮定する
（前段落により、このように仮定してよい）。補題
\ref{algebra-lemma-finite-presentation-flat-CM-locus-open} の開集合を
$W \subset \Spec(S)$ とする。$R \to S$ は忠実平坦なので、写像
$\Spec(S) \to \Spec(R)$ は全射である（補題 \ref{algebra-lemma-ff-rings} を参照）。
補題 \ref{algebra-lemma-generic-CM-flat-finite-presentation} により、写像
$W \to \Spec(R)$ も全射である。従って $S$ を積
$S_{g_1} \times \ldots \times S_{g_m}$ で置き換えることにより、
$W = \Spec(S)$ と仮定してよい。ここでは $\Spec(R)$ が準コンパクトであること
（補題 \ref{algebra-lemma-quasi-compact}）と、写像 $\Spec(S) \to \Spec(R)$ が開写像で
あること（命題 \ref{algebra-proposition-fppf-open}）を用いる。
$\mathfrak p \subset R$ を素イデアルとする。$\mathfrak p$ の上にあり、
ファイバー環 $S \otimes_R \kappa(\mathfrak p)$ の極大イデアルに対応する素イデアル
$\mathfrak q \subset S$ を選ぶ。ネーター局所環
$\overline{S}_{\mathfrak q} = S_{\mathfrak q}/\mathfrak pS_{\mathfrak q}$ は
Cohen--Macaulay であるとし、その次元を $d$ とする。$S_{\mathfrak q}$ の
極大イデアルに属し、$\overline{S}_{\mathfrak q}$ における像が正則列となる
$f_1, \ldots, f_d$ を選べる。$f_1, \ldots, f_d$ の共通分母
$g \in S$, $g \not \in \mathfrak q$ を選び、$R$-代数
$$
S' = S_g/(f_1, \ldots, f_d).
$$
を考える。構成により、$\mathfrak p$ の上にあり、$\mathfrak q$ に対応する
（$S_g \to S'_g$ を介して）素イデアル $\mathfrak q' \subset S'$ が存在する。
% P01-E1038: frozen source names the prime of S rather than the corresponding prime of S'.
また構成により、局所環
% P01-E1039: frozen source leaves the denominator sum unparenthesized; preserved verbatim.
$$
S'_{\mathfrak q'}/\mathfrak pS'_{\mathfrak q'} =
S_{\mathfrak q}/(f_1, \ldots, f_d) + \mathfrak pS_{\mathfrak q} =
\overline{S}_{\mathfrak q}/(\overline{f}_1, \ldots, \overline{f}_d)
$$
が次元零であるため（補題 \ref{algebra-lemma-isolated-point-fibre} を参照）、環準同型
$R \to S'$ は $\mathfrak q$ において準有限である。また構成により、
$R \to S'$ は有限表示である。最後に、補題
\ref{algebra-lemma-grothendieck-regular-sequence} により、局所環準同型
$R_{\mathfrak p} \to S'_{\mathfrak q'}$ は平坦である（ここで $R$ がネーターで
あることを用いる）。従って、平坦性が開であること
（定理 \ref{algebra-theorem-openness-flatness}）と準有限性が開であること
（補題 \ref{algebra-lemma-quasi-finite-open}）により、適当な $g' \in S$,
$g' \not \in \mathfrak q$ に対して $g$ を $gg'$ で置き換えた後、
$R \to S'$ は平坦かつ準有限であると仮定してよい。
像 $\Spec(S') \to \Spec(R)$ は開であり、$\mathfrak p$ を含む。
言い換えれば、任意に与えられた素イデアルを像に含む、補題の主張にあるような
環 $S'$ が（忠実性を除けば）存在することを示した。$\Spec(R)$ の
準コンパクト性をもう一度用いれば、このような環の有限個の積が求める条件を
満たすことが分かる。
\end{proof}
